\subsection{Benchmarking methodology}
\label{sec:benchmark-results}

The following tables present results comparins BioseqDB and minimap2 on one Macos machine and 4 AWS machines.

\begin{itemize}
  \item \textbf{Index.} A persisted index is constructed before measurement.
  Each repeat starts a fresh process and measures reading the index and mapping
  the queries. Construction of the persisted index is excluded. Operating-system page caches are not purged, so
  this is process-cold execution, but not a guaranteed cold disk read (I decided to measure it this way, because minimap2 still has a chance to use system page cache).
  \item \textbf{No index.} No pre-existing user index is available.
  \item \textbf{Two measurements per cell.} The first duration, A, is the
  nearest-rank P50 of the complete measured interval. The second, B, is the
  nearest-rank P50 after subtracting the union of instrumented CIGAR
  and MAPQ in each repeat
\end{itemize}

By default chunk cache for decompressed columns is disabled.
Each repeat runs 1,024 queries against one reference sequence:
\begin{itemize}
  \item References contain exactly
  $2\times10^9$ or $4\times10^9$ base pairs
  The parameter pairs are $(k,w)=(15,10)$ and $(20,15)$,
  \item Each query and target is randomly generated. DNA4 uses same probability for each base (25\%). The DNA5
  generator assigns probability 1\% to N and 24.75\% to each base.
  \item Query samples length is 20,000 bp and standard deviation 500 bp (before applying insertions and deletions), it's reverse-complemented with probability 0.5.
  \item Substitution probability is 0.001 and for insert/deletions it's 0.0005.
  \item Each query originates from one known sampled
  reference interval and later substitutions/insertions/deletions are applied to it.
\end{itemize}

BioseqDB uses task-parallel (most parallel-friendly mode) execution with prefiltering disabled and both CIGAR
and MAPQ projection enabled. 

AWS machines were missing GPUs so all Vulkan rows are grouped together and marked as not available.

\paragraph{Reading comparisons and highlights.}
Each cell places the A duration above its comparison, followed by the B duration
and its own comparison. Durations are rounded to tenths of a second for display.
All comparisons and minima use the original integer nanoseconds.
\begin{itemize}
  \item Minimap2 (light-gray) as the baseline is first row in each table. 
  \item All of the rows apart from minimap2 are ordered top-to-bottom by the minimum A time for the given "index" configuration.
  \item A BioseqDB value is compared with minimap2 for the same configuration (green/red triangles), while the minimap2 row compares with the fastest
  measured eligible BioseqDB effective path in that same column
  \item Pale yellow
  identifies the preferred BioseqDB configuration: MPF / Bounded / SIMD-DP /
  Block. This is the author's (my) selected configuration.
  \item A dot means there is no measurement, while a horizontal dash with \textsuperscript{g} identifies that GPU configuration wasn't available (no GPU on AWS machines)
\end{itemize}


% Generated by papers/scripts/generate_benchmark_results.py. Do not edit.
\definecolor{BenchBest}{HTML}{DDECF8}
\definecolor{BenchBaseline}{HTML}{F4F4F4}
\definecolor{BenchPreferred}{HTML}{FFFBEA}
\newgeometry{left=10mm,right=10mm,top=12mm,bottom=12mm}
\begin{landscape}
\subsection{Complete raw results}
\begin{table}[H]\centering\small
\caption{Machine specs}\label{tab:benchmark-machines}
\begin{tabularx}{\linewidth}{@{}p{24mm}*{5}{>{\raggedright\arraybackslash}X}@{}}\toprule
Property & Local & M7A & M7I & M4 & C7G\\\midrule
Machine & MacBook Pro & m7a.xlarge & m7i.xlarge & m4.xlarge & c7g.xlarge\\[3pt]
CPU & Apple M1 Pro & AMD EPYC 9R14 & Xeon Platinum 8488C & Xeon E5-2686 v4 @ 2.30GHz & Neoverse-V1\\[3pt]
Architecture & arm64 & x86\_64 & x86\_64 & x86\_64 & aarch64\\[3pt]
OS version / kernel & macOS 15.6 / Darwin 24.6.0 & Linux / 7.0.0-1012-aws & Linux / 7.0.0-1012-aws & Linux / 7.0.0-1012-aws & Linux / 7.0.0-1012-aws\\[3pt]
OS build / distribution & 24G5054d & Distribution not recorded & Distribution not recorded & Distribution not recorded & Distribution not recorded\\[3pt]
Logical / physical cores & 8 / 8 (6P + 2E) & 4 / 4 & 4 / 2 & 4 / 2 & 4 / 4\\[3pt]
Physical RAM & 16 GiB unified & 16 GiB & 16 GiB & 16 GiB & 8 GiB\\[3pt]
GPU & Apple M1 Pro, 14 cores & None available & None available & None available & None available\\[3pt]
Storage & APPLE SSD AP0512R & EBS gp3 & EBS gp3 & EBS gp3 & EBS gp3\\[3pt]
Disk capacity & 500.28 GB & 500 GiB & 500 GiB & 500 GiB & 500 GiB\\[3pt]
Provisioned disk I/O & Not recorded & 6,000 IOPS / 500 MiB/s & 6,000 IOPS / 500 MiB/s & 6,000 IOPS / 500 MiB/s & 6,000 IOPS / 500 MiB/s\\[3pt]
\bottomrule\end{tabularx}\end{table}
\clearpage
\refstepcounter{table}\label{tab:bench:local}
\noindent{\large\bfseries Table \thetable: (A): Local machine\hfill 1/12}\par\smallskip
{\fontsize{7.1}{8.1}\selectfont\setlength{\tabcolsep}{1pt}\setlength{\fboxsep}{2pt}\renewcommand{\arraystretch}{1.12}
\noindent\begin{tabular}{@{}p{14mm}p{27mm}|p{13.55mm}|p{13.55mm}|p{13.55mm}|p{13.55mm}|p{13.55mm}|p{13.55mm}|p{13.55mm}|p{13.55mm}|p{13.55mm}|p{13.55mm}|p{13.55mm}|p{13.55mm}|p{13.55mm}|p{13.55mm}|p{13.55mm}|p{13.55mm}|@{}}\hline
Engine & Execution & \multicolumn{8}{c|}{\textbf{Index}} & \multicolumn{8}{c|}{\textbf{No index}}\\
 &  & \multicolumn{4}{c|}{\textbf{4 Gbp}} & \multicolumn{4}{c|}{\textbf{2 Gbp}} & \multicolumn{4}{c|}{\textbf{4 Gbp}} & \multicolumn{4}{c|}{\textbf{2 Gbp}}\\
 &  & \multicolumn{2}{c|}{\textbf{k=15, w=10}} & \multicolumn{2}{c|}{\textbf{k=20, w=15}} & \multicolumn{2}{c|}{\textbf{k=15, w=10}} & \multicolumn{2}{c|}{\textbf{k=20, w=15}} & \multicolumn{2}{c|}{\textbf{k=15, w=10}} & \multicolumn{2}{c|}{\textbf{k=20, w=15}} & \multicolumn{2}{c|}{\textbf{k=15, w=10}} & \multicolumn{2}{c|}{\textbf{k=20, w=15}}\\
 &  & \multicolumn{1}{c|}{\textbf{DNA4}} & \multicolumn{1}{c|}{\textbf{DNA5}} & \multicolumn{1}{c|}{\textbf{DNA4}} & \multicolumn{1}{c|}{\textbf{DNA5}} & \multicolumn{1}{c|}{\textbf{DNA4}} & \multicolumn{1}{c|}{\textbf{DNA5}} & \multicolumn{1}{c|}{\textbf{DNA4}} & \multicolumn{1}{c|}{\textbf{DNA5}} & \multicolumn{1}{c|}{\textbf{DNA4}} & \multicolumn{1}{c|}{\textbf{DNA5}} & \multicolumn{1}{c|}{\textbf{DNA4}} & \multicolumn{1}{c|}{\textbf{DNA5}} & \multicolumn{1}{c|}{\textbf{DNA4}} & \multicolumn{1}{c|}{\textbf{DNA5}} & \multicolumn{1}{c|}{\textbf{DNA4}} & \multicolumn{1}{c|}{\textbf{DNA5}}\\
 & & \BenchmarkSourceLabel{bench:local:column:index-4-k15-w10-dna4-local}{C1}{\color{gray}C1} & \BenchmarkSourceLabel{bench:local:column:index-4-k15-w10-dna5-local}{C2}{\color{gray}C2} & \BenchmarkSourceLabel{bench:local:column:index-4-k20-w15-dna4-local}{C3}{\color{gray}C3} & \BenchmarkSourceLabel{bench:local:column:index-4-k20-w15-dna5-local}{C4}{\color{gray}C4} & \BenchmarkSourceLabel{bench:local:column:index-2-k15-w10-dna4-local}{C5}{\color{gray}C5} & \BenchmarkSourceLabel{bench:local:column:index-2-k15-w10-dna5-local}{C6}{\color{gray}C6} & \BenchmarkSourceLabel{bench:local:column:index-2-k20-w15-dna4-local}{C7}{\color{gray}C7} & \BenchmarkSourceLabel{bench:local:column:index-2-k20-w15-dna5-local}{C8}{\color{gray}C8} & \BenchmarkSourceLabel{bench:local:column:noindex-4-k15-w10-dna4-local}{C9}{\color{gray}C9} & \BenchmarkSourceLabel{bench:local:column:noindex-4-k15-w10-dna5-local}{C10}{\color{gray}C10} & \BenchmarkSourceLabel{bench:local:column:noindex-4-k20-w15-dna4-local}{C11}{\color{gray}C11} & \BenchmarkSourceLabel{bench:local:column:noindex-4-k20-w15-dna5-local}{C12}{\color{gray}C12} & \BenchmarkSourceLabel{bench:local:column:noindex-2-k15-w10-dna4-local}{C13}{\color{gray}C13} & \BenchmarkSourceLabel{bench:local:column:noindex-2-k15-w10-dna5-local}{C14}{\color{gray}C14} & \BenchmarkSourceLabel{bench:local:column:noindex-2-k20-w15-dna4-local}{C15}{\color{gray}C15} & \BenchmarkSourceLabel{bench:local:column:noindex-2-k20-w15-dna5-local}{C16}{\color{gray}C16}\\
\hline
\rowcolor{BenchBaseline}[0.5pt][0.5pt]\BenchmarkSourceLabel{bench:local:row:minimap2}{R1}\shortstack[l]{{\color{gray}R1}\\\textbf{minimap2}} & \textbf{Baseline} & \hspace*{1pt}\begin{tabular}[t]{@{}c@{}}\rule{0pt}{12pt}39.7s\\[-.2pt]{\fontsize{5.8}{6.4}\selectfont \textcolor{BrickRed}{$ \blacktriangledown\,2.2\times $}}\\[-.2pt]36.8s\\[-.2pt]{\fontsize{5.8}{6.4}\selectfont \textcolor{BrickRed}{$ \blacktriangledown\,2.2\times $}}\end{tabular} & \hspace*{1pt}\begin{tabular}[t]{@{}c@{}}\rule{0pt}{12pt}25.2s\\[-.2pt]{\fontsize{5.8}{6.4}\selectfont \textcolor{BrickRed}{$ \blacktriangledown\,1.5\times $}}\\[-.2pt]22.7s\\[-.2pt]{\fontsize{5.8}{6.4}\selectfont \textcolor{BrickRed}{$ \blacktriangledown\,1.5\times $}}\end{tabular} & \hspace*{1pt}\begin{tabular}[t]{@{}c@{}}\rule{0pt}{12pt}1m\,56.9s\\[-.2pt]{\fontsize{5.8}{6.4}\selectfont \textcolor{BrickRed}{$ \blacktriangledown\,29.1\times $}}\\[-.2pt]1m\,53.2s\\[-.2pt]{\fontsize{5.8}{6.4}\selectfont \textcolor{BrickRed}{$ \blacktriangledown\,32.9\times $}}\end{tabular} & \hspace*{1pt}\begin{tabular}[t]{@{}c@{}}\rule{0pt}{12pt}24.9s\\[-.2pt]{\fontsize{5.8}{6.4}\selectfont \textcolor{BrickRed}{$ \blacktriangledown\,1.6\times $}}\\[-.2pt]22.2s\\[-.2pt]{\fontsize{5.8}{6.4}\selectfont \textcolor{BrickRed}{$ \blacktriangledown\,1.6\times $}}\end{tabular} & \hspace*{1pt}\begin{tabular}[t]{@{}c@{}}\rule{0pt}{12pt}28.6s\\[-.2pt]{\fontsize{5.8}{6.4}\selectfont \textcolor{BrickRed}{$ \blacktriangledown\,11.5\times $}}\\[-.2pt]25.8s\\[-.2pt]{\fontsize{5.8}{6.4}\selectfont \textcolor{BrickRed}{$ \blacktriangledown\,15.6\times $}}\end{tabular} & \hspace*{1pt}\begin{tabular}[t]{@{}c@{}}\rule{0pt}{12pt}15.3s\\[-.2pt]{\fontsize{5.8}{6.4}\selectfont \textcolor{BrickRed}{$ \blacktriangledown\,2.8\times $}}\\[-.2pt]12.7s\\[-.2pt]{\fontsize{5.8}{6.4}\selectfont \textcolor{BrickRed}{$ \blacktriangledown\,3.3\times $}}\end{tabular} & \hspace*{1pt}\begin{tabular}[t]{@{}c@{}}\rule{0pt}{12pt}29.3s\\[-.2pt]{\fontsize{5.8}{6.4}\selectfont \textcolor{BrickRed}{$ \blacktriangledown\,15.0\times $}}\\[-.2pt]26.2s\\[-.2pt]{\fontsize{5.8}{6.4}\selectfont \textcolor{BrickRed}{$ \blacktriangledown\,22.9\times $}}\end{tabular} & \hspace*{1pt}\begin{tabular}[t]{@{}c@{}}\rule{0pt}{12pt}8.7s\\[-.2pt]{\fontsize{5.8}{6.4}\selectfont \textcolor{BrickRed}{$ \blacktriangledown\,1.7\times $}}\\[-.2pt]8.1s\\[-.2pt]{\fontsize{5.8}{6.4}\selectfont \textcolor{BrickRed}{$ \blacktriangledown\,2.3\times $}}\end{tabular} & \hspace*{1pt}\begin{tabular}[t]{@{}c@{}}\rule{0pt}{12pt}\colorbox{BenchBest}{\strut 2m\,52.4s}\\[-.2pt]{\fontsize{5.8}{6.4}\selectfont \textcolor{ForestGreen}{$ \blacktriangle\,19\% $}}\\[-.2pt]\colorbox{BenchBest}{\strut 2m\,48.5s}\\[-.2pt]{\fontsize{5.8}{6.4}\selectfont \textcolor{ForestGreen}{$ \blacktriangle\,15\% $}}\end{tabular} & \hspace*{1pt}\begin{tabular}[t]{@{}c@{}}\rule{0pt}{12pt}\colorbox{BenchBest}{\strut 2m\,17.1s}\\[-.2pt]{\fontsize{5.8}{6.4}\selectfont \textcolor{ForestGreen}{$ \blacktriangle\,3.9\times $}}\\[-.2pt]\colorbox{BenchBest}{\strut 2m\,15.9s}\\[-.2pt]{\fontsize{5.8}{6.4}\selectfont \textcolor{ForestGreen}{$ \blacktriangle\,1.7\times $}}\end{tabular} & \hspace*{1pt}\begin{tabular}[t]{@{}c@{}}\rule{0pt}{12pt}2m\,55s\\[-.2pt]{\fontsize{5.8}{6.4}\selectfont \textcolor{BrickRed}{$ \blacktriangledown\,6\% $}}\\[-.2pt]2m\,51s\\[-.2pt]{\fontsize{5.8}{6.4}\selectfont \textcolor{BrickRed}{$ \blacktriangledown\,12\% $}}\end{tabular} & \hspace*{1pt}\begin{tabular}[t]{@{}c@{}}\rule{0pt}{12pt}\colorbox{BenchBest}{\strut 55.9s}\\[-.2pt]{\fontsize{5.8}{6.4}\selectfont \textcolor{ForestGreen}{$ \blacktriangle\,8.6\times $}}\\[-.2pt]\colorbox{BenchBest}{\strut 55.1s}\\[-.2pt]{\fontsize{5.8}{6.4}\selectfont \textcolor{ForestGreen}{$ \blacktriangle\,3.1\times $}}\end{tabular} & \hspace*{1pt}\begin{tabular}[t]{@{}c@{}}\rule{0pt}{12pt}\colorbox{BenchBest}{\strut 43.3s}\\[-.2pt]{\fontsize{5.8}{6.4}\selectfont \textcolor{ForestGreen}{$ \blacktriangle\,2.6\times $}}\\[-.2pt]\colorbox{BenchBest}{\strut 42.5s}\\[-.2pt]{\fontsize{5.8}{6.4}\selectfont \textcolor{ForestGreen}{$ \blacktriangle\,2.4\times $}}\end{tabular} & \hspace*{1pt}\begin{tabular}[t]{@{}c@{}}\rule{0pt}{12pt}\colorbox{BenchBest}{\strut 38.9s}\\[-.2pt]{\fontsize{5.8}{6.4}\selectfont \textcolor{ForestGreen}{$ \blacktriangle\,8.0\times $}}\\[-.2pt]\colorbox{BenchBest}{\strut 38.2s}\\[-.2pt]{\fontsize{5.8}{6.4}\selectfont \textcolor{ForestGreen}{$ \blacktriangle\,2.8\times $}}\end{tabular} & \hspace*{1pt}\begin{tabular}[t]{@{}c@{}}\rule{0pt}{12pt}\colorbox{BenchBest}{\strut 53.1s}\\[-.2pt]{\fontsize{5.8}{6.4}\selectfont \textcolor{ForestGreen}{$ \blacktriangle\,1.6\times $}}\\[-.2pt]\colorbox{BenchBest}{\strut 48.2s}\\[-.2pt]{\fontsize{5.8}{6.4}\selectfont \textcolor{ForestGreen}{$ \blacktriangle\,1.5\times $}}\end{tabular} & \hspace*{1pt}\begin{tabular}[t]{@{}c@{}}\rule{0pt}{12pt}\colorbox{BenchBest}{\strut 28.7s}\\[-.2pt]{\fontsize{5.8}{6.4}\selectfont \textcolor{ForestGreen}{$ \blacktriangle\,9.8\times $}}\\[-.2pt]\colorbox{BenchBest}{\strut 28.2s}\\[-.2pt]{\fontsize{5.8}{6.4}\selectfont \textcolor{ForestGreen}{$ \blacktriangle\,2.9\times $}}\end{tabular}\\[3pt]\hline
\rowcolor{BenchPreferred}[0.5pt][0.5pt]\BenchmarkSourceLabel{bench:local:row:mpf-bounded-simd-block}{R2}\shortstack[l]{{\color{gray}R2}\\\textbf{BioseqDB}} & \begin{tabular}[t]{@{}l@{}}\rule{0pt}{12pt}MPF / Bounded\\SIMD-DP / Block\end{tabular} & \hspace*{1pt}\begin{tabular}[t]{@{}c@{}}\rule{0pt}{12pt}\colorbox{BenchBest}{\strut 17.7s}\\[-.2pt]{\fontsize{5.8}{6.4}\selectfont \textcolor{ForestGreen}{$ \blacktriangle\,2.2\times $}}\\[-.2pt]\colorbox{BenchBest}{\strut 17s}\\[-.2pt]{\fontsize{5.8}{6.4}\selectfont \textcolor{ForestGreen}{$ \blacktriangle\,2.2\times $}}\end{tabular} & \hspace*{1pt}\begin{tabular}[t]{@{}c@{}}\rule{0pt}{12pt}21.5s\\[-.2pt]{\fontsize{5.8}{6.4}\selectfont \textcolor{ForestGreen}{$ \blacktriangle\,15\% $}}\\[-.2pt]20.2s\\[-.2pt]{\fontsize{5.8}{6.4}\selectfont \textcolor{ForestGreen}{$ \blacktriangle\,11\% $}}\end{tabular} & \hspace*{1pt}\begin{tabular}[t]{@{}c@{}}\rule{0pt}{12pt}14.6s\\[-.2pt]{\fontsize{5.8}{6.4}\selectfont \textcolor{ForestGreen}{$ \blacktriangle\,8.0\times $}}\\[-.2pt]13.8s\\[-.2pt]{\fontsize{5.8}{6.4}\selectfont \textcolor{ForestGreen}{$ \blacktriangle\,8.2\times $}}\end{tabular} & \hspace*{1pt}\begin{tabular}[t]{@{}c@{}}\rule{0pt}{12pt}21s\\[-.2pt]{\fontsize{5.8}{6.4}\selectfont \textcolor{ForestGreen}{$ \blacktriangle\,16\% $}}\\[-.2pt]19.8s\\[-.2pt]{\fontsize{5.8}{6.4}\selectfont \textcolor{ForestGreen}{$ \blacktriangle\,11\% $}}\end{tabular} & \hspace*{1pt}\begin{tabular}[t]{@{}c@{}}\rule{0pt}{12pt}\colorbox{BenchBest}{\strut 2.5s}\\[-.2pt]{\fontsize{5.8}{6.4}\selectfont \textcolor{ForestGreen}{$ \blacktriangle\,11.5\times $}}\\[-.2pt]\colorbox{BenchBest}{\strut 1.7s}\\[-.2pt]{\fontsize{5.8}{6.4}\selectfont \textcolor{ForestGreen}{$ \blacktriangle\,15.6\times $}}\end{tabular} & \hspace*{1pt}\begin{tabular}[t]{@{}c@{}}\rule{0pt}{12pt}5.4s\\[-.2pt]{\fontsize{5.8}{6.4}\selectfont \textcolor{ForestGreen}{$ \blacktriangle\,2.8\times $}}\\[-.2pt]4.2s\\[-.2pt]{\fontsize{5.8}{6.4}\selectfont \textcolor{ForestGreen}{$ \blacktriangle\,3.1\times $}}\end{tabular} & \hspace*{1pt}\begin{tabular}[t]{@{}c@{}}\rule{0pt}{12pt}\colorbox{BenchBest}{\strut 2s}\\[-.2pt]{\fontsize{5.8}{6.4}\selectfont \textcolor{ForestGreen}{$ \blacktriangle\,15.0\times $}}\\[-.2pt]1.3s\\[-.2pt]{\fontsize{5.8}{6.4}\selectfont \textcolor{ForestGreen}{$ \blacktriangle\,20.4\times $}}\end{tabular} & \hspace*{1pt}\begin{tabular}[t]{@{}c@{}}\rule{0pt}{12pt}5.3s\\[-.2pt]{\fontsize{5.8}{6.4}\selectfont \textcolor{ForestGreen}{$ \blacktriangle\,1.6\times $}}\\[-.2pt]3.9s\\[-.2pt]{\fontsize{5.8}{6.4}\selectfont \textcolor{ForestGreen}{$ \blacktriangle\,2.1\times $}}\end{tabular} & \hspace*{1pt}\begin{tabular}[t]{@{}c@{}}\rule{0pt}{12pt}3m\,48.1s\\[-.2pt]{\fontsize{5.8}{6.4}\selectfont \textcolor{BrickRed}{$ \blacktriangledown\,1.3\times $}}\\[-.2pt]3m\,35s\\[-.2pt]{\fontsize{5.8}{6.4}\selectfont \textcolor{BrickRed}{$ \blacktriangledown\,28\% $}}\end{tabular} & \hspace*{1pt}\begin{tabular}[t]{@{}c@{}}\rule{0pt}{12pt}8m\,57.6s\\[-.2pt]{\fontsize{5.8}{6.4}\selectfont \textcolor{BrickRed}{$ \blacktriangledown\,3.9\times $}}\\[-.2pt]3m\,49.2s\\[-.2pt]{\fontsize{5.8}{6.4}\selectfont \textcolor{BrickRed}{$ \blacktriangledown\,1.7\times $}}\end{tabular} & \hspace*{1pt}\begin{tabular}[t]{@{}c@{}}\rule{0pt}{12pt}2m\,52s\\[-.2pt]{\fontsize{5.8}{6.4}\selectfont \textcolor{gray}{$ \blacktriangleright\,2\% $}}\\[-.2pt]2m\,39.3s\\[-.2pt]{\fontsize{5.8}{6.4}\selectfont \textcolor{ForestGreen}{$ \blacktriangle\,7\% $}}\end{tabular} & \hspace*{1pt}\begin{tabular}[t]{@{}c@{}}\rule{0pt}{12pt}8m\,13.1s\\[-.2pt]{\fontsize{5.8}{6.4}\selectfont \textcolor{BrickRed}{$ \blacktriangledown\,8.8\times $}}\\[-.2pt]2m\,54.5s\\[-.2pt]{\fontsize{5.8}{6.4}\selectfont \textcolor{BrickRed}{$ \blacktriangledown\,3.2\times $}}\end{tabular} & \hspace*{1pt}\begin{tabular}[t]{@{}c@{}}\rule{0pt}{12pt}1m\,54.3s\\[-.2pt]{\fontsize{5.8}{6.4}\selectfont \textcolor{BrickRed}{$ \blacktriangledown\,2.6\times $}}\\[-.2pt]1m\,45.5s\\[-.2pt]{\fontsize{5.8}{6.4}\selectfont \textcolor{BrickRed}{$ \blacktriangledown\,2.5\times $}}\end{tabular} & \hspace*{1pt}\begin{tabular}[t]{@{}c@{}}\rule{0pt}{12pt}5m\,17.3s\\[-.2pt]{\fontsize{5.8}{6.4}\selectfont \textcolor{BrickRed}{$ \blacktriangledown\,8.2\times $}}\\[-.2pt]1m\,54.9s\\[-.2pt]{\fontsize{5.8}{6.4}\selectfont \textcolor{BrickRed}{$ \blacktriangledown\,3.0\times $}}\end{tabular} & \hspace*{1pt}\begin{tabular}[t]{@{}c@{}}\rule{0pt}{12pt}1m\,28.3s\\[-.2pt]{\fontsize{5.8}{6.4}\selectfont \textcolor{BrickRed}{$ \blacktriangledown\,1.7\times $}}\\[-.2pt]1m\,16.8s\\[-.2pt]{\fontsize{5.8}{6.4}\selectfont \textcolor{BrickRed}{$ \blacktriangledown\,1.6\times $}}\end{tabular} & \hspace*{1pt}\begin{tabular}[t]{@{}c@{}}\rule{0pt}{12pt}4m\,45.7s\\[-.2pt]{\fontsize{5.8}{6.4}\selectfont \textcolor{BrickRed}{$ \blacktriangledown\,9.9\times $}}\\[-.2pt]1m\,26.6s\\[-.2pt]{\fontsize{5.8}{6.4}\selectfont \textcolor{BrickRed}{$ \blacktriangledown\,3.1\times $}}\end{tabular}\\[3pt]\hline
\BenchmarkSourceLabel{bench:local:row:mpf-bounded-scalar-ksw2}{R3}\shortstack[l]{{\color{gray}R3}} & \begin{tabular}[t]{@{}l@{}}\rule{0pt}{12pt}MPF / Bounded\\Scalar-DP / KSW2\end{tabular} & \textcolor{gray}{\ensuremath{\cdot}} & \hspace*{1pt}\begin{tabular}[t]{@{}c@{}}\rule{0pt}{12pt}36.5s\\[-.2pt]{\fontsize{5.8}{6.4}\selectfont \textcolor{BrickRed}{$ \blacktriangledown\,1.4\times $}}\\[-.2pt]34.8s\\[-.2pt]{\fontsize{5.8}{6.4}\selectfont \textcolor{BrickRed}{$ \blacktriangledown\,1.5\times $}}\end{tabular} & \hspace*{1pt}\begin{tabular}[t]{@{}c@{}}\rule{0pt}{12pt}27.1s\\[-.2pt]{\fontsize{5.8}{6.4}\selectfont \textcolor{ForestGreen}{$ \blacktriangle\,4.3\times $}}\\[-.2pt]25.6s\\[-.2pt]{\fontsize{5.8}{6.4}\selectfont \textcolor{ForestGreen}{$ \blacktriangle\,4.4\times $}}\end{tabular} & \hspace*{1pt}\begin{tabular}[t]{@{}c@{}}\rule{0pt}{12pt}30.3s\\[-.2pt]{\fontsize{5.8}{6.4}\selectfont \textcolor{BrickRed}{$ \blacktriangledown\,22\% $}}\\[-.2pt]28.7s\\[-.2pt]{\fontsize{5.8}{6.4}\selectfont \textcolor{BrickRed}{$ \blacktriangledown\,29\% $}}\end{tabular} & \textcolor{gray}{\ensuremath{\cdot}} & \hspace*{1pt}\begin{tabular}[t]{@{}c@{}}\rule{0pt}{12pt}5.8s\\[-.2pt]{\fontsize{5.8}{6.4}\selectfont \textcolor{ForestGreen}{$ \blacktriangle\,2.6\times $}}\\[-.2pt]4.2s\\[-.2pt]{\fontsize{5.8}{6.4}\selectfont \textcolor{ForestGreen}{$ \blacktriangle\,3.1\times $}}\end{tabular} & \hspace*{1pt}\begin{tabular}[t]{@{}c@{}}\rule{0pt}{12pt}2.4s\\[-.2pt]{\fontsize{5.8}{6.4}\selectfont \textcolor{ForestGreen}{$ \blacktriangle\,12.3\times $}}\\[-.2pt]\colorbox{BenchBest}{\strut 1.1s}\\[-.2pt]{\fontsize{5.8}{6.4}\selectfont \textcolor{ForestGreen}{$ \blacktriangle\,22.9\times $}}\end{tabular} & \hspace*{1pt}\begin{tabular}[t]{@{}c@{}}\rule{0pt}{12pt}5.4s\\[-.2pt]{\fontsize{5.8}{6.4}\selectfont \textcolor{ForestGreen}{$ \blacktriangle\,1.6\times $}}\\[-.2pt]3.8s\\[-.2pt]{\fontsize{5.8}{6.4}\selectfont \textcolor{ForestGreen}{$ \blacktriangle\,2.2\times $}}\end{tabular} & \textcolor{gray}{\ensuremath{\cdot}} & \textcolor{gray}{\ensuremath{\cdot}} & \textcolor{gray}{\ensuremath{\cdot}} & \textcolor{gray}{\ensuremath{\cdot}} & \textcolor{gray}{\ensuremath{\cdot}} & \textcolor{gray}{\ensuremath{\cdot}} & \textcolor{gray}{\ensuremath{\cdot}} & \textcolor{gray}{\ensuremath{\cdot}}\\[3pt]\hline
\BenchmarkSourceLabel{bench:local:row:mpf-bounded-scalar-block}{R4}\shortstack[l]{{\color{gray}R4}} & \begin{tabular}[t]{@{}l@{}}\rule{0pt}{12pt}MPF / Bounded\\Scalar-DP / Block\end{tabular} & \textcolor{gray}{\ensuremath{\cdot}} & \hspace*{1pt}\begin{tabular}[t]{@{}c@{}}\rule{0pt}{12pt}34.8s\\[-.2pt]{\fontsize{5.8}{6.4}\selectfont \textcolor{BrickRed}{$ \blacktriangledown\,1.4\times $}}\\[-.2pt]33.5s\\[-.2pt]{\fontsize{5.8}{6.4}\selectfont \textcolor{BrickRed}{$ \blacktriangledown\,1.5\times $}}\end{tabular} & \hspace*{1pt}\begin{tabular}[t]{@{}c@{}}\rule{0pt}{12pt}26s\\[-.2pt]{\fontsize{5.8}{6.4}\selectfont \textcolor{ForestGreen}{$ \blacktriangle\,4.5\times $}}\\[-.2pt]25.4s\\[-.2pt]{\fontsize{5.8}{6.4}\selectfont \textcolor{ForestGreen}{$ \blacktriangle\,4.5\times $}}\end{tabular} & \hspace*{1pt}\begin{tabular}[t]{@{}c@{}}\rule{0pt}{12pt}30.7s\\[-.2pt]{\fontsize{5.8}{6.4}\selectfont \textcolor{BrickRed}{$ \blacktriangledown\,23\% $}}\\[-.2pt]29.5s\\[-.2pt]{\fontsize{5.8}{6.4}\selectfont \textcolor{BrickRed}{$ \blacktriangledown\,1.3\times $}}\end{tabular} & \textcolor{gray}{\ensuremath{\cdot}} & \hspace*{1pt}\begin{tabular}[t]{@{}c@{}}\rule{0pt}{12pt}5.4s\\[-.2pt]{\fontsize{5.8}{6.4}\selectfont \textcolor{ForestGreen}{$ \blacktriangle\,2.8\times $}}\\[-.2pt]4.3s\\[-.2pt]{\fontsize{5.8}{6.4}\selectfont \textcolor{ForestGreen}{$ \blacktriangle\,3.0\times $}}\end{tabular} & \hspace*{1pt}\begin{tabular}[t]{@{}c@{}}\rule{0pt}{12pt}2.4s\\[-.2pt]{\fontsize{5.8}{6.4}\selectfont \textcolor{ForestGreen}{$ \blacktriangle\,12.0\times $}}\\[-.2pt]1.9s\\[-.2pt]{\fontsize{5.8}{6.4}\selectfont \textcolor{ForestGreen}{$ \blacktriangle\,13.8\times $}}\end{tabular} & \hspace*{1pt}\begin{tabular}[t]{@{}c@{}}\rule{0pt}{12pt}5.1s\\[-.2pt]{\fontsize{5.8}{6.4}\selectfont \textcolor{ForestGreen}{$ \blacktriangle\,1.7\times $}}\\[-.2pt]3.9s\\[-.2pt]{\fontsize{5.8}{6.4}\selectfont \textcolor{ForestGreen}{$ \blacktriangle\,2.1\times $}}\end{tabular} & \textcolor{gray}{\ensuremath{\cdot}} & \textcolor{gray}{\ensuremath{\cdot}} & \textcolor{gray}{\ensuremath{\cdot}} & \textcolor{gray}{\ensuremath{\cdot}} & \textcolor{gray}{\ensuremath{\cdot}} & \textcolor{gray}{\ensuremath{\cdot}} & \textcolor{gray}{\ensuremath{\cdot}} & \textcolor{gray}{\ensuremath{\cdot}}\\[3pt]\hline
\BenchmarkSourceLabel{bench:local:row:mpf-bounded-simd-ksw2}{R5}\shortstack[l]{{\color{gray}R5}} & \begin{tabular}[t]{@{}l@{}}\rule{0pt}{12pt}MPF / Bounded\\SIMD-DP / KSW2\end{tabular} & \hspace*{1pt}\begin{tabular}[t]{@{}c@{}}\rule{0pt}{12pt}20.8s\\[-.2pt]{\fontsize{5.8}{6.4}\selectfont \textcolor{ForestGreen}{$ \blacktriangle\,1.9\times $}}\\[-.2pt]19.1s\\[-.2pt]{\fontsize{5.8}{6.4}\selectfont \textcolor{ForestGreen}{$ \blacktriangle\,1.9\times $}}\end{tabular} & \hspace*{1pt}\begin{tabular}[t]{@{}c@{}}\rule{0pt}{12pt}\colorbox{BenchBest}{\strut 16.9s}\\[-.2pt]{\fontsize{5.8}{6.4}\selectfont \textcolor{ForestGreen}{$ \blacktriangle\,1.5\times $}}\\[-.2pt]\colorbox{BenchBest}{\strut 15.3s}\\[-.2pt]{\fontsize{5.8}{6.4}\selectfont \textcolor{ForestGreen}{$ \blacktriangle\,1.5\times $}}\end{tabular} & \hspace*{1pt}\begin{tabular}[t]{@{}c@{}}\rule{0pt}{12pt}19.8s\\[-.2pt]{\fontsize{5.8}{6.4}\selectfont \textcolor{ForestGreen}{$ \blacktriangle\,5.9\times $}}\\[-.2pt]18.2s\\[-.2pt]{\fontsize{5.8}{6.4}\selectfont \textcolor{ForestGreen}{$ \blacktriangle\,6.2\times $}}\end{tabular} & \hspace*{1pt}\begin{tabular}[t]{@{}c@{}}\rule{0pt}{12pt}30.1s\\[-.2pt]{\fontsize{5.8}{6.4}\selectfont \textcolor{BrickRed}{$ \blacktriangledown\,21\% $}}\\[-.2pt]28.5s\\[-.2pt]{\fontsize{5.8}{6.4}\selectfont \textcolor{BrickRed}{$ \blacktriangledown\,28\% $}}\end{tabular} & \hspace*{1pt}\begin{tabular}[t]{@{}c@{}}\rule{0pt}{12pt}3.7s\\[-.2pt]{\fontsize{5.8}{6.4}\selectfont \textcolor{ForestGreen}{$ \blacktriangle\,7.8\times $}}\\[-.2pt]2.4s\\[-.2pt]{\fontsize{5.8}{6.4}\selectfont \textcolor{ForestGreen}{$ \blacktriangle\,10.7\times $}}\end{tabular} & \hspace*{1pt}\begin{tabular}[t]{@{}c@{}}\rule{0pt}{12pt}6s\\[-.2pt]{\fontsize{5.8}{6.4}\selectfont \textcolor{ForestGreen}{$ \blacktriangle\,2.5\times $}}\\[-.2pt]4.5s\\[-.2pt]{\fontsize{5.8}{6.4}\selectfont \textcolor{ForestGreen}{$ \blacktriangle\,2.9\times $}}\end{tabular} & \hspace*{1pt}\begin{tabular}[t]{@{}c@{}}\rule{0pt}{12pt}2.5s\\[-.2pt]{\fontsize{5.8}{6.4}\selectfont \textcolor{ForestGreen}{$ \blacktriangle\,11.7\times $}}\\[-.2pt]1.2s\\[-.2pt]{\fontsize{5.8}{6.4}\selectfont \textcolor{ForestGreen}{$ \blacktriangle\,22.6\times $}}\end{tabular} & \hspace*{1pt}\begin{tabular}[t]{@{}c@{}}\rule{0pt}{12pt}\colorbox{BenchBest}{\strut 5s}\\[-.2pt]{\fontsize{5.8}{6.4}\selectfont \textcolor{ForestGreen}{$ \blacktriangle\,1.7\times $}}\\[-.2pt]\colorbox{BenchBest}{\strut 3.6s}\\[-.2pt]{\fontsize{5.8}{6.4}\selectfont \textcolor{ForestGreen}{$ \blacktriangle\,2.3\times $}}\end{tabular} & \hspace*{1pt}\begin{tabular}[t]{@{}c@{}}\rule{0pt}{12pt}3m\,32.7s\\[-.2pt]{\fontsize{5.8}{6.4}\selectfont \textcolor{BrickRed}{$ \blacktriangledown\,23\% $}}\\[-.2pt]3m\,17.7s\\[-.2pt]{\fontsize{5.8}{6.4}\selectfont \textcolor{BrickRed}{$ \blacktriangledown\,17\% $}}\end{tabular} & \textcolor{gray}{\ensuremath{\cdot}} & \hspace*{1pt}\begin{tabular}[t]{@{}c@{}}\rule{0pt}{12pt}\colorbox{BenchBest}{\strut 2m\,44.9s}\\[-.2pt]{\fontsize{5.8}{6.4}\selectfont \textcolor{ForestGreen}{$ \blacktriangle\,6\% $}}\\[-.2pt]\colorbox{BenchBest}{\strut 2m\,32.2s}\\[-.2pt]{\fontsize{5.8}{6.4}\selectfont \textcolor{ForestGreen}{$ \blacktriangle\,11\% $}}\end{tabular} & \textcolor{gray}{\ensuremath{\cdot}} & \hspace*{1pt}\begin{tabular}[t]{@{}c@{}}\rule{0pt}{12pt}1m\,52.7s\\[-.2pt]{\fontsize{5.8}{6.4}\selectfont \textcolor{BrickRed}{$ \blacktriangledown\,2.6\times $}}\\[-.2pt]1m\,40s\\[-.2pt]{\fontsize{5.8}{6.4}\selectfont \textcolor{BrickRed}{$ \blacktriangledown\,2.4\times $}}\end{tabular} & \hspace*{1pt}\begin{tabular}[t]{@{}c@{}}\rule{0pt}{12pt}5m\,12.1s\\[-.2pt]{\fontsize{5.8}{6.4}\selectfont \textcolor{BrickRed}{$ \blacktriangledown\,8.0\times $}}\\[-.2pt]1m\,46.7s\\[-.2pt]{\fontsize{5.8}{6.4}\selectfont \textcolor{BrickRed}{$ \blacktriangledown\,2.8\times $}}\end{tabular} & \hspace*{1pt}\begin{tabular}[t]{@{}c@{}}\rule{0pt}{12pt}1m\,26.1s\\[-.2pt]{\fontsize{5.8}{6.4}\selectfont \textcolor{BrickRed}{$ \blacktriangledown\,1.6\times $}}\\[-.2pt]1m\,13.9s\\[-.2pt]{\fontsize{5.8}{6.4}\selectfont \textcolor{BrickRed}{$ \blacktriangledown\,1.5\times $}}\end{tabular} & \hspace*{1pt}\begin{tabular}[t]{@{}c@{}}\rule{0pt}{12pt}4m\,41.4s\\[-.2pt]{\fontsize{5.8}{6.4}\selectfont \textcolor{BrickRed}{$ \blacktriangledown\,9.8\times $}}\\[-.2pt]1m\,21.9s\\[-.2pt]{\fontsize{5.8}{6.4}\selectfont \textcolor{BrickRed}{$ \blacktriangledown\,2.9\times $}}\end{tabular}\\[3pt]\hline
\BenchmarkSourceLabel{bench:local:row:mpf-bounded-rmq-ksw2}{R6}\shortstack[l]{{\color{gray}R6}} & \begin{tabular}[t]{@{}l@{}}\rule{0pt}{12pt}MPF / Bounded\\RMQ / KSW2\end{tabular} & \textcolor{gray}{\ensuremath{\cdot}} & \hspace*{1pt}\begin{tabular}[t]{@{}c@{}}\rule{0pt}{12pt}34.8s\\[-.2pt]{\fontsize{5.8}{6.4}\selectfont \textcolor{BrickRed}{$ \blacktriangledown\,1.4\times $}}\\[-.2pt]33.2s\\[-.2pt]{\fontsize{5.8}{6.4}\selectfont \textcolor{BrickRed}{$ \blacktriangledown\,1.5\times $}}\end{tabular} & \hspace*{1pt}\begin{tabular}[t]{@{}c@{}}\rule{0pt}{12pt}27.4s\\[-.2pt]{\fontsize{5.8}{6.4}\selectfont \textcolor{ForestGreen}{$ \blacktriangle\,4.3\times $}}\\[-.2pt]25.9s\\[-.2pt]{\fontsize{5.8}{6.4}\selectfont \textcolor{ForestGreen}{$ \blacktriangle\,4.4\times $}}\end{tabular} & \hspace*{1pt}\begin{tabular}[t]{@{}c@{}}\rule{0pt}{12pt}31.7s\\[-.2pt]{\fontsize{5.8}{6.4}\selectfont \textcolor{BrickRed}{$ \blacktriangledown\,27\% $}}\\[-.2pt]30.1s\\[-.2pt]{\fontsize{5.8}{6.4}\selectfont \textcolor{BrickRed}{$ \blacktriangledown\,1.4\times $}}\end{tabular} & \textcolor{gray}{\ensuremath{\cdot}} & \hspace*{1pt}\begin{tabular}[t]{@{}c@{}}\rule{0pt}{12pt}5.8s\\[-.2pt]{\fontsize{5.8}{6.4}\selectfont \textcolor{ForestGreen}{$ \blacktriangle\,2.6\times $}}\\[-.2pt]\colorbox{BenchBest}{\strut 3.9s}\\[-.2pt]{\fontsize{5.8}{6.4}\selectfont \textcolor{ForestGreen}{$ \blacktriangle\,3.3\times $}}\end{tabular} & \hspace*{1pt}\begin{tabular}[t]{@{}c@{}}\rule{0pt}{12pt}2.8s\\[-.2pt]{\fontsize{5.8}{6.4}\selectfont \textcolor{ForestGreen}{$ \blacktriangle\,10.6\times $}}\\[-.2pt]1.2s\\[-.2pt]{\fontsize{5.8}{6.4}\selectfont \textcolor{ForestGreen}{$ \blacktriangle\,21.1\times $}}\end{tabular} & \hspace*{1pt}\begin{tabular}[t]{@{}c@{}}\rule{0pt}{12pt}5.4s\\[-.2pt]{\fontsize{5.8}{6.4}\selectfont \textcolor{ForestGreen}{$ \blacktriangle\,1.6\times $}}\\[-.2pt]3.7s\\[-.2pt]{\fontsize{5.8}{6.4}\selectfont \textcolor{ForestGreen}{$ \blacktriangle\,2.2\times $}}\end{tabular} & \textcolor{gray}{\ensuremath{\cdot}} & \textcolor{gray}{\ensuremath{\cdot}} & \textcolor{gray}{\ensuremath{\cdot}} & \textcolor{gray}{\ensuremath{\cdot}} & \textcolor{gray}{\ensuremath{\cdot}} & \textcolor{gray}{\ensuremath{\cdot}} & \textcolor{gray}{\ensuremath{\cdot}} & \textcolor{gray}{\ensuremath{\cdot}}\\[3pt]\hline
\BenchmarkSourceLabel{bench:local:row:mpf-bounded-rmq-block}{R7}\shortstack[l]{{\color{gray}R7}} & \begin{tabular}[t]{@{}l@{}}\rule{0pt}{12pt}MPF / Bounded\\RMQ / Block\end{tabular} & \textcolor{gray}{\ensuremath{\cdot}} & \hspace*{1pt}\begin{tabular}[t]{@{}c@{}}\rule{0pt}{12pt}34s\\[-.2pt]{\fontsize{5.8}{6.4}\selectfont \textcolor{BrickRed}{$ \blacktriangledown\,1.3\times $}}\\[-.2pt]32.7s\\[-.2pt]{\fontsize{5.8}{6.4}\selectfont \textcolor{BrickRed}{$ \blacktriangledown\,1.4\times $}}\end{tabular} & \hspace*{1pt}\begin{tabular}[t]{@{}c@{}}\rule{0pt}{12pt}\colorbox{BenchBest}{\strut 4s}\\[-.2pt]{\fontsize{5.8}{6.4}\selectfont \textcolor{ForestGreen}{$ \blacktriangle\,29.1\times $}}\\[-.2pt]\colorbox{BenchBest}{\strut 3.4s}\\[-.2pt]{\fontsize{5.8}{6.4}\selectfont \textcolor{ForestGreen}{$ \blacktriangle\,32.9\times $}}\end{tabular} & \hspace*{1pt}\begin{tabular}[t]{@{}c@{}}\rule{0pt}{12pt}\colorbox{BenchBest}{\strut 15.4s}\\[-.2pt]{\fontsize{5.8}{6.4}\selectfont \textcolor{ForestGreen}{$ \blacktriangle\,1.6\times $}}\\[-.2pt]\colorbox{BenchBest}{\strut 14.2s}\\[-.2pt]{\fontsize{5.8}{6.4}\selectfont \textcolor{ForestGreen}{$ \blacktriangle\,1.6\times $}}\end{tabular} & \textcolor{gray}{\ensuremath{\cdot}} & \hspace*{1pt}\begin{tabular}[t]{@{}c@{}}\rule{0pt}{12pt}5.9s\\[-.2pt]{\fontsize{5.8}{6.4}\selectfont \textcolor{ForestGreen}{$ \blacktriangle\,2.6\times $}}\\[-.2pt]4.6s\\[-.2pt]{\fontsize{5.8}{6.4}\selectfont \textcolor{ForestGreen}{$ \blacktriangle\,2.8\times $}}\end{tabular} & \hspace*{1pt}\begin{tabular}[t]{@{}c@{}}\rule{0pt}{12pt}9.5s\\[-.2pt]{\fontsize{5.8}{6.4}\selectfont \textcolor{ForestGreen}{$ \blacktriangle\,3.1\times $}}\\[-.2pt]9s\\[-.2pt]{\fontsize{5.8}{6.4}\selectfont \textcolor{ForestGreen}{$ \blacktriangle\,2.9\times $}}\end{tabular} & \hspace*{1pt}\begin{tabular}[t]{@{}c@{}}\rule{0pt}{12pt}14s\\[-.2pt]{\fontsize{5.8}{6.4}\selectfont \textcolor{BrickRed}{$ \blacktriangledown\,1.6\times $}}\\[-.2pt]12.7s\\[-.2pt]{\fontsize{5.8}{6.4}\selectfont \textcolor{BrickRed}{$ \blacktriangledown\,1.6\times $}}\end{tabular} & \textcolor{gray}{\ensuremath{\cdot}} & \textcolor{gray}{\ensuremath{\cdot}} & \textcolor{gray}{\ensuremath{\cdot}} & \textcolor{gray}{\ensuremath{\cdot}} & \textcolor{gray}{\ensuremath{\cdot}} & \textcolor{gray}{\ensuremath{\cdot}} & \textcolor{gray}{\ensuremath{\cdot}} & \textcolor{gray}{\ensuremath{\cdot}}\\[3pt]\hline
\end{tabular}\par}
\smallskip{\footnotesize\noindent Each cell: A and B times, $\cdot$: no measurement. ---\textsuperscript{g}: GPU unavailable. n/a: no comparison.\par
\noindent Blue: one lowest A and one lowest B measurement for each column for each configuration.\par
\noindent Gray: minimap2. Yellow: my preferred MPF / Bounded / SIMD-DP / Block configuration.\par
\noindent MPF/MPS: minimizers per fragment/sequence. SIMD-ref: naive SIMD scan. Bounded: bounded SIMD algorithm.\par
}
\clearpage
\noindent{\large\bfseries Table \thetable: (A): Local machine\hfill 2/12}\par\smallskip
{\fontsize{7.1}{8.1}\selectfont\setlength{\tabcolsep}{1pt}\setlength{\fboxsep}{2pt}\renewcommand{\arraystretch}{1.12}
\noindent\begin{tabular}{@{}p{14mm}p{27mm}|p{13.55mm}|p{13.55mm}|p{13.55mm}|p{13.55mm}|p{13.55mm}|p{13.55mm}|p{13.55mm}|p{13.55mm}|p{13.55mm}|p{13.55mm}|p{13.55mm}|p{13.55mm}|p{13.55mm}|p{13.55mm}|p{13.55mm}|p{13.55mm}|@{}}\hline
Engine & Execution & \multicolumn{8}{c|}{\textbf{Index}} & \multicolumn{8}{c|}{\textbf{No index}}\\
 &  & \multicolumn{4}{c|}{\textbf{4 Gbp}} & \multicolumn{4}{c|}{\textbf{2 Gbp}} & \multicolumn{4}{c|}{\textbf{4 Gbp}} & \multicolumn{4}{c|}{\textbf{2 Gbp}}\\
 &  & \multicolumn{2}{c|}{\textbf{k=15, w=10}} & \multicolumn{2}{c|}{\textbf{k=20, w=15}} & \multicolumn{2}{c|}{\textbf{k=15, w=10}} & \multicolumn{2}{c|}{\textbf{k=20, w=15}} & \multicolumn{2}{c|}{\textbf{k=15, w=10}} & \multicolumn{2}{c|}{\textbf{k=20, w=15}} & \multicolumn{2}{c|}{\textbf{k=15, w=10}} & \multicolumn{2}{c|}{\textbf{k=20, w=15}}\\
 &  & \multicolumn{1}{c|}{\textbf{DNA4}} & \multicolumn{1}{c|}{\textbf{DNA5}} & \multicolumn{1}{c|}{\textbf{DNA4}} & \multicolumn{1}{c|}{\textbf{DNA5}} & \multicolumn{1}{c|}{\textbf{DNA4}} & \multicolumn{1}{c|}{\textbf{DNA5}} & \multicolumn{1}{c|}{\textbf{DNA4}} & \multicolumn{1}{c|}{\textbf{DNA5}} & \multicolumn{1}{c|}{\textbf{DNA4}} & \multicolumn{1}{c|}{\textbf{DNA5}} & \multicolumn{1}{c|}{\textbf{DNA4}} & \multicolumn{1}{c|}{\textbf{DNA5}} & \multicolumn{1}{c|}{\textbf{DNA4}} & \multicolumn{1}{c|}{\textbf{DNA5}} & \multicolumn{1}{c|}{\textbf{DNA4}} & \multicolumn{1}{c|}{\textbf{DNA5}}\\
 & & {\color{gray}C1} & {\color{gray}C2} & {\color{gray}C3} & {\color{gray}C4} & {\color{gray}C5} & {\color{gray}C6} & {\color{gray}C7} & {\color{gray}C8} & {\color{gray}C9} & {\color{gray}C10} & {\color{gray}C11} & {\color{gray}C12} & {\color{gray}C13} & {\color{gray}C14} & {\color{gray}C15} & {\color{gray}C16}\\
\hline
\rowcolor{BenchBaseline}[0.5pt][0.5pt]\shortstack[l]{{\color{gray}R1}\\\textbf{minimap2}} & \textbf{Baseline} & \hspace*{1pt}\begin{tabular}[t]{@{}c@{}}\rule{0pt}{12pt}39.7s\\[-.2pt]{\fontsize{5.8}{6.4}\selectfont \textcolor{BrickRed}{$ \blacktriangledown\,2.2\times $}}\\[-.2pt]36.8s\\[-.2pt]{\fontsize{5.8}{6.4}\selectfont \textcolor{BrickRed}{$ \blacktriangledown\,2.2\times $}}\end{tabular} & \hspace*{1pt}\begin{tabular}[t]{@{}c@{}}\rule{0pt}{12pt}25.2s\\[-.2pt]{\fontsize{5.8}{6.4}\selectfont \textcolor{BrickRed}{$ \blacktriangledown\,1.5\times $}}\\[-.2pt]22.7s\\[-.2pt]{\fontsize{5.8}{6.4}\selectfont \textcolor{BrickRed}{$ \blacktriangledown\,1.5\times $}}\end{tabular} & \hspace*{1pt}\begin{tabular}[t]{@{}c@{}}\rule{0pt}{12pt}1m\,56.9s\\[-.2pt]{\fontsize{5.8}{6.4}\selectfont \textcolor{BrickRed}{$ \blacktriangledown\,29.1\times $}}\\[-.2pt]1m\,53.2s\\[-.2pt]{\fontsize{5.8}{6.4}\selectfont \textcolor{BrickRed}{$ \blacktriangledown\,32.9\times $}}\end{tabular} & \hspace*{1pt}\begin{tabular}[t]{@{}c@{}}\rule{0pt}{12pt}24.9s\\[-.2pt]{\fontsize{5.8}{6.4}\selectfont \textcolor{BrickRed}{$ \blacktriangledown\,1.6\times $}}\\[-.2pt]22.2s\\[-.2pt]{\fontsize{5.8}{6.4}\selectfont \textcolor{BrickRed}{$ \blacktriangledown\,1.6\times $}}\end{tabular} & \hspace*{1pt}\begin{tabular}[t]{@{}c@{}}\rule{0pt}{12pt}28.6s\\[-.2pt]{\fontsize{5.8}{6.4}\selectfont \textcolor{BrickRed}{$ \blacktriangledown\,11.5\times $}}\\[-.2pt]25.8s\\[-.2pt]{\fontsize{5.8}{6.4}\selectfont \textcolor{BrickRed}{$ \blacktriangledown\,15.6\times $}}\end{tabular} & \hspace*{1pt}\begin{tabular}[t]{@{}c@{}}\rule{0pt}{12pt}15.3s\\[-.2pt]{\fontsize{5.8}{6.4}\selectfont \textcolor{BrickRed}{$ \blacktriangledown\,2.8\times $}}\\[-.2pt]12.7s\\[-.2pt]{\fontsize{5.8}{6.4}\selectfont \textcolor{BrickRed}{$ \blacktriangledown\,3.3\times $}}\end{tabular} & \hspace*{1pt}\begin{tabular}[t]{@{}c@{}}\rule{0pt}{12pt}29.3s\\[-.2pt]{\fontsize{5.8}{6.4}\selectfont \textcolor{BrickRed}{$ \blacktriangledown\,15.0\times $}}\\[-.2pt]26.2s\\[-.2pt]{\fontsize{5.8}{6.4}\selectfont \textcolor{BrickRed}{$ \blacktriangledown\,22.9\times $}}\end{tabular} & \hspace*{1pt}\begin{tabular}[t]{@{}c@{}}\rule{0pt}{12pt}8.7s\\[-.2pt]{\fontsize{5.8}{6.4}\selectfont \textcolor{BrickRed}{$ \blacktriangledown\,1.7\times $}}\\[-.2pt]8.1s\\[-.2pt]{\fontsize{5.8}{6.4}\selectfont \textcolor{BrickRed}{$ \blacktriangledown\,2.3\times $}}\end{tabular} & \hspace*{1pt}\begin{tabular}[t]{@{}c@{}}\rule{0pt}{12pt}2m\,52.4s\\[-.2pt]{\fontsize{5.8}{6.4}\selectfont \textcolor{ForestGreen}{$ \blacktriangle\,19\% $}}\\[-.2pt]2m\,48.5s\\[-.2pt]{\fontsize{5.8}{6.4}\selectfont \textcolor{ForestGreen}{$ \blacktriangle\,15\% $}}\end{tabular} & \hspace*{1pt}\begin{tabular}[t]{@{}c@{}}\rule{0pt}{12pt}2m\,17.1s\\[-.2pt]{\fontsize{5.8}{6.4}\selectfont \textcolor{ForestGreen}{$ \blacktriangle\,3.9\times $}}\\[-.2pt]2m\,15.9s\\[-.2pt]{\fontsize{5.8}{6.4}\selectfont \textcolor{ForestGreen}{$ \blacktriangle\,1.7\times $}}\end{tabular} & \hspace*{1pt}\begin{tabular}[t]{@{}c@{}}\rule{0pt}{12pt}2m\,55s\\[-.2pt]{\fontsize{5.8}{6.4}\selectfont \textcolor{BrickRed}{$ \blacktriangledown\,6\% $}}\\[-.2pt]2m\,51s\\[-.2pt]{\fontsize{5.8}{6.4}\selectfont \textcolor{BrickRed}{$ \blacktriangledown\,12\% $}}\end{tabular} & \hspace*{1pt}\begin{tabular}[t]{@{}c@{}}\rule{0pt}{12pt}55.9s\\[-.2pt]{\fontsize{5.8}{6.4}\selectfont \textcolor{ForestGreen}{$ \blacktriangle\,8.6\times $}}\\[-.2pt]55.1s\\[-.2pt]{\fontsize{5.8}{6.4}\selectfont \textcolor{ForestGreen}{$ \blacktriangle\,3.1\times $}}\end{tabular} & \hspace*{1pt}\begin{tabular}[t]{@{}c@{}}\rule{0pt}{12pt}43.3s\\[-.2pt]{\fontsize{5.8}{6.4}\selectfont \textcolor{ForestGreen}{$ \blacktriangle\,2.6\times $}}\\[-.2pt]42.5s\\[-.2pt]{\fontsize{5.8}{6.4}\selectfont \textcolor{ForestGreen}{$ \blacktriangle\,2.4\times $}}\end{tabular} & \hspace*{1pt}\begin{tabular}[t]{@{}c@{}}\rule{0pt}{12pt}38.9s\\[-.2pt]{\fontsize{5.8}{6.4}\selectfont \textcolor{ForestGreen}{$ \blacktriangle\,8.0\times $}}\\[-.2pt]38.2s\\[-.2pt]{\fontsize{5.8}{6.4}\selectfont \textcolor{ForestGreen}{$ \blacktriangle\,2.8\times $}}\end{tabular} & \hspace*{1pt}\begin{tabular}[t]{@{}c@{}}\rule{0pt}{12pt}53.1s\\[-.2pt]{\fontsize{5.8}{6.4}\selectfont \textcolor{ForestGreen}{$ \blacktriangle\,1.6\times $}}\\[-.2pt]48.2s\\[-.2pt]{\fontsize{5.8}{6.4}\selectfont \textcolor{ForestGreen}{$ \blacktriangle\,1.5\times $}}\end{tabular} & \hspace*{1pt}\begin{tabular}[t]{@{}c@{}}\rule{0pt}{12pt}28.7s\\[-.2pt]{\fontsize{5.8}{6.4}\selectfont \textcolor{ForestGreen}{$ \blacktriangle\,9.8\times $}}\\[-.2pt]28.2s\\[-.2pt]{\fontsize{5.8}{6.4}\selectfont \textcolor{ForestGreen}{$ \blacktriangle\,2.9\times $}}\end{tabular}\\[3pt]\hline
\BenchmarkSourceLabel{bench:local:row:mps-bounded-scalar-block}{R8}\shortstack[l]{{\color{gray}R8}\\\textbf{BioseqDB}} & \begin{tabular}[t]{@{}l@{}}\rule{0pt}{12pt}MPS / Bounded\\Scalar-DP / Block\end{tabular} & \textcolor{gray}{\ensuremath{\cdot}} & \textcolor{gray}{\ensuremath{\cdot}} & \textcolor{gray}{\ensuremath{\cdot}} & \textcolor{gray}{\ensuremath{\cdot}} & \hspace*{1pt}\begin{tabular}[t]{@{}c@{}}\rule{0pt}{12pt}4.1s\\[-.2pt]{\fontsize{5.8}{6.4}\selectfont \textcolor{ForestGreen}{$ \blacktriangle\,7.0\times $}}\\[-.2pt]3.2s\\[-.2pt]{\fontsize{5.8}{6.4}\selectfont \textcolor{ForestGreen}{$ \blacktriangle\,8.1\times $}}\end{tabular} & \hspace*{1pt}\begin{tabular}[t]{@{}c@{}}\rule{0pt}{12pt}3m\,47.5s\\[-.2pt]{\fontsize{5.8}{6.4}\selectfont \textcolor{BrickRed}{$ \blacktriangledown\,14.9\times $}}\\[-.2pt]21.8s\\[-.2pt]{\fontsize{5.8}{6.4}\selectfont \textcolor{BrickRed}{$ \blacktriangledown\,1.7\times $}}\end{tabular} & \hspace*{1pt}\begin{tabular}[t]{@{}c@{}}\rule{0pt}{12pt}4.5s\\[-.2pt]{\fontsize{5.8}{6.4}\selectfont \textcolor{ForestGreen}{$ \blacktriangle\,6.5\times $}}\\[-.2pt]3.7s\\[-.2pt]{\fontsize{5.8}{6.4}\selectfont \textcolor{ForestGreen}{$ \blacktriangle\,7.1\times $}}\end{tabular} & \hspace*{1pt}\begin{tabular}[t]{@{}c@{}}\rule{0pt}{12pt}3m\,29.2s\\[-.2pt]{\fontsize{5.8}{6.4}\selectfont \textcolor{BrickRed}{$ \blacktriangledown\,24.1\times $}}\\[-.2pt]13.3s\\[-.2pt]{\fontsize{5.8}{6.4}\selectfont \textcolor{BrickRed}{$ \blacktriangledown\,1.6\times $}}\end{tabular} & \textcolor{gray}{\ensuremath{\cdot}} & \textcolor{gray}{\ensuremath{\cdot}} & \textcolor{gray}{\ensuremath{\cdot}} & \textcolor{gray}{\ensuremath{\cdot}} & \textcolor{gray}{\ensuremath{\cdot}} & \textcolor{gray}{\ensuremath{\cdot}} & \textcolor{gray}{\ensuremath{\cdot}} & \textcolor{gray}{\ensuremath{\cdot}}\\[3pt]\hline
\BenchmarkSourceLabel{bench:local:row:mps-bounded-rmq-block}{R9}\shortstack[l]{{\color{gray}R9}} & \begin{tabular}[t]{@{}l@{}}\rule{0pt}{12pt}MPS / Bounded\\RMQ / Block\end{tabular} & \textcolor{gray}{\ensuremath{\cdot}} & \textcolor{gray}{\ensuremath{\cdot}} & \textcolor{gray}{\ensuremath{\cdot}} & \textcolor{gray}{\ensuremath{\cdot}} & \hspace*{1pt}\begin{tabular}[t]{@{}c@{}}\rule{0pt}{12pt}4.3s\\[-.2pt]{\fontsize{5.8}{6.4}\selectfont \textcolor{ForestGreen}{$ \blacktriangle\,6.7\times $}}\\[-.2pt]3.2s\\[-.2pt]{\fontsize{5.8}{6.4}\selectfont \textcolor{ForestGreen}{$ \blacktriangle\,8.0\times $}}\end{tabular} & \hspace*{1pt}\begin{tabular}[t]{@{}c@{}}\rule{0pt}{12pt}3m\,45.3s\\[-.2pt]{\fontsize{5.8}{6.4}\selectfont \textcolor{BrickRed}{$ \blacktriangledown\,14.8\times $}}\\[-.2pt]21.8s\\[-.2pt]{\fontsize{5.8}{6.4}\selectfont \textcolor{BrickRed}{$ \blacktriangledown\,1.7\times $}}\end{tabular} & \hspace*{1pt}\begin{tabular}[t]{@{}c@{}}\rule{0pt}{12pt}4.4s\\[-.2pt]{\fontsize{5.8}{6.4}\selectfont \textcolor{ForestGreen}{$ \blacktriangle\,6.6\times $}}\\[-.2pt]3.8s\\[-.2pt]{\fontsize{5.8}{6.4}\selectfont \textcolor{ForestGreen}{$ \blacktriangle\,7.0\times $}}\end{tabular} & \hspace*{1pt}\begin{tabular}[t]{@{}c@{}}\rule{0pt}{12pt}3m\,30s\\[-.2pt]{\fontsize{5.8}{6.4}\selectfont \textcolor{BrickRed}{$ \blacktriangledown\,24.2\times $}}\\[-.2pt]13.4s\\[-.2pt]{\fontsize{5.8}{6.4}\selectfont \textcolor{BrickRed}{$ \blacktriangledown\,1.6\times $}}\end{tabular} & \textcolor{gray}{\ensuremath{\cdot}} & \textcolor{gray}{\ensuremath{\cdot}} & \textcolor{gray}{\ensuremath{\cdot}} & \textcolor{gray}{\ensuremath{\cdot}} & \textcolor{gray}{\ensuremath{\cdot}} & \textcolor{gray}{\ensuremath{\cdot}} & \textcolor{gray}{\ensuremath{\cdot}} & \textcolor{gray}{\ensuremath{\cdot}}\\[3pt]\hline
\BenchmarkSourceLabel{bench:local:row:mps-bounded-simd-block}{R10}\shortstack[l]{{\color{gray}R10}} & \begin{tabular}[t]{@{}l@{}}\rule{0pt}{12pt}MPS / Bounded\\SIMD-DP / Block\end{tabular} & \textcolor{gray}{\ensuremath{\cdot}} & \textcolor{gray}{\ensuremath{\cdot}} & \textcolor{gray}{\ensuremath{\cdot}} & \textcolor{gray}{\ensuremath{\cdot}} & \hspace*{1pt}\begin{tabular}[t]{@{}c@{}}\rule{0pt}{12pt}4.3s\\[-.2pt]{\fontsize{5.8}{6.4}\selectfont \textcolor{ForestGreen}{$ \blacktriangle\,6.7\times $}}\\[-.2pt]3.3s\\[-.2pt]{\fontsize{5.8}{6.4}\selectfont \textcolor{ForestGreen}{$ \blacktriangle\,7.9\times $}}\end{tabular} & \hspace*{1pt}\begin{tabular}[t]{@{}c@{}}\rule{0pt}{12pt}3m\,45.2s\\[-.2pt]{\fontsize{5.8}{6.4}\selectfont \textcolor{BrickRed}{$ \blacktriangledown\,14.7\times $}}\\[-.2pt]22.3s\\[-.2pt]{\fontsize{5.8}{6.4}\selectfont \textcolor{BrickRed}{$ \blacktriangledown\,1.8\times $}}\end{tabular} & \hspace*{1pt}\begin{tabular}[t]{@{}c@{}}\rule{0pt}{12pt}4.5s\\[-.2pt]{\fontsize{5.8}{6.4}\selectfont \textcolor{ForestGreen}{$ \blacktriangle\,6.6\times $}}\\[-.2pt]3.8s\\[-.2pt]{\fontsize{5.8}{6.4}\selectfont \textcolor{ForestGreen}{$ \blacktriangle\,6.9\times $}}\end{tabular} & \hspace*{1pt}\begin{tabular}[t]{@{}c@{}}\rule{0pt}{12pt}3m\,27.4s\\[-.2pt]{\fontsize{5.8}{6.4}\selectfont \textcolor{BrickRed}{$ \blacktriangledown\,23.9\times $}}\\[-.2pt]14.1s\\[-.2pt]{\fontsize{5.8}{6.4}\selectfont \textcolor{BrickRed}{$ \blacktriangledown\,1.7\times $}}\end{tabular} & \hspace*{1pt}\begin{tabular}[t]{@{}c@{}}\rule{0pt}{12pt}3m\,53.6s\\[-.2pt]{\fontsize{5.8}{6.4}\selectfont \textcolor{BrickRed}{$ \blacktriangledown\,1.4\times $}}\\[-.2pt]3m\,40.3s\\[-.2pt]{\fontsize{5.8}{6.4}\selectfont \textcolor{BrickRed}{$ \blacktriangledown\,1.3\times $}}\end{tabular} & \hspace*{1pt}\begin{tabular}[t]{@{}c@{}}\rule{0pt}{12pt}8m\,50.7s\\[-.2pt]{\fontsize{5.8}{6.4}\selectfont \textcolor{BrickRed}{$ \blacktriangledown\,3.9\times $}}\\[-.2pt]3m\,46s\\[-.2pt]{\fontsize{5.8}{6.4}\selectfont \textcolor{BrickRed}{$ \blacktriangledown\,1.7\times $}}\end{tabular} & \hspace*{1pt}\begin{tabular}[t]{@{}c@{}}\rule{0pt}{12pt}2m\,51.3s\\[-.2pt]{\fontsize{5.8}{6.4}\selectfont \textcolor{gray}{$ \blacktriangleright\,2\% $}}\\[-.2pt]2m\,38.1s\\[-.2pt]{\fontsize{5.8}{6.4}\selectfont \textcolor{ForestGreen}{$ \blacktriangle\,8\% $}}\end{tabular} & \hspace*{1pt}\begin{tabular}[t]{@{}c@{}}\rule{0pt}{12pt}7m\,58.5s\\[-.2pt]{\fontsize{5.8}{6.4}\selectfont \textcolor{BrickRed}{$ \blacktriangledown\,8.6\times $}}\\[-.2pt]2m\,48.9s\\[-.2pt]{\fontsize{5.8}{6.4}\selectfont \textcolor{BrickRed}{$ \blacktriangledown\,3.1\times $}}\end{tabular} & \hspace*{1pt}\begin{tabular}[t]{@{}c@{}}\rule{0pt}{12pt}1m\,59.6s\\[-.2pt]{\fontsize{5.8}{6.4}\selectfont \textcolor{BrickRed}{$ \blacktriangledown\,2.8\times $}}\\[-.2pt]1m\,49s\\[-.2pt]{\fontsize{5.8}{6.4}\selectfont \textcolor{BrickRed}{$ \blacktriangledown\,2.6\times $}}\end{tabular} & \hspace*{1pt}\begin{tabular}[t]{@{}c@{}}\rule{0pt}{12pt}5m\,15.7s\\[-.2pt]{\fontsize{5.8}{6.4}\selectfont \textcolor{BrickRed}{$ \blacktriangledown\,8.1\times $}}\\[-.2pt]1m\,54.7s\\[-.2pt]{\fontsize{5.8}{6.4}\selectfont \textcolor{BrickRed}{$ \blacktriangledown\,3.0\times $}}\end{tabular} & \hspace*{1pt}\begin{tabular}[t]{@{}c@{}}\rule{0pt}{12pt}1m\,28s\\[-.2pt]{\fontsize{5.8}{6.4}\selectfont \textcolor{BrickRed}{$ \blacktriangledown\,1.7\times $}}\\[-.2pt]1m\,18.3s\\[-.2pt]{\fontsize{5.8}{6.4}\selectfont \textcolor{BrickRed}{$ \blacktriangledown\,1.6\times $}}\end{tabular} & \hspace*{1pt}\begin{tabular}[t]{@{}c@{}}\rule{0pt}{12pt}4m\,46.9s\\[-.2pt]{\fontsize{5.8}{6.4}\selectfont \textcolor{BrickRed}{$ \blacktriangledown\,10.0\times $}}\\[-.2pt]1m\,25.3s\\[-.2pt]{\fontsize{5.8}{6.4}\selectfont \textcolor{BrickRed}{$ \blacktriangledown\,3.0\times $}}\end{tabular}\\[3pt]\hline
\BenchmarkSourceLabel{bench:local:row:mps-twostack-scalar-block}{R11}\shortstack[l]{{\color{gray}R11}} & \begin{tabular}[t]{@{}l@{}}\rule{0pt}{12pt}MPS / Two-stack\\Scalar-DP / Block\end{tabular} & \textcolor{gray}{\ensuremath{\cdot}} & \textcolor{gray}{\ensuremath{\cdot}} & \textcolor{gray}{\ensuremath{\cdot}} & \textcolor{gray}{\ensuremath{\cdot}} & \hspace*{1pt}\begin{tabular}[t]{@{}c@{}}\rule{0pt}{12pt}4.3s\\[-.2pt]{\fontsize{5.8}{6.4}\selectfont \textcolor{ForestGreen}{$ \blacktriangle\,6.6\times $}}\\[-.2pt]3.4s\\[-.2pt]{\fontsize{5.8}{6.4}\selectfont \textcolor{ForestGreen}{$ \blacktriangle\,7.7\times $}}\end{tabular} & \hspace*{1pt}\begin{tabular}[t]{@{}c@{}}\rule{0pt}{12pt}3m\,32.4s\\[-.2pt]{\fontsize{5.8}{6.4}\selectfont \textcolor{BrickRed}{$ \blacktriangledown\,13.9\times $}}\\[-.2pt]16.6s\\[-.2pt]{\fontsize{5.8}{6.4}\selectfont \textcolor{BrickRed}{$ \blacktriangledown\,1.3\times $}}\end{tabular} & \hspace*{1pt}\begin{tabular}[t]{@{}c@{}}\rule{0pt}{12pt}11.8s\\[-.2pt]{\fontsize{5.8}{6.4}\selectfont \textcolor{ForestGreen}{$ \blacktriangle\,2.5\times $}}\\[-.2pt]4.5s\\[-.2pt]{\fontsize{5.8}{6.4}\selectfont \textcolor{ForestGreen}{$ \blacktriangle\,5.8\times $}}\end{tabular} & \hspace*{1pt}\begin{tabular}[t]{@{}c@{}}\rule{0pt}{12pt}3m\,21.7s\\[-.2pt]{\fontsize{5.8}{6.4}\selectfont \textcolor{BrickRed}{$ \blacktriangledown\,23.2\times $}}\\[-.2pt]8.6s\\[-.2pt]{\fontsize{5.8}{6.4}\selectfont \textcolor{BrickRed}{$ \blacktriangledown\,5\% $}}\end{tabular} & \textcolor{gray}{\ensuremath{\cdot}} & \textcolor{gray}{\ensuremath{\cdot}} & \textcolor{gray}{\ensuremath{\cdot}} & \textcolor{gray}{\ensuremath{\cdot}} & \textcolor{gray}{\ensuremath{\cdot}} & \textcolor{gray}{\ensuremath{\cdot}} & \textcolor{gray}{\ensuremath{\cdot}} & \textcolor{gray}{\ensuremath{\cdot}}\\[3pt]\hline
\BenchmarkSourceLabel{bench:local:row:mps-twostack-rmq-block}{R12}\shortstack[l]{{\color{gray}R12}} & \begin{tabular}[t]{@{}l@{}}\rule{0pt}{12pt}MPS / Two-stack\\RMQ / Block\end{tabular} & \textcolor{gray}{\ensuremath{\cdot}} & \textcolor{gray}{\ensuremath{\cdot}} & \textcolor{gray}{\ensuremath{\cdot}} & \textcolor{gray}{\ensuremath{\cdot}} & \hspace*{1pt}\begin{tabular}[t]{@{}c@{}}\rule{0pt}{12pt}4.3s\\[-.2pt]{\fontsize{5.8}{6.4}\selectfont \textcolor{ForestGreen}{$ \blacktriangle\,6.6\times $}}\\[-.2pt]3.4s\\[-.2pt]{\fontsize{5.8}{6.4}\selectfont \textcolor{ForestGreen}{$ \blacktriangle\,7.7\times $}}\end{tabular} & \hspace*{1pt}\begin{tabular}[t]{@{}c@{}}\rule{0pt}{12pt}3m\,35.1s\\[-.2pt]{\fontsize{5.8}{6.4}\selectfont \textcolor{BrickRed}{$ \blacktriangledown\,14.1\times $}}\\[-.2pt]17.4s\\[-.2pt]{\fontsize{5.8}{6.4}\selectfont \textcolor{BrickRed}{$ \blacktriangledown\,1.4\times $}}\end{tabular} & \hspace*{1pt}\begin{tabular}[t]{@{}c@{}}\rule{0pt}{12pt}11.8s\\[-.2pt]{\fontsize{5.8}{6.4}\selectfont \textcolor{ForestGreen}{$ \blacktriangle\,2.5\times $}}\\[-.2pt]4.6s\\[-.2pt]{\fontsize{5.8}{6.4}\selectfont \textcolor{ForestGreen}{$ \blacktriangle\,5.7\times $}}\end{tabular} & \hspace*{1pt}\begin{tabular}[t]{@{}c@{}}\rule{0pt}{12pt}3m\,22s\\[-.2pt]{\fontsize{5.8}{6.4}\selectfont \textcolor{BrickRed}{$ \blacktriangledown\,23.3\times $}}\\[-.2pt]8.6s\\[-.2pt]{\fontsize{5.8}{6.4}\selectfont \textcolor{BrickRed}{$ \blacktriangledown\,6\% $}}\end{tabular} & \textcolor{gray}{\ensuremath{\cdot}} & \textcolor{gray}{\ensuremath{\cdot}} & \textcolor{gray}{\ensuremath{\cdot}} & \textcolor{gray}{\ensuremath{\cdot}} & \textcolor{gray}{\ensuremath{\cdot}} & \textcolor{gray}{\ensuremath{\cdot}} & \textcolor{gray}{\ensuremath{\cdot}} & \textcolor{gray}{\ensuremath{\cdot}}\\[3pt]\hline
\BenchmarkSourceLabel{bench:local:row:mps-simdref-rmq-block}{R13}\shortstack[l]{{\color{gray}R13}} & \begin{tabular}[t]{@{}l@{}}\rule{0pt}{12pt}MPS / SIMD-ref\\RMQ / Block\end{tabular} & \textcolor{gray}{\ensuremath{\cdot}} & \textcolor{gray}{\ensuremath{\cdot}} & \textcolor{gray}{\ensuremath{\cdot}} & \textcolor{gray}{\ensuremath{\cdot}} & \hspace*{1pt}\begin{tabular}[t]{@{}c@{}}\rule{0pt}{12pt}4.3s\\[-.2pt]{\fontsize{5.8}{6.4}\selectfont \textcolor{ForestGreen}{$ \blacktriangle\,6.6\times $}}\\[-.2pt]3.5s\\[-.2pt]{\fontsize{5.8}{6.4}\selectfont \textcolor{ForestGreen}{$ \blacktriangle\,7.4\times $}}\end{tabular} & \hspace*{1pt}\begin{tabular}[t]{@{}c@{}}\rule{0pt}{12pt}3m\,51.3s\\[-.2pt]{\fontsize{5.8}{6.4}\selectfont \textcolor{BrickRed}{$ \blacktriangledown\,15.1\times $}}\\[-.2pt]27.3s\\[-.2pt]{\fontsize{5.8}{6.4}\selectfont \textcolor{BrickRed}{$ \blacktriangledown\,2.1\times $}}\end{tabular} & \hspace*{1pt}\begin{tabular}[t]{@{}c@{}}\rule{0pt}{12pt}12.9s\\[-.2pt]{\fontsize{5.8}{6.4}\selectfont \textcolor{ForestGreen}{$ \blacktriangle\,2.3\times $}}\\[-.2pt]5.8s\\[-.2pt]{\fontsize{5.8}{6.4}\selectfont \textcolor{ForestGreen}{$ \blacktriangle\,4.5\times $}}\end{tabular} & \hspace*{1pt}\begin{tabular}[t]{@{}c@{}}\rule{0pt}{12pt}3m\,38.6s\\[-.2pt]{\fontsize{5.8}{6.4}\selectfont \textcolor{BrickRed}{$ \blacktriangledown\,25.2\times $}}\\[-.2pt]20.6s\\[-.2pt]{\fontsize{5.8}{6.4}\selectfont \textcolor{BrickRed}{$ \blacktriangledown\,2.5\times $}}\end{tabular} & \textcolor{gray}{\ensuremath{\cdot}} & \textcolor{gray}{\ensuremath{\cdot}} & \textcolor{gray}{\ensuremath{\cdot}} & \textcolor{gray}{\ensuremath{\cdot}} & \textcolor{gray}{\ensuremath{\cdot}} & \textcolor{gray}{\ensuremath{\cdot}} & \textcolor{gray}{\ensuremath{\cdot}} & \textcolor{gray}{\ensuremath{\cdot}}\\[3pt]\hline
\BenchmarkSourceLabel{bench:local:row:mps-simdref-scalar-block}{R14}\shortstack[l]{{\color{gray}R14}} & \begin{tabular}[t]{@{}l@{}}\rule{0pt}{12pt}MPS / SIMD-ref\\Scalar-DP / Block\end{tabular} & \textcolor{gray}{\ensuremath{\cdot}} & \textcolor{gray}{\ensuremath{\cdot}} & \textcolor{gray}{\ensuremath{\cdot}} & \textcolor{gray}{\ensuremath{\cdot}} & \hspace*{1pt}\begin{tabular}[t]{@{}c@{}}\rule{0pt}{12pt}4.4s\\[-.2pt]{\fontsize{5.8}{6.4}\selectfont \textcolor{ForestGreen}{$ \blacktriangle\,6.6\times $}}\\[-.2pt]3.5s\\[-.2pt]{\fontsize{5.8}{6.4}\selectfont \textcolor{ForestGreen}{$ \blacktriangle\,7.5\times $}}\end{tabular} & \hspace*{1pt}\begin{tabular}[t]{@{}c@{}}\rule{0pt}{12pt}3m\,50.5s\\[-.2pt]{\fontsize{5.8}{6.4}\selectfont \textcolor{BrickRed}{$ \blacktriangledown\,15.1\times $}}\\[-.2pt]26.1s\\[-.2pt]{\fontsize{5.8}{6.4}\selectfont \textcolor{BrickRed}{$ \blacktriangledown\,2.1\times $}}\end{tabular} & \hspace*{1pt}\begin{tabular}[t]{@{}c@{}}\rule{0pt}{12pt}12.5s\\[-.2pt]{\fontsize{5.8}{6.4}\selectfont \textcolor{ForestGreen}{$ \blacktriangle\,2.3\times $}}\\[-.2pt]5.6s\\[-.2pt]{\fontsize{5.8}{6.4}\selectfont \textcolor{ForestGreen}{$ \blacktriangle\,4.7\times $}}\end{tabular} & \hspace*{1pt}\begin{tabular}[t]{@{}c@{}}\rule{0pt}{12pt}3m\,39.6s\\[-.2pt]{\fontsize{5.8}{6.4}\selectfont \textcolor{BrickRed}{$ \blacktriangledown\,25.3\times $}}\\[-.2pt]20.2s\\[-.2pt]{\fontsize{5.8}{6.4}\selectfont \textcolor{BrickRed}{$ \blacktriangledown\,2.5\times $}}\end{tabular} & \textcolor{gray}{\ensuremath{\cdot}} & \textcolor{gray}{\ensuremath{\cdot}} & \textcolor{gray}{\ensuremath{\cdot}} & \textcolor{gray}{\ensuremath{\cdot}} & \textcolor{gray}{\ensuremath{\cdot}} & \textcolor{gray}{\ensuremath{\cdot}} & \textcolor{gray}{\ensuremath{\cdot}} & \textcolor{gray}{\ensuremath{\cdot}}\\[3pt]\hline
\end{tabular}\par}
\smallskip{\footnotesize\noindent Each cell: A and B times, $\cdot$: no measurement. ---\textsuperscript{g}: GPU unavailable. n/a: no comparison.\par
\noindent Blue: one lowest A and one lowest B measurement for each column for each configuration.\par
\noindent Gray: minimap2. Yellow: my preferred MPF / Bounded / SIMD-DP / Block configuration.\par
\noindent MPF/MPS: minimizers per fragment/sequence. SIMD-ref: naive SIMD scan. Bounded: bounded SIMD algorithm.\par
}
\clearpage
\noindent{\large\bfseries Table \thetable: (A): Local machine\hfill 3/12}\par\smallskip
{\fontsize{7.1}{8.1}\selectfont\setlength{\tabcolsep}{1pt}\setlength{\fboxsep}{2pt}\renewcommand{\arraystretch}{1.12}
\noindent\begin{tabular}{@{}p{14mm}p{27mm}|p{13.55mm}|p{13.55mm}|p{13.55mm}|p{13.55mm}|p{13.55mm}|p{13.55mm}|p{13.55mm}|p{13.55mm}|p{13.55mm}|p{13.55mm}|p{13.55mm}|p{13.55mm}|p{13.55mm}|p{13.55mm}|p{13.55mm}|p{13.55mm}|@{}}\hline
Engine & Execution & \multicolumn{8}{c|}{\textbf{Index}} & \multicolumn{8}{c|}{\textbf{No index}}\\
 &  & \multicolumn{4}{c|}{\textbf{4 Gbp}} & \multicolumn{4}{c|}{\textbf{2 Gbp}} & \multicolumn{4}{c|}{\textbf{4 Gbp}} & \multicolumn{4}{c|}{\textbf{2 Gbp}}\\
 &  & \multicolumn{2}{c|}{\textbf{k=15, w=10}} & \multicolumn{2}{c|}{\textbf{k=20, w=15}} & \multicolumn{2}{c|}{\textbf{k=15, w=10}} & \multicolumn{2}{c|}{\textbf{k=20, w=15}} & \multicolumn{2}{c|}{\textbf{k=15, w=10}} & \multicolumn{2}{c|}{\textbf{k=20, w=15}} & \multicolumn{2}{c|}{\textbf{k=15, w=10}} & \multicolumn{2}{c|}{\textbf{k=20, w=15}}\\
 &  & \multicolumn{1}{c|}{\textbf{DNA4}} & \multicolumn{1}{c|}{\textbf{DNA5}} & \multicolumn{1}{c|}{\textbf{DNA4}} & \multicolumn{1}{c|}{\textbf{DNA5}} & \multicolumn{1}{c|}{\textbf{DNA4}} & \multicolumn{1}{c|}{\textbf{DNA5}} & \multicolumn{1}{c|}{\textbf{DNA4}} & \multicolumn{1}{c|}{\textbf{DNA5}} & \multicolumn{1}{c|}{\textbf{DNA4}} & \multicolumn{1}{c|}{\textbf{DNA5}} & \multicolumn{1}{c|}{\textbf{DNA4}} & \multicolumn{1}{c|}{\textbf{DNA5}} & \multicolumn{1}{c|}{\textbf{DNA4}} & \multicolumn{1}{c|}{\textbf{DNA5}} & \multicolumn{1}{c|}{\textbf{DNA4}} & \multicolumn{1}{c|}{\textbf{DNA5}}\\
 & & {\color{gray}C1} & {\color{gray}C2} & {\color{gray}C3} & {\color{gray}C4} & {\color{gray}C5} & {\color{gray}C6} & {\color{gray}C7} & {\color{gray}C8} & {\color{gray}C9} & {\color{gray}C10} & {\color{gray}C11} & {\color{gray}C12} & {\color{gray}C13} & {\color{gray}C14} & {\color{gray}C15} & {\color{gray}C16}\\
\hline
\rowcolor{BenchBaseline}[0.5pt][0.5pt]\shortstack[l]{{\color{gray}R1}\\\textbf{minimap2}} & \textbf{Baseline} & \hspace*{1pt}\begin{tabular}[t]{@{}c@{}}\rule{0pt}{12pt}39.7s\\[-.2pt]{\fontsize{5.8}{6.4}\selectfont \textcolor{BrickRed}{$ \blacktriangledown\,2.2\times $}}\\[-.2pt]36.8s\\[-.2pt]{\fontsize{5.8}{6.4}\selectfont \textcolor{BrickRed}{$ \blacktriangledown\,2.2\times $}}\end{tabular} & \hspace*{1pt}\begin{tabular}[t]{@{}c@{}}\rule{0pt}{12pt}25.2s\\[-.2pt]{\fontsize{5.8}{6.4}\selectfont \textcolor{BrickRed}{$ \blacktriangledown\,1.5\times $}}\\[-.2pt]22.7s\\[-.2pt]{\fontsize{5.8}{6.4}\selectfont \textcolor{BrickRed}{$ \blacktriangledown\,1.5\times $}}\end{tabular} & \hspace*{1pt}\begin{tabular}[t]{@{}c@{}}\rule{0pt}{12pt}1m\,56.9s\\[-.2pt]{\fontsize{5.8}{6.4}\selectfont \textcolor{BrickRed}{$ \blacktriangledown\,29.1\times $}}\\[-.2pt]1m\,53.2s\\[-.2pt]{\fontsize{5.8}{6.4}\selectfont \textcolor{BrickRed}{$ \blacktriangledown\,32.9\times $}}\end{tabular} & \hspace*{1pt}\begin{tabular}[t]{@{}c@{}}\rule{0pt}{12pt}24.9s\\[-.2pt]{\fontsize{5.8}{6.4}\selectfont \textcolor{BrickRed}{$ \blacktriangledown\,1.6\times $}}\\[-.2pt]22.2s\\[-.2pt]{\fontsize{5.8}{6.4}\selectfont \textcolor{BrickRed}{$ \blacktriangledown\,1.6\times $}}\end{tabular} & \hspace*{1pt}\begin{tabular}[t]{@{}c@{}}\rule{0pt}{12pt}28.6s\\[-.2pt]{\fontsize{5.8}{6.4}\selectfont \textcolor{BrickRed}{$ \blacktriangledown\,11.5\times $}}\\[-.2pt]25.8s\\[-.2pt]{\fontsize{5.8}{6.4}\selectfont \textcolor{BrickRed}{$ \blacktriangledown\,15.6\times $}}\end{tabular} & \hspace*{1pt}\begin{tabular}[t]{@{}c@{}}\rule{0pt}{12pt}15.3s\\[-.2pt]{\fontsize{5.8}{6.4}\selectfont \textcolor{BrickRed}{$ \blacktriangledown\,2.8\times $}}\\[-.2pt]12.7s\\[-.2pt]{\fontsize{5.8}{6.4}\selectfont \textcolor{BrickRed}{$ \blacktriangledown\,3.3\times $}}\end{tabular} & \hspace*{1pt}\begin{tabular}[t]{@{}c@{}}\rule{0pt}{12pt}29.3s\\[-.2pt]{\fontsize{5.8}{6.4}\selectfont \textcolor{BrickRed}{$ \blacktriangledown\,15.0\times $}}\\[-.2pt]26.2s\\[-.2pt]{\fontsize{5.8}{6.4}\selectfont \textcolor{BrickRed}{$ \blacktriangledown\,22.9\times $}}\end{tabular} & \hspace*{1pt}\begin{tabular}[t]{@{}c@{}}\rule{0pt}{12pt}8.7s\\[-.2pt]{\fontsize{5.8}{6.4}\selectfont \textcolor{BrickRed}{$ \blacktriangledown\,1.7\times $}}\\[-.2pt]8.1s\\[-.2pt]{\fontsize{5.8}{6.4}\selectfont \textcolor{BrickRed}{$ \blacktriangledown\,2.3\times $}}\end{tabular} & \hspace*{1pt}\begin{tabular}[t]{@{}c@{}}\rule{0pt}{12pt}2m\,52.4s\\[-.2pt]{\fontsize{5.8}{6.4}\selectfont \textcolor{ForestGreen}{$ \blacktriangle\,19\% $}}\\[-.2pt]2m\,48.5s\\[-.2pt]{\fontsize{5.8}{6.4}\selectfont \textcolor{ForestGreen}{$ \blacktriangle\,15\% $}}\end{tabular} & \hspace*{1pt}\begin{tabular}[t]{@{}c@{}}\rule{0pt}{12pt}2m\,17.1s\\[-.2pt]{\fontsize{5.8}{6.4}\selectfont \textcolor{ForestGreen}{$ \blacktriangle\,3.9\times $}}\\[-.2pt]2m\,15.9s\\[-.2pt]{\fontsize{5.8}{6.4}\selectfont \textcolor{ForestGreen}{$ \blacktriangle\,1.7\times $}}\end{tabular} & \hspace*{1pt}\begin{tabular}[t]{@{}c@{}}\rule{0pt}{12pt}2m\,55s\\[-.2pt]{\fontsize{5.8}{6.4}\selectfont \textcolor{BrickRed}{$ \blacktriangledown\,6\% $}}\\[-.2pt]2m\,51s\\[-.2pt]{\fontsize{5.8}{6.4}\selectfont \textcolor{BrickRed}{$ \blacktriangledown\,12\% $}}\end{tabular} & \hspace*{1pt}\begin{tabular}[t]{@{}c@{}}\rule{0pt}{12pt}55.9s\\[-.2pt]{\fontsize{5.8}{6.4}\selectfont \textcolor{ForestGreen}{$ \blacktriangle\,8.6\times $}}\\[-.2pt]55.1s\\[-.2pt]{\fontsize{5.8}{6.4}\selectfont \textcolor{ForestGreen}{$ \blacktriangle\,3.1\times $}}\end{tabular} & \hspace*{1pt}\begin{tabular}[t]{@{}c@{}}\rule{0pt}{12pt}43.3s\\[-.2pt]{\fontsize{5.8}{6.4}\selectfont \textcolor{ForestGreen}{$ \blacktriangle\,2.6\times $}}\\[-.2pt]42.5s\\[-.2pt]{\fontsize{5.8}{6.4}\selectfont \textcolor{ForestGreen}{$ \blacktriangle\,2.4\times $}}\end{tabular} & \hspace*{1pt}\begin{tabular}[t]{@{}c@{}}\rule{0pt}{12pt}38.9s\\[-.2pt]{\fontsize{5.8}{6.4}\selectfont \textcolor{ForestGreen}{$ \blacktriangle\,8.0\times $}}\\[-.2pt]38.2s\\[-.2pt]{\fontsize{5.8}{6.4}\selectfont \textcolor{ForestGreen}{$ \blacktriangle\,2.8\times $}}\end{tabular} & \hspace*{1pt}\begin{tabular}[t]{@{}c@{}}\rule{0pt}{12pt}53.1s\\[-.2pt]{\fontsize{5.8}{6.4}\selectfont \textcolor{ForestGreen}{$ \blacktriangle\,1.6\times $}}\\[-.2pt]48.2s\\[-.2pt]{\fontsize{5.8}{6.4}\selectfont \textcolor{ForestGreen}{$ \blacktriangle\,1.5\times $}}\end{tabular} & \hspace*{1pt}\begin{tabular}[t]{@{}c@{}}\rule{0pt}{12pt}28.7s\\[-.2pt]{\fontsize{5.8}{6.4}\selectfont \textcolor{ForestGreen}{$ \blacktriangle\,9.8\times $}}\\[-.2pt]28.2s\\[-.2pt]{\fontsize{5.8}{6.4}\selectfont \textcolor{ForestGreen}{$ \blacktriangle\,2.9\times $}}\end{tabular}\\[3pt]\hline
\BenchmarkSourceLabel{bench:local:row:mps-scalar-simd-block}{R15}\shortstack[l]{{\color{gray}R15}\\\textbf{BioseqDB}} & \begin{tabular}[t]{@{}l@{}}\rule{0pt}{12pt}MPS / Scalar\\SIMD-DP / Block\end{tabular} & \textcolor{gray}{\ensuremath{\cdot}} & \textcolor{gray}{\ensuremath{\cdot}} & \textcolor{gray}{\ensuremath{\cdot}} & \textcolor{gray}{\ensuremath{\cdot}} & \hspace*{1pt}\begin{tabular}[t]{@{}c@{}}\rule{0pt}{12pt}4.4s\\[-.2pt]{\fontsize{5.8}{6.4}\selectfont \textcolor{ForestGreen}{$ \blacktriangle\,6.6\times $}}\\[-.2pt]3.5s\\[-.2pt]{\fontsize{5.8}{6.4}\selectfont \textcolor{ForestGreen}{$ \blacktriangle\,7.5\times $}}\end{tabular} & \hspace*{1pt}\begin{tabular}[t]{@{}c@{}}\rule{0pt}{12pt}3m\,46.2s\\[-.2pt]{\fontsize{5.8}{6.4}\selectfont \textcolor{BrickRed}{$ \blacktriangledown\,14.8\times $}}\\[-.2pt]22.6s\\[-.2pt]{\fontsize{5.8}{6.4}\selectfont \textcolor{BrickRed}{$ \blacktriangledown\,1.8\times $}}\end{tabular} & \hspace*{1pt}\begin{tabular}[t]{@{}c@{}}\rule{0pt}{12pt}13.3s\\[-.2pt]{\fontsize{5.8}{6.4}\selectfont \textcolor{ForestGreen}{$ \blacktriangle\,2.2\times $}}\\[-.2pt]5.5s\\[-.2pt]{\fontsize{5.8}{6.4}\selectfont \textcolor{ForestGreen}{$ \blacktriangle\,4.7\times $}}\end{tabular} & \hspace*{1pt}\begin{tabular}[t]{@{}c@{}}\rule{0pt}{12pt}3m\,35s\\[-.2pt]{\fontsize{5.8}{6.4}\selectfont \textcolor{BrickRed}{$ \blacktriangledown\,24.8\times $}}\\[-.2pt]16.3s\\[-.2pt]{\fontsize{5.8}{6.4}\selectfont \textcolor{BrickRed}{$ \blacktriangledown\,2.0\times $}}\end{tabular} & \textcolor{gray}{\ensuremath{\cdot}} & \textcolor{gray}{\ensuremath{\cdot}} & \textcolor{gray}{\ensuremath{\cdot}} & \textcolor{gray}{\ensuremath{\cdot}} & \textcolor{gray}{\ensuremath{\cdot}} & \textcolor{gray}{\ensuremath{\cdot}} & \textcolor{gray}{\ensuremath{\cdot}} & \textcolor{gray}{\ensuremath{\cdot}}\\[3pt]\hline
\BenchmarkSourceLabel{bench:local:row:mps-simdref-simd-block}{R16}\shortstack[l]{{\color{gray}R16}} & \begin{tabular}[t]{@{}l@{}}\rule{0pt}{12pt}MPS / SIMD-ref\\SIMD-DP / Block\end{tabular} & \textcolor{gray}{\ensuremath{\cdot}} & \textcolor{gray}{\ensuremath{\cdot}} & \textcolor{gray}{\ensuremath{\cdot}} & \textcolor{gray}{\ensuremath{\cdot}} & \hspace*{1pt}\begin{tabular}[t]{@{}c@{}}\rule{0pt}{12pt}4.4s\\[-.2pt]{\fontsize{5.8}{6.4}\selectfont \textcolor{ForestGreen}{$ \blacktriangle\,6.6\times $}}\\[-.2pt]3.5s\\[-.2pt]{\fontsize{5.8}{6.4}\selectfont \textcolor{ForestGreen}{$ \blacktriangle\,7.4\times $}}\end{tabular} & \hspace*{1pt}\begin{tabular}[t]{@{}c@{}}\rule{0pt}{12pt}3m\,53.3s\\[-.2pt]{\fontsize{5.8}{6.4}\selectfont \textcolor{BrickRed}{$ \blacktriangledown\,15.3\times $}}\\[-.2pt]27s\\[-.2pt]{\fontsize{5.8}{6.4}\selectfont \textcolor{BrickRed}{$ \blacktriangledown\,2.1\times $}}\end{tabular} & \hspace*{1pt}\begin{tabular}[t]{@{}c@{}}\rule{0pt}{12pt}12.7s\\[-.2pt]{\fontsize{5.8}{6.4}\selectfont \textcolor{ForestGreen}{$ \blacktriangle\,2.3\times $}}\\[-.2pt]5.6s\\[-.2pt]{\fontsize{5.8}{6.4}\selectfont \textcolor{ForestGreen}{$ \blacktriangle\,4.7\times $}}\end{tabular} & \hspace*{1pt}\begin{tabular}[t]{@{}c@{}}\rule{0pt}{12pt}3m\,41.1s\\[-.2pt]{\fontsize{5.8}{6.4}\selectfont \textcolor{BrickRed}{$ \blacktriangledown\,25.5\times $}}\\[-.2pt]21s\\[-.2pt]{\fontsize{5.8}{6.4}\selectfont \textcolor{BrickRed}{$ \blacktriangledown\,2.6\times $}}\end{tabular} & \textcolor{gray}{\ensuremath{\cdot}} & \textcolor{gray}{\ensuremath{\cdot}} & \textcolor{gray}{\ensuremath{\cdot}} & \textcolor{gray}{\ensuremath{\cdot}} & \textcolor{gray}{\ensuremath{\cdot}} & \textcolor{gray}{\ensuremath{\cdot}} & \textcolor{gray}{\ensuremath{\cdot}} & \textcolor{gray}{\ensuremath{\cdot}}\\[3pt]\hline
\BenchmarkSourceLabel{bench:local:row:mps-twostack-simd-block}{R17}\shortstack[l]{{\color{gray}R17}} & \begin{tabular}[t]{@{}l@{}}\rule{0pt}{12pt}MPS / Two-stack\\SIMD-DP / Block\end{tabular} & \textcolor{gray}{\ensuremath{\cdot}} & \textcolor{gray}{\ensuremath{\cdot}} & \textcolor{gray}{\ensuremath{\cdot}} & \textcolor{gray}{\ensuremath{\cdot}} & \hspace*{1pt}\begin{tabular}[t]{@{}c@{}}\rule{0pt}{12pt}4.4s\\[-.2pt]{\fontsize{5.8}{6.4}\selectfont \textcolor{ForestGreen}{$ \blacktriangle\,6.5\times $}}\\[-.2pt]3.5s\\[-.2pt]{\fontsize{5.8}{6.4}\selectfont \textcolor{ForestGreen}{$ \blacktriangle\,7.5\times $}}\end{tabular} & \hspace*{1pt}\begin{tabular}[t]{@{}c@{}}\rule{0pt}{12pt}3m\,39s\\[-.2pt]{\fontsize{5.8}{6.4}\selectfont \textcolor{BrickRed}{$ \blacktriangledown\,14.3\times $}}\\[-.2pt]18.2s\\[-.2pt]{\fontsize{5.8}{6.4}\selectfont \textcolor{BrickRed}{$ \blacktriangledown\,1.4\times $}}\end{tabular} & \hspace*{1pt}\begin{tabular}[t]{@{}c@{}}\rule{0pt}{12pt}13.3s\\[-.2pt]{\fontsize{5.8}{6.4}\selectfont \textcolor{ForestGreen}{$ \blacktriangle\,2.2\times $}}\\[-.2pt]4.5s\\[-.2pt]{\fontsize{5.8}{6.4}\selectfont \textcolor{ForestGreen}{$ \blacktriangle\,5.9\times $}}\end{tabular} & \hspace*{1pt}\begin{tabular}[t]{@{}c@{}}\rule{0pt}{12pt}3m\,26s\\[-.2pt]{\fontsize{5.8}{6.4}\selectfont \textcolor{BrickRed}{$ \blacktriangledown\,23.7\times $}}\\[-.2pt]9.6s\\[-.2pt]{\fontsize{5.8}{6.4}\selectfont \textcolor{BrickRed}{$ \blacktriangledown\,18\% $}}\end{tabular} & \textcolor{gray}{\ensuremath{\cdot}} & \textcolor{gray}{\ensuremath{\cdot}} & \textcolor{gray}{\ensuremath{\cdot}} & \textcolor{gray}{\ensuremath{\cdot}} & \textcolor{gray}{\ensuremath{\cdot}} & \textcolor{gray}{\ensuremath{\cdot}} & \textcolor{gray}{\ensuremath{\cdot}} & \textcolor{gray}{\ensuremath{\cdot}}\\[3pt]\hline
\BenchmarkSourceLabel{bench:local:row:mps-scalar-scalar-block}{R18}\shortstack[l]{{\color{gray}R18}} & \begin{tabular}[t]{@{}l@{}}\rule{0pt}{12pt}MPS / Scalar\\Scalar-DP / Block\end{tabular} & \textcolor{gray}{\ensuremath{\cdot}} & \textcolor{gray}{\ensuremath{\cdot}} & \textcolor{gray}{\ensuremath{\cdot}} & \textcolor{gray}{\ensuremath{\cdot}} & \hspace*{1pt}\begin{tabular}[t]{@{}c@{}}\rule{0pt}{12pt}4.4s\\[-.2pt]{\fontsize{5.8}{6.4}\selectfont \textcolor{ForestGreen}{$ \blacktriangle\,6.5\times $}}\\[-.2pt]3.5s\\[-.2pt]{\fontsize{5.8}{6.4}\selectfont \textcolor{ForestGreen}{$ \blacktriangle\,7.5\times $}}\end{tabular} & \hspace*{1pt}\begin{tabular}[t]{@{}c@{}}\rule{0pt}{12pt}3m\,47.4s\\[-.2pt]{\fontsize{5.8}{6.4}\selectfont \textcolor{BrickRed}{$ \blacktriangledown\,14.9\times $}}\\[-.2pt]22.4s\\[-.2pt]{\fontsize{5.8}{6.4}\selectfont \textcolor{BrickRed}{$ \blacktriangledown\,1.8\times $}}\end{tabular} & \hspace*{1pt}\begin{tabular}[t]{@{}c@{}}\rule{0pt}{12pt}12.5s\\[-.2pt]{\fontsize{5.8}{6.4}\selectfont \textcolor{ForestGreen}{$ \blacktriangle\,2.4\times $}}\\[-.2pt]5.6s\\[-.2pt]{\fontsize{5.8}{6.4}\selectfont \textcolor{ForestGreen}{$ \blacktriangle\,4.7\times $}}\end{tabular} & \hspace*{1pt}\begin{tabular}[t]{@{}c@{}}\rule{0pt}{12pt}3m\,34.3s\\[-.2pt]{\fontsize{5.8}{6.4}\selectfont \textcolor{BrickRed}{$ \blacktriangledown\,24.7\times $}}\\[-.2pt]14.9s\\[-.2pt]{\fontsize{5.8}{6.4}\selectfont \textcolor{BrickRed}{$ \blacktriangledown\,1.8\times $}}\end{tabular} & \textcolor{gray}{\ensuremath{\cdot}} & \textcolor{gray}{\ensuremath{\cdot}} & \textcolor{gray}{\ensuremath{\cdot}} & \textcolor{gray}{\ensuremath{\cdot}} & \textcolor{gray}{\ensuremath{\cdot}} & \textcolor{gray}{\ensuremath{\cdot}} & \textcolor{gray}{\ensuremath{\cdot}} & \textcolor{gray}{\ensuremath{\cdot}}\\[3pt]\hline
\BenchmarkSourceLabel{bench:local:row:mps-scalar-rmq-block}{R19}\shortstack[l]{{\color{gray}R19}} & \begin{tabular}[t]{@{}l@{}}\rule{0pt}{12pt}MPS / Scalar\\RMQ / Block\end{tabular} & \textcolor{gray}{\ensuremath{\cdot}} & \textcolor{gray}{\ensuremath{\cdot}} & \textcolor{gray}{\ensuremath{\cdot}} & \textcolor{gray}{\ensuremath{\cdot}} & \hspace*{1pt}\begin{tabular}[t]{@{}c@{}}\rule{0pt}{12pt}4.4s\\[-.2pt]{\fontsize{5.8}{6.4}\selectfont \textcolor{ForestGreen}{$ \blacktriangle\,6.5\times $}}\\[-.2pt]3.5s\\[-.2pt]{\fontsize{5.8}{6.4}\selectfont \textcolor{ForestGreen}{$ \blacktriangle\,7.5\times $}}\end{tabular} & \hspace*{1pt}\begin{tabular}[t]{@{}c@{}}\rule{0pt}{12pt}3m\,43.8s\\[-.2pt]{\fontsize{5.8}{6.4}\selectfont \textcolor{BrickRed}{$ \blacktriangledown\,14.6\times $}}\\[-.2pt]21.9s\\[-.2pt]{\fontsize{5.8}{6.4}\selectfont \textcolor{BrickRed}{$ \blacktriangledown\,1.7\times $}}\end{tabular} & \hspace*{1pt}\begin{tabular}[t]{@{}c@{}}\rule{0pt}{12pt}12.9s\\[-.2pt]{\fontsize{5.8}{6.4}\selectfont \textcolor{ForestGreen}{$ \blacktriangle\,2.3\times $}}\\[-.2pt]5.8s\\[-.2pt]{\fontsize{5.8}{6.4}\selectfont \textcolor{ForestGreen}{$ \blacktriangle\,4.5\times $}}\end{tabular} & \hspace*{1pt}\begin{tabular}[t]{@{}c@{}}\rule{0pt}{12pt}3m\,35.1s\\[-.2pt]{\fontsize{5.8}{6.4}\selectfont \textcolor{BrickRed}{$ \blacktriangledown\,24.8\times $}}\\[-.2pt]15.1s\\[-.2pt]{\fontsize{5.8}{6.4}\selectfont \textcolor{BrickRed}{$ \blacktriangledown\,1.9\times $}}\end{tabular} & \textcolor{gray}{\ensuremath{\cdot}} & \textcolor{gray}{\ensuremath{\cdot}} & \textcolor{gray}{\ensuremath{\cdot}} & \textcolor{gray}{\ensuremath{\cdot}} & \textcolor{gray}{\ensuremath{\cdot}} & \textcolor{gray}{\ensuremath{\cdot}} & \textcolor{gray}{\ensuremath{\cdot}} & \textcolor{gray}{\ensuremath{\cdot}}\\[3pt]\hline
\BenchmarkSourceLabel{bench:local:row:mps-bounded-scalar-ksw2}{R20}\shortstack[l]{{\color{gray}R20}} & \begin{tabular}[t]{@{}l@{}}\rule{0pt}{12pt}MPS / Bounded\\Scalar-DP / KSW2\end{tabular} & \textcolor{gray}{\ensuremath{\cdot}} & \textcolor{gray}{\ensuremath{\cdot}} & \textcolor{gray}{\ensuremath{\cdot}} & \textcolor{gray}{\ensuremath{\cdot}} & \hspace*{1pt}\begin{tabular}[t]{@{}c@{}}\rule{0pt}{12pt}5s\\[-.2pt]{\fontsize{5.8}{6.4}\selectfont \textcolor{ForestGreen}{$ \blacktriangle\,5.7\times $}}\\[-.2pt]3.2s\\[-.2pt]{\fontsize{5.8}{6.4}\selectfont \textcolor{ForestGreen}{$ \blacktriangle\,8.1\times $}}\end{tabular} & \hspace*{1pt}\begin{tabular}[t]{@{}c@{}}\rule{0pt}{12pt}3m\,48.5s\\[-.2pt]{\fontsize{5.8}{6.4}\selectfont \textcolor{BrickRed}{$ \blacktriangledown\,15.0\times $}}\\[-.2pt]22s\\[-.2pt]{\fontsize{5.8}{6.4}\selectfont \textcolor{BrickRed}{$ \blacktriangledown\,1.7\times $}}\end{tabular} & \hspace*{1pt}\begin{tabular}[t]{@{}c@{}}\rule{0pt}{12pt}5.5s\\[-.2pt]{\fontsize{5.8}{6.4}\selectfont \textcolor{ForestGreen}{$ \blacktriangle\,5.3\times $}}\\[-.2pt]3.9s\\[-.2pt]{\fontsize{5.8}{6.4}\selectfont \textcolor{ForestGreen}{$ \blacktriangle\,6.7\times $}}\end{tabular} & \hspace*{1pt}\begin{tabular}[t]{@{}c@{}}\rule{0pt}{12pt}3m\,32s\\[-.2pt]{\fontsize{5.8}{6.4}\selectfont \textcolor{BrickRed}{$ \blacktriangledown\,24.4\times $}}\\[-.2pt]13.3s\\[-.2pt]{\fontsize{5.8}{6.4}\selectfont \textcolor{BrickRed}{$ \blacktriangledown\,1.6\times $}}\end{tabular} & \textcolor{gray}{\ensuremath{\cdot}} & \textcolor{gray}{\ensuremath{\cdot}} & \textcolor{gray}{\ensuremath{\cdot}} & \textcolor{gray}{\ensuremath{\cdot}} & \textcolor{gray}{\ensuremath{\cdot}} & \textcolor{gray}{\ensuremath{\cdot}} & \textcolor{gray}{\ensuremath{\cdot}} & \textcolor{gray}{\ensuremath{\cdot}}\\[3pt]\hline
\end{tabular}\par}
\smallskip{\footnotesize\noindent Each cell: A and B times, $\cdot$: no measurement. ---\textsuperscript{g}: GPU unavailable. n/a: no comparison.\par
\noindent Blue: one lowest A and one lowest B measurement for each column for each configuration.\par
\noindent Gray: minimap2. Yellow: my preferred MPF / Bounded / SIMD-DP / Block configuration.\par
\noindent MPF/MPS: minimizers per fragment/sequence. SIMD-ref: naive SIMD scan. Bounded: bounded SIMD algorithm.\par
}
\clearpage
\noindent{\large\bfseries Table \thetable: (A): Local machine\hfill 4/12}\par\smallskip
{\fontsize{7.1}{8.1}\selectfont\setlength{\tabcolsep}{1pt}\setlength{\fboxsep}{2pt}\renewcommand{\arraystretch}{1.12}
\noindent\begin{tabular}{@{}p{14mm}p{27mm}|p{13.55mm}|p{13.55mm}|p{13.55mm}|p{13.55mm}|p{13.55mm}|p{13.55mm}|p{13.55mm}|p{13.55mm}|p{13.55mm}|p{13.55mm}|p{13.55mm}|p{13.55mm}|p{13.55mm}|p{13.55mm}|p{13.55mm}|p{13.55mm}|@{}}\hline
Engine & Execution & \multicolumn{8}{c|}{\textbf{Index}} & \multicolumn{8}{c|}{\textbf{No index}}\\
 &  & \multicolumn{4}{c|}{\textbf{4 Gbp}} & \multicolumn{4}{c|}{\textbf{2 Gbp}} & \multicolumn{4}{c|}{\textbf{4 Gbp}} & \multicolumn{4}{c|}{\textbf{2 Gbp}}\\
 &  & \multicolumn{2}{c|}{\textbf{k=15, w=10}} & \multicolumn{2}{c|}{\textbf{k=20, w=15}} & \multicolumn{2}{c|}{\textbf{k=15, w=10}} & \multicolumn{2}{c|}{\textbf{k=20, w=15}} & \multicolumn{2}{c|}{\textbf{k=15, w=10}} & \multicolumn{2}{c|}{\textbf{k=20, w=15}} & \multicolumn{2}{c|}{\textbf{k=15, w=10}} & \multicolumn{2}{c|}{\textbf{k=20, w=15}}\\
 &  & \multicolumn{1}{c|}{\textbf{DNA4}} & \multicolumn{1}{c|}{\textbf{DNA5}} & \multicolumn{1}{c|}{\textbf{DNA4}} & \multicolumn{1}{c|}{\textbf{DNA5}} & \multicolumn{1}{c|}{\textbf{DNA4}} & \multicolumn{1}{c|}{\textbf{DNA5}} & \multicolumn{1}{c|}{\textbf{DNA4}} & \multicolumn{1}{c|}{\textbf{DNA5}} & \multicolumn{1}{c|}{\textbf{DNA4}} & \multicolumn{1}{c|}{\textbf{DNA5}} & \multicolumn{1}{c|}{\textbf{DNA4}} & \multicolumn{1}{c|}{\textbf{DNA5}} & \multicolumn{1}{c|}{\textbf{DNA4}} & \multicolumn{1}{c|}{\textbf{DNA5}} & \multicolumn{1}{c|}{\textbf{DNA4}} & \multicolumn{1}{c|}{\textbf{DNA5}}\\
 & & {\color{gray}C1} & {\color{gray}C2} & {\color{gray}C3} & {\color{gray}C4} & {\color{gray}C5} & {\color{gray}C6} & {\color{gray}C7} & {\color{gray}C8} & {\color{gray}C9} & {\color{gray}C10} & {\color{gray}C11} & {\color{gray}C12} & {\color{gray}C13} & {\color{gray}C14} & {\color{gray}C15} & {\color{gray}C16}\\
\hline
\rowcolor{BenchBaseline}[0.5pt][0.5pt]\shortstack[l]{{\color{gray}R1}\\\textbf{minimap2}} & \textbf{Baseline} & \hspace*{1pt}\begin{tabular}[t]{@{}c@{}}\rule{0pt}{12pt}39.7s\\[-.2pt]{\fontsize{5.8}{6.4}\selectfont \textcolor{BrickRed}{$ \blacktriangledown\,2.2\times $}}\\[-.2pt]36.8s\\[-.2pt]{\fontsize{5.8}{6.4}\selectfont \textcolor{BrickRed}{$ \blacktriangledown\,2.2\times $}}\end{tabular} & \hspace*{1pt}\begin{tabular}[t]{@{}c@{}}\rule{0pt}{12pt}25.2s\\[-.2pt]{\fontsize{5.8}{6.4}\selectfont \textcolor{BrickRed}{$ \blacktriangledown\,1.5\times $}}\\[-.2pt]22.7s\\[-.2pt]{\fontsize{5.8}{6.4}\selectfont \textcolor{BrickRed}{$ \blacktriangledown\,1.5\times $}}\end{tabular} & \hspace*{1pt}\begin{tabular}[t]{@{}c@{}}\rule{0pt}{12pt}1m\,56.9s\\[-.2pt]{\fontsize{5.8}{6.4}\selectfont \textcolor{BrickRed}{$ \blacktriangledown\,29.1\times $}}\\[-.2pt]1m\,53.2s\\[-.2pt]{\fontsize{5.8}{6.4}\selectfont \textcolor{BrickRed}{$ \blacktriangledown\,32.9\times $}}\end{tabular} & \hspace*{1pt}\begin{tabular}[t]{@{}c@{}}\rule{0pt}{12pt}24.9s\\[-.2pt]{\fontsize{5.8}{6.4}\selectfont \textcolor{BrickRed}{$ \blacktriangledown\,1.6\times $}}\\[-.2pt]22.2s\\[-.2pt]{\fontsize{5.8}{6.4}\selectfont \textcolor{BrickRed}{$ \blacktriangledown\,1.6\times $}}\end{tabular} & \hspace*{1pt}\begin{tabular}[t]{@{}c@{}}\rule{0pt}{12pt}28.6s\\[-.2pt]{\fontsize{5.8}{6.4}\selectfont \textcolor{BrickRed}{$ \blacktriangledown\,11.5\times $}}\\[-.2pt]25.8s\\[-.2pt]{\fontsize{5.8}{6.4}\selectfont \textcolor{BrickRed}{$ \blacktriangledown\,15.6\times $}}\end{tabular} & \hspace*{1pt}\begin{tabular}[t]{@{}c@{}}\rule{0pt}{12pt}15.3s\\[-.2pt]{\fontsize{5.8}{6.4}\selectfont \textcolor{BrickRed}{$ \blacktriangledown\,2.8\times $}}\\[-.2pt]12.7s\\[-.2pt]{\fontsize{5.8}{6.4}\selectfont \textcolor{BrickRed}{$ \blacktriangledown\,3.3\times $}}\end{tabular} & \hspace*{1pt}\begin{tabular}[t]{@{}c@{}}\rule{0pt}{12pt}29.3s\\[-.2pt]{\fontsize{5.8}{6.4}\selectfont \textcolor{BrickRed}{$ \blacktriangledown\,15.0\times $}}\\[-.2pt]26.2s\\[-.2pt]{\fontsize{5.8}{6.4}\selectfont \textcolor{BrickRed}{$ \blacktriangledown\,22.9\times $}}\end{tabular} & \hspace*{1pt}\begin{tabular}[t]{@{}c@{}}\rule{0pt}{12pt}8.7s\\[-.2pt]{\fontsize{5.8}{6.4}\selectfont \textcolor{BrickRed}{$ \blacktriangledown\,1.7\times $}}\\[-.2pt]8.1s\\[-.2pt]{\fontsize{5.8}{6.4}\selectfont \textcolor{BrickRed}{$ \blacktriangledown\,2.3\times $}}\end{tabular} & \hspace*{1pt}\begin{tabular}[t]{@{}c@{}}\rule{0pt}{12pt}2m\,52.4s\\[-.2pt]{\fontsize{5.8}{6.4}\selectfont \textcolor{ForestGreen}{$ \blacktriangle\,19\% $}}\\[-.2pt]2m\,48.5s\\[-.2pt]{\fontsize{5.8}{6.4}\selectfont \textcolor{ForestGreen}{$ \blacktriangle\,15\% $}}\end{tabular} & \hspace*{1pt}\begin{tabular}[t]{@{}c@{}}\rule{0pt}{12pt}2m\,17.1s\\[-.2pt]{\fontsize{5.8}{6.4}\selectfont \textcolor{ForestGreen}{$ \blacktriangle\,3.9\times $}}\\[-.2pt]2m\,15.9s\\[-.2pt]{\fontsize{5.8}{6.4}\selectfont \textcolor{ForestGreen}{$ \blacktriangle\,1.7\times $}}\end{tabular} & \hspace*{1pt}\begin{tabular}[t]{@{}c@{}}\rule{0pt}{12pt}2m\,55s\\[-.2pt]{\fontsize{5.8}{6.4}\selectfont \textcolor{BrickRed}{$ \blacktriangledown\,6\% $}}\\[-.2pt]2m\,51s\\[-.2pt]{\fontsize{5.8}{6.4}\selectfont \textcolor{BrickRed}{$ \blacktriangledown\,12\% $}}\end{tabular} & \hspace*{1pt}\begin{tabular}[t]{@{}c@{}}\rule{0pt}{12pt}55.9s\\[-.2pt]{\fontsize{5.8}{6.4}\selectfont \textcolor{ForestGreen}{$ \blacktriangle\,8.6\times $}}\\[-.2pt]55.1s\\[-.2pt]{\fontsize{5.8}{6.4}\selectfont \textcolor{ForestGreen}{$ \blacktriangle\,3.1\times $}}\end{tabular} & \hspace*{1pt}\begin{tabular}[t]{@{}c@{}}\rule{0pt}{12pt}43.3s\\[-.2pt]{\fontsize{5.8}{6.4}\selectfont \textcolor{ForestGreen}{$ \blacktriangle\,2.6\times $}}\\[-.2pt]42.5s\\[-.2pt]{\fontsize{5.8}{6.4}\selectfont \textcolor{ForestGreen}{$ \blacktriangle\,2.4\times $}}\end{tabular} & \hspace*{1pt}\begin{tabular}[t]{@{}c@{}}\rule{0pt}{12pt}38.9s\\[-.2pt]{\fontsize{5.8}{6.4}\selectfont \textcolor{ForestGreen}{$ \blacktriangle\,8.0\times $}}\\[-.2pt]38.2s\\[-.2pt]{\fontsize{5.8}{6.4}\selectfont \textcolor{ForestGreen}{$ \blacktriangle\,2.8\times $}}\end{tabular} & \hspace*{1pt}\begin{tabular}[t]{@{}c@{}}\rule{0pt}{12pt}53.1s\\[-.2pt]{\fontsize{5.8}{6.4}\selectfont \textcolor{ForestGreen}{$ \blacktriangle\,1.6\times $}}\\[-.2pt]48.2s\\[-.2pt]{\fontsize{5.8}{6.4}\selectfont \textcolor{ForestGreen}{$ \blacktriangle\,1.5\times $}}\end{tabular} & \hspace*{1pt}\begin{tabular}[t]{@{}c@{}}\rule{0pt}{12pt}28.7s\\[-.2pt]{\fontsize{5.8}{6.4}\selectfont \textcolor{ForestGreen}{$ \blacktriangle\,9.8\times $}}\\[-.2pt]28.2s\\[-.2pt]{\fontsize{5.8}{6.4}\selectfont \textcolor{ForestGreen}{$ \blacktriangle\,2.9\times $}}\end{tabular}\\[3pt]\hline
\BenchmarkSourceLabel{bench:local:row:mps-bounded-rmq-ksw2}{R21}\shortstack[l]{{\color{gray}R21}\\\textbf{BioseqDB}} & \begin{tabular}[t]{@{}l@{}}\rule{0pt}{12pt}MPS / Bounded\\RMQ / KSW2\end{tabular} & \textcolor{gray}{\ensuremath{\cdot}} & \textcolor{gray}{\ensuremath{\cdot}} & \textcolor{gray}{\ensuremath{\cdot}} & \textcolor{gray}{\ensuremath{\cdot}} & \hspace*{1pt}\begin{tabular}[t]{@{}c@{}}\rule{0pt}{12pt}5.1s\\[-.2pt]{\fontsize{5.8}{6.4}\selectfont \textcolor{ForestGreen}{$ \blacktriangle\,5.6\times $}}\\[-.2pt]3.2s\\[-.2pt]{\fontsize{5.8}{6.4}\selectfont \textcolor{ForestGreen}{$ \blacktriangle\,8.1\times $}}\end{tabular} & \hspace*{1pt}\begin{tabular}[t]{@{}c@{}}\rule{0pt}{12pt}3m\,47.8s\\[-.2pt]{\fontsize{5.8}{6.4}\selectfont \textcolor{BrickRed}{$ \blacktriangledown\,14.9\times $}}\\[-.2pt]22.2s\\[-.2pt]{\fontsize{5.8}{6.4}\selectfont \textcolor{BrickRed}{$ \blacktriangledown\,1.7\times $}}\end{tabular} & \hspace*{1pt}\begin{tabular}[t]{@{}c@{}}\rule{0pt}{12pt}5.3s\\[-.2pt]{\fontsize{5.8}{6.4}\selectfont \textcolor{ForestGreen}{$ \blacktriangle\,5.6\times $}}\\[-.2pt]3.7s\\[-.2pt]{\fontsize{5.8}{6.4}\selectfont \textcolor{ForestGreen}{$ \blacktriangle\,7.1\times $}}\end{tabular} & \hspace*{1pt}\begin{tabular}[t]{@{}c@{}}\rule{0pt}{12pt}3m\,31.9s\\[-.2pt]{\fontsize{5.8}{6.4}\selectfont \textcolor{BrickRed}{$ \blacktriangledown\,24.4\times $}}\\[-.2pt]13.4s\\[-.2pt]{\fontsize{5.8}{6.4}\selectfont \textcolor{BrickRed}{$ \blacktriangledown\,1.6\times $}}\end{tabular} & \textcolor{gray}{\ensuremath{\cdot}} & \textcolor{gray}{\ensuremath{\cdot}} & \textcolor{gray}{\ensuremath{\cdot}} & \textcolor{gray}{\ensuremath{\cdot}} & \textcolor{gray}{\ensuremath{\cdot}} & \textcolor{gray}{\ensuremath{\cdot}} & \textcolor{gray}{\ensuremath{\cdot}} & \textcolor{gray}{\ensuremath{\cdot}}\\[3pt]\hline
\BenchmarkSourceLabel{bench:local:row:mps-simdref-scalar-ksw2}{R22}\shortstack[l]{{\color{gray}R22}} & \begin{tabular}[t]{@{}l@{}}\rule{0pt}{12pt}MPS / SIMD-ref\\Scalar-DP / KSW2\end{tabular} & \textcolor{gray}{\ensuremath{\cdot}} & \textcolor{gray}{\ensuremath{\cdot}} & \textcolor{gray}{\ensuremath{\cdot}} & \textcolor{gray}{\ensuremath{\cdot}} & \hspace*{1pt}\begin{tabular}[t]{@{}c@{}}\rule{0pt}{12pt}5.1s\\[-.2pt]{\fontsize{5.8}{6.4}\selectfont \textcolor{ForestGreen}{$ \blacktriangle\,5.6\times $}}\\[-.2pt]3.4s\\[-.2pt]{\fontsize{5.8}{6.4}\selectfont \textcolor{ForestGreen}{$ \blacktriangle\,7.5\times $}}\end{tabular} & \hspace*{1pt}\begin{tabular}[t]{@{}c@{}}\rule{0pt}{12pt}3m\,51.5s\\[-.2pt]{\fontsize{5.8}{6.4}\selectfont \textcolor{BrickRed}{$ \blacktriangledown\,15.2\times $}}\\[-.2pt]25.9s\\[-.2pt]{\fontsize{5.8}{6.4}\selectfont \textcolor{BrickRed}{$ \blacktriangledown\,2.0\times $}}\end{tabular} & \hspace*{1pt}\begin{tabular}[t]{@{}c@{}}\rule{0pt}{12pt}14.7s\\[-.2pt]{\fontsize{5.8}{6.4}\selectfont \textcolor{ForestGreen}{$ \blacktriangle\,2.0\times $}}\\[-.2pt]5.7s\\[-.2pt]{\fontsize{5.8}{6.4}\selectfont \textcolor{ForestGreen}{$ \blacktriangle\,4.6\times $}}\end{tabular} & \hspace*{1pt}\begin{tabular}[t]{@{}c@{}}\rule{0pt}{12pt}3m\,44.9s\\[-.2pt]{\fontsize{5.8}{6.4}\selectfont \textcolor{BrickRed}{$ \blacktriangledown\,25.9\times $}}\\[-.2pt]21.1s\\[-.2pt]{\fontsize{5.8}{6.4}\selectfont \textcolor{BrickRed}{$ \blacktriangledown\,2.6\times $}}\end{tabular} & \textcolor{gray}{\ensuremath{\cdot}} & \textcolor{gray}{\ensuremath{\cdot}} & \textcolor{gray}{\ensuremath{\cdot}} & \textcolor{gray}{\ensuremath{\cdot}} & \textcolor{gray}{\ensuremath{\cdot}} & \textcolor{gray}{\ensuremath{\cdot}} & \textcolor{gray}{\ensuremath{\cdot}} & \textcolor{gray}{\ensuremath{\cdot}}\\[3pt]\hline
\BenchmarkSourceLabel{bench:local:row:mps-bounded-simd-ksw2}{R23}\shortstack[l]{{\color{gray}R23}} & \begin{tabular}[t]{@{}l@{}}\rule{0pt}{12pt}MPS / Bounded\\SIMD-DP / KSW2\end{tabular} & \textcolor{gray}{\ensuremath{\cdot}} & \textcolor{gray}{\ensuremath{\cdot}} & \textcolor{gray}{\ensuremath{\cdot}} & \textcolor{gray}{\ensuremath{\cdot}} & \hspace*{1pt}\begin{tabular}[t]{@{}c@{}}\rule{0pt}{12pt}5.2s\\[-.2pt]{\fontsize{5.8}{6.4}\selectfont \textcolor{ForestGreen}{$ \blacktriangle\,5.5\times $}}\\[-.2pt]3.2s\\[-.2pt]{\fontsize{5.8}{6.4}\selectfont \textcolor{ForestGreen}{$ \blacktriangle\,8.1\times $}}\end{tabular} & \hspace*{1pt}\begin{tabular}[t]{@{}c@{}}\rule{0pt}{12pt}3m\,47s\\[-.2pt]{\fontsize{5.8}{6.4}\selectfont \textcolor{BrickRed}{$ \blacktriangledown\,14.9\times $}}\\[-.2pt]21.5s\\[-.2pt]{\fontsize{5.8}{6.4}\selectfont \textcolor{BrickRed}{$ \blacktriangledown\,1.7\times $}}\end{tabular} & \hspace*{1pt}\begin{tabular}[t]{@{}c@{}}\rule{0pt}{12pt}5.3s\\[-.2pt]{\fontsize{5.8}{6.4}\selectfont \textcolor{ForestGreen}{$ \blacktriangle\,5.5\times $}}\\[-.2pt]3.8s\\[-.2pt]{\fontsize{5.8}{6.4}\selectfont \textcolor{ForestGreen}{$ \blacktriangle\,7.0\times $}}\end{tabular} & \hspace*{1pt}\begin{tabular}[t]{@{}c@{}}\rule{0pt}{12pt}3m\,32.1s\\[-.2pt]{\fontsize{5.8}{6.4}\selectfont \textcolor{BrickRed}{$ \blacktriangledown\,24.4\times $}}\\[-.2pt]13.2s\\[-.2pt]{\fontsize{5.8}{6.4}\selectfont \textcolor{BrickRed}{$ \blacktriangledown\,1.6\times $}}\end{tabular} & \hspace*{1pt}\begin{tabular}[t]{@{}c@{}}\rule{0pt}{12pt}3m\,54s\\[-.2pt]{\fontsize{5.8}{6.4}\selectfont \textcolor{BrickRed}{$ \blacktriangledown\,1.4\times $}}\\[-.2pt]3m\,39s\\[-.2pt]{\fontsize{5.8}{6.4}\selectfont \textcolor{BrickRed}{$ \blacktriangledown\,30\% $}}\end{tabular} & \textcolor{gray}{\ensuremath{\cdot}} & \hspace*{1pt}\begin{tabular}[t]{@{}c@{}}\rule{0pt}{12pt}2m\,52.6s\\[-.2pt]{\fontsize{5.8}{6.4}\selectfont \textcolor{gray}{$ \blacktriangleright\,1\% $}}\\[-.2pt]2m\,36.8s\\[-.2pt]{\fontsize{5.8}{6.4}\selectfont \textcolor{ForestGreen}{$ \blacktriangle\,8\% $}}\end{tabular} & \textcolor{gray}{\ensuremath{\cdot}} & \hspace*{1pt}\begin{tabular}[t]{@{}c@{}}\rule{0pt}{12pt}2m\,3.1s\\[-.2pt]{\fontsize{5.8}{6.4}\selectfont \textcolor{BrickRed}{$ \blacktriangledown\,2.8\times $}}\\[-.2pt]1m\,50.4s\\[-.2pt]{\fontsize{5.8}{6.4}\selectfont \textcolor{BrickRed}{$ \blacktriangledown\,2.6\times $}}\end{tabular} & \hspace*{1pt}\begin{tabular}[t]{@{}c@{}}\rule{0pt}{12pt}5m\,13.3s\\[-.2pt]{\fontsize{5.8}{6.4}\selectfont \textcolor{BrickRed}{$ \blacktriangledown\,8.1\times $}}\\[-.2pt]1m\,50.9s\\[-.2pt]{\fontsize{5.8}{6.4}\selectfont \textcolor{BrickRed}{$ \blacktriangledown\,2.9\times $}}\end{tabular} & \hspace*{1pt}\begin{tabular}[t]{@{}c@{}}\rule{0pt}{12pt}1m\,30.6s\\[-.2pt]{\fontsize{5.8}{6.4}\selectfont \textcolor{BrickRed}{$ \blacktriangledown\,1.7\times $}}\\[-.2pt]1m\,17.7s\\[-.2pt]{\fontsize{5.8}{6.4}\selectfont \textcolor{BrickRed}{$ \blacktriangledown\,1.6\times $}}\end{tabular} & \hspace*{1pt}\begin{tabular}[t]{@{}c@{}}\rule{0pt}{12pt}4m\,44.7s\\[-.2pt]{\fontsize{5.8}{6.4}\selectfont \textcolor{BrickRed}{$ \blacktriangledown\,9.9\times $}}\\[-.2pt]1m\,25.3s\\[-.2pt]{\fontsize{5.8}{6.4}\selectfont \textcolor{BrickRed}{$ \blacktriangledown\,3.0\times $}}\end{tabular}\\[3pt]\hline
\BenchmarkSourceLabel{bench:local:row:mps-twostack-rmq-ksw2}{R24}\shortstack[l]{{\color{gray}R24}} & \begin{tabular}[t]{@{}l@{}}\rule{0pt}{12pt}MPS / Two-stack\\RMQ / KSW2\end{tabular} & \textcolor{gray}{\ensuremath{\cdot}} & \textcolor{gray}{\ensuremath{\cdot}} & \textcolor{gray}{\ensuremath{\cdot}} & \textcolor{gray}{\ensuremath{\cdot}} & \hspace*{1pt}\begin{tabular}[t]{@{}c@{}}\rule{0pt}{12pt}5.2s\\[-.2pt]{\fontsize{5.8}{6.4}\selectfont \textcolor{ForestGreen}{$ \blacktriangle\,5.5\times $}}\\[-.2pt]3.4s\\[-.2pt]{\fontsize{5.8}{6.4}\selectfont \textcolor{ForestGreen}{$ \blacktriangle\,7.6\times $}}\end{tabular} & \hspace*{1pt}\begin{tabular}[t]{@{}c@{}}\rule{0pt}{12pt}3m\,34.7s\\[-.2pt]{\fontsize{5.8}{6.4}\selectfont \textcolor{BrickRed}{$ \blacktriangledown\,14.1\times $}}\\[-.2pt]16.7s\\[-.2pt]{\fontsize{5.8}{6.4}\selectfont \textcolor{BrickRed}{$ \blacktriangledown\,1.3\times $}}\end{tabular} & \hspace*{1pt}\begin{tabular}[t]{@{}c@{}}\rule{0pt}{12pt}16.7s\\[-.2pt]{\fontsize{5.8}{6.4}\selectfont \textcolor{ForestGreen}{$ \blacktriangle\,1.8\times $}}\\[-.2pt]4.6s\\[-.2pt]{\fontsize{5.8}{6.4}\selectfont \textcolor{ForestGreen}{$ \blacktriangle\,5.7\times $}}\end{tabular} & \hspace*{1pt}\begin{tabular}[t]{@{}c@{}}\rule{0pt}{12pt}3m\,27.7s\\[-.2pt]{\fontsize{5.8}{6.4}\selectfont \textcolor{BrickRed}{$ \blacktriangledown\,23.9\times $}}\\[-.2pt]8.4s\\[-.2pt]{\fontsize{5.8}{6.4}\selectfont \textcolor{gray}{$ \blacktriangleright\,3\% $}}\end{tabular} & \textcolor{gray}{\ensuremath{\cdot}} & \textcolor{gray}{\ensuremath{\cdot}} & \textcolor{gray}{\ensuremath{\cdot}} & \textcolor{gray}{\ensuremath{\cdot}} & \textcolor{gray}{\ensuremath{\cdot}} & \textcolor{gray}{\ensuremath{\cdot}} & \textcolor{gray}{\ensuremath{\cdot}} & \textcolor{gray}{\ensuremath{\cdot}}\\[3pt]\hline
\BenchmarkSourceLabel{bench:local:row:mps-twostack-scalar-ksw2}{R25}\shortstack[l]{{\color{gray}R25}} & \begin{tabular}[t]{@{}l@{}}\rule{0pt}{12pt}MPS / Two-stack\\Scalar-DP / KSW2\end{tabular} & \textcolor{gray}{\ensuremath{\cdot}} & \textcolor{gray}{\ensuremath{\cdot}} & \textcolor{gray}{\ensuremath{\cdot}} & \textcolor{gray}{\ensuremath{\cdot}} & \hspace*{1pt}\begin{tabular}[t]{@{}c@{}}\rule{0pt}{12pt}5.2s\\[-.2pt]{\fontsize{5.8}{6.4}\selectfont \textcolor{ForestGreen}{$ \blacktriangle\,5.5\times $}}\\[-.2pt]3.4s\\[-.2pt]{\fontsize{5.8}{6.4}\selectfont \textcolor{ForestGreen}{$ \blacktriangle\,7.6\times $}}\end{tabular} & \hspace*{1pt}\begin{tabular}[t]{@{}c@{}}\rule{0pt}{12pt}3m\,34.2s\\[-.2pt]{\fontsize{5.8}{6.4}\selectfont \textcolor{BrickRed}{$ \blacktriangledown\,14.0\times $}}\\[-.2pt]16.5s\\[-.2pt]{\fontsize{5.8}{6.4}\selectfont \textcolor{BrickRed}{$ \blacktriangledown\,30\% $}}\end{tabular} & \hspace*{1pt}\begin{tabular}[t]{@{}c@{}}\rule{0pt}{12pt}21.8s\\[-.2pt]{\fontsize{5.8}{6.4}\selectfont \textcolor{ForestGreen}{$ \blacktriangle\,26\% $}}\\[-.2pt]4.8s\\[-.2pt]{\fontsize{5.8}{6.4}\selectfont \textcolor{ForestGreen}{$ \blacktriangle\,5.5\times $}}\end{tabular} & \hspace*{1pt}\begin{tabular}[t]{@{}c@{}}\rule{0pt}{12pt}3m\,28.2s\\[-.2pt]{\fontsize{5.8}{6.4}\selectfont \textcolor{BrickRed}{$ \blacktriangledown\,24.0\times $}}\\[-.2pt]8.6s\\[-.2pt]{\fontsize{5.8}{6.4}\selectfont \textcolor{BrickRed}{$ \blacktriangledown\,6\% $}}\end{tabular} & \textcolor{gray}{\ensuremath{\cdot}} & \textcolor{gray}{\ensuremath{\cdot}} & \textcolor{gray}{\ensuremath{\cdot}} & \textcolor{gray}{\ensuremath{\cdot}} & \textcolor{gray}{\ensuremath{\cdot}} & \textcolor{gray}{\ensuremath{\cdot}} & \textcolor{gray}{\ensuremath{\cdot}} & \textcolor{gray}{\ensuremath{\cdot}}\\[3pt]\hline
\BenchmarkSourceLabel{bench:local:row:mps-twostack-simd-ksw2}{R26}\shortstack[l]{{\color{gray}R26}} & \begin{tabular}[t]{@{}l@{}}\rule{0pt}{12pt}MPS / Two-stack\\SIMD-DP / KSW2\end{tabular} & \textcolor{gray}{\ensuremath{\cdot}} & \textcolor{gray}{\ensuremath{\cdot}} & \textcolor{gray}{\ensuremath{\cdot}} & \textcolor{gray}{\ensuremath{\cdot}} & \hspace*{1pt}\begin{tabular}[t]{@{}c@{}}\rule{0pt}{12pt}5.2s\\[-.2pt]{\fontsize{5.8}{6.4}\selectfont \textcolor{ForestGreen}{$ \blacktriangle\,5.5\times $}}\\[-.2pt]3.4s\\[-.2pt]{\fontsize{5.8}{6.4}\selectfont \textcolor{ForestGreen}{$ \blacktriangle\,7.6\times $}}\end{tabular} & \hspace*{1pt}\begin{tabular}[t]{@{}c@{}}\rule{0pt}{12pt}3m\,40.4s\\[-.2pt]{\fontsize{5.8}{6.4}\selectfont \textcolor{BrickRed}{$ \blacktriangledown\,14.4\times $}}\\[-.2pt]17.5s\\[-.2pt]{\fontsize{5.8}{6.4}\selectfont \textcolor{BrickRed}{$ \blacktriangledown\,1.4\times $}}\end{tabular} & \hspace*{1pt}\begin{tabular}[t]{@{}c@{}}\rule{0pt}{12pt}16.6s\\[-.2pt]{\fontsize{5.8}{6.4}\selectfont \textcolor{ForestGreen}{$ \blacktriangle\,1.8\times $}}\\[-.2pt]4.5s\\[-.2pt]{\fontsize{5.8}{6.4}\selectfont \textcolor{ForestGreen}{$ \blacktriangle\,5.8\times $}}\end{tabular} & \hspace*{1pt}\begin{tabular}[t]{@{}c@{}}\rule{0pt}{12pt}3m\,28s\\[-.2pt]{\fontsize{5.8}{6.4}\selectfont \textcolor{BrickRed}{$ \blacktriangledown\,24.0\times $}}\\[-.2pt]8.5s\\[-.2pt]{\fontsize{5.8}{6.4}\selectfont \textcolor{BrickRed}{$ \blacktriangledown\,5\% $}}\end{tabular} & \textcolor{gray}{\ensuremath{\cdot}} & \textcolor{gray}{\ensuremath{\cdot}} & \textcolor{gray}{\ensuremath{\cdot}} & \textcolor{gray}{\ensuremath{\cdot}} & \textcolor{gray}{\ensuremath{\cdot}} & \textcolor{gray}{\ensuremath{\cdot}} & \textcolor{gray}{\ensuremath{\cdot}} & \textcolor{gray}{\ensuremath{\cdot}}\\[3pt]\hline
\BenchmarkSourceLabel{bench:local:row:mps-simdref-rmq-ksw2}{R27}\shortstack[l]{{\color{gray}R27}} & \begin{tabular}[t]{@{}l@{}}\rule{0pt}{12pt}MPS / SIMD-ref\\RMQ / KSW2\end{tabular} & \textcolor{gray}{\ensuremath{\cdot}} & \textcolor{gray}{\ensuremath{\cdot}} & \textcolor{gray}{\ensuremath{\cdot}} & \textcolor{gray}{\ensuremath{\cdot}} & \hspace*{1pt}\begin{tabular}[t]{@{}c@{}}\rule{0pt}{12pt}5.2s\\[-.2pt]{\fontsize{5.8}{6.4}\selectfont \textcolor{ForestGreen}{$ \blacktriangle\,5.5\times $}}\\[-.2pt]3.5s\\[-.2pt]{\fontsize{5.8}{6.4}\selectfont \textcolor{ForestGreen}{$ \blacktriangle\,7.3\times $}}\end{tabular} & \hspace*{1pt}\begin{tabular}[t]{@{}c@{}}\rule{0pt}{12pt}3m\,52.6s\\[-.2pt]{\fontsize{5.8}{6.4}\selectfont \textcolor{BrickRed}{$ \blacktriangledown\,15.2\times $}}\\[-.2pt]27.3s\\[-.2pt]{\fontsize{5.8}{6.4}\selectfont \textcolor{BrickRed}{$ \blacktriangledown\,2.2\times $}}\end{tabular} & \hspace*{1pt}\begin{tabular}[t]{@{}c@{}}\rule{0pt}{12pt}14.8s\\[-.2pt]{\fontsize{5.8}{6.4}\selectfont \textcolor{ForestGreen}{$ \blacktriangle\,2.0\times $}}\\[-.2pt]5.7s\\[-.2pt]{\fontsize{5.8}{6.4}\selectfont \textcolor{ForestGreen}{$ \blacktriangle\,4.6\times $}}\end{tabular} & \hspace*{1pt}\begin{tabular}[t]{@{}c@{}}\rule{0pt}{12pt}3m\,42.9s\\[-.2pt]{\fontsize{5.8}{6.4}\selectfont \textcolor{BrickRed}{$ \blacktriangledown\,25.7\times $}}\\[-.2pt]20.5s\\[-.2pt]{\fontsize{5.8}{6.4}\selectfont \textcolor{BrickRed}{$ \blacktriangledown\,2.5\times $}}\end{tabular} & \textcolor{gray}{\ensuremath{\cdot}} & \textcolor{gray}{\ensuremath{\cdot}} & \textcolor{gray}{\ensuremath{\cdot}} & \textcolor{gray}{\ensuremath{\cdot}} & \textcolor{gray}{\ensuremath{\cdot}} & \textcolor{gray}{\ensuremath{\cdot}} & \textcolor{gray}{\ensuremath{\cdot}} & \textcolor{gray}{\ensuremath{\cdot}}\\[3pt]\hline
\end{tabular}\par}
\smallskip{\footnotesize\noindent Each cell: A and B times, $\cdot$: no measurement. ---\textsuperscript{g}: GPU unavailable. n/a: no comparison.\par
\noindent Blue: one lowest A and one lowest B measurement for each column for each configuration.\par
\noindent Gray: minimap2. Yellow: my preferred MPF / Bounded / SIMD-DP / Block configuration.\par
\noindent MPF/MPS: minimizers per fragment/sequence. SIMD-ref: naive SIMD scan. Bounded: bounded SIMD algorithm.\par
}
\clearpage
\noindent{\large\bfseries Table \thetable: (A): Local machine\hfill 5/12}\par\smallskip
{\fontsize{7.1}{8.1}\selectfont\setlength{\tabcolsep}{1pt}\setlength{\fboxsep}{2pt}\renewcommand{\arraystretch}{1.12}
\noindent\begin{tabular}{@{}p{14mm}p{27mm}|p{13.55mm}|p{13.55mm}|p{13.55mm}|p{13.55mm}|p{13.55mm}|p{13.55mm}|p{13.55mm}|p{13.55mm}|p{13.55mm}|p{13.55mm}|p{13.55mm}|p{13.55mm}|p{13.55mm}|p{13.55mm}|p{13.55mm}|p{13.55mm}|@{}}\hline
Engine & Execution & \multicolumn{8}{c|}{\textbf{Index}} & \multicolumn{8}{c|}{\textbf{No index}}\\
 &  & \multicolumn{4}{c|}{\textbf{4 Gbp}} & \multicolumn{4}{c|}{\textbf{2 Gbp}} & \multicolumn{4}{c|}{\textbf{4 Gbp}} & \multicolumn{4}{c|}{\textbf{2 Gbp}}\\
 &  & \multicolumn{2}{c|}{\textbf{k=15, w=10}} & \multicolumn{2}{c|}{\textbf{k=20, w=15}} & \multicolumn{2}{c|}{\textbf{k=15, w=10}} & \multicolumn{2}{c|}{\textbf{k=20, w=15}} & \multicolumn{2}{c|}{\textbf{k=15, w=10}} & \multicolumn{2}{c|}{\textbf{k=20, w=15}} & \multicolumn{2}{c|}{\textbf{k=15, w=10}} & \multicolumn{2}{c|}{\textbf{k=20, w=15}}\\
 &  & \multicolumn{1}{c|}{\textbf{DNA4}} & \multicolumn{1}{c|}{\textbf{DNA5}} & \multicolumn{1}{c|}{\textbf{DNA4}} & \multicolumn{1}{c|}{\textbf{DNA5}} & \multicolumn{1}{c|}{\textbf{DNA4}} & \multicolumn{1}{c|}{\textbf{DNA5}} & \multicolumn{1}{c|}{\textbf{DNA4}} & \multicolumn{1}{c|}{\textbf{DNA5}} & \multicolumn{1}{c|}{\textbf{DNA4}} & \multicolumn{1}{c|}{\textbf{DNA5}} & \multicolumn{1}{c|}{\textbf{DNA4}} & \multicolumn{1}{c|}{\textbf{DNA5}} & \multicolumn{1}{c|}{\textbf{DNA4}} & \multicolumn{1}{c|}{\textbf{DNA5}} & \multicolumn{1}{c|}{\textbf{DNA4}} & \multicolumn{1}{c|}{\textbf{DNA5}}\\
 & & {\color{gray}C1} & {\color{gray}C2} & {\color{gray}C3} & {\color{gray}C4} & {\color{gray}C5} & {\color{gray}C6} & {\color{gray}C7} & {\color{gray}C8} & {\color{gray}C9} & {\color{gray}C10} & {\color{gray}C11} & {\color{gray}C12} & {\color{gray}C13} & {\color{gray}C14} & {\color{gray}C15} & {\color{gray}C16}\\
\hline
\rowcolor{BenchBaseline}[0.5pt][0.5pt]\shortstack[l]{{\color{gray}R1}\\\textbf{minimap2}} & \textbf{Baseline} & \hspace*{1pt}\begin{tabular}[t]{@{}c@{}}\rule{0pt}{12pt}39.7s\\[-.2pt]{\fontsize{5.8}{6.4}\selectfont \textcolor{BrickRed}{$ \blacktriangledown\,2.2\times $}}\\[-.2pt]36.8s\\[-.2pt]{\fontsize{5.8}{6.4}\selectfont \textcolor{BrickRed}{$ \blacktriangledown\,2.2\times $}}\end{tabular} & \hspace*{1pt}\begin{tabular}[t]{@{}c@{}}\rule{0pt}{12pt}25.2s\\[-.2pt]{\fontsize{5.8}{6.4}\selectfont \textcolor{BrickRed}{$ \blacktriangledown\,1.5\times $}}\\[-.2pt]22.7s\\[-.2pt]{\fontsize{5.8}{6.4}\selectfont \textcolor{BrickRed}{$ \blacktriangledown\,1.5\times $}}\end{tabular} & \hspace*{1pt}\begin{tabular}[t]{@{}c@{}}\rule{0pt}{12pt}1m\,56.9s\\[-.2pt]{\fontsize{5.8}{6.4}\selectfont \textcolor{BrickRed}{$ \blacktriangledown\,29.1\times $}}\\[-.2pt]1m\,53.2s\\[-.2pt]{\fontsize{5.8}{6.4}\selectfont \textcolor{BrickRed}{$ \blacktriangledown\,32.9\times $}}\end{tabular} & \hspace*{1pt}\begin{tabular}[t]{@{}c@{}}\rule{0pt}{12pt}24.9s\\[-.2pt]{\fontsize{5.8}{6.4}\selectfont \textcolor{BrickRed}{$ \blacktriangledown\,1.6\times $}}\\[-.2pt]22.2s\\[-.2pt]{\fontsize{5.8}{6.4}\selectfont \textcolor{BrickRed}{$ \blacktriangledown\,1.6\times $}}\end{tabular} & \hspace*{1pt}\begin{tabular}[t]{@{}c@{}}\rule{0pt}{12pt}28.6s\\[-.2pt]{\fontsize{5.8}{6.4}\selectfont \textcolor{BrickRed}{$ \blacktriangledown\,11.5\times $}}\\[-.2pt]25.8s\\[-.2pt]{\fontsize{5.8}{6.4}\selectfont \textcolor{BrickRed}{$ \blacktriangledown\,15.6\times $}}\end{tabular} & \hspace*{1pt}\begin{tabular}[t]{@{}c@{}}\rule{0pt}{12pt}15.3s\\[-.2pt]{\fontsize{5.8}{6.4}\selectfont \textcolor{BrickRed}{$ \blacktriangledown\,2.8\times $}}\\[-.2pt]12.7s\\[-.2pt]{\fontsize{5.8}{6.4}\selectfont \textcolor{BrickRed}{$ \blacktriangledown\,3.3\times $}}\end{tabular} & \hspace*{1pt}\begin{tabular}[t]{@{}c@{}}\rule{0pt}{12pt}29.3s\\[-.2pt]{\fontsize{5.8}{6.4}\selectfont \textcolor{BrickRed}{$ \blacktriangledown\,15.0\times $}}\\[-.2pt]26.2s\\[-.2pt]{\fontsize{5.8}{6.4}\selectfont \textcolor{BrickRed}{$ \blacktriangledown\,22.9\times $}}\end{tabular} & \hspace*{1pt}\begin{tabular}[t]{@{}c@{}}\rule{0pt}{12pt}8.7s\\[-.2pt]{\fontsize{5.8}{6.4}\selectfont \textcolor{BrickRed}{$ \blacktriangledown\,1.7\times $}}\\[-.2pt]8.1s\\[-.2pt]{\fontsize{5.8}{6.4}\selectfont \textcolor{BrickRed}{$ \blacktriangledown\,2.3\times $}}\end{tabular} & \hspace*{1pt}\begin{tabular}[t]{@{}c@{}}\rule{0pt}{12pt}2m\,52.4s\\[-.2pt]{\fontsize{5.8}{6.4}\selectfont \textcolor{ForestGreen}{$ \blacktriangle\,19\% $}}\\[-.2pt]2m\,48.5s\\[-.2pt]{\fontsize{5.8}{6.4}\selectfont \textcolor{ForestGreen}{$ \blacktriangle\,15\% $}}\end{tabular} & \hspace*{1pt}\begin{tabular}[t]{@{}c@{}}\rule{0pt}{12pt}2m\,17.1s\\[-.2pt]{\fontsize{5.8}{6.4}\selectfont \textcolor{ForestGreen}{$ \blacktriangle\,3.9\times $}}\\[-.2pt]2m\,15.9s\\[-.2pt]{\fontsize{5.8}{6.4}\selectfont \textcolor{ForestGreen}{$ \blacktriangle\,1.7\times $}}\end{tabular} & \hspace*{1pt}\begin{tabular}[t]{@{}c@{}}\rule{0pt}{12pt}2m\,55s\\[-.2pt]{\fontsize{5.8}{6.4}\selectfont \textcolor{BrickRed}{$ \blacktriangledown\,6\% $}}\\[-.2pt]2m\,51s\\[-.2pt]{\fontsize{5.8}{6.4}\selectfont \textcolor{BrickRed}{$ \blacktriangledown\,12\% $}}\end{tabular} & \hspace*{1pt}\begin{tabular}[t]{@{}c@{}}\rule{0pt}{12pt}55.9s\\[-.2pt]{\fontsize{5.8}{6.4}\selectfont \textcolor{ForestGreen}{$ \blacktriangle\,8.6\times $}}\\[-.2pt]55.1s\\[-.2pt]{\fontsize{5.8}{6.4}\selectfont \textcolor{ForestGreen}{$ \blacktriangle\,3.1\times $}}\end{tabular} & \hspace*{1pt}\begin{tabular}[t]{@{}c@{}}\rule{0pt}{12pt}43.3s\\[-.2pt]{\fontsize{5.8}{6.4}\selectfont \textcolor{ForestGreen}{$ \blacktriangle\,2.6\times $}}\\[-.2pt]42.5s\\[-.2pt]{\fontsize{5.8}{6.4}\selectfont \textcolor{ForestGreen}{$ \blacktriangle\,2.4\times $}}\end{tabular} & \hspace*{1pt}\begin{tabular}[t]{@{}c@{}}\rule{0pt}{12pt}38.9s\\[-.2pt]{\fontsize{5.8}{6.4}\selectfont \textcolor{ForestGreen}{$ \blacktriangle\,8.0\times $}}\\[-.2pt]38.2s\\[-.2pt]{\fontsize{5.8}{6.4}\selectfont \textcolor{ForestGreen}{$ \blacktriangle\,2.8\times $}}\end{tabular} & \hspace*{1pt}\begin{tabular}[t]{@{}c@{}}\rule{0pt}{12pt}53.1s\\[-.2pt]{\fontsize{5.8}{6.4}\selectfont \textcolor{ForestGreen}{$ \blacktriangle\,1.6\times $}}\\[-.2pt]48.2s\\[-.2pt]{\fontsize{5.8}{6.4}\selectfont \textcolor{ForestGreen}{$ \blacktriangle\,1.5\times $}}\end{tabular} & \hspace*{1pt}\begin{tabular}[t]{@{}c@{}}\rule{0pt}{12pt}28.7s\\[-.2pt]{\fontsize{5.8}{6.4}\selectfont \textcolor{ForestGreen}{$ \blacktriangle\,9.8\times $}}\\[-.2pt]28.2s\\[-.2pt]{\fontsize{5.8}{6.4}\selectfont \textcolor{ForestGreen}{$ \blacktriangle\,2.9\times $}}\end{tabular}\\[3pt]\hline
\BenchmarkSourceLabel{bench:local:row:mps-scalar-simd-ksw2}{R28}\shortstack[l]{{\color{gray}R28}\\\textbf{BioseqDB}} & \begin{tabular}[t]{@{}l@{}}\rule{0pt}{12pt}MPS / Scalar\\SIMD-DP / KSW2\end{tabular} & \textcolor{gray}{\ensuremath{\cdot}} & \textcolor{gray}{\ensuremath{\cdot}} & \textcolor{gray}{\ensuremath{\cdot}} & \textcolor{gray}{\ensuremath{\cdot}} & \hspace*{1pt}\begin{tabular}[t]{@{}c@{}}\rule{0pt}{12pt}5.3s\\[-.2pt]{\fontsize{5.8}{6.4}\selectfont \textcolor{ForestGreen}{$ \blacktriangle\,5.4\times $}}\\[-.2pt]3.4s\\[-.2pt]{\fontsize{5.8}{6.4}\selectfont \textcolor{ForestGreen}{$ \blacktriangle\,7.6\times $}}\end{tabular} & \hspace*{1pt}\begin{tabular}[t]{@{}c@{}}\rule{0pt}{12pt}3m\,46.1s\\[-.2pt]{\fontsize{5.8}{6.4}\selectfont \textcolor{BrickRed}{$ \blacktriangledown\,14.8\times $}}\\[-.2pt]21.5s\\[-.2pt]{\fontsize{5.8}{6.4}\selectfont \textcolor{BrickRed}{$ \blacktriangledown\,1.7\times $}}\end{tabular} & \hspace*{1pt}\begin{tabular}[t]{@{}c@{}}\rule{0pt}{12pt}17.6s\\[-.2pt]{\fontsize{5.8}{6.4}\selectfont \textcolor{ForestGreen}{$ \blacktriangle\,1.7\times $}}\\[-.2pt]5.5s\\[-.2pt]{\fontsize{5.8}{6.4}\selectfont \textcolor{ForestGreen}{$ \blacktriangle\,4.7\times $}}\end{tabular} & \hspace*{1pt}\begin{tabular}[t]{@{}c@{}}\rule{0pt}{12pt}3m\,36.6s\\[-.2pt]{\fontsize{5.8}{6.4}\selectfont \textcolor{BrickRed}{$ \blacktriangledown\,24.9\times $}}\\[-.2pt]15.3s\\[-.2pt]{\fontsize{5.8}{6.4}\selectfont \textcolor{BrickRed}{$ \blacktriangledown\,1.9\times $}}\end{tabular} & \textcolor{gray}{\ensuremath{\cdot}} & \textcolor{gray}{\ensuremath{\cdot}} & \textcolor{gray}{\ensuremath{\cdot}} & \textcolor{gray}{\ensuremath{\cdot}} & \textcolor{gray}{\ensuremath{\cdot}} & \textcolor{gray}{\ensuremath{\cdot}} & \textcolor{gray}{\ensuremath{\cdot}} & \textcolor{gray}{\ensuremath{\cdot}}\\[3pt]\hline
\BenchmarkSourceLabel{bench:local:row:mps-simdref-simd-ksw2}{R29}\shortstack[l]{{\color{gray}R29}} & \begin{tabular}[t]{@{}l@{}}\rule{0pt}{12pt}MPS / SIMD-ref\\SIMD-DP / KSW2\end{tabular} & \textcolor{gray}{\ensuremath{\cdot}} & \textcolor{gray}{\ensuremath{\cdot}} & \textcolor{gray}{\ensuremath{\cdot}} & \textcolor{gray}{\ensuremath{\cdot}} & \hspace*{1pt}\begin{tabular}[t]{@{}c@{}}\rule{0pt}{12pt}5.3s\\[-.2pt]{\fontsize{5.8}{6.4}\selectfont \textcolor{ForestGreen}{$ \blacktriangle\,5.4\times $}}\\[-.2pt]3.5s\\[-.2pt]{\fontsize{5.8}{6.4}\selectfont \textcolor{ForestGreen}{$ \blacktriangle\,7.4\times $}}\end{tabular} & \hspace*{1pt}\begin{tabular}[t]{@{}c@{}}\rule{0pt}{12pt}3m\,52.8s\\[-.2pt]{\fontsize{5.8}{6.4}\selectfont \textcolor{BrickRed}{$ \blacktriangledown\,15.2\times $}}\\[-.2pt]26.4s\\[-.2pt]{\fontsize{5.8}{6.4}\selectfont \textcolor{BrickRed}{$ \blacktriangledown\,2.1\times $}}\end{tabular} & \hspace*{1pt}\begin{tabular}[t]{@{}c@{}}\rule{0pt}{12pt}14.7s\\[-.2pt]{\fontsize{5.8}{6.4}\selectfont \textcolor{ForestGreen}{$ \blacktriangle\,2.0\times $}}\\[-.2pt]5.6s\\[-.2pt]{\fontsize{5.8}{6.4}\selectfont \textcolor{ForestGreen}{$ \blacktriangle\,4.7\times $}}\end{tabular} & \hspace*{1pt}\begin{tabular}[t]{@{}c@{}}\rule{0pt}{12pt}3m\,42.9s\\[-.2pt]{\fontsize{5.8}{6.4}\selectfont \textcolor{BrickRed}{$ \blacktriangledown\,25.7\times $}}\\[-.2pt]20.3s\\[-.2pt]{\fontsize{5.8}{6.4}\selectfont \textcolor{BrickRed}{$ \blacktriangledown\,2.5\times $}}\end{tabular} & \textcolor{gray}{\ensuremath{\cdot}} & \textcolor{gray}{\ensuremath{\cdot}} & \textcolor{gray}{\ensuremath{\cdot}} & \textcolor{gray}{\ensuremath{\cdot}} & \textcolor{gray}{\ensuremath{\cdot}} & \textcolor{gray}{\ensuremath{\cdot}} & \textcolor{gray}{\ensuremath{\cdot}} & \textcolor{gray}{\ensuremath{\cdot}}\\[3pt]\hline
\BenchmarkSourceLabel{bench:local:row:mps-scalar-rmq-ksw2}{R30}\shortstack[l]{{\color{gray}R30}} & \begin{tabular}[t]{@{}l@{}}\rule{0pt}{12pt}MPS / Scalar\\RMQ / KSW2\end{tabular} & \textcolor{gray}{\ensuremath{\cdot}} & \textcolor{gray}{\ensuremath{\cdot}} & \textcolor{gray}{\ensuremath{\cdot}} & \textcolor{gray}{\ensuremath{\cdot}} & \hspace*{1pt}\begin{tabular}[t]{@{}c@{}}\rule{0pt}{12pt}5.3s\\[-.2pt]{\fontsize{5.8}{6.4}\selectfont \textcolor{ForestGreen}{$ \blacktriangle\,5.4\times $}}\\[-.2pt]3.4s\\[-.2pt]{\fontsize{5.8}{6.4}\selectfont \textcolor{ForestGreen}{$ \blacktriangle\,7.6\times $}}\end{tabular} & \hspace*{1pt}\begin{tabular}[t]{@{}c@{}}\rule{0pt}{12pt}3m\,49.4s\\[-.2pt]{\fontsize{5.8}{6.4}\selectfont \textcolor{BrickRed}{$ \blacktriangledown\,15.0\times $}}\\[-.2pt]21.8s\\[-.2pt]{\fontsize{5.8}{6.4}\selectfont \textcolor{BrickRed}{$ \blacktriangledown\,1.7\times $}}\end{tabular} & \hspace*{1pt}\begin{tabular}[t]{@{}c@{}}\rule{0pt}{12pt}14.6s\\[-.2pt]{\fontsize{5.8}{6.4}\selectfont \textcolor{ForestGreen}{$ \blacktriangle\,2.0\times $}}\\[-.2pt]5.7s\\[-.2pt]{\fontsize{5.8}{6.4}\selectfont \textcolor{ForestGreen}{$ \blacktriangle\,4.6\times $}}\end{tabular} & \hspace*{1pt}\begin{tabular}[t]{@{}c@{}}\rule{0pt}{12pt}3m\,35.5s\\[-.2pt]{\fontsize{5.8}{6.4}\selectfont \textcolor{BrickRed}{$ \blacktriangledown\,24.8\times $}}\\[-.2pt]15.1s\\[-.2pt]{\fontsize{5.8}{6.4}\selectfont \textcolor{BrickRed}{$ \blacktriangledown\,1.9\times $}}\end{tabular} & \textcolor{gray}{\ensuremath{\cdot}} & \textcolor{gray}{\ensuremath{\cdot}} & \textcolor{gray}{\ensuremath{\cdot}} & \textcolor{gray}{\ensuremath{\cdot}} & \textcolor{gray}{\ensuremath{\cdot}} & \textcolor{gray}{\ensuremath{\cdot}} & \textcolor{gray}{\ensuremath{\cdot}} & \textcolor{gray}{\ensuremath{\cdot}}\\[3pt]\hline
\BenchmarkSourceLabel{bench:local:row:mps-scalar-scalar-ksw2}{R31}\shortstack[l]{{\color{gray}R31}} & \begin{tabular}[t]{@{}l@{}}\rule{0pt}{12pt}MPS / Scalar\\Scalar-DP / KSW2\end{tabular} & \textcolor{gray}{\ensuremath{\cdot}} & \textcolor{gray}{\ensuremath{\cdot}} & \textcolor{gray}{\ensuremath{\cdot}} & \textcolor{gray}{\ensuremath{\cdot}} & \hspace*{1pt}\begin{tabular}[t]{@{}c@{}}\rule{0pt}{12pt}5.3s\\[-.2pt]{\fontsize{5.8}{6.4}\selectfont \textcolor{ForestGreen}{$ \blacktriangle\,5.4\times $}}\\[-.2pt]3.4s\\[-.2pt]{\fontsize{5.8}{6.4}\selectfont \textcolor{ForestGreen}{$ \blacktriangle\,7.6\times $}}\end{tabular} & \hspace*{1pt}\begin{tabular}[t]{@{}c@{}}\rule{0pt}{12pt}3m\,47s\\[-.2pt]{\fontsize{5.8}{6.4}\selectfont \textcolor{BrickRed}{$ \blacktriangledown\,14.9\times $}}\\[-.2pt]22s\\[-.2pt]{\fontsize{5.8}{6.4}\selectfont \textcolor{BrickRed}{$ \blacktriangledown\,1.7\times $}}\end{tabular} & \hspace*{1pt}\begin{tabular}[t]{@{}c@{}}\rule{0pt}{12pt}16.3s\\[-.2pt]{\fontsize{5.8}{6.4}\selectfont \textcolor{ForestGreen}{$ \blacktriangle\,1.8\times $}}\\[-.2pt]5.6s\\[-.2pt]{\fontsize{5.8}{6.4}\selectfont \textcolor{ForestGreen}{$ \blacktriangle\,4.7\times $}}\end{tabular} & \hspace*{1pt}\begin{tabular}[t]{@{}c@{}}\rule{0pt}{12pt}3m\,35.8s\\[-.2pt]{\fontsize{5.8}{6.4}\selectfont \textcolor{BrickRed}{$ \blacktriangledown\,24.9\times $}}\\[-.2pt]15s\\[-.2pt]{\fontsize{5.8}{6.4}\selectfont \textcolor{BrickRed}{$ \blacktriangledown\,1.8\times $}}\end{tabular} & \textcolor{gray}{\ensuremath{\cdot}} & \textcolor{gray}{\ensuremath{\cdot}} & \textcolor{gray}{\ensuremath{\cdot}} & \textcolor{gray}{\ensuremath{\cdot}} & \textcolor{gray}{\ensuremath{\cdot}} & \textcolor{gray}{\ensuremath{\cdot}} & \textcolor{gray}{\ensuremath{\cdot}} & \textcolor{gray}{\ensuremath{\cdot}}\\[3pt]\hline
\BenchmarkSourceLabel{bench:local:row:mpf-twostack-simd-block}{R32}\shortstack[l]{{\color{gray}R32}} & \begin{tabular}[t]{@{}l@{}}\rule{0pt}{12pt}MPF / Two-stack\\SIMD-DP / Block\end{tabular} & \textcolor{gray}{\ensuremath{\cdot}} & \textcolor{gray}{\ensuremath{\cdot}} & \textcolor{gray}{\ensuremath{\cdot}} & \textcolor{gray}{\ensuremath{\cdot}} & \textcolor{gray}{\ensuremath{\cdot}} & \hspace*{1pt}\begin{tabular}[t]{@{}c@{}}\rule{0pt}{12pt}\colorbox{BenchBest}{\strut 5.4s}\\[-.2pt]{\fontsize{5.8}{6.4}\selectfont \textcolor{ForestGreen}{$ \blacktriangle\,2.8\times $}}\\[-.2pt]4.2s\\[-.2pt]{\fontsize{5.8}{6.4}\selectfont \textcolor{ForestGreen}{$ \blacktriangle\,3.0\times $}}\end{tabular} & \hspace*{1pt}\begin{tabular}[t]{@{}c@{}}\rule{0pt}{12pt}1m\,45s\\[-.2pt]{\fontsize{5.8}{6.4}\selectfont \textcolor{BrickRed}{$ \blacktriangledown\,3.6\times $}}\\[-.2pt]1m\,44.4s\\[-.2pt]{\fontsize{5.8}{6.4}\selectfont \textcolor{BrickRed}{$ \blacktriangledown\,4.0\times $}}\end{tabular} & \hspace*{1pt}\begin{tabular}[t]{@{}c@{}}\rule{0pt}{12pt}59.1s\\[-.2pt]{\fontsize{5.8}{6.4}\selectfont \textcolor{BrickRed}{$ \blacktriangledown\,6.8\times $}}\\[-.2pt]58.5s\\[-.2pt]{\fontsize{5.8}{6.4}\selectfont \textcolor{BrickRed}{$ \blacktriangledown\,7.2\times $}}\end{tabular} & \textcolor{gray}{\ensuremath{\cdot}} & \textcolor{gray}{\ensuremath{\cdot}} & \textcolor{gray}{\ensuremath{\cdot}} & \textcolor{gray}{\ensuremath{\cdot}} & \textcolor{gray}{\ensuremath{\cdot}} & \textcolor{gray}{\ensuremath{\cdot}} & \textcolor{gray}{\ensuremath{\cdot}} & \textcolor{gray}{\ensuremath{\cdot}}\\[3pt]\hline
\BenchmarkSourceLabel{bench:local:row:mpf-scalar-simd-block}{R33}\shortstack[l]{{\color{gray}R33}} & \begin{tabular}[t]{@{}l@{}}\rule{0pt}{12pt}MPF / Scalar\\SIMD-DP / Block\end{tabular} & \textcolor{gray}{\ensuremath{\cdot}} & \textcolor{gray}{\ensuremath{\cdot}} & \textcolor{gray}{\ensuremath{\cdot}} & \textcolor{gray}{\ensuremath{\cdot}} & \textcolor{gray}{\ensuremath{\cdot}} & \hspace*{1pt}\begin{tabular}[t]{@{}c@{}}\rule{0pt}{12pt}6.1s\\[-.2pt]{\fontsize{5.8}{6.4}\selectfont \textcolor{ForestGreen}{$ \blacktriangle\,2.5\times $}}\\[-.2pt]4.7s\\[-.2pt]{\fontsize{5.8}{6.4}\selectfont \textcolor{ForestGreen}{$ \blacktriangle\,2.7\times $}}\end{tabular} & \hspace*{1pt}\begin{tabular}[t]{@{}c@{}}\rule{0pt}{12pt}1m\,52.3s\\[-.2pt]{\fontsize{5.8}{6.4}\selectfont \textcolor{BrickRed}{$ \blacktriangledown\,3.8\times $}}\\[-.2pt]1m\,51.9s\\[-.2pt]{\fontsize{5.8}{6.4}\selectfont \textcolor{BrickRed}{$ \blacktriangledown\,4.3\times $}}\end{tabular} & \hspace*{1pt}\begin{tabular}[t]{@{}c@{}}\rule{0pt}{12pt}3m\,2.6s\\[-.2pt]{\fontsize{5.8}{6.4}\selectfont \textcolor{BrickRed}{$ \blacktriangledown\,21.0\times $}}\\[-.2pt]3m\,1.6s\\[-.2pt]{\fontsize{5.8}{6.4}\selectfont \textcolor{BrickRed}{$ \blacktriangledown\,22.3\times $}}\end{tabular} & \textcolor{gray}{\ensuremath{\cdot}} & \textcolor{gray}{\ensuremath{\cdot}} & \textcolor{gray}{\ensuremath{\cdot}} & \textcolor{gray}{\ensuremath{\cdot}} & \textcolor{gray}{\ensuremath{\cdot}} & \textcolor{gray}{\ensuremath{\cdot}} & \textcolor{gray}{\ensuremath{\cdot}} & \textcolor{gray}{\ensuremath{\cdot}}\\[3pt]\hline
\end{tabular}\par}
\smallskip{\footnotesize\noindent Each cell: A and B times, $\cdot$: no measurement. ---\textsuperscript{g}: GPU unavailable. n/a: no comparison.\par
\noindent Blue: one lowest A and one lowest B measurement for each column for each configuration.\par
\noindent Gray: minimap2. Yellow: my preferred MPF / Bounded / SIMD-DP / Block configuration.\par
\noindent MPF/MPS: minimizers per fragment/sequence. SIMD-ref: naive SIMD scan. Bounded: bounded SIMD algorithm.\par
}
\clearpage
\noindent{\large\bfseries Table \thetable: (A): Local machine\hfill 6/12}\par\smallskip
{\fontsize{7.1}{8.1}\selectfont\setlength{\tabcolsep}{1pt}\setlength{\fboxsep}{2pt}\renewcommand{\arraystretch}{1.12}
\noindent\begin{tabular}{@{}p{14mm}p{27mm}|p{13.55mm}|p{13.55mm}|p{13.55mm}|p{13.55mm}|p{13.55mm}|p{13.55mm}|p{13.55mm}|p{13.55mm}|p{13.55mm}|p{13.55mm}|p{13.55mm}|p{13.55mm}|p{13.55mm}|p{13.55mm}|p{13.55mm}|p{13.55mm}|@{}}\hline
Engine & Execution & \multicolumn{8}{c|}{\textbf{Index}} & \multicolumn{8}{c|}{\textbf{No index}}\\
 &  & \multicolumn{4}{c|}{\textbf{4 Gbp}} & \multicolumn{4}{c|}{\textbf{2 Gbp}} & \multicolumn{4}{c|}{\textbf{4 Gbp}} & \multicolumn{4}{c|}{\textbf{2 Gbp}}\\
 &  & \multicolumn{2}{c|}{\textbf{k=15, w=10}} & \multicolumn{2}{c|}{\textbf{k=20, w=15}} & \multicolumn{2}{c|}{\textbf{k=15, w=10}} & \multicolumn{2}{c|}{\textbf{k=20, w=15}} & \multicolumn{2}{c|}{\textbf{k=15, w=10}} & \multicolumn{2}{c|}{\textbf{k=20, w=15}} & \multicolumn{2}{c|}{\textbf{k=15, w=10}} & \multicolumn{2}{c|}{\textbf{k=20, w=15}}\\
 &  & \multicolumn{1}{c|}{\textbf{DNA4}} & \multicolumn{1}{c|}{\textbf{DNA5}} & \multicolumn{1}{c|}{\textbf{DNA4}} & \multicolumn{1}{c|}{\textbf{DNA5}} & \multicolumn{1}{c|}{\textbf{DNA4}} & \multicolumn{1}{c|}{\textbf{DNA5}} & \multicolumn{1}{c|}{\textbf{DNA4}} & \multicolumn{1}{c|}{\textbf{DNA5}} & \multicolumn{1}{c|}{\textbf{DNA4}} & \multicolumn{1}{c|}{\textbf{DNA5}} & \multicolumn{1}{c|}{\textbf{DNA4}} & \multicolumn{1}{c|}{\textbf{DNA5}} & \multicolumn{1}{c|}{\textbf{DNA4}} & \multicolumn{1}{c|}{\textbf{DNA5}} & \multicolumn{1}{c|}{\textbf{DNA4}} & \multicolumn{1}{c|}{\textbf{DNA5}}\\
 & & {\color{gray}C1} & {\color{gray}C2} & {\color{gray}C3} & {\color{gray}C4} & {\color{gray}C5} & {\color{gray}C6} & {\color{gray}C7} & {\color{gray}C8} & {\color{gray}C9} & {\color{gray}C10} & {\color{gray}C11} & {\color{gray}C12} & {\color{gray}C13} & {\color{gray}C14} & {\color{gray}C15} & {\color{gray}C16}\\
\hline
\rowcolor{BenchBaseline}[0.5pt][0.5pt]\shortstack[l]{{\color{gray}R1}\\\textbf{minimap2}} & \textbf{Baseline} & \hspace*{1pt}\begin{tabular}[t]{@{}c@{}}\rule{0pt}{12pt}39.7s\\[-.2pt]{\fontsize{5.8}{6.4}\selectfont \textcolor{BrickRed}{$ \blacktriangledown\,2.2\times $}}\\[-.2pt]36.8s\\[-.2pt]{\fontsize{5.8}{6.4}\selectfont \textcolor{BrickRed}{$ \blacktriangledown\,2.2\times $}}\end{tabular} & \hspace*{1pt}\begin{tabular}[t]{@{}c@{}}\rule{0pt}{12pt}25.2s\\[-.2pt]{\fontsize{5.8}{6.4}\selectfont \textcolor{BrickRed}{$ \blacktriangledown\,1.5\times $}}\\[-.2pt]22.7s\\[-.2pt]{\fontsize{5.8}{6.4}\selectfont \textcolor{BrickRed}{$ \blacktriangledown\,1.5\times $}}\end{tabular} & \hspace*{1pt}\begin{tabular}[t]{@{}c@{}}\rule{0pt}{12pt}1m\,56.9s\\[-.2pt]{\fontsize{5.8}{6.4}\selectfont \textcolor{BrickRed}{$ \blacktriangledown\,29.1\times $}}\\[-.2pt]1m\,53.2s\\[-.2pt]{\fontsize{5.8}{6.4}\selectfont \textcolor{BrickRed}{$ \blacktriangledown\,32.9\times $}}\end{tabular} & \hspace*{1pt}\begin{tabular}[t]{@{}c@{}}\rule{0pt}{12pt}24.9s\\[-.2pt]{\fontsize{5.8}{6.4}\selectfont \textcolor{BrickRed}{$ \blacktriangledown\,1.6\times $}}\\[-.2pt]22.2s\\[-.2pt]{\fontsize{5.8}{6.4}\selectfont \textcolor{BrickRed}{$ \blacktriangledown\,1.6\times $}}\end{tabular} & \hspace*{1pt}\begin{tabular}[t]{@{}c@{}}\rule{0pt}{12pt}28.6s\\[-.2pt]{\fontsize{5.8}{6.4}\selectfont \textcolor{BrickRed}{$ \blacktriangledown\,11.5\times $}}\\[-.2pt]25.8s\\[-.2pt]{\fontsize{5.8}{6.4}\selectfont \textcolor{BrickRed}{$ \blacktriangledown\,15.6\times $}}\end{tabular} & \hspace*{1pt}\begin{tabular}[t]{@{}c@{}}\rule{0pt}{12pt}15.3s\\[-.2pt]{\fontsize{5.8}{6.4}\selectfont \textcolor{BrickRed}{$ \blacktriangledown\,2.8\times $}}\\[-.2pt]12.7s\\[-.2pt]{\fontsize{5.8}{6.4}\selectfont \textcolor{BrickRed}{$ \blacktriangledown\,3.3\times $}}\end{tabular} & \hspace*{1pt}\begin{tabular}[t]{@{}c@{}}\rule{0pt}{12pt}29.3s\\[-.2pt]{\fontsize{5.8}{6.4}\selectfont \textcolor{BrickRed}{$ \blacktriangledown\,15.0\times $}}\\[-.2pt]26.2s\\[-.2pt]{\fontsize{5.8}{6.4}\selectfont \textcolor{BrickRed}{$ \blacktriangledown\,22.9\times $}}\end{tabular} & \hspace*{1pt}\begin{tabular}[t]{@{}c@{}}\rule{0pt}{12pt}8.7s\\[-.2pt]{\fontsize{5.8}{6.4}\selectfont \textcolor{BrickRed}{$ \blacktriangledown\,1.7\times $}}\\[-.2pt]8.1s\\[-.2pt]{\fontsize{5.8}{6.4}\selectfont \textcolor{BrickRed}{$ \blacktriangledown\,2.3\times $}}\end{tabular} & \hspace*{1pt}\begin{tabular}[t]{@{}c@{}}\rule{0pt}{12pt}2m\,52.4s\\[-.2pt]{\fontsize{5.8}{6.4}\selectfont \textcolor{ForestGreen}{$ \blacktriangle\,19\% $}}\\[-.2pt]2m\,48.5s\\[-.2pt]{\fontsize{5.8}{6.4}\selectfont \textcolor{ForestGreen}{$ \blacktriangle\,15\% $}}\end{tabular} & \hspace*{1pt}\begin{tabular}[t]{@{}c@{}}\rule{0pt}{12pt}2m\,17.1s\\[-.2pt]{\fontsize{5.8}{6.4}\selectfont \textcolor{ForestGreen}{$ \blacktriangle\,3.9\times $}}\\[-.2pt]2m\,15.9s\\[-.2pt]{\fontsize{5.8}{6.4}\selectfont \textcolor{ForestGreen}{$ \blacktriangle\,1.7\times $}}\end{tabular} & \hspace*{1pt}\begin{tabular}[t]{@{}c@{}}\rule{0pt}{12pt}2m\,55s\\[-.2pt]{\fontsize{5.8}{6.4}\selectfont \textcolor{BrickRed}{$ \blacktriangledown\,6\% $}}\\[-.2pt]2m\,51s\\[-.2pt]{\fontsize{5.8}{6.4}\selectfont \textcolor{BrickRed}{$ \blacktriangledown\,12\% $}}\end{tabular} & \hspace*{1pt}\begin{tabular}[t]{@{}c@{}}\rule{0pt}{12pt}55.9s\\[-.2pt]{\fontsize{5.8}{6.4}\selectfont \textcolor{ForestGreen}{$ \blacktriangle\,8.6\times $}}\\[-.2pt]55.1s\\[-.2pt]{\fontsize{5.8}{6.4}\selectfont \textcolor{ForestGreen}{$ \blacktriangle\,3.1\times $}}\end{tabular} & \hspace*{1pt}\begin{tabular}[t]{@{}c@{}}\rule{0pt}{12pt}43.3s\\[-.2pt]{\fontsize{5.8}{6.4}\selectfont \textcolor{ForestGreen}{$ \blacktriangle\,2.6\times $}}\\[-.2pt]42.5s\\[-.2pt]{\fontsize{5.8}{6.4}\selectfont \textcolor{ForestGreen}{$ \blacktriangle\,2.4\times $}}\end{tabular} & \hspace*{1pt}\begin{tabular}[t]{@{}c@{}}\rule{0pt}{12pt}38.9s\\[-.2pt]{\fontsize{5.8}{6.4}\selectfont \textcolor{ForestGreen}{$ \blacktriangle\,8.0\times $}}\\[-.2pt]38.2s\\[-.2pt]{\fontsize{5.8}{6.4}\selectfont \textcolor{ForestGreen}{$ \blacktriangle\,2.8\times $}}\end{tabular} & \hspace*{1pt}\begin{tabular}[t]{@{}c@{}}\rule{0pt}{12pt}53.1s\\[-.2pt]{\fontsize{5.8}{6.4}\selectfont \textcolor{ForestGreen}{$ \blacktriangle\,1.6\times $}}\\[-.2pt]48.2s\\[-.2pt]{\fontsize{5.8}{6.4}\selectfont \textcolor{ForestGreen}{$ \blacktriangle\,1.5\times $}}\end{tabular} & \hspace*{1pt}\begin{tabular}[t]{@{}c@{}}\rule{0pt}{12pt}28.7s\\[-.2pt]{\fontsize{5.8}{6.4}\selectfont \textcolor{ForestGreen}{$ \blacktriangle\,9.8\times $}}\\[-.2pt]28.2s\\[-.2pt]{\fontsize{5.8}{6.4}\selectfont \textcolor{ForestGreen}{$ \blacktriangle\,2.9\times $}}\end{tabular}\\[3pt]\hline
\BenchmarkSourceLabel{bench:local:row:mpf-scalar-rmq-block}{R34}\shortstack[l]{{\color{gray}R34}\\\textbf{BioseqDB}} & \begin{tabular}[t]{@{}l@{}}\rule{0pt}{12pt}MPF / Scalar\\RMQ / Block\end{tabular} & \textcolor{gray}{\ensuremath{\cdot}} & \textcolor{gray}{\ensuremath{\cdot}} & \textcolor{gray}{\ensuremath{\cdot}} & \textcolor{gray}{\ensuremath{\cdot}} & \textcolor{gray}{\ensuremath{\cdot}} & \hspace*{1pt}\begin{tabular}[t]{@{}c@{}}\rule{0pt}{12pt}8s\\[-.2pt]{\fontsize{5.8}{6.4}\selectfont \textcolor{ForestGreen}{$ \blacktriangle\,1.9\times $}}\\[-.2pt]6.7s\\[-.2pt]{\fontsize{5.8}{6.4}\selectfont \textcolor{ForestGreen}{$ \blacktriangle\,1.9\times $}}\end{tabular} & \hspace*{1pt}\begin{tabular}[t]{@{}c@{}}\rule{0pt}{12pt}1m\,50s\\[-.2pt]{\fontsize{5.8}{6.4}\selectfont \textcolor{BrickRed}{$ \blacktriangledown\,3.7\times $}}\\[-.2pt]1m\,49.6s\\[-.2pt]{\fontsize{5.8}{6.4}\selectfont \textcolor{BrickRed}{$ \blacktriangledown\,4.2\times $}}\end{tabular} & \hspace*{1pt}\begin{tabular}[t]{@{}c@{}}\rule{0pt}{12pt}3m\,3.4s\\[-.2pt]{\fontsize{5.8}{6.4}\selectfont \textcolor{BrickRed}{$ \blacktriangledown\,21.1\times $}}\\[-.2pt]3m\,2.4s\\[-.2pt]{\fontsize{5.8}{6.4}\selectfont \textcolor{BrickRed}{$ \blacktriangledown\,22.4\times $}}\end{tabular} & \textcolor{gray}{\ensuremath{\cdot}} & \textcolor{gray}{\ensuremath{\cdot}} & \textcolor{gray}{\ensuremath{\cdot}} & \textcolor{gray}{\ensuremath{\cdot}} & \textcolor{gray}{\ensuremath{\cdot}} & \textcolor{gray}{\ensuremath{\cdot}} & \textcolor{gray}{\ensuremath{\cdot}} & \textcolor{gray}{\ensuremath{\cdot}}\\[3pt]\hline
\BenchmarkSourceLabel{bench:local:row:mpf-scalar-scalar-ksw2}{R35}\shortstack[l]{{\color{gray}R35}} & \begin{tabular}[t]{@{}l@{}}\rule{0pt}{12pt}MPF / Scalar\\Scalar-DP / KSW2\end{tabular} & \textcolor{gray}{\ensuremath{\cdot}} & \textcolor{gray}{\ensuremath{\cdot}} & \textcolor{gray}{\ensuremath{\cdot}} & \textcolor{gray}{\ensuremath{\cdot}} & \textcolor{gray}{\ensuremath{\cdot}} & \hspace*{1pt}\begin{tabular}[t]{@{}c@{}}\rule{0pt}{12pt}8.5s\\[-.2pt]{\fontsize{5.8}{6.4}\selectfont \textcolor{ForestGreen}{$ \blacktriangle\,1.8\times $}}\\[-.2pt]7s\\[-.2pt]{\fontsize{5.8}{6.4}\selectfont \textcolor{ForestGreen}{$ \blacktriangle\,1.8\times $}}\end{tabular} & \hspace*{1pt}\begin{tabular}[t]{@{}c@{}}\rule{0pt}{12pt}1m\,43.2s\\[-.2pt]{\fontsize{5.8}{6.4}\selectfont \textcolor{BrickRed}{$ \blacktriangledown\,3.5\times $}}\\[-.2pt]1m\,41.9s\\[-.2pt]{\fontsize{5.8}{6.4}\selectfont \textcolor{BrickRed}{$ \blacktriangledown\,3.9\times $}}\end{tabular} & \hspace*{1pt}\begin{tabular}[t]{@{}c@{}}\rule{0pt}{12pt}3m\,8.6s\\[-.2pt]{\fontsize{5.8}{6.4}\selectfont \textcolor{BrickRed}{$ \blacktriangledown\,21.7\times $}}\\[-.2pt]3m\,7.3s\\[-.2pt]{\fontsize{5.8}{6.4}\selectfont \textcolor{BrickRed}{$ \blacktriangledown\,23.0\times $}}\end{tabular} & \textcolor{gray}{\ensuremath{\cdot}} & \textcolor{gray}{\ensuremath{\cdot}} & \textcolor{gray}{\ensuremath{\cdot}} & \textcolor{gray}{\ensuremath{\cdot}} & \textcolor{gray}{\ensuremath{\cdot}} & \textcolor{gray}{\ensuremath{\cdot}} & \textcolor{gray}{\ensuremath{\cdot}} & \textcolor{gray}{\ensuremath{\cdot}}\\[3pt]\hline
\BenchmarkSourceLabel{bench:local:row:mpf-scalar-simd-ksw2}{R36}\shortstack[l]{{\color{gray}R36}} & \begin{tabular}[t]{@{}l@{}}\rule{0pt}{12pt}MPF / Scalar\\SIMD-DP / KSW2\end{tabular} & \textcolor{gray}{\ensuremath{\cdot}} & \textcolor{gray}{\ensuremath{\cdot}} & \textcolor{gray}{\ensuremath{\cdot}} & \textcolor{gray}{\ensuremath{\cdot}} & \textcolor{gray}{\ensuremath{\cdot}} & \hspace*{1pt}\begin{tabular}[t]{@{}c@{}}\rule{0pt}{12pt}9s\\[-.2pt]{\fontsize{5.8}{6.4}\selectfont \textcolor{ForestGreen}{$ \blacktriangle\,1.7\times $}}\\[-.2pt]7.5s\\[-.2pt]{\fontsize{5.8}{6.4}\selectfont \textcolor{ForestGreen}{$ \blacktriangle\,1.7\times $}}\end{tabular} & \hspace*{1pt}\begin{tabular}[t]{@{}c@{}}\rule{0pt}{12pt}1m\,43s\\[-.2pt]{\fontsize{5.8}{6.4}\selectfont \textcolor{BrickRed}{$ \blacktriangledown\,3.5\times $}}\\[-.2pt]1m\,41.7s\\[-.2pt]{\fontsize{5.8}{6.4}\selectfont \textcolor{BrickRed}{$ \blacktriangledown\,3.9\times $}}\end{tabular} & \hspace*{1pt}\begin{tabular}[t]{@{}c@{}}\rule{0pt}{12pt}3m\,8.9s\\[-.2pt]{\fontsize{5.8}{6.4}\selectfont \textcolor{BrickRed}{$ \blacktriangledown\,21.8\times $}}\\[-.2pt]3m\,7.6s\\[-.2pt]{\fontsize{5.8}{6.4}\selectfont \textcolor{BrickRed}{$ \blacktriangledown\,23.1\times $}}\end{tabular} & \textcolor{gray}{\ensuremath{\cdot}} & \textcolor{gray}{\ensuremath{\cdot}} & \textcolor{gray}{\ensuremath{\cdot}} & \textcolor{gray}{\ensuremath{\cdot}} & \textcolor{gray}{\ensuremath{\cdot}} & \textcolor{gray}{\ensuremath{\cdot}} & \textcolor{gray}{\ensuremath{\cdot}} & \textcolor{gray}{\ensuremath{\cdot}}\\[3pt]\hline
\BenchmarkSourceLabel{bench:local:row:mpf-scalar-scalar-block}{R37}\shortstack[l]{{\color{gray}R37}} & \begin{tabular}[t]{@{}l@{}}\rule{0pt}{12pt}MPF / Scalar\\Scalar-DP / Block\end{tabular} & \textcolor{gray}{\ensuremath{\cdot}} & \textcolor{gray}{\ensuremath{\cdot}} & \textcolor{gray}{\ensuremath{\cdot}} & \textcolor{gray}{\ensuremath{\cdot}} & \textcolor{gray}{\ensuremath{\cdot}} & \hspace*{1pt}\begin{tabular}[t]{@{}c@{}}\rule{0pt}{12pt}10.1s\\[-.2pt]{\fontsize{5.8}{6.4}\selectfont \textcolor{ForestGreen}{$ \blacktriangle\,1.5\times $}}\\[-.2pt]8.9s\\[-.2pt]{\fontsize{5.8}{6.4}\selectfont \textcolor{ForestGreen}{$ \blacktriangle\,30\% $}}\end{tabular} & \hspace*{1pt}\begin{tabular}[t]{@{}c@{}}\rule{0pt}{12pt}1m\,48.8s\\[-.2pt]{\fontsize{5.8}{6.4}\selectfont \textcolor{BrickRed}{$ \blacktriangledown\,3.7\times $}}\\[-.2pt]1m\,48.4s\\[-.2pt]{\fontsize{5.8}{6.4}\selectfont \textcolor{BrickRed}{$ \blacktriangledown\,4.1\times $}}\end{tabular} & \hspace*{1pt}\begin{tabular}[t]{@{}c@{}}\rule{0pt}{12pt}3m\,4.8s\\[-.2pt]{\fontsize{5.8}{6.4}\selectfont \textcolor{BrickRed}{$ \blacktriangledown\,21.3\times $}}\\[-.2pt]3m\,3.8s\\[-.2pt]{\fontsize{5.8}{6.4}\selectfont \textcolor{BrickRed}{$ \blacktriangledown\,22.6\times $}}\end{tabular} & \textcolor{gray}{\ensuremath{\cdot}} & \textcolor{gray}{\ensuremath{\cdot}} & \textcolor{gray}{\ensuremath{\cdot}} & \textcolor{gray}{\ensuremath{\cdot}} & \textcolor{gray}{\ensuremath{\cdot}} & \textcolor{gray}{\ensuremath{\cdot}} & \textcolor{gray}{\ensuremath{\cdot}} & \textcolor{gray}{\ensuremath{\cdot}}\\[3pt]\hline
\BenchmarkSourceLabel{bench:local:row:mpf-twostack-simd-ksw2}{R38}\shortstack[l]{{\color{gray}R38}} & \begin{tabular}[t]{@{}l@{}}\rule{0pt}{12pt}MPF / Two-stack\\SIMD-DP / KSW2\end{tabular} & \textcolor{gray}{\ensuremath{\cdot}} & \textcolor{gray}{\ensuremath{\cdot}} & \textcolor{gray}{\ensuremath{\cdot}} & \textcolor{gray}{\ensuremath{\cdot}} & \textcolor{gray}{\ensuremath{\cdot}} & \hspace*{1pt}\begin{tabular}[t]{@{}c@{}}\rule{0pt}{12pt}10.4s\\[-.2pt]{\fontsize{5.8}{6.4}\selectfont \textcolor{ForestGreen}{$ \blacktriangle\,1.5\times $}}\\[-.2pt]8.8s\\[-.2pt]{\fontsize{5.8}{6.4}\selectfont \textcolor{ForestGreen}{$ \blacktriangle\,1.4\times $}}\end{tabular} & \hspace*{1pt}\begin{tabular}[t]{@{}c@{}}\rule{0pt}{12pt}1m\,21.3s\\[-.2pt]{\fontsize{5.8}{6.4}\selectfont \textcolor{BrickRed}{$ \blacktriangledown\,2.8\times $}}\\[-.2pt]1m\,19.5s\\[-.2pt]{\fontsize{5.8}{6.4}\selectfont \textcolor{BrickRed}{$ \blacktriangledown\,3.0\times $}}\end{tabular} & \hspace*{1pt}\begin{tabular}[t]{@{}c@{}}\rule{0pt}{12pt}1m\,1.1s\\[-.2pt]{\fontsize{5.8}{6.4}\selectfont \textcolor{BrickRed}{$ \blacktriangledown\,7.0\times $}}\\[-.2pt]58.6s\\[-.2pt]{\fontsize{5.8}{6.4}\selectfont \textcolor{BrickRed}{$ \blacktriangledown\,7.2\times $}}\end{tabular} & \textcolor{gray}{\ensuremath{\cdot}} & \textcolor{gray}{\ensuremath{\cdot}} & \textcolor{gray}{\ensuremath{\cdot}} & \textcolor{gray}{\ensuremath{\cdot}} & \textcolor{gray}{\ensuremath{\cdot}} & \textcolor{gray}{\ensuremath{\cdot}} & \textcolor{gray}{\ensuremath{\cdot}} & \textcolor{gray}{\ensuremath{\cdot}}\\[3pt]\hline
\BenchmarkSourceLabel{bench:local:row:mpf-simdref-rmq-block}{R39}\shortstack[l]{{\color{gray}R39}} & \begin{tabular}[t]{@{}l@{}}\rule{0pt}{12pt}MPF / SIMD-ref\\RMQ / Block\end{tabular} & \textcolor{gray}{\ensuremath{\cdot}} & \textcolor{gray}{\ensuremath{\cdot}} & \textcolor{gray}{\ensuremath{\cdot}} & \textcolor{gray}{\ensuremath{\cdot}} & \textcolor{gray}{\ensuremath{\cdot}} & \hspace*{1pt}\begin{tabular}[t]{@{}c@{}}\rule{0pt}{12pt}10.5s\\[-.2pt]{\fontsize{5.8}{6.4}\selectfont \textcolor{ForestGreen}{$ \blacktriangle\,1.4\times $}}\\[-.2pt]9.3s\\[-.2pt]{\fontsize{5.8}{6.4}\selectfont \textcolor{ForestGreen}{$ \blacktriangle\,26\% $}}\end{tabular} & \hspace*{1pt}\begin{tabular}[t]{@{}c@{}}\rule{0pt}{12pt}1m\,40.4s\\[-.2pt]{\fontsize{5.8}{6.4}\selectfont \textcolor{BrickRed}{$ \blacktriangledown\,3.4\times $}}\\[-.2pt]1m\,40s\\[-.2pt]{\fontsize{5.8}{6.4}\selectfont \textcolor{BrickRed}{$ \blacktriangledown\,3.8\times $}}\end{tabular} & \hspace*{1pt}\begin{tabular}[t]{@{}c@{}}\rule{0pt}{12pt}1m\,36.1s\\[-.2pt]{\fontsize{5.8}{6.4}\selectfont \textcolor{BrickRed}{$ \blacktriangledown\,11.1\times $}}\\[-.2pt]1m\,35.1s\\[-.2pt]{\fontsize{5.8}{6.4}\selectfont \textcolor{BrickRed}{$ \blacktriangledown\,11.7\times $}}\end{tabular} & \textcolor{gray}{\ensuremath{\cdot}} & \textcolor{gray}{\ensuremath{\cdot}} & \textcolor{gray}{\ensuremath{\cdot}} & \textcolor{gray}{\ensuremath{\cdot}} & \textcolor{gray}{\ensuremath{\cdot}} & \textcolor{gray}{\ensuremath{\cdot}} & \textcolor{gray}{\ensuremath{\cdot}} & \textcolor{gray}{\ensuremath{\cdot}}\\[3pt]\hline
\BenchmarkSourceLabel{bench:local:row:mpf-simdref-simd-block}{R40}\shortstack[l]{{\color{gray}R40}} & \begin{tabular}[t]{@{}l@{}}\rule{0pt}{12pt}MPF / SIMD-ref\\SIMD-DP / Block\end{tabular} & \textcolor{gray}{\ensuremath{\cdot}} & \textcolor{gray}{\ensuremath{\cdot}} & \textcolor{gray}{\ensuremath{\cdot}} & \textcolor{gray}{\ensuremath{\cdot}} & \textcolor{gray}{\ensuremath{\cdot}} & \hspace*{1pt}\begin{tabular}[t]{@{}c@{}}\rule{0pt}{12pt}10.8s\\[-.2pt]{\fontsize{5.8}{6.4}\selectfont \textcolor{ForestGreen}{$ \blacktriangle\,29\% $}}\\[-.2pt]9.6s\\[-.2pt]{\fontsize{5.8}{6.4}\selectfont \textcolor{ForestGreen}{$ \blacktriangle\,25\% $}}\end{tabular} & \hspace*{1pt}\begin{tabular}[t]{@{}c@{}}\rule{0pt}{12pt}1m\,42.9s\\[-.2pt]{\fontsize{5.8}{6.4}\selectfont \textcolor{BrickRed}{$ \blacktriangledown\,3.5\times $}}\\[-.2pt]1m\,42.6s\\[-.2pt]{\fontsize{5.8}{6.4}\selectfont \textcolor{BrickRed}{$ \blacktriangledown\,3.9\times $}}\end{tabular} & \hspace*{1pt}\begin{tabular}[t]{@{}c@{}}\rule{0pt}{12pt}1m\,57.8s\\[-.2pt]{\fontsize{5.8}{6.4}\selectfont \textcolor{BrickRed}{$ \blacktriangledown\,13.6\times $}}\\[-.2pt]1m\,56.7s\\[-.2pt]{\fontsize{5.8}{6.4}\selectfont \textcolor{BrickRed}{$ \blacktriangledown\,14.4\times $}}\end{tabular} & \textcolor{gray}{\ensuremath{\cdot}} & \textcolor{gray}{\ensuremath{\cdot}} & \textcolor{gray}{\ensuremath{\cdot}} & \textcolor{gray}{\ensuremath{\cdot}} & \textcolor{gray}{\ensuremath{\cdot}} & \textcolor{gray}{\ensuremath{\cdot}} & \textcolor{gray}{\ensuremath{\cdot}} & \textcolor{gray}{\ensuremath{\cdot}}\\[3pt]\hline
\end{tabular}\par}
\smallskip{\footnotesize\noindent Each cell: A and B times, $\cdot$: no measurement. ---\textsuperscript{g}: GPU unavailable. n/a: no comparison.\par
\noindent Blue: one lowest A and one lowest B measurement for each column for each configuration.\par
\noindent Gray: minimap2. Yellow: my preferred MPF / Bounded / SIMD-DP / Block configuration.\par
\noindent MPF/MPS: minimizers per fragment/sequence. SIMD-ref: naive SIMD scan. Bounded: bounded SIMD algorithm.\par
}
\clearpage
\noindent{\large\bfseries Table \thetable: (A): Local machine\hfill 7/12}\par\smallskip
{\fontsize{7.1}{8.1}\selectfont\setlength{\tabcolsep}{1pt}\setlength{\fboxsep}{2pt}\renewcommand{\arraystretch}{1.12}
\noindent\begin{tabular}{@{}p{14mm}p{27mm}|p{13.55mm}|p{13.55mm}|p{13.55mm}|p{13.55mm}|p{13.55mm}|p{13.55mm}|p{13.55mm}|p{13.55mm}|p{13.55mm}|p{13.55mm}|p{13.55mm}|p{13.55mm}|p{13.55mm}|p{13.55mm}|p{13.55mm}|p{13.55mm}|@{}}\hline
Engine & Execution & \multicolumn{8}{c|}{\textbf{Index}} & \multicolumn{8}{c|}{\textbf{No index}}\\
 &  & \multicolumn{4}{c|}{\textbf{4 Gbp}} & \multicolumn{4}{c|}{\textbf{2 Gbp}} & \multicolumn{4}{c|}{\textbf{4 Gbp}} & \multicolumn{4}{c|}{\textbf{2 Gbp}}\\
 &  & \multicolumn{2}{c|}{\textbf{k=15, w=10}} & \multicolumn{2}{c|}{\textbf{k=20, w=15}} & \multicolumn{2}{c|}{\textbf{k=15, w=10}} & \multicolumn{2}{c|}{\textbf{k=20, w=15}} & \multicolumn{2}{c|}{\textbf{k=15, w=10}} & \multicolumn{2}{c|}{\textbf{k=20, w=15}} & \multicolumn{2}{c|}{\textbf{k=15, w=10}} & \multicolumn{2}{c|}{\textbf{k=20, w=15}}\\
 &  & \multicolumn{1}{c|}{\textbf{DNA4}} & \multicolumn{1}{c|}{\textbf{DNA5}} & \multicolumn{1}{c|}{\textbf{DNA4}} & \multicolumn{1}{c|}{\textbf{DNA5}} & \multicolumn{1}{c|}{\textbf{DNA4}} & \multicolumn{1}{c|}{\textbf{DNA5}} & \multicolumn{1}{c|}{\textbf{DNA4}} & \multicolumn{1}{c|}{\textbf{DNA5}} & \multicolumn{1}{c|}{\textbf{DNA4}} & \multicolumn{1}{c|}{\textbf{DNA5}} & \multicolumn{1}{c|}{\textbf{DNA4}} & \multicolumn{1}{c|}{\textbf{DNA5}} & \multicolumn{1}{c|}{\textbf{DNA4}} & \multicolumn{1}{c|}{\textbf{DNA5}} & \multicolumn{1}{c|}{\textbf{DNA4}} & \multicolumn{1}{c|}{\textbf{DNA5}}\\
 & & {\color{gray}C1} & {\color{gray}C2} & {\color{gray}C3} & {\color{gray}C4} & {\color{gray}C5} & {\color{gray}C6} & {\color{gray}C7} & {\color{gray}C8} & {\color{gray}C9} & {\color{gray}C10} & {\color{gray}C11} & {\color{gray}C12} & {\color{gray}C13} & {\color{gray}C14} & {\color{gray}C15} & {\color{gray}C16}\\
\hline
\rowcolor{BenchBaseline}[0.5pt][0.5pt]\shortstack[l]{{\color{gray}R1}\\\textbf{minimap2}} & \textbf{Baseline} & \hspace*{1pt}\begin{tabular}[t]{@{}c@{}}\rule{0pt}{12pt}39.7s\\[-.2pt]{\fontsize{5.8}{6.4}\selectfont \textcolor{BrickRed}{$ \blacktriangledown\,2.2\times $}}\\[-.2pt]36.8s\\[-.2pt]{\fontsize{5.8}{6.4}\selectfont \textcolor{BrickRed}{$ \blacktriangledown\,2.2\times $}}\end{tabular} & \hspace*{1pt}\begin{tabular}[t]{@{}c@{}}\rule{0pt}{12pt}25.2s\\[-.2pt]{\fontsize{5.8}{6.4}\selectfont \textcolor{BrickRed}{$ \blacktriangledown\,1.5\times $}}\\[-.2pt]22.7s\\[-.2pt]{\fontsize{5.8}{6.4}\selectfont \textcolor{BrickRed}{$ \blacktriangledown\,1.5\times $}}\end{tabular} & \hspace*{1pt}\begin{tabular}[t]{@{}c@{}}\rule{0pt}{12pt}1m\,56.9s\\[-.2pt]{\fontsize{5.8}{6.4}\selectfont \textcolor{BrickRed}{$ \blacktriangledown\,29.1\times $}}\\[-.2pt]1m\,53.2s\\[-.2pt]{\fontsize{5.8}{6.4}\selectfont \textcolor{BrickRed}{$ \blacktriangledown\,32.9\times $}}\end{tabular} & \hspace*{1pt}\begin{tabular}[t]{@{}c@{}}\rule{0pt}{12pt}24.9s\\[-.2pt]{\fontsize{5.8}{6.4}\selectfont \textcolor{BrickRed}{$ \blacktriangledown\,1.6\times $}}\\[-.2pt]22.2s\\[-.2pt]{\fontsize{5.8}{6.4}\selectfont \textcolor{BrickRed}{$ \blacktriangledown\,1.6\times $}}\end{tabular} & \hspace*{1pt}\begin{tabular}[t]{@{}c@{}}\rule{0pt}{12pt}28.6s\\[-.2pt]{\fontsize{5.8}{6.4}\selectfont \textcolor{BrickRed}{$ \blacktriangledown\,11.5\times $}}\\[-.2pt]25.8s\\[-.2pt]{\fontsize{5.8}{6.4}\selectfont \textcolor{BrickRed}{$ \blacktriangledown\,15.6\times $}}\end{tabular} & \hspace*{1pt}\begin{tabular}[t]{@{}c@{}}\rule{0pt}{12pt}15.3s\\[-.2pt]{\fontsize{5.8}{6.4}\selectfont \textcolor{BrickRed}{$ \blacktriangledown\,2.8\times $}}\\[-.2pt]12.7s\\[-.2pt]{\fontsize{5.8}{6.4}\selectfont \textcolor{BrickRed}{$ \blacktriangledown\,3.3\times $}}\end{tabular} & \hspace*{1pt}\begin{tabular}[t]{@{}c@{}}\rule{0pt}{12pt}29.3s\\[-.2pt]{\fontsize{5.8}{6.4}\selectfont \textcolor{BrickRed}{$ \blacktriangledown\,15.0\times $}}\\[-.2pt]26.2s\\[-.2pt]{\fontsize{5.8}{6.4}\selectfont \textcolor{BrickRed}{$ \blacktriangledown\,22.9\times $}}\end{tabular} & \hspace*{1pt}\begin{tabular}[t]{@{}c@{}}\rule{0pt}{12pt}8.7s\\[-.2pt]{\fontsize{5.8}{6.4}\selectfont \textcolor{BrickRed}{$ \blacktriangledown\,1.7\times $}}\\[-.2pt]8.1s\\[-.2pt]{\fontsize{5.8}{6.4}\selectfont \textcolor{BrickRed}{$ \blacktriangledown\,2.3\times $}}\end{tabular} & \hspace*{1pt}\begin{tabular}[t]{@{}c@{}}\rule{0pt}{12pt}2m\,52.4s\\[-.2pt]{\fontsize{5.8}{6.4}\selectfont \textcolor{ForestGreen}{$ \blacktriangle\,19\% $}}\\[-.2pt]2m\,48.5s\\[-.2pt]{\fontsize{5.8}{6.4}\selectfont \textcolor{ForestGreen}{$ \blacktriangle\,15\% $}}\end{tabular} & \hspace*{1pt}\begin{tabular}[t]{@{}c@{}}\rule{0pt}{12pt}2m\,17.1s\\[-.2pt]{\fontsize{5.8}{6.4}\selectfont \textcolor{ForestGreen}{$ \blacktriangle\,3.9\times $}}\\[-.2pt]2m\,15.9s\\[-.2pt]{\fontsize{5.8}{6.4}\selectfont \textcolor{ForestGreen}{$ \blacktriangle\,1.7\times $}}\end{tabular} & \hspace*{1pt}\begin{tabular}[t]{@{}c@{}}\rule{0pt}{12pt}2m\,55s\\[-.2pt]{\fontsize{5.8}{6.4}\selectfont \textcolor{BrickRed}{$ \blacktriangledown\,6\% $}}\\[-.2pt]2m\,51s\\[-.2pt]{\fontsize{5.8}{6.4}\selectfont \textcolor{BrickRed}{$ \blacktriangledown\,12\% $}}\end{tabular} & \hspace*{1pt}\begin{tabular}[t]{@{}c@{}}\rule{0pt}{12pt}55.9s\\[-.2pt]{\fontsize{5.8}{6.4}\selectfont \textcolor{ForestGreen}{$ \blacktriangle\,8.6\times $}}\\[-.2pt]55.1s\\[-.2pt]{\fontsize{5.8}{6.4}\selectfont \textcolor{ForestGreen}{$ \blacktriangle\,3.1\times $}}\end{tabular} & \hspace*{1pt}\begin{tabular}[t]{@{}c@{}}\rule{0pt}{12pt}43.3s\\[-.2pt]{\fontsize{5.8}{6.4}\selectfont \textcolor{ForestGreen}{$ \blacktriangle\,2.6\times $}}\\[-.2pt]42.5s\\[-.2pt]{\fontsize{5.8}{6.4}\selectfont \textcolor{ForestGreen}{$ \blacktriangle\,2.4\times $}}\end{tabular} & \hspace*{1pt}\begin{tabular}[t]{@{}c@{}}\rule{0pt}{12pt}38.9s\\[-.2pt]{\fontsize{5.8}{6.4}\selectfont \textcolor{ForestGreen}{$ \blacktriangle\,8.0\times $}}\\[-.2pt]38.2s\\[-.2pt]{\fontsize{5.8}{6.4}\selectfont \textcolor{ForestGreen}{$ \blacktriangle\,2.8\times $}}\end{tabular} & \hspace*{1pt}\begin{tabular}[t]{@{}c@{}}\rule{0pt}{12pt}53.1s\\[-.2pt]{\fontsize{5.8}{6.4}\selectfont \textcolor{ForestGreen}{$ \blacktriangle\,1.6\times $}}\\[-.2pt]48.2s\\[-.2pt]{\fontsize{5.8}{6.4}\selectfont \textcolor{ForestGreen}{$ \blacktriangle\,1.5\times $}}\end{tabular} & \hspace*{1pt}\begin{tabular}[t]{@{}c@{}}\rule{0pt}{12pt}28.7s\\[-.2pt]{\fontsize{5.8}{6.4}\selectfont \textcolor{ForestGreen}{$ \blacktriangle\,9.8\times $}}\\[-.2pt]28.2s\\[-.2pt]{\fontsize{5.8}{6.4}\selectfont \textcolor{ForestGreen}{$ \blacktriangle\,2.9\times $}}\end{tabular}\\[3pt]\hline
\BenchmarkSourceLabel{bench:local:row:mpf-simdref-simd-ksw2}{R41}\shortstack[l]{{\color{gray}R41}\\\textbf{BioseqDB}} & \begin{tabular}[t]{@{}l@{}}\rule{0pt}{12pt}MPF / SIMD-ref\\SIMD-DP / KSW2\end{tabular} & \textcolor{gray}{\ensuremath{\cdot}} & \textcolor{gray}{\ensuremath{\cdot}} & \textcolor{gray}{\ensuremath{\cdot}} & \textcolor{gray}{\ensuremath{\cdot}} & \textcolor{gray}{\ensuremath{\cdot}} & \hspace*{1pt}\begin{tabular}[t]{@{}c@{}}\rule{0pt}{12pt}10.9s\\[-.2pt]{\fontsize{5.8}{6.4}\selectfont \textcolor{ForestGreen}{$ \blacktriangle\,29\% $}}\\[-.2pt]9.3s\\[-.2pt]{\fontsize{5.8}{6.4}\selectfont \textcolor{ForestGreen}{$ \blacktriangle\,27\% $}}\end{tabular} & \hspace*{1pt}\begin{tabular}[t]{@{}c@{}}\rule{0pt}{12pt}1m\,44.7s\\[-.2pt]{\fontsize{5.8}{6.4}\selectfont \textcolor{BrickRed}{$ \blacktriangledown\,3.6\times $}}\\[-.2pt]1m\,43.3s\\[-.2pt]{\fontsize{5.8}{6.4}\selectfont \textcolor{BrickRed}{$ \blacktriangledown\,3.9\times $}}\end{tabular} & \hspace*{1pt}\begin{tabular}[t]{@{}c@{}}\rule{0pt}{12pt}1m\,37.2s\\[-.2pt]{\fontsize{5.8}{6.4}\selectfont \textcolor{BrickRed}{$ \blacktriangledown\,11.2\times $}}\\[-.2pt]1m\,35.9s\\[-.2pt]{\fontsize{5.8}{6.4}\selectfont \textcolor{BrickRed}{$ \blacktriangledown\,11.8\times $}}\end{tabular} & \textcolor{gray}{\ensuremath{\cdot}} & \textcolor{gray}{\ensuremath{\cdot}} & \textcolor{gray}{\ensuremath{\cdot}} & \textcolor{gray}{\ensuremath{\cdot}} & \textcolor{gray}{\ensuremath{\cdot}} & \textcolor{gray}{\ensuremath{\cdot}} & \textcolor{gray}{\ensuremath{\cdot}} & \textcolor{gray}{\ensuremath{\cdot}}\\[3pt]\hline
\BenchmarkSourceLabel{bench:local:row:mpf-twostack-scalar-block}{R42}\shortstack[l]{{\color{gray}R42}} & \begin{tabular}[t]{@{}l@{}}\rule{0pt}{12pt}MPF / Two-stack\\Scalar-DP / Block\end{tabular} & \textcolor{gray}{\ensuremath{\cdot}} & \textcolor{gray}{\ensuremath{\cdot}} & \textcolor{gray}{\ensuremath{\cdot}} & \textcolor{gray}{\ensuremath{\cdot}} & \textcolor{gray}{\ensuremath{\cdot}} & \hspace*{1pt}\begin{tabular}[t]{@{}c@{}}\rule{0pt}{12pt}10.9s\\[-.2pt]{\fontsize{5.8}{6.4}\selectfont \textcolor{ForestGreen}{$ \blacktriangle\,28\% $}}\\[-.2pt]9.7s\\[-.2pt]{\fontsize{5.8}{6.4}\selectfont \textcolor{ForestGreen}{$ \blacktriangle\,24\% $}}\end{tabular} & \hspace*{1pt}\begin{tabular}[t]{@{}c@{}}\rule{0pt}{12pt}1m\,21s\\[-.2pt]{\fontsize{5.8}{6.4}\selectfont \textcolor{BrickRed}{$ \blacktriangledown\,2.8\times $}}\\[-.2pt]1m\,20.6s\\[-.2pt]{\fontsize{5.8}{6.4}\selectfont \textcolor{BrickRed}{$ \blacktriangledown\,3.1\times $}}\end{tabular} & \hspace*{1pt}\begin{tabular}[t]{@{}c@{}}\rule{0pt}{12pt}57.4s\\[-.2pt]{\fontsize{5.8}{6.4}\selectfont \textcolor{BrickRed}{$ \blacktriangledown\,6.6\times $}}\\[-.2pt]56.7s\\[-.2pt]{\fontsize{5.8}{6.4}\selectfont \textcolor{BrickRed}{$ \blacktriangledown\,7.0\times $}}\end{tabular} & \textcolor{gray}{\ensuremath{\cdot}} & \textcolor{gray}{\ensuremath{\cdot}} & \textcolor{gray}{\ensuremath{\cdot}} & \textcolor{gray}{\ensuremath{\cdot}} & \textcolor{gray}{\ensuremath{\cdot}} & \textcolor{gray}{\ensuremath{\cdot}} & \textcolor{gray}{\ensuremath{\cdot}} & \textcolor{gray}{\ensuremath{\cdot}}\\[3pt]\hline
\BenchmarkSourceLabel{bench:local:row:mpf-twostack-scalar-ksw2}{R43}\shortstack[l]{{\color{gray}R43}} & \begin{tabular}[t]{@{}l@{}}\rule{0pt}{12pt}MPF / Two-stack\\Scalar-DP / KSW2\end{tabular} & \textcolor{gray}{\ensuremath{\cdot}} & \textcolor{gray}{\ensuremath{\cdot}} & \textcolor{gray}{\ensuremath{\cdot}} & \textcolor{gray}{\ensuremath{\cdot}} & \textcolor{gray}{\ensuremath{\cdot}} & \hspace*{1pt}\begin{tabular}[t]{@{}c@{}}\rule{0pt}{12pt}11.4s\\[-.2pt]{\fontsize{5.8}{6.4}\selectfont \textcolor{ForestGreen}{$ \blacktriangle\,25\% $}}\\[-.2pt]9.9s\\[-.2pt]{\fontsize{5.8}{6.4}\selectfont \textcolor{ForestGreen}{$ \blacktriangle\,22\% $}}\end{tabular} & \hspace*{1pt}\begin{tabular}[t]{@{}c@{}}\rule{0pt}{12pt}1m\,21.1s\\[-.2pt]{\fontsize{5.8}{6.4}\selectfont \textcolor{BrickRed}{$ \blacktriangledown\,2.8\times $}}\\[-.2pt]1m\,19.3s\\[-.2pt]{\fontsize{5.8}{6.4}\selectfont \textcolor{BrickRed}{$ \blacktriangledown\,3.0\times $}}\end{tabular} & \hspace*{1pt}\begin{tabular}[t]{@{}c@{}}\rule{0pt}{12pt}58.7s\\[-.2pt]{\fontsize{5.8}{6.4}\selectfont \textcolor{BrickRed}{$ \blacktriangledown\,6.8\times $}}\\[-.2pt]56.2s\\[-.2pt]{\fontsize{5.8}{6.4}\selectfont \textcolor{BrickRed}{$ \blacktriangledown\,6.9\times $}}\end{tabular} & \textcolor{gray}{\ensuremath{\cdot}} & \textcolor{gray}{\ensuremath{\cdot}} & \textcolor{gray}{\ensuremath{\cdot}} & \textcolor{gray}{\ensuremath{\cdot}} & \textcolor{gray}{\ensuremath{\cdot}} & \textcolor{gray}{\ensuremath{\cdot}} & \textcolor{gray}{\ensuremath{\cdot}} & \textcolor{gray}{\ensuremath{\cdot}}\\[3pt]\hline
\BenchmarkSourceLabel{bench:local:row:mpf-scalar-rmq-ksw2}{R44}\shortstack[l]{{\color{gray}R44}} & \begin{tabular}[t]{@{}l@{}}\rule{0pt}{12pt}MPF / Scalar\\RMQ / KSW2\end{tabular} & \textcolor{gray}{\ensuremath{\cdot}} & \textcolor{gray}{\ensuremath{\cdot}} & \textcolor{gray}{\ensuremath{\cdot}} & \textcolor{gray}{\ensuremath{\cdot}} & \textcolor{gray}{\ensuremath{\cdot}} & \hspace*{1pt}\begin{tabular}[t]{@{}c@{}}\rule{0pt}{12pt}12s\\[-.2pt]{\fontsize{5.8}{6.4}\selectfont \textcolor{ForestGreen}{$ \blacktriangle\,22\% $}}\\[-.2pt]10.5s\\[-.2pt]{\fontsize{5.8}{6.4}\selectfont \textcolor{ForestGreen}{$ \blacktriangle\,17\% $}}\end{tabular} & \hspace*{1pt}\begin{tabular}[t]{@{}c@{}}\rule{0pt}{12pt}1m\,44.5s\\[-.2pt]{\fontsize{5.8}{6.4}\selectfont \textcolor{BrickRed}{$ \blacktriangledown\,3.6\times $}}\\[-.2pt]1m\,43.2s\\[-.2pt]{\fontsize{5.8}{6.4}\selectfont \textcolor{BrickRed}{$ \blacktriangledown\,3.9\times $}}\end{tabular} & \hspace*{1pt}\begin{tabular}[t]{@{}c@{}}\rule{0pt}{12pt}3m\,0.8s\\[-.2pt]{\fontsize{5.8}{6.4}\selectfont \textcolor{BrickRed}{$ \blacktriangledown\,20.8\times $}}\\[-.2pt]2m\,59.5s\\[-.2pt]{\fontsize{5.8}{6.4}\selectfont \textcolor{BrickRed}{$ \blacktriangledown\,22.1\times $}}\end{tabular} & \textcolor{gray}{\ensuremath{\cdot}} & \textcolor{gray}{\ensuremath{\cdot}} & \textcolor{gray}{\ensuremath{\cdot}} & \textcolor{gray}{\ensuremath{\cdot}} & \textcolor{gray}{\ensuremath{\cdot}} & \textcolor{gray}{\ensuremath{\cdot}} & \textcolor{gray}{\ensuremath{\cdot}} & \textcolor{gray}{\ensuremath{\cdot}}\\[3pt]\hline
\BenchmarkSourceLabel{bench:local:row:mpf-twostack-rmq-ksw2}{R45}\shortstack[l]{{\color{gray}R45}} & \begin{tabular}[t]{@{}l@{}}\rule{0pt}{12pt}MPF / Two-stack\\RMQ / KSW2\end{tabular} & \textcolor{gray}{\ensuremath{\cdot}} & \textcolor{gray}{\ensuremath{\cdot}} & \textcolor{gray}{\ensuremath{\cdot}} & \textcolor{gray}{\ensuremath{\cdot}} & \textcolor{gray}{\ensuremath{\cdot}} & \hspace*{1pt}\begin{tabular}[t]{@{}c@{}}\rule{0pt}{12pt}12.4s\\[-.2pt]{\fontsize{5.8}{6.4}\selectfont \textcolor{ForestGreen}{$ \blacktriangle\,19\% $}}\\[-.2pt]10.8s\\[-.2pt]{\fontsize{5.8}{6.4}\selectfont \textcolor{ForestGreen}{$ \blacktriangle\,15\% $}}\end{tabular} & \hspace*{1pt}\begin{tabular}[t]{@{}c@{}}\rule{0pt}{12pt}1m\,12.2s\\[-.2pt]{\fontsize{5.8}{6.4}\selectfont \textcolor{BrickRed}{$ \blacktriangledown\,2.5\times $}}\\[-.2pt]1m\,10.5s\\[-.2pt]{\fontsize{5.8}{6.4}\selectfont \textcolor{BrickRed}{$ \blacktriangledown\,2.7\times $}}\end{tabular} & \hspace*{1pt}\begin{tabular}[t]{@{}c@{}}\rule{0pt}{12pt}59.7s\\[-.2pt]{\fontsize{5.8}{6.4}\selectfont \textcolor{BrickRed}{$ \blacktriangledown\,6.9\times $}}\\[-.2pt]57.2s\\[-.2pt]{\fontsize{5.8}{6.4}\selectfont \textcolor{BrickRed}{$ \blacktriangledown\,7.0\times $}}\end{tabular} & \textcolor{gray}{\ensuremath{\cdot}} & \textcolor{gray}{\ensuremath{\cdot}} & \textcolor{gray}{\ensuremath{\cdot}} & \textcolor{gray}{\ensuremath{\cdot}} & \textcolor{gray}{\ensuremath{\cdot}} & \textcolor{gray}{\ensuremath{\cdot}} & \textcolor{gray}{\ensuremath{\cdot}} & \textcolor{gray}{\ensuremath{\cdot}}\\[3pt]\hline
\BenchmarkSourceLabel{bench:local:row:mpf-twostack-rmq-block}{R46}\shortstack[l]{{\color{gray}R46}} & \begin{tabular}[t]{@{}l@{}}\rule{0pt}{12pt}MPF / Two-stack\\RMQ / Block\end{tabular} & \textcolor{gray}{\ensuremath{\cdot}} & \textcolor{gray}{\ensuremath{\cdot}} & \textcolor{gray}{\ensuremath{\cdot}} & \textcolor{gray}{\ensuremath{\cdot}} & \textcolor{gray}{\ensuremath{\cdot}} & \hspace*{1pt}\begin{tabular}[t]{@{}c@{}}\rule{0pt}{12pt}12.7s\\[-.2pt]{\fontsize{5.8}{6.4}\selectfont \textcolor{ForestGreen}{$ \blacktriangle\,17\% $}}\\[-.2pt]11.5s\\[-.2pt]{\fontsize{5.8}{6.4}\selectfont \textcolor{ForestGreen}{$ \blacktriangle\,10\% $}}\end{tabular} & \hspace*{1pt}\begin{tabular}[t]{@{}c@{}}\rule{0pt}{12pt}1m\,35.2s\\[-.2pt]{\fontsize{5.8}{6.4}\selectfont \textcolor{BrickRed}{$ \blacktriangledown\,3.2\times $}}\\[-.2pt]1m\,34.5s\\[-.2pt]{\fontsize{5.8}{6.4}\selectfont \textcolor{BrickRed}{$ \blacktriangledown\,3.6\times $}}\end{tabular} & \hspace*{1pt}\begin{tabular}[t]{@{}c@{}}\rule{0pt}{12pt}58s\\[-.2pt]{\fontsize{5.8}{6.4}\selectfont \textcolor{BrickRed}{$ \blacktriangledown\,6.7\times $}}\\[-.2pt]57.2s\\[-.2pt]{\fontsize{5.8}{6.4}\selectfont \textcolor{BrickRed}{$ \blacktriangledown\,7.0\times $}}\end{tabular} & \textcolor{gray}{\ensuremath{\cdot}} & \textcolor{gray}{\ensuremath{\cdot}} & \textcolor{gray}{\ensuremath{\cdot}} & \textcolor{gray}{\ensuremath{\cdot}} & \textcolor{gray}{\ensuremath{\cdot}} & \textcolor{gray}{\ensuremath{\cdot}} & \textcolor{gray}{\ensuremath{\cdot}} & \textcolor{gray}{\ensuremath{\cdot}}\\[3pt]\hline
\BenchmarkSourceLabel{bench:local:row:mps-twostack-vulkan-block}{R47}\shortstack[l]{{\color{gray}R47}} & \begin{tabular}[t]{@{}l@{}}\rule{0pt}{12pt}MPS / Two-stack\\Vulkan-DP / Block\end{tabular} & \textcolor{gray}{\ensuremath{\cdot}} & \textcolor{gray}{\ensuremath{\cdot}} & \textcolor{gray}{\ensuremath{\cdot}} & \textcolor{gray}{\ensuremath{\cdot}} & \hspace*{1pt}\begin{tabular}[t]{@{}c@{}}\rule{0pt}{12pt}26.2s\\[-.2pt]{\fontsize{5.8}{6.4}\selectfont \textcolor{ForestGreen}{$ \blacktriangle\,8\% $}}\\[-.2pt]19s\\[-.2pt]{\fontsize{5.8}{6.4}\selectfont \textcolor{ForestGreen}{$ \blacktriangle\,26\% $}}\end{tabular} & \hspace*{1pt}\begin{tabular}[t]{@{}c@{}}\rule{0pt}{12pt}3m\,49.5s\\[-.2pt]{\fontsize{5.8}{6.4}\selectfont \textcolor{BrickRed}{$ \blacktriangledown\,15.0\times $}}\\[-.2pt]33.7s\\[-.2pt]{\fontsize{5.8}{6.4}\selectfont \textcolor{BrickRed}{$ \blacktriangledown\,2.6\times $}}\end{tabular} & \hspace*{1pt}\begin{tabular}[t]{@{}c@{}}\rule{0pt}{12pt}13.1s\\[-.2pt]{\fontsize{5.8}{6.4}\selectfont \textcolor{ForestGreen}{$ \blacktriangle\,2.2\times $}}\\[-.2pt]5.8s\\[-.2pt]{\fontsize{5.8}{6.4}\selectfont \textcolor{ForestGreen}{$ \blacktriangle\,4.5\times $}}\end{tabular} & \hspace*{1pt}\begin{tabular}[t]{@{}c@{}}\rule{0pt}{12pt}3m\,23.1s\\[-.2pt]{\fontsize{5.8}{6.4}\selectfont \textcolor{BrickRed}{$ \blacktriangledown\,23.4\times $}}\\[-.2pt]9.9s\\[-.2pt]{\fontsize{5.8}{6.4}\selectfont \textcolor{BrickRed}{$ \blacktriangledown\,22\% $}}\end{tabular} & \textcolor{gray}{\ensuremath{\cdot}} & \textcolor{gray}{\ensuremath{\cdot}} & \textcolor{gray}{\ensuremath{\cdot}} & \textcolor{gray}{\ensuremath{\cdot}} & \textcolor{gray}{\ensuremath{\cdot}} & \textcolor{gray}{\ensuremath{\cdot}} & \textcolor{gray}{\ensuremath{\cdot}} & \textcolor{gray}{\ensuremath{\cdot}}\\[3pt]\hline
\end{tabular}\par}
\smallskip{\footnotesize\noindent Each cell: A and B times, $\cdot$: no measurement. ---\textsuperscript{g}: GPU unavailable. n/a: no comparison.\par
\noindent Blue: one lowest A and one lowest B measurement for each column for each configuration.\par
\noindent Gray: minimap2. Yellow: my preferred MPF / Bounded / SIMD-DP / Block configuration.\par
\noindent MPF/MPS: minimizers per fragment/sequence. SIMD-ref: naive SIMD scan. Bounded: bounded SIMD algorithm.\par
}
\clearpage
\noindent{\large\bfseries Table \thetable: (A): Local machine\hfill 8/12}\par\smallskip
{\fontsize{7.1}{8.1}\selectfont\setlength{\tabcolsep}{1pt}\setlength{\fboxsep}{2pt}\renewcommand{\arraystretch}{1.12}
\noindent\begin{tabular}{@{}p{14mm}p{27mm}|p{13.55mm}|p{13.55mm}|p{13.55mm}|p{13.55mm}|p{13.55mm}|p{13.55mm}|p{13.55mm}|p{13.55mm}|p{13.55mm}|p{13.55mm}|p{13.55mm}|p{13.55mm}|p{13.55mm}|p{13.55mm}|p{13.55mm}|p{13.55mm}|@{}}\hline
Engine & Execution & \multicolumn{8}{c|}{\textbf{Index}} & \multicolumn{8}{c|}{\textbf{No index}}\\
 &  & \multicolumn{4}{c|}{\textbf{4 Gbp}} & \multicolumn{4}{c|}{\textbf{2 Gbp}} & \multicolumn{4}{c|}{\textbf{4 Gbp}} & \multicolumn{4}{c|}{\textbf{2 Gbp}}\\
 &  & \multicolumn{2}{c|}{\textbf{k=15, w=10}} & \multicolumn{2}{c|}{\textbf{k=20, w=15}} & \multicolumn{2}{c|}{\textbf{k=15, w=10}} & \multicolumn{2}{c|}{\textbf{k=20, w=15}} & \multicolumn{2}{c|}{\textbf{k=15, w=10}} & \multicolumn{2}{c|}{\textbf{k=20, w=15}} & \multicolumn{2}{c|}{\textbf{k=15, w=10}} & \multicolumn{2}{c|}{\textbf{k=20, w=15}}\\
 &  & \multicolumn{1}{c|}{\textbf{DNA4}} & \multicolumn{1}{c|}{\textbf{DNA5}} & \multicolumn{1}{c|}{\textbf{DNA4}} & \multicolumn{1}{c|}{\textbf{DNA5}} & \multicolumn{1}{c|}{\textbf{DNA4}} & \multicolumn{1}{c|}{\textbf{DNA5}} & \multicolumn{1}{c|}{\textbf{DNA4}} & \multicolumn{1}{c|}{\textbf{DNA5}} & \multicolumn{1}{c|}{\textbf{DNA4}} & \multicolumn{1}{c|}{\textbf{DNA5}} & \multicolumn{1}{c|}{\textbf{DNA4}} & \multicolumn{1}{c|}{\textbf{DNA5}} & \multicolumn{1}{c|}{\textbf{DNA4}} & \multicolumn{1}{c|}{\textbf{DNA5}} & \multicolumn{1}{c|}{\textbf{DNA4}} & \multicolumn{1}{c|}{\textbf{DNA5}}\\
 & & {\color{gray}C1} & {\color{gray}C2} & {\color{gray}C3} & {\color{gray}C4} & {\color{gray}C5} & {\color{gray}C6} & {\color{gray}C7} & {\color{gray}C8} & {\color{gray}C9} & {\color{gray}C10} & {\color{gray}C11} & {\color{gray}C12} & {\color{gray}C13} & {\color{gray}C14} & {\color{gray}C15} & {\color{gray}C16}\\
\hline
\rowcolor{BenchBaseline}[0.5pt][0.5pt]\shortstack[l]{{\color{gray}R1}\\\textbf{minimap2}} & \textbf{Baseline} & \hspace*{1pt}\begin{tabular}[t]{@{}c@{}}\rule{0pt}{12pt}39.7s\\[-.2pt]{\fontsize{5.8}{6.4}\selectfont \textcolor{BrickRed}{$ \blacktriangledown\,2.2\times $}}\\[-.2pt]36.8s\\[-.2pt]{\fontsize{5.8}{6.4}\selectfont \textcolor{BrickRed}{$ \blacktriangledown\,2.2\times $}}\end{tabular} & \hspace*{1pt}\begin{tabular}[t]{@{}c@{}}\rule{0pt}{12pt}25.2s\\[-.2pt]{\fontsize{5.8}{6.4}\selectfont \textcolor{BrickRed}{$ \blacktriangledown\,1.5\times $}}\\[-.2pt]22.7s\\[-.2pt]{\fontsize{5.8}{6.4}\selectfont \textcolor{BrickRed}{$ \blacktriangledown\,1.5\times $}}\end{tabular} & \hspace*{1pt}\begin{tabular}[t]{@{}c@{}}\rule{0pt}{12pt}1m\,56.9s\\[-.2pt]{\fontsize{5.8}{6.4}\selectfont \textcolor{BrickRed}{$ \blacktriangledown\,29.1\times $}}\\[-.2pt]1m\,53.2s\\[-.2pt]{\fontsize{5.8}{6.4}\selectfont \textcolor{BrickRed}{$ \blacktriangledown\,32.9\times $}}\end{tabular} & \hspace*{1pt}\begin{tabular}[t]{@{}c@{}}\rule{0pt}{12pt}24.9s\\[-.2pt]{\fontsize{5.8}{6.4}\selectfont \textcolor{BrickRed}{$ \blacktriangledown\,1.6\times $}}\\[-.2pt]22.2s\\[-.2pt]{\fontsize{5.8}{6.4}\selectfont \textcolor{BrickRed}{$ \blacktriangledown\,1.6\times $}}\end{tabular} & \hspace*{1pt}\begin{tabular}[t]{@{}c@{}}\rule{0pt}{12pt}28.6s\\[-.2pt]{\fontsize{5.8}{6.4}\selectfont \textcolor{BrickRed}{$ \blacktriangledown\,11.5\times $}}\\[-.2pt]25.8s\\[-.2pt]{\fontsize{5.8}{6.4}\selectfont \textcolor{BrickRed}{$ \blacktriangledown\,15.6\times $}}\end{tabular} & \hspace*{1pt}\begin{tabular}[t]{@{}c@{}}\rule{0pt}{12pt}15.3s\\[-.2pt]{\fontsize{5.8}{6.4}\selectfont \textcolor{BrickRed}{$ \blacktriangledown\,2.8\times $}}\\[-.2pt]12.7s\\[-.2pt]{\fontsize{5.8}{6.4}\selectfont \textcolor{BrickRed}{$ \blacktriangledown\,3.3\times $}}\end{tabular} & \hspace*{1pt}\begin{tabular}[t]{@{}c@{}}\rule{0pt}{12pt}29.3s\\[-.2pt]{\fontsize{5.8}{6.4}\selectfont \textcolor{BrickRed}{$ \blacktriangledown\,15.0\times $}}\\[-.2pt]26.2s\\[-.2pt]{\fontsize{5.8}{6.4}\selectfont \textcolor{BrickRed}{$ \blacktriangledown\,22.9\times $}}\end{tabular} & \hspace*{1pt}\begin{tabular}[t]{@{}c@{}}\rule{0pt}{12pt}8.7s\\[-.2pt]{\fontsize{5.8}{6.4}\selectfont \textcolor{BrickRed}{$ \blacktriangledown\,1.7\times $}}\\[-.2pt]8.1s\\[-.2pt]{\fontsize{5.8}{6.4}\selectfont \textcolor{BrickRed}{$ \blacktriangledown\,2.3\times $}}\end{tabular} & \hspace*{1pt}\begin{tabular}[t]{@{}c@{}}\rule{0pt}{12pt}2m\,52.4s\\[-.2pt]{\fontsize{5.8}{6.4}\selectfont \textcolor{ForestGreen}{$ \blacktriangle\,19\% $}}\\[-.2pt]2m\,48.5s\\[-.2pt]{\fontsize{5.8}{6.4}\selectfont \textcolor{ForestGreen}{$ \blacktriangle\,15\% $}}\end{tabular} & \hspace*{1pt}\begin{tabular}[t]{@{}c@{}}\rule{0pt}{12pt}2m\,17.1s\\[-.2pt]{\fontsize{5.8}{6.4}\selectfont \textcolor{ForestGreen}{$ \blacktriangle\,3.9\times $}}\\[-.2pt]2m\,15.9s\\[-.2pt]{\fontsize{5.8}{6.4}\selectfont \textcolor{ForestGreen}{$ \blacktriangle\,1.7\times $}}\end{tabular} & \hspace*{1pt}\begin{tabular}[t]{@{}c@{}}\rule{0pt}{12pt}2m\,55s\\[-.2pt]{\fontsize{5.8}{6.4}\selectfont \textcolor{BrickRed}{$ \blacktriangledown\,6\% $}}\\[-.2pt]2m\,51s\\[-.2pt]{\fontsize{5.8}{6.4}\selectfont \textcolor{BrickRed}{$ \blacktriangledown\,12\% $}}\end{tabular} & \hspace*{1pt}\begin{tabular}[t]{@{}c@{}}\rule{0pt}{12pt}55.9s\\[-.2pt]{\fontsize{5.8}{6.4}\selectfont \textcolor{ForestGreen}{$ \blacktriangle\,8.6\times $}}\\[-.2pt]55.1s\\[-.2pt]{\fontsize{5.8}{6.4}\selectfont \textcolor{ForestGreen}{$ \blacktriangle\,3.1\times $}}\end{tabular} & \hspace*{1pt}\begin{tabular}[t]{@{}c@{}}\rule{0pt}{12pt}43.3s\\[-.2pt]{\fontsize{5.8}{6.4}\selectfont \textcolor{ForestGreen}{$ \blacktriangle\,2.6\times $}}\\[-.2pt]42.5s\\[-.2pt]{\fontsize{5.8}{6.4}\selectfont \textcolor{ForestGreen}{$ \blacktriangle\,2.4\times $}}\end{tabular} & \hspace*{1pt}\begin{tabular}[t]{@{}c@{}}\rule{0pt}{12pt}38.9s\\[-.2pt]{\fontsize{5.8}{6.4}\selectfont \textcolor{ForestGreen}{$ \blacktriangle\,8.0\times $}}\\[-.2pt]38.2s\\[-.2pt]{\fontsize{5.8}{6.4}\selectfont \textcolor{ForestGreen}{$ \blacktriangle\,2.8\times $}}\end{tabular} & \hspace*{1pt}\begin{tabular}[t]{@{}c@{}}\rule{0pt}{12pt}53.1s\\[-.2pt]{\fontsize{5.8}{6.4}\selectfont \textcolor{ForestGreen}{$ \blacktriangle\,1.6\times $}}\\[-.2pt]48.2s\\[-.2pt]{\fontsize{5.8}{6.4}\selectfont \textcolor{ForestGreen}{$ \blacktriangle\,1.5\times $}}\end{tabular} & \hspace*{1pt}\begin{tabular}[t]{@{}c@{}}\rule{0pt}{12pt}28.7s\\[-.2pt]{\fontsize{5.8}{6.4}\selectfont \textcolor{ForestGreen}{$ \blacktriangle\,9.8\times $}}\\[-.2pt]28.2s\\[-.2pt]{\fontsize{5.8}{6.4}\selectfont \textcolor{ForestGreen}{$ \blacktriangle\,2.9\times $}}\end{tabular}\\[3pt]\hline
\BenchmarkSourceLabel{bench:local:row:mps-twostack-vulkan-ksw2}{R48}\shortstack[l]{{\color{gray}R48}\\\textbf{BioseqDB}} & \begin{tabular}[t]{@{}l@{}}\rule{0pt}{12pt}MPS / Two-stack\\Vulkan-DP / KSW2\end{tabular} & \textcolor{gray}{\ensuremath{\cdot}} & \textcolor{gray}{\ensuremath{\cdot}} & \textcolor{gray}{\ensuremath{\cdot}} & \textcolor{gray}{\ensuremath{\cdot}} & \hspace*{1pt}\begin{tabular}[t]{@{}c@{}}\rule{0pt}{12pt}27.6s\\[-.2pt]{\fontsize{5.8}{6.4}\selectfont \textcolor{gray}{$ \blacktriangleright\,4\% $}}\\[-.2pt]18.6s\\[-.2pt]{\fontsize{5.8}{6.4}\selectfont \textcolor{ForestGreen}{$ \blacktriangle\,28\% $}}\end{tabular} & \hspace*{1pt}\begin{tabular}[t]{@{}c@{}}\rule{0pt}{12pt}3m\,51.3s\\[-.2pt]{\fontsize{5.8}{6.4}\selectfont \textcolor{BrickRed}{$ \blacktriangledown\,15.1\times $}}\\[-.2pt]33.5s\\[-.2pt]{\fontsize{5.8}{6.4}\selectfont \textcolor{BrickRed}{$ \blacktriangledown\,2.6\times $}}\end{tabular} & \hspace*{1pt}\begin{tabular}[t]{@{}c@{}}\rule{0pt}{12pt}18.1s\\[-.2pt]{\fontsize{5.8}{6.4}\selectfont \textcolor{ForestGreen}{$ \blacktriangle\,1.6\times $}}\\[-.2pt]5.8s\\[-.2pt]{\fontsize{5.8}{6.4}\selectfont \textcolor{ForestGreen}{$ \blacktriangle\,4.5\times $}}\end{tabular} & \hspace*{1pt}\begin{tabular}[t]{@{}c@{}}\rule{0pt}{12pt}3m\,29.5s\\[-.2pt]{\fontsize{5.8}{6.4}\selectfont \textcolor{BrickRed}{$ \blacktriangledown\,24.1\times $}}\\[-.2pt]9.9s\\[-.2pt]{\fontsize{5.8}{6.4}\selectfont \textcolor{BrickRed}{$ \blacktriangledown\,21\% $}}\end{tabular} & \textcolor{gray}{\ensuremath{\cdot}} & \textcolor{gray}{\ensuremath{\cdot}} & \textcolor{gray}{\ensuremath{\cdot}} & \textcolor{gray}{\ensuremath{\cdot}} & \textcolor{gray}{\ensuremath{\cdot}} & \textcolor{gray}{\ensuremath{\cdot}} & \textcolor{gray}{\ensuremath{\cdot}} & \textcolor{gray}{\ensuremath{\cdot}}\\[3pt]\hline
\BenchmarkSourceLabel{bench:local:row:mps-bounded-vulkan-block}{R49}\shortstack[l]{{\color{gray}R49}} & \begin{tabular}[t]{@{}l@{}}\rule{0pt}{12pt}MPS / Bounded\\Vulkan-DP / Block\end{tabular} & \textcolor{gray}{\ensuremath{\cdot}} & \textcolor{gray}{\ensuremath{\cdot}} & \textcolor{gray}{\ensuremath{\cdot}} & \textcolor{gray}{\ensuremath{\cdot}} & \hspace*{1pt}\begin{tabular}[t]{@{}c@{}}\rule{0pt}{12pt}25.5s\\[-.2pt]{\fontsize{5.8}{6.4}\selectfont \textcolor{ForestGreen}{$ \blacktriangle\,11\% $}}\\[-.2pt]18.2s\\[-.2pt]{\fontsize{5.8}{6.4}\selectfont \textcolor{ForestGreen}{$ \blacktriangle\,29\% $}}\end{tabular} & \hspace*{1pt}\begin{tabular}[t]{@{}c@{}}\rule{0pt}{12pt}4m\,7.7s\\[-.2pt]{\fontsize{5.8}{6.4}\selectfont \textcolor{BrickRed}{$ \blacktriangledown\,16.2\times $}}\\[-.2pt]43.2s\\[-.2pt]{\fontsize{5.8}{6.4}\selectfont \textcolor{BrickRed}{$ \blacktriangledown\,3.4\times $}}\end{tabular} & \hspace*{1pt}\begin{tabular}[t]{@{}c@{}}\rule{0pt}{12pt}19.8s\\[-.2pt]{\fontsize{5.8}{6.4}\selectfont \textcolor{ForestGreen}{$ \blacktriangle\,1.5\times $}}\\[-.2pt]12.9s\\[-.2pt]{\fontsize{5.8}{6.4}\selectfont \textcolor{ForestGreen}{$ \blacktriangle\,2.0\times $}}\end{tabular} & \hspace*{1pt}\begin{tabular}[t]{@{}c@{}}\rule{0pt}{12pt}3m\,47.2s\\[-.2pt]{\fontsize{5.8}{6.4}\selectfont \textcolor{BrickRed}{$ \blacktriangledown\,26.2\times $}}\\[-.2pt]29s\\[-.2pt]{\fontsize{5.8}{6.4}\selectfont \textcolor{BrickRed}{$ \blacktriangledown\,3.6\times $}}\end{tabular} & \textcolor{gray}{\ensuremath{\cdot}} & \textcolor{gray}{\ensuremath{\cdot}} & \textcolor{gray}{\ensuremath{\cdot}} & \textcolor{gray}{\ensuremath{\cdot}} & \textcolor{gray}{\ensuremath{\cdot}} & \textcolor{gray}{\ensuremath{\cdot}} & \textcolor{gray}{\ensuremath{\cdot}} & \textcolor{gray}{\ensuremath{\cdot}}\\[3pt]\hline
\BenchmarkSourceLabel{bench:local:row:mps-scalar-vulkan-block}{R50}\shortstack[l]{{\color{gray}R50}} & \begin{tabular}[t]{@{}l@{}}\rule{0pt}{12pt}MPS / Scalar\\Vulkan-DP / Block\end{tabular} & \textcolor{gray}{\ensuremath{\cdot}} & \textcolor{gray}{\ensuremath{\cdot}} & \textcolor{gray}{\ensuremath{\cdot}} & \textcolor{gray}{\ensuremath{\cdot}} & \hspace*{1pt}\begin{tabular}[t]{@{}c@{}}\rule{0pt}{12pt}28.8s\\[-.2pt]{\fontsize{5.8}{6.4}\selectfont \textcolor{gray}{$ \blacktriangleright\,1\% $}}\\[-.2pt]18.6s\\[-.2pt]{\fontsize{5.8}{6.4}\selectfont \textcolor{ForestGreen}{$ \blacktriangle\,28\% $}}\end{tabular} & \hspace*{1pt}\begin{tabular}[t]{@{}c@{}}\rule{0pt}{12pt}4m\,7s\\[-.2pt]{\fontsize{5.8}{6.4}\selectfont \textcolor{BrickRed}{$ \blacktriangledown\,16.2\times $}}\\[-.2pt]44.1s\\[-.2pt]{\fontsize{5.8}{6.4}\selectfont \textcolor{BrickRed}{$ \blacktriangledown\,3.5\times $}}\end{tabular} & \hspace*{1pt}\begin{tabular}[t]{@{}c@{}}\rule{0pt}{12pt}20.9s\\[-.2pt]{\fontsize{5.8}{6.4}\selectfont \textcolor{ForestGreen}{$ \blacktriangle\,29\% $}}\\[-.2pt]13.8s\\[-.2pt]{\fontsize{5.8}{6.4}\selectfont \textcolor{ForestGreen}{$ \blacktriangle\,1.9\times $}}\end{tabular} & \hspace*{1pt}\begin{tabular}[t]{@{}c@{}}\rule{0pt}{12pt}3m\,50.9s\\[-.2pt]{\fontsize{5.8}{6.4}\selectfont \textcolor{BrickRed}{$ \blacktriangledown\,26.6\times $}}\\[-.2pt]31.3s\\[-.2pt]{\fontsize{5.8}{6.4}\selectfont \textcolor{BrickRed}{$ \blacktriangledown\,3.8\times $}}\end{tabular} & \textcolor{gray}{\ensuremath{\cdot}} & \textcolor{gray}{\ensuremath{\cdot}} & \textcolor{gray}{\ensuremath{\cdot}} & \textcolor{gray}{\ensuremath{\cdot}} & \textcolor{gray}{\ensuremath{\cdot}} & \textcolor{gray}{\ensuremath{\cdot}} & \textcolor{gray}{\ensuremath{\cdot}} & \textcolor{gray}{\ensuremath{\cdot}}\\[3pt]\hline
\BenchmarkSourceLabel{bench:local:row:mps-simdref-vulkan-block}{R51}\shortstack[l]{{\color{gray}R51}} & \begin{tabular}[t]{@{}l@{}}\rule{0pt}{12pt}MPS / SIMD-ref\\Vulkan-DP / Block\end{tabular} & \textcolor{gray}{\ensuremath{\cdot}} & \textcolor{gray}{\ensuremath{\cdot}} & \textcolor{gray}{\ensuremath{\cdot}} & \textcolor{gray}{\ensuremath{\cdot}} & \hspace*{1pt}\begin{tabular}[t]{@{}c@{}}\rule{0pt}{12pt}25.2s\\[-.2pt]{\fontsize{5.8}{6.4}\selectfont \textcolor{ForestGreen}{$ \blacktriangle\,12\% $}}\\[-.2pt]17.9s\\[-.2pt]{\fontsize{5.8}{6.4}\selectfont \textcolor{ForestGreen}{$ \blacktriangle\,1.4\times $}}\end{tabular} & \hspace*{1pt}\begin{tabular}[t]{@{}c@{}}\rule{0pt}{12pt}3m\,59.3s\\[-.2pt]{\fontsize{5.8}{6.4}\selectfont \textcolor{BrickRed}{$ \blacktriangledown\,15.7\times $}}\\[-.2pt]35.3s\\[-.2pt]{\fontsize{5.8}{6.4}\selectfont \textcolor{BrickRed}{$ \blacktriangledown\,2.8\times $}}\end{tabular} & \hspace*{1pt}\begin{tabular}[t]{@{}c@{}}\rule{0pt}{12pt}21.4s\\[-.2pt]{\fontsize{5.8}{6.4}\selectfont \textcolor{ForestGreen}{$ \blacktriangle\,27\% $}}\\[-.2pt]14.2s\\[-.2pt]{\fontsize{5.8}{6.4}\selectfont \textcolor{ForestGreen}{$ \blacktriangle\,1.8\times $}}\end{tabular} & \hspace*{1pt}\begin{tabular}[t]{@{}c@{}}\rule{0pt}{12pt}3m\,57.1s\\[-.2pt]{\fontsize{5.8}{6.4}\selectfont \textcolor{BrickRed}{$ \blacktriangledown\,27.3\times $}}\\[-.2pt]37.5s\\[-.2pt]{\fontsize{5.8}{6.4}\selectfont \textcolor{BrickRed}{$ \blacktriangledown\,4.6\times $}}\end{tabular} & \textcolor{gray}{\ensuremath{\cdot}} & \textcolor{gray}{\ensuremath{\cdot}} & \textcolor{gray}{\ensuremath{\cdot}} & \textcolor{gray}{\ensuremath{\cdot}} & \textcolor{gray}{\ensuremath{\cdot}} & \textcolor{gray}{\ensuremath{\cdot}} & \textcolor{gray}{\ensuremath{\cdot}} & \textcolor{gray}{\ensuremath{\cdot}}\\[3pt]\hline
\BenchmarkSourceLabel{bench:local:row:mps-bounded-vulkan-ksw2}{R52}\shortstack[l]{{\color{gray}R52}} & \begin{tabular}[t]{@{}l@{}}\rule{0pt}{12pt}MPS / Bounded\\Vulkan-DP / KSW2\end{tabular} & \textcolor{gray}{\ensuremath{\cdot}} & \textcolor{gray}{\ensuremath{\cdot}} & \textcolor{gray}{\ensuremath{\cdot}} & \textcolor{gray}{\ensuremath{\cdot}} & \hspace*{1pt}\begin{tabular}[t]{@{}c@{}}\rule{0pt}{12pt}28.4s\\[-.2pt]{\fontsize{5.8}{6.4}\selectfont \textcolor{gray}{$ \blacktriangleright\,0\% $}}\\[-.2pt]19.2s\\[-.2pt]{\fontsize{5.8}{6.4}\selectfont \textcolor{ForestGreen}{$ \blacktriangle\,26\% $}}\end{tabular} & \hspace*{1pt}\begin{tabular}[t]{@{}c@{}}\rule{0pt}{12pt}4m\,4.9s\\[-.2pt]{\fontsize{5.8}{6.4}\selectfont \textcolor{BrickRed}{$ \blacktriangledown\,16.0\times $}}\\[-.2pt]42.5s\\[-.2pt]{\fontsize{5.8}{6.4}\selectfont \textcolor{BrickRed}{$ \blacktriangledown\,3.3\times $}}\end{tabular} & \hspace*{1pt}\begin{tabular}[t]{@{}c@{}}\rule{0pt}{12pt}23s\\[-.2pt]{\fontsize{5.8}{6.4}\selectfont \textcolor{ForestGreen}{$ \blacktriangle\,21\% $}}\\[-.2pt]14s\\[-.2pt]{\fontsize{5.8}{6.4}\selectfont \textcolor{ForestGreen}{$ \blacktriangle\,1.9\times $}}\end{tabular} & \hspace*{1pt}\begin{tabular}[t]{@{}c@{}}\rule{0pt}{12pt}3m\,50.4s\\[-.2pt]{\fontsize{5.8}{6.4}\selectfont \textcolor{BrickRed}{$ \blacktriangledown\,26.5\times $}}\\[-.2pt]28.8s\\[-.2pt]{\fontsize{5.8}{6.4}\selectfont \textcolor{BrickRed}{$ \blacktriangledown\,3.5\times $}}\end{tabular} & \textcolor{gray}{\ensuremath{\cdot}} & \textcolor{gray}{\ensuremath{\cdot}} & \textcolor{gray}{\ensuremath{\cdot}} & \textcolor{gray}{\ensuremath{\cdot}} & \textcolor{gray}{\ensuremath{\cdot}} & \textcolor{gray}{\ensuremath{\cdot}} & \textcolor{gray}{\ensuremath{\cdot}} & \textcolor{gray}{\ensuremath{\cdot}}\\[3pt]\hline
\BenchmarkSourceLabel{bench:local:row:mps-simdref-vulkan-ksw2}{R53}\shortstack[l]{{\color{gray}R53}} & \begin{tabular}[t]{@{}l@{}}\rule{0pt}{12pt}MPS / SIMD-ref\\Vulkan-DP / KSW2\end{tabular} & \textcolor{gray}{\ensuremath{\cdot}} & \textcolor{gray}{\ensuremath{\cdot}} & \textcolor{gray}{\ensuremath{\cdot}} & \textcolor{gray}{\ensuremath{\cdot}} & \hspace*{1pt}\begin{tabular}[t]{@{}c@{}}\rule{0pt}{12pt}28.4s\\[-.2pt]{\fontsize{5.8}{6.4}\selectfont \textcolor{gray}{$ \blacktriangleright\,0\% $}}\\[-.2pt]18.2s\\[-.2pt]{\fontsize{5.8}{6.4}\selectfont \textcolor{ForestGreen}{$ \blacktriangle\,30\% $}}\end{tabular} & \hspace*{1pt}\begin{tabular}[t]{@{}c@{}}\rule{0pt}{12pt}4m\,1.2s\\[-.2pt]{\fontsize{5.8}{6.4}\selectfont \textcolor{BrickRed}{$ \blacktriangledown\,15.8\times $}}\\[-.2pt]35.5s\\[-.2pt]{\fontsize{5.8}{6.4}\selectfont \textcolor{BrickRed}{$ \blacktriangledown\,2.8\times $}}\end{tabular} & \hspace*{1pt}\begin{tabular}[t]{@{}c@{}}\rule{0pt}{12pt}23.8s\\[-.2pt]{\fontsize{5.8}{6.4}\selectfont \textcolor{ForestGreen}{$ \blacktriangle\,19\% $}}\\[-.2pt]14.5s\\[-.2pt]{\fontsize{5.8}{6.4}\selectfont \textcolor{ForestGreen}{$ \blacktriangle\,1.8\times $}}\end{tabular} & \hspace*{1pt}\begin{tabular}[t]{@{}c@{}}\rule{0pt}{12pt}3m\,57.5s\\[-.2pt]{\fontsize{5.8}{6.4}\selectfont \textcolor{BrickRed}{$ \blacktriangledown\,27.4\times $}}\\[-.2pt]37.7s\\[-.2pt]{\fontsize{5.8}{6.4}\selectfont \textcolor{BrickRed}{$ \blacktriangledown\,4.6\times $}}\end{tabular} & \textcolor{gray}{\ensuremath{\cdot}} & \textcolor{gray}{\ensuremath{\cdot}} & \textcolor{gray}{\ensuremath{\cdot}} & \textcolor{gray}{\ensuremath{\cdot}} & \textcolor{gray}{\ensuremath{\cdot}} & \textcolor{gray}{\ensuremath{\cdot}} & \textcolor{gray}{\ensuremath{\cdot}} & \textcolor{gray}{\ensuremath{\cdot}}\\[3pt]\hline
\end{tabular}\par}
\smallskip{\footnotesize\noindent Each cell: A and B times, $\cdot$: no measurement. ---\textsuperscript{g}: GPU unavailable. n/a: no comparison.\par
\noindent Blue: one lowest A and one lowest B measurement for each column for each configuration.\par
\noindent Gray: minimap2. Yellow: my preferred MPF / Bounded / SIMD-DP / Block configuration.\par
\noindent MPF/MPS: minimizers per fragment/sequence. SIMD-ref: naive SIMD scan. Bounded: bounded SIMD algorithm.\par
}
\clearpage
\noindent{\large\bfseries Table \thetable: (A): Local machine\hfill 9/12}\par\smallskip
{\fontsize{7.1}{8.1}\selectfont\setlength{\tabcolsep}{1pt}\setlength{\fboxsep}{2pt}\renewcommand{\arraystretch}{1.12}
\noindent\begin{tabular}{@{}p{14mm}p{27mm}|p{13.55mm}|p{13.55mm}|p{13.55mm}|p{13.55mm}|p{13.55mm}|p{13.55mm}|p{13.55mm}|p{13.55mm}|p{13.55mm}|p{13.55mm}|p{13.55mm}|p{13.55mm}|p{13.55mm}|p{13.55mm}|p{13.55mm}|p{13.55mm}|@{}}\hline
Engine & Execution & \multicolumn{8}{c|}{\textbf{Index}} & \multicolumn{8}{c|}{\textbf{No index}}\\
 &  & \multicolumn{4}{c|}{\textbf{4 Gbp}} & \multicolumn{4}{c|}{\textbf{2 Gbp}} & \multicolumn{4}{c|}{\textbf{4 Gbp}} & \multicolumn{4}{c|}{\textbf{2 Gbp}}\\
 &  & \multicolumn{2}{c|}{\textbf{k=15, w=10}} & \multicolumn{2}{c|}{\textbf{k=20, w=15}} & \multicolumn{2}{c|}{\textbf{k=15, w=10}} & \multicolumn{2}{c|}{\textbf{k=20, w=15}} & \multicolumn{2}{c|}{\textbf{k=15, w=10}} & \multicolumn{2}{c|}{\textbf{k=20, w=15}} & \multicolumn{2}{c|}{\textbf{k=15, w=10}} & \multicolumn{2}{c|}{\textbf{k=20, w=15}}\\
 &  & \multicolumn{1}{c|}{\textbf{DNA4}} & \multicolumn{1}{c|}{\textbf{DNA5}} & \multicolumn{1}{c|}{\textbf{DNA4}} & \multicolumn{1}{c|}{\textbf{DNA5}} & \multicolumn{1}{c|}{\textbf{DNA4}} & \multicolumn{1}{c|}{\textbf{DNA5}} & \multicolumn{1}{c|}{\textbf{DNA4}} & \multicolumn{1}{c|}{\textbf{DNA5}} & \multicolumn{1}{c|}{\textbf{DNA4}} & \multicolumn{1}{c|}{\textbf{DNA5}} & \multicolumn{1}{c|}{\textbf{DNA4}} & \multicolumn{1}{c|}{\textbf{DNA5}} & \multicolumn{1}{c|}{\textbf{DNA4}} & \multicolumn{1}{c|}{\textbf{DNA5}} & \multicolumn{1}{c|}{\textbf{DNA4}} & \multicolumn{1}{c|}{\textbf{DNA5}}\\
 & & {\color{gray}C1} & {\color{gray}C2} & {\color{gray}C3} & {\color{gray}C4} & {\color{gray}C5} & {\color{gray}C6} & {\color{gray}C7} & {\color{gray}C8} & {\color{gray}C9} & {\color{gray}C10} & {\color{gray}C11} & {\color{gray}C12} & {\color{gray}C13} & {\color{gray}C14} & {\color{gray}C15} & {\color{gray}C16}\\
\hline
\rowcolor{BenchBaseline}[0.5pt][0.5pt]\shortstack[l]{{\color{gray}R1}\\\textbf{minimap2}} & \textbf{Baseline} & \hspace*{1pt}\begin{tabular}[t]{@{}c@{}}\rule{0pt}{12pt}39.7s\\[-.2pt]{\fontsize{5.8}{6.4}\selectfont \textcolor{BrickRed}{$ \blacktriangledown\,2.2\times $}}\\[-.2pt]36.8s\\[-.2pt]{\fontsize{5.8}{6.4}\selectfont \textcolor{BrickRed}{$ \blacktriangledown\,2.2\times $}}\end{tabular} & \hspace*{1pt}\begin{tabular}[t]{@{}c@{}}\rule{0pt}{12pt}25.2s\\[-.2pt]{\fontsize{5.8}{6.4}\selectfont \textcolor{BrickRed}{$ \blacktriangledown\,1.5\times $}}\\[-.2pt]22.7s\\[-.2pt]{\fontsize{5.8}{6.4}\selectfont \textcolor{BrickRed}{$ \blacktriangledown\,1.5\times $}}\end{tabular} & \hspace*{1pt}\begin{tabular}[t]{@{}c@{}}\rule{0pt}{12pt}1m\,56.9s\\[-.2pt]{\fontsize{5.8}{6.4}\selectfont \textcolor{BrickRed}{$ \blacktriangledown\,29.1\times $}}\\[-.2pt]1m\,53.2s\\[-.2pt]{\fontsize{5.8}{6.4}\selectfont \textcolor{BrickRed}{$ \blacktriangledown\,32.9\times $}}\end{tabular} & \hspace*{1pt}\begin{tabular}[t]{@{}c@{}}\rule{0pt}{12pt}24.9s\\[-.2pt]{\fontsize{5.8}{6.4}\selectfont \textcolor{BrickRed}{$ \blacktriangledown\,1.6\times $}}\\[-.2pt]22.2s\\[-.2pt]{\fontsize{5.8}{6.4}\selectfont \textcolor{BrickRed}{$ \blacktriangledown\,1.6\times $}}\end{tabular} & \hspace*{1pt}\begin{tabular}[t]{@{}c@{}}\rule{0pt}{12pt}28.6s\\[-.2pt]{\fontsize{5.8}{6.4}\selectfont \textcolor{BrickRed}{$ \blacktriangledown\,11.5\times $}}\\[-.2pt]25.8s\\[-.2pt]{\fontsize{5.8}{6.4}\selectfont \textcolor{BrickRed}{$ \blacktriangledown\,15.6\times $}}\end{tabular} & \hspace*{1pt}\begin{tabular}[t]{@{}c@{}}\rule{0pt}{12pt}15.3s\\[-.2pt]{\fontsize{5.8}{6.4}\selectfont \textcolor{BrickRed}{$ \blacktriangledown\,2.8\times $}}\\[-.2pt]12.7s\\[-.2pt]{\fontsize{5.8}{6.4}\selectfont \textcolor{BrickRed}{$ \blacktriangledown\,3.3\times $}}\end{tabular} & \hspace*{1pt}\begin{tabular}[t]{@{}c@{}}\rule{0pt}{12pt}29.3s\\[-.2pt]{\fontsize{5.8}{6.4}\selectfont \textcolor{BrickRed}{$ \blacktriangledown\,15.0\times $}}\\[-.2pt]26.2s\\[-.2pt]{\fontsize{5.8}{6.4}\selectfont \textcolor{BrickRed}{$ \blacktriangledown\,22.9\times $}}\end{tabular} & \hspace*{1pt}\begin{tabular}[t]{@{}c@{}}\rule{0pt}{12pt}8.7s\\[-.2pt]{\fontsize{5.8}{6.4}\selectfont \textcolor{BrickRed}{$ \blacktriangledown\,1.7\times $}}\\[-.2pt]8.1s\\[-.2pt]{\fontsize{5.8}{6.4}\selectfont \textcolor{BrickRed}{$ \blacktriangledown\,2.3\times $}}\end{tabular} & \hspace*{1pt}\begin{tabular}[t]{@{}c@{}}\rule{0pt}{12pt}2m\,52.4s\\[-.2pt]{\fontsize{5.8}{6.4}\selectfont \textcolor{ForestGreen}{$ \blacktriangle\,19\% $}}\\[-.2pt]2m\,48.5s\\[-.2pt]{\fontsize{5.8}{6.4}\selectfont \textcolor{ForestGreen}{$ \blacktriangle\,15\% $}}\end{tabular} & \hspace*{1pt}\begin{tabular}[t]{@{}c@{}}\rule{0pt}{12pt}2m\,17.1s\\[-.2pt]{\fontsize{5.8}{6.4}\selectfont \textcolor{ForestGreen}{$ \blacktriangle\,3.9\times $}}\\[-.2pt]2m\,15.9s\\[-.2pt]{\fontsize{5.8}{6.4}\selectfont \textcolor{ForestGreen}{$ \blacktriangle\,1.7\times $}}\end{tabular} & \hspace*{1pt}\begin{tabular}[t]{@{}c@{}}\rule{0pt}{12pt}2m\,55s\\[-.2pt]{\fontsize{5.8}{6.4}\selectfont \textcolor{BrickRed}{$ \blacktriangledown\,6\% $}}\\[-.2pt]2m\,51s\\[-.2pt]{\fontsize{5.8}{6.4}\selectfont \textcolor{BrickRed}{$ \blacktriangledown\,12\% $}}\end{tabular} & \hspace*{1pt}\begin{tabular}[t]{@{}c@{}}\rule{0pt}{12pt}55.9s\\[-.2pt]{\fontsize{5.8}{6.4}\selectfont \textcolor{ForestGreen}{$ \blacktriangle\,8.6\times $}}\\[-.2pt]55.1s\\[-.2pt]{\fontsize{5.8}{6.4}\selectfont \textcolor{ForestGreen}{$ \blacktriangle\,3.1\times $}}\end{tabular} & \hspace*{1pt}\begin{tabular}[t]{@{}c@{}}\rule{0pt}{12pt}43.3s\\[-.2pt]{\fontsize{5.8}{6.4}\selectfont \textcolor{ForestGreen}{$ \blacktriangle\,2.6\times $}}\\[-.2pt]42.5s\\[-.2pt]{\fontsize{5.8}{6.4}\selectfont \textcolor{ForestGreen}{$ \blacktriangle\,2.4\times $}}\end{tabular} & \hspace*{1pt}\begin{tabular}[t]{@{}c@{}}\rule{0pt}{12pt}38.9s\\[-.2pt]{\fontsize{5.8}{6.4}\selectfont \textcolor{ForestGreen}{$ \blacktriangle\,8.0\times $}}\\[-.2pt]38.2s\\[-.2pt]{\fontsize{5.8}{6.4}\selectfont \textcolor{ForestGreen}{$ \blacktriangle\,2.8\times $}}\end{tabular} & \hspace*{1pt}\begin{tabular}[t]{@{}c@{}}\rule{0pt}{12pt}53.1s\\[-.2pt]{\fontsize{5.8}{6.4}\selectfont \textcolor{ForestGreen}{$ \blacktriangle\,1.6\times $}}\\[-.2pt]48.2s\\[-.2pt]{\fontsize{5.8}{6.4}\selectfont \textcolor{ForestGreen}{$ \blacktriangle\,1.5\times $}}\end{tabular} & \hspace*{1pt}\begin{tabular}[t]{@{}c@{}}\rule{0pt}{12pt}28.7s\\[-.2pt]{\fontsize{5.8}{6.4}\selectfont \textcolor{ForestGreen}{$ \blacktriangle\,9.8\times $}}\\[-.2pt]28.2s\\[-.2pt]{\fontsize{5.8}{6.4}\selectfont \textcolor{ForestGreen}{$ \blacktriangle\,2.9\times $}}\end{tabular}\\[3pt]\hline
\BenchmarkSourceLabel{bench:local:row:mps-scalar-vulkan-ksw2}{R54}\shortstack[l]{{\color{gray}R54}\\\textbf{BioseqDB}} & \begin{tabular}[t]{@{}l@{}}\rule{0pt}{12pt}MPS / Scalar\\Vulkan-DP / KSW2\end{tabular} & \textcolor{gray}{\ensuremath{\cdot}} & \textcolor{gray}{\ensuremath{\cdot}} & \textcolor{gray}{\ensuremath{\cdot}} & \textcolor{gray}{\ensuremath{\cdot}} & \hspace*{1pt}\begin{tabular}[t]{@{}c@{}}\rule{0pt}{12pt}31s\\[-.2pt]{\fontsize{5.8}{6.4}\selectfont \textcolor{BrickRed}{$ \blacktriangledown\,9\% $}}\\[-.2pt]18.4s\\[-.2pt]{\fontsize{5.8}{6.4}\selectfont \textcolor{ForestGreen}{$ \blacktriangle\,29\% $}}\end{tabular} & \hspace*{1pt}\begin{tabular}[t]{@{}c@{}}\rule{0pt}{12pt}4m\,7.9s\\[-.2pt]{\fontsize{5.8}{6.4}\selectfont \textcolor{BrickRed}{$ \blacktriangledown\,16.2\times $}}\\[-.2pt]43.3s\\[-.2pt]{\fontsize{5.8}{6.4}\selectfont \textcolor{BrickRed}{$ \blacktriangledown\,3.4\times $}}\end{tabular} & \hspace*{1pt}\begin{tabular}[t]{@{}c@{}}\rule{0pt}{12pt}24.2s\\[-.2pt]{\fontsize{5.8}{6.4}\selectfont \textcolor{ForestGreen}{$ \blacktriangle\,18\% $}}\\[-.2pt]15s\\[-.2pt]{\fontsize{5.8}{6.4}\selectfont \textcolor{ForestGreen}{$ \blacktriangle\,1.7\times $}}\end{tabular} & \hspace*{1pt}\begin{tabular}[t]{@{}c@{}}\rule{0pt}{12pt}3m\,55.7s\\[-.2pt]{\fontsize{5.8}{6.4}\selectfont \textcolor{BrickRed}{$ \blacktriangledown\,27.1\times $}}\\[-.2pt]32.7s\\[-.2pt]{\fontsize{5.8}{6.4}\selectfont \textcolor{BrickRed}{$ \blacktriangledown\,4.0\times $}}\end{tabular} & \textcolor{gray}{\ensuremath{\cdot}} & \textcolor{gray}{\ensuremath{\cdot}} & \textcolor{gray}{\ensuremath{\cdot}} & \textcolor{gray}{\ensuremath{\cdot}} & \textcolor{gray}{\ensuremath{\cdot}} & \textcolor{gray}{\ensuremath{\cdot}} & \textcolor{gray}{\ensuremath{\cdot}} & \textcolor{gray}{\ensuremath{\cdot}}\\[3pt]\hline
\BenchmarkSourceLabel{bench:local:row:mpf-bounded-vulkan-ksw2}{R55}\shortstack[l]{{\color{gray}R55}} & \begin{tabular}[t]{@{}l@{}}\rule{0pt}{12pt}MPF / Bounded\\Vulkan-DP / KSW2\end{tabular} & \hspace*{1pt}\begin{tabular}[t]{@{}c@{}}\rule{0pt}{12pt}57.1s\\[-.2pt]{\fontsize{5.8}{6.4}\selectfont \textcolor{BrickRed}{$ \blacktriangledown\,1.4\times $}}\\[-.2pt]40.2s\\[-.2pt]{\fontsize{5.8}{6.4}\selectfont \textcolor{BrickRed}{$ \blacktriangledown\,9\% $}}\end{tabular} & \hspace*{1pt}\begin{tabular}[t]{@{}c@{}}\rule{0pt}{12pt}5m\,53.3s\\[-.2pt]{\fontsize{5.8}{6.4}\selectfont \textcolor{BrickRed}{$ \blacktriangledown\,14.0\times $}}\\[-.2pt]51.2s\\[-.2pt]{\fontsize{5.8}{6.4}\selectfont \textcolor{BrickRed}{$ \blacktriangledown\,2.3\times $}}\end{tabular} & \hspace*{1pt}\begin{tabular}[t]{@{}c@{}}\rule{0pt}{12pt}1m\,3.5s\\[-.2pt]{\fontsize{5.8}{6.4}\selectfont \textcolor{ForestGreen}{$ \blacktriangle\,1.8\times $}}\\[-.2pt]46.7s\\[-.2pt]{\fontsize{5.8}{6.4}\selectfont \textcolor{ForestGreen}{$ \blacktriangle\,2.4\times $}}\end{tabular} & \hspace*{1pt}\begin{tabular}[t]{@{}c@{}}\rule{0pt}{12pt}6m\,14.8s\\[-.2pt]{\fontsize{5.8}{6.4}\selectfont \textcolor{BrickRed}{$ \blacktriangledown\,15.1\times $}}\\[-.2pt]1m\,3.1s\\[-.2pt]{\fontsize{5.8}{6.4}\selectfont \textcolor{BrickRed}{$ \blacktriangledown\,2.8\times $}}\end{tabular} & \hspace*{1pt}\begin{tabular}[t]{@{}c@{}}\rule{0pt}{12pt}42s\\[-.2pt]{\fontsize{5.8}{6.4}\selectfont \textcolor{BrickRed}{$ \blacktriangledown\,1.5\times $}}\\[-.2pt]27.9s\\[-.2pt]{\fontsize{5.8}{6.4}\selectfont \textcolor{BrickRed}{$ \blacktriangledown\,8\% $}}\end{tabular} & \hspace*{1pt}\begin{tabular}[t]{@{}c@{}}\rule{0pt}{12pt}3m\,59.2s\\[-.2pt]{\fontsize{5.8}{6.4}\selectfont \textcolor{BrickRed}{$ \blacktriangledown\,15.7\times $}}\\[-.2pt]30.8s\\[-.2pt]{\fontsize{5.8}{6.4}\selectfont \textcolor{BrickRed}{$ \blacktriangledown\,2.4\times $}}\end{tabular} & \hspace*{1pt}\begin{tabular}[t]{@{}c@{}}\rule{0pt}{12pt}32.3s\\[-.2pt]{\fontsize{5.8}{6.4}\selectfont \textcolor{BrickRed}{$ \blacktriangledown\,10\% $}}\\[-.2pt]17.8s\\[-.2pt]{\fontsize{5.8}{6.4}\selectfont \textcolor{ForestGreen}{$ \blacktriangle\,1.5\times $}}\end{tabular} & \hspace*{1pt}\begin{tabular}[t]{@{}c@{}}\rule{0pt}{12pt}3m\,47.2s\\[-.2pt]{\fontsize{5.8}{6.4}\selectfont \textcolor{BrickRed}{$ \blacktriangledown\,26.2\times $}}\\[-.2pt]23.7s\\[-.2pt]{\fontsize{5.8}{6.4}\selectfont \textcolor{BrickRed}{$ \blacktriangledown\,2.9\times $}}\end{tabular} & \textcolor{gray}{\ensuremath{\cdot}} & \textcolor{gray}{\ensuremath{\cdot}} & \textcolor{gray}{\ensuremath{\cdot}} & \textcolor{gray}{\ensuremath{\cdot}} & \textcolor{gray}{\ensuremath{\cdot}} & \textcolor{gray}{\ensuremath{\cdot}} & \textcolor{gray}{\ensuremath{\cdot}} & \textcolor{gray}{\ensuremath{\cdot}}\\[3pt]\hline
\BenchmarkSourceLabel{bench:local:row:mpf-bounded-vulkan-block}{R56}\shortstack[l]{{\color{gray}R56}} & \begin{tabular}[t]{@{}l@{}}\rule{0pt}{12pt}MPF / Bounded\\Vulkan-DP / Block\end{tabular} & \hspace*{1pt}\begin{tabular}[t]{@{}c@{}}\rule{0pt}{12pt}44.5s\\[-.2pt]{\fontsize{5.8}{6.4}\selectfont \textcolor{BrickRed}{$ \blacktriangledown\,12\% $}}\\[-.2pt]29.5s\\[-.2pt]{\fontsize{5.8}{6.4}\selectfont \textcolor{ForestGreen}{$ \blacktriangle\,20\% $}}\end{tabular} & \hspace*{1pt}\begin{tabular}[t]{@{}c@{}}\rule{0pt}{12pt}6m\,10.7s\\[-.2pt]{\fontsize{5.8}{6.4}\selectfont \textcolor{BrickRed}{$ \blacktriangledown\,14.7\times $}}\\[-.2pt]1m\,8.5s\\[-.2pt]{\fontsize{5.8}{6.4}\selectfont \textcolor{BrickRed}{$ \blacktriangledown\,3.0\times $}}\end{tabular} & \hspace*{1pt}\begin{tabular}[t]{@{}c@{}}\rule{0pt}{12pt}1m\,0.8s\\[-.2pt]{\fontsize{5.8}{6.4}\selectfont \textcolor{ForestGreen}{$ \blacktriangle\,1.9\times $}}\\[-.2pt]45.9s\\[-.2pt]{\fontsize{5.8}{6.4}\selectfont \textcolor{ForestGreen}{$ \blacktriangle\,2.5\times $}}\end{tabular} & \hspace*{1pt}\begin{tabular}[t]{@{}c@{}}\rule{0pt}{12pt}6m\,29.4s\\[-.2pt]{\fontsize{5.8}{6.4}\selectfont \textcolor{BrickRed}{$ \blacktriangledown\,15.6\times $}}\\[-.2pt]1m\,21.6s\\[-.2pt]{\fontsize{5.8}{6.4}\selectfont \textcolor{BrickRed}{$ \blacktriangledown\,3.7\times $}}\end{tabular} & \hspace*{1pt}\begin{tabular}[t]{@{}c@{}}\rule{0pt}{12pt}37.9s\\[-.2pt]{\fontsize{5.8}{6.4}\selectfont \textcolor{BrickRed}{$ \blacktriangledown\,1.3\times $}}\\[-.2pt]25.3s\\[-.2pt]{\fontsize{5.8}{6.4}\selectfont \textcolor{gray}{$ \blacktriangleright\,2\% $}}\end{tabular} & \hspace*{1pt}\begin{tabular}[t]{@{}c@{}}\rule{0pt}{12pt}4m\,4.2s\\[-.2pt]{\fontsize{5.8}{6.4}\selectfont \textcolor{BrickRed}{$ \blacktriangledown\,16.0\times $}}\\[-.2pt]38.6s\\[-.2pt]{\fontsize{5.8}{6.4}\selectfont \textcolor{BrickRed}{$ \blacktriangledown\,3.0\times $}}\end{tabular} & \hspace*{1pt}\begin{tabular}[t]{@{}c@{}}\rule{0pt}{12pt}37s\\[-.2pt]{\fontsize{5.8}{6.4}\selectfont \textcolor{BrickRed}{$ \blacktriangledown\,26\% $}}\\[-.2pt]24.7s\\[-.2pt]{\fontsize{5.8}{6.4}\selectfont \textcolor{ForestGreen}{$ \blacktriangle\,6\% $}}\end{tabular} & \hspace*{1pt}\begin{tabular}[t]{@{}c@{}}\rule{0pt}{12pt}3m\,48.2s\\[-.2pt]{\fontsize{5.8}{6.4}\selectfont \textcolor{BrickRed}{$ \blacktriangledown\,26.3\times $}}\\[-.2pt]23.4s\\[-.2pt]{\fontsize{5.8}{6.4}\selectfont \textcolor{BrickRed}{$ \blacktriangledown\,2.9\times $}}\end{tabular} & \textcolor{gray}{\ensuremath{\cdot}} & \textcolor{gray}{\ensuremath{\cdot}} & \textcolor{gray}{\ensuremath{\cdot}} & \textcolor{gray}{\ensuremath{\cdot}} & \textcolor{gray}{\ensuremath{\cdot}} & \textcolor{gray}{\ensuremath{\cdot}} & \textcolor{gray}{\ensuremath{\cdot}} & \textcolor{gray}{\ensuremath{\cdot}}\\[3pt]\hline
\BenchmarkSourceLabel{bench:local:row:mpf-simdref-vulkan-block}{R57}\shortstack[l]{{\color{gray}R57}} & \begin{tabular}[t]{@{}l@{}}\rule{0pt}{12pt}MPF / SIMD-ref\\Vulkan-DP / Block\end{tabular} & \textcolor{gray}{\ensuremath{\cdot}} & \textcolor{gray}{\ensuremath{\cdot}} & \textcolor{gray}{\ensuremath{\cdot}} & \textcolor{gray}{\ensuremath{\cdot}} & \hspace*{1pt}\begin{tabular}[t]{@{}c@{}}\rule{0pt}{12pt}39s\\[-.2pt]{\fontsize{5.8}{6.4}\selectfont \textcolor{BrickRed}{$ \blacktriangledown\,1.4\times $}}\\[-.2pt]26.4s\\[-.2pt]{\fontsize{5.8}{6.4}\selectfont \textcolor{gray}{$ \blacktriangleright\,2\% $}}\end{tabular} & \hspace*{1pt}\begin{tabular}[t]{@{}c@{}}\rule{0pt}{12pt}3m\,49.5s\\[-.2pt]{\fontsize{5.8}{6.4}\selectfont \textcolor{BrickRed}{$ \blacktriangledown\,15.0\times $}}\\[-.2pt]24s\\[-.2pt]{\fontsize{5.8}{6.4}\selectfont \textcolor{BrickRed}{$ \blacktriangledown\,1.9\times $}}\end{tabular} & \hspace*{1pt}\begin{tabular}[t]{@{}c@{}}\rule{0pt}{12pt}13m\,49.2s\\[-.2pt]{\fontsize{5.8}{6.4}\selectfont \textcolor{BrickRed}{$ \blacktriangledown\,28.3\times $}}\\[-.2pt]13m\,37.6s\\[-.2pt]{\fontsize{5.8}{6.4}\selectfont \textcolor{BrickRed}{$ \blacktriangledown\,31.2\times $}}\end{tabular} & \hspace*{1pt}\begin{tabular}[t]{@{}c@{}}\rule{0pt}{12pt}14m\,54.7s\\[-.2pt]{\fontsize{5.8}{6.4}\selectfont \textcolor{BrickRed}{$ \blacktriangledown\,103.0\times $}}\\[-.2pt]11m\,32.8s\\[-.2pt]{\fontsize{5.8}{6.4}\selectfont \textcolor{BrickRed}{$ \blacktriangledown\,85.2\times $}}\end{tabular} & \textcolor{gray}{\ensuremath{\cdot}} & \textcolor{gray}{\ensuremath{\cdot}} & \textcolor{gray}{\ensuremath{\cdot}} & \textcolor{gray}{\ensuremath{\cdot}} & \textcolor{gray}{\ensuremath{\cdot}} & \textcolor{gray}{\ensuremath{\cdot}} & \textcolor{gray}{\ensuremath{\cdot}} & \textcolor{gray}{\ensuremath{\cdot}}\\[3pt]\hline
\BenchmarkSourceLabel{bench:local:row:mpf-scalar-vulkan-block}{R58}\shortstack[l]{{\color{gray}R58}} & \begin{tabular}[t]{@{}l@{}}\rule{0pt}{12pt}MPF / Scalar\\Vulkan-DP / Block\end{tabular} & \hspace*{1pt}\begin{tabular}[t]{@{}c@{}}\rule{0pt}{12pt}43.3s\\[-.2pt]{\fontsize{5.8}{6.4}\selectfont \textcolor{BrickRed}{$ \blacktriangledown\,9\% $}}\\[-.2pt]28.9s\\[-.2pt]{\fontsize{5.8}{6.4}\selectfont \textcolor{ForestGreen}{$ \blacktriangle\,21\% $}}\end{tabular} & \textcolor{gray}{\ensuremath{\cdot}} & \textcolor{gray}{\ensuremath{\cdot}} & \textcolor{gray}{\ensuremath{\cdot}} & \hspace*{1pt}\begin{tabular}[t]{@{}c@{}}\rule{0pt}{12pt}39s\\[-.2pt]{\fontsize{5.8}{6.4}\selectfont \textcolor{BrickRed}{$ \blacktriangledown\,1.4\times $}}\\[-.2pt]26.5s\\[-.2pt]{\fontsize{5.8}{6.4}\selectfont \textcolor{gray}{$ \blacktriangleright\,3\% $}}\end{tabular} & \hspace*{1pt}\begin{tabular}[t]{@{}c@{}}\rule{0pt}{12pt}4m\,5.4s\\[-.2pt]{\fontsize{5.8}{6.4}\selectfont \textcolor{BrickRed}{$ \blacktriangledown\,16.1\times $}}\\[-.2pt]39.5s\\[-.2pt]{\fontsize{5.8}{6.4}\selectfont \textcolor{BrickRed}{$ \blacktriangledown\,3.1\times $}}\end{tabular} & \hspace*{1pt}\begin{tabular}[t]{@{}c@{}}\rule{0pt}{12pt}12m\,37.6s\\[-.2pt]{\fontsize{5.8}{6.4}\selectfont \textcolor{BrickRed}{$ \blacktriangledown\,25.8\times $}}\\[-.2pt]12m\,26.5s\\[-.2pt]{\fontsize{5.8}{6.4}\selectfont \textcolor{BrickRed}{$ \blacktriangledown\,28.5\times $}}\end{tabular} & \hspace*{1pt}\begin{tabular}[t]{@{}c@{}}\rule{0pt}{12pt}27m\,27.3s\\[-.2pt]{\fontsize{5.8}{6.4}\selectfont \textcolor{BrickRed}{$ \blacktriangledown\,189.7\times $}}\\[-.2pt]24m\,4.7s\\[-.2pt]{\fontsize{5.8}{6.4}\selectfont \textcolor{BrickRed}{$ \blacktriangledown\,177.7\times $}}\end{tabular} & \textcolor{gray}{\ensuremath{\cdot}} & \textcolor{gray}{\ensuremath{\cdot}} & \textcolor{gray}{\ensuremath{\cdot}} & \textcolor{gray}{\ensuremath{\cdot}} & \textcolor{gray}{\ensuremath{\cdot}} & \textcolor{gray}{\ensuremath{\cdot}} & \textcolor{gray}{\ensuremath{\cdot}} & \textcolor{gray}{\ensuremath{\cdot}}\\[3pt]\hline
\BenchmarkSourceLabel{bench:local:row:mpf-twostack-vulkan-ksw2}{R59}\shortstack[l]{{\color{gray}R59}} & \begin{tabular}[t]{@{}l@{}}\rule{0pt}{12pt}MPF / Two-stack\\Vulkan-DP / KSW2\end{tabular} & \textcolor{gray}{\ensuremath{\cdot}} & \textcolor{gray}{\ensuremath{\cdot}} & \textcolor{gray}{\ensuremath{\cdot}} & \textcolor{gray}{\ensuremath{\cdot}} & \hspace*{1pt}\begin{tabular}[t]{@{}c@{}}\rule{0pt}{12pt}40.9s\\[-.2pt]{\fontsize{5.8}{6.4}\selectfont \textcolor{BrickRed}{$ \blacktriangledown\,1.4\times $}}\\[-.2pt]27.2s\\[-.2pt]{\fontsize{5.8}{6.4}\selectfont \textcolor{BrickRed}{$ \blacktriangledown\,5\% $}}\end{tabular} & \hspace*{1pt}\begin{tabular}[t]{@{}c@{}}\rule{0pt}{12pt}4m\,6.8s\\[-.2pt]{\fontsize{5.8}{6.4}\selectfont \textcolor{BrickRed}{$ \blacktriangledown\,16.2\times $}}\\[-.2pt]36.8s\\[-.2pt]{\fontsize{5.8}{6.4}\selectfont \textcolor{BrickRed}{$ \blacktriangledown\,2.9\times $}}\end{tabular} & \hspace*{1pt}\begin{tabular}[t]{@{}c@{}}\rule{0pt}{12pt}13m\,56.9s\\[-.2pt]{\fontsize{5.8}{6.4}\selectfont \textcolor{BrickRed}{$ \blacktriangledown\,28.5\times $}}\\[-.2pt]13m\,25.1s\\[-.2pt]{\fontsize{5.8}{6.4}\selectfont \textcolor{BrickRed}{$ \blacktriangledown\,30.7\times $}}\end{tabular} & \hspace*{1pt}\begin{tabular}[t]{@{}c@{}}\rule{0pt}{12pt}10m\,19.9s\\[-.2pt]{\fontsize{5.8}{6.4}\selectfont \textcolor{BrickRed}{$ \blacktriangledown\,71.4\times $}}\\[-.2pt]6m\,50.6s\\[-.2pt]{\fontsize{5.8}{6.4}\selectfont \textcolor{BrickRed}{$ \blacktriangledown\,50.5\times $}}\end{tabular} & \textcolor{gray}{\ensuremath{\cdot}} & \textcolor{gray}{\ensuremath{\cdot}} & \textcolor{gray}{\ensuremath{\cdot}} & \textcolor{gray}{\ensuremath{\cdot}} & \textcolor{gray}{\ensuremath{\cdot}} & \textcolor{gray}{\ensuremath{\cdot}} & \textcolor{gray}{\ensuremath{\cdot}} & \textcolor{gray}{\ensuremath{\cdot}}\\[3pt]\hline
\BenchmarkSourceLabel{bench:local:row:mpf-scalar-vulkan-ksw2}{R60}\shortstack[l]{{\color{gray}R60}} & \begin{tabular}[t]{@{}l@{}}\rule{0pt}{12pt}MPF / Scalar\\Vulkan-DP / KSW2\end{tabular} & \hspace*{1pt}\begin{tabular}[t]{@{}c@{}}\rule{0pt}{12pt}59.4s\\[-.2pt]{\fontsize{5.8}{6.4}\selectfont \textcolor{BrickRed}{$ \blacktriangledown\,1.5\times $}}\\[-.2pt]42s\\[-.2pt]{\fontsize{5.8}{6.4}\selectfont \textcolor{BrickRed}{$ \blacktriangledown\,14\% $}}\end{tabular} & \textcolor{gray}{\ensuremath{\cdot}} & \textcolor{gray}{\ensuremath{\cdot}} & \textcolor{gray}{\ensuremath{\cdot}} & \hspace*{1pt}\begin{tabular}[t]{@{}c@{}}\rule{0pt}{12pt}40.9s\\[-.2pt]{\fontsize{5.8}{6.4}\selectfont \textcolor{BrickRed}{$ \blacktriangledown\,1.4\times $}}\\[-.2pt]26.8s\\[-.2pt]{\fontsize{5.8}{6.4}\selectfont \textcolor{gray}{$ \blacktriangleright\,4\% $}}\end{tabular} & \hspace*{1pt}\begin{tabular}[t]{@{}c@{}}\rule{0pt}{12pt}3m\,59.5s\\[-.2pt]{\fontsize{5.8}{6.4}\selectfont \textcolor{BrickRed}{$ \blacktriangledown\,15.7\times $}}\\[-.2pt]31.6s\\[-.2pt]{\fontsize{5.8}{6.4}\selectfont \textcolor{BrickRed}{$ \blacktriangledown\,2.5\times $}}\end{tabular} & \hspace*{1pt}\begin{tabular}[t]{@{}c@{}}\rule{0pt}{12pt}13m\,53.7s\\[-.2pt]{\fontsize{5.8}{6.4}\selectfont \textcolor{BrickRed}{$ \blacktriangledown\,28.4\times $}}\\[-.2pt]13m\,39.4s\\[-.2pt]{\fontsize{5.8}{6.4}\selectfont \textcolor{BrickRed}{$ \blacktriangledown\,31.3\times $}}\end{tabular} & \textcolor{gray}{\ensuremath{\cdot}} & \textcolor{gray}{\ensuremath{\cdot}} & \textcolor{gray}{\ensuremath{\cdot}} & \textcolor{gray}{\ensuremath{\cdot}} & \textcolor{gray}{\ensuremath{\cdot}} & \textcolor{gray}{\ensuremath{\cdot}} & \textcolor{gray}{\ensuremath{\cdot}} & \textcolor{gray}{\ensuremath{\cdot}} & \textcolor{gray}{\ensuremath{\cdot}}\\[3pt]\hline
\end{tabular}\par}
\smallskip{\footnotesize\noindent Each cell: A and B times, $\cdot$: no measurement. ---\textsuperscript{g}: GPU unavailable. n/a: no comparison.\par
\noindent Blue: one lowest A and one lowest B measurement for each column for each configuration.\par
\noindent Gray: minimap2. Yellow: my preferred MPF / Bounded / SIMD-DP / Block configuration.\par
\noindent MPF/MPS: minimizers per fragment/sequence. SIMD-ref: naive SIMD scan. Bounded: bounded SIMD algorithm.\par
}
\clearpage
\noindent{\large\bfseries Table \thetable: (A): Local machine\hfill 10/12}\par\smallskip
{\fontsize{7.1}{8.1}\selectfont\setlength{\tabcolsep}{1pt}\setlength{\fboxsep}{2pt}\renewcommand{\arraystretch}{1.12}
\noindent\begin{tabular}{@{}p{14mm}p{27mm}|p{13.55mm}|p{13.55mm}|p{13.55mm}|p{13.55mm}|p{13.55mm}|p{13.55mm}|p{13.55mm}|p{13.55mm}|p{13.55mm}|p{13.55mm}|p{13.55mm}|p{13.55mm}|p{13.55mm}|p{13.55mm}|p{13.55mm}|p{13.55mm}|@{}}\hline
Engine & Execution & \multicolumn{8}{c|}{\textbf{Index}} & \multicolumn{8}{c|}{\textbf{No index}}\\
 &  & \multicolumn{4}{c|}{\textbf{4 Gbp}} & \multicolumn{4}{c|}{\textbf{2 Gbp}} & \multicolumn{4}{c|}{\textbf{4 Gbp}} & \multicolumn{4}{c|}{\textbf{2 Gbp}}\\
 &  & \multicolumn{2}{c|}{\textbf{k=15, w=10}} & \multicolumn{2}{c|}{\textbf{k=20, w=15}} & \multicolumn{2}{c|}{\textbf{k=15, w=10}} & \multicolumn{2}{c|}{\textbf{k=20, w=15}} & \multicolumn{2}{c|}{\textbf{k=15, w=10}} & \multicolumn{2}{c|}{\textbf{k=20, w=15}} & \multicolumn{2}{c|}{\textbf{k=15, w=10}} & \multicolumn{2}{c|}{\textbf{k=20, w=15}}\\
 &  & \multicolumn{1}{c|}{\textbf{DNA4}} & \multicolumn{1}{c|}{\textbf{DNA5}} & \multicolumn{1}{c|}{\textbf{DNA4}} & \multicolumn{1}{c|}{\textbf{DNA5}} & \multicolumn{1}{c|}{\textbf{DNA4}} & \multicolumn{1}{c|}{\textbf{DNA5}} & \multicolumn{1}{c|}{\textbf{DNA4}} & \multicolumn{1}{c|}{\textbf{DNA5}} & \multicolumn{1}{c|}{\textbf{DNA4}} & \multicolumn{1}{c|}{\textbf{DNA5}} & \multicolumn{1}{c|}{\textbf{DNA4}} & \multicolumn{1}{c|}{\textbf{DNA5}} & \multicolumn{1}{c|}{\textbf{DNA4}} & \multicolumn{1}{c|}{\textbf{DNA5}} & \multicolumn{1}{c|}{\textbf{DNA4}} & \multicolumn{1}{c|}{\textbf{DNA5}}\\
 & & {\color{gray}C1} & {\color{gray}C2} & {\color{gray}C3} & {\color{gray}C4} & {\color{gray}C5} & {\color{gray}C6} & {\color{gray}C7} & {\color{gray}C8} & {\color{gray}C9} & {\color{gray}C10} & {\color{gray}C11} & {\color{gray}C12} & {\color{gray}C13} & {\color{gray}C14} & {\color{gray}C15} & {\color{gray}C16}\\
\hline
\rowcolor{BenchBaseline}[0.5pt][0.5pt]\shortstack[l]{{\color{gray}R1}\\\textbf{minimap2}} & \textbf{Baseline} & \hspace*{1pt}\begin{tabular}[t]{@{}c@{}}\rule{0pt}{12pt}39.7s\\[-.2pt]{\fontsize{5.8}{6.4}\selectfont \textcolor{BrickRed}{$ \blacktriangledown\,2.2\times $}}\\[-.2pt]36.8s\\[-.2pt]{\fontsize{5.8}{6.4}\selectfont \textcolor{BrickRed}{$ \blacktriangledown\,2.2\times $}}\end{tabular} & \hspace*{1pt}\begin{tabular}[t]{@{}c@{}}\rule{0pt}{12pt}25.2s\\[-.2pt]{\fontsize{5.8}{6.4}\selectfont \textcolor{BrickRed}{$ \blacktriangledown\,1.5\times $}}\\[-.2pt]22.7s\\[-.2pt]{\fontsize{5.8}{6.4}\selectfont \textcolor{BrickRed}{$ \blacktriangledown\,1.5\times $}}\end{tabular} & \hspace*{1pt}\begin{tabular}[t]{@{}c@{}}\rule{0pt}{12pt}1m\,56.9s\\[-.2pt]{\fontsize{5.8}{6.4}\selectfont \textcolor{BrickRed}{$ \blacktriangledown\,29.1\times $}}\\[-.2pt]1m\,53.2s\\[-.2pt]{\fontsize{5.8}{6.4}\selectfont \textcolor{BrickRed}{$ \blacktriangledown\,32.9\times $}}\end{tabular} & \hspace*{1pt}\begin{tabular}[t]{@{}c@{}}\rule{0pt}{12pt}24.9s\\[-.2pt]{\fontsize{5.8}{6.4}\selectfont \textcolor{BrickRed}{$ \blacktriangledown\,1.6\times $}}\\[-.2pt]22.2s\\[-.2pt]{\fontsize{5.8}{6.4}\selectfont \textcolor{BrickRed}{$ \blacktriangledown\,1.6\times $}}\end{tabular} & \hspace*{1pt}\begin{tabular}[t]{@{}c@{}}\rule{0pt}{12pt}28.6s\\[-.2pt]{\fontsize{5.8}{6.4}\selectfont \textcolor{BrickRed}{$ \blacktriangledown\,11.5\times $}}\\[-.2pt]25.8s\\[-.2pt]{\fontsize{5.8}{6.4}\selectfont \textcolor{BrickRed}{$ \blacktriangledown\,15.6\times $}}\end{tabular} & \hspace*{1pt}\begin{tabular}[t]{@{}c@{}}\rule{0pt}{12pt}15.3s\\[-.2pt]{\fontsize{5.8}{6.4}\selectfont \textcolor{BrickRed}{$ \blacktriangledown\,2.8\times $}}\\[-.2pt]12.7s\\[-.2pt]{\fontsize{5.8}{6.4}\selectfont \textcolor{BrickRed}{$ \blacktriangledown\,3.3\times $}}\end{tabular} & \hspace*{1pt}\begin{tabular}[t]{@{}c@{}}\rule{0pt}{12pt}29.3s\\[-.2pt]{\fontsize{5.8}{6.4}\selectfont \textcolor{BrickRed}{$ \blacktriangledown\,15.0\times $}}\\[-.2pt]26.2s\\[-.2pt]{\fontsize{5.8}{6.4}\selectfont \textcolor{BrickRed}{$ \blacktriangledown\,22.9\times $}}\end{tabular} & \hspace*{1pt}\begin{tabular}[t]{@{}c@{}}\rule{0pt}{12pt}8.7s\\[-.2pt]{\fontsize{5.8}{6.4}\selectfont \textcolor{BrickRed}{$ \blacktriangledown\,1.7\times $}}\\[-.2pt]8.1s\\[-.2pt]{\fontsize{5.8}{6.4}\selectfont \textcolor{BrickRed}{$ \blacktriangledown\,2.3\times $}}\end{tabular} & \hspace*{1pt}\begin{tabular}[t]{@{}c@{}}\rule{0pt}{12pt}2m\,52.4s\\[-.2pt]{\fontsize{5.8}{6.4}\selectfont \textcolor{ForestGreen}{$ \blacktriangle\,19\% $}}\\[-.2pt]2m\,48.5s\\[-.2pt]{\fontsize{5.8}{6.4}\selectfont \textcolor{ForestGreen}{$ \blacktriangle\,15\% $}}\end{tabular} & \hspace*{1pt}\begin{tabular}[t]{@{}c@{}}\rule{0pt}{12pt}2m\,17.1s\\[-.2pt]{\fontsize{5.8}{6.4}\selectfont \textcolor{ForestGreen}{$ \blacktriangle\,3.9\times $}}\\[-.2pt]2m\,15.9s\\[-.2pt]{\fontsize{5.8}{6.4}\selectfont \textcolor{ForestGreen}{$ \blacktriangle\,1.7\times $}}\end{tabular} & \hspace*{1pt}\begin{tabular}[t]{@{}c@{}}\rule{0pt}{12pt}2m\,55s\\[-.2pt]{\fontsize{5.8}{6.4}\selectfont \textcolor{BrickRed}{$ \blacktriangledown\,6\% $}}\\[-.2pt]2m\,51s\\[-.2pt]{\fontsize{5.8}{6.4}\selectfont \textcolor{BrickRed}{$ \blacktriangledown\,12\% $}}\end{tabular} & \hspace*{1pt}\begin{tabular}[t]{@{}c@{}}\rule{0pt}{12pt}55.9s\\[-.2pt]{\fontsize{5.8}{6.4}\selectfont \textcolor{ForestGreen}{$ \blacktriangle\,8.6\times $}}\\[-.2pt]55.1s\\[-.2pt]{\fontsize{5.8}{6.4}\selectfont \textcolor{ForestGreen}{$ \blacktriangle\,3.1\times $}}\end{tabular} & \hspace*{1pt}\begin{tabular}[t]{@{}c@{}}\rule{0pt}{12pt}43.3s\\[-.2pt]{\fontsize{5.8}{6.4}\selectfont \textcolor{ForestGreen}{$ \blacktriangle\,2.6\times $}}\\[-.2pt]42.5s\\[-.2pt]{\fontsize{5.8}{6.4}\selectfont \textcolor{ForestGreen}{$ \blacktriangle\,2.4\times $}}\end{tabular} & \hspace*{1pt}\begin{tabular}[t]{@{}c@{}}\rule{0pt}{12pt}38.9s\\[-.2pt]{\fontsize{5.8}{6.4}\selectfont \textcolor{ForestGreen}{$ \blacktriangle\,8.0\times $}}\\[-.2pt]38.2s\\[-.2pt]{\fontsize{5.8}{6.4}\selectfont \textcolor{ForestGreen}{$ \blacktriangle\,2.8\times $}}\end{tabular} & \hspace*{1pt}\begin{tabular}[t]{@{}c@{}}\rule{0pt}{12pt}53.1s\\[-.2pt]{\fontsize{5.8}{6.4}\selectfont \textcolor{ForestGreen}{$ \blacktriangle\,1.6\times $}}\\[-.2pt]48.2s\\[-.2pt]{\fontsize{5.8}{6.4}\selectfont \textcolor{ForestGreen}{$ \blacktriangle\,1.5\times $}}\end{tabular} & \hspace*{1pt}\begin{tabular}[t]{@{}c@{}}\rule{0pt}{12pt}28.7s\\[-.2pt]{\fontsize{5.8}{6.4}\selectfont \textcolor{ForestGreen}{$ \blacktriangle\,9.8\times $}}\\[-.2pt]28.2s\\[-.2pt]{\fontsize{5.8}{6.4}\selectfont \textcolor{ForestGreen}{$ \blacktriangle\,2.9\times $}}\end{tabular}\\[3pt]\hline
\BenchmarkSourceLabel{bench:local:row:mpf-simdref-vulkan-ksw2}{R61}\shortstack[l]{{\color{gray}R61}\\\textbf{BioseqDB}} & \begin{tabular}[t]{@{}l@{}}\rule{0pt}{12pt}MPF / SIMD-ref\\Vulkan-DP / KSW2\end{tabular} & \textcolor{gray}{\ensuremath{\cdot}} & \textcolor{gray}{\ensuremath{\cdot}} & \textcolor{gray}{\ensuremath{\cdot}} & \textcolor{gray}{\ensuremath{\cdot}} & \hspace*{1pt}\begin{tabular}[t]{@{}c@{}}\rule{0pt}{12pt}41.8s\\[-.2pt]{\fontsize{5.8}{6.4}\selectfont \textcolor{BrickRed}{$ \blacktriangledown\,1.5\times $}}\\[-.2pt]27.5s\\[-.2pt]{\fontsize{5.8}{6.4}\selectfont \textcolor{BrickRed}{$ \blacktriangledown\,6\% $}}\end{tabular} & \hspace*{1pt}\begin{tabular}[t]{@{}c@{}}\rule{0pt}{12pt}3m\,50.7s\\[-.2pt]{\fontsize{5.8}{6.4}\selectfont \textcolor{BrickRed}{$ \blacktriangledown\,15.1\times $}}\\[-.2pt]23.9s\\[-.2pt]{\fontsize{5.8}{6.4}\selectfont \textcolor{BrickRed}{$ \blacktriangledown\,1.9\times $}}\end{tabular} & \hspace*{1pt}\begin{tabular}[t]{@{}c@{}}\rule{0pt}{12pt}13m\,53s\\[-.2pt]{\fontsize{5.8}{6.4}\selectfont \textcolor{BrickRed}{$ \blacktriangledown\,28.4\times $}}\\[-.2pt]13m\,39.1s\\[-.2pt]{\fontsize{5.8}{6.4}\selectfont \textcolor{BrickRed}{$ \blacktriangledown\,31.2\times $}}\end{tabular} & \hspace*{1pt}\begin{tabular}[t]{@{}c@{}}\rule{0pt}{12pt}15m\,18.3s\\[-.2pt]{\fontsize{5.8}{6.4}\selectfont \textcolor{BrickRed}{$ \blacktriangledown\,105.7\times $}}\\[-.2pt]11m\,51.5s\\[-.2pt]{\fontsize{5.8}{6.4}\selectfont \textcolor{BrickRed}{$ \blacktriangledown\,87.5\times $}}\end{tabular} & \textcolor{gray}{\ensuremath{\cdot}} & \textcolor{gray}{\ensuremath{\cdot}} & \textcolor{gray}{\ensuremath{\cdot}} & \textcolor{gray}{\ensuremath{\cdot}} & \textcolor{gray}{\ensuremath{\cdot}} & \textcolor{gray}{\ensuremath{\cdot}} & \textcolor{gray}{\ensuremath{\cdot}} & \textcolor{gray}{\ensuremath{\cdot}}\\[3pt]\hline
\BenchmarkSourceLabel{bench:local:row:mpf-twostack-vulkan-block}{R62}\shortstack[l]{{\color{gray}R62}} & \begin{tabular}[t]{@{}l@{}}\rule{0pt}{12pt}MPF / Two-stack\\Vulkan-DP / Block\end{tabular} & \textcolor{gray}{\ensuremath{\cdot}} & \textcolor{gray}{\ensuremath{\cdot}} & \textcolor{gray}{\ensuremath{\cdot}} & \textcolor{gray}{\ensuremath{\cdot}} & \hspace*{1pt}\begin{tabular}[t]{@{}c@{}}\rule{0pt}{12pt}43.6s\\[-.2pt]{\fontsize{5.8}{6.4}\selectfont \textcolor{BrickRed}{$ \blacktriangledown\,1.5\times $}}\\[-.2pt]31.6s\\[-.2pt]{\fontsize{5.8}{6.4}\selectfont \textcolor{BrickRed}{$ \blacktriangledown\,22\% $}}\end{tabular} & \hspace*{1pt}\begin{tabular}[t]{@{}c@{}}\rule{0pt}{12pt}4m\,5.1s\\[-.2pt]{\fontsize{5.8}{6.4}\selectfont \textcolor{BrickRed}{$ \blacktriangledown\,16.0\times $}}\\[-.2pt]36.3s\\[-.2pt]{\fontsize{5.8}{6.4}\selectfont \textcolor{BrickRed}{$ \blacktriangledown\,2.9\times $}}\end{tabular} & \hspace*{1pt}\begin{tabular}[t]{@{}c@{}}\rule{0pt}{12pt}12m\,53.5s\\[-.2pt]{\fontsize{5.8}{6.4}\selectfont \textcolor{BrickRed}{$ \blacktriangledown\,26.4\times $}}\\[-.2pt]12m\,25.3s\\[-.2pt]{\fontsize{5.8}{6.4}\selectfont \textcolor{BrickRed}{$ \blacktriangledown\,28.4\times $}}\end{tabular} & \hspace*{1pt}\begin{tabular}[t]{@{}c@{}}\rule{0pt}{12pt}10m\,19.7s\\[-.2pt]{\fontsize{5.8}{6.4}\selectfont \textcolor{BrickRed}{$ \blacktriangledown\,71.4\times $}}\\[-.2pt]6m\,54.6s\\[-.2pt]{\fontsize{5.8}{6.4}\selectfont \textcolor{BrickRed}{$ \blacktriangledown\,51.0\times $}}\end{tabular} & \textcolor{gray}{\ensuremath{\cdot}} & \textcolor{gray}{\ensuremath{\cdot}} & \textcolor{gray}{\ensuremath{\cdot}} & \textcolor{gray}{\ensuremath{\cdot}} & \textcolor{gray}{\ensuremath{\cdot}} & \textcolor{gray}{\ensuremath{\cdot}} & \textcolor{gray}{\ensuremath{\cdot}} & \textcolor{gray}{\ensuremath{\cdot}}\\[3pt]\hline
\BenchmarkSourceLabel{bench:local:row:mpf-simdref-scalar-block}{R63}\shortstack[l]{{\color{gray}R63}} & \begin{tabular}[t]{@{}l@{}}\rule{0pt}{12pt}MPF / SIMD-ref\\Scalar-DP / Block\end{tabular} & \textcolor{gray}{\ensuremath{\cdot}} & \textcolor{gray}{\ensuremath{\cdot}} & \textcolor{gray}{\ensuremath{\cdot}} & \textcolor{gray}{\ensuremath{\cdot}} & \textcolor{gray}{\ensuremath{\cdot}} & \textcolor{gray}{\ensuremath{\cdot}} & \hspace*{1pt}\begin{tabular}[t]{@{}c@{}}\rule{0pt}{12pt}1m\,34.2s\\[-.2pt]{\fontsize{5.8}{6.4}\selectfont \textcolor{BrickRed}{$ \blacktriangledown\,3.2\times $}}\\[-.2pt]1m\,34s\\[-.2pt]{\fontsize{5.8}{6.4}\selectfont \textcolor{BrickRed}{$ \blacktriangledown\,3.6\times $}}\end{tabular} & \hspace*{1pt}\begin{tabular}[t]{@{}c@{}}\rule{0pt}{12pt}1m\,37.5s\\[-.2pt]{\fontsize{5.8}{6.4}\selectfont \textcolor{BrickRed}{$ \blacktriangledown\,11.2\times $}}\\[-.2pt]1m\,36.6s\\[-.2pt]{\fontsize{5.8}{6.4}\selectfont \textcolor{BrickRed}{$ \blacktriangledown\,11.9\times $}}\end{tabular} & \textcolor{gray}{\ensuremath{\cdot}} & \textcolor{gray}{\ensuremath{\cdot}} & \textcolor{gray}{\ensuremath{\cdot}} & \textcolor{gray}{\ensuremath{\cdot}} & \textcolor{gray}{\ensuremath{\cdot}} & \textcolor{gray}{\ensuremath{\cdot}} & \textcolor{gray}{\ensuremath{\cdot}} & \textcolor{gray}{\ensuremath{\cdot}}\\[3pt]\hline
\BenchmarkSourceLabel{bench:local:row:mpf-simdref-rmq-ksw2}{R64}\shortstack[l]{{\color{gray}R64}} & \begin{tabular}[t]{@{}l@{}}\rule{0pt}{12pt}MPF / SIMD-ref\\RMQ / KSW2\end{tabular} & \textcolor{gray}{\ensuremath{\cdot}} & \textcolor{gray}{\ensuremath{\cdot}} & \textcolor{gray}{\ensuremath{\cdot}} & \textcolor{gray}{\ensuremath{\cdot}} & \textcolor{gray}{\ensuremath{\cdot}} & \textcolor{gray}{\ensuremath{\cdot}} & \hspace*{1pt}\begin{tabular}[t]{@{}c@{}}\rule{0pt}{12pt}1m\,44s\\[-.2pt]{\fontsize{5.8}{6.4}\selectfont \textcolor{BrickRed}{$ \blacktriangledown\,3.5\times $}}\\[-.2pt]1m\,42.7s\\[-.2pt]{\fontsize{5.8}{6.4}\selectfont \textcolor{BrickRed}{$ \blacktriangledown\,3.9\times $}}\end{tabular} & \hspace*{1pt}\begin{tabular}[t]{@{}c@{}}\rule{0pt}{12pt}1m\,34.8s\\[-.2pt]{\fontsize{5.8}{6.4}\selectfont \textcolor{BrickRed}{$ \blacktriangledown\,10.9\times $}}\\[-.2pt]1m\,33.5s\\[-.2pt]{\fontsize{5.8}{6.4}\selectfont \textcolor{BrickRed}{$ \blacktriangledown\,11.5\times $}}\end{tabular} & \textcolor{gray}{\ensuremath{\cdot}} & \textcolor{gray}{\ensuremath{\cdot}} & \textcolor{gray}{\ensuremath{\cdot}} & \textcolor{gray}{\ensuremath{\cdot}} & \textcolor{gray}{\ensuremath{\cdot}} & \textcolor{gray}{\ensuremath{\cdot}} & \textcolor{gray}{\ensuremath{\cdot}} & \textcolor{gray}{\ensuremath{\cdot}}\\[3pt]\hline
\BenchmarkSourceLabel{bench:local:row:mpf-simdref-scalar-ksw2}{R65}\shortstack[l]{{\color{gray}R65}} & \begin{tabular}[t]{@{}l@{}}\rule{0pt}{12pt}MPF / SIMD-ref\\Scalar-DP / KSW2\end{tabular} & \textcolor{gray}{\ensuremath{\cdot}} & \textcolor{gray}{\ensuremath{\cdot}} & \textcolor{gray}{\ensuremath{\cdot}} & \textcolor{gray}{\ensuremath{\cdot}} & \textcolor{gray}{\ensuremath{\cdot}} & \textcolor{gray}{\ensuremath{\cdot}} & \hspace*{1pt}\begin{tabular}[t]{@{}c@{}}\rule{0pt}{12pt}1m\,44.3s\\[-.2pt]{\fontsize{5.8}{6.4}\selectfont \textcolor{BrickRed}{$ \blacktriangledown\,3.6\times $}}\\[-.2pt]1m\,43s\\[-.2pt]{\fontsize{5.8}{6.4}\selectfont \textcolor{BrickRed}{$ \blacktriangledown\,3.9\times $}}\end{tabular} & \hspace*{1pt}\begin{tabular}[t]{@{}c@{}}\rule{0pt}{12pt}1m\,39.6s\\[-.2pt]{\fontsize{5.8}{6.4}\selectfont \textcolor{BrickRed}{$ \blacktriangledown\,11.5\times $}}\\[-.2pt]1m\,38.4s\\[-.2pt]{\fontsize{5.8}{6.4}\selectfont \textcolor{BrickRed}{$ \blacktriangledown\,12.1\times $}}\end{tabular} & \textcolor{gray}{\ensuremath{\cdot}} & \textcolor{gray}{\ensuremath{\cdot}} & \textcolor{gray}{\ensuremath{\cdot}} & \textcolor{gray}{\ensuremath{\cdot}} & \textcolor{gray}{\ensuremath{\cdot}} & \textcolor{gray}{\ensuremath{\cdot}} & \textcolor{gray}{\ensuremath{\cdot}} & \textcolor{gray}{\ensuremath{\cdot}}\\[3pt]\hline
\BenchmarkSourceLabel{bench:local:row:mpf-vulkan-rmq-block}{R66}\shortstack[l]{{\color{gray}R66}} & \begin{tabular}[t]{@{}l@{}}\rule{0pt}{12pt}MPF / Vulkan\\RMQ / Block\end{tabular} & \textcolor{gray}{\ensuremath{\cdot}} & \textcolor{gray}{\ensuremath{\cdot}} & \textcolor{gray}{\ensuremath{\cdot}} & \textcolor{gray}{\ensuremath{\cdot}} & \hspace*{1pt}\begin{tabular}[t]{@{}c@{}}\rule{0pt}{12pt}5m\,41.1s\\[-.2pt]{\fontsize{5.8}{6.4}\selectfont \textcolor{BrickRed}{$ \blacktriangledown\,11.9\times $}}\\[-.2pt]5m\,31.7s\\[-.2pt]{\fontsize{5.8}{6.4}\selectfont \textcolor{BrickRed}{$ \blacktriangledown\,12.8\times $}}\end{tabular} & \hspace*{1pt}\begin{tabular}[t]{@{}c@{}}\rule{0pt}{12pt}8m\,50.1s\\[-.2pt]{\fontsize{5.8}{6.4}\selectfont \textcolor{BrickRed}{$ \blacktriangledown\,34.7\times $}}\\[-.2pt]5m\,27s\\[-.2pt]{\fontsize{5.8}{6.4}\selectfont \textcolor{BrickRed}{$ \blacktriangledown\,25.7\times $}}\end{tabular} & \hspace*{1pt}\begin{tabular}[t]{@{}c@{}}\rule{0pt}{12pt}3m\,9.8s\\[-.2pt]{\fontsize{5.8}{6.4}\selectfont \textcolor{BrickRed}{$ \blacktriangledown\,6.5\times $}}\\[-.2pt]3m\,1.4s\\[-.2pt]{\fontsize{5.8}{6.4}\selectfont \textcolor{BrickRed}{$ \blacktriangledown\,6.9\times $}}\end{tabular} & \hspace*{1pt}\begin{tabular}[t]{@{}c@{}}\rule{0pt}{12pt}6m\,26.4s\\[-.2pt]{\fontsize{5.8}{6.4}\selectfont \textcolor{BrickRed}{$ \blacktriangledown\,44.5\times $}}\\[-.2pt]3m\,9.2s\\[-.2pt]{\fontsize{5.8}{6.4}\selectfont \textcolor{BrickRed}{$ \blacktriangledown\,23.3\times $}}\end{tabular} & \textcolor{gray}{\ensuremath{\cdot}} & \textcolor{gray}{\ensuremath{\cdot}} & \textcolor{gray}{\ensuremath{\cdot}} & \textcolor{gray}{\ensuremath{\cdot}} & \textcolor{gray}{\ensuremath{\cdot}} & \textcolor{gray}{\ensuremath{\cdot}} & \textcolor{gray}{\ensuremath{\cdot}} & \textcolor{gray}{\ensuremath{\cdot}}\\[3pt]\hline
\end{tabular}\par}
\smallskip{\footnotesize\noindent Each cell: A and B times, $\cdot$: no measurement. ---\textsuperscript{g}: GPU unavailable. n/a: no comparison.\par
\noindent Blue: one lowest A and one lowest B measurement for each column for each configuration.\par
\noindent Gray: minimap2. Yellow: my preferred MPF / Bounded / SIMD-DP / Block configuration.\par
\noindent MPF/MPS: minimizers per fragment/sequence. SIMD-ref: naive SIMD scan. Bounded: bounded SIMD algorithm.\par
}
\clearpage
\noindent{\large\bfseries Table \thetable: (A): Local machine\hfill 11/12}\par\smallskip
{\fontsize{7.1}{8.1}\selectfont\setlength{\tabcolsep}{1pt}\setlength{\fboxsep}{2pt}\renewcommand{\arraystretch}{1.12}
\noindent\begin{tabular}{@{}p{14mm}p{27mm}|p{13.55mm}|p{13.55mm}|p{13.55mm}|p{13.55mm}|p{13.55mm}|p{13.55mm}|p{13.55mm}|p{13.55mm}|p{13.55mm}|p{13.55mm}|p{13.55mm}|p{13.55mm}|p{13.55mm}|p{13.55mm}|p{13.55mm}|p{13.55mm}|@{}}\hline
Engine & Execution & \multicolumn{8}{c|}{\textbf{Index}} & \multicolumn{8}{c|}{\textbf{No index}}\\
 &  & \multicolumn{4}{c|}{\textbf{4 Gbp}} & \multicolumn{4}{c|}{\textbf{2 Gbp}} & \multicolumn{4}{c|}{\textbf{4 Gbp}} & \multicolumn{4}{c|}{\textbf{2 Gbp}}\\
 &  & \multicolumn{2}{c|}{\textbf{k=15, w=10}} & \multicolumn{2}{c|}{\textbf{k=20, w=15}} & \multicolumn{2}{c|}{\textbf{k=15, w=10}} & \multicolumn{2}{c|}{\textbf{k=20, w=15}} & \multicolumn{2}{c|}{\textbf{k=15, w=10}} & \multicolumn{2}{c|}{\textbf{k=20, w=15}} & \multicolumn{2}{c|}{\textbf{k=15, w=10}} & \multicolumn{2}{c|}{\textbf{k=20, w=15}}\\
 &  & \multicolumn{1}{c|}{\textbf{DNA4}} & \multicolumn{1}{c|}{\textbf{DNA5}} & \multicolumn{1}{c|}{\textbf{DNA4}} & \multicolumn{1}{c|}{\textbf{DNA5}} & \multicolumn{1}{c|}{\textbf{DNA4}} & \multicolumn{1}{c|}{\textbf{DNA5}} & \multicolumn{1}{c|}{\textbf{DNA4}} & \multicolumn{1}{c|}{\textbf{DNA5}} & \multicolumn{1}{c|}{\textbf{DNA4}} & \multicolumn{1}{c|}{\textbf{DNA5}} & \multicolumn{1}{c|}{\textbf{DNA4}} & \multicolumn{1}{c|}{\textbf{DNA5}} & \multicolumn{1}{c|}{\textbf{DNA4}} & \multicolumn{1}{c|}{\textbf{DNA5}} & \multicolumn{1}{c|}{\textbf{DNA4}} & \multicolumn{1}{c|}{\textbf{DNA5}}\\
 & & {\color{gray}C1} & {\color{gray}C2} & {\color{gray}C3} & {\color{gray}C4} & {\color{gray}C5} & {\color{gray}C6} & {\color{gray}C7} & {\color{gray}C8} & {\color{gray}C9} & {\color{gray}C10} & {\color{gray}C11} & {\color{gray}C12} & {\color{gray}C13} & {\color{gray}C14} & {\color{gray}C15} & {\color{gray}C16}\\
\hline
\rowcolor{BenchBaseline}[0.5pt][0.5pt]\shortstack[l]{{\color{gray}R1}\\\textbf{minimap2}} & \textbf{Baseline} & \hspace*{1pt}\begin{tabular}[t]{@{}c@{}}\rule{0pt}{12pt}39.7s\\[-.2pt]{\fontsize{5.8}{6.4}\selectfont \textcolor{BrickRed}{$ \blacktriangledown\,2.2\times $}}\\[-.2pt]36.8s\\[-.2pt]{\fontsize{5.8}{6.4}\selectfont \textcolor{BrickRed}{$ \blacktriangledown\,2.2\times $}}\end{tabular} & \hspace*{1pt}\begin{tabular}[t]{@{}c@{}}\rule{0pt}{12pt}25.2s\\[-.2pt]{\fontsize{5.8}{6.4}\selectfont \textcolor{BrickRed}{$ \blacktriangledown\,1.5\times $}}\\[-.2pt]22.7s\\[-.2pt]{\fontsize{5.8}{6.4}\selectfont \textcolor{BrickRed}{$ \blacktriangledown\,1.5\times $}}\end{tabular} & \hspace*{1pt}\begin{tabular}[t]{@{}c@{}}\rule{0pt}{12pt}1m\,56.9s\\[-.2pt]{\fontsize{5.8}{6.4}\selectfont \textcolor{BrickRed}{$ \blacktriangledown\,29.1\times $}}\\[-.2pt]1m\,53.2s\\[-.2pt]{\fontsize{5.8}{6.4}\selectfont \textcolor{BrickRed}{$ \blacktriangledown\,32.9\times $}}\end{tabular} & \hspace*{1pt}\begin{tabular}[t]{@{}c@{}}\rule{0pt}{12pt}24.9s\\[-.2pt]{\fontsize{5.8}{6.4}\selectfont \textcolor{BrickRed}{$ \blacktriangledown\,1.6\times $}}\\[-.2pt]22.2s\\[-.2pt]{\fontsize{5.8}{6.4}\selectfont \textcolor{BrickRed}{$ \blacktriangledown\,1.6\times $}}\end{tabular} & \hspace*{1pt}\begin{tabular}[t]{@{}c@{}}\rule{0pt}{12pt}28.6s\\[-.2pt]{\fontsize{5.8}{6.4}\selectfont \textcolor{BrickRed}{$ \blacktriangledown\,11.5\times $}}\\[-.2pt]25.8s\\[-.2pt]{\fontsize{5.8}{6.4}\selectfont \textcolor{BrickRed}{$ \blacktriangledown\,15.6\times $}}\end{tabular} & \hspace*{1pt}\begin{tabular}[t]{@{}c@{}}\rule{0pt}{12pt}15.3s\\[-.2pt]{\fontsize{5.8}{6.4}\selectfont \textcolor{BrickRed}{$ \blacktriangledown\,2.8\times $}}\\[-.2pt]12.7s\\[-.2pt]{\fontsize{5.8}{6.4}\selectfont \textcolor{BrickRed}{$ \blacktriangledown\,3.3\times $}}\end{tabular} & \hspace*{1pt}\begin{tabular}[t]{@{}c@{}}\rule{0pt}{12pt}29.3s\\[-.2pt]{\fontsize{5.8}{6.4}\selectfont \textcolor{BrickRed}{$ \blacktriangledown\,15.0\times $}}\\[-.2pt]26.2s\\[-.2pt]{\fontsize{5.8}{6.4}\selectfont \textcolor{BrickRed}{$ \blacktriangledown\,22.9\times $}}\end{tabular} & \hspace*{1pt}\begin{tabular}[t]{@{}c@{}}\rule{0pt}{12pt}8.7s\\[-.2pt]{\fontsize{5.8}{6.4}\selectfont \textcolor{BrickRed}{$ \blacktriangledown\,1.7\times $}}\\[-.2pt]8.1s\\[-.2pt]{\fontsize{5.8}{6.4}\selectfont \textcolor{BrickRed}{$ \blacktriangledown\,2.3\times $}}\end{tabular} & \hspace*{1pt}\begin{tabular}[t]{@{}c@{}}\rule{0pt}{12pt}2m\,52.4s\\[-.2pt]{\fontsize{5.8}{6.4}\selectfont \textcolor{ForestGreen}{$ \blacktriangle\,19\% $}}\\[-.2pt]2m\,48.5s\\[-.2pt]{\fontsize{5.8}{6.4}\selectfont \textcolor{ForestGreen}{$ \blacktriangle\,15\% $}}\end{tabular} & \hspace*{1pt}\begin{tabular}[t]{@{}c@{}}\rule{0pt}{12pt}2m\,17.1s\\[-.2pt]{\fontsize{5.8}{6.4}\selectfont \textcolor{ForestGreen}{$ \blacktriangle\,3.9\times $}}\\[-.2pt]2m\,15.9s\\[-.2pt]{\fontsize{5.8}{6.4}\selectfont \textcolor{ForestGreen}{$ \blacktriangle\,1.7\times $}}\end{tabular} & \hspace*{1pt}\begin{tabular}[t]{@{}c@{}}\rule{0pt}{12pt}2m\,55s\\[-.2pt]{\fontsize{5.8}{6.4}\selectfont \textcolor{BrickRed}{$ \blacktriangledown\,6\% $}}\\[-.2pt]2m\,51s\\[-.2pt]{\fontsize{5.8}{6.4}\selectfont \textcolor{BrickRed}{$ \blacktriangledown\,12\% $}}\end{tabular} & \hspace*{1pt}\begin{tabular}[t]{@{}c@{}}\rule{0pt}{12pt}55.9s\\[-.2pt]{\fontsize{5.8}{6.4}\selectfont \textcolor{ForestGreen}{$ \blacktriangle\,8.6\times $}}\\[-.2pt]55.1s\\[-.2pt]{\fontsize{5.8}{6.4}\selectfont \textcolor{ForestGreen}{$ \blacktriangle\,3.1\times $}}\end{tabular} & \hspace*{1pt}\begin{tabular}[t]{@{}c@{}}\rule{0pt}{12pt}43.3s\\[-.2pt]{\fontsize{5.8}{6.4}\selectfont \textcolor{ForestGreen}{$ \blacktriangle\,2.6\times $}}\\[-.2pt]42.5s\\[-.2pt]{\fontsize{5.8}{6.4}\selectfont \textcolor{ForestGreen}{$ \blacktriangle\,2.4\times $}}\end{tabular} & \hspace*{1pt}\begin{tabular}[t]{@{}c@{}}\rule{0pt}{12pt}38.9s\\[-.2pt]{\fontsize{5.8}{6.4}\selectfont \textcolor{ForestGreen}{$ \blacktriangle\,8.0\times $}}\\[-.2pt]38.2s\\[-.2pt]{\fontsize{5.8}{6.4}\selectfont \textcolor{ForestGreen}{$ \blacktriangle\,2.8\times $}}\end{tabular} & \hspace*{1pt}\begin{tabular}[t]{@{}c@{}}\rule{0pt}{12pt}53.1s\\[-.2pt]{\fontsize{5.8}{6.4}\selectfont \textcolor{ForestGreen}{$ \blacktriangle\,1.6\times $}}\\[-.2pt]48.2s\\[-.2pt]{\fontsize{5.8}{6.4}\selectfont \textcolor{ForestGreen}{$ \blacktriangle\,1.5\times $}}\end{tabular} & \hspace*{1pt}\begin{tabular}[t]{@{}c@{}}\rule{0pt}{12pt}28.7s\\[-.2pt]{\fontsize{5.8}{6.4}\selectfont \textcolor{ForestGreen}{$ \blacktriangle\,9.8\times $}}\\[-.2pt]28.2s\\[-.2pt]{\fontsize{5.8}{6.4}\selectfont \textcolor{ForestGreen}{$ \blacktriangle\,2.9\times $}}\end{tabular}\\[3pt]\hline
\BenchmarkSourceLabel{bench:local:row:mpf-vulkan-vulkan-block}{R67}\shortstack[l]{{\color{gray}R67}\\\textbf{BioseqDB}} & \begin{tabular}[t]{@{}l@{}}\rule{0pt}{12pt}MPF / Vulkan\\Vulkan-DP / Block\end{tabular} & \hspace*{1pt}\begin{tabular}[t]{@{}c@{}}\rule{0pt}{12pt}9m\,36.5s\\[-.2pt]{\fontsize{5.8}{6.4}\selectfont \textcolor{BrickRed}{$ \blacktriangledown\,14.5\times $}}\\[-.2pt]9m\,22.3s\\[-.2pt]{\fontsize{5.8}{6.4}\selectfont \textcolor{BrickRed}{$ \blacktriangledown\,15.3\times $}}\end{tabular} & \hspace*{1pt}\begin{tabular}[t]{@{}c@{}}\rule{0pt}{12pt}13m\,49.8s\\[-.2pt]{\fontsize{5.8}{6.4}\selectfont \textcolor{BrickRed}{$ \blacktriangledown\,32.9\times $}}\\[-.2pt]8m\,45.2s\\[-.2pt]{\fontsize{5.8}{6.4}\selectfont \textcolor{BrickRed}{$ \blacktriangledown\,23.2\times $}}\end{tabular} & \hspace*{1pt}\begin{tabular}[t]{@{}c@{}}\rule{0pt}{12pt}4m\,57.7s\\[-.2pt]{\fontsize{5.8}{6.4}\selectfont \textcolor{BrickRed}{$ \blacktriangledown\,2.5\times $}}\\[-.2pt]4m\,42.7s\\[-.2pt]{\fontsize{5.8}{6.4}\selectfont \textcolor{BrickRed}{$ \blacktriangledown\,2.5\times $}}\end{tabular} & \hspace*{1pt}\begin{tabular}[t]{@{}c@{}}\rule{0pt}{12pt}10m\,50.3s\\[-.2pt]{\fontsize{5.8}{6.4}\selectfont \textcolor{BrickRed}{$ \blacktriangledown\,26.1\times $}}\\[-.2pt]5m\,42.9s\\[-.2pt]{\fontsize{5.8}{6.4}\selectfont \textcolor{BrickRed}{$ \blacktriangledown\,15.4\times $}}\end{tabular} & \hspace*{1pt}\begin{tabular}[t]{@{}c@{}}\rule{0pt}{12pt}6m\,46.4s\\[-.2pt]{\fontsize{5.8}{6.4}\selectfont \textcolor{BrickRed}{$ \blacktriangledown\,14.2\times $}}\\[-.2pt]6m\,34.1s\\[-.2pt]{\fontsize{5.8}{6.4}\selectfont \textcolor{BrickRed}{$ \blacktriangledown\,15.3\times $}}\end{tabular} & \hspace*{1pt}\begin{tabular}[t]{@{}c@{}}\rule{0pt}{12pt}8m\,46.4s\\[-.2pt]{\fontsize{5.8}{6.4}\selectfont \textcolor{BrickRed}{$ \blacktriangledown\,34.5\times $}}\\[-.2pt]5m\,22.9s\\[-.2pt]{\fontsize{5.8}{6.4}\selectfont \textcolor{BrickRed}{$ \blacktriangledown\,25.4\times $}}\end{tabular} & \hspace*{1pt}\begin{tabular}[t]{@{}c@{}}\rule{0pt}{12pt}3m\,11.5s\\[-.2pt]{\fontsize{5.8}{6.4}\selectfont \textcolor{BrickRed}{$ \blacktriangledown\,6.5\times $}}\\[-.2pt]3m\,2.2s\\[-.2pt]{\fontsize{5.8}{6.4}\selectfont \textcolor{BrickRed}{$ \blacktriangledown\,6.9\times $}}\end{tabular} & \hspace*{1pt}\begin{tabular}[t]{@{}c@{}}\rule{0pt}{12pt}6m\,30.8s\\[-.2pt]{\fontsize{5.8}{6.4}\selectfont \textcolor{BrickRed}{$ \blacktriangledown\,45.0\times $}}\\[-.2pt]3m\,13.4s\\[-.2pt]{\fontsize{5.8}{6.4}\selectfont \textcolor{BrickRed}{$ \blacktriangledown\,23.8\times $}}\end{tabular} & \hspace*{1pt}\begin{tabular}[t]{@{}c@{}}\rule{0pt}{12pt}9m\,50.5s\\[-.2pt]{\fontsize{5.8}{6.4}\selectfont \textcolor{BrickRed}{$ \blacktriangledown\,3.4\times $}}\\[-.2pt]9m\,37.9s\\[-.2pt]{\fontsize{5.8}{6.4}\selectfont \textcolor{BrickRed}{$ \blacktriangledown\,3.4\times $}}\end{tabular} & \hspace*{1pt}\begin{tabular}[t]{@{}c@{}}\rule{0pt}{12pt}14m\,21.8s\\[-.2pt]{\fontsize{5.8}{6.4}\selectfont \textcolor{BrickRed}{$ \blacktriangledown\,6.3\times $}}\\[-.2pt]9m\,10.7s\\[-.2pt]{\fontsize{5.8}{6.4}\selectfont \textcolor{BrickRed}{$ \blacktriangledown\,4.1\times $}}\end{tabular} & \hspace*{1pt}\begin{tabular}[t]{@{}c@{}}\rule{0pt}{12pt}4m\,10.1s\\[-.2pt]{\fontsize{5.8}{6.4}\selectfont \textcolor{BrickRed}{$ \blacktriangledown\,1.4\times $}}\\[-.2pt]4m\,1s\\[-.2pt]{\fontsize{5.8}{6.4}\selectfont \textcolor{BrickRed}{$ \blacktriangledown\,1.4\times $}}\end{tabular} & \hspace*{1pt}\begin{tabular}[t]{@{}c@{}}\rule{0pt}{12pt}9m\,51.7s\\[-.2pt]{\fontsize{5.8}{6.4}\selectfont \textcolor{BrickRed}{$ \blacktriangledown\,10.6\times $}}\\[-.2pt]4m\,35.5s\\[-.2pt]{\fontsize{5.8}{6.4}\selectfont \textcolor{BrickRed}{$ \blacktriangledown\,5.0\times $}}\end{tabular} & \hspace*{1pt}\begin{tabular}[t]{@{}c@{}}\rule{0pt}{12pt}6m\,16.9s\\[-.2pt]{\fontsize{5.8}{6.4}\selectfont \textcolor{BrickRed}{$ \blacktriangledown\,8.7\times $}}\\[-.2pt]6m\,6.9s\\[-.2pt]{\fontsize{5.8}{6.4}\selectfont \textcolor{BrickRed}{$ \blacktriangledown\,8.6\times $}}\end{tabular} & \hspace*{1pt}\begin{tabular}[t]{@{}c@{}}\rule{0pt}{12pt}9m\,20.4s\\[-.2pt]{\fontsize{5.8}{6.4}\selectfont \textcolor{BrickRed}{$ \blacktriangledown\,14.4\times $}}\\[-.2pt]5m\,57.2s\\[-.2pt]{\fontsize{5.8}{6.4}\selectfont \textcolor{BrickRed}{$ \blacktriangledown\,9.4\times $}}\end{tabular} & \hspace*{1pt}\begin{tabular}[t]{@{}c@{}}\rule{0pt}{12pt}3m\,10.7s\\[-.2pt]{\fontsize{5.8}{6.4}\selectfont \textcolor{BrickRed}{$ \blacktriangledown\,3.6\times $}}\\[-.2pt]3m\,0.2s\\[-.2pt]{\fontsize{5.8}{6.4}\selectfont \textcolor{BrickRed}{$ \blacktriangledown\,3.7\times $}}\end{tabular} & \hspace*{1pt}\begin{tabular}[t]{@{}c@{}}\rule{0pt}{12pt}6m\,41.2s\\[-.2pt]{\fontsize{5.8}{6.4}\selectfont \textcolor{BrickRed}{$ \blacktriangledown\,14.0\times $}}\\[-.2pt]3m\,21.5s\\[-.2pt]{\fontsize{5.8}{6.4}\selectfont \textcolor{BrickRed}{$ \blacktriangledown\,7.1\times $}}\end{tabular}\\[3pt]\hline
\BenchmarkSourceLabel{bench:local:row:mpf-vulkan-rmq-ksw2}{R68}\shortstack[l]{{\color{gray}R68}} & \begin{tabular}[t]{@{}l@{}}\rule{0pt}{12pt}MPF / Vulkan\\RMQ / KSW2\end{tabular} & \textcolor{gray}{\ensuremath{\cdot}} & \textcolor{gray}{\ensuremath{\cdot}} & \textcolor{gray}{\ensuremath{\cdot}} & \textcolor{gray}{\ensuremath{\cdot}} & \hspace*{1pt}\begin{tabular}[t]{@{}c@{}}\rule{0pt}{12pt}5m\,51.9s\\[-.2pt]{\fontsize{5.8}{6.4}\selectfont \textcolor{BrickRed}{$ \blacktriangledown\,12.3\times $}}\\[-.2pt]5m\,40.4s\\[-.2pt]{\fontsize{5.8}{6.4}\selectfont \textcolor{BrickRed}{$ \blacktriangledown\,13.2\times $}}\end{tabular} & \hspace*{1pt}\begin{tabular}[t]{@{}c@{}}\rule{0pt}{12pt}8m\,46.4s\\[-.2pt]{\fontsize{5.8}{6.4}\selectfont \textcolor{BrickRed}{$ \blacktriangledown\,34.5\times $}}\\[-.2pt]5m\,22.2s\\[-.2pt]{\fontsize{5.8}{6.4}\selectfont \textcolor{BrickRed}{$ \blacktriangledown\,25.3\times $}}\end{tabular} & \hspace*{1pt}\begin{tabular}[t]{@{}c@{}}\rule{0pt}{12pt}3m\,12.6s\\[-.2pt]{\fontsize{5.8}{6.4}\selectfont \textcolor{BrickRed}{$ \blacktriangledown\,6.6\times $}}\\[-.2pt]2m\,58.9s\\[-.2pt]{\fontsize{5.8}{6.4}\selectfont \textcolor{BrickRed}{$ \blacktriangledown\,6.8\times $}}\end{tabular} & \hspace*{1pt}\begin{tabular}[t]{@{}c@{}}\rule{0pt}{12pt}6m\,29.8s\\[-.2pt]{\fontsize{5.8}{6.4}\selectfont \textcolor{BrickRed}{$ \blacktriangledown\,44.9\times $}}\\[-.2pt]3m\,8.7s\\[-.2pt]{\fontsize{5.8}{6.4}\selectfont \textcolor{BrickRed}{$ \blacktriangledown\,23.2\times $}}\end{tabular} & \textcolor{gray}{\ensuremath{\cdot}} & \textcolor{gray}{\ensuremath{\cdot}} & \textcolor{gray}{\ensuremath{\cdot}} & \textcolor{gray}{\ensuremath{\cdot}} & \textcolor{gray}{\ensuremath{\cdot}} & \textcolor{gray}{\ensuremath{\cdot}} & \textcolor{gray}{\ensuremath{\cdot}} & \textcolor{gray}{\ensuremath{\cdot}}\\[3pt]\hline
\BenchmarkSourceLabel{bench:local:row:mpf-vulkan-scalar-ksw2}{R69}\shortstack[l]{{\color{gray}R69}} & \begin{tabular}[t]{@{}l@{}}\rule{0pt}{12pt}MPF / Vulkan\\Scalar-DP / KSW2\end{tabular} & \textcolor{gray}{\ensuremath{\cdot}} & \textcolor{gray}{\ensuremath{\cdot}} & \textcolor{gray}{\ensuremath{\cdot}} & \textcolor{gray}{\ensuremath{\cdot}} & \hspace*{1pt}\begin{tabular}[t]{@{}c@{}}\rule{0pt}{12pt}5m\,38.2s\\[-.2pt]{\fontsize{5.8}{6.4}\selectfont \textcolor{BrickRed}{$ \blacktriangledown\,11.8\times $}}\\[-.2pt]5m\,28s\\[-.2pt]{\fontsize{5.8}{6.4}\selectfont \textcolor{BrickRed}{$ \blacktriangledown\,12.7\times $}}\end{tabular} & \hspace*{1pt}\begin{tabular}[t]{@{}c@{}}\rule{0pt}{12pt}8m\,37.9s\\[-.2pt]{\fontsize{5.8}{6.4}\selectfont \textcolor{BrickRed}{$ \blacktriangledown\,33.9\times $}}\\[-.2pt]5m\,16s\\[-.2pt]{\fontsize{5.8}{6.4}\selectfont \textcolor{BrickRed}{$ \blacktriangledown\,24.9\times $}}\end{tabular} & \hspace*{1pt}\begin{tabular}[t]{@{}c@{}}\rule{0pt}{12pt}3m\,25.8s\\[-.2pt]{\fontsize{5.8}{6.4}\selectfont \textcolor{BrickRed}{$ \blacktriangledown\,7.0\times $}}\\[-.2pt]3m\,14.4s\\[-.2pt]{\fontsize{5.8}{6.4}\selectfont \textcolor{BrickRed}{$ \blacktriangledown\,7.4\times $}}\end{tabular} & \hspace*{1pt}\begin{tabular}[t]{@{}c@{}}\rule{0pt}{12pt}6m\,24.9s\\[-.2pt]{\fontsize{5.8}{6.4}\selectfont \textcolor{BrickRed}{$ \blacktriangledown\,44.3\times $}}\\[-.2pt]3m\,8.7s\\[-.2pt]{\fontsize{5.8}{6.4}\selectfont \textcolor{BrickRed}{$ \blacktriangledown\,23.2\times $}}\end{tabular} & \textcolor{gray}{\ensuremath{\cdot}} & \textcolor{gray}{\ensuremath{\cdot}} & \textcolor{gray}{\ensuremath{\cdot}} & \textcolor{gray}{\ensuremath{\cdot}} & \textcolor{gray}{\ensuremath{\cdot}} & \textcolor{gray}{\ensuremath{\cdot}} & \textcolor{gray}{\ensuremath{\cdot}} & \textcolor{gray}{\ensuremath{\cdot}}\\[3pt]\hline
\BenchmarkSourceLabel{bench:local:row:mpf-vulkan-simd-ksw2}{R70}\shortstack[l]{{\color{gray}R70}} & \begin{tabular}[t]{@{}l@{}}\rule{0pt}{12pt}MPF / Vulkan\\SIMD-DP / KSW2\end{tabular} & \textcolor{gray}{\ensuremath{\cdot}} & \textcolor{gray}{\ensuremath{\cdot}} & \textcolor{gray}{\ensuremath{\cdot}} & \textcolor{gray}{\ensuremath{\cdot}} & \hspace*{1pt}\begin{tabular}[t]{@{}c@{}}\rule{0pt}{12pt}5m\,44.6s\\[-.2pt]{\fontsize{5.8}{6.4}\selectfont \textcolor{BrickRed}{$ \blacktriangledown\,12.1\times $}}\\[-.2pt]5m\,34.2s\\[-.2pt]{\fontsize{5.8}{6.4}\selectfont \textcolor{BrickRed}{$ \blacktriangledown\,12.9\times $}}\end{tabular} & \hspace*{1pt}\begin{tabular}[t]{@{}c@{}}\rule{0pt}{12pt}8m\,41.2s\\[-.2pt]{\fontsize{5.8}{6.4}\selectfont \textcolor{BrickRed}{$ \blacktriangledown\,34.1\times $}}\\[-.2pt]5m\,20s\\[-.2pt]{\fontsize{5.8}{6.4}\selectfont \textcolor{BrickRed}{$ \blacktriangledown\,25.2\times $}}\end{tabular} & \hspace*{1pt}\begin{tabular}[t]{@{}c@{}}\rule{0pt}{12pt}3m\,33.4s\\[-.2pt]{\fontsize{5.8}{6.4}\selectfont \textcolor{BrickRed}{$ \blacktriangledown\,7.3\times $}}\\[-.2pt]3m\,19s\\[-.2pt]{\fontsize{5.8}{6.4}\selectfont \textcolor{BrickRed}{$ \blacktriangledown\,7.6\times $}}\end{tabular} & \hspace*{1pt}\begin{tabular}[t]{@{}c@{}}\rule{0pt}{12pt}6m\,39.3s\\[-.2pt]{\fontsize{5.8}{6.4}\selectfont \textcolor{BrickRed}{$ \blacktriangledown\,46.0\times $}}\\[-.2pt]3m\,17.5s\\[-.2pt]{\fontsize{5.8}{6.4}\selectfont \textcolor{BrickRed}{$ \blacktriangledown\,24.3\times $}}\end{tabular} & \textcolor{gray}{\ensuremath{\cdot}} & \textcolor{gray}{\ensuremath{\cdot}} & \textcolor{gray}{\ensuremath{\cdot}} & \textcolor{gray}{\ensuremath{\cdot}} & \textcolor{gray}{\ensuremath{\cdot}} & \textcolor{gray}{\ensuremath{\cdot}} & \textcolor{gray}{\ensuremath{\cdot}} & \textcolor{gray}{\ensuremath{\cdot}}\\[3pt]\hline
\BenchmarkSourceLabel{bench:local:row:mpf-vulkan-scalar-block}{R71}\shortstack[l]{{\color{gray}R71}} & \begin{tabular}[t]{@{}l@{}}\rule{0pt}{12pt}MPF / Vulkan\\Scalar-DP / Block\end{tabular} & \textcolor{gray}{\ensuremath{\cdot}} & \textcolor{gray}{\ensuremath{\cdot}} & \textcolor{gray}{\ensuremath{\cdot}} & \textcolor{gray}{\ensuremath{\cdot}} & \hspace*{1pt}\begin{tabular}[t]{@{}c@{}}\rule{0pt}{12pt}5m\,40.3s\\[-.2pt]{\fontsize{5.8}{6.4}\selectfont \textcolor{BrickRed}{$ \blacktriangledown\,11.9\times $}}\\[-.2pt]5m\,29.9s\\[-.2pt]{\fontsize{5.8}{6.4}\selectfont \textcolor{BrickRed}{$ \blacktriangledown\,12.8\times $}}\end{tabular} & \hspace*{1pt}\begin{tabular}[t]{@{}c@{}}\rule{0pt}{12pt}8m\,42.9s\\[-.2pt]{\fontsize{5.8}{6.4}\selectfont \textcolor{BrickRed}{$ \blacktriangledown\,34.2\times $}}\\[-.2pt]5m\,22.1s\\[-.2pt]{\fontsize{5.8}{6.4}\selectfont \textcolor{BrickRed}{$ \blacktriangledown\,25.3\times $}}\end{tabular} & \hspace*{1pt}\begin{tabular}[t]{@{}c@{}}\rule{0pt}{12pt}3m\,41.9s\\[-.2pt]{\fontsize{5.8}{6.4}\selectfont \textcolor{BrickRed}{$ \blacktriangledown\,7.6\times $}}\\[-.2pt]3m\,31.8s\\[-.2pt]{\fontsize{5.8}{6.4}\selectfont \textcolor{BrickRed}{$ \blacktriangledown\,8.1\times $}}\end{tabular} & \hspace*{1pt}\begin{tabular}[t]{@{}c@{}}\rule{0pt}{12pt}6m\,25.4s\\[-.2pt]{\fontsize{5.8}{6.4}\selectfont \textcolor{BrickRed}{$ \blacktriangledown\,44.4\times $}}\\[-.2pt]3m\,9.3s\\[-.2pt]{\fontsize{5.8}{6.4}\selectfont \textcolor{BrickRed}{$ \blacktriangledown\,23.3\times $}}\end{tabular} & \textcolor{gray}{\ensuremath{\cdot}} & \textcolor{gray}{\ensuremath{\cdot}} & \textcolor{gray}{\ensuremath{\cdot}} & \textcolor{gray}{\ensuremath{\cdot}} & \textcolor{gray}{\ensuremath{\cdot}} & \textcolor{gray}{\ensuremath{\cdot}} & \textcolor{gray}{\ensuremath{\cdot}} & \textcolor{gray}{\ensuremath{\cdot}}\\[3pt]\hline
\BenchmarkSourceLabel{bench:local:row:mpf-vulkan-vulkan-ksw2}{R72}\shortstack[l]{{\color{gray}R72}} & \begin{tabular}[t]{@{}l@{}}\rule{0pt}{12pt}MPF / Vulkan\\Vulkan-DP / KSW2\end{tabular} & \textcolor{gray}{\ensuremath{\cdot}} & \textcolor{gray}{\ensuremath{\cdot}} & \textcolor{gray}{\ensuremath{\cdot}} & \textcolor{gray}{\ensuremath{\cdot}} & \hspace*{1pt}\begin{tabular}[t]{@{}c@{}}\rule{0pt}{12pt}6m\,5.5s\\[-.2pt]{\fontsize{5.8}{6.4}\selectfont \textcolor{BrickRed}{$ \blacktriangledown\,12.8\times $}}\\[-.2pt]5m\,53.5s\\[-.2pt]{\fontsize{5.8}{6.4}\selectfont \textcolor{BrickRed}{$ \blacktriangledown\,13.7\times $}}\end{tabular} & \hspace*{1pt}\begin{tabular}[t]{@{}c@{}}\rule{0pt}{12pt}8m\,57.9s\\[-.2pt]{\fontsize{5.8}{6.4}\selectfont \textcolor{BrickRed}{$ \blacktriangledown\,35.2\times $}}\\[-.2pt]5m\,38.3s\\[-.2pt]{\fontsize{5.8}{6.4}\selectfont \textcolor{BrickRed}{$ \blacktriangledown\,26.6\times $}}\end{tabular} & \hspace*{1pt}\begin{tabular}[t]{@{}c@{}}\rule{0pt}{12pt}3m\,46.3s\\[-.2pt]{\fontsize{5.8}{6.4}\selectfont \textcolor{BrickRed}{$ \blacktriangledown\,7.7\times $}}\\[-.2pt]3m\,33.1s\\[-.2pt]{\fontsize{5.8}{6.4}\selectfont \textcolor{BrickRed}{$ \blacktriangledown\,8.1\times $}}\end{tabular} & \hspace*{1pt}\begin{tabular}[t]{@{}c@{}}\rule{0pt}{12pt}6m\,39s\\[-.2pt]{\fontsize{5.8}{6.4}\selectfont \textcolor{BrickRed}{$ \blacktriangledown\,45.9\times $}}\\[-.2pt]3m\,23.4s\\[-.2pt]{\fontsize{5.8}{6.4}\selectfont \textcolor{BrickRed}{$ \blacktriangledown\,25.0\times $}}\end{tabular} & \textcolor{gray}{\ensuremath{\cdot}} & \textcolor{gray}{\ensuremath{\cdot}} & \textcolor{gray}{\ensuremath{\cdot}} & \textcolor{gray}{\ensuremath{\cdot}} & \textcolor{gray}{\ensuremath{\cdot}} & \textcolor{gray}{\ensuremath{\cdot}} & \textcolor{gray}{\ensuremath{\cdot}} & \textcolor{gray}{\ensuremath{\cdot}}\\[3pt]\hline
\BenchmarkSourceLabel{bench:local:row:mpf-vulkan-simd-block}{R73}\shortstack[l]{{\color{gray}R73}} & \begin{tabular}[t]{@{}l@{}}\rule{0pt}{12pt}MPF / Vulkan\\SIMD-DP / Block\end{tabular} & \textcolor{gray}{\ensuremath{\cdot}} & \textcolor{gray}{\ensuremath{\cdot}} & \textcolor{gray}{\ensuremath{\cdot}} & \textcolor{gray}{\ensuremath{\cdot}} & \hspace*{1pt}\begin{tabular}[t]{@{}c@{}}\rule{0pt}{12pt}5m\,43.9s\\[-.2pt]{\fontsize{5.8}{6.4}\selectfont \textcolor{BrickRed}{$ \blacktriangledown\,12.0\times $}}\\[-.2pt]5m\,34.1s\\[-.2pt]{\fontsize{5.8}{6.4}\selectfont \textcolor{BrickRed}{$ \blacktriangledown\,12.9\times $}}\end{tabular} & \hspace*{1pt}\begin{tabular}[t]{@{}c@{}}\rule{0pt}{12pt}8m\,37.5s\\[-.2pt]{\fontsize{5.8}{6.4}\selectfont \textcolor{BrickRed}{$ \blacktriangledown\,33.9\times $}}\\[-.2pt]5m\,14.9s\\[-.2pt]{\fontsize{5.8}{6.4}\selectfont \textcolor{BrickRed}{$ \blacktriangledown\,24.8\times $}}\end{tabular} & \hspace*{1pt}\begin{tabular}[t]{@{}c@{}}\rule{0pt}{12pt}4m\,42.2s\\[-.2pt]{\fontsize{5.8}{6.4}\selectfont \textcolor{BrickRed}{$ \blacktriangledown\,9.6\times $}}\\[-.2pt]4m\,30.5s\\[-.2pt]{\fontsize{5.8}{6.4}\selectfont \textcolor{BrickRed}{$ \blacktriangledown\,10.3\times $}}\end{tabular} & \hspace*{1pt}\begin{tabular}[t]{@{}c@{}}\rule{0pt}{12pt}6m\,35.3s\\[-.2pt]{\fontsize{5.8}{6.4}\selectfont \textcolor{BrickRed}{$ \blacktriangledown\,45.5\times $}}\\[-.2pt]3m\,19.5s\\[-.2pt]{\fontsize{5.8}{6.4}\selectfont \textcolor{BrickRed}{$ \blacktriangledown\,24.5\times $}}\end{tabular} & \textcolor{gray}{\ensuremath{\cdot}} & \textcolor{gray}{\ensuremath{\cdot}} & \textcolor{gray}{\ensuremath{\cdot}} & \textcolor{gray}{\ensuremath{\cdot}} & \textcolor{gray}{\ensuremath{\cdot}} & \textcolor{gray}{\ensuremath{\cdot}} & \textcolor{gray}{\ensuremath{\cdot}} & \textcolor{gray}{\ensuremath{\cdot}}\\[3pt]\hline
\end{tabular}\par}
\smallskip{\footnotesize\noindent Each cell: A and B times, $\cdot$: no measurement. ---\textsuperscript{g}: GPU unavailable. n/a: no comparison.\par
\noindent Blue: one lowest A and one lowest B measurement for each column for each configuration.\par
\noindent Gray: minimap2. Yellow: my preferred MPF / Bounded / SIMD-DP / Block configuration.\par
\noindent MPF/MPS: minimizers per fragment/sequence. SIMD-ref: naive SIMD scan. Bounded: bounded SIMD algorithm.\par
}
\clearpage
\noindent{\large\bfseries Table \thetable: (A): Local machine\hfill 12/12}\par\smallskip
{\fontsize{7.1}{8.1}\selectfont\setlength{\tabcolsep}{1pt}\setlength{\fboxsep}{2pt}\renewcommand{\arraystretch}{1.12}
\noindent\begin{tabular}{@{}p{14mm}p{27mm}|p{13.55mm}|p{13.55mm}|p{13.55mm}|p{13.55mm}|p{13.55mm}|p{13.55mm}|p{13.55mm}|p{13.55mm}|p{13.55mm}|p{13.55mm}|p{13.55mm}|p{13.55mm}|p{13.55mm}|p{13.55mm}|p{13.55mm}|p{13.55mm}|@{}}\hline
Engine & Execution & \multicolumn{8}{c|}{\textbf{Index}} & \multicolumn{8}{c|}{\textbf{No index}}\\
 &  & \multicolumn{4}{c|}{\textbf{4 Gbp}} & \multicolumn{4}{c|}{\textbf{2 Gbp}} & \multicolumn{4}{c|}{\textbf{4 Gbp}} & \multicolumn{4}{c|}{\textbf{2 Gbp}}\\
 &  & \multicolumn{2}{c|}{\textbf{k=15, w=10}} & \multicolumn{2}{c|}{\textbf{k=20, w=15}} & \multicolumn{2}{c|}{\textbf{k=15, w=10}} & \multicolumn{2}{c|}{\textbf{k=20, w=15}} & \multicolumn{2}{c|}{\textbf{k=15, w=10}} & \multicolumn{2}{c|}{\textbf{k=20, w=15}} & \multicolumn{2}{c|}{\textbf{k=15, w=10}} & \multicolumn{2}{c|}{\textbf{k=20, w=15}}\\
 &  & \multicolumn{1}{c|}{\textbf{DNA4}} & \multicolumn{1}{c|}{\textbf{DNA5}} & \multicolumn{1}{c|}{\textbf{DNA4}} & \multicolumn{1}{c|}{\textbf{DNA5}} & \multicolumn{1}{c|}{\textbf{DNA4}} & \multicolumn{1}{c|}{\textbf{DNA5}} & \multicolumn{1}{c|}{\textbf{DNA4}} & \multicolumn{1}{c|}{\textbf{DNA5}} & \multicolumn{1}{c|}{\textbf{DNA4}} & \multicolumn{1}{c|}{\textbf{DNA5}} & \multicolumn{1}{c|}{\textbf{DNA4}} & \multicolumn{1}{c|}{\textbf{DNA5}} & \multicolumn{1}{c|}{\textbf{DNA4}} & \multicolumn{1}{c|}{\textbf{DNA5}} & \multicolumn{1}{c|}{\textbf{DNA4}} & \multicolumn{1}{c|}{\textbf{DNA5}}\\
 & & {\color{gray}C1} & {\color{gray}C2} & {\color{gray}C3} & {\color{gray}C4} & {\color{gray}C5} & {\color{gray}C6} & {\color{gray}C7} & {\color{gray}C8} & {\color{gray}C9} & {\color{gray}C10} & {\color{gray}C11} & {\color{gray}C12} & {\color{gray}C13} & {\color{gray}C14} & {\color{gray}C15} & {\color{gray}C16}\\
\hline
\rowcolor{BenchBaseline}[0.5pt][0.5pt]\shortstack[l]{{\color{gray}R1}\\\textbf{minimap2}} & \textbf{Baseline} & \hspace*{1pt}\begin{tabular}[t]{@{}c@{}}\rule{0pt}{12pt}39.7s\\[-.2pt]{\fontsize{5.8}{6.4}\selectfont \textcolor{BrickRed}{$ \blacktriangledown\,2.2\times $}}\\[-.2pt]36.8s\\[-.2pt]{\fontsize{5.8}{6.4}\selectfont \textcolor{BrickRed}{$ \blacktriangledown\,2.2\times $}}\end{tabular} & \hspace*{1pt}\begin{tabular}[t]{@{}c@{}}\rule{0pt}{12pt}25.2s\\[-.2pt]{\fontsize{5.8}{6.4}\selectfont \textcolor{BrickRed}{$ \blacktriangledown\,1.5\times $}}\\[-.2pt]22.7s\\[-.2pt]{\fontsize{5.8}{6.4}\selectfont \textcolor{BrickRed}{$ \blacktriangledown\,1.5\times $}}\end{tabular} & \hspace*{1pt}\begin{tabular}[t]{@{}c@{}}\rule{0pt}{12pt}1m\,56.9s\\[-.2pt]{\fontsize{5.8}{6.4}\selectfont \textcolor{BrickRed}{$ \blacktriangledown\,29.1\times $}}\\[-.2pt]1m\,53.2s\\[-.2pt]{\fontsize{5.8}{6.4}\selectfont \textcolor{BrickRed}{$ \blacktriangledown\,32.9\times $}}\end{tabular} & \hspace*{1pt}\begin{tabular}[t]{@{}c@{}}\rule{0pt}{12pt}24.9s\\[-.2pt]{\fontsize{5.8}{6.4}\selectfont \textcolor{BrickRed}{$ \blacktriangledown\,1.6\times $}}\\[-.2pt]22.2s\\[-.2pt]{\fontsize{5.8}{6.4}\selectfont \textcolor{BrickRed}{$ \blacktriangledown\,1.6\times $}}\end{tabular} & \hspace*{1pt}\begin{tabular}[t]{@{}c@{}}\rule{0pt}{12pt}28.6s\\[-.2pt]{\fontsize{5.8}{6.4}\selectfont \textcolor{BrickRed}{$ \blacktriangledown\,11.5\times $}}\\[-.2pt]25.8s\\[-.2pt]{\fontsize{5.8}{6.4}\selectfont \textcolor{BrickRed}{$ \blacktriangledown\,15.6\times $}}\end{tabular} & \hspace*{1pt}\begin{tabular}[t]{@{}c@{}}\rule{0pt}{12pt}15.3s\\[-.2pt]{\fontsize{5.8}{6.4}\selectfont \textcolor{BrickRed}{$ \blacktriangledown\,2.8\times $}}\\[-.2pt]12.7s\\[-.2pt]{\fontsize{5.8}{6.4}\selectfont \textcolor{BrickRed}{$ \blacktriangledown\,3.3\times $}}\end{tabular} & \hspace*{1pt}\begin{tabular}[t]{@{}c@{}}\rule{0pt}{12pt}29.3s\\[-.2pt]{\fontsize{5.8}{6.4}\selectfont \textcolor{BrickRed}{$ \blacktriangledown\,15.0\times $}}\\[-.2pt]26.2s\\[-.2pt]{\fontsize{5.8}{6.4}\selectfont \textcolor{BrickRed}{$ \blacktriangledown\,22.9\times $}}\end{tabular} & \hspace*{1pt}\begin{tabular}[t]{@{}c@{}}\rule{0pt}{12pt}8.7s\\[-.2pt]{\fontsize{5.8}{6.4}\selectfont \textcolor{BrickRed}{$ \blacktriangledown\,1.7\times $}}\\[-.2pt]8.1s\\[-.2pt]{\fontsize{5.8}{6.4}\selectfont \textcolor{BrickRed}{$ \blacktriangledown\,2.3\times $}}\end{tabular} & \hspace*{1pt}\begin{tabular}[t]{@{}c@{}}\rule{0pt}{12pt}2m\,52.4s\\[-.2pt]{\fontsize{5.8}{6.4}\selectfont \textcolor{ForestGreen}{$ \blacktriangle\,19\% $}}\\[-.2pt]2m\,48.5s\\[-.2pt]{\fontsize{5.8}{6.4}\selectfont \textcolor{ForestGreen}{$ \blacktriangle\,15\% $}}\end{tabular} & \hspace*{1pt}\begin{tabular}[t]{@{}c@{}}\rule{0pt}{12pt}2m\,17.1s\\[-.2pt]{\fontsize{5.8}{6.4}\selectfont \textcolor{ForestGreen}{$ \blacktriangle\,3.9\times $}}\\[-.2pt]2m\,15.9s\\[-.2pt]{\fontsize{5.8}{6.4}\selectfont \textcolor{ForestGreen}{$ \blacktriangle\,1.7\times $}}\end{tabular} & \hspace*{1pt}\begin{tabular}[t]{@{}c@{}}\rule{0pt}{12pt}2m\,55s\\[-.2pt]{\fontsize{5.8}{6.4}\selectfont \textcolor{BrickRed}{$ \blacktriangledown\,6\% $}}\\[-.2pt]2m\,51s\\[-.2pt]{\fontsize{5.8}{6.4}\selectfont \textcolor{BrickRed}{$ \blacktriangledown\,12\% $}}\end{tabular} & \hspace*{1pt}\begin{tabular}[t]{@{}c@{}}\rule{0pt}{12pt}55.9s\\[-.2pt]{\fontsize{5.8}{6.4}\selectfont \textcolor{ForestGreen}{$ \blacktriangle\,8.6\times $}}\\[-.2pt]55.1s\\[-.2pt]{\fontsize{5.8}{6.4}\selectfont \textcolor{ForestGreen}{$ \blacktriangle\,3.1\times $}}\end{tabular} & \hspace*{1pt}\begin{tabular}[t]{@{}c@{}}\rule{0pt}{12pt}43.3s\\[-.2pt]{\fontsize{5.8}{6.4}\selectfont \textcolor{ForestGreen}{$ \blacktriangle\,2.6\times $}}\\[-.2pt]42.5s\\[-.2pt]{\fontsize{5.8}{6.4}\selectfont \textcolor{ForestGreen}{$ \blacktriangle\,2.4\times $}}\end{tabular} & \hspace*{1pt}\begin{tabular}[t]{@{}c@{}}\rule{0pt}{12pt}38.9s\\[-.2pt]{\fontsize{5.8}{6.4}\selectfont \textcolor{ForestGreen}{$ \blacktriangle\,8.0\times $}}\\[-.2pt]38.2s\\[-.2pt]{\fontsize{5.8}{6.4}\selectfont \textcolor{ForestGreen}{$ \blacktriangle\,2.8\times $}}\end{tabular} & \hspace*{1pt}\begin{tabular}[t]{@{}c@{}}\rule{0pt}{12pt}53.1s\\[-.2pt]{\fontsize{5.8}{6.4}\selectfont \textcolor{ForestGreen}{$ \blacktriangle\,1.6\times $}}\\[-.2pt]48.2s\\[-.2pt]{\fontsize{5.8}{6.4}\selectfont \textcolor{ForestGreen}{$ \blacktriangle\,1.5\times $}}\end{tabular} & \hspace*{1pt}\begin{tabular}[t]{@{}c@{}}\rule{0pt}{12pt}28.7s\\[-.2pt]{\fontsize{5.8}{6.4}\selectfont \textcolor{ForestGreen}{$ \blacktriangle\,9.8\times $}}\\[-.2pt]28.2s\\[-.2pt]{\fontsize{5.8}{6.4}\selectfont \textcolor{ForestGreen}{$ \blacktriangle\,2.9\times $}}\end{tabular}\\[3pt]\hline
\BenchmarkSourceLabel{bench:local:row:mps-vulkan-rmq-ksw2}{R74}\shortstack[l]{{\color{gray}R74}\\\textbf{BioseqDB}} & \begin{tabular}[t]{@{}l@{}}\rule{0pt}{12pt}MPS / Vulkan\\RMQ / KSW2\end{tabular} & \textcolor{gray}{\ensuremath{\cdot}} & \textcolor{gray}{\ensuremath{\cdot}} & \textcolor{gray}{\ensuremath{\cdot}} & \textcolor{gray}{\ensuremath{\cdot}} & \hspace*{1pt}\begin{tabular}[t]{@{}c@{}}\rule{0pt}{12pt}{\fontsize{5.6}{6.4}\selectfont 1h\,56m\,53.5s}\\[-.2pt]{\fontsize{5.8}{6.4}\selectfont \textcolor{BrickRed}{$ \blacktriangledown\,245.5\times $}}\\[-.2pt]{\fontsize{5.6}{6.4}\selectfont 1h\,56m\,41.7s}\\[-.2pt]{\fontsize{5.8}{6.4}\selectfont \textcolor{BrickRed}{$ \blacktriangledown\,271.1\times $}}\end{tabular} & \textcolor{gray}{\ensuremath{\cdot}} & \textcolor{gray}{\ensuremath{\cdot}} & \textcolor{gray}{\ensuremath{\cdot}} & \textcolor{gray}{\ensuremath{\cdot}} & \textcolor{gray}{\ensuremath{\cdot}} & \textcolor{gray}{\ensuremath{\cdot}} & \textcolor{gray}{\ensuremath{\cdot}} & \textcolor{gray}{\ensuremath{\cdot}} & \textcolor{gray}{\ensuremath{\cdot}} & \textcolor{gray}{\ensuremath{\cdot}} & \textcolor{gray}{\ensuremath{\cdot}}\\[3pt]\hline
\BenchmarkSourceLabel{bench:local:row:mps-vulkan-rmq-block}{R75}\shortstack[l]{{\color{gray}R75}} & \begin{tabular}[t]{@{}l@{}}\rule{0pt}{12pt}MPS / Vulkan\\RMQ / Block\end{tabular} & \textcolor{gray}{\ensuremath{\cdot}} & \textcolor{gray}{\ensuremath{\cdot}} & \textcolor{gray}{\ensuremath{\cdot}} & \textcolor{gray}{\ensuremath{\cdot}} & \hspace*{1pt}\begin{tabular}[t]{@{}c@{}}\rule{0pt}{12pt}{\fontsize{5.6}{6.4}\selectfont 1h\,56m\,56.6s}\\[-.2pt]{\fontsize{5.8}{6.4}\selectfont \textcolor{BrickRed}{$ \blacktriangledown\,245.6\times $}}\\[-.2pt]{\fontsize{5.6}{6.4}\selectfont 1h\,56m\,45.8s}\\[-.2pt]{\fontsize{5.8}{6.4}\selectfont \textcolor{BrickRed}{$ \blacktriangledown\,271.2\times $}}\end{tabular} & \textcolor{gray}{\ensuremath{\cdot}} & \textcolor{gray}{\ensuremath{\cdot}} & \textcolor{gray}{\ensuremath{\cdot}} & \textcolor{gray}{\ensuremath{\cdot}} & \textcolor{gray}{\ensuremath{\cdot}} & \textcolor{gray}{\ensuremath{\cdot}} & \textcolor{gray}{\ensuremath{\cdot}} & \textcolor{gray}{\ensuremath{\cdot}} & \textcolor{gray}{\ensuremath{\cdot}} & \textcolor{gray}{\ensuremath{\cdot}} & \textcolor{gray}{\ensuremath{\cdot}}\\[3pt]\hline
\BenchmarkSourceLabel{bench:local:row:mps-vulkan-simd-ksw2}{R76}\shortstack[l]{{\color{gray}R76}} & \begin{tabular}[t]{@{}l@{}}\rule{0pt}{12pt}MPS / Vulkan\\SIMD-DP / KSW2\end{tabular} & \textcolor{gray}{\ensuremath{\cdot}} & \textcolor{gray}{\ensuremath{\cdot}} & \textcolor{gray}{\ensuremath{\cdot}} & \textcolor{gray}{\ensuremath{\cdot}} & \hspace*{1pt}\begin{tabular}[t]{@{}c@{}}\rule{0pt}{12pt}{\fontsize{5.6}{6.4}\selectfont 1h\,56m\,58.7s}\\[-.2pt]{\fontsize{5.8}{6.4}\selectfont \textcolor{BrickRed}{$ \blacktriangledown\,245.7\times $}}\\[-.2pt]{\fontsize{5.6}{6.4}\selectfont 1h\,56m\,45.6s}\\[-.2pt]{\fontsize{5.8}{6.4}\selectfont \textcolor{BrickRed}{$ \blacktriangledown\,271.2\times $}}\end{tabular} & \textcolor{gray}{\ensuremath{\cdot}} & \textcolor{gray}{\ensuremath{\cdot}} & \textcolor{gray}{\ensuremath{\cdot}} & \textcolor{gray}{\ensuremath{\cdot}} & \textcolor{gray}{\ensuremath{\cdot}} & \textcolor{gray}{\ensuremath{\cdot}} & \textcolor{gray}{\ensuremath{\cdot}} & \textcolor{gray}{\ensuremath{\cdot}} & \textcolor{gray}{\ensuremath{\cdot}} & \textcolor{gray}{\ensuremath{\cdot}} & \textcolor{gray}{\ensuremath{\cdot}}\\[3pt]\hline
\BenchmarkSourceLabel{bench:local:row:mps-vulkan-scalar-block}{R77}\shortstack[l]{{\color{gray}R77}} & \begin{tabular}[t]{@{}l@{}}\rule{0pt}{12pt}MPS / Vulkan\\Scalar-DP / Block\end{tabular} & \textcolor{gray}{\ensuremath{\cdot}} & \textcolor{gray}{\ensuremath{\cdot}} & \textcolor{gray}{\ensuremath{\cdot}} & \textcolor{gray}{\ensuremath{\cdot}} & \hspace*{1pt}\begin{tabular}[t]{@{}c@{}}\rule{0pt}{12pt}{\fontsize{5.6}{6.4}\selectfont 1h\,57m\,18.3s}\\[-.2pt]{\fontsize{5.8}{6.4}\selectfont \textcolor{BrickRed}{$ \blacktriangledown\,246.4\times $}}\\[-.2pt]{\fontsize{5.6}{6.4}\selectfont 1h\,57m\,8.7s}\\[-.2pt]{\fontsize{5.8}{6.4}\selectfont \textcolor{BrickRed}{$ \blacktriangledown\,272.1\times $}}\end{tabular} & \textcolor{gray}{\ensuremath{\cdot}} & \textcolor{gray}{\ensuremath{\cdot}} & \textcolor{gray}{\ensuremath{\cdot}} & \textcolor{gray}{\ensuremath{\cdot}} & \textcolor{gray}{\ensuremath{\cdot}} & \textcolor{gray}{\ensuremath{\cdot}} & \textcolor{gray}{\ensuremath{\cdot}} & \textcolor{gray}{\ensuremath{\cdot}} & \textcolor{gray}{\ensuremath{\cdot}} & \textcolor{gray}{\ensuremath{\cdot}} & \textcolor{gray}{\ensuremath{\cdot}}\\[3pt]\hline
\BenchmarkSourceLabel{bench:local:row:mps-vulkan-simd-block}{R78}\shortstack[l]{{\color{gray}R78}} & \begin{tabular}[t]{@{}l@{}}\rule{0pt}{12pt}MPS / Vulkan\\SIMD-DP / Block\end{tabular} & \textcolor{gray}{\ensuremath{\cdot}} & \textcolor{gray}{\ensuremath{\cdot}} & \textcolor{gray}{\ensuremath{\cdot}} & \textcolor{gray}{\ensuremath{\cdot}} & \hspace*{1pt}\begin{tabular}[t]{@{}c@{}}\rule{0pt}{12pt}{\fontsize{5.6}{6.4}\selectfont 1h\,57m\,20.2s}\\[-.2pt]{\fontsize{5.8}{6.4}\selectfont \textcolor{BrickRed}{$ \blacktriangledown\,246.4\times $}}\\[-.2pt]{\fontsize{5.6}{6.4}\selectfont 1h\,57m\,9.8s}\\[-.2pt]{\fontsize{5.8}{6.4}\selectfont \textcolor{BrickRed}{$ \blacktriangledown\,272.2\times $}}\end{tabular} & \textcolor{gray}{\ensuremath{\cdot}} & \textcolor{gray}{\ensuremath{\cdot}} & \textcolor{gray}{\ensuremath{\cdot}} & \textcolor{gray}{\ensuremath{\cdot}} & \textcolor{gray}{\ensuremath{\cdot}} & \textcolor{gray}{\ensuremath{\cdot}} & \textcolor{gray}{\ensuremath{\cdot}} & \textcolor{gray}{\ensuremath{\cdot}} & \textcolor{gray}{\ensuremath{\cdot}} & \textcolor{gray}{\ensuremath{\cdot}} & \textcolor{gray}{\ensuremath{\cdot}}\\[3pt]\hline
\BenchmarkSourceLabel{bench:local:row:mps-vulkan-scalar-ksw2}{R79}\shortstack[l]{{\color{gray}R79}} & \begin{tabular}[t]{@{}l@{}}\rule{0pt}{12pt}MPS / Vulkan\\Scalar-DP / KSW2\end{tabular} & \textcolor{gray}{\ensuremath{\cdot}} & \textcolor{gray}{\ensuremath{\cdot}} & \textcolor{gray}{\ensuremath{\cdot}} & \textcolor{gray}{\ensuremath{\cdot}} & \hspace*{1pt}\begin{tabular}[t]{@{}c@{}}\rule{0pt}{12pt}{\fontsize{5.6}{6.4}\selectfont 1h\,57m\,21.3s}\\[-.2pt]{\fontsize{5.8}{6.4}\selectfont \textcolor{BrickRed}{$ \blacktriangledown\,246.5\times $}}\\[-.2pt]{\fontsize{5.6}{6.4}\selectfont 1h\,57m\,7.9s}\\[-.2pt]{\fontsize{5.8}{6.4}\selectfont \textcolor{BrickRed}{$ \blacktriangledown\,272.1\times $}}\end{tabular} & \textcolor{gray}{\ensuremath{\cdot}} & \textcolor{gray}{\ensuremath{\cdot}} & \textcolor{gray}{\ensuremath{\cdot}} & \textcolor{gray}{\ensuremath{\cdot}} & \textcolor{gray}{\ensuremath{\cdot}} & \textcolor{gray}{\ensuremath{\cdot}} & \textcolor{gray}{\ensuremath{\cdot}} & \textcolor{gray}{\ensuremath{\cdot}} & \textcolor{gray}{\ensuremath{\cdot}} & \textcolor{gray}{\ensuremath{\cdot}} & \textcolor{gray}{\ensuremath{\cdot}}\\[3pt]\hline
\BenchmarkSourceLabel{bench:local:row:mps-vulkan-vulkan-block}{R80}\shortstack[l]{{\color{gray}R80}} & \begin{tabular}[t]{@{}l@{}}\rule{0pt}{12pt}MPS / Vulkan\\Vulkan-DP / Block\end{tabular} & \textcolor{gray}{\ensuremath{\cdot}} & \textcolor{gray}{\ensuremath{\cdot}} & \textcolor{gray}{\ensuremath{\cdot}} & \textcolor{gray}{\ensuremath{\cdot}} & \hspace*{1pt}\begin{tabular}[t]{@{}c@{}}\rule{0pt}{12pt}{\fontsize{5.6}{6.4}\selectfont 1h\,57m\,33.2s}\\[-.2pt]{\fontsize{5.8}{6.4}\selectfont \textcolor{BrickRed}{$ \blacktriangledown\,246.9\times $}}\\[-.2pt]{\fontsize{5.6}{6.4}\selectfont 1h\,57m\,22.2s}\\[-.2pt]{\fontsize{5.8}{6.4}\selectfont \textcolor{BrickRed}{$ \blacktriangledown\,272.6\times $}}\end{tabular} & \textcolor{gray}{\ensuremath{\cdot}} & \textcolor{gray}{\ensuremath{\cdot}} & \textcolor{gray}{\ensuremath{\cdot}} & \textcolor{gray}{\ensuremath{\cdot}} & \textcolor{gray}{\ensuremath{\cdot}} & \textcolor{gray}{\ensuremath{\cdot}} & \textcolor{gray}{\ensuremath{\cdot}} & \textcolor{gray}{\ensuremath{\cdot}} & \textcolor{gray}{\ensuremath{\cdot}} & \textcolor{gray}{\ensuremath{\cdot}} & \textcolor{gray}{\ensuremath{\cdot}}\\[3pt]\hline
\end{tabular}\par}
\smallskip{\footnotesize\noindent Each cell: A and B times, $\cdot$: no measurement. ---\textsuperscript{g}: GPU unavailable. n/a: no comparison.\par
\noindent Blue: one lowest A and one lowest B measurement for each column for each configuration.\par
\noindent Gray: minimap2. Yellow: my preferred MPF / Bounded / SIMD-DP / Block configuration.\par
\noindent MPF/MPS: minimizers per fragment/sequence. SIMD-ref: naive SIMD scan. Bounded: bounded SIMD algorithm.\par
}
\clearpage
\refstepcounter{table}\label{tab:bench:aws-4-k15-w10}
\noindent{\large\bfseries Table \thetable: (B): AWS: 4 Gbp, k=15, w=10\hfill 1/7}\par\smallskip
{\fontsize{7.1}{8.1}\selectfont\setlength{\tabcolsep}{1pt}\setlength{\fboxsep}{2pt}\renewcommand{\arraystretch}{1.12}
\noindent\begin{tabular}{@{}p{14mm}p{27mm}|p{13.55mm}|p{13.55mm}|p{13.55mm}|p{13.55mm}|p{13.55mm}|p{13.55mm}|p{13.55mm}|p{13.55mm}|p{13.55mm}|p{13.55mm}|p{13.55mm}|p{13.55mm}|p{13.55mm}|p{13.55mm}|p{13.55mm}|p{13.55mm}|@{}}\hline
Engine & Execution & \multicolumn{8}{c|}{\textbf{Index}} & \multicolumn{8}{c|}{\textbf{No index}}\\
 &  & \multicolumn{8}{c|}{\textbf{4 Gbp}} & \multicolumn{8}{c|}{\textbf{4 Gbp}}\\
 &  & \multicolumn{8}{c|}{\textbf{k=15, w=10}} & \multicolumn{8}{c|}{\textbf{k=15, w=10}}\\
 &  & \multicolumn{4}{c|}{\textbf{DNA4}} & \multicolumn{4}{c|}{\textbf{DNA5}} & \multicolumn{4}{c|}{\textbf{DNA4}} & \multicolumn{4}{c|}{\textbf{DNA5}}\\
 &  & \multicolumn{1}{c|}{\textbf{M7A}} & \multicolumn{1}{c|}{\textbf{M7I}} & \multicolumn{1}{c|}{\textbf{M4}} & \multicolumn{1}{c|}{\textbf{C7G}} & \multicolumn{1}{c|}{\textbf{M7A}} & \multicolumn{1}{c|}{\textbf{M7I}} & \multicolumn{1}{c|}{\textbf{M4}} & \multicolumn{1}{c|}{\textbf{C7G}} & \multicolumn{1}{c|}{\textbf{M7A}} & \multicolumn{1}{c|}{\textbf{M7I}} & \multicolumn{1}{c|}{\textbf{M4}} & \multicolumn{1}{c|}{\textbf{C7G}} & \multicolumn{1}{c|}{\textbf{M7A}} & \multicolumn{1}{c|}{\textbf{M7I}} & \multicolumn{1}{c|}{\textbf{M4}} & \multicolumn{1}{c|}{\textbf{C7G}}\\
 & & \BenchmarkSourceLabel{bench:aws-4-k15-w10:column:index-4-k15-w10-dna4-m7a}{C1}{\color{gray}C1} & \BenchmarkSourceLabel{bench:aws-4-k15-w10:column:index-4-k15-w10-dna4-m7i}{C2}{\color{gray}C2} & \BenchmarkSourceLabel{bench:aws-4-k15-w10:column:index-4-k15-w10-dna4-m4}{C3}{\color{gray}C3} & \BenchmarkSourceLabel{bench:aws-4-k15-w10:column:index-4-k15-w10-dna4-c7g}{C4}{\color{gray}C4} & \BenchmarkSourceLabel{bench:aws-4-k15-w10:column:index-4-k15-w10-dna5-m7a}{C5}{\color{gray}C5} & \BenchmarkSourceLabel{bench:aws-4-k15-w10:column:index-4-k15-w10-dna5-m7i}{C6}{\color{gray}C6} & \BenchmarkSourceLabel{bench:aws-4-k15-w10:column:index-4-k15-w10-dna5-m4}{C7}{\color{gray}C7} & \BenchmarkSourceLabel{bench:aws-4-k15-w10:column:index-4-k15-w10-dna5-c7g}{C8}{\color{gray}C8} & \BenchmarkSourceLabel{bench:aws-4-k15-w10:column:noindex-4-k15-w10-dna4-m7a}{C9}{\color{gray}C9} & \BenchmarkSourceLabel{bench:aws-4-k15-w10:column:noindex-4-k15-w10-dna4-m7i}{C10}{\color{gray}C10} & \BenchmarkSourceLabel{bench:aws-4-k15-w10:column:noindex-4-k15-w10-dna4-m4}{C11}{\color{gray}C11} & \BenchmarkSourceLabel{bench:aws-4-k15-w10:column:noindex-4-k15-w10-dna4-c7g}{C12}{\color{gray}C12} & \BenchmarkSourceLabel{bench:aws-4-k15-w10:column:noindex-4-k15-w10-dna5-m7a}{C13}{\color{gray}C13} & \BenchmarkSourceLabel{bench:aws-4-k15-w10:column:noindex-4-k15-w10-dna5-m7i}{C14}{\color{gray}C14} & \BenchmarkSourceLabel{bench:aws-4-k15-w10:column:noindex-4-k15-w10-dna5-m4}{C15}{\color{gray}C15} & \BenchmarkSourceLabel{bench:aws-4-k15-w10:column:noindex-4-k15-w10-dna5-c7g}{C16}{\color{gray}C16}\\
\hline
\rowcolor{BenchBaseline}[0.5pt][0.5pt]\BenchmarkSourceLabel{bench:aws-4-k15-w10:row:minimap2}{R1}\shortstack[l]{{\color{gray}R1}\\\textbf{minimap2}} & \textbf{Baseline} & \hspace*{1pt}\begin{tabular}[t]{@{}c@{}}\rule{0pt}{12pt}36.2s\\[-.2pt]{\fontsize{5.8}{6.4}\selectfont \textcolor{BrickRed}{$ \blacktriangledown\,6.9\times $}}\\[-.2pt]35s\\[-.2pt]{\fontsize{5.8}{6.4}\selectfont \textcolor{BrickRed}{$ \blacktriangledown\,8.0\times $}}\end{tabular} & \hspace*{1pt}\begin{tabular}[t]{@{}c@{}}\rule{0pt}{12pt}38.1s\\[-.2pt]{\fontsize{5.8}{6.4}\selectfont \textcolor{BrickRed}{$ \blacktriangledown\,7.8\times $}}\\[-.2pt]36s\\[-.2pt]{\fontsize{5.8}{6.4}\selectfont \textcolor{BrickRed}{$ \blacktriangledown\,9.3\times $}}\end{tabular} & \hspace*{1pt}\begin{tabular}[t]{@{}c@{}}\rule{0pt}{12pt}1m\,32.4s\\[-.2pt]{\fontsize{5.8}{6.4}\selectfont \textcolor{BrickRed}{$ \blacktriangledown\,10.7\times $}}\\[-.2pt]1m\,28.5s\\[-.2pt]{\fontsize{5.8}{6.4}\selectfont \textcolor{BrickRed}{$ \blacktriangledown\,13.5\times $}}\end{tabular} & \hspace*{1pt}\begin{tabular}[t]{@{}c@{}}\rule{0pt}{12pt}6m\,31.5s\\[-.2pt]{\fontsize{5.8}{6.4}\selectfont \textcolor{BrickRed}{$ \blacktriangledown\,11.5\times $}}\\[-.2pt]6m\,26.3s\\[-.2pt]{\fontsize{5.8}{6.4}\selectfont \textcolor{BrickRed}{$ \blacktriangledown\,12.7\times $}}\end{tabular} & \hspace*{1pt}\begin{tabular}[t]{@{}c@{}}\rule{0pt}{12pt}31.6s\\[-.2pt]{\fontsize{5.8}{6.4}\selectfont \textcolor{BrickRed}{$ \blacktriangledown\,2.1\times $}}\\[-.2pt]30.3s\\[-.2pt]{\fontsize{5.8}{6.4}\selectfont \textcolor{BrickRed}{$ \blacktriangledown\,3.0\times $}}\end{tabular} & \hspace*{1pt}\begin{tabular}[t]{@{}c@{}}\rule{0pt}{12pt}34.1s\\[-.2pt]{\fontsize{5.8}{6.4}\selectfont \textcolor{BrickRed}{$ \blacktriangledown\,2.9\times $}}\\[-.2pt]32s\\[-.2pt]{\fontsize{5.8}{6.4}\selectfont \textcolor{BrickRed}{$ \blacktriangledown\,3.7\times $}}\end{tabular} & \hspace*{1pt}\begin{tabular}[t]{@{}c@{}}\rule{0pt}{12pt}1m\,9.5s\\[-.2pt]{\fontsize{5.8}{6.4}\selectfont \textcolor{BrickRed}{$ \blacktriangledown\,3.0\times $}}\\[-.2pt]1m\,6.1s\\[-.2pt]{\fontsize{5.8}{6.4}\selectfont \textcolor{BrickRed}{$ \blacktriangledown\,4.3\times $}}\end{tabular} & \hspace*{1pt}\begin{tabular}[t]{@{}c@{}}\rule{0pt}{12pt}3m\,1s\\[-.2pt]{\fontsize{5.8}{6.4}\selectfont \textcolor{BrickRed}{$ \blacktriangledown\,7.9\times $}}\\[-.2pt]2m\,51.5s\\[-.2pt]{\fontsize{5.8}{6.4}\selectfont \textcolor{BrickRed}{$ \blacktriangledown\,8.6\times $}}\end{tabular} & \hspace*{1pt}\begin{tabular}[t]{@{}c@{}}\rule{0pt}{12pt}\colorbox{BenchBest}{\strut 2m\,55.5s}\\[-.2pt]{\fontsize{5.8}{6.4}\selectfont \textcolor{ForestGreen}{$ \blacktriangle\,20\% $}}\\[-.2pt]\colorbox{BenchBest}{\strut 2m\,54s}\\[-.2pt]{\fontsize{5.8}{6.4}\selectfont \textcolor{ForestGreen}{$ \blacktriangle\,17\% $}}\end{tabular} & \hspace*{1pt}\begin{tabular}[t]{@{}c@{}}\rule{0pt}{12pt}\colorbox{BenchBest}{\strut 2m\,50.4s}\\[-.2pt]{\fontsize{5.8}{6.4}\selectfont \textcolor{ForestGreen}{$ \blacktriangle\,23\% $}}\\[-.2pt]\colorbox{BenchBest}{\strut 2m\,48.2s}\\[-.2pt]{\fontsize{5.8}{6.4}\selectfont \textcolor{ForestGreen}{$ \blacktriangle\,20\% $}}\end{tabular} & \hspace*{1pt}\begin{tabular}[t]{@{}c@{}}\rule{0pt}{12pt}\colorbox{BenchBest}{\strut 5m\,35.7s}\\[-.2pt]{\fontsize{5.8}{6.4}\selectfont \textcolor{ForestGreen}{$ \blacktriangle\,18\% $}}\\[-.2pt]\colorbox{BenchBest}{\strut 5m\,31.9s}\\[-.2pt]{\fontsize{5.8}{6.4}\selectfont \textcolor{ForestGreen}{$ \blacktriangle\,14\% $}}\end{tabular} & \hspace*{1pt}\begin{tabular}[t]{@{}c@{}}\rule{0pt}{12pt}\colorbox{BenchBest}{\strut 44m\,20.5s}\\[-.2pt]{\fontsize{5.8}{6.4}\selectfont \textcolor{gray}{n/a}}\\[-.2pt]\colorbox{BenchBest}{\strut 44m\,15.6s}\\[-.2pt]{\fontsize{5.8}{6.4}\selectfont \textcolor{gray}{n/a}}\end{tabular} & \hspace*{1pt}\begin{tabular}[t]{@{}c@{}}\rule{0pt}{12pt}\colorbox{BenchBest}{\strut 1m\,47.1s}\\[-.2pt]{\fontsize{5.8}{6.4}\selectfont \textcolor{ForestGreen}{$ \blacktriangle\,5.5\times $}}\\[-.2pt]\colorbox{BenchBest}{\strut 1m\,45.8s}\\[-.2pt]{\fontsize{5.8}{6.4}\selectfont \textcolor{ForestGreen}{$ \blacktriangle\,2.1\times $}}\end{tabular} & \hspace*{1pt}\begin{tabular}[t]{@{}c@{}}\rule{0pt}{12pt}\colorbox{BenchBest}{\strut 1m\,47.2s}\\[-.2pt]{\fontsize{5.8}{6.4}\selectfont \textcolor{ForestGreen}{$ \blacktriangle\,4.6\times $}}\\[-.2pt]\colorbox{BenchBest}{\strut 1m\,45.1s}\\[-.2pt]{\fontsize{5.8}{6.4}\selectfont \textcolor{ForestGreen}{$ \blacktriangle\,2.1\times $}}\end{tabular} & \hspace*{1pt}\begin{tabular}[t]{@{}c@{}}\rule{0pt}{12pt}\colorbox{BenchBest}{\strut 3m\,40.3s}\\[-.2pt]{\fontsize{5.8}{6.4}\selectfont \textcolor{ForestGreen}{$ \blacktriangle\,4.2\times $}}\\[-.2pt]\colorbox{BenchBest}{\strut 3m\,36.6s}\\[-.2pt]{\fontsize{5.8}{6.4}\selectfont \textcolor{ForestGreen}{$ \blacktriangle\,1.7\times $}}\end{tabular} & \hspace*{1pt}\begin{tabular}[t]{@{}c@{}}\rule{0pt}{12pt}\colorbox{BenchBest}{\strut 29m\,43.4s}\\[-.2pt]{\fontsize{5.8}{6.4}\selectfont \textcolor{gray}{n/a}}\\[-.2pt]\colorbox{BenchBest}{\strut 29m\,38.5s}\\[-.2pt]{\fontsize{5.8}{6.4}\selectfont \textcolor{gray}{n/a}}\end{tabular}\\[3pt]\hline
\BenchmarkSourceLabel{bench:aws-4-k15-w10:row:mpf-simdref-scalar-block}{R2}\shortstack[l]{{\color{gray}R2}\\\textbf{BioseqDB}} & \begin{tabular}[t]{@{}l@{}}\rule{0pt}{12pt}MPF / SIMD-ref\\Scalar-DP / Block\end{tabular} & \hspace*{1pt}\begin{tabular}[t]{@{}c@{}}\rule{0pt}{12pt}\colorbox{BenchBest}{\strut 5.3s}\\[-.2pt]{\fontsize{5.8}{6.4}\selectfont \textcolor{ForestGreen}{$ \blacktriangle\,6.9\times $}}\\[-.2pt]4.5s\\[-.2pt]{\fontsize{5.8}{6.4}\selectfont \textcolor{ForestGreen}{$ \blacktriangle\,7.8\times $}}\end{tabular} & \hspace*{1pt}\begin{tabular}[t]{@{}c@{}}\rule{0pt}{12pt}\colorbox{BenchBest}{\strut 4.9s}\\[-.2pt]{\fontsize{5.8}{6.4}\selectfont \textcolor{ForestGreen}{$ \blacktriangle\,7.8\times $}}\\[-.2pt]4s\\[-.2pt]{\fontsize{5.8}{6.4}\selectfont \textcolor{ForestGreen}{$ \blacktriangle\,9.0\times $}}\end{tabular} & \hspace*{1pt}\begin{tabular}[t]{@{}c@{}}\rule{0pt}{12pt}11.6s\\[-.2pt]{\fontsize{5.8}{6.4}\selectfont \textcolor{ForestGreen}{$ \blacktriangle\,7.9\times $}}\\[-.2pt]9s\\[-.2pt]{\fontsize{5.8}{6.4}\selectfont \textcolor{ForestGreen}{$ \blacktriangle\,9.9\times $}}\end{tabular} & \hspace*{1pt}\begin{tabular}[t]{@{}c@{}}\rule{0pt}{12pt}35.1s\\[-.2pt]{\fontsize{5.8}{6.4}\selectfont \textcolor{ForestGreen}{$ \blacktriangle\,11.2\times $}}\\[-.2pt]34s\\[-.2pt]{\fontsize{5.8}{6.4}\selectfont \textcolor{ForestGreen}{$ \blacktriangle\,11.4\times $}}\end{tabular} & \hspace*{1pt}\begin{tabular}[t]{@{}c@{}}\rule{0pt}{12pt}14.8s\\[-.2pt]{\fontsize{5.8}{6.4}\selectfont \textcolor{ForestGreen}{$ \blacktriangle\,2.1\times $}}\\[-.2pt]10.3s\\[-.2pt]{\fontsize{5.8}{6.4}\selectfont \textcolor{ForestGreen}{$ \blacktriangle\,2.9\times $}}\end{tabular} & \hspace*{1pt}\begin{tabular}[t]{@{}c@{}}\rule{0pt}{12pt}15s\\[-.2pt]{\fontsize{5.8}{6.4}\selectfont \textcolor{ForestGreen}{$ \blacktriangle\,2.3\times $}}\\[-.2pt]10.9s\\[-.2pt]{\fontsize{5.8}{6.4}\selectfont \textcolor{ForestGreen}{$ \blacktriangle\,2.9\times $}}\end{tabular} & \hspace*{1pt}\begin{tabular}[t]{@{}c@{}}\rule{0pt}{12pt}26.1s\\[-.2pt]{\fontsize{5.8}{6.4}\selectfont \textcolor{ForestGreen}{$ \blacktriangle\,2.7\times $}}\\[-.2pt]17.6s\\[-.2pt]{\fontsize{5.8}{6.4}\selectfont \textcolor{ForestGreen}{$ \blacktriangle\,3.8\times $}}\end{tabular} & \hspace*{1pt}\begin{tabular}[t]{@{}c@{}}\rule{0pt}{12pt}\colorbox{BenchBest}{\strut 22.9s}\\[-.2pt]{\fontsize{5.8}{6.4}\selectfont \textcolor{ForestGreen}{$ \blacktriangle\,7.9\times $}}\\[-.2pt]20.4s\\[-.2pt]{\fontsize{5.8}{6.4}\selectfont \textcolor{ForestGreen}{$ \blacktriangle\,8.4\times $}}\end{tabular} & \hspace*{1pt}\begin{tabular}[t]{@{}c@{}}\rule{0pt}{12pt}4m\,50.9s\\[-.2pt]{\fontsize{5.8}{6.4}\selectfont \textcolor{BrickRed}{$ \blacktriangledown\,1.7\times $}}\\[-.2pt]4m\,39.8s\\[-.2pt]{\fontsize{5.8}{6.4}\selectfont \textcolor{BrickRed}{$ \blacktriangledown\,1.6\times $}}\end{tabular} & \hspace*{1pt}\begin{tabular}[t]{@{}c@{}}\rule{0pt}{12pt}4m\,54.9s\\[-.2pt]{\fontsize{5.8}{6.4}\selectfont \textcolor{BrickRed}{$ \blacktriangledown\,1.7\times $}}\\[-.2pt]4m\,41.8s\\[-.2pt]{\fontsize{5.8}{6.4}\selectfont \textcolor{BrickRed}{$ \blacktriangledown\,1.7\times $}}\end{tabular} & \textcolor{gray}{\ensuremath{\cdot}} & \textcolor{gray}{\ensuremath{\cdot}} & \textcolor{gray}{\ensuremath{\cdot}} & \textcolor{gray}{\ensuremath{\cdot}} & \textcolor{gray}{\ensuremath{\cdot}} & \textcolor{gray}{\ensuremath{\cdot}}\\[3pt]\hline
\BenchmarkSourceLabel{bench:aws-4-k15-w10:row:mpf-twostack-scalar-block}{R3}\shortstack[l]{{\color{gray}R3}} & \begin{tabular}[t]{@{}l@{}}\rule{0pt}{12pt}MPF / Two-stack\\Scalar-DP / Block\end{tabular} & \hspace*{1pt}\begin{tabular}[t]{@{}c@{}}\rule{0pt}{12pt}5.3s\\[-.2pt]{\fontsize{5.8}{6.4}\selectfont \textcolor{ForestGreen}{$ \blacktriangle\,6.8\times $}}\\[-.2pt]4.5s\\[-.2pt]{\fontsize{5.8}{6.4}\selectfont \textcolor{ForestGreen}{$ \blacktriangle\,7.7\times $}}\end{tabular} & \hspace*{1pt}\begin{tabular}[t]{@{}c@{}}\rule{0pt}{12pt}4.9s\\[-.2pt]{\fontsize{5.8}{6.4}\selectfont \textcolor{ForestGreen}{$ \blacktriangle\,7.8\times $}}\\[-.2pt]4s\\[-.2pt]{\fontsize{5.8}{6.4}\selectfont \textcolor{ForestGreen}{$ \blacktriangle\,8.9\times $}}\end{tabular} & \hspace*{1pt}\begin{tabular}[t]{@{}c@{}}\rule{0pt}{12pt}10.6s\\[-.2pt]{\fontsize{5.8}{6.4}\selectfont \textcolor{ForestGreen}{$ \blacktriangle\,8.7\times $}}\\[-.2pt]7.4s\\[-.2pt]{\fontsize{5.8}{6.4}\selectfont \textcolor{ForestGreen}{$ \blacktriangle\,11.9\times $}}\end{tabular} & \hspace*{1pt}\begin{tabular}[t]{@{}c@{}}\rule{0pt}{12pt}34.5s\\[-.2pt]{\fontsize{5.8}{6.4}\selectfont \textcolor{ForestGreen}{$ \blacktriangle\,11.3\times $}}\\[-.2pt]33.4s\\[-.2pt]{\fontsize{5.8}{6.4}\selectfont \textcolor{ForestGreen}{$ \blacktriangle\,11.6\times $}}\end{tabular} & \hspace*{1pt}\begin{tabular}[t]{@{}c@{}}\rule{0pt}{12pt}15.4s\\[-.2pt]{\fontsize{5.8}{6.4}\selectfont \textcolor{ForestGreen}{$ \blacktriangle\,2.0\times $}}\\[-.2pt]11.3s\\[-.2pt]{\fontsize{5.8}{6.4}\selectfont \textcolor{ForestGreen}{$ \blacktriangle\,2.7\times $}}\end{tabular} & \hspace*{1pt}\begin{tabular}[t]{@{}c@{}}\rule{0pt}{12pt}\colorbox{BenchBest}{\strut 11.8s}\\[-.2pt]{\fontsize{5.8}{6.4}\selectfont \textcolor{ForestGreen}{$ \blacktriangle\,2.9\times $}}\\[-.2pt]\colorbox{BenchBest}{\strut 8.8s}\\[-.2pt]{\fontsize{5.8}{6.4}\selectfont \textcolor{ForestGreen}{$ \blacktriangle\,3.7\times $}}\end{tabular} & \hspace*{1pt}\begin{tabular}[t]{@{}c@{}}\rule{0pt}{12pt}23.4s\\[-.2pt]{\fontsize{5.8}{6.4}\selectfont \textcolor{ForestGreen}{$ \blacktriangle\,3.0\times $}}\\[-.2pt]16.1s\\[-.2pt]{\fontsize{5.8}{6.4}\selectfont \textcolor{ForestGreen}{$ \blacktriangle\,4.1\times $}}\end{tabular} & \textcolor{gray}{\ensuremath{\cdot}} & \hspace*{1pt}\begin{tabular}[t]{@{}c@{}}\rule{0pt}{12pt}5m\,6.4s\\[-.2pt]{\fontsize{5.8}{6.4}\selectfont \textcolor{BrickRed}{$ \blacktriangledown\,1.7\times $}}\\[-.2pt]4m\,54.3s\\[-.2pt]{\fontsize{5.8}{6.4}\selectfont \textcolor{BrickRed}{$ \blacktriangledown\,1.7\times $}}\end{tabular} & \hspace*{1pt}\begin{tabular}[t]{@{}c@{}}\rule{0pt}{12pt}4m\,10.4s\\[-.2pt]{\fontsize{5.8}{6.4}\selectfont \textcolor{BrickRed}{$ \blacktriangledown\,1.5\times $}}\\[-.2pt]3m\,58.2s\\[-.2pt]{\fontsize{5.8}{6.4}\selectfont \textcolor{BrickRed}{$ \blacktriangledown\,1.4\times $}}\end{tabular} & \textcolor{gray}{\ensuremath{\cdot}} & \textcolor{gray}{\ensuremath{\cdot}} & \hspace*{1pt}\begin{tabular}[t]{@{}c@{}}\rule{0pt}{12pt}10m\,12.8s\\[-.2pt]{\fontsize{5.8}{6.4}\selectfont \textcolor{BrickRed}{$ \blacktriangledown\,5.7\times $}}\\[-.2pt]4m\,6.6s\\[-.2pt]{\fontsize{5.8}{6.4}\selectfont \textcolor{BrickRed}{$ \blacktriangledown\,2.3\times $}}\end{tabular} & \hspace*{1pt}\begin{tabular}[t]{@{}c@{}}\rule{0pt}{12pt}8m\,47.2s\\[-.2pt]{\fontsize{5.8}{6.4}\selectfont \textcolor{BrickRed}{$ \blacktriangledown\,4.9\times $}}\\[-.2pt]3m\,42.1s\\[-.2pt]{\fontsize{5.8}{6.4}\selectfont \textcolor{BrickRed}{$ \blacktriangledown\,2.1\times $}}\end{tabular} & \textcolor{gray}{\ensuremath{\cdot}} & \textcolor{gray}{\ensuremath{\cdot}}\\[3pt]\hline
\BenchmarkSourceLabel{bench:aws-4-k15-w10:row:mpf-scalar-scalar-block}{R4}\shortstack[l]{{\color{gray}R4}} & \begin{tabular}[t]{@{}l@{}}\rule{0pt}{12pt}MPF / Scalar\\Scalar-DP / Block\end{tabular} & \hspace*{1pt}\begin{tabular}[t]{@{}c@{}}\rule{0pt}{12pt}5.4s\\[-.2pt]{\fontsize{5.8}{6.4}\selectfont \textcolor{ForestGreen}{$ \blacktriangle\,6.7\times $}}\\[-.2pt]4.6s\\[-.2pt]{\fontsize{5.8}{6.4}\selectfont \textcolor{ForestGreen}{$ \blacktriangle\,7.6\times $}}\end{tabular} & \hspace*{1pt}\begin{tabular}[t]{@{}c@{}}\rule{0pt}{12pt}4.9s\\[-.2pt]{\fontsize{5.8}{6.4}\selectfont \textcolor{ForestGreen}{$ \blacktriangle\,7.7\times $}}\\[-.2pt]4s\\[-.2pt]{\fontsize{5.8}{6.4}\selectfont \textcolor{ForestGreen}{$ \blacktriangle\,8.9\times $}}\end{tabular} & \hspace*{1pt}\begin{tabular}[t]{@{}c@{}}\rule{0pt}{12pt}11.6s\\[-.2pt]{\fontsize{5.8}{6.4}\selectfont \textcolor{ForestGreen}{$ \blacktriangle\,8.0\times $}}\\[-.2pt]8.6s\\[-.2pt]{\fontsize{5.8}{6.4}\selectfont \textcolor{ForestGreen}{$ \blacktriangle\,10.3\times $}}\end{tabular} & \hspace*{1pt}\begin{tabular}[t]{@{}c@{}}\rule{0pt}{12pt}37.7s\\[-.2pt]{\fontsize{5.8}{6.4}\selectfont \textcolor{ForestGreen}{$ \blacktriangle\,10.4\times $}}\\[-.2pt]36.6s\\[-.2pt]{\fontsize{5.8}{6.4}\selectfont \textcolor{ForestGreen}{$ \blacktriangle\,10.6\times $}}\end{tabular} & \hspace*{1pt}\begin{tabular}[t]{@{}c@{}}\rule{0pt}{12pt}19.1s\\[-.2pt]{\fontsize{5.8}{6.4}\selectfont \textcolor{ForestGreen}{$ \blacktriangle\,1.7\times $}}\\[-.2pt]14.6s\\[-.2pt]{\fontsize{5.8}{6.4}\selectfont \textcolor{ForestGreen}{$ \blacktriangle\,2.1\times $}}\end{tabular} & \hspace*{1pt}\begin{tabular}[t]{@{}c@{}}\rule{0pt}{12pt}16.9s\\[-.2pt]{\fontsize{5.8}{6.4}\selectfont \textcolor{ForestGreen}{$ \blacktriangle\,2.0\times $}}\\[-.2pt]13.7s\\[-.2pt]{\fontsize{5.8}{6.4}\selectfont \textcolor{ForestGreen}{$ \blacktriangle\,2.3\times $}}\end{tabular} & \hspace*{1pt}\begin{tabular}[t]{@{}c@{}}\rule{0pt}{12pt}28s\\[-.2pt]{\fontsize{5.8}{6.4}\selectfont \textcolor{ForestGreen}{$ \blacktriangle\,2.5\times $}}\\[-.2pt]20.6s\\[-.2pt]{\fontsize{5.8}{6.4}\selectfont \textcolor{ForestGreen}{$ \blacktriangle\,3.2\times $}}\end{tabular} & \hspace*{1pt}\begin{tabular}[t]{@{}c@{}}\rule{0pt}{12pt}49.5s\\[-.2pt]{\fontsize{5.8}{6.4}\selectfont \textcolor{ForestGreen}{$ \blacktriangle\,3.7\times $}}\\[-.2pt]45.9s\\[-.2pt]{\fontsize{5.8}{6.4}\selectfont \textcolor{ForestGreen}{$ \blacktriangle\,3.7\times $}}\end{tabular} & \hspace*{1pt}\begin{tabular}[t]{@{}c@{}}\rule{0pt}{12pt}5m\,20.4s\\[-.2pt]{\fontsize{5.8}{6.4}\selectfont \textcolor{BrickRed}{$ \blacktriangledown\,1.8\times $}}\\[-.2pt]5m\,8s\\[-.2pt]{\fontsize{5.8}{6.4}\selectfont \textcolor{BrickRed}{$ \blacktriangledown\,1.8\times $}}\end{tabular} & \hspace*{1pt}\begin{tabular}[t]{@{}c@{}}\rule{0pt}{12pt}5m\,8.8s\\[-.2pt]{\fontsize{5.8}{6.4}\selectfont \textcolor{BrickRed}{$ \blacktriangledown\,1.8\times $}}\\[-.2pt]4m\,57s\\[-.2pt]{\fontsize{5.8}{6.4}\selectfont \textcolor{BrickRed}{$ \blacktriangledown\,1.8\times $}}\end{tabular} & \textcolor{gray}{\ensuremath{\cdot}} & \textcolor{gray}{\ensuremath{\cdot}} & \hspace*{1pt}\begin{tabular}[t]{@{}c@{}}\rule{0pt}{12pt}11m\,9.4s\\[-.2pt]{\fontsize{5.8}{6.4}\selectfont \textcolor{BrickRed}{$ \blacktriangledown\,6.3\times $}}\\[-.2pt]5m\,3.9s\\[-.2pt]{\fontsize{5.8}{6.4}\selectfont \textcolor{BrickRed}{$ \blacktriangledown\,2.9\times $}}\end{tabular} & \hspace*{1pt}\begin{tabular}[t]{@{}c@{}}\rule{0pt}{12pt}9m\,52.5s\\[-.2pt]{\fontsize{5.8}{6.4}\selectfont \textcolor{BrickRed}{$ \blacktriangledown\,5.5\times $}}\\[-.2pt]5m\,21.2s\\[-.2pt]{\fontsize{5.8}{6.4}\selectfont \textcolor{BrickRed}{$ \blacktriangledown\,3.1\times $}}\end{tabular} & \textcolor{gray}{\ensuremath{\cdot}} & \textcolor{gray}{\ensuremath{\cdot}}\\[3pt]\hline
\BenchmarkSourceLabel{bench:aws-4-k15-w10:row:mpf-simdref-rmq-block}{R5}\shortstack[l]{{\color{gray}R5}} & \begin{tabular}[t]{@{}l@{}}\rule{0pt}{12pt}MPF / SIMD-ref\\RMQ / Block\end{tabular} & \hspace*{1pt}\begin{tabular}[t]{@{}c@{}}\rule{0pt}{12pt}5.4s\\[-.2pt]{\fontsize{5.8}{6.4}\selectfont \textcolor{ForestGreen}{$ \blacktriangle\,6.7\times $}}\\[-.2pt]4.5s\\[-.2pt]{\fontsize{5.8}{6.4}\selectfont \textcolor{ForestGreen}{$ \blacktriangle\,7.8\times $}}\end{tabular} & \hspace*{1pt}\begin{tabular}[t]{@{}c@{}}\rule{0pt}{12pt}5s\\[-.2pt]{\fontsize{5.8}{6.4}\selectfont \textcolor{ForestGreen}{$ \blacktriangle\,7.7\times $}}\\[-.2pt]4.1s\\[-.2pt]{\fontsize{5.8}{6.4}\selectfont \textcolor{ForestGreen}{$ \blacktriangle\,8.8\times $}}\end{tabular} & \hspace*{1pt}\begin{tabular}[t]{@{}c@{}}\rule{0pt}{12pt}1m\,11.9s\\[-.2pt]{\fontsize{5.8}{6.4}\selectfont \textcolor{ForestGreen}{$ \blacktriangle\,22\% $}}\\[-.2pt]1m\,8.9s\\[-.2pt]{\fontsize{5.8}{6.4}\selectfont \textcolor{ForestGreen}{$ \blacktriangle\,22\% $}}\end{tabular} & \hspace*{1pt}\begin{tabular}[t]{@{}c@{}}\rule{0pt}{12pt}36s\\[-.2pt]{\fontsize{5.8}{6.4}\selectfont \textcolor{ForestGreen}{$ \blacktriangle\,10.9\times $}}\\[-.2pt]34.6s\\[-.2pt]{\fontsize{5.8}{6.4}\selectfont \textcolor{ForestGreen}{$ \blacktriangle\,11.2\times $}}\end{tabular} & \hspace*{1pt}\begin{tabular}[t]{@{}c@{}}\rule{0pt}{12pt}26.6s\\[-.2pt]{\fontsize{5.8}{6.4}\selectfont \textcolor{ForestGreen}{$ \blacktriangle\,16\% $}}\\[-.2pt]21.2s\\[-.2pt]{\fontsize{5.8}{6.4}\selectfont \textcolor{ForestGreen}{$ \blacktriangle\,30\% $}}\end{tabular} & \hspace*{1pt}\begin{tabular}[t]{@{}c@{}}\rule{0pt}{12pt}22.9s\\[-.2pt]{\fontsize{5.8}{6.4}\selectfont \textcolor{ForestGreen}{$ \blacktriangle\,1.5\times $}}\\[-.2pt]17.8s\\[-.2pt]{\fontsize{5.8}{6.4}\selectfont \textcolor{ForestGreen}{$ \blacktriangle\,1.8\times $}}\end{tabular} & \hspace*{1pt}\begin{tabular}[t]{@{}c@{}}\rule{0pt}{12pt}27.7s\\[-.2pt]{\fontsize{5.8}{6.4}\selectfont \textcolor{ForestGreen}{$ \blacktriangle\,2.5\times $}}\\[-.2pt]17.3s\\[-.2pt]{\fontsize{5.8}{6.4}\selectfont \textcolor{ForestGreen}{$ \blacktriangle\,3.8\times $}}\end{tabular} & \hspace*{1pt}\begin{tabular}[t]{@{}c@{}}\rule{0pt}{12pt}50.2s\\[-.2pt]{\fontsize{5.8}{6.4}\selectfont \textcolor{ForestGreen}{$ \blacktriangle\,3.6\times $}}\\[-.2pt]47.7s\\[-.2pt]{\fontsize{5.8}{6.4}\selectfont \textcolor{ForestGreen}{$ \blacktriangle\,3.6\times $}}\end{tabular} & \hspace*{1pt}\begin{tabular}[t]{@{}c@{}}\rule{0pt}{12pt}5m\,0.9s\\[-.2pt]{\fontsize{5.8}{6.4}\selectfont \textcolor{BrickRed}{$ \blacktriangledown\,1.7\times $}}\\[-.2pt]4m\,49.9s\\[-.2pt]{\fontsize{5.8}{6.4}\selectfont \textcolor{BrickRed}{$ \blacktriangledown\,1.7\times $}}\end{tabular} & \hspace*{1pt}\begin{tabular}[t]{@{}c@{}}\rule{0pt}{12pt}5m\,32s\\[-.2pt]{\fontsize{5.8}{6.4}\selectfont \textcolor{BrickRed}{$ \blacktriangledown\,1.9\times $}}\\[-.2pt]5m\,18.9s\\[-.2pt]{\fontsize{5.8}{6.4}\selectfont \textcolor{BrickRed}{$ \blacktriangledown\,1.9\times $}}\end{tabular} & \textcolor{gray}{\ensuremath{\cdot}} & \textcolor{gray}{\ensuremath{\cdot}} & \textcolor{gray}{\ensuremath{\cdot}} & \textcolor{gray}{\ensuremath{\cdot}} & \textcolor{gray}{\ensuremath{\cdot}} & \textcolor{gray}{\ensuremath{\cdot}}\\[3pt]\hline
\BenchmarkSourceLabel{bench:aws-4-k15-w10:row:mpf-twostack-rmq-block}{R6}\shortstack[l]{{\color{gray}R6}} & \begin{tabular}[t]{@{}l@{}}\rule{0pt}{12pt}MPF / Two-stack\\RMQ / Block\end{tabular} & \hspace*{1pt}\begin{tabular}[t]{@{}c@{}}\rule{0pt}{12pt}5.3s\\[-.2pt]{\fontsize{5.8}{6.4}\selectfont \textcolor{ForestGreen}{$ \blacktriangle\,6.8\times $}}\\[-.2pt]4.4s\\[-.2pt]{\fontsize{5.8}{6.4}\selectfont \textcolor{ForestGreen}{$ \blacktriangle\,7.9\times $}}\end{tabular} & \hspace*{1pt}\begin{tabular}[t]{@{}c@{}}\rule{0pt}{12pt}5s\\[-.2pt]{\fontsize{5.8}{6.4}\selectfont \textcolor{ForestGreen}{$ \blacktriangle\,7.7\times $}}\\[-.2pt]4.1s\\[-.2pt]{\fontsize{5.8}{6.4}\selectfont \textcolor{ForestGreen}{$ \blacktriangle\,8.9\times $}}\end{tabular} & \hspace*{1pt}\begin{tabular}[t]{@{}c@{}}\rule{0pt}{12pt}1m\,21s\\[-.2pt]{\fontsize{5.8}{6.4}\selectfont \textcolor{ForestGreen}{$ \blacktriangle\,12\% $}}\\[-.2pt]1m\,17.7s\\[-.2pt]{\fontsize{5.8}{6.4}\selectfont \textcolor{ForestGreen}{$ \blacktriangle\,12\% $}}\end{tabular} & \hspace*{1pt}\begin{tabular}[t]{@{}c@{}}\rule{0pt}{12pt}35s\\[-.2pt]{\fontsize{5.8}{6.4}\selectfont \textcolor{ForestGreen}{$ \blacktriangle\,11.2\times $}}\\[-.2pt]33.7s\\[-.2pt]{\fontsize{5.8}{6.4}\selectfont \textcolor{ForestGreen}{$ \blacktriangle\,11.5\times $}}\end{tabular} & \hspace*{1pt}\begin{tabular}[t]{@{}c@{}}\rule{0pt}{12pt}21.2s\\[-.2pt]{\fontsize{5.8}{6.4}\selectfont \textcolor{ForestGreen}{$ \blacktriangle\,1.5\times $}}\\[-.2pt]16.9s\\[-.2pt]{\fontsize{5.8}{6.4}\selectfont \textcolor{ForestGreen}{$ \blacktriangle\,1.8\times $}}\end{tabular} & \hspace*{1pt}\begin{tabular}[t]{@{}c@{}}\rule{0pt}{12pt}21.7s\\[-.2pt]{\fontsize{5.8}{6.4}\selectfont \textcolor{ForestGreen}{$ \blacktriangle\,1.6\times $}}\\[-.2pt]18.6s\\[-.2pt]{\fontsize{5.8}{6.4}\selectfont \textcolor{ForestGreen}{$ \blacktriangle\,1.7\times $}}\end{tabular} & \hspace*{1pt}\begin{tabular}[t]{@{}c@{}}\rule{0pt}{12pt}30.6s\\[-.2pt]{\fontsize{5.8}{6.4}\selectfont \textcolor{ForestGreen}{$ \blacktriangle\,2.3\times $}}\\[-.2pt]23s\\[-.2pt]{\fontsize{5.8}{6.4}\selectfont \textcolor{ForestGreen}{$ \blacktriangle\,2.9\times $}}\end{tabular} & \textcolor{gray}{\ensuremath{\cdot}} & \hspace*{1pt}\begin{tabular}[t]{@{}c@{}}\rule{0pt}{12pt}4m\,31.5s\\[-.2pt]{\fontsize{5.8}{6.4}\selectfont \textcolor{BrickRed}{$ \blacktriangledown\,1.5\times $}}\\[-.2pt]4m\,22.4s\\[-.2pt]{\fontsize{5.8}{6.4}\selectfont \textcolor{BrickRed}{$ \blacktriangledown\,1.5\times $}}\end{tabular} & \hspace*{1pt}\begin{tabular}[t]{@{}c@{}}\rule{0pt}{12pt}4m\,45s\\[-.2pt]{\fontsize{5.8}{6.4}\selectfont \textcolor{BrickRed}{$ \blacktriangledown\,1.7\times $}}\\[-.2pt]4m\,33.2s\\[-.2pt]{\fontsize{5.8}{6.4}\selectfont \textcolor{BrickRed}{$ \blacktriangledown\,1.6\times $}}\end{tabular} & \textcolor{gray}{\ensuremath{\cdot}} & \textcolor{gray}{\ensuremath{\cdot}} & \hspace*{1pt}\begin{tabular}[t]{@{}c@{}}\rule{0pt}{12pt}10m\,15.4s\\[-.2pt]{\fontsize{5.8}{6.4}\selectfont \textcolor{BrickRed}{$ \blacktriangledown\,5.7\times $}}\\[-.2pt]4m\,7.2s\\[-.2pt]{\fontsize{5.8}{6.4}\selectfont \textcolor{BrickRed}{$ \blacktriangledown\,2.3\times $}}\end{tabular} & \hspace*{1pt}\begin{tabular}[t]{@{}c@{}}\rule{0pt}{12pt}8m\,43.3s\\[-.2pt]{\fontsize{5.8}{6.4}\selectfont \textcolor{BrickRed}{$ \blacktriangledown\,4.9\times $}}\\[-.2pt]3m\,39s\\[-.2pt]{\fontsize{5.8}{6.4}\selectfont \textcolor{BrickRed}{$ \blacktriangledown\,2.1\times $}}\end{tabular} & \textcolor{gray}{\ensuremath{\cdot}} & \textcolor{gray}{\ensuremath{\cdot}}\\[3pt]\hline
\BenchmarkSourceLabel{bench:aws-4-k15-w10:row:mpf-scalar-rmq-block}{R7}\shortstack[l]{{\color{gray}R7}} & \begin{tabular}[t]{@{}l@{}}\rule{0pt}{12pt}MPF / Scalar\\RMQ / Block\end{tabular} & \hspace*{1pt}\begin{tabular}[t]{@{}c@{}}\rule{0pt}{12pt}5.5s\\[-.2pt]{\fontsize{5.8}{6.4}\selectfont \textcolor{ForestGreen}{$ \blacktriangle\,6.6\times $}}\\[-.2pt]4.6s\\[-.2pt]{\fontsize{5.8}{6.4}\selectfont \textcolor{ForestGreen}{$ \blacktriangle\,7.6\times $}}\end{tabular} & \hspace*{1pt}\begin{tabular}[t]{@{}c@{}}\rule{0pt}{12pt}5s\\[-.2pt]{\fontsize{5.8}{6.4}\selectfont \textcolor{ForestGreen}{$ \blacktriangle\,7.6\times $}}\\[-.2pt]4.1s\\[-.2pt]{\fontsize{5.8}{6.4}\selectfont \textcolor{ForestGreen}{$ \blacktriangle\,8.8\times $}}\end{tabular} & \hspace*{1pt}\begin{tabular}[t]{@{}c@{}}\rule{0pt}{12pt}22.8s\\[-.2pt]{\fontsize{5.8}{6.4}\selectfont \textcolor{ForestGreen}{$ \blacktriangle\,4.0\times $}}\\[-.2pt]19.7s\\[-.2pt]{\fontsize{5.8}{6.4}\selectfont \textcolor{ForestGreen}{$ \blacktriangle\,4.5\times $}}\end{tabular} & \hspace*{1pt}\begin{tabular}[t]{@{}c@{}}\rule{0pt}{12pt}35.2s\\[-.2pt]{\fontsize{5.8}{6.4}\selectfont \textcolor{ForestGreen}{$ \blacktriangle\,11.1\times $}}\\[-.2pt]34s\\[-.2pt]{\fontsize{5.8}{6.4}\selectfont \textcolor{ForestGreen}{$ \blacktriangle\,11.4\times $}}\end{tabular} & \hspace*{1pt}\begin{tabular}[t]{@{}c@{}}\rule{0pt}{12pt}35s\\[-.2pt]{\fontsize{5.8}{6.4}\selectfont \textcolor{BrickRed}{$ \blacktriangledown\,11\% $}}\\[-.2pt]30.5s\\[-.2pt]{\fontsize{5.8}{6.4}\selectfont \textcolor{gray}{$ \blacktriangleright\,1\% $}}\end{tabular} & \hspace*{1pt}\begin{tabular}[t]{@{}c@{}}\rule{0pt}{12pt}35.6s\\[-.2pt]{\fontsize{5.8}{6.4}\selectfont \textcolor{gray}{$ \blacktriangleright\,4\% $}}\\[-.2pt]32.4s\\[-.2pt]{\fontsize{5.8}{6.4}\selectfont \textcolor{gray}{$ \blacktriangleright\,1\% $}}\end{tabular} & \hspace*{1pt}\begin{tabular}[t]{@{}c@{}}\rule{0pt}{12pt}34s\\[-.2pt]{\fontsize{5.8}{6.4}\selectfont \textcolor{ForestGreen}{$ \blacktriangle\,2.0\times $}}\\[-.2pt]26.4s\\[-.2pt]{\fontsize{5.8}{6.4}\selectfont \textcolor{ForestGreen}{$ \blacktriangle\,2.5\times $}}\end{tabular} & \hspace*{1pt}\begin{tabular}[t]{@{}c@{}}\rule{0pt}{12pt}56.3s\\[-.2pt]{\fontsize{5.8}{6.4}\selectfont \textcolor{ForestGreen}{$ \blacktriangle\,3.2\times $}}\\[-.2pt]53s\\[-.2pt]{\fontsize{5.8}{6.4}\selectfont \textcolor{ForestGreen}{$ \blacktriangle\,3.2\times $}}\end{tabular} & \hspace*{1pt}\begin{tabular}[t]{@{}c@{}}\rule{0pt}{12pt}5m\,26.3s\\[-.2pt]{\fontsize{5.8}{6.4}\selectfont \textcolor{BrickRed}{$ \blacktriangledown\,1.9\times $}}\\[-.2pt]5m\,17.1s\\[-.2pt]{\fontsize{5.8}{6.4}\selectfont \textcolor{BrickRed}{$ \blacktriangledown\,1.8\times $}}\end{tabular} & \hspace*{1pt}\begin{tabular}[t]{@{}c@{}}\rule{0pt}{12pt}5m\,17.7s\\[-.2pt]{\fontsize{5.8}{6.4}\selectfont \textcolor{BrickRed}{$ \blacktriangledown\,1.9\times $}}\\[-.2pt]5m\,5.9s\\[-.2pt]{\fontsize{5.8}{6.4}\selectfont \textcolor{BrickRed}{$ \blacktriangledown\,1.8\times $}}\end{tabular} & \textcolor{gray}{\ensuremath{\cdot}} & \textcolor{gray}{\ensuremath{\cdot}} & \hspace*{1pt}\begin{tabular}[t]{@{}c@{}}\rule{0pt}{12pt}11m\,28.1s\\[-.2pt]{\fontsize{5.8}{6.4}\selectfont \textcolor{BrickRed}{$ \blacktriangledown\,6.4\times $}}\\[-.2pt]5m\,21.3s\\[-.2pt]{\fontsize{5.8}{6.4}\selectfont \textcolor{BrickRed}{$ \blacktriangledown\,3.0\times $}}\end{tabular} & \hspace*{1pt}\begin{tabular}[t]{@{}c@{}}\rule{0pt}{12pt}10m\,39.7s\\[-.2pt]{\fontsize{5.8}{6.4}\selectfont \textcolor{BrickRed}{$ \blacktriangledown\,6.0\times $}}\\[-.2pt]5m\,34.9s\\[-.2pt]{\fontsize{5.8}{6.4}\selectfont \textcolor{BrickRed}{$ \blacktriangledown\,3.2\times $}}\end{tabular} & \textcolor{gray}{\ensuremath{\cdot}} & \textcolor{gray}{\ensuremath{\cdot}}\\[3pt]\hline
\BenchmarkSourceLabel{bench:aws-4-k15-w10:row:mpf-twostack-simd-block}{R8}\shortstack[l]{{\color{gray}R8}} & \begin{tabular}[t]{@{}l@{}}\rule{0pt}{12pt}MPF / Two-stack\\SIMD-DP / Block\end{tabular} & \hspace*{1pt}\begin{tabular}[t]{@{}c@{}}\rule{0pt}{12pt}5.4s\\[-.2pt]{\fontsize{5.8}{6.4}\selectfont \textcolor{ForestGreen}{$ \blacktriangle\,6.7\times $}}\\[-.2pt]4.5s\\[-.2pt]{\fontsize{5.8}{6.4}\selectfont \textcolor{ForestGreen}{$ \blacktriangle\,7.8\times $}}\end{tabular} & \hspace*{1pt}\begin{tabular}[t]{@{}c@{}}\rule{0pt}{12pt}5s\\[-.2pt]{\fontsize{5.8}{6.4}\selectfont \textcolor{ForestGreen}{$ \blacktriangle\,7.6\times $}}\\[-.2pt]4s\\[-.2pt]{\fontsize{5.8}{6.4}\selectfont \textcolor{ForestGreen}{$ \blacktriangle\,8.9\times $}}\end{tabular} & \hspace*{1pt}\begin{tabular}[t]{@{}c@{}}\rule{0pt}{12pt}43.8s\\[-.2pt]{\fontsize{5.8}{6.4}\selectfont \textcolor{ForestGreen}{$ \blacktriangle\,2.1\times $}}\\[-.2pt]40.6s\\[-.2pt]{\fontsize{5.8}{6.4}\selectfont \textcolor{ForestGreen}{$ \blacktriangle\,2.2\times $}}\end{tabular} & \hspace*{1pt}\begin{tabular}[t]{@{}c@{}}\rule{0pt}{12pt}35.7s\\[-.2pt]{\fontsize{5.8}{6.4}\selectfont \textcolor{ForestGreen}{$ \blacktriangle\,11.0\times $}}\\[-.2pt]34.6s\\[-.2pt]{\fontsize{5.8}{6.4}\selectfont \textcolor{ForestGreen}{$ \blacktriangle\,11.2\times $}}\end{tabular} & \hspace*{1pt}\begin{tabular}[t]{@{}c@{}}\rule{0pt}{12pt}17.7s\\[-.2pt]{\fontsize{5.8}{6.4}\selectfont \textcolor{ForestGreen}{$ \blacktriangle\,1.8\times $}}\\[-.2pt]13.5s\\[-.2pt]{\fontsize{5.8}{6.4}\selectfont \textcolor{ForestGreen}{$ \blacktriangle\,2.2\times $}}\end{tabular} & \hspace*{1pt}\begin{tabular}[t]{@{}c@{}}\rule{0pt}{12pt}22s\\[-.2pt]{\fontsize{5.8}{6.4}\selectfont \textcolor{ForestGreen}{$ \blacktriangle\,1.6\times $}}\\[-.2pt]18.6s\\[-.2pt]{\fontsize{5.8}{6.4}\selectfont \textcolor{ForestGreen}{$ \blacktriangle\,1.7\times $}}\end{tabular} & \hspace*{1pt}\begin{tabular}[t]{@{}c@{}}\rule{0pt}{12pt}\colorbox{BenchBest}{\strut 23.1s}\\[-.2pt]{\fontsize{5.8}{6.4}\selectfont \textcolor{ForestGreen}{$ \blacktriangle\,3.0\times $}}\\[-.2pt]15.6s\\[-.2pt]{\fontsize{5.8}{6.4}\selectfont \textcolor{ForestGreen}{$ \blacktriangle\,4.2\times $}}\end{tabular} & \hspace*{1pt}\begin{tabular}[t]{@{}c@{}}\rule{0pt}{12pt}36.6s\\[-.2pt]{\fontsize{5.8}{6.4}\selectfont \textcolor{ForestGreen}{$ \blacktriangle\,5.0\times $}}\\[-.2pt]33.2s\\[-.2pt]{\fontsize{5.8}{6.4}\selectfont \textcolor{ForestGreen}{$ \blacktriangle\,5.2\times $}}\end{tabular} & \hspace*{1pt}\begin{tabular}[t]{@{}c@{}}\rule{0pt}{12pt}4m\,54.2s\\[-.2pt]{\fontsize{5.8}{6.4}\selectfont \textcolor{BrickRed}{$ \blacktriangledown\,1.7\times $}}\\[-.2pt]4m\,41.8s\\[-.2pt]{\fontsize{5.8}{6.4}\selectfont \textcolor{BrickRed}{$ \blacktriangledown\,1.6\times $}}\end{tabular} & \hspace*{1pt}\begin{tabular}[t]{@{}c@{}}\rule{0pt}{12pt}4m\,16.1s\\[-.2pt]{\fontsize{5.8}{6.4}\selectfont \textcolor{BrickRed}{$ \blacktriangledown\,1.5\times $}}\\[-.2pt]4m\,3.7s\\[-.2pt]{\fontsize{5.8}{6.4}\selectfont \textcolor{BrickRed}{$ \blacktriangledown\,1.4\times $}}\end{tabular} & \textcolor{gray}{\ensuremath{\cdot}} & \textcolor{gray}{\ensuremath{\cdot}} & \hspace*{1pt}\begin{tabular}[t]{@{}c@{}}\rule{0pt}{12pt}10m\,20.1s\\[-.2pt]{\fontsize{5.8}{6.4}\selectfont \textcolor{BrickRed}{$ \blacktriangledown\,5.8\times $}}\\[-.2pt]4m\,12.8s\\[-.2pt]{\fontsize{5.8}{6.4}\selectfont \textcolor{BrickRed}{$ \blacktriangledown\,2.4\times $}}\end{tabular} & \hspace*{1pt}\begin{tabular}[t]{@{}c@{}}\rule{0pt}{12pt}8m\,41.2s\\[-.2pt]{\fontsize{5.8}{6.4}\selectfont \textcolor{BrickRed}{$ \blacktriangledown\,4.9\times $}}\\[-.2pt]3m\,38.4s\\[-.2pt]{\fontsize{5.8}{6.4}\selectfont \textcolor{BrickRed}{$ \blacktriangledown\,2.1\times $}}\end{tabular} & \textcolor{gray}{\ensuremath{\cdot}} & \textcolor{gray}{\ensuremath{\cdot}}\\[3pt]\hline
\end{tabular}\par}
\smallskip{\footnotesize\noindent Each cell: A and B times, $\cdot$: no measurement. ---\textsuperscript{g}: GPU unavailable. n/a: no comparison.\par
\noindent Blue: one lowest A and one lowest B measurement for each column for each configuration.\par
\noindent Gray: minimap2. Yellow: my preferred MPF / Bounded / SIMD-DP / Block configuration.\par
\noindent MPF/MPS: minimizers per fragment/sequence. SIMD-ref: naive SIMD scan. Bounded: bounded SIMD algorithm.\par
}
\clearpage
\noindent{\large\bfseries Table \thetable: (B): AWS: 4 Gbp, k=15, w=10\hfill 2/7}\par\smallskip
{\fontsize{7.1}{8.1}\selectfont\setlength{\tabcolsep}{1pt}\setlength{\fboxsep}{2pt}\renewcommand{\arraystretch}{1.12}
\noindent\begin{tabular}{@{}p{14mm}p{27mm}|p{13.55mm}|p{13.55mm}|p{13.55mm}|p{13.55mm}|p{13.55mm}|p{13.55mm}|p{13.55mm}|p{13.55mm}|p{13.55mm}|p{13.55mm}|p{13.55mm}|p{13.55mm}|p{13.55mm}|p{13.55mm}|p{13.55mm}|p{13.55mm}|@{}}\hline
Engine & Execution & \multicolumn{8}{c|}{\textbf{Index}} & \multicolumn{8}{c|}{\textbf{No index}}\\
 &  & \multicolumn{8}{c|}{\textbf{4 Gbp}} & \multicolumn{8}{c|}{\textbf{4 Gbp}}\\
 &  & \multicolumn{8}{c|}{\textbf{k=15, w=10}} & \multicolumn{8}{c|}{\textbf{k=15, w=10}}\\
 &  & \multicolumn{4}{c|}{\textbf{DNA4}} & \multicolumn{4}{c|}{\textbf{DNA5}} & \multicolumn{4}{c|}{\textbf{DNA4}} & \multicolumn{4}{c|}{\textbf{DNA5}}\\
 &  & \multicolumn{1}{c|}{\textbf{M7A}} & \multicolumn{1}{c|}{\textbf{M7I}} & \multicolumn{1}{c|}{\textbf{M4}} & \multicolumn{1}{c|}{\textbf{C7G}} & \multicolumn{1}{c|}{\textbf{M7A}} & \multicolumn{1}{c|}{\textbf{M7I}} & \multicolumn{1}{c|}{\textbf{M4}} & \multicolumn{1}{c|}{\textbf{C7G}} & \multicolumn{1}{c|}{\textbf{M7A}} & \multicolumn{1}{c|}{\textbf{M7I}} & \multicolumn{1}{c|}{\textbf{M4}} & \multicolumn{1}{c|}{\textbf{C7G}} & \multicolumn{1}{c|}{\textbf{M7A}} & \multicolumn{1}{c|}{\textbf{M7I}} & \multicolumn{1}{c|}{\textbf{M4}} & \multicolumn{1}{c|}{\textbf{C7G}}\\
 & & {\color{gray}C1} & {\color{gray}C2} & {\color{gray}C3} & {\color{gray}C4} & {\color{gray}C5} & {\color{gray}C6} & {\color{gray}C7} & {\color{gray}C8} & {\color{gray}C9} & {\color{gray}C10} & {\color{gray}C11} & {\color{gray}C12} & {\color{gray}C13} & {\color{gray}C14} & {\color{gray}C15} & {\color{gray}C16}\\
\hline
\rowcolor{BenchBaseline}[0.5pt][0.5pt]\shortstack[l]{{\color{gray}R1}\\\textbf{minimap2}} & \textbf{Baseline} & \hspace*{1pt}\begin{tabular}[t]{@{}c@{}}\rule{0pt}{12pt}36.2s\\[-.2pt]{\fontsize{5.8}{6.4}\selectfont \textcolor{BrickRed}{$ \blacktriangledown\,6.9\times $}}\\[-.2pt]35s\\[-.2pt]{\fontsize{5.8}{6.4}\selectfont \textcolor{BrickRed}{$ \blacktriangledown\,8.0\times $}}\end{tabular} & \hspace*{1pt}\begin{tabular}[t]{@{}c@{}}\rule{0pt}{12pt}38.1s\\[-.2pt]{\fontsize{5.8}{6.4}\selectfont \textcolor{BrickRed}{$ \blacktriangledown\,7.8\times $}}\\[-.2pt]36s\\[-.2pt]{\fontsize{5.8}{6.4}\selectfont \textcolor{BrickRed}{$ \blacktriangledown\,9.3\times $}}\end{tabular} & \hspace*{1pt}\begin{tabular}[t]{@{}c@{}}\rule{0pt}{12pt}1m\,32.4s\\[-.2pt]{\fontsize{5.8}{6.4}\selectfont \textcolor{BrickRed}{$ \blacktriangledown\,10.7\times $}}\\[-.2pt]1m\,28.5s\\[-.2pt]{\fontsize{5.8}{6.4}\selectfont \textcolor{BrickRed}{$ \blacktriangledown\,13.5\times $}}\end{tabular} & \hspace*{1pt}\begin{tabular}[t]{@{}c@{}}\rule{0pt}{12pt}6m\,31.5s\\[-.2pt]{\fontsize{5.8}{6.4}\selectfont \textcolor{BrickRed}{$ \blacktriangledown\,11.5\times $}}\\[-.2pt]6m\,26.3s\\[-.2pt]{\fontsize{5.8}{6.4}\selectfont \textcolor{BrickRed}{$ \blacktriangledown\,12.7\times $}}\end{tabular} & \hspace*{1pt}\begin{tabular}[t]{@{}c@{}}\rule{0pt}{12pt}31.6s\\[-.2pt]{\fontsize{5.8}{6.4}\selectfont \textcolor{BrickRed}{$ \blacktriangledown\,2.1\times $}}\\[-.2pt]30.3s\\[-.2pt]{\fontsize{5.8}{6.4}\selectfont \textcolor{BrickRed}{$ \blacktriangledown\,3.0\times $}}\end{tabular} & \hspace*{1pt}\begin{tabular}[t]{@{}c@{}}\rule{0pt}{12pt}34.1s\\[-.2pt]{\fontsize{5.8}{6.4}\selectfont \textcolor{BrickRed}{$ \blacktriangledown\,2.9\times $}}\\[-.2pt]32s\\[-.2pt]{\fontsize{5.8}{6.4}\selectfont \textcolor{BrickRed}{$ \blacktriangledown\,3.7\times $}}\end{tabular} & \hspace*{1pt}\begin{tabular}[t]{@{}c@{}}\rule{0pt}{12pt}1m\,9.5s\\[-.2pt]{\fontsize{5.8}{6.4}\selectfont \textcolor{BrickRed}{$ \blacktriangledown\,3.0\times $}}\\[-.2pt]1m\,6.1s\\[-.2pt]{\fontsize{5.8}{6.4}\selectfont \textcolor{BrickRed}{$ \blacktriangledown\,4.3\times $}}\end{tabular} & \hspace*{1pt}\begin{tabular}[t]{@{}c@{}}\rule{0pt}{12pt}3m\,1s\\[-.2pt]{\fontsize{5.8}{6.4}\selectfont \textcolor{BrickRed}{$ \blacktriangledown\,7.9\times $}}\\[-.2pt]2m\,51.5s\\[-.2pt]{\fontsize{5.8}{6.4}\selectfont \textcolor{BrickRed}{$ \blacktriangledown\,8.6\times $}}\end{tabular} & \hspace*{1pt}\begin{tabular}[t]{@{}c@{}}\rule{0pt}{12pt}2m\,55.5s\\[-.2pt]{\fontsize{5.8}{6.4}\selectfont \textcolor{ForestGreen}{$ \blacktriangle\,20\% $}}\\[-.2pt]2m\,54s\\[-.2pt]{\fontsize{5.8}{6.4}\selectfont \textcolor{ForestGreen}{$ \blacktriangle\,17\% $}}\end{tabular} & \hspace*{1pt}\begin{tabular}[t]{@{}c@{}}\rule{0pt}{12pt}2m\,50.4s\\[-.2pt]{\fontsize{5.8}{6.4}\selectfont \textcolor{ForestGreen}{$ \blacktriangle\,23\% $}}\\[-.2pt]2m\,48.2s\\[-.2pt]{\fontsize{5.8}{6.4}\selectfont \textcolor{ForestGreen}{$ \blacktriangle\,20\% $}}\end{tabular} & \hspace*{1pt}\begin{tabular}[t]{@{}c@{}}\rule{0pt}{12pt}5m\,35.7s\\[-.2pt]{\fontsize{5.8}{6.4}\selectfont \textcolor{ForestGreen}{$ \blacktriangle\,18\% $}}\\[-.2pt]5m\,31.9s\\[-.2pt]{\fontsize{5.8}{6.4}\selectfont \textcolor{ForestGreen}{$ \blacktriangle\,14\% $}}\end{tabular} & \hspace*{1pt}\begin{tabular}[t]{@{}c@{}}\rule{0pt}{12pt}44m\,20.5s\\[-.2pt]{\fontsize{5.8}{6.4}\selectfont \textcolor{gray}{n/a}}\\[-.2pt]44m\,15.6s\\[-.2pt]{\fontsize{5.8}{6.4}\selectfont \textcolor{gray}{n/a}}\end{tabular} & \hspace*{1pt}\begin{tabular}[t]{@{}c@{}}\rule{0pt}{12pt}1m\,47.1s\\[-.2pt]{\fontsize{5.8}{6.4}\selectfont \textcolor{ForestGreen}{$ \blacktriangle\,5.5\times $}}\\[-.2pt]1m\,45.8s\\[-.2pt]{\fontsize{5.8}{6.4}\selectfont \textcolor{ForestGreen}{$ \blacktriangle\,2.1\times $}}\end{tabular} & \hspace*{1pt}\begin{tabular}[t]{@{}c@{}}\rule{0pt}{12pt}1m\,47.2s\\[-.2pt]{\fontsize{5.8}{6.4}\selectfont \textcolor{ForestGreen}{$ \blacktriangle\,4.6\times $}}\\[-.2pt]1m\,45.1s\\[-.2pt]{\fontsize{5.8}{6.4}\selectfont \textcolor{ForestGreen}{$ \blacktriangle\,2.1\times $}}\end{tabular} & \hspace*{1pt}\begin{tabular}[t]{@{}c@{}}\rule{0pt}{12pt}3m\,40.3s\\[-.2pt]{\fontsize{5.8}{6.4}\selectfont \textcolor{ForestGreen}{$ \blacktriangle\,4.2\times $}}\\[-.2pt]3m\,36.6s\\[-.2pt]{\fontsize{5.8}{6.4}\selectfont \textcolor{ForestGreen}{$ \blacktriangle\,1.7\times $}}\end{tabular} & \hspace*{1pt}\begin{tabular}[t]{@{}c@{}}\rule{0pt}{12pt}29m\,43.4s\\[-.2pt]{\fontsize{5.8}{6.4}\selectfont \textcolor{gray}{n/a}}\\[-.2pt]29m\,38.5s\\[-.2pt]{\fontsize{5.8}{6.4}\selectfont \textcolor{gray}{n/a}}\end{tabular}\\[3pt]\hline
\BenchmarkSourceLabel{bench:aws-4-k15-w10:row:mpf-scalar-simd-block}{R9}\shortstack[l]{{\color{gray}R9}\\\textbf{BioseqDB}} & \begin{tabular}[t]{@{}l@{}}\rule{0pt}{12pt}MPF / Scalar\\SIMD-DP / Block\end{tabular} & \hspace*{1pt}\begin{tabular}[t]{@{}c@{}}\rule{0pt}{12pt}5.5s\\[-.2pt]{\fontsize{5.8}{6.4}\selectfont \textcolor{ForestGreen}{$ \blacktriangle\,6.6\times $}}\\[-.2pt]4.6s\\[-.2pt]{\fontsize{5.8}{6.4}\selectfont \textcolor{ForestGreen}{$ \blacktriangle\,7.6\times $}}\end{tabular} & \hspace*{1pt}\begin{tabular}[t]{@{}c@{}}\rule{0pt}{12pt}5.1s\\[-.2pt]{\fontsize{5.8}{6.4}\selectfont \textcolor{ForestGreen}{$ \blacktriangle\,7.5\times $}}\\[-.2pt]4.1s\\[-.2pt]{\fontsize{5.8}{6.4}\selectfont \textcolor{ForestGreen}{$ \blacktriangle\,8.8\times $}}\end{tabular} & \hspace*{1pt}\begin{tabular}[t]{@{}c@{}}\rule{0pt}{12pt}32.2s\\[-.2pt]{\fontsize{5.8}{6.4}\selectfont \textcolor{ForestGreen}{$ \blacktriangle\,2.9\times $}}\\[-.2pt]29.2s\\[-.2pt]{\fontsize{5.8}{6.4}\selectfont \textcolor{ForestGreen}{$ \blacktriangle\,3.0\times $}}\end{tabular} & \hspace*{1pt}\begin{tabular}[t]{@{}c@{}}\rule{0pt}{12pt}35.7s\\[-.2pt]{\fontsize{5.8}{6.4}\selectfont \textcolor{ForestGreen}{$ \blacktriangle\,11.0\times $}}\\[-.2pt]34.6s\\[-.2pt]{\fontsize{5.8}{6.4}\selectfont \textcolor{ForestGreen}{$ \blacktriangle\,11.2\times $}}\end{tabular} & \hspace*{1pt}\begin{tabular}[t]{@{}c@{}}\rule{0pt}{12pt}23.2s\\[-.2pt]{\fontsize{5.8}{6.4}\selectfont \textcolor{ForestGreen}{$ \blacktriangle\,27\% $}}\\[-.2pt]18.8s\\[-.2pt]{\fontsize{5.8}{6.4}\selectfont \textcolor{ForestGreen}{$ \blacktriangle\,1.6\times $}}\end{tabular} & \hspace*{1pt}\begin{tabular}[t]{@{}c@{}}\rule{0pt}{12pt}32.3s\\[-.2pt]{\fontsize{5.8}{6.4}\selectfont \textcolor{ForestGreen}{$ \blacktriangle\,5\% $}}\\[-.2pt]28.5s\\[-.2pt]{\fontsize{5.8}{6.4}\selectfont \textcolor{ForestGreen}{$ \blacktriangle\,11\% $}}\end{tabular} & \hspace*{1pt}\begin{tabular}[t]{@{}c@{}}\rule{0pt}{12pt}1m\,28.5s\\[-.2pt]{\fontsize{5.8}{6.4}\selectfont \textcolor{BrickRed}{$ \blacktriangledown\,27\% $}}\\[-.2pt]1m\,20.8s\\[-.2pt]{\fontsize{5.8}{6.4}\selectfont \textcolor{BrickRed}{$ \blacktriangledown\,22\% $}}\end{tabular} & \hspace*{1pt}\begin{tabular}[t]{@{}c@{}}\rule{0pt}{12pt}57.3s\\[-.2pt]{\fontsize{5.8}{6.4}\selectfont \textcolor{ForestGreen}{$ \blacktriangle\,3.2\times $}}\\[-.2pt]53.9s\\[-.2pt]{\fontsize{5.8}{6.4}\selectfont \textcolor{ForestGreen}{$ \blacktriangle\,3.2\times $}}\end{tabular} & \hspace*{1pt}\begin{tabular}[t]{@{}c@{}}\rule{0pt}{12pt}5m\,24.9s\\[-.2pt]{\fontsize{5.8}{6.4}\selectfont \textcolor{BrickRed}{$ \blacktriangledown\,1.9\times $}}\\[-.2pt]5m\,16s\\[-.2pt]{\fontsize{5.8}{6.4}\selectfont \textcolor{BrickRed}{$ \blacktriangledown\,1.8\times $}}\end{tabular} & \hspace*{1pt}\begin{tabular}[t]{@{}c@{}}\rule{0pt}{12pt}5m\,22s\\[-.2pt]{\fontsize{5.8}{6.4}\selectfont \textcolor{BrickRed}{$ \blacktriangledown\,1.9\times $}}\\[-.2pt]5m\,9.2s\\[-.2pt]{\fontsize{5.8}{6.4}\selectfont \textcolor{BrickRed}{$ \blacktriangledown\,1.8\times $}}\end{tabular} & \textcolor{gray}{\ensuremath{\cdot}} & \textcolor{gray}{\ensuremath{\cdot}} & \hspace*{1pt}\begin{tabular}[t]{@{}c@{}}\rule{0pt}{12pt}11m\,28.2s\\[-.2pt]{\fontsize{5.8}{6.4}\selectfont \textcolor{BrickRed}{$ \blacktriangledown\,6.4\times $}}\\[-.2pt]5m\,22.9s\\[-.2pt]{\fontsize{5.8}{6.4}\selectfont \textcolor{BrickRed}{$ \blacktriangledown\,3.1\times $}}\end{tabular} & \hspace*{1pt}\begin{tabular}[t]{@{}c@{}}\rule{0pt}{12pt}9m\,48.5s\\[-.2pt]{\fontsize{5.8}{6.4}\selectfont \textcolor{BrickRed}{$ \blacktriangledown\,5.5\times $}}\\[-.2pt]5m\,17.7s\\[-.2pt]{\fontsize{5.8}{6.4}\selectfont \textcolor{BrickRed}{$ \blacktriangledown\,3.0\times $}}\end{tabular} & \textcolor{gray}{\ensuremath{\cdot}} & \textcolor{gray}{\ensuremath{\cdot}}\\[3pt]\hline
\BenchmarkSourceLabel{bench:aws-4-k15-w10:row:mpf-bounded-scalar-block}{R10}\shortstack[l]{{\color{gray}R10}} & \begin{tabular}[t]{@{}l@{}}\rule{0pt}{12pt}MPF / Bounded\\Scalar-DP / Block\end{tabular} & \hspace*{1pt}\begin{tabular}[t]{@{}c@{}}\rule{0pt}{12pt}5.3s\\[-.2pt]{\fontsize{5.8}{6.4}\selectfont \textcolor{ForestGreen}{$ \blacktriangle\,6.8\times $}}\\[-.2pt]4.6s\\[-.2pt]{\fontsize{5.8}{6.4}\selectfont \textcolor{ForestGreen}{$ \blacktriangle\,7.7\times $}}\end{tabular} & \hspace*{1pt}\begin{tabular}[t]{@{}c@{}}\rule{0pt}{12pt}5.5s\\[-.2pt]{\fontsize{5.8}{6.4}\selectfont \textcolor{ForestGreen}{$ \blacktriangle\,6.9\times $}}\\[-.2pt]4.6s\\[-.2pt]{\fontsize{5.8}{6.4}\selectfont \textcolor{ForestGreen}{$ \blacktriangle\,7.8\times $}}\end{tabular} & \hspace*{1pt}\begin{tabular}[t]{@{}c@{}}\rule{0pt}{12pt}10.8s\\[-.2pt]{\fontsize{5.8}{6.4}\selectfont \textcolor{ForestGreen}{$ \blacktriangle\,8.5\times $}}\\[-.2pt]7.8s\\[-.2pt]{\fontsize{5.8}{6.4}\selectfont \textcolor{ForestGreen}{$ \blacktriangle\,11.4\times $}}\end{tabular} & \hspace*{1pt}\begin{tabular}[t]{@{}c@{}}\rule{0pt}{12pt}\colorbox{BenchBest}{\strut 34.1s}\\[-.2pt]{\fontsize{5.8}{6.4}\selectfont \textcolor{ForestGreen}{$ \blacktriangle\,11.5\times $}}\\[-.2pt]33s\\[-.2pt]{\fontsize{5.8}{6.4}\selectfont \textcolor{ForestGreen}{$ \blacktriangle\,11.7\times $}}\end{tabular} & \hspace*{1pt}\begin{tabular}[t]{@{}c@{}}\rule{0pt}{12pt}18.7s\\[-.2pt]{\fontsize{5.8}{6.4}\selectfont \textcolor{ForestGreen}{$ \blacktriangle\,1.7\times $}}\\[-.2pt]14.3s\\[-.2pt]{\fontsize{5.8}{6.4}\selectfont \textcolor{ForestGreen}{$ \blacktriangle\,2.1\times $}}\end{tabular} & \hspace*{1pt}\begin{tabular}[t]{@{}c@{}}\rule{0pt}{12pt}18.3s\\[-.2pt]{\fontsize{5.8}{6.4}\selectfont \textcolor{ForestGreen}{$ \blacktriangle\,1.9\times $}}\\[-.2pt]14.8s\\[-.2pt]{\fontsize{5.8}{6.4}\selectfont \textcolor{ForestGreen}{$ \blacktriangle\,2.2\times $}}\end{tabular} & \hspace*{1pt}\begin{tabular}[t]{@{}c@{}}\rule{0pt}{12pt}23.6s\\[-.2pt]{\fontsize{5.8}{6.4}\selectfont \textcolor{ForestGreen}{$ \blacktriangle\,3.0\times $}}\\[-.2pt]16s\\[-.2pt]{\fontsize{5.8}{6.4}\selectfont \textcolor{ForestGreen}{$ \blacktriangle\,4.1\times $}}\end{tabular} & \hspace*{1pt}\begin{tabular}[t]{@{}c@{}}\rule{0pt}{12pt}54s\\[-.2pt]{\fontsize{5.8}{6.4}\selectfont \textcolor{ForestGreen}{$ \blacktriangle\,3.3\times $}}\\[-.2pt]50.6s\\[-.2pt]{\fontsize{5.8}{6.4}\selectfont \textcolor{ForestGreen}{$ \blacktriangle\,3.4\times $}}\end{tabular} & \hspace*{1pt}\begin{tabular}[t]{@{}c@{}}\rule{0pt}{12pt}3m\,40.6s\\[-.2pt]{\fontsize{5.8}{6.4}\selectfont \textcolor{BrickRed}{$ \blacktriangledown\,26\% $}}\\[-.2pt]3m\,31.8s\\[-.2pt]{\fontsize{5.8}{6.4}\selectfont \textcolor{BrickRed}{$ \blacktriangledown\,22\% $}}\end{tabular} & \hspace*{1pt}\begin{tabular}[t]{@{}c@{}}\rule{0pt}{12pt}4m\,25.9s\\[-.2pt]{\fontsize{5.8}{6.4}\selectfont \textcolor{BrickRed}{$ \blacktriangledown\,1.6\times $}}\\[-.2pt]4m\,13.9s\\[-.2pt]{\fontsize{5.8}{6.4}\selectfont \textcolor{BrickRed}{$ \blacktriangledown\,1.5\times $}}\end{tabular} & \textcolor{gray}{\ensuremath{\cdot}} & \textcolor{gray}{\ensuremath{\cdot}} & \hspace*{1pt}\begin{tabular}[t]{@{}c@{}}\rule{0pt}{12pt}10m\,4.6s\\[-.2pt]{\fontsize{5.8}{6.4}\selectfont \textcolor{BrickRed}{$ \blacktriangledown\,5.6\times $}}\\[-.2pt]3m\,56.7s\\[-.2pt]{\fontsize{5.8}{6.4}\selectfont \textcolor{BrickRed}{$ \blacktriangledown\,2.2\times $}}\end{tabular} & \hspace*{1pt}\begin{tabular}[t]{@{}c@{}}\rule{0pt}{12pt}14m\,49s\\[-.2pt]{\fontsize{5.8}{6.4}\selectfont \textcolor{BrickRed}{$ \blacktriangledown\,8.3\times $}}\\[-.2pt]8m\,18.7s\\[-.2pt]{\fontsize{5.8}{6.4}\selectfont \textcolor{BrickRed}{$ \blacktriangledown\,4.7\times $}}\end{tabular} & \hspace*{1pt}\begin{tabular}[t]{@{}c@{}}\rule{0pt}{12pt}17m\,59.1s\\[-.2pt]{\fontsize{5.8}{6.4}\selectfont \textcolor{BrickRed}{$ \blacktriangledown\,4.9\times $}}\\[-.2pt]9m\,9.4s\\[-.2pt]{\fontsize{5.8}{6.4}\selectfont \textcolor{BrickRed}{$ \blacktriangledown\,2.5\times $}}\end{tabular} & \textcolor{gray}{\ensuremath{\cdot}}\\[3pt]\hline
\BenchmarkSourceLabel{bench:aws-4-k15-w10:row:mpf-bounded-rmq-block}{R11}\shortstack[l]{{\color{gray}R11}} & \begin{tabular}[t]{@{}l@{}}\rule{0pt}{12pt}MPF / Bounded\\RMQ / Block\end{tabular} & \hspace*{1pt}\begin{tabular}[t]{@{}c@{}}\rule{0pt}{12pt}5.4s\\[-.2pt]{\fontsize{5.8}{6.4}\selectfont \textcolor{ForestGreen}{$ \blacktriangle\,6.8\times $}}\\[-.2pt]4.6s\\[-.2pt]{\fontsize{5.8}{6.4}\selectfont \textcolor{ForestGreen}{$ \blacktriangle\,7.6\times $}}\end{tabular} & \hspace*{1pt}\begin{tabular}[t]{@{}c@{}}\rule{0pt}{12pt}9.9s\\[-.2pt]{\fontsize{5.8}{6.4}\selectfont \textcolor{ForestGreen}{$ \blacktriangle\,3.9\times $}}\\[-.2pt]9.1s\\[-.2pt]{\fontsize{5.8}{6.4}\selectfont \textcolor{ForestGreen}{$ \blacktriangle\,4.0\times $}}\end{tabular} & \hspace*{1pt}\begin{tabular}[t]{@{}c@{}}\rule{0pt}{12pt}1m\,5.5s\\[-.2pt]{\fontsize{5.8}{6.4}\selectfont \textcolor{ForestGreen}{$ \blacktriangle\,29\% $}}\\[-.2pt]1m\,2.5s\\[-.2pt]{\fontsize{5.8}{6.4}\selectfont \textcolor{ForestGreen}{$ \blacktriangle\,29\% $}}\end{tabular} & \hspace*{1pt}\begin{tabular}[t]{@{}c@{}}\rule{0pt}{12pt}37.3s\\[-.2pt]{\fontsize{5.8}{6.4}\selectfont \textcolor{ForestGreen}{$ \blacktriangle\,10.5\times $}}\\[-.2pt]36.1s\\[-.2pt]{\fontsize{5.8}{6.4}\selectfont \textcolor{ForestGreen}{$ \blacktriangle\,10.7\times $}}\end{tabular} & \hspace*{1pt}\begin{tabular}[t]{@{}c@{}}\rule{0pt}{12pt}25.7s\\[-.2pt]{\fontsize{5.8}{6.4}\selectfont \textcolor{ForestGreen}{$ \blacktriangle\,19\% $}}\\[-.2pt]21.3s\\[-.2pt]{\fontsize{5.8}{6.4}\selectfont \textcolor{ForestGreen}{$ \blacktriangle\,30\% $}}\end{tabular} & \hspace*{1pt}\begin{tabular}[t]{@{}c@{}}\rule{0pt}{12pt}25.7s\\[-.2pt]{\fontsize{5.8}{6.4}\selectfont \textcolor{ForestGreen}{$ \blacktriangle\,25\% $}}\\[-.2pt]22.1s\\[-.2pt]{\fontsize{5.8}{6.4}\selectfont \textcolor{ForestGreen}{$ \blacktriangle\,1.5\times $}}\end{tabular} & \hspace*{1pt}\begin{tabular}[t]{@{}c@{}}\rule{0pt}{12pt}23.6s\\[-.2pt]{\fontsize{5.8}{6.4}\selectfont \textcolor{ForestGreen}{$ \blacktriangle\,2.9\times $}}\\[-.2pt]15.9s\\[-.2pt]{\fontsize{5.8}{6.4}\selectfont \textcolor{ForestGreen}{$ \blacktriangle\,4.2\times $}}\end{tabular} & \hspace*{1pt}\begin{tabular}[t]{@{}c@{}}\rule{0pt}{12pt}53s\\[-.2pt]{\fontsize{5.8}{6.4}\selectfont \textcolor{ForestGreen}{$ \blacktriangle\,3.4\times $}}\\[-.2pt]49.5s\\[-.2pt]{\fontsize{5.8}{6.4}\selectfont \textcolor{ForestGreen}{$ \blacktriangle\,3.5\times $}}\end{tabular} & \hspace*{1pt}\begin{tabular}[t]{@{}c@{}}\rule{0pt}{12pt}4m\,2.1s\\[-.2pt]{\fontsize{5.8}{6.4}\selectfont \textcolor{BrickRed}{$ \blacktriangledown\,1.4\times $}}\\[-.2pt]3m\,53.3s\\[-.2pt]{\fontsize{5.8}{6.4}\selectfont \textcolor{BrickRed}{$ \blacktriangledown\,1.3\times $}}\end{tabular} & \hspace*{1pt}\begin{tabular}[t]{@{}c@{}}\rule{0pt}{12pt}3m\,42.6s\\[-.2pt]{\fontsize{5.8}{6.4}\selectfont \textcolor{BrickRed}{$ \blacktriangledown\,1.3\times $}}\\[-.2pt]3m\,31.1s\\[-.2pt]{\fontsize{5.8}{6.4}\selectfont \textcolor{BrickRed}{$ \blacktriangledown\,25\% $}}\end{tabular} & \textcolor{gray}{\ensuremath{\cdot}} & \textcolor{gray}{\ensuremath{\cdot}} & \hspace*{1pt}\begin{tabular}[t]{@{}c@{}}\rule{0pt}{12pt}10m\,3.7s\\[-.2pt]{\fontsize{5.8}{6.4}\selectfont \textcolor{BrickRed}{$ \blacktriangledown\,5.6\times $}}\\[-.2pt]3m\,53.5s\\[-.2pt]{\fontsize{5.8}{6.4}\selectfont \textcolor{BrickRed}{$ \blacktriangledown\,2.2\times $}}\end{tabular} & \hspace*{1pt}\begin{tabular}[t]{@{}c@{}}\rule{0pt}{12pt}14m\,40.9s\\[-.2pt]{\fontsize{5.8}{6.4}\selectfont \textcolor{BrickRed}{$ \blacktriangledown\,8.2\times $}}\\[-.2pt]8m\,18.9s\\[-.2pt]{\fontsize{5.8}{6.4}\selectfont \textcolor{BrickRed}{$ \blacktriangledown\,4.7\times $}}\end{tabular} & \textcolor{gray}{\ensuremath{\cdot}} & \textcolor{gray}{\ensuremath{\cdot}}\\[3pt]\hline
\BenchmarkSourceLabel{bench:aws-4-k15-w10:row:mpf-simdref-simd-block}{R12}\shortstack[l]{{\color{gray}R12}} & \begin{tabular}[t]{@{}l@{}}\rule{0pt}{12pt}MPF / SIMD-ref\\SIMD-DP / Block\end{tabular} & \hspace*{1pt}\begin{tabular}[t]{@{}c@{}}\rule{0pt}{12pt}5.4s\\[-.2pt]{\fontsize{5.8}{6.4}\selectfont \textcolor{ForestGreen}{$ \blacktriangle\,6.7\times $}}\\[-.2pt]4.5s\\[-.2pt]{\fontsize{5.8}{6.4}\selectfont \textcolor{ForestGreen}{$ \blacktriangle\,7.8\times $}}\end{tabular} & \hspace*{1pt}\begin{tabular}[t]{@{}c@{}}\rule{0pt}{12pt}17.6s\\[-.2pt]{\fontsize{5.8}{6.4}\selectfont \textcolor{ForestGreen}{$ \blacktriangle\,2.2\times $}}\\[-.2pt]16.5s\\[-.2pt]{\fontsize{5.8}{6.4}\selectfont \textcolor{ForestGreen}{$ \blacktriangle\,2.2\times $}}\end{tabular} & \hspace*{1pt}\begin{tabular}[t]{@{}c@{}}\rule{0pt}{12pt}\colorbox{BenchBest}{\strut 8.7s}\\[-.2pt]{\fontsize{5.8}{6.4}\selectfont \textcolor{ForestGreen}{$ \blacktriangle\,10.7\times $}}\\[-.2pt]\colorbox{BenchBest}{\strut 6.5s}\\[-.2pt]{\fontsize{5.8}{6.4}\selectfont \textcolor{ForestGreen}{$ \blacktriangle\,13.5\times $}}\end{tabular} & \hspace*{1pt}\begin{tabular}[t]{@{}c@{}}\rule{0pt}{12pt}35.9s\\[-.2pt]{\fontsize{5.8}{6.4}\selectfont \textcolor{ForestGreen}{$ \blacktriangle\,10.9\times $}}\\[-.2pt]34.7s\\[-.2pt]{\fontsize{5.8}{6.4}\selectfont \textcolor{ForestGreen}{$ \blacktriangle\,11.1\times $}}\end{tabular} & \hspace*{1pt}\begin{tabular}[t]{@{}c@{}}\rule{0pt}{12pt}29.2s\\[-.2pt]{\fontsize{5.8}{6.4}\selectfont \textcolor{ForestGreen}{$ \blacktriangle\,8\% $}}\\[-.2pt]22.8s\\[-.2pt]{\fontsize{5.8}{6.4}\selectfont \textcolor{ForestGreen}{$ \blacktriangle\,25\% $}}\end{tabular} & \hspace*{1pt}\begin{tabular}[t]{@{}c@{}}\rule{0pt}{12pt}33s\\[-.2pt]{\fontsize{5.8}{6.4}\selectfont \textcolor{gray}{$ \blacktriangleright\,3\% $}}\\[-.2pt]27s\\[-.2pt]{\fontsize{5.8}{6.4}\selectfont \textcolor{ForestGreen}{$ \blacktriangle\,16\% $}}\end{tabular} & \hspace*{1pt}\begin{tabular}[t]{@{}c@{}}\rule{0pt}{12pt}54.5s\\[-.2pt]{\fontsize{5.8}{6.4}\selectfont \textcolor{ForestGreen}{$ \blacktriangle\,22\% $}}\\[-.2pt]43.2s\\[-.2pt]{\fontsize{5.8}{6.4}\selectfont \textcolor{ForestGreen}{$ \blacktriangle\,1.5\times $}}\end{tabular} & \hspace*{1pt}\begin{tabular}[t]{@{}c@{}}\rule{0pt}{12pt}37.7s\\[-.2pt]{\fontsize{5.8}{6.4}\selectfont \textcolor{ForestGreen}{$ \blacktriangle\,4.8\times $}}\\[-.2pt]35.3s\\[-.2pt]{\fontsize{5.8}{6.4}\selectfont \textcolor{ForestGreen}{$ \blacktriangle\,4.9\times $}}\end{tabular} & \hspace*{1pt}\begin{tabular}[t]{@{}c@{}}\rule{0pt}{12pt}4m\,54.9s\\[-.2pt]{\fontsize{5.8}{6.4}\selectfont \textcolor{BrickRed}{$ \blacktriangledown\,1.7\times $}}\\[-.2pt]4m\,43.5s\\[-.2pt]{\fontsize{5.8}{6.4}\selectfont \textcolor{BrickRed}{$ \blacktriangledown\,1.6\times $}}\end{tabular} & \hspace*{1pt}\begin{tabular}[t]{@{}c@{}}\rule{0pt}{12pt}4m\,42.7s\\[-.2pt]{\fontsize{5.8}{6.4}\selectfont \textcolor{BrickRed}{$ \blacktriangledown\,1.7\times $}}\\[-.2pt]4m\,29.7s\\[-.2pt]{\fontsize{5.8}{6.4}\selectfont \textcolor{BrickRed}{$ \blacktriangledown\,1.6\times $}}\end{tabular} & \textcolor{gray}{\ensuremath{\cdot}} & \textcolor{gray}{\ensuremath{\cdot}} & \textcolor{gray}{\ensuremath{\cdot}} & \textcolor{gray}{\ensuremath{\cdot}} & \textcolor{gray}{\ensuremath{\cdot}} & \textcolor{gray}{\ensuremath{\cdot}}\\[3pt]\hline
\BenchmarkSourceLabel{bench:aws-4-k15-w10:row:mpf-simdref-scalar-ksw2}{R13}\shortstack[l]{{\color{gray}R13}} & \begin{tabular}[t]{@{}l@{}}\rule{0pt}{12pt}MPF / SIMD-ref\\Scalar-DP / KSW2\end{tabular} & \hspace*{1pt}\begin{tabular}[t]{@{}c@{}}\rule{0pt}{12pt}13.8s\\[-.2pt]{\fontsize{5.8}{6.4}\selectfont \textcolor{ForestGreen}{$ \blacktriangle\,2.6\times $}}\\[-.2pt]11.3s\\[-.2pt]{\fontsize{5.8}{6.4}\selectfont \textcolor{ForestGreen}{$ \blacktriangle\,3.1\times $}}\end{tabular} & \hspace*{1pt}\begin{tabular}[t]{@{}c@{}}\rule{0pt}{12pt}6.3s\\[-.2pt]{\fontsize{5.8}{6.4}\selectfont \textcolor{ForestGreen}{$ \blacktriangle\,6.0\times $}}\\[-.2pt]3.9s\\[-.2pt]{\fontsize{5.8}{6.4}\selectfont \textcolor{ForestGreen}{$ \blacktriangle\,9.3\times $}}\end{tabular} & \hspace*{1pt}\begin{tabular}[t]{@{}c@{}}\rule{0pt}{12pt}17.9s\\[-.2pt]{\fontsize{5.8}{6.4}\selectfont \textcolor{ForestGreen}{$ \blacktriangle\,5.2\times $}}\\[-.2pt]10.1s\\[-.2pt]{\fontsize{5.8}{6.4}\selectfont \textcolor{ForestGreen}{$ \blacktriangle\,8.8\times $}}\end{tabular} & \hspace*{1pt}\begin{tabular}[t]{@{}c@{}}\rule{0pt}{12pt}39s\\[-.2pt]{\fontsize{5.8}{6.4}\selectfont \textcolor{ForestGreen}{$ \blacktriangle\,10.0\times $}}\\[-.2pt]35.1s\\[-.2pt]{\fontsize{5.8}{6.4}\selectfont \textcolor{ForestGreen}{$ \blacktriangle\,11.0\times $}}\end{tabular} & \hspace*{1pt}\begin{tabular}[t]{@{}c@{}}\rule{0pt}{12pt}\colorbox{BenchBest}{\strut 14.8s}\\[-.2pt]{\fontsize{5.8}{6.4}\selectfont \textcolor{ForestGreen}{$ \blacktriangle\,2.1\times $}}\\[-.2pt]\colorbox{BenchBest}{\strut 10.2s}\\[-.2pt]{\fontsize{5.8}{6.4}\selectfont \textcolor{ForestGreen}{$ \blacktriangle\,3.0\times $}}\end{tabular} & \hspace*{1pt}\begin{tabular}[t]{@{}c@{}}\rule{0pt}{12pt}16s\\[-.2pt]{\fontsize{5.8}{6.4}\selectfont \textcolor{ForestGreen}{$ \blacktriangle\,2.1\times $}}\\[-.2pt]11.5s\\[-.2pt]{\fontsize{5.8}{6.4}\selectfont \textcolor{ForestGreen}{$ \blacktriangle\,2.8\times $}}\end{tabular} & \hspace*{1pt}\begin{tabular}[t]{@{}c@{}}\rule{0pt}{12pt}26.6s\\[-.2pt]{\fontsize{5.8}{6.4}\selectfont \textcolor{ForestGreen}{$ \blacktriangle\,2.6\times $}}\\[-.2pt]17.2s\\[-.2pt]{\fontsize{5.8}{6.4}\selectfont \textcolor{ForestGreen}{$ \blacktriangle\,3.8\times $}}\end{tabular} & \hspace*{1pt}\begin{tabular}[t]{@{}c@{}}\rule{0pt}{12pt}23.6s\\[-.2pt]{\fontsize{5.8}{6.4}\selectfont \textcolor{ForestGreen}{$ \blacktriangle\,7.7\times $}}\\[-.2pt]20.4s\\[-.2pt]{\fontsize{5.8}{6.4}\selectfont \textcolor{ForestGreen}{$ \blacktriangle\,8.4\times $}}\end{tabular} & \hspace*{1pt}\begin{tabular}[t]{@{}c@{}}\rule{0pt}{12pt}4m\,52.7s\\[-.2pt]{\fontsize{5.8}{6.4}\selectfont \textcolor{BrickRed}{$ \blacktriangledown\,1.7\times $}}\\[-.2pt]4m\,39.7s\\[-.2pt]{\fontsize{5.8}{6.4}\selectfont \textcolor{BrickRed}{$ \blacktriangledown\,1.6\times $}}\end{tabular} & \hspace*{1pt}\begin{tabular}[t]{@{}c@{}}\rule{0pt}{12pt}5m\,1.6s\\[-.2pt]{\fontsize{5.8}{6.4}\selectfont \textcolor{BrickRed}{$ \blacktriangledown\,1.8\times $}}\\[-.2pt]4m\,46.5s\\[-.2pt]{\fontsize{5.8}{6.4}\selectfont \textcolor{BrickRed}{$ \blacktriangledown\,1.7\times $}}\end{tabular} & \textcolor{gray}{\ensuremath{\cdot}} & \textcolor{gray}{\ensuremath{\cdot}} & \textcolor{gray}{\ensuremath{\cdot}} & \textcolor{gray}{\ensuremath{\cdot}} & \textcolor{gray}{\ensuremath{\cdot}} & \textcolor{gray}{\ensuremath{\cdot}}\\[3pt]\hline
\BenchmarkSourceLabel{bench:aws-4-k15-w10:row:mpf-twostack-rmq-ksw2}{R14}\shortstack[l]{{\color{gray}R14}} & \begin{tabular}[t]{@{}l@{}}\rule{0pt}{12pt}MPF / Two-stack\\RMQ / KSW2\end{tabular} & \hspace*{1pt}\begin{tabular}[t]{@{}c@{}}\rule{0pt}{12pt}6.9s\\[-.2pt]{\fontsize{5.8}{6.4}\selectfont \textcolor{ForestGreen}{$ \blacktriangle\,5.2\times $}}\\[-.2pt]4.4s\\[-.2pt]{\fontsize{5.8}{6.4}\selectfont \textcolor{ForestGreen}{$ \blacktriangle\,8.0\times $}}\end{tabular} & \hspace*{1pt}\begin{tabular}[t]{@{}c@{}}\rule{0pt}{12pt}6.4s\\[-.2pt]{\fontsize{5.8}{6.4}\selectfont \textcolor{ForestGreen}{$ \blacktriangle\,5.9\times $}}\\[-.2pt]\colorbox{BenchBest}{\strut 3.9s}\\[-.2pt]{\fontsize{5.8}{6.4}\selectfont \textcolor{ForestGreen}{$ \blacktriangle\,9.3\times $}}\end{tabular} & \hspace*{1pt}\begin{tabular}[t]{@{}c@{}}\rule{0pt}{12pt}32.1s\\[-.2pt]{\fontsize{5.8}{6.4}\selectfont \textcolor{ForestGreen}{$ \blacktriangle\,2.9\times $}}\\[-.2pt]23.8s\\[-.2pt]{\fontsize{5.8}{6.4}\selectfont \textcolor{ForestGreen}{$ \blacktriangle\,3.7\times $}}\end{tabular} & \hspace*{1pt}\begin{tabular}[t]{@{}c@{}}\rule{0pt}{12pt}35.5s\\[-.2pt]{\fontsize{5.8}{6.4}\selectfont \textcolor{ForestGreen}{$ \blacktriangle\,11.0\times $}}\\[-.2pt]31.4s\\[-.2pt]{\fontsize{5.8}{6.4}\selectfont \textcolor{ForestGreen}{$ \blacktriangle\,12.3\times $}}\end{tabular} & \hspace*{1pt}\begin{tabular}[t]{@{}c@{}}\rule{0pt}{12pt}17.2s\\[-.2pt]{\fontsize{5.8}{6.4}\selectfont \textcolor{ForestGreen}{$ \blacktriangle\,1.8\times $}}\\[-.2pt]13s\\[-.2pt]{\fontsize{5.8}{6.4}\selectfont \textcolor{ForestGreen}{$ \blacktriangle\,2.3\times $}}\end{tabular} & \hspace*{1pt}\begin{tabular}[t]{@{}c@{}}\rule{0pt}{12pt}17.3s\\[-.2pt]{\fontsize{5.8}{6.4}\selectfont \textcolor{ForestGreen}{$ \blacktriangle\,2.0\times $}}\\[-.2pt]13.9s\\[-.2pt]{\fontsize{5.8}{6.4}\selectfont \textcolor{ForestGreen}{$ \blacktriangle\,2.3\times $}}\end{tabular} & \hspace*{1pt}\begin{tabular}[t]{@{}c@{}}\rule{0pt}{12pt}24.4s\\[-.2pt]{\fontsize{5.8}{6.4}\selectfont \textcolor{ForestGreen}{$ \blacktriangle\,2.8\times $}}\\[-.2pt]16.4s\\[-.2pt]{\fontsize{5.8}{6.4}\selectfont \textcolor{ForestGreen}{$ \blacktriangle\,4.0\times $}}\end{tabular} & \textcolor{gray}{\ensuremath{\cdot}} & \hspace*{1pt}\begin{tabular}[t]{@{}c@{}}\rule{0pt}{12pt}4m\,33.3s\\[-.2pt]{\fontsize{5.8}{6.4}\selectfont \textcolor{BrickRed}{$ \blacktriangledown\,1.6\times $}}\\[-.2pt]4m\,19.4s\\[-.2pt]{\fontsize{5.8}{6.4}\selectfont \textcolor{BrickRed}{$ \blacktriangledown\,1.5\times $}}\end{tabular} & \hspace*{1pt}\begin{tabular}[t]{@{}c@{}}\rule{0pt}{12pt}4m\,13.5s\\[-.2pt]{\fontsize{5.8}{6.4}\selectfont \textcolor{BrickRed}{$ \blacktriangledown\,1.5\times $}}\\[-.2pt]4m\,0s\\[-.2pt]{\fontsize{5.8}{6.4}\selectfont \textcolor{BrickRed}{$ \blacktriangledown\,1.4\times $}}\end{tabular} & \textcolor{gray}{\ensuremath{\cdot}} & \textcolor{gray}{\ensuremath{\cdot}} & \hspace*{1pt}\begin{tabular}[t]{@{}c@{}}\rule{0pt}{12pt}10m\,11.1s\\[-.2pt]{\fontsize{5.8}{6.4}\selectfont \textcolor{BrickRed}{$ \blacktriangledown\,5.7\times $}}\\[-.2pt]4m\,1.3s\\[-.2pt]{\fontsize{5.8}{6.4}\selectfont \textcolor{BrickRed}{$ \blacktriangledown\,2.3\times $}}\end{tabular} & \hspace*{1pt}\begin{tabular}[t]{@{}c@{}}\rule{0pt}{12pt}8m\,47.2s\\[-.2pt]{\fontsize{5.8}{6.4}\selectfont \textcolor{BrickRed}{$ \blacktriangledown\,4.9\times $}}\\[-.2pt]3m\,42.1s\\[-.2pt]{\fontsize{5.8}{6.4}\selectfont \textcolor{BrickRed}{$ \blacktriangledown\,2.1\times $}}\end{tabular} & \textcolor{gray}{\ensuremath{\cdot}} & \textcolor{gray}{\ensuremath{\cdot}}\\[3pt]\hline
\BenchmarkSourceLabel{bench:aws-4-k15-w10:row:mpf-simdref-rmq-ksw2}{R15}\shortstack[l]{{\color{gray}R15}} & \begin{tabular}[t]{@{}l@{}}\rule{0pt}{12pt}MPF / SIMD-ref\\RMQ / KSW2\end{tabular} & \hspace*{1pt}\begin{tabular}[t]{@{}c@{}}\rule{0pt}{12pt}6.9s\\[-.2pt]{\fontsize{5.8}{6.4}\selectfont \textcolor{ForestGreen}{$ \blacktriangle\,5.3\times $}}\\[-.2pt]4.4s\\[-.2pt]{\fontsize{5.8}{6.4}\selectfont \textcolor{ForestGreen}{$ \blacktriangle\,8.0\times $}}\end{tabular} & \hspace*{1pt}\begin{tabular}[t]{@{}c@{}}\rule{0pt}{12pt}6.5s\\[-.2pt]{\fontsize{5.8}{6.4}\selectfont \textcolor{ForestGreen}{$ \blacktriangle\,5.9\times $}}\\[-.2pt]3.9s\\[-.2pt]{\fontsize{5.8}{6.4}\selectfont \textcolor{ForestGreen}{$ \blacktriangle\,9.3\times $}}\end{tabular} & \hspace*{1pt}\begin{tabular}[t]{@{}c@{}}\rule{0pt}{12pt}19.3s\\[-.2pt]{\fontsize{5.8}{6.4}\selectfont \textcolor{ForestGreen}{$ \blacktriangle\,4.8\times $}}\\[-.2pt]11.4s\\[-.2pt]{\fontsize{5.8}{6.4}\selectfont \textcolor{ForestGreen}{$ \blacktriangle\,7.8\times $}}\end{tabular} & \hspace*{1pt}\begin{tabular}[t]{@{}c@{}}\rule{0pt}{12pt}39.9s\\[-.2pt]{\fontsize{5.8}{6.4}\selectfont \textcolor{ForestGreen}{$ \blacktriangle\,9.8\times $}}\\[-.2pt]35.7s\\[-.2pt]{\fontsize{5.8}{6.4}\selectfont \textcolor{ForestGreen}{$ \blacktriangle\,10.8\times $}}\end{tabular} & \hspace*{1pt}\begin{tabular}[t]{@{}c@{}}\rule{0pt}{12pt}15.9s\\[-.2pt]{\fontsize{5.8}{6.4}\selectfont \textcolor{ForestGreen}{$ \blacktriangle\,2.0\times $}}\\[-.2pt]10.4s\\[-.2pt]{\fontsize{5.8}{6.4}\selectfont \textcolor{ForestGreen}{$ \blacktriangle\,2.9\times $}}\end{tabular} & \hspace*{1pt}\begin{tabular}[t]{@{}c@{}}\rule{0pt}{12pt}17.9s\\[-.2pt]{\fontsize{5.8}{6.4}\selectfont \textcolor{ForestGreen}{$ \blacktriangle\,1.9\times $}}\\[-.2pt]12.3s\\[-.2pt]{\fontsize{5.8}{6.4}\selectfont \textcolor{ForestGreen}{$ \blacktriangle\,2.6\times $}}\end{tabular} & \hspace*{1pt}\begin{tabular}[t]{@{}c@{}}\rule{0pt}{12pt}29.4s\\[-.2pt]{\fontsize{5.8}{6.4}\selectfont \textcolor{ForestGreen}{$ \blacktriangle\,2.4\times $}}\\[-.2pt]17.7s\\[-.2pt]{\fontsize{5.8}{6.4}\selectfont \textcolor{ForestGreen}{$ \blacktriangle\,3.7\times $}}\end{tabular} & \hspace*{1pt}\begin{tabular}[t]{@{}c@{}}\rule{0pt}{12pt}23.4s\\[-.2pt]{\fontsize{5.8}{6.4}\selectfont \textcolor{ForestGreen}{$ \blacktriangle\,7.7\times $}}\\[-.2pt]\colorbox{BenchBest}{\strut 20s}\\[-.2pt]{\fontsize{5.8}{6.4}\selectfont \textcolor{ForestGreen}{$ \blacktriangle\,8.6\times $}}\end{tabular} & \hspace*{1pt}\begin{tabular}[t]{@{}c@{}}\rule{0pt}{12pt}4m\,58.1s\\[-.2pt]{\fontsize{5.8}{6.4}\selectfont \textcolor{BrickRed}{$ \blacktriangledown\,1.7\times $}}\\[-.2pt]4m\,45.3s\\[-.2pt]{\fontsize{5.8}{6.4}\selectfont \textcolor{BrickRed}{$ \blacktriangledown\,1.6\times $}}\end{tabular} & \hspace*{1pt}\begin{tabular}[t]{@{}c@{}}\rule{0pt}{12pt}5m\,2.1s\\[-.2pt]{\fontsize{5.8}{6.4}\selectfont \textcolor{BrickRed}{$ \blacktriangledown\,1.8\times $}}\\[-.2pt]4m\,47.2s\\[-.2pt]{\fontsize{5.8}{6.4}\selectfont \textcolor{BrickRed}{$ \blacktriangledown\,1.7\times $}}\end{tabular} & \textcolor{gray}{\ensuremath{\cdot}} & \textcolor{gray}{\ensuremath{\cdot}} & \textcolor{gray}{\ensuremath{\cdot}} & \textcolor{gray}{\ensuremath{\cdot}} & \textcolor{gray}{\ensuremath{\cdot}} & \textcolor{gray}{\ensuremath{\cdot}}\\[3pt]\hline
\end{tabular}\par}
\smallskip{\footnotesize\noindent Each cell: A and B times, $\cdot$: no measurement. ---\textsuperscript{g}: GPU unavailable. n/a: no comparison.\par
\noindent Blue: one lowest A and one lowest B measurement for each column for each configuration.\par
\noindent Gray: minimap2. Yellow: my preferred MPF / Bounded / SIMD-DP / Block configuration.\par
\noindent MPF/MPS: minimizers per fragment/sequence. SIMD-ref: naive SIMD scan. Bounded: bounded SIMD algorithm.\par
}
\clearpage
\noindent{\large\bfseries Table \thetable: (B): AWS: 4 Gbp, k=15, w=10\hfill 3/7}\par\smallskip
{\fontsize{7.1}{8.1}\selectfont\setlength{\tabcolsep}{1pt}\setlength{\fboxsep}{2pt}\renewcommand{\arraystretch}{1.12}
\noindent\begin{tabular}{@{}p{14mm}p{27mm}|p{13.55mm}|p{13.55mm}|p{13.55mm}|p{13.55mm}|p{13.55mm}|p{13.55mm}|p{13.55mm}|p{13.55mm}|p{13.55mm}|p{13.55mm}|p{13.55mm}|p{13.55mm}|p{13.55mm}|p{13.55mm}|p{13.55mm}|p{13.55mm}|@{}}\hline
Engine & Execution & \multicolumn{8}{c|}{\textbf{Index}} & \multicolumn{8}{c|}{\textbf{No index}}\\
 &  & \multicolumn{8}{c|}{\textbf{4 Gbp}} & \multicolumn{8}{c|}{\textbf{4 Gbp}}\\
 &  & \multicolumn{8}{c|}{\textbf{k=15, w=10}} & \multicolumn{8}{c|}{\textbf{k=15, w=10}}\\
 &  & \multicolumn{4}{c|}{\textbf{DNA4}} & \multicolumn{4}{c|}{\textbf{DNA5}} & \multicolumn{4}{c|}{\textbf{DNA4}} & \multicolumn{4}{c|}{\textbf{DNA5}}\\
 &  & \multicolumn{1}{c|}{\textbf{M7A}} & \multicolumn{1}{c|}{\textbf{M7I}} & \multicolumn{1}{c|}{\textbf{M4}} & \multicolumn{1}{c|}{\textbf{C7G}} & \multicolumn{1}{c|}{\textbf{M7A}} & \multicolumn{1}{c|}{\textbf{M7I}} & \multicolumn{1}{c|}{\textbf{M4}} & \multicolumn{1}{c|}{\textbf{C7G}} & \multicolumn{1}{c|}{\textbf{M7A}} & \multicolumn{1}{c|}{\textbf{M7I}} & \multicolumn{1}{c|}{\textbf{M4}} & \multicolumn{1}{c|}{\textbf{C7G}} & \multicolumn{1}{c|}{\textbf{M7A}} & \multicolumn{1}{c|}{\textbf{M7I}} & \multicolumn{1}{c|}{\textbf{M4}} & \multicolumn{1}{c|}{\textbf{C7G}}\\
 & & {\color{gray}C1} & {\color{gray}C2} & {\color{gray}C3} & {\color{gray}C4} & {\color{gray}C5} & {\color{gray}C6} & {\color{gray}C7} & {\color{gray}C8} & {\color{gray}C9} & {\color{gray}C10} & {\color{gray}C11} & {\color{gray}C12} & {\color{gray}C13} & {\color{gray}C14} & {\color{gray}C15} & {\color{gray}C16}\\
\hline
\rowcolor{BenchBaseline}[0.5pt][0.5pt]\shortstack[l]{{\color{gray}R1}\\\textbf{minimap2}} & \textbf{Baseline} & \hspace*{1pt}\begin{tabular}[t]{@{}c@{}}\rule{0pt}{12pt}36.2s\\[-.2pt]{\fontsize{5.8}{6.4}\selectfont \textcolor{BrickRed}{$ \blacktriangledown\,6.9\times $}}\\[-.2pt]35s\\[-.2pt]{\fontsize{5.8}{6.4}\selectfont \textcolor{BrickRed}{$ \blacktriangledown\,8.0\times $}}\end{tabular} & \hspace*{1pt}\begin{tabular}[t]{@{}c@{}}\rule{0pt}{12pt}38.1s\\[-.2pt]{\fontsize{5.8}{6.4}\selectfont \textcolor{BrickRed}{$ \blacktriangledown\,7.8\times $}}\\[-.2pt]36s\\[-.2pt]{\fontsize{5.8}{6.4}\selectfont \textcolor{BrickRed}{$ \blacktriangledown\,9.3\times $}}\end{tabular} & \hspace*{1pt}\begin{tabular}[t]{@{}c@{}}\rule{0pt}{12pt}1m\,32.4s\\[-.2pt]{\fontsize{5.8}{6.4}\selectfont \textcolor{BrickRed}{$ \blacktriangledown\,10.7\times $}}\\[-.2pt]1m\,28.5s\\[-.2pt]{\fontsize{5.8}{6.4}\selectfont \textcolor{BrickRed}{$ \blacktriangledown\,13.5\times $}}\end{tabular} & \hspace*{1pt}\begin{tabular}[t]{@{}c@{}}\rule{0pt}{12pt}6m\,31.5s\\[-.2pt]{\fontsize{5.8}{6.4}\selectfont \textcolor{BrickRed}{$ \blacktriangledown\,11.5\times $}}\\[-.2pt]6m\,26.3s\\[-.2pt]{\fontsize{5.8}{6.4}\selectfont \textcolor{BrickRed}{$ \blacktriangledown\,12.7\times $}}\end{tabular} & \hspace*{1pt}\begin{tabular}[t]{@{}c@{}}\rule{0pt}{12pt}31.6s\\[-.2pt]{\fontsize{5.8}{6.4}\selectfont \textcolor{BrickRed}{$ \blacktriangledown\,2.1\times $}}\\[-.2pt]30.3s\\[-.2pt]{\fontsize{5.8}{6.4}\selectfont \textcolor{BrickRed}{$ \blacktriangledown\,3.0\times $}}\end{tabular} & \hspace*{1pt}\begin{tabular}[t]{@{}c@{}}\rule{0pt}{12pt}34.1s\\[-.2pt]{\fontsize{5.8}{6.4}\selectfont \textcolor{BrickRed}{$ \blacktriangledown\,2.9\times $}}\\[-.2pt]32s\\[-.2pt]{\fontsize{5.8}{6.4}\selectfont \textcolor{BrickRed}{$ \blacktriangledown\,3.7\times $}}\end{tabular} & \hspace*{1pt}\begin{tabular}[t]{@{}c@{}}\rule{0pt}{12pt}1m\,9.5s\\[-.2pt]{\fontsize{5.8}{6.4}\selectfont \textcolor{BrickRed}{$ \blacktriangledown\,3.0\times $}}\\[-.2pt]1m\,6.1s\\[-.2pt]{\fontsize{5.8}{6.4}\selectfont \textcolor{BrickRed}{$ \blacktriangledown\,4.3\times $}}\end{tabular} & \hspace*{1pt}\begin{tabular}[t]{@{}c@{}}\rule{0pt}{12pt}3m\,1s\\[-.2pt]{\fontsize{5.8}{6.4}\selectfont \textcolor{BrickRed}{$ \blacktriangledown\,7.9\times $}}\\[-.2pt]2m\,51.5s\\[-.2pt]{\fontsize{5.8}{6.4}\selectfont \textcolor{BrickRed}{$ \blacktriangledown\,8.6\times $}}\end{tabular} & \hspace*{1pt}\begin{tabular}[t]{@{}c@{}}\rule{0pt}{12pt}2m\,55.5s\\[-.2pt]{\fontsize{5.8}{6.4}\selectfont \textcolor{ForestGreen}{$ \blacktriangle\,20\% $}}\\[-.2pt]2m\,54s\\[-.2pt]{\fontsize{5.8}{6.4}\selectfont \textcolor{ForestGreen}{$ \blacktriangle\,17\% $}}\end{tabular} & \hspace*{1pt}\begin{tabular}[t]{@{}c@{}}\rule{0pt}{12pt}2m\,50.4s\\[-.2pt]{\fontsize{5.8}{6.4}\selectfont \textcolor{ForestGreen}{$ \blacktriangle\,23\% $}}\\[-.2pt]2m\,48.2s\\[-.2pt]{\fontsize{5.8}{6.4}\selectfont \textcolor{ForestGreen}{$ \blacktriangle\,20\% $}}\end{tabular} & \hspace*{1pt}\begin{tabular}[t]{@{}c@{}}\rule{0pt}{12pt}5m\,35.7s\\[-.2pt]{\fontsize{5.8}{6.4}\selectfont \textcolor{ForestGreen}{$ \blacktriangle\,18\% $}}\\[-.2pt]5m\,31.9s\\[-.2pt]{\fontsize{5.8}{6.4}\selectfont \textcolor{ForestGreen}{$ \blacktriangle\,14\% $}}\end{tabular} & \hspace*{1pt}\begin{tabular}[t]{@{}c@{}}\rule{0pt}{12pt}44m\,20.5s\\[-.2pt]{\fontsize{5.8}{6.4}\selectfont \textcolor{gray}{n/a}}\\[-.2pt]44m\,15.6s\\[-.2pt]{\fontsize{5.8}{6.4}\selectfont \textcolor{gray}{n/a}}\end{tabular} & \hspace*{1pt}\begin{tabular}[t]{@{}c@{}}\rule{0pt}{12pt}1m\,47.1s\\[-.2pt]{\fontsize{5.8}{6.4}\selectfont \textcolor{ForestGreen}{$ \blacktriangle\,5.5\times $}}\\[-.2pt]1m\,45.8s\\[-.2pt]{\fontsize{5.8}{6.4}\selectfont \textcolor{ForestGreen}{$ \blacktriangle\,2.1\times $}}\end{tabular} & \hspace*{1pt}\begin{tabular}[t]{@{}c@{}}\rule{0pt}{12pt}1m\,47.2s\\[-.2pt]{\fontsize{5.8}{6.4}\selectfont \textcolor{ForestGreen}{$ \blacktriangle\,4.6\times $}}\\[-.2pt]1m\,45.1s\\[-.2pt]{\fontsize{5.8}{6.4}\selectfont \textcolor{ForestGreen}{$ \blacktriangle\,2.1\times $}}\end{tabular} & \hspace*{1pt}\begin{tabular}[t]{@{}c@{}}\rule{0pt}{12pt}3m\,40.3s\\[-.2pt]{\fontsize{5.8}{6.4}\selectfont \textcolor{ForestGreen}{$ \blacktriangle\,4.2\times $}}\\[-.2pt]3m\,36.6s\\[-.2pt]{\fontsize{5.8}{6.4}\selectfont \textcolor{ForestGreen}{$ \blacktriangle\,1.7\times $}}\end{tabular} & \hspace*{1pt}\begin{tabular}[t]{@{}c@{}}\rule{0pt}{12pt}29m\,43.4s\\[-.2pt]{\fontsize{5.8}{6.4}\selectfont \textcolor{gray}{n/a}}\\[-.2pt]29m\,38.5s\\[-.2pt]{\fontsize{5.8}{6.4}\selectfont \textcolor{gray}{n/a}}\end{tabular}\\[3pt]\hline
\BenchmarkSourceLabel{bench:aws-4-k15-w10:row:mpf-scalar-simd-ksw2}{R16}\shortstack[l]{{\color{gray}R16}\\\textbf{BioseqDB}} & \begin{tabular}[t]{@{}l@{}}\rule{0pt}{12pt}MPF / Scalar\\SIMD-DP / KSW2\end{tabular} & \hspace*{1pt}\begin{tabular}[t]{@{}c@{}}\rule{0pt}{12pt}8.2s\\[-.2pt]{\fontsize{5.8}{6.4}\selectfont \textcolor{ForestGreen}{$ \blacktriangle\,4.4\times $}}\\[-.2pt]5.6s\\[-.2pt]{\fontsize{5.8}{6.4}\selectfont \textcolor{ForestGreen}{$ \blacktriangle\,6.2\times $}}\end{tabular} & \hspace*{1pt}\begin{tabular}[t]{@{}c@{}}\rule{0pt}{12pt}6.5s\\[-.2pt]{\fontsize{5.8}{6.4}\selectfont \textcolor{ForestGreen}{$ \blacktriangle\,5.9\times $}}\\[-.2pt]3.9s\\[-.2pt]{\fontsize{5.8}{6.4}\selectfont \textcolor{ForestGreen}{$ \blacktriangle\,9.1\times $}}\end{tabular} & \hspace*{1pt}\begin{tabular}[t]{@{}c@{}}\rule{0pt}{12pt}2m\,15.5s\\[-.2pt]{\fontsize{5.8}{6.4}\selectfont \textcolor{BrickRed}{$ \blacktriangledown\,1.5\times $}}\\[-.2pt]2m\,7.7s\\[-.2pt]{\fontsize{5.8}{6.4}\selectfont \textcolor{BrickRed}{$ \blacktriangledown\,1.4\times $}}\end{tabular} & \hspace*{1pt}\begin{tabular}[t]{@{}c@{}}\rule{0pt}{12pt}37.8s\\[-.2pt]{\fontsize{5.8}{6.4}\selectfont \textcolor{ForestGreen}{$ \blacktriangle\,10.3\times $}}\\[-.2pt]33.9s\\[-.2pt]{\fontsize{5.8}{6.4}\selectfont \textcolor{ForestGreen}{$ \blacktriangle\,11.4\times $}}\end{tabular} & \hspace*{1pt}\begin{tabular}[t]{@{}c@{}}\rule{0pt}{12pt}42.2s\\[-.2pt]{\fontsize{5.8}{6.4}\selectfont \textcolor{BrickRed}{$ \blacktriangledown\,1.3\times $}}\\[-.2pt]37.6s\\[-.2pt]{\fontsize{5.8}{6.4}\selectfont \textcolor{BrickRed}{$ \blacktriangledown\,24\% $}}\end{tabular} & \hspace*{1pt}\begin{tabular}[t]{@{}c@{}}\rule{0pt}{12pt}26.9s\\[-.2pt]{\fontsize{5.8}{6.4}\selectfont \textcolor{ForestGreen}{$ \blacktriangle\,21\% $}}\\[-.2pt]23.3s\\[-.2pt]{\fontsize{5.8}{6.4}\selectfont \textcolor{ForestGreen}{$ \blacktriangle\,27\% $}}\end{tabular} & \hspace*{1pt}\begin{tabular}[t]{@{}c@{}}\rule{0pt}{12pt}1m\,59.3s\\[-.2pt]{\fontsize{5.8}{6.4}\selectfont \textcolor{BrickRed}{$ \blacktriangledown\,1.7\times $}}\\[-.2pt]1m\,51.3s\\[-.2pt]{\fontsize{5.8}{6.4}\selectfont \textcolor{BrickRed}{$ \blacktriangledown\,1.7\times $}}\end{tabular} & \hspace*{1pt}\begin{tabular}[t]{@{}c@{}}\rule{0pt}{12pt}55s\\[-.2pt]{\fontsize{5.8}{6.4}\selectfont \textcolor{ForestGreen}{$ \blacktriangle\,3.3\times $}}\\[-.2pt]50.9s\\[-.2pt]{\fontsize{5.8}{6.4}\selectfont \textcolor{ForestGreen}{$ \blacktriangle\,3.4\times $}}\end{tabular} & \hspace*{1pt}\begin{tabular}[t]{@{}c@{}}\rule{0pt}{12pt}5m\,27.2s\\[-.2pt]{\fontsize{5.8}{6.4}\selectfont \textcolor{BrickRed}{$ \blacktriangledown\,1.9\times $}}\\[-.2pt]5m\,16.6s\\[-.2pt]{\fontsize{5.8}{6.4}\selectfont \textcolor{BrickRed}{$ \blacktriangledown\,1.8\times $}}\end{tabular} & \hspace*{1pt}\begin{tabular}[t]{@{}c@{}}\rule{0pt}{12pt}5m\,50s\\[-.2pt]{\fontsize{5.8}{6.4}\selectfont \textcolor{BrickRed}{$ \blacktriangledown\,2.1\times $}}\\[-.2pt]5m\,36.1s\\[-.2pt]{\fontsize{5.8}{6.4}\selectfont \textcolor{BrickRed}{$ \blacktriangledown\,2.0\times $}}\end{tabular} & \textcolor{gray}{\ensuremath{\cdot}} & \textcolor{gray}{\ensuremath{\cdot}} & \hspace*{1pt}\begin{tabular}[t]{@{}c@{}}\rule{0pt}{12pt}11m\,14.3s\\[-.2pt]{\fontsize{5.8}{6.4}\selectfont \textcolor{BrickRed}{$ \blacktriangledown\,6.3\times $}}\\[-.2pt]5m\,8.1s\\[-.2pt]{\fontsize{5.8}{6.4}\selectfont \textcolor{BrickRed}{$ \blacktriangledown\,2.9\times $}}\end{tabular} & \hspace*{1pt}\begin{tabular}[t]{@{}c@{}}\rule{0pt}{12pt}9m\,48.4s\\[-.2pt]{\fontsize{5.8}{6.4}\selectfont \textcolor{BrickRed}{$ \blacktriangledown\,5.5\times $}}\\[-.2pt]5m\,16.4s\\[-.2pt]{\fontsize{5.8}{6.4}\selectfont \textcolor{BrickRed}{$ \blacktriangledown\,3.0\times $}}\end{tabular} & \textcolor{gray}{\ensuremath{\cdot}} & \textcolor{gray}{\ensuremath{\cdot}}\\[3pt]\hline
\BenchmarkSourceLabel{bench:aws-4-k15-w10:row:mpf-simdref-simd-ksw2}{R17}\shortstack[l]{{\color{gray}R17}} & \begin{tabular}[t]{@{}l@{}}\rule{0pt}{12pt}MPF / SIMD-ref\\SIMD-DP / KSW2\end{tabular} & \hspace*{1pt}\begin{tabular}[t]{@{}c@{}}\rule{0pt}{12pt}6.9s\\[-.2pt]{\fontsize{5.8}{6.4}\selectfont \textcolor{ForestGreen}{$ \blacktriangle\,5.2\times $}}\\[-.2pt]4.4s\\[-.2pt]{\fontsize{5.8}{6.4}\selectfont \textcolor{ForestGreen}{$ \blacktriangle\,8.0\times $}}\end{tabular} & \hspace*{1pt}\begin{tabular}[t]{@{}c@{}}\rule{0pt}{12pt}6.5s\\[-.2pt]{\fontsize{5.8}{6.4}\selectfont \textcolor{ForestGreen}{$ \blacktriangle\,5.9\times $}}\\[-.2pt]3.9s\\[-.2pt]{\fontsize{5.8}{6.4}\selectfont \textcolor{ForestGreen}{$ \blacktriangle\,9.2\times $}}\end{tabular} & \hspace*{1pt}\begin{tabular}[t]{@{}c@{}}\rule{0pt}{12pt}53.6s\\[-.2pt]{\fontsize{5.8}{6.4}\selectfont \textcolor{ForestGreen}{$ \blacktriangle\,1.7\times $}}\\[-.2pt]45.3s\\[-.2pt]{\fontsize{5.8}{6.4}\selectfont \textcolor{ForestGreen}{$ \blacktriangle\,2.0\times $}}\end{tabular} & \hspace*{1pt}\begin{tabular}[t]{@{}c@{}}\rule{0pt}{12pt}39.8s\\[-.2pt]{\fontsize{5.8}{6.4}\selectfont \textcolor{ForestGreen}{$ \blacktriangle\,9.8\times $}}\\[-.2pt]35.9s\\[-.2pt]{\fontsize{5.8}{6.4}\selectfont \textcolor{ForestGreen}{$ \blacktriangle\,10.8\times $}}\end{tabular} & \hspace*{1pt}\begin{tabular}[t]{@{}c@{}}\rule{0pt}{12pt}23.3s\\[-.2pt]{\fontsize{5.8}{6.4}\selectfont \textcolor{ForestGreen}{$ \blacktriangle\,26\% $}}\\[-.2pt]17.1s\\[-.2pt]{\fontsize{5.8}{6.4}\selectfont \textcolor{ForestGreen}{$ \blacktriangle\,1.8\times $}}\end{tabular} & \hspace*{1pt}\begin{tabular}[t]{@{}c@{}}\rule{0pt}{12pt}38.6s\\[-.2pt]{\fontsize{5.8}{6.4}\selectfont \textcolor{BrickRed}{$ \blacktriangledown\,13\% $}}\\[-.2pt]32.2s\\[-.2pt]{\fontsize{5.8}{6.4}\selectfont \textcolor{gray}{$ \blacktriangleright\,1\% $}}\end{tabular} & \hspace*{1pt}\begin{tabular}[t]{@{}c@{}}\rule{0pt}{12pt}46.7s\\[-.2pt]{\fontsize{5.8}{6.4}\selectfont \textcolor{ForestGreen}{$ \blacktriangle\,1.5\times $}}\\[-.2pt]34.1s\\[-.2pt]{\fontsize{5.8}{6.4}\selectfont \textcolor{ForestGreen}{$ \blacktriangle\,1.9\times $}}\end{tabular} & \hspace*{1pt}\begin{tabular}[t]{@{}c@{}}\rule{0pt}{12pt}51.3s\\[-.2pt]{\fontsize{5.8}{6.4}\selectfont \textcolor{ForestGreen}{$ \blacktriangle\,3.5\times $}}\\[-.2pt]47.9s\\[-.2pt]{\fontsize{5.8}{6.4}\selectfont \textcolor{ForestGreen}{$ \blacktriangle\,3.6\times $}}\end{tabular} & \hspace*{1pt}\begin{tabular}[t]{@{}c@{}}\rule{0pt}{12pt}4m\,57.4s\\[-.2pt]{\fontsize{5.8}{6.4}\selectfont \textcolor{BrickRed}{$ \blacktriangledown\,1.7\times $}}\\[-.2pt]4m\,44.6s\\[-.2pt]{\fontsize{5.8}{6.4}\selectfont \textcolor{BrickRed}{$ \blacktriangledown\,1.6\times $}}\end{tabular} & \hspace*{1pt}\begin{tabular}[t]{@{}c@{}}\rule{0pt}{12pt}5m\,23.4s\\[-.2pt]{\fontsize{5.8}{6.4}\selectfont \textcolor{BrickRed}{$ \blacktriangledown\,1.9\times $}}\\[-.2pt]5m\,7.5s\\[-.2pt]{\fontsize{5.8}{6.4}\selectfont \textcolor{BrickRed}{$ \blacktriangledown\,1.8\times $}}\end{tabular} & \textcolor{gray}{\ensuremath{\cdot}} & \textcolor{gray}{\ensuremath{\cdot}} & \textcolor{gray}{\ensuremath{\cdot}} & \textcolor{gray}{\ensuremath{\cdot}} & \textcolor{gray}{\ensuremath{\cdot}} & \textcolor{gray}{\ensuremath{\cdot}}\\[3pt]\hline
\BenchmarkSourceLabel{bench:aws-4-k15-w10:row:mpf-scalar-rmq-ksw2}{R18}\shortstack[l]{{\color{gray}R18}} & \begin{tabular}[t]{@{}l@{}}\rule{0pt}{12pt}MPF / Scalar\\RMQ / KSW2\end{tabular} & \hspace*{1pt}\begin{tabular}[t]{@{}c@{}}\rule{0pt}{12pt}7s\\[-.2pt]{\fontsize{5.8}{6.4}\selectfont \textcolor{ForestGreen}{$ \blacktriangle\,5.1\times $}}\\[-.2pt]4.5s\\[-.2pt]{\fontsize{5.8}{6.4}\selectfont \textcolor{ForestGreen}{$ \blacktriangle\,7.7\times $}}\end{tabular} & \hspace*{1pt}\begin{tabular}[t]{@{}c@{}}\rule{0pt}{12pt}6.5s\\[-.2pt]{\fontsize{5.8}{6.4}\selectfont \textcolor{ForestGreen}{$ \blacktriangle\,5.9\times $}}\\[-.2pt]3.9s\\[-.2pt]{\fontsize{5.8}{6.4}\selectfont \textcolor{ForestGreen}{$ \blacktriangle\,9.1\times $}}\end{tabular} & \hspace*{1pt}\begin{tabular}[t]{@{}c@{}}\rule{0pt}{12pt}17.7s\\[-.2pt]{\fontsize{5.8}{6.4}\selectfont \textcolor{ForestGreen}{$ \blacktriangle\,5.2\times $}}\\[-.2pt]9.5s\\[-.2pt]{\fontsize{5.8}{6.4}\selectfont \textcolor{ForestGreen}{$ \blacktriangle\,9.3\times $}}\end{tabular} & \hspace*{1pt}\begin{tabular}[t]{@{}c@{}}\rule{0pt}{12pt}36.3s\\[-.2pt]{\fontsize{5.8}{6.4}\selectfont \textcolor{ForestGreen}{$ \blacktriangle\,10.8\times $}}\\[-.2pt]32.3s\\[-.2pt]{\fontsize{5.8}{6.4}\selectfont \textcolor{ForestGreen}{$ \blacktriangle\,12.0\times $}}\end{tabular} & \hspace*{1pt}\begin{tabular}[t]{@{}c@{}}\rule{0pt}{12pt}26.4s\\[-.2pt]{\fontsize{5.8}{6.4}\selectfont \textcolor{ForestGreen}{$ \blacktriangle\,17\% $}}\\[-.2pt]21.9s\\[-.2pt]{\fontsize{5.8}{6.4}\selectfont \textcolor{ForestGreen}{$ \blacktriangle\,28\% $}}\end{tabular} & \hspace*{1pt}\begin{tabular}[t]{@{}c@{}}\rule{0pt}{12pt}21s\\[-.2pt]{\fontsize{5.8}{6.4}\selectfont \textcolor{ForestGreen}{$ \blacktriangle\,1.6\times $}}\\[-.2pt]17.6s\\[-.2pt]{\fontsize{5.8}{6.4}\selectfont \textcolor{ForestGreen}{$ \blacktriangle\,1.8\times $}}\end{tabular} & \hspace*{1pt}\begin{tabular}[t]{@{}c@{}}\rule{0pt}{12pt}35.4s\\[-.2pt]{\fontsize{5.8}{6.4}\selectfont \textcolor{ForestGreen}{$ \blacktriangle\,2.0\times $}}\\[-.2pt]27.1s\\[-.2pt]{\fontsize{5.8}{6.4}\selectfont \textcolor{ForestGreen}{$ \blacktriangle\,2.4\times $}}\end{tabular} & \hspace*{1pt}\begin{tabular}[t]{@{}c@{}}\rule{0pt}{12pt}51.9s\\[-.2pt]{\fontsize{5.8}{6.4}\selectfont \textcolor{ForestGreen}{$ \blacktriangle\,3.5\times $}}\\[-.2pt]48s\\[-.2pt]{\fontsize{5.8}{6.4}\selectfont \textcolor{ForestGreen}{$ \blacktriangle\,3.6\times $}}\end{tabular} & \hspace*{1pt}\begin{tabular}[t]{@{}c@{}}\rule{0pt}{12pt}5m\,15.8s\\[-.2pt]{\fontsize{5.8}{6.4}\selectfont \textcolor{BrickRed}{$ \blacktriangledown\,1.8\times $}}\\[-.2pt]5m\,2.5s\\[-.2pt]{\fontsize{5.8}{6.4}\selectfont \textcolor{BrickRed}{$ \blacktriangledown\,1.7\times $}}\end{tabular} & \hspace*{1pt}\begin{tabular}[t]{@{}c@{}}\rule{0pt}{12pt}5m\,19.7s\\[-.2pt]{\fontsize{5.8}{6.4}\selectfont \textcolor{BrickRed}{$ \blacktriangledown\,1.9\times $}}\\[-.2pt]5m\,5.9s\\[-.2pt]{\fontsize{5.8}{6.4}\selectfont \textcolor{BrickRed}{$ \blacktriangledown\,1.8\times $}}\end{tabular} & \textcolor{gray}{\ensuremath{\cdot}} & \textcolor{gray}{\ensuremath{\cdot}} & \hspace*{1pt}\begin{tabular}[t]{@{}c@{}}\rule{0pt}{12pt}11m\,18.9s\\[-.2pt]{\fontsize{5.8}{6.4}\selectfont \textcolor{BrickRed}{$ \blacktriangledown\,6.3\times $}}\\[-.2pt]5m\,12.1s\\[-.2pt]{\fontsize{5.8}{6.4}\selectfont \textcolor{BrickRed}{$ \blacktriangledown\,2.9\times $}}\end{tabular} & \hspace*{1pt}\begin{tabular}[t]{@{}c@{}}\rule{0pt}{12pt}9m\,52.6s\\[-.2pt]{\fontsize{5.8}{6.4}\selectfont \textcolor{BrickRed}{$ \blacktriangledown\,5.5\times $}}\\[-.2pt]5m\,20s\\[-.2pt]{\fontsize{5.8}{6.4}\selectfont \textcolor{BrickRed}{$ \blacktriangledown\,3.0\times $}}\end{tabular} & \textcolor{gray}{\ensuremath{\cdot}} & \textcolor{gray}{\ensuremath{\cdot}}\\[3pt]\hline
\BenchmarkSourceLabel{bench:aws-4-k15-w10:row:mpf-bounded-scalar-ksw2}{R19}\shortstack[l]{{\color{gray}R19}} & \begin{tabular}[t]{@{}l@{}}\rule{0pt}{12pt}MPF / Bounded\\Scalar-DP / KSW2\end{tabular} & \hspace*{1pt}\begin{tabular}[t]{@{}c@{}}\rule{0pt}{12pt}6.9s\\[-.2pt]{\fontsize{5.8}{6.4}\selectfont \textcolor{ForestGreen}{$ \blacktriangle\,5.3\times $}}\\[-.2pt]4.4s\\[-.2pt]{\fontsize{5.8}{6.4}\selectfont \textcolor{ForestGreen}{$ \blacktriangle\,7.9\times $}}\end{tabular} & \hspace*{1pt}\begin{tabular}[t]{@{}c@{}}\rule{0pt}{12pt}7.6s\\[-.2pt]{\fontsize{5.8}{6.4}\selectfont \textcolor{ForestGreen}{$ \blacktriangle\,5.0\times $}}\\[-.2pt]5.2s\\[-.2pt]{\fontsize{5.8}{6.4}\selectfont \textcolor{ForestGreen}{$ \blacktriangle\,6.9\times $}}\end{tabular} & \hspace*{1pt}\begin{tabular}[t]{@{}c@{}}\rule{0pt}{12pt}14.7s\\[-.2pt]{\fontsize{5.8}{6.4}\selectfont \textcolor{ForestGreen}{$ \blacktriangle\,6.3\times $}}\\[-.2pt]7s\\[-.2pt]{\fontsize{5.8}{6.4}\selectfont \textcolor{ForestGreen}{$ \blacktriangle\,12.7\times $}}\end{tabular} & \hspace*{1pt}\begin{tabular}[t]{@{}c@{}}\rule{0pt}{12pt}39.1s\\[-.2pt]{\fontsize{5.8}{6.4}\selectfont \textcolor{ForestGreen}{$ \blacktriangle\,10.0\times $}}\\[-.2pt]35.3s\\[-.2pt]{\fontsize{5.8}{6.4}\selectfont \textcolor{ForestGreen}{$ \blacktriangle\,11.0\times $}}\end{tabular} & \hspace*{1pt}\begin{tabular}[t]{@{}c@{}}\rule{0pt}{12pt}22.4s\\[-.2pt]{\fontsize{5.8}{6.4}\selectfont \textcolor{ForestGreen}{$ \blacktriangle\,29\% $}}\\[-.2pt]18.1s\\[-.2pt]{\fontsize{5.8}{6.4}\selectfont \textcolor{ForestGreen}{$ \blacktriangle\,1.7\times $}}\end{tabular} & \hspace*{1pt}\begin{tabular}[t]{@{}c@{}}\rule{0pt}{12pt}14.1s\\[-.2pt]{\fontsize{5.8}{6.4}\selectfont \textcolor{ForestGreen}{$ \blacktriangle\,2.4\times $}}\\[-.2pt]10s\\[-.2pt]{\fontsize{5.8}{6.4}\selectfont \textcolor{ForestGreen}{$ \blacktriangle\,3.2\times $}}\end{tabular} & \hspace*{1pt}\begin{tabular}[t]{@{}c@{}}\rule{0pt}{12pt}24.1s\\[-.2pt]{\fontsize{5.8}{6.4}\selectfont \textcolor{ForestGreen}{$ \blacktriangle\,2.9\times $}}\\[-.2pt]16s\\[-.2pt]{\fontsize{5.8}{6.4}\selectfont \textcolor{ForestGreen}{$ \blacktriangle\,4.1\times $}}\end{tabular} & \hspace*{1pt}\begin{tabular}[t]{@{}c@{}}\rule{0pt}{12pt}52s\\[-.2pt]{\fontsize{5.8}{6.4}\selectfont \textcolor{ForestGreen}{$ \blacktriangle\,3.5\times $}}\\[-.2pt]48.1s\\[-.2pt]{\fontsize{5.8}{6.4}\selectfont \textcolor{ForestGreen}{$ \blacktriangle\,3.6\times $}}\end{tabular} & \hspace*{1pt}\begin{tabular}[t]{@{}c@{}}\rule{0pt}{12pt}3m\,56.3s\\[-.2pt]{\fontsize{5.8}{6.4}\selectfont \textcolor{BrickRed}{$ \blacktriangledown\,1.3\times $}}\\[-.2pt]3m\,44.9s\\[-.2pt]{\fontsize{5.8}{6.4}\selectfont \textcolor{BrickRed}{$ \blacktriangledown\,29\% $}}\end{tabular} & \hspace*{1pt}\begin{tabular}[t]{@{}c@{}}\rule{0pt}{12pt}4m\,6.8s\\[-.2pt]{\fontsize{5.8}{6.4}\selectfont \textcolor{BrickRed}{$ \blacktriangledown\,1.4\times $}}\\[-.2pt]3m\,52.3s\\[-.2pt]{\fontsize{5.8}{6.4}\selectfont \textcolor{BrickRed}{$ \blacktriangledown\,1.4\times $}}\end{tabular} & \textcolor{gray}{\ensuremath{\cdot}} & \textcolor{gray}{\ensuremath{\cdot}} & \hspace*{1pt}\begin{tabular}[t]{@{}c@{}}\rule{0pt}{12pt}10m\,13.7s\\[-.2pt]{\fontsize{5.8}{6.4}\selectfont \textcolor{BrickRed}{$ \blacktriangledown\,5.7\times $}}\\[-.2pt]4m\,6.2s\\[-.2pt]{\fontsize{5.8}{6.4}\selectfont \textcolor{BrickRed}{$ \blacktriangledown\,2.3\times $}}\end{tabular} & \hspace*{1pt}\begin{tabular}[t]{@{}c@{}}\rule{0pt}{12pt}14m\,38.1s\\[-.2pt]{\fontsize{5.8}{6.4}\selectfont \textcolor{BrickRed}{$ \blacktriangledown\,8.2\times $}}\\[-.2pt]8m\,47s\\[-.2pt]{\fontsize{5.8}{6.4}\selectfont \textcolor{BrickRed}{$ \blacktriangledown\,5.0\times $}}\end{tabular} & \textcolor{gray}{\ensuremath{\cdot}} & \textcolor{gray}{\ensuremath{\cdot}}\\[3pt]\hline
\BenchmarkSourceLabel{bench:aws-4-k15-w10:row:mpf-twostack-simd-ksw2}{R20}\shortstack[l]{{\color{gray}R20}} & \begin{tabular}[t]{@{}l@{}}\rule{0pt}{12pt}MPF / Two-stack\\SIMD-DP / KSW2\end{tabular} & \hspace*{1pt}\begin{tabular}[t]{@{}c@{}}\rule{0pt}{12pt}6.9s\\[-.2pt]{\fontsize{5.8}{6.4}\selectfont \textcolor{ForestGreen}{$ \blacktriangle\,5.3\times $}}\\[-.2pt]\colorbox{BenchBest}{\strut 4.4s}\\[-.2pt]{\fontsize{5.8}{6.4}\selectfont \textcolor{ForestGreen}{$ \blacktriangle\,8.0\times $}}\end{tabular} & \hspace*{1pt}\begin{tabular}[t]{@{}c@{}}\rule{0pt}{12pt}7.3s\\[-.2pt]{\fontsize{5.8}{6.4}\selectfont \textcolor{ForestGreen}{$ \blacktriangle\,5.2\times $}}\\[-.2pt]4.8s\\[-.2pt]{\fontsize{5.8}{6.4}\selectfont \textcolor{ForestGreen}{$ \blacktriangle\,7.5\times $}}\end{tabular} & \hspace*{1pt}\begin{tabular}[t]{@{}c@{}}\rule{0pt}{12pt}12.5s\\[-.2pt]{\fontsize{5.8}{6.4}\selectfont \textcolor{ForestGreen}{$ \blacktriangle\,7.4\times $}}\\[-.2pt]6.8s\\[-.2pt]{\fontsize{5.8}{6.4}\selectfont \textcolor{ForestGreen}{$ \blacktriangle\,13.1\times $}}\end{tabular} & \hspace*{1pt}\begin{tabular}[t]{@{}c@{}}\rule{0pt}{12pt}41.9s\\[-.2pt]{\fontsize{5.8}{6.4}\selectfont \textcolor{ForestGreen}{$ \blacktriangle\,9.4\times $}}\\[-.2pt]37.8s\\[-.2pt]{\fontsize{5.8}{6.4}\selectfont \textcolor{ForestGreen}{$ \blacktriangle\,10.2\times $}}\end{tabular} & \hspace*{1pt}\begin{tabular}[t]{@{}c@{}}\rule{0pt}{12pt}28.6s\\[-.2pt]{\fontsize{5.8}{6.4}\selectfont \textcolor{ForestGreen}{$ \blacktriangle\,9\% $}}\\[-.2pt]24.4s\\[-.2pt]{\fontsize{5.8}{6.4}\selectfont \textcolor{ForestGreen}{$ \blacktriangle\,19\% $}}\end{tabular} & \hspace*{1pt}\begin{tabular}[t]{@{}c@{}}\rule{0pt}{12pt}24s\\[-.2pt]{\fontsize{5.8}{6.4}\selectfont \textcolor{ForestGreen}{$ \blacktriangle\,30\% $}}\\[-.2pt]20.6s\\[-.2pt]{\fontsize{5.8}{6.4}\selectfont \textcolor{ForestGreen}{$ \blacktriangle\,1.6\times $}}\end{tabular} & \hspace*{1pt}\begin{tabular}[t]{@{}c@{}}\rule{0pt}{12pt}57.1s\\[-.2pt]{\fontsize{5.8}{6.4}\selectfont \textcolor{ForestGreen}{$ \blacktriangle\,18\% $}}\\[-.2pt]49.1s\\[-.2pt]{\fontsize{5.8}{6.4}\selectfont \textcolor{ForestGreen}{$ \blacktriangle\,26\% $}}\end{tabular} & \textcolor{gray}{\ensuremath{\cdot}} & \hspace*{1pt}\begin{tabular}[t]{@{}c@{}}\rule{0pt}{12pt}4m\,37.5s\\[-.2pt]{\fontsize{5.8}{6.4}\selectfont \textcolor{BrickRed}{$ \blacktriangledown\,1.6\times $}}\\[-.2pt]4m\,25.9s\\[-.2pt]{\fontsize{5.8}{6.4}\selectfont \textcolor{BrickRed}{$ \blacktriangledown\,1.5\times $}}\end{tabular} & \hspace*{1pt}\begin{tabular}[t]{@{}c@{}}\rule{0pt}{12pt}4m\,48.3s\\[-.2pt]{\fontsize{5.8}{6.4}\selectfont \textcolor{BrickRed}{$ \blacktriangledown\,1.7\times $}}\\[-.2pt]4m\,34.5s\\[-.2pt]{\fontsize{5.8}{6.4}\selectfont \textcolor{BrickRed}{$ \blacktriangledown\,1.6\times $}}\end{tabular} & \textcolor{gray}{\ensuremath{\cdot}} & \textcolor{gray}{\ensuremath{\cdot}} & \hspace*{1pt}\begin{tabular}[t]{@{}c@{}}\rule{0pt}{12pt}10m\,6.4s\\[-.2pt]{\fontsize{5.8}{6.4}\selectfont \textcolor{BrickRed}{$ \blacktriangledown\,5.7\times $}}\\[-.2pt]3m\,55.9s\\[-.2pt]{\fontsize{5.8}{6.4}\selectfont \textcolor{BrickRed}{$ \blacktriangledown\,2.2\times $}}\end{tabular} & \hspace*{1pt}\begin{tabular}[t]{@{}c@{}}\rule{0pt}{12pt}8m\,51.6s\\[-.2pt]{\fontsize{5.8}{6.4}\selectfont \textcolor{BrickRed}{$ \blacktriangledown\,5.0\times $}}\\[-.2pt]3m\,47s\\[-.2pt]{\fontsize{5.8}{6.4}\selectfont \textcolor{BrickRed}{$ \blacktriangledown\,2.2\times $}}\end{tabular} & \textcolor{gray}{\ensuremath{\cdot}} & \textcolor{gray}{\ensuremath{\cdot}}\\[3pt]\hline
\BenchmarkSourceLabel{bench:aws-4-k15-w10:row:mpf-bounded-rmq-ksw2}{R21}\shortstack[l]{{\color{gray}R21}} & \begin{tabular}[t]{@{}l@{}}\rule{0pt}{12pt}MPF / Bounded\\RMQ / KSW2\end{tabular} & \hspace*{1pt}\begin{tabular}[t]{@{}c@{}}\rule{0pt}{12pt}6.9s\\[-.2pt]{\fontsize{5.8}{6.4}\selectfont \textcolor{ForestGreen}{$ \blacktriangle\,5.2\times $}}\\[-.2pt]4.5s\\[-.2pt]{\fontsize{5.8}{6.4}\selectfont \textcolor{ForestGreen}{$ \blacktriangle\,7.9\times $}}\end{tabular} & \hspace*{1pt}\begin{tabular}[t]{@{}c@{}}\rule{0pt}{12pt}9.6s\\[-.2pt]{\fontsize{5.8}{6.4}\selectfont \textcolor{ForestGreen}{$ \blacktriangle\,3.9\times $}}\\[-.2pt]7.2s\\[-.2pt]{\fontsize{5.8}{6.4}\selectfont \textcolor{ForestGreen}{$ \blacktriangle\,5.0\times $}}\end{tabular} & \hspace*{1pt}\begin{tabular}[t]{@{}c@{}}\rule{0pt}{12pt}14.6s\\[-.2pt]{\fontsize{5.8}{6.4}\selectfont \textcolor{ForestGreen}{$ \blacktriangle\,6.3\times $}}\\[-.2pt]6.7s\\[-.2pt]{\fontsize{5.8}{6.4}\selectfont \textcolor{ForestGreen}{$ \blacktriangle\,13.1\times $}}\end{tabular} & \hspace*{1pt}\begin{tabular}[t]{@{}c@{}}\rule{0pt}{12pt}39.9s\\[-.2pt]{\fontsize{5.8}{6.4}\selectfont \textcolor{ForestGreen}{$ \blacktriangle\,9.8\times $}}\\[-.2pt]35.8s\\[-.2pt]{\fontsize{5.8}{6.4}\selectfont \textcolor{ForestGreen}{$ \blacktriangle\,10.8\times $}}\end{tabular} & \hspace*{1pt}\begin{tabular}[t]{@{}c@{}}\rule{0pt}{12pt}30.9s\\[-.2pt]{\fontsize{5.8}{6.4}\selectfont \textcolor{gray}{$ \blacktriangleright\,2\% $}}\\[-.2pt]26.5s\\[-.2pt]{\fontsize{5.8}{6.4}\selectfont \textcolor{ForestGreen}{$ \blacktriangle\,12\% $}}\end{tabular} & \hspace*{1pt}\begin{tabular}[t]{@{}c@{}}\rule{0pt}{12pt}37.1s\\[-.2pt]{\fontsize{5.8}{6.4}\selectfont \textcolor{BrickRed}{$ \blacktriangledown\,9\% $}}\\[-.2pt]33.3s\\[-.2pt]{\fontsize{5.8}{6.4}\selectfont \textcolor{gray}{$ \blacktriangleright\,4\% $}}\end{tabular} & \hspace*{1pt}\begin{tabular}[t]{@{}c@{}}\rule{0pt}{12pt}24.1s\\[-.2pt]{\fontsize{5.8}{6.4}\selectfont \textcolor{ForestGreen}{$ \blacktriangle\,2.9\times $}}\\[-.2pt]15.8s\\[-.2pt]{\fontsize{5.8}{6.4}\selectfont \textcolor{ForestGreen}{$ \blacktriangle\,4.2\times $}}\end{tabular} & \hspace*{1pt}\begin{tabular}[t]{@{}c@{}}\rule{0pt}{12pt}58.6s\\[-.2pt]{\fontsize{5.8}{6.4}\selectfont \textcolor{ForestGreen}{$ \blacktriangle\,3.1\times $}}\\[-.2pt]54.5s\\[-.2pt]{\fontsize{5.8}{6.4}\selectfont \textcolor{ForestGreen}{$ \blacktriangle\,3.1\times $}}\end{tabular} & \hspace*{1pt}\begin{tabular}[t]{@{}c@{}}\rule{0pt}{12pt}3m\,41.3s\\[-.2pt]{\fontsize{5.8}{6.4}\selectfont \textcolor{BrickRed}{$ \blacktriangledown\,26\% $}}\\[-.2pt]3m\,30.7s\\[-.2pt]{\fontsize{5.8}{6.4}\selectfont \textcolor{BrickRed}{$ \blacktriangledown\,21\% $}}\end{tabular} & \hspace*{1pt}\begin{tabular}[t]{@{}c@{}}\rule{0pt}{12pt}3m\,59.5s\\[-.2pt]{\fontsize{5.8}{6.4}\selectfont \textcolor{BrickRed}{$ \blacktriangledown\,1.4\times $}}\\[-.2pt]3m\,45.5s\\[-.2pt]{\fontsize{5.8}{6.4}\selectfont \textcolor{BrickRed}{$ \blacktriangledown\,1.3\times $}}\end{tabular} & \textcolor{gray}{\ensuremath{\cdot}} & \textcolor{gray}{\ensuremath{\cdot}} & \hspace*{1pt}\begin{tabular}[t]{@{}c@{}}\rule{0pt}{12pt}10m\,2.7s\\[-.2pt]{\fontsize{5.8}{6.4}\selectfont \textcolor{BrickRed}{$ \blacktriangledown\,5.6\times $}}\\[-.2pt]3m\,55.6s\\[-.2pt]{\fontsize{5.8}{6.4}\selectfont \textcolor{BrickRed}{$ \blacktriangledown\,2.2\times $}}\end{tabular} & \hspace*{1pt}\begin{tabular}[t]{@{}c@{}}\rule{0pt}{12pt}15m\,4.3s\\[-.2pt]{\fontsize{5.8}{6.4}\selectfont \textcolor{BrickRed}{$ \blacktriangledown\,8.4\times $}}\\[-.2pt]8m\,30.7s\\[-.2pt]{\fontsize{5.8}{6.4}\selectfont \textcolor{BrickRed}{$ \blacktriangledown\,4.9\times $}}\end{tabular} & \textcolor{gray}{\ensuremath{\cdot}} & \textcolor{gray}{\ensuremath{\cdot}}\\[3pt]\hline
\BenchmarkSourceLabel{bench:aws-4-k15-w10:row:mpf-bounded-simd-ksw2}{R22}\shortstack[l]{{\color{gray}R22}} & \begin{tabular}[t]{@{}l@{}}\rule{0pt}{12pt}MPF / Bounded\\SIMD-DP / KSW2\end{tabular} & \hspace*{1pt}\begin{tabular}[t]{@{}c@{}}\rule{0pt}{12pt}14.7s\\[-.2pt]{\fontsize{5.8}{6.4}\selectfont \textcolor{ForestGreen}{$ \blacktriangle\,2.5\times $}}\\[-.2pt]12.1s\\[-.2pt]{\fontsize{5.8}{6.4}\selectfont \textcolor{ForestGreen}{$ \blacktriangle\,2.9\times $}}\end{tabular} & \hspace*{1pt}\begin{tabular}[t]{@{}c@{}}\rule{0pt}{12pt}8.1s\\[-.2pt]{\fontsize{5.8}{6.4}\selectfont \textcolor{ForestGreen}{$ \blacktriangle\,4.7\times $}}\\[-.2pt]5.5s\\[-.2pt]{\fontsize{5.8}{6.4}\selectfont \textcolor{ForestGreen}{$ \blacktriangle\,6.5\times $}}\end{tabular} & \hspace*{1pt}\begin{tabular}[t]{@{}c@{}}\rule{0pt}{12pt}24.9s\\[-.2pt]{\fontsize{5.8}{6.4}\selectfont \textcolor{ForestGreen}{$ \blacktriangle\,3.7\times $}}\\[-.2pt]16.9s\\[-.2pt]{\fontsize{5.8}{6.4}\selectfont \textcolor{ForestGreen}{$ \blacktriangle\,5.2\times $}}\end{tabular} & \hspace*{1pt}\begin{tabular}[t]{@{}c@{}}\rule{0pt}{12pt}34.6s\\[-.2pt]{\fontsize{5.8}{6.4}\selectfont \textcolor{ForestGreen}{$ \blacktriangle\,11.3\times $}}\\[-.2pt]\colorbox{BenchBest}{\strut 30.4s}\\[-.2pt]{\fontsize{5.8}{6.4}\selectfont \textcolor{ForestGreen}{$ \blacktriangle\,12.7\times $}}\end{tabular} & \hspace*{1pt}\begin{tabular}[t]{@{}c@{}}\rule{0pt}{12pt}30s\\[-.2pt]{\fontsize{5.8}{6.4}\selectfont \textcolor{ForestGreen}{$ \blacktriangle\,5\% $}}\\[-.2pt]25.5s\\[-.2pt]{\fontsize{5.8}{6.4}\selectfont \textcolor{ForestGreen}{$ \blacktriangle\,16\% $}}\end{tabular} & \hspace*{1pt}\begin{tabular}[t]{@{}c@{}}\rule{0pt}{12pt}38.5s\\[-.2pt]{\fontsize{5.8}{6.4}\selectfont \textcolor{BrickRed}{$ \blacktriangledown\,13\% $}}\\[-.2pt]34.3s\\[-.2pt]{\fontsize{5.8}{6.4}\selectfont \textcolor{BrickRed}{$ \blacktriangledown\,7\% $}}\end{tabular} & \hspace*{1pt}\begin{tabular}[t]{@{}c@{}}\rule{0pt}{12pt}51.7s\\[-.2pt]{\fontsize{5.8}{6.4}\selectfont \textcolor{ForestGreen}{$ \blacktriangle\,26\% $}}\\[-.2pt]43.6s\\[-.2pt]{\fontsize{5.8}{6.4}\selectfont \textcolor{ForestGreen}{$ \blacktriangle\,1.5\times $}}\end{tabular} & \hspace*{1pt}\begin{tabular}[t]{@{}c@{}}\rule{0pt}{12pt}55.6s\\[-.2pt]{\fontsize{5.8}{6.4}\selectfont \textcolor{ForestGreen}{$ \blacktriangle\,3.3\times $}}\\[-.2pt]51.7s\\[-.2pt]{\fontsize{5.8}{6.4}\selectfont \textcolor{ForestGreen}{$ \blacktriangle\,3.3\times $}}\end{tabular} & \hspace*{1pt}\begin{tabular}[t]{@{}c@{}}\rule{0pt}{12pt}3m\,57.7s\\[-.2pt]{\fontsize{5.8}{6.4}\selectfont \textcolor{BrickRed}{$ \blacktriangledown\,1.4\times $}}\\[-.2pt]3m\,46.9s\\[-.2pt]{\fontsize{5.8}{6.4}\selectfont \textcolor{BrickRed}{$ \blacktriangledown\,1.3\times $}}\end{tabular} & \hspace*{1pt}\begin{tabular}[t]{@{}c@{}}\rule{0pt}{12pt}4m\,0.1s\\[-.2pt]{\fontsize{5.8}{6.4}\selectfont \textcolor{BrickRed}{$ \blacktriangledown\,1.4\times $}}\\[-.2pt]3m\,46s\\[-.2pt]{\fontsize{5.8}{6.4}\selectfont \textcolor{BrickRed}{$ \blacktriangledown\,1.3\times $}}\end{tabular} & \textcolor{gray}{\ensuremath{\cdot}} & \textcolor{gray}{\ensuremath{\cdot}} & \hspace*{1pt}\begin{tabular}[t]{@{}c@{}}\rule{0pt}{12pt}10m\,21.9s\\[-.2pt]{\fontsize{5.8}{6.4}\selectfont \textcolor{BrickRed}{$ \blacktriangledown\,5.8\times $}}\\[-.2pt]4m\,12s\\[-.2pt]{\fontsize{5.8}{6.4}\selectfont \textcolor{BrickRed}{$ \blacktriangledown\,2.4\times $}}\end{tabular} & \hspace*{1pt}\begin{tabular}[t]{@{}c@{}}\rule{0pt}{12pt}11m\,12.6s\\[-.2pt]{\fontsize{5.8}{6.4}\selectfont \textcolor{BrickRed}{$ \blacktriangledown\,6.3\times $}}\\[-.2pt]5m\,0.9s\\[-.2pt]{\fontsize{5.8}{6.4}\selectfont \textcolor{BrickRed}{$ \blacktriangledown\,2.9\times $}}\end{tabular} & \textcolor{gray}{\ensuremath{\cdot}} & \textcolor{gray}{\ensuremath{\cdot}}\\[3pt]\hline
\end{tabular}\par}
\smallskip{\footnotesize\noindent Each cell: A and B times, $\cdot$: no measurement. ---\textsuperscript{g}: GPU unavailable. n/a: no comparison.\par
\noindent Blue: one lowest A and one lowest B measurement for each column for each configuration.\par
\noindent Gray: minimap2. Yellow: my preferred MPF / Bounded / SIMD-DP / Block configuration.\par
\noindent MPF/MPS: minimizers per fragment/sequence. SIMD-ref: naive SIMD scan. Bounded: bounded SIMD algorithm.\par
}
\clearpage
\noindent{\large\bfseries Table \thetable: (B): AWS: 4 Gbp, k=15, w=10\hfill 4/7}\par\smallskip
{\fontsize{7.1}{8.1}\selectfont\setlength{\tabcolsep}{1pt}\setlength{\fboxsep}{2pt}\renewcommand{\arraystretch}{1.12}
\noindent\begin{tabular}{@{}p{14mm}p{27mm}|p{13.55mm}|p{13.55mm}|p{13.55mm}|p{13.55mm}|p{13.55mm}|p{13.55mm}|p{13.55mm}|p{13.55mm}|p{13.55mm}|p{13.55mm}|p{13.55mm}|p{13.55mm}|p{13.55mm}|p{13.55mm}|p{13.55mm}|p{13.55mm}|@{}}\hline
Engine & Execution & \multicolumn{8}{c|}{\textbf{Index}} & \multicolumn{8}{c|}{\textbf{No index}}\\
 &  & \multicolumn{8}{c|}{\textbf{4 Gbp}} & \multicolumn{8}{c|}{\textbf{4 Gbp}}\\
 &  & \multicolumn{8}{c|}{\textbf{k=15, w=10}} & \multicolumn{8}{c|}{\textbf{k=15, w=10}}\\
 &  & \multicolumn{4}{c|}{\textbf{DNA4}} & \multicolumn{4}{c|}{\textbf{DNA5}} & \multicolumn{4}{c|}{\textbf{DNA4}} & \multicolumn{4}{c|}{\textbf{DNA5}}\\
 &  & \multicolumn{1}{c|}{\textbf{M7A}} & \multicolumn{1}{c|}{\textbf{M7I}} & \multicolumn{1}{c|}{\textbf{M4}} & \multicolumn{1}{c|}{\textbf{C7G}} & \multicolumn{1}{c|}{\textbf{M7A}} & \multicolumn{1}{c|}{\textbf{M7I}} & \multicolumn{1}{c|}{\textbf{M4}} & \multicolumn{1}{c|}{\textbf{C7G}} & \multicolumn{1}{c|}{\textbf{M7A}} & \multicolumn{1}{c|}{\textbf{M7I}} & \multicolumn{1}{c|}{\textbf{M4}} & \multicolumn{1}{c|}{\textbf{C7G}} & \multicolumn{1}{c|}{\textbf{M7A}} & \multicolumn{1}{c|}{\textbf{M7I}} & \multicolumn{1}{c|}{\textbf{M4}} & \multicolumn{1}{c|}{\textbf{C7G}}\\
 & & {\color{gray}C1} & {\color{gray}C2} & {\color{gray}C3} & {\color{gray}C4} & {\color{gray}C5} & {\color{gray}C6} & {\color{gray}C7} & {\color{gray}C8} & {\color{gray}C9} & {\color{gray}C10} & {\color{gray}C11} & {\color{gray}C12} & {\color{gray}C13} & {\color{gray}C14} & {\color{gray}C15} & {\color{gray}C16}\\
\hline
\rowcolor{BenchBaseline}[0.5pt][0.5pt]\shortstack[l]{{\color{gray}R1}\\\textbf{minimap2}} & \textbf{Baseline} & \hspace*{1pt}\begin{tabular}[t]{@{}c@{}}\rule{0pt}{12pt}36.2s\\[-.2pt]{\fontsize{5.8}{6.4}\selectfont \textcolor{BrickRed}{$ \blacktriangledown\,6.9\times $}}\\[-.2pt]35s\\[-.2pt]{\fontsize{5.8}{6.4}\selectfont \textcolor{BrickRed}{$ \blacktriangledown\,8.0\times $}}\end{tabular} & \hspace*{1pt}\begin{tabular}[t]{@{}c@{}}\rule{0pt}{12pt}38.1s\\[-.2pt]{\fontsize{5.8}{6.4}\selectfont \textcolor{BrickRed}{$ \blacktriangledown\,7.8\times $}}\\[-.2pt]36s\\[-.2pt]{\fontsize{5.8}{6.4}\selectfont \textcolor{BrickRed}{$ \blacktriangledown\,9.3\times $}}\end{tabular} & \hspace*{1pt}\begin{tabular}[t]{@{}c@{}}\rule{0pt}{12pt}1m\,32.4s\\[-.2pt]{\fontsize{5.8}{6.4}\selectfont \textcolor{BrickRed}{$ \blacktriangledown\,10.7\times $}}\\[-.2pt]1m\,28.5s\\[-.2pt]{\fontsize{5.8}{6.4}\selectfont \textcolor{BrickRed}{$ \blacktriangledown\,13.5\times $}}\end{tabular} & \hspace*{1pt}\begin{tabular}[t]{@{}c@{}}\rule{0pt}{12pt}6m\,31.5s\\[-.2pt]{\fontsize{5.8}{6.4}\selectfont \textcolor{BrickRed}{$ \blacktriangledown\,11.5\times $}}\\[-.2pt]6m\,26.3s\\[-.2pt]{\fontsize{5.8}{6.4}\selectfont \textcolor{BrickRed}{$ \blacktriangledown\,12.7\times $}}\end{tabular} & \hspace*{1pt}\begin{tabular}[t]{@{}c@{}}\rule{0pt}{12pt}31.6s\\[-.2pt]{\fontsize{5.8}{6.4}\selectfont \textcolor{BrickRed}{$ \blacktriangledown\,2.1\times $}}\\[-.2pt]30.3s\\[-.2pt]{\fontsize{5.8}{6.4}\selectfont \textcolor{BrickRed}{$ \blacktriangledown\,3.0\times $}}\end{tabular} & \hspace*{1pt}\begin{tabular}[t]{@{}c@{}}\rule{0pt}{12pt}34.1s\\[-.2pt]{\fontsize{5.8}{6.4}\selectfont \textcolor{BrickRed}{$ \blacktriangledown\,2.9\times $}}\\[-.2pt]32s\\[-.2pt]{\fontsize{5.8}{6.4}\selectfont \textcolor{BrickRed}{$ \blacktriangledown\,3.7\times $}}\end{tabular} & \hspace*{1pt}\begin{tabular}[t]{@{}c@{}}\rule{0pt}{12pt}1m\,9.5s\\[-.2pt]{\fontsize{5.8}{6.4}\selectfont \textcolor{BrickRed}{$ \blacktriangledown\,3.0\times $}}\\[-.2pt]1m\,6.1s\\[-.2pt]{\fontsize{5.8}{6.4}\selectfont \textcolor{BrickRed}{$ \blacktriangledown\,4.3\times $}}\end{tabular} & \hspace*{1pt}\begin{tabular}[t]{@{}c@{}}\rule{0pt}{12pt}3m\,1s\\[-.2pt]{\fontsize{5.8}{6.4}\selectfont \textcolor{BrickRed}{$ \blacktriangledown\,7.9\times $}}\\[-.2pt]2m\,51.5s\\[-.2pt]{\fontsize{5.8}{6.4}\selectfont \textcolor{BrickRed}{$ \blacktriangledown\,8.6\times $}}\end{tabular} & \hspace*{1pt}\begin{tabular}[t]{@{}c@{}}\rule{0pt}{12pt}2m\,55.5s\\[-.2pt]{\fontsize{5.8}{6.4}\selectfont \textcolor{ForestGreen}{$ \blacktriangle\,20\% $}}\\[-.2pt]2m\,54s\\[-.2pt]{\fontsize{5.8}{6.4}\selectfont \textcolor{ForestGreen}{$ \blacktriangle\,17\% $}}\end{tabular} & \hspace*{1pt}\begin{tabular}[t]{@{}c@{}}\rule{0pt}{12pt}2m\,50.4s\\[-.2pt]{\fontsize{5.8}{6.4}\selectfont \textcolor{ForestGreen}{$ \blacktriangle\,23\% $}}\\[-.2pt]2m\,48.2s\\[-.2pt]{\fontsize{5.8}{6.4}\selectfont \textcolor{ForestGreen}{$ \blacktriangle\,20\% $}}\end{tabular} & \hspace*{1pt}\begin{tabular}[t]{@{}c@{}}\rule{0pt}{12pt}5m\,35.7s\\[-.2pt]{\fontsize{5.8}{6.4}\selectfont \textcolor{ForestGreen}{$ \blacktriangle\,18\% $}}\\[-.2pt]5m\,31.9s\\[-.2pt]{\fontsize{5.8}{6.4}\selectfont \textcolor{ForestGreen}{$ \blacktriangle\,14\% $}}\end{tabular} & \hspace*{1pt}\begin{tabular}[t]{@{}c@{}}\rule{0pt}{12pt}44m\,20.5s\\[-.2pt]{\fontsize{5.8}{6.4}\selectfont \textcolor{gray}{n/a}}\\[-.2pt]44m\,15.6s\\[-.2pt]{\fontsize{5.8}{6.4}\selectfont \textcolor{gray}{n/a}}\end{tabular} & \hspace*{1pt}\begin{tabular}[t]{@{}c@{}}\rule{0pt}{12pt}1m\,47.1s\\[-.2pt]{\fontsize{5.8}{6.4}\selectfont \textcolor{ForestGreen}{$ \blacktriangle\,5.5\times $}}\\[-.2pt]1m\,45.8s\\[-.2pt]{\fontsize{5.8}{6.4}\selectfont \textcolor{ForestGreen}{$ \blacktriangle\,2.1\times $}}\end{tabular} & \hspace*{1pt}\begin{tabular}[t]{@{}c@{}}\rule{0pt}{12pt}1m\,47.2s\\[-.2pt]{\fontsize{5.8}{6.4}\selectfont \textcolor{ForestGreen}{$ \blacktriangle\,4.6\times $}}\\[-.2pt]1m\,45.1s\\[-.2pt]{\fontsize{5.8}{6.4}\selectfont \textcolor{ForestGreen}{$ \blacktriangle\,2.1\times $}}\end{tabular} & \hspace*{1pt}\begin{tabular}[t]{@{}c@{}}\rule{0pt}{12pt}3m\,40.3s\\[-.2pt]{\fontsize{5.8}{6.4}\selectfont \textcolor{ForestGreen}{$ \blacktriangle\,4.2\times $}}\\[-.2pt]3m\,36.6s\\[-.2pt]{\fontsize{5.8}{6.4}\selectfont \textcolor{ForestGreen}{$ \blacktriangle\,1.7\times $}}\end{tabular} & \hspace*{1pt}\begin{tabular}[t]{@{}c@{}}\rule{0pt}{12pt}29m\,43.4s\\[-.2pt]{\fontsize{5.8}{6.4}\selectfont \textcolor{gray}{n/a}}\\[-.2pt]29m\,38.5s\\[-.2pt]{\fontsize{5.8}{6.4}\selectfont \textcolor{gray}{n/a}}\end{tabular}\\[3pt]\hline
\rowcolor{BenchPreferred}[0.5pt][0.5pt]\BenchmarkSourceLabel{bench:aws-4-k15-w10:row:mpf-bounded-simd-block}{R23}\shortstack[l]{{\color{gray}R23}\\\textbf{BioseqDB}} & \begin{tabular}[t]{@{}l@{}}\rule{0pt}{12pt}MPF / Bounded\\SIMD-DP / Block\end{tabular} & \hspace*{1pt}\begin{tabular}[t]{@{}c@{}}\rule{0pt}{12pt}9.9s\\[-.2pt]{\fontsize{5.8}{6.4}\selectfont \textcolor{ForestGreen}{$ \blacktriangle\,3.6\times $}}\\[-.2pt]9s\\[-.2pt]{\fontsize{5.8}{6.4}\selectfont \textcolor{ForestGreen}{$ \blacktriangle\,3.9\times $}}\end{tabular} & \hspace*{1pt}\begin{tabular}[t]{@{}c@{}}\rule{0pt}{12pt}13s\\[-.2pt]{\fontsize{5.8}{6.4}\selectfont \textcolor{ForestGreen}{$ \blacktriangle\,2.9\times $}}\\[-.2pt]11.9s\\[-.2pt]{\fontsize{5.8}{6.4}\selectfont \textcolor{ForestGreen}{$ \blacktriangle\,3.0\times $}}\end{tabular} & \hspace*{1pt}\begin{tabular}[t]{@{}c@{}}\rule{0pt}{12pt}9.4s\\[-.2pt]{\fontsize{5.8}{6.4}\selectfont \textcolor{ForestGreen}{$ \blacktriangle\,9.8\times $}}\\[-.2pt]7s\\[-.2pt]{\fontsize{5.8}{6.4}\selectfont \textcolor{ForestGreen}{$ \blacktriangle\,12.7\times $}}\end{tabular} & \hspace*{1pt}\begin{tabular}[t]{@{}c@{}}\rule{0pt}{12pt}35s\\[-.2pt]{\fontsize{5.8}{6.4}\selectfont \textcolor{ForestGreen}{$ \blacktriangle\,11.2\times $}}\\[-.2pt]33.8s\\[-.2pt]{\fontsize{5.8}{6.4}\selectfont \textcolor{ForestGreen}{$ \blacktriangle\,11.4\times $}}\end{tabular} & \hspace*{1pt}\begin{tabular}[t]{@{}c@{}}\rule{0pt}{12pt}23.1s\\[-.2pt]{\fontsize{5.8}{6.4}\selectfont \textcolor{ForestGreen}{$ \blacktriangle\,27\% $}}\\[-.2pt]18.6s\\[-.2pt]{\fontsize{5.8}{6.4}\selectfont \textcolor{ForestGreen}{$ \blacktriangle\,1.6\times $}}\end{tabular} & \hspace*{1pt}\begin{tabular}[t]{@{}c@{}}\rule{0pt}{12pt}25.2s\\[-.2pt]{\fontsize{5.8}{6.4}\selectfont \textcolor{ForestGreen}{$ \blacktriangle\,26\% $}}\\[-.2pt]21.2s\\[-.2pt]{\fontsize{5.8}{6.4}\selectfont \textcolor{ForestGreen}{$ \blacktriangle\,1.5\times $}}\end{tabular} & \hspace*{1pt}\begin{tabular}[t]{@{}c@{}}\rule{0pt}{12pt}54.4s\\[-.2pt]{\fontsize{5.8}{6.4}\selectfont \textcolor{ForestGreen}{$ \blacktriangle\,22\% $}}\\[-.2pt]46.7s\\[-.2pt]{\fontsize{5.8}{6.4}\selectfont \textcolor{ForestGreen}{$ \blacktriangle\,29\% $}}\end{tabular} & \hspace*{1pt}\begin{tabular}[t]{@{}c@{}}\rule{0pt}{12pt}49.8s\\[-.2pt]{\fontsize{5.8}{6.4}\selectfont \textcolor{ForestGreen}{$ \blacktriangle\,3.6\times $}}\\[-.2pt]46.4s\\[-.2pt]{\fontsize{5.8}{6.4}\selectfont \textcolor{ForestGreen}{$ \blacktriangle\,3.7\times $}}\end{tabular} & \hspace*{1pt}\begin{tabular}[t]{@{}c@{}}\rule{0pt}{12pt}3m\,44.5s\\[-.2pt]{\fontsize{5.8}{6.4}\selectfont \textcolor{BrickRed}{$ \blacktriangledown\,28\% $}}\\[-.2pt]3m\,35.9s\\[-.2pt]{\fontsize{5.8}{6.4}\selectfont \textcolor{BrickRed}{$ \blacktriangledown\,24\% $}}\end{tabular} & \hspace*{1pt}\begin{tabular}[t]{@{}c@{}}\rule{0pt}{12pt}3m\,57.4s\\[-.2pt]{\fontsize{5.8}{6.4}\selectfont \textcolor{BrickRed}{$ \blacktriangledown\,1.4\times $}}\\[-.2pt]3m\,45.7s\\[-.2pt]{\fontsize{5.8}{6.4}\selectfont \textcolor{BrickRed}{$ \blacktriangledown\,1.3\times $}}\end{tabular} & \hspace*{1pt}\begin{tabular}[t]{@{}c@{}}\rule{0pt}{12pt}6m\,50.6s\\[-.2pt]{\fontsize{5.8}{6.4}\selectfont \textcolor{BrickRed}{$ \blacktriangledown\,22\% $}}\\[-.2pt]6m\,30.7s\\[-.2pt]{\fontsize{5.8}{6.4}\selectfont \textcolor{BrickRed}{$ \blacktriangledown\,18\% $}}\end{tabular} & \textcolor{gray}{\ensuremath{\cdot}} & \hspace*{1pt}\begin{tabular}[t]{@{}c@{}}\rule{0pt}{12pt}9m\,45.5s\\[-.2pt]{\fontsize{5.8}{6.4}\selectfont \textcolor{BrickRed}{$ \blacktriangledown\,5.5\times $}}\\[-.2pt]3m\,41.5s\\[-.2pt]{\fontsize{5.8}{6.4}\selectfont \textcolor{BrickRed}{$ \blacktriangledown\,2.1\times $}}\end{tabular} & \hspace*{1pt}\begin{tabular}[t]{@{}c@{}}\rule{0pt}{12pt}8m\,25.6s\\[-.2pt]{\fontsize{5.8}{6.4}\selectfont \textcolor{BrickRed}{$ \blacktriangledown\,4.7\times $}}\\[-.2pt]3m\,56.7s\\[-.2pt]{\fontsize{5.8}{6.4}\selectfont \textcolor{BrickRed}{$ \blacktriangledown\,2.3\times $}}\end{tabular} & \hspace*{1pt}\begin{tabular}[t]{@{}c@{}}\rule{0pt}{12pt}15m\,24.9s\\[-.2pt]{\fontsize{5.8}{6.4}\selectfont \textcolor{BrickRed}{$ \blacktriangledown\,4.2\times $}}\\[-.2pt]6m\,0.2s\\[-.2pt]{\fontsize{5.8}{6.4}\selectfont \textcolor{BrickRed}{$ \blacktriangledown\,1.7\times $}}\end{tabular} & \textcolor{gray}{\ensuremath{\cdot}}\\[3pt]\hline
\BenchmarkSourceLabel{bench:aws-4-k15-w10:row:mpf-scalar-scalar-ksw2}{R24}\shortstack[l]{{\color{gray}R24}} & \begin{tabular}[t]{@{}l@{}}\rule{0pt}{12pt}MPF / Scalar\\Scalar-DP / KSW2\end{tabular} & \hspace*{1pt}\begin{tabular}[t]{@{}c@{}}\rule{0pt}{12pt}15s\\[-.2pt]{\fontsize{5.8}{6.4}\selectfont \textcolor{ForestGreen}{$ \blacktriangle\,2.4\times $}}\\[-.2pt]12.5s\\[-.2pt]{\fontsize{5.8}{6.4}\selectfont \textcolor{ForestGreen}{$ \blacktriangle\,2.8\times $}}\end{tabular} & \hspace*{1pt}\begin{tabular}[t]{@{}c@{}}\rule{0pt}{12pt}10.6s\\[-.2pt]{\fontsize{5.8}{6.4}\selectfont \textcolor{ForestGreen}{$ \blacktriangle\,3.6\times $}}\\[-.2pt]8.2s\\[-.2pt]{\fontsize{5.8}{6.4}\selectfont \textcolor{ForestGreen}{$ \blacktriangle\,4.4\times $}}\end{tabular} & \hspace*{1pt}\begin{tabular}[t]{@{}c@{}}\rule{0pt}{12pt}15.3s\\[-.2pt]{\fontsize{5.8}{6.4}\selectfont \textcolor{ForestGreen}{$ \blacktriangle\,6.1\times $}}\\[-.2pt]7.3s\\[-.2pt]{\fontsize{5.8}{6.4}\selectfont \textcolor{ForestGreen}{$ \blacktriangle\,12.1\times $}}\end{tabular} & \hspace*{1pt}\begin{tabular}[t]{@{}c@{}}\rule{0pt}{12pt}36.2s\\[-.2pt]{\fontsize{5.8}{6.4}\selectfont \textcolor{ForestGreen}{$ \blacktriangle\,10.8\times $}}\\[-.2pt]32.2s\\[-.2pt]{\fontsize{5.8}{6.4}\selectfont \textcolor{ForestGreen}{$ \blacktriangle\,12.0\times $}}\end{tabular} & \hspace*{1pt}\begin{tabular}[t]{@{}c@{}}\rule{0pt}{12pt}20s\\[-.2pt]{\fontsize{5.8}{6.4}\selectfont \textcolor{ForestGreen}{$ \blacktriangle\,1.6\times $}}\\[-.2pt]15.6s\\[-.2pt]{\fontsize{5.8}{6.4}\selectfont \textcolor{ForestGreen}{$ \blacktriangle\,1.9\times $}}\end{tabular} & \hspace*{1pt}\begin{tabular}[t]{@{}c@{}}\rule{0pt}{12pt}12.4s\\[-.2pt]{\fontsize{5.8}{6.4}\selectfont \textcolor{ForestGreen}{$ \blacktriangle\,2.8\times $}}\\[-.2pt]9.1s\\[-.2pt]{\fontsize{5.8}{6.4}\selectfont \textcolor{ForestGreen}{$ \blacktriangle\,3.5\times $}}\end{tabular} & \hspace*{1pt}\begin{tabular}[t]{@{}c@{}}\rule{0pt}{12pt}25.3s\\[-.2pt]{\fontsize{5.8}{6.4}\selectfont \textcolor{ForestGreen}{$ \blacktriangle\,2.7\times $}}\\[-.2pt]17.1s\\[-.2pt]{\fontsize{5.8}{6.4}\selectfont \textcolor{ForestGreen}{$ \blacktriangle\,3.9\times $}}\end{tabular} & \hspace*{1pt}\begin{tabular}[t]{@{}c@{}}\rule{0pt}{12pt}57.6s\\[-.2pt]{\fontsize{5.8}{6.4}\selectfont \textcolor{ForestGreen}{$ \blacktriangle\,3.1\times $}}\\[-.2pt]53.6s\\[-.2pt]{\fontsize{5.8}{6.4}\selectfont \textcolor{ForestGreen}{$ \blacktriangle\,3.2\times $}}\end{tabular} & \hspace*{1pt}\begin{tabular}[t]{@{}c@{}}\rule{0pt}{12pt}5m\,25.6s\\[-.2pt]{\fontsize{5.8}{6.4}\selectfont \textcolor{BrickRed}{$ \blacktriangledown\,1.9\times $}}\\[-.2pt]5m\,12.1s\\[-.2pt]{\fontsize{5.8}{6.4}\selectfont \textcolor{BrickRed}{$ \blacktriangledown\,1.8\times $}}\end{tabular} & \hspace*{1pt}\begin{tabular}[t]{@{}c@{}}\rule{0pt}{12pt}5m\,5.9s\\[-.2pt]{\fontsize{5.8}{6.4}\selectfont \textcolor{BrickRed}{$ \blacktriangledown\,1.8\times $}}\\[-.2pt]4m\,52.6s\\[-.2pt]{\fontsize{5.8}{6.4}\selectfont \textcolor{BrickRed}{$ \blacktriangledown\,1.7\times $}}\end{tabular} & \textcolor{gray}{\ensuremath{\cdot}} & \textcolor{gray}{\ensuremath{\cdot}} & \hspace*{1pt}\begin{tabular}[t]{@{}c@{}}\rule{0pt}{12pt}11m\,14.1s\\[-.2pt]{\fontsize{5.8}{6.4}\selectfont \textcolor{BrickRed}{$ \blacktriangledown\,6.3\times $}}\\[-.2pt]5m\,5.9s\\[-.2pt]{\fontsize{5.8}{6.4}\selectfont \textcolor{BrickRed}{$ \blacktriangledown\,2.9\times $}}\end{tabular} & \hspace*{1pt}\begin{tabular}[t]{@{}c@{}}\rule{0pt}{12pt}9m\,54.4s\\[-.2pt]{\fontsize{5.8}{6.4}\selectfont \textcolor{BrickRed}{$ \blacktriangledown\,5.5\times $}}\\[-.2pt]5m\,20.4s\\[-.2pt]{\fontsize{5.8}{6.4}\selectfont \textcolor{BrickRed}{$ \blacktriangledown\,3.0\times $}}\end{tabular} & \textcolor{gray}{\ensuremath{\cdot}} & \textcolor{gray}{\ensuremath{\cdot}}\\[3pt]\hline
\BenchmarkSourceLabel{bench:aws-4-k15-w10:row:mps-bounded-scalar-block}{R25}\shortstack[l]{{\color{gray}R25}} & \begin{tabular}[t]{@{}l@{}}\rule{0pt}{12pt}MPS / Bounded\\Scalar-DP / Block\end{tabular} & \hspace*{1pt}\begin{tabular}[t]{@{}c@{}}\rule{0pt}{12pt}10.9s\\[-.2pt]{\fontsize{5.8}{6.4}\selectfont \textcolor{ForestGreen}{$ \blacktriangle\,3.3\times $}}\\[-.2pt]10.3s\\[-.2pt]{\fontsize{5.8}{6.4}\selectfont \textcolor{ForestGreen}{$ \blacktriangle\,3.4\times $}}\end{tabular} & \hspace*{1pt}\begin{tabular}[t]{@{}c@{}}\rule{0pt}{12pt}11.1s\\[-.2pt]{\fontsize{5.8}{6.4}\selectfont \textcolor{ForestGreen}{$ \blacktriangle\,3.4\times $}}\\[-.2pt]10.3s\\[-.2pt]{\fontsize{5.8}{6.4}\selectfont \textcolor{ForestGreen}{$ \blacktriangle\,3.5\times $}}\end{tabular} & \hspace*{1pt}\begin{tabular}[t]{@{}c@{}}\rule{0pt}{12pt}21.5s\\[-.2pt]{\fontsize{5.8}{6.4}\selectfont \textcolor{ForestGreen}{$ \blacktriangle\,4.3\times $}}\\[-.2pt]18.9s\\[-.2pt]{\fontsize{5.8}{6.4}\selectfont \textcolor{ForestGreen}{$ \blacktriangle\,4.7\times $}}\end{tabular} & \textcolor{gray}{\ensuremath{\cdot}} & \hspace*{1pt}\begin{tabular}[t]{@{}c@{}}\rule{0pt}{12pt}8m\,24.2s\\[-.2pt]{\fontsize{5.8}{6.4}\selectfont \textcolor{BrickRed}{$ \blacktriangledown\,16.0\times $}}\\[-.2pt]59.7s\\[-.2pt]{\fontsize{5.8}{6.4}\selectfont \textcolor{BrickRed}{$ \blacktriangledown\,2.0\times $}}\end{tabular} & \hspace*{1pt}\begin{tabular}[t]{@{}c@{}}\rule{0pt}{12pt}6m\,45.3s\\[-.2pt]{\fontsize{5.8}{6.4}\selectfont \textcolor{BrickRed}{$ \blacktriangledown\,11.9\times $}}\\[-.2pt]1m\,3.2s\\[-.2pt]{\fontsize{5.8}{6.4}\selectfont \textcolor{BrickRed}{$ \blacktriangledown\,2.0\times $}}\end{tabular} & \textcolor{gray}{\ensuremath{\cdot}} & \textcolor{gray}{\ensuremath{\cdot}} & \hspace*{1pt}\begin{tabular}[t]{@{}c@{}}\rule{0pt}{12pt}4m\,17.5s\\[-.2pt]{\fontsize{5.8}{6.4}\selectfont \textcolor{BrickRed}{$ \blacktriangledown\,1.5\times $}}\\[-.2pt]4m\,8.2s\\[-.2pt]{\fontsize{5.8}{6.4}\selectfont \textcolor{BrickRed}{$ \blacktriangledown\,1.4\times $}}\end{tabular} & \hspace*{1pt}\begin{tabular}[t]{@{}c@{}}\rule{0pt}{12pt}4m\,7.1s\\[-.2pt]{\fontsize{5.8}{6.4}\selectfont \textcolor{BrickRed}{$ \blacktriangledown\,1.4\times $}}\\[-.2pt]3m\,55.3s\\[-.2pt]{\fontsize{5.8}{6.4}\selectfont \textcolor{BrickRed}{$ \blacktriangledown\,1.4\times $}}\end{tabular} & \textcolor{gray}{\ensuremath{\cdot}} & \textcolor{gray}{\ensuremath{\cdot}} & \hspace*{1pt}\begin{tabular}[t]{@{}c@{}}\rule{0pt}{12pt}10m\,38s\\[-.2pt]{\fontsize{5.8}{6.4}\selectfont \textcolor{BrickRed}{$ \blacktriangledown\,6.0\times $}}\\[-.2pt]4m\,13.2s\\[-.2pt]{\fontsize{5.8}{6.4}\selectfont \textcolor{BrickRed}{$ \blacktriangledown\,2.4\times $}}\end{tabular} & \hspace*{1pt}\begin{tabular}[t]{@{}c@{}}\rule{0pt}{12pt}9m\,41.8s\\[-.2pt]{\fontsize{5.8}{6.4}\selectfont \textcolor{BrickRed}{$ \blacktriangledown\,5.4\times $}}\\[-.2pt]4m\,18.5s\\[-.2pt]{\fontsize{5.8}{6.4}\selectfont \textcolor{BrickRed}{$ \blacktriangledown\,2.5\times $}}\end{tabular} & \textcolor{gray}{\ensuremath{\cdot}} & \textcolor{gray}{\ensuremath{\cdot}}\\[3pt]\hline
\BenchmarkSourceLabel{bench:aws-4-k15-w10:row:mps-simdref-scalar-block}{R26}\shortstack[l]{{\color{gray}R26}} & \begin{tabular}[t]{@{}l@{}}\rule{0pt}{12pt}MPS / SIMD-ref\\Scalar-DP / Block\end{tabular} & \hspace*{1pt}\begin{tabular}[t]{@{}c@{}}\rule{0pt}{12pt}11.9s\\[-.2pt]{\fontsize{5.8}{6.4}\selectfont \textcolor{ForestGreen}{$ \blacktriangle\,3.0\times $}}\\[-.2pt]10.9s\\[-.2pt]{\fontsize{5.8}{6.4}\selectfont \textcolor{ForestGreen}{$ \blacktriangle\,3.2\times $}}\end{tabular} & \hspace*{1pt}\begin{tabular}[t]{@{}c@{}}\rule{0pt}{12pt}10.9s\\[-.2pt]{\fontsize{5.8}{6.4}\selectfont \textcolor{ForestGreen}{$ \blacktriangle\,3.5\times $}}\\[-.2pt]10s\\[-.2pt]{\fontsize{5.8}{6.4}\selectfont \textcolor{ForestGreen}{$ \blacktriangle\,3.6\times $}}\end{tabular} & \textcolor{gray}{\ensuremath{\cdot}} & \textcolor{gray}{\ensuremath{\cdot}} & \hspace*{1pt}\begin{tabular}[t]{@{}c@{}}\rule{0pt}{12pt}6m\,13.5s\\[-.2pt]{\fontsize{5.8}{6.4}\selectfont \textcolor{BrickRed}{$ \blacktriangledown\,11.8\times $}}\\[-.2pt]39.8s\\[-.2pt]{\fontsize{5.8}{6.4}\selectfont \textcolor{BrickRed}{$ \blacktriangledown\,1.3\times $}}\end{tabular} & \hspace*{1pt}\begin{tabular}[t]{@{}c@{}}\rule{0pt}{12pt}5m\,21s\\[-.2pt]{\fontsize{5.8}{6.4}\selectfont \textcolor{BrickRed}{$ \blacktriangledown\,9.4\times $}}\\[-.2pt]35.2s\\[-.2pt]{\fontsize{5.8}{6.4}\selectfont \textcolor{BrickRed}{$ \blacktriangledown\,10\% $}}\end{tabular} & \textcolor{gray}{\ensuremath{\cdot}} & \textcolor{gray}{\ensuremath{\cdot}} & \textcolor{gray}{\ensuremath{\cdot}} & \textcolor{gray}{\ensuremath{\cdot}} & \textcolor{gray}{\ensuremath{\cdot}} & \textcolor{gray}{\ensuremath{\cdot}} & \textcolor{gray}{\ensuremath{\cdot}} & \textcolor{gray}{\ensuremath{\cdot}} & \textcolor{gray}{\ensuremath{\cdot}} & \textcolor{gray}{\ensuremath{\cdot}}\\[3pt]\hline
\BenchmarkSourceLabel{bench:aws-4-k15-w10:row:mps-simdref-rmq-block}{R27}\shortstack[l]{{\color{gray}R27}} & \begin{tabular}[t]{@{}l@{}}\rule{0pt}{12pt}MPS / SIMD-ref\\RMQ / Block\end{tabular} & \hspace*{1pt}\begin{tabular}[t]{@{}c@{}}\rule{0pt}{12pt}12s\\[-.2pt]{\fontsize{5.8}{6.4}\selectfont \textcolor{ForestGreen}{$ \blacktriangle\,3.0\times $}}\\[-.2pt]11s\\[-.2pt]{\fontsize{5.8}{6.4}\selectfont \textcolor{ForestGreen}{$ \blacktriangle\,3.2\times $}}\end{tabular} & \hspace*{1pt}\begin{tabular}[t]{@{}c@{}}\rule{0pt}{12pt}11s\\[-.2pt]{\fontsize{5.8}{6.4}\selectfont \textcolor{ForestGreen}{$ \blacktriangle\,3.5\times $}}\\[-.2pt]10s\\[-.2pt]{\fontsize{5.8}{6.4}\selectfont \textcolor{ForestGreen}{$ \blacktriangle\,3.6\times $}}\end{tabular} & \textcolor{gray}{\ensuremath{\cdot}} & \textcolor{gray}{\ensuremath{\cdot}} & \hspace*{1pt}\begin{tabular}[t]{@{}c@{}}\rule{0pt}{12pt}6m\,15.3s\\[-.2pt]{\fontsize{5.8}{6.4}\selectfont \textcolor{BrickRed}{$ \blacktriangledown\,11.9\times $}}\\[-.2pt]41.2s\\[-.2pt]{\fontsize{5.8}{6.4}\selectfont \textcolor{BrickRed}{$ \blacktriangledown\,1.4\times $}}\end{tabular} & \hspace*{1pt}\begin{tabular}[t]{@{}c@{}}\rule{0pt}{12pt}5m\,22.7s\\[-.2pt]{\fontsize{5.8}{6.4}\selectfont \textcolor{BrickRed}{$ \blacktriangledown\,9.5\times $}}\\[-.2pt]37s\\[-.2pt]{\fontsize{5.8}{6.4}\selectfont \textcolor{BrickRed}{$ \blacktriangledown\,15\% $}}\end{tabular} & \textcolor{gray}{\ensuremath{\cdot}} & \textcolor{gray}{\ensuremath{\cdot}} & \textcolor{gray}{\ensuremath{\cdot}} & \textcolor{gray}{\ensuremath{\cdot}} & \textcolor{gray}{\ensuremath{\cdot}} & \textcolor{gray}{\ensuremath{\cdot}} & \textcolor{gray}{\ensuremath{\cdot}} & \textcolor{gray}{\ensuremath{\cdot}} & \textcolor{gray}{\ensuremath{\cdot}} & \textcolor{gray}{\ensuremath{\cdot}}\\[3pt]\hline
\BenchmarkSourceLabel{bench:aws-4-k15-w10:row:mps-bounded-rmq-block}{R28}\shortstack[l]{{\color{gray}R28}} & \begin{tabular}[t]{@{}l@{}}\rule{0pt}{12pt}MPS / Bounded\\RMQ / Block\end{tabular} & \hspace*{1pt}\begin{tabular}[t]{@{}c@{}}\rule{0pt}{12pt}11s\\[-.2pt]{\fontsize{5.8}{6.4}\selectfont \textcolor{ForestGreen}{$ \blacktriangle\,3.3\times $}}\\[-.2pt]10.3s\\[-.2pt]{\fontsize{5.8}{6.4}\selectfont \textcolor{ForestGreen}{$ \blacktriangle\,3.4\times $}}\end{tabular} & \hspace*{1pt}\begin{tabular}[t]{@{}c@{}}\rule{0pt}{12pt}11.5s\\[-.2pt]{\fontsize{5.8}{6.4}\selectfont \textcolor{ForestGreen}{$ \blacktriangle\,3.3\times $}}\\[-.2pt]10.7s\\[-.2pt]{\fontsize{5.8}{6.4}\selectfont \textcolor{ForestGreen}{$ \blacktriangle\,3.4\times $}}\end{tabular} & \hspace*{1pt}\begin{tabular}[t]{@{}c@{}}\rule{0pt}{12pt}21.8s\\[-.2pt]{\fontsize{5.8}{6.4}\selectfont \textcolor{ForestGreen}{$ \blacktriangle\,4.2\times $}}\\[-.2pt]18.8s\\[-.2pt]{\fontsize{5.8}{6.4}\selectfont \textcolor{ForestGreen}{$ \blacktriangle\,4.7\times $}}\end{tabular} & \textcolor{gray}{\ensuremath{\cdot}} & \hspace*{1pt}\begin{tabular}[t]{@{}c@{}}\rule{0pt}{12pt}8m\,23.3s\\[-.2pt]{\fontsize{5.8}{6.4}\selectfont \textcolor{BrickRed}{$ \blacktriangledown\,15.9\times $}}\\[-.2pt]59.1s\\[-.2pt]{\fontsize{5.8}{6.4}\selectfont \textcolor{BrickRed}{$ \blacktriangledown\,1.9\times $}}\end{tabular} & \hspace*{1pt}\begin{tabular}[t]{@{}c@{}}\rule{0pt}{12pt}7m\,3.6s\\[-.2pt]{\fontsize{5.8}{6.4}\selectfont \textcolor{BrickRed}{$ \blacktriangledown\,12.4\times $}}\\[-.2pt]1m\,6.2s\\[-.2pt]{\fontsize{5.8}{6.4}\selectfont \textcolor{BrickRed}{$ \blacktriangledown\,2.1\times $}}\end{tabular} & \textcolor{gray}{\ensuremath{\cdot}} & \textcolor{gray}{\ensuremath{\cdot}} & \hspace*{1pt}\begin{tabular}[t]{@{}c@{}}\rule{0pt}{12pt}4m\,10.2s\\[-.2pt]{\fontsize{5.8}{6.4}\selectfont \textcolor{BrickRed}{$ \blacktriangledown\,1.4\times $}}\\[-.2pt]4m\,1s\\[-.2pt]{\fontsize{5.8}{6.4}\selectfont \textcolor{BrickRed}{$ \blacktriangledown\,1.4\times $}}\end{tabular} & \hspace*{1pt}\begin{tabular}[t]{@{}c@{}}\rule{0pt}{12pt}4m\,11.6s\\[-.2pt]{\fontsize{5.8}{6.4}\selectfont \textcolor{BrickRed}{$ \blacktriangledown\,1.5\times $}}\\[-.2pt]3m\,59.9s\\[-.2pt]{\fontsize{5.8}{6.4}\selectfont \textcolor{BrickRed}{$ \blacktriangledown\,1.4\times $}}\end{tabular} & \hspace*{1pt}\begin{tabular}[t]{@{}c@{}}\rule{0pt}{12pt}6m\,51.1s\\[-.2pt]{\fontsize{5.8}{6.4}\selectfont \textcolor{BrickRed}{$ \blacktriangledown\,22\% $}}\\[-.2pt]6m\,31.1s\\[-.2pt]{\fontsize{5.8}{6.4}\selectfont \textcolor{BrickRed}{$ \blacktriangledown\,18\% $}}\end{tabular} & \textcolor{gray}{\ensuremath{\cdot}} & \hspace*{1pt}\begin{tabular}[t]{@{}c@{}}\rule{0pt}{12pt}10m\,27.7s\\[-.2pt]{\fontsize{5.8}{6.4}\selectfont \textcolor{BrickRed}{$ \blacktriangledown\,5.9\times $}}\\[-.2pt]4m\,6s\\[-.2pt]{\fontsize{5.8}{6.4}\selectfont \textcolor{BrickRed}{$ \blacktriangledown\,2.3\times $}}\end{tabular} & \hspace*{1pt}\begin{tabular}[t]{@{}c@{}}\rule{0pt}{12pt}9m\,42.6s\\[-.2pt]{\fontsize{5.8}{6.4}\selectfont \textcolor{BrickRed}{$ \blacktriangledown\,5.4\times $}}\\[-.2pt]4m\,14.7s\\[-.2pt]{\fontsize{5.8}{6.4}\selectfont \textcolor{BrickRed}{$ \blacktriangledown\,2.4\times $}}\end{tabular} & \textcolor{gray}{\ensuremath{\cdot}} & \textcolor{gray}{\ensuremath{\cdot}}\\[3pt]\hline
\BenchmarkSourceLabel{bench:aws-4-k15-w10:row:mps-bounded-simd-block}{R29}\shortstack[l]{{\color{gray}R29}} & \begin{tabular}[t]{@{}l@{}}\rule{0pt}{12pt}MPS / Bounded\\SIMD-DP / Block\end{tabular} & \hspace*{1pt}\begin{tabular}[t]{@{}c@{}}\rule{0pt}{12pt}11.1s\\[-.2pt]{\fontsize{5.8}{6.4}\selectfont \textcolor{ForestGreen}{$ \blacktriangle\,3.3\times $}}\\[-.2pt]10.5s\\[-.2pt]{\fontsize{5.8}{6.4}\selectfont \textcolor{ForestGreen}{$ \blacktriangle\,3.3\times $}}\end{tabular} & \hspace*{1pt}\begin{tabular}[t]{@{}c@{}}\rule{0pt}{12pt}11.4s\\[-.2pt]{\fontsize{5.8}{6.4}\selectfont \textcolor{ForestGreen}{$ \blacktriangle\,3.4\times $}}\\[-.2pt]10.5s\\[-.2pt]{\fontsize{5.8}{6.4}\selectfont \textcolor{ForestGreen}{$ \blacktriangle\,3.4\times $}}\end{tabular} & \hspace*{1pt}\begin{tabular}[t]{@{}c@{}}\rule{0pt}{12pt}30.1s\\[-.2pt]{\fontsize{5.8}{6.4}\selectfont \textcolor{ForestGreen}{$ \blacktriangle\,3.1\times $}}\\[-.2pt]27.5s\\[-.2pt]{\fontsize{5.8}{6.4}\selectfont \textcolor{ForestGreen}{$ \blacktriangle\,3.2\times $}}\end{tabular} & \textcolor{gray}{\ensuremath{\cdot}} & \hspace*{1pt}\begin{tabular}[t]{@{}c@{}}\rule{0pt}{12pt}8m\,12s\\[-.2pt]{\fontsize{5.8}{6.4}\selectfont \textcolor{BrickRed}{$ \blacktriangledown\,15.6\times $}}\\[-.2pt]1m\,0.9s\\[-.2pt]{\fontsize{5.8}{6.4}\selectfont \textcolor{BrickRed}{$ \blacktriangledown\,2.0\times $}}\end{tabular} & \hspace*{1pt}\begin{tabular}[t]{@{}c@{}}\rule{0pt}{12pt}6m\,45.6s\\[-.2pt]{\fontsize{5.8}{6.4}\selectfont \textcolor{BrickRed}{$ \blacktriangledown\,11.9\times $}}\\[-.2pt]1m\,1.4s\\[-.2pt]{\fontsize{5.8}{6.4}\selectfont \textcolor{BrickRed}{$ \blacktriangledown\,1.9\times $}}\end{tabular} & \textcolor{gray}{\ensuremath{\cdot}} & \textcolor{gray}{\ensuremath{\cdot}} & \hspace*{1pt}\begin{tabular}[t]{@{}c@{}}\rule{0pt}{12pt}4m\,19.3s\\[-.2pt]{\fontsize{5.8}{6.4}\selectfont \textcolor{BrickRed}{$ \blacktriangledown\,1.5\times $}}\\[-.2pt]4m\,10.3s\\[-.2pt]{\fontsize{5.8}{6.4}\selectfont \textcolor{BrickRed}{$ \blacktriangledown\,1.4\times $}}\end{tabular} & \hspace*{1pt}\begin{tabular}[t]{@{}c@{}}\rule{0pt}{12pt}4m\,29.2s\\[-.2pt]{\fontsize{5.8}{6.4}\selectfont \textcolor{BrickRed}{$ \blacktriangledown\,1.6\times $}}\\[-.2pt]4m\,17.1s\\[-.2pt]{\fontsize{5.8}{6.4}\selectfont \textcolor{BrickRed}{$ \blacktriangledown\,1.5\times $}}\end{tabular} & \textcolor{gray}{\ensuremath{\cdot}} & \textcolor{gray}{\ensuremath{\cdot}} & \hspace*{1pt}\begin{tabular}[t]{@{}c@{}}\rule{0pt}{12pt}10m\,25.7s\\[-.2pt]{\fontsize{5.8}{6.4}\selectfont \textcolor{BrickRed}{$ \blacktriangledown\,5.8\times $}}\\[-.2pt]4m\,1.1s\\[-.2pt]{\fontsize{5.8}{6.4}\selectfont \textcolor{BrickRed}{$ \blacktriangledown\,2.3\times $}}\end{tabular} & \hspace*{1pt}\begin{tabular}[t]{@{}c@{}}\rule{0pt}{12pt}9m\,26.4s\\[-.2pt]{\fontsize{5.8}{6.4}\selectfont \textcolor{BrickRed}{$ \blacktriangledown\,5.3\times $}}\\[-.2pt]4m\,6.2s\\[-.2pt]{\fontsize{5.8}{6.4}\selectfont \textcolor{BrickRed}{$ \blacktriangledown\,2.3\times $}}\end{tabular} & \textcolor{gray}{\ensuremath{\cdot}} & \textcolor{gray}{\ensuremath{\cdot}}\\[3pt]\hline
\end{tabular}\par}
\smallskip{\footnotesize\noindent Each cell: A and B times, $\cdot$: no measurement. ---\textsuperscript{g}: GPU unavailable. n/a: no comparison.\par
\noindent Blue: one lowest A and one lowest B measurement for each column for each configuration.\par
\noindent Gray: minimap2. Yellow: my preferred MPF / Bounded / SIMD-DP / Block configuration.\par
\noindent MPF/MPS: minimizers per fragment/sequence. SIMD-ref: naive SIMD scan. Bounded: bounded SIMD algorithm.\par
}
\clearpage
\noindent{\large\bfseries Table \thetable: (B): AWS: 4 Gbp, k=15, w=10\hfill 5/7}\par\smallskip
{\fontsize{7.1}{8.1}\selectfont\setlength{\tabcolsep}{1pt}\setlength{\fboxsep}{2pt}\renewcommand{\arraystretch}{1.12}
\noindent\begin{tabular}{@{}p{14mm}p{27mm}|p{13.55mm}|p{13.55mm}|p{13.55mm}|p{13.55mm}|p{13.55mm}|p{13.55mm}|p{13.55mm}|p{13.55mm}|p{13.55mm}|p{13.55mm}|p{13.55mm}|p{13.55mm}|p{13.55mm}|p{13.55mm}|p{13.55mm}|p{13.55mm}|@{}}\hline
Engine & Execution & \multicolumn{8}{c|}{\textbf{Index}} & \multicolumn{8}{c|}{\textbf{No index}}\\
 &  & \multicolumn{8}{c|}{\textbf{4 Gbp}} & \multicolumn{8}{c|}{\textbf{4 Gbp}}\\
 &  & \multicolumn{8}{c|}{\textbf{k=15, w=10}} & \multicolumn{8}{c|}{\textbf{k=15, w=10}}\\
 &  & \multicolumn{4}{c|}{\textbf{DNA4}} & \multicolumn{4}{c|}{\textbf{DNA5}} & \multicolumn{4}{c|}{\textbf{DNA4}} & \multicolumn{4}{c|}{\textbf{DNA5}}\\
 &  & \multicolumn{1}{c|}{\textbf{M7A}} & \multicolumn{1}{c|}{\textbf{M7I}} & \multicolumn{1}{c|}{\textbf{M4}} & \multicolumn{1}{c|}{\textbf{C7G}} & \multicolumn{1}{c|}{\textbf{M7A}} & \multicolumn{1}{c|}{\textbf{M7I}} & \multicolumn{1}{c|}{\textbf{M4}} & \multicolumn{1}{c|}{\textbf{C7G}} & \multicolumn{1}{c|}{\textbf{M7A}} & \multicolumn{1}{c|}{\textbf{M7I}} & \multicolumn{1}{c|}{\textbf{M4}} & \multicolumn{1}{c|}{\textbf{C7G}} & \multicolumn{1}{c|}{\textbf{M7A}} & \multicolumn{1}{c|}{\textbf{M7I}} & \multicolumn{1}{c|}{\textbf{M4}} & \multicolumn{1}{c|}{\textbf{C7G}}\\
 & & {\color{gray}C1} & {\color{gray}C2} & {\color{gray}C3} & {\color{gray}C4} & {\color{gray}C5} & {\color{gray}C6} & {\color{gray}C7} & {\color{gray}C8} & {\color{gray}C9} & {\color{gray}C10} & {\color{gray}C11} & {\color{gray}C12} & {\color{gray}C13} & {\color{gray}C14} & {\color{gray}C15} & {\color{gray}C16}\\
\hline
\rowcolor{BenchBaseline}[0.5pt][0.5pt]\shortstack[l]{{\color{gray}R1}\\\textbf{minimap2}} & \textbf{Baseline} & \hspace*{1pt}\begin{tabular}[t]{@{}c@{}}\rule{0pt}{12pt}36.2s\\[-.2pt]{\fontsize{5.8}{6.4}\selectfont \textcolor{BrickRed}{$ \blacktriangledown\,6.9\times $}}\\[-.2pt]35s\\[-.2pt]{\fontsize{5.8}{6.4}\selectfont \textcolor{BrickRed}{$ \blacktriangledown\,8.0\times $}}\end{tabular} & \hspace*{1pt}\begin{tabular}[t]{@{}c@{}}\rule{0pt}{12pt}38.1s\\[-.2pt]{\fontsize{5.8}{6.4}\selectfont \textcolor{BrickRed}{$ \blacktriangledown\,7.8\times $}}\\[-.2pt]36s\\[-.2pt]{\fontsize{5.8}{6.4}\selectfont \textcolor{BrickRed}{$ \blacktriangledown\,9.3\times $}}\end{tabular} & \hspace*{1pt}\begin{tabular}[t]{@{}c@{}}\rule{0pt}{12pt}1m\,32.4s\\[-.2pt]{\fontsize{5.8}{6.4}\selectfont \textcolor{BrickRed}{$ \blacktriangledown\,10.7\times $}}\\[-.2pt]1m\,28.5s\\[-.2pt]{\fontsize{5.8}{6.4}\selectfont \textcolor{BrickRed}{$ \blacktriangledown\,13.5\times $}}\end{tabular} & \hspace*{1pt}\begin{tabular}[t]{@{}c@{}}\rule{0pt}{12pt}6m\,31.5s\\[-.2pt]{\fontsize{5.8}{6.4}\selectfont \textcolor{BrickRed}{$ \blacktriangledown\,11.5\times $}}\\[-.2pt]6m\,26.3s\\[-.2pt]{\fontsize{5.8}{6.4}\selectfont \textcolor{BrickRed}{$ \blacktriangledown\,12.7\times $}}\end{tabular} & \hspace*{1pt}\begin{tabular}[t]{@{}c@{}}\rule{0pt}{12pt}31.6s\\[-.2pt]{\fontsize{5.8}{6.4}\selectfont \textcolor{BrickRed}{$ \blacktriangledown\,2.1\times $}}\\[-.2pt]30.3s\\[-.2pt]{\fontsize{5.8}{6.4}\selectfont \textcolor{BrickRed}{$ \blacktriangledown\,3.0\times $}}\end{tabular} & \hspace*{1pt}\begin{tabular}[t]{@{}c@{}}\rule{0pt}{12pt}34.1s\\[-.2pt]{\fontsize{5.8}{6.4}\selectfont \textcolor{BrickRed}{$ \blacktriangledown\,2.9\times $}}\\[-.2pt]32s\\[-.2pt]{\fontsize{5.8}{6.4}\selectfont \textcolor{BrickRed}{$ \blacktriangledown\,3.7\times $}}\end{tabular} & \hspace*{1pt}\begin{tabular}[t]{@{}c@{}}\rule{0pt}{12pt}1m\,9.5s\\[-.2pt]{\fontsize{5.8}{6.4}\selectfont \textcolor{BrickRed}{$ \blacktriangledown\,3.0\times $}}\\[-.2pt]1m\,6.1s\\[-.2pt]{\fontsize{5.8}{6.4}\selectfont \textcolor{BrickRed}{$ \blacktriangledown\,4.3\times $}}\end{tabular} & \hspace*{1pt}\begin{tabular}[t]{@{}c@{}}\rule{0pt}{12pt}3m\,1s\\[-.2pt]{\fontsize{5.8}{6.4}\selectfont \textcolor{BrickRed}{$ \blacktriangledown\,7.9\times $}}\\[-.2pt]2m\,51.5s\\[-.2pt]{\fontsize{5.8}{6.4}\selectfont \textcolor{BrickRed}{$ \blacktriangledown\,8.6\times $}}\end{tabular} & \hspace*{1pt}\begin{tabular}[t]{@{}c@{}}\rule{0pt}{12pt}2m\,55.5s\\[-.2pt]{\fontsize{5.8}{6.4}\selectfont \textcolor{ForestGreen}{$ \blacktriangle\,20\% $}}\\[-.2pt]2m\,54s\\[-.2pt]{\fontsize{5.8}{6.4}\selectfont \textcolor{ForestGreen}{$ \blacktriangle\,17\% $}}\end{tabular} & \hspace*{1pt}\begin{tabular}[t]{@{}c@{}}\rule{0pt}{12pt}2m\,50.4s\\[-.2pt]{\fontsize{5.8}{6.4}\selectfont \textcolor{ForestGreen}{$ \blacktriangle\,23\% $}}\\[-.2pt]2m\,48.2s\\[-.2pt]{\fontsize{5.8}{6.4}\selectfont \textcolor{ForestGreen}{$ \blacktriangle\,20\% $}}\end{tabular} & \hspace*{1pt}\begin{tabular}[t]{@{}c@{}}\rule{0pt}{12pt}5m\,35.7s\\[-.2pt]{\fontsize{5.8}{6.4}\selectfont \textcolor{ForestGreen}{$ \blacktriangle\,18\% $}}\\[-.2pt]5m\,31.9s\\[-.2pt]{\fontsize{5.8}{6.4}\selectfont \textcolor{ForestGreen}{$ \blacktriangle\,14\% $}}\end{tabular} & \hspace*{1pt}\begin{tabular}[t]{@{}c@{}}\rule{0pt}{12pt}44m\,20.5s\\[-.2pt]{\fontsize{5.8}{6.4}\selectfont \textcolor{gray}{n/a}}\\[-.2pt]44m\,15.6s\\[-.2pt]{\fontsize{5.8}{6.4}\selectfont \textcolor{gray}{n/a}}\end{tabular} & \hspace*{1pt}\begin{tabular}[t]{@{}c@{}}\rule{0pt}{12pt}1m\,47.1s\\[-.2pt]{\fontsize{5.8}{6.4}\selectfont \textcolor{ForestGreen}{$ \blacktriangle\,5.5\times $}}\\[-.2pt]1m\,45.8s\\[-.2pt]{\fontsize{5.8}{6.4}\selectfont \textcolor{ForestGreen}{$ \blacktriangle\,2.1\times $}}\end{tabular} & \hspace*{1pt}\begin{tabular}[t]{@{}c@{}}\rule{0pt}{12pt}1m\,47.2s\\[-.2pt]{\fontsize{5.8}{6.4}\selectfont \textcolor{ForestGreen}{$ \blacktriangle\,4.6\times $}}\\[-.2pt]1m\,45.1s\\[-.2pt]{\fontsize{5.8}{6.4}\selectfont \textcolor{ForestGreen}{$ \blacktriangle\,2.1\times $}}\end{tabular} & \hspace*{1pt}\begin{tabular}[t]{@{}c@{}}\rule{0pt}{12pt}3m\,40.3s\\[-.2pt]{\fontsize{5.8}{6.4}\selectfont \textcolor{ForestGreen}{$ \blacktriangle\,4.2\times $}}\\[-.2pt]3m\,36.6s\\[-.2pt]{\fontsize{5.8}{6.4}\selectfont \textcolor{ForestGreen}{$ \blacktriangle\,1.7\times $}}\end{tabular} & \hspace*{1pt}\begin{tabular}[t]{@{}c@{}}\rule{0pt}{12pt}29m\,43.4s\\[-.2pt]{\fontsize{5.8}{6.4}\selectfont \textcolor{gray}{n/a}}\\[-.2pt]29m\,38.5s\\[-.2pt]{\fontsize{5.8}{6.4}\selectfont \textcolor{gray}{n/a}}\end{tabular}\\[3pt]\hline
\BenchmarkSourceLabel{bench:aws-4-k15-w10:row:mps-twostack-simd-block}{R30}\shortstack[l]{{\color{gray}R30}\\\textbf{BioseqDB}} & \begin{tabular}[t]{@{}l@{}}\rule{0pt}{12pt}MPS / Two-stack\\SIMD-DP / Block\end{tabular} & \hspace*{1pt}\begin{tabular}[t]{@{}c@{}}\rule{0pt}{12pt}11.5s\\[-.2pt]{\fontsize{5.8}{6.4}\selectfont \textcolor{ForestGreen}{$ \blacktriangle\,3.1\times $}}\\[-.2pt]10.9s\\[-.2pt]{\fontsize{5.8}{6.4}\selectfont \textcolor{ForestGreen}{$ \blacktriangle\,3.2\times $}}\end{tabular} & \hspace*{1pt}\begin{tabular}[t]{@{}c@{}}\rule{0pt}{12pt}11.2s\\[-.2pt]{\fontsize{5.8}{6.4}\selectfont \textcolor{ForestGreen}{$ \blacktriangle\,3.4\times $}}\\[-.2pt]10.4s\\[-.2pt]{\fontsize{5.8}{6.4}\selectfont \textcolor{ForestGreen}{$ \blacktriangle\,3.5\times $}}\end{tabular} & \textcolor{gray}{\ensuremath{\cdot}} & \textcolor{gray}{\ensuremath{\cdot}} & \hspace*{1pt}\begin{tabular}[t]{@{}c@{}}\rule{0pt}{12pt}8m\,2s\\[-.2pt]{\fontsize{5.8}{6.4}\selectfont \textcolor{BrickRed}{$ \blacktriangledown\,15.3\times $}}\\[-.2pt]48.3s\\[-.2pt]{\fontsize{5.8}{6.4}\selectfont \textcolor{BrickRed}{$ \blacktriangledown\,1.6\times $}}\end{tabular} & \hspace*{1pt}\begin{tabular}[t]{@{}c@{}}\rule{0pt}{12pt}7m\,1.6s\\[-.2pt]{\fontsize{5.8}{6.4}\selectfont \textcolor{BrickRed}{$ \blacktriangledown\,12.4\times $}}\\[-.2pt]48s\\[-.2pt]{\fontsize{5.8}{6.4}\selectfont \textcolor{BrickRed}{$ \blacktriangledown\,1.5\times $}}\end{tabular} & \textcolor{gray}{\ensuremath{\cdot}} & \textcolor{gray}{\ensuremath{\cdot}} & \hspace*{1pt}\begin{tabular}[t]{@{}c@{}}\rule{0pt}{12pt}4m\,39.5s\\[-.2pt]{\fontsize{5.8}{6.4}\selectfont \textcolor{BrickRed}{$ \blacktriangledown\,1.6\times $}}\\[-.2pt]4m\,30.1s\\[-.2pt]{\fontsize{5.8}{6.4}\selectfont \textcolor{BrickRed}{$ \blacktriangledown\,1.6\times $}}\end{tabular} & \hspace*{1pt}\begin{tabular}[t]{@{}c@{}}\rule{0pt}{12pt}4m\,24.7s\\[-.2pt]{\fontsize{5.8}{6.4}\selectfont \textcolor{BrickRed}{$ \blacktriangledown\,1.6\times $}}\\[-.2pt]4m\,13s\\[-.2pt]{\fontsize{5.8}{6.4}\selectfont \textcolor{BrickRed}{$ \blacktriangledown\,1.5\times $}}\end{tabular} & \textcolor{gray}{\ensuremath{\cdot}} & \textcolor{gray}{\ensuremath{\cdot}} & \textcolor{gray}{\ensuremath{\cdot}} & \textcolor{gray}{\ensuremath{\cdot}} & \textcolor{gray}{\ensuremath{\cdot}} & \textcolor{gray}{\ensuremath{\cdot}}\\[3pt]\hline
\BenchmarkSourceLabel{bench:aws-4-k15-w10:row:mps-scalar-simd-block}{R31}\shortstack[l]{{\color{gray}R31}} & \begin{tabular}[t]{@{}l@{}}\rule{0pt}{12pt}MPS / Scalar\\SIMD-DP / Block\end{tabular} & \hspace*{1pt}\begin{tabular}[t]{@{}c@{}}\rule{0pt}{12pt}11.3s\\[-.2pt]{\fontsize{5.8}{6.4}\selectfont \textcolor{ForestGreen}{$ \blacktriangle\,3.2\times $}}\\[-.2pt]10.7s\\[-.2pt]{\fontsize{5.8}{6.4}\selectfont \textcolor{ForestGreen}{$ \blacktriangle\,3.3\times $}}\end{tabular} & \hspace*{1pt}\begin{tabular}[t]{@{}c@{}}\rule{0pt}{12pt}11.4s\\[-.2pt]{\fontsize{5.8}{6.4}\selectfont \textcolor{ForestGreen}{$ \blacktriangle\,3.3\times $}}\\[-.2pt]10.7s\\[-.2pt]{\fontsize{5.8}{6.4}\selectfont \textcolor{ForestGreen}{$ \blacktriangle\,3.4\times $}}\end{tabular} & \textcolor{gray}{\ensuremath{\cdot}} & \textcolor{gray}{\ensuremath{\cdot}} & \hspace*{1pt}\begin{tabular}[t]{@{}c@{}}\rule{0pt}{12pt}8m\,11.6s\\[-.2pt]{\fontsize{5.8}{6.4}\selectfont \textcolor{BrickRed}{$ \blacktriangledown\,15.6\times $}}\\[-.2pt]1m\,1.8s\\[-.2pt]{\fontsize{5.8}{6.4}\selectfont \textcolor{BrickRed}{$ \blacktriangledown\,2.0\times $}}\end{tabular} & \hspace*{1pt}\begin{tabular}[t]{@{}c@{}}\rule{0pt}{12pt}6m\,34.1s\\[-.2pt]{\fontsize{5.8}{6.4}\selectfont \textcolor{BrickRed}{$ \blacktriangledown\,11.6\times $}}\\[-.2pt]58.8s\\[-.2pt]{\fontsize{5.8}{6.4}\selectfont \textcolor{BrickRed}{$ \blacktriangledown\,1.8\times $}}\end{tabular} & \textcolor{gray}{\ensuremath{\cdot}} & \textcolor{gray}{\ensuremath{\cdot}} & \hspace*{1pt}\begin{tabular}[t]{@{}c@{}}\rule{0pt}{12pt}5m\,48.1s\\[-.2pt]{\fontsize{5.8}{6.4}\selectfont \textcolor{BrickRed}{$ \blacktriangledown\,2.0\times $}}\\[-.2pt]5m\,36.4s\\[-.2pt]{\fontsize{5.8}{6.4}\selectfont \textcolor{BrickRed}{$ \blacktriangledown\,1.9\times $}}\end{tabular} & \hspace*{1pt}\begin{tabular}[t]{@{}c@{}}\rule{0pt}{12pt}5m\,51.7s\\[-.2pt]{\fontsize{5.8}{6.4}\selectfont \textcolor{BrickRed}{$ \blacktriangledown\,2.1\times $}}\\[-.2pt]5m\,39.9s\\[-.2pt]{\fontsize{5.8}{6.4}\selectfont \textcolor{BrickRed}{$ \blacktriangledown\,2.0\times $}}\end{tabular} & \textcolor{gray}{\ensuremath{\cdot}} & \textcolor{gray}{\ensuremath{\cdot}} & \hspace*{1pt}\begin{tabular}[t]{@{}c@{}}\rule{0pt}{12pt}11m\,45.1s\\[-.2pt]{\fontsize{5.8}{6.4}\selectfont \textcolor{BrickRed}{$ \blacktriangledown\,6.6\times $}}\\[-.2pt]5m\,20.1s\\[-.2pt]{\fontsize{5.8}{6.4}\selectfont \textcolor{BrickRed}{$ \blacktriangledown\,3.0\times $}}\end{tabular} & \hspace*{1pt}\begin{tabular}[t]{@{}c@{}}\rule{0pt}{12pt}10m\,10.3s\\[-.2pt]{\fontsize{5.8}{6.4}\selectfont \textcolor{BrickRed}{$ \blacktriangledown\,5.7\times $}}\\[-.2pt]5m\,30.2s\\[-.2pt]{\fontsize{5.8}{6.4}\selectfont \textcolor{BrickRed}{$ \blacktriangledown\,3.1\times $}}\end{tabular} & \textcolor{gray}{\ensuremath{\cdot}} & \textcolor{gray}{\ensuremath{\cdot}}\\[3pt]\hline
\BenchmarkSourceLabel{bench:aws-4-k15-w10:row:mps-twostack-rmq-block}{R32}\shortstack[l]{{\color{gray}R32}} & \begin{tabular}[t]{@{}l@{}}\rule{0pt}{12pt}MPS / Two-stack\\RMQ / Block\end{tabular} & \hspace*{1pt}\begin{tabular}[t]{@{}c@{}}\rule{0pt}{12pt}11.4s\\[-.2pt]{\fontsize{5.8}{6.4}\selectfont \textcolor{ForestGreen}{$ \blacktriangle\,3.2\times $}}\\[-.2pt]10.9s\\[-.2pt]{\fontsize{5.8}{6.4}\selectfont \textcolor{ForestGreen}{$ \blacktriangle\,3.2\times $}}\end{tabular} & \hspace*{1pt}\begin{tabular}[t]{@{}c@{}}\rule{0pt}{12pt}11.8s\\[-.2pt]{\fontsize{5.8}{6.4}\selectfont \textcolor{ForestGreen}{$ \blacktriangle\,3.2\times $}}\\[-.2pt]10.9s\\[-.2pt]{\fontsize{5.8}{6.4}\selectfont \textcolor{ForestGreen}{$ \blacktriangle\,3.3\times $}}\end{tabular} & \textcolor{gray}{\ensuremath{\cdot}} & \textcolor{gray}{\ensuremath{\cdot}} & \hspace*{1pt}\begin{tabular}[t]{@{}c@{}}\rule{0pt}{12pt}8m\,8.4s\\[-.2pt]{\fontsize{5.8}{6.4}\selectfont \textcolor{BrickRed}{$ \blacktriangledown\,15.5\times $}}\\[-.2pt]45.9s\\[-.2pt]{\fontsize{5.8}{6.4}\selectfont \textcolor{BrickRed}{$ \blacktriangledown\,1.5\times $}}\end{tabular} & \hspace*{1pt}\begin{tabular}[t]{@{}c@{}}\rule{0pt}{12pt}7m\,22.7s\\[-.2pt]{\fontsize{5.8}{6.4}\selectfont \textcolor{BrickRed}{$ \blacktriangledown\,13.0\times $}}\\[-.2pt]47.8s\\[-.2pt]{\fontsize{5.8}{6.4}\selectfont \textcolor{BrickRed}{$ \blacktriangledown\,1.5\times $}}\end{tabular} & \textcolor{gray}{\ensuremath{\cdot}} & \textcolor{gray}{\ensuremath{\cdot}} & \hspace*{1pt}\begin{tabular}[t]{@{}c@{}}\rule{0pt}{12pt}4m\,44.1s\\[-.2pt]{\fontsize{5.8}{6.4}\selectfont \textcolor{BrickRed}{$ \blacktriangledown\,1.6\times $}}\\[-.2pt]4m\,35.1s\\[-.2pt]{\fontsize{5.8}{6.4}\selectfont \textcolor{BrickRed}{$ \blacktriangledown\,1.6\times $}}\end{tabular} & \hspace*{1pt}\begin{tabular}[t]{@{}c@{}}\rule{0pt}{12pt}4m\,42.8s\\[-.2pt]{\fontsize{5.8}{6.4}\selectfont \textcolor{BrickRed}{$ \blacktriangledown\,1.7\times $}}\\[-.2pt]4m\,30.9s\\[-.2pt]{\fontsize{5.8}{6.4}\selectfont \textcolor{BrickRed}{$ \blacktriangledown\,1.6\times $}}\end{tabular} & \textcolor{gray}{\ensuremath{\cdot}} & \textcolor{gray}{\ensuremath{\cdot}} & \textcolor{gray}{\ensuremath{\cdot}} & \textcolor{gray}{\ensuremath{\cdot}} & \textcolor{gray}{\ensuremath{\cdot}} & \textcolor{gray}{\ensuremath{\cdot}}\\[3pt]\hline
\BenchmarkSourceLabel{bench:aws-4-k15-w10:row:mps-scalar-rmq-block}{R33}\shortstack[l]{{\color{gray}R33}} & \begin{tabular}[t]{@{}l@{}}\rule{0pt}{12pt}MPS / Scalar\\RMQ / Block\end{tabular} & \hspace*{1pt}\begin{tabular}[t]{@{}c@{}}\rule{0pt}{12pt}11.5s\\[-.2pt]{\fontsize{5.8}{6.4}\selectfont \textcolor{ForestGreen}{$ \blacktriangle\,3.2\times $}}\\[-.2pt]10.9s\\[-.2pt]{\fontsize{5.8}{6.4}\selectfont \textcolor{ForestGreen}{$ \blacktriangle\,3.2\times $}}\end{tabular} & \hspace*{1pt}\begin{tabular}[t]{@{}c@{}}\rule{0pt}{12pt}11.9s\\[-.2pt]{\fontsize{5.8}{6.4}\selectfont \textcolor{ForestGreen}{$ \blacktriangle\,3.2\times $}}\\[-.2pt]11s\\[-.2pt]{\fontsize{5.8}{6.4}\selectfont \textcolor{ForestGreen}{$ \blacktriangle\,3.3\times $}}\end{tabular} & \textcolor{gray}{\ensuremath{\cdot}} & \textcolor{gray}{\ensuremath{\cdot}} & \hspace*{1pt}\begin{tabular}[t]{@{}c@{}}\rule{0pt}{12pt}7m\,41.6s\\[-.2pt]{\fontsize{5.8}{6.4}\selectfont \textcolor{BrickRed}{$ \blacktriangledown\,14.6\times $}}\\[-.2pt]59.8s\\[-.2pt]{\fontsize{5.8}{6.4}\selectfont \textcolor{BrickRed}{$ \blacktriangledown\,2.0\times $}}\end{tabular} & \hspace*{1pt}\begin{tabular}[t]{@{}c@{}}\rule{0pt}{12pt}6m\,41.2s\\[-.2pt]{\fontsize{5.8}{6.4}\selectfont \textcolor{BrickRed}{$ \blacktriangledown\,11.8\times $}}\\[-.2pt]59.4s\\[-.2pt]{\fontsize{5.8}{6.4}\selectfont \textcolor{BrickRed}{$ \blacktriangledown\,1.9\times $}}\end{tabular} & \textcolor{gray}{\ensuremath{\cdot}} & \textcolor{gray}{\ensuremath{\cdot}} & \hspace*{1pt}\begin{tabular}[t]{@{}c@{}}\rule{0pt}{12pt}5m\,40.1s\\[-.2pt]{\fontsize{5.8}{6.4}\selectfont \textcolor{BrickRed}{$ \blacktriangledown\,1.9\times $}}\\[-.2pt]5m\,31s\\[-.2pt]{\fontsize{5.8}{6.4}\selectfont \textcolor{BrickRed}{$ \blacktriangledown\,1.9\times $}}\end{tabular} & \hspace*{1pt}\begin{tabular}[t]{@{}c@{}}\rule{0pt}{12pt}6m\,11.3s\\[-.2pt]{\fontsize{5.8}{6.4}\selectfont \textcolor{BrickRed}{$ \blacktriangledown\,2.2\times $}}\\[-.2pt]5m\,59.5s\\[-.2pt]{\fontsize{5.8}{6.4}\selectfont \textcolor{BrickRed}{$ \blacktriangledown\,2.1\times $}}\end{tabular} & \textcolor{gray}{\ensuremath{\cdot}} & \textcolor{gray}{\ensuremath{\cdot}} & \hspace*{1pt}\begin{tabular}[t]{@{}c@{}}\rule{0pt}{12pt}11m\,59.3s\\[-.2pt]{\fontsize{5.8}{6.4}\selectfont \textcolor{BrickRed}{$ \blacktriangledown\,6.7\times $}}\\[-.2pt]5m\,36.2s\\[-.2pt]{\fontsize{5.8}{6.4}\selectfont \textcolor{BrickRed}{$ \blacktriangledown\,3.2\times $}}\end{tabular} & \textcolor{gray}{\ensuremath{\cdot}} & \textcolor{gray}{\ensuremath{\cdot}} & \textcolor{gray}{\ensuremath{\cdot}}\\[3pt]\hline
\BenchmarkSourceLabel{bench:aws-4-k15-w10:row:mps-twostack-scalar-block}{R34}\shortstack[l]{{\color{gray}R34}} & \begin{tabular}[t]{@{}l@{}}\rule{0pt}{12pt}MPS / Two-stack\\Scalar-DP / Block\end{tabular} & \hspace*{1pt}\begin{tabular}[t]{@{}c@{}}\rule{0pt}{12pt}12s\\[-.2pt]{\fontsize{5.8}{6.4}\selectfont \textcolor{ForestGreen}{$ \blacktriangle\,3.0\times $}}\\[-.2pt]11s\\[-.2pt]{\fontsize{5.8}{6.4}\selectfont \textcolor{ForestGreen}{$ \blacktriangle\,3.2\times $}}\end{tabular} & \hspace*{1pt}\begin{tabular}[t]{@{}c@{}}\rule{0pt}{12pt}11.6s\\[-.2pt]{\fontsize{5.8}{6.4}\selectfont \textcolor{ForestGreen}{$ \blacktriangle\,3.3\times $}}\\[-.2pt]10.6s\\[-.2pt]{\fontsize{5.8}{6.4}\selectfont \textcolor{ForestGreen}{$ \blacktriangle\,3.4\times $}}\end{tabular} & \textcolor{gray}{\ensuremath{\cdot}} & \textcolor{gray}{\ensuremath{\cdot}} & \hspace*{1pt}\begin{tabular}[t]{@{}c@{}}\rule{0pt}{12pt}8m\,0.3s\\[-.2pt]{\fontsize{5.8}{6.4}\selectfont \textcolor{BrickRed}{$ \blacktriangledown\,15.2\times $}}\\[-.2pt]45.7s\\[-.2pt]{\fontsize{5.8}{6.4}\selectfont \textcolor{BrickRed}{$ \blacktriangledown\,1.5\times $}}\end{tabular} & \hspace*{1pt}\begin{tabular}[t]{@{}c@{}}\rule{0pt}{12pt}6m\,34.2s\\[-.2pt]{\fontsize{5.8}{6.4}\selectfont \textcolor{BrickRed}{$ \blacktriangledown\,11.6\times $}}\\[-.2pt]45.6s\\[-.2pt]{\fontsize{5.8}{6.4}\selectfont \textcolor{BrickRed}{$ \blacktriangledown\,1.4\times $}}\end{tabular} & \textcolor{gray}{\ensuremath{\cdot}} & \textcolor{gray}{\ensuremath{\cdot}} & \textcolor{gray}{\ensuremath{\cdot}} & \textcolor{gray}{\ensuremath{\cdot}} & \textcolor{gray}{\ensuremath{\cdot}} & \textcolor{gray}{\ensuremath{\cdot}} & \textcolor{gray}{\ensuremath{\cdot}} & \textcolor{gray}{\ensuremath{\cdot}} & \textcolor{gray}{\ensuremath{\cdot}} & \textcolor{gray}{\ensuremath{\cdot}}\\[3pt]\hline
\BenchmarkSourceLabel{bench:aws-4-k15-w10:row:mps-scalar-scalar-block}{R35}\shortstack[l]{{\color{gray}R35}} & \begin{tabular}[t]{@{}l@{}}\rule{0pt}{12pt}MPS / Scalar\\Scalar-DP / Block\end{tabular} & \hspace*{1pt}\begin{tabular}[t]{@{}c@{}}\rule{0pt}{12pt}15.6s\\[-.2pt]{\fontsize{5.8}{6.4}\selectfont \textcolor{ForestGreen}{$ \blacktriangle\,2.3\times $}}\\[-.2pt]14.7s\\[-.2pt]{\fontsize{5.8}{6.4}\selectfont \textcolor{ForestGreen}{$ \blacktriangle\,2.4\times $}}\end{tabular} & \hspace*{1pt}\begin{tabular}[t]{@{}c@{}}\rule{0pt}{12pt}12.1s\\[-.2pt]{\fontsize{5.8}{6.4}\selectfont \textcolor{ForestGreen}{$ \blacktriangle\,3.1\times $}}\\[-.2pt]11s\\[-.2pt]{\fontsize{5.8}{6.4}\selectfont \textcolor{ForestGreen}{$ \blacktriangle\,3.3\times $}}\end{tabular} & \textcolor{gray}{\ensuremath{\cdot}} & \textcolor{gray}{\ensuremath{\cdot}} & \hspace*{1pt}\begin{tabular}[t]{@{}c@{}}\rule{0pt}{12pt}7m\,42.1s\\[-.2pt]{\fontsize{5.8}{6.4}\selectfont \textcolor{BrickRed}{$ \blacktriangledown\,14.6\times $}}\\[-.2pt]58.9s\\[-.2pt]{\fontsize{5.8}{6.4}\selectfont \textcolor{BrickRed}{$ \blacktriangledown\,1.9\times $}}\end{tabular} & \hspace*{1pt}\begin{tabular}[t]{@{}c@{}}\rule{0pt}{12pt}6m\,19.9s\\[-.2pt]{\fontsize{5.8}{6.4}\selectfont \textcolor{BrickRed}{$ \blacktriangledown\,11.1\times $}}\\[-.2pt]54.2s\\[-.2pt]{\fontsize{5.8}{6.4}\selectfont \textcolor{BrickRed}{$ \blacktriangledown\,1.7\times $}}\end{tabular} & \textcolor{gray}{\ensuremath{\cdot}} & \textcolor{gray}{\ensuremath{\cdot}} & \hspace*{1pt}\begin{tabular}[t]{@{}c@{}}\rule{0pt}{12pt}5m\,51.9s\\[-.2pt]{\fontsize{5.8}{6.4}\selectfont \textcolor{BrickRed}{$ \blacktriangledown\,2.0\times $}}\\[-.2pt]5m\,39.9s\\[-.2pt]{\fontsize{5.8}{6.4}\selectfont \textcolor{BrickRed}{$ \blacktriangledown\,2.0\times $}}\end{tabular} & \textcolor{gray}{\ensuremath{\cdot}} & \textcolor{gray}{\ensuremath{\cdot}} & \textcolor{gray}{\ensuremath{\cdot}} & \hspace*{1pt}\begin{tabular}[t]{@{}c@{}}\rule{0pt}{12pt}11m\,44.1s\\[-.2pt]{\fontsize{5.8}{6.4}\selectfont \textcolor{BrickRed}{$ \blacktriangledown\,6.6\times $}}\\[-.2pt]5m\,20.6s\\[-.2pt]{\fontsize{5.8}{6.4}\selectfont \textcolor{BrickRed}{$ \blacktriangledown\,3.0\times $}}\end{tabular} & \hspace*{1pt}\begin{tabular}[t]{@{}c@{}}\rule{0pt}{12pt}10m\,13.1s\\[-.2pt]{\fontsize{5.8}{6.4}\selectfont \textcolor{BrickRed}{$ \blacktriangledown\,5.7\times $}}\\[-.2pt]5m\,32.9s\\[-.2pt]{\fontsize{5.8}{6.4}\selectfont \textcolor{BrickRed}{$ \blacktriangledown\,3.2\times $}}\end{tabular} & \textcolor{gray}{\ensuremath{\cdot}} & \textcolor{gray}{\ensuremath{\cdot}}\\[3pt]\hline
\BenchmarkSourceLabel{bench:aws-4-k15-w10:row:mps-bounded-simd-ksw2}{R36}\shortstack[l]{{\color{gray}R36}} & \begin{tabular}[t]{@{}l@{}}\rule{0pt}{12pt}MPS / Bounded\\SIMD-DP / KSW2\end{tabular} & \hspace*{1pt}\begin{tabular}[t]{@{}c@{}}\rule{0pt}{12pt}12.3s\\[-.2pt]{\fontsize{5.8}{6.4}\selectfont \textcolor{ForestGreen}{$ \blacktriangle\,3.0\times $}}\\[-.2pt]9.7s\\[-.2pt]{\fontsize{5.8}{6.4}\selectfont \textcolor{ForestGreen}{$ \blacktriangle\,3.6\times $}}\end{tabular} & \hspace*{1pt}\begin{tabular}[t]{@{}c@{}}\rule{0pt}{12pt}12.1s\\[-.2pt]{\fontsize{5.8}{6.4}\selectfont \textcolor{ForestGreen}{$ \blacktriangle\,3.1\times $}}\\[-.2pt]9.1s\\[-.2pt]{\fontsize{5.8}{6.4}\selectfont \textcolor{ForestGreen}{$ \blacktriangle\,4.0\times $}}\end{tabular} & \hspace*{1pt}\begin{tabular}[t]{@{}c@{}}\rule{0pt}{12pt}26.1s\\[-.2pt]{\fontsize{5.8}{6.4}\selectfont \textcolor{ForestGreen}{$ \blacktriangle\,3.5\times $}}\\[-.2pt]16.6s\\[-.2pt]{\fontsize{5.8}{6.4}\selectfont \textcolor{ForestGreen}{$ \blacktriangle\,5.3\times $}}\end{tabular} & \textcolor{gray}{\ensuremath{\cdot}} & \hspace*{1pt}\begin{tabular}[t]{@{}c@{}}\rule{0pt}{12pt}8m\,21.9s\\[-.2pt]{\fontsize{5.8}{6.4}\selectfont \textcolor{BrickRed}{$ \blacktriangledown\,15.9\times $}}\\[-.2pt]57.7s\\[-.2pt]{\fontsize{5.8}{6.4}\selectfont \textcolor{BrickRed}{$ \blacktriangledown\,1.9\times $}}\end{tabular} & \hspace*{1pt}\begin{tabular}[t]{@{}c@{}}\rule{0pt}{12pt}7m\,3s\\[-.2pt]{\fontsize{5.8}{6.4}\selectfont \textcolor{BrickRed}{$ \blacktriangledown\,12.4\times $}}\\[-.2pt]1m\,0.6s\\[-.2pt]{\fontsize{5.8}{6.4}\selectfont \textcolor{BrickRed}{$ \blacktriangledown\,1.9\times $}}\end{tabular} & \textcolor{gray}{\ensuremath{\cdot}} & \textcolor{gray}{\ensuremath{\cdot}} & \hspace*{1pt}\begin{tabular}[t]{@{}c@{}}\rule{0pt}{12pt}4m\,19.7s\\[-.2pt]{\fontsize{5.8}{6.4}\selectfont \textcolor{BrickRed}{$ \blacktriangledown\,1.5\times $}}\\[-.2pt]4m\,8.9s\\[-.2pt]{\fontsize{5.8}{6.4}\selectfont \textcolor{BrickRed}{$ \blacktriangledown\,1.4\times $}}\end{tabular} & \hspace*{1pt}\begin{tabular}[t]{@{}c@{}}\rule{0pt}{12pt}3m\,58.7s\\[-.2pt]{\fontsize{5.8}{6.4}\selectfont \textcolor{BrickRed}{$ \blacktriangledown\,1.4\times $}}\\[-.2pt]3m\,45.4s\\[-.2pt]{\fontsize{5.8}{6.4}\selectfont \textcolor{BrickRed}{$ \blacktriangledown\,1.3\times $}}\end{tabular} & \hspace*{1pt}\begin{tabular}[t]{@{}c@{}}\rule{0pt}{12pt}10m\,17.4s\\[-.2pt]{\fontsize{5.8}{6.4}\selectfont \textcolor{BrickRed}{$ \blacktriangledown\,1.8\times $}}\\[-.2pt]9m\,56s\\[-.2pt]{\fontsize{5.8}{6.4}\selectfont \textcolor{BrickRed}{$ \blacktriangledown\,1.8\times $}}\end{tabular} & \textcolor{gray}{\ensuremath{\cdot}} & \hspace*{1pt}\begin{tabular}[t]{@{}c@{}}\rule{0pt}{12pt}10m\,43.5s\\[-.2pt]{\fontsize{5.8}{6.4}\selectfont \textcolor{BrickRed}{$ \blacktriangledown\,6.0\times $}}\\[-.2pt]4m\,18.3s\\[-.2pt]{\fontsize{5.8}{6.4}\selectfont \textcolor{BrickRed}{$ \blacktriangledown\,2.4\times $}}\end{tabular} & \hspace*{1pt}\begin{tabular}[t]{@{}c@{}}\rule{0pt}{12pt}9m\,31.6s\\[-.2pt]{\fontsize{5.8}{6.4}\selectfont \textcolor{BrickRed}{$ \blacktriangledown\,5.3\times $}}\\[-.2pt]4m\,7.6s\\[-.2pt]{\fontsize{5.8}{6.4}\selectfont \textcolor{BrickRed}{$ \blacktriangledown\,2.4\times $}}\end{tabular} & \textcolor{gray}{\ensuremath{\cdot}} & \textcolor{gray}{\ensuremath{\cdot}}\\[3pt]\hline
\end{tabular}\par}
\smallskip{\footnotesize\noindent Each cell: A and B times, $\cdot$: no measurement. ---\textsuperscript{g}: GPU unavailable. n/a: no comparison.\par
\noindent Blue: one lowest A and one lowest B measurement for each column for each configuration.\par
\noindent Gray: minimap2. Yellow: my preferred MPF / Bounded / SIMD-DP / Block configuration.\par
\noindent MPF/MPS: minimizers per fragment/sequence. SIMD-ref: naive SIMD scan. Bounded: bounded SIMD algorithm.\par
}
\clearpage
\noindent{\large\bfseries Table \thetable: (B): AWS: 4 Gbp, k=15, w=10\hfill 6/7}\par\smallskip
{\fontsize{7.1}{8.1}\selectfont\setlength{\tabcolsep}{1pt}\setlength{\fboxsep}{2pt}\renewcommand{\arraystretch}{1.12}
\noindent\begin{tabular}{@{}p{14mm}p{27mm}|p{13.55mm}|p{13.55mm}|p{13.55mm}|p{13.55mm}|p{13.55mm}|p{13.55mm}|p{13.55mm}|p{13.55mm}|p{13.55mm}|p{13.55mm}|p{13.55mm}|p{13.55mm}|p{13.55mm}|p{13.55mm}|p{13.55mm}|p{13.55mm}|@{}}\hline
Engine & Execution & \multicolumn{8}{c|}{\textbf{Index}} & \multicolumn{8}{c|}{\textbf{No index}}\\
 &  & \multicolumn{8}{c|}{\textbf{4 Gbp}} & \multicolumn{8}{c|}{\textbf{4 Gbp}}\\
 &  & \multicolumn{8}{c|}{\textbf{k=15, w=10}} & \multicolumn{8}{c|}{\textbf{k=15, w=10}}\\
 &  & \multicolumn{4}{c|}{\textbf{DNA4}} & \multicolumn{4}{c|}{\textbf{DNA5}} & \multicolumn{4}{c|}{\textbf{DNA4}} & \multicolumn{4}{c|}{\textbf{DNA5}}\\
 &  & \multicolumn{1}{c|}{\textbf{M7A}} & \multicolumn{1}{c|}{\textbf{M7I}} & \multicolumn{1}{c|}{\textbf{M4}} & \multicolumn{1}{c|}{\textbf{C7G}} & \multicolumn{1}{c|}{\textbf{M7A}} & \multicolumn{1}{c|}{\textbf{M7I}} & \multicolumn{1}{c|}{\textbf{M4}} & \multicolumn{1}{c|}{\textbf{C7G}} & \multicolumn{1}{c|}{\textbf{M7A}} & \multicolumn{1}{c|}{\textbf{M7I}} & \multicolumn{1}{c|}{\textbf{M4}} & \multicolumn{1}{c|}{\textbf{C7G}} & \multicolumn{1}{c|}{\textbf{M7A}} & \multicolumn{1}{c|}{\textbf{M7I}} & \multicolumn{1}{c|}{\textbf{M4}} & \multicolumn{1}{c|}{\textbf{C7G}}\\
 & & {\color{gray}C1} & {\color{gray}C2} & {\color{gray}C3} & {\color{gray}C4} & {\color{gray}C5} & {\color{gray}C6} & {\color{gray}C7} & {\color{gray}C8} & {\color{gray}C9} & {\color{gray}C10} & {\color{gray}C11} & {\color{gray}C12} & {\color{gray}C13} & {\color{gray}C14} & {\color{gray}C15} & {\color{gray}C16}\\
\hline
\rowcolor{BenchBaseline}[0.5pt][0.5pt]\shortstack[l]{{\color{gray}R1}\\\textbf{minimap2}} & \textbf{Baseline} & \hspace*{1pt}\begin{tabular}[t]{@{}c@{}}\rule{0pt}{12pt}36.2s\\[-.2pt]{\fontsize{5.8}{6.4}\selectfont \textcolor{BrickRed}{$ \blacktriangledown\,6.9\times $}}\\[-.2pt]35s\\[-.2pt]{\fontsize{5.8}{6.4}\selectfont \textcolor{BrickRed}{$ \blacktriangledown\,8.0\times $}}\end{tabular} & \hspace*{1pt}\begin{tabular}[t]{@{}c@{}}\rule{0pt}{12pt}38.1s\\[-.2pt]{\fontsize{5.8}{6.4}\selectfont \textcolor{BrickRed}{$ \blacktriangledown\,7.8\times $}}\\[-.2pt]36s\\[-.2pt]{\fontsize{5.8}{6.4}\selectfont \textcolor{BrickRed}{$ \blacktriangledown\,9.3\times $}}\end{tabular} & \hspace*{1pt}\begin{tabular}[t]{@{}c@{}}\rule{0pt}{12pt}1m\,32.4s\\[-.2pt]{\fontsize{5.8}{6.4}\selectfont \textcolor{BrickRed}{$ \blacktriangledown\,10.7\times $}}\\[-.2pt]1m\,28.5s\\[-.2pt]{\fontsize{5.8}{6.4}\selectfont \textcolor{BrickRed}{$ \blacktriangledown\,13.5\times $}}\end{tabular} & \hspace*{1pt}\begin{tabular}[t]{@{}c@{}}\rule{0pt}{12pt}6m\,31.5s\\[-.2pt]{\fontsize{5.8}{6.4}\selectfont \textcolor{BrickRed}{$ \blacktriangledown\,11.5\times $}}\\[-.2pt]6m\,26.3s\\[-.2pt]{\fontsize{5.8}{6.4}\selectfont \textcolor{BrickRed}{$ \blacktriangledown\,12.7\times $}}\end{tabular} & \hspace*{1pt}\begin{tabular}[t]{@{}c@{}}\rule{0pt}{12pt}31.6s\\[-.2pt]{\fontsize{5.8}{6.4}\selectfont \textcolor{BrickRed}{$ \blacktriangledown\,2.1\times $}}\\[-.2pt]30.3s\\[-.2pt]{\fontsize{5.8}{6.4}\selectfont \textcolor{BrickRed}{$ \blacktriangledown\,3.0\times $}}\end{tabular} & \hspace*{1pt}\begin{tabular}[t]{@{}c@{}}\rule{0pt}{12pt}34.1s\\[-.2pt]{\fontsize{5.8}{6.4}\selectfont \textcolor{BrickRed}{$ \blacktriangledown\,2.9\times $}}\\[-.2pt]32s\\[-.2pt]{\fontsize{5.8}{6.4}\selectfont \textcolor{BrickRed}{$ \blacktriangledown\,3.7\times $}}\end{tabular} & \hspace*{1pt}\begin{tabular}[t]{@{}c@{}}\rule{0pt}{12pt}1m\,9.5s\\[-.2pt]{\fontsize{5.8}{6.4}\selectfont \textcolor{BrickRed}{$ \blacktriangledown\,3.0\times $}}\\[-.2pt]1m\,6.1s\\[-.2pt]{\fontsize{5.8}{6.4}\selectfont \textcolor{BrickRed}{$ \blacktriangledown\,4.3\times $}}\end{tabular} & \hspace*{1pt}\begin{tabular}[t]{@{}c@{}}\rule{0pt}{12pt}3m\,1s\\[-.2pt]{\fontsize{5.8}{6.4}\selectfont \textcolor{BrickRed}{$ \blacktriangledown\,7.9\times $}}\\[-.2pt]2m\,51.5s\\[-.2pt]{\fontsize{5.8}{6.4}\selectfont \textcolor{BrickRed}{$ \blacktriangledown\,8.6\times $}}\end{tabular} & \hspace*{1pt}\begin{tabular}[t]{@{}c@{}}\rule{0pt}{12pt}2m\,55.5s\\[-.2pt]{\fontsize{5.8}{6.4}\selectfont \textcolor{ForestGreen}{$ \blacktriangle\,20\% $}}\\[-.2pt]2m\,54s\\[-.2pt]{\fontsize{5.8}{6.4}\selectfont \textcolor{ForestGreen}{$ \blacktriangle\,17\% $}}\end{tabular} & \hspace*{1pt}\begin{tabular}[t]{@{}c@{}}\rule{0pt}{12pt}2m\,50.4s\\[-.2pt]{\fontsize{5.8}{6.4}\selectfont \textcolor{ForestGreen}{$ \blacktriangle\,23\% $}}\\[-.2pt]2m\,48.2s\\[-.2pt]{\fontsize{5.8}{6.4}\selectfont \textcolor{ForestGreen}{$ \blacktriangle\,20\% $}}\end{tabular} & \hspace*{1pt}\begin{tabular}[t]{@{}c@{}}\rule{0pt}{12pt}5m\,35.7s\\[-.2pt]{\fontsize{5.8}{6.4}\selectfont \textcolor{ForestGreen}{$ \blacktriangle\,18\% $}}\\[-.2pt]5m\,31.9s\\[-.2pt]{\fontsize{5.8}{6.4}\selectfont \textcolor{ForestGreen}{$ \blacktriangle\,14\% $}}\end{tabular} & \hspace*{1pt}\begin{tabular}[t]{@{}c@{}}\rule{0pt}{12pt}44m\,20.5s\\[-.2pt]{\fontsize{5.8}{6.4}\selectfont \textcolor{gray}{n/a}}\\[-.2pt]44m\,15.6s\\[-.2pt]{\fontsize{5.8}{6.4}\selectfont \textcolor{gray}{n/a}}\end{tabular} & \hspace*{1pt}\begin{tabular}[t]{@{}c@{}}\rule{0pt}{12pt}1m\,47.1s\\[-.2pt]{\fontsize{5.8}{6.4}\selectfont \textcolor{ForestGreen}{$ \blacktriangle\,5.5\times $}}\\[-.2pt]1m\,45.8s\\[-.2pt]{\fontsize{5.8}{6.4}\selectfont \textcolor{ForestGreen}{$ \blacktriangle\,2.1\times $}}\end{tabular} & \hspace*{1pt}\begin{tabular}[t]{@{}c@{}}\rule{0pt}{12pt}1m\,47.2s\\[-.2pt]{\fontsize{5.8}{6.4}\selectfont \textcolor{ForestGreen}{$ \blacktriangle\,4.6\times $}}\\[-.2pt]1m\,45.1s\\[-.2pt]{\fontsize{5.8}{6.4}\selectfont \textcolor{ForestGreen}{$ \blacktriangle\,2.1\times $}}\end{tabular} & \hspace*{1pt}\begin{tabular}[t]{@{}c@{}}\rule{0pt}{12pt}3m\,40.3s\\[-.2pt]{\fontsize{5.8}{6.4}\selectfont \textcolor{ForestGreen}{$ \blacktriangle\,4.2\times $}}\\[-.2pt]3m\,36.6s\\[-.2pt]{\fontsize{5.8}{6.4}\selectfont \textcolor{ForestGreen}{$ \blacktriangle\,1.7\times $}}\end{tabular} & \hspace*{1pt}\begin{tabular}[t]{@{}c@{}}\rule{0pt}{12pt}29m\,43.4s\\[-.2pt]{\fontsize{5.8}{6.4}\selectfont \textcolor{gray}{n/a}}\\[-.2pt]29m\,38.5s\\[-.2pt]{\fontsize{5.8}{6.4}\selectfont \textcolor{gray}{n/a}}\end{tabular}\\[3pt]\hline
\BenchmarkSourceLabel{bench:aws-4-k15-w10:row:mps-bounded-scalar-ksw2}{R37}\shortstack[l]{{\color{gray}R37}\\\textbf{BioseqDB}} & \begin{tabular}[t]{@{}l@{}}\rule{0pt}{12pt}MPS / Bounded\\Scalar-DP / KSW2\end{tabular} & \hspace*{1pt}\begin{tabular}[t]{@{}c@{}}\rule{0pt}{12pt}12.3s\\[-.2pt]{\fontsize{5.8}{6.4}\selectfont \textcolor{ForestGreen}{$ \blacktriangle\,2.9\times $}}\\[-.2pt]9.7s\\[-.2pt]{\fontsize{5.8}{6.4}\selectfont \textcolor{ForestGreen}{$ \blacktriangle\,3.6\times $}}\end{tabular} & \hspace*{1pt}\begin{tabular}[t]{@{}c@{}}\rule{0pt}{12pt}13.5s\\[-.2pt]{\fontsize{5.8}{6.4}\selectfont \textcolor{ForestGreen}{$ \blacktriangle\,2.8\times $}}\\[-.2pt]9.6s\\[-.2pt]{\fontsize{5.8}{6.4}\selectfont \textcolor{ForestGreen}{$ \blacktriangle\,3.7\times $}}\end{tabular} & \hspace*{1pt}\begin{tabular}[t]{@{}c@{}}\rule{0pt}{12pt}26.1s\\[-.2pt]{\fontsize{5.8}{6.4}\selectfont \textcolor{ForestGreen}{$ \blacktriangle\,3.5\times $}}\\[-.2pt]16.5s\\[-.2pt]{\fontsize{5.8}{6.4}\selectfont \textcolor{ForestGreen}{$ \blacktriangle\,5.4\times $}}\end{tabular} & \textcolor{gray}{\ensuremath{\cdot}} & \hspace*{1pt}\begin{tabular}[t]{@{}c@{}}\rule{0pt}{12pt}8m\,23.3s\\[-.2pt]{\fontsize{5.8}{6.4}\selectfont \textcolor{BrickRed}{$ \blacktriangledown\,15.9\times $}}\\[-.2pt]58.5s\\[-.2pt]{\fontsize{5.8}{6.4}\selectfont \textcolor{BrickRed}{$ \blacktriangledown\,1.9\times $}}\end{tabular} & \hspace*{1pt}\begin{tabular}[t]{@{}c@{}}\rule{0pt}{12pt}6m\,59.2s\\[-.2pt]{\fontsize{5.8}{6.4}\selectfont \textcolor{BrickRed}{$ \blacktriangledown\,12.3\times $}}\\[-.2pt]1m\,4.7s\\[-.2pt]{\fontsize{5.8}{6.4}\selectfont \textcolor{BrickRed}{$ \blacktriangledown\,2.0\times $}}\end{tabular} & \textcolor{gray}{\ensuremath{\cdot}} & \textcolor{gray}{\ensuremath{\cdot}} & \hspace*{1pt}\begin{tabular}[t]{@{}c@{}}\rule{0pt}{12pt}4m\,19.1s\\[-.2pt]{\fontsize{5.8}{6.4}\selectfont \textcolor{BrickRed}{$ \blacktriangledown\,1.5\times $}}\\[-.2pt]4m\,8.7s\\[-.2pt]{\fontsize{5.8}{6.4}\selectfont \textcolor{BrickRed}{$ \blacktriangledown\,1.4\times $}}\end{tabular} & \hspace*{1pt}\begin{tabular}[t]{@{}c@{}}\rule{0pt}{12pt}4m\,29.2s\\[-.2pt]{\fontsize{5.8}{6.4}\selectfont \textcolor{BrickRed}{$ \blacktriangledown\,1.6\times $}}\\[-.2pt]4m\,15s\\[-.2pt]{\fontsize{5.8}{6.4}\selectfont \textcolor{BrickRed}{$ \blacktriangledown\,1.5\times $}}\end{tabular} & \textcolor{gray}{\ensuremath{\cdot}} & \textcolor{gray}{\ensuremath{\cdot}} & \hspace*{1pt}\begin{tabular}[t]{@{}c@{}}\rule{0pt}{12pt}10m\,32.3s\\[-.2pt]{\fontsize{5.8}{6.4}\selectfont \textcolor{BrickRed}{$ \blacktriangledown\,5.9\times $}}\\[-.2pt]4m\,7.7s\\[-.2pt]{\fontsize{5.8}{6.4}\selectfont \textcolor{BrickRed}{$ \blacktriangledown\,2.3\times $}}\end{tabular} & \hspace*{1pt}\begin{tabular}[t]{@{}c@{}}\rule{0pt}{12pt}9m\,53.4s\\[-.2pt]{\fontsize{5.8}{6.4}\selectfont \textcolor{BrickRed}{$ \blacktriangledown\,5.5\times $}}\\[-.2pt]4m\,23.7s\\[-.2pt]{\fontsize{5.8}{6.4}\selectfont \textcolor{BrickRed}{$ \blacktriangledown\,2.5\times $}}\end{tabular} & \textcolor{gray}{\ensuremath{\cdot}} & \textcolor{gray}{\ensuremath{\cdot}}\\[3pt]\hline
\BenchmarkSourceLabel{bench:aws-4-k15-w10:row:mps-simdref-simd-ksw2}{R38}\shortstack[l]{{\color{gray}R38}} & \begin{tabular}[t]{@{}l@{}}\rule{0pt}{12pt}MPS / SIMD-ref\\SIMD-DP / KSW2\end{tabular} & \hspace*{1pt}\begin{tabular}[t]{@{}c@{}}\rule{0pt}{12pt}14.6s\\[-.2pt]{\fontsize{5.8}{6.4}\selectfont \textcolor{ForestGreen}{$ \blacktriangle\,2.5\times $}}\\[-.2pt]9.8s\\[-.2pt]{\fontsize{5.8}{6.4}\selectfont \textcolor{ForestGreen}{$ \blacktriangle\,3.6\times $}}\end{tabular} & \hspace*{1pt}\begin{tabular}[t]{@{}c@{}}\rule{0pt}{12pt}12.6s\\[-.2pt]{\fontsize{5.8}{6.4}\selectfont \textcolor{ForestGreen}{$ \blacktriangle\,3.0\times $}}\\[-.2pt]8.7s\\[-.2pt]{\fontsize{5.8}{6.4}\selectfont \textcolor{ForestGreen}{$ \blacktriangle\,4.1\times $}}\end{tabular} & \textcolor{gray}{\ensuremath{\cdot}} & \textcolor{gray}{\ensuremath{\cdot}} & \hspace*{1pt}\begin{tabular}[t]{@{}c@{}}\rule{0pt}{12pt}6m\,16.2s\\[-.2pt]{\fontsize{5.8}{6.4}\selectfont \textcolor{BrickRed}{$ \blacktriangledown\,11.9\times $}}\\[-.2pt]40.5s\\[-.2pt]{\fontsize{5.8}{6.4}\selectfont \textcolor{BrickRed}{$ \blacktriangledown\,1.3\times $}}\end{tabular} & \hspace*{1pt}\begin{tabular}[t]{@{}c@{}}\rule{0pt}{12pt}5m\,23.1s\\[-.2pt]{\fontsize{5.8}{6.4}\selectfont \textcolor{BrickRed}{$ \blacktriangledown\,9.5\times $}}\\[-.2pt]36.2s\\[-.2pt]{\fontsize{5.8}{6.4}\selectfont \textcolor{BrickRed}{$ \blacktriangledown\,13\% $}}\end{tabular} & \textcolor{gray}{\ensuremath{\cdot}} & \textcolor{gray}{\ensuremath{\cdot}} & \textcolor{gray}{\ensuremath{\cdot}} & \textcolor{gray}{\ensuremath{\cdot}} & \textcolor{gray}{\ensuremath{\cdot}} & \textcolor{gray}{\ensuremath{\cdot}} & \textcolor{gray}{\ensuremath{\cdot}} & \textcolor{gray}{\ensuremath{\cdot}} & \textcolor{gray}{\ensuremath{\cdot}} & \textcolor{gray}{\ensuremath{\cdot}}\\[3pt]\hline
\BenchmarkSourceLabel{bench:aws-4-k15-w10:row:mps-bounded-rmq-ksw2}{R39}\shortstack[l]{{\color{gray}R39}} & \begin{tabular}[t]{@{}l@{}}\rule{0pt}{12pt}MPS / Bounded\\RMQ / KSW2\end{tabular} & \hspace*{1pt}\begin{tabular}[t]{@{}c@{}}\rule{0pt}{12pt}12.6s\\[-.2pt]{\fontsize{5.8}{6.4}\selectfont \textcolor{ForestGreen}{$ \blacktriangle\,2.9\times $}}\\[-.2pt]9.7s\\[-.2pt]{\fontsize{5.8}{6.4}\selectfont \textcolor{ForestGreen}{$ \blacktriangle\,3.6\times $}}\end{tabular} & \hspace*{1pt}\begin{tabular}[t]{@{}c@{}}\rule{0pt}{12pt}13.7s\\[-.2pt]{\fontsize{5.8}{6.4}\selectfont \textcolor{ForestGreen}{$ \blacktriangle\,2.8\times $}}\\[-.2pt]9.5s\\[-.2pt]{\fontsize{5.8}{6.4}\selectfont \textcolor{ForestGreen}{$ \blacktriangle\,3.8\times $}}\end{tabular} & \hspace*{1pt}\begin{tabular}[t]{@{}c@{}}\rule{0pt}{12pt}26.4s\\[-.2pt]{\fontsize{5.8}{6.4}\selectfont \textcolor{ForestGreen}{$ \blacktriangle\,3.5\times $}}\\[-.2pt]16.6s\\[-.2pt]{\fontsize{5.8}{6.4}\selectfont \textcolor{ForestGreen}{$ \blacktriangle\,5.3\times $}}\end{tabular} & \textcolor{gray}{\ensuremath{\cdot}} & \hspace*{1pt}\begin{tabular}[t]{@{}c@{}}\rule{0pt}{12pt}8m\,35.7s\\[-.2pt]{\fontsize{5.8}{6.4}\selectfont \textcolor{BrickRed}{$ \blacktriangledown\,16.3\times $}}\\[-.2pt]58.8s\\[-.2pt]{\fontsize{5.8}{6.4}\selectfont \textcolor{BrickRed}{$ \blacktriangledown\,1.9\times $}}\end{tabular} & \hspace*{1pt}\begin{tabular}[t]{@{}c@{}}\rule{0pt}{12pt}6m\,49.4s\\[-.2pt]{\fontsize{5.8}{6.4}\selectfont \textcolor{BrickRed}{$ \blacktriangledown\,12.0\times $}}\\[-.2pt]1m\,4.5s\\[-.2pt]{\fontsize{5.8}{6.4}\selectfont \textcolor{BrickRed}{$ \blacktriangledown\,2.0\times $}}\end{tabular} & \textcolor{gray}{\ensuremath{\cdot}} & \textcolor{gray}{\ensuremath{\cdot}} & \hspace*{1pt}\begin{tabular}[t]{@{}c@{}}\rule{0pt}{12pt}4m\,20.8s\\[-.2pt]{\fontsize{5.8}{6.4}\selectfont \textcolor{BrickRed}{$ \blacktriangledown\,1.5\times $}}\\[-.2pt]4m\,9.9s\\[-.2pt]{\fontsize{5.8}{6.4}\selectfont \textcolor{BrickRed}{$ \blacktriangledown\,1.4\times $}}\end{tabular} & \hspace*{1pt}\begin{tabular}[t]{@{}c@{}}\rule{0pt}{12pt}3m\,58.2s\\[-.2pt]{\fontsize{5.8}{6.4}\selectfont \textcolor{BrickRed}{$ \blacktriangledown\,1.4\times $}}\\[-.2pt]3m\,44.9s\\[-.2pt]{\fontsize{5.8}{6.4}\selectfont \textcolor{BrickRed}{$ \blacktriangledown\,1.3\times $}}\end{tabular} & \hspace*{1pt}\begin{tabular}[t]{@{}c@{}}\rule{0pt}{12pt}6m\,49.2s\\[-.2pt]{\fontsize{5.8}{6.4}\selectfont \textcolor{BrickRed}{$ \blacktriangledown\,22\% $}}\\[-.2pt]6m\,25.6s\\[-.2pt]{\fontsize{5.8}{6.4}\selectfont \textcolor{BrickRed}{$ \blacktriangledown\,16\% $}}\end{tabular} & \textcolor{gray}{\ensuremath{\cdot}} & \hspace*{1pt}\begin{tabular}[t]{@{}c@{}}\rule{0pt}{12pt}10m\,29.6s\\[-.2pt]{\fontsize{5.8}{6.4}\selectfont \textcolor{BrickRed}{$ \blacktriangledown\,5.9\times $}}\\[-.2pt]4m\,5.4s\\[-.2pt]{\fontsize{5.8}{6.4}\selectfont \textcolor{BrickRed}{$ \blacktriangledown\,2.3\times $}}\end{tabular} & \hspace*{1pt}\begin{tabular}[t]{@{}c@{}}\rule{0pt}{12pt}9m\,40.1s\\[-.2pt]{\fontsize{5.8}{6.4}\selectfont \textcolor{BrickRed}{$ \blacktriangledown\,5.4\times $}}\\[-.2pt]4m\,14.5s\\[-.2pt]{\fontsize{5.8}{6.4}\selectfont \textcolor{BrickRed}{$ \blacktriangledown\,2.4\times $}}\end{tabular} & \textcolor{gray}{\ensuremath{\cdot}} & \textcolor{gray}{\ensuremath{\cdot}}\\[3pt]\hline
\BenchmarkSourceLabel{bench:aws-4-k15-w10:row:mps-simdref-scalar-ksw2}{R40}\shortstack[l]{{\color{gray}R40}} & \begin{tabular}[t]{@{}l@{}}\rule{0pt}{12pt}MPS / SIMD-ref\\Scalar-DP / KSW2\end{tabular} & \hspace*{1pt}\begin{tabular}[t]{@{}c@{}}\rule{0pt}{12pt}14.5s\\[-.2pt]{\fontsize{5.8}{6.4}\selectfont \textcolor{ForestGreen}{$ \blacktriangle\,2.5\times $}}\\[-.2pt]10s\\[-.2pt]{\fontsize{5.8}{6.4}\selectfont \textcolor{ForestGreen}{$ \blacktriangle\,3.5\times $}}\end{tabular} & \hspace*{1pt}\begin{tabular}[t]{@{}c@{}}\rule{0pt}{12pt}12.7s\\[-.2pt]{\fontsize{5.8}{6.4}\selectfont \textcolor{ForestGreen}{$ \blacktriangle\,3.0\times $}}\\[-.2pt]8.7s\\[-.2pt]{\fontsize{5.8}{6.4}\selectfont \textcolor{ForestGreen}{$ \blacktriangle\,4.1\times $}}\end{tabular} & \textcolor{gray}{\ensuremath{\cdot}} & \textcolor{gray}{\ensuremath{\cdot}} & \hspace*{1pt}\begin{tabular}[t]{@{}c@{}}\rule{0pt}{12pt}6m\,13.5s\\[-.2pt]{\fontsize{5.8}{6.4}\selectfont \textcolor{BrickRed}{$ \blacktriangledown\,11.8\times $}}\\[-.2pt]39.6s\\[-.2pt]{\fontsize{5.8}{6.4}\selectfont \textcolor{BrickRed}{$ \blacktriangledown\,1.3\times $}}\end{tabular} & \hspace*{1pt}\begin{tabular}[t]{@{}c@{}}\rule{0pt}{12pt}5m\,21.7s\\[-.2pt]{\fontsize{5.8}{6.4}\selectfont \textcolor{BrickRed}{$ \blacktriangledown\,9.4\times $}}\\[-.2pt]35.3s\\[-.2pt]{\fontsize{5.8}{6.4}\selectfont \textcolor{BrickRed}{$ \blacktriangledown\,10\% $}}\end{tabular} & \textcolor{gray}{\ensuremath{\cdot}} & \textcolor{gray}{\ensuremath{\cdot}} & \textcolor{gray}{\ensuremath{\cdot}} & \textcolor{gray}{\ensuremath{\cdot}} & \textcolor{gray}{\ensuremath{\cdot}} & \textcolor{gray}{\ensuremath{\cdot}} & \textcolor{gray}{\ensuremath{\cdot}} & \textcolor{gray}{\ensuremath{\cdot}} & \textcolor{gray}{\ensuremath{\cdot}} & \textcolor{gray}{\ensuremath{\cdot}}\\[3pt]\hline
\BenchmarkSourceLabel{bench:aws-4-k15-w10:row:mps-scalar-simd-ksw2}{R41}\shortstack[l]{{\color{gray}R41}} & \begin{tabular}[t]{@{}l@{}}\rule{0pt}{12pt}MPS / Scalar\\SIMD-DP / KSW2\end{tabular} & \hspace*{1pt}\begin{tabular}[t]{@{}c@{}}\rule{0pt}{12pt}12.7s\\[-.2pt]{\fontsize{5.8}{6.4}\selectfont \textcolor{ForestGreen}{$ \blacktriangle\,2.9\times $}}\\[-.2pt]9.9s\\[-.2pt]{\fontsize{5.8}{6.4}\selectfont \textcolor{ForestGreen}{$ \blacktriangle\,3.6\times $}}\end{tabular} & \hspace*{1pt}\begin{tabular}[t]{@{}c@{}}\rule{0pt}{12pt}13.2s\\[-.2pt]{\fontsize{5.8}{6.4}\selectfont \textcolor{ForestGreen}{$ \blacktriangle\,2.9\times $}}\\[-.2pt]9.7s\\[-.2pt]{\fontsize{5.8}{6.4}\selectfont \textcolor{ForestGreen}{$ \blacktriangle\,3.7\times $}}\end{tabular} & \textcolor{gray}{\ensuremath{\cdot}} & \textcolor{gray}{\ensuremath{\cdot}} & \hspace*{1pt}\begin{tabular}[t]{@{}c@{}}\rule{0pt}{12pt}8m\,31.6s\\[-.2pt]{\fontsize{5.8}{6.4}\selectfont \textcolor{BrickRed}{$ \blacktriangledown\,16.2\times $}}\\[-.2pt]1m\,3.7s\\[-.2pt]{\fontsize{5.8}{6.4}\selectfont \textcolor{BrickRed}{$ \blacktriangledown\,2.1\times $}}\end{tabular} & \hspace*{1pt}\begin{tabular}[t]{@{}c@{}}\rule{0pt}{12pt}6m\,40.8s\\[-.2pt]{\fontsize{5.8}{6.4}\selectfont \textcolor{BrickRed}{$ \blacktriangledown\,11.8\times $}}\\[-.2pt]56.5s\\[-.2pt]{\fontsize{5.8}{6.4}\selectfont \textcolor{BrickRed}{$ \blacktriangledown\,1.8\times $}}\end{tabular} & \textcolor{gray}{\ensuremath{\cdot}} & \textcolor{gray}{\ensuremath{\cdot}} & \hspace*{1pt}\begin{tabular}[t]{@{}c@{}}\rule{0pt}{12pt}5m\,39.7s\\[-.2pt]{\fontsize{5.8}{6.4}\selectfont \textcolor{BrickRed}{$ \blacktriangledown\,1.9\times $}}\\[-.2pt]5m\,28.2s\\[-.2pt]{\fontsize{5.8}{6.4}\selectfont \textcolor{BrickRed}{$ \blacktriangledown\,1.9\times $}}\end{tabular} & \hspace*{1pt}\begin{tabular}[t]{@{}c@{}}\rule{0pt}{12pt}5m\,42.5s\\[-.2pt]{\fontsize{5.8}{6.4}\selectfont \textcolor{BrickRed}{$ \blacktriangledown\,2.0\times $}}\\[-.2pt]5m\,29s\\[-.2pt]{\fontsize{5.8}{6.4}\selectfont \textcolor{BrickRed}{$ \blacktriangledown\,2.0\times $}}\end{tabular} & \textcolor{gray}{\ensuremath{\cdot}} & \textcolor{gray}{\ensuremath{\cdot}} & \hspace*{1pt}\begin{tabular}[t]{@{}c@{}}\rule{0pt}{12pt}11m\,45s\\[-.2pt]{\fontsize{5.8}{6.4}\selectfont \textcolor{BrickRed}{$ \blacktriangledown\,6.6\times $}}\\[-.2pt]5m\,19.1s\\[-.2pt]{\fontsize{5.8}{6.4}\selectfont \textcolor{BrickRed}{$ \blacktriangledown\,3.0\times $}}\end{tabular} & \hspace*{1pt}\begin{tabular}[t]{@{}c@{}}\rule{0pt}{12pt}10m\,13.5s\\[-.2pt]{\fontsize{5.8}{6.4}\selectfont \textcolor{BrickRed}{$ \blacktriangledown\,5.7\times $}}\\[-.2pt]5m\,31.5s\\[-.2pt]{\fontsize{5.8}{6.4}\selectfont \textcolor{BrickRed}{$ \blacktriangledown\,3.2\times $}}\end{tabular} & \textcolor{gray}{\ensuremath{\cdot}} & \textcolor{gray}{\ensuremath{\cdot}}\\[3pt]\hline
\BenchmarkSourceLabel{bench:aws-4-k15-w10:row:mps-simdref-rmq-ksw2}{R42}\shortstack[l]{{\color{gray}R42}} & \begin{tabular}[t]{@{}l@{}}\rule{0pt}{12pt}MPS / SIMD-ref\\RMQ / KSW2\end{tabular} & \hspace*{1pt}\begin{tabular}[t]{@{}c@{}}\rule{0pt}{12pt}14.5s\\[-.2pt]{\fontsize{5.8}{6.4}\selectfont \textcolor{ForestGreen}{$ \blacktriangle\,2.5\times $}}\\[-.2pt]10s\\[-.2pt]{\fontsize{5.8}{6.4}\selectfont \textcolor{ForestGreen}{$ \blacktriangle\,3.5\times $}}\end{tabular} & \hspace*{1pt}\begin{tabular}[t]{@{}c@{}}\rule{0pt}{12pt}12.8s\\[-.2pt]{\fontsize{5.8}{6.4}\selectfont \textcolor{ForestGreen}{$ \blacktriangle\,3.0\times $}}\\[-.2pt]8.7s\\[-.2pt]{\fontsize{5.8}{6.4}\selectfont \textcolor{ForestGreen}{$ \blacktriangle\,4.1\times $}}\end{tabular} & \textcolor{gray}{\ensuremath{\cdot}} & \textcolor{gray}{\ensuremath{\cdot}} & \hspace*{1pt}\begin{tabular}[t]{@{}c@{}}\rule{0pt}{12pt}6m\,21.1s\\[-.2pt]{\fontsize{5.8}{6.4}\selectfont \textcolor{BrickRed}{$ \blacktriangledown\,12.1\times $}}\\[-.2pt]41.4s\\[-.2pt]{\fontsize{5.8}{6.4}\selectfont \textcolor{BrickRed}{$ \blacktriangledown\,1.4\times $}}\end{tabular} & \hspace*{1pt}\begin{tabular}[t]{@{}c@{}}\rule{0pt}{12pt}5m\,23.8s\\[-.2pt]{\fontsize{5.8}{6.4}\selectfont \textcolor{BrickRed}{$ \blacktriangledown\,9.5\times $}}\\[-.2pt]37s\\[-.2pt]{\fontsize{5.8}{6.4}\selectfont \textcolor{BrickRed}{$ \blacktriangledown\,16\% $}}\end{tabular} & \textcolor{gray}{\ensuremath{\cdot}} & \textcolor{gray}{\ensuremath{\cdot}} & \textcolor{gray}{\ensuremath{\cdot}} & \textcolor{gray}{\ensuremath{\cdot}} & \textcolor{gray}{\ensuremath{\cdot}} & \textcolor{gray}{\ensuremath{\cdot}} & \textcolor{gray}{\ensuremath{\cdot}} & \textcolor{gray}{\ensuremath{\cdot}} & \textcolor{gray}{\ensuremath{\cdot}} & \textcolor{gray}{\ensuremath{\cdot}}\\[3pt]\hline
\BenchmarkSourceLabel{bench:aws-4-k15-w10:row:mps-twostack-simd-ksw2}{R43}\shortstack[l]{{\color{gray}R43}} & \begin{tabular}[t]{@{}l@{}}\rule{0pt}{12pt}MPS / Two-stack\\SIMD-DP / KSW2\end{tabular} & \hspace*{1pt}\begin{tabular}[t]{@{}c@{}}\rule{0pt}{12pt}13.1s\\[-.2pt]{\fontsize{5.8}{6.4}\selectfont \textcolor{ForestGreen}{$ \blacktriangle\,2.8\times $}}\\[-.2pt]10.5s\\[-.2pt]{\fontsize{5.8}{6.4}\selectfont \textcolor{ForestGreen}{$ \blacktriangle\,3.3\times $}}\end{tabular} & \hspace*{1pt}\begin{tabular}[t]{@{}c@{}}\rule{0pt}{12pt}13.4s\\[-.2pt]{\fontsize{5.8}{6.4}\selectfont \textcolor{ForestGreen}{$ \blacktriangle\,2.8\times $}}\\[-.2pt]9.7s\\[-.2pt]{\fontsize{5.8}{6.4}\selectfont \textcolor{ForestGreen}{$ \blacktriangle\,3.7\times $}}\end{tabular} & \textcolor{gray}{\ensuremath{\cdot}} & \textcolor{gray}{\ensuremath{\cdot}} & \hspace*{1pt}\begin{tabular}[t]{@{}c@{}}\rule{0pt}{12pt}8m\,20.3s\\[-.2pt]{\fontsize{5.8}{6.4}\selectfont \textcolor{BrickRed}{$ \blacktriangledown\,15.8\times $}}\\[-.2pt]45.6s\\[-.2pt]{\fontsize{5.8}{6.4}\selectfont \textcolor{BrickRed}{$ \blacktriangledown\,1.5\times $}}\end{tabular} & \hspace*{1pt}\begin{tabular}[t]{@{}c@{}}\rule{0pt}{12pt}6m\,45s\\[-.2pt]{\fontsize{5.8}{6.4}\selectfont \textcolor{BrickRed}{$ \blacktriangledown\,11.9\times $}}\\[-.2pt]45.4s\\[-.2pt]{\fontsize{5.8}{6.4}\selectfont \textcolor{BrickRed}{$ \blacktriangledown\,1.4\times $}}\end{tabular} & \textcolor{gray}{\ensuremath{\cdot}} & \textcolor{gray}{\ensuremath{\cdot}} & \hspace*{1pt}\begin{tabular}[t]{@{}c@{}}\rule{0pt}{12pt}4m\,54.4s\\[-.2pt]{\fontsize{5.8}{6.4}\selectfont \textcolor{BrickRed}{$ \blacktriangledown\,1.7\times $}}\\[-.2pt]4m\,44.1s\\[-.2pt]{\fontsize{5.8}{6.4}\selectfont \textcolor{BrickRed}{$ \blacktriangledown\,1.6\times $}}\end{tabular} & \hspace*{1pt}\begin{tabular}[t]{@{}c@{}}\rule{0pt}{12pt}4m\,39.9s\\[-.2pt]{\fontsize{5.8}{6.4}\selectfont \textcolor{BrickRed}{$ \blacktriangledown\,1.6\times $}}\\[-.2pt]4m\,26.4s\\[-.2pt]{\fontsize{5.8}{6.4}\selectfont \textcolor{BrickRed}{$ \blacktriangledown\,1.6\times $}}\end{tabular} & \textcolor{gray}{\ensuremath{\cdot}} & \textcolor{gray}{\ensuremath{\cdot}} & \textcolor{gray}{\ensuremath{\cdot}} & \textcolor{gray}{\ensuremath{\cdot}} & \textcolor{gray}{\ensuremath{\cdot}} & \textcolor{gray}{\ensuremath{\cdot}}\\[3pt]\hline
\end{tabular}\par}
\smallskip{\footnotesize\noindent Each cell: A and B times, $\cdot$: no measurement. ---\textsuperscript{g}: GPU unavailable. n/a: no comparison.\par
\noindent Blue: one lowest A and one lowest B measurement for each column for each configuration.\par
\noindent Gray: minimap2. Yellow: my preferred MPF / Bounded / SIMD-DP / Block configuration.\par
\noindent MPF/MPS: minimizers per fragment/sequence. SIMD-ref: naive SIMD scan. Bounded: bounded SIMD algorithm.\par
}
\clearpage
\noindent{\large\bfseries Table \thetable: (B): AWS: 4 Gbp, k=15, w=10\hfill 7/7}\par\smallskip
{\fontsize{7.1}{8.1}\selectfont\setlength{\tabcolsep}{1pt}\setlength{\fboxsep}{2pt}\renewcommand{\arraystretch}{1.12}
\noindent\begin{tabular}{@{}p{14mm}p{27mm}|p{13.55mm}|p{13.55mm}|p{13.55mm}|p{13.55mm}|p{13.55mm}|p{13.55mm}|p{13.55mm}|p{13.55mm}|p{13.55mm}|p{13.55mm}|p{13.55mm}|p{13.55mm}|p{13.55mm}|p{13.55mm}|p{13.55mm}|p{13.55mm}|@{}}\hline
Engine & Execution & \multicolumn{8}{c|}{\textbf{Index}} & \multicolumn{8}{c|}{\textbf{No index}}\\
 &  & \multicolumn{8}{c|}{\textbf{4 Gbp}} & \multicolumn{8}{c|}{\textbf{4 Gbp}}\\
 &  & \multicolumn{8}{c|}{\textbf{k=15, w=10}} & \multicolumn{8}{c|}{\textbf{k=15, w=10}}\\
 &  & \multicolumn{4}{c|}{\textbf{DNA4}} & \multicolumn{4}{c|}{\textbf{DNA5}} & \multicolumn{4}{c|}{\textbf{DNA4}} & \multicolumn{4}{c|}{\textbf{DNA5}}\\
 &  & \multicolumn{1}{c|}{\textbf{M7A}} & \multicolumn{1}{c|}{\textbf{M7I}} & \multicolumn{1}{c|}{\textbf{M4}} & \multicolumn{1}{c|}{\textbf{C7G}} & \multicolumn{1}{c|}{\textbf{M7A}} & \multicolumn{1}{c|}{\textbf{M7I}} & \multicolumn{1}{c|}{\textbf{M4}} & \multicolumn{1}{c|}{\textbf{C7G}} & \multicolumn{1}{c|}{\textbf{M7A}} & \multicolumn{1}{c|}{\textbf{M7I}} & \multicolumn{1}{c|}{\textbf{M4}} & \multicolumn{1}{c|}{\textbf{C7G}} & \multicolumn{1}{c|}{\textbf{M7A}} & \multicolumn{1}{c|}{\textbf{M7I}} & \multicolumn{1}{c|}{\textbf{M4}} & \multicolumn{1}{c|}{\textbf{C7G}}\\
 & & {\color{gray}C1} & {\color{gray}C2} & {\color{gray}C3} & {\color{gray}C4} & {\color{gray}C5} & {\color{gray}C6} & {\color{gray}C7} & {\color{gray}C8} & {\color{gray}C9} & {\color{gray}C10} & {\color{gray}C11} & {\color{gray}C12} & {\color{gray}C13} & {\color{gray}C14} & {\color{gray}C15} & {\color{gray}C16}\\
\hline
\rowcolor{BenchBaseline}[0.5pt][0.5pt]\shortstack[l]{{\color{gray}R1}\\\textbf{minimap2}} & \textbf{Baseline} & \hspace*{1pt}\begin{tabular}[t]{@{}c@{}}\rule{0pt}{12pt}36.2s\\[-.2pt]{\fontsize{5.8}{6.4}\selectfont \textcolor{BrickRed}{$ \blacktriangledown\,6.9\times $}}\\[-.2pt]35s\\[-.2pt]{\fontsize{5.8}{6.4}\selectfont \textcolor{BrickRed}{$ \blacktriangledown\,8.0\times $}}\end{tabular} & \hspace*{1pt}\begin{tabular}[t]{@{}c@{}}\rule{0pt}{12pt}38.1s\\[-.2pt]{\fontsize{5.8}{6.4}\selectfont \textcolor{BrickRed}{$ \blacktriangledown\,7.8\times $}}\\[-.2pt]36s\\[-.2pt]{\fontsize{5.8}{6.4}\selectfont \textcolor{BrickRed}{$ \blacktriangledown\,9.3\times $}}\end{tabular} & \hspace*{1pt}\begin{tabular}[t]{@{}c@{}}\rule{0pt}{12pt}1m\,32.4s\\[-.2pt]{\fontsize{5.8}{6.4}\selectfont \textcolor{BrickRed}{$ \blacktriangledown\,10.7\times $}}\\[-.2pt]1m\,28.5s\\[-.2pt]{\fontsize{5.8}{6.4}\selectfont \textcolor{BrickRed}{$ \blacktriangledown\,13.5\times $}}\end{tabular} & \hspace*{1pt}\begin{tabular}[t]{@{}c@{}}\rule{0pt}{12pt}6m\,31.5s\\[-.2pt]{\fontsize{5.8}{6.4}\selectfont \textcolor{BrickRed}{$ \blacktriangledown\,11.5\times $}}\\[-.2pt]6m\,26.3s\\[-.2pt]{\fontsize{5.8}{6.4}\selectfont \textcolor{BrickRed}{$ \blacktriangledown\,12.7\times $}}\end{tabular} & \hspace*{1pt}\begin{tabular}[t]{@{}c@{}}\rule{0pt}{12pt}31.6s\\[-.2pt]{\fontsize{5.8}{6.4}\selectfont \textcolor{BrickRed}{$ \blacktriangledown\,2.1\times $}}\\[-.2pt]30.3s\\[-.2pt]{\fontsize{5.8}{6.4}\selectfont \textcolor{BrickRed}{$ \blacktriangledown\,3.0\times $}}\end{tabular} & \hspace*{1pt}\begin{tabular}[t]{@{}c@{}}\rule{0pt}{12pt}34.1s\\[-.2pt]{\fontsize{5.8}{6.4}\selectfont \textcolor{BrickRed}{$ \blacktriangledown\,2.9\times $}}\\[-.2pt]32s\\[-.2pt]{\fontsize{5.8}{6.4}\selectfont \textcolor{BrickRed}{$ \blacktriangledown\,3.7\times $}}\end{tabular} & \hspace*{1pt}\begin{tabular}[t]{@{}c@{}}\rule{0pt}{12pt}1m\,9.5s\\[-.2pt]{\fontsize{5.8}{6.4}\selectfont \textcolor{BrickRed}{$ \blacktriangledown\,3.0\times $}}\\[-.2pt]1m\,6.1s\\[-.2pt]{\fontsize{5.8}{6.4}\selectfont \textcolor{BrickRed}{$ \blacktriangledown\,4.3\times $}}\end{tabular} & \hspace*{1pt}\begin{tabular}[t]{@{}c@{}}\rule{0pt}{12pt}3m\,1s\\[-.2pt]{\fontsize{5.8}{6.4}\selectfont \textcolor{BrickRed}{$ \blacktriangledown\,7.9\times $}}\\[-.2pt]2m\,51.5s\\[-.2pt]{\fontsize{5.8}{6.4}\selectfont \textcolor{BrickRed}{$ \blacktriangledown\,8.6\times $}}\end{tabular} & \hspace*{1pt}\begin{tabular}[t]{@{}c@{}}\rule{0pt}{12pt}2m\,55.5s\\[-.2pt]{\fontsize{5.8}{6.4}\selectfont \textcolor{ForestGreen}{$ \blacktriangle\,20\% $}}\\[-.2pt]2m\,54s\\[-.2pt]{\fontsize{5.8}{6.4}\selectfont \textcolor{ForestGreen}{$ \blacktriangle\,17\% $}}\end{tabular} & \hspace*{1pt}\begin{tabular}[t]{@{}c@{}}\rule{0pt}{12pt}2m\,50.4s\\[-.2pt]{\fontsize{5.8}{6.4}\selectfont \textcolor{ForestGreen}{$ \blacktriangle\,23\% $}}\\[-.2pt]2m\,48.2s\\[-.2pt]{\fontsize{5.8}{6.4}\selectfont \textcolor{ForestGreen}{$ \blacktriangle\,20\% $}}\end{tabular} & \hspace*{1pt}\begin{tabular}[t]{@{}c@{}}\rule{0pt}{12pt}5m\,35.7s\\[-.2pt]{\fontsize{5.8}{6.4}\selectfont \textcolor{ForestGreen}{$ \blacktriangle\,18\% $}}\\[-.2pt]5m\,31.9s\\[-.2pt]{\fontsize{5.8}{6.4}\selectfont \textcolor{ForestGreen}{$ \blacktriangle\,14\% $}}\end{tabular} & \hspace*{1pt}\begin{tabular}[t]{@{}c@{}}\rule{0pt}{12pt}44m\,20.5s\\[-.2pt]{\fontsize{5.8}{6.4}\selectfont \textcolor{gray}{n/a}}\\[-.2pt]44m\,15.6s\\[-.2pt]{\fontsize{5.8}{6.4}\selectfont \textcolor{gray}{n/a}}\end{tabular} & \hspace*{1pt}\begin{tabular}[t]{@{}c@{}}\rule{0pt}{12pt}1m\,47.1s\\[-.2pt]{\fontsize{5.8}{6.4}\selectfont \textcolor{ForestGreen}{$ \blacktriangle\,5.5\times $}}\\[-.2pt]1m\,45.8s\\[-.2pt]{\fontsize{5.8}{6.4}\selectfont \textcolor{ForestGreen}{$ \blacktriangle\,2.1\times $}}\end{tabular} & \hspace*{1pt}\begin{tabular}[t]{@{}c@{}}\rule{0pt}{12pt}1m\,47.2s\\[-.2pt]{\fontsize{5.8}{6.4}\selectfont \textcolor{ForestGreen}{$ \blacktriangle\,4.6\times $}}\\[-.2pt]1m\,45.1s\\[-.2pt]{\fontsize{5.8}{6.4}\selectfont \textcolor{ForestGreen}{$ \blacktriangle\,2.1\times $}}\end{tabular} & \hspace*{1pt}\begin{tabular}[t]{@{}c@{}}\rule{0pt}{12pt}3m\,40.3s\\[-.2pt]{\fontsize{5.8}{6.4}\selectfont \textcolor{ForestGreen}{$ \blacktriangle\,4.2\times $}}\\[-.2pt]3m\,36.6s\\[-.2pt]{\fontsize{5.8}{6.4}\selectfont \textcolor{ForestGreen}{$ \blacktriangle\,1.7\times $}}\end{tabular} & \hspace*{1pt}\begin{tabular}[t]{@{}c@{}}\rule{0pt}{12pt}29m\,43.4s\\[-.2pt]{\fontsize{5.8}{6.4}\selectfont \textcolor{gray}{n/a}}\\[-.2pt]29m\,38.5s\\[-.2pt]{\fontsize{5.8}{6.4}\selectfont \textcolor{gray}{n/a}}\end{tabular}\\[3pt]\hline
\BenchmarkSourceLabel{bench:aws-4-k15-w10:row:mps-twostack-scalar-ksw2}{R44}\shortstack[l]{{\color{gray}R44}\\\textbf{BioseqDB}} & \begin{tabular}[t]{@{}l@{}}\rule{0pt}{12pt}MPS / Two-stack\\Scalar-DP / KSW2\end{tabular} & \hspace*{1pt}\begin{tabular}[t]{@{}c@{}}\rule{0pt}{12pt}14.6s\\[-.2pt]{\fontsize{5.8}{6.4}\selectfont \textcolor{ForestGreen}{$ \blacktriangle\,2.5\times $}}\\[-.2pt]9.9s\\[-.2pt]{\fontsize{5.8}{6.4}\selectfont \textcolor{ForestGreen}{$ \blacktriangle\,3.5\times $}}\end{tabular} & \hspace*{1pt}\begin{tabular}[t]{@{}c@{}}\rule{0pt}{12pt}13.2s\\[-.2pt]{\fontsize{5.8}{6.4}\selectfont \textcolor{ForestGreen}{$ \blacktriangle\,2.9\times $}}\\[-.2pt]9.2s\\[-.2pt]{\fontsize{5.8}{6.4}\selectfont \textcolor{ForestGreen}{$ \blacktriangle\,3.9\times $}}\end{tabular} & \textcolor{gray}{\ensuremath{\cdot}} & \textcolor{gray}{\ensuremath{\cdot}} & \hspace*{1pt}\begin{tabular}[t]{@{}c@{}}\rule{0pt}{12pt}8m\,22s\\[-.2pt]{\fontsize{5.8}{6.4}\selectfont \textcolor{BrickRed}{$ \blacktriangledown\,15.9\times $}}\\[-.2pt]46.5s\\[-.2pt]{\fontsize{5.8}{6.4}\selectfont \textcolor{BrickRed}{$ \blacktriangledown\,1.5\times $}}\end{tabular} & \hspace*{1pt}\begin{tabular}[t]{@{}c@{}}\rule{0pt}{12pt}6m\,54.6s\\[-.2pt]{\fontsize{5.8}{6.4}\selectfont \textcolor{BrickRed}{$ \blacktriangledown\,12.2\times $}}\\[-.2pt]45.8s\\[-.2pt]{\fontsize{5.8}{6.4}\selectfont \textcolor{BrickRed}{$ \blacktriangledown\,1.4\times $}}\end{tabular} & \textcolor{gray}{\ensuremath{\cdot}} & \textcolor{gray}{\ensuremath{\cdot}} & \textcolor{gray}{\ensuremath{\cdot}} & \textcolor{gray}{\ensuremath{\cdot}} & \textcolor{gray}{\ensuremath{\cdot}} & \textcolor{gray}{\ensuremath{\cdot}} & \hspace*{1pt}\begin{tabular}[t]{@{}c@{}}\rule{0pt}{12pt}10m\,24.5s\\[-.2pt]{\fontsize{5.8}{6.4}\selectfont \textcolor{BrickRed}{$ \blacktriangledown\,5.8\times $}}\\[-.2pt]4m\,4.1s\\[-.2pt]{\fontsize{5.8}{6.4}\selectfont \textcolor{BrickRed}{$ \blacktriangledown\,2.3\times $}}\end{tabular} & \hspace*{1pt}\begin{tabular}[t]{@{}c@{}}\rule{0pt}{12pt}8m\,30.6s\\[-.2pt]{\fontsize{5.8}{6.4}\selectfont \textcolor{BrickRed}{$ \blacktriangledown\,4.8\times $}}\\[-.2pt]3m\,55.3s\\[-.2pt]{\fontsize{5.8}{6.4}\selectfont \textcolor{BrickRed}{$ \blacktriangledown\,2.2\times $}}\end{tabular} & \textcolor{gray}{\ensuremath{\cdot}} & \textcolor{gray}{\ensuremath{\cdot}}\\[3pt]\hline
\BenchmarkSourceLabel{bench:aws-4-k15-w10:row:mps-scalar-rmq-ksw2}{R45}\shortstack[l]{{\color{gray}R45}} & \begin{tabular}[t]{@{}l@{}}\rule{0pt}{12pt}MPS / Scalar\\RMQ / KSW2\end{tabular} & \hspace*{1pt}\begin{tabular}[t]{@{}c@{}}\rule{0pt}{12pt}13.3s\\[-.2pt]{\fontsize{5.8}{6.4}\selectfont \textcolor{ForestGreen}{$ \blacktriangle\,2.7\times $}}\\[-.2pt]10.8s\\[-.2pt]{\fontsize{5.8}{6.4}\selectfont \textcolor{ForestGreen}{$ \blacktriangle\,3.3\times $}}\end{tabular} & \hspace*{1pt}\begin{tabular}[t]{@{}c@{}}\rule{0pt}{12pt}18.6s\\[-.2pt]{\fontsize{5.8}{6.4}\selectfont \textcolor{ForestGreen}{$ \blacktriangle\,2.0\times $}}\\[-.2pt]13.7s\\[-.2pt]{\fontsize{5.8}{6.4}\selectfont \textcolor{ForestGreen}{$ \blacktriangle\,2.6\times $}}\end{tabular} & \textcolor{gray}{\ensuremath{\cdot}} & \textcolor{gray}{\ensuremath{\cdot}} & \hspace*{1pt}\begin{tabular}[t]{@{}c@{}}\rule{0pt}{12pt}7m\,41.4s\\[-.2pt]{\fontsize{5.8}{6.4}\selectfont \textcolor{BrickRed}{$ \blacktriangledown\,14.6\times $}}\\[-.2pt]59.5s\\[-.2pt]{\fontsize{5.8}{6.4}\selectfont \textcolor{BrickRed}{$ \blacktriangledown\,2.0\times $}}\end{tabular} & \hspace*{1pt}\begin{tabular}[t]{@{}c@{}}\rule{0pt}{12pt}6m\,30.2s\\[-.2pt]{\fontsize{5.8}{6.4}\selectfont \textcolor{BrickRed}{$ \blacktriangledown\,11.5\times $}}\\[-.2pt]54.7s\\[-.2pt]{\fontsize{5.8}{6.4}\selectfont \textcolor{BrickRed}{$ \blacktriangledown\,1.7\times $}}\end{tabular} & \textcolor{gray}{\ensuremath{\cdot}} & \textcolor{gray}{\ensuremath{\cdot}} & \hspace*{1pt}\begin{tabular}[t]{@{}c@{}}\rule{0pt}{12pt}5m\,34.3s\\[-.2pt]{\fontsize{5.8}{6.4}\selectfont \textcolor{BrickRed}{$ \blacktriangledown\,1.9\times $}}\\[-.2pt]5m\,20.2s\\[-.2pt]{\fontsize{5.8}{6.4}\selectfont \textcolor{BrickRed}{$ \blacktriangledown\,1.8\times $}}\end{tabular} & \textcolor{gray}{\ensuremath{\cdot}} & \textcolor{gray}{\ensuremath{\cdot}} & \textcolor{gray}{\ensuremath{\cdot}} & \hspace*{1pt}\begin{tabular}[t]{@{}c@{}}\rule{0pt}{12pt}12m\,14s\\[-.2pt]{\fontsize{5.8}{6.4}\selectfont \textcolor{BrickRed}{$ \blacktriangledown\,6.9\times $}}\\[-.2pt]5m\,48.2s\\[-.2pt]{\fontsize{5.8}{6.4}\selectfont \textcolor{BrickRed}{$ \blacktriangledown\,3.3\times $}}\end{tabular} & \textcolor{gray}{\ensuremath{\cdot}} & \textcolor{gray}{\ensuremath{\cdot}} & \textcolor{gray}{\ensuremath{\cdot}}\\[3pt]\hline
\BenchmarkSourceLabel{bench:aws-4-k15-w10:row:mps-twostack-rmq-ksw2}{R46}\shortstack[l]{{\color{gray}R46}} & \begin{tabular}[t]{@{}l@{}}\rule{0pt}{12pt}MPS / Two-stack\\RMQ / KSW2\end{tabular} & \hspace*{1pt}\begin{tabular}[t]{@{}c@{}}\rule{0pt}{12pt}13.3s\\[-.2pt]{\fontsize{5.8}{6.4}\selectfont \textcolor{ForestGreen}{$ \blacktriangle\,2.7\times $}}\\[-.2pt]10.8s\\[-.2pt]{\fontsize{5.8}{6.4}\selectfont \textcolor{ForestGreen}{$ \blacktriangle\,3.2\times $}}\end{tabular} & \hspace*{1pt}\begin{tabular}[t]{@{}c@{}}\rule{0pt}{12pt}17.9s\\[-.2pt]{\fontsize{5.8}{6.4}\selectfont \textcolor{ForestGreen}{$ \blacktriangle\,2.1\times $}}\\[-.2pt]13.5s\\[-.2pt]{\fontsize{5.8}{6.4}\selectfont \textcolor{ForestGreen}{$ \blacktriangle\,2.7\times $}}\end{tabular} & \textcolor{gray}{\ensuremath{\cdot}} & \textcolor{gray}{\ensuremath{\cdot}} & \hspace*{1pt}\begin{tabular}[t]{@{}c@{}}\rule{0pt}{12pt}8m\,11.3s\\[-.2pt]{\fontsize{5.8}{6.4}\selectfont \textcolor{BrickRed}{$ \blacktriangledown\,15.6\times $}}\\[-.2pt]46s\\[-.2pt]{\fontsize{5.8}{6.4}\selectfont \textcolor{BrickRed}{$ \blacktriangledown\,1.5\times $}}\end{tabular} & \hspace*{1pt}\begin{tabular}[t]{@{}c@{}}\rule{0pt}{12pt}6m\,46s\\[-.2pt]{\fontsize{5.8}{6.4}\selectfont \textcolor{BrickRed}{$ \blacktriangledown\,11.9\times $}}\\[-.2pt]45.8s\\[-.2pt]{\fontsize{5.8}{6.4}\selectfont \textcolor{BrickRed}{$ \blacktriangledown\,1.4\times $}}\end{tabular} & \textcolor{gray}{\ensuremath{\cdot}} & \textcolor{gray}{\ensuremath{\cdot}} & \textcolor{gray}{\ensuremath{\cdot}} & \textcolor{gray}{\ensuremath{\cdot}} & \textcolor{gray}{\ensuremath{\cdot}} & \textcolor{gray}{\ensuremath{\cdot}} & \textcolor{gray}{\ensuremath{\cdot}} & \textcolor{gray}{\ensuremath{\cdot}} & \textcolor{gray}{\ensuremath{\cdot}} & \textcolor{gray}{\ensuremath{\cdot}}\\[3pt]\hline
\BenchmarkSourceLabel{bench:aws-4-k15-w10:row:mpf-twostack-scalar-ksw2}{R47}\shortstack[l]{{\color{gray}R47}} & \begin{tabular}[t]{@{}l@{}}\rule{0pt}{12pt}MPF / Two-stack\\Scalar-DP / KSW2\end{tabular} & \hspace*{1pt}\begin{tabular}[t]{@{}c@{}}\rule{0pt}{12pt}20.1s\\[-.2pt]{\fontsize{5.8}{6.4}\selectfont \textcolor{ForestGreen}{$ \blacktriangle\,1.8\times $}}\\[-.2pt]17.6s\\[-.2pt]{\fontsize{5.8}{6.4}\selectfont \textcolor{ForestGreen}{$ \blacktriangle\,2.0\times $}}\end{tabular} & \hspace*{1pt}\begin{tabular}[t]{@{}c@{}}\rule{0pt}{12pt}16.8s\\[-.2pt]{\fontsize{5.8}{6.4}\selectfont \textcolor{ForestGreen}{$ \blacktriangle\,2.3\times $}}\\[-.2pt]14.4s\\[-.2pt]{\fontsize{5.8}{6.4}\selectfont \textcolor{ForestGreen}{$ \blacktriangle\,2.5\times $}}\end{tabular} & \hspace*{1pt}\begin{tabular}[t]{@{}c@{}}\rule{0pt}{12pt}17.5s\\[-.2pt]{\fontsize{5.8}{6.4}\selectfont \textcolor{ForestGreen}{$ \blacktriangle\,5.3\times $}}\\[-.2pt]9.2s\\[-.2pt]{\fontsize{5.8}{6.4}\selectfont \textcolor{ForestGreen}{$ \blacktriangle\,9.6\times $}}\end{tabular} & \hspace*{1pt}\begin{tabular}[t]{@{}c@{}}\rule{0pt}{12pt}35.7s\\[-.2pt]{\fontsize{5.8}{6.4}\selectfont \textcolor{ForestGreen}{$ \blacktriangle\,11.0\times $}}\\[-.2pt]31.7s\\[-.2pt]{\fontsize{5.8}{6.4}\selectfont \textcolor{ForestGreen}{$ \blacktriangle\,12.2\times $}}\end{tabular} & \hspace*{1pt}\begin{tabular}[t]{@{}c@{}}\rule{0pt}{12pt}15.2s\\[-.2pt]{\fontsize{5.8}{6.4}\selectfont \textcolor{ForestGreen}{$ \blacktriangle\,2.1\times $}}\\[-.2pt]11s\\[-.2pt]{\fontsize{5.8}{6.4}\selectfont \textcolor{ForestGreen}{$ \blacktriangle\,2.7\times $}}\end{tabular} & \hspace*{1pt}\begin{tabular}[t]{@{}c@{}}\rule{0pt}{12pt}14.2s\\[-.2pt]{\fontsize{5.8}{6.4}\selectfont \textcolor{ForestGreen}{$ \blacktriangle\,2.4\times $}}\\[-.2pt]10.7s\\[-.2pt]{\fontsize{5.8}{6.4}\selectfont \textcolor{ForestGreen}{$ \blacktriangle\,3.0\times $}}\end{tabular} & \hspace*{1pt}\begin{tabular}[t]{@{}c@{}}\rule{0pt}{12pt}23.3s\\[-.2pt]{\fontsize{5.8}{6.4}\selectfont \textcolor{ForestGreen}{$ \blacktriangle\,3.0\times $}}\\[-.2pt]\colorbox{BenchBest}{\strut 15.5s}\\[-.2pt]{\fontsize{5.8}{6.4}\selectfont \textcolor{ForestGreen}{$ \blacktriangle\,4.3\times $}}\end{tabular} & \textcolor{gray}{\ensuremath{\cdot}} & \hspace*{1pt}\begin{tabular}[t]{@{}c@{}}\rule{0pt}{12pt}4m\,31.8s\\[-.2pt]{\fontsize{5.8}{6.4}\selectfont \textcolor{BrickRed}{$ \blacktriangledown\,1.5\times $}}\\[-.2pt]4m\,18.1s\\[-.2pt]{\fontsize{5.8}{6.4}\selectfont \textcolor{BrickRed}{$ \blacktriangledown\,1.5\times $}}\end{tabular} & \hspace*{1pt}\begin{tabular}[t]{@{}c@{}}\rule{0pt}{12pt}4m\,18s\\[-.2pt]{\fontsize{5.8}{6.4}\selectfont \textcolor{BrickRed}{$ \blacktriangledown\,1.5\times $}}\\[-.2pt]4m\,3.8s\\[-.2pt]{\fontsize{5.8}{6.4}\selectfont \textcolor{BrickRed}{$ \blacktriangledown\,1.4\times $}}\end{tabular} & \textcolor{gray}{\ensuremath{\cdot}} & \textcolor{gray}{\ensuremath{\cdot}} & \hspace*{1pt}\begin{tabular}[t]{@{}c@{}}\rule{0pt}{12pt}10m\,1.7s\\[-.2pt]{\fontsize{5.8}{6.4}\selectfont \textcolor{BrickRed}{$ \blacktriangledown\,5.6\times $}}\\[-.2pt]3m\,53.1s\\[-.2pt]{\fontsize{5.8}{6.4}\selectfont \textcolor{BrickRed}{$ \blacktriangledown\,2.2\times $}}\end{tabular} & \hspace*{1pt}\begin{tabular}[t]{@{}c@{}}\rule{0pt}{12pt}8m\,16.7s\\[-.2pt]{\fontsize{5.8}{6.4}\selectfont \textcolor{BrickRed}{$ \blacktriangledown\,4.6\times $}}\\[-.2pt]3m\,46.2s\\[-.2pt]{\fontsize{5.8}{6.4}\selectfont \textcolor{BrickRed}{$ \blacktriangledown\,2.2\times $}}\end{tabular} & \textcolor{gray}{\ensuremath{\cdot}} & \textcolor{gray}{\ensuremath{\cdot}}\\[3pt]\hline
\BenchmarkSourceLabel{bench:aws-4-k15-w10:row:mps-scalar-scalar-ksw2}{R48}\shortstack[l]{{\color{gray}R48}} & \begin{tabular}[t]{@{}l@{}}\rule{0pt}{12pt}MPS / Scalar\\Scalar-DP / KSW2\end{tabular} & \hspace*{1pt}\begin{tabular}[t]{@{}c@{}}\rule{0pt}{12pt}14.5s\\[-.2pt]{\fontsize{5.8}{6.4}\selectfont \textcolor{ForestGreen}{$ \blacktriangle\,2.5\times $}}\\[-.2pt]10.1s\\[-.2pt]{\fontsize{5.8}{6.4}\selectfont \textcolor{ForestGreen}{$ \blacktriangle\,3.5\times $}}\end{tabular} & \hspace*{1pt}\begin{tabular}[t]{@{}c@{}}\rule{0pt}{12pt}14.3s\\[-.2pt]{\fontsize{5.8}{6.4}\selectfont \textcolor{ForestGreen}{$ \blacktriangle\,2.7\times $}}\\[-.2pt]9.5s\\[-.2pt]{\fontsize{5.8}{6.4}\selectfont \textcolor{ForestGreen}{$ \blacktriangle\,3.8\times $}}\end{tabular} & \textcolor{gray}{\ensuremath{\cdot}} & \textcolor{gray}{\ensuremath{\cdot}} & \hspace*{1pt}\begin{tabular}[t]{@{}c@{}}\rule{0pt}{12pt}7m\,42s\\[-.2pt]{\fontsize{5.8}{6.4}\selectfont \textcolor{BrickRed}{$ \blacktriangledown\,14.6\times $}}\\[-.2pt]58.8s\\[-.2pt]{\fontsize{5.8}{6.4}\selectfont \textcolor{BrickRed}{$ \blacktriangledown\,1.9\times $}}\end{tabular} & \hspace*{1pt}\begin{tabular}[t]{@{}c@{}}\rule{0pt}{12pt}6m\,31.7s\\[-.2pt]{\fontsize{5.8}{6.4}\selectfont \textcolor{BrickRed}{$ \blacktriangledown\,11.5\times $}}\\[-.2pt]54.5s\\[-.2pt]{\fontsize{5.8}{6.4}\selectfont \textcolor{BrickRed}{$ \blacktriangledown\,1.7\times $}}\end{tabular} & \textcolor{gray}{\ensuremath{\cdot}} & \textcolor{gray}{\ensuremath{\cdot}} & \hspace*{1pt}\begin{tabular}[t]{@{}c@{}}\rule{0pt}{12pt}5m\,38.2s\\[-.2pt]{\fontsize{5.8}{6.4}\selectfont \textcolor{BrickRed}{$ \blacktriangledown\,1.9\times $}}\\[-.2pt]5m\,24.8s\\[-.2pt]{\fontsize{5.8}{6.4}\selectfont \textcolor{BrickRed}{$ \blacktriangledown\,1.9\times $}}\end{tabular} & \textcolor{gray}{\ensuremath{\cdot}} & \textcolor{gray}{\ensuremath{\cdot}} & \textcolor{gray}{\ensuremath{\cdot}} & \hspace*{1pt}\begin{tabular}[t]{@{}c@{}}\rule{0pt}{12pt}12m\,12s\\[-.2pt]{\fontsize{5.8}{6.4}\selectfont \textcolor{BrickRed}{$ \blacktriangledown\,6.8\times $}}\\[-.2pt]5m\,47s\\[-.2pt]{\fontsize{5.8}{6.4}\selectfont \textcolor{BrickRed}{$ \blacktriangledown\,3.3\times $}}\end{tabular} & \textcolor{gray}{\ensuremath{\cdot}} & \textcolor{gray}{\ensuremath{\cdot}} & \textcolor{gray}{\ensuremath{\cdot}}\\[3pt]\hline
\BenchmarkSourceLabel{bench:aws-4-k15-w10:row:mps-simdref-simd-block}{R49}\shortstack[l]{{\color{gray}R49}} & \begin{tabular}[t]{@{}l@{}}\rule{0pt}{12pt}MPS / SIMD-ref\\SIMD-DP / Block\end{tabular} & \hspace*{1pt}\begin{tabular}[t]{@{}c@{}}\rule{0pt}{12pt}15.4s\\[-.2pt]{\fontsize{5.8}{6.4}\selectfont \textcolor{ForestGreen}{$ \blacktriangle\,2.3\times $}}\\[-.2pt]14.5s\\[-.2pt]{\fontsize{5.8}{6.4}\selectfont \textcolor{ForestGreen}{$ \blacktriangle\,2.4\times $}}\end{tabular} & \hspace*{1pt}\begin{tabular}[t]{@{}c@{}}\rule{0pt}{12pt}15.1s\\[-.2pt]{\fontsize{5.8}{6.4}\selectfont \textcolor{ForestGreen}{$ \blacktriangle\,2.5\times $}}\\[-.2pt]14.2s\\[-.2pt]{\fontsize{5.8}{6.4}\selectfont \textcolor{ForestGreen}{$ \blacktriangle\,2.5\times $}}\end{tabular} & \textcolor{gray}{\ensuremath{\cdot}} & \textcolor{gray}{\ensuremath{\cdot}} & \hspace*{1pt}\begin{tabular}[t]{@{}c@{}}\rule{0pt}{12pt}6m\,17.1s\\[-.2pt]{\fontsize{5.8}{6.4}\selectfont \textcolor{BrickRed}{$ \blacktriangledown\,11.9\times $}}\\[-.2pt]42.5s\\[-.2pt]{\fontsize{5.8}{6.4}\selectfont \textcolor{BrickRed}{$ \blacktriangledown\,1.4\times $}}\end{tabular} & \hspace*{1pt}\begin{tabular}[t]{@{}c@{}}\rule{0pt}{12pt}5m\,21.5s\\[-.2pt]{\fontsize{5.8}{6.4}\selectfont \textcolor{BrickRed}{$ \blacktriangledown\,9.4\times $}}\\[-.2pt]37.9s\\[-.2pt]{\fontsize{5.8}{6.4}\selectfont \textcolor{BrickRed}{$ \blacktriangledown\,18\% $}}\end{tabular} & \textcolor{gray}{\ensuremath{\cdot}} & \textcolor{gray}{\ensuremath{\cdot}} & \textcolor{gray}{\ensuremath{\cdot}} & \textcolor{gray}{\ensuremath{\cdot}} & \hspace*{1pt}\begin{tabular}[t]{@{}c@{}}\rule{0pt}{12pt}8m\,12.5s\\[-.2pt]{\fontsize{5.8}{6.4}\selectfont \textcolor{BrickRed}{$ \blacktriangledown\,1.5\times $}}\\[-.2pt]7m\,55.2s\\[-.2pt]{\fontsize{5.8}{6.4}\selectfont \textcolor{BrickRed}{$ \blacktriangledown\,1.4\times $}}\end{tabular} & \textcolor{gray}{\ensuremath{\cdot}} & \textcolor{gray}{\ensuremath{\cdot}} & \textcolor{gray}{\ensuremath{\cdot}} & \textcolor{gray}{\ensuremath{\cdot}} & \textcolor{gray}{\ensuremath{\cdot}}\\[3pt]\hline
\BenchmarkSourceLabel{bench:aws-4-k15-w10:row:vulkan-unavailable}{R50}\shortstack[l]{{\color{gray}R50}\\\rule{0pt}{12pt}\textbf{BioseqDB}} & Vulkan (all modes) & \textcolor{gray}{---\textsuperscript{g}} & \textcolor{gray}{---\textsuperscript{g}} & \textcolor{gray}{---\textsuperscript{g}} & \textcolor{gray}{---\textsuperscript{g}} & \textcolor{gray}{---\textsuperscript{g}} & \textcolor{gray}{---\textsuperscript{g}} & \textcolor{gray}{---\textsuperscript{g}} & \textcolor{gray}{---\textsuperscript{g}} & \textcolor{gray}{---\textsuperscript{g}} & \textcolor{gray}{---\textsuperscript{g}} & \textcolor{gray}{---\textsuperscript{g}} & \textcolor{gray}{---\textsuperscript{g}} & \textcolor{gray}{---\textsuperscript{g}} & \textcolor{gray}{---\textsuperscript{g}} & \textcolor{gray}{---\textsuperscript{g}} & \textcolor{gray}{---\textsuperscript{g}}\\[3pt]\hline
\end{tabular}\par}
\smallskip{\footnotesize\noindent Each cell: A and B times, $\cdot$: no measurement. ---\textsuperscript{g}: GPU unavailable. n/a: no comparison.\par
\noindent Blue: one lowest A and one lowest B measurement for each column for each configuration.\par
\noindent Gray: minimap2. Yellow: my preferred MPF / Bounded / SIMD-DP / Block configuration.\par
\noindent MPF/MPS: minimizers per fragment/sequence. SIMD-ref: naive SIMD scan. Bounded: bounded SIMD algorithm.\par
}
\clearpage
\refstepcounter{table}\label{tab:bench:aws-4-k20-w15}
\noindent{\large\bfseries Table \thetable: (C): AWS: 4 Gbp, k=20, w=15\hfill 1/7}\par\smallskip
{\fontsize{7.1}{8.1}\selectfont\setlength{\tabcolsep}{1pt}\setlength{\fboxsep}{2pt}\renewcommand{\arraystretch}{1.12}
\noindent\begin{tabular}{@{}p{14mm}p{27mm}|p{13.55mm}|p{13.55mm}|p{13.55mm}|p{13.55mm}|p{13.55mm}|p{13.55mm}|p{13.55mm}|p{13.55mm}|p{13.55mm}|p{13.55mm}|p{13.55mm}|p{13.55mm}|p{13.55mm}|p{13.55mm}|p{13.55mm}|p{13.55mm}|@{}}\hline
Engine & Execution & \multicolumn{8}{c|}{\textbf{Index}} & \multicolumn{8}{c|}{\textbf{No index}}\\
 &  & \multicolumn{8}{c|}{\textbf{4 Gbp}} & \multicolumn{8}{c|}{\textbf{4 Gbp}}\\
 &  & \multicolumn{8}{c|}{\textbf{k=20, w=15}} & \multicolumn{8}{c|}{\textbf{k=20, w=15}}\\
 &  & \multicolumn{4}{c|}{\textbf{DNA4}} & \multicolumn{4}{c|}{\textbf{DNA5}} & \multicolumn{4}{c|}{\textbf{DNA4}} & \multicolumn{4}{c|}{\textbf{DNA5}}\\
 &  & \multicolumn{1}{c|}{\textbf{M7A}} & \multicolumn{1}{c|}{\textbf{M7I}} & \multicolumn{1}{c|}{\textbf{M4}} & \multicolumn{1}{c|}{\textbf{C7G}} & \multicolumn{1}{c|}{\textbf{M7A}} & \multicolumn{1}{c|}{\textbf{M7I}} & \multicolumn{1}{c|}{\textbf{M4}} & \multicolumn{1}{c|}{\textbf{C7G}} & \multicolumn{1}{c|}{\textbf{M7A}} & \multicolumn{1}{c|}{\textbf{M7I}} & \multicolumn{1}{c|}{\textbf{M4}} & \multicolumn{1}{c|}{\textbf{C7G}} & \multicolumn{1}{c|}{\textbf{M7A}} & \multicolumn{1}{c|}{\textbf{M7I}} & \multicolumn{1}{c|}{\textbf{M4}} & \multicolumn{1}{c|}{\textbf{C7G}}\\
 & & \BenchmarkSourceLabel{bench:aws-4-k20-w15:column:index-4-k20-w15-dna4-m7a}{C1}{\color{gray}C1} & \BenchmarkSourceLabel{bench:aws-4-k20-w15:column:index-4-k20-w15-dna4-m7i}{C2}{\color{gray}C2} & \BenchmarkSourceLabel{bench:aws-4-k20-w15:column:index-4-k20-w15-dna4-m4}{C3}{\color{gray}C3} & \BenchmarkSourceLabel{bench:aws-4-k20-w15:column:index-4-k20-w15-dna4-c7g}{C4}{\color{gray}C4} & \BenchmarkSourceLabel{bench:aws-4-k20-w15:column:index-4-k20-w15-dna5-m7a}{C5}{\color{gray}C5} & \BenchmarkSourceLabel{bench:aws-4-k20-w15:column:index-4-k20-w15-dna5-m7i}{C6}{\color{gray}C6} & \BenchmarkSourceLabel{bench:aws-4-k20-w15:column:index-4-k20-w15-dna5-m4}{C7}{\color{gray}C7} & \BenchmarkSourceLabel{bench:aws-4-k20-w15:column:index-4-k20-w15-dna5-c7g}{C8}{\color{gray}C8} & \BenchmarkSourceLabel{bench:aws-4-k20-w15:column:noindex-4-k20-w15-dna4-m7a}{C9}{\color{gray}C9} & \BenchmarkSourceLabel{bench:aws-4-k20-w15:column:noindex-4-k20-w15-dna4-m7i}{C10}{\color{gray}C10} & \BenchmarkSourceLabel{bench:aws-4-k20-w15:column:noindex-4-k20-w15-dna4-m4}{C11}{\color{gray}C11} & \BenchmarkSourceLabel{bench:aws-4-k20-w15:column:noindex-4-k20-w15-dna4-c7g}{C12}{\color{gray}C12} & \BenchmarkSourceLabel{bench:aws-4-k20-w15:column:noindex-4-k20-w15-dna5-m7a}{C13}{\color{gray}C13} & \BenchmarkSourceLabel{bench:aws-4-k20-w15:column:noindex-4-k20-w15-dna5-m7i}{C14}{\color{gray}C14} & \BenchmarkSourceLabel{bench:aws-4-k20-w15:column:noindex-4-k20-w15-dna5-m4}{C15}{\color{gray}C15} & \BenchmarkSourceLabel{bench:aws-4-k20-w15:column:noindex-4-k20-w15-dna5-c7g}{C16}{\color{gray}C16}\\
\hline
\rowcolor{BenchBaseline}[0.5pt][0.5pt]\BenchmarkSourceLabel{bench:aws-4-k20-w15:row:minimap2}{R1}\shortstack[l]{{\color{gray}R1}\\\textbf{minimap2}} & \textbf{Baseline} & \hspace*{1pt}\begin{tabular}[t]{@{}c@{}}\rule{0pt}{12pt}2m\,48.1s\\[-.2pt]{\fontsize{5.8}{6.4}\selectfont \textcolor{BrickRed}{$ \blacktriangledown\,20.4\times $}}\\[-.2pt]2m\,41.7s\\[-.2pt]{\fontsize{5.8}{6.4}\selectfont \textcolor{BrickRed}{$ \blacktriangledown\,22.6\times $}}\end{tabular} & \hspace*{1pt}\begin{tabular}[t]{@{}c@{}}\rule{0pt}{12pt}2m\,50.3s\\[-.2pt]{\fontsize{5.8}{6.4}\selectfont \textcolor{BrickRed}{$ \blacktriangledown\,19.8\times $}}\\[-.2pt]2m\,43s\\[-.2pt]{\fontsize{5.8}{6.4}\selectfont \textcolor{BrickRed}{$ \blacktriangledown\,22.2\times $}}\end{tabular} & \hspace*{1pt}\begin{tabular}[t]{@{}c@{}}\rule{0pt}{12pt}4m\,54.3s\\[-.2pt]{\fontsize{5.8}{6.4}\selectfont \textcolor{BrickRed}{$ \blacktriangledown\,43.2\times $}}\\[-.2pt]4m\,39s\\[-.2pt]{\fontsize{5.8}{6.4}\selectfont \textcolor{BrickRed}{$ \blacktriangledown\,54.6\times $}}\end{tabular} & \hspace*{1pt}\begin{tabular}[t]{@{}c@{}}\rule{0pt}{12pt}25m\,5.5s\\[-.2pt]{\fontsize{5.8}{6.4}\selectfont \textcolor{BrickRed}{$ \blacktriangledown\,41.8\times $}}\\[-.2pt]25m\,0.7s\\[-.2pt]{\fontsize{5.8}{6.4}\selectfont \textcolor{BrickRed}{$ \blacktriangledown\,46.6\times $}}\end{tabular} & \hspace*{1pt}\begin{tabular}[t]{@{}c@{}}\rule{0pt}{12pt}33.9s\\[-.2pt]{\fontsize{5.8}{6.4}\selectfont \textcolor{BrickRed}{$ \blacktriangledown\,2.3\times $}}\\[-.2pt]32.9s\\[-.2pt]{\fontsize{5.8}{6.4}\selectfont \textcolor{BrickRed}{$ \blacktriangledown\,3.1\times $}}\end{tabular} & \hspace*{1pt}\begin{tabular}[t]{@{}c@{}}\rule{0pt}{12pt}34.4s\\[-.2pt]{\fontsize{5.8}{6.4}\selectfont \textcolor{BrickRed}{$ \blacktriangledown\,2.6\times $}}\\[-.2pt]32.7s\\[-.2pt]{\fontsize{5.8}{6.4}\selectfont \textcolor{BrickRed}{$ \blacktriangledown\,2.9\times $}}\end{tabular} & \hspace*{1pt}\begin{tabular}[t]{@{}c@{}}\rule{0pt}{12pt}1m\,8.2s\\[-.2pt]{\fontsize{5.8}{6.4}\selectfont \textcolor{BrickRed}{$ \blacktriangledown\,3.7\times $}}\\[-.2pt]1m\,5.1s\\[-.2pt]{\fontsize{5.8}{6.4}\selectfont \textcolor{BrickRed}{$ \blacktriangledown\,4.6\times $}}\end{tabular} & \hspace*{1pt}\begin{tabular}[t]{@{}c@{}}\rule{0pt}{12pt}3m\,9.8s\\[-.2pt]{\fontsize{5.8}{6.4}\selectfont \textcolor{BrickRed}{$ \blacktriangledown\,3.9\times $}}\\[-.2pt]3m\,3.7s\\[-.2pt]{\fontsize{5.8}{6.4}\selectfont \textcolor{BrickRed}{$ \blacktriangledown\,4.1\times $}}\end{tabular} & \hspace*{1pt}\begin{tabular}[t]{@{}c@{}}\rule{0pt}{12pt}3m\,28.3s\\[-.2pt]{\fontsize{5.8}{6.4}\selectfont \textcolor{BrickRed}{$ \blacktriangledown\,1.3\times $}}\\[-.2pt]3m\,22.8s\\[-.2pt]{\fontsize{5.8}{6.4}\selectfont \textcolor{BrickRed}{$ \blacktriangledown\,1.4\times $}}\end{tabular} & \hspace*{1pt}\begin{tabular}[t]{@{}c@{}}\rule{0pt}{12pt}3m\,53.2s\\[-.2pt]{\fontsize{5.8}{6.4}\selectfont \textcolor{BrickRed}{$ \blacktriangledown\,1.4\times $}}\\[-.2pt]3m\,43.5s\\[-.2pt]{\fontsize{5.8}{6.4}\selectfont \textcolor{BrickRed}{$ \blacktriangledown\,1.5\times $}}\end{tabular} & \hspace*{1pt}\begin{tabular}[t]{@{}c@{}}\rule{0pt}{12pt}7m\,14.6s\\[-.2pt]{\fontsize{5.8}{6.4}\selectfont \textcolor{BrickRed}{$ \blacktriangledown\,1.5\times $}}\\[-.2pt]6m\,57.2s\\[-.2pt]{\fontsize{5.8}{6.4}\selectfont \textcolor{BrickRed}{$ \blacktriangledown\,1.6\times $}}\end{tabular} & \hspace*{1pt}\begin{tabular}[t]{@{}c@{}}\rule{0pt}{12pt}\colorbox{BenchBest}{\strut 51m\,22s}\\[-.2pt]{\fontsize{5.8}{6.4}\selectfont \textcolor{gray}{n/a}}\\[-.2pt]\colorbox{BenchBest}{\strut 51m\,17.2s}\\[-.2pt]{\fontsize{5.8}{6.4}\selectfont \textcolor{gray}{n/a}}\end{tabular} & \hspace*{1pt}\begin{tabular}[t]{@{}c@{}}\rule{0pt}{12pt}\colorbox{BenchBest}{\strut 1m\,10.3s}\\[-.2pt]{\fontsize{5.8}{6.4}\selectfont \textcolor{ForestGreen}{$ \blacktriangle\,7.6\times $}}\\[-.2pt]\colorbox{BenchBest}{\strut 1m\,9s}\\[-.2pt]{\fontsize{5.8}{6.4}\selectfont \textcolor{ForestGreen}{$ \blacktriangle\,2.4\times $}}\end{tabular} & \hspace*{1pt}\begin{tabular}[t]{@{}c@{}}\rule{0pt}{12pt}\colorbox{BenchBest}{\strut 1m\,8.6s}\\[-.2pt]{\fontsize{5.8}{6.4}\selectfont \textcolor{ForestGreen}{$ \blacktriangle\,6.4\times $}}\\[-.2pt]\colorbox{BenchBest}{\strut 1m\,6.9s}\\[-.2pt]{\fontsize{5.8}{6.4}\selectfont \textcolor{ForestGreen}{$ \blacktriangle\,2.3\times $}}\end{tabular} & \hspace*{1pt}\begin{tabular}[t]{@{}c@{}}\rule{0pt}{12pt}\colorbox{BenchBest}{\strut 2m\,21.5s}\\[-.2pt]{\fontsize{5.8}{6.4}\selectfont \textcolor{ForestGreen}{$ \blacktriangle\,6.0\times $}}\\[-.2pt]\colorbox{BenchBest}{\strut 2m\,18.4s}\\[-.2pt]{\fontsize{5.8}{6.4}\selectfont \textcolor{ForestGreen}{$ \blacktriangle\,2.2\times $}}\end{tabular} & \hspace*{1pt}\begin{tabular}[t]{@{}c@{}}\rule{0pt}{12pt}\colorbox{BenchBest}{\strut 15m\,56.8s}\\[-.2pt]{\fontsize{5.8}{6.4}\selectfont \textcolor{gray}{n/a}}\\[-.2pt]\colorbox{BenchBest}{\strut 15m\,51.9s}\\[-.2pt]{\fontsize{5.8}{6.4}\selectfont \textcolor{gray}{n/a}}\end{tabular}\\[3pt]\hline
\rowcolor{BenchPreferred}[0.5pt][0.5pt]\BenchmarkSourceLabel{bench:aws-4-k20-w15:row:mpf-bounded-simd-block}{R2}\shortstack[l]{{\color{gray}R2}\\\textbf{BioseqDB}} & \begin{tabular}[t]{@{}l@{}}\rule{0pt}{12pt}MPF / Bounded\\SIMD-DP / Block\end{tabular} & \hspace*{1pt}\begin{tabular}[t]{@{}c@{}}\rule{0pt}{12pt}9.2s\\[-.2pt]{\fontsize{5.8}{6.4}\selectfont \textcolor{ForestGreen}{$ \blacktriangle\,18.3\times $}}\\[-.2pt]8.5s\\[-.2pt]{\fontsize{5.8}{6.4}\selectfont \textcolor{ForestGreen}{$ \blacktriangle\,19.0\times $}}\end{tabular} & \hspace*{1pt}\begin{tabular}[t]{@{}c@{}}\rule{0pt}{12pt}11.4s\\[-.2pt]{\fontsize{5.8}{6.4}\selectfont \textcolor{ForestGreen}{$ \blacktriangle\,15.0\times $}}\\[-.2pt]10.7s\\[-.2pt]{\fontsize{5.8}{6.4}\selectfont \textcolor{ForestGreen}{$ \blacktriangle\,15.3\times $}}\end{tabular} & \hspace*{1pt}\begin{tabular}[t]{@{}c@{}}\rule{0pt}{12pt}\colorbox{BenchBest}{\strut 6.8s}\\[-.2pt]{\fontsize{5.8}{6.4}\selectfont \textcolor{ForestGreen}{$ \blacktriangle\,43.2\times $}}\\[-.2pt]\colorbox{BenchBest}{\strut 5.1s}\\[-.2pt]{\fontsize{5.8}{6.4}\selectfont \textcolor{ForestGreen}{$ \blacktriangle\,54.6\times $}}\end{tabular} & \hspace*{1pt}\begin{tabular}[t]{@{}c@{}}\rule{0pt}{12pt}39.1s\\[-.2pt]{\fontsize{5.8}{6.4}\selectfont \textcolor{ForestGreen}{$ \blacktriangle\,38.5\times $}}\\[-.2pt]38s\\[-.2pt]{\fontsize{5.8}{6.4}\selectfont \textcolor{ForestGreen}{$ \blacktriangle\,39.5\times $}}\end{tabular} & \hspace*{1pt}\begin{tabular}[t]{@{}c@{}}\rule{0pt}{12pt}18.2s\\[-.2pt]{\fontsize{5.8}{6.4}\selectfont \textcolor{ForestGreen}{$ \blacktriangle\,1.9\times $}}\\[-.2pt]16s\\[-.2pt]{\fontsize{5.8}{6.4}\selectfont \textcolor{ForestGreen}{$ \blacktriangle\,2.1\times $}}\end{tabular} & \hspace*{1pt}\begin{tabular}[t]{@{}c@{}}\rule{0pt}{12pt}\colorbox{BenchBest}{\strut 13.2s}\\[-.2pt]{\fontsize{5.8}{6.4}\selectfont \textcolor{ForestGreen}{$ \blacktriangle\,2.6\times $}}\\[-.2pt]\colorbox{BenchBest}{\strut 11.4s}\\[-.2pt]{\fontsize{5.8}{6.4}\selectfont \textcolor{ForestGreen}{$ \blacktriangle\,2.9\times $}}\end{tabular} & \hspace*{1pt}\begin{tabular}[t]{@{}c@{}}\rule{0pt}{12pt}\colorbox{BenchBest}{\strut 18.6s}\\[-.2pt]{\fontsize{5.8}{6.4}\selectfont \textcolor{ForestGreen}{$ \blacktriangle\,3.7\times $}}\\[-.2pt]\colorbox{BenchBest}{\strut 14.3s}\\[-.2pt]{\fontsize{5.8}{6.4}\selectfont \textcolor{ForestGreen}{$ \blacktriangle\,4.6\times $}}\end{tabular} & \hspace*{1pt}\begin{tabular}[t]{@{}c@{}}\rule{0pt}{12pt}57.2s\\[-.2pt]{\fontsize{5.8}{6.4}\selectfont \textcolor{ForestGreen}{$ \blacktriangle\,3.3\times $}}\\[-.2pt]53.9s\\[-.2pt]{\fontsize{5.8}{6.4}\selectfont \textcolor{ForestGreen}{$ \blacktriangle\,3.4\times $}}\end{tabular} & \hspace*{1pt}\begin{tabular}[t]{@{}c@{}}\rule{0pt}{12pt}\colorbox{BenchBest}{\strut 2m\,38.5s}\\[-.2pt]{\fontsize{5.8}{6.4}\selectfont \textcolor{ForestGreen}{$ \blacktriangle\,24\% $}}\\[-.2pt]\colorbox{BenchBest}{\strut 2m\,29.8s}\\[-.2pt]{\fontsize{5.8}{6.4}\selectfont \textcolor{ForestGreen}{$ \blacktriangle\,26\% $}}\end{tabular} & \hspace*{1pt}\begin{tabular}[t]{@{}c@{}}\rule{0pt}{12pt}2m\,49s\\[-.2pt]{\fontsize{5.8}{6.4}\selectfont \textcolor{ForestGreen}{$ \blacktriangle\,28\% $}}\\[-.2pt]2m\,37s\\[-.2pt]{\fontsize{5.8}{6.4}\selectfont \textcolor{ForestGreen}{$ \blacktriangle\,30\% $}}\end{tabular} & \hspace*{1pt}\begin{tabular}[t]{@{}c@{}}\rule{0pt}{12pt}5m\,3.3s\\[-.2pt]{\fontsize{5.8}{6.4}\selectfont \textcolor{ForestGreen}{$ \blacktriangle\,1.4\times $}}\\[-.2pt]4m\,42.7s\\[-.2pt]{\fontsize{5.8}{6.4}\selectfont \textcolor{ForestGreen}{$ \blacktriangle\,1.5\times $}}\end{tabular} & \textcolor{gray}{\ensuremath{\cdot}} & \hspace*{1pt}\begin{tabular}[t]{@{}c@{}}\rule{0pt}{12pt}8m\,52.3s\\[-.2pt]{\fontsize{5.8}{6.4}\selectfont \textcolor{BrickRed}{$ \blacktriangledown\,7.6\times $}}\\[-.2pt]2m\,45.5s\\[-.2pt]{\fontsize{5.8}{6.4}\selectfont \textcolor{BrickRed}{$ \blacktriangledown\,2.4\times $}}\end{tabular} & \hspace*{1pt}\begin{tabular}[t]{@{}c@{}}\rule{0pt}{12pt}7m\,22.2s\\[-.2pt]{\fontsize{5.8}{6.4}\selectfont \textcolor{BrickRed}{$ \blacktriangledown\,6.4\times $}}\\[-.2pt]2m\,48.2s\\[-.2pt]{\fontsize{5.8}{6.4}\selectfont \textcolor{BrickRed}{$ \blacktriangledown\,2.5\times $}}\end{tabular} & \hspace*{1pt}\begin{tabular}[t]{@{}c@{}}\rule{0pt}{12pt}14m\,43.3s\\[-.2pt]{\fontsize{5.8}{6.4}\selectfont \textcolor{BrickRed}{$ \blacktriangledown\,6.2\times $}}\\[-.2pt]5m\,9.5s\\[-.2pt]{\fontsize{5.8}{6.4}\selectfont \textcolor{BrickRed}{$ \blacktriangledown\,2.2\times $}}\end{tabular} & \textcolor{gray}{\ensuremath{\cdot}}\\[3pt]\hline
\BenchmarkSourceLabel{bench:aws-4-k20-w15:row:mpf-bounded-scalar-block}{R3}\shortstack[l]{{\color{gray}R3}} & \begin{tabular}[t]{@{}l@{}}\rule{0pt}{12pt}MPF / Bounded\\Scalar-DP / Block\end{tabular} & \hspace*{1pt}\begin{tabular}[t]{@{}c@{}}\rule{0pt}{12pt}10.7s\\[-.2pt]{\fontsize{5.8}{6.4}\selectfont \textcolor{ForestGreen}{$ \blacktriangle\,15.7\times $}}\\[-.2pt]10.1s\\[-.2pt]{\fontsize{5.8}{6.4}\selectfont \textcolor{ForestGreen}{$ \blacktriangle\,16.1\times $}}\end{tabular} & \hspace*{1pt}\begin{tabular}[t]{@{}c@{}}\rule{0pt}{12pt}11.5s\\[-.2pt]{\fontsize{5.8}{6.4}\selectfont \textcolor{ForestGreen}{$ \blacktriangle\,14.8\times $}}\\[-.2pt]10.8s\\[-.2pt]{\fontsize{5.8}{6.4}\selectfont \textcolor{ForestGreen}{$ \blacktriangle\,15.2\times $}}\end{tabular} & \hspace*{1pt}\begin{tabular}[t]{@{}c@{}}\rule{0pt}{12pt}7.6s\\[-.2pt]{\fontsize{5.8}{6.4}\selectfont \textcolor{ForestGreen}{$ \blacktriangle\,38.5\times $}}\\[-.2pt]5.4s\\[-.2pt]{\fontsize{5.8}{6.4}\selectfont \textcolor{ForestGreen}{$ \blacktriangle\,51.4\times $}}\end{tabular} & \hspace*{1pt}\begin{tabular}[t]{@{}c@{}}\rule{0pt}{12pt}36.6s\\[-.2pt]{\fontsize{5.8}{6.4}\selectfont \textcolor{ForestGreen}{$ \blacktriangle\,41.1\times $}}\\[-.2pt]35.5s\\[-.2pt]{\fontsize{5.8}{6.4}\selectfont \textcolor{ForestGreen}{$ \blacktriangle\,42.2\times $}}\end{tabular} & \hspace*{1pt}\begin{tabular}[t]{@{}c@{}}\rule{0pt}{12pt}17s\\[-.2pt]{\fontsize{5.8}{6.4}\selectfont \textcolor{ForestGreen}{$ \blacktriangle\,2.0\times $}}\\[-.2pt]12.8s\\[-.2pt]{\fontsize{5.8}{6.4}\selectfont \textcolor{ForestGreen}{$ \blacktriangle\,2.6\times $}}\end{tabular} & \hspace*{1pt}\begin{tabular}[t]{@{}c@{}}\rule{0pt}{12pt}17.3s\\[-.2pt]{\fontsize{5.8}{6.4}\selectfont \textcolor{ForestGreen}{$ \blacktriangle\,2.0\times $}}\\[-.2pt]14.1s\\[-.2pt]{\fontsize{5.8}{6.4}\selectfont \textcolor{ForestGreen}{$ \blacktriangle\,2.3\times $}}\end{tabular} & \hspace*{1pt}\begin{tabular}[t]{@{}c@{}}\rule{0pt}{12pt}23s\\[-.2pt]{\fontsize{5.8}{6.4}\selectfont \textcolor{ForestGreen}{$ \blacktriangle\,3.0\times $}}\\[-.2pt]15.8s\\[-.2pt]{\fontsize{5.8}{6.4}\selectfont \textcolor{ForestGreen}{$ \blacktriangle\,4.1\times $}}\end{tabular} & \hspace*{1pt}\begin{tabular}[t]{@{}c@{}}\rule{0pt}{12pt}51.2s\\[-.2pt]{\fontsize{5.8}{6.4}\selectfont \textcolor{ForestGreen}{$ \blacktriangle\,3.7\times $}}\\[-.2pt]47.8s\\[-.2pt]{\fontsize{5.8}{6.4}\selectfont \textcolor{ForestGreen}{$ \blacktriangle\,3.8\times $}}\end{tabular} & \hspace*{1pt}\begin{tabular}[t]{@{}c@{}}\rule{0pt}{12pt}2m\,50s\\[-.2pt]{\fontsize{5.8}{6.4}\selectfont \textcolor{ForestGreen}{$ \blacktriangle\,18\% $}}\\[-.2pt]2m\,41s\\[-.2pt]{\fontsize{5.8}{6.4}\selectfont \textcolor{ForestGreen}{$ \blacktriangle\,21\% $}}\end{tabular} & \hspace*{1pt}\begin{tabular}[t]{@{}c@{}}\rule{0pt}{12pt}\colorbox{BenchBest}{\strut 2m\,48.3s}\\[-.2pt]{\fontsize{5.8}{6.4}\selectfont \textcolor{ForestGreen}{$ \blacktriangle\,28\% $}}\\[-.2pt]2m\,35.9s\\[-.2pt]{\fontsize{5.8}{6.4}\selectfont \textcolor{ForestGreen}{$ \blacktriangle\,1.4\times $}}\end{tabular} & \hspace*{1pt}\begin{tabular}[t]{@{}c@{}}\rule{0pt}{12pt}\colorbox{BenchBest}{\strut 4m\,47.5s}\\[-.2pt]{\fontsize{5.8}{6.4}\selectfont \textcolor{ForestGreen}{$ \blacktriangle\,1.5\times $}}\\[-.2pt]\colorbox{BenchBest}{\strut 4m\,26.9s}\\[-.2pt]{\fontsize{5.8}{6.4}\selectfont \textcolor{ForestGreen}{$ \blacktriangle\,1.6\times $}}\end{tabular} & \textcolor{gray}{\ensuremath{\cdot}} & \hspace*{1pt}\begin{tabular}[t]{@{}c@{}}\rule{0pt}{12pt}9m\,21.4s\\[-.2pt]{\fontsize{5.8}{6.4}\selectfont \textcolor{BrickRed}{$ \blacktriangledown\,8.0\times $}}\\[-.2pt]3m\,11.2s\\[-.2pt]{\fontsize{5.8}{6.4}\selectfont \textcolor{BrickRed}{$ \blacktriangledown\,2.8\times $}}\end{tabular} & \hspace*{1pt}\begin{tabular}[t]{@{}c@{}}\rule{0pt}{12pt}7m\,45.7s\\[-.2pt]{\fontsize{5.8}{6.4}\selectfont \textcolor{BrickRed}{$ \blacktriangledown\,6.8\times $}}\\[-.2pt]2m\,38.2s\\[-.2pt]{\fontsize{5.8}{6.4}\selectfont \textcolor{BrickRed}{$ \blacktriangledown\,2.4\times $}}\end{tabular} & \textcolor{gray}{\ensuremath{\cdot}} & \textcolor{gray}{\ensuremath{\cdot}}\\[3pt]\hline
\BenchmarkSourceLabel{bench:aws-4-k20-w15:row:mpf-bounded-rmq-block}{R4}\shortstack[l]{{\color{gray}R4}} & \begin{tabular}[t]{@{}l@{}}\rule{0pt}{12pt}MPF / Bounded\\RMQ / Block\end{tabular} & \hspace*{1pt}\begin{tabular}[t]{@{}c@{}}\rule{0pt}{12pt}\colorbox{BenchBest}{\strut 8.3s}\\[-.2pt]{\fontsize{5.8}{6.4}\selectfont \textcolor{ForestGreen}{$ \blacktriangle\,20.4\times $}}\\[-.2pt]7.6s\\[-.2pt]{\fontsize{5.8}{6.4}\selectfont \textcolor{ForestGreen}{$ \blacktriangle\,21.3\times $}}\end{tabular} & \hspace*{1pt}\begin{tabular}[t]{@{}c@{}}\rule{0pt}{12pt}\colorbox{BenchBest}{\strut 8.6s}\\[-.2pt]{\fontsize{5.8}{6.4}\selectfont \textcolor{ForestGreen}{$ \blacktriangle\,19.8\times $}}\\[-.2pt]7.8s\\[-.2pt]{\fontsize{5.8}{6.4}\selectfont \textcolor{ForestGreen}{$ \blacktriangle\,20.8\times $}}\end{tabular} & \hspace*{1pt}\begin{tabular}[t]{@{}c@{}}\rule{0pt}{12pt}35s\\[-.2pt]{\fontsize{5.8}{6.4}\selectfont \textcolor{ForestGreen}{$ \blacktriangle\,8.4\times $}}\\[-.2pt]32.9s\\[-.2pt]{\fontsize{5.8}{6.4}\selectfont \textcolor{ForestGreen}{$ \blacktriangle\,8.5\times $}}\end{tabular} & \hspace*{1pt}\begin{tabular}[t]{@{}c@{}}\rule{0pt}{12pt}38.3s\\[-.2pt]{\fontsize{5.8}{6.4}\selectfont \textcolor{ForestGreen}{$ \blacktriangle\,39.3\times $}}\\[-.2pt]37.2s\\[-.2pt]{\fontsize{5.8}{6.4}\selectfont \textcolor{ForestGreen}{$ \blacktriangle\,40.3\times $}}\end{tabular} & \hspace*{1pt}\begin{tabular}[t]{@{}c@{}}\rule{0pt}{12pt}35.9s\\[-.2pt]{\fontsize{5.8}{6.4}\selectfont \textcolor{BrickRed}{$ \blacktriangledown\,6\% $}}\\[-.2pt]31.8s\\[-.2pt]{\fontsize{5.8}{6.4}\selectfont \textcolor{gray}{$ \blacktriangleright\,3\% $}}\end{tabular} & \hspace*{1pt}\begin{tabular}[t]{@{}c@{}}\rule{0pt}{12pt}33.2s\\[-.2pt]{\fontsize{5.8}{6.4}\selectfont \textcolor{gray}{$ \blacktriangleright\,4\% $}}\\[-.2pt]30.1s\\[-.2pt]{\fontsize{5.8}{6.4}\selectfont \textcolor{ForestGreen}{$ \blacktriangle\,8\% $}}\end{tabular} & \hspace*{1pt}\begin{tabular}[t]{@{}c@{}}\rule{0pt}{12pt}46.2s\\[-.2pt]{\fontsize{5.8}{6.4}\selectfont \textcolor{ForestGreen}{$ \blacktriangle\,1.5\times $}}\\[-.2pt]39.4s\\[-.2pt]{\fontsize{5.8}{6.4}\selectfont \textcolor{ForestGreen}{$ \blacktriangle\,1.7\times $}}\end{tabular} & \hspace*{1pt}\begin{tabular}[t]{@{}c@{}}\rule{0pt}{12pt}49.2s\\[-.2pt]{\fontsize{5.8}{6.4}\selectfont \textcolor{ForestGreen}{$ \blacktriangle\,3.9\times $}}\\[-.2pt]45.8s\\[-.2pt]{\fontsize{5.8}{6.4}\selectfont \textcolor{ForestGreen}{$ \blacktriangle\,4.0\times $}}\end{tabular} & \hspace*{1pt}\begin{tabular}[t]{@{}c@{}}\rule{0pt}{12pt}2m\,53.8s\\[-.2pt]{\fontsize{5.8}{6.4}\selectfont \textcolor{ForestGreen}{$ \blacktriangle\,17\% $}}\\[-.2pt]2m\,43.7s\\[-.2pt]{\fontsize{5.8}{6.4}\selectfont \textcolor{ForestGreen}{$ \blacktriangle\,19\% $}}\end{tabular} & \hspace*{1pt}\begin{tabular}[t]{@{}c@{}}\rule{0pt}{12pt}3m\,4.5s\\[-.2pt]{\fontsize{5.8}{6.4}\selectfont \textcolor{ForestGreen}{$ \blacktriangle\,21\% $}}\\[-.2pt]2m\,52.1s\\[-.2pt]{\fontsize{5.8}{6.4}\selectfont \textcolor{ForestGreen}{$ \blacktriangle\,23\% $}}\end{tabular} & \textcolor{gray}{\ensuremath{\cdot}} & \textcolor{gray}{\ensuremath{\cdot}} & \hspace*{1pt}\begin{tabular}[t]{@{}c@{}}\rule{0pt}{12pt}9m\,6.6s\\[-.2pt]{\fontsize{5.8}{6.4}\selectfont \textcolor{BrickRed}{$ \blacktriangledown\,7.8\times $}}\\[-.2pt]2m\,56.5s\\[-.2pt]{\fontsize{5.8}{6.4}\selectfont \textcolor{BrickRed}{$ \blacktriangledown\,2.6\times $}}\end{tabular} & \hspace*{1pt}\begin{tabular}[t]{@{}c@{}}\rule{0pt}{12pt}7m\,59.6s\\[-.2pt]{\fontsize{5.8}{6.4}\selectfont \textcolor{BrickRed}{$ \blacktriangledown\,7.0\times $}}\\[-.2pt]2m\,51.2s\\[-.2pt]{\fontsize{5.8}{6.4}\selectfont \textcolor{BrickRed}{$ \blacktriangledown\,2.6\times $}}\end{tabular} & \textcolor{gray}{\ensuremath{\cdot}} & \textcolor{gray}{\ensuremath{\cdot}}\\[3pt]\hline
\BenchmarkSourceLabel{bench:aws-4-k20-w15:row:mpf-bounded-simd-ksw2}{R5}\shortstack[l]{{\color{gray}R5}} & \begin{tabular}[t]{@{}l@{}}\rule{0pt}{12pt}MPF / Bounded\\SIMD-DP / KSW2\end{tabular} & \hspace*{1pt}\begin{tabular}[t]{@{}c@{}}\rule{0pt}{12pt}9.6s\\[-.2pt]{\fontsize{5.8}{6.4}\selectfont \textcolor{ForestGreen}{$ \blacktriangle\,17.6\times $}}\\[-.2pt]\colorbox{BenchBest}{\strut 7.1s}\\[-.2pt]{\fontsize{5.8}{6.4}\selectfont \textcolor{ForestGreen}{$ \blacktriangle\,22.6\times $}}\end{tabular} & \hspace*{1pt}\begin{tabular}[t]{@{}c@{}}\rule{0pt}{12pt}9.6s\\[-.2pt]{\fontsize{5.8}{6.4}\selectfont \textcolor{ForestGreen}{$ \blacktriangle\,17.7\times $}}\\[-.2pt]\colorbox{BenchBest}{\strut 7.3s}\\[-.2pt]{\fontsize{5.8}{6.4}\selectfont \textcolor{ForestGreen}{$ \blacktriangle\,22.2\times $}}\end{tabular} & \hspace*{1pt}\begin{tabular}[t]{@{}c@{}}\rule{0pt}{12pt}38.3s\\[-.2pt]{\fontsize{5.8}{6.4}\selectfont \textcolor{ForestGreen}{$ \blacktriangle\,7.7\times $}}\\[-.2pt]30.9s\\[-.2pt]{\fontsize{5.8}{6.4}\selectfont \textcolor{ForestGreen}{$ \blacktriangle\,9.0\times $}}\end{tabular} & \hspace*{1pt}\begin{tabular}[t]{@{}c@{}}\rule{0pt}{12pt}40.5s\\[-.2pt]{\fontsize{5.8}{6.4}\selectfont \textcolor{ForestGreen}{$ \blacktriangle\,37.2\times $}}\\[-.2pt]36.4s\\[-.2pt]{\fontsize{5.8}{6.4}\selectfont \textcolor{ForestGreen}{$ \blacktriangle\,41.2\times $}}\end{tabular} & \hspace*{1pt}\begin{tabular}[t]{@{}c@{}}\rule{0pt}{12pt}34.6s\\[-.2pt]{\fontsize{5.8}{6.4}\selectfont \textcolor{gray}{$ \blacktriangleright\,2\% $}}\\[-.2pt]30.4s\\[-.2pt]{\fontsize{5.8}{6.4}\selectfont \textcolor{ForestGreen}{$ \blacktriangle\,8\% $}}\end{tabular} & \hspace*{1pt}\begin{tabular}[t]{@{}c@{}}\rule{0pt}{12pt}21.3s\\[-.2pt]{\fontsize{5.8}{6.4}\selectfont \textcolor{ForestGreen}{$ \blacktriangle\,1.6\times $}}\\[-.2pt]17.7s\\[-.2pt]{\fontsize{5.8}{6.4}\selectfont \textcolor{ForestGreen}{$ \blacktriangle\,1.8\times $}}\end{tabular} & \hspace*{1pt}\begin{tabular}[t]{@{}c@{}}\rule{0pt}{12pt}37s\\[-.2pt]{\fontsize{5.8}{6.4}\selectfont \textcolor{ForestGreen}{$ \blacktriangle\,1.8\times $}}\\[-.2pt]29.3s\\[-.2pt]{\fontsize{5.8}{6.4}\selectfont \textcolor{ForestGreen}{$ \blacktriangle\,2.2\times $}}\end{tabular} & \hspace*{1pt}\begin{tabular}[t]{@{}c@{}}\rule{0pt}{12pt}55.9s\\[-.2pt]{\fontsize{5.8}{6.4}\selectfont \textcolor{ForestGreen}{$ \blacktriangle\,3.4\times $}}\\[-.2pt]52s\\[-.2pt]{\fontsize{5.8}{6.4}\selectfont \textcolor{ForestGreen}{$ \blacktriangle\,3.5\times $}}\end{tabular} & \hspace*{1pt}\begin{tabular}[t]{@{}c@{}}\rule{0pt}{12pt}2m\,51.2s\\[-.2pt]{\fontsize{5.8}{6.4}\selectfont \textcolor{ForestGreen}{$ \blacktriangle\,18\% $}}\\[-.2pt]2m\,39.2s\\[-.2pt]{\fontsize{5.8}{6.4}\selectfont \textcolor{ForestGreen}{$ \blacktriangle\,21\% $}}\end{tabular} & \hspace*{1pt}\begin{tabular}[t]{@{}c@{}}\rule{0pt}{12pt}3m\,0s\\[-.2pt]{\fontsize{5.8}{6.4}\selectfont \textcolor{ForestGreen}{$ \blacktriangle\,23\% $}}\\[-.2pt]2m\,45.6s\\[-.2pt]{\fontsize{5.8}{6.4}\selectfont \textcolor{ForestGreen}{$ \blacktriangle\,26\% $}}\end{tabular} & \textcolor{gray}{\ensuremath{\cdot}} & \textcolor{gray}{\ensuremath{\cdot}} & \hspace*{1pt}\begin{tabular}[t]{@{}c@{}}\rule{0pt}{12pt}9m\,4.2s\\[-.2pt]{\fontsize{5.8}{6.4}\selectfont \textcolor{BrickRed}{$ \blacktriangledown\,7.7\times $}}\\[-.2pt]2m\,55.2s\\[-.2pt]{\fontsize{5.8}{6.4}\selectfont \textcolor{BrickRed}{$ \blacktriangledown\,2.5\times $}}\end{tabular} & \hspace*{1pt}\begin{tabular}[t]{@{}c@{}}\rule{0pt}{12pt}7m\,44.4s\\[-.2pt]{\fontsize{5.8}{6.4}\selectfont \textcolor{BrickRed}{$ \blacktriangledown\,6.8\times $}}\\[-.2pt]2m\,37.2s\\[-.2pt]{\fontsize{5.8}{6.4}\selectfont \textcolor{BrickRed}{$ \blacktriangledown\,2.3\times $}}\end{tabular} & \textcolor{gray}{\ensuremath{\cdot}} & \textcolor{gray}{\ensuremath{\cdot}}\\[3pt]\hline
\BenchmarkSourceLabel{bench:aws-4-k20-w15:row:mpf-bounded-rmq-ksw2}{R6}\shortstack[l]{{\color{gray}R6}} & \begin{tabular}[t]{@{}l@{}}\rule{0pt}{12pt}MPF / Bounded\\RMQ / KSW2\end{tabular} & \hspace*{1pt}\begin{tabular}[t]{@{}c@{}}\rule{0pt}{12pt}10.6s\\[-.2pt]{\fontsize{5.8}{6.4}\selectfont \textcolor{ForestGreen}{$ \blacktriangle\,15.9\times $}}\\[-.2pt]8.3s\\[-.2pt]{\fontsize{5.8}{6.4}\selectfont \textcolor{ForestGreen}{$ \blacktriangle\,19.5\times $}}\end{tabular} & \hspace*{1pt}\begin{tabular}[t]{@{}c@{}}\rule{0pt}{12pt}14.6s\\[-.2pt]{\fontsize{5.8}{6.4}\selectfont \textcolor{ForestGreen}{$ \blacktriangle\,11.7\times $}}\\[-.2pt]12.5s\\[-.2pt]{\fontsize{5.8}{6.4}\selectfont \textcolor{ForestGreen}{$ \blacktriangle\,13.0\times $}}\end{tabular} & \hspace*{1pt}\begin{tabular}[t]{@{}c@{}}\rule{0pt}{12pt}12.5s\\[-.2pt]{\fontsize{5.8}{6.4}\selectfont \textcolor{ForestGreen}{$ \blacktriangle\,23.5\times $}}\\[-.2pt]5.1s\\[-.2pt]{\fontsize{5.8}{6.4}\selectfont \textcolor{ForestGreen}{$ \blacktriangle\,54.4\times $}}\end{tabular} & \hspace*{1pt}\begin{tabular}[t]{@{}c@{}}\rule{0pt}{12pt}\colorbox{BenchBest}{\strut 36s}\\[-.2pt]{\fontsize{5.8}{6.4}\selectfont \textcolor{ForestGreen}{$ \blacktriangle\,41.8\times $}}\\[-.2pt]\colorbox{BenchBest}{\strut 32.2s}\\[-.2pt]{\fontsize{5.8}{6.4}\selectfont \textcolor{ForestGreen}{$ \blacktriangle\,46.6\times $}}\end{tabular} & \hspace*{1pt}\begin{tabular}[t]{@{}c@{}}\rule{0pt}{12pt}26.8s\\[-.2pt]{\fontsize{5.8}{6.4}\selectfont \textcolor{ForestGreen}{$ \blacktriangle\,21\% $}}\\[-.2pt]22.7s\\[-.2pt]{\fontsize{5.8}{6.4}\selectfont \textcolor{ForestGreen}{$ \blacktriangle\,1.5\times $}}\end{tabular} & \hspace*{1pt}\begin{tabular}[t]{@{}c@{}}\rule{0pt}{12pt}34.6s\\[-.2pt]{\fontsize{5.8}{6.4}\selectfont \textcolor{gray}{$ \blacktriangleright\,1\% $}}\\[-.2pt]31.2s\\[-.2pt]{\fontsize{5.8}{6.4}\selectfont \textcolor{gray}{$ \blacktriangleright\,4\% $}}\end{tabular} & \hspace*{1pt}\begin{tabular}[t]{@{}c@{}}\rule{0pt}{12pt}26s\\[-.2pt]{\fontsize{5.8}{6.4}\selectfont \textcolor{ForestGreen}{$ \blacktriangle\,2.6\times $}}\\[-.2pt]18.5s\\[-.2pt]{\fontsize{5.8}{6.4}\selectfont \textcolor{ForestGreen}{$ \blacktriangle\,3.5\times $}}\end{tabular} & \hspace*{1pt}\begin{tabular}[t]{@{}c@{}}\rule{0pt}{12pt}53.3s\\[-.2pt]{\fontsize{5.8}{6.4}\selectfont \textcolor{ForestGreen}{$ \blacktriangle\,3.6\times $}}\\[-.2pt]49.4s\\[-.2pt]{\fontsize{5.8}{6.4}\selectfont \textcolor{ForestGreen}{$ \blacktriangle\,3.7\times $}}\end{tabular} & \hspace*{1pt}\begin{tabular}[t]{@{}c@{}}\rule{0pt}{12pt}2m\,51.8s\\[-.2pt]{\fontsize{5.8}{6.4}\selectfont \textcolor{ForestGreen}{$ \blacktriangle\,18\% $}}\\[-.2pt]2m\,40.6s\\[-.2pt]{\fontsize{5.8}{6.4}\selectfont \textcolor{ForestGreen}{$ \blacktriangle\,21\% $}}\end{tabular} & \hspace*{1pt}\begin{tabular}[t]{@{}c@{}}\rule{0pt}{12pt}2m\,50.2s\\[-.2pt]{\fontsize{5.8}{6.4}\selectfont \textcolor{ForestGreen}{$ \blacktriangle\,27\% $}}\\[-.2pt]2m\,36.4s\\[-.2pt]{\fontsize{5.8}{6.4}\selectfont \textcolor{ForestGreen}{$ \blacktriangle\,1.4\times $}}\end{tabular} & \textcolor{gray}{\ensuremath{\cdot}} & \textcolor{gray}{\ensuremath{\cdot}} & \hspace*{1pt}\begin{tabular}[t]{@{}c@{}}\rule{0pt}{12pt}9m\,12.4s\\[-.2pt]{\fontsize{5.8}{6.4}\selectfont \textcolor{BrickRed}{$ \blacktriangledown\,7.9\times $}}\\[-.2pt]2m\,59.4s\\[-.2pt]{\fontsize{5.8}{6.4}\selectfont \textcolor{BrickRed}{$ \blacktriangledown\,2.6\times $}}\end{tabular} & \hspace*{1pt}\begin{tabular}[t]{@{}c@{}}\rule{0pt}{12pt}7m\,51.1s\\[-.2pt]{\fontsize{5.8}{6.4}\selectfont \textcolor{BrickRed}{$ \blacktriangledown\,6.9\times $}}\\[-.2pt]2m\,41.8s\\[-.2pt]{\fontsize{5.8}{6.4}\selectfont \textcolor{BrickRed}{$ \blacktriangledown\,2.4\times $}}\end{tabular} & \textcolor{gray}{\ensuremath{\cdot}} & \textcolor{gray}{\ensuremath{\cdot}}\\[3pt]\hline
\BenchmarkSourceLabel{bench:aws-4-k20-w15:row:mps-bounded-simd-block}{R7}\shortstack[l]{{\color{gray}R7}} & \begin{tabular}[t]{@{}l@{}}\rule{0pt}{12pt}MPS / Bounded\\SIMD-DP / Block\end{tabular} & \hspace*{1pt}\begin{tabular}[t]{@{}c@{}}\rule{0pt}{12pt}11.7s\\[-.2pt]{\fontsize{5.8}{6.4}\selectfont \textcolor{ForestGreen}{$ \blacktriangle\,14.3\times $}}\\[-.2pt]11s\\[-.2pt]{\fontsize{5.8}{6.4}\selectfont \textcolor{ForestGreen}{$ \blacktriangle\,14.7\times $}}\end{tabular} & \hspace*{1pt}\begin{tabular}[t]{@{}c@{}}\rule{0pt}{12pt}11.9s\\[-.2pt]{\fontsize{5.8}{6.4}\selectfont \textcolor{ForestGreen}{$ \blacktriangle\,14.3\times $}}\\[-.2pt]11.2s\\[-.2pt]{\fontsize{5.8}{6.4}\selectfont \textcolor{ForestGreen}{$ \blacktriangle\,14.6\times $}}\end{tabular} & \textcolor{gray}{\ensuremath{\cdot}} & \textcolor{gray}{\ensuremath{\cdot}} & \hspace*{1pt}\begin{tabular}[t]{@{}c@{}}\rule{0pt}{12pt}8m\,5.8s\\[-.2pt]{\fontsize{5.8}{6.4}\selectfont \textcolor{BrickRed}{$ \blacktriangledown\,14.3\times $}}\\[-.2pt]38.1s\\[-.2pt]{\fontsize{5.8}{6.4}\selectfont \textcolor{BrickRed}{$ \blacktriangledown\,16\% $}}\end{tabular} & \hspace*{1pt}\begin{tabular}[t]{@{}c@{}}\rule{0pt}{12pt}5m\,56.7s\\[-.2pt]{\fontsize{5.8}{6.4}\selectfont \textcolor{BrickRed}{$ \blacktriangledown\,10.4\times $}}\\[-.2pt]35.4s\\[-.2pt]{\fontsize{5.8}{6.4}\selectfont \textcolor{BrickRed}{$ \blacktriangledown\,8\% $}}\end{tabular} & \textcolor{gray}{\ensuremath{\cdot}} & \textcolor{gray}{\ensuremath{\cdot}} & \hspace*{1pt}\begin{tabular}[t]{@{}c@{}}\rule{0pt}{12pt}2m\,43.8s\\[-.2pt]{\fontsize{5.8}{6.4}\selectfont \textcolor{ForestGreen}{$ \blacktriangle\,21\% $}}\\[-.2pt]2m\,34.7s\\[-.2pt]{\fontsize{5.8}{6.4}\selectfont \textcolor{ForestGreen}{$ \blacktriangle\,24\% $}}\end{tabular} & \hspace*{1pt}\begin{tabular}[t]{@{}c@{}}\rule{0pt}{12pt}2m\,58.9s\\[-.2pt]{\fontsize{5.8}{6.4}\selectfont \textcolor{ForestGreen}{$ \blacktriangle\,23\% $}}\\[-.2pt]2m\,46.3s\\[-.2pt]{\fontsize{5.8}{6.4}\selectfont \textcolor{ForestGreen}{$ \blacktriangle\,26\% $}}\end{tabular} & \textcolor{gray}{\ensuremath{\cdot}} & \textcolor{gray}{\ensuremath{\cdot}} & \hspace*{1pt}\begin{tabular}[t]{@{}c@{}}\rule{0pt}{12pt}9m\,21.2s\\[-.2pt]{\fontsize{5.8}{6.4}\selectfont \textcolor{BrickRed}{$ \blacktriangledown\,8.0\times $}}\\[-.2pt]3m\,2.3s\\[-.2pt]{\fontsize{5.8}{6.4}\selectfont \textcolor{BrickRed}{$ \blacktriangledown\,2.6\times $}}\end{tabular} & \hspace*{1pt}\begin{tabular}[t]{@{}c@{}}\rule{0pt}{12pt}8m\,1.2s\\[-.2pt]{\fontsize{5.8}{6.4}\selectfont \textcolor{BrickRed}{$ \blacktriangledown\,7.0\times $}}\\[-.2pt]2m\,44.6s\\[-.2pt]{\fontsize{5.8}{6.4}\selectfont \textcolor{BrickRed}{$ \blacktriangledown\,2.5\times $}}\end{tabular} & \textcolor{gray}{\ensuremath{\cdot}} & \textcolor{gray}{\ensuremath{\cdot}}\\[3pt]\hline
\BenchmarkSourceLabel{bench:aws-4-k20-w15:row:mps-bounded-rmq-block}{R8}\shortstack[l]{{\color{gray}R8}} & \begin{tabular}[t]{@{}l@{}}\rule{0pt}{12pt}MPS / Bounded\\RMQ / Block\end{tabular} & \hspace*{1pt}\begin{tabular}[t]{@{}c@{}}\rule{0pt}{12pt}11.7s\\[-.2pt]{\fontsize{5.8}{6.4}\selectfont \textcolor{ForestGreen}{$ \blacktriangle\,14.3\times $}}\\[-.2pt]11s\\[-.2pt]{\fontsize{5.8}{6.4}\selectfont \textcolor{ForestGreen}{$ \blacktriangle\,14.6\times $}}\end{tabular} & \hspace*{1pt}\begin{tabular}[t]{@{}c@{}}\rule{0pt}{12pt}11.8s\\[-.2pt]{\fontsize{5.8}{6.4}\selectfont \textcolor{ForestGreen}{$ \blacktriangle\,14.5\times $}}\\[-.2pt]11.1s\\[-.2pt]{\fontsize{5.8}{6.4}\selectfont \textcolor{ForestGreen}{$ \blacktriangle\,14.7\times $}}\end{tabular} & \textcolor{gray}{\ensuremath{\cdot}} & \textcolor{gray}{\ensuremath{\cdot}} & \hspace*{1pt}\begin{tabular}[t]{@{}c@{}}\rule{0pt}{12pt}7m\,41.1s\\[-.2pt]{\fontsize{5.8}{6.4}\selectfont \textcolor{BrickRed}{$ \blacktriangledown\,13.6\times $}}\\[-.2pt]35.7s\\[-.2pt]{\fontsize{5.8}{6.4}\selectfont \textcolor{BrickRed}{$ \blacktriangledown\,8\% $}}\end{tabular} & \hspace*{1pt}\begin{tabular}[t]{@{}c@{}}\rule{0pt}{12pt}5m\,57.8s\\[-.2pt]{\fontsize{5.8}{6.4}\selectfont \textcolor{BrickRed}{$ \blacktriangledown\,10.4\times $}}\\[-.2pt]34.4s\\[-.2pt]{\fontsize{5.8}{6.4}\selectfont \textcolor{BrickRed}{$ \blacktriangledown\,5\% $}}\end{tabular} & \textcolor{gray}{\ensuremath{\cdot}} & \textcolor{gray}{\ensuremath{\cdot}} & \hspace*{1pt}\begin{tabular}[t]{@{}c@{}}\rule{0pt}{12pt}2m\,43.8s\\[-.2pt]{\fontsize{5.8}{6.4}\selectfont \textcolor{ForestGreen}{$ \blacktriangle\,21\% $}}\\[-.2pt]2m\,34.7s\\[-.2pt]{\fontsize{5.8}{6.4}\selectfont \textcolor{ForestGreen}{$ \blacktriangle\,24\% $}}\end{tabular} & \hspace*{1pt}\begin{tabular}[t]{@{}c@{}}\rule{0pt}{12pt}2m\,52.1s\\[-.2pt]{\fontsize{5.8}{6.4}\selectfont \textcolor{ForestGreen}{$ \blacktriangle\,26\% $}}\\[-.2pt]2m\,39.6s\\[-.2pt]{\fontsize{5.8}{6.4}\selectfont \textcolor{ForestGreen}{$ \blacktriangle\,29\% $}}\end{tabular} & \textcolor{gray}{\ensuremath{\cdot}} & \textcolor{gray}{\ensuremath{\cdot}} & \hspace*{1pt}\begin{tabular}[t]{@{}c@{}}\rule{0pt}{12pt}9m\,17.3s\\[-.2pt]{\fontsize{5.8}{6.4}\selectfont \textcolor{BrickRed}{$ \blacktriangledown\,7.9\times $}}\\[-.2pt]2m\,59.8s\\[-.2pt]{\fontsize{5.8}{6.4}\selectfont \textcolor{BrickRed}{$ \blacktriangledown\,2.6\times $}}\end{tabular} & \textcolor{gray}{\ensuremath{\cdot}} & \hspace*{1pt}\begin{tabular}[t]{@{}c@{}}\rule{0pt}{12pt}14m\,16s\\[-.2pt]{\fontsize{5.8}{6.4}\selectfont \textcolor{BrickRed}{$ \blacktriangledown\,6.0\times $}}\\[-.2pt]5m\,26.4s\\[-.2pt]{\fontsize{5.8}{6.4}\selectfont \textcolor{BrickRed}{$ \blacktriangledown\,2.4\times $}}\end{tabular} & \textcolor{gray}{\ensuremath{\cdot}}\\[3pt]\hline
\end{tabular}\par}
\smallskip{\footnotesize\noindent Each cell: A and B times, $\cdot$: no measurement. ---\textsuperscript{g}: GPU unavailable. n/a: no comparison.\par
\noindent Blue: one lowest A and one lowest B measurement for each column for each configuration.\par
\noindent Gray: minimap2. Yellow: my preferred MPF / Bounded / SIMD-DP / Block configuration.\par
\noindent MPF/MPS: minimizers per fragment/sequence. SIMD-ref: naive SIMD scan. Bounded: bounded SIMD algorithm.\par
}
\clearpage
\noindent{\large\bfseries Table \thetable: (C): AWS: 4 Gbp, k=20, w=15\hfill 2/7}\par\smallskip
{\fontsize{7.1}{8.1}\selectfont\setlength{\tabcolsep}{1pt}\setlength{\fboxsep}{2pt}\renewcommand{\arraystretch}{1.12}
\noindent\begin{tabular}{@{}p{14mm}p{27mm}|p{13.55mm}|p{13.55mm}|p{13.55mm}|p{13.55mm}|p{13.55mm}|p{13.55mm}|p{13.55mm}|p{13.55mm}|p{13.55mm}|p{13.55mm}|p{13.55mm}|p{13.55mm}|p{13.55mm}|p{13.55mm}|p{13.55mm}|p{13.55mm}|@{}}\hline
Engine & Execution & \multicolumn{8}{c|}{\textbf{Index}} & \multicolumn{8}{c|}{\textbf{No index}}\\
 &  & \multicolumn{8}{c|}{\textbf{4 Gbp}} & \multicolumn{8}{c|}{\textbf{4 Gbp}}\\
 &  & \multicolumn{8}{c|}{\textbf{k=20, w=15}} & \multicolumn{8}{c|}{\textbf{k=20, w=15}}\\
 &  & \multicolumn{4}{c|}{\textbf{DNA4}} & \multicolumn{4}{c|}{\textbf{DNA5}} & \multicolumn{4}{c|}{\textbf{DNA4}} & \multicolumn{4}{c|}{\textbf{DNA5}}\\
 &  & \multicolumn{1}{c|}{\textbf{M7A}} & \multicolumn{1}{c|}{\textbf{M7I}} & \multicolumn{1}{c|}{\textbf{M4}} & \multicolumn{1}{c|}{\textbf{C7G}} & \multicolumn{1}{c|}{\textbf{M7A}} & \multicolumn{1}{c|}{\textbf{M7I}} & \multicolumn{1}{c|}{\textbf{M4}} & \multicolumn{1}{c|}{\textbf{C7G}} & \multicolumn{1}{c|}{\textbf{M7A}} & \multicolumn{1}{c|}{\textbf{M7I}} & \multicolumn{1}{c|}{\textbf{M4}} & \multicolumn{1}{c|}{\textbf{C7G}} & \multicolumn{1}{c|}{\textbf{M7A}} & \multicolumn{1}{c|}{\textbf{M7I}} & \multicolumn{1}{c|}{\textbf{M4}} & \multicolumn{1}{c|}{\textbf{C7G}}\\
 & & {\color{gray}C1} & {\color{gray}C2} & {\color{gray}C3} & {\color{gray}C4} & {\color{gray}C5} & {\color{gray}C6} & {\color{gray}C7} & {\color{gray}C8} & {\color{gray}C9} & {\color{gray}C10} & {\color{gray}C11} & {\color{gray}C12} & {\color{gray}C13} & {\color{gray}C14} & {\color{gray}C15} & {\color{gray}C16}\\
\hline
\rowcolor{BenchBaseline}[0.5pt][0.5pt]\shortstack[l]{{\color{gray}R1}\\\textbf{minimap2}} & \textbf{Baseline} & \hspace*{1pt}\begin{tabular}[t]{@{}c@{}}\rule{0pt}{12pt}2m\,48.1s\\[-.2pt]{\fontsize{5.8}{6.4}\selectfont \textcolor{BrickRed}{$ \blacktriangledown\,20.4\times $}}\\[-.2pt]2m\,41.7s\\[-.2pt]{\fontsize{5.8}{6.4}\selectfont \textcolor{BrickRed}{$ \blacktriangledown\,22.6\times $}}\end{tabular} & \hspace*{1pt}\begin{tabular}[t]{@{}c@{}}\rule{0pt}{12pt}2m\,50.3s\\[-.2pt]{\fontsize{5.8}{6.4}\selectfont \textcolor{BrickRed}{$ \blacktriangledown\,19.8\times $}}\\[-.2pt]2m\,43s\\[-.2pt]{\fontsize{5.8}{6.4}\selectfont \textcolor{BrickRed}{$ \blacktriangledown\,22.2\times $}}\end{tabular} & \hspace*{1pt}\begin{tabular}[t]{@{}c@{}}\rule{0pt}{12pt}4m\,54.3s\\[-.2pt]{\fontsize{5.8}{6.4}\selectfont \textcolor{BrickRed}{$ \blacktriangledown\,43.2\times $}}\\[-.2pt]4m\,39s\\[-.2pt]{\fontsize{5.8}{6.4}\selectfont \textcolor{BrickRed}{$ \blacktriangledown\,54.6\times $}}\end{tabular} & \hspace*{1pt}\begin{tabular}[t]{@{}c@{}}\rule{0pt}{12pt}25m\,5.5s\\[-.2pt]{\fontsize{5.8}{6.4}\selectfont \textcolor{BrickRed}{$ \blacktriangledown\,41.8\times $}}\\[-.2pt]25m\,0.7s\\[-.2pt]{\fontsize{5.8}{6.4}\selectfont \textcolor{BrickRed}{$ \blacktriangledown\,46.6\times $}}\end{tabular} & \hspace*{1pt}\begin{tabular}[t]{@{}c@{}}\rule{0pt}{12pt}33.9s\\[-.2pt]{\fontsize{5.8}{6.4}\selectfont \textcolor{BrickRed}{$ \blacktriangledown\,2.3\times $}}\\[-.2pt]32.9s\\[-.2pt]{\fontsize{5.8}{6.4}\selectfont \textcolor{BrickRed}{$ \blacktriangledown\,3.1\times $}}\end{tabular} & \hspace*{1pt}\begin{tabular}[t]{@{}c@{}}\rule{0pt}{12pt}34.4s\\[-.2pt]{\fontsize{5.8}{6.4}\selectfont \textcolor{BrickRed}{$ \blacktriangledown\,2.6\times $}}\\[-.2pt]32.7s\\[-.2pt]{\fontsize{5.8}{6.4}\selectfont \textcolor{BrickRed}{$ \blacktriangledown\,2.9\times $}}\end{tabular} & \hspace*{1pt}\begin{tabular}[t]{@{}c@{}}\rule{0pt}{12pt}1m\,8.2s\\[-.2pt]{\fontsize{5.8}{6.4}\selectfont \textcolor{BrickRed}{$ \blacktriangledown\,3.7\times $}}\\[-.2pt]1m\,5.1s\\[-.2pt]{\fontsize{5.8}{6.4}\selectfont \textcolor{BrickRed}{$ \blacktriangledown\,4.6\times $}}\end{tabular} & \hspace*{1pt}\begin{tabular}[t]{@{}c@{}}\rule{0pt}{12pt}3m\,9.8s\\[-.2pt]{\fontsize{5.8}{6.4}\selectfont \textcolor{BrickRed}{$ \blacktriangledown\,3.9\times $}}\\[-.2pt]3m\,3.7s\\[-.2pt]{\fontsize{5.8}{6.4}\selectfont \textcolor{BrickRed}{$ \blacktriangledown\,4.1\times $}}\end{tabular} & \hspace*{1pt}\begin{tabular}[t]{@{}c@{}}\rule{0pt}{12pt}3m\,28.3s\\[-.2pt]{\fontsize{5.8}{6.4}\selectfont \textcolor{BrickRed}{$ \blacktriangledown\,1.3\times $}}\\[-.2pt]3m\,22.8s\\[-.2pt]{\fontsize{5.8}{6.4}\selectfont \textcolor{BrickRed}{$ \blacktriangledown\,1.4\times $}}\end{tabular} & \hspace*{1pt}\begin{tabular}[t]{@{}c@{}}\rule{0pt}{12pt}3m\,53.2s\\[-.2pt]{\fontsize{5.8}{6.4}\selectfont \textcolor{BrickRed}{$ \blacktriangledown\,1.4\times $}}\\[-.2pt]3m\,43.5s\\[-.2pt]{\fontsize{5.8}{6.4}\selectfont \textcolor{BrickRed}{$ \blacktriangledown\,1.5\times $}}\end{tabular} & \hspace*{1pt}\begin{tabular}[t]{@{}c@{}}\rule{0pt}{12pt}7m\,14.6s\\[-.2pt]{\fontsize{5.8}{6.4}\selectfont \textcolor{BrickRed}{$ \blacktriangledown\,1.5\times $}}\\[-.2pt]6m\,57.2s\\[-.2pt]{\fontsize{5.8}{6.4}\selectfont \textcolor{BrickRed}{$ \blacktriangledown\,1.6\times $}}\end{tabular} & \hspace*{1pt}\begin{tabular}[t]{@{}c@{}}\rule{0pt}{12pt}51m\,22s\\[-.2pt]{\fontsize{5.8}{6.4}\selectfont \textcolor{gray}{n/a}}\\[-.2pt]51m\,17.2s\\[-.2pt]{\fontsize{5.8}{6.4}\selectfont \textcolor{gray}{n/a}}\end{tabular} & \hspace*{1pt}\begin{tabular}[t]{@{}c@{}}\rule{0pt}{12pt}1m\,10.3s\\[-.2pt]{\fontsize{5.8}{6.4}\selectfont \textcolor{ForestGreen}{$ \blacktriangle\,7.6\times $}}\\[-.2pt]1m\,9s\\[-.2pt]{\fontsize{5.8}{6.4}\selectfont \textcolor{ForestGreen}{$ \blacktriangle\,2.4\times $}}\end{tabular} & \hspace*{1pt}\begin{tabular}[t]{@{}c@{}}\rule{0pt}{12pt}1m\,8.6s\\[-.2pt]{\fontsize{5.8}{6.4}\selectfont \textcolor{ForestGreen}{$ \blacktriangle\,6.4\times $}}\\[-.2pt]1m\,6.9s\\[-.2pt]{\fontsize{5.8}{6.4}\selectfont \textcolor{ForestGreen}{$ \blacktriangle\,2.3\times $}}\end{tabular} & \hspace*{1pt}\begin{tabular}[t]{@{}c@{}}\rule{0pt}{12pt}2m\,21.5s\\[-.2pt]{\fontsize{5.8}{6.4}\selectfont \textcolor{ForestGreen}{$ \blacktriangle\,6.0\times $}}\\[-.2pt]2m\,18.4s\\[-.2pt]{\fontsize{5.8}{6.4}\selectfont \textcolor{ForestGreen}{$ \blacktriangle\,2.2\times $}}\end{tabular} & \hspace*{1pt}\begin{tabular}[t]{@{}c@{}}\rule{0pt}{12pt}15m\,56.8s\\[-.2pt]{\fontsize{5.8}{6.4}\selectfont \textcolor{gray}{n/a}}\\[-.2pt]15m\,51.9s\\[-.2pt]{\fontsize{5.8}{6.4}\selectfont \textcolor{gray}{n/a}}\end{tabular}\\[3pt]\hline
\BenchmarkSourceLabel{bench:aws-4-k20-w15:row:mps-bounded-scalar-block}{R9}\shortstack[l]{{\color{gray}R9}\\\textbf{BioseqDB}} & \begin{tabular}[t]{@{}l@{}}\rule{0pt}{12pt}MPS / Bounded\\Scalar-DP / Block\end{tabular} & \hspace*{1pt}\begin{tabular}[t]{@{}c@{}}\rule{0pt}{12pt}11.8s\\[-.2pt]{\fontsize{5.8}{6.4}\selectfont \textcolor{ForestGreen}{$ \blacktriangle\,14.2\times $}}\\[-.2pt]11.1s\\[-.2pt]{\fontsize{5.8}{6.4}\selectfont \textcolor{ForestGreen}{$ \blacktriangle\,14.6\times $}}\end{tabular} & \hspace*{1pt}\begin{tabular}[t]{@{}c@{}}\rule{0pt}{12pt}12.7s\\[-.2pt]{\fontsize{5.8}{6.4}\selectfont \textcolor{ForestGreen}{$ \blacktriangle\,13.4\times $}}\\[-.2pt]11.9s\\[-.2pt]{\fontsize{5.8}{6.4}\selectfont \textcolor{ForestGreen}{$ \blacktriangle\,13.7\times $}}\end{tabular} & \textcolor{gray}{\ensuremath{\cdot}} & \textcolor{gray}{\ensuremath{\cdot}} & \hspace*{1pt}\begin{tabular}[t]{@{}c@{}}\rule{0pt}{12pt}7m\,27.3s\\[-.2pt]{\fontsize{5.8}{6.4}\selectfont \textcolor{BrickRed}{$ \blacktriangledown\,13.2\times $}}\\[-.2pt]36.3s\\[-.2pt]{\fontsize{5.8}{6.4}\selectfont \textcolor{BrickRed}{$ \blacktriangledown\,10\% $}}\end{tabular} & \hspace*{1pt}\begin{tabular}[t]{@{}c@{}}\rule{0pt}{12pt}5m\,57.5s\\[-.2pt]{\fontsize{5.8}{6.4}\selectfont \textcolor{BrickRed}{$ \blacktriangledown\,10.4\times $}}\\[-.2pt]34.3s\\[-.2pt]{\fontsize{5.8}{6.4}\selectfont \textcolor{gray}{$ \blacktriangleright\,5\% $}}\end{tabular} & \textcolor{gray}{\ensuremath{\cdot}} & \textcolor{gray}{\ensuremath{\cdot}} & \hspace*{1pt}\begin{tabular}[t]{@{}c@{}}\rule{0pt}{12pt}2m\,52.6s\\[-.2pt]{\fontsize{5.8}{6.4}\selectfont \textcolor{ForestGreen}{$ \blacktriangle\,17\% $}}\\[-.2pt]2m\,42.8s\\[-.2pt]{\fontsize{5.8}{6.4}\selectfont \textcolor{ForestGreen}{$ \blacktriangle\,20\% $}}\end{tabular} & \hspace*{1pt}\begin{tabular}[t]{@{}c@{}}\rule{0pt}{12pt}2m\,58.9s\\[-.2pt]{\fontsize{5.8}{6.4}\selectfont \textcolor{ForestGreen}{$ \blacktriangle\,23\% $}}\\[-.2pt]2m\,46.9s\\[-.2pt]{\fontsize{5.8}{6.4}\selectfont \textcolor{ForestGreen}{$ \blacktriangle\,25\% $}}\end{tabular} & \textcolor{gray}{\ensuremath{\cdot}} & \textcolor{gray}{\ensuremath{\cdot}} & \hspace*{1pt}\begin{tabular}[t]{@{}c@{}}\rule{0pt}{12pt}9m\,15.5s\\[-.2pt]{\fontsize{5.8}{6.4}\selectfont \textcolor{BrickRed}{$ \blacktriangledown\,7.9\times $}}\\[-.2pt]3m\,0.5s\\[-.2pt]{\fontsize{5.8}{6.4}\selectfont \textcolor{BrickRed}{$ \blacktriangledown\,2.6\times $}}\end{tabular} & \textcolor{gray}{\ensuremath{\cdot}} & \textcolor{gray}{\ensuremath{\cdot}} & \textcolor{gray}{\ensuremath{\cdot}}\\[3pt]\hline
\BenchmarkSourceLabel{bench:aws-4-k20-w15:row:mpf-bounded-scalar-ksw2}{R10}\shortstack[l]{{\color{gray}R10}} & \begin{tabular}[t]{@{}l@{}}\rule{0pt}{12pt}MPF / Bounded\\Scalar-DP / KSW2\end{tabular} & \hspace*{1pt}\begin{tabular}[t]{@{}c@{}}\rule{0pt}{12pt}14.5s\\[-.2pt]{\fontsize{5.8}{6.4}\selectfont \textcolor{ForestGreen}{$ \blacktriangle\,11.6\times $}}\\[-.2pt]12.2s\\[-.2pt]{\fontsize{5.8}{6.4}\selectfont \textcolor{ForestGreen}{$ \blacktriangle\,13.3\times $}}\end{tabular} & \hspace*{1pt}\begin{tabular}[t]{@{}c@{}}\rule{0pt}{12pt}14.2s\\[-.2pt]{\fontsize{5.8}{6.4}\selectfont \textcolor{ForestGreen}{$ \blacktriangle\,12.0\times $}}\\[-.2pt]11.9s\\[-.2pt]{\fontsize{5.8}{6.4}\selectfont \textcolor{ForestGreen}{$ \blacktriangle\,13.7\times $}}\end{tabular} & \hspace*{1pt}\begin{tabular}[t]{@{}c@{}}\rule{0pt}{12pt}12.7s\\[-.2pt]{\fontsize{5.8}{6.4}\selectfont \textcolor{ForestGreen}{$ \blacktriangle\,23.2\times $}}\\[-.2pt]5.2s\\[-.2pt]{\fontsize{5.8}{6.4}\selectfont \textcolor{ForestGreen}{$ \blacktriangle\,53.2\times $}}\end{tabular} & \hspace*{1pt}\begin{tabular}[t]{@{}c@{}}\rule{0pt}{12pt}36.8s\\[-.2pt]{\fontsize{5.8}{6.4}\selectfont \textcolor{ForestGreen}{$ \blacktriangle\,40.9\times $}}\\[-.2pt]32.8s\\[-.2pt]{\fontsize{5.8}{6.4}\selectfont \textcolor{ForestGreen}{$ \blacktriangle\,45.7\times $}}\end{tabular} & \hspace*{1pt}\begin{tabular}[t]{@{}c@{}}\rule{0pt}{12pt}\colorbox{BenchBest}{\strut 14.8s}\\[-.2pt]{\fontsize{5.8}{6.4}\selectfont \textcolor{ForestGreen}{$ \blacktriangle\,2.3\times $}}\\[-.2pt]\colorbox{BenchBest}{\strut 10.7s}\\[-.2pt]{\fontsize{5.8}{6.4}\selectfont \textcolor{ForestGreen}{$ \blacktriangle\,3.1\times $}}\end{tabular} & \hspace*{1pt}\begin{tabular}[t]{@{}c@{}}\rule{0pt}{12pt}17.7s\\[-.2pt]{\fontsize{5.8}{6.4}\selectfont \textcolor{ForestGreen}{$ \blacktriangle\,1.9\times $}}\\[-.2pt]14.3s\\[-.2pt]{\fontsize{5.8}{6.4}\selectfont \textcolor{ForestGreen}{$ \blacktriangle\,2.3\times $}}\end{tabular} & \hspace*{1pt}\begin{tabular}[t]{@{}c@{}}\rule{0pt}{12pt}27s\\[-.2pt]{\fontsize{5.8}{6.4}\selectfont \textcolor{ForestGreen}{$ \blacktriangle\,2.5\times $}}\\[-.2pt]19.4s\\[-.2pt]{\fontsize{5.8}{6.4}\selectfont \textcolor{ForestGreen}{$ \blacktriangle\,3.4\times $}}\end{tabular} & \hspace*{1pt}\begin{tabular}[t]{@{}c@{}}\rule{0pt}{12pt}\colorbox{BenchBest}{\strut 48.6s}\\[-.2pt]{\fontsize{5.8}{6.4}\selectfont \textcolor{ForestGreen}{$ \blacktriangle\,3.9\times $}}\\[-.2pt]\colorbox{BenchBest}{\strut 44.8s}\\[-.2pt]{\fontsize{5.8}{6.4}\selectfont \textcolor{ForestGreen}{$ \blacktriangle\,4.1\times $}}\end{tabular} & \hspace*{1pt}\begin{tabular}[t]{@{}c@{}}\rule{0pt}{12pt}2m\,40.6s\\[-.2pt]{\fontsize{5.8}{6.4}\selectfont \textcolor{ForestGreen}{$ \blacktriangle\,23\% $}}\\[-.2pt]2m\,29.9s\\[-.2pt]{\fontsize{5.8}{6.4}\selectfont \textcolor{ForestGreen}{$ \blacktriangle\,26\% $}}\end{tabular} & \hspace*{1pt}\begin{tabular}[t]{@{}c@{}}\rule{0pt}{12pt}3m\,2.9s\\[-.2pt]{\fontsize{5.8}{6.4}\selectfont \textcolor{ForestGreen}{$ \blacktriangle\,22\% $}}\\[-.2pt]2m\,48.3s\\[-.2pt]{\fontsize{5.8}{6.4}\selectfont \textcolor{ForestGreen}{$ \blacktriangle\,25\% $}}\end{tabular} & \textcolor{gray}{\ensuremath{\cdot}} & \textcolor{gray}{\ensuremath{\cdot}} & \hspace*{1pt}\begin{tabular}[t]{@{}c@{}}\rule{0pt}{12pt}9m\,11s\\[-.2pt]{\fontsize{5.8}{6.4}\selectfont \textcolor{BrickRed}{$ \blacktriangledown\,7.8\times $}}\\[-.2pt]2m\,56.6s\\[-.2pt]{\fontsize{5.8}{6.4}\selectfont \textcolor{BrickRed}{$ \blacktriangledown\,2.6\times $}}\end{tabular} & \hspace*{1pt}\begin{tabular}[t]{@{}c@{}}\rule{0pt}{12pt}7m\,47.4s\\[-.2pt]{\fontsize{5.8}{6.4}\selectfont \textcolor{BrickRed}{$ \blacktriangledown\,6.8\times $}}\\[-.2pt]2m\,37.9s\\[-.2pt]{\fontsize{5.8}{6.4}\selectfont \textcolor{BrickRed}{$ \blacktriangledown\,2.4\times $}}\end{tabular} & \textcolor{gray}{\ensuremath{\cdot}} & \textcolor{gray}{\ensuremath{\cdot}}\\[3pt]\hline
\BenchmarkSourceLabel{bench:aws-4-k20-w15:row:mps-bounded-rmq-ksw2}{R11}\shortstack[l]{{\color{gray}R11}} & \begin{tabular}[t]{@{}l@{}}\rule{0pt}{12pt}MPS / Bounded\\RMQ / KSW2\end{tabular} & \hspace*{1pt}\begin{tabular}[t]{@{}c@{}}\rule{0pt}{12pt}13.5s\\[-.2pt]{\fontsize{5.8}{6.4}\selectfont \textcolor{ForestGreen}{$ \blacktriangle\,12.4\times $}}\\[-.2pt]11s\\[-.2pt]{\fontsize{5.8}{6.4}\selectfont \textcolor{ForestGreen}{$ \blacktriangle\,14.7\times $}}\end{tabular} & \hspace*{1pt}\begin{tabular}[t]{@{}c@{}}\rule{0pt}{12pt}12.8s\\[-.2pt]{\fontsize{5.8}{6.4}\selectfont \textcolor{ForestGreen}{$ \blacktriangle\,13.3\times $}}\\[-.2pt]10.6s\\[-.2pt]{\fontsize{5.8}{6.4}\selectfont \textcolor{ForestGreen}{$ \blacktriangle\,15.4\times $}}\end{tabular} & \textcolor{gray}{\ensuremath{\cdot}} & \textcolor{gray}{\ensuremath{\cdot}} & \hspace*{1pt}\begin{tabular}[t]{@{}c@{}}\rule{0pt}{12pt}7m\,20.4s\\[-.2pt]{\fontsize{5.8}{6.4}\selectfont \textcolor{BrickRed}{$ \blacktriangledown\,13.0\times $}}\\[-.2pt]35.7s\\[-.2pt]{\fontsize{5.8}{6.4}\selectfont \textcolor{BrickRed}{$ \blacktriangledown\,8\% $}}\end{tabular} & \hspace*{1pt}\begin{tabular}[t]{@{}c@{}}\rule{0pt}{12pt}6m\,8.8s\\[-.2pt]{\fontsize{5.8}{6.4}\selectfont \textcolor{BrickRed}{$ \blacktriangledown\,10.7\times $}}\\[-.2pt]34.6s\\[-.2pt]{\fontsize{5.8}{6.4}\selectfont \textcolor{BrickRed}{$ \blacktriangledown\,6\% $}}\end{tabular} & \textcolor{gray}{\ensuremath{\cdot}} & \textcolor{gray}{\ensuremath{\cdot}} & \hspace*{1pt}\begin{tabular}[t]{@{}c@{}}\rule{0pt}{12pt}2m\,56.8s\\[-.2pt]{\fontsize{5.8}{6.4}\selectfont \textcolor{ForestGreen}{$ \blacktriangle\,15\% $}}\\[-.2pt]2m\,44.8s\\[-.2pt]{\fontsize{5.8}{6.4}\selectfont \textcolor{ForestGreen}{$ \blacktriangle\,19\% $}}\end{tabular} & \hspace*{1pt}\begin{tabular}[t]{@{}c@{}}\rule{0pt}{12pt}3m\,11.4s\\[-.2pt]{\fontsize{5.8}{6.4}\selectfont \textcolor{ForestGreen}{$ \blacktriangle\,18\% $}}\\[-.2pt]2m\,56.8s\\[-.2pt]{\fontsize{5.8}{6.4}\selectfont \textcolor{ForestGreen}{$ \blacktriangle\,21\% $}}\end{tabular} & \textcolor{gray}{\ensuremath{\cdot}} & \textcolor{gray}{\ensuremath{\cdot}} & \hspace*{1pt}\begin{tabular}[t]{@{}c@{}}\rule{0pt}{12pt}9m\,19.1s\\[-.2pt]{\fontsize{5.8}{6.4}\selectfont \textcolor{BrickRed}{$ \blacktriangledown\,8.0\times $}}\\[-.2pt]3m\,4s\\[-.2pt]{\fontsize{5.8}{6.4}\selectfont \textcolor{BrickRed}{$ \blacktriangledown\,2.7\times $}}\end{tabular} & \textcolor{gray}{\ensuremath{\cdot}} & \textcolor{gray}{\ensuremath{\cdot}} & \textcolor{gray}{\ensuremath{\cdot}}\\[3pt]\hline
\BenchmarkSourceLabel{bench:aws-4-k20-w15:row:mps-bounded-scalar-ksw2}{R12}\shortstack[l]{{\color{gray}R12}} & \begin{tabular}[t]{@{}l@{}}\rule{0pt}{12pt}MPS / Bounded\\Scalar-DP / KSW2\end{tabular} & \hspace*{1pt}\begin{tabular}[t]{@{}c@{}}\rule{0pt}{12pt}13.4s\\[-.2pt]{\fontsize{5.8}{6.4}\selectfont \textcolor{ForestGreen}{$ \blacktriangle\,12.6\times $}}\\[-.2pt]10.9s\\[-.2pt]{\fontsize{5.8}{6.4}\selectfont \textcolor{ForestGreen}{$ \blacktriangle\,14.9\times $}}\end{tabular} & \hspace*{1pt}\begin{tabular}[t]{@{}c@{}}\rule{0pt}{12pt}13s\\[-.2pt]{\fontsize{5.8}{6.4}\selectfont \textcolor{ForestGreen}{$ \blacktriangle\,13.1\times $}}\\[-.2pt]10.7s\\[-.2pt]{\fontsize{5.8}{6.4}\selectfont \textcolor{ForestGreen}{$ \blacktriangle\,15.2\times $}}\end{tabular} & \textcolor{gray}{\ensuremath{\cdot}} & \textcolor{gray}{\ensuremath{\cdot}} & \hspace*{1pt}\begin{tabular}[t]{@{}c@{}}\rule{0pt}{12pt}7m\,44.9s\\[-.2pt]{\fontsize{5.8}{6.4}\selectfont \textcolor{BrickRed}{$ \blacktriangledown\,13.7\times $}}\\[-.2pt]35.9s\\[-.2pt]{\fontsize{5.8}{6.4}\selectfont \textcolor{BrickRed}{$ \blacktriangledown\,9\% $}}\end{tabular} & \hspace*{1pt}\begin{tabular}[t]{@{}c@{}}\rule{0pt}{12pt}5m\,59.4s\\[-.2pt]{\fontsize{5.8}{6.4}\selectfont \textcolor{BrickRed}{$ \blacktriangledown\,10.4\times $}}\\[-.2pt]34s\\[-.2pt]{\fontsize{5.8}{6.4}\selectfont \textcolor{gray}{$ \blacktriangleright\,4\% $}}\end{tabular} & \textcolor{gray}{\ensuremath{\cdot}} & \textcolor{gray}{\ensuremath{\cdot}} & \hspace*{1pt}\begin{tabular}[t]{@{}c@{}}\rule{0pt}{12pt}2m\,59.1s\\[-.2pt]{\fontsize{5.8}{6.4}\selectfont \textcolor{ForestGreen}{$ \blacktriangle\,14\% $}}\\[-.2pt]2m\,47.7s\\[-.2pt]{\fontsize{5.8}{6.4}\selectfont \textcolor{ForestGreen}{$ \blacktriangle\,17\% $}}\end{tabular} & \hspace*{1pt}\begin{tabular}[t]{@{}c@{}}\rule{0pt}{12pt}2m\,50.8s\\[-.2pt]{\fontsize{5.8}{6.4}\selectfont \textcolor{ForestGreen}{$ \blacktriangle\,27\% $}}\\[-.2pt]2m\,36.9s\\[-.2pt]{\fontsize{5.8}{6.4}\selectfont \textcolor{ForestGreen}{$ \blacktriangle\,30\% $}}\end{tabular} & \textcolor{gray}{\ensuremath{\cdot}} & \textcolor{gray}{\ensuremath{\cdot}} & \hspace*{1pt}\begin{tabular}[t]{@{}c@{}}\rule{0pt}{12pt}9m\,18.6s\\[-.2pt]{\fontsize{5.8}{6.4}\selectfont \textcolor{BrickRed}{$ \blacktriangledown\,8.0\times $}}\\[-.2pt]3m\,3.8s\\[-.2pt]{\fontsize{5.8}{6.4}\selectfont \textcolor{BrickRed}{$ \blacktriangledown\,2.7\times $}}\end{tabular} & \textcolor{gray}{\ensuremath{\cdot}} & \textcolor{gray}{\ensuremath{\cdot}} & \textcolor{gray}{\ensuremath{\cdot}}\\[3pt]\hline
\BenchmarkSourceLabel{bench:aws-4-k20-w15:row:mps-bounded-simd-ksw2}{R13}\shortstack[l]{{\color{gray}R13}} & \begin{tabular}[t]{@{}l@{}}\rule{0pt}{12pt}MPS / Bounded\\SIMD-DP / KSW2\end{tabular} & \hspace*{1pt}\begin{tabular}[t]{@{}c@{}}\rule{0pt}{12pt}13.3s\\[-.2pt]{\fontsize{5.8}{6.4}\selectfont \textcolor{ForestGreen}{$ \blacktriangle\,12.7\times $}}\\[-.2pt]10.8s\\[-.2pt]{\fontsize{5.8}{6.4}\selectfont \textcolor{ForestGreen}{$ \blacktriangle\,15.0\times $}}\end{tabular} & \hspace*{1pt}\begin{tabular}[t]{@{}c@{}}\rule{0pt}{12pt}13.6s\\[-.2pt]{\fontsize{5.8}{6.4}\selectfont \textcolor{ForestGreen}{$ \blacktriangle\,12.5\times $}}\\[-.2pt]11.1s\\[-.2pt]{\fontsize{5.8}{6.4}\selectfont \textcolor{ForestGreen}{$ \blacktriangle\,14.7\times $}}\end{tabular} & \textcolor{gray}{\ensuremath{\cdot}} & \textcolor{gray}{\ensuremath{\cdot}} & \hspace*{1pt}\begin{tabular}[t]{@{}c@{}}\rule{0pt}{12pt}7m\,20.3s\\[-.2pt]{\fontsize{5.8}{6.4}\selectfont \textcolor{BrickRed}{$ \blacktriangledown\,13.0\times $}}\\[-.2pt]35.3s\\[-.2pt]{\fontsize{5.8}{6.4}\selectfont \textcolor{BrickRed}{$ \blacktriangledown\,7\% $}}\end{tabular} & \hspace*{1pt}\begin{tabular}[t]{@{}c@{}}\rule{0pt}{12pt}5m\,56.2s\\[-.2pt]{\fontsize{5.8}{6.4}\selectfont \textcolor{BrickRed}{$ \blacktriangledown\,10.3\times $}}\\[-.2pt]34s\\[-.2pt]{\fontsize{5.8}{6.4}\selectfont \textcolor{gray}{$ \blacktriangleright\,4\% $}}\end{tabular} & \textcolor{gray}{\ensuremath{\cdot}} & \textcolor{gray}{\ensuremath{\cdot}} & \hspace*{1pt}\begin{tabular}[t]{@{}c@{}}\rule{0pt}{12pt}2m\,59s\\[-.2pt]{\fontsize{5.8}{6.4}\selectfont \textcolor{ForestGreen}{$ \blacktriangle\,14\% $}}\\[-.2pt]2m\,48.5s\\[-.2pt]{\fontsize{5.8}{6.4}\selectfont \textcolor{ForestGreen}{$ \blacktriangle\,17\% $}}\end{tabular} & \hspace*{1pt}\begin{tabular}[t]{@{}c@{}}\rule{0pt}{12pt}2m\,52.5s\\[-.2pt]{\fontsize{5.8}{6.4}\selectfont \textcolor{ForestGreen}{$ \blacktriangle\,26\% $}}\\[-.2pt]2m\,38.8s\\[-.2pt]{\fontsize{5.8}{6.4}\selectfont \textcolor{ForestGreen}{$ \blacktriangle\,29\% $}}\end{tabular} & \textcolor{gray}{\ensuremath{\cdot}} & \textcolor{gray}{\ensuremath{\cdot}} & \hspace*{1pt}\begin{tabular}[t]{@{}c@{}}\rule{0pt}{12pt}9m\,21.1s\\[-.2pt]{\fontsize{5.8}{6.4}\selectfont \textcolor{BrickRed}{$ \blacktriangledown\,8.0\times $}}\\[-.2pt]3m\,3.5s\\[-.2pt]{\fontsize{5.8}{6.4}\selectfont \textcolor{BrickRed}{$ \blacktriangledown\,2.7\times $}}\end{tabular} & \hspace*{1pt}\begin{tabular}[t]{@{}c@{}}\rule{0pt}{12pt}8m\,2.9s\\[-.2pt]{\fontsize{5.8}{6.4}\selectfont \textcolor{BrickRed}{$ \blacktriangledown\,7.0\times $}}\\[-.2pt]2m\,44s\\[-.2pt]{\fontsize{5.8}{6.4}\selectfont \textcolor{BrickRed}{$ \blacktriangledown\,2.5\times $}}\end{tabular} & \textcolor{gray}{\ensuremath{\cdot}} & \textcolor{gray}{\ensuremath{\cdot}}\\[3pt]\hline
\BenchmarkSourceLabel{bench:aws-4-k20-w15:row:mps-twostack-simd-block}{R14}\shortstack[l]{{\color{gray}R14}} & \begin{tabular}[t]{@{}l@{}}\rule{0pt}{12pt}MPS / Two-stack\\SIMD-DP / Block\end{tabular} & \hspace*{1pt}\begin{tabular}[t]{@{}c@{}}\rule{0pt}{12pt}51.6s\\[-.2pt]{\fontsize{5.8}{6.4}\selectfont \textcolor{ForestGreen}{$ \blacktriangle\,3.3\times $}}\\[-.2pt]9.8s\\[-.2pt]{\fontsize{5.8}{6.4}\selectfont \textcolor{ForestGreen}{$ \blacktriangle\,16.6\times $}}\end{tabular} & \hspace*{1pt}\begin{tabular}[t]{@{}c@{}}\rule{0pt}{12pt}44.7s\\[-.2pt]{\fontsize{5.8}{6.4}\selectfont \textcolor{ForestGreen}{$ \blacktriangle\,3.8\times $}}\\[-.2pt]9.6s\\[-.2pt]{\fontsize{5.8}{6.4}\selectfont \textcolor{ForestGreen}{$ \blacktriangle\,17.0\times $}}\end{tabular} & \hspace*{1pt}\begin{tabular}[t]{@{}c@{}}\rule{0pt}{12pt}1m\,16.1s\\[-.2pt]{\fontsize{5.8}{6.4}\selectfont \textcolor{ForestGreen}{$ \blacktriangle\,3.9\times $}}\\[-.2pt]15.7s\\[-.2pt]{\fontsize{5.8}{6.4}\selectfont \textcolor{ForestGreen}{$ \blacktriangle\,17.7\times $}}\end{tabular} & \textcolor{gray}{\ensuremath{\cdot}} & \textcolor{gray}{\ensuremath{\cdot}} & \textcolor{gray}{\ensuremath{\cdot}} & \textcolor{gray}{\ensuremath{\cdot}} & \textcolor{gray}{\ensuremath{\cdot}} & \textcolor{gray}{\ensuremath{\cdot}} & \textcolor{gray}{\ensuremath{\cdot}} & \textcolor{gray}{\ensuremath{\cdot}} & \textcolor{gray}{\ensuremath{\cdot}} & \textcolor{gray}{\ensuremath{\cdot}} & \textcolor{gray}{\ensuremath{\cdot}} & \textcolor{gray}{\ensuremath{\cdot}} & \textcolor{gray}{\ensuremath{\cdot}}\\[3pt]\hline
\BenchmarkSourceLabel{bench:aws-4-k20-w15:row:mps-twostack-rmq-block}{R15}\shortstack[l]{{\color{gray}R15}} & \begin{tabular}[t]{@{}l@{}}\rule{0pt}{12pt}MPS / Two-stack\\RMQ / Block\end{tabular} & \hspace*{1pt}\begin{tabular}[t]{@{}c@{}}\rule{0pt}{12pt}50.7s\\[-.2pt]{\fontsize{5.8}{6.4}\selectfont \textcolor{ForestGreen}{$ \blacktriangle\,3.3\times $}}\\[-.2pt]9.4s\\[-.2pt]{\fontsize{5.8}{6.4}\selectfont \textcolor{ForestGreen}{$ \blacktriangle\,17.2\times $}}\end{tabular} & \hspace*{1pt}\begin{tabular}[t]{@{}c@{}}\rule{0pt}{12pt}44.9s\\[-.2pt]{\fontsize{5.8}{6.4}\selectfont \textcolor{ForestGreen}{$ \blacktriangle\,3.8\times $}}\\[-.2pt]9.9s\\[-.2pt]{\fontsize{5.8}{6.4}\selectfont \textcolor{ForestGreen}{$ \blacktriangle\,16.5\times $}}\end{tabular} & \textcolor{gray}{\ensuremath{\cdot}} & \textcolor{gray}{\ensuremath{\cdot}} & \textcolor{gray}{\ensuremath{\cdot}} & \textcolor{gray}{\ensuremath{\cdot}} & \textcolor{gray}{\ensuremath{\cdot}} & \textcolor{gray}{\ensuremath{\cdot}} & \textcolor{gray}{\ensuremath{\cdot}} & \textcolor{gray}{\ensuremath{\cdot}} & \textcolor{gray}{\ensuremath{\cdot}} & \textcolor{gray}{\ensuremath{\cdot}} & \textcolor{gray}{\ensuremath{\cdot}} & \textcolor{gray}{\ensuremath{\cdot}} & \textcolor{gray}{\ensuremath{\cdot}} & \textcolor{gray}{\ensuremath{\cdot}}\\[3pt]\hline
\end{tabular}\par}
\smallskip{\footnotesize\noindent Each cell: A and B times, $\cdot$: no measurement. ---\textsuperscript{g}: GPU unavailable. n/a: no comparison.\par
\noindent Blue: one lowest A and one lowest B measurement for each column for each configuration.\par
\noindent Gray: minimap2. Yellow: my preferred MPF / Bounded / SIMD-DP / Block configuration.\par
\noindent MPF/MPS: minimizers per fragment/sequence. SIMD-ref: naive SIMD scan. Bounded: bounded SIMD algorithm.\par
}
\clearpage
\noindent{\large\bfseries Table \thetable: (C): AWS: 4 Gbp, k=20, w=15\hfill 3/7}\par\smallskip
{\fontsize{7.1}{8.1}\selectfont\setlength{\tabcolsep}{1pt}\setlength{\fboxsep}{2pt}\renewcommand{\arraystretch}{1.12}
\noindent\begin{tabular}{@{}p{14mm}p{27mm}|p{13.55mm}|p{13.55mm}|p{13.55mm}|p{13.55mm}|p{13.55mm}|p{13.55mm}|p{13.55mm}|p{13.55mm}|p{13.55mm}|p{13.55mm}|p{13.55mm}|p{13.55mm}|p{13.55mm}|p{13.55mm}|p{13.55mm}|p{13.55mm}|@{}}\hline
Engine & Execution & \multicolumn{8}{c|}{\textbf{Index}} & \multicolumn{8}{c|}{\textbf{No index}}\\
 &  & \multicolumn{8}{c|}{\textbf{4 Gbp}} & \multicolumn{8}{c|}{\textbf{4 Gbp}}\\
 &  & \multicolumn{8}{c|}{\textbf{k=20, w=15}} & \multicolumn{8}{c|}{\textbf{k=20, w=15}}\\
 &  & \multicolumn{4}{c|}{\textbf{DNA4}} & \multicolumn{4}{c|}{\textbf{DNA5}} & \multicolumn{4}{c|}{\textbf{DNA4}} & \multicolumn{4}{c|}{\textbf{DNA5}}\\
 &  & \multicolumn{1}{c|}{\textbf{M7A}} & \multicolumn{1}{c|}{\textbf{M7I}} & \multicolumn{1}{c|}{\textbf{M4}} & \multicolumn{1}{c|}{\textbf{C7G}} & \multicolumn{1}{c|}{\textbf{M7A}} & \multicolumn{1}{c|}{\textbf{M7I}} & \multicolumn{1}{c|}{\textbf{M4}} & \multicolumn{1}{c|}{\textbf{C7G}} & \multicolumn{1}{c|}{\textbf{M7A}} & \multicolumn{1}{c|}{\textbf{M7I}} & \multicolumn{1}{c|}{\textbf{M4}} & \multicolumn{1}{c|}{\textbf{C7G}} & \multicolumn{1}{c|}{\textbf{M7A}} & \multicolumn{1}{c|}{\textbf{M7I}} & \multicolumn{1}{c|}{\textbf{M4}} & \multicolumn{1}{c|}{\textbf{C7G}}\\
 & & {\color{gray}C1} & {\color{gray}C2} & {\color{gray}C3} & {\color{gray}C4} & {\color{gray}C5} & {\color{gray}C6} & {\color{gray}C7} & {\color{gray}C8} & {\color{gray}C9} & {\color{gray}C10} & {\color{gray}C11} & {\color{gray}C12} & {\color{gray}C13} & {\color{gray}C14} & {\color{gray}C15} & {\color{gray}C16}\\
\hline
\rowcolor{BenchBaseline}[0.5pt][0.5pt]\shortstack[l]{{\color{gray}R1}\\\textbf{minimap2}} & \textbf{Baseline} & \hspace*{1pt}\begin{tabular}[t]{@{}c@{}}\rule{0pt}{12pt}2m\,48.1s\\[-.2pt]{\fontsize{5.8}{6.4}\selectfont \textcolor{BrickRed}{$ \blacktriangledown\,20.4\times $}}\\[-.2pt]2m\,41.7s\\[-.2pt]{\fontsize{5.8}{6.4}\selectfont \textcolor{BrickRed}{$ \blacktriangledown\,22.6\times $}}\end{tabular} & \hspace*{1pt}\begin{tabular}[t]{@{}c@{}}\rule{0pt}{12pt}2m\,50.3s\\[-.2pt]{\fontsize{5.8}{6.4}\selectfont \textcolor{BrickRed}{$ \blacktriangledown\,19.8\times $}}\\[-.2pt]2m\,43s\\[-.2pt]{\fontsize{5.8}{6.4}\selectfont \textcolor{BrickRed}{$ \blacktriangledown\,22.2\times $}}\end{tabular} & \hspace*{1pt}\begin{tabular}[t]{@{}c@{}}\rule{0pt}{12pt}4m\,54.3s\\[-.2pt]{\fontsize{5.8}{6.4}\selectfont \textcolor{BrickRed}{$ \blacktriangledown\,43.2\times $}}\\[-.2pt]4m\,39s\\[-.2pt]{\fontsize{5.8}{6.4}\selectfont \textcolor{BrickRed}{$ \blacktriangledown\,54.6\times $}}\end{tabular} & \hspace*{1pt}\begin{tabular}[t]{@{}c@{}}\rule{0pt}{12pt}25m\,5.5s\\[-.2pt]{\fontsize{5.8}{6.4}\selectfont \textcolor{BrickRed}{$ \blacktriangledown\,41.8\times $}}\\[-.2pt]25m\,0.7s\\[-.2pt]{\fontsize{5.8}{6.4}\selectfont \textcolor{BrickRed}{$ \blacktriangledown\,46.6\times $}}\end{tabular} & \hspace*{1pt}\begin{tabular}[t]{@{}c@{}}\rule{0pt}{12pt}33.9s\\[-.2pt]{\fontsize{5.8}{6.4}\selectfont \textcolor{BrickRed}{$ \blacktriangledown\,2.3\times $}}\\[-.2pt]32.9s\\[-.2pt]{\fontsize{5.8}{6.4}\selectfont \textcolor{BrickRed}{$ \blacktriangledown\,3.1\times $}}\end{tabular} & \hspace*{1pt}\begin{tabular}[t]{@{}c@{}}\rule{0pt}{12pt}34.4s\\[-.2pt]{\fontsize{5.8}{6.4}\selectfont \textcolor{BrickRed}{$ \blacktriangledown\,2.6\times $}}\\[-.2pt]32.7s\\[-.2pt]{\fontsize{5.8}{6.4}\selectfont \textcolor{BrickRed}{$ \blacktriangledown\,2.9\times $}}\end{tabular} & \hspace*{1pt}\begin{tabular}[t]{@{}c@{}}\rule{0pt}{12pt}1m\,8.2s\\[-.2pt]{\fontsize{5.8}{6.4}\selectfont \textcolor{BrickRed}{$ \blacktriangledown\,3.7\times $}}\\[-.2pt]1m\,5.1s\\[-.2pt]{\fontsize{5.8}{6.4}\selectfont \textcolor{BrickRed}{$ \blacktriangledown\,4.6\times $}}\end{tabular} & \hspace*{1pt}\begin{tabular}[t]{@{}c@{}}\rule{0pt}{12pt}3m\,9.8s\\[-.2pt]{\fontsize{5.8}{6.4}\selectfont \textcolor{BrickRed}{$ \blacktriangledown\,3.9\times $}}\\[-.2pt]3m\,3.7s\\[-.2pt]{\fontsize{5.8}{6.4}\selectfont \textcolor{BrickRed}{$ \blacktriangledown\,4.1\times $}}\end{tabular} & \hspace*{1pt}\begin{tabular}[t]{@{}c@{}}\rule{0pt}{12pt}3m\,28.3s\\[-.2pt]{\fontsize{5.8}{6.4}\selectfont \textcolor{BrickRed}{$ \blacktriangledown\,1.3\times $}}\\[-.2pt]3m\,22.8s\\[-.2pt]{\fontsize{5.8}{6.4}\selectfont \textcolor{BrickRed}{$ \blacktriangledown\,1.4\times $}}\end{tabular} & \hspace*{1pt}\begin{tabular}[t]{@{}c@{}}\rule{0pt}{12pt}3m\,53.2s\\[-.2pt]{\fontsize{5.8}{6.4}\selectfont \textcolor{BrickRed}{$ \blacktriangledown\,1.4\times $}}\\[-.2pt]3m\,43.5s\\[-.2pt]{\fontsize{5.8}{6.4}\selectfont \textcolor{BrickRed}{$ \blacktriangledown\,1.5\times $}}\end{tabular} & \hspace*{1pt}\begin{tabular}[t]{@{}c@{}}\rule{0pt}{12pt}7m\,14.6s\\[-.2pt]{\fontsize{5.8}{6.4}\selectfont \textcolor{BrickRed}{$ \blacktriangledown\,1.5\times $}}\\[-.2pt]6m\,57.2s\\[-.2pt]{\fontsize{5.8}{6.4}\selectfont \textcolor{BrickRed}{$ \blacktriangledown\,1.6\times $}}\end{tabular} & \hspace*{1pt}\begin{tabular}[t]{@{}c@{}}\rule{0pt}{12pt}51m\,22s\\[-.2pt]{\fontsize{5.8}{6.4}\selectfont \textcolor{gray}{n/a}}\\[-.2pt]51m\,17.2s\\[-.2pt]{\fontsize{5.8}{6.4}\selectfont \textcolor{gray}{n/a}}\end{tabular} & \hspace*{1pt}\begin{tabular}[t]{@{}c@{}}\rule{0pt}{12pt}1m\,10.3s\\[-.2pt]{\fontsize{5.8}{6.4}\selectfont \textcolor{ForestGreen}{$ \blacktriangle\,7.6\times $}}\\[-.2pt]1m\,9s\\[-.2pt]{\fontsize{5.8}{6.4}\selectfont \textcolor{ForestGreen}{$ \blacktriangle\,2.4\times $}}\end{tabular} & \hspace*{1pt}\begin{tabular}[t]{@{}c@{}}\rule{0pt}{12pt}1m\,8.6s\\[-.2pt]{\fontsize{5.8}{6.4}\selectfont \textcolor{ForestGreen}{$ \blacktriangle\,6.4\times $}}\\[-.2pt]1m\,6.9s\\[-.2pt]{\fontsize{5.8}{6.4}\selectfont \textcolor{ForestGreen}{$ \blacktriangle\,2.3\times $}}\end{tabular} & \hspace*{1pt}\begin{tabular}[t]{@{}c@{}}\rule{0pt}{12pt}2m\,21.5s\\[-.2pt]{\fontsize{5.8}{6.4}\selectfont \textcolor{ForestGreen}{$ \blacktriangle\,6.0\times $}}\\[-.2pt]2m\,18.4s\\[-.2pt]{\fontsize{5.8}{6.4}\selectfont \textcolor{ForestGreen}{$ \blacktriangle\,2.2\times $}}\end{tabular} & \hspace*{1pt}\begin{tabular}[t]{@{}c@{}}\rule{0pt}{12pt}15m\,56.8s\\[-.2pt]{\fontsize{5.8}{6.4}\selectfont \textcolor{gray}{n/a}}\\[-.2pt]15m\,51.9s\\[-.2pt]{\fontsize{5.8}{6.4}\selectfont \textcolor{gray}{n/a}}\end{tabular}\\[3pt]\hline
\BenchmarkSourceLabel{bench:aws-4-k20-w15:row:mps-twostack-rmq-ksw2}{R16}\shortstack[l]{{\color{gray}R16}\\\textbf{BioseqDB}} & \begin{tabular}[t]{@{}l@{}}\rule{0pt}{12pt}MPS / Two-stack\\RMQ / KSW2\end{tabular} & \hspace*{1pt}\begin{tabular}[t]{@{}c@{}}\rule{0pt}{12pt}53.5s\\[-.2pt]{\fontsize{5.8}{6.4}\selectfont \textcolor{ForestGreen}{$ \blacktriangle\,3.1\times $}}\\[-.2pt]9.3s\\[-.2pt]{\fontsize{5.8}{6.4}\selectfont \textcolor{ForestGreen}{$ \blacktriangle\,17.4\times $}}\end{tabular} & \hspace*{1pt}\begin{tabular}[t]{@{}c@{}}\rule{0pt}{12pt}47.1s\\[-.2pt]{\fontsize{5.8}{6.4}\selectfont \textcolor{ForestGreen}{$ \blacktriangle\,3.6\times $}}\\[-.2pt]9s\\[-.2pt]{\fontsize{5.8}{6.4}\selectfont \textcolor{ForestGreen}{$ \blacktriangle\,18.2\times $}}\end{tabular} & \textcolor{gray}{\ensuremath{\cdot}} & \textcolor{gray}{\ensuremath{\cdot}} & \textcolor{gray}{\ensuremath{\cdot}} & \textcolor{gray}{\ensuremath{\cdot}} & \textcolor{gray}{\ensuremath{\cdot}} & \textcolor{gray}{\ensuremath{\cdot}} & \textcolor{gray}{\ensuremath{\cdot}} & \textcolor{gray}{\ensuremath{\cdot}} & \textcolor{gray}{\ensuremath{\cdot}} & \textcolor{gray}{\ensuremath{\cdot}} & \textcolor{gray}{\ensuremath{\cdot}} & \textcolor{gray}{\ensuremath{\cdot}} & \textcolor{gray}{\ensuremath{\cdot}} & \textcolor{gray}{\ensuremath{\cdot}}\\[3pt]\hline
\BenchmarkSourceLabel{bench:aws-4-k20-w15:row:mps-twostack-simd-ksw2}{R17}\shortstack[l]{{\color{gray}R17}} & \begin{tabular}[t]{@{}l@{}}\rule{0pt}{12pt}MPS / Two-stack\\SIMD-DP / KSW2\end{tabular} & \hspace*{1pt}\begin{tabular}[t]{@{}c@{}}\rule{0pt}{12pt}55.3s\\[-.2pt]{\fontsize{5.8}{6.4}\selectfont \textcolor{ForestGreen}{$ \blacktriangle\,3.0\times $}}\\[-.2pt]9.6s\\[-.2pt]{\fontsize{5.8}{6.4}\selectfont \textcolor{ForestGreen}{$ \blacktriangle\,16.8\times $}}\end{tabular} & \hspace*{1pt}\begin{tabular}[t]{@{}c@{}}\rule{0pt}{12pt}48.6s\\[-.2pt]{\fontsize{5.8}{6.4}\selectfont \textcolor{ForestGreen}{$ \blacktriangle\,3.5\times $}}\\[-.2pt]9.6s\\[-.2pt]{\fontsize{5.8}{6.4}\selectfont \textcolor{ForestGreen}{$ \blacktriangle\,17.0\times $}}\end{tabular} & \textcolor{gray}{\ensuremath{\cdot}} & \textcolor{gray}{\ensuremath{\cdot}} & \textcolor{gray}{\ensuremath{\cdot}} & \textcolor{gray}{\ensuremath{\cdot}} & \textcolor{gray}{\ensuremath{\cdot}} & \textcolor{gray}{\ensuremath{\cdot}} & \textcolor{gray}{\ensuremath{\cdot}} & \textcolor{gray}{\ensuremath{\cdot}} & \textcolor{gray}{\ensuremath{\cdot}} & \textcolor{gray}{\ensuremath{\cdot}} & \textcolor{gray}{\ensuremath{\cdot}} & \textcolor{gray}{\ensuremath{\cdot}} & \textcolor{gray}{\ensuremath{\cdot}} & \textcolor{gray}{\ensuremath{\cdot}}\\[3pt]\hline
\BenchmarkSourceLabel{bench:aws-4-k20-w15:row:mps-twostack-scalar-ksw2}{R18}\shortstack[l]{{\color{gray}R18}} & \begin{tabular}[t]{@{}l@{}}\rule{0pt}{12pt}MPS / Two-stack\\Scalar-DP / KSW2\end{tabular} & \hspace*{1pt}\begin{tabular}[t]{@{}c@{}}\rule{0pt}{12pt}57.2s\\[-.2pt]{\fontsize{5.8}{6.4}\selectfont \textcolor{ForestGreen}{$ \blacktriangle\,2.9\times $}}\\[-.2pt]9.8s\\[-.2pt]{\fontsize{5.8}{6.4}\selectfont \textcolor{ForestGreen}{$ \blacktriangle\,16.5\times $}}\end{tabular} & \hspace*{1pt}\begin{tabular}[t]{@{}c@{}}\rule{0pt}{12pt}49.4s\\[-.2pt]{\fontsize{5.8}{6.4}\selectfont \textcolor{ForestGreen}{$ \blacktriangle\,3.4\times $}}\\[-.2pt]9.2s\\[-.2pt]{\fontsize{5.8}{6.4}\selectfont \textcolor{ForestGreen}{$ \blacktriangle\,17.8\times $}}\end{tabular} & \textcolor{gray}{\ensuremath{\cdot}} & \textcolor{gray}{\ensuremath{\cdot}} & \textcolor{gray}{\ensuremath{\cdot}} & \textcolor{gray}{\ensuremath{\cdot}} & \textcolor{gray}{\ensuremath{\cdot}} & \textcolor{gray}{\ensuremath{\cdot}} & \textcolor{gray}{\ensuremath{\cdot}} & \textcolor{gray}{\ensuremath{\cdot}} & \textcolor{gray}{\ensuremath{\cdot}} & \textcolor{gray}{\ensuremath{\cdot}} & \textcolor{gray}{\ensuremath{\cdot}} & \textcolor{gray}{\ensuremath{\cdot}} & \textcolor{gray}{\ensuremath{\cdot}} & \textcolor{gray}{\ensuremath{\cdot}}\\[3pt]\hline
\BenchmarkSourceLabel{bench:aws-4-k20-w15:row:mps-scalar-simd-block}{R19}\shortstack[l]{{\color{gray}R19}} & \begin{tabular}[t]{@{}l@{}}\rule{0pt}{12pt}MPS / Scalar\\SIMD-DP / Block\end{tabular} & \hspace*{1pt}\begin{tabular}[t]{@{}c@{}}\rule{0pt}{12pt}57s\\[-.2pt]{\fontsize{5.8}{6.4}\selectfont \textcolor{ForestGreen}{$ \blacktriangle\,3.0\times $}}\\[-.2pt]14.3s\\[-.2pt]{\fontsize{5.8}{6.4}\selectfont \textcolor{ForestGreen}{$ \blacktriangle\,11.3\times $}}\end{tabular} & \hspace*{1pt}\begin{tabular}[t]{@{}c@{}}\rule{0pt}{12pt}50.6s\\[-.2pt]{\fontsize{5.8}{6.4}\selectfont \textcolor{ForestGreen}{$ \blacktriangle\,3.4\times $}}\\[-.2pt]15.1s\\[-.2pt]{\fontsize{5.8}{6.4}\selectfont \textcolor{ForestGreen}{$ \blacktriangle\,10.8\times $}}\end{tabular} & \textcolor{gray}{\ensuremath{\cdot}} & \textcolor{gray}{\ensuremath{\cdot}} & \hspace*{1pt}\begin{tabular}[t]{@{}c@{}}\rule{0pt}{12pt}7m\,14.5s\\[-.2pt]{\fontsize{5.8}{6.4}\selectfont \textcolor{BrickRed}{$ \blacktriangledown\,12.8\times $}}\\[-.2pt]41.9s\\[-.2pt]{\fontsize{5.8}{6.4}\selectfont \textcolor{BrickRed}{$ \blacktriangledown\,27\% $}}\end{tabular} & \hspace*{1pt}\begin{tabular}[t]{@{}c@{}}\rule{0pt}{12pt}6m\,19.3s\\[-.2pt]{\fontsize{5.8}{6.4}\selectfont \textcolor{BrickRed}{$ \blacktriangledown\,11.0\times $}}\\[-.2pt]41.3s\\[-.2pt]{\fontsize{5.8}{6.4}\selectfont \textcolor{BrickRed}{$ \blacktriangledown\,26\% $}}\end{tabular} & \textcolor{gray}{\ensuremath{\cdot}} & \textcolor{gray}{\ensuremath{\cdot}} & \hspace*{1pt}\begin{tabular}[t]{@{}c@{}}\rule{0pt}{12pt}5m\,10.1s\\[-.2pt]{\fontsize{5.8}{6.4}\selectfont \textcolor{BrickRed}{$ \blacktriangledown\,1.5\times $}}\\[-.2pt]4m\,27.6s\\[-.2pt]{\fontsize{5.8}{6.4}\selectfont \textcolor{BrickRed}{$ \blacktriangledown\,1.3\times $}}\end{tabular} & \hspace*{1pt}\begin{tabular}[t]{@{}c@{}}\rule{0pt}{12pt}5m\,3.5s\\[-.2pt]{\fontsize{5.8}{6.4}\selectfont \textcolor{BrickRed}{$ \blacktriangledown\,1.3\times $}}\\[-.2pt]4m\,28.9s\\[-.2pt]{\fontsize{5.8}{6.4}\selectfont \textcolor{BrickRed}{$ \blacktriangledown\,20\% $}}\end{tabular} & \textcolor{gray}{\ensuremath{\cdot}} & \textcolor{gray}{\ensuremath{\cdot}} & \textcolor{gray}{\ensuremath{\cdot}} & \textcolor{gray}{\ensuremath{\cdot}} & \textcolor{gray}{\ensuremath{\cdot}} & \textcolor{gray}{\ensuremath{\cdot}}\\[3pt]\hline
\BenchmarkSourceLabel{bench:aws-4-k20-w15:row:mps-simdref-scalar-block}{R20}\shortstack[l]{{\color{gray}R20}} & \begin{tabular}[t]{@{}l@{}}\rule{0pt}{12pt}MPS / SIMD-ref\\Scalar-DP / Block\end{tabular} & \hspace*{1pt}\begin{tabular}[t]{@{}c@{}}\rule{0pt}{12pt}1m\,1.1s\\[-.2pt]{\fontsize{5.8}{6.4}\selectfont \textcolor{ForestGreen}{$ \blacktriangle\,2.7\times $}}\\[-.2pt]14.9s\\[-.2pt]{\fontsize{5.8}{6.4}\selectfont \textcolor{ForestGreen}{$ \blacktriangle\,10.8\times $}}\end{tabular} & \hspace*{1pt}\begin{tabular}[t]{@{}c@{}}\rule{0pt}{12pt}50.8s\\[-.2pt]{\fontsize{5.8}{6.4}\selectfont \textcolor{ForestGreen}{$ \blacktriangle\,3.4\times $}}\\[-.2pt]14.6s\\[-.2pt]{\fontsize{5.8}{6.4}\selectfont \textcolor{ForestGreen}{$ \blacktriangle\,11.2\times $}}\end{tabular} & \textcolor{gray}{\ensuremath{\cdot}} & \textcolor{gray}{\ensuremath{\cdot}} & \hspace*{1pt}\begin{tabular}[t]{@{}c@{}}\rule{0pt}{12pt}7m\,13.1s\\[-.2pt]{\fontsize{5.8}{6.4}\selectfont \textcolor{BrickRed}{$ \blacktriangledown\,12.8\times $}}\\[-.2pt]52.7s\\[-.2pt]{\fontsize{5.8}{6.4}\selectfont \textcolor{BrickRed}{$ \blacktriangledown\,1.6\times $}}\end{tabular} & \hspace*{1pt}\begin{tabular}[t]{@{}c@{}}\rule{0pt}{12pt}6m\,7.7s\\[-.2pt]{\fontsize{5.8}{6.4}\selectfont \textcolor{BrickRed}{$ \blacktriangledown\,10.7\times $}}\\[-.2pt]50.4s\\[-.2pt]{\fontsize{5.8}{6.4}\selectfont \textcolor{BrickRed}{$ \blacktriangledown\,1.5\times $}}\end{tabular} & \textcolor{gray}{\ensuremath{\cdot}} & \textcolor{gray}{\ensuremath{\cdot}} & \textcolor{gray}{\ensuremath{\cdot}} & \textcolor{gray}{\ensuremath{\cdot}} & \textcolor{gray}{\ensuremath{\cdot}} & \textcolor{gray}{\ensuremath{\cdot}} & \textcolor{gray}{\ensuremath{\cdot}} & \textcolor{gray}{\ensuremath{\cdot}} & \textcolor{gray}{\ensuremath{\cdot}} & \textcolor{gray}{\ensuremath{\cdot}}\\[3pt]\hline
\BenchmarkSourceLabel{bench:aws-4-k20-w15:row:mps-simdref-rmq-block}{R21}\shortstack[l]{{\color{gray}R21}} & \begin{tabular}[t]{@{}l@{}}\rule{0pt}{12pt}MPS / SIMD-ref\\RMQ / Block\end{tabular} & \hspace*{1pt}\begin{tabular}[t]{@{}c@{}}\rule{0pt}{12pt}1m\,0.8s\\[-.2pt]{\fontsize{5.8}{6.4}\selectfont \textcolor{ForestGreen}{$ \blacktriangle\,2.8\times $}}\\[-.2pt]15s\\[-.2pt]{\fontsize{5.8}{6.4}\selectfont \textcolor{ForestGreen}{$ \blacktriangle\,10.8\times $}}\end{tabular} & \hspace*{1pt}\begin{tabular}[t]{@{}c@{}}\rule{0pt}{12pt}50.9s\\[-.2pt]{\fontsize{5.8}{6.4}\selectfont \textcolor{ForestGreen}{$ \blacktriangle\,3.3\times $}}\\[-.2pt]14.7s\\[-.2pt]{\fontsize{5.8}{6.4}\selectfont \textcolor{ForestGreen}{$ \blacktriangle\,11.1\times $}}\end{tabular} & \textcolor{gray}{\ensuremath{\cdot}} & \textcolor{gray}{\ensuremath{\cdot}} & \hspace*{1pt}\begin{tabular}[t]{@{}c@{}}\rule{0pt}{12pt}7m\,13.3s\\[-.2pt]{\fontsize{5.8}{6.4}\selectfont \textcolor{BrickRed}{$ \blacktriangledown\,12.8\times $}}\\[-.2pt]53.4s\\[-.2pt]{\fontsize{5.8}{6.4}\selectfont \textcolor{BrickRed}{$ \blacktriangledown\,1.6\times $}}\end{tabular} & \hspace*{1pt}\begin{tabular}[t]{@{}c@{}}\rule{0pt}{12pt}6m\,11.5s\\[-.2pt]{\fontsize{5.8}{6.4}\selectfont \textcolor{BrickRed}{$ \blacktriangledown\,10.8\times $}}\\[-.2pt]54s\\[-.2pt]{\fontsize{5.8}{6.4}\selectfont \textcolor{BrickRed}{$ \blacktriangledown\,1.7\times $}}\end{tabular} & \textcolor{gray}{\ensuremath{\cdot}} & \textcolor{gray}{\ensuremath{\cdot}} & \textcolor{gray}{\ensuremath{\cdot}} & \textcolor{gray}{\ensuremath{\cdot}} & \textcolor{gray}{\ensuremath{\cdot}} & \textcolor{gray}{\ensuremath{\cdot}} & \textcolor{gray}{\ensuremath{\cdot}} & \textcolor{gray}{\ensuremath{\cdot}} & \textcolor{gray}{\ensuremath{\cdot}} & \textcolor{gray}{\ensuremath{\cdot}}\\[3pt]\hline
\BenchmarkSourceLabel{bench:aws-4-k20-w15:row:mps-scalar-rmq-block}{R22}\shortstack[l]{{\color{gray}R22}} & \begin{tabular}[t]{@{}l@{}}\rule{0pt}{12pt}MPS / Scalar\\RMQ / Block\end{tabular} & \hspace*{1pt}\begin{tabular}[t]{@{}c@{}}\rule{0pt}{12pt}57.4s\\[-.2pt]{\fontsize{5.8}{6.4}\selectfont \textcolor{ForestGreen}{$ \blacktriangle\,2.9\times $}}\\[-.2pt]14.6s\\[-.2pt]{\fontsize{5.8}{6.4}\selectfont \textcolor{ForestGreen}{$ \blacktriangle\,11.1\times $}}\end{tabular} & \hspace*{1pt}\begin{tabular}[t]{@{}c@{}}\rule{0pt}{12pt}51.1s\\[-.2pt]{\fontsize{5.8}{6.4}\selectfont \textcolor{ForestGreen}{$ \blacktriangle\,3.3\times $}}\\[-.2pt]15.9s\\[-.2pt]{\fontsize{5.8}{6.4}\selectfont \textcolor{ForestGreen}{$ \blacktriangle\,10.3\times $}}\end{tabular} & \textcolor{gray}{\ensuremath{\cdot}} & \textcolor{gray}{\ensuremath{\cdot}} & \hspace*{1pt}\begin{tabular}[t]{@{}c@{}}\rule{0pt}{12pt}7m\,2.9s\\[-.2pt]{\fontsize{5.8}{6.4}\selectfont \textcolor{BrickRed}{$ \blacktriangledown\,12.5\times $}}\\[-.2pt]38.9s\\[-.2pt]{\fontsize{5.8}{6.4}\selectfont \textcolor{BrickRed}{$ \blacktriangledown\,18\% $}}\end{tabular} & \hspace*{1pt}\begin{tabular}[t]{@{}c@{}}\rule{0pt}{12pt}6m\,25.9s\\[-.2pt]{\fontsize{5.8}{6.4}\selectfont \textcolor{BrickRed}{$ \blacktriangledown\,11.2\times $}}\\[-.2pt]44s\\[-.2pt]{\fontsize{5.8}{6.4}\selectfont \textcolor{BrickRed}{$ \blacktriangledown\,1.3\times $}}\end{tabular} & \textcolor{gray}{\ensuremath{\cdot}} & \textcolor{gray}{\ensuremath{\cdot}} & \hspace*{1pt}\begin{tabular}[t]{@{}c@{}}\rule{0pt}{12pt}5m\,10.6s\\[-.2pt]{\fontsize{5.8}{6.4}\selectfont \textcolor{BrickRed}{$ \blacktriangledown\,1.5\times $}}\\[-.2pt]4m\,28.1s\\[-.2pt]{\fontsize{5.8}{6.4}\selectfont \textcolor{BrickRed}{$ \blacktriangledown\,1.3\times $}}\end{tabular} & \hspace*{1pt}\begin{tabular}[t]{@{}c@{}}\rule{0pt}{12pt}5m\,37s\\[-.2pt]{\fontsize{5.8}{6.4}\selectfont \textcolor{BrickRed}{$ \blacktriangledown\,1.4\times $}}\\[-.2pt]5m\,0.7s\\[-.2pt]{\fontsize{5.8}{6.4}\selectfont \textcolor{BrickRed}{$ \blacktriangledown\,1.3\times $}}\end{tabular} & \textcolor{gray}{\ensuremath{\cdot}} & \textcolor{gray}{\ensuremath{\cdot}} & \textcolor{gray}{\ensuremath{\cdot}} & \textcolor{gray}{\ensuremath{\cdot}} & \textcolor{gray}{\ensuremath{\cdot}} & \textcolor{gray}{\ensuremath{\cdot}}\\[3pt]\hline
\end{tabular}\par}
\smallskip{\footnotesize\noindent Each cell: A and B times, $\cdot$: no measurement. ---\textsuperscript{g}: GPU unavailable. n/a: no comparison.\par
\noindent Blue: one lowest A and one lowest B measurement for each column for each configuration.\par
\noindent Gray: minimap2. Yellow: my preferred MPF / Bounded / SIMD-DP / Block configuration.\par
\noindent MPF/MPS: minimizers per fragment/sequence. SIMD-ref: naive SIMD scan. Bounded: bounded SIMD algorithm.\par
}
\clearpage
\noindent{\large\bfseries Table \thetable: (C): AWS: 4 Gbp, k=20, w=15\hfill 4/7}\par\smallskip
{\fontsize{7.1}{8.1}\selectfont\setlength{\tabcolsep}{1pt}\setlength{\fboxsep}{2pt}\renewcommand{\arraystretch}{1.12}
\noindent\begin{tabular}{@{}p{14mm}p{27mm}|p{13.55mm}|p{13.55mm}|p{13.55mm}|p{13.55mm}|p{13.55mm}|p{13.55mm}|p{13.55mm}|p{13.55mm}|p{13.55mm}|p{13.55mm}|p{13.55mm}|p{13.55mm}|p{13.55mm}|p{13.55mm}|p{13.55mm}|p{13.55mm}|@{}}\hline
Engine & Execution & \multicolumn{8}{c|}{\textbf{Index}} & \multicolumn{8}{c|}{\textbf{No index}}\\
 &  & \multicolumn{8}{c|}{\textbf{4 Gbp}} & \multicolumn{8}{c|}{\textbf{4 Gbp}}\\
 &  & \multicolumn{8}{c|}{\textbf{k=20, w=15}} & \multicolumn{8}{c|}{\textbf{k=20, w=15}}\\
 &  & \multicolumn{4}{c|}{\textbf{DNA4}} & \multicolumn{4}{c|}{\textbf{DNA5}} & \multicolumn{4}{c|}{\textbf{DNA4}} & \multicolumn{4}{c|}{\textbf{DNA5}}\\
 &  & \multicolumn{1}{c|}{\textbf{M7A}} & \multicolumn{1}{c|}{\textbf{M7I}} & \multicolumn{1}{c|}{\textbf{M4}} & \multicolumn{1}{c|}{\textbf{C7G}} & \multicolumn{1}{c|}{\textbf{M7A}} & \multicolumn{1}{c|}{\textbf{M7I}} & \multicolumn{1}{c|}{\textbf{M4}} & \multicolumn{1}{c|}{\textbf{C7G}} & \multicolumn{1}{c|}{\textbf{M7A}} & \multicolumn{1}{c|}{\textbf{M7I}} & \multicolumn{1}{c|}{\textbf{M4}} & \multicolumn{1}{c|}{\textbf{C7G}} & \multicolumn{1}{c|}{\textbf{M7A}} & \multicolumn{1}{c|}{\textbf{M7I}} & \multicolumn{1}{c|}{\textbf{M4}} & \multicolumn{1}{c|}{\textbf{C7G}}\\
 & & {\color{gray}C1} & {\color{gray}C2} & {\color{gray}C3} & {\color{gray}C4} & {\color{gray}C5} & {\color{gray}C6} & {\color{gray}C7} & {\color{gray}C8} & {\color{gray}C9} & {\color{gray}C10} & {\color{gray}C11} & {\color{gray}C12} & {\color{gray}C13} & {\color{gray}C14} & {\color{gray}C15} & {\color{gray}C16}\\
\hline
\rowcolor{BenchBaseline}[0.5pt][0.5pt]\shortstack[l]{{\color{gray}R1}\\\textbf{minimap2}} & \textbf{Baseline} & \hspace*{1pt}\begin{tabular}[t]{@{}c@{}}\rule{0pt}{12pt}2m\,48.1s\\[-.2pt]{\fontsize{5.8}{6.4}\selectfont \textcolor{BrickRed}{$ \blacktriangledown\,20.4\times $}}\\[-.2pt]2m\,41.7s\\[-.2pt]{\fontsize{5.8}{6.4}\selectfont \textcolor{BrickRed}{$ \blacktriangledown\,22.6\times $}}\end{tabular} & \hspace*{1pt}\begin{tabular}[t]{@{}c@{}}\rule{0pt}{12pt}2m\,50.3s\\[-.2pt]{\fontsize{5.8}{6.4}\selectfont \textcolor{BrickRed}{$ \blacktriangledown\,19.8\times $}}\\[-.2pt]2m\,43s\\[-.2pt]{\fontsize{5.8}{6.4}\selectfont \textcolor{BrickRed}{$ \blacktriangledown\,22.2\times $}}\end{tabular} & \hspace*{1pt}\begin{tabular}[t]{@{}c@{}}\rule{0pt}{12pt}4m\,54.3s\\[-.2pt]{\fontsize{5.8}{6.4}\selectfont \textcolor{BrickRed}{$ \blacktriangledown\,43.2\times $}}\\[-.2pt]4m\,39s\\[-.2pt]{\fontsize{5.8}{6.4}\selectfont \textcolor{BrickRed}{$ \blacktriangledown\,54.6\times $}}\end{tabular} & \hspace*{1pt}\begin{tabular}[t]{@{}c@{}}\rule{0pt}{12pt}25m\,5.5s\\[-.2pt]{\fontsize{5.8}{6.4}\selectfont \textcolor{BrickRed}{$ \blacktriangledown\,41.8\times $}}\\[-.2pt]25m\,0.7s\\[-.2pt]{\fontsize{5.8}{6.4}\selectfont \textcolor{BrickRed}{$ \blacktriangledown\,46.6\times $}}\end{tabular} & \hspace*{1pt}\begin{tabular}[t]{@{}c@{}}\rule{0pt}{12pt}33.9s\\[-.2pt]{\fontsize{5.8}{6.4}\selectfont \textcolor{BrickRed}{$ \blacktriangledown\,2.3\times $}}\\[-.2pt]32.9s\\[-.2pt]{\fontsize{5.8}{6.4}\selectfont \textcolor{BrickRed}{$ \blacktriangledown\,3.1\times $}}\end{tabular} & \hspace*{1pt}\begin{tabular}[t]{@{}c@{}}\rule{0pt}{12pt}34.4s\\[-.2pt]{\fontsize{5.8}{6.4}\selectfont \textcolor{BrickRed}{$ \blacktriangledown\,2.6\times $}}\\[-.2pt]32.7s\\[-.2pt]{\fontsize{5.8}{6.4}\selectfont \textcolor{BrickRed}{$ \blacktriangledown\,2.9\times $}}\end{tabular} & \hspace*{1pt}\begin{tabular}[t]{@{}c@{}}\rule{0pt}{12pt}1m\,8.2s\\[-.2pt]{\fontsize{5.8}{6.4}\selectfont \textcolor{BrickRed}{$ \blacktriangledown\,3.7\times $}}\\[-.2pt]1m\,5.1s\\[-.2pt]{\fontsize{5.8}{6.4}\selectfont \textcolor{BrickRed}{$ \blacktriangledown\,4.6\times $}}\end{tabular} & \hspace*{1pt}\begin{tabular}[t]{@{}c@{}}\rule{0pt}{12pt}3m\,9.8s\\[-.2pt]{\fontsize{5.8}{6.4}\selectfont \textcolor{BrickRed}{$ \blacktriangledown\,3.9\times $}}\\[-.2pt]3m\,3.7s\\[-.2pt]{\fontsize{5.8}{6.4}\selectfont \textcolor{BrickRed}{$ \blacktriangledown\,4.1\times $}}\end{tabular} & \hspace*{1pt}\begin{tabular}[t]{@{}c@{}}\rule{0pt}{12pt}3m\,28.3s\\[-.2pt]{\fontsize{5.8}{6.4}\selectfont \textcolor{BrickRed}{$ \blacktriangledown\,1.3\times $}}\\[-.2pt]3m\,22.8s\\[-.2pt]{\fontsize{5.8}{6.4}\selectfont \textcolor{BrickRed}{$ \blacktriangledown\,1.4\times $}}\end{tabular} & \hspace*{1pt}\begin{tabular}[t]{@{}c@{}}\rule{0pt}{12pt}3m\,53.2s\\[-.2pt]{\fontsize{5.8}{6.4}\selectfont \textcolor{BrickRed}{$ \blacktriangledown\,1.4\times $}}\\[-.2pt]3m\,43.5s\\[-.2pt]{\fontsize{5.8}{6.4}\selectfont \textcolor{BrickRed}{$ \blacktriangledown\,1.5\times $}}\end{tabular} & \hspace*{1pt}\begin{tabular}[t]{@{}c@{}}\rule{0pt}{12pt}7m\,14.6s\\[-.2pt]{\fontsize{5.8}{6.4}\selectfont \textcolor{BrickRed}{$ \blacktriangledown\,1.5\times $}}\\[-.2pt]6m\,57.2s\\[-.2pt]{\fontsize{5.8}{6.4}\selectfont \textcolor{BrickRed}{$ \blacktriangledown\,1.6\times $}}\end{tabular} & \hspace*{1pt}\begin{tabular}[t]{@{}c@{}}\rule{0pt}{12pt}51m\,22s\\[-.2pt]{\fontsize{5.8}{6.4}\selectfont \textcolor{gray}{n/a}}\\[-.2pt]51m\,17.2s\\[-.2pt]{\fontsize{5.8}{6.4}\selectfont \textcolor{gray}{n/a}}\end{tabular} & \hspace*{1pt}\begin{tabular}[t]{@{}c@{}}\rule{0pt}{12pt}1m\,10.3s\\[-.2pt]{\fontsize{5.8}{6.4}\selectfont \textcolor{ForestGreen}{$ \blacktriangle\,7.6\times $}}\\[-.2pt]1m\,9s\\[-.2pt]{\fontsize{5.8}{6.4}\selectfont \textcolor{ForestGreen}{$ \blacktriangle\,2.4\times $}}\end{tabular} & \hspace*{1pt}\begin{tabular}[t]{@{}c@{}}\rule{0pt}{12pt}1m\,8.6s\\[-.2pt]{\fontsize{5.8}{6.4}\selectfont \textcolor{ForestGreen}{$ \blacktriangle\,6.4\times $}}\\[-.2pt]1m\,6.9s\\[-.2pt]{\fontsize{5.8}{6.4}\selectfont \textcolor{ForestGreen}{$ \blacktriangle\,2.3\times $}}\end{tabular} & \hspace*{1pt}\begin{tabular}[t]{@{}c@{}}\rule{0pt}{12pt}2m\,21.5s\\[-.2pt]{\fontsize{5.8}{6.4}\selectfont \textcolor{ForestGreen}{$ \blacktriangle\,6.0\times $}}\\[-.2pt]2m\,18.4s\\[-.2pt]{\fontsize{5.8}{6.4}\selectfont \textcolor{ForestGreen}{$ \blacktriangle\,2.2\times $}}\end{tabular} & \hspace*{1pt}\begin{tabular}[t]{@{}c@{}}\rule{0pt}{12pt}15m\,56.8s\\[-.2pt]{\fontsize{5.8}{6.4}\selectfont \textcolor{gray}{n/a}}\\[-.2pt]15m\,51.9s\\[-.2pt]{\fontsize{5.8}{6.4}\selectfont \textcolor{gray}{n/a}}\end{tabular}\\[3pt]\hline
\BenchmarkSourceLabel{bench:aws-4-k20-w15:row:mps-simdref-simd-block}{R23}\shortstack[l]{{\color{gray}R23}\\\textbf{BioseqDB}} & \begin{tabular}[t]{@{}l@{}}\rule{0pt}{12pt}MPS / SIMD-ref\\SIMD-DP / Block\end{tabular} & \hspace*{1pt}\begin{tabular}[t]{@{}c@{}}\rule{0pt}{12pt}1m\,5.1s\\[-.2pt]{\fontsize{5.8}{6.4}\selectfont \textcolor{ForestGreen}{$ \blacktriangle\,2.6\times $}}\\[-.2pt]15s\\[-.2pt]{\fontsize{5.8}{6.4}\selectfont \textcolor{ForestGreen}{$ \blacktriangle\,10.8\times $}}\end{tabular} & \hspace*{1pt}\begin{tabular}[t]{@{}c@{}}\rule{0pt}{12pt}51.5s\\[-.2pt]{\fontsize{5.8}{6.4}\selectfont \textcolor{ForestGreen}{$ \blacktriangle\,3.3\times $}}\\[-.2pt]14.6s\\[-.2pt]{\fontsize{5.8}{6.4}\selectfont \textcolor{ForestGreen}{$ \blacktriangle\,11.1\times $}}\end{tabular} & \textcolor{gray}{\ensuremath{\cdot}} & \textcolor{gray}{\ensuremath{\cdot}} & \hspace*{1pt}\begin{tabular}[t]{@{}c@{}}\rule{0pt}{12pt}7m\,16s\\[-.2pt]{\fontsize{5.8}{6.4}\selectfont \textcolor{BrickRed}{$ \blacktriangledown\,12.8\times $}}\\[-.2pt]56.4s\\[-.2pt]{\fontsize{5.8}{6.4}\selectfont \textcolor{BrickRed}{$ \blacktriangledown\,1.7\times $}}\end{tabular} & \hspace*{1pt}\begin{tabular}[t]{@{}c@{}}\rule{0pt}{12pt}6m\,11.8s\\[-.2pt]{\fontsize{5.8}{6.4}\selectfont \textcolor{BrickRed}{$ \blacktriangledown\,10.8\times $}}\\[-.2pt]53s\\[-.2pt]{\fontsize{5.8}{6.4}\selectfont \textcolor{BrickRed}{$ \blacktriangledown\,1.6\times $}}\end{tabular} & \textcolor{gray}{\ensuremath{\cdot}} & \textcolor{gray}{\ensuremath{\cdot}} & \textcolor{gray}{\ensuremath{\cdot}} & \textcolor{gray}{\ensuremath{\cdot}} & \hspace*{1pt}\begin{tabular}[t]{@{}c@{}}\rule{0pt}{12pt}6m\,45.8s\\[-.2pt]{\fontsize{5.8}{6.4}\selectfont \textcolor{ForestGreen}{$ \blacktriangle\,7\% $}}\\[-.2pt]5m\,46.6s\\[-.2pt]{\fontsize{5.8}{6.4}\selectfont \textcolor{ForestGreen}{$ \blacktriangle\,17\% $}}\end{tabular} & \textcolor{gray}{\ensuremath{\cdot}} & \textcolor{gray}{\ensuremath{\cdot}} & \textcolor{gray}{\ensuremath{\cdot}} & \textcolor{gray}{\ensuremath{\cdot}} & \textcolor{gray}{\ensuremath{\cdot}}\\[3pt]\hline
\BenchmarkSourceLabel{bench:aws-4-k20-w15:row:mps-scalar-scalar-block}{R24}\shortstack[l]{{\color{gray}R24}} & \begin{tabular}[t]{@{}l@{}}\rule{0pt}{12pt}MPS / Scalar\\Scalar-DP / Block\end{tabular} & \hspace*{1pt}\begin{tabular}[t]{@{}c@{}}\rule{0pt}{12pt}1m\,0s\\[-.2pt]{\fontsize{5.8}{6.4}\selectfont \textcolor{ForestGreen}{$ \blacktriangle\,2.8\times $}}\\[-.2pt]14.7s\\[-.2pt]{\fontsize{5.8}{6.4}\selectfont \textcolor{ForestGreen}{$ \blacktriangle\,11.0\times $}}\end{tabular} & \hspace*{1pt}\begin{tabular}[t]{@{}c@{}}\rule{0pt}{12pt}51.7s\\[-.2pt]{\fontsize{5.8}{6.4}\selectfont \textcolor{ForestGreen}{$ \blacktriangle\,3.3\times $}}\\[-.2pt]14.7s\\[-.2pt]{\fontsize{5.8}{6.4}\selectfont \textcolor{ForestGreen}{$ \blacktriangle\,11.1\times $}}\end{tabular} & \textcolor{gray}{\ensuremath{\cdot}} & \textcolor{gray}{\ensuremath{\cdot}} & \hspace*{1pt}\begin{tabular}[t]{@{}c@{}}\rule{0pt}{12pt}7m\,3s\\[-.2pt]{\fontsize{5.8}{6.4}\selectfont \textcolor{BrickRed}{$ \blacktriangledown\,12.5\times $}}\\[-.2pt]38.5s\\[-.2pt]{\fontsize{5.8}{6.4}\selectfont \textcolor{BrickRed}{$ \blacktriangledown\,17\% $}}\end{tabular} & \hspace*{1pt}\begin{tabular}[t]{@{}c@{}}\rule{0pt}{12pt}6m\,18s\\[-.2pt]{\fontsize{5.8}{6.4}\selectfont \textcolor{BrickRed}{$ \blacktriangledown\,11.0\times $}}\\[-.2pt]41.1s\\[-.2pt]{\fontsize{5.8}{6.4}\selectfont \textcolor{BrickRed}{$ \blacktriangledown\,26\% $}}\end{tabular} & \textcolor{gray}{\ensuremath{\cdot}} & \textcolor{gray}{\ensuremath{\cdot}} & \hspace*{1pt}\begin{tabular}[t]{@{}c@{}}\rule{0pt}{12pt}5m\,22.6s\\[-.2pt]{\fontsize{5.8}{6.4}\selectfont \textcolor{BrickRed}{$ \blacktriangledown\,1.5\times $}}\\[-.2pt]4m\,36.5s\\[-.2pt]{\fontsize{5.8}{6.4}\selectfont \textcolor{BrickRed}{$ \blacktriangledown\,1.4\times $}}\end{tabular} & \textcolor{gray}{\ensuremath{\cdot}} & \textcolor{gray}{\ensuremath{\cdot}} & \textcolor{gray}{\ensuremath{\cdot}} & \textcolor{gray}{\ensuremath{\cdot}} & \textcolor{gray}{\ensuremath{\cdot}} & \textcolor{gray}{\ensuremath{\cdot}} & \textcolor{gray}{\ensuremath{\cdot}}\\[3pt]\hline
\BenchmarkSourceLabel{bench:aws-4-k20-w15:row:mps-twostack-scalar-block}{R25}\shortstack[l]{{\color{gray}R25}} & \begin{tabular}[t]{@{}l@{}}\rule{0pt}{12pt}MPS / Two-stack\\Scalar-DP / Block\end{tabular} & \hspace*{1pt}\begin{tabular}[t]{@{}c@{}}\rule{0pt}{12pt}56.5s\\[-.2pt]{\fontsize{5.8}{6.4}\selectfont \textcolor{ForestGreen}{$ \blacktriangle\,3.0\times $}}\\[-.2pt]9.7s\\[-.2pt]{\fontsize{5.8}{6.4}\selectfont \textcolor{ForestGreen}{$ \blacktriangle\,16.7\times $}}\end{tabular} & \hspace*{1pt}\begin{tabular}[t]{@{}c@{}}\rule{0pt}{12pt}51.7s\\[-.2pt]{\fontsize{5.8}{6.4}\selectfont \textcolor{ForestGreen}{$ \blacktriangle\,3.3\times $}}\\[-.2pt]9.9s\\[-.2pt]{\fontsize{5.8}{6.4}\selectfont \textcolor{ForestGreen}{$ \blacktriangle\,16.5\times $}}\end{tabular} & \textcolor{gray}{\ensuremath{\cdot}} & \textcolor{gray}{\ensuremath{\cdot}} & \textcolor{gray}{\ensuremath{\cdot}} & \textcolor{gray}{\ensuremath{\cdot}} & \textcolor{gray}{\ensuremath{\cdot}} & \textcolor{gray}{\ensuremath{\cdot}} & \textcolor{gray}{\ensuremath{\cdot}} & \textcolor{gray}{\ensuremath{\cdot}} & \textcolor{gray}{\ensuremath{\cdot}} & \textcolor{gray}{\ensuremath{\cdot}} & \textcolor{gray}{\ensuremath{\cdot}} & \textcolor{gray}{\ensuremath{\cdot}} & \textcolor{gray}{\ensuremath{\cdot}} & \textcolor{gray}{\ensuremath{\cdot}}\\[3pt]\hline
\BenchmarkSourceLabel{bench:aws-4-k20-w15:row:mps-simdref-simd-ksw2}{R26}\shortstack[l]{{\color{gray}R26}} & \begin{tabular}[t]{@{}l@{}}\rule{0pt}{12pt}MPS / SIMD-ref\\SIMD-DP / KSW2\end{tabular} & \hspace*{1pt}\begin{tabular}[t]{@{}c@{}}\rule{0pt}{12pt}1m\,2.9s\\[-.2pt]{\fontsize{5.8}{6.4}\selectfont \textcolor{ForestGreen}{$ \blacktriangle\,2.7\times $}}\\[-.2pt]14.9s\\[-.2pt]{\fontsize{5.8}{6.4}\selectfont \textcolor{ForestGreen}{$ \blacktriangle\,10.8\times $}}\end{tabular} & \hspace*{1pt}\begin{tabular}[t]{@{}c@{}}\rule{0pt}{12pt}52s\\[-.2pt]{\fontsize{5.8}{6.4}\selectfont \textcolor{ForestGreen}{$ \blacktriangle\,3.3\times $}}\\[-.2pt]14.4s\\[-.2pt]{\fontsize{5.8}{6.4}\selectfont \textcolor{ForestGreen}{$ \blacktriangle\,11.3\times $}}\end{tabular} & \textcolor{gray}{\ensuremath{\cdot}} & \textcolor{gray}{\ensuremath{\cdot}} & \hspace*{1pt}\begin{tabular}[t]{@{}c@{}}\rule{0pt}{12pt}7m\,15.8s\\[-.2pt]{\fontsize{5.8}{6.4}\selectfont \textcolor{BrickRed}{$ \blacktriangledown\,12.8\times $}}\\[-.2pt]53.3s\\[-.2pt]{\fontsize{5.8}{6.4}\selectfont \textcolor{BrickRed}{$ \blacktriangledown\,1.6\times $}}\end{tabular} & \hspace*{1pt}\begin{tabular}[t]{@{}c@{}}\rule{0pt}{12pt}6m\,10.8s\\[-.2pt]{\fontsize{5.8}{6.4}\selectfont \textcolor{BrickRed}{$ \blacktriangledown\,10.8\times $}}\\[-.2pt]51.3s\\[-.2pt]{\fontsize{5.8}{6.4}\selectfont \textcolor{BrickRed}{$ \blacktriangledown\,1.6\times $}}\end{tabular} & \textcolor{gray}{\ensuremath{\cdot}} & \textcolor{gray}{\ensuremath{\cdot}} & \textcolor{gray}{\ensuremath{\cdot}} & \textcolor{gray}{\ensuremath{\cdot}} & \textcolor{gray}{\ensuremath{\cdot}} & \textcolor{gray}{\ensuremath{\cdot}} & \textcolor{gray}{\ensuremath{\cdot}} & \textcolor{gray}{\ensuremath{\cdot}} & \textcolor{gray}{\ensuremath{\cdot}} & \textcolor{gray}{\ensuremath{\cdot}}\\[3pt]\hline
\BenchmarkSourceLabel{bench:aws-4-k20-w15:row:mps-scalar-simd-ksw2}{R27}\shortstack[l]{{\color{gray}R27}} & \begin{tabular}[t]{@{}l@{}}\rule{0pt}{12pt}MPS / Scalar\\SIMD-DP / KSW2\end{tabular} & \hspace*{1pt}\begin{tabular}[t]{@{}c@{}}\rule{0pt}{12pt}58.6s\\[-.2pt]{\fontsize{5.8}{6.4}\selectfont \textcolor{ForestGreen}{$ \blacktriangle\,2.9\times $}}\\[-.2pt]14.4s\\[-.2pt]{\fontsize{5.8}{6.4}\selectfont \textcolor{ForestGreen}{$ \blacktriangle\,11.2\times $}}\end{tabular} & \hspace*{1pt}\begin{tabular}[t]{@{}c@{}}\rule{0pt}{12pt}52.3s\\[-.2pt]{\fontsize{5.8}{6.4}\selectfont \textcolor{ForestGreen}{$ \blacktriangle\,3.3\times $}}\\[-.2pt]15.4s\\[-.2pt]{\fontsize{5.8}{6.4}\selectfont \textcolor{ForestGreen}{$ \blacktriangle\,10.6\times $}}\end{tabular} & \textcolor{gray}{\ensuremath{\cdot}} & \textcolor{gray}{\ensuremath{\cdot}} & \hspace*{1pt}\begin{tabular}[t]{@{}c@{}}\rule{0pt}{12pt}7m\,14.6s\\[-.2pt]{\fontsize{5.8}{6.4}\selectfont \textcolor{BrickRed}{$ \blacktriangledown\,12.8\times $}}\\[-.2pt]39s\\[-.2pt]{\fontsize{5.8}{6.4}\selectfont \textcolor{BrickRed}{$ \blacktriangledown\,19\% $}}\end{tabular} & \hspace*{1pt}\begin{tabular}[t]{@{}c@{}}\rule{0pt}{12pt}6m\,25.1s\\[-.2pt]{\fontsize{5.8}{6.4}\selectfont \textcolor{BrickRed}{$ \blacktriangledown\,11.2\times $}}\\[-.2pt]39s\\[-.2pt]{\fontsize{5.8}{6.4}\selectfont \textcolor{BrickRed}{$ \blacktriangledown\,19\% $}}\end{tabular} & \textcolor{gray}{\ensuremath{\cdot}} & \textcolor{gray}{\ensuremath{\cdot}} & \hspace*{1pt}\begin{tabular}[t]{@{}c@{}}\rule{0pt}{12pt}5m\,10.5s\\[-.2pt]{\fontsize{5.8}{6.4}\selectfont \textcolor{BrickRed}{$ \blacktriangledown\,1.5\times $}}\\[-.2pt]4m\,25.5s\\[-.2pt]{\fontsize{5.8}{6.4}\selectfont \textcolor{BrickRed}{$ \blacktriangledown\,1.3\times $}}\end{tabular} & \hspace*{1pt}\begin{tabular}[t]{@{}c@{}}\rule{0pt}{12pt}4m\,45.4s\\[-.2pt]{\fontsize{5.8}{6.4}\selectfont \textcolor{BrickRed}{$ \blacktriangledown\,22\% $}}\\[-.2pt]4m\,9.3s\\[-.2pt]{\fontsize{5.8}{6.4}\selectfont \textcolor{BrickRed}{$ \blacktriangledown\,12\% $}}\end{tabular} & \textcolor{gray}{\ensuremath{\cdot}} & \textcolor{gray}{\ensuremath{\cdot}} & \textcolor{gray}{\ensuremath{\cdot}} & \textcolor{gray}{\ensuremath{\cdot}} & \textcolor{gray}{\ensuremath{\cdot}} & \textcolor{gray}{\ensuremath{\cdot}}\\[3pt]\hline
\BenchmarkSourceLabel{bench:aws-4-k20-w15:row:mps-scalar-scalar-ksw2}{R28}\shortstack[l]{{\color{gray}R28}} & \begin{tabular}[t]{@{}l@{}}\rule{0pt}{12pt}MPS / Scalar\\Scalar-DP / KSW2\end{tabular} & \hspace*{1pt}\begin{tabular}[t]{@{}c@{}}\rule{0pt}{12pt}1m\,1.8s\\[-.2pt]{\fontsize{5.8}{6.4}\selectfont \textcolor{ForestGreen}{$ \blacktriangle\,2.7\times $}}\\[-.2pt]14.7s\\[-.2pt]{\fontsize{5.8}{6.4}\selectfont \textcolor{ForestGreen}{$ \blacktriangle\,11.0\times $}}\end{tabular} & \hspace*{1pt}\begin{tabular}[t]{@{}c@{}}\rule{0pt}{12pt}52.7s\\[-.2pt]{\fontsize{5.8}{6.4}\selectfont \textcolor{ForestGreen}{$ \blacktriangle\,3.2\times $}}\\[-.2pt]14.6s\\[-.2pt]{\fontsize{5.8}{6.4}\selectfont \textcolor{ForestGreen}{$ \blacktriangle\,11.2\times $}}\end{tabular} & \textcolor{gray}{\ensuremath{\cdot}} & \textcolor{gray}{\ensuremath{\cdot}} & \hspace*{1pt}\begin{tabular}[t]{@{}c@{}}\rule{0pt}{12pt}7m\,1.1s\\[-.2pt]{\fontsize{5.8}{6.4}\selectfont \textcolor{BrickRed}{$ \blacktriangledown\,12.4\times $}}\\[-.2pt]38.9s\\[-.2pt]{\fontsize{5.8}{6.4}\selectfont \textcolor{BrickRed}{$ \blacktriangledown\,18\% $}}\end{tabular} & \hspace*{1pt}\begin{tabular}[t]{@{}c@{}}\rule{0pt}{12pt}6m\,4.9s\\[-.2pt]{\fontsize{5.8}{6.4}\selectfont \textcolor{BrickRed}{$ \blacktriangledown\,10.6\times $}}\\[-.2pt]39.9s\\[-.2pt]{\fontsize{5.8}{6.4}\selectfont \textcolor{BrickRed}{$ \blacktriangledown\,22\% $}}\end{tabular} & \textcolor{gray}{\ensuremath{\cdot}} & \textcolor{gray}{\ensuremath{\cdot}} & \hspace*{1pt}\begin{tabular}[t]{@{}c@{}}\rule{0pt}{12pt}5m\,9.5s\\[-.2pt]{\fontsize{5.8}{6.4}\selectfont \textcolor{BrickRed}{$ \blacktriangledown\,1.5\times $}}\\[-.2pt]4m\,23s\\[-.2pt]{\fontsize{5.8}{6.4}\selectfont \textcolor{BrickRed}{$ \blacktriangledown\,30\% $}}\end{tabular} & \textcolor{gray}{\ensuremath{\cdot}} & \textcolor{gray}{\ensuremath{\cdot}} & \textcolor{gray}{\ensuremath{\cdot}} & \textcolor{gray}{\ensuremath{\cdot}} & \textcolor{gray}{\ensuremath{\cdot}} & \textcolor{gray}{\ensuremath{\cdot}} & \textcolor{gray}{\ensuremath{\cdot}}\\[3pt]\hline
\BenchmarkSourceLabel{bench:aws-4-k20-w15:row:mps-simdref-rmq-ksw2}{R29}\shortstack[l]{{\color{gray}R29}} & \begin{tabular}[t]{@{}l@{}}\rule{0pt}{12pt}MPS / SIMD-ref\\RMQ / KSW2\end{tabular} & \hspace*{1pt}\begin{tabular}[t]{@{}c@{}}\rule{0pt}{12pt}1m\,2.5s\\[-.2pt]{\fontsize{5.8}{6.4}\selectfont \textcolor{ForestGreen}{$ \blacktriangle\,2.7\times $}}\\[-.2pt]15s\\[-.2pt]{\fontsize{5.8}{6.4}\selectfont \textcolor{ForestGreen}{$ \blacktriangle\,10.8\times $}}\end{tabular} & \hspace*{1pt}\begin{tabular}[t]{@{}c@{}}\rule{0pt}{12pt}53.6s\\[-.2pt]{\fontsize{5.8}{6.4}\selectfont \textcolor{ForestGreen}{$ \blacktriangle\,3.2\times $}}\\[-.2pt]15.3s\\[-.2pt]{\fontsize{5.8}{6.4}\selectfont \textcolor{ForestGreen}{$ \blacktriangle\,10.7\times $}}\end{tabular} & \textcolor{gray}{\ensuremath{\cdot}} & \textcolor{gray}{\ensuremath{\cdot}} & \hspace*{1pt}\begin{tabular}[t]{@{}c@{}}\rule{0pt}{12pt}7m\,16s\\[-.2pt]{\fontsize{5.8}{6.4}\selectfont \textcolor{BrickRed}{$ \blacktriangledown\,12.8\times $}}\\[-.2pt]53.4s\\[-.2pt]{\fontsize{5.8}{6.4}\selectfont \textcolor{BrickRed}{$ \blacktriangledown\,1.6\times $}}\end{tabular} & \hspace*{1pt}\begin{tabular}[t]{@{}c@{}}\rule{0pt}{12pt}6m\,11.5s\\[-.2pt]{\fontsize{5.8}{6.4}\selectfont \textcolor{BrickRed}{$ \blacktriangledown\,10.8\times $}}\\[-.2pt]52.1s\\[-.2pt]{\fontsize{5.8}{6.4}\selectfont \textcolor{BrickRed}{$ \blacktriangledown\,1.6\times $}}\end{tabular} & \textcolor{gray}{\ensuremath{\cdot}} & \textcolor{gray}{\ensuremath{\cdot}} & \textcolor{gray}{\ensuremath{\cdot}} & \textcolor{gray}{\ensuremath{\cdot}} & \textcolor{gray}{\ensuremath{\cdot}} & \textcolor{gray}{\ensuremath{\cdot}} & \textcolor{gray}{\ensuremath{\cdot}} & \textcolor{gray}{\ensuremath{\cdot}} & \textcolor{gray}{\ensuremath{\cdot}} & \textcolor{gray}{\ensuremath{\cdot}}\\[3pt]\hline
\end{tabular}\par}
\smallskip{\footnotesize\noindent Each cell: A and B times, $\cdot$: no measurement. ---\textsuperscript{g}: GPU unavailable. n/a: no comparison.\par
\noindent Blue: one lowest A and one lowest B measurement for each column for each configuration.\par
\noindent Gray: minimap2. Yellow: my preferred MPF / Bounded / SIMD-DP / Block configuration.\par
\noindent MPF/MPS: minimizers per fragment/sequence. SIMD-ref: naive SIMD scan. Bounded: bounded SIMD algorithm.\par
}
\clearpage
\noindent{\large\bfseries Table \thetable: (C): AWS: 4 Gbp, k=20, w=15\hfill 5/7}\par\smallskip
{\fontsize{7.1}{8.1}\selectfont\setlength{\tabcolsep}{1pt}\setlength{\fboxsep}{2pt}\renewcommand{\arraystretch}{1.12}
\noindent\begin{tabular}{@{}p{14mm}p{27mm}|p{13.55mm}|p{13.55mm}|p{13.55mm}|p{13.55mm}|p{13.55mm}|p{13.55mm}|p{13.55mm}|p{13.55mm}|p{13.55mm}|p{13.55mm}|p{13.55mm}|p{13.55mm}|p{13.55mm}|p{13.55mm}|p{13.55mm}|p{13.55mm}|@{}}\hline
Engine & Execution & \multicolumn{8}{c|}{\textbf{Index}} & \multicolumn{8}{c|}{\textbf{No index}}\\
 &  & \multicolumn{8}{c|}{\textbf{4 Gbp}} & \multicolumn{8}{c|}{\textbf{4 Gbp}}\\
 &  & \multicolumn{8}{c|}{\textbf{k=20, w=15}} & \multicolumn{8}{c|}{\textbf{k=20, w=15}}\\
 &  & \multicolumn{4}{c|}{\textbf{DNA4}} & \multicolumn{4}{c|}{\textbf{DNA5}} & \multicolumn{4}{c|}{\textbf{DNA4}} & \multicolumn{4}{c|}{\textbf{DNA5}}\\
 &  & \multicolumn{1}{c|}{\textbf{M7A}} & \multicolumn{1}{c|}{\textbf{M7I}} & \multicolumn{1}{c|}{\textbf{M4}} & \multicolumn{1}{c|}{\textbf{C7G}} & \multicolumn{1}{c|}{\textbf{M7A}} & \multicolumn{1}{c|}{\textbf{M7I}} & \multicolumn{1}{c|}{\textbf{M4}} & \multicolumn{1}{c|}{\textbf{C7G}} & \multicolumn{1}{c|}{\textbf{M7A}} & \multicolumn{1}{c|}{\textbf{M7I}} & \multicolumn{1}{c|}{\textbf{M4}} & \multicolumn{1}{c|}{\textbf{C7G}} & \multicolumn{1}{c|}{\textbf{M7A}} & \multicolumn{1}{c|}{\textbf{M7I}} & \multicolumn{1}{c|}{\textbf{M4}} & \multicolumn{1}{c|}{\textbf{C7G}}\\
 & & {\color{gray}C1} & {\color{gray}C2} & {\color{gray}C3} & {\color{gray}C4} & {\color{gray}C5} & {\color{gray}C6} & {\color{gray}C7} & {\color{gray}C8} & {\color{gray}C9} & {\color{gray}C10} & {\color{gray}C11} & {\color{gray}C12} & {\color{gray}C13} & {\color{gray}C14} & {\color{gray}C15} & {\color{gray}C16}\\
\hline
\rowcolor{BenchBaseline}[0.5pt][0.5pt]\shortstack[l]{{\color{gray}R1}\\\textbf{minimap2}} & \textbf{Baseline} & \hspace*{1pt}\begin{tabular}[t]{@{}c@{}}\rule{0pt}{12pt}2m\,48.1s\\[-.2pt]{\fontsize{5.8}{6.4}\selectfont \textcolor{BrickRed}{$ \blacktriangledown\,20.4\times $}}\\[-.2pt]2m\,41.7s\\[-.2pt]{\fontsize{5.8}{6.4}\selectfont \textcolor{BrickRed}{$ \blacktriangledown\,22.6\times $}}\end{tabular} & \hspace*{1pt}\begin{tabular}[t]{@{}c@{}}\rule{0pt}{12pt}2m\,50.3s\\[-.2pt]{\fontsize{5.8}{6.4}\selectfont \textcolor{BrickRed}{$ \blacktriangledown\,19.8\times $}}\\[-.2pt]2m\,43s\\[-.2pt]{\fontsize{5.8}{6.4}\selectfont \textcolor{BrickRed}{$ \blacktriangledown\,22.2\times $}}\end{tabular} & \hspace*{1pt}\begin{tabular}[t]{@{}c@{}}\rule{0pt}{12pt}4m\,54.3s\\[-.2pt]{\fontsize{5.8}{6.4}\selectfont \textcolor{BrickRed}{$ \blacktriangledown\,43.2\times $}}\\[-.2pt]4m\,39s\\[-.2pt]{\fontsize{5.8}{6.4}\selectfont \textcolor{BrickRed}{$ \blacktriangledown\,54.6\times $}}\end{tabular} & \hspace*{1pt}\begin{tabular}[t]{@{}c@{}}\rule{0pt}{12pt}25m\,5.5s\\[-.2pt]{\fontsize{5.8}{6.4}\selectfont \textcolor{BrickRed}{$ \blacktriangledown\,41.8\times $}}\\[-.2pt]25m\,0.7s\\[-.2pt]{\fontsize{5.8}{6.4}\selectfont \textcolor{BrickRed}{$ \blacktriangledown\,46.6\times $}}\end{tabular} & \hspace*{1pt}\begin{tabular}[t]{@{}c@{}}\rule{0pt}{12pt}33.9s\\[-.2pt]{\fontsize{5.8}{6.4}\selectfont \textcolor{BrickRed}{$ \blacktriangledown\,2.3\times $}}\\[-.2pt]32.9s\\[-.2pt]{\fontsize{5.8}{6.4}\selectfont \textcolor{BrickRed}{$ \blacktriangledown\,3.1\times $}}\end{tabular} & \hspace*{1pt}\begin{tabular}[t]{@{}c@{}}\rule{0pt}{12pt}34.4s\\[-.2pt]{\fontsize{5.8}{6.4}\selectfont \textcolor{BrickRed}{$ \blacktriangledown\,2.6\times $}}\\[-.2pt]32.7s\\[-.2pt]{\fontsize{5.8}{6.4}\selectfont \textcolor{BrickRed}{$ \blacktriangledown\,2.9\times $}}\end{tabular} & \hspace*{1pt}\begin{tabular}[t]{@{}c@{}}\rule{0pt}{12pt}1m\,8.2s\\[-.2pt]{\fontsize{5.8}{6.4}\selectfont \textcolor{BrickRed}{$ \blacktriangledown\,3.7\times $}}\\[-.2pt]1m\,5.1s\\[-.2pt]{\fontsize{5.8}{6.4}\selectfont \textcolor{BrickRed}{$ \blacktriangledown\,4.6\times $}}\end{tabular} & \hspace*{1pt}\begin{tabular}[t]{@{}c@{}}\rule{0pt}{12pt}3m\,9.8s\\[-.2pt]{\fontsize{5.8}{6.4}\selectfont \textcolor{BrickRed}{$ \blacktriangledown\,3.9\times $}}\\[-.2pt]3m\,3.7s\\[-.2pt]{\fontsize{5.8}{6.4}\selectfont \textcolor{BrickRed}{$ \blacktriangledown\,4.1\times $}}\end{tabular} & \hspace*{1pt}\begin{tabular}[t]{@{}c@{}}\rule{0pt}{12pt}3m\,28.3s\\[-.2pt]{\fontsize{5.8}{6.4}\selectfont \textcolor{BrickRed}{$ \blacktriangledown\,1.3\times $}}\\[-.2pt]3m\,22.8s\\[-.2pt]{\fontsize{5.8}{6.4}\selectfont \textcolor{BrickRed}{$ \blacktriangledown\,1.4\times $}}\end{tabular} & \hspace*{1pt}\begin{tabular}[t]{@{}c@{}}\rule{0pt}{12pt}3m\,53.2s\\[-.2pt]{\fontsize{5.8}{6.4}\selectfont \textcolor{BrickRed}{$ \blacktriangledown\,1.4\times $}}\\[-.2pt]3m\,43.5s\\[-.2pt]{\fontsize{5.8}{6.4}\selectfont \textcolor{BrickRed}{$ \blacktriangledown\,1.5\times $}}\end{tabular} & \hspace*{1pt}\begin{tabular}[t]{@{}c@{}}\rule{0pt}{12pt}7m\,14.6s\\[-.2pt]{\fontsize{5.8}{6.4}\selectfont \textcolor{BrickRed}{$ \blacktriangledown\,1.5\times $}}\\[-.2pt]6m\,57.2s\\[-.2pt]{\fontsize{5.8}{6.4}\selectfont \textcolor{BrickRed}{$ \blacktriangledown\,1.6\times $}}\end{tabular} & \hspace*{1pt}\begin{tabular}[t]{@{}c@{}}\rule{0pt}{12pt}51m\,22s\\[-.2pt]{\fontsize{5.8}{6.4}\selectfont \textcolor{gray}{n/a}}\\[-.2pt]51m\,17.2s\\[-.2pt]{\fontsize{5.8}{6.4}\selectfont \textcolor{gray}{n/a}}\end{tabular} & \hspace*{1pt}\begin{tabular}[t]{@{}c@{}}\rule{0pt}{12pt}1m\,10.3s\\[-.2pt]{\fontsize{5.8}{6.4}\selectfont \textcolor{ForestGreen}{$ \blacktriangle\,7.6\times $}}\\[-.2pt]1m\,9s\\[-.2pt]{\fontsize{5.8}{6.4}\selectfont \textcolor{ForestGreen}{$ \blacktriangle\,2.4\times $}}\end{tabular} & \hspace*{1pt}\begin{tabular}[t]{@{}c@{}}\rule{0pt}{12pt}1m\,8.6s\\[-.2pt]{\fontsize{5.8}{6.4}\selectfont \textcolor{ForestGreen}{$ \blacktriangle\,6.4\times $}}\\[-.2pt]1m\,6.9s\\[-.2pt]{\fontsize{5.8}{6.4}\selectfont \textcolor{ForestGreen}{$ \blacktriangle\,2.3\times $}}\end{tabular} & \hspace*{1pt}\begin{tabular}[t]{@{}c@{}}\rule{0pt}{12pt}2m\,21.5s\\[-.2pt]{\fontsize{5.8}{6.4}\selectfont \textcolor{ForestGreen}{$ \blacktriangle\,6.0\times $}}\\[-.2pt]2m\,18.4s\\[-.2pt]{\fontsize{5.8}{6.4}\selectfont \textcolor{ForestGreen}{$ \blacktriangle\,2.2\times $}}\end{tabular} & \hspace*{1pt}\begin{tabular}[t]{@{}c@{}}\rule{0pt}{12pt}15m\,56.8s\\[-.2pt]{\fontsize{5.8}{6.4}\selectfont \textcolor{gray}{n/a}}\\[-.2pt]15m\,51.9s\\[-.2pt]{\fontsize{5.8}{6.4}\selectfont \textcolor{gray}{n/a}}\end{tabular}\\[3pt]\hline
\BenchmarkSourceLabel{bench:aws-4-k20-w15:row:mps-simdref-scalar-ksw2}{R30}\shortstack[l]{{\color{gray}R30}\\\textbf{BioseqDB}} & \begin{tabular}[t]{@{}l@{}}\rule{0pt}{12pt}MPS / SIMD-ref\\Scalar-DP / KSW2\end{tabular} & \hspace*{1pt}\begin{tabular}[t]{@{}c@{}}\rule{0pt}{12pt}1m\,2.8s\\[-.2pt]{\fontsize{5.8}{6.4}\selectfont \textcolor{ForestGreen}{$ \blacktriangle\,2.7\times $}}\\[-.2pt]15s\\[-.2pt]{\fontsize{5.8}{6.4}\selectfont \textcolor{ForestGreen}{$ \blacktriangle\,10.8\times $}}\end{tabular} & \hspace*{1pt}\begin{tabular}[t]{@{}c@{}}\rule{0pt}{12pt}53.9s\\[-.2pt]{\fontsize{5.8}{6.4}\selectfont \textcolor{ForestGreen}{$ \blacktriangle\,3.2\times $}}\\[-.2pt]15.2s\\[-.2pt]{\fontsize{5.8}{6.4}\selectfont \textcolor{ForestGreen}{$ \blacktriangle\,10.7\times $}}\end{tabular} & \textcolor{gray}{\ensuremath{\cdot}} & \textcolor{gray}{\ensuremath{\cdot}} & \hspace*{1pt}\begin{tabular}[t]{@{}c@{}}\rule{0pt}{12pt}7m\,17.8s\\[-.2pt]{\fontsize{5.8}{6.4}\selectfont \textcolor{BrickRed}{$ \blacktriangledown\,12.9\times $}}\\[-.2pt]53.3s\\[-.2pt]{\fontsize{5.8}{6.4}\selectfont \textcolor{BrickRed}{$ \blacktriangledown\,1.6\times $}}\end{tabular} & \hspace*{1pt}\begin{tabular}[t]{@{}c@{}}\rule{0pt}{12pt}6m\,9.3s\\[-.2pt]{\fontsize{5.8}{6.4}\selectfont \textcolor{BrickRed}{$ \blacktriangledown\,10.7\times $}}\\[-.2pt]50.2s\\[-.2pt]{\fontsize{5.8}{6.4}\selectfont \textcolor{BrickRed}{$ \blacktriangledown\,1.5\times $}}\end{tabular} & \textcolor{gray}{\ensuremath{\cdot}} & \textcolor{gray}{\ensuremath{\cdot}} & \textcolor{gray}{\ensuremath{\cdot}} & \textcolor{gray}{\ensuremath{\cdot}} & \textcolor{gray}{\ensuremath{\cdot}} & \textcolor{gray}{\ensuremath{\cdot}} & \textcolor{gray}{\ensuremath{\cdot}} & \textcolor{gray}{\ensuremath{\cdot}} & \textcolor{gray}{\ensuremath{\cdot}} & \textcolor{gray}{\ensuremath{\cdot}}\\[3pt]\hline
\BenchmarkSourceLabel{bench:aws-4-k20-w15:row:mps-scalar-rmq-ksw2}{R31}\shortstack[l]{{\color{gray}R31}} & \begin{tabular}[t]{@{}l@{}}\rule{0pt}{12pt}MPS / Scalar\\RMQ / KSW2\end{tabular} & \hspace*{1pt}\begin{tabular}[t]{@{}c@{}}\rule{0pt}{12pt}59.3s\\[-.2pt]{\fontsize{5.8}{6.4}\selectfont \textcolor{ForestGreen}{$ \blacktriangle\,2.8\times $}}\\[-.2pt]14.7s\\[-.2pt]{\fontsize{5.8}{6.4}\selectfont \textcolor{ForestGreen}{$ \blacktriangle\,11.0\times $}}\end{tabular} & \hspace*{1pt}\begin{tabular}[t]{@{}c@{}}\rule{0pt}{12pt}56.8s\\[-.2pt]{\fontsize{5.8}{6.4}\selectfont \textcolor{ForestGreen}{$ \blacktriangle\,3.0\times $}}\\[-.2pt]14.7s\\[-.2pt]{\fontsize{5.8}{6.4}\selectfont \textcolor{ForestGreen}{$ \blacktriangle\,11.1\times $}}\end{tabular} & \textcolor{gray}{\ensuremath{\cdot}} & \textcolor{gray}{\ensuremath{\cdot}} & \hspace*{1pt}\begin{tabular}[t]{@{}c@{}}\rule{0pt}{12pt}7m\,5s\\[-.2pt]{\fontsize{5.8}{6.4}\selectfont \textcolor{BrickRed}{$ \blacktriangledown\,12.5\times $}}\\[-.2pt]39.2s\\[-.2pt]{\fontsize{5.8}{6.4}\selectfont \textcolor{BrickRed}{$ \blacktriangledown\,19\% $}}\end{tabular} & \hspace*{1pt}\begin{tabular}[t]{@{}c@{}}\rule{0pt}{12pt}6m\,26.9s\\[-.2pt]{\fontsize{5.8}{6.4}\selectfont \textcolor{BrickRed}{$ \blacktriangledown\,11.2\times $}}\\[-.2pt]41.2s\\[-.2pt]{\fontsize{5.8}{6.4}\selectfont \textcolor{BrickRed}{$ \blacktriangledown\,26\% $}}\end{tabular} & \textcolor{gray}{\ensuremath{\cdot}} & \textcolor{gray}{\ensuremath{\cdot}} & \hspace*{1pt}\begin{tabular}[t]{@{}c@{}}\rule{0pt}{12pt}5m\,14.2s\\[-.2pt]{\fontsize{5.8}{6.4}\selectfont \textcolor{BrickRed}{$ \blacktriangledown\,1.5\times $}}\\[-.2pt]4m\,26.6s\\[-.2pt]{\fontsize{5.8}{6.4}\selectfont \textcolor{BrickRed}{$ \blacktriangledown\,1.3\times $}}\end{tabular} & \textcolor{gray}{\ensuremath{\cdot}} & \textcolor{gray}{\ensuremath{\cdot}} & \textcolor{gray}{\ensuremath{\cdot}} & \textcolor{gray}{\ensuremath{\cdot}} & \textcolor{gray}{\ensuremath{\cdot}} & \textcolor{gray}{\ensuremath{\cdot}} & \textcolor{gray}{\ensuremath{\cdot}}\\[3pt]\hline
\BenchmarkSourceLabel{bench:aws-4-k20-w15:row:mpf-twostack-rmq-ksw2}{R32}\shortstack[l]{{\color{gray}R32}} & \begin{tabular}[t]{@{}l@{}}\rule{0pt}{12pt}MPF / Two-stack\\RMQ / KSW2\end{tabular} & \hspace*{1pt}\begin{tabular}[t]{@{}c@{}}\rule{0pt}{12pt}9m\,9.8s\\[-.2pt]{\fontsize{5.8}{6.4}\selectfont \textcolor{BrickRed}{$ \blacktriangledown\,3.3\times $}}\\[-.2pt]9m\,7.7s\\[-.2pt]{\fontsize{5.8}{6.4}\selectfont \textcolor{BrickRed}{$ \blacktriangledown\,3.4\times $}}\end{tabular} & \hspace*{1pt}\begin{tabular}[t]{@{}c@{}}\rule{0pt}{12pt}15m\,43.9s\\[-.2pt]{\fontsize{5.8}{6.4}\selectfont \textcolor{BrickRed}{$ \blacktriangledown\,5.5\times $}}\\[-.2pt]15m\,40.9s\\[-.2pt]{\fontsize{5.8}{6.4}\selectfont \textcolor{BrickRed}{$ \blacktriangledown\,5.8\times $}}\end{tabular} & \hspace*{1pt}\begin{tabular}[t]{@{}c@{}}\rule{0pt}{12pt}33m\,26.8s\\[-.2pt]{\fontsize{5.8}{6.4}\selectfont \textcolor{BrickRed}{$ \blacktriangledown\,6.8\times $}}\\[-.2pt]33m\,20.7s\\[-.2pt]{\fontsize{5.8}{6.4}\selectfont \textcolor{BrickRed}{$ \blacktriangledown\,7.2\times $}}\end{tabular} & \textcolor{gray}{\ensuremath{\cdot}} & \textcolor{gray}{\ensuremath{\cdot}} & \textcolor{gray}{\ensuremath{\cdot}} & \textcolor{gray}{\ensuremath{\cdot}} & \textcolor{gray}{\ensuremath{\cdot}} & \hspace*{1pt}\begin{tabular}[t]{@{}c@{}}\rule{0pt}{12pt}3m\,22.2s\\[-.2pt]{\fontsize{5.8}{6.4}\selectfont \textcolor{gray}{$ \blacktriangleright\,3\% $}}\\[-.2pt]2m\,35.6s\\[-.2pt]{\fontsize{5.8}{6.4}\selectfont \textcolor{ForestGreen}{$ \blacktriangle\,23\% $}}\end{tabular} & \hspace*{1pt}\begin{tabular}[t]{@{}c@{}}\rule{0pt}{12pt}3m\,3.8s\\[-.2pt]{\fontsize{5.8}{6.4}\selectfont \textcolor{ForestGreen}{$ \blacktriangle\,21\% $}}\\[-.2pt]\colorbox{BenchBest}{\strut 2m\,25s}\\[-.2pt]{\fontsize{5.8}{6.4}\selectfont \textcolor{ForestGreen}{$ \blacktriangle\,1.5\times $}}\end{tabular} & \textcolor{gray}{\ensuremath{\cdot}} & \textcolor{gray}{\ensuremath{\cdot}} & \textcolor{gray}{\ensuremath{\cdot}} & \textcolor{gray}{\ensuremath{\cdot}} & \textcolor{gray}{\ensuremath{\cdot}} & \textcolor{gray}{\ensuremath{\cdot}}\\[3pt]\hline
\BenchmarkSourceLabel{bench:aws-4-k20-w15:row:mpf-twostack-rmq-block}{R33}\shortstack[l]{{\color{gray}R33}} & \begin{tabular}[t]{@{}l@{}}\rule{0pt}{12pt}MPF / Two-stack\\RMQ / Block\end{tabular} & \hspace*{1pt}\begin{tabular}[t]{@{}c@{}}\rule{0pt}{12pt}9m\,11s\\[-.2pt]{\fontsize{5.8}{6.4}\selectfont \textcolor{BrickRed}{$ \blacktriangledown\,3.3\times $}}\\[-.2pt]9m\,10.5s\\[-.2pt]{\fontsize{5.8}{6.4}\selectfont \textcolor{BrickRed}{$ \blacktriangledown\,3.4\times $}}\end{tabular} & \hspace*{1pt}\begin{tabular}[t]{@{}c@{}}\rule{0pt}{12pt}13m\,20.9s\\[-.2pt]{\fontsize{5.8}{6.4}\selectfont \textcolor{BrickRed}{$ \blacktriangledown\,4.7\times $}}\\[-.2pt]13m\,20.1s\\[-.2pt]{\fontsize{5.8}{6.4}\selectfont \textcolor{BrickRed}{$ \blacktriangledown\,4.9\times $}}\end{tabular} & \hspace*{1pt}\begin{tabular}[t]{@{}c@{}}\rule{0pt}{12pt}33m\,23s\\[-.2pt]{\fontsize{5.8}{6.4}\selectfont \textcolor{BrickRed}{$ \blacktriangledown\,6.8\times $}}\\[-.2pt]33m\,21.5s\\[-.2pt]{\fontsize{5.8}{6.4}\selectfont \textcolor{BrickRed}{$ \blacktriangledown\,7.2\times $}}\end{tabular} & \textcolor{gray}{\ensuremath{\cdot}} & \textcolor{gray}{\ensuremath{\cdot}} & \textcolor{gray}{\ensuremath{\cdot}} & \textcolor{gray}{\ensuremath{\cdot}} & \textcolor{gray}{\ensuremath{\cdot}} & \hspace*{1pt}\begin{tabular}[t]{@{}c@{}}\rule{0pt}{12pt}3m\,21.7s\\[-.2pt]{\fontsize{5.8}{6.4}\selectfont \textcolor{gray}{$ \blacktriangleright\,3\% $}}\\[-.2pt]2m\,38.8s\\[-.2pt]{\fontsize{5.8}{6.4}\selectfont \textcolor{ForestGreen}{$ \blacktriangle\,22\% $}}\end{tabular} & \hspace*{1pt}\begin{tabular}[t]{@{}c@{}}\rule{0pt}{12pt}3m\,0.8s\\[-.2pt]{\fontsize{5.8}{6.4}\selectfont \textcolor{ForestGreen}{$ \blacktriangle\,22\% $}}\\[-.2pt]2m\,25s\\[-.2pt]{\fontsize{5.8}{6.4}\selectfont \textcolor{ForestGreen}{$ \blacktriangle\,1.5\times $}}\end{tabular} & \textcolor{gray}{\ensuremath{\cdot}} & \textcolor{gray}{\ensuremath{\cdot}} & \textcolor{gray}{\ensuremath{\cdot}} & \textcolor{gray}{\ensuremath{\cdot}} & \textcolor{gray}{\ensuremath{\cdot}} & \textcolor{gray}{\ensuremath{\cdot}}\\[3pt]\hline
\BenchmarkSourceLabel{bench:aws-4-k20-w15:row:mpf-twostack-simd-block}{R34}\shortstack[l]{{\color{gray}R34}} & \begin{tabular}[t]{@{}l@{}}\rule{0pt}{12pt}MPF / Two-stack\\SIMD-DP / Block\end{tabular} & \hspace*{1pt}\begin{tabular}[t]{@{}c@{}}\rule{0pt}{12pt}9m\,11.2s\\[-.2pt]{\fontsize{5.8}{6.4}\selectfont \textcolor{BrickRed}{$ \blacktriangledown\,3.3\times $}}\\[-.2pt]9m\,10.7s\\[-.2pt]{\fontsize{5.8}{6.4}\selectfont \textcolor{BrickRed}{$ \blacktriangledown\,3.4\times $}}\end{tabular} & \hspace*{1pt}\begin{tabular}[t]{@{}c@{}}\rule{0pt}{12pt}15m\,34.3s\\[-.2pt]{\fontsize{5.8}{6.4}\selectfont \textcolor{BrickRed}{$ \blacktriangledown\,5.5\times $}}\\[-.2pt]15m\,33.4s\\[-.2pt]{\fontsize{5.8}{6.4}\selectfont \textcolor{BrickRed}{$ \blacktriangledown\,5.7\times $}}\end{tabular} & \hspace*{1pt}\begin{tabular}[t]{@{}c@{}}\rule{0pt}{12pt}33m\,48.6s\\[-.2pt]{\fontsize{5.8}{6.4}\selectfont \textcolor{BrickRed}{$ \blacktriangledown\,6.9\times $}}\\[-.2pt]33m\,47.1s\\[-.2pt]{\fontsize{5.8}{6.4}\selectfont \textcolor{BrickRed}{$ \blacktriangledown\,7.3\times $}}\end{tabular} & \textcolor{gray}{\ensuremath{\cdot}} & \textcolor{gray}{\ensuremath{\cdot}} & \textcolor{gray}{\ensuremath{\cdot}} & \textcolor{gray}{\ensuremath{\cdot}} & \textcolor{gray}{\ensuremath{\cdot}} & \hspace*{1pt}\begin{tabular}[t]{@{}c@{}}\rule{0pt}{12pt}3m\,19s\\[-.2pt]{\fontsize{5.8}{6.4}\selectfont \textcolor{gray}{$ \blacktriangleright\,4\% $}}\\[-.2pt]2m\,36.1s\\[-.2pt]{\fontsize{5.8}{6.4}\selectfont \textcolor{ForestGreen}{$ \blacktriangle\,23\% $}}\end{tabular} & \hspace*{1pt}\begin{tabular}[t]{@{}c@{}}\rule{0pt}{12pt}3m\,6.1s\\[-.2pt]{\fontsize{5.8}{6.4}\selectfont \textcolor{ForestGreen}{$ \blacktriangle\,20\% $}}\\[-.2pt]2m\,30.7s\\[-.2pt]{\fontsize{5.8}{6.4}\selectfont \textcolor{ForestGreen}{$ \blacktriangle\,1.5\times $}}\end{tabular} & \textcolor{gray}{\ensuremath{\cdot}} & \textcolor{gray}{\ensuremath{\cdot}} & \textcolor{gray}{\ensuremath{\cdot}} & \textcolor{gray}{\ensuremath{\cdot}} & \textcolor{gray}{\ensuremath{\cdot}} & \textcolor{gray}{\ensuremath{\cdot}}\\[3pt]\hline
\BenchmarkSourceLabel{bench:aws-4-k20-w15:row:mpf-twostack-simd-ksw2}{R35}\shortstack[l]{{\color{gray}R35}} & \begin{tabular}[t]{@{}l@{}}\rule{0pt}{12pt}MPF / Two-stack\\SIMD-DP / KSW2\end{tabular} & \hspace*{1pt}\begin{tabular}[t]{@{}c@{}}\rule{0pt}{12pt}9m\,11.4s\\[-.2pt]{\fontsize{5.8}{6.4}\selectfont \textcolor{BrickRed}{$ \blacktriangledown\,3.3\times $}}\\[-.2pt]9m\,9.5s\\[-.2pt]{\fontsize{5.8}{6.4}\selectfont \textcolor{BrickRed}{$ \blacktriangledown\,3.4\times $}}\end{tabular} & \hspace*{1pt}\begin{tabular}[t]{@{}c@{}}\rule{0pt}{12pt}13m\,22s\\[-.2pt]{\fontsize{5.8}{6.4}\selectfont \textcolor{BrickRed}{$ \blacktriangledown\,4.7\times $}}\\[-.2pt]13m\,19.3s\\[-.2pt]{\fontsize{5.8}{6.4}\selectfont \textcolor{BrickRed}{$ \blacktriangledown\,4.9\times $}}\end{tabular} & \hspace*{1pt}\begin{tabular}[t]{@{}c@{}}\rule{0pt}{12pt}33m\,35.9s\\[-.2pt]{\fontsize{5.8}{6.4}\selectfont \textcolor{BrickRed}{$ \blacktriangledown\,6.9\times $}}\\[-.2pt]33m\,29.9s\\[-.2pt]{\fontsize{5.8}{6.4}\selectfont \textcolor{BrickRed}{$ \blacktriangledown\,7.2\times $}}\end{tabular} & \textcolor{gray}{\ensuremath{\cdot}} & \textcolor{gray}{\ensuremath{\cdot}} & \textcolor{gray}{\ensuremath{\cdot}} & \textcolor{gray}{\ensuremath{\cdot}} & \textcolor{gray}{\ensuremath{\cdot}} & \hspace*{1pt}\begin{tabular}[t]{@{}c@{}}\rule{0pt}{12pt}3m\,21.5s\\[-.2pt]{\fontsize{5.8}{6.4}\selectfont \textcolor{gray}{$ \blacktriangleright\,3\% $}}\\[-.2pt]2m\,35s\\[-.2pt]{\fontsize{5.8}{6.4}\selectfont \textcolor{ForestGreen}{$ \blacktriangle\,24\% $}}\end{tabular} & \hspace*{1pt}\begin{tabular}[t]{@{}c@{}}\rule{0pt}{12pt}3m\,4.9s\\[-.2pt]{\fontsize{5.8}{6.4}\selectfont \textcolor{ForestGreen}{$ \blacktriangle\,21\% $}}\\[-.2pt]2m\,25.8s\\[-.2pt]{\fontsize{5.8}{6.4}\selectfont \textcolor{ForestGreen}{$ \blacktriangle\,1.5\times $}}\end{tabular} & \textcolor{gray}{\ensuremath{\cdot}} & \textcolor{gray}{\ensuremath{\cdot}} & \textcolor{gray}{\ensuremath{\cdot}} & \textcolor{gray}{\ensuremath{\cdot}} & \textcolor{gray}{\ensuremath{\cdot}} & \textcolor{gray}{\ensuremath{\cdot}}\\[3pt]\hline
\BenchmarkSourceLabel{bench:aws-4-k20-w15:row:mpf-twostack-scalar-block}{R36}\shortstack[l]{{\color{gray}R36}} & \begin{tabular}[t]{@{}l@{}}\rule{0pt}{12pt}MPF / Two-stack\\Scalar-DP / Block\end{tabular} & \hspace*{1pt}\begin{tabular}[t]{@{}c@{}}\rule{0pt}{12pt}9m\,15.8s\\[-.2pt]{\fontsize{5.8}{6.4}\selectfont \textcolor{BrickRed}{$ \blacktriangledown\,3.3\times $}}\\[-.2pt]9m\,15s\\[-.2pt]{\fontsize{5.8}{6.4}\selectfont \textcolor{BrickRed}{$ \blacktriangledown\,3.4\times $}}\end{tabular} & \hspace*{1pt}\begin{tabular}[t]{@{}c@{}}\rule{0pt}{12pt}15m\,26.5s\\[-.2pt]{\fontsize{5.8}{6.4}\selectfont \textcolor{BrickRed}{$ \blacktriangledown\,5.4\times $}}\\[-.2pt]15m\,25.5s\\[-.2pt]{\fontsize{5.8}{6.4}\selectfont \textcolor{BrickRed}{$ \blacktriangledown\,5.7\times $}}\end{tabular} & \hspace*{1pt}\begin{tabular}[t]{@{}c@{}}\rule{0pt}{12pt}33m\,22.2s\\[-.2pt]{\fontsize{5.8}{6.4}\selectfont \textcolor{BrickRed}{$ \blacktriangledown\,6.8\times $}}\\[-.2pt]33m\,20.6s\\[-.2pt]{\fontsize{5.8}{6.4}\selectfont \textcolor{BrickRed}{$ \blacktriangledown\,7.2\times $}}\end{tabular} & \textcolor{gray}{\ensuremath{\cdot}} & \textcolor{gray}{\ensuremath{\cdot}} & \textcolor{gray}{\ensuremath{\cdot}} & \textcolor{gray}{\ensuremath{\cdot}} & \textcolor{gray}{\ensuremath{\cdot}} & \hspace*{1pt}\begin{tabular}[t]{@{}c@{}}\rule{0pt}{12pt}3m\,18.4s\\[-.2pt]{\fontsize{5.8}{6.4}\selectfont \textcolor{gray}{$ \blacktriangleright\,5\% $}}\\[-.2pt]2m\,36s\\[-.2pt]{\fontsize{5.8}{6.4}\selectfont \textcolor{ForestGreen}{$ \blacktriangle\,23\% $}}\end{tabular} & \hspace*{1pt}\begin{tabular}[t]{@{}c@{}}\rule{0pt}{12pt}3m\,4.1s\\[-.2pt]{\fontsize{5.8}{6.4}\selectfont \textcolor{ForestGreen}{$ \blacktriangle\,21\% $}}\\[-.2pt]2m\,28.5s\\[-.2pt]{\fontsize{5.8}{6.4}\selectfont \textcolor{ForestGreen}{$ \blacktriangle\,1.5\times $}}\end{tabular} & \textcolor{gray}{\ensuremath{\cdot}} & \textcolor{gray}{\ensuremath{\cdot}} & \textcolor{gray}{\ensuremath{\cdot}} & \textcolor{gray}{\ensuremath{\cdot}} & \textcolor{gray}{\ensuremath{\cdot}} & \textcolor{gray}{\ensuremath{\cdot}}\\[3pt]\hline
\end{tabular}\par}
\smallskip{\footnotesize\noindent Each cell: A and B times, $\cdot$: no measurement. ---\textsuperscript{g}: GPU unavailable. n/a: no comparison.\par
\noindent Blue: one lowest A and one lowest B measurement for each column for each configuration.\par
\noindent Gray: minimap2. Yellow: my preferred MPF / Bounded / SIMD-DP / Block configuration.\par
\noindent MPF/MPS: minimizers per fragment/sequence. SIMD-ref: naive SIMD scan. Bounded: bounded SIMD algorithm.\par
}
\clearpage
\noindent{\large\bfseries Table \thetable: (C): AWS: 4 Gbp, k=20, w=15\hfill 6/7}\par\smallskip
{\fontsize{7.1}{8.1}\selectfont\setlength{\tabcolsep}{1pt}\setlength{\fboxsep}{2pt}\renewcommand{\arraystretch}{1.12}
\noindent\begin{tabular}{@{}p{14mm}p{27mm}|p{13.55mm}|p{13.55mm}|p{13.55mm}|p{13.55mm}|p{13.55mm}|p{13.55mm}|p{13.55mm}|p{13.55mm}|p{13.55mm}|p{13.55mm}|p{13.55mm}|p{13.55mm}|p{13.55mm}|p{13.55mm}|p{13.55mm}|p{13.55mm}|@{}}\hline
Engine & Execution & \multicolumn{8}{c|}{\textbf{Index}} & \multicolumn{8}{c|}{\textbf{No index}}\\
 &  & \multicolumn{8}{c|}{\textbf{4 Gbp}} & \multicolumn{8}{c|}{\textbf{4 Gbp}}\\
 &  & \multicolumn{8}{c|}{\textbf{k=20, w=15}} & \multicolumn{8}{c|}{\textbf{k=20, w=15}}\\
 &  & \multicolumn{4}{c|}{\textbf{DNA4}} & \multicolumn{4}{c|}{\textbf{DNA5}} & \multicolumn{4}{c|}{\textbf{DNA4}} & \multicolumn{4}{c|}{\textbf{DNA5}}\\
 &  & \multicolumn{1}{c|}{\textbf{M7A}} & \multicolumn{1}{c|}{\textbf{M7I}} & \multicolumn{1}{c|}{\textbf{M4}} & \multicolumn{1}{c|}{\textbf{C7G}} & \multicolumn{1}{c|}{\textbf{M7A}} & \multicolumn{1}{c|}{\textbf{M7I}} & \multicolumn{1}{c|}{\textbf{M4}} & \multicolumn{1}{c|}{\textbf{C7G}} & \multicolumn{1}{c|}{\textbf{M7A}} & \multicolumn{1}{c|}{\textbf{M7I}} & \multicolumn{1}{c|}{\textbf{M4}} & \multicolumn{1}{c|}{\textbf{C7G}} & \multicolumn{1}{c|}{\textbf{M7A}} & \multicolumn{1}{c|}{\textbf{M7I}} & \multicolumn{1}{c|}{\textbf{M4}} & \multicolumn{1}{c|}{\textbf{C7G}}\\
 & & {\color{gray}C1} & {\color{gray}C2} & {\color{gray}C3} & {\color{gray}C4} & {\color{gray}C5} & {\color{gray}C6} & {\color{gray}C7} & {\color{gray}C8} & {\color{gray}C9} & {\color{gray}C10} & {\color{gray}C11} & {\color{gray}C12} & {\color{gray}C13} & {\color{gray}C14} & {\color{gray}C15} & {\color{gray}C16}\\
\hline
\rowcolor{BenchBaseline}[0.5pt][0.5pt]\shortstack[l]{{\color{gray}R1}\\\textbf{minimap2}} & \textbf{Baseline} & \hspace*{1pt}\begin{tabular}[t]{@{}c@{}}\rule{0pt}{12pt}2m\,48.1s\\[-.2pt]{\fontsize{5.8}{6.4}\selectfont \textcolor{BrickRed}{$ \blacktriangledown\,20.4\times $}}\\[-.2pt]2m\,41.7s\\[-.2pt]{\fontsize{5.8}{6.4}\selectfont \textcolor{BrickRed}{$ \blacktriangledown\,22.6\times $}}\end{tabular} & \hspace*{1pt}\begin{tabular}[t]{@{}c@{}}\rule{0pt}{12pt}2m\,50.3s\\[-.2pt]{\fontsize{5.8}{6.4}\selectfont \textcolor{BrickRed}{$ \blacktriangledown\,19.8\times $}}\\[-.2pt]2m\,43s\\[-.2pt]{\fontsize{5.8}{6.4}\selectfont \textcolor{BrickRed}{$ \blacktriangledown\,22.2\times $}}\end{tabular} & \hspace*{1pt}\begin{tabular}[t]{@{}c@{}}\rule{0pt}{12pt}4m\,54.3s\\[-.2pt]{\fontsize{5.8}{6.4}\selectfont \textcolor{BrickRed}{$ \blacktriangledown\,43.2\times $}}\\[-.2pt]4m\,39s\\[-.2pt]{\fontsize{5.8}{6.4}\selectfont \textcolor{BrickRed}{$ \blacktriangledown\,54.6\times $}}\end{tabular} & \hspace*{1pt}\begin{tabular}[t]{@{}c@{}}\rule{0pt}{12pt}25m\,5.5s\\[-.2pt]{\fontsize{5.8}{6.4}\selectfont \textcolor{BrickRed}{$ \blacktriangledown\,41.8\times $}}\\[-.2pt]25m\,0.7s\\[-.2pt]{\fontsize{5.8}{6.4}\selectfont \textcolor{BrickRed}{$ \blacktriangledown\,46.6\times $}}\end{tabular} & \hspace*{1pt}\begin{tabular}[t]{@{}c@{}}\rule{0pt}{12pt}33.9s\\[-.2pt]{\fontsize{5.8}{6.4}\selectfont \textcolor{BrickRed}{$ \blacktriangledown\,2.3\times $}}\\[-.2pt]32.9s\\[-.2pt]{\fontsize{5.8}{6.4}\selectfont \textcolor{BrickRed}{$ \blacktriangledown\,3.1\times $}}\end{tabular} & \hspace*{1pt}\begin{tabular}[t]{@{}c@{}}\rule{0pt}{12pt}34.4s\\[-.2pt]{\fontsize{5.8}{6.4}\selectfont \textcolor{BrickRed}{$ \blacktriangledown\,2.6\times $}}\\[-.2pt]32.7s\\[-.2pt]{\fontsize{5.8}{6.4}\selectfont \textcolor{BrickRed}{$ \blacktriangledown\,2.9\times $}}\end{tabular} & \hspace*{1pt}\begin{tabular}[t]{@{}c@{}}\rule{0pt}{12pt}1m\,8.2s\\[-.2pt]{\fontsize{5.8}{6.4}\selectfont \textcolor{BrickRed}{$ \blacktriangledown\,3.7\times $}}\\[-.2pt]1m\,5.1s\\[-.2pt]{\fontsize{5.8}{6.4}\selectfont \textcolor{BrickRed}{$ \blacktriangledown\,4.6\times $}}\end{tabular} & \hspace*{1pt}\begin{tabular}[t]{@{}c@{}}\rule{0pt}{12pt}3m\,9.8s\\[-.2pt]{\fontsize{5.8}{6.4}\selectfont \textcolor{BrickRed}{$ \blacktriangledown\,3.9\times $}}\\[-.2pt]3m\,3.7s\\[-.2pt]{\fontsize{5.8}{6.4}\selectfont \textcolor{BrickRed}{$ \blacktriangledown\,4.1\times $}}\end{tabular} & \hspace*{1pt}\begin{tabular}[t]{@{}c@{}}\rule{0pt}{12pt}3m\,28.3s\\[-.2pt]{\fontsize{5.8}{6.4}\selectfont \textcolor{BrickRed}{$ \blacktriangledown\,1.3\times $}}\\[-.2pt]3m\,22.8s\\[-.2pt]{\fontsize{5.8}{6.4}\selectfont \textcolor{BrickRed}{$ \blacktriangledown\,1.4\times $}}\end{tabular} & \hspace*{1pt}\begin{tabular}[t]{@{}c@{}}\rule{0pt}{12pt}3m\,53.2s\\[-.2pt]{\fontsize{5.8}{6.4}\selectfont \textcolor{BrickRed}{$ \blacktriangledown\,1.4\times $}}\\[-.2pt]3m\,43.5s\\[-.2pt]{\fontsize{5.8}{6.4}\selectfont \textcolor{BrickRed}{$ \blacktriangledown\,1.5\times $}}\end{tabular} & \hspace*{1pt}\begin{tabular}[t]{@{}c@{}}\rule{0pt}{12pt}7m\,14.6s\\[-.2pt]{\fontsize{5.8}{6.4}\selectfont \textcolor{BrickRed}{$ \blacktriangledown\,1.5\times $}}\\[-.2pt]6m\,57.2s\\[-.2pt]{\fontsize{5.8}{6.4}\selectfont \textcolor{BrickRed}{$ \blacktriangledown\,1.6\times $}}\end{tabular} & \hspace*{1pt}\begin{tabular}[t]{@{}c@{}}\rule{0pt}{12pt}51m\,22s\\[-.2pt]{\fontsize{5.8}{6.4}\selectfont \textcolor{gray}{n/a}}\\[-.2pt]51m\,17.2s\\[-.2pt]{\fontsize{5.8}{6.4}\selectfont \textcolor{gray}{n/a}}\end{tabular} & \hspace*{1pt}\begin{tabular}[t]{@{}c@{}}\rule{0pt}{12pt}1m\,10.3s\\[-.2pt]{\fontsize{5.8}{6.4}\selectfont \textcolor{ForestGreen}{$ \blacktriangle\,7.6\times $}}\\[-.2pt]1m\,9s\\[-.2pt]{\fontsize{5.8}{6.4}\selectfont \textcolor{ForestGreen}{$ \blacktriangle\,2.4\times $}}\end{tabular} & \hspace*{1pt}\begin{tabular}[t]{@{}c@{}}\rule{0pt}{12pt}1m\,8.6s\\[-.2pt]{\fontsize{5.8}{6.4}\selectfont \textcolor{ForestGreen}{$ \blacktriangle\,6.4\times $}}\\[-.2pt]1m\,6.9s\\[-.2pt]{\fontsize{5.8}{6.4}\selectfont \textcolor{ForestGreen}{$ \blacktriangle\,2.3\times $}}\end{tabular} & \hspace*{1pt}\begin{tabular}[t]{@{}c@{}}\rule{0pt}{12pt}2m\,21.5s\\[-.2pt]{\fontsize{5.8}{6.4}\selectfont \textcolor{ForestGreen}{$ \blacktriangle\,6.0\times $}}\\[-.2pt]2m\,18.4s\\[-.2pt]{\fontsize{5.8}{6.4}\selectfont \textcolor{ForestGreen}{$ \blacktriangle\,2.2\times $}}\end{tabular} & \hspace*{1pt}\begin{tabular}[t]{@{}c@{}}\rule{0pt}{12pt}15m\,56.8s\\[-.2pt]{\fontsize{5.8}{6.4}\selectfont \textcolor{gray}{n/a}}\\[-.2pt]15m\,51.9s\\[-.2pt]{\fontsize{5.8}{6.4}\selectfont \textcolor{gray}{n/a}}\end{tabular}\\[3pt]\hline
\BenchmarkSourceLabel{bench:aws-4-k20-w15:row:mpf-twostack-scalar-ksw2}{R37}\shortstack[l]{{\color{gray}R37}\\\textbf{BioseqDB}} & \begin{tabular}[t]{@{}l@{}}\rule{0pt}{12pt}MPF / Two-stack\\Scalar-DP / KSW2\end{tabular} & \hspace*{1pt}\begin{tabular}[t]{@{}c@{}}\rule{0pt}{12pt}9m\,16.3s\\[-.2pt]{\fontsize{5.8}{6.4}\selectfont \textcolor{BrickRed}{$ \blacktriangledown\,3.3\times $}}\\[-.2pt]9m\,13.5s\\[-.2pt]{\fontsize{5.8}{6.4}\selectfont \textcolor{BrickRed}{$ \blacktriangledown\,3.4\times $}}\end{tabular} & \hspace*{1pt}\begin{tabular}[t]{@{}c@{}}\rule{0pt}{12pt}16m\,8.5s\\[-.2pt]{\fontsize{5.8}{6.4}\selectfont \textcolor{BrickRed}{$ \blacktriangledown\,5.7\times $}}\\[-.2pt]16m\,5.1s\\[-.2pt]{\fontsize{5.8}{6.4}\selectfont \textcolor{BrickRed}{$ \blacktriangledown\,5.9\times $}}\end{tabular} & \hspace*{1pt}\begin{tabular}[t]{@{}c@{}}\rule{0pt}{12pt}33m\,26.6s\\[-.2pt]{\fontsize{5.8}{6.4}\selectfont \textcolor{BrickRed}{$ \blacktriangledown\,6.8\times $}}\\[-.2pt]33m\,20.5s\\[-.2pt]{\fontsize{5.8}{6.4}\selectfont \textcolor{BrickRed}{$ \blacktriangledown\,7.2\times $}}\end{tabular} & \textcolor{gray}{\ensuremath{\cdot}} & \textcolor{gray}{\ensuremath{\cdot}} & \textcolor{gray}{\ensuremath{\cdot}} & \textcolor{gray}{\ensuremath{\cdot}} & \textcolor{gray}{\ensuremath{\cdot}} & \hspace*{1pt}\begin{tabular}[t]{@{}c@{}}\rule{0pt}{12pt}3m\,21.4s\\[-.2pt]{\fontsize{5.8}{6.4}\selectfont \textcolor{gray}{$ \blacktriangleright\,3\% $}}\\[-.2pt]2m\,35s\\[-.2pt]{\fontsize{5.8}{6.4}\selectfont \textcolor{ForestGreen}{$ \blacktriangle\,24\% $}}\end{tabular} & \hspace*{1pt}\begin{tabular}[t]{@{}c@{}}\rule{0pt}{12pt}3m\,4.9s\\[-.2pt]{\fontsize{5.8}{6.4}\selectfont \textcolor{ForestGreen}{$ \blacktriangle\,21\% $}}\\[-.2pt]2m\,25.7s\\[-.2pt]{\fontsize{5.8}{6.4}\selectfont \textcolor{ForestGreen}{$ \blacktriangle\,1.5\times $}}\end{tabular} & \textcolor{gray}{\ensuremath{\cdot}} & \textcolor{gray}{\ensuremath{\cdot}} & \textcolor{gray}{\ensuremath{\cdot}} & \textcolor{gray}{\ensuremath{\cdot}} & \textcolor{gray}{\ensuremath{\cdot}} & \textcolor{gray}{\ensuremath{\cdot}}\\[3pt]\hline
\BenchmarkSourceLabel{bench:aws-4-k20-w15:row:mpf-scalar-scalar-ksw2}{R38}\shortstack[l]{{\color{gray}R38}} & \begin{tabular}[t]{@{}l@{}}\rule{0pt}{12pt}MPF / Scalar\\Scalar-DP / KSW2\end{tabular} & \hspace*{1pt}\begin{tabular}[t]{@{}c@{}}\rule{0pt}{12pt}11m\,11.3s\\[-.2pt]{\fontsize{5.8}{6.4}\selectfont \textcolor{BrickRed}{$ \blacktriangledown\,4.0\times $}}\\[-.2pt]11m\,7.6s\\[-.2pt]{\fontsize{5.8}{6.4}\selectfont \textcolor{BrickRed}{$ \blacktriangledown\,4.1\times $}}\end{tabular} & \hspace*{1pt}\begin{tabular}[t]{@{}c@{}}\rule{0pt}{12pt}24m\,35.8s\\[-.2pt]{\fontsize{5.8}{6.4}\selectfont \textcolor{BrickRed}{$ \blacktriangledown\,8.7\times $}}\\[-.2pt]24m\,32.7s\\[-.2pt]{\fontsize{5.8}{6.4}\selectfont \textcolor{BrickRed}{$ \blacktriangledown\,9.0\times $}}\end{tabular} & \hspace*{1pt}\begin{tabular}[t]{@{}c@{}}\rule{0pt}{12pt}49m\,58.6s\\[-.2pt]{\fontsize{5.8}{6.4}\selectfont \textcolor{BrickRed}{$ \blacktriangledown\,10.2\times $}}\\[-.2pt]49m\,52.4s\\[-.2pt]{\fontsize{5.8}{6.4}\selectfont \textcolor{BrickRed}{$ \blacktriangledown\,10.7\times $}}\end{tabular} & \hspace*{1pt}\begin{tabular}[t]{@{}c@{}}\rule{0pt}{12pt}{\fontsize{5.6}{6.4}\selectfont 3h\,11m\,57.8s}\\[-.2pt]{\fontsize{5.8}{6.4}\selectfont \textcolor{BrickRed}{$ \blacktriangledown\,7.7\times $}}\\[-.2pt]{\fontsize{5.6}{6.4}\selectfont 3h\,11m\,53.5s}\\[-.2pt]{\fontsize{5.8}{6.4}\selectfont \textcolor{BrickRed}{$ \blacktriangledown\,7.7\times $}}\end{tabular} & \hspace*{1pt}\begin{tabular}[t]{@{}c@{}}\rule{0pt}{12pt}16m\,14.9s\\[-.2pt]{\fontsize{5.8}{6.4}\selectfont \textcolor{BrickRed}{$ \blacktriangledown\,28.7\times $}}\\[-.2pt]16m\,11.2s\\[-.2pt]{\fontsize{5.8}{6.4}\selectfont \textcolor{BrickRed}{$ \blacktriangledown\,29.5\times $}}\end{tabular} & \hspace*{1pt}\begin{tabular}[t]{@{}c@{}}\rule{0pt}{12pt}38m\,2.4s\\[-.2pt]{\fontsize{5.8}{6.4}\selectfont \textcolor{BrickRed}{$ \blacktriangledown\,66.3\times $}}\\[-.2pt]37m\,59.5s\\[-.2pt]{\fontsize{5.8}{6.4}\selectfont \textcolor{BrickRed}{$ \blacktriangledown\,69.7\times $}}\end{tabular} & \hspace*{1pt}\begin{tabular}[t]{@{}c@{}}\rule{0pt}{12pt}{\fontsize{5.6}{6.4}\selectfont 1h\,11m\,44.1s}\\[-.2pt]{\fontsize{5.8}{6.4}\selectfont \textcolor{BrickRed}{$ \blacktriangledown\,63.1\times $}}\\[-.2pt]{\fontsize{5.6}{6.4}\selectfont 1h\,11m\,37.9s}\\[-.2pt]{\fontsize{5.8}{6.4}\selectfont \textcolor{BrickRed}{$ \blacktriangledown\,66.0\times $}}\end{tabular} & \hspace*{1pt}\begin{tabular}[t]{@{}c@{}}\rule{0pt}{12pt}{\fontsize{5.6}{6.4}\selectfont 2h\,23m\,50.8s}\\[-.2pt]{\fontsize{5.8}{6.4}\selectfont \textcolor{BrickRed}{$ \blacktriangledown\,45.5\times $}}\\[-.2pt]{\fontsize{5.6}{6.4}\selectfont 2h\,23m\,46.4s}\\[-.2pt]{\fontsize{5.8}{6.4}\selectfont \textcolor{BrickRed}{$ \blacktriangledown\,47.0\times $}}\end{tabular} & \hspace*{1pt}\begin{tabular}[t]{@{}c@{}}\rule{0pt}{12pt}5m\,8.2s\\[-.2pt]{\fontsize{5.8}{6.4}\selectfont \textcolor{BrickRed}{$ \blacktriangledown\,1.5\times $}}\\[-.2pt]4m\,20.4s\\[-.2pt]{\fontsize{5.8}{6.4}\selectfont \textcolor{BrickRed}{$ \blacktriangledown\,28\% $}}\end{tabular} & \hspace*{1pt}\begin{tabular}[t]{@{}c@{}}\rule{0pt}{12pt}4m\,17.9s\\[-.2pt]{\fontsize{5.8}{6.4}\selectfont \textcolor{BrickRed}{$ \blacktriangledown\,11\% $}}\\[-.2pt]3m\,42.6s\\[-.2pt]{\fontsize{5.8}{6.4}\selectfont \textcolor{gray}{$ \blacktriangleright\,0\% $}}\end{tabular} & \textcolor{gray}{\ensuremath{\cdot}} & \textcolor{gray}{\ensuremath{\cdot}} & \hspace*{1pt}\begin{tabular}[t]{@{}c@{}}\rule{0pt}{12pt}12m\,31.3s\\[-.2pt]{\fontsize{5.8}{6.4}\selectfont \textcolor{BrickRed}{$ \blacktriangledown\,10.7\times $}}\\[-.2pt]4m\,43.5s\\[-.2pt]{\fontsize{5.8}{6.4}\selectfont \textcolor{BrickRed}{$ \blacktriangledown\,4.1\times $}}\end{tabular} & \hspace*{1pt}\begin{tabular}[t]{@{}c@{}}\rule{0pt}{12pt}11m\,42.9s\\[-.2pt]{\fontsize{5.8}{6.4}\selectfont \textcolor{BrickRed}{$ \blacktriangledown\,10.2\times $}}\\[-.2pt]5m\,21.5s\\[-.2pt]{\fontsize{5.8}{6.4}\selectfont \textcolor{BrickRed}{$ \blacktriangledown\,4.8\times $}}\end{tabular} & \textcolor{gray}{\ensuremath{\cdot}} & \textcolor{gray}{\ensuremath{\cdot}}\\[3pt]\hline
\BenchmarkSourceLabel{bench:aws-4-k20-w15:row:mpf-scalar-scalar-block}{R39}\shortstack[l]{{\color{gray}R39}} & \begin{tabular}[t]{@{}l@{}}\rule{0pt}{12pt}MPF / Scalar\\Scalar-DP / Block\end{tabular} & \hspace*{1pt}\begin{tabular}[t]{@{}c@{}}\rule{0pt}{12pt}11m\,13.8s\\[-.2pt]{\fontsize{5.8}{6.4}\selectfont \textcolor{BrickRed}{$ \blacktriangledown\,4.0\times $}}\\[-.2pt]11m\,13s\\[-.2pt]{\fontsize{5.8}{6.4}\selectfont \textcolor{BrickRed}{$ \blacktriangledown\,4.2\times $}}\end{tabular} & \hspace*{1pt}\begin{tabular}[t]{@{}c@{}}\rule{0pt}{12pt}26m\,53.1s\\[-.2pt]{\fontsize{5.8}{6.4}\selectfont \textcolor{BrickRed}{$ \blacktriangledown\,9.5\times $}}\\[-.2pt]26m\,52.3s\\[-.2pt]{\fontsize{5.8}{6.4}\selectfont \textcolor{BrickRed}{$ \blacktriangledown\,9.9\times $}}\end{tabular} & \hspace*{1pt}\begin{tabular}[t]{@{}c@{}}\rule{0pt}{12pt}50m\,56.1s\\[-.2pt]{\fontsize{5.8}{6.4}\selectfont \textcolor{BrickRed}{$ \blacktriangledown\,10.4\times $}}\\[-.2pt]50m\,54.6s\\[-.2pt]{\fontsize{5.8}{6.4}\selectfont \textcolor{BrickRed}{$ \blacktriangledown\,10.9\times $}}\end{tabular} & \hspace*{1pt}\begin{tabular}[t]{@{}c@{}}\rule{0pt}{12pt}{\fontsize{5.6}{6.4}\selectfont 3h\,12m\,16.6s}\\[-.2pt]{\fontsize{5.8}{6.4}\selectfont \textcolor{BrickRed}{$ \blacktriangledown\,7.7\times $}}\\[-.2pt]{\fontsize{5.6}{6.4}\selectfont 3h\,12m\,14.4s}\\[-.2pt]{\fontsize{5.8}{6.4}\selectfont \textcolor{BrickRed}{$ \blacktriangledown\,7.7\times $}}\end{tabular} & \hspace*{1pt}\begin{tabular}[t]{@{}c@{}}\rule{0pt}{12pt}14m\,50.9s\\[-.2pt]{\fontsize{5.8}{6.4}\selectfont \textcolor{BrickRed}{$ \blacktriangledown\,26.2\times $}}\\[-.2pt]14m\,47.2s\\[-.2pt]{\fontsize{5.8}{6.4}\selectfont \textcolor{BrickRed}{$ \blacktriangledown\,26.9\times $}}\end{tabular} & \hspace*{1pt}\begin{tabular}[t]{@{}c@{}}\rule{0pt}{12pt}29m\,53.2s\\[-.2pt]{\fontsize{5.8}{6.4}\selectfont \textcolor{BrickRed}{$ \blacktriangledown\,52.1\times $}}\\[-.2pt]29m\,50.6s\\[-.2pt]{\fontsize{5.8}{6.4}\selectfont \textcolor{BrickRed}{$ \blacktriangledown\,54.7\times $}}\end{tabular} & \hspace*{1pt}\begin{tabular}[t]{@{}c@{}}\rule{0pt}{12pt}{\fontsize{5.6}{6.4}\selectfont 1h\,12m\,1.2s}\\[-.2pt]{\fontsize{5.8}{6.4}\selectfont \textcolor{BrickRed}{$ \blacktriangledown\,63.4\times $}}\\[-.2pt]{\fontsize{5.6}{6.4}\selectfont 1h\,11m\,55.6s}\\[-.2pt]{\fontsize{5.8}{6.4}\selectfont \textcolor{BrickRed}{$ \blacktriangledown\,66.3\times $}}\end{tabular} & \hspace*{1pt}\begin{tabular}[t]{@{}c@{}}\rule{0pt}{12pt}{\fontsize{5.6}{6.4}\selectfont 2h\,25m\,29.9s}\\[-.2pt]{\fontsize{5.8}{6.4}\selectfont \textcolor{BrickRed}{$ \blacktriangledown\,46.0\times $}}\\[-.2pt]{\fontsize{5.6}{6.4}\selectfont 2h\,25m\,25.6s}\\[-.2pt]{\fontsize{5.8}{6.4}\selectfont \textcolor{BrickRed}{$ \blacktriangledown\,47.5\times $}}\end{tabular} & \hspace*{1pt}\begin{tabular}[t]{@{}c@{}}\rule{0pt}{12pt}5m\,21.8s\\[-.2pt]{\fontsize{5.8}{6.4}\selectfont \textcolor{BrickRed}{$ \blacktriangledown\,1.5\times $}}\\[-.2pt]4m\,34.8s\\[-.2pt]{\fontsize{5.8}{6.4}\selectfont \textcolor{BrickRed}{$ \blacktriangledown\,1.4\times $}}\end{tabular} & \hspace*{1pt}\begin{tabular}[t]{@{}c@{}}\rule{0pt}{12pt}4m\,17.1s\\[-.2pt]{\fontsize{5.8}{6.4}\selectfont \textcolor{BrickRed}{$ \blacktriangledown\,10\% $}}\\[-.2pt]3m\,43.8s\\[-.2pt]{\fontsize{5.8}{6.4}\selectfont \textcolor{gray}{$ \blacktriangleright\,0\% $}}\end{tabular} & \textcolor{gray}{\ensuremath{\cdot}} & \textcolor{gray}{\ensuremath{\cdot}} & \hspace*{1pt}\begin{tabular}[t]{@{}c@{}}\rule{0pt}{12pt}12m\,28.5s\\[-.2pt]{\fontsize{5.8}{6.4}\selectfont \textcolor{BrickRed}{$ \blacktriangledown\,10.7\times $}}\\[-.2pt]4m\,39.8s\\[-.2pt]{\fontsize{5.8}{6.4}\selectfont \textcolor{BrickRed}{$ \blacktriangledown\,4.1\times $}}\end{tabular} & \hspace*{1pt}\begin{tabular}[t]{@{}c@{}}\rule{0pt}{12pt}11m\,56.7s\\[-.2pt]{\fontsize{5.8}{6.4}\selectfont \textcolor{BrickRed}{$ \blacktriangledown\,10.4\times $}}\\[-.2pt]5m\,31.3s\\[-.2pt]{\fontsize{5.8}{6.4}\selectfont \textcolor{BrickRed}{$ \blacktriangledown\,5.0\times $}}\end{tabular} & \textcolor{gray}{\ensuremath{\cdot}} & \textcolor{gray}{\ensuremath{\cdot}}\\[3pt]\hline
\BenchmarkSourceLabel{bench:aws-4-k20-w15:row:mpf-scalar-simd-ksw2}{R40}\shortstack[l]{{\color{gray}R40}} & \begin{tabular}[t]{@{}l@{}}\rule{0pt}{12pt}MPF / Scalar\\SIMD-DP / KSW2\end{tabular} & \hspace*{1pt}\begin{tabular}[t]{@{}c@{}}\rule{0pt}{12pt}11m\,14.3s\\[-.2pt]{\fontsize{5.8}{6.4}\selectfont \textcolor{BrickRed}{$ \blacktriangledown\,4.0\times $}}\\[-.2pt]11m\,12.2s\\[-.2pt]{\fontsize{5.8}{6.4}\selectfont \textcolor{BrickRed}{$ \blacktriangledown\,4.2\times $}}\end{tabular} & \hspace*{1pt}\begin{tabular}[t]{@{}c@{}}\rule{0pt}{12pt}24m\,37.5s\\[-.2pt]{\fontsize{5.8}{6.4}\selectfont \textcolor{BrickRed}{$ \blacktriangledown\,8.7\times $}}\\[-.2pt]24m\,35.6s\\[-.2pt]{\fontsize{5.8}{6.4}\selectfont \textcolor{BrickRed}{$ \blacktriangledown\,9.1\times $}}\end{tabular} & \hspace*{1pt}\begin{tabular}[t]{@{}c@{}}\rule{0pt}{12pt}44m\,15.3s\\[-.2pt]{\fontsize{5.8}{6.4}\selectfont \textcolor{BrickRed}{$ \blacktriangledown\,9.0\times $}}\\[-.2pt]44m\,9.6s\\[-.2pt]{\fontsize{5.8}{6.4}\selectfont \textcolor{BrickRed}{$ \blacktriangledown\,9.5\times $}}\end{tabular} & \hspace*{1pt}\begin{tabular}[t]{@{}c@{}}\rule{0pt}{12pt}{\fontsize{5.6}{6.4}\selectfont 1h\,34m\,25.1s}\\[-.2pt]{\fontsize{5.8}{6.4}\selectfont \textcolor{BrickRed}{$ \blacktriangledown\,3.8\times $}}\\[-.2pt]{\fontsize{5.6}{6.4}\selectfont 1h\,34m\,20.6s}\\[-.2pt]{\fontsize{5.8}{6.4}\selectfont \textcolor{BrickRed}{$ \blacktriangledown\,3.8\times $}}\end{tabular} & \hspace*{1pt}\begin{tabular}[t]{@{}c@{}}\rule{0pt}{12pt}14m\,52.8s\\[-.2pt]{\fontsize{5.8}{6.4}\selectfont \textcolor{BrickRed}{$ \blacktriangledown\,26.3\times $}}\\[-.2pt]14m\,49.1s\\[-.2pt]{\fontsize{5.8}{6.4}\selectfont \textcolor{BrickRed}{$ \blacktriangledown\,27.0\times $}}\end{tabular} & \hspace*{1pt}\begin{tabular}[t]{@{}c@{}}\rule{0pt}{12pt}32m\,52.9s\\[-.2pt]{\fontsize{5.8}{6.4}\selectfont \textcolor{BrickRed}{$ \blacktriangledown\,57.3\times $}}\\[-.2pt]32m\,50s\\[-.2pt]{\fontsize{5.8}{6.4}\selectfont \textcolor{BrickRed}{$ \blacktriangledown\,60.2\times $}}\end{tabular} & \hspace*{1pt}\begin{tabular}[t]{@{}c@{}}\rule{0pt}{12pt}{\fontsize{5.6}{6.4}\selectfont 1h\,8m\,24.8s}\\[-.2pt]{\fontsize{5.8}{6.4}\selectfont \textcolor{BrickRed}{$ \blacktriangledown\,60.2\times $}}\\[-.2pt]{\fontsize{5.6}{6.4}\selectfont 1h\,8m\,18.7s}\\[-.2pt]{\fontsize{5.8}{6.4}\selectfont \textcolor{BrickRed}{$ \blacktriangledown\,62.9\times $}}\end{tabular} & \hspace*{1pt}\begin{tabular}[t]{@{}c@{}}\rule{0pt}{12pt}{\fontsize{5.6}{6.4}\selectfont 1h\,26m\,14.3s}\\[-.2pt]{\fontsize{5.8}{6.4}\selectfont \textcolor{BrickRed}{$ \blacktriangledown\,27.3\times $}}\\[-.2pt]{\fontsize{5.6}{6.4}\selectfont 1h\,26m\,9.8s}\\[-.2pt]{\fontsize{5.8}{6.4}\selectfont \textcolor{BrickRed}{$ \blacktriangledown\,28.1\times $}}\end{tabular} & \hspace*{1pt}\begin{tabular}[t]{@{}c@{}}\rule{0pt}{12pt}4m\,53s\\[-.2pt]{\fontsize{5.8}{6.4}\selectfont \textcolor{BrickRed}{$ \blacktriangledown\,1.4\times $}}\\[-.2pt]4m\,8.2s\\[-.2pt]{\fontsize{5.8}{6.4}\selectfont \textcolor{BrickRed}{$ \blacktriangledown\,22\% $}}\end{tabular} & \hspace*{1pt}\begin{tabular}[t]{@{}c@{}}\rule{0pt}{12pt}5m\,9.3s\\[-.2pt]{\fontsize{5.8}{6.4}\selectfont \textcolor{BrickRed}{$ \blacktriangledown\,1.3\times $}}\\[-.2pt]4m\,31.7s\\[-.2pt]{\fontsize{5.8}{6.4}\selectfont \textcolor{BrickRed}{$ \blacktriangledown\,22\% $}}\end{tabular} & \textcolor{gray}{\ensuremath{\cdot}} & \textcolor{gray}{\ensuremath{\cdot}} & \hspace*{1pt}\begin{tabular}[t]{@{}c@{}}\rule{0pt}{12pt}12m\,34.5s\\[-.2pt]{\fontsize{5.8}{6.4}\selectfont \textcolor{BrickRed}{$ \blacktriangledown\,10.7\times $}}\\[-.2pt]4m\,44.3s\\[-.2pt]{\fontsize{5.8}{6.4}\selectfont \textcolor{BrickRed}{$ \blacktriangledown\,4.1\times $}}\end{tabular} & \hspace*{1pt}\begin{tabular}[t]{@{}c@{}}\rule{0pt}{12pt}11m\,58.3s\\[-.2pt]{\fontsize{5.8}{6.4}\selectfont \textcolor{BrickRed}{$ \blacktriangledown\,10.5\times $}}\\[-.2pt]5m\,31.6s\\[-.2pt]{\fontsize{5.8}{6.4}\selectfont \textcolor{BrickRed}{$ \blacktriangledown\,5.0\times $}}\end{tabular} & \textcolor{gray}{\ensuremath{\cdot}} & \textcolor{gray}{\ensuremath{\cdot}}\\[3pt]\hline
\BenchmarkSourceLabel{bench:aws-4-k20-w15:row:mpf-scalar-rmq-ksw2}{R41}\shortstack[l]{{\color{gray}R41}} & \begin{tabular}[t]{@{}l@{}}\rule{0pt}{12pt}MPF / Scalar\\RMQ / KSW2\end{tabular} & \hspace*{1pt}\begin{tabular}[t]{@{}c@{}}\rule{0pt}{12pt}11m\,14.6s\\[-.2pt]{\fontsize{5.8}{6.4}\selectfont \textcolor{BrickRed}{$ \blacktriangledown\,4.0\times $}}\\[-.2pt]11m\,12.5s\\[-.2pt]{\fontsize{5.8}{6.4}\selectfont \textcolor{BrickRed}{$ \blacktriangledown\,4.2\times $}}\end{tabular} & \hspace*{1pt}\begin{tabular}[t]{@{}c@{}}\rule{0pt}{12pt}25m\,12.4s\\[-.2pt]{\fontsize{5.8}{6.4}\selectfont \textcolor{BrickRed}{$ \blacktriangledown\,8.9\times $}}\\[-.2pt]25m\,10.5s\\[-.2pt]{\fontsize{5.8}{6.4}\selectfont \textcolor{BrickRed}{$ \blacktriangledown\,9.3\times $}}\end{tabular} & \hspace*{1pt}\begin{tabular}[t]{@{}c@{}}\rule{0pt}{12pt}53m\,45.8s\\[-.2pt]{\fontsize{5.8}{6.4}\selectfont \textcolor{BrickRed}{$ \blacktriangledown\,11.0\times $}}\\[-.2pt]53m\,39.5s\\[-.2pt]{\fontsize{5.8}{6.4}\selectfont \textcolor{BrickRed}{$ \blacktriangledown\,11.5\times $}}\end{tabular} & \hspace*{1pt}\begin{tabular}[t]{@{}c@{}}\rule{0pt}{12pt}{\fontsize{5.6}{6.4}\selectfont 3h\,11m\,29.7s}\\[-.2pt]{\fontsize{5.8}{6.4}\selectfont \textcolor{BrickRed}{$ \blacktriangledown\,7.6\times $}}\\[-.2pt]{\fontsize{5.6}{6.4}\selectfont 3h\,11m\,25.2s}\\[-.2pt]{\fontsize{5.8}{6.4}\selectfont \textcolor{BrickRed}{$ \blacktriangledown\,7.7\times $}}\end{tabular} & \hspace*{1pt}\begin{tabular}[t]{@{}c@{}}\rule{0pt}{12pt}15m\,59.1s\\[-.2pt]{\fontsize{5.8}{6.4}\selectfont \textcolor{BrickRed}{$ \blacktriangledown\,28.3\times $}}\\[-.2pt]15m\,55.5s\\[-.2pt]{\fontsize{5.8}{6.4}\selectfont \textcolor{BrickRed}{$ \blacktriangledown\,29.0\times $}}\end{tabular} & \hspace*{1pt}\begin{tabular}[t]{@{}c@{}}\rule{0pt}{12pt}29m\,55.2s\\[-.2pt]{\fontsize{5.8}{6.4}\selectfont \textcolor{BrickRed}{$ \blacktriangledown\,52.2\times $}}\\[-.2pt]29m\,52.6s\\[-.2pt]{\fontsize{5.8}{6.4}\selectfont \textcolor{BrickRed}{$ \blacktriangledown\,54.8\times $}}\end{tabular} & \hspace*{1pt}\begin{tabular}[t]{@{}c@{}}\rule{0pt}{12pt}{\fontsize{5.6}{6.4}\selectfont 1h\,11m\,28.9s}\\[-.2pt]{\fontsize{5.8}{6.4}\selectfont \textcolor{BrickRed}{$ \blacktriangledown\,62.9\times $}}\\[-.2pt]{\fontsize{5.6}{6.4}\selectfont 1h\,11m\,22.8s}\\[-.2pt]{\fontsize{5.8}{6.4}\selectfont \textcolor{BrickRed}{$ \blacktriangledown\,65.8\times $}}\end{tabular} & \hspace*{1pt}\begin{tabular}[t]{@{}c@{}}\rule{0pt}{12pt}{\fontsize{5.6}{6.4}\selectfont 1h\,56m\,30.8s}\\[-.2pt]{\fontsize{5.8}{6.4}\selectfont \textcolor{BrickRed}{$ \blacktriangledown\,36.8\times $}}\\[-.2pt]{\fontsize{5.6}{6.4}\selectfont 1h\,56m\,26.1s}\\[-.2pt]{\fontsize{5.8}{6.4}\selectfont \textcolor{BrickRed}{$ \blacktriangledown\,38.0\times $}}\end{tabular} & \hspace*{1pt}\begin{tabular}[t]{@{}c@{}}\rule{0pt}{12pt}5m\,17.4s\\[-.2pt]{\fontsize{5.8}{6.4}\selectfont \textcolor{BrickRed}{$ \blacktriangledown\,1.5\times $}}\\[-.2pt]4m\,30.6s\\[-.2pt]{\fontsize{5.8}{6.4}\selectfont \textcolor{BrickRed}{$ \blacktriangledown\,1.3\times $}}\end{tabular} & \hspace*{1pt}\begin{tabular}[t]{@{}c@{}}\rule{0pt}{12pt}4m\,16.1s\\[-.2pt]{\fontsize{5.8}{6.4}\selectfont \textcolor{BrickRed}{$ \blacktriangledown\,10\% $}}\\[-.2pt]3m\,41.1s\\[-.2pt]{\fontsize{5.8}{6.4}\selectfont \textcolor{gray}{$ \blacktriangleright\,1\% $}}\end{tabular} & \textcolor{gray}{\ensuremath{\cdot}} & \textcolor{gray}{\ensuremath{\cdot}} & \hspace*{1pt}\begin{tabular}[t]{@{}c@{}}\rule{0pt}{12pt}12m\,25.5s\\[-.2pt]{\fontsize{5.8}{6.4}\selectfont \textcolor{BrickRed}{$ \blacktriangledown\,10.6\times $}}\\[-.2pt]4m\,33.9s\\[-.2pt]{\fontsize{5.8}{6.4}\selectfont \textcolor{BrickRed}{$ \blacktriangledown\,4.0\times $}}\end{tabular} & \hspace*{1pt}\begin{tabular}[t]{@{}c@{}}\rule{0pt}{12pt}12m\,4.7s\\[-.2pt]{\fontsize{5.8}{6.4}\selectfont \textcolor{BrickRed}{$ \blacktriangledown\,10.6\times $}}\\[-.2pt]5m\,35.2s\\[-.2pt]{\fontsize{5.8}{6.4}\selectfont \textcolor{BrickRed}{$ \blacktriangledown\,5.0\times $}}\end{tabular} & \textcolor{gray}{\ensuremath{\cdot}} & \textcolor{gray}{\ensuremath{\cdot}}\\[3pt]\hline
\BenchmarkSourceLabel{bench:aws-4-k20-w15:row:mpf-simdref-rmq-block}{R42}\shortstack[l]{{\color{gray}R42}} & \begin{tabular}[t]{@{}l@{}}\rule{0pt}{12pt}MPF / SIMD-ref\\RMQ / Block\end{tabular} & \hspace*{1pt}\begin{tabular}[t]{@{}c@{}}\rule{0pt}{12pt}11m\,19.5s\\[-.2pt]{\fontsize{5.8}{6.4}\selectfont \textcolor{BrickRed}{$ \blacktriangledown\,4.0\times $}}\\[-.2pt]11m\,18.8s\\[-.2pt]{\fontsize{5.8}{6.4}\selectfont \textcolor{BrickRed}{$ \blacktriangledown\,4.2\times $}}\end{tabular} & \hspace*{1pt}\begin{tabular}[t]{@{}c@{}}\rule{0pt}{12pt}23m\,45.6s\\[-.2pt]{\fontsize{5.8}{6.4}\selectfont \textcolor{BrickRed}{$ \blacktriangledown\,8.4\times $}}\\[-.2pt]23m\,44.9s\\[-.2pt]{\fontsize{5.8}{6.4}\selectfont \textcolor{BrickRed}{$ \blacktriangledown\,8.7\times $}}\end{tabular} & \hspace*{1pt}\begin{tabular}[t]{@{}c@{}}\rule{0pt}{12pt}45m\,51.5s\\[-.2pt]{\fontsize{5.8}{6.4}\selectfont \textcolor{BrickRed}{$ \blacktriangledown\,9.3\times $}}\\[-.2pt]45m\,50.2s\\[-.2pt]{\fontsize{5.8}{6.4}\selectfont \textcolor{BrickRed}{$ \blacktriangledown\,9.9\times $}}\end{tabular} & \hspace*{1pt}\begin{tabular}[t]{@{}c@{}}\rule{0pt}{12pt}{\fontsize{5.6}{6.4}\selectfont 2h\,47m\,50.9s}\\[-.2pt]{\fontsize{5.8}{6.4}\selectfont \textcolor{BrickRed}{$ \blacktriangledown\,6.7\times $}}\\[-.2pt]{\fontsize{5.6}{6.4}\selectfont 2h\,47m\,48.7s}\\[-.2pt]{\fontsize{5.8}{6.4}\selectfont \textcolor{BrickRed}{$ \blacktriangledown\,6.7\times $}}\end{tabular} & \hspace*{1pt}\begin{tabular}[t]{@{}c@{}}\rule{0pt}{12pt}12m\,10.8s\\[-.2pt]{\fontsize{5.8}{6.4}\selectfont \textcolor{BrickRed}{$ \blacktriangledown\,21.5\times $}}\\[-.2pt]12m\,7.1s\\[-.2pt]{\fontsize{5.8}{6.4}\selectfont \textcolor{BrickRed}{$ \blacktriangledown\,22.1\times $}}\end{tabular} & \hspace*{1pt}\begin{tabular}[t]{@{}c@{}}\rule{0pt}{12pt}31m\,7.5s\\[-.2pt]{\fontsize{5.8}{6.4}\selectfont \textcolor{BrickRed}{$ \blacktriangledown\,54.3\times $}}\\[-.2pt]31m\,4.5s\\[-.2pt]{\fontsize{5.8}{6.4}\selectfont \textcolor{BrickRed}{$ \blacktriangledown\,57.0\times $}}\end{tabular} & \hspace*{1pt}\begin{tabular}[t]{@{}c@{}}\rule{0pt}{12pt}47m\,1.1s\\[-.2pt]{\fontsize{5.8}{6.4}\selectfont \textcolor{BrickRed}{$ \blacktriangledown\,41.4\times $}}\\[-.2pt]46m\,55.3s\\[-.2pt]{\fontsize{5.8}{6.4}\selectfont \textcolor{BrickRed}{$ \blacktriangledown\,43.2\times $}}\end{tabular} & \textcolor{gray}{\ensuremath{\cdot}} & \hspace*{1pt}\begin{tabular}[t]{@{}c@{}}\rule{0pt}{12pt}4m\,39.8s\\[-.2pt]{\fontsize{5.8}{6.4}\selectfont \textcolor{BrickRed}{$ \blacktriangledown\,1.3\times $}}\\[-.2pt]3m\,55.2s\\[-.2pt]{\fontsize{5.8}{6.4}\selectfont \textcolor{BrickRed}{$ \blacktriangledown\,16\% $}}\end{tabular} & \hspace*{1pt}\begin{tabular}[t]{@{}c@{}}\rule{0pt}{12pt}4m\,28.8s\\[-.2pt]{\fontsize{5.8}{6.4}\selectfont \textcolor{BrickRed}{$ \blacktriangledown\,15\% $}}\\[-.2pt]3m\,52s\\[-.2pt]{\fontsize{5.8}{6.4}\selectfont \textcolor{gray}{$ \blacktriangleright\,4\% $}}\end{tabular} & \textcolor{gray}{\ensuremath{\cdot}} & \textcolor{gray}{\ensuremath{\cdot}} & \hspace*{1pt}\begin{tabular}[t]{@{}c@{}}\rule{0pt}{12pt}10m\,49.3s\\[-.2pt]{\fontsize{5.8}{6.4}\selectfont \textcolor{BrickRed}{$ \blacktriangledown\,9.2\times $}}\\[-.2pt]4m\,37.1s\\[-.2pt]{\fontsize{5.8}{6.4}\selectfont \textcolor{BrickRed}{$ \blacktriangledown\,4.0\times $}}\end{tabular} & \hspace*{1pt}\begin{tabular}[t]{@{}c@{}}\rule{0pt}{12pt}10m\,38.4s\\[-.2pt]{\fontsize{5.8}{6.4}\selectfont \textcolor{BrickRed}{$ \blacktriangledown\,9.3\times $}}\\[-.2pt]5m\,1.8s\\[-.2pt]{\fontsize{5.8}{6.4}\selectfont \textcolor{BrickRed}{$ \blacktriangledown\,4.5\times $}}\end{tabular} & \textcolor{gray}{\ensuremath{\cdot}} & \textcolor{gray}{\ensuremath{\cdot}}\\[3pt]\hline
\BenchmarkSourceLabel{bench:aws-4-k20-w15:row:mpf-simdref-rmq-ksw2}{R43}\shortstack[l]{{\color{gray}R43}} & \begin{tabular}[t]{@{}l@{}}\rule{0pt}{12pt}MPF / SIMD-ref\\RMQ / KSW2\end{tabular} & \hspace*{1pt}\begin{tabular}[t]{@{}c@{}}\rule{0pt}{12pt}11m\,21.2s\\[-.2pt]{\fontsize{5.8}{6.4}\selectfont \textcolor{BrickRed}{$ \blacktriangledown\,4.1\times $}}\\[-.2pt]11m\,17.6s\\[-.2pt]{\fontsize{5.8}{6.4}\selectfont \textcolor{BrickRed}{$ \blacktriangledown\,4.2\times $}}\end{tabular} & \hspace*{1pt}\begin{tabular}[t]{@{}c@{}}\rule{0pt}{12pt}23m\,50.3s\\[-.2pt]{\fontsize{5.8}{6.4}\selectfont \textcolor{BrickRed}{$ \blacktriangledown\,8.4\times $}}\\[-.2pt]23m\,47.4s\\[-.2pt]{\fontsize{5.8}{6.4}\selectfont \textcolor{BrickRed}{$ \blacktriangledown\,8.8\times $}}\end{tabular} & \hspace*{1pt}\begin{tabular}[t]{@{}c@{}}\rule{0pt}{12pt}47m\,6.3s\\[-.2pt]{\fontsize{5.8}{6.4}\selectfont \textcolor{BrickRed}{$ \blacktriangledown\,9.6\times $}}\\[-.2pt]47m\,0.3s\\[-.2pt]{\fontsize{5.8}{6.4}\selectfont \textcolor{BrickRed}{$ \blacktriangledown\,10.1\times $}}\end{tabular} & \hspace*{1pt}\begin{tabular}[t]{@{}c@{}}\rule{0pt}{12pt}{\fontsize{5.6}{6.4}\selectfont 3h\,11m\,12.9s}\\[-.2pt]{\fontsize{5.8}{6.4}\selectfont \textcolor{BrickRed}{$ \blacktriangledown\,7.6\times $}}\\[-.2pt]{\fontsize{5.6}{6.4}\selectfont 3h\,11m\,8.5s}\\[-.2pt]{\fontsize{5.8}{6.4}\selectfont \textcolor{BrickRed}{$ \blacktriangledown\,7.6\times $}}\end{tabular} & \hspace*{1pt}\begin{tabular}[t]{@{}c@{}}\rule{0pt}{12pt}12m\,7.2s\\[-.2pt]{\fontsize{5.8}{6.4}\selectfont \textcolor{BrickRed}{$ \blacktriangledown\,21.4\times $}}\\[-.2pt]12m\,3.6s\\[-.2pt]{\fontsize{5.8}{6.4}\selectfont \textcolor{BrickRed}{$ \blacktriangledown\,22.0\times $}}\end{tabular} & \hspace*{1pt}\begin{tabular}[t]{@{}c@{}}\rule{0pt}{12pt}25m\,37s\\[-.2pt]{\fontsize{5.8}{6.4}\selectfont \textcolor{BrickRed}{$ \blacktriangledown\,44.7\times $}}\\[-.2pt]25m\,34.2s\\[-.2pt]{\fontsize{5.8}{6.4}\selectfont \textcolor{BrickRed}{$ \blacktriangledown\,46.9\times $}}\end{tabular} & \hspace*{1pt}\begin{tabular}[t]{@{}c@{}}\rule{0pt}{12pt}47m\,19.6s\\[-.2pt]{\fontsize{5.8}{6.4}\selectfont \textcolor{BrickRed}{$ \blacktriangledown\,41.6\times $}}\\[-.2pt]47m\,13.3s\\[-.2pt]{\fontsize{5.8}{6.4}\selectfont \textcolor{BrickRed}{$ \blacktriangledown\,43.5\times $}}\end{tabular} & \textcolor{gray}{\ensuremath{\cdot}} & \hspace*{1pt}\begin{tabular}[t]{@{}c@{}}\rule{0pt}{12pt}4m\,43.9s\\[-.2pt]{\fontsize{5.8}{6.4}\selectfont \textcolor{BrickRed}{$ \blacktriangledown\,1.4\times $}}\\[-.2pt]3m\,57.8s\\[-.2pt]{\fontsize{5.8}{6.4}\selectfont \textcolor{BrickRed}{$ \blacktriangledown\,17\% $}}\end{tabular} & \hspace*{1pt}\begin{tabular}[t]{@{}c@{}}\rule{0pt}{12pt}4m\,55.1s\\[-.2pt]{\fontsize{5.8}{6.4}\selectfont \textcolor{BrickRed}{$ \blacktriangledown\,27\% $}}\\[-.2pt]4m\,16.8s\\[-.2pt]{\fontsize{5.8}{6.4}\selectfont \textcolor{BrickRed}{$ \blacktriangledown\,15\% $}}\end{tabular} & \textcolor{gray}{\ensuremath{\cdot}} & \textcolor{gray}{\ensuremath{\cdot}} & \hspace*{1pt}\begin{tabular}[t]{@{}c@{}}\rule{0pt}{12pt}10m\,52.3s\\[-.2pt]{\fontsize{5.8}{6.4}\selectfont \textcolor{BrickRed}{$ \blacktriangledown\,9.3\times $}}\\[-.2pt]4m\,28.1s\\[-.2pt]{\fontsize{5.8}{6.4}\selectfont \textcolor{BrickRed}{$ \blacktriangledown\,3.9\times $}}\end{tabular} & \hspace*{1pt}\begin{tabular}[t]{@{}c@{}}\rule{0pt}{12pt}10m\,29.8s\\[-.2pt]{\fontsize{5.8}{6.4}\selectfont \textcolor{BrickRed}{$ \blacktriangledown\,9.2\times $}}\\[-.2pt]4m\,56.7s\\[-.2pt]{\fontsize{5.8}{6.4}\selectfont \textcolor{BrickRed}{$ \blacktriangledown\,4.4\times $}}\end{tabular} & \textcolor{gray}{\ensuremath{\cdot}} & \textcolor{gray}{\ensuremath{\cdot}}\\[3pt]\hline
\end{tabular}\par}
\smallskip{\footnotesize\noindent Each cell: A and B times, $\cdot$: no measurement. ---\textsuperscript{g}: GPU unavailable. n/a: no comparison.\par
\noindent Blue: one lowest A and one lowest B measurement for each column for each configuration.\par
\noindent Gray: minimap2. Yellow: my preferred MPF / Bounded / SIMD-DP / Block configuration.\par
\noindent MPF/MPS: minimizers per fragment/sequence. SIMD-ref: naive SIMD scan. Bounded: bounded SIMD algorithm.\par
}
\clearpage
\noindent{\large\bfseries Table \thetable: (C): AWS: 4 Gbp, k=20, w=15\hfill 7/7}\par\smallskip
{\fontsize{7.1}{8.1}\selectfont\setlength{\tabcolsep}{1pt}\setlength{\fboxsep}{2pt}\renewcommand{\arraystretch}{1.12}
\noindent\begin{tabular}{@{}p{14mm}p{27mm}|p{13.55mm}|p{13.55mm}|p{13.55mm}|p{13.55mm}|p{13.55mm}|p{13.55mm}|p{13.55mm}|p{13.55mm}|p{13.55mm}|p{13.55mm}|p{13.55mm}|p{13.55mm}|p{13.55mm}|p{13.55mm}|p{13.55mm}|p{13.55mm}|@{}}\hline
Engine & Execution & \multicolumn{8}{c|}{\textbf{Index}} & \multicolumn{8}{c|}{\textbf{No index}}\\
 &  & \multicolumn{8}{c|}{\textbf{4 Gbp}} & \multicolumn{8}{c|}{\textbf{4 Gbp}}\\
 &  & \multicolumn{8}{c|}{\textbf{k=20, w=15}} & \multicolumn{8}{c|}{\textbf{k=20, w=15}}\\
 &  & \multicolumn{4}{c|}{\textbf{DNA4}} & \multicolumn{4}{c|}{\textbf{DNA5}} & \multicolumn{4}{c|}{\textbf{DNA4}} & \multicolumn{4}{c|}{\textbf{DNA5}}\\
 &  & \multicolumn{1}{c|}{\textbf{M7A}} & \multicolumn{1}{c|}{\textbf{M7I}} & \multicolumn{1}{c|}{\textbf{M4}} & \multicolumn{1}{c|}{\textbf{C7G}} & \multicolumn{1}{c|}{\textbf{M7A}} & \multicolumn{1}{c|}{\textbf{M7I}} & \multicolumn{1}{c|}{\textbf{M4}} & \multicolumn{1}{c|}{\textbf{C7G}} & \multicolumn{1}{c|}{\textbf{M7A}} & \multicolumn{1}{c|}{\textbf{M7I}} & \multicolumn{1}{c|}{\textbf{M4}} & \multicolumn{1}{c|}{\textbf{C7G}} & \multicolumn{1}{c|}{\textbf{M7A}} & \multicolumn{1}{c|}{\textbf{M7I}} & \multicolumn{1}{c|}{\textbf{M4}} & \multicolumn{1}{c|}{\textbf{C7G}}\\
 & & {\color{gray}C1} & {\color{gray}C2} & {\color{gray}C3} & {\color{gray}C4} & {\color{gray}C5} & {\color{gray}C6} & {\color{gray}C7} & {\color{gray}C8} & {\color{gray}C9} & {\color{gray}C10} & {\color{gray}C11} & {\color{gray}C12} & {\color{gray}C13} & {\color{gray}C14} & {\color{gray}C15} & {\color{gray}C16}\\
\hline
\rowcolor{BenchBaseline}[0.5pt][0.5pt]\shortstack[l]{{\color{gray}R1}\\\textbf{minimap2}} & \textbf{Baseline} & \hspace*{1pt}\begin{tabular}[t]{@{}c@{}}\rule{0pt}{12pt}2m\,48.1s\\[-.2pt]{\fontsize{5.8}{6.4}\selectfont \textcolor{BrickRed}{$ \blacktriangledown\,20.4\times $}}\\[-.2pt]2m\,41.7s\\[-.2pt]{\fontsize{5.8}{6.4}\selectfont \textcolor{BrickRed}{$ \blacktriangledown\,22.6\times $}}\end{tabular} & \hspace*{1pt}\begin{tabular}[t]{@{}c@{}}\rule{0pt}{12pt}2m\,50.3s\\[-.2pt]{\fontsize{5.8}{6.4}\selectfont \textcolor{BrickRed}{$ \blacktriangledown\,19.8\times $}}\\[-.2pt]2m\,43s\\[-.2pt]{\fontsize{5.8}{6.4}\selectfont \textcolor{BrickRed}{$ \blacktriangledown\,22.2\times $}}\end{tabular} & \hspace*{1pt}\begin{tabular}[t]{@{}c@{}}\rule{0pt}{12pt}4m\,54.3s\\[-.2pt]{\fontsize{5.8}{6.4}\selectfont \textcolor{BrickRed}{$ \blacktriangledown\,43.2\times $}}\\[-.2pt]4m\,39s\\[-.2pt]{\fontsize{5.8}{6.4}\selectfont \textcolor{BrickRed}{$ \blacktriangledown\,54.6\times $}}\end{tabular} & \hspace*{1pt}\begin{tabular}[t]{@{}c@{}}\rule{0pt}{12pt}25m\,5.5s\\[-.2pt]{\fontsize{5.8}{6.4}\selectfont \textcolor{BrickRed}{$ \blacktriangledown\,41.8\times $}}\\[-.2pt]25m\,0.7s\\[-.2pt]{\fontsize{5.8}{6.4}\selectfont \textcolor{BrickRed}{$ \blacktriangledown\,46.6\times $}}\end{tabular} & \hspace*{1pt}\begin{tabular}[t]{@{}c@{}}\rule{0pt}{12pt}33.9s\\[-.2pt]{\fontsize{5.8}{6.4}\selectfont \textcolor{BrickRed}{$ \blacktriangledown\,2.3\times $}}\\[-.2pt]32.9s\\[-.2pt]{\fontsize{5.8}{6.4}\selectfont \textcolor{BrickRed}{$ \blacktriangledown\,3.1\times $}}\end{tabular} & \hspace*{1pt}\begin{tabular}[t]{@{}c@{}}\rule{0pt}{12pt}34.4s\\[-.2pt]{\fontsize{5.8}{6.4}\selectfont \textcolor{BrickRed}{$ \blacktriangledown\,2.6\times $}}\\[-.2pt]32.7s\\[-.2pt]{\fontsize{5.8}{6.4}\selectfont \textcolor{BrickRed}{$ \blacktriangledown\,2.9\times $}}\end{tabular} & \hspace*{1pt}\begin{tabular}[t]{@{}c@{}}\rule{0pt}{12pt}1m\,8.2s\\[-.2pt]{\fontsize{5.8}{6.4}\selectfont \textcolor{BrickRed}{$ \blacktriangledown\,3.7\times $}}\\[-.2pt]1m\,5.1s\\[-.2pt]{\fontsize{5.8}{6.4}\selectfont \textcolor{BrickRed}{$ \blacktriangledown\,4.6\times $}}\end{tabular} & \hspace*{1pt}\begin{tabular}[t]{@{}c@{}}\rule{0pt}{12pt}3m\,9.8s\\[-.2pt]{\fontsize{5.8}{6.4}\selectfont \textcolor{BrickRed}{$ \blacktriangledown\,3.9\times $}}\\[-.2pt]3m\,3.7s\\[-.2pt]{\fontsize{5.8}{6.4}\selectfont \textcolor{BrickRed}{$ \blacktriangledown\,4.1\times $}}\end{tabular} & \hspace*{1pt}\begin{tabular}[t]{@{}c@{}}\rule{0pt}{12pt}3m\,28.3s\\[-.2pt]{\fontsize{5.8}{6.4}\selectfont \textcolor{BrickRed}{$ \blacktriangledown\,1.3\times $}}\\[-.2pt]3m\,22.8s\\[-.2pt]{\fontsize{5.8}{6.4}\selectfont \textcolor{BrickRed}{$ \blacktriangledown\,1.4\times $}}\end{tabular} & \hspace*{1pt}\begin{tabular}[t]{@{}c@{}}\rule{0pt}{12pt}3m\,53.2s\\[-.2pt]{\fontsize{5.8}{6.4}\selectfont \textcolor{BrickRed}{$ \blacktriangledown\,1.4\times $}}\\[-.2pt]3m\,43.5s\\[-.2pt]{\fontsize{5.8}{6.4}\selectfont \textcolor{BrickRed}{$ \blacktriangledown\,1.5\times $}}\end{tabular} & \hspace*{1pt}\begin{tabular}[t]{@{}c@{}}\rule{0pt}{12pt}7m\,14.6s\\[-.2pt]{\fontsize{5.8}{6.4}\selectfont \textcolor{BrickRed}{$ \blacktriangledown\,1.5\times $}}\\[-.2pt]6m\,57.2s\\[-.2pt]{\fontsize{5.8}{6.4}\selectfont \textcolor{BrickRed}{$ \blacktriangledown\,1.6\times $}}\end{tabular} & \hspace*{1pt}\begin{tabular}[t]{@{}c@{}}\rule{0pt}{12pt}51m\,22s\\[-.2pt]{\fontsize{5.8}{6.4}\selectfont \textcolor{gray}{n/a}}\\[-.2pt]51m\,17.2s\\[-.2pt]{\fontsize{5.8}{6.4}\selectfont \textcolor{gray}{n/a}}\end{tabular} & \hspace*{1pt}\begin{tabular}[t]{@{}c@{}}\rule{0pt}{12pt}1m\,10.3s\\[-.2pt]{\fontsize{5.8}{6.4}\selectfont \textcolor{ForestGreen}{$ \blacktriangle\,7.6\times $}}\\[-.2pt]1m\,9s\\[-.2pt]{\fontsize{5.8}{6.4}\selectfont \textcolor{ForestGreen}{$ \blacktriangle\,2.4\times $}}\end{tabular} & \hspace*{1pt}\begin{tabular}[t]{@{}c@{}}\rule{0pt}{12pt}1m\,8.6s\\[-.2pt]{\fontsize{5.8}{6.4}\selectfont \textcolor{ForestGreen}{$ \blacktriangle\,6.4\times $}}\\[-.2pt]1m\,6.9s\\[-.2pt]{\fontsize{5.8}{6.4}\selectfont \textcolor{ForestGreen}{$ \blacktriangle\,2.3\times $}}\end{tabular} & \hspace*{1pt}\begin{tabular}[t]{@{}c@{}}\rule{0pt}{12pt}2m\,21.5s\\[-.2pt]{\fontsize{5.8}{6.4}\selectfont \textcolor{ForestGreen}{$ \blacktriangle\,6.0\times $}}\\[-.2pt]2m\,18.4s\\[-.2pt]{\fontsize{5.8}{6.4}\selectfont \textcolor{ForestGreen}{$ \blacktriangle\,2.2\times $}}\end{tabular} & \hspace*{1pt}\begin{tabular}[t]{@{}c@{}}\rule{0pt}{12pt}15m\,56.8s\\[-.2pt]{\fontsize{5.8}{6.4}\selectfont \textcolor{gray}{n/a}}\\[-.2pt]15m\,51.9s\\[-.2pt]{\fontsize{5.8}{6.4}\selectfont \textcolor{gray}{n/a}}\end{tabular}\\[3pt]\hline
\BenchmarkSourceLabel{bench:aws-4-k20-w15:row:mpf-simdref-simd-ksw2}{R44}\shortstack[l]{{\color{gray}R44}\\\textbf{BioseqDB}} & \begin{tabular}[t]{@{}l@{}}\rule{0pt}{12pt}MPF / SIMD-ref\\SIMD-DP / KSW2\end{tabular} & \hspace*{1pt}\begin{tabular}[t]{@{}c@{}}\rule{0pt}{12pt}11m\,21.3s\\[-.2pt]{\fontsize{5.8}{6.4}\selectfont \textcolor{BrickRed}{$ \blacktriangledown\,4.1\times $}}\\[-.2pt]11m\,17.7s\\[-.2pt]{\fontsize{5.8}{6.4}\selectfont \textcolor{BrickRed}{$ \blacktriangledown\,4.2\times $}}\end{tabular} & \hspace*{1pt}\begin{tabular}[t]{@{}c@{}}\rule{0pt}{12pt}24m\,2.7s\\[-.2pt]{\fontsize{5.8}{6.4}\selectfont \textcolor{BrickRed}{$ \blacktriangledown\,8.5\times $}}\\[-.2pt]23m\,59.7s\\[-.2pt]{\fontsize{5.8}{6.4}\selectfont \textcolor{BrickRed}{$ \blacktriangledown\,8.8\times $}}\end{tabular} & \hspace*{1pt}\begin{tabular}[t]{@{}c@{}}\rule{0pt}{12pt}46m\,8s\\[-.2pt]{\fontsize{5.8}{6.4}\selectfont \textcolor{BrickRed}{$ \blacktriangledown\,9.4\times $}}\\[-.2pt]46m\,1.9s\\[-.2pt]{\fontsize{5.8}{6.4}\selectfont \textcolor{BrickRed}{$ \blacktriangledown\,9.9\times $}}\end{tabular} & \hspace*{1pt}\begin{tabular}[t]{@{}c@{}}\rule{0pt}{12pt}{\fontsize{5.6}{6.4}\selectfont 1h\,35m\,48.4s}\\[-.2pt]{\fontsize{5.8}{6.4}\selectfont \textcolor{BrickRed}{$ \blacktriangledown\,3.8\times $}}\\[-.2pt]{\fontsize{5.6}{6.4}\selectfont 1h\,35m\,44.2s}\\[-.2pt]{\fontsize{5.8}{6.4}\selectfont \textcolor{BrickRed}{$ \blacktriangledown\,3.8\times $}}\end{tabular} & \hspace*{1pt}\begin{tabular}[t]{@{}c@{}}\rule{0pt}{12pt}12m\,1.2s\\[-.2pt]{\fontsize{5.8}{6.4}\selectfont \textcolor{BrickRed}{$ \blacktriangledown\,21.2\times $}}\\[-.2pt]11m\,57.5s\\[-.2pt]{\fontsize{5.8}{6.4}\selectfont \textcolor{BrickRed}{$ \blacktriangledown\,21.8\times $}}\end{tabular} & \hspace*{1pt}\begin{tabular}[t]{@{}c@{}}\rule{0pt}{12pt}35m\,19s\\[-.2pt]{\fontsize{5.8}{6.4}\selectfont \textcolor{BrickRed}{$ \blacktriangledown\,61.6\times $}}\\[-.2pt]35m\,15.4s\\[-.2pt]{\fontsize{5.8}{6.4}\selectfont \textcolor{BrickRed}{$ \blacktriangledown\,64.7\times $}}\end{tabular} & \hspace*{1pt}\begin{tabular}[t]{@{}c@{}}\rule{0pt}{12pt}46m\,59.7s\\[-.2pt]{\fontsize{5.8}{6.4}\selectfont \textcolor{BrickRed}{$ \blacktriangledown\,41.3\times $}}\\[-.2pt]46m\,53.4s\\[-.2pt]{\fontsize{5.8}{6.4}\selectfont \textcolor{BrickRed}{$ \blacktriangledown\,43.2\times $}}\end{tabular} & \textcolor{gray}{\ensuremath{\cdot}} & \hspace*{1pt}\begin{tabular}[t]{@{}c@{}}\rule{0pt}{12pt}4m\,42.6s\\[-.2pt]{\fontsize{5.8}{6.4}\selectfont \textcolor{BrickRed}{$ \blacktriangledown\,1.4\times $}}\\[-.2pt]3m\,56.3s\\[-.2pt]{\fontsize{5.8}{6.4}\selectfont \textcolor{BrickRed}{$ \blacktriangledown\,17\% $}}\end{tabular} & \hspace*{1pt}\begin{tabular}[t]{@{}c@{}}\rule{0pt}{12pt}4m\,36.4s\\[-.2pt]{\fontsize{5.8}{6.4}\selectfont \textcolor{BrickRed}{$ \blacktriangledown\,19\% $}}\\[-.2pt]3m\,58.1s\\[-.2pt]{\fontsize{5.8}{6.4}\selectfont \textcolor{BrickRed}{$ \blacktriangledown\,7\% $}}\end{tabular} & \textcolor{gray}{\ensuremath{\cdot}} & \textcolor{gray}{\ensuremath{\cdot}} & \hspace*{1pt}\begin{tabular}[t]{@{}c@{}}\rule{0pt}{12pt}11m\,3.4s\\[-.2pt]{\fontsize{5.8}{6.4}\selectfont \textcolor{BrickRed}{$ \blacktriangledown\,9.4\times $}}\\[-.2pt]4m\,43.3s\\[-.2pt]{\fontsize{5.8}{6.4}\selectfont \textcolor{BrickRed}{$ \blacktriangledown\,4.1\times $}}\end{tabular} & \hspace*{1pt}\begin{tabular}[t]{@{}c@{}}\rule{0pt}{12pt}10m\,45.6s\\[-.2pt]{\fontsize{5.8}{6.4}\selectfont \textcolor{BrickRed}{$ \blacktriangledown\,9.4\times $}}\\[-.2pt]5m\,9.2s\\[-.2pt]{\fontsize{5.8}{6.4}\selectfont \textcolor{BrickRed}{$ \blacktriangledown\,4.6\times $}}\end{tabular} & \textcolor{gray}{\ensuremath{\cdot}} & \textcolor{gray}{\ensuremath{\cdot}}\\[3pt]\hline
\BenchmarkSourceLabel{bench:aws-4-k20-w15:row:mpf-simdref-scalar-block}{R45}\shortstack[l]{{\color{gray}R45}} & \begin{tabular}[t]{@{}l@{}}\rule{0pt}{12pt}MPF / SIMD-ref\\Scalar-DP / Block\end{tabular} & \hspace*{1pt}\begin{tabular}[t]{@{}c@{}}\rule{0pt}{12pt}11m\,22.1s\\[-.2pt]{\fontsize{5.8}{6.4}\selectfont \textcolor{BrickRed}{$ \blacktriangledown\,4.1\times $}}\\[-.2pt]11m\,21.3s\\[-.2pt]{\fontsize{5.8}{6.4}\selectfont \textcolor{BrickRed}{$ \blacktriangledown\,4.2\times $}}\end{tabular} & \hspace*{1pt}\begin{tabular}[t]{@{}c@{}}\rule{0pt}{12pt}24m\,22.6s\\[-.2pt]{\fontsize{5.8}{6.4}\selectfont \textcolor{BrickRed}{$ \blacktriangledown\,8.6\times $}}\\[-.2pt]24m\,21.8s\\[-.2pt]{\fontsize{5.8}{6.4}\selectfont \textcolor{BrickRed}{$ \blacktriangledown\,9.0\times $}}\end{tabular} & \hspace*{1pt}\begin{tabular}[t]{@{}c@{}}\rule{0pt}{12pt}46m\,17.2s\\[-.2pt]{\fontsize{5.8}{6.4}\selectfont \textcolor{BrickRed}{$ \blacktriangledown\,9.4\times $}}\\[-.2pt]46m\,15.7s\\[-.2pt]{\fontsize{5.8}{6.4}\selectfont \textcolor{BrickRed}{$ \blacktriangledown\,9.9\times $}}\end{tabular} & \hspace*{1pt}\begin{tabular}[t]{@{}c@{}}\rule{0pt}{12pt}{\fontsize{5.6}{6.4}\selectfont 2h\,59m\,10.5s}\\[-.2pt]{\fontsize{5.8}{6.4}\selectfont \textcolor{BrickRed}{$ \blacktriangledown\,7.1\times $}}\\[-.2pt]{\fontsize{5.6}{6.4}\selectfont 2h\,59m\,8.2s}\\[-.2pt]{\fontsize{5.8}{6.4}\selectfont \textcolor{BrickRed}{$ \blacktriangledown\,7.2\times $}}\end{tabular} & \hspace*{1pt}\begin{tabular}[t]{@{}c@{}}\rule{0pt}{12pt}12m\,25.5s\\[-.2pt]{\fontsize{5.8}{6.4}\selectfont \textcolor{BrickRed}{$ \blacktriangledown\,22.0\times $}}\\[-.2pt]12m\,21.9s\\[-.2pt]{\fontsize{5.8}{6.4}\selectfont \textcolor{BrickRed}{$ \blacktriangledown\,22.5\times $}}\end{tabular} & \hspace*{1pt}\begin{tabular}[t]{@{}c@{}}\rule{0pt}{12pt}25m\,32.6s\\[-.2pt]{\fontsize{5.8}{6.4}\selectfont \textcolor{BrickRed}{$ \blacktriangledown\,44.5\times $}}\\[-.2pt]25m\,30s\\[-.2pt]{\fontsize{5.8}{6.4}\selectfont \textcolor{BrickRed}{$ \blacktriangledown\,46.8\times $}}\end{tabular} & \hspace*{1pt}\begin{tabular}[t]{@{}c@{}}\rule{0pt}{12pt}47m\,21.3s\\[-.2pt]{\fontsize{5.8}{6.4}\selectfont \textcolor{BrickRed}{$ \blacktriangledown\,41.7\times $}}\\[-.2pt]47m\,15.6s\\[-.2pt]{\fontsize{5.8}{6.4}\selectfont \textcolor{BrickRed}{$ \blacktriangledown\,43.5\times $}}\end{tabular} & \textcolor{gray}{\ensuremath{\cdot}} & \hspace*{1pt}\begin{tabular}[t]{@{}c@{}}\rule{0pt}{12pt}4m\,45.1s\\[-.2pt]{\fontsize{5.8}{6.4}\selectfont \textcolor{BrickRed}{$ \blacktriangledown\,1.4\times $}}\\[-.2pt]4m\,0.5s\\[-.2pt]{\fontsize{5.8}{6.4}\selectfont \textcolor{BrickRed}{$ \blacktriangledown\,19\% $}}\end{tabular} & \hspace*{1pt}\begin{tabular}[t]{@{}c@{}}\rule{0pt}{12pt}4m\,33.4s\\[-.2pt]{\fontsize{5.8}{6.4}\selectfont \textcolor{BrickRed}{$ \blacktriangledown\,17\% $}}\\[-.2pt]3m\,57.3s\\[-.2pt]{\fontsize{5.8}{6.4}\selectfont \textcolor{BrickRed}{$ \blacktriangledown\,6\% $}}\end{tabular} & \textcolor{gray}{\ensuremath{\cdot}} & \textcolor{gray}{\ensuremath{\cdot}} & \hspace*{1pt}\begin{tabular}[t]{@{}c@{}}\rule{0pt}{12pt}11m\,4.3s\\[-.2pt]{\fontsize{5.8}{6.4}\selectfont \textcolor{BrickRed}{$ \blacktriangledown\,9.5\times $}}\\[-.2pt]4m\,42.4s\\[-.2pt]{\fontsize{5.8}{6.4}\selectfont \textcolor{BrickRed}{$ \blacktriangledown\,4.1\times $}}\end{tabular} & \hspace*{1pt}\begin{tabular}[t]{@{}c@{}}\rule{0pt}{12pt}10m\,27.9s\\[-.2pt]{\fontsize{5.8}{6.4}\selectfont \textcolor{BrickRed}{$ \blacktriangledown\,9.2\times $}}\\[-.2pt]4m\,57.1s\\[-.2pt]{\fontsize{5.8}{6.4}\selectfont \textcolor{BrickRed}{$ \blacktriangledown\,4.4\times $}}\end{tabular} & \textcolor{gray}{\ensuremath{\cdot}} & \textcolor{gray}{\ensuremath{\cdot}}\\[3pt]\hline
\BenchmarkSourceLabel{bench:aws-4-k20-w15:row:mpf-scalar-rmq-block}{R46}\shortstack[l]{{\color{gray}R46}} & \begin{tabular}[t]{@{}l@{}}\rule{0pt}{12pt}MPF / Scalar\\RMQ / Block\end{tabular} & \hspace*{1pt}\begin{tabular}[t]{@{}c@{}}\rule{0pt}{12pt}11m\,22.5s\\[-.2pt]{\fontsize{5.8}{6.4}\selectfont \textcolor{BrickRed}{$ \blacktriangledown\,4.1\times $}}\\[-.2pt]11m\,22s\\[-.2pt]{\fontsize{5.8}{6.4}\selectfont \textcolor{BrickRed}{$ \blacktriangledown\,4.2\times $}}\end{tabular} & \hspace*{1pt}\begin{tabular}[t]{@{}c@{}}\rule{0pt}{12pt}24m\,52.9s\\[-.2pt]{\fontsize{5.8}{6.4}\selectfont \textcolor{BrickRed}{$ \blacktriangledown\,8.8\times $}}\\[-.2pt]24m\,52.4s\\[-.2pt]{\fontsize{5.8}{6.4}\selectfont \textcolor{BrickRed}{$ \blacktriangledown\,9.2\times $}}\end{tabular} & \hspace*{1pt}\begin{tabular}[t]{@{}c@{}}\rule{0pt}{12pt}46m\,1.7s\\[-.2pt]{\fontsize{5.8}{6.4}\selectfont \textcolor{BrickRed}{$ \blacktriangledown\,9.4\times $}}\\[-.2pt]46m\,0.4s\\[-.2pt]{\fontsize{5.8}{6.4}\selectfont \textcolor{BrickRed}{$ \blacktriangledown\,9.9\times $}}\end{tabular} & \hspace*{1pt}\begin{tabular}[t]{@{}c@{}}\rule{0pt}{12pt}{\fontsize{5.6}{6.4}\selectfont 2h\,57m\,27.5s}\\[-.2pt]{\fontsize{5.8}{6.4}\selectfont \textcolor{BrickRed}{$ \blacktriangledown\,7.1\times $}}\\[-.2pt]{\fontsize{5.6}{6.4}\selectfont 2h\,57m\,25.3s}\\[-.2pt]{\fontsize{5.8}{6.4}\selectfont \textcolor{BrickRed}{$ \blacktriangledown\,7.1\times $}}\end{tabular} & \hspace*{1pt}\begin{tabular}[t]{@{}c@{}}\rule{0pt}{12pt}14m\,45.6s\\[-.2pt]{\fontsize{5.8}{6.4}\selectfont \textcolor{BrickRed}{$ \blacktriangledown\,26.1\times $}}\\[-.2pt]14m\,42s\\[-.2pt]{\fontsize{5.8}{6.4}\selectfont \textcolor{BrickRed}{$ \blacktriangledown\,26.8\times $}}\end{tabular} & \hspace*{1pt}\begin{tabular}[t]{@{}c@{}}\rule{0pt}{12pt}31m\,31s\\[-.2pt]{\fontsize{5.8}{6.4}\selectfont \textcolor{BrickRed}{$ \blacktriangledown\,54.9\times $}}\\[-.2pt]31m\,28.4s\\[-.2pt]{\fontsize{5.8}{6.4}\selectfont \textcolor{BrickRed}{$ \blacktriangledown\,57.7\times $}}\end{tabular} & \hspace*{1pt}\begin{tabular}[t]{@{}c@{}}\rule{0pt}{12pt}{\fontsize{5.6}{6.4}\selectfont 1h\,7m\,40.8s}\\[-.2pt]{\fontsize{5.8}{6.4}\selectfont \textcolor{BrickRed}{$ \blacktriangledown\,59.5\times $}}\\[-.2pt]{\fontsize{5.6}{6.4}\selectfont 1h\,7m\,35.3s}\\[-.2pt]{\fontsize{5.8}{6.4}\selectfont \textcolor{BrickRed}{$ \blacktriangledown\,62.3\times $}}\end{tabular} & \hspace*{1pt}\begin{tabular}[t]{@{}c@{}}\rule{0pt}{12pt}{\fontsize{5.6}{6.4}\selectfont 1h\,26m\,25.6s}\\[-.2pt]{\fontsize{5.8}{6.4}\selectfont \textcolor{BrickRed}{$ \blacktriangledown\,27.3\times $}}\\[-.2pt]{\fontsize{5.6}{6.4}\selectfont 1h\,26m\,20.7s}\\[-.2pt]{\fontsize{5.8}{6.4}\selectfont \textcolor{BrickRed}{$ \blacktriangledown\,28.2\times $}}\end{tabular} & \hspace*{1pt}\begin{tabular}[t]{@{}c@{}}\rule{0pt}{12pt}5m\,16.6s\\[-.2pt]{\fontsize{5.8}{6.4}\selectfont \textcolor{BrickRed}{$ \blacktriangledown\,1.5\times $}}\\[-.2pt]4m\,33.9s\\[-.2pt]{\fontsize{5.8}{6.4}\selectfont \textcolor{BrickRed}{$ \blacktriangledown\,1.4\times $}}\end{tabular} & \hspace*{1pt}\begin{tabular}[t]{@{}c@{}}\rule{0pt}{12pt}4m\,53.5s\\[-.2pt]{\fontsize{5.8}{6.4}\selectfont \textcolor{BrickRed}{$ \blacktriangledown\,26\% $}}\\[-.2pt]4m\,17.9s\\[-.2pt]{\fontsize{5.8}{6.4}\selectfont \textcolor{BrickRed}{$ \blacktriangledown\,15\% $}}\end{tabular} & \textcolor{gray}{\ensuremath{\cdot}} & \textcolor{gray}{\ensuremath{\cdot}} & \hspace*{1pt}\begin{tabular}[t]{@{}c@{}}\rule{0pt}{12pt}12m\,47.3s\\[-.2pt]{\fontsize{5.8}{6.4}\selectfont \textcolor{BrickRed}{$ \blacktriangledown\,10.9\times $}}\\[-.2pt]4m\,55.6s\\[-.2pt]{\fontsize{5.8}{6.4}\selectfont \textcolor{BrickRed}{$ \blacktriangledown\,4.3\times $}}\end{tabular} & \hspace*{1pt}\begin{tabular}[t]{@{}c@{}}\rule{0pt}{12pt}12m\,11.3s\\[-.2pt]{\fontsize{5.8}{6.4}\selectfont \textcolor{BrickRed}{$ \blacktriangledown\,10.7\times $}}\\[-.2pt]5m\,42.7s\\[-.2pt]{\fontsize{5.8}{6.4}\selectfont \textcolor{BrickRed}{$ \blacktriangledown\,5.1\times $}}\end{tabular} & \textcolor{gray}{\ensuremath{\cdot}} & \textcolor{gray}{\ensuremath{\cdot}}\\[3pt]\hline
\BenchmarkSourceLabel{bench:aws-4-k20-w15:row:mpf-simdref-scalar-ksw2}{R47}\shortstack[l]{{\color{gray}R47}} & \begin{tabular}[t]{@{}l@{}}\rule{0pt}{12pt}MPF / SIMD-ref\\Scalar-DP / KSW2\end{tabular} & \hspace*{1pt}\begin{tabular}[t]{@{}c@{}}\rule{0pt}{12pt}11m\,28.4s\\[-.2pt]{\fontsize{5.8}{6.4}\selectfont \textcolor{BrickRed}{$ \blacktriangledown\,4.1\times $}}\\[-.2pt]11m\,24.7s\\[-.2pt]{\fontsize{5.8}{6.4}\selectfont \textcolor{BrickRed}{$ \blacktriangledown\,4.2\times $}}\end{tabular} & \hspace*{1pt}\begin{tabular}[t]{@{}c@{}}\rule{0pt}{12pt}24m\,37.6s\\[-.2pt]{\fontsize{5.8}{6.4}\selectfont \textcolor{BrickRed}{$ \blacktriangledown\,8.7\times $}}\\[-.2pt]24m\,34.9s\\[-.2pt]{\fontsize{5.8}{6.4}\selectfont \textcolor{BrickRed}{$ \blacktriangledown\,9.0\times $}}\end{tabular} & \hspace*{1pt}\begin{tabular}[t]{@{}c@{}}\rule{0pt}{12pt}46m\,9.8s\\[-.2pt]{\fontsize{5.8}{6.4}\selectfont \textcolor{BrickRed}{$ \blacktriangledown\,9.4\times $}}\\[-.2pt]46m\,3.7s\\[-.2pt]{\fontsize{5.8}{6.4}\selectfont \textcolor{BrickRed}{$ \blacktriangledown\,9.9\times $}}\end{tabular} & \textcolor{gray}{\ensuremath{\cdot}} & \hspace*{1pt}\begin{tabular}[t]{@{}c@{}}\rule{0pt}{12pt}12m\,9.4s\\[-.2pt]{\fontsize{5.8}{6.4}\selectfont \textcolor{BrickRed}{$ \blacktriangledown\,21.5\times $}}\\[-.2pt]12m\,5.8s\\[-.2pt]{\fontsize{5.8}{6.4}\selectfont \textcolor{BrickRed}{$ \blacktriangledown\,22.0\times $}}\end{tabular} & \hspace*{1pt}\begin{tabular}[t]{@{}c@{}}\rule{0pt}{12pt}25m\,52.5s\\[-.2pt]{\fontsize{5.8}{6.4}\selectfont \textcolor{BrickRed}{$ \blacktriangledown\,45.1\times $}}\\[-.2pt]25m\,49.6s\\[-.2pt]{\fontsize{5.8}{6.4}\selectfont \textcolor{BrickRed}{$ \blacktriangledown\,47.4\times $}}\end{tabular} & \hspace*{1pt}\begin{tabular}[t]{@{}c@{}}\rule{0pt}{12pt}47m\,13.9s\\[-.2pt]{\fontsize{5.8}{6.4}\selectfont \textcolor{BrickRed}{$ \blacktriangledown\,41.6\times $}}\\[-.2pt]47m\,7.6s\\[-.2pt]{\fontsize{5.8}{6.4}\selectfont \textcolor{BrickRed}{$ \blacktriangledown\,43.4\times $}}\end{tabular} & \textcolor{gray}{\ensuremath{\cdot}} & \hspace*{1pt}\begin{tabular}[t]{@{}c@{}}\rule{0pt}{12pt}4m\,48.9s\\[-.2pt]{\fontsize{5.8}{6.4}\selectfont \textcolor{BrickRed}{$ \blacktriangledown\,1.4\times $}}\\[-.2pt]4m\,2.5s\\[-.2pt]{\fontsize{5.8}{6.4}\selectfont \textcolor{BrickRed}{$ \blacktriangledown\,20\% $}}\end{tabular} & \hspace*{1pt}\begin{tabular}[t]{@{}c@{}}\rule{0pt}{12pt}4m\,34s\\[-.2pt]{\fontsize{5.8}{6.4}\selectfont \textcolor{BrickRed}{$ \blacktriangledown\,17\% $}}\\[-.2pt]3m\,55.9s\\[-.2pt]{\fontsize{5.8}{6.4}\selectfont \textcolor{BrickRed}{$ \blacktriangledown\,6\% $}}\end{tabular} & \textcolor{gray}{\ensuremath{\cdot}} & \textcolor{gray}{\ensuremath{\cdot}} & \hspace*{1pt}\begin{tabular}[t]{@{}c@{}}\rule{0pt}{12pt}10m\,57s\\[-.2pt]{\fontsize{5.8}{6.4}\selectfont \textcolor{BrickRed}{$ \blacktriangledown\,9.4\times $}}\\[-.2pt]4m\,39.3s\\[-.2pt]{\fontsize{5.8}{6.4}\selectfont \textcolor{BrickRed}{$ \blacktriangledown\,4.0\times $}}\end{tabular} & \hspace*{1pt}\begin{tabular}[t]{@{}c@{}}\rule{0pt}{12pt}10m\,35.2s\\[-.2pt]{\fontsize{5.8}{6.4}\selectfont \textcolor{BrickRed}{$ \blacktriangledown\,9.3\times $}}\\[-.2pt]5m\,2s\\[-.2pt]{\fontsize{5.8}{6.4}\selectfont \textcolor{BrickRed}{$ \blacktriangledown\,4.5\times $}}\end{tabular} & \textcolor{gray}{\ensuremath{\cdot}} & \textcolor{gray}{\ensuremath{\cdot}}\\[3pt]\hline
\BenchmarkSourceLabel{bench:aws-4-k20-w15:row:mpf-scalar-simd-block}{R48}\shortstack[l]{{\color{gray}R48}} & \begin{tabular}[t]{@{}l@{}}\rule{0pt}{12pt}MPF / Scalar\\SIMD-DP / Block\end{tabular} & \hspace*{1pt}\begin{tabular}[t]{@{}c@{}}\rule{0pt}{12pt}11m\,30.6s\\[-.2pt]{\fontsize{5.8}{6.4}\selectfont \textcolor{BrickRed}{$ \blacktriangledown\,4.1\times $}}\\[-.2pt]11m\,29.9s\\[-.2pt]{\fontsize{5.8}{6.4}\selectfont \textcolor{BrickRed}{$ \blacktriangledown\,4.3\times $}}\end{tabular} & \hspace*{1pt}\begin{tabular}[t]{@{}c@{}}\rule{0pt}{12pt}25m\,2.6s\\[-.2pt]{\fontsize{5.8}{6.4}\selectfont \textcolor{BrickRed}{$ \blacktriangledown\,8.8\times $}}\\[-.2pt]25m\,1.8s\\[-.2pt]{\fontsize{5.8}{6.4}\selectfont \textcolor{BrickRed}{$ \blacktriangledown\,9.2\times $}}\end{tabular} & \hspace*{1pt}\begin{tabular}[t]{@{}c@{}}\rule{0pt}{12pt}45m\,36s\\[-.2pt]{\fontsize{5.8}{6.4}\selectfont \textcolor{BrickRed}{$ \blacktriangledown\,9.3\times $}}\\[-.2pt]45m\,34.7s\\[-.2pt]{\fontsize{5.8}{6.4}\selectfont \textcolor{BrickRed}{$ \blacktriangledown\,9.8\times $}}\end{tabular} & \hspace*{1pt}\begin{tabular}[t]{@{}c@{}}\rule{0pt}{12pt}{\fontsize{5.6}{6.4}\selectfont 1h\,36m\,16.5s}\\[-.2pt]{\fontsize{5.8}{6.4}\selectfont \textcolor{BrickRed}{$ \blacktriangledown\,3.8\times $}}\\[-.2pt]{\fontsize{5.6}{6.4}\selectfont 1h\,36m\,14.3s}\\[-.2pt]{\fontsize{5.8}{6.4}\selectfont \textcolor{BrickRed}{$ \blacktriangledown\,3.8\times $}}\end{tabular} & \hspace*{1pt}\begin{tabular}[t]{@{}c@{}}\rule{0pt}{12pt}15m\,49.6s\\[-.2pt]{\fontsize{5.8}{6.4}\selectfont \textcolor{BrickRed}{$ \blacktriangledown\,28.0\times $}}\\[-.2pt]15m\,45.8s\\[-.2pt]{\fontsize{5.8}{6.4}\selectfont \textcolor{BrickRed}{$ \blacktriangledown\,28.7\times $}}\end{tabular} & \hspace*{1pt}\begin{tabular}[t]{@{}c@{}}\rule{0pt}{12pt}32m\,50.4s\\[-.2pt]{\fontsize{5.8}{6.4}\selectfont \textcolor{BrickRed}{$ \blacktriangledown\,57.2\times $}}\\[-.2pt]32m\,47.7s\\[-.2pt]{\fontsize{5.8}{6.4}\selectfont \textcolor{BrickRed}{$ \blacktriangledown\,60.2\times $}}\end{tabular} & \hspace*{1pt}\begin{tabular}[t]{@{}c@{}}\rule{0pt}{12pt}{\fontsize{5.6}{6.4}\selectfont 1h\,6m\,46.8s}\\[-.2pt]{\fontsize{5.8}{6.4}\selectfont \textcolor{BrickRed}{$ \blacktriangledown\,58.8\times $}}\\[-.2pt]{\fontsize{5.6}{6.4}\selectfont 1h\,6m\,41.3s}\\[-.2pt]{\fontsize{5.8}{6.4}\selectfont \textcolor{BrickRed}{$ \blacktriangledown\,61.4\times $}}\end{tabular} & \hspace*{1pt}\begin{tabular}[t]{@{}c@{}}\rule{0pt}{12pt}{\fontsize{5.6}{6.4}\selectfont 1h\,30m\,32.9s}\\[-.2pt]{\fontsize{5.8}{6.4}\selectfont \textcolor{BrickRed}{$ \blacktriangledown\,28.6\times $}}\\[-.2pt]{\fontsize{5.6}{6.4}\selectfont 1h\,30m\,28.4s}\\[-.2pt]{\fontsize{5.8}{6.4}\selectfont \textcolor{BrickRed}{$ \blacktriangledown\,29.6\times $}}\end{tabular} & \hspace*{1pt}\begin{tabular}[t]{@{}c@{}}\rule{0pt}{12pt}4m\,51s\\[-.2pt]{\fontsize{5.8}{6.4}\selectfont \textcolor{BrickRed}{$ \blacktriangledown\,1.4\times $}}\\[-.2pt]4m\,8.2s\\[-.2pt]{\fontsize{5.8}{6.4}\selectfont \textcolor{BrickRed}{$ \blacktriangledown\,22\% $}}\end{tabular} & \hspace*{1pt}\begin{tabular}[t]{@{}c@{}}\rule{0pt}{12pt}4m\,40s\\[-.2pt]{\fontsize{5.8}{6.4}\selectfont \textcolor{BrickRed}{$ \blacktriangledown\,20\% $}}\\[-.2pt]4m\,5.7s\\[-.2pt]{\fontsize{5.8}{6.4}\selectfont \textcolor{BrickRed}{$ \blacktriangledown\,10\% $}}\end{tabular} & \textcolor{gray}{\ensuremath{\cdot}} & \textcolor{gray}{\ensuremath{\cdot}} & \hspace*{1pt}\begin{tabular}[t]{@{}c@{}}\rule{0pt}{12pt}12m\,25.7s\\[-.2pt]{\fontsize{5.8}{6.4}\selectfont \textcolor{BrickRed}{$ \blacktriangledown\,10.6\times $}}\\[-.2pt]4m\,36.4s\\[-.2pt]{\fontsize{5.8}{6.4}\selectfont \textcolor{BrickRed}{$ \blacktriangledown\,4.0\times $}}\end{tabular} & \hspace*{1pt}\begin{tabular}[t]{@{}c@{}}\rule{0pt}{12pt}11m\,44.4s\\[-.2pt]{\fontsize{5.8}{6.4}\selectfont \textcolor{BrickRed}{$ \blacktriangledown\,10.3\times $}}\\[-.2pt]5m\,24.7s\\[-.2pt]{\fontsize{5.8}{6.4}\selectfont \textcolor{BrickRed}{$ \blacktriangledown\,4.9\times $}}\end{tabular} & \textcolor{gray}{\ensuremath{\cdot}} & \textcolor{gray}{\ensuremath{\cdot}}\\[3pt]\hline
\BenchmarkSourceLabel{bench:aws-4-k20-w15:row:mpf-simdref-simd-block}{R49}\shortstack[l]{{\color{gray}R49}} & \begin{tabular}[t]{@{}l@{}}\rule{0pt}{12pt}MPF / SIMD-ref\\SIMD-DP / Block\end{tabular} & \hspace*{1pt}\begin{tabular}[t]{@{}c@{}}\rule{0pt}{12pt}12m\,1.2s\\[-.2pt]{\fontsize{5.8}{6.4}\selectfont \textcolor{BrickRed}{$ \blacktriangledown\,4.3\times $}}\\[-.2pt]12m\,0.4s\\[-.2pt]{\fontsize{5.8}{6.4}\selectfont \textcolor{BrickRed}{$ \blacktriangledown\,4.5\times $}}\end{tabular} & \hspace*{1pt}\begin{tabular}[t]{@{}c@{}}\rule{0pt}{12pt}24m\,25.1s\\[-.2pt]{\fontsize{5.8}{6.4}\selectfont \textcolor{BrickRed}{$ \blacktriangledown\,8.6\times $}}\\[-.2pt]24m\,24.2s\\[-.2pt]{\fontsize{5.8}{6.4}\selectfont \textcolor{BrickRed}{$ \blacktriangledown\,9.0\times $}}\end{tabular} & \hspace*{1pt}\begin{tabular}[t]{@{}c@{}}\rule{0pt}{12pt}46m\,40.1s\\[-.2pt]{\fontsize{5.8}{6.4}\selectfont \textcolor{BrickRed}{$ \blacktriangledown\,9.5\times $}}\\[-.2pt]46m\,38.8s\\[-.2pt]{\fontsize{5.8}{6.4}\selectfont \textcolor{BrickRed}{$ \blacktriangledown\,10.0\times $}}\end{tabular} & \hspace*{1pt}\begin{tabular}[t]{@{}c@{}}\rule{0pt}{12pt}{\fontsize{5.6}{6.4}\selectfont 1h\,35m\,35.8s}\\[-.2pt]{\fontsize{5.8}{6.4}\selectfont \textcolor{BrickRed}{$ \blacktriangledown\,3.8\times $}}\\[-.2pt]{\fontsize{5.6}{6.4}\selectfont 1h\,35m\,33.7s}\\[-.2pt]{\fontsize{5.8}{6.4}\selectfont \textcolor{BrickRed}{$ \blacktriangledown\,3.8\times $}}\end{tabular} & \hspace*{1pt}\begin{tabular}[t]{@{}c@{}}\rule{0pt}{12pt}12m\,28.9s\\[-.2pt]{\fontsize{5.8}{6.4}\selectfont \textcolor{BrickRed}{$ \blacktriangledown\,22.1\times $}}\\[-.2pt]12m\,25.2s\\[-.2pt]{\fontsize{5.8}{6.4}\selectfont \textcolor{BrickRed}{$ \blacktriangledown\,22.6\times $}}\end{tabular} & \hspace*{1pt}\begin{tabular}[t]{@{}c@{}}\rule{0pt}{12pt}34m\,56.4s\\[-.2pt]{\fontsize{5.8}{6.4}\selectfont \textcolor{BrickRed}{$ \blacktriangledown\,60.9\times $}}\\[-.2pt]34m\,53.2s\\[-.2pt]{\fontsize{5.8}{6.4}\selectfont \textcolor{BrickRed}{$ \blacktriangledown\,64.0\times $}}\end{tabular} & \hspace*{1pt}\begin{tabular}[t]{@{}c@{}}\rule{0pt}{12pt}49m\,6.1s\\[-.2pt]{\fontsize{5.8}{6.4}\selectfont \textcolor{BrickRed}{$ \blacktriangledown\,43.2\times $}}\\[-.2pt]49m\,0.3s\\[-.2pt]{\fontsize{5.8}{6.4}\selectfont \textcolor{BrickRed}{$ \blacktriangledown\,45.2\times $}}\end{tabular} & \textcolor{gray}{\ensuremath{\cdot}} & \hspace*{1pt}\begin{tabular}[t]{@{}c@{}}\rule{0pt}{12pt}4m\,45.5s\\[-.2pt]{\fontsize{5.8}{6.4}\selectfont \textcolor{BrickRed}{$ \blacktriangledown\,1.4\times $}}\\[-.2pt]4m\,1.2s\\[-.2pt]{\fontsize{5.8}{6.4}\selectfont \textcolor{BrickRed}{$ \blacktriangledown\,19\% $}}\end{tabular} & \hspace*{1pt}\begin{tabular}[t]{@{}c@{}}\rule{0pt}{12pt}4m\,32.2s\\[-.2pt]{\fontsize{5.8}{6.4}\selectfont \textcolor{BrickRed}{$ \blacktriangledown\,17\% $}}\\[-.2pt]3m\,55.6s\\[-.2pt]{\fontsize{5.8}{6.4}\selectfont \textcolor{BrickRed}{$ \blacktriangledown\,5\% $}}\end{tabular} & \textcolor{gray}{\ensuremath{\cdot}} & \textcolor{gray}{\ensuremath{\cdot}} & \hspace*{1pt}\begin{tabular}[t]{@{}c@{}}\rule{0pt}{12pt}10m\,51.7s\\[-.2pt]{\fontsize{5.8}{6.4}\selectfont \textcolor{BrickRed}{$ \blacktriangledown\,9.3\times $}}\\[-.2pt]4m\,28.8s\\[-.2pt]{\fontsize{5.8}{6.4}\selectfont \textcolor{BrickRed}{$ \blacktriangledown\,3.9\times $}}\end{tabular} & \hspace*{1pt}\begin{tabular}[t]{@{}c@{}}\rule{0pt}{12pt}10m\,57.3s\\[-.2pt]{\fontsize{5.8}{6.4}\selectfont \textcolor{BrickRed}{$ \blacktriangledown\,9.6\times $}}\\[-.2pt]5m\,24.8s\\[-.2pt]{\fontsize{5.8}{6.4}\selectfont \textcolor{BrickRed}{$ \blacktriangledown\,4.9\times $}}\end{tabular} & \textcolor{gray}{\ensuremath{\cdot}} & \textcolor{gray}{\ensuremath{\cdot}}\\[3pt]\hline
\BenchmarkSourceLabel{bench:aws-4-k20-w15:row:vulkan-unavailable}{R50}\shortstack[l]{{\color{gray}R50}\\\rule{0pt}{12pt}\textbf{BioseqDB}} & Vulkan (all modes) & \textcolor{gray}{---\textsuperscript{g}} & \textcolor{gray}{---\textsuperscript{g}} & \textcolor{gray}{---\textsuperscript{g}} & \textcolor{gray}{---\textsuperscript{g}} & \textcolor{gray}{---\textsuperscript{g}} & \textcolor{gray}{---\textsuperscript{g}} & \textcolor{gray}{---\textsuperscript{g}} & \textcolor{gray}{---\textsuperscript{g}} & \textcolor{gray}{---\textsuperscript{g}} & \textcolor{gray}{---\textsuperscript{g}} & \textcolor{gray}{---\textsuperscript{g}} & \textcolor{gray}{---\textsuperscript{g}} & \textcolor{gray}{---\textsuperscript{g}} & \textcolor{gray}{---\textsuperscript{g}} & \textcolor{gray}{---\textsuperscript{g}} & \textcolor{gray}{---\textsuperscript{g}}\\[3pt]\hline
\end{tabular}\par}
\smallskip{\footnotesize\noindent Each cell: A and B times, $\cdot$: no measurement. ---\textsuperscript{g}: GPU unavailable. n/a: no comparison.\par
\noindent Blue: one lowest A and one lowest B measurement for each column for each configuration.\par
\noindent Gray: minimap2. Yellow: my preferred MPF / Bounded / SIMD-DP / Block configuration.\par
\noindent MPF/MPS: minimizers per fragment/sequence. SIMD-ref: naive SIMD scan. Bounded: bounded SIMD algorithm.\par
}
\clearpage
\refstepcounter{table}\label{tab:bench:aws-2-k15-w10}
\noindent{\large\bfseries Table \thetable: (D): AWS: 2 Gbp, k=15, w=10\hfill 1/7}\par\smallskip
{\fontsize{7.1}{8.1}\selectfont\setlength{\tabcolsep}{1pt}\setlength{\fboxsep}{2pt}\renewcommand{\arraystretch}{1.12}
\noindent\begin{tabular}{@{}p{14mm}p{27mm}|p{13.55mm}|p{13.55mm}|p{13.55mm}|p{13.55mm}|p{13.55mm}|p{13.55mm}|p{13.55mm}|p{13.55mm}|p{13.55mm}|p{13.55mm}|p{13.55mm}|p{13.55mm}|p{13.55mm}|p{13.55mm}|p{13.55mm}|p{13.55mm}|@{}}\hline
Engine & Execution & \multicolumn{8}{c|}{\textbf{Index}} & \multicolumn{8}{c|}{\textbf{No index}}\\
 &  & \multicolumn{8}{c|}{\textbf{2 Gbp}} & \multicolumn{8}{c|}{\textbf{2 Gbp}}\\
 &  & \multicolumn{8}{c|}{\textbf{k=15, w=10}} & \multicolumn{8}{c|}{\textbf{k=15, w=10}}\\
 &  & \multicolumn{4}{c|}{\textbf{DNA4}} & \multicolumn{4}{c|}{\textbf{DNA5}} & \multicolumn{4}{c|}{\textbf{DNA4}} & \multicolumn{4}{c|}{\textbf{DNA5}}\\
 &  & \multicolumn{1}{c|}{\textbf{M7A}} & \multicolumn{1}{c|}{\textbf{M7I}} & \multicolumn{1}{c|}{\textbf{M4}} & \multicolumn{1}{c|}{\textbf{C7G}} & \multicolumn{1}{c|}{\textbf{M7A}} & \multicolumn{1}{c|}{\textbf{M7I}} & \multicolumn{1}{c|}{\textbf{M4}} & \multicolumn{1}{c|}{\textbf{C7G}} & \multicolumn{1}{c|}{\textbf{M7A}} & \multicolumn{1}{c|}{\textbf{M7I}} & \multicolumn{1}{c|}{\textbf{M4}} & \multicolumn{1}{c|}{\textbf{C7G}} & \multicolumn{1}{c|}{\textbf{M7A}} & \multicolumn{1}{c|}{\textbf{M7I}} & \multicolumn{1}{c|}{\textbf{M4}} & \multicolumn{1}{c|}{\textbf{C7G}}\\
 & & \BenchmarkSourceLabel{bench:aws-2-k15-w10:column:index-2-k15-w10-dna4-m7a}{C1}{\color{gray}C1} & \BenchmarkSourceLabel{bench:aws-2-k15-w10:column:index-2-k15-w10-dna4-m7i}{C2}{\color{gray}C2} & \BenchmarkSourceLabel{bench:aws-2-k15-w10:column:index-2-k15-w10-dna4-m4}{C3}{\color{gray}C3} & \BenchmarkSourceLabel{bench:aws-2-k15-w10:column:index-2-k15-w10-dna4-c7g}{C4}{\color{gray}C4} & \BenchmarkSourceLabel{bench:aws-2-k15-w10:column:index-2-k15-w10-dna5-m7a}{C5}{\color{gray}C5} & \BenchmarkSourceLabel{bench:aws-2-k15-w10:column:index-2-k15-w10-dna5-m7i}{C6}{\color{gray}C6} & \BenchmarkSourceLabel{bench:aws-2-k15-w10:column:index-2-k15-w10-dna5-m4}{C7}{\color{gray}C7} & \BenchmarkSourceLabel{bench:aws-2-k15-w10:column:index-2-k15-w10-dna5-c7g}{C8}{\color{gray}C8} & \BenchmarkSourceLabel{bench:aws-2-k15-w10:column:noindex-2-k15-w10-dna4-m7a}{C9}{\color{gray}C9} & \BenchmarkSourceLabel{bench:aws-2-k15-w10:column:noindex-2-k15-w10-dna4-m7i}{C10}{\color{gray}C10} & \BenchmarkSourceLabel{bench:aws-2-k15-w10:column:noindex-2-k15-w10-dna4-m4}{C11}{\color{gray}C11} & \BenchmarkSourceLabel{bench:aws-2-k15-w10:column:noindex-2-k15-w10-dna4-c7g}{C12}{\color{gray}C12} & \BenchmarkSourceLabel{bench:aws-2-k15-w10:column:noindex-2-k15-w10-dna5-m7a}{C13}{\color{gray}C13} & \BenchmarkSourceLabel{bench:aws-2-k15-w10:column:noindex-2-k15-w10-dna5-m7i}{C14}{\color{gray}C14} & \BenchmarkSourceLabel{bench:aws-2-k15-w10:column:noindex-2-k15-w10-dna5-m4}{C15}{\color{gray}C15} & \BenchmarkSourceLabel{bench:aws-2-k15-w10:column:noindex-2-k15-w10-dna5-c7g}{C16}{\color{gray}C16}\\
\hline
\rowcolor{BenchBaseline}[0.5pt][0.5pt]\BenchmarkSourceLabel{bench:aws-2-k15-w10:row:minimap2}{R1}\shortstack[l]{{\color{gray}R1}\\\textbf{minimap2}} & \textbf{Baseline} & \hspace*{1pt}\begin{tabular}[t]{@{}c@{}}\rule{0pt}{12pt}11.2s\\[-.2pt]{\fontsize{5.8}{6.4}\selectfont \textcolor{BrickRed}{$ \blacktriangledown\,3.1\times $}}\\[-.2pt]10s\\[-.2pt]{\fontsize{5.8}{6.4}\selectfont \textcolor{BrickRed}{$ \blacktriangledown\,3.6\times $}}\end{tabular} & \hspace*{1pt}\begin{tabular}[t]{@{}c@{}}\rule{0pt}{12pt}10.7s\\[-.2pt]{\fontsize{5.8}{6.4}\selectfont \textcolor{BrickRed}{$ \blacktriangledown\,3.2\times $}}\\[-.2pt]8.6s\\[-.2pt]{\fontsize{5.8}{6.4}\selectfont \textcolor{BrickRed}{$ \blacktriangledown\,3.5\times $}}\end{tabular} & \hspace*{1pt}\begin{tabular}[t]{@{}c@{}}\rule{0pt}{12pt}22.5s\\[-.2pt]{\fontsize{5.8}{6.4}\selectfont \textcolor{BrickRed}{$ \blacktriangledown\,3.6\times $}}\\[-.2pt]18.8s\\[-.2pt]{\fontsize{5.8}{6.4}\selectfont \textcolor{BrickRed}{$ \blacktriangledown\,4.4\times $}}\end{tabular} & \hspace*{1pt}\begin{tabular}[t]{@{}c@{}}\rule{0pt}{12pt}1m\,13.1s\\[-.2pt]{\fontsize{5.8}{6.4}\selectfont \textcolor{BrickRed}{$ \blacktriangledown\,22.5\times $}}\\[-.2pt]1m\,5.7s\\[-.2pt]{\fontsize{5.8}{6.4}\selectfont \textcolor{BrickRed}{$ \blacktriangledown\,27.7\times $}}\end{tabular} & \hspace*{1pt}\begin{tabular}[t]{@{}c@{}}\rule{0pt}{12pt}9.9s\\[-.2pt]{\fontsize{5.8}{6.4}\selectfont \textcolor{BrickRed}{$ \blacktriangledown\,21\% $}}\\[-.2pt]8.8s\\[-.2pt]{\fontsize{5.8}{6.4}\selectfont \textcolor{BrickRed}{$ \blacktriangledown\,1.6\times $}}\end{tabular} & \hspace*{1pt}\begin{tabular}[t]{@{}c@{}}\rule{0pt}{12pt}9.5s\\[-.2pt]{\fontsize{5.8}{6.4}\selectfont \textcolor{BrickRed}{$ \blacktriangledown\,1.4\times $}}\\[-.2pt]7.6s\\[-.2pt]{\fontsize{5.8}{6.4}\selectfont \textcolor{BrickRed}{$ \blacktriangledown\,1.6\times $}}\end{tabular} & \hspace*{1pt}\begin{tabular}[t]{@{}c@{}}\rule{0pt}{12pt}16.5s\\[-.2pt]{\fontsize{5.8}{6.4}\selectfont \textcolor{BrickRed}{$ \blacktriangledown\,1.7\times $}}\\[-.2pt]12.9s\\[-.2pt]{\fontsize{5.8}{6.4}\selectfont \textcolor{BrickRed}{$ \blacktriangledown\,2.0\times $}}\end{tabular} & \hspace*{1pt}\begin{tabular}[t]{@{}c@{}}\rule{0pt}{12pt}28.2s\\[-.2pt]{\fontsize{5.8}{6.4}\selectfont \textcolor{BrickRed}{$ \blacktriangledown\,2.3\times $}}\\[-.2pt]26.2s\\[-.2pt]{\fontsize{5.8}{6.4}\selectfont \textcolor{BrickRed}{$ \blacktriangledown\,3.0\times $}}\end{tabular} & \hspace*{1pt}\begin{tabular}[t]{@{}c@{}}\rule{0pt}{12pt}\colorbox{BenchBest}{\strut 53s}\\[-.2pt]{\fontsize{5.8}{6.4}\selectfont \textcolor{ForestGreen}{$ \blacktriangle\,2.2\times $}}\\[-.2pt]\colorbox{BenchBest}{\strut 51.9s}\\[-.2pt]{\fontsize{5.8}{6.4}\selectfont \textcolor{ForestGreen}{$ \blacktriangle\,2.1\times $}}\end{tabular} & \hspace*{1pt}\begin{tabular}[t]{@{}c@{}}\rule{0pt}{12pt}\colorbox{BenchBest}{\strut 56.6s}\\[-.2pt]{\fontsize{5.8}{6.4}\selectfont \textcolor{ForestGreen}{$ \blacktriangle\,2.1\times $}}\\[-.2pt]\colorbox{BenchBest}{\strut 54.6s}\\[-.2pt]{\fontsize{5.8}{6.4}\selectfont \textcolor{ForestGreen}{$ \blacktriangle\,1.9\times $}}\end{tabular} & \hspace*{1pt}\begin{tabular}[t]{@{}c@{}}\rule{0pt}{12pt}\colorbox{BenchBest}{\strut 1m\,55.5s}\\[-.2pt]{\fontsize{5.8}{6.4}\selectfont \textcolor{ForestGreen}{$ \blacktriangle\,1.8\times $}}\\[-.2pt]\colorbox{BenchBest}{\strut 1m\,51.8s}\\[-.2pt]{\fontsize{5.8}{6.4}\selectfont \textcolor{ForestGreen}{$ \blacktriangle\,1.5\times $}}\end{tabular} & \hspace*{1pt}\begin{tabular}[t]{@{}c@{}}\rule{0pt}{12pt}6m\,1.3s\\[-.2pt]{\fontsize{5.8}{6.4}\selectfont \textcolor{BrickRed}{$ \blacktriangledown\,3.0\times $}}\\[-.2pt]5m\,55.2s\\[-.2pt]{\fontsize{5.8}{6.4}\selectfont \textcolor{BrickRed}{$ \blacktriangledown\,3.1\times $}}\end{tabular} & \hspace*{1pt}\begin{tabular}[t]{@{}c@{}}\rule{0pt}{12pt}\colorbox{BenchBest}{\strut 48.1s}\\[-.2pt]{\fontsize{5.8}{6.4}\selectfont \textcolor{ForestGreen}{$ \blacktriangle\,7.3\times $}}\\[-.2pt]\colorbox{BenchBest}{\strut 47s}\\[-.2pt]{\fontsize{5.8}{6.4}\selectfont \textcolor{ForestGreen}{$ \blacktriangle\,2.4\times $}}\end{tabular} & \hspace*{1pt}\begin{tabular}[t]{@{}c@{}}\rule{0pt}{12pt}\colorbox{BenchBest}{\strut 44.5s}\\[-.2pt]{\fontsize{5.8}{6.4}\selectfont \textcolor{ForestGreen}{$ \blacktriangle\,6.5\times $}}\\[-.2pt]\colorbox{BenchBest}{\strut 42.6s}\\[-.2pt]{\fontsize{5.8}{6.4}\selectfont \textcolor{ForestGreen}{$ \blacktriangle\,2.7\times $}}\end{tabular} & \hspace*{1pt}\begin{tabular}[t]{@{}c@{}}\rule{0pt}{12pt}\colorbox{BenchBest}{\strut 1m\,30.2s}\\[-.2pt]{\fontsize{5.8}{6.4}\selectfont \textcolor{ForestGreen}{$ \blacktriangle\,6.4\times $}}\\[-.2pt]\colorbox{BenchBest}{\strut 1m\,26.7s}\\[-.2pt]{\fontsize{5.8}{6.4}\selectfont \textcolor{ForestGreen}{$ \blacktriangle\,2.2\times $}}\end{tabular} & \hspace*{1pt}\begin{tabular}[t]{@{}c@{}}\rule{0pt}{12pt}\colorbox{BenchBest}{\strut 2m\,46.4s}\\[-.2pt]{\fontsize{5.8}{6.4}\selectfont \textcolor{ForestGreen}{$ \blacktriangle\,4.0\times $}}\\[-.2pt]2m\,42.7s\\[-.2pt]{\fontsize{5.8}{6.4}\selectfont \textcolor{BrickRed}{$ \blacktriangledown\,27\% $}}\end{tabular}\\[3pt]\hline
\BenchmarkSourceLabel{bench:aws-2-k15-w10:row:mpf-bounded-scalar-block}{R2}\shortstack[l]{{\color{gray}R2}\\\textbf{BioseqDB}} & \begin{tabular}[t]{@{}l@{}}\rule{0pt}{12pt}MPF / Bounded\\Scalar-DP / Block\end{tabular} & \hspace*{1pt}\begin{tabular}[t]{@{}c@{}}\rule{0pt}{12pt}3.6s\\[-.2pt]{\fontsize{5.8}{6.4}\selectfont \textcolor{ForestGreen}{$ \blacktriangle\,3.1\times $}}\\[-.2pt]2.8s\\[-.2pt]{\fontsize{5.8}{6.4}\selectfont \textcolor{ForestGreen}{$ \blacktriangle\,3.5\times $}}\end{tabular} & \hspace*{1pt}\begin{tabular}[t]{@{}c@{}}\rule{0pt}{12pt}3.4s\\[-.2pt]{\fontsize{5.8}{6.4}\selectfont \textcolor{ForestGreen}{$ \blacktriangle\,3.1\times $}}\\[-.2pt]2.6s\\[-.2pt]{\fontsize{5.8}{6.4}\selectfont \textcolor{ForestGreen}{$ \blacktriangle\,3.3\times $}}\end{tabular} & \hspace*{1pt}\begin{tabular}[t]{@{}c@{}}\rule{0pt}{12pt}6.4s\\[-.2pt]{\fontsize{5.8}{6.4}\selectfont \textcolor{ForestGreen}{$ \blacktriangle\,3.5\times $}}\\[-.2pt]4.4s\\[-.2pt]{\fontsize{5.8}{6.4}\selectfont \textcolor{ForestGreen}{$ \blacktriangle\,4.2\times $}}\end{tabular} & \hspace*{1pt}\begin{tabular}[t]{@{}c@{}}\rule{0pt}{12pt}\colorbox{BenchBest}{\strut 3.2s}\\[-.2pt]{\fontsize{5.8}{6.4}\selectfont \textcolor{ForestGreen}{$ \blacktriangle\,22.5\times $}}\\[-.2pt]\colorbox{BenchBest}{\strut 2.4s}\\[-.2pt]{\fontsize{5.8}{6.4}\selectfont \textcolor{ForestGreen}{$ \blacktriangle\,27.7\times $}}\end{tabular} & \hspace*{1pt}\begin{tabular}[t]{@{}c@{}}\rule{0pt}{12pt}10.2s\\[-.2pt]{\fontsize{5.8}{6.4}\selectfont \textcolor{gray}{$ \blacktriangleright\,3\% $}}\\[-.2pt]6s\\[-.2pt]{\fontsize{5.8}{6.4}\selectfont \textcolor{ForestGreen}{$ \blacktriangle\,1.5\times $}}\end{tabular} & \hspace*{1pt}\begin{tabular}[t]{@{}c@{}}\rule{0pt}{12pt}8.6s\\[-.2pt]{\fontsize{5.8}{6.4}\selectfont \textcolor{ForestGreen}{$ \blacktriangle\,10\% $}}\\[-.2pt]5.5s\\[-.2pt]{\fontsize{5.8}{6.4}\selectfont \textcolor{ForestGreen}{$ \blacktriangle\,28\% $}}\end{tabular} & \hspace*{1pt}\begin{tabular}[t]{@{}c@{}}\rule{0pt}{12pt}14.6s\\[-.2pt]{\fontsize{5.8}{6.4}\selectfont \textcolor{ForestGreen}{$ \blacktriangle\,11\% $}}\\[-.2pt]8.4s\\[-.2pt]{\fontsize{5.8}{6.4}\selectfont \textcolor{ForestGreen}{$ \blacktriangle\,1.5\times $}}\end{tabular} & \hspace*{1pt}\begin{tabular}[t]{@{}c@{}}\rule{0pt}{12pt}12.4s\\[-.2pt]{\fontsize{5.8}{6.4}\selectfont \textcolor{ForestGreen}{$ \blacktriangle\,2.3\times $}}\\[-.2pt]9.2s\\[-.2pt]{\fontsize{5.8}{6.4}\selectfont \textcolor{ForestGreen}{$ \blacktriangle\,2.9\times $}}\end{tabular} & \hspace*{1pt}\begin{tabular}[t]{@{}c@{}}\rule{0pt}{12pt}1m\,56.2s\\[-.2pt]{\fontsize{5.8}{6.4}\selectfont \textcolor{BrickRed}{$ \blacktriangledown\,2.2\times $}}\\[-.2pt]1m\,47.8s\\[-.2pt]{\fontsize{5.8}{6.4}\selectfont \textcolor{BrickRed}{$ \blacktriangledown\,2.1\times $}}\end{tabular} & \hspace*{1pt}\begin{tabular}[t]{@{}c@{}}\rule{0pt}{12pt}2m\,7.4s\\[-.2pt]{\fontsize{5.8}{6.4}\selectfont \textcolor{BrickRed}{$ \blacktriangledown\,2.3\times $}}\\[-.2pt]1m\,56.8s\\[-.2pt]{\fontsize{5.8}{6.4}\selectfont \textcolor{BrickRed}{$ \blacktriangledown\,2.1\times $}}\end{tabular} & \textcolor{gray}{\ensuremath{\cdot}} & \hspace*{1pt}\begin{tabular}[t]{@{}c@{}}\rule{0pt}{12pt}\colorbox{BenchBest}{\strut 2m\,1.9s}\\[-.2pt]{\fontsize{5.8}{6.4}\selectfont \textcolor{ForestGreen}{$ \blacktriangle\,3.0\times $}}\\[-.2pt]\colorbox{BenchBest}{\strut 1m\,56.1s}\\[-.2pt]{\fontsize{5.8}{6.4}\selectfont \textcolor{ForestGreen}{$ \blacktriangle\,3.1\times $}}\end{tabular} & \hspace*{1pt}\begin{tabular}[t]{@{}c@{}}\rule{0pt}{12pt}6m\,7.7s\\[-.2pt]{\fontsize{5.8}{6.4}\selectfont \textcolor{BrickRed}{$ \blacktriangledown\,7.7\times $}}\\[-.2pt]2m\,6.1s\\[-.2pt]{\fontsize{5.8}{6.4}\selectfont \textcolor{BrickRed}{$ \blacktriangledown\,2.7\times $}}\end{tabular} & \hspace*{1pt}\begin{tabular}[t]{@{}c@{}}\rule{0pt}{12pt}6m\,48.2s\\[-.2pt]{\fontsize{5.8}{6.4}\selectfont \textcolor{BrickRed}{$ \blacktriangledown\,9.2\times $}}\\[-.2pt]3m\,14.7s\\[-.2pt]{\fontsize{5.8}{6.4}\selectfont \textcolor{BrickRed}{$ \blacktriangledown\,4.6\times $}}\end{tabular} & \textcolor{gray}{\ensuremath{\cdot}} & \hspace*{1pt}\begin{tabular}[t]{@{}c@{}}\rule{0pt}{12pt}11m\,10.4s\\[-.2pt]{\fontsize{5.8}{6.4}\selectfont \textcolor{BrickRed}{$ \blacktriangledown\,4.0\times $}}\\[-.2pt]\colorbox{BenchBest}{\strut 2m\,8s}\\[-.2pt]{\fontsize{5.8}{6.4}\selectfont \textcolor{ForestGreen}{$ \blacktriangle\,21\% $}}\end{tabular}\\[3pt]\hline
\BenchmarkSourceLabel{bench:aws-2-k15-w10:row:mpf-simdref-scalar-block}{R3}\shortstack[l]{{\color{gray}R3}} & \begin{tabular}[t]{@{}l@{}}\rule{0pt}{12pt}MPF / SIMD-ref\\Scalar-DP / Block\end{tabular} & \hspace*{1pt}\begin{tabular}[t]{@{}c@{}}\rule{0pt}{12pt}\colorbox{BenchBest}{\strut 3.6s}\\[-.2pt]{\fontsize{5.8}{6.4}\selectfont \textcolor{ForestGreen}{$ \blacktriangle\,3.1\times $}}\\[-.2pt]2.9s\\[-.2pt]{\fontsize{5.8}{6.4}\selectfont \textcolor{ForestGreen}{$ \blacktriangle\,3.5\times $}}\end{tabular} & \hspace*{1pt}\begin{tabular}[t]{@{}c@{}}\rule{0pt}{12pt}\colorbox{BenchBest}{\strut 3.4s}\\[-.2pt]{\fontsize{5.8}{6.4}\selectfont \textcolor{ForestGreen}{$ \blacktriangle\,3.2\times $}}\\[-.2pt]2.6s\\[-.2pt]{\fontsize{5.8}{6.4}\selectfont \textcolor{ForestGreen}{$ \blacktriangle\,3.3\times $}}\end{tabular} & \hspace*{1pt}\begin{tabular}[t]{@{}c@{}}\rule{0pt}{12pt}6.5s\\[-.2pt]{\fontsize{5.8}{6.4}\selectfont \textcolor{ForestGreen}{$ \blacktriangle\,3.5\times $}}\\[-.2pt]4.4s\\[-.2pt]{\fontsize{5.8}{6.4}\selectfont \textcolor{ForestGreen}{$ \blacktriangle\,4.2\times $}}\end{tabular} & \hspace*{1pt}\begin{tabular}[t]{@{}c@{}}\rule{0pt}{12pt}3.3s\\[-.2pt]{\fontsize{5.8}{6.4}\selectfont \textcolor{ForestGreen}{$ \blacktriangle\,22.1\times $}}\\[-.2pt]2.5s\\[-.2pt]{\fontsize{5.8}{6.4}\selectfont \textcolor{ForestGreen}{$ \blacktriangle\,26.2\times $}}\end{tabular} & \hspace*{1pt}\begin{tabular}[t]{@{}c@{}}\rule{0pt}{12pt}12.7s\\[-.2pt]{\fontsize{5.8}{6.4}\selectfont \textcolor{BrickRed}{$ \blacktriangledown\,29\% $}}\\[-.2pt]7.4s\\[-.2pt]{\fontsize{5.8}{6.4}\selectfont \textcolor{ForestGreen}{$ \blacktriangle\,15\% $}}\end{tabular} & \hspace*{1pt}\begin{tabular}[t]{@{}c@{}}\rule{0pt}{12pt}13s\\[-.2pt]{\fontsize{5.8}{6.4}\selectfont \textcolor{BrickRed}{$ \blacktriangledown\,1.4\times $}}\\[-.2pt]7.4s\\[-.2pt]{\fontsize{5.8}{6.4}\selectfont \textcolor{gray}{$ \blacktriangleright\,3\% $}}\end{tabular} & \hspace*{1pt}\begin{tabular}[t]{@{}c@{}}\rule{0pt}{12pt}24s\\[-.2pt]{\fontsize{5.8}{6.4}\selectfont \textcolor{BrickRed}{$ \blacktriangledown\,1.5\times $}}\\[-.2pt]13.6s\\[-.2pt]{\fontsize{5.8}{6.4}\selectfont \textcolor{BrickRed}{$ \blacktriangledown\,5\% $}}\end{tabular} & \hspace*{1pt}\begin{tabular}[t]{@{}c@{}}\rule{0pt}{12pt}16.5s\\[-.2pt]{\fontsize{5.8}{6.4}\selectfont \textcolor{ForestGreen}{$ \blacktriangle\,1.7\times $}}\\[-.2pt]13.6s\\[-.2pt]{\fontsize{5.8}{6.4}\selectfont \textcolor{ForestGreen}{$ \blacktriangle\,1.9\times $}}\end{tabular} & \hspace*{1pt}\begin{tabular}[t]{@{}c@{}}\rule{0pt}{12pt}2m\,51.6s\\[-.2pt]{\fontsize{5.8}{6.4}\selectfont \textcolor{BrickRed}{$ \blacktriangledown\,3.2\times $}}\\[-.2pt]2m\,26.6s\\[-.2pt]{\fontsize{5.8}{6.4}\selectfont \textcolor{BrickRed}{$ \blacktriangledown\,2.8\times $}}\end{tabular} & \hspace*{1pt}\begin{tabular}[t]{@{}c@{}}\rule{0pt}{12pt}2m\,23.6s\\[-.2pt]{\fontsize{5.8}{6.4}\selectfont \textcolor{BrickRed}{$ \blacktriangledown\,2.5\times $}}\\[-.2pt]2m\,2.9s\\[-.2pt]{\fontsize{5.8}{6.4}\selectfont \textcolor{BrickRed}{$ \blacktriangledown\,2.3\times $}}\end{tabular} & \textcolor{gray}{\ensuremath{\cdot}} & \textcolor{gray}{\ensuremath{\cdot}} & \textcolor{gray}{\ensuremath{\cdot}} & \textcolor{gray}{\ensuremath{\cdot}} & \textcolor{gray}{\ensuremath{\cdot}} & \textcolor{gray}{\ensuremath{\cdot}}\\[3pt]\hline
\BenchmarkSourceLabel{bench:aws-2-k15-w10:row:mpf-scalar-scalar-block}{R4}\shortstack[l]{{\color{gray}R4}} & \begin{tabular}[t]{@{}l@{}}\rule{0pt}{12pt}MPF / Scalar\\Scalar-DP / Block\end{tabular} & \hspace*{1pt}\begin{tabular}[t]{@{}c@{}}\rule{0pt}{12pt}3.8s\\[-.2pt]{\fontsize{5.8}{6.4}\selectfont \textcolor{ForestGreen}{$ \blacktriangle\,3.0\times $}}\\[-.2pt]3.1s\\[-.2pt]{\fontsize{5.8}{6.4}\selectfont \textcolor{ForestGreen}{$ \blacktriangle\,3.3\times $}}\end{tabular} & \hspace*{1pt}\begin{tabular}[t]{@{}c@{}}\rule{0pt}{12pt}3.5s\\[-.2pt]{\fontsize{5.8}{6.4}\selectfont \textcolor{ForestGreen}{$ \blacktriangle\,3.1\times $}}\\[-.2pt]2.7s\\[-.2pt]{\fontsize{5.8}{6.4}\selectfont \textcolor{ForestGreen}{$ \blacktriangle\,3.2\times $}}\end{tabular} & \hspace*{1pt}\begin{tabular}[t]{@{}c@{}}\rule{0pt}{12pt}6.5s\\[-.2pt]{\fontsize{5.8}{6.4}\selectfont \textcolor{ForestGreen}{$ \blacktriangle\,3.4\times $}}\\[-.2pt]4.8s\\[-.2pt]{\fontsize{5.8}{6.4}\selectfont \textcolor{ForestGreen}{$ \blacktriangle\,3.9\times $}}\end{tabular} & \hspace*{1pt}\begin{tabular}[t]{@{}c@{}}\rule{0pt}{12pt}3.4s\\[-.2pt]{\fontsize{5.8}{6.4}\selectfont \textcolor{ForestGreen}{$ \blacktriangle\,21.8\times $}}\\[-.2pt]2.5s\\[-.2pt]{\fontsize{5.8}{6.4}\selectfont \textcolor{ForestGreen}{$ \blacktriangle\,26.2\times $}}\end{tabular} & \hspace*{1pt}\begin{tabular}[t]{@{}c@{}}\rule{0pt}{12pt}10.1s\\[-.2pt]{\fontsize{5.8}{6.4}\selectfont \textcolor{gray}{$ \blacktriangleright\,2\% $}}\\[-.2pt]5.9s\\[-.2pt]{\fontsize{5.8}{6.4}\selectfont \textcolor{ForestGreen}{$ \blacktriangle\,1.5\times $}}\end{tabular} & \hspace*{1pt}\begin{tabular}[t]{@{}c@{}}\rule{0pt}{12pt}8.3s\\[-.2pt]{\fontsize{5.8}{6.4}\selectfont \textcolor{ForestGreen}{$ \blacktriangle\,13\% $}}\\[-.2pt]5.2s\\[-.2pt]{\fontsize{5.8}{6.4}\selectfont \textcolor{ForestGreen}{$ \blacktriangle\,1.5\times $}}\end{tabular} & \hspace*{1pt}\begin{tabular}[t]{@{}c@{}}\rule{0pt}{12pt}15.9s\\[-.2pt]{\fontsize{5.8}{6.4}\selectfont \textcolor{gray}{$ \blacktriangleright\,4\% $}}\\[-.2pt]9.6s\\[-.2pt]{\fontsize{5.8}{6.4}\selectfont \textcolor{ForestGreen}{$ \blacktriangle\,26\% $}}\end{tabular} & \hspace*{1pt}\begin{tabular}[t]{@{}c@{}}\rule{0pt}{12pt}14.5s\\[-.2pt]{\fontsize{5.8}{6.4}\selectfont \textcolor{ForestGreen}{$ \blacktriangle\,1.9\times $}}\\[-.2pt]11.3s\\[-.2pt]{\fontsize{5.8}{6.4}\selectfont \textcolor{ForestGreen}{$ \blacktriangle\,2.3\times $}}\end{tabular} & \hspace*{1pt}\begin{tabular}[t]{@{}c@{}}\rule{0pt}{12pt}2m\,51.4s\\[-.2pt]{\fontsize{5.8}{6.4}\selectfont \textcolor{BrickRed}{$ \blacktriangledown\,3.2\times $}}\\[-.2pt]2m\,29.1s\\[-.2pt]{\fontsize{5.8}{6.4}\selectfont \textcolor{BrickRed}{$ \blacktriangledown\,2.9\times $}}\end{tabular} & \hspace*{1pt}\begin{tabular}[t]{@{}c@{}}\rule{0pt}{12pt}2m\,59.3s\\[-.2pt]{\fontsize{5.8}{6.4}\selectfont \textcolor{BrickRed}{$ \blacktriangledown\,3.2\times $}}\\[-.2pt]2m\,39.2s\\[-.2pt]{\fontsize{5.8}{6.4}\selectfont \textcolor{BrickRed}{$ \blacktriangledown\,2.9\times $}}\end{tabular} & \textcolor{gray}{\ensuremath{\cdot}} & \textcolor{gray}{\ensuremath{\cdot}} & \hspace*{1pt}\begin{tabular}[t]{@{}c@{}}\rule{0pt}{12pt}7m\,9.6s\\[-.2pt]{\fontsize{5.8}{6.4}\selectfont \textcolor{BrickRed}{$ \blacktriangledown\,8.9\times $}}\\[-.2pt]2m\,48s\\[-.2pt]{\fontsize{5.8}{6.4}\selectfont \textcolor{BrickRed}{$ \blacktriangledown\,3.6\times $}}\end{tabular} & \hspace*{1pt}\begin{tabular}[t]{@{}c@{}}\rule{0pt}{12pt}6m\,28.6s\\[-.2pt]{\fontsize{5.8}{6.4}\selectfont \textcolor{BrickRed}{$ \blacktriangledown\,8.7\times $}}\\[-.2pt]2m\,48.4s\\[-.2pt]{\fontsize{5.8}{6.4}\selectfont \textcolor{BrickRed}{$ \blacktriangledown\,4.0\times $}}\end{tabular} & \textcolor{gray}{\ensuremath{\cdot}} & \textcolor{gray}{\ensuremath{\cdot}}\\[3pt]\hline
\BenchmarkSourceLabel{bench:aws-2-k15-w10:row:mpf-twostack-scalar-block}{R5}\shortstack[l]{{\color{gray}R5}} & \begin{tabular}[t]{@{}l@{}}\rule{0pt}{12pt}MPF / Two-stack\\Scalar-DP / Block\end{tabular} & \hspace*{1pt}\begin{tabular}[t]{@{}c@{}}\rule{0pt}{12pt}3.6s\\[-.2pt]{\fontsize{5.8}{6.4}\selectfont \textcolor{ForestGreen}{$ \blacktriangle\,3.1\times $}}\\[-.2pt]2.9s\\[-.2pt]{\fontsize{5.8}{6.4}\selectfont \textcolor{ForestGreen}{$ \blacktriangle\,3.4\times $}}\end{tabular} & \hspace*{1pt}\begin{tabular}[t]{@{}c@{}}\rule{0pt}{12pt}3.4s\\[-.2pt]{\fontsize{5.8}{6.4}\selectfont \textcolor{ForestGreen}{$ \blacktriangle\,3.2\times $}}\\[-.2pt]2.6s\\[-.2pt]{\fontsize{5.8}{6.4}\selectfont \textcolor{ForestGreen}{$ \blacktriangle\,3.3\times $}}\end{tabular} & \hspace*{1pt}\begin{tabular}[t]{@{}c@{}}\rule{0pt}{12pt}\colorbox{BenchBest}{\strut 6.3s}\\[-.2pt]{\fontsize{5.8}{6.4}\selectfont \textcolor{ForestGreen}{$ \blacktriangle\,3.6\times $}}\\[-.2pt]4.6s\\[-.2pt]{\fontsize{5.8}{6.4}\selectfont \textcolor{ForestGreen}{$ \blacktriangle\,4.1\times $}}\end{tabular} & \hspace*{1pt}\begin{tabular}[t]{@{}c@{}}\rule{0pt}{12pt}7.9s\\[-.2pt]{\fontsize{5.8}{6.4}\selectfont \textcolor{ForestGreen}{$ \blacktriangle\,9.2\times $}}\\[-.2pt]7.1s\\[-.2pt]{\fontsize{5.8}{6.4}\selectfont \textcolor{ForestGreen}{$ \blacktriangle\,9.3\times $}}\end{tabular} & \hspace*{1pt}\begin{tabular}[t]{@{}c@{}}\rule{0pt}{12pt}9.7s\\[-.2pt]{\fontsize{5.8}{6.4}\selectfont \textcolor{gray}{$ \blacktriangleright\,3\% $}}\\[-.2pt]5.5s\\[-.2pt]{\fontsize{5.8}{6.4}\selectfont \textcolor{ForestGreen}{$ \blacktriangle\,1.6\times $}}\end{tabular} & \hspace*{1pt}\begin{tabular}[t]{@{}c@{}}\rule{0pt}{12pt}7.6s\\[-.2pt]{\fontsize{5.8}{6.4}\selectfont \textcolor{ForestGreen}{$ \blacktriangle\,20\% $}}\\[-.2pt]4.7s\\[-.2pt]{\fontsize{5.8}{6.4}\selectfont \textcolor{ForestGreen}{$ \blacktriangle\,1.6\times $}}\end{tabular} & \hspace*{1pt}\begin{tabular}[t]{@{}c@{}}\rule{0pt}{12pt}15.8s\\[-.2pt]{\fontsize{5.8}{6.4}\selectfont \textcolor{gray}{$ \blacktriangleright\,4\% $}}\\[-.2pt]8.8s\\[-.2pt]{\fontsize{5.8}{6.4}\selectfont \textcolor{ForestGreen}{$ \blacktriangle\,1.5\times $}}\end{tabular} & \hspace*{1pt}\begin{tabular}[t]{@{}c@{}}\rule{0pt}{12pt}\colorbox{BenchBest}{\strut 12.4s}\\[-.2pt]{\fontsize{5.8}{6.4}\selectfont \textcolor{ForestGreen}{$ \blacktriangle\,2.3\times $}}\\[-.2pt]8.9s\\[-.2pt]{\fontsize{5.8}{6.4}\selectfont \textcolor{ForestGreen}{$ \blacktriangle\,3.0\times $}}\end{tabular} & \hspace*{1pt}\begin{tabular}[t]{@{}c@{}}\rule{0pt}{12pt}2m\,41.5s\\[-.2pt]{\fontsize{5.8}{6.4}\selectfont \textcolor{BrickRed}{$ \blacktriangledown\,3.0\times $}}\\[-.2pt]2m\,17s\\[-.2pt]{\fontsize{5.8}{6.4}\selectfont \textcolor{BrickRed}{$ \blacktriangledown\,2.6\times $}}\end{tabular} & \hspace*{1pt}\begin{tabular}[t]{@{}c@{}}\rule{0pt}{12pt}2m\,37.7s\\[-.2pt]{\fontsize{5.8}{6.4}\selectfont \textcolor{BrickRed}{$ \blacktriangledown\,2.8\times $}}\\[-.2pt]2m\,16.7s\\[-.2pt]{\fontsize{5.8}{6.4}\selectfont \textcolor{BrickRed}{$ \blacktriangledown\,2.5\times $}}\end{tabular} & \textcolor{gray}{\ensuremath{\cdot}} & \textcolor{gray}{\ensuremath{\cdot}} & \hspace*{1pt}\begin{tabular}[t]{@{}c@{}}\rule{0pt}{12pt}6m\,32.3s\\[-.2pt]{\fontsize{5.8}{6.4}\selectfont \textcolor{BrickRed}{$ \blacktriangledown\,8.2\times $}}\\[-.2pt]2m\,11.7s\\[-.2pt]{\fontsize{5.8}{6.4}\selectfont \textcolor{BrickRed}{$ \blacktriangledown\,2.8\times $}}\end{tabular} & \hspace*{1pt}\begin{tabular}[t]{@{}c@{}}\rule{0pt}{12pt}5m\,45.5s\\[-.2pt]{\fontsize{5.8}{6.4}\selectfont \textcolor{BrickRed}{$ \blacktriangledown\,7.8\times $}}\\[-.2pt]2m\,3.6s\\[-.2pt]{\fontsize{5.8}{6.4}\selectfont \textcolor{BrickRed}{$ \blacktriangledown\,2.9\times $}}\end{tabular} & \textcolor{gray}{\ensuremath{\cdot}} & \textcolor{gray}{\ensuremath{\cdot}}\\[3pt]\hline
\BenchmarkSourceLabel{bench:aws-2-k15-w10:row:mpf-twostack-simd-block}{R6}\shortstack[l]{{\color{gray}R6}} & \begin{tabular}[t]{@{}l@{}}\rule{0pt}{12pt}MPF / Two-stack\\SIMD-DP / Block\end{tabular} & \hspace*{1pt}\begin{tabular}[t]{@{}c@{}}\rule{0pt}{12pt}3.7s\\[-.2pt]{\fontsize{5.8}{6.4}\selectfont \textcolor{ForestGreen}{$ \blacktriangle\,3.0\times $}}\\[-.2pt]3s\\[-.2pt]{\fontsize{5.8}{6.4}\selectfont \textcolor{ForestGreen}{$ \blacktriangle\,3.4\times $}}\end{tabular} & \hspace*{1pt}\begin{tabular}[t]{@{}c@{}}\rule{0pt}{12pt}3.5s\\[-.2pt]{\fontsize{5.8}{6.4}\selectfont \textcolor{ForestGreen}{$ \blacktriangle\,3.0\times $}}\\[-.2pt]2.6s\\[-.2pt]{\fontsize{5.8}{6.4}\selectfont \textcolor{ForestGreen}{$ \blacktriangle\,3.3\times $}}\end{tabular} & \hspace*{1pt}\begin{tabular}[t]{@{}c@{}}\rule{0pt}{12pt}6.7s\\[-.2pt]{\fontsize{5.8}{6.4}\selectfont \textcolor{ForestGreen}{$ \blacktriangle\,3.4\times $}}\\[-.2pt]4.7s\\[-.2pt]{\fontsize{5.8}{6.4}\selectfont \textcolor{ForestGreen}{$ \blacktriangle\,4.0\times $}}\end{tabular} & \hspace*{1pt}\begin{tabular}[t]{@{}c@{}}\rule{0pt}{12pt}3.4s\\[-.2pt]{\fontsize{5.8}{6.4}\selectfont \textcolor{ForestGreen}{$ \blacktriangle\,21.5\times $}}\\[-.2pt]2.5s\\[-.2pt]{\fontsize{5.8}{6.4}\selectfont \textcolor{ForestGreen}{$ \blacktriangle\,25.8\times $}}\end{tabular} & \hspace*{1pt}\begin{tabular}[t]{@{}c@{}}\rule{0pt}{12pt}11.8s\\[-.2pt]{\fontsize{5.8}{6.4}\selectfont \textcolor{BrickRed}{$ \blacktriangledown\,19\% $}}\\[-.2pt]7.6s\\[-.2pt]{\fontsize{5.8}{6.4}\selectfont \textcolor{ForestGreen}{$ \blacktriangle\,14\% $}}\end{tabular} & \hspace*{1pt}\begin{tabular}[t]{@{}c@{}}\rule{0pt}{12pt}8.6s\\[-.2pt]{\fontsize{5.8}{6.4}\selectfont \textcolor{ForestGreen}{$ \blacktriangle\,10\% $}}\\[-.2pt]5s\\[-.2pt]{\fontsize{5.8}{6.4}\selectfont \textcolor{ForestGreen}{$ \blacktriangle\,1.5\times $}}\end{tabular} & \hspace*{1pt}\begin{tabular}[t]{@{}c@{}}\rule{0pt}{12pt}16s\\[-.2pt]{\fontsize{5.8}{6.4}\selectfont \textcolor{gray}{$ \blacktriangleright\,3\% $}}\\[-.2pt]9.1s\\[-.2pt]{\fontsize{5.8}{6.4}\selectfont \textcolor{ForestGreen}{$ \blacktriangle\,30\% $}}\end{tabular} & \hspace*{1pt}\begin{tabular}[t]{@{}c@{}}\rule{0pt}{12pt}13.5s\\[-.2pt]{\fontsize{5.8}{6.4}\selectfont \textcolor{ForestGreen}{$ \blacktriangle\,2.1\times $}}\\[-.2pt]11s\\[-.2pt]{\fontsize{5.8}{6.4}\selectfont \textcolor{ForestGreen}{$ \blacktriangle\,2.4\times $}}\end{tabular} & \hspace*{1pt}\begin{tabular}[t]{@{}c@{}}\rule{0pt}{12pt}2m\,38.5s\\[-.2pt]{\fontsize{5.8}{6.4}\selectfont \textcolor{BrickRed}{$ \blacktriangledown\,3.0\times $}}\\[-.2pt]2m\,16s\\[-.2pt]{\fontsize{5.8}{6.4}\selectfont \textcolor{BrickRed}{$ \blacktriangledown\,2.6\times $}}\end{tabular} & \hspace*{1pt}\begin{tabular}[t]{@{}c@{}}\rule{0pt}{12pt}2m\,34.1s\\[-.2pt]{\fontsize{5.8}{6.4}\selectfont \textcolor{BrickRed}{$ \blacktriangledown\,2.7\times $}}\\[-.2pt]2m\,13.1s\\[-.2pt]{\fontsize{5.8}{6.4}\selectfont \textcolor{BrickRed}{$ \blacktriangledown\,2.4\times $}}\end{tabular} & \textcolor{gray}{\ensuremath{\cdot}} & \hspace*{1pt}\begin{tabular}[t]{@{}c@{}}\rule{0pt}{12pt}3m\,52.6s\\[-.2pt]{\fontsize{5.8}{6.4}\selectfont \textcolor{ForestGreen}{$ \blacktriangle\,1.6\times $}}\\[-.2pt]3m\,35.8s\\[-.2pt]{\fontsize{5.8}{6.4}\selectfont \textcolor{ForestGreen}{$ \blacktriangle\,1.6\times $}}\end{tabular} & \hspace*{1pt}\begin{tabular}[t]{@{}c@{}}\rule{0pt}{12pt}6m\,24.3s\\[-.2pt]{\fontsize{5.8}{6.4}\selectfont \textcolor{BrickRed}{$ \blacktriangledown\,8.0\times $}}\\[-.2pt]2m\,4.9s\\[-.2pt]{\fontsize{5.8}{6.4}\selectfont \textcolor{BrickRed}{$ \blacktriangledown\,2.7\times $}}\end{tabular} & \hspace*{1pt}\begin{tabular}[t]{@{}c@{}}\rule{0pt}{12pt}5m\,40.8s\\[-.2pt]{\fontsize{5.8}{6.4}\selectfont \textcolor{BrickRed}{$ \blacktriangledown\,7.7\times $}}\\[-.2pt]1m\,59.2s\\[-.2pt]{\fontsize{5.8}{6.4}\selectfont \textcolor{BrickRed}{$ \blacktriangledown\,2.8\times $}}\end{tabular} & \textcolor{gray}{\ensuremath{\cdot}} & \hspace*{1pt}\begin{tabular}[t]{@{}c@{}}\rule{0pt}{12pt}13m\,14.7s\\[-.2pt]{\fontsize{5.8}{6.4}\selectfont \textcolor{BrickRed}{$ \blacktriangledown\,4.8\times $}}\\[-.2pt]3m\,54.6s\\[-.2pt]{\fontsize{5.8}{6.4}\selectfont \textcolor{BrickRed}{$ \blacktriangledown\,1.4\times $}}\end{tabular}\\[3pt]\hline
\BenchmarkSourceLabel{bench:aws-2-k15-w10:row:mpf-simdref-simd-block}{R7}\shortstack[l]{{\color{gray}R7}} & \begin{tabular}[t]{@{}l@{}}\rule{0pt}{12pt}MPF / SIMD-ref\\SIMD-DP / Block\end{tabular} & \hspace*{1pt}\begin{tabular}[t]{@{}c@{}}\rule{0pt}{12pt}3.7s\\[-.2pt]{\fontsize{5.8}{6.4}\selectfont \textcolor{ForestGreen}{$ \blacktriangle\,3.0\times $}}\\[-.2pt]2.9s\\[-.2pt]{\fontsize{5.8}{6.4}\selectfont \textcolor{ForestGreen}{$ \blacktriangle\,3.4\times $}}\end{tabular} & \hspace*{1pt}\begin{tabular}[t]{@{}c@{}}\rule{0pt}{12pt}3.5s\\[-.2pt]{\fontsize{5.8}{6.4}\selectfont \textcolor{ForestGreen}{$ \blacktriangle\,3.0\times $}}\\[-.2pt]2.7s\\[-.2pt]{\fontsize{5.8}{6.4}\selectfont \textcolor{ForestGreen}{$ \blacktriangle\,3.2\times $}}\end{tabular} & \hspace*{1pt}\begin{tabular}[t]{@{}c@{}}\rule{0pt}{12pt}6.9s\\[-.2pt]{\fontsize{5.8}{6.4}\selectfont \textcolor{ForestGreen}{$ \blacktriangle\,3.3\times $}}\\[-.2pt]4.5s\\[-.2pt]{\fontsize{5.8}{6.4}\selectfont \textcolor{ForestGreen}{$ \blacktriangle\,4.2\times $}}\end{tabular} & \hspace*{1pt}\begin{tabular}[t]{@{}c@{}}\rule{0pt}{12pt}3.4s\\[-.2pt]{\fontsize{5.8}{6.4}\selectfont \textcolor{ForestGreen}{$ \blacktriangle\,21.4\times $}}\\[-.2pt]2.5s\\[-.2pt]{\fontsize{5.8}{6.4}\selectfont \textcolor{ForestGreen}{$ \blacktriangle\,26.3\times $}}\end{tabular} & \hspace*{1pt}\begin{tabular}[t]{@{}c@{}}\rule{0pt}{12pt}13.2s\\[-.2pt]{\fontsize{5.8}{6.4}\selectfont \textcolor{BrickRed}{$ \blacktriangledown\,1.3\times $}}\\[-.2pt]7.7s\\[-.2pt]{\fontsize{5.8}{6.4}\selectfont \textcolor{ForestGreen}{$ \blacktriangle\,13\% $}}\end{tabular} & \hspace*{1pt}\begin{tabular}[t]{@{}c@{}}\rule{0pt}{12pt}13.9s\\[-.2pt]{\fontsize{5.8}{6.4}\selectfont \textcolor{BrickRed}{$ \blacktriangledown\,1.5\times $}}\\[-.2pt]8s\\[-.2pt]{\fontsize{5.8}{6.4}\selectfont \textcolor{gray}{$ \blacktriangleright\,5\% $}}\end{tabular} & \hspace*{1pt}\begin{tabular}[t]{@{}c@{}}\rule{0pt}{12pt}28.7s\\[-.2pt]{\fontsize{5.8}{6.4}\selectfont \textcolor{BrickRed}{$ \blacktriangledown\,1.7\times $}}\\[-.2pt]14.5s\\[-.2pt]{\fontsize{5.8}{6.4}\selectfont \textcolor{BrickRed}{$ \blacktriangledown\,13\% $}}\end{tabular} & \hspace*{1pt}\begin{tabular}[t]{@{}c@{}}\rule{0pt}{12pt}26.2s\\[-.2pt]{\fontsize{5.8}{6.4}\selectfont \textcolor{ForestGreen}{$ \blacktriangle\,7\% $}}\\[-.2pt]24.4s\\[-.2pt]{\fontsize{5.8}{6.4}\selectfont \textcolor{ForestGreen}{$ \blacktriangle\,7\% $}}\end{tabular} & \hspace*{1pt}\begin{tabular}[t]{@{}c@{}}\rule{0pt}{12pt}2m\,47.2s\\[-.2pt]{\fontsize{5.8}{6.4}\selectfont \textcolor{BrickRed}{$ \blacktriangledown\,3.2\times $}}\\[-.2pt]2m\,22.2s\\[-.2pt]{\fontsize{5.8}{6.4}\selectfont \textcolor{BrickRed}{$ \blacktriangledown\,2.7\times $}}\end{tabular} & \hspace*{1pt}\begin{tabular}[t]{@{}c@{}}\rule{0pt}{12pt}2m\,32.6s\\[-.2pt]{\fontsize{5.8}{6.4}\selectfont \textcolor{BrickRed}{$ \blacktriangledown\,2.7\times $}}\\[-.2pt]2m\,12.9s\\[-.2pt]{\fontsize{5.8}{6.4}\selectfont \textcolor{BrickRed}{$ \blacktriangledown\,2.4\times $}}\end{tabular} & \textcolor{gray}{\ensuremath{\cdot}} & \hspace*{1pt}\begin{tabular}[t]{@{}c@{}}\rule{0pt}{12pt}3m\,6.8s\\[-.2pt]{\fontsize{5.8}{6.4}\selectfont \textcolor{ForestGreen}{$ \blacktriangle\,1.9\times $}}\\[-.2pt]2m\,49.2s\\[-.2pt]{\fontsize{5.8}{6.4}\selectfont \textcolor{ForestGreen}{$ \blacktriangle\,2.1\times $}}\end{tabular} & \textcolor{gray}{\ensuremath{\cdot}} & \textcolor{gray}{\ensuremath{\cdot}} & \textcolor{gray}{\ensuremath{\cdot}} & \textcolor{gray}{\ensuremath{\cdot}}\\[3pt]\hline
\BenchmarkSourceLabel{bench:aws-2-k15-w10:row:mpf-bounded-rmq-block}{R8}\shortstack[l]{{\color{gray}R8}} & \begin{tabular}[t]{@{}l@{}}\rule{0pt}{12pt}MPF / Bounded\\RMQ / Block\end{tabular} & \hspace*{1pt}\begin{tabular}[t]{@{}c@{}}\rule{0pt}{12pt}3.7s\\[-.2pt]{\fontsize{5.8}{6.4}\selectfont \textcolor{ForestGreen}{$ \blacktriangle\,3.0\times $}}\\[-.2pt]2.9s\\[-.2pt]{\fontsize{5.8}{6.4}\selectfont \textcolor{ForestGreen}{$ \blacktriangle\,3.4\times $}}\end{tabular} & \hspace*{1pt}\begin{tabular}[t]{@{}c@{}}\rule{0pt}{12pt}3.5s\\[-.2pt]{\fontsize{5.8}{6.4}\selectfont \textcolor{ForestGreen}{$ \blacktriangle\,3.0\times $}}\\[-.2pt]2.7s\\[-.2pt]{\fontsize{5.8}{6.4}\selectfont \textcolor{ForestGreen}{$ \blacktriangle\,3.2\times $}}\end{tabular} & \hspace*{1pt}\begin{tabular}[t]{@{}c@{}}\rule{0pt}{12pt}6.5s\\[-.2pt]{\fontsize{5.8}{6.4}\selectfont \textcolor{ForestGreen}{$ \blacktriangle\,3.5\times $}}\\[-.2pt]4.5s\\[-.2pt]{\fontsize{5.8}{6.4}\selectfont \textcolor{ForestGreen}{$ \blacktriangle\,4.2\times $}}\end{tabular} & \hspace*{1pt}\begin{tabular}[t]{@{}c@{}}\rule{0pt}{12pt}3.4s\\[-.2pt]{\fontsize{5.8}{6.4}\selectfont \textcolor{ForestGreen}{$ \blacktriangle\,21.2\times $}}\\[-.2pt]2.5s\\[-.2pt]{\fontsize{5.8}{6.4}\selectfont \textcolor{ForestGreen}{$ \blacktriangle\,26.1\times $}}\end{tabular} & \hspace*{1pt}\begin{tabular}[t]{@{}c@{}}\rule{0pt}{12pt}18.7s\\[-.2pt]{\fontsize{5.8}{6.4}\selectfont \textcolor{BrickRed}{$ \blacktriangledown\,1.9\times $}}\\[-.2pt]14.5s\\[-.2pt]{\fontsize{5.8}{6.4}\selectfont \textcolor{BrickRed}{$ \blacktriangledown\,1.6\times $}}\end{tabular} & \hspace*{1pt}\begin{tabular}[t]{@{}c@{}}\rule{0pt}{12pt}20s\\[-.2pt]{\fontsize{5.8}{6.4}\selectfont \textcolor{BrickRed}{$ \blacktriangledown\,2.1\times $}}\\[-.2pt]16.6s\\[-.2pt]{\fontsize{5.8}{6.4}\selectfont \textcolor{BrickRed}{$ \blacktriangledown\,2.2\times $}}\end{tabular} & \hspace*{1pt}\begin{tabular}[t]{@{}c@{}}\rule{0pt}{12pt}21.3s\\[-.2pt]{\fontsize{5.8}{6.4}\selectfont \textcolor{BrickRed}{$ \blacktriangledown\,29\% $}}\\[-.2pt]14.6s\\[-.2pt]{\fontsize{5.8}{6.4}\selectfont \textcolor{BrickRed}{$ \blacktriangledown\,13\% $}}\end{tabular} & \hspace*{1pt}\begin{tabular}[t]{@{}c@{}}\rule{0pt}{12pt}12.4s\\[-.2pt]{\fontsize{5.8}{6.4}\selectfont \textcolor{ForestGreen}{$ \blacktriangle\,2.3\times $}}\\[-.2pt]9s\\[-.2pt]{\fontsize{5.8}{6.4}\selectfont \textcolor{ForestGreen}{$ \blacktriangle\,2.9\times $}}\end{tabular} & \hspace*{1pt}\begin{tabular}[t]{@{}c@{}}\rule{0pt}{12pt}1m\,55.5s\\[-.2pt]{\fontsize{5.8}{6.4}\selectfont \textcolor{BrickRed}{$ \blacktriangledown\,2.2\times $}}\\[-.2pt]1m\,47.1s\\[-.2pt]{\fontsize{5.8}{6.4}\selectfont \textcolor{BrickRed}{$ \blacktriangledown\,2.1\times $}}\end{tabular} & \hspace*{1pt}\begin{tabular}[t]{@{}c@{}}\rule{0pt}{12pt}2m\,21s\\[-.2pt]{\fontsize{5.8}{6.4}\selectfont \textcolor{BrickRed}{$ \blacktriangledown\,2.5\times $}}\\[-.2pt]2m\,10.3s\\[-.2pt]{\fontsize{5.8}{6.4}\selectfont \textcolor{BrickRed}{$ \blacktriangledown\,2.4\times $}}\end{tabular} & \textcolor{gray}{\ensuremath{\cdot}} & \hspace*{1pt}\begin{tabular}[t]{@{}c@{}}\rule{0pt}{12pt}2m\,4s\\[-.2pt]{\fontsize{5.8}{6.4}\selectfont \textcolor{ForestGreen}{$ \blacktriangle\,2.9\times $}}\\[-.2pt]1m\,58.1s\\[-.2pt]{\fontsize{5.8}{6.4}\selectfont \textcolor{ForestGreen}{$ \blacktriangle\,3.0\times $}}\end{tabular} & \hspace*{1pt}\begin{tabular}[t]{@{}c@{}}\rule{0pt}{12pt}6m\,8.6s\\[-.2pt]{\fontsize{5.8}{6.4}\selectfont \textcolor{BrickRed}{$ \blacktriangledown\,7.7\times $}}\\[-.2pt]2m\,7s\\[-.2pt]{\fontsize{5.8}{6.4}\selectfont \textcolor{BrickRed}{$ \blacktriangledown\,2.7\times $}}\end{tabular} & \hspace*{1pt}\begin{tabular}[t]{@{}c@{}}\rule{0pt}{12pt}5m\,25.3s\\[-.2pt]{\fontsize{5.8}{6.4}\selectfont \textcolor{BrickRed}{$ \blacktriangledown\,7.3\times $}}\\[-.2pt]1m\,58.4s\\[-.2pt]{\fontsize{5.8}{6.4}\selectfont \textcolor{BrickRed}{$ \blacktriangledown\,2.8\times $}}\end{tabular} & \textcolor{gray}{\ensuremath{\cdot}} & \hspace*{1pt}\begin{tabular}[t]{@{}c@{}}\rule{0pt}{12pt}11m\,12.6s\\[-.2pt]{\fontsize{5.8}{6.4}\selectfont \textcolor{BrickRed}{$ \blacktriangledown\,4.0\times $}}\\[-.2pt]2m\,10.3s\\[-.2pt]{\fontsize{5.8}{6.4}\selectfont \textcolor{ForestGreen}{$ \blacktriangle\,20\% $}}\end{tabular}\\[3pt]\hline
\end{tabular}\par}
\smallskip{\footnotesize\noindent Each cell: A and B times, $\cdot$: no measurement. ---\textsuperscript{g}: GPU unavailable. n/a: no comparison.\par
\noindent Blue: one lowest A and one lowest B measurement for each column for each configuration.\par
\noindent Gray: minimap2. Yellow: my preferred MPF / Bounded / SIMD-DP / Block configuration.\par
\noindent MPF/MPS: minimizers per fragment/sequence. SIMD-ref: naive SIMD scan. Bounded: bounded SIMD algorithm.\par
}
\clearpage
\noindent{\large\bfseries Table \thetable: (D): AWS: 2 Gbp, k=15, w=10\hfill 2/7}\par\smallskip
{\fontsize{7.1}{8.1}\selectfont\setlength{\tabcolsep}{1pt}\setlength{\fboxsep}{2pt}\renewcommand{\arraystretch}{1.12}
\noindent\begin{tabular}{@{}p{14mm}p{27mm}|p{13.55mm}|p{13.55mm}|p{13.55mm}|p{13.55mm}|p{13.55mm}|p{13.55mm}|p{13.55mm}|p{13.55mm}|p{13.55mm}|p{13.55mm}|p{13.55mm}|p{13.55mm}|p{13.55mm}|p{13.55mm}|p{13.55mm}|p{13.55mm}|@{}}\hline
Engine & Execution & \multicolumn{8}{c|}{\textbf{Index}} & \multicolumn{8}{c|}{\textbf{No index}}\\
 &  & \multicolumn{8}{c|}{\textbf{2 Gbp}} & \multicolumn{8}{c|}{\textbf{2 Gbp}}\\
 &  & \multicolumn{8}{c|}{\textbf{k=15, w=10}} & \multicolumn{8}{c|}{\textbf{k=15, w=10}}\\
 &  & \multicolumn{4}{c|}{\textbf{DNA4}} & \multicolumn{4}{c|}{\textbf{DNA5}} & \multicolumn{4}{c|}{\textbf{DNA4}} & \multicolumn{4}{c|}{\textbf{DNA5}}\\
 &  & \multicolumn{1}{c|}{\textbf{M7A}} & \multicolumn{1}{c|}{\textbf{M7I}} & \multicolumn{1}{c|}{\textbf{M4}} & \multicolumn{1}{c|}{\textbf{C7G}} & \multicolumn{1}{c|}{\textbf{M7A}} & \multicolumn{1}{c|}{\textbf{M7I}} & \multicolumn{1}{c|}{\textbf{M4}} & \multicolumn{1}{c|}{\textbf{C7G}} & \multicolumn{1}{c|}{\textbf{M7A}} & \multicolumn{1}{c|}{\textbf{M7I}} & \multicolumn{1}{c|}{\textbf{M4}} & \multicolumn{1}{c|}{\textbf{C7G}} & \multicolumn{1}{c|}{\textbf{M7A}} & \multicolumn{1}{c|}{\textbf{M7I}} & \multicolumn{1}{c|}{\textbf{M4}} & \multicolumn{1}{c|}{\textbf{C7G}}\\
 & & {\color{gray}C1} & {\color{gray}C2} & {\color{gray}C3} & {\color{gray}C4} & {\color{gray}C5} & {\color{gray}C6} & {\color{gray}C7} & {\color{gray}C8} & {\color{gray}C9} & {\color{gray}C10} & {\color{gray}C11} & {\color{gray}C12} & {\color{gray}C13} & {\color{gray}C14} & {\color{gray}C15} & {\color{gray}C16}\\
\hline
\rowcolor{BenchBaseline}[0.5pt][0.5pt]\shortstack[l]{{\color{gray}R1}\\\textbf{minimap2}} & \textbf{Baseline} & \hspace*{1pt}\begin{tabular}[t]{@{}c@{}}\rule{0pt}{12pt}11.2s\\[-.2pt]{\fontsize{5.8}{6.4}\selectfont \textcolor{BrickRed}{$ \blacktriangledown\,3.1\times $}}\\[-.2pt]10s\\[-.2pt]{\fontsize{5.8}{6.4}\selectfont \textcolor{BrickRed}{$ \blacktriangledown\,3.6\times $}}\end{tabular} & \hspace*{1pt}\begin{tabular}[t]{@{}c@{}}\rule{0pt}{12pt}10.7s\\[-.2pt]{\fontsize{5.8}{6.4}\selectfont \textcolor{BrickRed}{$ \blacktriangledown\,3.2\times $}}\\[-.2pt]8.6s\\[-.2pt]{\fontsize{5.8}{6.4}\selectfont \textcolor{BrickRed}{$ \blacktriangledown\,3.5\times $}}\end{tabular} & \hspace*{1pt}\begin{tabular}[t]{@{}c@{}}\rule{0pt}{12pt}22.5s\\[-.2pt]{\fontsize{5.8}{6.4}\selectfont \textcolor{BrickRed}{$ \blacktriangledown\,3.6\times $}}\\[-.2pt]18.8s\\[-.2pt]{\fontsize{5.8}{6.4}\selectfont \textcolor{BrickRed}{$ \blacktriangledown\,4.4\times $}}\end{tabular} & \hspace*{1pt}\begin{tabular}[t]{@{}c@{}}\rule{0pt}{12pt}1m\,13.1s\\[-.2pt]{\fontsize{5.8}{6.4}\selectfont \textcolor{BrickRed}{$ \blacktriangledown\,22.5\times $}}\\[-.2pt]1m\,5.7s\\[-.2pt]{\fontsize{5.8}{6.4}\selectfont \textcolor{BrickRed}{$ \blacktriangledown\,27.7\times $}}\end{tabular} & \hspace*{1pt}\begin{tabular}[t]{@{}c@{}}\rule{0pt}{12pt}9.9s\\[-.2pt]{\fontsize{5.8}{6.4}\selectfont \textcolor{BrickRed}{$ \blacktriangledown\,21\% $}}\\[-.2pt]8.8s\\[-.2pt]{\fontsize{5.8}{6.4}\selectfont \textcolor{BrickRed}{$ \blacktriangledown\,1.6\times $}}\end{tabular} & \hspace*{1pt}\begin{tabular}[t]{@{}c@{}}\rule{0pt}{12pt}9.5s\\[-.2pt]{\fontsize{5.8}{6.4}\selectfont \textcolor{BrickRed}{$ \blacktriangledown\,1.4\times $}}\\[-.2pt]7.6s\\[-.2pt]{\fontsize{5.8}{6.4}\selectfont \textcolor{BrickRed}{$ \blacktriangledown\,1.6\times $}}\end{tabular} & \hspace*{1pt}\begin{tabular}[t]{@{}c@{}}\rule{0pt}{12pt}16.5s\\[-.2pt]{\fontsize{5.8}{6.4}\selectfont \textcolor{BrickRed}{$ \blacktriangledown\,1.7\times $}}\\[-.2pt]12.9s\\[-.2pt]{\fontsize{5.8}{6.4}\selectfont \textcolor{BrickRed}{$ \blacktriangledown\,2.0\times $}}\end{tabular} & \hspace*{1pt}\begin{tabular}[t]{@{}c@{}}\rule{0pt}{12pt}28.2s\\[-.2pt]{\fontsize{5.8}{6.4}\selectfont \textcolor{BrickRed}{$ \blacktriangledown\,2.3\times $}}\\[-.2pt]26.2s\\[-.2pt]{\fontsize{5.8}{6.4}\selectfont \textcolor{BrickRed}{$ \blacktriangledown\,3.0\times $}}\end{tabular} & \hspace*{1pt}\begin{tabular}[t]{@{}c@{}}\rule{0pt}{12pt}53s\\[-.2pt]{\fontsize{5.8}{6.4}\selectfont \textcolor{ForestGreen}{$ \blacktriangle\,2.2\times $}}\\[-.2pt]51.9s\\[-.2pt]{\fontsize{5.8}{6.4}\selectfont \textcolor{ForestGreen}{$ \blacktriangle\,2.1\times $}}\end{tabular} & \hspace*{1pt}\begin{tabular}[t]{@{}c@{}}\rule{0pt}{12pt}56.6s\\[-.2pt]{\fontsize{5.8}{6.4}\selectfont \textcolor{ForestGreen}{$ \blacktriangle\,2.1\times $}}\\[-.2pt]54.6s\\[-.2pt]{\fontsize{5.8}{6.4}\selectfont \textcolor{ForestGreen}{$ \blacktriangle\,1.9\times $}}\end{tabular} & \hspace*{1pt}\begin{tabular}[t]{@{}c@{}}\rule{0pt}{12pt}1m\,55.5s\\[-.2pt]{\fontsize{5.8}{6.4}\selectfont \textcolor{ForestGreen}{$ \blacktriangle\,1.8\times $}}\\[-.2pt]1m\,51.8s\\[-.2pt]{\fontsize{5.8}{6.4}\selectfont \textcolor{ForestGreen}{$ \blacktriangle\,1.5\times $}}\end{tabular} & \hspace*{1pt}\begin{tabular}[t]{@{}c@{}}\rule{0pt}{12pt}6m\,1.3s\\[-.2pt]{\fontsize{5.8}{6.4}\selectfont \textcolor{BrickRed}{$ \blacktriangledown\,3.0\times $}}\\[-.2pt]5m\,55.2s\\[-.2pt]{\fontsize{5.8}{6.4}\selectfont \textcolor{BrickRed}{$ \blacktriangledown\,3.1\times $}}\end{tabular} & \hspace*{1pt}\begin{tabular}[t]{@{}c@{}}\rule{0pt}{12pt}48.1s\\[-.2pt]{\fontsize{5.8}{6.4}\selectfont \textcolor{ForestGreen}{$ \blacktriangle\,7.3\times $}}\\[-.2pt]47s\\[-.2pt]{\fontsize{5.8}{6.4}\selectfont \textcolor{ForestGreen}{$ \blacktriangle\,2.4\times $}}\end{tabular} & \hspace*{1pt}\begin{tabular}[t]{@{}c@{}}\rule{0pt}{12pt}44.5s\\[-.2pt]{\fontsize{5.8}{6.4}\selectfont \textcolor{ForestGreen}{$ \blacktriangle\,6.5\times $}}\\[-.2pt]42.6s\\[-.2pt]{\fontsize{5.8}{6.4}\selectfont \textcolor{ForestGreen}{$ \blacktriangle\,2.7\times $}}\end{tabular} & \hspace*{1pt}\begin{tabular}[t]{@{}c@{}}\rule{0pt}{12pt}1m\,30.2s\\[-.2pt]{\fontsize{5.8}{6.4}\selectfont \textcolor{ForestGreen}{$ \blacktriangle\,6.4\times $}}\\[-.2pt]1m\,26.7s\\[-.2pt]{\fontsize{5.8}{6.4}\selectfont \textcolor{ForestGreen}{$ \blacktriangle\,2.2\times $}}\end{tabular} & \hspace*{1pt}\begin{tabular}[t]{@{}c@{}}\rule{0pt}{12pt}2m\,46.4s\\[-.2pt]{\fontsize{5.8}{6.4}\selectfont \textcolor{ForestGreen}{$ \blacktriangle\,4.0\times $}}\\[-.2pt]2m\,42.7s\\[-.2pt]{\fontsize{5.8}{6.4}\selectfont \textcolor{BrickRed}{$ \blacktriangledown\,27\% $}}\end{tabular}\\[3pt]\hline
\BenchmarkSourceLabel{bench:aws-2-k15-w10:row:mpf-simdref-rmq-block}{R9}\shortstack[l]{{\color{gray}R9}\\\textbf{BioseqDB}} & \begin{tabular}[t]{@{}l@{}}\rule{0pt}{12pt}MPF / SIMD-ref\\RMQ / Block\end{tabular} & \hspace*{1pt}\begin{tabular}[t]{@{}c@{}}\rule{0pt}{12pt}3.7s\\[-.2pt]{\fontsize{5.8}{6.4}\selectfont \textcolor{ForestGreen}{$ \blacktriangle\,3.0\times $}}\\[-.2pt]3s\\[-.2pt]{\fontsize{5.8}{6.4}\selectfont \textcolor{ForestGreen}{$ \blacktriangle\,3.4\times $}}\end{tabular} & \hspace*{1pt}\begin{tabular}[t]{@{}c@{}}\rule{0pt}{12pt}3.5s\\[-.2pt]{\fontsize{5.8}{6.4}\selectfont \textcolor{ForestGreen}{$ \blacktriangle\,3.0\times $}}\\[-.2pt]2.7s\\[-.2pt]{\fontsize{5.8}{6.4}\selectfont \textcolor{ForestGreen}{$ \blacktriangle\,3.2\times $}}\end{tabular} & \hspace*{1pt}\begin{tabular}[t]{@{}c@{}}\rule{0pt}{12pt}6.7s\\[-.2pt]{\fontsize{5.8}{6.4}\selectfont \textcolor{ForestGreen}{$ \blacktriangle\,3.4\times $}}\\[-.2pt]4.6s\\[-.2pt]{\fontsize{5.8}{6.4}\selectfont \textcolor{ForestGreen}{$ \blacktriangle\,4.1\times $}}\end{tabular} & \hspace*{1pt}\begin{tabular}[t]{@{}c@{}}\rule{0pt}{12pt}3.5s\\[-.2pt]{\fontsize{5.8}{6.4}\selectfont \textcolor{ForestGreen}{$ \blacktriangle\,21.0\times $}}\\[-.2pt]2.6s\\[-.2pt]{\fontsize{5.8}{6.4}\selectfont \textcolor{ForestGreen}{$ \blacktriangle\,25.5\times $}}\end{tabular} & \hspace*{1pt}\begin{tabular}[t]{@{}c@{}}\rule{0pt}{12pt}13.6s\\[-.2pt]{\fontsize{5.8}{6.4}\selectfont \textcolor{BrickRed}{$ \blacktriangledown\,1.4\times $}}\\[-.2pt]7.8s\\[-.2pt]{\fontsize{5.8}{6.4}\selectfont \textcolor{ForestGreen}{$ \blacktriangle\,12\% $}}\end{tabular} & \hspace*{1pt}\begin{tabular}[t]{@{}c@{}}\rule{0pt}{12pt}13.2s\\[-.2pt]{\fontsize{5.8}{6.4}\selectfont \textcolor{BrickRed}{$ \blacktriangledown\,1.4\times $}}\\[-.2pt]8.1s\\[-.2pt]{\fontsize{5.8}{6.4}\selectfont \textcolor{BrickRed}{$ \blacktriangledown\,6\% $}}\end{tabular} & \hspace*{1pt}\begin{tabular}[t]{@{}c@{}}\rule{0pt}{12pt}25.7s\\[-.2pt]{\fontsize{5.8}{6.4}\selectfont \textcolor{BrickRed}{$ \blacktriangledown\,1.6\times $}}\\[-.2pt]14s\\[-.2pt]{\fontsize{5.8}{6.4}\selectfont \textcolor{BrickRed}{$ \blacktriangledown\,9\% $}}\end{tabular} & \hspace*{1pt}\begin{tabular}[t]{@{}c@{}}\rule{0pt}{12pt}18.3s\\[-.2pt]{\fontsize{5.8}{6.4}\selectfont \textcolor{ForestGreen}{$ \blacktriangle\,1.5\times $}}\\[-.2pt]15.5s\\[-.2pt]{\fontsize{5.8}{6.4}\selectfont \textcolor{ForestGreen}{$ \blacktriangle\,1.7\times $}}\end{tabular} & \hspace*{1pt}\begin{tabular}[t]{@{}c@{}}\rule{0pt}{12pt}2m\,42.5s\\[-.2pt]{\fontsize{5.8}{6.4}\selectfont \textcolor{BrickRed}{$ \blacktriangledown\,3.1\times $}}\\[-.2pt]2m\,17.6s\\[-.2pt]{\fontsize{5.8}{6.4}\selectfont \textcolor{BrickRed}{$ \blacktriangledown\,2.7\times $}}\end{tabular} & \hspace*{1pt}\begin{tabular}[t]{@{}c@{}}\rule{0pt}{12pt}2m\,24.5s\\[-.2pt]{\fontsize{5.8}{6.4}\selectfont \textcolor{BrickRed}{$ \blacktriangledown\,2.6\times $}}\\[-.2pt]2m\,4.8s\\[-.2pt]{\fontsize{5.8}{6.4}\selectfont \textcolor{BrickRed}{$ \blacktriangledown\,2.3\times $}}\end{tabular} & \textcolor{gray}{\ensuremath{\cdot}} & \textcolor{gray}{\ensuremath{\cdot}} & \textcolor{gray}{\ensuremath{\cdot}} & \textcolor{gray}{\ensuremath{\cdot}} & \textcolor{gray}{\ensuremath{\cdot}} & \textcolor{gray}{\ensuremath{\cdot}}\\[3pt]\hline
\rowcolor{BenchPreferred}[0.5pt][0.5pt]\BenchmarkSourceLabel{bench:aws-2-k15-w10:row:mpf-bounded-simd-block}{R10}\shortstack[l]{{\color{gray}R10}} & \begin{tabular}[t]{@{}l@{}}\rule{0pt}{12pt}MPF / Bounded\\SIMD-DP / Block\end{tabular} & \hspace*{1pt}\begin{tabular}[t]{@{}c@{}}\rule{0pt}{12pt}3.7s\\[-.2pt]{\fontsize{5.8}{6.4}\selectfont \textcolor{ForestGreen}{$ \blacktriangle\,3.0\times $}}\\[-.2pt]3s\\[-.2pt]{\fontsize{5.8}{6.4}\selectfont \textcolor{ForestGreen}{$ \blacktriangle\,3.4\times $}}\end{tabular} & \hspace*{1pt}\begin{tabular}[t]{@{}c@{}}\rule{0pt}{12pt}3.5s\\[-.2pt]{\fontsize{5.8}{6.4}\selectfont \textcolor{ForestGreen}{$ \blacktriangle\,3.0\times $}}\\[-.2pt]2.7s\\[-.2pt]{\fontsize{5.8}{6.4}\selectfont \textcolor{ForestGreen}{$ \blacktriangle\,3.2\times $}}\end{tabular} & \hspace*{1pt}\begin{tabular}[t]{@{}c@{}}\rule{0pt}{12pt}6.7s\\[-.2pt]{\fontsize{5.8}{6.4}\selectfont \textcolor{ForestGreen}{$ \blacktriangle\,3.3\times $}}\\[-.2pt]4.5s\\[-.2pt]{\fontsize{5.8}{6.4}\selectfont \textcolor{ForestGreen}{$ \blacktriangle\,4.2\times $}}\end{tabular} & \hspace*{1pt}\begin{tabular}[t]{@{}c@{}}\rule{0pt}{12pt}8.8s\\[-.2pt]{\fontsize{5.8}{6.4}\selectfont \textcolor{ForestGreen}{$ \blacktriangle\,8.3\times $}}\\[-.2pt]7.9s\\[-.2pt]{\fontsize{5.8}{6.4}\selectfont \textcolor{ForestGreen}{$ \blacktriangle\,8.3\times $}}\end{tabular} & \hspace*{1pt}\begin{tabular}[t]{@{}c@{}}\rule{0pt}{12pt}\colorbox{BenchBest}{\strut 8.2s}\\[-.2pt]{\fontsize{5.8}{6.4}\selectfont \textcolor{ForestGreen}{$ \blacktriangle\,17\% $}}\\[-.2pt]5.8s\\[-.2pt]{\fontsize{5.8}{6.4}\selectfont \textcolor{ForestGreen}{$ \blacktriangle\,1.5\times $}}\end{tabular} & \hspace*{1pt}\begin{tabular}[t]{@{}c@{}}\rule{0pt}{12pt}\colorbox{BenchBest}{\strut 6.6s}\\[-.2pt]{\fontsize{5.8}{6.4}\selectfont \textcolor{ForestGreen}{$ \blacktriangle\,1.4\times $}}\\[-.2pt]\colorbox{BenchBest}{\strut 4.6s}\\[-.2pt]{\fontsize{5.8}{6.4}\selectfont \textcolor{ForestGreen}{$ \blacktriangle\,1.6\times $}}\end{tabular} & \hspace*{1pt}\begin{tabular}[t]{@{}c@{}}\rule{0pt}{12pt}13.7s\\[-.2pt]{\fontsize{5.8}{6.4}\selectfont \textcolor{ForestGreen}{$ \blacktriangle\,17\% $}}\\[-.2pt]8.8s\\[-.2pt]{\fontsize{5.8}{6.4}\selectfont \textcolor{ForestGreen}{$ \blacktriangle\,1.5\times $}}\end{tabular} & \hspace*{1pt}\begin{tabular}[t]{@{}c@{}}\rule{0pt}{12pt}20.2s\\[-.2pt]{\fontsize{5.8}{6.4}\selectfont \textcolor{ForestGreen}{$ \blacktriangle\,28\% $}}\\[-.2pt]16.8s\\[-.2pt]{\fontsize{5.8}{6.4}\selectfont \textcolor{ForestGreen}{$ \blacktriangle\,1.6\times $}}\end{tabular} & \hspace*{1pt}\begin{tabular}[t]{@{}c@{}}\rule{0pt}{12pt}1m\,56.4s\\[-.2pt]{\fontsize{5.8}{6.4}\selectfont \textcolor{BrickRed}{$ \blacktriangledown\,2.2\times $}}\\[-.2pt]1m\,48.1s\\[-.2pt]{\fontsize{5.8}{6.4}\selectfont \textcolor{BrickRed}{$ \blacktriangledown\,2.1\times $}}\end{tabular} & \hspace*{1pt}\begin{tabular}[t]{@{}c@{}}\rule{0pt}{12pt}1m\,56.2s\\[-.2pt]{\fontsize{5.8}{6.4}\selectfont \textcolor{BrickRed}{$ \blacktriangledown\,2.1\times $}}\\[-.2pt]1m\,46.2s\\[-.2pt]{\fontsize{5.8}{6.4}\selectfont \textcolor{BrickRed}{$ \blacktriangledown\,1.9\times $}}\end{tabular} & \hspace*{1pt}\begin{tabular}[t]{@{}c@{}}\rule{0pt}{12pt}3m\,42.9s\\[-.2pt]{\fontsize{5.8}{6.4}\selectfont \textcolor{BrickRed}{$ \blacktriangledown\,1.9\times $}}\\[-.2pt]3m\,25s\\[-.2pt]{\fontsize{5.8}{6.4}\selectfont \textcolor{BrickRed}{$ \blacktriangledown\,1.8\times $}}\end{tabular} & \hspace*{1pt}\begin{tabular}[t]{@{}c@{}}\rule{0pt}{12pt}2m\,5.6s\\[-.2pt]{\fontsize{5.8}{6.4}\selectfont \textcolor{ForestGreen}{$ \blacktriangle\,2.9\times $}}\\[-.2pt]1m\,59.7s\\[-.2pt]{\fontsize{5.8}{6.4}\selectfont \textcolor{ForestGreen}{$ \blacktriangle\,3.0\times $}}\end{tabular} & \hspace*{1pt}\begin{tabular}[t]{@{}c@{}}\rule{0pt}{12pt}5m\,52.1s\\[-.2pt]{\fontsize{5.8}{6.4}\selectfont \textcolor{BrickRed}{$ \blacktriangledown\,7.3\times $}}\\[-.2pt]1m\,53.5s\\[-.2pt]{\fontsize{5.8}{6.4}\selectfont \textcolor{BrickRed}{$ \blacktriangledown\,2.4\times $}}\end{tabular} & \hspace*{1pt}\begin{tabular}[t]{@{}c@{}}\rule{0pt}{12pt}4m\,49.2s\\[-.2pt]{\fontsize{5.8}{6.4}\selectfont \textcolor{BrickRed}{$ \blacktriangledown\,6.5\times $}}\\[-.2pt]1m\,56.8s\\[-.2pt]{\fontsize{5.8}{6.4}\selectfont \textcolor{BrickRed}{$ \blacktriangledown\,2.7\times $}}\end{tabular} & \hspace*{1pt}\begin{tabular}[t]{@{}c@{}}\rule{0pt}{12pt}9m\,39.5s\\[-.2pt]{\fontsize{5.8}{6.4}\selectfont \textcolor{BrickRed}{$ \blacktriangledown\,6.4\times $}}\\[-.2pt]3m\,18.1s\\[-.2pt]{\fontsize{5.8}{6.4}\selectfont \textcolor{BrickRed}{$ \blacktriangledown\,2.3\times $}}\end{tabular} & \hspace*{1pt}\begin{tabular}[t]{@{}c@{}}\rule{0pt}{12pt}11m\,24.8s\\[-.2pt]{\fontsize{5.8}{6.4}\selectfont \textcolor{BrickRed}{$ \blacktriangledown\,4.1\times $}}\\[-.2pt]2m\,20.4s\\[-.2pt]{\fontsize{5.8}{6.4}\selectfont \textcolor{ForestGreen}{$ \blacktriangle\,14\% $}}\end{tabular}\\[3pt]\hline
\BenchmarkSourceLabel{bench:aws-2-k15-w10:row:mpf-scalar-rmq-block}{R11}\shortstack[l]{{\color{gray}R11}} & \begin{tabular}[t]{@{}l@{}}\rule{0pt}{12pt}MPF / Scalar\\RMQ / Block\end{tabular} & \hspace*{1pt}\begin{tabular}[t]{@{}c@{}}\rule{0pt}{12pt}4s\\[-.2pt]{\fontsize{5.8}{6.4}\selectfont \textcolor{ForestGreen}{$ \blacktriangle\,2.8\times $}}\\[-.2pt]3.1s\\[-.2pt]{\fontsize{5.8}{6.4}\selectfont \textcolor{ForestGreen}{$ \blacktriangle\,3.2\times $}}\end{tabular} & \hspace*{1pt}\begin{tabular}[t]{@{}c@{}}\rule{0pt}{12pt}3.7s\\[-.2pt]{\fontsize{5.8}{6.4}\selectfont \textcolor{ForestGreen}{$ \blacktriangle\,2.9\times $}}\\[-.2pt]2.8s\\[-.2pt]{\fontsize{5.8}{6.4}\selectfont \textcolor{ForestGreen}{$ \blacktriangle\,3.1\times $}}\end{tabular} & \hspace*{1pt}\begin{tabular}[t]{@{}c@{}}\rule{0pt}{12pt}7.1s\\[-.2pt]{\fontsize{5.8}{6.4}\selectfont \textcolor{ForestGreen}{$ \blacktriangle\,3.2\times $}}\\[-.2pt]4.8s\\[-.2pt]{\fontsize{5.8}{6.4}\selectfont \textcolor{ForestGreen}{$ \blacktriangle\,3.9\times $}}\end{tabular} & \hspace*{1pt}\begin{tabular}[t]{@{}c@{}}\rule{0pt}{12pt}3.5s\\[-.2pt]{\fontsize{5.8}{6.4}\selectfont \textcolor{ForestGreen}{$ \blacktriangle\,20.6\times $}}\\[-.2pt]2.6s\\[-.2pt]{\fontsize{5.8}{6.4}\selectfont \textcolor{ForestGreen}{$ \blacktriangle\,25.3\times $}}\end{tabular} & \hspace*{1pt}\begin{tabular}[t]{@{}c@{}}\rule{0pt}{12pt}10s\\[-.2pt]{\fontsize{5.8}{6.4}\selectfont \textcolor{gray}{$ \blacktriangleright\,1\% $}}\\[-.2pt]5.9s\\[-.2pt]{\fontsize{5.8}{6.4}\selectfont \textcolor{ForestGreen}{$ \blacktriangle\,1.5\times $}}\end{tabular} & \hspace*{1pt}\begin{tabular}[t]{@{}c@{}}\rule{0pt}{12pt}8.7s\\[-.2pt]{\fontsize{5.8}{6.4}\selectfont \textcolor{ForestGreen}{$ \blacktriangle\,8\% $}}\\[-.2pt]5.7s\\[-.2pt]{\fontsize{5.8}{6.4}\selectfont \textcolor{ForestGreen}{$ \blacktriangle\,25\% $}}\end{tabular} & \hspace*{1pt}\begin{tabular}[t]{@{}c@{}}\rule{0pt}{12pt}\colorbox{BenchBest}{\strut 9.5s}\\[-.2pt]{\fontsize{5.8}{6.4}\selectfont \textcolor{ForestGreen}{$ \blacktriangle\,1.7\times $}}\\[-.2pt]\colorbox{BenchBest}{\strut 6.5s}\\[-.2pt]{\fontsize{5.8}{6.4}\selectfont \textcolor{ForestGreen}{$ \blacktriangle\,2.0\times $}}\end{tabular} & \hspace*{1pt}\begin{tabular}[t]{@{}c@{}}\rule{0pt}{12pt}23s\\[-.2pt]{\fontsize{5.8}{6.4}\selectfont \textcolor{ForestGreen}{$ \blacktriangle\,18\% $}}\\[-.2pt]19.8s\\[-.2pt]{\fontsize{5.8}{6.4}\selectfont \textcolor{ForestGreen}{$ \blacktriangle\,25\% $}}\end{tabular} & \hspace*{1pt}\begin{tabular}[t]{@{}c@{}}\rule{0pt}{12pt}2m\,51.5s\\[-.2pt]{\fontsize{5.8}{6.4}\selectfont \textcolor{BrickRed}{$ \blacktriangledown\,3.2\times $}}\\[-.2pt]2m\,29.1s\\[-.2pt]{\fontsize{5.8}{6.4}\selectfont \textcolor{BrickRed}{$ \blacktriangledown\,2.9\times $}}\end{tabular} & \hspace*{1pt}\begin{tabular}[t]{@{}c@{}}\rule{0pt}{12pt}3m\,4.8s\\[-.2pt]{\fontsize{5.8}{6.4}\selectfont \textcolor{BrickRed}{$ \blacktriangledown\,3.3\times $}}\\[-.2pt]2m\,44.5s\\[-.2pt]{\fontsize{5.8}{6.4}\selectfont \textcolor{BrickRed}{$ \blacktriangledown\,3.0\times $}}\end{tabular} & \textcolor{gray}{\ensuremath{\cdot}} & \textcolor{gray}{\ensuremath{\cdot}} & \hspace*{1pt}\begin{tabular}[t]{@{}c@{}}\rule{0pt}{12pt}7m\,6.3s\\[-.2pt]{\fontsize{5.8}{6.4}\selectfont \textcolor{BrickRed}{$ \blacktriangledown\,8.9\times $}}\\[-.2pt]2m\,44.6s\\[-.2pt]{\fontsize{5.8}{6.4}\selectfont \textcolor{BrickRed}{$ \blacktriangledown\,3.5\times $}}\end{tabular} & \hspace*{1pt}\begin{tabular}[t]{@{}c@{}}\rule{0pt}{12pt}6m\,52.5s\\[-.2pt]{\fontsize{5.8}{6.4}\selectfont \textcolor{BrickRed}{$ \blacktriangledown\,9.3\times $}}\\[-.2pt]3m\,7.3s\\[-.2pt]{\fontsize{5.8}{6.4}\selectfont \textcolor{BrickRed}{$ \blacktriangledown\,4.4\times $}}\end{tabular} & \textcolor{gray}{\ensuremath{\cdot}} & \textcolor{gray}{\ensuremath{\cdot}}\\[3pt]\hline
\BenchmarkSourceLabel{bench:aws-2-k15-w10:row:mpf-twostack-rmq-block}{R12}\shortstack[l]{{\color{gray}R12}} & \begin{tabular}[t]{@{}l@{}}\rule{0pt}{12pt}MPF / Two-stack\\RMQ / Block\end{tabular} & \hspace*{1pt}\begin{tabular}[t]{@{}c@{}}\rule{0pt}{12pt}3.8s\\[-.2pt]{\fontsize{5.8}{6.4}\selectfont \textcolor{ForestGreen}{$ \blacktriangle\,2.9\times $}}\\[-.2pt]3s\\[-.2pt]{\fontsize{5.8}{6.4}\selectfont \textcolor{ForestGreen}{$ \blacktriangle\,3.3\times $}}\end{tabular} & \hspace*{1pt}\begin{tabular}[t]{@{}c@{}}\rule{0pt}{12pt}3.6s\\[-.2pt]{\fontsize{5.8}{6.4}\selectfont \textcolor{ForestGreen}{$ \blacktriangle\,3.0\times $}}\\[-.2pt]2.7s\\[-.2pt]{\fontsize{5.8}{6.4}\selectfont \textcolor{ForestGreen}{$ \blacktriangle\,3.2\times $}}\end{tabular} & \hspace*{1pt}\begin{tabular}[t]{@{}c@{}}\rule{0pt}{12pt}6.8s\\[-.2pt]{\fontsize{5.8}{6.4}\selectfont \textcolor{ForestGreen}{$ \blacktriangle\,3.3\times $}}\\[-.2pt]4.7s\\[-.2pt]{\fontsize{5.8}{6.4}\selectfont \textcolor{ForestGreen}{$ \blacktriangle\,4.0\times $}}\end{tabular} & \hspace*{1pt}\begin{tabular}[t]{@{}c@{}}\rule{0pt}{12pt}3.6s\\[-.2pt]{\fontsize{5.8}{6.4}\selectfont \textcolor{ForestGreen}{$ \blacktriangle\,20.1\times $}}\\[-.2pt]2.7s\\[-.2pt]{\fontsize{5.8}{6.4}\selectfont \textcolor{ForestGreen}{$ \blacktriangle\,24.2\times $}}\end{tabular} & \hspace*{1pt}\begin{tabular}[t]{@{}c@{}}\rule{0pt}{12pt}9.6s\\[-.2pt]{\fontsize{5.8}{6.4}\selectfont \textcolor{gray}{$ \blacktriangleright\,3\% $}}\\[-.2pt]5.5s\\[-.2pt]{\fontsize{5.8}{6.4}\selectfont \textcolor{ForestGreen}{$ \blacktriangle\,1.6\times $}}\end{tabular} & \hspace*{1pt}\begin{tabular}[t]{@{}c@{}}\rule{0pt}{12pt}8.2s\\[-.2pt]{\fontsize{5.8}{6.4}\selectfont \textcolor{ForestGreen}{$ \blacktriangle\,14\% $}}\\[-.2pt]5s\\[-.2pt]{\fontsize{5.8}{6.4}\selectfont \textcolor{ForestGreen}{$ \blacktriangle\,1.5\times $}}\end{tabular} & \hspace*{1pt}\begin{tabular}[t]{@{}c@{}}\rule{0pt}{12pt}12.8s\\[-.2pt]{\fontsize{5.8}{6.4}\selectfont \textcolor{ForestGreen}{$ \blacktriangle\,23\% $}}\\[-.2pt]8.6s\\[-.2pt]{\fontsize{5.8}{6.4}\selectfont \textcolor{ForestGreen}{$ \blacktriangle\,1.5\times $}}\end{tabular} & \hspace*{1pt}\begin{tabular}[t]{@{}c@{}}\rule{0pt}{12pt}14.3s\\[-.2pt]{\fontsize{5.8}{6.4}\selectfont \textcolor{ForestGreen}{$ \blacktriangle\,2.0\times $}}\\[-.2pt]11.1s\\[-.2pt]{\fontsize{5.8}{6.4}\selectfont \textcolor{ForestGreen}{$ \blacktriangle\,2.4\times $}}\end{tabular} & \hspace*{1pt}\begin{tabular}[t]{@{}c@{}}\rule{0pt}{12pt}2m\,38.4s\\[-.2pt]{\fontsize{5.8}{6.4}\selectfont \textcolor{BrickRed}{$ \blacktriangledown\,3.0\times $}}\\[-.2pt]2m\,16.3s\\[-.2pt]{\fontsize{5.8}{6.4}\selectfont \textcolor{BrickRed}{$ \blacktriangledown\,2.6\times $}}\end{tabular} & \hspace*{1pt}\begin{tabular}[t]{@{}c@{}}\rule{0pt}{12pt}2m\,38s\\[-.2pt]{\fontsize{5.8}{6.4}\selectfont \textcolor{BrickRed}{$ \blacktriangledown\,2.8\times $}}\\[-.2pt]2m\,17.8s\\[-.2pt]{\fontsize{5.8}{6.4}\selectfont \textcolor{BrickRed}{$ \blacktriangledown\,2.5\times $}}\end{tabular} & \textcolor{gray}{\ensuremath{\cdot}} & \textcolor{gray}{\ensuremath{\cdot}} & \hspace*{1pt}\begin{tabular}[t]{@{}c@{}}\rule{0pt}{12pt}6m\,27.6s\\[-.2pt]{\fontsize{5.8}{6.4}\selectfont \textcolor{BrickRed}{$ \blacktriangledown\,8.1\times $}}\\[-.2pt]2m\,7.9s\\[-.2pt]{\fontsize{5.8}{6.4}\selectfont \textcolor{BrickRed}{$ \blacktriangledown\,2.7\times $}}\end{tabular} & \hspace*{1pt}\begin{tabular}[t]{@{}c@{}}\rule{0pt}{12pt}6m\,2.1s\\[-.2pt]{\fontsize{5.8}{6.4}\selectfont \textcolor{BrickRed}{$ \blacktriangledown\,8.1\times $}}\\[-.2pt]2m\,19.9s\\[-.2pt]{\fontsize{5.8}{6.4}\selectfont \textcolor{BrickRed}{$ \blacktriangledown\,3.3\times $}}\end{tabular} & \textcolor{gray}{\ensuremath{\cdot}} & \textcolor{gray}{\ensuremath{\cdot}}\\[3pt]\hline
\BenchmarkSourceLabel{bench:aws-2-k15-w10:row:mpf-scalar-simd-block}{R13}\shortstack[l]{{\color{gray}R13}} & \begin{tabular}[t]{@{}l@{}}\rule{0pt}{12pt}MPF / Scalar\\SIMD-DP / Block\end{tabular} & \hspace*{1pt}\begin{tabular}[t]{@{}c@{}}\rule{0pt}{12pt}3.9s\\[-.2pt]{\fontsize{5.8}{6.4}\selectfont \textcolor{ForestGreen}{$ \blacktriangle\,2.9\times $}}\\[-.2pt]3.1s\\[-.2pt]{\fontsize{5.8}{6.4}\selectfont \textcolor{ForestGreen}{$ \blacktriangle\,3.3\times $}}\end{tabular} & \hspace*{1pt}\begin{tabular}[t]{@{}c@{}}\rule{0pt}{12pt}4.5s\\[-.2pt]{\fontsize{5.8}{6.4}\selectfont \textcolor{ForestGreen}{$ \blacktriangle\,2.4\times $}}\\[-.2pt]3.6s\\[-.2pt]{\fontsize{5.8}{6.4}\selectfont \textcolor{ForestGreen}{$ \blacktriangle\,2.4\times $}}\end{tabular} & \hspace*{1pt}\begin{tabular}[t]{@{}c@{}}\rule{0pt}{12pt}6.6s\\[-.2pt]{\fontsize{5.8}{6.4}\selectfont \textcolor{ForestGreen}{$ \blacktriangle\,3.4\times $}}\\[-.2pt]4.7s\\[-.2pt]{\fontsize{5.8}{6.4}\selectfont \textcolor{ForestGreen}{$ \blacktriangle\,4.0\times $}}\end{tabular} & \hspace*{1pt}\begin{tabular}[t]{@{}c@{}}\rule{0pt}{12pt}9.1s\\[-.2pt]{\fontsize{5.8}{6.4}\selectfont \textcolor{ForestGreen}{$ \blacktriangle\,8.1\times $}}\\[-.2pt]8.1s\\[-.2pt]{\fontsize{5.8}{6.4}\selectfont \textcolor{ForestGreen}{$ \blacktriangle\,8.1\times $}}\end{tabular} & \hspace*{1pt}\begin{tabular}[t]{@{}c@{}}\rule{0pt}{12pt}10.4s\\[-.2pt]{\fontsize{5.8}{6.4}\selectfont \textcolor{BrickRed}{$ \blacktriangledown\,5\% $}}\\[-.2pt]6s\\[-.2pt]{\fontsize{5.8}{6.4}\selectfont \textcolor{ForestGreen}{$ \blacktriangle\,1.5\times $}}\end{tabular} & \hspace*{1pt}\begin{tabular}[t]{@{}c@{}}\rule{0pt}{12pt}11.1s\\[-.2pt]{\fontsize{5.8}{6.4}\selectfont \textcolor{BrickRed}{$ \blacktriangledown\,17\% $}}\\[-.2pt]7.8s\\[-.2pt]{\fontsize{5.8}{6.4}\selectfont \textcolor{gray}{$ \blacktriangleright\,2\% $}}\end{tabular} & \hspace*{1pt}\begin{tabular}[t]{@{}c@{}}\rule{0pt}{12pt}16.6s\\[-.2pt]{\fontsize{5.8}{6.4}\selectfont \textcolor{gray}{$ \blacktriangleright\,1\% $}}\\[-.2pt]10.1s\\[-.2pt]{\fontsize{5.8}{6.4}\selectfont \textcolor{ForestGreen}{$ \blacktriangle\,22\% $}}\end{tabular} & \hspace*{1pt}\begin{tabular}[t]{@{}c@{}}\rule{0pt}{12pt}14.3s\\[-.2pt]{\fontsize{5.8}{6.4}\selectfont \textcolor{ForestGreen}{$ \blacktriangle\,2.0\times $}}\\[-.2pt]12.1s\\[-.2pt]{\fontsize{5.8}{6.4}\selectfont \textcolor{ForestGreen}{$ \blacktriangle\,2.2\times $}}\end{tabular} & \hspace*{1pt}\begin{tabular}[t]{@{}c@{}}\rule{0pt}{12pt}3m\,8.5s\\[-.2pt]{\fontsize{5.8}{6.4}\selectfont \textcolor{BrickRed}{$ \blacktriangledown\,3.6\times $}}\\[-.2pt]2m\,46.3s\\[-.2pt]{\fontsize{5.8}{6.4}\selectfont \textcolor{BrickRed}{$ \blacktriangledown\,3.2\times $}}\end{tabular} & \hspace*{1pt}\begin{tabular}[t]{@{}c@{}}\rule{0pt}{12pt}3m\,1.2s\\[-.2pt]{\fontsize{5.8}{6.4}\selectfont \textcolor{BrickRed}{$ \blacktriangledown\,3.2\times $}}\\[-.2pt]2m\,41.3s\\[-.2pt]{\fontsize{5.8}{6.4}\selectfont \textcolor{BrickRed}{$ \blacktriangledown\,3.0\times $}}\end{tabular} & \textcolor{gray}{\ensuremath{\cdot}} & \hspace*{1pt}\begin{tabular}[t]{@{}c@{}}\rule{0pt}{12pt}3m\,2.9s\\[-.2pt]{\fontsize{5.8}{6.4}\selectfont \textcolor{ForestGreen}{$ \blacktriangle\,2.0\times $}}\\[-.2pt]2m\,46s\\[-.2pt]{\fontsize{5.8}{6.4}\selectfont \textcolor{ForestGreen}{$ \blacktriangle\,2.1\times $}}\end{tabular} & \hspace*{1pt}\begin{tabular}[t]{@{}c@{}}\rule{0pt}{12pt}7m\,10.9s\\[-.2pt]{\fontsize{5.8}{6.4}\selectfont \textcolor{BrickRed}{$ \blacktriangledown\,9.0\times $}}\\[-.2pt]2m\,49.7s\\[-.2pt]{\fontsize{5.8}{6.4}\selectfont \textcolor{BrickRed}{$ \blacktriangledown\,3.6\times $}}\end{tabular} & \hspace*{1pt}\begin{tabular}[t]{@{}c@{}}\rule{0pt}{12pt}6m\,59s\\[-.2pt]{\fontsize{5.8}{6.4}\selectfont \textcolor{BrickRed}{$ \blacktriangledown\,9.4\times $}}\\[-.2pt]3m\,11.3s\\[-.2pt]{\fontsize{5.8}{6.4}\selectfont \textcolor{BrickRed}{$ \blacktriangledown\,4.5\times $}}\end{tabular} & \textcolor{gray}{\ensuremath{\cdot}} & \hspace*{1pt}\begin{tabular}[t]{@{}c@{}}\rule{0pt}{12pt}12m\,21.6s\\[-.2pt]{\fontsize{5.8}{6.4}\selectfont \textcolor{BrickRed}{$ \blacktriangledown\,4.5\times $}}\\[-.2pt]3m\,1s\\[-.2pt]{\fontsize{5.8}{6.4}\selectfont \textcolor{BrickRed}{$ \blacktriangledown\,11\% $}}\end{tabular}\\[3pt]\hline
\BenchmarkSourceLabel{bench:aws-2-k15-w10:row:mpf-bounded-simd-ksw2}{R14}\shortstack[l]{{\color{gray}R14}} & \begin{tabular}[t]{@{}l@{}}\rule{0pt}{12pt}MPF / Bounded\\SIMD-DP / KSW2\end{tabular} & \hspace*{1pt}\begin{tabular}[t]{@{}c@{}}\rule{0pt}{12pt}5.2s\\[-.2pt]{\fontsize{5.8}{6.4}\selectfont \textcolor{ForestGreen}{$ \blacktriangle\,2.1\times $}}\\[-.2pt]\colorbox{BenchBest}{\strut 2.8s}\\[-.2pt]{\fontsize{5.8}{6.4}\selectfont \textcolor{ForestGreen}{$ \blacktriangle\,3.6\times $}}\end{tabular} & \hspace*{1pt}\begin{tabular}[t]{@{}c@{}}\rule{0pt}{12pt}4.7s\\[-.2pt]{\fontsize{5.8}{6.4}\selectfont \textcolor{ForestGreen}{$ \blacktriangle\,2.3\times $}}\\[-.2pt]2.5s\\[-.2pt]{\fontsize{5.8}{6.4}\selectfont \textcolor{ForestGreen}{$ \blacktriangle\,3.5\times $}}\end{tabular} & \hspace*{1pt}\begin{tabular}[t]{@{}c@{}}\rule{0pt}{12pt}9.9s\\[-.2pt]{\fontsize{5.8}{6.4}\selectfont \textcolor{ForestGreen}{$ \blacktriangle\,2.3\times $}}\\[-.2pt]4.3s\\[-.2pt]{\fontsize{5.8}{6.4}\selectfont \textcolor{ForestGreen}{$ \blacktriangle\,4.4\times $}}\end{tabular} & \hspace*{1pt}\begin{tabular}[t]{@{}c@{}}\rule{0pt}{12pt}7.4s\\[-.2pt]{\fontsize{5.8}{6.4}\selectfont \textcolor{ForestGreen}{$ \blacktriangle\,9.9\times $}}\\[-.2pt]4.4s\\[-.2pt]{\fontsize{5.8}{6.4}\selectfont \textcolor{ForestGreen}{$ \blacktriangle\,15.1\times $}}\end{tabular} & \hspace*{1pt}\begin{tabular}[t]{@{}c@{}}\rule{0pt}{12pt}10.2s\\[-.2pt]{\fontsize{5.8}{6.4}\selectfont \textcolor{gray}{$ \blacktriangleright\,3\% $}}\\[-.2pt]5.9s\\[-.2pt]{\fontsize{5.8}{6.4}\selectfont \textcolor{ForestGreen}{$ \blacktriangle\,1.5\times $}}\end{tabular} & \hspace*{1pt}\begin{tabular}[t]{@{}c@{}}\rule{0pt}{12pt}12.5s\\[-.2pt]{\fontsize{5.8}{6.4}\selectfont \textcolor{BrickRed}{$ \blacktriangledown\,1.3\times $}}\\[-.2pt]8.7s\\[-.2pt]{\fontsize{5.8}{6.4}\selectfont \textcolor{BrickRed}{$ \blacktriangledown\,14\% $}}\end{tabular} & \hspace*{1pt}\begin{tabular}[t]{@{}c@{}}\rule{0pt}{12pt}14.1s\\[-.2pt]{\fontsize{5.8}{6.4}\selectfont \textcolor{ForestGreen}{$ \blacktriangle\,14\% $}}\\[-.2pt]8.5s\\[-.2pt]{\fontsize{5.8}{6.4}\selectfont \textcolor{ForestGreen}{$ \blacktriangle\,1.5\times $}}\end{tabular} & \hspace*{1pt}\begin{tabular}[t]{@{}c@{}}\rule{0pt}{12pt}21.5s\\[-.2pt]{\fontsize{5.8}{6.4}\selectfont \textcolor{ForestGreen}{$ \blacktriangle\,24\% $}}\\[-.2pt]17.6s\\[-.2pt]{\fontsize{5.8}{6.4}\selectfont \textcolor{ForestGreen}{$ \blacktriangle\,1.5\times $}}\end{tabular} & \hspace*{1pt}\begin{tabular}[t]{@{}c@{}}\rule{0pt}{12pt}2m\,11.6s\\[-.2pt]{\fontsize{5.8}{6.4}\selectfont \textcolor{BrickRed}{$ \blacktriangledown\,2.5\times $}}\\[-.2pt]2m\,1.1s\\[-.2pt]{\fontsize{5.8}{6.4}\selectfont \textcolor{BrickRed}{$ \blacktriangledown\,2.3\times $}}\end{tabular} & \hspace*{1pt}\begin{tabular}[t]{@{}c@{}}\rule{0pt}{12pt}2m\,7.5s\\[-.2pt]{\fontsize{5.8}{6.4}\selectfont \textcolor{BrickRed}{$ \blacktriangledown\,2.3\times $}}\\[-.2pt]1m\,54.8s\\[-.2pt]{\fontsize{5.8}{6.4}\selectfont \textcolor{BrickRed}{$ \blacktriangledown\,2.1\times $}}\end{tabular} & \textcolor{gray}{\ensuremath{\cdot}} & \hspace*{1pt}\begin{tabular}[t]{@{}c@{}}\rule{0pt}{12pt}2m\,20.7s\\[-.2pt]{\fontsize{5.8}{6.4}\selectfont \textcolor{ForestGreen}{$ \blacktriangle\,2.6\times $}}\\[-.2pt]2m\,10.4s\\[-.2pt]{\fontsize{5.8}{6.4}\selectfont \textcolor{ForestGreen}{$ \blacktriangle\,2.7\times $}}\end{tabular} & \hspace*{1pt}\begin{tabular}[t]{@{}c@{}}\rule{0pt}{12pt}6m\,9.1s\\[-.2pt]{\fontsize{5.8}{6.4}\selectfont \textcolor{BrickRed}{$ \blacktriangledown\,7.7\times $}}\\[-.2pt]2m\,5.6s\\[-.2pt]{\fontsize{5.8}{6.4}\selectfont \textcolor{BrickRed}{$ \blacktriangledown\,2.7\times $}}\end{tabular} & \hspace*{1pt}\begin{tabular}[t]{@{}c@{}}\rule{0pt}{12pt}5m\,32.3s\\[-.2pt]{\fontsize{5.8}{6.4}\selectfont \textcolor{BrickRed}{$ \blacktriangledown\,7.5\times $}}\\[-.2pt]1m\,56.1s\\[-.2pt]{\fontsize{5.8}{6.4}\selectfont \textcolor{BrickRed}{$ \blacktriangledown\,2.7\times $}}\end{tabular} & \textcolor{gray}{\ensuremath{\cdot}} & \hspace*{1pt}\begin{tabular}[t]{@{}c@{}}\rule{0pt}{12pt}11m\,23.6s\\[-.2pt]{\fontsize{5.8}{6.4}\selectfont \textcolor{BrickRed}{$ \blacktriangledown\,4.1\times $}}\\[-.2pt]2m\,22.7s\\[-.2pt]{\fontsize{5.8}{6.4}\selectfont \textcolor{ForestGreen}{$ \blacktriangle\,12\% $}}\end{tabular}\\[3pt]\hline
\BenchmarkSourceLabel{bench:aws-2-k15-w10:row:mpf-simdref-scalar-ksw2}{R15}\shortstack[l]{{\color{gray}R15}} & \begin{tabular}[t]{@{}l@{}}\rule{0pt}{12pt}MPF / SIMD-ref\\Scalar-DP / KSW2\end{tabular} & \hspace*{1pt}\begin{tabular}[t]{@{}c@{}}\rule{0pt}{12pt}5.1s\\[-.2pt]{\fontsize{5.8}{6.4}\selectfont \textcolor{ForestGreen}{$ \blacktriangle\,2.2\times $}}\\[-.2pt]2.8s\\[-.2pt]{\fontsize{5.8}{6.4}\selectfont \textcolor{ForestGreen}{$ \blacktriangle\,3.6\times $}}\end{tabular} & \hspace*{1pt}\begin{tabular}[t]{@{}c@{}}\rule{0pt}{12pt}4.7s\\[-.2pt]{\fontsize{5.8}{6.4}\selectfont \textcolor{ForestGreen}{$ \blacktriangle\,2.3\times $}}\\[-.2pt]2.5s\\[-.2pt]{\fontsize{5.8}{6.4}\selectfont \textcolor{ForestGreen}{$ \blacktriangle\,3.4\times $}}\end{tabular} & \hspace*{1pt}\begin{tabular}[t]{@{}c@{}}\rule{0pt}{12pt}9.8s\\[-.2pt]{\fontsize{5.8}{6.4}\selectfont \textcolor{ForestGreen}{$ \blacktriangle\,2.3\times $}}\\[-.2pt]4.3s\\[-.2pt]{\fontsize{5.8}{6.4}\selectfont \textcolor{ForestGreen}{$ \blacktriangle\,4.3\times $}}\end{tabular} & \hspace*{1pt}\begin{tabular}[t]{@{}c@{}}\rule{0pt}{12pt}7.2s\\[-.2pt]{\fontsize{5.8}{6.4}\selectfont \textcolor{ForestGreen}{$ \blacktriangle\,10.1\times $}}\\[-.2pt]4.4s\\[-.2pt]{\fontsize{5.8}{6.4}\selectfont \textcolor{ForestGreen}{$ \blacktriangle\,15.1\times $}}\end{tabular} & \hspace*{1pt}\begin{tabular}[t]{@{}c@{}}\rule{0pt}{12pt}12.8s\\[-.2pt]{\fontsize{5.8}{6.4}\selectfont \textcolor{BrickRed}{$ \blacktriangledown\,29\% $}}\\[-.2pt]7.3s\\[-.2pt]{\fontsize{5.8}{6.4}\selectfont \textcolor{ForestGreen}{$ \blacktriangle\,17\% $}}\end{tabular} & \hspace*{1pt}\begin{tabular}[t]{@{}c@{}}\rule{0pt}{12pt}13s\\[-.2pt]{\fontsize{5.8}{6.4}\selectfont \textcolor{BrickRed}{$ \blacktriangledown\,1.4\times $}}\\[-.2pt]7.3s\\[-.2pt]{\fontsize{5.8}{6.4}\selectfont \textcolor{gray}{$ \blacktriangleright\,4\% $}}\end{tabular} & \hspace*{1pt}\begin{tabular}[t]{@{}c@{}}\rule{0pt}{12pt}24.5s\\[-.2pt]{\fontsize{5.8}{6.4}\selectfont \textcolor{BrickRed}{$ \blacktriangledown\,1.5\times $}}\\[-.2pt]13.4s\\[-.2pt]{\fontsize{5.8}{6.4}\selectfont \textcolor{gray}{$ \blacktriangleright\,4\% $}}\end{tabular} & \hspace*{1pt}\begin{tabular}[t]{@{}c@{}}\rule{0pt}{12pt}16.2s\\[-.2pt]{\fontsize{5.8}{6.4}\selectfont \textcolor{ForestGreen}{$ \blacktriangle\,1.7\times $}}\\[-.2pt]12.9s\\[-.2pt]{\fontsize{5.8}{6.4}\selectfont \textcolor{ForestGreen}{$ \blacktriangle\,2.0\times $}}\end{tabular} & \hspace*{1pt}\begin{tabular}[t]{@{}c@{}}\rule{0pt}{12pt}2m\,52.8s\\[-.2pt]{\fontsize{5.8}{6.4}\selectfont \textcolor{BrickRed}{$ \blacktriangledown\,3.3\times $}}\\[-.2pt]2m\,25.9s\\[-.2pt]{\fontsize{5.8}{6.4}\selectfont \textcolor{BrickRed}{$ \blacktriangledown\,2.8\times $}}\end{tabular} & \hspace*{1pt}\begin{tabular}[t]{@{}c@{}}\rule{0pt}{12pt}2m\,31.6s\\[-.2pt]{\fontsize{5.8}{6.4}\selectfont \textcolor{BrickRed}{$ \blacktriangledown\,2.7\times $}}\\[-.2pt]2m\,9.7s\\[-.2pt]{\fontsize{5.8}{6.4}\selectfont \textcolor{BrickRed}{$ \blacktriangledown\,2.4\times $}}\end{tabular} & \textcolor{gray}{\ensuremath{\cdot}} & \textcolor{gray}{\ensuremath{\cdot}} & \textcolor{gray}{\ensuremath{\cdot}} & \textcolor{gray}{\ensuremath{\cdot}} & \textcolor{gray}{\ensuremath{\cdot}} & \textcolor{gray}{\ensuremath{\cdot}}\\[3pt]\hline
\end{tabular}\par}
\smallskip{\footnotesize\noindent Each cell: A and B times, $\cdot$: no measurement. ---\textsuperscript{g}: GPU unavailable. n/a: no comparison.\par
\noindent Blue: one lowest A and one lowest B measurement for each column for each configuration.\par
\noindent Gray: minimap2. Yellow: my preferred MPF / Bounded / SIMD-DP / Block configuration.\par
\noindent MPF/MPS: minimizers per fragment/sequence. SIMD-ref: naive SIMD scan. Bounded: bounded SIMD algorithm.\par
}
\clearpage
\noindent{\large\bfseries Table \thetable: (D): AWS: 2 Gbp, k=15, w=10\hfill 3/7}\par\smallskip
{\fontsize{7.1}{8.1}\selectfont\setlength{\tabcolsep}{1pt}\setlength{\fboxsep}{2pt}\renewcommand{\arraystretch}{1.12}
\noindent\begin{tabular}{@{}p{14mm}p{27mm}|p{13.55mm}|p{13.55mm}|p{13.55mm}|p{13.55mm}|p{13.55mm}|p{13.55mm}|p{13.55mm}|p{13.55mm}|p{13.55mm}|p{13.55mm}|p{13.55mm}|p{13.55mm}|p{13.55mm}|p{13.55mm}|p{13.55mm}|p{13.55mm}|@{}}\hline
Engine & Execution & \multicolumn{8}{c|}{\textbf{Index}} & \multicolumn{8}{c|}{\textbf{No index}}\\
 &  & \multicolumn{8}{c|}{\textbf{2 Gbp}} & \multicolumn{8}{c|}{\textbf{2 Gbp}}\\
 &  & \multicolumn{8}{c|}{\textbf{k=15, w=10}} & \multicolumn{8}{c|}{\textbf{k=15, w=10}}\\
 &  & \multicolumn{4}{c|}{\textbf{DNA4}} & \multicolumn{4}{c|}{\textbf{DNA5}} & \multicolumn{4}{c|}{\textbf{DNA4}} & \multicolumn{4}{c|}{\textbf{DNA5}}\\
 &  & \multicolumn{1}{c|}{\textbf{M7A}} & \multicolumn{1}{c|}{\textbf{M7I}} & \multicolumn{1}{c|}{\textbf{M4}} & \multicolumn{1}{c|}{\textbf{C7G}} & \multicolumn{1}{c|}{\textbf{M7A}} & \multicolumn{1}{c|}{\textbf{M7I}} & \multicolumn{1}{c|}{\textbf{M4}} & \multicolumn{1}{c|}{\textbf{C7G}} & \multicolumn{1}{c|}{\textbf{M7A}} & \multicolumn{1}{c|}{\textbf{M7I}} & \multicolumn{1}{c|}{\textbf{M4}} & \multicolumn{1}{c|}{\textbf{C7G}} & \multicolumn{1}{c|}{\textbf{M7A}} & \multicolumn{1}{c|}{\textbf{M7I}} & \multicolumn{1}{c|}{\textbf{M4}} & \multicolumn{1}{c|}{\textbf{C7G}}\\
 & & {\color{gray}C1} & {\color{gray}C2} & {\color{gray}C3} & {\color{gray}C4} & {\color{gray}C5} & {\color{gray}C6} & {\color{gray}C7} & {\color{gray}C8} & {\color{gray}C9} & {\color{gray}C10} & {\color{gray}C11} & {\color{gray}C12} & {\color{gray}C13} & {\color{gray}C14} & {\color{gray}C15} & {\color{gray}C16}\\
\hline
\rowcolor{BenchBaseline}[0.5pt][0.5pt]\shortstack[l]{{\color{gray}R1}\\\textbf{minimap2}} & \textbf{Baseline} & \hspace*{1pt}\begin{tabular}[t]{@{}c@{}}\rule{0pt}{12pt}11.2s\\[-.2pt]{\fontsize{5.8}{6.4}\selectfont \textcolor{BrickRed}{$ \blacktriangledown\,3.1\times $}}\\[-.2pt]10s\\[-.2pt]{\fontsize{5.8}{6.4}\selectfont \textcolor{BrickRed}{$ \blacktriangledown\,3.6\times $}}\end{tabular} & \hspace*{1pt}\begin{tabular}[t]{@{}c@{}}\rule{0pt}{12pt}10.7s\\[-.2pt]{\fontsize{5.8}{6.4}\selectfont \textcolor{BrickRed}{$ \blacktriangledown\,3.2\times $}}\\[-.2pt]8.6s\\[-.2pt]{\fontsize{5.8}{6.4}\selectfont \textcolor{BrickRed}{$ \blacktriangledown\,3.5\times $}}\end{tabular} & \hspace*{1pt}\begin{tabular}[t]{@{}c@{}}\rule{0pt}{12pt}22.5s\\[-.2pt]{\fontsize{5.8}{6.4}\selectfont \textcolor{BrickRed}{$ \blacktriangledown\,3.6\times $}}\\[-.2pt]18.8s\\[-.2pt]{\fontsize{5.8}{6.4}\selectfont \textcolor{BrickRed}{$ \blacktriangledown\,4.4\times $}}\end{tabular} & \hspace*{1pt}\begin{tabular}[t]{@{}c@{}}\rule{0pt}{12pt}1m\,13.1s\\[-.2pt]{\fontsize{5.8}{6.4}\selectfont \textcolor{BrickRed}{$ \blacktriangledown\,22.5\times $}}\\[-.2pt]1m\,5.7s\\[-.2pt]{\fontsize{5.8}{6.4}\selectfont \textcolor{BrickRed}{$ \blacktriangledown\,27.7\times $}}\end{tabular} & \hspace*{1pt}\begin{tabular}[t]{@{}c@{}}\rule{0pt}{12pt}9.9s\\[-.2pt]{\fontsize{5.8}{6.4}\selectfont \textcolor{BrickRed}{$ \blacktriangledown\,21\% $}}\\[-.2pt]8.8s\\[-.2pt]{\fontsize{5.8}{6.4}\selectfont \textcolor{BrickRed}{$ \blacktriangledown\,1.6\times $}}\end{tabular} & \hspace*{1pt}\begin{tabular}[t]{@{}c@{}}\rule{0pt}{12pt}9.5s\\[-.2pt]{\fontsize{5.8}{6.4}\selectfont \textcolor{BrickRed}{$ \blacktriangledown\,1.4\times $}}\\[-.2pt]7.6s\\[-.2pt]{\fontsize{5.8}{6.4}\selectfont \textcolor{BrickRed}{$ \blacktriangledown\,1.6\times $}}\end{tabular} & \hspace*{1pt}\begin{tabular}[t]{@{}c@{}}\rule{0pt}{12pt}16.5s\\[-.2pt]{\fontsize{5.8}{6.4}\selectfont \textcolor{BrickRed}{$ \blacktriangledown\,1.7\times $}}\\[-.2pt]12.9s\\[-.2pt]{\fontsize{5.8}{6.4}\selectfont \textcolor{BrickRed}{$ \blacktriangledown\,2.0\times $}}\end{tabular} & \hspace*{1pt}\begin{tabular}[t]{@{}c@{}}\rule{0pt}{12pt}28.2s\\[-.2pt]{\fontsize{5.8}{6.4}\selectfont \textcolor{BrickRed}{$ \blacktriangledown\,2.3\times $}}\\[-.2pt]26.2s\\[-.2pt]{\fontsize{5.8}{6.4}\selectfont \textcolor{BrickRed}{$ \blacktriangledown\,3.0\times $}}\end{tabular} & \hspace*{1pt}\begin{tabular}[t]{@{}c@{}}\rule{0pt}{12pt}53s\\[-.2pt]{\fontsize{5.8}{6.4}\selectfont \textcolor{ForestGreen}{$ \blacktriangle\,2.2\times $}}\\[-.2pt]51.9s\\[-.2pt]{\fontsize{5.8}{6.4}\selectfont \textcolor{ForestGreen}{$ \blacktriangle\,2.1\times $}}\end{tabular} & \hspace*{1pt}\begin{tabular}[t]{@{}c@{}}\rule{0pt}{12pt}56.6s\\[-.2pt]{\fontsize{5.8}{6.4}\selectfont \textcolor{ForestGreen}{$ \blacktriangle\,2.1\times $}}\\[-.2pt]54.6s\\[-.2pt]{\fontsize{5.8}{6.4}\selectfont \textcolor{ForestGreen}{$ \blacktriangle\,1.9\times $}}\end{tabular} & \hspace*{1pt}\begin{tabular}[t]{@{}c@{}}\rule{0pt}{12pt}1m\,55.5s\\[-.2pt]{\fontsize{5.8}{6.4}\selectfont \textcolor{ForestGreen}{$ \blacktriangle\,1.8\times $}}\\[-.2pt]1m\,51.8s\\[-.2pt]{\fontsize{5.8}{6.4}\selectfont \textcolor{ForestGreen}{$ \blacktriangle\,1.5\times $}}\end{tabular} & \hspace*{1pt}\begin{tabular}[t]{@{}c@{}}\rule{0pt}{12pt}6m\,1.3s\\[-.2pt]{\fontsize{5.8}{6.4}\selectfont \textcolor{BrickRed}{$ \blacktriangledown\,3.0\times $}}\\[-.2pt]5m\,55.2s\\[-.2pt]{\fontsize{5.8}{6.4}\selectfont \textcolor{BrickRed}{$ \blacktriangledown\,3.1\times $}}\end{tabular} & \hspace*{1pt}\begin{tabular}[t]{@{}c@{}}\rule{0pt}{12pt}48.1s\\[-.2pt]{\fontsize{5.8}{6.4}\selectfont \textcolor{ForestGreen}{$ \blacktriangle\,7.3\times $}}\\[-.2pt]47s\\[-.2pt]{\fontsize{5.8}{6.4}\selectfont \textcolor{ForestGreen}{$ \blacktriangle\,2.4\times $}}\end{tabular} & \hspace*{1pt}\begin{tabular}[t]{@{}c@{}}\rule{0pt}{12pt}44.5s\\[-.2pt]{\fontsize{5.8}{6.4}\selectfont \textcolor{ForestGreen}{$ \blacktriangle\,6.5\times $}}\\[-.2pt]42.6s\\[-.2pt]{\fontsize{5.8}{6.4}\selectfont \textcolor{ForestGreen}{$ \blacktriangle\,2.7\times $}}\end{tabular} & \hspace*{1pt}\begin{tabular}[t]{@{}c@{}}\rule{0pt}{12pt}1m\,30.2s\\[-.2pt]{\fontsize{5.8}{6.4}\selectfont \textcolor{ForestGreen}{$ \blacktriangle\,6.4\times $}}\\[-.2pt]1m\,26.7s\\[-.2pt]{\fontsize{5.8}{6.4}\selectfont \textcolor{ForestGreen}{$ \blacktriangle\,2.2\times $}}\end{tabular} & \hspace*{1pt}\begin{tabular}[t]{@{}c@{}}\rule{0pt}{12pt}2m\,46.4s\\[-.2pt]{\fontsize{5.8}{6.4}\selectfont \textcolor{ForestGreen}{$ \blacktriangle\,4.0\times $}}\\[-.2pt]2m\,42.7s\\[-.2pt]{\fontsize{5.8}{6.4}\selectfont \textcolor{BrickRed}{$ \blacktriangledown\,27\% $}}\end{tabular}\\[3pt]\hline
\BenchmarkSourceLabel{bench:aws-2-k15-w10:row:mpf-bounded-scalar-ksw2}{R16}\shortstack[l]{{\color{gray}R16}\\\textbf{BioseqDB}} & \begin{tabular}[t]{@{}l@{}}\rule{0pt}{12pt}MPF / Bounded\\Scalar-DP / KSW2\end{tabular} & \hspace*{1pt}\begin{tabular}[t]{@{}c@{}}\rule{0pt}{12pt}5.2s\\[-.2pt]{\fontsize{5.8}{6.4}\selectfont \textcolor{ForestGreen}{$ \blacktriangle\,2.1\times $}}\\[-.2pt]2.8s\\[-.2pt]{\fontsize{5.8}{6.4}\selectfont \textcolor{ForestGreen}{$ \blacktriangle\,3.6\times $}}\end{tabular} & \hspace*{1pt}\begin{tabular}[t]{@{}c@{}}\rule{0pt}{12pt}4.7s\\[-.2pt]{\fontsize{5.8}{6.4}\selectfont \textcolor{ForestGreen}{$ \blacktriangle\,2.3\times $}}\\[-.2pt]\colorbox{BenchBest}{\strut 2.5s}\\[-.2pt]{\fontsize{5.8}{6.4}\selectfont \textcolor{ForestGreen}{$ \blacktriangle\,3.5\times $}}\end{tabular} & \hspace*{1pt}\begin{tabular}[t]{@{}c@{}}\rule{0pt}{12pt}9.7s\\[-.2pt]{\fontsize{5.8}{6.4}\selectfont \textcolor{ForestGreen}{$ \blacktriangle\,2.3\times $}}\\[-.2pt]4.3s\\[-.2pt]{\fontsize{5.8}{6.4}\selectfont \textcolor{ForestGreen}{$ \blacktriangle\,4.4\times $}}\end{tabular} & \hspace*{1pt}\begin{tabular}[t]{@{}c@{}}\rule{0pt}{12pt}7.2s\\[-.2pt]{\fontsize{5.8}{6.4}\selectfont \textcolor{ForestGreen}{$ \blacktriangle\,10.1\times $}}\\[-.2pt]4.3s\\[-.2pt]{\fontsize{5.8}{6.4}\selectfont \textcolor{ForestGreen}{$ \blacktriangle\,15.4\times $}}\end{tabular} & \hspace*{1pt}\begin{tabular}[t]{@{}c@{}}\rule{0pt}{12pt}9.9s\\[-.2pt]{\fontsize{5.8}{6.4}\selectfont \textcolor{gray}{$ \blacktriangleright\,0\% $}}\\[-.2pt]5.8s\\[-.2pt]{\fontsize{5.8}{6.4}\selectfont \textcolor{ForestGreen}{$ \blacktriangle\,1.5\times $}}\end{tabular} & \hspace*{1pt}\begin{tabular}[t]{@{}c@{}}\rule{0pt}{12pt}8.5s\\[-.2pt]{\fontsize{5.8}{6.4}\selectfont \textcolor{ForestGreen}{$ \blacktriangle\,11\% $}}\\[-.2pt]5.1s\\[-.2pt]{\fontsize{5.8}{6.4}\selectfont \textcolor{ForestGreen}{$ \blacktriangle\,1.5\times $}}\end{tabular} & \hspace*{1pt}\begin{tabular}[t]{@{}c@{}}\rule{0pt}{12pt}15.2s\\[-.2pt]{\fontsize{5.8}{6.4}\selectfont \textcolor{ForestGreen}{$ \blacktriangle\,8\% $}}\\[-.2pt]8.5s\\[-.2pt]{\fontsize{5.8}{6.4}\selectfont \textcolor{ForestGreen}{$ \blacktriangle\,1.5\times $}}\end{tabular} & \hspace*{1pt}\begin{tabular}[t]{@{}c@{}}\rule{0pt}{12pt}13s\\[-.2pt]{\fontsize{5.8}{6.4}\selectfont \textcolor{ForestGreen}{$ \blacktriangle\,2.2\times $}}\\[-.2pt]9.4s\\[-.2pt]{\fontsize{5.8}{6.4}\selectfont \textcolor{ForestGreen}{$ \blacktriangle\,2.8\times $}}\end{tabular} & \hspace*{1pt}\begin{tabular}[t]{@{}c@{}}\rule{0pt}{12pt}2m\,8.2s\\[-.2pt]{\fontsize{5.8}{6.4}\selectfont \textcolor{BrickRed}{$ \blacktriangledown\,2.4\times $}}\\[-.2pt]1m\,58.3s\\[-.2pt]{\fontsize{5.8}{6.4}\selectfont \textcolor{BrickRed}{$ \blacktriangledown\,2.3\times $}}\end{tabular} & \hspace*{1pt}\begin{tabular}[t]{@{}c@{}}\rule{0pt}{12pt}2m\,3.7s\\[-.2pt]{\fontsize{5.8}{6.4}\selectfont \textcolor{BrickRed}{$ \blacktriangledown\,2.2\times $}}\\[-.2pt]1m\,51.2s\\[-.2pt]{\fontsize{5.8}{6.4}\selectfont \textcolor{BrickRed}{$ \blacktriangledown\,2.0\times $}}\end{tabular} & \textcolor{gray}{\ensuremath{\cdot}} & \hspace*{1pt}\begin{tabular}[t]{@{}c@{}}\rule{0pt}{12pt}2m\,25.1s\\[-.2pt]{\fontsize{5.8}{6.4}\selectfont \textcolor{ForestGreen}{$ \blacktriangle\,2.5\times $}}\\[-.2pt]2m\,14.9s\\[-.2pt]{\fontsize{5.8}{6.4}\selectfont \textcolor{ForestGreen}{$ \blacktriangle\,2.6\times $}}\end{tabular} & \hspace*{1pt}\begin{tabular}[t]{@{}c@{}}\rule{0pt}{12pt}6m\,7.9s\\[-.2pt]{\fontsize{5.8}{6.4}\selectfont \textcolor{BrickRed}{$ \blacktriangledown\,7.7\times $}}\\[-.2pt]2m\,3.4s\\[-.2pt]{\fontsize{5.8}{6.4}\selectfont \textcolor{BrickRed}{$ \blacktriangledown\,2.6\times $}}\end{tabular} & \hspace*{1pt}\begin{tabular}[t]{@{}c@{}}\rule{0pt}{12pt}5m\,36.8s\\[-.2pt]{\fontsize{5.8}{6.4}\selectfont \textcolor{BrickRed}{$ \blacktriangledown\,7.6\times $}}\\[-.2pt]2m\,6.4s\\[-.2pt]{\fontsize{5.8}{6.4}\selectfont \textcolor{BrickRed}{$ \blacktriangledown\,3.0\times $}}\end{tabular} & \hspace*{1pt}\begin{tabular}[t]{@{}c@{}}\rule{0pt}{12pt}9m\,39.3s\\[-.2pt]{\fontsize{5.8}{6.4}\selectfont \textcolor{BrickRed}{$ \blacktriangledown\,6.4\times $}}\\[-.2pt]3m\,13.2s\\[-.2pt]{\fontsize{5.8}{6.4}\selectfont \textcolor{BrickRed}{$ \blacktriangledown\,2.2\times $}}\end{tabular} & \hspace*{1pt}\begin{tabular}[t]{@{}c@{}}\rule{0pt}{12pt}11m\,21s\\[-.2pt]{\fontsize{5.8}{6.4}\selectfont \textcolor{BrickRed}{$ \blacktriangledown\,4.1\times $}}\\[-.2pt]2m\,19.3s\\[-.2pt]{\fontsize{5.8}{6.4}\selectfont \textcolor{ForestGreen}{$ \blacktriangle\,14\% $}}\end{tabular}\\[3pt]\hline
\BenchmarkSourceLabel{bench:aws-2-k15-w10:row:mpf-twostack-scalar-ksw2}{R17}\shortstack[l]{{\color{gray}R17}} & \begin{tabular}[t]{@{}l@{}}\rule{0pt}{12pt}MPF / Two-stack\\Scalar-DP / KSW2\end{tabular} & \hspace*{1pt}\begin{tabular}[t]{@{}c@{}}\rule{0pt}{12pt}5.2s\\[-.2pt]{\fontsize{5.8}{6.4}\selectfont \textcolor{ForestGreen}{$ \blacktriangle\,2.1\times $}}\\[-.2pt]2.8s\\[-.2pt]{\fontsize{5.8}{6.4}\selectfont \textcolor{ForestGreen}{$ \blacktriangle\,3.5\times $}}\end{tabular} & \hspace*{1pt}\begin{tabular}[t]{@{}c@{}}\rule{0pt}{12pt}4.7s\\[-.2pt]{\fontsize{5.8}{6.4}\selectfont \textcolor{ForestGreen}{$ \blacktriangle\,2.3\times $}}\\[-.2pt]2.5s\\[-.2pt]{\fontsize{5.8}{6.4}\selectfont \textcolor{ForestGreen}{$ \blacktriangle\,3.4\times $}}\end{tabular} & \hspace*{1pt}\begin{tabular}[t]{@{}c@{}}\rule{0pt}{12pt}9.5s\\[-.2pt]{\fontsize{5.8}{6.4}\selectfont \textcolor{ForestGreen}{$ \blacktriangle\,2.4\times $}}\\[-.2pt]4.5s\\[-.2pt]{\fontsize{5.8}{6.4}\selectfont \textcolor{ForestGreen}{$ \blacktriangle\,4.2\times $}}\end{tabular} & \hspace*{1pt}\begin{tabular}[t]{@{}c@{}}\rule{0pt}{12pt}7.2s\\[-.2pt]{\fontsize{5.8}{6.4}\selectfont \textcolor{ForestGreen}{$ \blacktriangle\,10.1\times $}}\\[-.2pt]4.4s\\[-.2pt]{\fontsize{5.8}{6.4}\selectfont \textcolor{ForestGreen}{$ \blacktriangle\,14.9\times $}}\end{tabular} & \hspace*{1pt}\begin{tabular}[t]{@{}c@{}}\rule{0pt}{12pt}9.7s\\[-.2pt]{\fontsize{5.8}{6.4}\selectfont \textcolor{gray}{$ \blacktriangleright\,2\% $}}\\[-.2pt]5.6s\\[-.2pt]{\fontsize{5.8}{6.4}\selectfont \textcolor{ForestGreen}{$ \blacktriangle\,1.6\times $}}\end{tabular} & \hspace*{1pt}\begin{tabular}[t]{@{}c@{}}\rule{0pt}{12pt}7.8s\\[-.2pt]{\fontsize{5.8}{6.4}\selectfont \textcolor{ForestGreen}{$ \blacktriangle\,18\% $}}\\[-.2pt]4.7s\\[-.2pt]{\fontsize{5.8}{6.4}\selectfont \textcolor{ForestGreen}{$ \blacktriangle\,1.6\times $}}\end{tabular} & \hspace*{1pt}\begin{tabular}[t]{@{}c@{}}\rule{0pt}{12pt}16.5s\\[-.2pt]{\fontsize{5.8}{6.4}\selectfont \textcolor{gray}{$ \blacktriangleright\,0\% $}}\\[-.2pt]8.9s\\[-.2pt]{\fontsize{5.8}{6.4}\selectfont \textcolor{ForestGreen}{$ \blacktriangle\,1.5\times $}}\end{tabular} & \hspace*{1pt}\begin{tabular}[t]{@{}c@{}}\rule{0pt}{12pt}12.8s\\[-.2pt]{\fontsize{5.8}{6.4}\selectfont \textcolor{ForestGreen}{$ \blacktriangle\,2.2\times $}}\\[-.2pt]9s\\[-.2pt]{\fontsize{5.8}{6.4}\selectfont \textcolor{ForestGreen}{$ \blacktriangle\,2.9\times $}}\end{tabular} & \hspace*{1pt}\begin{tabular}[t]{@{}c@{}}\rule{0pt}{12pt}2m\,46.6s\\[-.2pt]{\fontsize{5.8}{6.4}\selectfont \textcolor{BrickRed}{$ \blacktriangledown\,3.1\times $}}\\[-.2pt]2m\,20.5s\\[-.2pt]{\fontsize{5.8}{6.4}\selectfont \textcolor{BrickRed}{$ \blacktriangledown\,2.7\times $}}\end{tabular} & \hspace*{1pt}\begin{tabular}[t]{@{}c@{}}\rule{0pt}{12pt}2m\,46.1s\\[-.2pt]{\fontsize{5.8}{6.4}\selectfont \textcolor{BrickRed}{$ \blacktriangledown\,2.9\times $}}\\[-.2pt]2m\,22.8s\\[-.2pt]{\fontsize{5.8}{6.4}\selectfont \textcolor{BrickRed}{$ \blacktriangledown\,2.6\times $}}\end{tabular} & \textcolor{gray}{\ensuremath{\cdot}} & \textcolor{gray}{\ensuremath{\cdot}} & \hspace*{1pt}\begin{tabular}[t]{@{}c@{}}\rule{0pt}{12pt}6m\,29s\\[-.2pt]{\fontsize{5.8}{6.4}\selectfont \textcolor{BrickRed}{$ \blacktriangledown\,8.1\times $}}\\[-.2pt]2m\,7.9s\\[-.2pt]{\fontsize{5.8}{6.4}\selectfont \textcolor{BrickRed}{$ \blacktriangledown\,2.7\times $}}\end{tabular} & \hspace*{1pt}\begin{tabular}[t]{@{}c@{}}\rule{0pt}{12pt}5m\,48.7s\\[-.2pt]{\fontsize{5.8}{6.4}\selectfont \textcolor{BrickRed}{$ \blacktriangledown\,7.8\times $}}\\[-.2pt]2m\,5.1s\\[-.2pt]{\fontsize{5.8}{6.4}\selectfont \textcolor{BrickRed}{$ \blacktriangledown\,2.9\times $}}\end{tabular} & \textcolor{gray}{\ensuremath{\cdot}} & \textcolor{gray}{\ensuremath{\cdot}}\\[3pt]\hline
\BenchmarkSourceLabel{bench:aws-2-k15-w10:row:mpf-bounded-rmq-ksw2}{R18}\shortstack[l]{{\color{gray}R18}} & \begin{tabular}[t]{@{}l@{}}\rule{0pt}{12pt}MPF / Bounded\\RMQ / KSW2\end{tabular} & \hspace*{1pt}\begin{tabular}[t]{@{}c@{}}\rule{0pt}{12pt}5.3s\\[-.2pt]{\fontsize{5.8}{6.4}\selectfont \textcolor{ForestGreen}{$ \blacktriangle\,2.1\times $}}\\[-.2pt]2.8s\\[-.2pt]{\fontsize{5.8}{6.4}\selectfont \textcolor{ForestGreen}{$ \blacktriangle\,3.6\times $}}\end{tabular} & \hspace*{1pt}\begin{tabular}[t]{@{}c@{}}\rule{0pt}{12pt}4.8s\\[-.2pt]{\fontsize{5.8}{6.4}\selectfont \textcolor{ForestGreen}{$ \blacktriangle\,2.2\times $}}\\[-.2pt]2.5s\\[-.2pt]{\fontsize{5.8}{6.4}\selectfont \textcolor{ForestGreen}{$ \blacktriangle\,3.5\times $}}\end{tabular} & \hspace*{1pt}\begin{tabular}[t]{@{}c@{}}\rule{0pt}{12pt}9.8s\\[-.2pt]{\fontsize{5.8}{6.4}\selectfont \textcolor{ForestGreen}{$ \blacktriangle\,2.3\times $}}\\[-.2pt]\colorbox{BenchBest}{\strut 4.2s}\\[-.2pt]{\fontsize{5.8}{6.4}\selectfont \textcolor{ForestGreen}{$ \blacktriangle\,4.4\times $}}\end{tabular} & \hspace*{1pt}\begin{tabular}[t]{@{}c@{}}\rule{0pt}{12pt}7.7s\\[-.2pt]{\fontsize{5.8}{6.4}\selectfont \textcolor{ForestGreen}{$ \blacktriangle\,9.5\times $}}\\[-.2pt]4.6s\\[-.2pt]{\fontsize{5.8}{6.4}\selectfont \textcolor{ForestGreen}{$ \blacktriangle\,14.2\times $}}\end{tabular} & \hspace*{1pt}\begin{tabular}[t]{@{}c@{}}\rule{0pt}{12pt}12s\\[-.2pt]{\fontsize{5.8}{6.4}\selectfont \textcolor{BrickRed}{$ \blacktriangledown\,21\% $}}\\[-.2pt]7.7s\\[-.2pt]{\fontsize{5.8}{6.4}\selectfont \textcolor{ForestGreen}{$ \blacktriangle\,13\% $}}\end{tabular} & \hspace*{1pt}\begin{tabular}[t]{@{}c@{}}\rule{0pt}{12pt}11.5s\\[-.2pt]{\fontsize{5.8}{6.4}\selectfont \textcolor{BrickRed}{$ \blacktriangledown\,21\% $}}\\[-.2pt]7.6s\\[-.2pt]{\fontsize{5.8}{6.4}\selectfont \textcolor{gray}{$ \blacktriangleright\,0\% $}}\end{tabular} & \hspace*{1pt}\begin{tabular}[t]{@{}c@{}}\rule{0pt}{12pt}17.2s\\[-.2pt]{\fontsize{5.8}{6.4}\selectfont \textcolor{gray}{$ \blacktriangleright\,4\% $}}\\[-.2pt]10s\\[-.2pt]{\fontsize{5.8}{6.4}\selectfont \textcolor{ForestGreen}{$ \blacktriangle\,22\% $}}\end{tabular} & \hspace*{1pt}\begin{tabular}[t]{@{}c@{}}\rule{0pt}{12pt}13.2s\\[-.2pt]{\fontsize{5.8}{6.4}\selectfont \textcolor{ForestGreen}{$ \blacktriangle\,2.1\times $}}\\[-.2pt]9.4s\\[-.2pt]{\fontsize{5.8}{6.4}\selectfont \textcolor{ForestGreen}{$ \blacktriangle\,2.8\times $}}\end{tabular} & \hspace*{1pt}\begin{tabular}[t]{@{}c@{}}\rule{0pt}{12pt}2m\,17.7s\\[-.2pt]{\fontsize{5.8}{6.4}\selectfont \textcolor{BrickRed}{$ \blacktriangledown\,2.6\times $}}\\[-.2pt]2m\,7s\\[-.2pt]{\fontsize{5.8}{6.4}\selectfont \textcolor{BrickRed}{$ \blacktriangledown\,2.4\times $}}\end{tabular} & \hspace*{1pt}\begin{tabular}[t]{@{}c@{}}\rule{0pt}{12pt}1m\,59.4s\\[-.2pt]{\fontsize{5.8}{6.4}\selectfont \textcolor{BrickRed}{$ \blacktriangledown\,2.1\times $}}\\[-.2pt]1m\,47.5s\\[-.2pt]{\fontsize{5.8}{6.4}\selectfont \textcolor{BrickRed}{$ \blacktriangledown\,2.0\times $}}\end{tabular} & \textcolor{gray}{\ensuremath{\cdot}} & \hspace*{1pt}\begin{tabular}[t]{@{}c@{}}\rule{0pt}{12pt}2m\,24s\\[-.2pt]{\fontsize{5.8}{6.4}\selectfont \textcolor{ForestGreen}{$ \blacktriangle\,2.5\times $}}\\[-.2pt]2m\,13.9s\\[-.2pt]{\fontsize{5.8}{6.4}\selectfont \textcolor{ForestGreen}{$ \blacktriangle\,2.7\times $}}\end{tabular} & \hspace*{1pt}\begin{tabular}[t]{@{}c@{}}\rule{0pt}{12pt}6m\,7.2s\\[-.2pt]{\fontsize{5.8}{6.4}\selectfont \textcolor{BrickRed}{$ \blacktriangledown\,7.6\times $}}\\[-.2pt]2m\,3.6s\\[-.2pt]{\fontsize{5.8}{6.4}\selectfont \textcolor{BrickRed}{$ \blacktriangledown\,2.6\times $}}\end{tabular} & \hspace*{1pt}\begin{tabular}[t]{@{}c@{}}\rule{0pt}{12pt}5m\,26.8s\\[-.2pt]{\fontsize{5.8}{6.4}\selectfont \textcolor{BrickRed}{$ \blacktriangledown\,7.3\times $}}\\[-.2pt]1m\,56.2s\\[-.2pt]{\fontsize{5.8}{6.4}\selectfont \textcolor{BrickRed}{$ \blacktriangledown\,2.7\times $}}\end{tabular} & \textcolor{gray}{\ensuremath{\cdot}} & \textcolor{gray}{\ensuremath{\cdot}}\\[3pt]\hline
\BenchmarkSourceLabel{bench:aws-2-k15-w10:row:mpf-scalar-scalar-ksw2}{R19}\shortstack[l]{{\color{gray}R19}} & \begin{tabular}[t]{@{}l@{}}\rule{0pt}{12pt}MPF / Scalar\\Scalar-DP / KSW2\end{tabular} & \hspace*{1pt}\begin{tabular}[t]{@{}c@{}}\rule{0pt}{12pt}5.2s\\[-.2pt]{\fontsize{5.8}{6.4}\selectfont \textcolor{ForestGreen}{$ \blacktriangle\,2.1\times $}}\\[-.2pt]2.9s\\[-.2pt]{\fontsize{5.8}{6.4}\selectfont \textcolor{ForestGreen}{$ \blacktriangle\,3.5\times $}}\end{tabular} & \hspace*{1pt}\begin{tabular}[t]{@{}c@{}}\rule{0pt}{12pt}4.8s\\[-.2pt]{\fontsize{5.8}{6.4}\selectfont \textcolor{ForestGreen}{$ \blacktriangle\,2.2\times $}}\\[-.2pt]2.6s\\[-.2pt]{\fontsize{5.8}{6.4}\selectfont \textcolor{ForestGreen}{$ \blacktriangle\,3.3\times $}}\end{tabular} & \hspace*{1pt}\begin{tabular}[t]{@{}c@{}}\rule{0pt}{12pt}9.5s\\[-.2pt]{\fontsize{5.8}{6.4}\selectfont \textcolor{ForestGreen}{$ \blacktriangle\,2.4\times $}}\\[-.2pt]4.5s\\[-.2pt]{\fontsize{5.8}{6.4}\selectfont \textcolor{ForestGreen}{$ \blacktriangle\,4.2\times $}}\end{tabular} & \hspace*{1pt}\begin{tabular}[t]{@{}c@{}}\rule{0pt}{12pt}8.6s\\[-.2pt]{\fontsize{5.8}{6.4}\selectfont \textcolor{ForestGreen}{$ \blacktriangle\,8.5\times $}}\\[-.2pt]5.7s\\[-.2pt]{\fontsize{5.8}{6.4}\selectfont \textcolor{ForestGreen}{$ \blacktriangle\,11.6\times $}}\end{tabular} & \hspace*{1pt}\begin{tabular}[t]{@{}c@{}}\rule{0pt}{12pt}10.1s\\[-.2pt]{\fontsize{5.8}{6.4}\selectfont \textcolor{gray}{$ \blacktriangleright\,2\% $}}\\[-.2pt]5.9s\\[-.2pt]{\fontsize{5.8}{6.4}\selectfont \textcolor{ForestGreen}{$ \blacktriangle\,1.5\times $}}\end{tabular} & \hspace*{1pt}\begin{tabular}[t]{@{}c@{}}\rule{0pt}{12pt}8.8s\\[-.2pt]{\fontsize{5.8}{6.4}\selectfont \textcolor{ForestGreen}{$ \blacktriangle\,8\% $}}\\[-.2pt]5.3s\\[-.2pt]{\fontsize{5.8}{6.4}\selectfont \textcolor{ForestGreen}{$ \blacktriangle\,1.4\times $}}\end{tabular} & \hspace*{1pt}\begin{tabular}[t]{@{}c@{}}\rule{0pt}{12pt}15.9s\\[-.2pt]{\fontsize{5.8}{6.4}\selectfont \textcolor{gray}{$ \blacktriangleright\,4\% $}}\\[-.2pt]9.2s\\[-.2pt]{\fontsize{5.8}{6.4}\selectfont \textcolor{ForestGreen}{$ \blacktriangle\,29\% $}}\end{tabular} & \hspace*{1pt}\begin{tabular}[t]{@{}c@{}}\rule{0pt}{12pt}13.1s\\[-.2pt]{\fontsize{5.8}{6.4}\selectfont \textcolor{ForestGreen}{$ \blacktriangle\,2.1\times $}}\\[-.2pt]9.4s\\[-.2pt]{\fontsize{5.8}{6.4}\selectfont \textcolor{ForestGreen}{$ \blacktriangle\,2.8\times $}}\end{tabular} & \hspace*{1pt}\begin{tabular}[t]{@{}c@{}}\rule{0pt}{12pt}2m\,52.3s\\[-.2pt]{\fontsize{5.8}{6.4}\selectfont \textcolor{BrickRed}{$ \blacktriangledown\,3.3\times $}}\\[-.2pt]2m\,28.6s\\[-.2pt]{\fontsize{5.8}{6.4}\selectfont \textcolor{BrickRed}{$ \blacktriangledown\,2.9\times $}}\end{tabular} & \hspace*{1pt}\begin{tabular}[t]{@{}c@{}}\rule{0pt}{12pt}3m\,4.9s\\[-.2pt]{\fontsize{5.8}{6.4}\selectfont \textcolor{BrickRed}{$ \blacktriangledown\,3.3\times $}}\\[-.2pt]2m\,43.4s\\[-.2pt]{\fontsize{5.8}{6.4}\selectfont \textcolor{BrickRed}{$ \blacktriangledown\,3.0\times $}}\end{tabular} & \textcolor{gray}{\ensuremath{\cdot}} & \textcolor{gray}{\ensuremath{\cdot}} & \hspace*{1pt}\begin{tabular}[t]{@{}c@{}}\rule{0pt}{12pt}7m\,3.4s\\[-.2pt]{\fontsize{5.8}{6.4}\selectfont \textcolor{BrickRed}{$ \blacktriangledown\,8.8\times $}}\\[-.2pt]2m\,40.4s\\[-.2pt]{\fontsize{5.8}{6.4}\selectfont \textcolor{BrickRed}{$ \blacktriangledown\,3.4\times $}}\end{tabular} & \hspace*{1pt}\begin{tabular}[t]{@{}c@{}}\rule{0pt}{12pt}6m\,32.5s\\[-.2pt]{\fontsize{5.8}{6.4}\selectfont \textcolor{BrickRed}{$ \blacktriangledown\,8.8\times $}}\\[-.2pt]2m\,49.9s\\[-.2pt]{\fontsize{5.8}{6.4}\selectfont \textcolor{BrickRed}{$ \blacktriangledown\,4.0\times $}}\end{tabular} & \textcolor{gray}{\ensuremath{\cdot}} & \textcolor{gray}{\ensuremath{\cdot}}\\[3pt]\hline
\BenchmarkSourceLabel{bench:aws-2-k15-w10:row:mpf-simdref-simd-ksw2}{R20}\shortstack[l]{{\color{gray}R20}} & \begin{tabular}[t]{@{}l@{}}\rule{0pt}{12pt}MPF / SIMD-ref\\SIMD-DP / KSW2\end{tabular} & \hspace*{1pt}\begin{tabular}[t]{@{}c@{}}\rule{0pt}{12pt}5.3s\\[-.2pt]{\fontsize{5.8}{6.4}\selectfont \textcolor{ForestGreen}{$ \blacktriangle\,2.1\times $}}\\[-.2pt]2.8s\\[-.2pt]{\fontsize{5.8}{6.4}\selectfont \textcolor{ForestGreen}{$ \blacktriangle\,3.6\times $}}\end{tabular} & \hspace*{1pt}\begin{tabular}[t]{@{}c@{}}\rule{0pt}{12pt}4.8s\\[-.2pt]{\fontsize{5.8}{6.4}\selectfont \textcolor{ForestGreen}{$ \blacktriangle\,2.2\times $}}\\[-.2pt]2.5s\\[-.2pt]{\fontsize{5.8}{6.4}\selectfont \textcolor{ForestGreen}{$ \blacktriangle\,3.4\times $}}\end{tabular} & \hspace*{1pt}\begin{tabular}[t]{@{}c@{}}\rule{0pt}{12pt}9.7s\\[-.2pt]{\fontsize{5.8}{6.4}\selectfont \textcolor{ForestGreen}{$ \blacktriangle\,2.3\times $}}\\[-.2pt]4.3s\\[-.2pt]{\fontsize{5.8}{6.4}\selectfont \textcolor{ForestGreen}{$ \blacktriangle\,4.4\times $}}\end{tabular} & \hspace*{1pt}\begin{tabular}[t]{@{}c@{}}\rule{0pt}{12pt}9s\\[-.2pt]{\fontsize{5.8}{6.4}\selectfont \textcolor{ForestGreen}{$ \blacktriangle\,8.1\times $}}\\[-.2pt]6.2s\\[-.2pt]{\fontsize{5.8}{6.4}\selectfont \textcolor{ForestGreen}{$ \blacktriangle\,10.6\times $}}\end{tabular} & \hspace*{1pt}\begin{tabular}[t]{@{}c@{}}\rule{0pt}{12pt}16.7s\\[-.2pt]{\fontsize{5.8}{6.4}\selectfont \textcolor{BrickRed}{$ \blacktriangledown\,1.7\times $}}\\[-.2pt]11.2s\\[-.2pt]{\fontsize{5.8}{6.4}\selectfont \textcolor{BrickRed}{$ \blacktriangledown\,27\% $}}\end{tabular} & \hspace*{1pt}\begin{tabular}[t]{@{}c@{}}\rule{0pt}{12pt}16.4s\\[-.2pt]{\fontsize{5.8}{6.4}\selectfont \textcolor{BrickRed}{$ \blacktriangledown\,1.7\times $}}\\[-.2pt]9.7s\\[-.2pt]{\fontsize{5.8}{6.4}\selectfont \textcolor{BrickRed}{$ \blacktriangledown\,27\% $}}\end{tabular} & \hspace*{1pt}\begin{tabular}[t]{@{}c@{}}\rule{0pt}{12pt}30.5s\\[-.2pt]{\fontsize{5.8}{6.4}\selectfont \textcolor{BrickRed}{$ \blacktriangledown\,1.8\times $}}\\[-.2pt]14.9s\\[-.2pt]{\fontsize{5.8}{6.4}\selectfont \textcolor{BrickRed}{$ \blacktriangledown\,15\% $}}\end{tabular} & \hspace*{1pt}\begin{tabular}[t]{@{}c@{}}\rule{0pt}{12pt}17.3s\\[-.2pt]{\fontsize{5.8}{6.4}\selectfont \textcolor{ForestGreen}{$ \blacktriangle\,1.6\times $}}\\[-.2pt]13.9s\\[-.2pt]{\fontsize{5.8}{6.4}\selectfont \textcolor{ForestGreen}{$ \blacktriangle\,1.9\times $}}\end{tabular} & \hspace*{1pt}\begin{tabular}[t]{@{}c@{}}\rule{0pt}{12pt}2m\,51.2s\\[-.2pt]{\fontsize{5.8}{6.4}\selectfont \textcolor{BrickRed}{$ \blacktriangledown\,3.2\times $}}\\[-.2pt]2m\,24.3s\\[-.2pt]{\fontsize{5.8}{6.4}\selectfont \textcolor{BrickRed}{$ \blacktriangledown\,2.8\times $}}\end{tabular} & \hspace*{1pt}\begin{tabular}[t]{@{}c@{}}\rule{0pt}{12pt}2m\,22.1s\\[-.2pt]{\fontsize{5.8}{6.4}\selectfont \textcolor{BrickRed}{$ \blacktriangledown\,2.5\times $}}\\[-.2pt]2m\,1.6s\\[-.2pt]{\fontsize{5.8}{6.4}\selectfont \textcolor{BrickRed}{$ \blacktriangledown\,2.2\times $}}\end{tabular} & \hspace*{1pt}\begin{tabular}[t]{@{}c@{}}\rule{0pt}{12pt}3m\,23.3s\\[-.2pt]{\fontsize{5.8}{6.4}\selectfont \textcolor{BrickRed}{$ \blacktriangledown\,1.8\times $}}\\[-.2pt]2m\,50.3s\\[-.2pt]{\fontsize{5.8}{6.4}\selectfont \textcolor{BrickRed}{$ \blacktriangledown\,1.5\times $}}\end{tabular} & \textcolor{gray}{\ensuremath{\cdot}} & \textcolor{gray}{\ensuremath{\cdot}} & \textcolor{gray}{\ensuremath{\cdot}} & \textcolor{gray}{\ensuremath{\cdot}} & \textcolor{gray}{\ensuremath{\cdot}}\\[3pt]\hline
\BenchmarkSourceLabel{bench:aws-2-k15-w10:row:mpf-twostack-simd-ksw2}{R21}\shortstack[l]{{\color{gray}R21}} & \begin{tabular}[t]{@{}l@{}}\rule{0pt}{12pt}MPF / Two-stack\\SIMD-DP / KSW2\end{tabular} & \hspace*{1pt}\begin{tabular}[t]{@{}c@{}}\rule{0pt}{12pt}5.3s\\[-.2pt]{\fontsize{5.8}{6.4}\selectfont \textcolor{ForestGreen}{$ \blacktriangle\,2.1\times $}}\\[-.2pt]2.9s\\[-.2pt]{\fontsize{5.8}{6.4}\selectfont \textcolor{ForestGreen}{$ \blacktriangle\,3.5\times $}}\end{tabular} & \hspace*{1pt}\begin{tabular}[t]{@{}c@{}}\rule{0pt}{12pt}4.8s\\[-.2pt]{\fontsize{5.8}{6.4}\selectfont \textcolor{ForestGreen}{$ \blacktriangle\,2.2\times $}}\\[-.2pt]2.5s\\[-.2pt]{\fontsize{5.8}{6.4}\selectfont \textcolor{ForestGreen}{$ \blacktriangle\,3.4\times $}}\end{tabular} & \hspace*{1pt}\begin{tabular}[t]{@{}c@{}}\rule{0pt}{12pt}9.8s\\[-.2pt]{\fontsize{5.8}{6.4}\selectfont \textcolor{ForestGreen}{$ \blacktriangle\,2.3\times $}}\\[-.2pt]4.5s\\[-.2pt]{\fontsize{5.8}{6.4}\selectfont \textcolor{ForestGreen}{$ \blacktriangle\,4.2\times $}}\end{tabular} & \hspace*{1pt}\begin{tabular}[t]{@{}c@{}}\rule{0pt}{12pt}7.1s\\[-.2pt]{\fontsize{5.8}{6.4}\selectfont \textcolor{ForestGreen}{$ \blacktriangle\,10.3\times $}}\\[-.2pt]4.2s\\[-.2pt]{\fontsize{5.8}{6.4}\selectfont \textcolor{ForestGreen}{$ \blacktriangle\,15.6\times $}}\end{tabular} & \hspace*{1pt}\begin{tabular}[t]{@{}c@{}}\rule{0pt}{12pt}9.7s\\[-.2pt]{\fontsize{5.8}{6.4}\selectfont \textcolor{gray}{$ \blacktriangleright\,2\% $}}\\[-.2pt]5.5s\\[-.2pt]{\fontsize{5.8}{6.4}\selectfont \textcolor{ForestGreen}{$ \blacktriangle\,1.6\times $}}\end{tabular} & \hspace*{1pt}\begin{tabular}[t]{@{}c@{}}\rule{0pt}{12pt}9.6s\\[-.2pt]{\fontsize{5.8}{6.4}\selectfont \textcolor{gray}{$ \blacktriangleright\,1\% $}}\\[-.2pt]5.2s\\[-.2pt]{\fontsize{5.8}{6.4}\selectfont \textcolor{ForestGreen}{$ \blacktriangle\,1.5\times $}}\end{tabular} & \hspace*{1pt}\begin{tabular}[t]{@{}c@{}}\rule{0pt}{12pt}16.4s\\[-.2pt]{\fontsize{5.8}{6.4}\selectfont \textcolor{gray}{$ \blacktriangleright\,0\% $}}\\[-.2pt]8.8s\\[-.2pt]{\fontsize{5.8}{6.4}\selectfont \textcolor{ForestGreen}{$ \blacktriangle\,1.5\times $}}\end{tabular} & \hspace*{1pt}\begin{tabular}[t]{@{}c@{}}\rule{0pt}{12pt}14.9s\\[-.2pt]{\fontsize{5.8}{6.4}\selectfont \textcolor{ForestGreen}{$ \blacktriangle\,1.9\times $}}\\[-.2pt]11s\\[-.2pt]{\fontsize{5.8}{6.4}\selectfont \textcolor{ForestGreen}{$ \blacktriangle\,2.4\times $}}\end{tabular} & \hspace*{1pt}\begin{tabular}[t]{@{}c@{}}\rule{0pt}{12pt}2m\,27.4s\\[-.2pt]{\fontsize{5.8}{6.4}\selectfont \textcolor{BrickRed}{$ \blacktriangledown\,2.8\times $}}\\[-.2pt]2m\,3.1s\\[-.2pt]{\fontsize{5.8}{6.4}\selectfont \textcolor{BrickRed}{$ \blacktriangledown\,2.4\times $}}\end{tabular} & \hspace*{1pt}\begin{tabular}[t]{@{}c@{}}\rule{0pt}{12pt}2m\,26.8s\\[-.2pt]{\fontsize{5.8}{6.4}\selectfont \textcolor{BrickRed}{$ \blacktriangledown\,2.6\times $}}\\[-.2pt]2m\,5.2s\\[-.2pt]{\fontsize{5.8}{6.4}\selectfont \textcolor{BrickRed}{$ \blacktriangledown\,2.3\times $}}\end{tabular} & \textcolor{gray}{\ensuremath{\cdot}} & \textcolor{gray}{\ensuremath{\cdot}} & \hspace*{1pt}\begin{tabular}[t]{@{}c@{}}\rule{0pt}{12pt}6m\,28.7s\\[-.2pt]{\fontsize{5.8}{6.4}\selectfont \textcolor{BrickRed}{$ \blacktriangledown\,8.1\times $}}\\[-.2pt]2m\,7.4s\\[-.2pt]{\fontsize{5.8}{6.4}\selectfont \textcolor{BrickRed}{$ \blacktriangledown\,2.7\times $}}\end{tabular} & \hspace*{1pt}\begin{tabular}[t]{@{}c@{}}\rule{0pt}{12pt}5m\,46.6s\\[-.2pt]{\fontsize{5.8}{6.4}\selectfont \textcolor{BrickRed}{$ \blacktriangledown\,7.8\times $}}\\[-.2pt]2m\,1.5s\\[-.2pt]{\fontsize{5.8}{6.4}\selectfont \textcolor{BrickRed}{$ \blacktriangledown\,2.9\times $}}\end{tabular} & \textcolor{gray}{\ensuremath{\cdot}} & \textcolor{gray}{\ensuremath{\cdot}}\\[3pt]\hline
\BenchmarkSourceLabel{bench:aws-2-k15-w10:row:mpf-simdref-rmq-ksw2}{R22}\shortstack[l]{{\color{gray}R22}} & \begin{tabular}[t]{@{}l@{}}\rule{0pt}{12pt}MPF / SIMD-ref\\RMQ / KSW2\end{tabular} & \hspace*{1pt}\begin{tabular}[t]{@{}c@{}}\rule{0pt}{12pt}5.3s\\[-.2pt]{\fontsize{5.8}{6.4}\selectfont \textcolor{ForestGreen}{$ \blacktriangle\,2.1\times $}}\\[-.2pt]2.8s\\[-.2pt]{\fontsize{5.8}{6.4}\selectfont \textcolor{ForestGreen}{$ \blacktriangle\,3.6\times $}}\end{tabular} & \hspace*{1pt}\begin{tabular}[t]{@{}c@{}}\rule{0pt}{12pt}4.9s\\[-.2pt]{\fontsize{5.8}{6.4}\selectfont \textcolor{ForestGreen}{$ \blacktriangle\,2.2\times $}}\\[-.2pt]2.5s\\[-.2pt]{\fontsize{5.8}{6.4}\selectfont \textcolor{ForestGreen}{$ \blacktriangle\,3.4\times $}}\end{tabular} & \hspace*{1pt}\begin{tabular}[t]{@{}c@{}}\rule{0pt}{12pt}9.9s\\[-.2pt]{\fontsize{5.8}{6.4}\selectfont \textcolor{ForestGreen}{$ \blacktriangle\,2.3\times $}}\\[-.2pt]4.3s\\[-.2pt]{\fontsize{5.8}{6.4}\selectfont \textcolor{ForestGreen}{$ \blacktriangle\,4.4\times $}}\end{tabular} & \hspace*{1pt}\begin{tabular}[t]{@{}c@{}}\rule{0pt}{12pt}5.5s\\[-.2pt]{\fontsize{5.8}{6.4}\selectfont \textcolor{ForestGreen}{$ \blacktriangle\,13.3\times $}}\\[-.2pt]2.4s\\[-.2pt]{\fontsize{5.8}{6.4}\selectfont \textcolor{ForestGreen}{$ \blacktriangle\,26.8\times $}}\end{tabular} & \hspace*{1pt}\begin{tabular}[t]{@{}c@{}}\rule{0pt}{12pt}13.6s\\[-.2pt]{\fontsize{5.8}{6.4}\selectfont \textcolor{BrickRed}{$ \blacktriangledown\,1.4\times $}}\\[-.2pt]7.5s\\[-.2pt]{\fontsize{5.8}{6.4}\selectfont \textcolor{ForestGreen}{$ \blacktriangle\,15\% $}}\end{tabular} & \hspace*{1pt}\begin{tabular}[t]{@{}c@{}}\rule{0pt}{12pt}13.8s\\[-.2pt]{\fontsize{5.8}{6.4}\selectfont \textcolor{BrickRed}{$ \blacktriangledown\,1.5\times $}}\\[-.2pt]7.7s\\[-.2pt]{\fontsize{5.8}{6.4}\selectfont \textcolor{gray}{$ \blacktriangleright\,0\% $}}\end{tabular} & \hspace*{1pt}\begin{tabular}[t]{@{}c@{}}\rule{0pt}{12pt}26.7s\\[-.2pt]{\fontsize{5.8}{6.4}\selectfont \textcolor{BrickRed}{$ \blacktriangledown\,1.6\times $}}\\[-.2pt]13.8s\\[-.2pt]{\fontsize{5.8}{6.4}\selectfont \textcolor{BrickRed}{$ \blacktriangledown\,7\% $}}\end{tabular} & \hspace*{1pt}\begin{tabular}[t]{@{}c@{}}\rule{0pt}{12pt}31.9s\\[-.2pt]{\fontsize{5.8}{6.4}\selectfont \textcolor{BrickRed}{$ \blacktriangledown\,13\% $}}\\[-.2pt]28.3s\\[-.2pt]{\fontsize{5.8}{6.4}\selectfont \textcolor{BrickRed}{$ \blacktriangledown\,8\% $}}\end{tabular} & \hspace*{1pt}\begin{tabular}[t]{@{}c@{}}\rule{0pt}{12pt}2m\,45.3s\\[-.2pt]{\fontsize{5.8}{6.4}\selectfont \textcolor{BrickRed}{$ \blacktriangledown\,3.1\times $}}\\[-.2pt]2m\,18.5s\\[-.2pt]{\fontsize{5.8}{6.4}\selectfont \textcolor{BrickRed}{$ \blacktriangledown\,2.7\times $}}\end{tabular} & \hspace*{1pt}\begin{tabular}[t]{@{}c@{}}\rule{0pt}{12pt}2m\,24.6s\\[-.2pt]{\fontsize{5.8}{6.4}\selectfont \textcolor{BrickRed}{$ \blacktriangledown\,2.6\times $}}\\[-.2pt]2m\,3.8s\\[-.2pt]{\fontsize{5.8}{6.4}\selectfont \textcolor{BrickRed}{$ \blacktriangledown\,2.3\times $}}\end{tabular} & \textcolor{gray}{\ensuremath{\cdot}} & \textcolor{gray}{\ensuremath{\cdot}} & \textcolor{gray}{\ensuremath{\cdot}} & \textcolor{gray}{\ensuremath{\cdot}} & \textcolor{gray}{\ensuremath{\cdot}} & \textcolor{gray}{\ensuremath{\cdot}}\\[3pt]\hline
\end{tabular}\par}
\smallskip{\footnotesize\noindent Each cell: A and B times, $\cdot$: no measurement. ---\textsuperscript{g}: GPU unavailable. n/a: no comparison.\par
\noindent Blue: one lowest A and one lowest B measurement for each column for each configuration.\par
\noindent Gray: minimap2. Yellow: my preferred MPF / Bounded / SIMD-DP / Block configuration.\par
\noindent MPF/MPS: minimizers per fragment/sequence. SIMD-ref: naive SIMD scan. Bounded: bounded SIMD algorithm.\par
}
\clearpage
\noindent{\large\bfseries Table \thetable: (D): AWS: 2 Gbp, k=15, w=10\hfill 4/7}\par\smallskip
{\fontsize{7.1}{8.1}\selectfont\setlength{\tabcolsep}{1pt}\setlength{\fboxsep}{2pt}\renewcommand{\arraystretch}{1.12}
\noindent\begin{tabular}{@{}p{14mm}p{27mm}|p{13.55mm}|p{13.55mm}|p{13.55mm}|p{13.55mm}|p{13.55mm}|p{13.55mm}|p{13.55mm}|p{13.55mm}|p{13.55mm}|p{13.55mm}|p{13.55mm}|p{13.55mm}|p{13.55mm}|p{13.55mm}|p{13.55mm}|p{13.55mm}|@{}}\hline
Engine & Execution & \multicolumn{8}{c|}{\textbf{Index}} & \multicolumn{8}{c|}{\textbf{No index}}\\
 &  & \multicolumn{8}{c|}{\textbf{2 Gbp}} & \multicolumn{8}{c|}{\textbf{2 Gbp}}\\
 &  & \multicolumn{8}{c|}{\textbf{k=15, w=10}} & \multicolumn{8}{c|}{\textbf{k=15, w=10}}\\
 &  & \multicolumn{4}{c|}{\textbf{DNA4}} & \multicolumn{4}{c|}{\textbf{DNA5}} & \multicolumn{4}{c|}{\textbf{DNA4}} & \multicolumn{4}{c|}{\textbf{DNA5}}\\
 &  & \multicolumn{1}{c|}{\textbf{M7A}} & \multicolumn{1}{c|}{\textbf{M7I}} & \multicolumn{1}{c|}{\textbf{M4}} & \multicolumn{1}{c|}{\textbf{C7G}} & \multicolumn{1}{c|}{\textbf{M7A}} & \multicolumn{1}{c|}{\textbf{M7I}} & \multicolumn{1}{c|}{\textbf{M4}} & \multicolumn{1}{c|}{\textbf{C7G}} & \multicolumn{1}{c|}{\textbf{M7A}} & \multicolumn{1}{c|}{\textbf{M7I}} & \multicolumn{1}{c|}{\textbf{M4}} & \multicolumn{1}{c|}{\textbf{C7G}} & \multicolumn{1}{c|}{\textbf{M7A}} & \multicolumn{1}{c|}{\textbf{M7I}} & \multicolumn{1}{c|}{\textbf{M4}} & \multicolumn{1}{c|}{\textbf{C7G}}\\
 & & {\color{gray}C1} & {\color{gray}C2} & {\color{gray}C3} & {\color{gray}C4} & {\color{gray}C5} & {\color{gray}C6} & {\color{gray}C7} & {\color{gray}C8} & {\color{gray}C9} & {\color{gray}C10} & {\color{gray}C11} & {\color{gray}C12} & {\color{gray}C13} & {\color{gray}C14} & {\color{gray}C15} & {\color{gray}C16}\\
\hline
\rowcolor{BenchBaseline}[0.5pt][0.5pt]\shortstack[l]{{\color{gray}R1}\\\textbf{minimap2}} & \textbf{Baseline} & \hspace*{1pt}\begin{tabular}[t]{@{}c@{}}\rule{0pt}{12pt}11.2s\\[-.2pt]{\fontsize{5.8}{6.4}\selectfont \textcolor{BrickRed}{$ \blacktriangledown\,3.1\times $}}\\[-.2pt]10s\\[-.2pt]{\fontsize{5.8}{6.4}\selectfont \textcolor{BrickRed}{$ \blacktriangledown\,3.6\times $}}\end{tabular} & \hspace*{1pt}\begin{tabular}[t]{@{}c@{}}\rule{0pt}{12pt}10.7s\\[-.2pt]{\fontsize{5.8}{6.4}\selectfont \textcolor{BrickRed}{$ \blacktriangledown\,3.2\times $}}\\[-.2pt]8.6s\\[-.2pt]{\fontsize{5.8}{6.4}\selectfont \textcolor{BrickRed}{$ \blacktriangledown\,3.5\times $}}\end{tabular} & \hspace*{1pt}\begin{tabular}[t]{@{}c@{}}\rule{0pt}{12pt}22.5s\\[-.2pt]{\fontsize{5.8}{6.4}\selectfont \textcolor{BrickRed}{$ \blacktriangledown\,3.6\times $}}\\[-.2pt]18.8s\\[-.2pt]{\fontsize{5.8}{6.4}\selectfont \textcolor{BrickRed}{$ \blacktriangledown\,4.4\times $}}\end{tabular} & \hspace*{1pt}\begin{tabular}[t]{@{}c@{}}\rule{0pt}{12pt}1m\,13.1s\\[-.2pt]{\fontsize{5.8}{6.4}\selectfont \textcolor{BrickRed}{$ \blacktriangledown\,22.5\times $}}\\[-.2pt]1m\,5.7s\\[-.2pt]{\fontsize{5.8}{6.4}\selectfont \textcolor{BrickRed}{$ \blacktriangledown\,27.7\times $}}\end{tabular} & \hspace*{1pt}\begin{tabular}[t]{@{}c@{}}\rule{0pt}{12pt}9.9s\\[-.2pt]{\fontsize{5.8}{6.4}\selectfont \textcolor{BrickRed}{$ \blacktriangledown\,21\% $}}\\[-.2pt]8.8s\\[-.2pt]{\fontsize{5.8}{6.4}\selectfont \textcolor{BrickRed}{$ \blacktriangledown\,1.6\times $}}\end{tabular} & \hspace*{1pt}\begin{tabular}[t]{@{}c@{}}\rule{0pt}{12pt}9.5s\\[-.2pt]{\fontsize{5.8}{6.4}\selectfont \textcolor{BrickRed}{$ \blacktriangledown\,1.4\times $}}\\[-.2pt]7.6s\\[-.2pt]{\fontsize{5.8}{6.4}\selectfont \textcolor{BrickRed}{$ \blacktriangledown\,1.6\times $}}\end{tabular} & \hspace*{1pt}\begin{tabular}[t]{@{}c@{}}\rule{0pt}{12pt}16.5s\\[-.2pt]{\fontsize{5.8}{6.4}\selectfont \textcolor{BrickRed}{$ \blacktriangledown\,1.7\times $}}\\[-.2pt]12.9s\\[-.2pt]{\fontsize{5.8}{6.4}\selectfont \textcolor{BrickRed}{$ \blacktriangledown\,2.0\times $}}\end{tabular} & \hspace*{1pt}\begin{tabular}[t]{@{}c@{}}\rule{0pt}{12pt}28.2s\\[-.2pt]{\fontsize{5.8}{6.4}\selectfont \textcolor{BrickRed}{$ \blacktriangledown\,2.3\times $}}\\[-.2pt]26.2s\\[-.2pt]{\fontsize{5.8}{6.4}\selectfont \textcolor{BrickRed}{$ \blacktriangledown\,3.0\times $}}\end{tabular} & \hspace*{1pt}\begin{tabular}[t]{@{}c@{}}\rule{0pt}{12pt}53s\\[-.2pt]{\fontsize{5.8}{6.4}\selectfont \textcolor{ForestGreen}{$ \blacktriangle\,2.2\times $}}\\[-.2pt]51.9s\\[-.2pt]{\fontsize{5.8}{6.4}\selectfont \textcolor{ForestGreen}{$ \blacktriangle\,2.1\times $}}\end{tabular} & \hspace*{1pt}\begin{tabular}[t]{@{}c@{}}\rule{0pt}{12pt}56.6s\\[-.2pt]{\fontsize{5.8}{6.4}\selectfont \textcolor{ForestGreen}{$ \blacktriangle\,2.1\times $}}\\[-.2pt]54.6s\\[-.2pt]{\fontsize{5.8}{6.4}\selectfont \textcolor{ForestGreen}{$ \blacktriangle\,1.9\times $}}\end{tabular} & \hspace*{1pt}\begin{tabular}[t]{@{}c@{}}\rule{0pt}{12pt}1m\,55.5s\\[-.2pt]{\fontsize{5.8}{6.4}\selectfont \textcolor{ForestGreen}{$ \blacktriangle\,1.8\times $}}\\[-.2pt]1m\,51.8s\\[-.2pt]{\fontsize{5.8}{6.4}\selectfont \textcolor{ForestGreen}{$ \blacktriangle\,1.5\times $}}\end{tabular} & \hspace*{1pt}\begin{tabular}[t]{@{}c@{}}\rule{0pt}{12pt}6m\,1.3s\\[-.2pt]{\fontsize{5.8}{6.4}\selectfont \textcolor{BrickRed}{$ \blacktriangledown\,3.0\times $}}\\[-.2pt]5m\,55.2s\\[-.2pt]{\fontsize{5.8}{6.4}\selectfont \textcolor{BrickRed}{$ \blacktriangledown\,3.1\times $}}\end{tabular} & \hspace*{1pt}\begin{tabular}[t]{@{}c@{}}\rule{0pt}{12pt}48.1s\\[-.2pt]{\fontsize{5.8}{6.4}\selectfont \textcolor{ForestGreen}{$ \blacktriangle\,7.3\times $}}\\[-.2pt]47s\\[-.2pt]{\fontsize{5.8}{6.4}\selectfont \textcolor{ForestGreen}{$ \blacktriangle\,2.4\times $}}\end{tabular} & \hspace*{1pt}\begin{tabular}[t]{@{}c@{}}\rule{0pt}{12pt}44.5s\\[-.2pt]{\fontsize{5.8}{6.4}\selectfont \textcolor{ForestGreen}{$ \blacktriangle\,6.5\times $}}\\[-.2pt]42.6s\\[-.2pt]{\fontsize{5.8}{6.4}\selectfont \textcolor{ForestGreen}{$ \blacktriangle\,2.7\times $}}\end{tabular} & \hspace*{1pt}\begin{tabular}[t]{@{}c@{}}\rule{0pt}{12pt}1m\,30.2s\\[-.2pt]{\fontsize{5.8}{6.4}\selectfont \textcolor{ForestGreen}{$ \blacktriangle\,6.4\times $}}\\[-.2pt]1m\,26.7s\\[-.2pt]{\fontsize{5.8}{6.4}\selectfont \textcolor{ForestGreen}{$ \blacktriangle\,2.2\times $}}\end{tabular} & \hspace*{1pt}\begin{tabular}[t]{@{}c@{}}\rule{0pt}{12pt}2m\,46.4s\\[-.2pt]{\fontsize{5.8}{6.4}\selectfont \textcolor{ForestGreen}{$ \blacktriangle\,4.0\times $}}\\[-.2pt]2m\,42.7s\\[-.2pt]{\fontsize{5.8}{6.4}\selectfont \textcolor{BrickRed}{$ \blacktriangledown\,27\% $}}\end{tabular}\\[3pt]\hline
\BenchmarkSourceLabel{bench:aws-2-k15-w10:row:mpf-scalar-simd-ksw2}{R23}\shortstack[l]{{\color{gray}R23}\\\textbf{BioseqDB}} & \begin{tabular}[t]{@{}l@{}}\rule{0pt}{12pt}MPF / Scalar\\SIMD-DP / KSW2\end{tabular} & \hspace*{1pt}\begin{tabular}[t]{@{}c@{}}\rule{0pt}{12pt}5.4s\\[-.2pt]{\fontsize{5.8}{6.4}\selectfont \textcolor{ForestGreen}{$ \blacktriangle\,2.1\times $}}\\[-.2pt]2.9s\\[-.2pt]{\fontsize{5.8}{6.4}\selectfont \textcolor{ForestGreen}{$ \blacktriangle\,3.5\times $}}\end{tabular} & \hspace*{1pt}\begin{tabular}[t]{@{}c@{}}\rule{0pt}{12pt}4.9s\\[-.2pt]{\fontsize{5.8}{6.4}\selectfont \textcolor{ForestGreen}{$ \blacktriangle\,2.2\times $}}\\[-.2pt]2.6s\\[-.2pt]{\fontsize{5.8}{6.4}\selectfont \textcolor{ForestGreen}{$ \blacktriangle\,3.3\times $}}\end{tabular} & \hspace*{1pt}\begin{tabular}[t]{@{}c@{}}\rule{0pt}{12pt}10.3s\\[-.2pt]{\fontsize{5.8}{6.4}\selectfont \textcolor{ForestGreen}{$ \blacktriangle\,2.2\times $}}\\[-.2pt]4.6s\\[-.2pt]{\fontsize{5.8}{6.4}\selectfont \textcolor{ForestGreen}{$ \blacktriangle\,4.1\times $}}\end{tabular} & \hspace*{1pt}\begin{tabular}[t]{@{}c@{}}\rule{0pt}{12pt}7.8s\\[-.2pt]{\fontsize{5.8}{6.4}\selectfont \textcolor{ForestGreen}{$ \blacktriangle\,9.3\times $}}\\[-.2pt]5s\\[-.2pt]{\fontsize{5.8}{6.4}\selectfont \textcolor{ForestGreen}{$ \blacktriangle\,13.2\times $}}\end{tabular} & \hspace*{1pt}\begin{tabular}[t]{@{}c@{}}\rule{0pt}{12pt}10.2s\\[-.2pt]{\fontsize{5.8}{6.4}\selectfont \textcolor{gray}{$ \blacktriangleright\,3\% $}}\\[-.2pt]5.9s\\[-.2pt]{\fontsize{5.8}{6.4}\selectfont \textcolor{ForestGreen}{$ \blacktriangle\,1.5\times $}}\end{tabular} & \hspace*{1pt}\begin{tabular}[t]{@{}c@{}}\rule{0pt}{12pt}8.9s\\[-.2pt]{\fontsize{5.8}{6.4}\selectfont \textcolor{ForestGreen}{$ \blacktriangle\,6\% $}}\\[-.2pt]5.3s\\[-.2pt]{\fontsize{5.8}{6.4}\selectfont \textcolor{ForestGreen}{$ \blacktriangle\,1.5\times $}}\end{tabular} & \hspace*{1pt}\begin{tabular}[t]{@{}c@{}}\rule{0pt}{12pt}17.8s\\[-.2pt]{\fontsize{5.8}{6.4}\selectfont \textcolor{BrickRed}{$ \blacktriangledown\,8\% $}}\\[-.2pt]10.8s\\[-.2pt]{\fontsize{5.8}{6.4}\selectfont \textcolor{ForestGreen}{$ \blacktriangle\,16\% $}}\end{tabular} & \hspace*{1pt}\begin{tabular}[t]{@{}c@{}}\rule{0pt}{12pt}15.6s\\[-.2pt]{\fontsize{5.8}{6.4}\selectfont \textcolor{ForestGreen}{$ \blacktriangle\,1.8\times $}}\\[-.2pt]12.5s\\[-.2pt]{\fontsize{5.8}{6.4}\selectfont \textcolor{ForestGreen}{$ \blacktriangle\,2.1\times $}}\end{tabular} & \hspace*{1pt}\begin{tabular}[t]{@{}c@{}}\rule{0pt}{12pt}3m\,4.9s\\[-.2pt]{\fontsize{5.8}{6.4}\selectfont \textcolor{BrickRed}{$ \blacktriangledown\,3.5\times $}}\\[-.2pt]2m\,40.5s\\[-.2pt]{\fontsize{5.8}{6.4}\selectfont \textcolor{BrickRed}{$ \blacktriangledown\,3.1\times $}}\end{tabular} & \hspace*{1pt}\begin{tabular}[t]{@{}c@{}}\rule{0pt}{12pt}3m\,7.6s\\[-.2pt]{\fontsize{5.8}{6.4}\selectfont \textcolor{BrickRed}{$ \blacktriangledown\,3.3\times $}}\\[-.2pt]2m\,46s\\[-.2pt]{\fontsize{5.8}{6.4}\selectfont \textcolor{BrickRed}{$ \blacktriangledown\,3.0\times $}}\end{tabular} & \textcolor{gray}{\ensuremath{\cdot}} & \textcolor{gray}{\ensuremath{\cdot}} & \hspace*{1pt}\begin{tabular}[t]{@{}c@{}}\rule{0pt}{12pt}7m\,6.6s\\[-.2pt]{\fontsize{5.8}{6.4}\selectfont \textcolor{BrickRed}{$ \blacktriangledown\,8.9\times $}}\\[-.2pt]2m\,43.2s\\[-.2pt]{\fontsize{5.8}{6.4}\selectfont \textcolor{BrickRed}{$ \blacktriangledown\,3.5\times $}}\end{tabular} & \hspace*{1pt}\begin{tabular}[t]{@{}c@{}}\rule{0pt}{12pt}6m\,54s\\[-.2pt]{\fontsize{5.8}{6.4}\selectfont \textcolor{BrickRed}{$ \blacktriangledown\,9.3\times $}}\\[-.2pt]3m\,6.6s\\[-.2pt]{\fontsize{5.8}{6.4}\selectfont \textcolor{BrickRed}{$ \blacktriangledown\,4.4\times $}}\end{tabular} & \textcolor{gray}{\ensuremath{\cdot}} & \textcolor{gray}{\ensuremath{\cdot}}\\[3pt]\hline
\BenchmarkSourceLabel{bench:aws-2-k15-w10:row:mpf-twostack-rmq-ksw2}{R24}\shortstack[l]{{\color{gray}R24}} & \begin{tabular}[t]{@{}l@{}}\rule{0pt}{12pt}MPF / Two-stack\\RMQ / KSW2\end{tabular} & \hspace*{1pt}\begin{tabular}[t]{@{}c@{}}\rule{0pt}{12pt}5.3s\\[-.2pt]{\fontsize{5.8}{6.4}\selectfont \textcolor{ForestGreen}{$ \blacktriangle\,2.1\times $}}\\[-.2pt]2.9s\\[-.2pt]{\fontsize{5.8}{6.4}\selectfont \textcolor{ForestGreen}{$ \blacktriangle\,3.5\times $}}\end{tabular} & \hspace*{1pt}\begin{tabular}[t]{@{}c@{}}\rule{0pt}{12pt}5s\\[-.2pt]{\fontsize{5.8}{6.4}\selectfont \textcolor{ForestGreen}{$ \blacktriangle\,2.1\times $}}\\[-.2pt]2.5s\\[-.2pt]{\fontsize{5.8}{6.4}\selectfont \textcolor{ForestGreen}{$ \blacktriangle\,3.4\times $}}\end{tabular} & \hspace*{1pt}\begin{tabular}[t]{@{}c@{}}\rule{0pt}{12pt}10.1s\\[-.2pt]{\fontsize{5.8}{6.4}\selectfont \textcolor{ForestGreen}{$ \blacktriangle\,2.2\times $}}\\[-.2pt]4.5s\\[-.2pt]{\fontsize{5.8}{6.4}\selectfont \textcolor{ForestGreen}{$ \blacktriangle\,4.2\times $}}\end{tabular} & \hspace*{1pt}\begin{tabular}[t]{@{}c@{}}\rule{0pt}{12pt}5.6s\\[-.2pt]{\fontsize{5.8}{6.4}\selectfont \textcolor{ForestGreen}{$ \blacktriangle\,12.9\times $}}\\[-.2pt]2.6s\\[-.2pt]{\fontsize{5.8}{6.4}\selectfont \textcolor{ForestGreen}{$ \blacktriangle\,25.5\times $}}\end{tabular} & \hspace*{1pt}\begin{tabular}[t]{@{}c@{}}\rule{0pt}{12pt}9.6s\\[-.2pt]{\fontsize{5.8}{6.4}\selectfont \textcolor{gray}{$ \blacktriangleright\,4\% $}}\\[-.2pt]\colorbox{BenchBest}{\strut 5.5s}\\[-.2pt]{\fontsize{5.8}{6.4}\selectfont \textcolor{ForestGreen}{$ \blacktriangle\,1.6\times $}}\end{tabular} & \hspace*{1pt}\begin{tabular}[t]{@{}c@{}}\rule{0pt}{12pt}8.9s\\[-.2pt]{\fontsize{5.8}{6.4}\selectfont \textcolor{ForestGreen}{$ \blacktriangle\,6\% $}}\\[-.2pt]5.1s\\[-.2pt]{\fontsize{5.8}{6.4}\selectfont \textcolor{ForestGreen}{$ \blacktriangle\,1.5\times $}}\end{tabular} & \hspace*{1pt}\begin{tabular}[t]{@{}c@{}}\rule{0pt}{12pt}16.4s\\[-.2pt]{\fontsize{5.8}{6.4}\selectfont \textcolor{gray}{$ \blacktriangleright\,1\% $}}\\[-.2pt]8.9s\\[-.2pt]{\fontsize{5.8}{6.4}\selectfont \textcolor{ForestGreen}{$ \blacktriangle\,1.5\times $}}\end{tabular} & \hspace*{1pt}\begin{tabular}[t]{@{}c@{}}\rule{0pt}{12pt}12.7s\\[-.2pt]{\fontsize{5.8}{6.4}\selectfont \textcolor{ForestGreen}{$ \blacktriangle\,2.2\times $}}\\[-.2pt]\colorbox{BenchBest}{\strut 8.8s}\\[-.2pt]{\fontsize{5.8}{6.4}\selectfont \textcolor{ForestGreen}{$ \blacktriangle\,3.0\times $}}\end{tabular} & \hspace*{1pt}\begin{tabular}[t]{@{}c@{}}\rule{0pt}{12pt}2m\,37s\\[-.2pt]{\fontsize{5.8}{6.4}\selectfont \textcolor{BrickRed}{$ \blacktriangledown\,3.0\times $}}\\[-.2pt]2m\,12.9s\\[-.2pt]{\fontsize{5.8}{6.4}\selectfont \textcolor{BrickRed}{$ \blacktriangledown\,2.6\times $}}\end{tabular} & \hspace*{1pt}\begin{tabular}[t]{@{}c@{}}\rule{0pt}{12pt}2m\,34.6s\\[-.2pt]{\fontsize{5.8}{6.4}\selectfont \textcolor{BrickRed}{$ \blacktriangledown\,2.7\times $}}\\[-.2pt]2m\,12.9s\\[-.2pt]{\fontsize{5.8}{6.4}\selectfont \textcolor{BrickRed}{$ \blacktriangledown\,2.4\times $}}\end{tabular} & \textcolor{gray}{\ensuremath{\cdot}} & \textcolor{gray}{\ensuremath{\cdot}} & \hspace*{1pt}\begin{tabular}[t]{@{}c@{}}\rule{0pt}{12pt}6m\,27.3s\\[-.2pt]{\fontsize{5.8}{6.4}\selectfont \textcolor{BrickRed}{$ \blacktriangledown\,8.1\times $}}\\[-.2pt]2m\,5.7s\\[-.2pt]{\fontsize{5.8}{6.4}\selectfont \textcolor{BrickRed}{$ \blacktriangledown\,2.7\times $}}\end{tabular} & \hspace*{1pt}\begin{tabular}[t]{@{}c@{}}\rule{0pt}{12pt}5m\,42s\\[-.2pt]{\fontsize{5.8}{6.4}\selectfont \textcolor{BrickRed}{$ \blacktriangledown\,7.7\times $}}\\[-.2pt]1m\,58.6s\\[-.2pt]{\fontsize{5.8}{6.4}\selectfont \textcolor{BrickRed}{$ \blacktriangledown\,2.8\times $}}\end{tabular} & \textcolor{gray}{\ensuremath{\cdot}} & \textcolor{gray}{\ensuremath{\cdot}}\\[3pt]\hline
\BenchmarkSourceLabel{bench:aws-2-k15-w10:row:mpf-scalar-rmq-ksw2}{R25}\shortstack[l]{{\color{gray}R25}} & \begin{tabular}[t]{@{}l@{}}\rule{0pt}{12pt}MPF / Scalar\\RMQ / KSW2\end{tabular} & \hspace*{1pt}\begin{tabular}[t]{@{}c@{}}\rule{0pt}{12pt}5.5s\\[-.2pt]{\fontsize{5.8}{6.4}\selectfont \textcolor{ForestGreen}{$ \blacktriangle\,2.0\times $}}\\[-.2pt]2.9s\\[-.2pt]{\fontsize{5.8}{6.4}\selectfont \textcolor{ForestGreen}{$ \blacktriangle\,3.4\times $}}\end{tabular} & \hspace*{1pt}\begin{tabular}[t]{@{}c@{}}\rule{0pt}{12pt}5.1s\\[-.2pt]{\fontsize{5.8}{6.4}\selectfont \textcolor{ForestGreen}{$ \blacktriangle\,2.1\times $}}\\[-.2pt]2.7s\\[-.2pt]{\fontsize{5.8}{6.4}\selectfont \textcolor{ForestGreen}{$ \blacktriangle\,3.2\times $}}\end{tabular} & \hspace*{1pt}\begin{tabular}[t]{@{}c@{}}\rule{0pt}{12pt}10.3s\\[-.2pt]{\fontsize{5.8}{6.4}\selectfont \textcolor{ForestGreen}{$ \blacktriangle\,2.2\times $}}\\[-.2pt]4.6s\\[-.2pt]{\fontsize{5.8}{6.4}\selectfont \textcolor{ForestGreen}{$ \blacktriangle\,4.1\times $}}\end{tabular} & \hspace*{1pt}\begin{tabular}[t]{@{}c@{}}\rule{0pt}{12pt}5.6s\\[-.2pt]{\fontsize{5.8}{6.4}\selectfont \textcolor{ForestGreen}{$ \blacktriangle\,13.1\times $}}\\[-.2pt]2.5s\\[-.2pt]{\fontsize{5.8}{6.4}\selectfont \textcolor{ForestGreen}{$ \blacktriangle\,26.3\times $}}\end{tabular} & \hspace*{1pt}\begin{tabular}[t]{@{}c@{}}\rule{0pt}{12pt}10s\\[-.2pt]{\fontsize{5.8}{6.4}\selectfont \textcolor{gray}{$ \blacktriangleright\,1\% $}}\\[-.2pt]5.9s\\[-.2pt]{\fontsize{5.8}{6.4}\selectfont \textcolor{ForestGreen}{$ \blacktriangle\,1.5\times $}}\end{tabular} & \hspace*{1pt}\begin{tabular}[t]{@{}c@{}}\rule{0pt}{12pt}8.7s\\[-.2pt]{\fontsize{5.8}{6.4}\selectfont \textcolor{ForestGreen}{$ \blacktriangle\,8\% $}}\\[-.2pt]5.3s\\[-.2pt]{\fontsize{5.8}{6.4}\selectfont \textcolor{ForestGreen}{$ \blacktriangle\,1.4\times $}}\end{tabular} & \hspace*{1pt}\begin{tabular}[t]{@{}c@{}}\rule{0pt}{12pt}16s\\[-.2pt]{\fontsize{5.8}{6.4}\selectfont \textcolor{gray}{$ \blacktriangleright\,3\% $}}\\[-.2pt]9.3s\\[-.2pt]{\fontsize{5.8}{6.4}\selectfont \textcolor{ForestGreen}{$ \blacktriangle\,28\% $}}\end{tabular} & \hspace*{1pt}\begin{tabular}[t]{@{}c@{}}\rule{0pt}{12pt}21.1s\\[-.2pt]{\fontsize{5.8}{6.4}\selectfont \textcolor{ForestGreen}{$ \blacktriangle\,25\% $}}\\[-.2pt]17.5s\\[-.2pt]{\fontsize{5.8}{6.4}\selectfont \textcolor{ForestGreen}{$ \blacktriangle\,1.5\times $}}\end{tabular} & \hspace*{1pt}\begin{tabular}[t]{@{}c@{}}\rule{0pt}{12pt}3m\,1.5s\\[-.2pt]{\fontsize{5.8}{6.4}\selectfont \textcolor{BrickRed}{$ \blacktriangledown\,3.4\times $}}\\[-.2pt]2m\,37.7s\\[-.2pt]{\fontsize{5.8}{6.4}\selectfont \textcolor{BrickRed}{$ \blacktriangledown\,3.0\times $}}\end{tabular} & \hspace*{1pt}\begin{tabular}[t]{@{}c@{}}\rule{0pt}{12pt}2m\,54.6s\\[-.2pt]{\fontsize{5.8}{6.4}\selectfont \textcolor{BrickRed}{$ \blacktriangledown\,3.1\times $}}\\[-.2pt]2m\,33.3s\\[-.2pt]{\fontsize{5.8}{6.4}\selectfont \textcolor{BrickRed}{$ \blacktriangledown\,2.8\times $}}\end{tabular} & \textcolor{gray}{\ensuremath{\cdot}} & \textcolor{gray}{\ensuremath{\cdot}} & \hspace*{1pt}\begin{tabular}[t]{@{}c@{}}\rule{0pt}{12pt}7m\,3s\\[-.2pt]{\fontsize{5.8}{6.4}\selectfont \textcolor{BrickRed}{$ \blacktriangledown\,8.8\times $}}\\[-.2pt]2m\,38.7s\\[-.2pt]{\fontsize{5.8}{6.4}\selectfont \textcolor{BrickRed}{$ \blacktriangledown\,3.4\times $}}\end{tabular} & \hspace*{1pt}\begin{tabular}[t]{@{}c@{}}\rule{0pt}{12pt}6m\,59.2s\\[-.2pt]{\fontsize{5.8}{6.4}\selectfont \textcolor{BrickRed}{$ \blacktriangledown\,9.4\times $}}\\[-.2pt]3m\,13.2s\\[-.2pt]{\fontsize{5.8}{6.4}\selectfont \textcolor{BrickRed}{$ \blacktriangledown\,4.5\times $}}\end{tabular} & \textcolor{gray}{\ensuremath{\cdot}} & \textcolor{gray}{\ensuremath{\cdot}}\\[3pt]\hline
\BenchmarkSourceLabel{bench:aws-2-k15-w10:row:mps-bounded-scalar-block}{R26}\shortstack[l]{{\color{gray}R26}} & \begin{tabular}[t]{@{}l@{}}\rule{0pt}{12pt}MPS / Bounded\\Scalar-DP / Block\end{tabular} & \hspace*{1pt}\begin{tabular}[t]{@{}c@{}}\rule{0pt}{12pt}6.4s\\[-.2pt]{\fontsize{5.8}{6.4}\selectfont \textcolor{ForestGreen}{$ \blacktriangle\,1.7\times $}}\\[-.2pt]5.7s\\[-.2pt]{\fontsize{5.8}{6.4}\selectfont \textcolor{ForestGreen}{$ \blacktriangle\,1.8\times $}}\end{tabular} & \hspace*{1pt}\begin{tabular}[t]{@{}c@{}}\rule{0pt}{12pt}6.6s\\[-.2pt]{\fontsize{5.8}{6.4}\selectfont \textcolor{ForestGreen}{$ \blacktriangle\,1.6\times $}}\\[-.2pt]5.7s\\[-.2pt]{\fontsize{5.8}{6.4}\selectfont \textcolor{ForestGreen}{$ \blacktriangle\,1.5\times $}}\end{tabular} & \hspace*{1pt}\begin{tabular}[t]{@{}c@{}}\rule{0pt}{12pt}12.5s\\[-.2pt]{\fontsize{5.8}{6.4}\selectfont \textcolor{ForestGreen}{$ \blacktriangle\,1.8\times $}}\\[-.2pt]10.1s\\[-.2pt]{\fontsize{5.8}{6.4}\selectfont \textcolor{ForestGreen}{$ \blacktriangle\,1.9\times $}}\end{tabular} & \hspace*{1pt}\begin{tabular}[t]{@{}c@{}}\rule{0pt}{12pt}5.9s\\[-.2pt]{\fontsize{5.8}{6.4}\selectfont \textcolor{ForestGreen}{$ \blacktriangle\,12.3\times $}}\\[-.2pt]4.9s\\[-.2pt]{\fontsize{5.8}{6.4}\selectfont \textcolor{ForestGreen}{$ \blacktriangle\,13.4\times $}}\end{tabular} & \hspace*{1pt}\begin{tabular}[t]{@{}c@{}}\rule{0pt}{12pt}4m\,54s\\[-.2pt]{\fontsize{5.8}{6.4}\selectfont \textcolor{BrickRed}{$ \blacktriangledown\,29.7\times $}}\\[-.2pt]32s\\[-.2pt]{\fontsize{5.8}{6.4}\selectfont \textcolor{BrickRed}{$ \blacktriangledown\,3.6\times $}}\end{tabular} & \hspace*{1pt}\begin{tabular}[t]{@{}c@{}}\rule{0pt}{12pt}4m\,0.6s\\[-.2pt]{\fontsize{5.8}{6.4}\selectfont \textcolor{BrickRed}{$ \blacktriangledown\,25.4\times $}}\\[-.2pt]28.5s\\[-.2pt]{\fontsize{5.8}{6.4}\selectfont \textcolor{BrickRed}{$ \blacktriangledown\,3.7\times $}}\end{tabular} & \hspace*{1pt}\begin{tabular}[t]{@{}c@{}}\rule{0pt}{12pt}7m\,33.3s\\[-.2pt]{\fontsize{5.8}{6.4}\selectfont \textcolor{BrickRed}{$ \blacktriangledown\,27.5\times $}}\\[-.2pt]53.3s\\[-.2pt]{\fontsize{5.8}{6.4}\selectfont \textcolor{BrickRed}{$ \blacktriangledown\,4.1\times $}}\end{tabular} & \textcolor{gray}{\ensuremath{\cdot}} & \hspace*{1pt}\begin{tabular}[t]{@{}c@{}}\rule{0pt}{12pt}2m\,24.9s\\[-.2pt]{\fontsize{5.8}{6.4}\selectfont \textcolor{BrickRed}{$ \blacktriangledown\,2.7\times $}}\\[-.2pt]2m\,2.4s\\[-.2pt]{\fontsize{5.8}{6.4}\selectfont \textcolor{BrickRed}{$ \blacktriangledown\,2.4\times $}}\end{tabular} & \hspace*{1pt}\begin{tabular}[t]{@{}c@{}}\rule{0pt}{12pt}2m\,39.7s\\[-.2pt]{\fontsize{5.8}{6.4}\selectfont \textcolor{BrickRed}{$ \blacktriangledown\,2.8\times $}}\\[-.2pt]2m\,18.8s\\[-.2pt]{\fontsize{5.8}{6.4}\selectfont \textcolor{BrickRed}{$ \blacktriangledown\,2.5\times $}}\end{tabular} & \textcolor{gray}{\ensuremath{\cdot}} & \hspace*{1pt}\begin{tabular}[t]{@{}c@{}}\rule{0pt}{12pt}2m\,43s\\[-.2pt]{\fontsize{5.8}{6.4}\selectfont \textcolor{ForestGreen}{$ \blacktriangle\,2.2\times $}}\\[-.2pt]2m\,26.5s\\[-.2pt]{\fontsize{5.8}{6.4}\selectfont \textcolor{ForestGreen}{$ \blacktriangle\,2.4\times $}}\end{tabular} & \hspace*{1pt}\begin{tabular}[t]{@{}c@{}}\rule{0pt}{12pt}6m\,32.4s\\[-.2pt]{\fontsize{5.8}{6.4}\selectfont \textcolor{BrickRed}{$ \blacktriangledown\,8.2\times $}}\\[-.2pt]2m\,10.2s\\[-.2pt]{\fontsize{5.8}{6.4}\selectfont \textcolor{BrickRed}{$ \blacktriangledown\,2.8\times $}}\end{tabular} & \hspace*{1pt}\begin{tabular}[t]{@{}c@{}}\rule{0pt}{12pt}5m\,33.1s\\[-.2pt]{\fontsize{5.8}{6.4}\selectfont \textcolor{BrickRed}{$ \blacktriangledown\,7.5\times $}}\\[-.2pt]1m\,56s\\[-.2pt]{\fontsize{5.8}{6.4}\selectfont \textcolor{BrickRed}{$ \blacktriangledown\,2.7\times $}}\end{tabular} & \textcolor{gray}{\ensuremath{\cdot}} & \hspace*{1pt}\begin{tabular}[t]{@{}c@{}}\rule{0pt}{12pt}11m\,41.8s\\[-.2pt]{\fontsize{5.8}{6.4}\selectfont \textcolor{BrickRed}{$ \blacktriangledown\,4.2\times $}}\\[-.2pt]2m\,22.1s\\[-.2pt]{\fontsize{5.8}{6.4}\selectfont \textcolor{ForestGreen}{$ \blacktriangle\,13\% $}}\end{tabular}\\[3pt]\hline
\BenchmarkSourceLabel{bench:aws-2-k15-w10:row:mps-bounded-simd-block}{R27}\shortstack[l]{{\color{gray}R27}} & \begin{tabular}[t]{@{}l@{}}\rule{0pt}{12pt}MPS / Bounded\\SIMD-DP / Block\end{tabular} & \hspace*{1pt}\begin{tabular}[t]{@{}c@{}}\rule{0pt}{12pt}6.7s\\[-.2pt]{\fontsize{5.8}{6.4}\selectfont \textcolor{ForestGreen}{$ \blacktriangle\,1.7\times $}}\\[-.2pt]6s\\[-.2pt]{\fontsize{5.8}{6.4}\selectfont \textcolor{ForestGreen}{$ \blacktriangle\,1.7\times $}}\end{tabular} & \hspace*{1pt}\begin{tabular}[t]{@{}c@{}}\rule{0pt}{12pt}6.9s\\[-.2pt]{\fontsize{5.8}{6.4}\selectfont \textcolor{ForestGreen}{$ \blacktriangle\,1.5\times $}}\\[-.2pt]6s\\[-.2pt]{\fontsize{5.8}{6.4}\selectfont \textcolor{ForestGreen}{$ \blacktriangle\,1.4\times $}}\end{tabular} & \hspace*{1pt}\begin{tabular}[t]{@{}c@{}}\rule{0pt}{12pt}12.5s\\[-.2pt]{\fontsize{5.8}{6.4}\selectfont \textcolor{ForestGreen}{$ \blacktriangle\,1.8\times $}}\\[-.2pt]10.4s\\[-.2pt]{\fontsize{5.8}{6.4}\selectfont \textcolor{ForestGreen}{$ \blacktriangle\,1.8\times $}}\end{tabular} & \hspace*{1pt}\begin{tabular}[t]{@{}c@{}}\rule{0pt}{12pt}6s\\[-.2pt]{\fontsize{5.8}{6.4}\selectfont \textcolor{ForestGreen}{$ \blacktriangle\,12.1\times $}}\\[-.2pt]5s\\[-.2pt]{\fontsize{5.8}{6.4}\selectfont \textcolor{ForestGreen}{$ \blacktriangle\,13.0\times $}}\end{tabular} & \hspace*{1pt}\begin{tabular}[t]{@{}c@{}}\rule{0pt}{12pt}4m\,54.5s\\[-.2pt]{\fontsize{5.8}{6.4}\selectfont \textcolor{BrickRed}{$ \blacktriangledown\,29.7\times $}}\\[-.2pt]33.3s\\[-.2pt]{\fontsize{5.8}{6.4}\selectfont \textcolor{BrickRed}{$ \blacktriangledown\,3.8\times $}}\end{tabular} & \hspace*{1pt}\begin{tabular}[t]{@{}c@{}}\rule{0pt}{12pt}4m\,26.3s\\[-.2pt]{\fontsize{5.8}{6.4}\selectfont \textcolor{BrickRed}{$ \blacktriangledown\,28.1\times $}}\\[-.2pt]33.8s\\[-.2pt]{\fontsize{5.8}{6.4}\selectfont \textcolor{BrickRed}{$ \blacktriangledown\,4.4\times $}}\end{tabular} & \hspace*{1pt}\begin{tabular}[t]{@{}c@{}}\rule{0pt}{12pt}8m\,17s\\[-.2pt]{\fontsize{5.8}{6.4}\selectfont \textcolor{BrickRed}{$ \blacktriangledown\,30.1\times $}}\\[-.2pt]54.4s\\[-.2pt]{\fontsize{5.8}{6.4}\selectfont \textcolor{BrickRed}{$ \blacktriangledown\,4.2\times $}}\end{tabular} & \textcolor{gray}{\ensuremath{\cdot}} & \hspace*{1pt}\begin{tabular}[t]{@{}c@{}}\rule{0pt}{12pt}2m\,20.8s\\[-.2pt]{\fontsize{5.8}{6.4}\selectfont \textcolor{BrickRed}{$ \blacktriangledown\,2.7\times $}}\\[-.2pt]1m\,57.4s\\[-.2pt]{\fontsize{5.8}{6.4}\selectfont \textcolor{BrickRed}{$ \blacktriangledown\,2.3\times $}}\end{tabular} & \hspace*{1pt}\begin{tabular}[t]{@{}c@{}}\rule{0pt}{12pt}2m\,25.9s\\[-.2pt]{\fontsize{5.8}{6.4}\selectfont \textcolor{BrickRed}{$ \blacktriangledown\,2.6\times $}}\\[-.2pt]2m\,5.5s\\[-.2pt]{\fontsize{5.8}{6.4}\selectfont \textcolor{BrickRed}{$ \blacktriangledown\,2.3\times $}}\end{tabular} & \textcolor{gray}{\ensuremath{\cdot}} & \hspace*{1pt}\begin{tabular}[t]{@{}c@{}}\rule{0pt}{12pt}2m\,42.2s\\[-.2pt]{\fontsize{5.8}{6.4}\selectfont \textcolor{ForestGreen}{$ \blacktriangle\,2.2\times $}}\\[-.2pt]2m\,23.6s\\[-.2pt]{\fontsize{5.8}{6.4}\selectfont \textcolor{ForestGreen}{$ \blacktriangle\,2.5\times $}}\end{tabular} & \hspace*{1pt}\begin{tabular}[t]{@{}c@{}}\rule{0pt}{12pt}6m\,30.1s\\[-.2pt]{\fontsize{5.8}{6.4}\selectfont \textcolor{BrickRed}{$ \blacktriangledown\,8.1\times $}}\\[-.2pt]2m\,18.5s\\[-.2pt]{\fontsize{5.8}{6.4}\selectfont \textcolor{BrickRed}{$ \blacktriangledown\,2.9\times $}}\end{tabular} & \hspace*{1pt}\begin{tabular}[t]{@{}c@{}}\rule{0pt}{12pt}5m\,34.2s\\[-.2pt]{\fontsize{5.8}{6.4}\selectfont \textcolor{BrickRed}{$ \blacktriangledown\,7.5\times $}}\\[-.2pt]1m\,56.9s\\[-.2pt]{\fontsize{5.8}{6.4}\selectfont \textcolor{BrickRed}{$ \blacktriangledown\,2.7\times $}}\end{tabular} & \textcolor{gray}{\ensuremath{\cdot}} & \hspace*{1pt}\begin{tabular}[t]{@{}c@{}}\rule{0pt}{12pt}11m\,44s\\[-.2pt]{\fontsize{5.8}{6.4}\selectfont \textcolor{BrickRed}{$ \blacktriangledown\,4.2\times $}}\\[-.2pt]2m\,21.6s\\[-.2pt]{\fontsize{5.8}{6.4}\selectfont \textcolor{ForestGreen}{$ \blacktriangle\,13\% $}}\end{tabular}\\[3pt]\hline
\BenchmarkSourceLabel{bench:aws-2-k15-w10:row:mps-bounded-rmq-block}{R28}\shortstack[l]{{\color{gray}R28}} & \begin{tabular}[t]{@{}l@{}}\rule{0pt}{12pt}MPS / Bounded\\RMQ / Block\end{tabular} & \hspace*{1pt}\begin{tabular}[t]{@{}c@{}}\rule{0pt}{12pt}6.5s\\[-.2pt]{\fontsize{5.8}{6.4}\selectfont \textcolor{ForestGreen}{$ \blacktriangle\,1.7\times $}}\\[-.2pt]5.7s\\[-.2pt]{\fontsize{5.8}{6.4}\selectfont \textcolor{ForestGreen}{$ \blacktriangle\,1.8\times $}}\end{tabular} & \hspace*{1pt}\begin{tabular}[t]{@{}c@{}}\rule{0pt}{12pt}6.5s\\[-.2pt]{\fontsize{5.8}{6.4}\selectfont \textcolor{ForestGreen}{$ \blacktriangle\,1.6\times $}}\\[-.2pt]5.8s\\[-.2pt]{\fontsize{5.8}{6.4}\selectfont \textcolor{ForestGreen}{$ \blacktriangle\,1.5\times $}}\end{tabular} & \hspace*{1pt}\begin{tabular}[t]{@{}c@{}}\rule{0pt}{12pt}12.7s\\[-.2pt]{\fontsize{5.8}{6.4}\selectfont \textcolor{ForestGreen}{$ \blacktriangle\,1.8\times $}}\\[-.2pt]10.2s\\[-.2pt]{\fontsize{5.8}{6.4}\selectfont \textcolor{ForestGreen}{$ \blacktriangle\,1.8\times $}}\end{tabular} & \hspace*{1pt}\begin{tabular}[t]{@{}c@{}}\rule{0pt}{12pt}6.1s\\[-.2pt]{\fontsize{5.8}{6.4}\selectfont \textcolor{ForestGreen}{$ \blacktriangle\,12.0\times $}}\\[-.2pt]5.1s\\[-.2pt]{\fontsize{5.8}{6.4}\selectfont \textcolor{ForestGreen}{$ \blacktriangle\,12.8\times $}}\end{tabular} & \hspace*{1pt}\begin{tabular}[t]{@{}c@{}}\rule{0pt}{12pt}4m\,55.9s\\[-.2pt]{\fontsize{5.8}{6.4}\selectfont \textcolor{BrickRed}{$ \blacktriangledown\,29.9\times $}}\\[-.2pt]33.3s\\[-.2pt]{\fontsize{5.8}{6.4}\selectfont \textcolor{BrickRed}{$ \blacktriangledown\,3.8\times $}}\end{tabular} & \hspace*{1pt}\begin{tabular}[t]{@{}c@{}}\rule{0pt}{12pt}4m\,1.2s\\[-.2pt]{\fontsize{5.8}{6.4}\selectfont \textcolor{BrickRed}{$ \blacktriangledown\,25.4\times $}}\\[-.2pt]29.1s\\[-.2pt]{\fontsize{5.8}{6.4}\selectfont \textcolor{BrickRed}{$ \blacktriangledown\,3.8\times $}}\end{tabular} & \hspace*{1pt}\begin{tabular}[t]{@{}c@{}}\rule{0pt}{12pt}7m\,32.9s\\[-.2pt]{\fontsize{5.8}{6.4}\selectfont \textcolor{BrickRed}{$ \blacktriangledown\,27.5\times $}}\\[-.2pt]53.3s\\[-.2pt]{\fontsize{5.8}{6.4}\selectfont \textcolor{BrickRed}{$ \blacktriangledown\,4.1\times $}}\end{tabular} & \textcolor{gray}{\ensuremath{\cdot}} & \hspace*{1pt}\begin{tabular}[t]{@{}c@{}}\rule{0pt}{12pt}2m\,31.3s\\[-.2pt]{\fontsize{5.8}{6.4}\selectfont \textcolor{BrickRed}{$ \blacktriangledown\,2.9\times $}}\\[-.2pt]2m\,9.3s\\[-.2pt]{\fontsize{5.8}{6.4}\selectfont \textcolor{BrickRed}{$ \blacktriangledown\,2.5\times $}}\end{tabular} & \hspace*{1pt}\begin{tabular}[t]{@{}c@{}}\rule{0pt}{12pt}2m\,16.4s\\[-.2pt]{\fontsize{5.8}{6.4}\selectfont \textcolor{BrickRed}{$ \blacktriangledown\,2.4\times $}}\\[-.2pt]1m\,56.6s\\[-.2pt]{\fontsize{5.8}{6.4}\selectfont \textcolor{BrickRed}{$ \blacktriangledown\,2.1\times $}}\end{tabular} & \textcolor{gray}{\ensuremath{\cdot}} & \hspace*{1pt}\begin{tabular}[t]{@{}c@{}}\rule{0pt}{12pt}2m\,44.7s\\[-.2pt]{\fontsize{5.8}{6.4}\selectfont \textcolor{ForestGreen}{$ \blacktriangle\,2.2\times $}}\\[-.2pt]2m\,26.9s\\[-.2pt]{\fontsize{5.8}{6.4}\selectfont \textcolor{ForestGreen}{$ \blacktriangle\,2.4\times $}}\end{tabular} & \hspace*{1pt}\begin{tabular}[t]{@{}c@{}}\rule{0pt}{12pt}6m\,30.9s\\[-.2pt]{\fontsize{5.8}{6.4}\selectfont \textcolor{BrickRed}{$ \blacktriangledown\,8.1\times $}}\\[-.2pt]2m\,9.4s\\[-.2pt]{\fontsize{5.8}{6.4}\selectfont \textcolor{BrickRed}{$ \blacktriangledown\,2.8\times $}}\end{tabular} & \hspace*{1pt}\begin{tabular}[t]{@{}c@{}}\rule{0pt}{12pt}5m\,39.3s\\[-.2pt]{\fontsize{5.8}{6.4}\selectfont \textcolor{BrickRed}{$ \blacktriangledown\,7.6\times $}}\\[-.2pt]2m\,0.7s\\[-.2pt]{\fontsize{5.8}{6.4}\selectfont \textcolor{BrickRed}{$ \blacktriangledown\,2.8\times $}}\end{tabular} & \textcolor{gray}{\ensuremath{\cdot}} & \hspace*{1pt}\begin{tabular}[t]{@{}c@{}}\rule{0pt}{12pt}11m\,44.3s\\[-.2pt]{\fontsize{5.8}{6.4}\selectfont \textcolor{BrickRed}{$ \blacktriangledown\,4.2\times $}}\\[-.2pt]2m\,23.5s\\[-.2pt]{\fontsize{5.8}{6.4}\selectfont \textcolor{ForestGreen}{$ \blacktriangle\,12\% $}}\end{tabular}\\[3pt]\hline
\BenchmarkSourceLabel{bench:aws-2-k15-w10:row:mps-simdref-simd-block}{R29}\shortstack[l]{{\color{gray}R29}} & \begin{tabular}[t]{@{}l@{}}\rule{0pt}{12pt}MPS / SIMD-ref\\SIMD-DP / Block\end{tabular} & \hspace*{1pt}\begin{tabular}[t]{@{}c@{}}\rule{0pt}{12pt}7.2s\\[-.2pt]{\fontsize{5.8}{6.4}\selectfont \textcolor{ForestGreen}{$ \blacktriangle\,1.5\times $}}\\[-.2pt]6.1s\\[-.2pt]{\fontsize{5.8}{6.4}\selectfont \textcolor{ForestGreen}{$ \blacktriangle\,1.6\times $}}\end{tabular} & \hspace*{1pt}\begin{tabular}[t]{@{}c@{}}\rule{0pt}{12pt}6.7s\\[-.2pt]{\fontsize{5.8}{6.4}\selectfont \textcolor{ForestGreen}{$ \blacktriangle\,1.6\times $}}\\[-.2pt]5.7s\\[-.2pt]{\fontsize{5.8}{6.4}\selectfont \textcolor{ForestGreen}{$ \blacktriangle\,1.5\times $}}\end{tabular} & \textcolor{gray}{\ensuremath{\cdot}} & \hspace*{1pt}\begin{tabular}[t]{@{}c@{}}\rule{0pt}{12pt}6.4s\\[-.2pt]{\fontsize{5.8}{6.4}\selectfont \textcolor{ForestGreen}{$ \blacktriangle\,11.4\times $}}\\[-.2pt]5.5s\\[-.2pt]{\fontsize{5.8}{6.4}\selectfont \textcolor{ForestGreen}{$ \blacktriangle\,12.0\times $}}\end{tabular} & \hspace*{1pt}\begin{tabular}[t]{@{}c@{}}\rule{0pt}{12pt}4m\,55.8s\\[-.2pt]{\fontsize{5.8}{6.4}\selectfont \textcolor{BrickRed}{$ \blacktriangledown\,29.9\times $}}\\[-.2pt]34.7s\\[-.2pt]{\fontsize{5.8}{6.4}\selectfont \textcolor{BrickRed}{$ \blacktriangledown\,3.9\times $}}\end{tabular} & \hspace*{1pt}\begin{tabular}[t]{@{}c@{}}\rule{0pt}{12pt}4m\,10.3s\\[-.2pt]{\fontsize{5.8}{6.4}\selectfont \textcolor{BrickRed}{$ \blacktriangledown\,26.4\times $}}\\[-.2pt]31.6s\\[-.2pt]{\fontsize{5.8}{6.4}\selectfont \textcolor{BrickRed}{$ \blacktriangledown\,4.1\times $}}\end{tabular} & \textcolor{gray}{\ensuremath{\cdot}} & \textcolor{gray}{\ensuremath{\cdot}} & \textcolor{gray}{\ensuremath{\cdot}} & \textcolor{gray}{\ensuremath{\cdot}} & \hspace*{1pt}\begin{tabular}[t]{@{}c@{}}\rule{0pt}{12pt}4m\,15.9s\\[-.2pt]{\fontsize{5.8}{6.4}\selectfont \textcolor{BrickRed}{$ \blacktriangledown\,2.2\times $}}\\[-.2pt]3m\,41.9s\\[-.2pt]{\fontsize{5.8}{6.4}\selectfont \textcolor{BrickRed}{$ \blacktriangledown\,2.0\times $}}\end{tabular} & \hspace*{1pt}\begin{tabular}[t]{@{}c@{}}\rule{0pt}{12pt}3m\,19.6s\\[-.2pt]{\fontsize{5.8}{6.4}\selectfont \textcolor{ForestGreen}{$ \blacktriangle\,1.8\times $}}\\[-.2pt]3m\,1.1s\\[-.2pt]{\fontsize{5.8}{6.4}\selectfont \textcolor{ForestGreen}{$ \blacktriangle\,2.0\times $}}\end{tabular} & \textcolor{gray}{\ensuremath{\cdot}} & \textcolor{gray}{\ensuremath{\cdot}} & \textcolor{gray}{\ensuremath{\cdot}} & \textcolor{gray}{\ensuremath{\cdot}}\\[3pt]\hline
\end{tabular}\par}
\smallskip{\footnotesize\noindent Each cell: A and B times, $\cdot$: no measurement. ---\textsuperscript{g}: GPU unavailable. n/a: no comparison.\par
\noindent Blue: one lowest A and one lowest B measurement for each column for each configuration.\par
\noindent Gray: minimap2. Yellow: my preferred MPF / Bounded / SIMD-DP / Block configuration.\par
\noindent MPF/MPS: minimizers per fragment/sequence. SIMD-ref: naive SIMD scan. Bounded: bounded SIMD algorithm.\par
}
\clearpage
\noindent{\large\bfseries Table \thetable: (D): AWS: 2 Gbp, k=15, w=10\hfill 5/7}\par\smallskip
{\fontsize{7.1}{8.1}\selectfont\setlength{\tabcolsep}{1pt}\setlength{\fboxsep}{2pt}\renewcommand{\arraystretch}{1.12}
\noindent\begin{tabular}{@{}p{14mm}p{27mm}|p{13.55mm}|p{13.55mm}|p{13.55mm}|p{13.55mm}|p{13.55mm}|p{13.55mm}|p{13.55mm}|p{13.55mm}|p{13.55mm}|p{13.55mm}|p{13.55mm}|p{13.55mm}|p{13.55mm}|p{13.55mm}|p{13.55mm}|p{13.55mm}|@{}}\hline
Engine & Execution & \multicolumn{8}{c|}{\textbf{Index}} & \multicolumn{8}{c|}{\textbf{No index}}\\
 &  & \multicolumn{8}{c|}{\textbf{2 Gbp}} & \multicolumn{8}{c|}{\textbf{2 Gbp}}\\
 &  & \multicolumn{8}{c|}{\textbf{k=15, w=10}} & \multicolumn{8}{c|}{\textbf{k=15, w=10}}\\
 &  & \multicolumn{4}{c|}{\textbf{DNA4}} & \multicolumn{4}{c|}{\textbf{DNA5}} & \multicolumn{4}{c|}{\textbf{DNA4}} & \multicolumn{4}{c|}{\textbf{DNA5}}\\
 &  & \multicolumn{1}{c|}{\textbf{M7A}} & \multicolumn{1}{c|}{\textbf{M7I}} & \multicolumn{1}{c|}{\textbf{M4}} & \multicolumn{1}{c|}{\textbf{C7G}} & \multicolumn{1}{c|}{\textbf{M7A}} & \multicolumn{1}{c|}{\textbf{M7I}} & \multicolumn{1}{c|}{\textbf{M4}} & \multicolumn{1}{c|}{\textbf{C7G}} & \multicolumn{1}{c|}{\textbf{M7A}} & \multicolumn{1}{c|}{\textbf{M7I}} & \multicolumn{1}{c|}{\textbf{M4}} & \multicolumn{1}{c|}{\textbf{C7G}} & \multicolumn{1}{c|}{\textbf{M7A}} & \multicolumn{1}{c|}{\textbf{M7I}} & \multicolumn{1}{c|}{\textbf{M4}} & \multicolumn{1}{c|}{\textbf{C7G}}\\
 & & {\color{gray}C1} & {\color{gray}C2} & {\color{gray}C3} & {\color{gray}C4} & {\color{gray}C5} & {\color{gray}C6} & {\color{gray}C7} & {\color{gray}C8} & {\color{gray}C9} & {\color{gray}C10} & {\color{gray}C11} & {\color{gray}C12} & {\color{gray}C13} & {\color{gray}C14} & {\color{gray}C15} & {\color{gray}C16}\\
\hline
\rowcolor{BenchBaseline}[0.5pt][0.5pt]\shortstack[l]{{\color{gray}R1}\\\textbf{minimap2}} & \textbf{Baseline} & \hspace*{1pt}\begin{tabular}[t]{@{}c@{}}\rule{0pt}{12pt}11.2s\\[-.2pt]{\fontsize{5.8}{6.4}\selectfont \textcolor{BrickRed}{$ \blacktriangledown\,3.1\times $}}\\[-.2pt]10s\\[-.2pt]{\fontsize{5.8}{6.4}\selectfont \textcolor{BrickRed}{$ \blacktriangledown\,3.6\times $}}\end{tabular} & \hspace*{1pt}\begin{tabular}[t]{@{}c@{}}\rule{0pt}{12pt}10.7s\\[-.2pt]{\fontsize{5.8}{6.4}\selectfont \textcolor{BrickRed}{$ \blacktriangledown\,3.2\times $}}\\[-.2pt]8.6s\\[-.2pt]{\fontsize{5.8}{6.4}\selectfont \textcolor{BrickRed}{$ \blacktriangledown\,3.5\times $}}\end{tabular} & \hspace*{1pt}\begin{tabular}[t]{@{}c@{}}\rule{0pt}{12pt}22.5s\\[-.2pt]{\fontsize{5.8}{6.4}\selectfont \textcolor{BrickRed}{$ \blacktriangledown\,3.6\times $}}\\[-.2pt]18.8s\\[-.2pt]{\fontsize{5.8}{6.4}\selectfont \textcolor{BrickRed}{$ \blacktriangledown\,4.4\times $}}\end{tabular} & \hspace*{1pt}\begin{tabular}[t]{@{}c@{}}\rule{0pt}{12pt}1m\,13.1s\\[-.2pt]{\fontsize{5.8}{6.4}\selectfont \textcolor{BrickRed}{$ \blacktriangledown\,22.5\times $}}\\[-.2pt]1m\,5.7s\\[-.2pt]{\fontsize{5.8}{6.4}\selectfont \textcolor{BrickRed}{$ \blacktriangledown\,27.7\times $}}\end{tabular} & \hspace*{1pt}\begin{tabular}[t]{@{}c@{}}\rule{0pt}{12pt}9.9s\\[-.2pt]{\fontsize{5.8}{6.4}\selectfont \textcolor{BrickRed}{$ \blacktriangledown\,21\% $}}\\[-.2pt]8.8s\\[-.2pt]{\fontsize{5.8}{6.4}\selectfont \textcolor{BrickRed}{$ \blacktriangledown\,1.6\times $}}\end{tabular} & \hspace*{1pt}\begin{tabular}[t]{@{}c@{}}\rule{0pt}{12pt}9.5s\\[-.2pt]{\fontsize{5.8}{6.4}\selectfont \textcolor{BrickRed}{$ \blacktriangledown\,1.4\times $}}\\[-.2pt]7.6s\\[-.2pt]{\fontsize{5.8}{6.4}\selectfont \textcolor{BrickRed}{$ \blacktriangledown\,1.6\times $}}\end{tabular} & \hspace*{1pt}\begin{tabular}[t]{@{}c@{}}\rule{0pt}{12pt}16.5s\\[-.2pt]{\fontsize{5.8}{6.4}\selectfont \textcolor{BrickRed}{$ \blacktriangledown\,1.7\times $}}\\[-.2pt]12.9s\\[-.2pt]{\fontsize{5.8}{6.4}\selectfont \textcolor{BrickRed}{$ \blacktriangledown\,2.0\times $}}\end{tabular} & \hspace*{1pt}\begin{tabular}[t]{@{}c@{}}\rule{0pt}{12pt}28.2s\\[-.2pt]{\fontsize{5.8}{6.4}\selectfont \textcolor{BrickRed}{$ \blacktriangledown\,2.3\times $}}\\[-.2pt]26.2s\\[-.2pt]{\fontsize{5.8}{6.4}\selectfont \textcolor{BrickRed}{$ \blacktriangledown\,3.0\times $}}\end{tabular} & \hspace*{1pt}\begin{tabular}[t]{@{}c@{}}\rule{0pt}{12pt}53s\\[-.2pt]{\fontsize{5.8}{6.4}\selectfont \textcolor{ForestGreen}{$ \blacktriangle\,2.2\times $}}\\[-.2pt]51.9s\\[-.2pt]{\fontsize{5.8}{6.4}\selectfont \textcolor{ForestGreen}{$ \blacktriangle\,2.1\times $}}\end{tabular} & \hspace*{1pt}\begin{tabular}[t]{@{}c@{}}\rule{0pt}{12pt}56.6s\\[-.2pt]{\fontsize{5.8}{6.4}\selectfont \textcolor{ForestGreen}{$ \blacktriangle\,2.1\times $}}\\[-.2pt]54.6s\\[-.2pt]{\fontsize{5.8}{6.4}\selectfont \textcolor{ForestGreen}{$ \blacktriangle\,1.9\times $}}\end{tabular} & \hspace*{1pt}\begin{tabular}[t]{@{}c@{}}\rule{0pt}{12pt}1m\,55.5s\\[-.2pt]{\fontsize{5.8}{6.4}\selectfont \textcolor{ForestGreen}{$ \blacktriangle\,1.8\times $}}\\[-.2pt]1m\,51.8s\\[-.2pt]{\fontsize{5.8}{6.4}\selectfont \textcolor{ForestGreen}{$ \blacktriangle\,1.5\times $}}\end{tabular} & \hspace*{1pt}\begin{tabular}[t]{@{}c@{}}\rule{0pt}{12pt}6m\,1.3s\\[-.2pt]{\fontsize{5.8}{6.4}\selectfont \textcolor{BrickRed}{$ \blacktriangledown\,3.0\times $}}\\[-.2pt]5m\,55.2s\\[-.2pt]{\fontsize{5.8}{6.4}\selectfont \textcolor{BrickRed}{$ \blacktriangledown\,3.1\times $}}\end{tabular} & \hspace*{1pt}\begin{tabular}[t]{@{}c@{}}\rule{0pt}{12pt}48.1s\\[-.2pt]{\fontsize{5.8}{6.4}\selectfont \textcolor{ForestGreen}{$ \blacktriangle\,7.3\times $}}\\[-.2pt]47s\\[-.2pt]{\fontsize{5.8}{6.4}\selectfont \textcolor{ForestGreen}{$ \blacktriangle\,2.4\times $}}\end{tabular} & \hspace*{1pt}\begin{tabular}[t]{@{}c@{}}\rule{0pt}{12pt}44.5s\\[-.2pt]{\fontsize{5.8}{6.4}\selectfont \textcolor{ForestGreen}{$ \blacktriangle\,6.5\times $}}\\[-.2pt]42.6s\\[-.2pt]{\fontsize{5.8}{6.4}\selectfont \textcolor{ForestGreen}{$ \blacktriangle\,2.7\times $}}\end{tabular} & \hspace*{1pt}\begin{tabular}[t]{@{}c@{}}\rule{0pt}{12pt}1m\,30.2s\\[-.2pt]{\fontsize{5.8}{6.4}\selectfont \textcolor{ForestGreen}{$ \blacktriangle\,6.4\times $}}\\[-.2pt]1m\,26.7s\\[-.2pt]{\fontsize{5.8}{6.4}\selectfont \textcolor{ForestGreen}{$ \blacktriangle\,2.2\times $}}\end{tabular} & \hspace*{1pt}\begin{tabular}[t]{@{}c@{}}\rule{0pt}{12pt}2m\,46.4s\\[-.2pt]{\fontsize{5.8}{6.4}\selectfont \textcolor{ForestGreen}{$ \blacktriangle\,4.0\times $}}\\[-.2pt]2m\,42.7s\\[-.2pt]{\fontsize{5.8}{6.4}\selectfont \textcolor{BrickRed}{$ \blacktriangledown\,27\% $}}\end{tabular}\\[3pt]\hline
\BenchmarkSourceLabel{bench:aws-2-k15-w10:row:mps-twostack-simd-block}{R30}\shortstack[l]{{\color{gray}R30}\\\textbf{BioseqDB}} & \begin{tabular}[t]{@{}l@{}}\rule{0pt}{12pt}MPS / Two-stack\\SIMD-DP / Block\end{tabular} & \hspace*{1pt}\begin{tabular}[t]{@{}c@{}}\rule{0pt}{12pt}6.8s\\[-.2pt]{\fontsize{5.8}{6.4}\selectfont \textcolor{ForestGreen}{$ \blacktriangle\,1.6\times $}}\\[-.2pt]6.1s\\[-.2pt]{\fontsize{5.8}{6.4}\selectfont \textcolor{ForestGreen}{$ \blacktriangle\,1.7\times $}}\end{tabular} & \hspace*{1pt}\begin{tabular}[t]{@{}c@{}}\rule{0pt}{12pt}6.8s\\[-.2pt]{\fontsize{5.8}{6.4}\selectfont \textcolor{ForestGreen}{$ \blacktriangle\,1.6\times $}}\\[-.2pt]6s\\[-.2pt]{\fontsize{5.8}{6.4}\selectfont \textcolor{ForestGreen}{$ \blacktriangle\,1.4\times $}}\end{tabular} & \hspace*{1pt}\begin{tabular}[t]{@{}c@{}}\rule{0pt}{12pt}12.6s\\[-.2pt]{\fontsize{5.8}{6.4}\selectfont \textcolor{ForestGreen}{$ \blacktriangle\,1.8\times $}}\\[-.2pt]10.6s\\[-.2pt]{\fontsize{5.8}{6.4}\selectfont \textcolor{ForestGreen}{$ \blacktriangle\,1.8\times $}}\end{tabular} & \hspace*{1pt}\begin{tabular}[t]{@{}c@{}}\rule{0pt}{12pt}6.5s\\[-.2pt]{\fontsize{5.8}{6.4}\selectfont \textcolor{ForestGreen}{$ \blacktriangle\,11.2\times $}}\\[-.2pt]5.7s\\[-.2pt]{\fontsize{5.8}{6.4}\selectfont \textcolor{ForestGreen}{$ \blacktriangle\,11.6\times $}}\end{tabular} & \hspace*{1pt}\begin{tabular}[t]{@{}c@{}}\rule{0pt}{12pt}4m\,47.8s\\[-.2pt]{\fontsize{5.8}{6.4}\selectfont \textcolor{BrickRed}{$ \blacktriangledown\,29.1\times $}}\\[-.2pt]26.2s\\[-.2pt]{\fontsize{5.8}{6.4}\selectfont \textcolor{BrickRed}{$ \blacktriangledown\,3.0\times $}}\end{tabular} & \hspace*{1pt}\begin{tabular}[t]{@{}c@{}}\rule{0pt}{12pt}3m\,58.5s\\[-.2pt]{\fontsize{5.8}{6.4}\selectfont \textcolor{BrickRed}{$ \blacktriangledown\,25.1\times $}}\\[-.2pt]24.4s\\[-.2pt]{\fontsize{5.8}{6.4}\selectfont \textcolor{BrickRed}{$ \blacktriangledown\,3.2\times $}}\end{tabular} & \textcolor{gray}{\ensuremath{\cdot}} & \textcolor{gray}{\ensuremath{\cdot}} & \hspace*{1pt}\begin{tabular}[t]{@{}c@{}}\rule{0pt}{12pt}2m\,46.7s\\[-.2pt]{\fontsize{5.8}{6.4}\selectfont \textcolor{BrickRed}{$ \blacktriangledown\,3.1\times $}}\\[-.2pt]2m\,24.6s\\[-.2pt]{\fontsize{5.8}{6.4}\selectfont \textcolor{BrickRed}{$ \blacktriangledown\,2.8\times $}}\end{tabular} & \hspace*{1pt}\begin{tabular}[t]{@{}c@{}}\rule{0pt}{12pt}2m\,47.4s\\[-.2pt]{\fontsize{5.8}{6.4}\selectfont \textcolor{BrickRed}{$ \blacktriangledown\,3.0\times $}}\\[-.2pt]2m\,27.3s\\[-.2pt]{\fontsize{5.8}{6.4}\selectfont \textcolor{BrickRed}{$ \blacktriangledown\,2.7\times $}}\end{tabular} & \textcolor{gray}{\ensuremath{\cdot}} & \hspace*{1pt}\begin{tabular}[t]{@{}c@{}}\rule{0pt}{12pt}3m\,50.1s\\[-.2pt]{\fontsize{5.8}{6.4}\selectfont \textcolor{ForestGreen}{$ \blacktriangle\,1.6\times $}}\\[-.2pt]3m\,32.7s\\[-.2pt]{\fontsize{5.8}{6.4}\selectfont \textcolor{ForestGreen}{$ \blacktriangle\,1.7\times $}}\end{tabular} & \textcolor{gray}{\ensuremath{\cdot}} & \textcolor{gray}{\ensuremath{\cdot}} & \textcolor{gray}{\ensuremath{\cdot}} & \textcolor{gray}{\ensuremath{\cdot}}\\[3pt]\hline
\BenchmarkSourceLabel{bench:aws-2-k15-w10:row:mps-scalar-simd-block}{R31}\shortstack[l]{{\color{gray}R31}} & \begin{tabular}[t]{@{}l@{}}\rule{0pt}{12pt}MPS / Scalar\\SIMD-DP / Block\end{tabular} & \hspace*{1pt}\begin{tabular}[t]{@{}c@{}}\rule{0pt}{12pt}6.8s\\[-.2pt]{\fontsize{5.8}{6.4}\selectfont \textcolor{ForestGreen}{$ \blacktriangle\,1.6\times $}}\\[-.2pt]6.2s\\[-.2pt]{\fontsize{5.8}{6.4}\selectfont \textcolor{ForestGreen}{$ \blacktriangle\,1.6\times $}}\end{tabular} & \hspace*{1pt}\begin{tabular}[t]{@{}c@{}}\rule{0pt}{12pt}7.1s\\[-.2pt]{\fontsize{5.8}{6.4}\selectfont \textcolor{ForestGreen}{$ \blacktriangle\,1.5\times $}}\\[-.2pt]6.2s\\[-.2pt]{\fontsize{5.8}{6.4}\selectfont \textcolor{ForestGreen}{$ \blacktriangle\,28\% $}}\end{tabular} & \textcolor{gray}{\ensuremath{\cdot}} & \hspace*{1pt}\begin{tabular}[t]{@{}c@{}}\rule{0pt}{12pt}6.6s\\[-.2pt]{\fontsize{5.8}{6.4}\selectfont \textcolor{ForestGreen}{$ \blacktriangle\,11.1\times $}}\\[-.2pt]5.6s\\[-.2pt]{\fontsize{5.8}{6.4}\selectfont \textcolor{ForestGreen}{$ \blacktriangle\,11.7\times $}}\end{tabular} & \hspace*{1pt}\begin{tabular}[t]{@{}c@{}}\rule{0pt}{12pt}4m\,53.5s\\[-.2pt]{\fontsize{5.8}{6.4}\selectfont \textcolor{BrickRed}{$ \blacktriangledown\,29.6\times $}}\\[-.2pt]32.6s\\[-.2pt]{\fontsize{5.8}{6.4}\selectfont \textcolor{BrickRed}{$ \blacktriangledown\,3.7\times $}}\end{tabular} & \hspace*{1pt}\begin{tabular}[t]{@{}c@{}}\rule{0pt}{12pt}4m\,10s\\[-.2pt]{\fontsize{5.8}{6.4}\selectfont \textcolor{BrickRed}{$ \blacktriangledown\,26.4\times $}}\\[-.2pt]31.4s\\[-.2pt]{\fontsize{5.8}{6.4}\selectfont \textcolor{BrickRed}{$ \blacktriangledown\,4.1\times $}}\end{tabular} & \textcolor{gray}{\ensuremath{\cdot}} & \textcolor{gray}{\ensuremath{\cdot}} & \hspace*{1pt}\begin{tabular}[t]{@{}c@{}}\rule{0pt}{12pt}3m\,9.8s\\[-.2pt]{\fontsize{5.8}{6.4}\selectfont \textcolor{BrickRed}{$ \blacktriangledown\,3.6\times $}}\\[-.2pt]2m\,47.5s\\[-.2pt]{\fontsize{5.8}{6.4}\selectfont \textcolor{BrickRed}{$ \blacktriangledown\,3.2\times $}}\end{tabular} & \hspace*{1pt}\begin{tabular}[t]{@{}c@{}}\rule{0pt}{12pt}3m\,6.4s\\[-.2pt]{\fontsize{5.8}{6.4}\selectfont \textcolor{BrickRed}{$ \blacktriangledown\,3.3\times $}}\\[-.2pt]2m\,46.4s\\[-.2pt]{\fontsize{5.8}{6.4}\selectfont \textcolor{BrickRed}{$ \blacktriangledown\,3.0\times $}}\end{tabular} & \textcolor{gray}{\ensuremath{\cdot}} & \hspace*{1pt}\begin{tabular}[t]{@{}c@{}}\rule{0pt}{12pt}3m\,9.9s\\[-.2pt]{\fontsize{5.8}{6.4}\selectfont \textcolor{ForestGreen}{$ \blacktriangle\,1.9\times $}}\\[-.2pt]2m\,52.7s\\[-.2pt]{\fontsize{5.8}{6.4}\selectfont \textcolor{ForestGreen}{$ \blacktriangle\,2.1\times $}}\end{tabular} & \hspace*{1pt}\begin{tabular}[t]{@{}c@{}}\rule{0pt}{12pt}7m\,6.1s\\[-.2pt]{\fontsize{5.8}{6.4}\selectfont \textcolor{BrickRed}{$ \blacktriangledown\,8.9\times $}}\\[-.2pt]2m\,43.9s\\[-.2pt]{\fontsize{5.8}{6.4}\selectfont \textcolor{BrickRed}{$ \blacktriangledown\,3.5\times $}}\end{tabular} & \textcolor{gray}{\ensuremath{\cdot}} & \textcolor{gray}{\ensuremath{\cdot}} & \hspace*{1pt}\begin{tabular}[t]{@{}c@{}}\rule{0pt}{12pt}12m\,35.2s\\[-.2pt]{\fontsize{5.8}{6.4}\selectfont \textcolor{BrickRed}{$ \blacktriangledown\,4.5\times $}}\\[-.2pt]3m\,14.5s\\[-.2pt]{\fontsize{5.8}{6.4}\selectfont \textcolor{BrickRed}{$ \blacktriangledown\,20\% $}}\end{tabular}\\[3pt]\hline
\BenchmarkSourceLabel{bench:aws-2-k15-w10:row:mps-simdref-scalar-block}{R32}\shortstack[l]{{\color{gray}R32}} & \begin{tabular}[t]{@{}l@{}}\rule{0pt}{12pt}MPS / SIMD-ref\\Scalar-DP / Block\end{tabular} & \hspace*{1pt}\begin{tabular}[t]{@{}c@{}}\rule{0pt}{12pt}7.2s\\[-.2pt]{\fontsize{5.8}{6.4}\selectfont \textcolor{ForestGreen}{$ \blacktriangle\,1.6\times $}}\\[-.2pt]6.1s\\[-.2pt]{\fontsize{5.8}{6.4}\selectfont \textcolor{ForestGreen}{$ \blacktriangle\,1.6\times $}}\end{tabular} & \hspace*{1pt}\begin{tabular}[t]{@{}c@{}}\rule{0pt}{12pt}6.8s\\[-.2pt]{\fontsize{5.8}{6.4}\selectfont \textcolor{ForestGreen}{$ \blacktriangle\,1.6\times $}}\\[-.2pt]5.8s\\[-.2pt]{\fontsize{5.8}{6.4}\selectfont \textcolor{ForestGreen}{$ \blacktriangle\,1.5\times $}}\end{tabular} & \textcolor{gray}{\ensuremath{\cdot}} & \textcolor{gray}{\ensuremath{\cdot}} & \hspace*{1pt}\begin{tabular}[t]{@{}c@{}}\rule{0pt}{12pt}4m\,54.2s\\[-.2pt]{\fontsize{5.8}{6.4}\selectfont \textcolor{BrickRed}{$ \blacktriangledown\,29.7\times $}}\\[-.2pt]32.4s\\[-.2pt]{\fontsize{5.8}{6.4}\selectfont \textcolor{BrickRed}{$ \blacktriangledown\,3.7\times $}}\end{tabular} & \hspace*{1pt}\begin{tabular}[t]{@{}c@{}}\rule{0pt}{12pt}4m\,9.1s\\[-.2pt]{\fontsize{5.8}{6.4}\selectfont \textcolor{BrickRed}{$ \blacktriangledown\,26.3\times $}}\\[-.2pt]31.3s\\[-.2pt]{\fontsize{5.8}{6.4}\selectfont \textcolor{BrickRed}{$ \blacktriangledown\,4.1\times $}}\end{tabular} & \textcolor{gray}{\ensuremath{\cdot}} & \textcolor{gray}{\ensuremath{\cdot}} & \textcolor{gray}{\ensuremath{\cdot}} & \textcolor{gray}{\ensuremath{\cdot}} & \textcolor{gray}{\ensuremath{\cdot}} & \textcolor{gray}{\ensuremath{\cdot}} & \textcolor{gray}{\ensuremath{\cdot}} & \textcolor{gray}{\ensuremath{\cdot}} & \textcolor{gray}{\ensuremath{\cdot}} & \textcolor{gray}{\ensuremath{\cdot}}\\[3pt]\hline
\BenchmarkSourceLabel{bench:aws-2-k15-w10:row:mps-simdref-rmq-block}{R33}\shortstack[l]{{\color{gray}R33}} & \begin{tabular}[t]{@{}l@{}}\rule{0pt}{12pt}MPS / SIMD-ref\\RMQ / Block\end{tabular} & \hspace*{1pt}\begin{tabular}[t]{@{}c@{}}\rule{0pt}{12pt}7.1s\\[-.2pt]{\fontsize{5.8}{6.4}\selectfont \textcolor{ForestGreen}{$ \blacktriangle\,1.6\times $}}\\[-.2pt]6s\\[-.2pt]{\fontsize{5.8}{6.4}\selectfont \textcolor{ForestGreen}{$ \blacktriangle\,1.7\times $}}\end{tabular} & \hspace*{1pt}\begin{tabular}[t]{@{}c@{}}\rule{0pt}{12pt}6.8s\\[-.2pt]{\fontsize{5.8}{6.4}\selectfont \textcolor{ForestGreen}{$ \blacktriangle\,1.6\times $}}\\[-.2pt]5.8s\\[-.2pt]{\fontsize{5.8}{6.4}\selectfont \textcolor{ForestGreen}{$ \blacktriangle\,1.5\times $}}\end{tabular} & \textcolor{gray}{\ensuremath{\cdot}} & \textcolor{gray}{\ensuremath{\cdot}} & \hspace*{1pt}\begin{tabular}[t]{@{}c@{}}\rule{0pt}{12pt}4m\,56.7s\\[-.2pt]{\fontsize{5.8}{6.4}\selectfont \textcolor{BrickRed}{$ \blacktriangledown\,29.9\times $}}\\[-.2pt]35.7s\\[-.2pt]{\fontsize{5.8}{6.4}\selectfont \textcolor{BrickRed}{$ \blacktriangledown\,4.1\times $}}\end{tabular} & \hspace*{1pt}\begin{tabular}[t]{@{}c@{}}\rule{0pt}{12pt}4m\,11.8s\\[-.2pt]{\fontsize{5.8}{6.4}\selectfont \textcolor{BrickRed}{$ \blacktriangledown\,26.6\times $}}\\[-.2pt]33.7s\\[-.2pt]{\fontsize{5.8}{6.4}\selectfont \textcolor{BrickRed}{$ \blacktriangledown\,4.4\times $}}\end{tabular} & \textcolor{gray}{\ensuremath{\cdot}} & \textcolor{gray}{\ensuremath{\cdot}} & \textcolor{gray}{\ensuremath{\cdot}} & \textcolor{gray}{\ensuremath{\cdot}} & \hspace*{1pt}\begin{tabular}[t]{@{}c@{}}\rule{0pt}{12pt}3m\,24.8s\\[-.2pt]{\fontsize{5.8}{6.4}\selectfont \textcolor{BrickRed}{$ \blacktriangledown\,1.8\times $}}\\[-.2pt]2m\,59.9s\\[-.2pt]{\fontsize{5.8}{6.4}\selectfont \textcolor{BrickRed}{$ \blacktriangledown\,1.6\times $}}\end{tabular} & \textcolor{gray}{\ensuremath{\cdot}} & \textcolor{gray}{\ensuremath{\cdot}} & \textcolor{gray}{\ensuremath{\cdot}} & \textcolor{gray}{\ensuremath{\cdot}} & \textcolor{gray}{\ensuremath{\cdot}}\\[3pt]\hline
\BenchmarkSourceLabel{bench:aws-2-k15-w10:row:mps-twostack-rmq-block}{R34}\shortstack[l]{{\color{gray}R34}} & \begin{tabular}[t]{@{}l@{}}\rule{0pt}{12pt}MPS / Two-stack\\RMQ / Block\end{tabular} & \hspace*{1pt}\begin{tabular}[t]{@{}c@{}}\rule{0pt}{12pt}6.9s\\[-.2pt]{\fontsize{5.8}{6.4}\selectfont \textcolor{ForestGreen}{$ \blacktriangle\,1.6\times $}}\\[-.2pt]6.2s\\[-.2pt]{\fontsize{5.8}{6.4}\selectfont \textcolor{ForestGreen}{$ \blacktriangle\,1.6\times $}}\end{tabular} & \hspace*{1pt}\begin{tabular}[t]{@{}c@{}}\rule{0pt}{12pt}7.5s\\[-.2pt]{\fontsize{5.8}{6.4}\selectfont \textcolor{ForestGreen}{$ \blacktriangle\,29\% $}}\\[-.2pt]6.6s\\[-.2pt]{\fontsize{5.8}{6.4}\selectfont \textcolor{ForestGreen}{$ \blacktriangle\,23\% $}}\end{tabular} & \textcolor{gray}{\ensuremath{\cdot}} & \textcolor{gray}{\ensuremath{\cdot}} & \hspace*{1pt}\begin{tabular}[t]{@{}c@{}}\rule{0pt}{12pt}4m\,45.1s\\[-.2pt]{\fontsize{5.8}{6.4}\selectfont \textcolor{BrickRed}{$ \blacktriangledown\,28.8\times $}}\\[-.2pt]25.1s\\[-.2pt]{\fontsize{5.8}{6.4}\selectfont \textcolor{BrickRed}{$ \blacktriangledown\,2.8\times $}}\end{tabular} & \hspace*{1pt}\begin{tabular}[t]{@{}c@{}}\rule{0pt}{12pt}3m\,56.8s\\[-.2pt]{\fontsize{5.8}{6.4}\selectfont \textcolor{BrickRed}{$ \blacktriangledown\,25.0\times $}}\\[-.2pt]23.9s\\[-.2pt]{\fontsize{5.8}{6.4}\selectfont \textcolor{BrickRed}{$ \blacktriangledown\,3.1\times $}}\end{tabular} & \textcolor{gray}{\ensuremath{\cdot}} & \textcolor{gray}{\ensuremath{\cdot}} & \hspace*{1pt}\begin{tabular}[t]{@{}c@{}}\rule{0pt}{12pt}2m\,47.3s\\[-.2pt]{\fontsize{5.8}{6.4}\selectfont \textcolor{BrickRed}{$ \blacktriangledown\,3.2\times $}}\\[-.2pt]2m\,24.6s\\[-.2pt]{\fontsize{5.8}{6.4}\selectfont \textcolor{BrickRed}{$ \blacktriangledown\,2.8\times $}}\end{tabular} & \hspace*{1pt}\begin{tabular}[t]{@{}c@{}}\rule{0pt}{12pt}2m\,30.3s\\[-.2pt]{\fontsize{5.8}{6.4}\selectfont \textcolor{BrickRed}{$ \blacktriangledown\,2.7\times $}}\\[-.2pt]2m\,10.8s\\[-.2pt]{\fontsize{5.8}{6.4}\selectfont \textcolor{BrickRed}{$ \blacktriangledown\,2.4\times $}}\end{tabular} & \hspace*{1pt}\begin{tabular}[t]{@{}c@{}}\rule{0pt}{12pt}4m\,45s\\[-.2pt]{\fontsize{5.8}{6.4}\selectfont \textcolor{BrickRed}{$ \blacktriangledown\,2.5\times $}}\\[-.2pt]4m\,11.3s\\[-.2pt]{\fontsize{5.8}{6.4}\selectfont \textcolor{BrickRed}{$ \blacktriangledown\,2.2\times $}}\end{tabular} & \textcolor{gray}{\ensuremath{\cdot}} & \textcolor{gray}{\ensuremath{\cdot}} & \textcolor{gray}{\ensuremath{\cdot}} & \textcolor{gray}{\ensuremath{\cdot}} & \textcolor{gray}{\ensuremath{\cdot}}\\[3pt]\hline
\BenchmarkSourceLabel{bench:aws-2-k15-w10:row:mps-twostack-scalar-block}{R35}\shortstack[l]{{\color{gray}R35}} & \begin{tabular}[t]{@{}l@{}}\rule{0pt}{12pt}MPS / Two-stack\\Scalar-DP / Block\end{tabular} & \hspace*{1pt}\begin{tabular}[t]{@{}c@{}}\rule{0pt}{12pt}6.9s\\[-.2pt]{\fontsize{5.8}{6.4}\selectfont \textcolor{ForestGreen}{$ \blacktriangle\,1.6\times $}}\\[-.2pt]6.2s\\[-.2pt]{\fontsize{5.8}{6.4}\selectfont \textcolor{ForestGreen}{$ \blacktriangle\,1.6\times $}}\end{tabular} & \hspace*{1pt}\begin{tabular}[t]{@{}c@{}}\rule{0pt}{12pt}7.2s\\[-.2pt]{\fontsize{5.8}{6.4}\selectfont \textcolor{ForestGreen}{$ \blacktriangle\,1.5\times $}}\\[-.2pt]6.3s\\[-.2pt]{\fontsize{5.8}{6.4}\selectfont \textcolor{ForestGreen}{$ \blacktriangle\,26\% $}}\end{tabular} & \textcolor{gray}{\ensuremath{\cdot}} & \textcolor{gray}{\ensuremath{\cdot}} & \hspace*{1pt}\begin{tabular}[t]{@{}c@{}}\rule{0pt}{12pt}4m\,44.1s\\[-.2pt]{\fontsize{5.8}{6.4}\selectfont \textcolor{BrickRed}{$ \blacktriangledown\,28.7\times $}}\\[-.2pt]24.5s\\[-.2pt]{\fontsize{5.8}{6.4}\selectfont \textcolor{BrickRed}{$ \blacktriangledown\,2.8\times $}}\end{tabular} & \hspace*{1pt}\begin{tabular}[t]{@{}c@{}}\rule{0pt}{12pt}3m\,57.3s\\[-.2pt]{\fontsize{5.8}{6.4}\selectfont \textcolor{BrickRed}{$ \blacktriangledown\,25.0\times $}}\\[-.2pt]23.7s\\[-.2pt]{\fontsize{5.8}{6.4}\selectfont \textcolor{BrickRed}{$ \blacktriangledown\,3.1\times $}}\end{tabular} & \textcolor{gray}{\ensuremath{\cdot}} & \textcolor{gray}{\ensuremath{\cdot}} & \textcolor{gray}{\ensuremath{\cdot}} & \textcolor{gray}{\ensuremath{\cdot}} & \textcolor{gray}{\ensuremath{\cdot}} & \textcolor{gray}{\ensuremath{\cdot}} & \textcolor{gray}{\ensuremath{\cdot}} & \textcolor{gray}{\ensuremath{\cdot}} & \textcolor{gray}{\ensuremath{\cdot}} & \textcolor{gray}{\ensuremath{\cdot}}\\[3pt]\hline
\BenchmarkSourceLabel{bench:aws-2-k15-w10:row:mps-scalar-rmq-block}{R36}\shortstack[l]{{\color{gray}R36}} & \begin{tabular}[t]{@{}l@{}}\rule{0pt}{12pt}MPS / Scalar\\RMQ / Block\end{tabular} & \hspace*{1pt}\begin{tabular}[t]{@{}c@{}}\rule{0pt}{12pt}6.9s\\[-.2pt]{\fontsize{5.8}{6.4}\selectfont \textcolor{ForestGreen}{$ \blacktriangle\,1.6\times $}}\\[-.2pt]6.2s\\[-.2pt]{\fontsize{5.8}{6.4}\selectfont \textcolor{ForestGreen}{$ \blacktriangle\,1.6\times $}}\end{tabular} & \hspace*{1pt}\begin{tabular}[t]{@{}c@{}}\rule{0pt}{12pt}7s\\[-.2pt]{\fontsize{5.8}{6.4}\selectfont \textcolor{ForestGreen}{$ \blacktriangle\,1.5\times $}}\\[-.2pt]6.2s\\[-.2pt]{\fontsize{5.8}{6.4}\selectfont \textcolor{ForestGreen}{$ \blacktriangle\,28\% $}}\end{tabular} & \textcolor{gray}{\ensuremath{\cdot}} & \textcolor{gray}{\ensuremath{\cdot}} & \hspace*{1pt}\begin{tabular}[t]{@{}c@{}}\rule{0pt}{12pt}4m\,52.3s\\[-.2pt]{\fontsize{5.8}{6.4}\selectfont \textcolor{BrickRed}{$ \blacktriangledown\,29.5\times $}}\\[-.2pt]31.8s\\[-.2pt]{\fontsize{5.8}{6.4}\selectfont \textcolor{BrickRed}{$ \blacktriangledown\,3.6\times $}}\end{tabular} & \hspace*{1pt}\begin{tabular}[t]{@{}c@{}}\rule{0pt}{12pt}4m\,9.2s\\[-.2pt]{\fontsize{5.8}{6.4}\selectfont \textcolor{BrickRed}{$ \blacktriangledown\,26.3\times $}}\\[-.2pt]30.9s\\[-.2pt]{\fontsize{5.8}{6.4}\selectfont \textcolor{BrickRed}{$ \blacktriangledown\,4.0\times $}}\end{tabular} & \textcolor{gray}{\ensuremath{\cdot}} & \textcolor{gray}{\ensuremath{\cdot}} & \hspace*{1pt}\begin{tabular}[t]{@{}c@{}}\rule{0pt}{12pt}3m\,15.8s\\[-.2pt]{\fontsize{5.8}{6.4}\selectfont \textcolor{BrickRed}{$ \blacktriangledown\,3.7\times $}}\\[-.2pt]2m\,53.2s\\[-.2pt]{\fontsize{5.8}{6.4}\selectfont \textcolor{BrickRed}{$ \blacktriangledown\,3.3\times $}}\end{tabular} & \hspace*{1pt}\begin{tabular}[t]{@{}c@{}}\rule{0pt}{12pt}3m\,13.2s\\[-.2pt]{\fontsize{5.8}{6.4}\selectfont \textcolor{BrickRed}{$ \blacktriangledown\,3.4\times $}}\\[-.2pt]2m\,53s\\[-.2pt]{\fontsize{5.8}{6.4}\selectfont \textcolor{BrickRed}{$ \blacktriangledown\,3.2\times $}}\end{tabular} & \textcolor{gray}{\ensuremath{\cdot}} & \textcolor{gray}{\ensuremath{\cdot}} & \hspace*{1pt}\begin{tabular}[t]{@{}c@{}}\rule{0pt}{12pt}7m\,7.4s\\[-.2pt]{\fontsize{5.8}{6.4}\selectfont \textcolor{BrickRed}{$ \blacktriangledown\,8.9\times $}}\\[-.2pt]2m\,45.6s\\[-.2pt]{\fontsize{5.8}{6.4}\selectfont \textcolor{BrickRed}{$ \blacktriangledown\,3.5\times $}}\end{tabular} & \textcolor{gray}{\ensuremath{\cdot}} & \textcolor{gray}{\ensuremath{\cdot}} & \textcolor{gray}{\ensuremath{\cdot}}\\[3pt]\hline
\end{tabular}\par}
\smallskip{\footnotesize\noindent Each cell: A and B times, $\cdot$: no measurement. ---\textsuperscript{g}: GPU unavailable. n/a: no comparison.\par
\noindent Blue: one lowest A and one lowest B measurement for each column for each configuration.\par
\noindent Gray: minimap2. Yellow: my preferred MPF / Bounded / SIMD-DP / Block configuration.\par
\noindent MPF/MPS: minimizers per fragment/sequence. SIMD-ref: naive SIMD scan. Bounded: bounded SIMD algorithm.\par
}
\clearpage
\noindent{\large\bfseries Table \thetable: (D): AWS: 2 Gbp, k=15, w=10\hfill 6/7}\par\smallskip
{\fontsize{7.1}{8.1}\selectfont\setlength{\tabcolsep}{1pt}\setlength{\fboxsep}{2pt}\renewcommand{\arraystretch}{1.12}
\noindent\begin{tabular}{@{}p{14mm}p{27mm}|p{13.55mm}|p{13.55mm}|p{13.55mm}|p{13.55mm}|p{13.55mm}|p{13.55mm}|p{13.55mm}|p{13.55mm}|p{13.55mm}|p{13.55mm}|p{13.55mm}|p{13.55mm}|p{13.55mm}|p{13.55mm}|p{13.55mm}|p{13.55mm}|@{}}\hline
Engine & Execution & \multicolumn{8}{c|}{\textbf{Index}} & \multicolumn{8}{c|}{\textbf{No index}}\\
 &  & \multicolumn{8}{c|}{\textbf{2 Gbp}} & \multicolumn{8}{c|}{\textbf{2 Gbp}}\\
 &  & \multicolumn{8}{c|}{\textbf{k=15, w=10}} & \multicolumn{8}{c|}{\textbf{k=15, w=10}}\\
 &  & \multicolumn{4}{c|}{\textbf{DNA4}} & \multicolumn{4}{c|}{\textbf{DNA5}} & \multicolumn{4}{c|}{\textbf{DNA4}} & \multicolumn{4}{c|}{\textbf{DNA5}}\\
 &  & \multicolumn{1}{c|}{\textbf{M7A}} & \multicolumn{1}{c|}{\textbf{M7I}} & \multicolumn{1}{c|}{\textbf{M4}} & \multicolumn{1}{c|}{\textbf{C7G}} & \multicolumn{1}{c|}{\textbf{M7A}} & \multicolumn{1}{c|}{\textbf{M7I}} & \multicolumn{1}{c|}{\textbf{M4}} & \multicolumn{1}{c|}{\textbf{C7G}} & \multicolumn{1}{c|}{\textbf{M7A}} & \multicolumn{1}{c|}{\textbf{M7I}} & \multicolumn{1}{c|}{\textbf{M4}} & \multicolumn{1}{c|}{\textbf{C7G}} & \multicolumn{1}{c|}{\textbf{M7A}} & \multicolumn{1}{c|}{\textbf{M7I}} & \multicolumn{1}{c|}{\textbf{M4}} & \multicolumn{1}{c|}{\textbf{C7G}}\\
 & & {\color{gray}C1} & {\color{gray}C2} & {\color{gray}C3} & {\color{gray}C4} & {\color{gray}C5} & {\color{gray}C6} & {\color{gray}C7} & {\color{gray}C8} & {\color{gray}C9} & {\color{gray}C10} & {\color{gray}C11} & {\color{gray}C12} & {\color{gray}C13} & {\color{gray}C14} & {\color{gray}C15} & {\color{gray}C16}\\
\hline
\rowcolor{BenchBaseline}[0.5pt][0.5pt]\shortstack[l]{{\color{gray}R1}\\\textbf{minimap2}} & \textbf{Baseline} & \hspace*{1pt}\begin{tabular}[t]{@{}c@{}}\rule{0pt}{12pt}11.2s\\[-.2pt]{\fontsize{5.8}{6.4}\selectfont \textcolor{BrickRed}{$ \blacktriangledown\,3.1\times $}}\\[-.2pt]10s\\[-.2pt]{\fontsize{5.8}{6.4}\selectfont \textcolor{BrickRed}{$ \blacktriangledown\,3.6\times $}}\end{tabular} & \hspace*{1pt}\begin{tabular}[t]{@{}c@{}}\rule{0pt}{12pt}10.7s\\[-.2pt]{\fontsize{5.8}{6.4}\selectfont \textcolor{BrickRed}{$ \blacktriangledown\,3.2\times $}}\\[-.2pt]8.6s\\[-.2pt]{\fontsize{5.8}{6.4}\selectfont \textcolor{BrickRed}{$ \blacktriangledown\,3.5\times $}}\end{tabular} & \hspace*{1pt}\begin{tabular}[t]{@{}c@{}}\rule{0pt}{12pt}22.5s\\[-.2pt]{\fontsize{5.8}{6.4}\selectfont \textcolor{BrickRed}{$ \blacktriangledown\,3.6\times $}}\\[-.2pt]18.8s\\[-.2pt]{\fontsize{5.8}{6.4}\selectfont \textcolor{BrickRed}{$ \blacktriangledown\,4.4\times $}}\end{tabular} & \hspace*{1pt}\begin{tabular}[t]{@{}c@{}}\rule{0pt}{12pt}1m\,13.1s\\[-.2pt]{\fontsize{5.8}{6.4}\selectfont \textcolor{BrickRed}{$ \blacktriangledown\,22.5\times $}}\\[-.2pt]1m\,5.7s\\[-.2pt]{\fontsize{5.8}{6.4}\selectfont \textcolor{BrickRed}{$ \blacktriangledown\,27.7\times $}}\end{tabular} & \hspace*{1pt}\begin{tabular}[t]{@{}c@{}}\rule{0pt}{12pt}9.9s\\[-.2pt]{\fontsize{5.8}{6.4}\selectfont \textcolor{BrickRed}{$ \blacktriangledown\,21\% $}}\\[-.2pt]8.8s\\[-.2pt]{\fontsize{5.8}{6.4}\selectfont \textcolor{BrickRed}{$ \blacktriangledown\,1.6\times $}}\end{tabular} & \hspace*{1pt}\begin{tabular}[t]{@{}c@{}}\rule{0pt}{12pt}9.5s\\[-.2pt]{\fontsize{5.8}{6.4}\selectfont \textcolor{BrickRed}{$ \blacktriangledown\,1.4\times $}}\\[-.2pt]7.6s\\[-.2pt]{\fontsize{5.8}{6.4}\selectfont \textcolor{BrickRed}{$ \blacktriangledown\,1.6\times $}}\end{tabular} & \hspace*{1pt}\begin{tabular}[t]{@{}c@{}}\rule{0pt}{12pt}16.5s\\[-.2pt]{\fontsize{5.8}{6.4}\selectfont \textcolor{BrickRed}{$ \blacktriangledown\,1.7\times $}}\\[-.2pt]12.9s\\[-.2pt]{\fontsize{5.8}{6.4}\selectfont \textcolor{BrickRed}{$ \blacktriangledown\,2.0\times $}}\end{tabular} & \hspace*{1pt}\begin{tabular}[t]{@{}c@{}}\rule{0pt}{12pt}28.2s\\[-.2pt]{\fontsize{5.8}{6.4}\selectfont \textcolor{BrickRed}{$ \blacktriangledown\,2.3\times $}}\\[-.2pt]26.2s\\[-.2pt]{\fontsize{5.8}{6.4}\selectfont \textcolor{BrickRed}{$ \blacktriangledown\,3.0\times $}}\end{tabular} & \hspace*{1pt}\begin{tabular}[t]{@{}c@{}}\rule{0pt}{12pt}53s\\[-.2pt]{\fontsize{5.8}{6.4}\selectfont \textcolor{ForestGreen}{$ \blacktriangle\,2.2\times $}}\\[-.2pt]51.9s\\[-.2pt]{\fontsize{5.8}{6.4}\selectfont \textcolor{ForestGreen}{$ \blacktriangle\,2.1\times $}}\end{tabular} & \hspace*{1pt}\begin{tabular}[t]{@{}c@{}}\rule{0pt}{12pt}56.6s\\[-.2pt]{\fontsize{5.8}{6.4}\selectfont \textcolor{ForestGreen}{$ \blacktriangle\,2.1\times $}}\\[-.2pt]54.6s\\[-.2pt]{\fontsize{5.8}{6.4}\selectfont \textcolor{ForestGreen}{$ \blacktriangle\,1.9\times $}}\end{tabular} & \hspace*{1pt}\begin{tabular}[t]{@{}c@{}}\rule{0pt}{12pt}1m\,55.5s\\[-.2pt]{\fontsize{5.8}{6.4}\selectfont \textcolor{ForestGreen}{$ \blacktriangle\,1.8\times $}}\\[-.2pt]1m\,51.8s\\[-.2pt]{\fontsize{5.8}{6.4}\selectfont \textcolor{ForestGreen}{$ \blacktriangle\,1.5\times $}}\end{tabular} & \hspace*{1pt}\begin{tabular}[t]{@{}c@{}}\rule{0pt}{12pt}6m\,1.3s\\[-.2pt]{\fontsize{5.8}{6.4}\selectfont \textcolor{BrickRed}{$ \blacktriangledown\,3.0\times $}}\\[-.2pt]5m\,55.2s\\[-.2pt]{\fontsize{5.8}{6.4}\selectfont \textcolor{BrickRed}{$ \blacktriangledown\,3.1\times $}}\end{tabular} & \hspace*{1pt}\begin{tabular}[t]{@{}c@{}}\rule{0pt}{12pt}48.1s\\[-.2pt]{\fontsize{5.8}{6.4}\selectfont \textcolor{ForestGreen}{$ \blacktriangle\,7.3\times $}}\\[-.2pt]47s\\[-.2pt]{\fontsize{5.8}{6.4}\selectfont \textcolor{ForestGreen}{$ \blacktriangle\,2.4\times $}}\end{tabular} & \hspace*{1pt}\begin{tabular}[t]{@{}c@{}}\rule{0pt}{12pt}44.5s\\[-.2pt]{\fontsize{5.8}{6.4}\selectfont \textcolor{ForestGreen}{$ \blacktriangle\,6.5\times $}}\\[-.2pt]42.6s\\[-.2pt]{\fontsize{5.8}{6.4}\selectfont \textcolor{ForestGreen}{$ \blacktriangle\,2.7\times $}}\end{tabular} & \hspace*{1pt}\begin{tabular}[t]{@{}c@{}}\rule{0pt}{12pt}1m\,30.2s\\[-.2pt]{\fontsize{5.8}{6.4}\selectfont \textcolor{ForestGreen}{$ \blacktriangle\,6.4\times $}}\\[-.2pt]1m\,26.7s\\[-.2pt]{\fontsize{5.8}{6.4}\selectfont \textcolor{ForestGreen}{$ \blacktriangle\,2.2\times $}}\end{tabular} & \hspace*{1pt}\begin{tabular}[t]{@{}c@{}}\rule{0pt}{12pt}2m\,46.4s\\[-.2pt]{\fontsize{5.8}{6.4}\selectfont \textcolor{ForestGreen}{$ \blacktriangle\,4.0\times $}}\\[-.2pt]2m\,42.7s\\[-.2pt]{\fontsize{5.8}{6.4}\selectfont \textcolor{BrickRed}{$ \blacktriangledown\,27\% $}}\end{tabular}\\[3pt]\hline
\BenchmarkSourceLabel{bench:aws-2-k15-w10:row:mps-scalar-scalar-block}{R37}\shortstack[l]{{\color{gray}R37}\\\textbf{BioseqDB}} & \begin{tabular}[t]{@{}l@{}}\rule{0pt}{12pt}MPS / Scalar\\Scalar-DP / Block\end{tabular} & \hspace*{1pt}\begin{tabular}[t]{@{}c@{}}\rule{0pt}{12pt}7.5s\\[-.2pt]{\fontsize{5.8}{6.4}\selectfont \textcolor{ForestGreen}{$ \blacktriangle\,1.5\times $}}\\[-.2pt]6.3s\\[-.2pt]{\fontsize{5.8}{6.4}\selectfont \textcolor{ForestGreen}{$ \blacktriangle\,1.6\times $}}\end{tabular} & \hspace*{1pt}\begin{tabular}[t]{@{}c@{}}\rule{0pt}{12pt}7.2s\\[-.2pt]{\fontsize{5.8}{6.4}\selectfont \textcolor{ForestGreen}{$ \blacktriangle\,1.5\times $}}\\[-.2pt]6.4s\\[-.2pt]{\fontsize{5.8}{6.4}\selectfont \textcolor{ForestGreen}{$ \blacktriangle\,26\% $}}\end{tabular} & \textcolor{gray}{\ensuremath{\cdot}} & \textcolor{gray}{\ensuremath{\cdot}} & \hspace*{1pt}\begin{tabular}[t]{@{}c@{}}\rule{0pt}{12pt}4m\,52s\\[-.2pt]{\fontsize{5.8}{6.4}\selectfont \textcolor{BrickRed}{$ \blacktriangledown\,29.5\times $}}\\[-.2pt]31.3s\\[-.2pt]{\fontsize{5.8}{6.4}\selectfont \textcolor{BrickRed}{$ \blacktriangledown\,3.6\times $}}\end{tabular} & \hspace*{1pt}\begin{tabular}[t]{@{}c@{}}\rule{0pt}{12pt}4m\,9.3s\\[-.2pt]{\fontsize{5.8}{6.4}\selectfont \textcolor{BrickRed}{$ \blacktriangledown\,26.3\times $}}\\[-.2pt]30.8s\\[-.2pt]{\fontsize{5.8}{6.4}\selectfont \textcolor{BrickRed}{$ \blacktriangledown\,4.0\times $}}\end{tabular} & \textcolor{gray}{\ensuremath{\cdot}} & \textcolor{gray}{\ensuremath{\cdot}} & \hspace*{1pt}\begin{tabular}[t]{@{}c@{}}\rule{0pt}{12pt}3m\,13.5s\\[-.2pt]{\fontsize{5.8}{6.4}\selectfont \textcolor{BrickRed}{$ \blacktriangledown\,3.6\times $}}\\[-.2pt]2m\,47.5s\\[-.2pt]{\fontsize{5.8}{6.4}\selectfont \textcolor{BrickRed}{$ \blacktriangledown\,3.2\times $}}\end{tabular} & \hspace*{1pt}\begin{tabular}[t]{@{}c@{}}\rule{0pt}{12pt}3m\,16.5s\\[-.2pt]{\fontsize{5.8}{6.4}\selectfont \textcolor{BrickRed}{$ \blacktriangledown\,3.5\times $}}\\[-.2pt]2m\,55.3s\\[-.2pt]{\fontsize{5.8}{6.4}\selectfont \textcolor{BrickRed}{$ \blacktriangledown\,3.2\times $}}\end{tabular} & \textcolor{gray}{\ensuremath{\cdot}} & \textcolor{gray}{\ensuremath{\cdot}} & \hspace*{1pt}\begin{tabular}[t]{@{}c@{}}\rule{0pt}{12pt}7m\,3.9s\\[-.2pt]{\fontsize{5.8}{6.4}\selectfont \textcolor{BrickRed}{$ \blacktriangledown\,8.8\times $}}\\[-.2pt]2m\,41.8s\\[-.2pt]{\fontsize{5.8}{6.4}\selectfont \textcolor{BrickRed}{$ \blacktriangledown\,3.4\times $}}\end{tabular} & \textcolor{gray}{\ensuremath{\cdot}} & \textcolor{gray}{\ensuremath{\cdot}} & \textcolor{gray}{\ensuremath{\cdot}}\\[3pt]\hline
\BenchmarkSourceLabel{bench:aws-2-k15-w10:row:mps-bounded-scalar-ksw2}{R38}\shortstack[l]{{\color{gray}R38}} & \begin{tabular}[t]{@{}l@{}}\rule{0pt}{12pt}MPS / Bounded\\Scalar-DP / KSW2\end{tabular} & \hspace*{1pt}\begin{tabular}[t]{@{}c@{}}\rule{0pt}{12pt}7.9s\\[-.2pt]{\fontsize{5.8}{6.4}\selectfont \textcolor{ForestGreen}{$ \blacktriangle\,29\% $}}\\[-.2pt]4.9s\\[-.2pt]{\fontsize{5.8}{6.4}\selectfont \textcolor{ForestGreen}{$ \blacktriangle\,2.0\times $}}\end{tabular} & \hspace*{1pt}\begin{tabular}[t]{@{}c@{}}\rule{0pt}{12pt}8.6s\\[-.2pt]{\fontsize{5.8}{6.4}\selectfont \textcolor{ForestGreen}{$ \blacktriangle\,20\% $}}\\[-.2pt]5s\\[-.2pt]{\fontsize{5.8}{6.4}\selectfont \textcolor{ForestGreen}{$ \blacktriangle\,1.7\times $}}\end{tabular} & \hspace*{1pt}\begin{tabular}[t]{@{}c@{}}\rule{0pt}{12pt}17.2s\\[-.2pt]{\fontsize{5.8}{6.4}\selectfont \textcolor{ForestGreen}{$ \blacktriangle\,24\% $}}\\[-.2pt]8.9s\\[-.2pt]{\fontsize{5.8}{6.4}\selectfont \textcolor{ForestGreen}{$ \blacktriangle\,2.1\times $}}\end{tabular} & \hspace*{1pt}\begin{tabular}[t]{@{}c@{}}\rule{0pt}{12pt}8s\\[-.2pt]{\fontsize{5.8}{6.4}\selectfont \textcolor{ForestGreen}{$ \blacktriangle\,9.2\times $}}\\[-.2pt]4.6s\\[-.2pt]{\fontsize{5.8}{6.4}\selectfont \textcolor{ForestGreen}{$ \blacktriangle\,14.2\times $}}\end{tabular} & \hspace*{1pt}\begin{tabular}[t]{@{}c@{}}\rule{0pt}{12pt}4m\,56.3s\\[-.2pt]{\fontsize{5.8}{6.4}\selectfont \textcolor{BrickRed}{$ \blacktriangledown\,29.9\times $}}\\[-.2pt]32.9s\\[-.2pt]{\fontsize{5.8}{6.4}\selectfont \textcolor{BrickRed}{$ \blacktriangledown\,3.7\times $}}\end{tabular} & \hspace*{1pt}\begin{tabular}[t]{@{}c@{}}\rule{0pt}{12pt}4m\,2.1s\\[-.2pt]{\fontsize{5.8}{6.4}\selectfont \textcolor{BrickRed}{$ \blacktriangledown\,25.5\times $}}\\[-.2pt]28.6s\\[-.2pt]{\fontsize{5.8}{6.4}\selectfont \textcolor{BrickRed}{$ \blacktriangledown\,3.7\times $}}\end{tabular} & \hspace*{1pt}\begin{tabular}[t]{@{}c@{}}\rule{0pt}{12pt}8m\,11.2s\\[-.2pt]{\fontsize{5.8}{6.4}\selectfont \textcolor{BrickRed}{$ \blacktriangledown\,29.8\times $}}\\[-.2pt]53.8s\\[-.2pt]{\fontsize{5.8}{6.4}\selectfont \textcolor{BrickRed}{$ \blacktriangledown\,4.2\times $}}\end{tabular} & \textcolor{gray}{\ensuremath{\cdot}} & \hspace*{1pt}\begin{tabular}[t]{@{}c@{}}\rule{0pt}{12pt}2m\,37s\\[-.2pt]{\fontsize{5.8}{6.4}\selectfont \textcolor{BrickRed}{$ \blacktriangledown\,3.0\times $}}\\[-.2pt]2m\,12s\\[-.2pt]{\fontsize{5.8}{6.4}\selectfont \textcolor{BrickRed}{$ \blacktriangledown\,2.5\times $}}\end{tabular} & \hspace*{1pt}\begin{tabular}[t]{@{}c@{}}\rule{0pt}{12pt}2m\,30.6s\\[-.2pt]{\fontsize{5.8}{6.4}\selectfont \textcolor{BrickRed}{$ \blacktriangledown\,2.7\times $}}\\[-.2pt]2m\,7.1s\\[-.2pt]{\fontsize{5.8}{6.4}\selectfont \textcolor{BrickRed}{$ \blacktriangledown\,2.3\times $}}\end{tabular} & \textcolor{gray}{\ensuremath{\cdot}} & \hspace*{1pt}\begin{tabular}[t]{@{}c@{}}\rule{0pt}{12pt}2m\,29s\\[-.2pt]{\fontsize{5.8}{6.4}\selectfont \textcolor{ForestGreen}{$ \blacktriangle\,2.4\times $}}\\[-.2pt]2m\,7.2s\\[-.2pt]{\fontsize{5.8}{6.4}\selectfont \textcolor{ForestGreen}{$ \blacktriangle\,2.8\times $}}\end{tabular} & \hspace*{1pt}\begin{tabular}[t]{@{}c@{}}\rule{0pt}{12pt}6m\,26.9s\\[-.2pt]{\fontsize{5.8}{6.4}\selectfont \textcolor{BrickRed}{$ \blacktriangledown\,8.1\times $}}\\[-.2pt]2m\,3.3s\\[-.2pt]{\fontsize{5.8}{6.4}\selectfont \textcolor{BrickRed}{$ \blacktriangledown\,2.6\times $}}\end{tabular} & \hspace*{1pt}\begin{tabular}[t]{@{}c@{}}\rule{0pt}{12pt}5m\,37.5s\\[-.2pt]{\fontsize{5.8}{6.4}\selectfont \textcolor{BrickRed}{$ \blacktriangledown\,7.6\times $}}\\[-.2pt]1m\,58.7s\\[-.2pt]{\fontsize{5.8}{6.4}\selectfont \textcolor{BrickRed}{$ \blacktriangledown\,2.8\times $}}\end{tabular} & \textcolor{gray}{\ensuremath{\cdot}} & \hspace*{1pt}\begin{tabular}[t]{@{}c@{}}\rule{0pt}{12pt}12m\,0.3s\\[-.2pt]{\fontsize{5.8}{6.4}\selectfont \textcolor{BrickRed}{$ \blacktriangledown\,4.3\times $}}\\[-.2pt]2m\,37s\\[-.2pt]{\fontsize{5.8}{6.4}\selectfont \textcolor{gray}{$ \blacktriangleright\,4\% $}}\end{tabular}\\[3pt]\hline
\BenchmarkSourceLabel{bench:aws-2-k15-w10:row:mps-bounded-rmq-ksw2}{R39}\shortstack[l]{{\color{gray}R39}} & \begin{tabular}[t]{@{}l@{}}\rule{0pt}{12pt}MPS / Bounded\\RMQ / KSW2\end{tabular} & \hspace*{1pt}\begin{tabular}[t]{@{}c@{}}\rule{0pt}{12pt}7.9s\\[-.2pt]{\fontsize{5.8}{6.4}\selectfont \textcolor{ForestGreen}{$ \blacktriangle\,29\% $}}\\[-.2pt]4.9s\\[-.2pt]{\fontsize{5.8}{6.4}\selectfont \textcolor{ForestGreen}{$ \blacktriangle\,2.0\times $}}\end{tabular} & \hspace*{1pt}\begin{tabular}[t]{@{}c@{}}\rule{0pt}{12pt}8.5s\\[-.2pt]{\fontsize{5.8}{6.4}\selectfont \textcolor{ForestGreen}{$ \blacktriangle\,20\% $}}\\[-.2pt]4.9s\\[-.2pt]{\fontsize{5.8}{6.4}\selectfont \textcolor{ForestGreen}{$ \blacktriangle\,1.7\times $}}\end{tabular} & \hspace*{1pt}\begin{tabular}[t]{@{}c@{}}\rule{0pt}{12pt}17.4s\\[-.2pt]{\fontsize{5.8}{6.4}\selectfont \textcolor{ForestGreen}{$ \blacktriangle\,23\% $}}\\[-.2pt]8.8s\\[-.2pt]{\fontsize{5.8}{6.4}\selectfont \textcolor{ForestGreen}{$ \blacktriangle\,2.1\times $}}\end{tabular} & \hspace*{1pt}\begin{tabular}[t]{@{}c@{}}\rule{0pt}{12pt}8.1s\\[-.2pt]{\fontsize{5.8}{6.4}\selectfont \textcolor{ForestGreen}{$ \blacktriangle\,9.1\times $}}\\[-.2pt]4.6s\\[-.2pt]{\fontsize{5.8}{6.4}\selectfont \textcolor{ForestGreen}{$ \blacktriangle\,14.2\times $}}\end{tabular} & \hspace*{1pt}\begin{tabular}[t]{@{}c@{}}\rule{0pt}{12pt}4m\,56.6s\\[-.2pt]{\fontsize{5.8}{6.4}\selectfont \textcolor{BrickRed}{$ \blacktriangledown\,29.9\times $}}\\[-.2pt]32.7s\\[-.2pt]{\fontsize{5.8}{6.4}\selectfont \textcolor{BrickRed}{$ \blacktriangledown\,3.7\times $}}\end{tabular} & \hspace*{1pt}\begin{tabular}[t]{@{}c@{}}\rule{0pt}{12pt}4m\,2.8s\\[-.2pt]{\fontsize{5.8}{6.4}\selectfont \textcolor{BrickRed}{$ \blacktriangledown\,25.6\times $}}\\[-.2pt]28.9s\\[-.2pt]{\fontsize{5.8}{6.4}\selectfont \textcolor{BrickRed}{$ \blacktriangledown\,3.8\times $}}\end{tabular} & \hspace*{1pt}\begin{tabular}[t]{@{}c@{}}\rule{0pt}{12pt}7m\,36.6s\\[-.2pt]{\fontsize{5.8}{6.4}\selectfont \textcolor{BrickRed}{$ \blacktriangledown\,27.7\times $}}\\[-.2pt]53.2s\\[-.2pt]{\fontsize{5.8}{6.4}\selectfont \textcolor{BrickRed}{$ \blacktriangledown\,4.1\times $}}\end{tabular} & \textcolor{gray}{\ensuremath{\cdot}} & \hspace*{1pt}\begin{tabular}[t]{@{}c@{}}\rule{0pt}{12pt}2m\,21.9s\\[-.2pt]{\fontsize{5.8}{6.4}\selectfont \textcolor{BrickRed}{$ \blacktriangledown\,2.7\times $}}\\[-.2pt]1m\,57.6s\\[-.2pt]{\fontsize{5.8}{6.4}\selectfont \textcolor{BrickRed}{$ \blacktriangledown\,2.3\times $}}\end{tabular} & \hspace*{1pt}\begin{tabular}[t]{@{}c@{}}\rule{0pt}{12pt}2m\,27.9s\\[-.2pt]{\fontsize{5.8}{6.4}\selectfont \textcolor{BrickRed}{$ \blacktriangledown\,2.6\times $}}\\[-.2pt]2m\,5.4s\\[-.2pt]{\fontsize{5.8}{6.4}\selectfont \textcolor{BrickRed}{$ \blacktriangledown\,2.3\times $}}\end{tabular} & \textcolor{gray}{\ensuremath{\cdot}} & \hspace*{1pt}\begin{tabular}[t]{@{}c@{}}\rule{0pt}{12pt}2m\,29.1s\\[-.2pt]{\fontsize{5.8}{6.4}\selectfont \textcolor{ForestGreen}{$ \blacktriangle\,2.4\times $}}\\[-.2pt]2m\,8.2s\\[-.2pt]{\fontsize{5.8}{6.4}\selectfont \textcolor{ForestGreen}{$ \blacktriangle\,2.8\times $}}\end{tabular} & \hspace*{1pt}\begin{tabular}[t]{@{}c@{}}\rule{0pt}{12pt}6m\,41.2s\\[-.2pt]{\fontsize{5.8}{6.4}\selectfont \textcolor{BrickRed}{$ \blacktriangledown\,8.3\times $}}\\[-.2pt]2m\,18.2s\\[-.2pt]{\fontsize{5.8}{6.4}\selectfont \textcolor{BrickRed}{$ \blacktriangledown\,2.9\times $}}\end{tabular} & \hspace*{1pt}\begin{tabular}[t]{@{}c@{}}\rule{0pt}{12pt}5m\,35.7s\\[-.2pt]{\fontsize{5.8}{6.4}\selectfont \textcolor{BrickRed}{$ \blacktriangledown\,7.5\times $}}\\[-.2pt]1m\,59.6s\\[-.2pt]{\fontsize{5.8}{6.4}\selectfont \textcolor{BrickRed}{$ \blacktriangledown\,2.8\times $}}\end{tabular} & \textcolor{gray}{\ensuremath{\cdot}} & \textcolor{gray}{\ensuremath{\cdot}}\\[3pt]\hline
\BenchmarkSourceLabel{bench:aws-2-k15-w10:row:mps-twostack-rmq-ksw2}{R40}\shortstack[l]{{\color{gray}R40}} & \begin{tabular}[t]{@{}l@{}}\rule{0pt}{12pt}MPS / Two-stack\\RMQ / KSW2\end{tabular} & \hspace*{1pt}\begin{tabular}[t]{@{}c@{}}\rule{0pt}{12pt}8.3s\\[-.2pt]{\fontsize{5.8}{6.4}\selectfont \textcolor{ForestGreen}{$ \blacktriangle\,26\% $}}\\[-.2pt]5.3s\\[-.2pt]{\fontsize{5.8}{6.4}\selectfont \textcolor{ForestGreen}{$ \blacktriangle\,1.9\times $}}\end{tabular} & \hspace*{1pt}\begin{tabular}[t]{@{}c@{}}\rule{0pt}{12pt}8.1s\\[-.2pt]{\fontsize{5.8}{6.4}\selectfont \textcolor{ForestGreen}{$ \blacktriangle\,23\% $}}\\[-.2pt]5.1s\\[-.2pt]{\fontsize{5.8}{6.4}\selectfont \textcolor{ForestGreen}{$ \blacktriangle\,1.7\times $}}\end{tabular} & \textcolor{gray}{\ensuremath{\cdot}} & \textcolor{gray}{\ensuremath{\cdot}} & \hspace*{1pt}\begin{tabular}[t]{@{}c@{}}\rule{0pt}{12pt}4m\,47.6s\\[-.2pt]{\fontsize{5.8}{6.4}\selectfont \textcolor{BrickRed}{$ \blacktriangledown\,29.0\times $}}\\[-.2pt]25.3s\\[-.2pt]{\fontsize{5.8}{6.4}\selectfont \textcolor{BrickRed}{$ \blacktriangledown\,2.9\times $}}\end{tabular} & \hspace*{1pt}\begin{tabular}[t]{@{}c@{}}\rule{0pt}{12pt}3m\,58.5s\\[-.2pt]{\fontsize{5.8}{6.4}\selectfont \textcolor{BrickRed}{$ \blacktriangledown\,25.2\times $}}\\[-.2pt]23.4s\\[-.2pt]{\fontsize{5.8}{6.4}\selectfont \textcolor{BrickRed}{$ \blacktriangledown\,3.1\times $}}\end{tabular} & \textcolor{gray}{\ensuremath{\cdot}} & \textcolor{gray}{\ensuremath{\cdot}} & \hspace*{1pt}\begin{tabular}[t]{@{}c@{}}\rule{0pt}{12pt}2m\,51.1s\\[-.2pt]{\fontsize{5.8}{6.4}\selectfont \textcolor{BrickRed}{$ \blacktriangledown\,3.2\times $}}\\[-.2pt]2m\,26.7s\\[-.2pt]{\fontsize{5.8}{6.4}\selectfont \textcolor{BrickRed}{$ \blacktriangledown\,2.8\times $}}\end{tabular} & \hspace*{1pt}\begin{tabular}[t]{@{}c@{}}\rule{0pt}{12pt}2m\,59.5s\\[-.2pt]{\fontsize{5.8}{6.4}\selectfont \textcolor{BrickRed}{$ \blacktriangledown\,3.2\times $}}\\[-.2pt]2m\,36.8s\\[-.2pt]{\fontsize{5.8}{6.4}\selectfont \textcolor{BrickRed}{$ \blacktriangledown\,2.9\times $}}\end{tabular} & \textcolor{gray}{\ensuremath{\cdot}} & \textcolor{gray}{\ensuremath{\cdot}} & \textcolor{gray}{\ensuremath{\cdot}} & \textcolor{gray}{\ensuremath{\cdot}} & \textcolor{gray}{\ensuremath{\cdot}} & \textcolor{gray}{\ensuremath{\cdot}}\\[3pt]\hline
\BenchmarkSourceLabel{bench:aws-2-k15-w10:row:mps-bounded-simd-ksw2}{R41}\shortstack[l]{{\color{gray}R41}} & \begin{tabular}[t]{@{}l@{}}\rule{0pt}{12pt}MPS / Bounded\\SIMD-DP / KSW2\end{tabular} & \hspace*{1pt}\begin{tabular}[t]{@{}c@{}}\rule{0pt}{12pt}8.2s\\[-.2pt]{\fontsize{5.8}{6.4}\selectfont \textcolor{ForestGreen}{$ \blacktriangle\,26\% $}}\\[-.2pt]5.4s\\[-.2pt]{\fontsize{5.8}{6.4}\selectfont \textcolor{ForestGreen}{$ \blacktriangle\,1.8\times $}}\end{tabular} & \hspace*{1pt}\begin{tabular}[t]{@{}c@{}}\rule{0pt}{12pt}8.4s\\[-.2pt]{\fontsize{5.8}{6.4}\selectfont \textcolor{ForestGreen}{$ \blacktriangle\,21\% $}}\\[-.2pt]4.9s\\[-.2pt]{\fontsize{5.8}{6.4}\selectfont \textcolor{ForestGreen}{$ \blacktriangle\,1.8\times $}}\end{tabular} & \hspace*{1pt}\begin{tabular}[t]{@{}c@{}}\rule{0pt}{12pt}17.8s\\[-.2pt]{\fontsize{5.8}{6.4}\selectfont \textcolor{ForestGreen}{$ \blacktriangle\,21\% $}}\\[-.2pt]8.9s\\[-.2pt]{\fontsize{5.8}{6.4}\selectfont \textcolor{ForestGreen}{$ \blacktriangle\,2.1\times $}}\end{tabular} & \hspace*{1pt}\begin{tabular}[t]{@{}c@{}}\rule{0pt}{12pt}9.6s\\[-.2pt]{\fontsize{5.8}{6.4}\selectfont \textcolor{ForestGreen}{$ \blacktriangle\,7.6\times $}}\\[-.2pt]6s\\[-.2pt]{\fontsize{5.8}{6.4}\selectfont \textcolor{ForestGreen}{$ \blacktriangle\,10.9\times $}}\end{tabular} & \hspace*{1pt}\begin{tabular}[t]{@{}c@{}}\rule{0pt}{12pt}4m\,56.1s\\[-.2pt]{\fontsize{5.8}{6.4}\selectfont \textcolor{BrickRed}{$ \blacktriangledown\,29.9\times $}}\\[-.2pt]32.4s\\[-.2pt]{\fontsize{5.8}{6.4}\selectfont \textcolor{BrickRed}{$ \blacktriangledown\,3.7\times $}}\end{tabular} & \hspace*{1pt}\begin{tabular}[t]{@{}c@{}}\rule{0pt}{12pt}4m\,12.1s\\[-.2pt]{\fontsize{5.8}{6.4}\selectfont \textcolor{BrickRed}{$ \blacktriangledown\,26.6\times $}}\\[-.2pt]30.4s\\[-.2pt]{\fontsize{5.8}{6.4}\selectfont \textcolor{BrickRed}{$ \blacktriangledown\,4.0\times $}}\end{tabular} & \hspace*{1pt}\begin{tabular}[t]{@{}c@{}}\rule{0pt}{12pt}8m\,4.9s\\[-.2pt]{\fontsize{5.8}{6.4}\selectfont \textcolor{BrickRed}{$ \blacktriangledown\,29.4\times $}}\\[-.2pt]52.6s\\[-.2pt]{\fontsize{5.8}{6.4}\selectfont \textcolor{BrickRed}{$ \blacktriangledown\,4.1\times $}}\end{tabular} & \textcolor{gray}{\ensuremath{\cdot}} & \hspace*{1pt}\begin{tabular}[t]{@{}c@{}}\rule{0pt}{12pt}2m\,20.7s\\[-.2pt]{\fontsize{5.8}{6.4}\selectfont \textcolor{BrickRed}{$ \blacktriangledown\,2.7\times $}}\\[-.2pt]1m\,56.8s\\[-.2pt]{\fontsize{5.8}{6.4}\selectfont \textcolor{BrickRed}{$ \blacktriangledown\,2.3\times $}}\end{tabular} & \hspace*{1pt}\begin{tabular}[t]{@{}c@{}}\rule{0pt}{12pt}2m\,19.7s\\[-.2pt]{\fontsize{5.8}{6.4}\selectfont \textcolor{BrickRed}{$ \blacktriangledown\,2.5\times $}}\\[-.2pt]1m\,57.9s\\[-.2pt]{\fontsize{5.8}{6.4}\selectfont \textcolor{BrickRed}{$ \blacktriangledown\,2.2\times $}}\end{tabular} & \textcolor{gray}{\ensuremath{\cdot}} & \hspace*{1pt}\begin{tabular}[t]{@{}c@{}}\rule{0pt}{12pt}2m\,42.4s\\[-.2pt]{\fontsize{5.8}{6.4}\selectfont \textcolor{ForestGreen}{$ \blacktriangle\,2.2\times $}}\\[-.2pt]2m\,21.1s\\[-.2pt]{\fontsize{5.8}{6.4}\selectfont \textcolor{ForestGreen}{$ \blacktriangle\,2.5\times $}}\end{tabular} & \hspace*{1pt}\begin{tabular}[t]{@{}c@{}}\rule{0pt}{12pt}6m\,37.2s\\[-.2pt]{\fontsize{5.8}{6.4}\selectfont \textcolor{BrickRed}{$ \blacktriangledown\,8.3\times $}}\\[-.2pt]2m\,13.5s\\[-.2pt]{\fontsize{5.8}{6.4}\selectfont \textcolor{BrickRed}{$ \blacktriangledown\,2.8\times $}}\end{tabular} & \hspace*{1pt}\begin{tabular}[t]{@{}c@{}}\rule{0pt}{12pt}5m\,36.2s\\[-.2pt]{\fontsize{5.8}{6.4}\selectfont \textcolor{BrickRed}{$ \blacktriangledown\,7.6\times $}}\\[-.2pt]1m\,59.7s\\[-.2pt]{\fontsize{5.8}{6.4}\selectfont \textcolor{BrickRed}{$ \blacktriangledown\,2.8\times $}}\end{tabular} & \textcolor{gray}{\ensuremath{\cdot}} & \hspace*{1pt}\begin{tabular}[t]{@{}c@{}}\rule{0pt}{12pt}11m\,54s\\[-.2pt]{\fontsize{5.8}{6.4}\selectfont \textcolor{BrickRed}{$ \blacktriangledown\,4.3\times $}}\\[-.2pt]2m\,27.8s\\[-.2pt]{\fontsize{5.8}{6.4}\selectfont \textcolor{ForestGreen}{$ \blacktriangle\,9\% $}}\end{tabular}\\[3pt]\hline
\BenchmarkSourceLabel{bench:aws-2-k15-w10:row:mps-scalar-rmq-ksw2}{R42}\shortstack[l]{{\color{gray}R42}} & \begin{tabular}[t]{@{}l@{}}\rule{0pt}{12pt}MPS / Scalar\\RMQ / KSW2\end{tabular} & \hspace*{1pt}\begin{tabular}[t]{@{}c@{}}\rule{0pt}{12pt}8.2s\\[-.2pt]{\fontsize{5.8}{6.4}\selectfont \textcolor{ForestGreen}{$ \blacktriangle\,26\% $}}\\[-.2pt]5.4s\\[-.2pt]{\fontsize{5.8}{6.4}\selectfont \textcolor{ForestGreen}{$ \blacktriangle\,1.9\times $}}\end{tabular} & \hspace*{1pt}\begin{tabular}[t]{@{}c@{}}\rule{0pt}{12pt}8.7s\\[-.2pt]{\fontsize{5.8}{6.4}\selectfont \textcolor{ForestGreen}{$ \blacktriangle\,18\% $}}\\[-.2pt]5.4s\\[-.2pt]{\fontsize{5.8}{6.4}\selectfont \textcolor{ForestGreen}{$ \blacktriangle\,1.6\times $}}\end{tabular} & \textcolor{gray}{\ensuremath{\cdot}} & \textcolor{gray}{\ensuremath{\cdot}} & \hspace*{1pt}\begin{tabular}[t]{@{}c@{}}\rule{0pt}{12pt}4m\,53.7s\\[-.2pt]{\fontsize{5.8}{6.4}\selectfont \textcolor{BrickRed}{$ \blacktriangledown\,29.7\times $}}\\[-.2pt]31.6s\\[-.2pt]{\fontsize{5.8}{6.4}\selectfont \textcolor{BrickRed}{$ \blacktriangledown\,3.6\times $}}\end{tabular} & \hspace*{1pt}\begin{tabular}[t]{@{}c@{}}\rule{0pt}{12pt}4m\,11.2s\\[-.2pt]{\fontsize{5.8}{6.4}\selectfont \textcolor{BrickRed}{$ \blacktriangledown\,26.5\times $}}\\[-.2pt]31.2s\\[-.2pt]{\fontsize{5.8}{6.4}\selectfont \textcolor{BrickRed}{$ \blacktriangledown\,4.1\times $}}\end{tabular} & \textcolor{gray}{\ensuremath{\cdot}} & \textcolor{gray}{\ensuremath{\cdot}} & \hspace*{1pt}\begin{tabular}[t]{@{}c@{}}\rule{0pt}{12pt}3m\,14.2s\\[-.2pt]{\fontsize{5.8}{6.4}\selectfont \textcolor{BrickRed}{$ \blacktriangledown\,3.7\times $}}\\[-.2pt]2m\,49.7s\\[-.2pt]{\fontsize{5.8}{6.4}\selectfont \textcolor{BrickRed}{$ \blacktriangledown\,3.3\times $}}\end{tabular} & \hspace*{1pt}\begin{tabular}[t]{@{}c@{}}\rule{0pt}{12pt}3m\,39.1s\\[-.2pt]{\fontsize{5.8}{6.4}\selectfont \textcolor{BrickRed}{$ \blacktriangledown\,3.9\times $}}\\[-.2pt]3m\,16.1s\\[-.2pt]{\fontsize{5.8}{6.4}\selectfont \textcolor{BrickRed}{$ \blacktriangledown\,3.6\times $}}\end{tabular} & \textcolor{gray}{\ensuremath{\cdot}} & \textcolor{gray}{\ensuremath{\cdot}} & \hspace*{1pt}\begin{tabular}[t]{@{}c@{}}\rule{0pt}{12pt}7m\,9s\\[-.2pt]{\fontsize{5.8}{6.4}\selectfont \textcolor{BrickRed}{$ \blacktriangledown\,8.9\times $}}\\[-.2pt]2m\,43.5s\\[-.2pt]{\fontsize{5.8}{6.4}\selectfont \textcolor{BrickRed}{$ \blacktriangledown\,3.5\times $}}\end{tabular} & \textcolor{gray}{\ensuremath{\cdot}} & \textcolor{gray}{\ensuremath{\cdot}} & \textcolor{gray}{\ensuremath{\cdot}}\\[3pt]\hline
\BenchmarkSourceLabel{bench:aws-2-k15-w10:row:mps-twostack-simd-ksw2}{R43}\shortstack[l]{{\color{gray}R43}} & \begin{tabular}[t]{@{}l@{}}\rule{0pt}{12pt}MPS / Two-stack\\SIMD-DP / KSW2\end{tabular} & \hspace*{1pt}\begin{tabular}[t]{@{}c@{}}\rule{0pt}{12pt}8.3s\\[-.2pt]{\fontsize{5.8}{6.4}\selectfont \textcolor{ForestGreen}{$ \blacktriangle\,26\% $}}\\[-.2pt]5.4s\\[-.2pt]{\fontsize{5.8}{6.4}\selectfont \textcolor{ForestGreen}{$ \blacktriangle\,1.8\times $}}\end{tabular} & \hspace*{1pt}\begin{tabular}[t]{@{}c@{}}\rule{0pt}{12pt}8.4s\\[-.2pt]{\fontsize{5.8}{6.4}\selectfont \textcolor{ForestGreen}{$ \blacktriangle\,21\% $}}\\[-.2pt]5.2s\\[-.2pt]{\fontsize{5.8}{6.4}\selectfont \textcolor{ForestGreen}{$ \blacktriangle\,1.7\times $}}\end{tabular} & \hspace*{1pt}\begin{tabular}[t]{@{}c@{}}\rule{0pt}{12pt}16.2s\\[-.2pt]{\fontsize{5.8}{6.4}\selectfont \textcolor{ForestGreen}{$ \blacktriangle\,28\% $}}\\[-.2pt]9.2s\\[-.2pt]{\fontsize{5.8}{6.4}\selectfont \textcolor{ForestGreen}{$ \blacktriangle\,2.0\times $}}\end{tabular} & \textcolor{gray}{\ensuremath{\cdot}} & \hspace*{1pt}\begin{tabular}[t]{@{}c@{}}\rule{0pt}{12pt}4m\,46.4s\\[-.2pt]{\fontsize{5.8}{6.4}\selectfont \textcolor{BrickRed}{$ \blacktriangledown\,28.9\times $}}\\[-.2pt]24.9s\\[-.2pt]{\fontsize{5.8}{6.4}\selectfont \textcolor{BrickRed}{$ \blacktriangledown\,2.8\times $}}\end{tabular} & \hspace*{1pt}\begin{tabular}[t]{@{}c@{}}\rule{0pt}{12pt}3m\,58.8s\\[-.2pt]{\fontsize{5.8}{6.4}\selectfont \textcolor{BrickRed}{$ \blacktriangledown\,25.2\times $}}\\[-.2pt]23.6s\\[-.2pt]{\fontsize{5.8}{6.4}\selectfont \textcolor{BrickRed}{$ \blacktriangledown\,3.1\times $}}\end{tabular} & \textcolor{gray}{\ensuremath{\cdot}} & \textcolor{gray}{\ensuremath{\cdot}} & \hspace*{1pt}\begin{tabular}[t]{@{}c@{}}\rule{0pt}{12pt}2m\,40.2s\\[-.2pt]{\fontsize{5.8}{6.4}\selectfont \textcolor{BrickRed}{$ \blacktriangledown\,3.0\times $}}\\[-.2pt]2m\,15.4s\\[-.2pt]{\fontsize{5.8}{6.4}\selectfont \textcolor{BrickRed}{$ \blacktriangledown\,2.6\times $}}\end{tabular} & \hspace*{1pt}\begin{tabular}[t]{@{}c@{}}\rule{0pt}{12pt}2m\,42.5s\\[-.2pt]{\fontsize{5.8}{6.4}\selectfont \textcolor{BrickRed}{$ \blacktriangledown\,2.9\times $}}\\[-.2pt]2m\,20.9s\\[-.2pt]{\fontsize{5.8}{6.4}\selectfont \textcolor{BrickRed}{$ \blacktriangledown\,2.6\times $}}\end{tabular} & \textcolor{gray}{\ensuremath{\cdot}} & \textcolor{gray}{\ensuremath{\cdot}} & \textcolor{gray}{\ensuremath{\cdot}} & \textcolor{gray}{\ensuremath{\cdot}} & \textcolor{gray}{\ensuremath{\cdot}} & \textcolor{gray}{\ensuremath{\cdot}}\\[3pt]\hline
\end{tabular}\par}
\smallskip{\footnotesize\noindent Each cell: A and B times, $\cdot$: no measurement. ---\textsuperscript{g}: GPU unavailable. n/a: no comparison.\par
\noindent Blue: one lowest A and one lowest B measurement for each column for each configuration.\par
\noindent Gray: minimap2. Yellow: my preferred MPF / Bounded / SIMD-DP / Block configuration.\par
\noindent MPF/MPS: minimizers per fragment/sequence. SIMD-ref: naive SIMD scan. Bounded: bounded SIMD algorithm.\par
}
\clearpage
\noindent{\large\bfseries Table \thetable: (D): AWS: 2 Gbp, k=15, w=10\hfill 7/7}\par\smallskip
{\fontsize{7.1}{8.1}\selectfont\setlength{\tabcolsep}{1pt}\setlength{\fboxsep}{2pt}\renewcommand{\arraystretch}{1.12}
\noindent\begin{tabular}{@{}p{14mm}p{27mm}|p{13.55mm}|p{13.55mm}|p{13.55mm}|p{13.55mm}|p{13.55mm}|p{13.55mm}|p{13.55mm}|p{13.55mm}|p{13.55mm}|p{13.55mm}|p{13.55mm}|p{13.55mm}|p{13.55mm}|p{13.55mm}|p{13.55mm}|p{13.55mm}|@{}}\hline
Engine & Execution & \multicolumn{8}{c|}{\textbf{Index}} & \multicolumn{8}{c|}{\textbf{No index}}\\
 &  & \multicolumn{8}{c|}{\textbf{2 Gbp}} & \multicolumn{8}{c|}{\textbf{2 Gbp}}\\
 &  & \multicolumn{8}{c|}{\textbf{k=15, w=10}} & \multicolumn{8}{c|}{\textbf{k=15, w=10}}\\
 &  & \multicolumn{4}{c|}{\textbf{DNA4}} & \multicolumn{4}{c|}{\textbf{DNA5}} & \multicolumn{4}{c|}{\textbf{DNA4}} & \multicolumn{4}{c|}{\textbf{DNA5}}\\
 &  & \multicolumn{1}{c|}{\textbf{M7A}} & \multicolumn{1}{c|}{\textbf{M7I}} & \multicolumn{1}{c|}{\textbf{M4}} & \multicolumn{1}{c|}{\textbf{C7G}} & \multicolumn{1}{c|}{\textbf{M7A}} & \multicolumn{1}{c|}{\textbf{M7I}} & \multicolumn{1}{c|}{\textbf{M4}} & \multicolumn{1}{c|}{\textbf{C7G}} & \multicolumn{1}{c|}{\textbf{M7A}} & \multicolumn{1}{c|}{\textbf{M7I}} & \multicolumn{1}{c|}{\textbf{M4}} & \multicolumn{1}{c|}{\textbf{C7G}} & \multicolumn{1}{c|}{\textbf{M7A}} & \multicolumn{1}{c|}{\textbf{M7I}} & \multicolumn{1}{c|}{\textbf{M4}} & \multicolumn{1}{c|}{\textbf{C7G}}\\
 & & {\color{gray}C1} & {\color{gray}C2} & {\color{gray}C3} & {\color{gray}C4} & {\color{gray}C5} & {\color{gray}C6} & {\color{gray}C7} & {\color{gray}C8} & {\color{gray}C9} & {\color{gray}C10} & {\color{gray}C11} & {\color{gray}C12} & {\color{gray}C13} & {\color{gray}C14} & {\color{gray}C15} & {\color{gray}C16}\\
\hline
\rowcolor{BenchBaseline}[0.5pt][0.5pt]\shortstack[l]{{\color{gray}R1}\\\textbf{minimap2}} & \textbf{Baseline} & \hspace*{1pt}\begin{tabular}[t]{@{}c@{}}\rule{0pt}{12pt}11.2s\\[-.2pt]{\fontsize{5.8}{6.4}\selectfont \textcolor{BrickRed}{$ \blacktriangledown\,3.1\times $}}\\[-.2pt]10s\\[-.2pt]{\fontsize{5.8}{6.4}\selectfont \textcolor{BrickRed}{$ \blacktriangledown\,3.6\times $}}\end{tabular} & \hspace*{1pt}\begin{tabular}[t]{@{}c@{}}\rule{0pt}{12pt}10.7s\\[-.2pt]{\fontsize{5.8}{6.4}\selectfont \textcolor{BrickRed}{$ \blacktriangledown\,3.2\times $}}\\[-.2pt]8.6s\\[-.2pt]{\fontsize{5.8}{6.4}\selectfont \textcolor{BrickRed}{$ \blacktriangledown\,3.5\times $}}\end{tabular} & \hspace*{1pt}\begin{tabular}[t]{@{}c@{}}\rule{0pt}{12pt}22.5s\\[-.2pt]{\fontsize{5.8}{6.4}\selectfont \textcolor{BrickRed}{$ \blacktriangledown\,3.6\times $}}\\[-.2pt]18.8s\\[-.2pt]{\fontsize{5.8}{6.4}\selectfont \textcolor{BrickRed}{$ \blacktriangledown\,4.4\times $}}\end{tabular} & \hspace*{1pt}\begin{tabular}[t]{@{}c@{}}\rule{0pt}{12pt}1m\,13.1s\\[-.2pt]{\fontsize{5.8}{6.4}\selectfont \textcolor{BrickRed}{$ \blacktriangledown\,22.5\times $}}\\[-.2pt]1m\,5.7s\\[-.2pt]{\fontsize{5.8}{6.4}\selectfont \textcolor{BrickRed}{$ \blacktriangledown\,27.7\times $}}\end{tabular} & \hspace*{1pt}\begin{tabular}[t]{@{}c@{}}\rule{0pt}{12pt}9.9s\\[-.2pt]{\fontsize{5.8}{6.4}\selectfont \textcolor{BrickRed}{$ \blacktriangledown\,21\% $}}\\[-.2pt]8.8s\\[-.2pt]{\fontsize{5.8}{6.4}\selectfont \textcolor{BrickRed}{$ \blacktriangledown\,1.6\times $}}\end{tabular} & \hspace*{1pt}\begin{tabular}[t]{@{}c@{}}\rule{0pt}{12pt}9.5s\\[-.2pt]{\fontsize{5.8}{6.4}\selectfont \textcolor{BrickRed}{$ \blacktriangledown\,1.4\times $}}\\[-.2pt]7.6s\\[-.2pt]{\fontsize{5.8}{6.4}\selectfont \textcolor{BrickRed}{$ \blacktriangledown\,1.6\times $}}\end{tabular} & \hspace*{1pt}\begin{tabular}[t]{@{}c@{}}\rule{0pt}{12pt}16.5s\\[-.2pt]{\fontsize{5.8}{6.4}\selectfont \textcolor{BrickRed}{$ \blacktriangledown\,1.7\times $}}\\[-.2pt]12.9s\\[-.2pt]{\fontsize{5.8}{6.4}\selectfont \textcolor{BrickRed}{$ \blacktriangledown\,2.0\times $}}\end{tabular} & \hspace*{1pt}\begin{tabular}[t]{@{}c@{}}\rule{0pt}{12pt}28.2s\\[-.2pt]{\fontsize{5.8}{6.4}\selectfont \textcolor{BrickRed}{$ \blacktriangledown\,2.3\times $}}\\[-.2pt]26.2s\\[-.2pt]{\fontsize{5.8}{6.4}\selectfont \textcolor{BrickRed}{$ \blacktriangledown\,3.0\times $}}\end{tabular} & \hspace*{1pt}\begin{tabular}[t]{@{}c@{}}\rule{0pt}{12pt}53s\\[-.2pt]{\fontsize{5.8}{6.4}\selectfont \textcolor{ForestGreen}{$ \blacktriangle\,2.2\times $}}\\[-.2pt]51.9s\\[-.2pt]{\fontsize{5.8}{6.4}\selectfont \textcolor{ForestGreen}{$ \blacktriangle\,2.1\times $}}\end{tabular} & \hspace*{1pt}\begin{tabular}[t]{@{}c@{}}\rule{0pt}{12pt}56.6s\\[-.2pt]{\fontsize{5.8}{6.4}\selectfont \textcolor{ForestGreen}{$ \blacktriangle\,2.1\times $}}\\[-.2pt]54.6s\\[-.2pt]{\fontsize{5.8}{6.4}\selectfont \textcolor{ForestGreen}{$ \blacktriangle\,1.9\times $}}\end{tabular} & \hspace*{1pt}\begin{tabular}[t]{@{}c@{}}\rule{0pt}{12pt}1m\,55.5s\\[-.2pt]{\fontsize{5.8}{6.4}\selectfont \textcolor{ForestGreen}{$ \blacktriangle\,1.8\times $}}\\[-.2pt]1m\,51.8s\\[-.2pt]{\fontsize{5.8}{6.4}\selectfont \textcolor{ForestGreen}{$ \blacktriangle\,1.5\times $}}\end{tabular} & \hspace*{1pt}\begin{tabular}[t]{@{}c@{}}\rule{0pt}{12pt}6m\,1.3s\\[-.2pt]{\fontsize{5.8}{6.4}\selectfont \textcolor{BrickRed}{$ \blacktriangledown\,3.0\times $}}\\[-.2pt]5m\,55.2s\\[-.2pt]{\fontsize{5.8}{6.4}\selectfont \textcolor{BrickRed}{$ \blacktriangledown\,3.1\times $}}\end{tabular} & \hspace*{1pt}\begin{tabular}[t]{@{}c@{}}\rule{0pt}{12pt}48.1s\\[-.2pt]{\fontsize{5.8}{6.4}\selectfont \textcolor{ForestGreen}{$ \blacktriangle\,7.3\times $}}\\[-.2pt]47s\\[-.2pt]{\fontsize{5.8}{6.4}\selectfont \textcolor{ForestGreen}{$ \blacktriangle\,2.4\times $}}\end{tabular} & \hspace*{1pt}\begin{tabular}[t]{@{}c@{}}\rule{0pt}{12pt}44.5s\\[-.2pt]{\fontsize{5.8}{6.4}\selectfont \textcolor{ForestGreen}{$ \blacktriangle\,6.5\times $}}\\[-.2pt]42.6s\\[-.2pt]{\fontsize{5.8}{6.4}\selectfont \textcolor{ForestGreen}{$ \blacktriangle\,2.7\times $}}\end{tabular} & \hspace*{1pt}\begin{tabular}[t]{@{}c@{}}\rule{0pt}{12pt}1m\,30.2s\\[-.2pt]{\fontsize{5.8}{6.4}\selectfont \textcolor{ForestGreen}{$ \blacktriangle\,6.4\times $}}\\[-.2pt]1m\,26.7s\\[-.2pt]{\fontsize{5.8}{6.4}\selectfont \textcolor{ForestGreen}{$ \blacktriangle\,2.2\times $}}\end{tabular} & \hspace*{1pt}\begin{tabular}[t]{@{}c@{}}\rule{0pt}{12pt}2m\,46.4s\\[-.2pt]{\fontsize{5.8}{6.4}\selectfont \textcolor{ForestGreen}{$ \blacktriangle\,4.0\times $}}\\[-.2pt]2m\,42.7s\\[-.2pt]{\fontsize{5.8}{6.4}\selectfont \textcolor{BrickRed}{$ \blacktriangledown\,27\% $}}\end{tabular}\\[3pt]\hline
\BenchmarkSourceLabel{bench:aws-2-k15-w10:row:mps-scalar-simd-ksw2}{R44}\shortstack[l]{{\color{gray}R44}\\\textbf{BioseqDB}} & \begin{tabular}[t]{@{}l@{}}\rule{0pt}{12pt}MPS / Scalar\\SIMD-DP / KSW2\end{tabular} & \hspace*{1pt}\begin{tabular}[t]{@{}c@{}}\rule{0pt}{12pt}8.3s\\[-.2pt]{\fontsize{5.8}{6.4}\selectfont \textcolor{ForestGreen}{$ \blacktriangle\,26\% $}}\\[-.2pt]5.5s\\[-.2pt]{\fontsize{5.8}{6.4}\selectfont \textcolor{ForestGreen}{$ \blacktriangle\,1.8\times $}}\end{tabular} & \hspace*{1pt}\begin{tabular}[t]{@{}c@{}}\rule{0pt}{12pt}8.3s\\[-.2pt]{\fontsize{5.8}{6.4}\selectfont \textcolor{ForestGreen}{$ \blacktriangle\,22\% $}}\\[-.2pt]5.2s\\[-.2pt]{\fontsize{5.8}{6.4}\selectfont \textcolor{ForestGreen}{$ \blacktriangle\,1.6\times $}}\end{tabular} & \textcolor{gray}{\ensuremath{\cdot}} & \textcolor{gray}{\ensuremath{\cdot}} & \hspace*{1pt}\begin{tabular}[t]{@{}c@{}}\rule{0pt}{12pt}4m\,52.9s\\[-.2pt]{\fontsize{5.8}{6.4}\selectfont \textcolor{BrickRed}{$ \blacktriangledown\,29.6\times $}}\\[-.2pt]31.1s\\[-.2pt]{\fontsize{5.8}{6.4}\selectfont \textcolor{BrickRed}{$ \blacktriangledown\,3.5\times $}}\end{tabular} & \hspace*{1pt}\begin{tabular}[t]{@{}c@{}}\rule{0pt}{12pt}4m\,10.6s\\[-.2pt]{\fontsize{5.8}{6.4}\selectfont \textcolor{BrickRed}{$ \blacktriangledown\,26.4\times $}}\\[-.2pt]30.7s\\[-.2pt]{\fontsize{5.8}{6.4}\selectfont \textcolor{BrickRed}{$ \blacktriangledown\,4.0\times $}}\end{tabular} & \textcolor{gray}{\ensuremath{\cdot}} & \textcolor{gray}{\ensuremath{\cdot}} & \hspace*{1pt}\begin{tabular}[t]{@{}c@{}}\rule{0pt}{12pt}3m\,1.5s\\[-.2pt]{\fontsize{5.8}{6.4}\selectfont \textcolor{BrickRed}{$ \blacktriangledown\,3.4\times $}}\\[-.2pt]2m\,37.6s\\[-.2pt]{\fontsize{5.8}{6.4}\selectfont \textcolor{BrickRed}{$ \blacktriangledown\,3.0\times $}}\end{tabular} & \hspace*{1pt}\begin{tabular}[t]{@{}c@{}}\rule{0pt}{12pt}3m\,21.1s\\[-.2pt]{\fontsize{5.8}{6.4}\selectfont \textcolor{BrickRed}{$ \blacktriangledown\,3.6\times $}}\\[-.2pt]2m\,58.7s\\[-.2pt]{\fontsize{5.8}{6.4}\selectfont \textcolor{BrickRed}{$ \blacktriangledown\,3.3\times $}}\end{tabular} & \textcolor{gray}{\ensuremath{\cdot}} & \textcolor{gray}{\ensuremath{\cdot}} & \hspace*{1pt}\begin{tabular}[t]{@{}c@{}}\rule{0pt}{12pt}7m\,6.4s\\[-.2pt]{\fontsize{5.8}{6.4}\selectfont \textcolor{BrickRed}{$ \blacktriangledown\,8.9\times $}}\\[-.2pt]2m\,43.1s\\[-.2pt]{\fontsize{5.8}{6.4}\selectfont \textcolor{BrickRed}{$ \blacktriangledown\,3.5\times $}}\end{tabular} & \textcolor{gray}{\ensuremath{\cdot}} & \textcolor{gray}{\ensuremath{\cdot}} & \textcolor{gray}{\ensuremath{\cdot}}\\[3pt]\hline
\BenchmarkSourceLabel{bench:aws-2-k15-w10:row:mps-scalar-scalar-ksw2}{R45}\shortstack[l]{{\color{gray}R45}} & \begin{tabular}[t]{@{}l@{}}\rule{0pt}{12pt}MPS / Scalar\\Scalar-DP / KSW2\end{tabular} & \hspace*{1pt}\begin{tabular}[t]{@{}c@{}}\rule{0pt}{12pt}8.3s\\[-.2pt]{\fontsize{5.8}{6.4}\selectfont \textcolor{ForestGreen}{$ \blacktriangle\,26\% $}}\\[-.2pt]5.4s\\[-.2pt]{\fontsize{5.8}{6.4}\selectfont \textcolor{ForestGreen}{$ \blacktriangle\,1.8\times $}}\end{tabular} & \hspace*{1pt}\begin{tabular}[t]{@{}c@{}}\rule{0pt}{12pt}8.7s\\[-.2pt]{\fontsize{5.8}{6.4}\selectfont \textcolor{ForestGreen}{$ \blacktriangle\,18\% $}}\\[-.2pt]5.3s\\[-.2pt]{\fontsize{5.8}{6.4}\selectfont \textcolor{ForestGreen}{$ \blacktriangle\,1.6\times $}}\end{tabular} & \textcolor{gray}{\ensuremath{\cdot}} & \textcolor{gray}{\ensuremath{\cdot}} & \hspace*{1pt}\begin{tabular}[t]{@{}c@{}}\rule{0pt}{12pt}4m\,54.2s\\[-.2pt]{\fontsize{5.8}{6.4}\selectfont \textcolor{BrickRed}{$ \blacktriangledown\,29.7\times $}}\\[-.2pt]31.5s\\[-.2pt]{\fontsize{5.8}{6.4}\selectfont \textcolor{BrickRed}{$ \blacktriangledown\,3.6\times $}}\end{tabular} & \hspace*{1pt}\begin{tabular}[t]{@{}c@{}}\rule{0pt}{12pt}4m\,11s\\[-.2pt]{\fontsize{5.8}{6.4}\selectfont \textcolor{BrickRed}{$ \blacktriangledown\,26.5\times $}}\\[-.2pt]31s\\[-.2pt]{\fontsize{5.8}{6.4}\selectfont \textcolor{BrickRed}{$ \blacktriangledown\,4.1\times $}}\end{tabular} & \textcolor{gray}{\ensuremath{\cdot}} & \textcolor{gray}{\ensuremath{\cdot}} & \hspace*{1pt}\begin{tabular}[t]{@{}c@{}}\rule{0pt}{12pt}3m\,21.3s\\[-.2pt]{\fontsize{5.8}{6.4}\selectfont \textcolor{BrickRed}{$ \blacktriangledown\,3.8\times $}}\\[-.2pt]2m\,53.7s\\[-.2pt]{\fontsize{5.8}{6.4}\selectfont \textcolor{BrickRed}{$ \blacktriangledown\,3.3\times $}}\end{tabular} & \textcolor{gray}{\ensuremath{\cdot}} & \textcolor{gray}{\ensuremath{\cdot}} & \textcolor{gray}{\ensuremath{\cdot}} & \hspace*{1pt}\begin{tabular}[t]{@{}c@{}}\rule{0pt}{12pt}7m\,19.5s\\[-.2pt]{\fontsize{5.8}{6.4}\selectfont \textcolor{BrickRed}{$ \blacktriangledown\,9.1\times $}}\\[-.2pt]2m\,56s\\[-.2pt]{\fontsize{5.8}{6.4}\selectfont \textcolor{BrickRed}{$ \blacktriangledown\,3.7\times $}}\end{tabular} & \textcolor{gray}{\ensuremath{\cdot}} & \textcolor{gray}{\ensuremath{\cdot}} & \textcolor{gray}{\ensuremath{\cdot}}\\[3pt]\hline
\BenchmarkSourceLabel{bench:aws-2-k15-w10:row:mps-twostack-scalar-ksw2}{R46}\shortstack[l]{{\color{gray}R46}} & \begin{tabular}[t]{@{}l@{}}\rule{0pt}{12pt}MPS / Two-stack\\Scalar-DP / KSW2\end{tabular} & \hspace*{1pt}\begin{tabular}[t]{@{}c@{}}\rule{0pt}{12pt}8.4s\\[-.2pt]{\fontsize{5.8}{6.4}\selectfont \textcolor{ForestGreen}{$ \blacktriangle\,25\% $}}\\[-.2pt]5.4s\\[-.2pt]{\fontsize{5.8}{6.4}\selectfont \textcolor{ForestGreen}{$ \blacktriangle\,1.8\times $}}\end{tabular} & \hspace*{1pt}\begin{tabular}[t]{@{}c@{}}\rule{0pt}{12pt}9.4s\\[-.2pt]{\fontsize{5.8}{6.4}\selectfont \textcolor{ForestGreen}{$ \blacktriangle\,12\% $}}\\[-.2pt]5.5s\\[-.2pt]{\fontsize{5.8}{6.4}\selectfont \textcolor{ForestGreen}{$ \blacktriangle\,1.6\times $}}\end{tabular} & \textcolor{gray}{\ensuremath{\cdot}} & \textcolor{gray}{\ensuremath{\cdot}} & \hspace*{1pt}\begin{tabular}[t]{@{}c@{}}\rule{0pt}{12pt}4m\,46.5s\\[-.2pt]{\fontsize{5.8}{6.4}\selectfont \textcolor{BrickRed}{$ \blacktriangledown\,28.9\times $}}\\[-.2pt]24.8s\\[-.2pt]{\fontsize{5.8}{6.4}\selectfont \textcolor{BrickRed}{$ \blacktriangledown\,2.8\times $}}\end{tabular} & \hspace*{1pt}\begin{tabular}[t]{@{}c@{}}\rule{0pt}{12pt}3m\,58s\\[-.2pt]{\fontsize{5.8}{6.4}\selectfont \textcolor{BrickRed}{$ \blacktriangledown\,25.1\times $}}\\[-.2pt]23s\\[-.2pt]{\fontsize{5.8}{6.4}\selectfont \textcolor{BrickRed}{$ \blacktriangledown\,3.0\times $}}\end{tabular} & \textcolor{gray}{\ensuremath{\cdot}} & \textcolor{gray}{\ensuremath{\cdot}} & \textcolor{gray}{\ensuremath{\cdot}} & \textcolor{gray}{\ensuremath{\cdot}} & \textcolor{gray}{\ensuremath{\cdot}} & \textcolor{gray}{\ensuremath{\cdot}} & \textcolor{gray}{\ensuremath{\cdot}} & \textcolor{gray}{\ensuremath{\cdot}} & \textcolor{gray}{\ensuremath{\cdot}} & \textcolor{gray}{\ensuremath{\cdot}}\\[3pt]\hline
\BenchmarkSourceLabel{bench:aws-2-k15-w10:row:mps-simdref-simd-ksw2}{R47}\shortstack[l]{{\color{gray}R47}} & \begin{tabular}[t]{@{}l@{}}\rule{0pt}{12pt}MPS / SIMD-ref\\SIMD-DP / KSW2\end{tabular} & \hspace*{1pt}\begin{tabular}[t]{@{}c@{}}\rule{0pt}{12pt}9.8s\\[-.2pt]{\fontsize{5.8}{6.4}\selectfont \textcolor{ForestGreen}{$ \blacktriangle\,13\% $}}\\[-.2pt]5.3s\\[-.2pt]{\fontsize{5.8}{6.4}\selectfont \textcolor{ForestGreen}{$ \blacktriangle\,1.9\times $}}\end{tabular} & \hspace*{1pt}\begin{tabular}[t]{@{}c@{}}\rule{0pt}{12pt}8.6s\\[-.2pt]{\fontsize{5.8}{6.4}\selectfont \textcolor{ForestGreen}{$ \blacktriangle\,20\% $}}\\[-.2pt]5s\\[-.2pt]{\fontsize{5.8}{6.4}\selectfont \textcolor{ForestGreen}{$ \blacktriangle\,1.7\times $}}\end{tabular} & \hspace*{1pt}\begin{tabular}[t]{@{}c@{}}\rule{0pt}{12pt}15.7s\\[-.2pt]{\fontsize{5.8}{6.4}\selectfont \textcolor{ForestGreen}{$ \blacktriangle\,1.4\times $}}\\[-.2pt]9.3s\\[-.2pt]{\fontsize{5.8}{6.4}\selectfont \textcolor{ForestGreen}{$ \blacktriangle\,2.0\times $}}\end{tabular} & \textcolor{gray}{\ensuremath{\cdot}} & \hspace*{1pt}\begin{tabular}[t]{@{}c@{}}\rule{0pt}{12pt}4m\,55.9s\\[-.2pt]{\fontsize{5.8}{6.4}\selectfont \textcolor{BrickRed}{$ \blacktriangledown\,29.9\times $}}\\[-.2pt]33.3s\\[-.2pt]{\fontsize{5.8}{6.4}\selectfont \textcolor{BrickRed}{$ \blacktriangledown\,3.8\times $}}\end{tabular} & \hspace*{1pt}\begin{tabular}[t]{@{}c@{}}\rule{0pt}{12pt}4m\,11.3s\\[-.2pt]{\fontsize{5.8}{6.4}\selectfont \textcolor{BrickRed}{$ \blacktriangledown\,26.5\times $}}\\[-.2pt]31.3s\\[-.2pt]{\fontsize{5.8}{6.4}\selectfont \textcolor{BrickRed}{$ \blacktriangledown\,4.1\times $}}\end{tabular} & \textcolor{gray}{\ensuremath{\cdot}} & \textcolor{gray}{\ensuremath{\cdot}} & \textcolor{gray}{\ensuremath{\cdot}} & \textcolor{gray}{\ensuremath{\cdot}} & \textcolor{gray}{\ensuremath{\cdot}} & \textcolor{gray}{\ensuremath{\cdot}} & \textcolor{gray}{\ensuremath{\cdot}} & \textcolor{gray}{\ensuremath{\cdot}} & \textcolor{gray}{\ensuremath{\cdot}} & \textcolor{gray}{\ensuremath{\cdot}}\\[3pt]\hline
\BenchmarkSourceLabel{bench:aws-2-k15-w10:row:mps-simdref-rmq-ksw2}{R48}\shortstack[l]{{\color{gray}R48}} & \begin{tabular}[t]{@{}l@{}}\rule{0pt}{12pt}MPS / SIMD-ref\\RMQ / KSW2\end{tabular} & \hspace*{1pt}\begin{tabular}[t]{@{}c@{}}\rule{0pt}{12pt}9.8s\\[-.2pt]{\fontsize{5.8}{6.4}\selectfont \textcolor{ForestGreen}{$ \blacktriangle\,12\% $}}\\[-.2pt]5.2s\\[-.2pt]{\fontsize{5.8}{6.4}\selectfont \textcolor{ForestGreen}{$ \blacktriangle\,1.9\times $}}\end{tabular} & \hspace*{1pt}\begin{tabular}[t]{@{}c@{}}\rule{0pt}{12pt}8.6s\\[-.2pt]{\fontsize{5.8}{6.4}\selectfont \textcolor{ForestGreen}{$ \blacktriangle\,19\% $}}\\[-.2pt]4.9s\\[-.2pt]{\fontsize{5.8}{6.4}\selectfont \textcolor{ForestGreen}{$ \blacktriangle\,1.7\times $}}\end{tabular} & \textcolor{gray}{\ensuremath{\cdot}} & \textcolor{gray}{\ensuremath{\cdot}} & \hspace*{1pt}\begin{tabular}[t]{@{}c@{}}\rule{0pt}{12pt}4m\,58.8s\\[-.2pt]{\fontsize{5.8}{6.4}\selectfont \textcolor{BrickRed}{$ \blacktriangledown\,30.2\times $}}\\[-.2pt]36.1s\\[-.2pt]{\fontsize{5.8}{6.4}\selectfont \textcolor{BrickRed}{$ \blacktriangledown\,4.1\times $}}\end{tabular} & \hspace*{1pt}\begin{tabular}[t]{@{}c@{}}\rule{0pt}{12pt}4m\,12.9s\\[-.2pt]{\fontsize{5.8}{6.4}\selectfont \textcolor{BrickRed}{$ \blacktriangledown\,26.7\times $}}\\[-.2pt]33.1s\\[-.2pt]{\fontsize{5.8}{6.4}\selectfont \textcolor{BrickRed}{$ \blacktriangledown\,4.3\times $}}\end{tabular} & \textcolor{gray}{\ensuremath{\cdot}} & \textcolor{gray}{\ensuremath{\cdot}} & \textcolor{gray}{\ensuremath{\cdot}} & \textcolor{gray}{\ensuremath{\cdot}} & \textcolor{gray}{\ensuremath{\cdot}} & \textcolor{gray}{\ensuremath{\cdot}} & \textcolor{gray}{\ensuremath{\cdot}} & \textcolor{gray}{\ensuremath{\cdot}} & \textcolor{gray}{\ensuremath{\cdot}} & \textcolor{gray}{\ensuremath{\cdot}}\\[3pt]\hline
\BenchmarkSourceLabel{bench:aws-2-k15-w10:row:mps-simdref-scalar-ksw2}{R49}\shortstack[l]{{\color{gray}R49}} & \begin{tabular}[t]{@{}l@{}}\rule{0pt}{12pt}MPS / SIMD-ref\\Scalar-DP / KSW2\end{tabular} & \hspace*{1pt}\begin{tabular}[t]{@{}c@{}}\rule{0pt}{12pt}9.8s\\[-.2pt]{\fontsize{5.8}{6.4}\selectfont \textcolor{ForestGreen}{$ \blacktriangle\,12\% $}}\\[-.2pt]5.2s\\[-.2pt]{\fontsize{5.8}{6.4}\selectfont \textcolor{ForestGreen}{$ \blacktriangle\,1.9\times $}}\end{tabular} & \hspace*{1pt}\begin{tabular}[t]{@{}c@{}}\rule{0pt}{12pt}8.7s\\[-.2pt]{\fontsize{5.8}{6.4}\selectfont \textcolor{ForestGreen}{$ \blacktriangle\,19\% $}}\\[-.2pt]5s\\[-.2pt]{\fontsize{5.8}{6.4}\selectfont \textcolor{ForestGreen}{$ \blacktriangle\,1.7\times $}}\end{tabular} & \hspace*{1pt}\begin{tabular}[t]{@{}c@{}}\rule{0pt}{12pt}16.4s\\[-.2pt]{\fontsize{5.8}{6.4}\selectfont \textcolor{ForestGreen}{$ \blacktriangle\,27\% $}}\\[-.2pt]8.9s\\[-.2pt]{\fontsize{5.8}{6.4}\selectfont \textcolor{ForestGreen}{$ \blacktriangle\,2.1\times $}}\end{tabular} & \textcolor{gray}{\ensuremath{\cdot}} & \hspace*{1pt}\begin{tabular}[t]{@{}c@{}}\rule{0pt}{12pt}4m\,55s\\[-.2pt]{\fontsize{5.8}{6.4}\selectfont \textcolor{BrickRed}{$ \blacktriangledown\,29.8\times $}}\\[-.2pt]32s\\[-.2pt]{\fontsize{5.8}{6.4}\selectfont \textcolor{BrickRed}{$ \blacktriangledown\,3.6\times $}}\end{tabular} & \hspace*{1pt}\begin{tabular}[t]{@{}c@{}}\rule{0pt}{12pt}4m\,9.4s\\[-.2pt]{\fontsize{5.8}{6.4}\selectfont \textcolor{BrickRed}{$ \blacktriangledown\,26.3\times $}}\\[-.2pt]29.7s\\[-.2pt]{\fontsize{5.8}{6.4}\selectfont \textcolor{BrickRed}{$ \blacktriangledown\,3.9\times $}}\end{tabular} & \textcolor{gray}{\ensuremath{\cdot}} & \textcolor{gray}{\ensuremath{\cdot}} & \textcolor{gray}{\ensuremath{\cdot}} & \textcolor{gray}{\ensuremath{\cdot}} & \textcolor{gray}{\ensuremath{\cdot}} & \textcolor{gray}{\ensuremath{\cdot}} & \textcolor{gray}{\ensuremath{\cdot}} & \textcolor{gray}{\ensuremath{\cdot}} & \textcolor{gray}{\ensuremath{\cdot}} & \textcolor{gray}{\ensuremath{\cdot}}\\[3pt]\hline
\BenchmarkSourceLabel{bench:aws-2-k15-w10:row:vulkan-unavailable}{R50}\shortstack[l]{{\color{gray}R50}\\\rule{0pt}{12pt}\textbf{BioseqDB}} & Vulkan (all modes) & \textcolor{gray}{---\textsuperscript{g}} & \textcolor{gray}{---\textsuperscript{g}} & \textcolor{gray}{---\textsuperscript{g}} & \textcolor{gray}{---\textsuperscript{g}} & \textcolor{gray}{---\textsuperscript{g}} & \textcolor{gray}{---\textsuperscript{g}} & \textcolor{gray}{---\textsuperscript{g}} & \textcolor{gray}{---\textsuperscript{g}} & \textcolor{gray}{---\textsuperscript{g}} & \textcolor{gray}{---\textsuperscript{g}} & \textcolor{gray}{---\textsuperscript{g}} & \textcolor{gray}{---\textsuperscript{g}} & \textcolor{gray}{---\textsuperscript{g}} & \textcolor{gray}{---\textsuperscript{g}} & \textcolor{gray}{---\textsuperscript{g}} & \textcolor{gray}{---\textsuperscript{g}}\\[3pt]\hline
\end{tabular}\par}
\smallskip{\footnotesize\noindent Each cell: A and B times, $\cdot$: no measurement. ---\textsuperscript{g}: GPU unavailable. n/a: no comparison.\par
\noindent Blue: one lowest A and one lowest B measurement for each column for each configuration.\par
\noindent Gray: minimap2. Yellow: my preferred MPF / Bounded / SIMD-DP / Block configuration.\par
\noindent MPF/MPS: minimizers per fragment/sequence. SIMD-ref: naive SIMD scan. Bounded: bounded SIMD algorithm.\par
}
\clearpage
\refstepcounter{table}\label{tab:bench:aws-2-k20-w15}
\noindent{\large\bfseries Table \thetable: (E): AWS: 2 Gbp, k=20, w=15\hfill 1/7}\par\smallskip
{\fontsize{7.1}{8.1}\selectfont\setlength{\tabcolsep}{1pt}\setlength{\fboxsep}{2pt}\renewcommand{\arraystretch}{1.12}
\noindent\begin{tabular}{@{}p{14mm}p{27mm}|p{13.55mm}|p{13.55mm}|p{13.55mm}|p{13.55mm}|p{13.55mm}|p{13.55mm}|p{13.55mm}|p{13.55mm}|p{13.55mm}|p{13.55mm}|p{13.55mm}|p{13.55mm}|p{13.55mm}|p{13.55mm}|p{13.55mm}|p{13.55mm}|@{}}\hline
Engine & Execution & \multicolumn{8}{c|}{\textbf{Index}} & \multicolumn{8}{c|}{\textbf{No index}}\\
 &  & \multicolumn{8}{c|}{\textbf{2 Gbp}} & \multicolumn{8}{c|}{\textbf{2 Gbp}}\\
 &  & \multicolumn{8}{c|}{\textbf{k=20, w=15}} & \multicolumn{8}{c|}{\textbf{k=20, w=15}}\\
 &  & \multicolumn{4}{c|}{\textbf{DNA4}} & \multicolumn{4}{c|}{\textbf{DNA5}} & \multicolumn{4}{c|}{\textbf{DNA4}} & \multicolumn{4}{c|}{\textbf{DNA5}}\\
 &  & \multicolumn{1}{c|}{\textbf{M7A}} & \multicolumn{1}{c|}{\textbf{M7I}} & \multicolumn{1}{c|}{\textbf{M4}} & \multicolumn{1}{c|}{\textbf{C7G}} & \multicolumn{1}{c|}{\textbf{M7A}} & \multicolumn{1}{c|}{\textbf{M7I}} & \multicolumn{1}{c|}{\textbf{M4}} & \multicolumn{1}{c|}{\textbf{C7G}} & \multicolumn{1}{c|}{\textbf{M7A}} & \multicolumn{1}{c|}{\textbf{M7I}} & \multicolumn{1}{c|}{\textbf{M4}} & \multicolumn{1}{c|}{\textbf{C7G}} & \multicolumn{1}{c|}{\textbf{M7A}} & \multicolumn{1}{c|}{\textbf{M7I}} & \multicolumn{1}{c|}{\textbf{M4}} & \multicolumn{1}{c|}{\textbf{C7G}}\\
 & & \BenchmarkSourceLabel{bench:aws-2-k20-w15:column:index-2-k20-w15-dna4-m7a}{C1}{\color{gray}C1} & \BenchmarkSourceLabel{bench:aws-2-k20-w15:column:index-2-k20-w15-dna4-m7i}{C2}{\color{gray}C2} & \BenchmarkSourceLabel{bench:aws-2-k20-w15:column:index-2-k20-w15-dna4-m4}{C3}{\color{gray}C3} & \BenchmarkSourceLabel{bench:aws-2-k20-w15:column:index-2-k20-w15-dna4-c7g}{C4}{\color{gray}C4} & \BenchmarkSourceLabel{bench:aws-2-k20-w15:column:index-2-k20-w15-dna5-m7a}{C5}{\color{gray}C5} & \BenchmarkSourceLabel{bench:aws-2-k20-w15:column:index-2-k20-w15-dna5-m7i}{C6}{\color{gray}C6} & \BenchmarkSourceLabel{bench:aws-2-k20-w15:column:index-2-k20-w15-dna5-m4}{C7}{\color{gray}C7} & \BenchmarkSourceLabel{bench:aws-2-k20-w15:column:index-2-k20-w15-dna5-c7g}{C8}{\color{gray}C8} & \BenchmarkSourceLabel{bench:aws-2-k20-w15:column:noindex-2-k20-w15-dna4-m7a}{C9}{\color{gray}C9} & \BenchmarkSourceLabel{bench:aws-2-k20-w15:column:noindex-2-k20-w15-dna4-m7i}{C10}{\color{gray}C10} & \BenchmarkSourceLabel{bench:aws-2-k20-w15:column:noindex-2-k20-w15-dna4-m4}{C11}{\color{gray}C11} & \BenchmarkSourceLabel{bench:aws-2-k20-w15:column:noindex-2-k20-w15-dna4-c7g}{C12}{\color{gray}C12} & \BenchmarkSourceLabel{bench:aws-2-k20-w15:column:noindex-2-k20-w15-dna5-m7a}{C13}{\color{gray}C13} & \BenchmarkSourceLabel{bench:aws-2-k20-w15:column:noindex-2-k20-w15-dna5-m7i}{C14}{\color{gray}C14} & \BenchmarkSourceLabel{bench:aws-2-k20-w15:column:noindex-2-k20-w15-dna5-m4}{C15}{\color{gray}C15} & \BenchmarkSourceLabel{bench:aws-2-k20-w15:column:noindex-2-k20-w15-dna5-c7g}{C16}{\color{gray}C16}\\
\hline
\rowcolor{BenchBaseline}[0.5pt][0.5pt]\BenchmarkSourceLabel{bench:aws-2-k20-w15:row:minimap2}{R1}\shortstack[l]{{\color{gray}R1}\\\textbf{minimap2}} & \textbf{Baseline} & \hspace*{1pt}\begin{tabular}[t]{@{}c@{}}\rule{0pt}{12pt}19.6s\\[-.2pt]{\fontsize{5.8}{6.4}\selectfont \textcolor{BrickRed}{$ \blacktriangledown\,6.7\times $}}\\[-.2pt]18.6s\\[-.2pt]{\fontsize{5.8}{6.4}\selectfont \textcolor{BrickRed}{$ \blacktriangledown\,8.4\times $}}\end{tabular} & \hspace*{1pt}\begin{tabular}[t]{@{}c@{}}\rule{0pt}{12pt}18.2s\\[-.2pt]{\fontsize{5.8}{6.4}\selectfont \textcolor{BrickRed}{$ \blacktriangledown\,6.9\times $}}\\[-.2pt]16.5s\\[-.2pt]{\fontsize{5.8}{6.4}\selectfont \textcolor{BrickRed}{$ \blacktriangledown\,8.6\times $}}\end{tabular} & \hspace*{1pt}\begin{tabular}[t]{@{}c@{}}\rule{0pt}{12pt}42.5s\\[-.2pt]{\fontsize{5.8}{6.4}\selectfont \textcolor{BrickRed}{$ \blacktriangledown\,8.2\times $}}\\[-.2pt]39.4s\\[-.2pt]{\fontsize{5.8}{6.4}\selectfont \textcolor{BrickRed}{$ \blacktriangledown\,12.6\times $}}\end{tabular} & \hspace*{1pt}\begin{tabular}[t]{@{}c@{}}\rule{0pt}{12pt}6m\,29.2s\\[-.2pt]{\fontsize{5.8}{6.4}\selectfont \textcolor{BrickRed}{$ \blacktriangledown\,92.8\times $}}\\[-.2pt]5m\,56.9s\\[-.2pt]{\fontsize{5.8}{6.4}\selectfont \textcolor{BrickRed}{$ \blacktriangledown\,184.9\times $}}\end{tabular} & \hspace*{1pt}\begin{tabular}[t]{@{}c@{}}\rule{0pt}{12pt}9.3s\\[-.2pt]{\fontsize{5.8}{6.4}\selectfont \textcolor{BrickRed}{$ \blacktriangledown\,24\% $}}\\[-.2pt]8.4s\\[-.2pt]{\fontsize{5.8}{6.4}\selectfont \textcolor{BrickRed}{$ \blacktriangledown\,1.6\times $}}\end{tabular} & \hspace*{1pt}\begin{tabular}[t]{@{}c@{}}\rule{0pt}{12pt}8.7s\\[-.2pt]{\fontsize{5.8}{6.4}\selectfont \textcolor{BrickRed}{$ \blacktriangledown\,1.5\times $}}\\[-.2pt]7s\\[-.2pt]{\fontsize{5.8}{6.4}\selectfont \textcolor{BrickRed}{$ \blacktriangledown\,1.7\times $}}\end{tabular} & \hspace*{1pt}\begin{tabular}[t]{@{}c@{}}\rule{0pt}{12pt}15.8s\\[-.2pt]{\fontsize{5.8}{6.4}\selectfont \textcolor{BrickRed}{$ \blacktriangledown\,26\% $}}\\[-.2pt]12.7s\\[-.2pt]{\fontsize{5.8}{6.4}\selectfont \textcolor{BrickRed}{$ \blacktriangledown\,1.7\times $}}\end{tabular} & \hspace*{1pt}\begin{tabular}[t]{@{}c@{}}\rule{0pt}{12pt}18.5s\\[-.2pt]{\fontsize{5.8}{6.4}\selectfont \textcolor{BrickRed}{$ \blacktriangledown\,1.6\times $}}\\[-.2pt]17s\\[-.2pt]{\fontsize{5.8}{6.4}\selectfont \textcolor{BrickRed}{$ \blacktriangledown\,1.8\times $}}\end{tabular} & \hspace*{1pt}\begin{tabular}[t]{@{}c@{}}\rule{0pt}{12pt}\colorbox{BenchBest}{\strut 42.4s}\\[-.2pt]{\fontsize{5.8}{6.4}\selectfont \textcolor{ForestGreen}{$ \blacktriangle\,2.0\times $}}\\[-.2pt]\colorbox{BenchBest}{\strut 41.4s}\\[-.2pt]{\fontsize{5.8}{6.4}\selectfont \textcolor{ForestGreen}{$ \blacktriangle\,1.9\times $}}\end{tabular} & \hspace*{1pt}\begin{tabular}[t]{@{}c@{}}\rule{0pt}{12pt}\colorbox{BenchBest}{\strut 42.9s}\\[-.2pt]{\fontsize{5.8}{6.4}\selectfont \textcolor{ForestGreen}{$ \blacktriangle\,2.0\times $}}\\[-.2pt]\colorbox{BenchBest}{\strut 41.1s}\\[-.2pt]{\fontsize{5.8}{6.4}\selectfont \textcolor{ForestGreen}{$ \blacktriangle\,1.8\times $}}\end{tabular} & \hspace*{1pt}\begin{tabular}[t]{@{}c@{}}\rule{0pt}{12pt}\colorbox{BenchBest}{\strut 1m\,26.9s}\\[-.2pt]{\fontsize{5.8}{6.4}\selectfont \textcolor{ForestGreen}{$ \blacktriangle\,1.9\times $}}\\[-.2pt]\colorbox{BenchBest}{\strut 1m\,23.7s}\\[-.2pt]{\fontsize{5.8}{6.4}\selectfont \textcolor{ForestGreen}{$ \blacktriangle\,1.7\times $}}\end{tabular} & \hspace*{1pt}\begin{tabular}[t]{@{}c@{}}\rule{0pt}{12pt}4m\,51.6s\\[-.2pt]{\fontsize{5.8}{6.4}\selectfont \textcolor{BrickRed}{$ \blacktriangledown\,2.9\times $}}\\[-.2pt]4m\,46.6s\\[-.2pt]{\fontsize{5.8}{6.4}\selectfont \textcolor{BrickRed}{$ \blacktriangledown\,3.0\times $}}\end{tabular} & \hspace*{1pt}\begin{tabular}[t]{@{}c@{}}\rule{0pt}{12pt}\colorbox{BenchBest}{\strut 35s}\\[-.2pt]{\fontsize{5.8}{6.4}\selectfont \textcolor{ForestGreen}{$ \blacktriangle\,8.6\times $}}\\[-.2pt]\colorbox{BenchBest}{\strut 34s}\\[-.2pt]{\fontsize{5.8}{6.4}\selectfont \textcolor{ForestGreen}{$ \blacktriangle\,2.3\times $}}\end{tabular} & \hspace*{1pt}\begin{tabular}[t]{@{}c@{}}\rule{0pt}{12pt}\colorbox{BenchBest}{\strut 35.9s}\\[-.2pt]{\fontsize{5.8}{6.4}\selectfont \textcolor{ForestGreen}{$ \blacktriangle\,6.7\times $}}\\[-.2pt]\colorbox{BenchBest}{\strut 34.2s}\\[-.2pt]{\fontsize{5.8}{6.4}\selectfont \textcolor{ForestGreen}{$ \blacktriangle\,2.1\times $}}\end{tabular} & \hspace*{1pt}\begin{tabular}[t]{@{}c@{}}\rule{0pt}{12pt}\colorbox{BenchBest}{\strut 1m\,14.1s}\\[-.2pt]{\fontsize{5.8}{6.4}\selectfont \textcolor{ForestGreen}{$ \blacktriangle\,6.9\times $}}\\[-.2pt]\colorbox{BenchBest}{\strut 1m\,11.1s}\\[-.2pt]{\fontsize{5.8}{6.4}\selectfont \textcolor{ForestGreen}{$ \blacktriangle\,2.2\times $}}\end{tabular} & \hspace*{1pt}\begin{tabular}[t]{@{}c@{}}\rule{0pt}{12pt}\colorbox{BenchBest}{\strut 1m\,0.1s}\\[-.2pt]{\fontsize{5.8}{6.4}\selectfont \textcolor{ForestGreen}{$ \blacktriangle\,10.4\times $}}\\[-.2pt]\colorbox{BenchBest}{\strut 58.6s}\\[-.2pt]{\fontsize{5.8}{6.4}\selectfont \textcolor{ForestGreen}{$ \blacktriangle\,1.8\times $}}\end{tabular}\\[3pt]\hline
\rowcolor{BenchPreferred}[0.5pt][0.5pt]\BenchmarkSourceLabel{bench:aws-2-k20-w15:row:mpf-bounded-simd-block}{R2}\shortstack[l]{{\color{gray}R2}\\\textbf{BioseqDB}} & \begin{tabular}[t]{@{}l@{}}\rule{0pt}{12pt}MPF / Bounded\\SIMD-DP / Block\end{tabular} & \hspace*{1pt}\begin{tabular}[t]{@{}c@{}}\rule{0pt}{12pt}3s\\[-.2pt]{\fontsize{5.8}{6.4}\selectfont \textcolor{ForestGreen}{$ \blacktriangle\,6.6\times $}}\\[-.2pt]2.3s\\[-.2pt]{\fontsize{5.8}{6.4}\selectfont \textcolor{ForestGreen}{$ \blacktriangle\,8.1\times $}}\end{tabular} & \hspace*{1pt}\begin{tabular}[t]{@{}c@{}}\rule{0pt}{12pt}\colorbox{BenchBest}{\strut 2.6s}\\[-.2pt]{\fontsize{5.8}{6.4}\selectfont \textcolor{ForestGreen}{$ \blacktriangle\,6.9\times $}}\\[-.2pt]\colorbox{BenchBest}{\strut 1.9s}\\[-.2pt]{\fontsize{5.8}{6.4}\selectfont \textcolor{ForestGreen}{$ \blacktriangle\,8.6\times $}}\end{tabular} & \hspace*{1pt}\begin{tabular}[t]{@{}c@{}}\rule{0pt}{12pt}5.4s\\[-.2pt]{\fontsize{5.8}{6.4}\selectfont \textcolor{ForestGreen}{$ \blacktriangle\,7.9\times $}}\\[-.2pt]3.4s\\[-.2pt]{\fontsize{5.8}{6.4}\selectfont \textcolor{ForestGreen}{$ \blacktriangle\,11.7\times $}}\end{tabular} & \hspace*{1pt}\begin{tabular}[t]{@{}c@{}}\rule{0pt}{12pt}4.5s\\[-.2pt]{\fontsize{5.8}{6.4}\selectfont \textcolor{ForestGreen}{$ \blacktriangle\,86.7\times $}}\\[-.2pt]3.6s\\[-.2pt]{\fontsize{5.8}{6.4}\selectfont \textcolor{ForestGreen}{$ \blacktriangle\,97.9\times $}}\end{tabular} & \hspace*{1pt}\begin{tabular}[t]{@{}c@{}}\rule{0pt}{12pt}\colorbox{BenchBest}{\strut 7.5s}\\[-.2pt]{\fontsize{5.8}{6.4}\selectfont \textcolor{ForestGreen}{$ \blacktriangle\,19\% $}}\\[-.2pt]\colorbox{BenchBest}{\strut 5.2s}\\[-.2pt]{\fontsize{5.8}{6.4}\selectfont \textcolor{ForestGreen}{$ \blacktriangle\,1.6\times $}}\end{tabular} & \hspace*{1pt}\begin{tabular}[t]{@{}c@{}}\rule{0pt}{12pt}\colorbox{BenchBest}{\strut 5.8s}\\[-.2pt]{\fontsize{5.8}{6.4}\selectfont \textcolor{ForestGreen}{$ \blacktriangle\,1.5\times $}}\\[-.2pt]\colorbox{BenchBest}{\strut 4s}\\[-.2pt]{\fontsize{5.8}{6.4}\selectfont \textcolor{ForestGreen}{$ \blacktriangle\,1.7\times $}}\end{tabular} & \hspace*{1pt}\begin{tabular}[t]{@{}c@{}}\rule{0pt}{12pt}\colorbox{BenchBest}{\strut 12.6s}\\[-.2pt]{\fontsize{5.8}{6.4}\selectfont \textcolor{ForestGreen}{$ \blacktriangle\,21\% $}}\\[-.2pt]7.9s\\[-.2pt]{\fontsize{5.8}{6.4}\selectfont \textcolor{ForestGreen}{$ \blacktriangle\,1.6\times $}}\end{tabular} & \hspace*{1pt}\begin{tabular}[t]{@{}c@{}}\rule{0pt}{12pt}12.9s\\[-.2pt]{\fontsize{5.8}{6.4}\selectfont \textcolor{ForestGreen}{$ \blacktriangle\,1.4\times $}}\\[-.2pt]10.6s\\[-.2pt]{\fontsize{5.8}{6.4}\selectfont \textcolor{ForestGreen}{$ \blacktriangle\,1.6\times $}}\end{tabular} & \hspace*{1pt}\begin{tabular}[t]{@{}c@{}}\rule{0pt}{12pt}1m\,24.8s\\[-.2pt]{\fontsize{5.8}{6.4}\selectfont \textcolor{BrickRed}{$ \blacktriangledown\,2.0\times $}}\\[-.2pt]1m\,16.9s\\[-.2pt]{\fontsize{5.8}{6.4}\selectfont \textcolor{BrickRed}{$ \blacktriangledown\,1.9\times $}}\end{tabular} & \hspace*{1pt}\begin{tabular}[t]{@{}c@{}}\rule{0pt}{12pt}1m\,25.4s\\[-.2pt]{\fontsize{5.8}{6.4}\selectfont \textcolor{BrickRed}{$ \blacktriangledown\,2.0\times $}}\\[-.2pt]1m\,15.7s\\[-.2pt]{\fontsize{5.8}{6.4}\selectfont \textcolor{BrickRed}{$ \blacktriangledown\,1.8\times $}}\end{tabular} & \hspace*{1pt}\begin{tabular}[t]{@{}c@{}}\rule{0pt}{12pt}2m\,44s\\[-.2pt]{\fontsize{5.8}{6.4}\selectfont \textcolor{BrickRed}{$ \blacktriangledown\,1.9\times $}}\\[-.2pt]2m\,26.2s\\[-.2pt]{\fontsize{5.8}{6.4}\selectfont \textcolor{BrickRed}{$ \blacktriangledown\,1.7\times $}}\end{tabular} & \hspace*{1pt}\begin{tabular}[t]{@{}c@{}}\rule{0pt}{12pt}\colorbox{BenchBest}{\strut 1m\,42s}\\[-.2pt]{\fontsize{5.8}{6.4}\selectfont \textcolor{ForestGreen}{$ \blacktriangle\,2.9\times $}}\\[-.2pt]1m\,35.7s\\[-.2pt]{\fontsize{5.8}{6.4}\selectfont \textcolor{ForestGreen}{$ \blacktriangle\,3.0\times $}}\end{tabular} & \hspace*{1pt}\begin{tabular}[t]{@{}c@{}}\rule{0pt}{12pt}5m\,1.5s\\[-.2pt]{\fontsize{5.8}{6.4}\selectfont \textcolor{BrickRed}{$ \blacktriangledown\,8.6\times $}}\\[-.2pt]1m\,24.2s\\[-.2pt]{\fontsize{5.8}{6.4}\selectfont \textcolor{BrickRed}{$ \blacktriangledown\,2.5\times $}}\end{tabular} & \hspace*{1pt}\begin{tabular}[t]{@{}c@{}}\rule{0pt}{12pt}4m\,0.6s\\[-.2pt]{\fontsize{5.8}{6.4}\selectfont \textcolor{BrickRed}{$ \blacktriangledown\,6.7\times $}}\\[-.2pt]1m\,24.3s\\[-.2pt]{\fontsize{5.8}{6.4}\selectfont \textcolor{BrickRed}{$ \blacktriangledown\,2.5\times $}}\end{tabular} & \hspace*{1pt}\begin{tabular}[t]{@{}c@{}}\rule{0pt}{12pt}8m\,31.7s\\[-.2pt]{\fontsize{5.8}{6.4}\selectfont \textcolor{BrickRed}{$ \blacktriangledown\,6.9\times $}}\\[-.2pt]2m\,34.1s\\[-.2pt]{\fontsize{5.8}{6.4}\selectfont \textcolor{BrickRed}{$ \blacktriangledown\,2.2\times $}}\end{tabular} & \hspace*{1pt}\begin{tabular}[t]{@{}c@{}}\rule{0pt}{12pt}10m\,27.6s\\[-.2pt]{\fontsize{5.8}{6.4}\selectfont \textcolor{BrickRed}{$ \blacktriangledown\,10.4\times $}}\\[-.2pt]1m\,51.4s\\[-.2pt]{\fontsize{5.8}{6.4}\selectfont \textcolor{BrickRed}{$ \blacktriangledown\,1.9\times $}}\end{tabular}\\[3pt]\hline
\BenchmarkSourceLabel{bench:aws-2-k20-w15:row:mpf-bounded-rmq-block}{R3}\shortstack[l]{{\color{gray}R3}} & \begin{tabular}[t]{@{}l@{}}\rule{0pt}{12pt}MPF / Bounded\\RMQ / Block\end{tabular} & \hspace*{1pt}\begin{tabular}[t]{@{}c@{}}\rule{0pt}{12pt}\colorbox{BenchBest}{\strut 2.9s}\\[-.2pt]{\fontsize{5.8}{6.4}\selectfont \textcolor{ForestGreen}{$ \blacktriangle\,6.7\times $}}\\[-.2pt]2.2s\\[-.2pt]{\fontsize{5.8}{6.4}\selectfont \textcolor{ForestGreen}{$ \blacktriangle\,8.3\times $}}\end{tabular} & \hspace*{1pt}\begin{tabular}[t]{@{}c@{}}\rule{0pt}{12pt}2.7s\\[-.2pt]{\fontsize{5.8}{6.4}\selectfont \textcolor{ForestGreen}{$ \blacktriangle\,6.7\times $}}\\[-.2pt]2s\\[-.2pt]{\fontsize{5.8}{6.4}\selectfont \textcolor{ForestGreen}{$ \blacktriangle\,8.3\times $}}\end{tabular} & \hspace*{1pt}\begin{tabular}[t]{@{}c@{}}\rule{0pt}{12pt}\colorbox{BenchBest}{\strut 5.2s}\\[-.2pt]{\fontsize{5.8}{6.4}\selectfont \textcolor{ForestGreen}{$ \blacktriangle\,8.2\times $}}\\[-.2pt]3.3s\\[-.2pt]{\fontsize{5.8}{6.4}\selectfont \textcolor{ForestGreen}{$ \blacktriangle\,12.1\times $}}\end{tabular} & \hspace*{1pt}\begin{tabular}[t]{@{}c@{}}\rule{0pt}{12pt}4.7s\\[-.2pt]{\fontsize{5.8}{6.4}\selectfont \textcolor{ForestGreen}{$ \blacktriangle\,82.6\times $}}\\[-.2pt]3.9s\\[-.2pt]{\fontsize{5.8}{6.4}\selectfont \textcolor{ForestGreen}{$ \blacktriangle\,92.2\times $}}\end{tabular} & \hspace*{1pt}\begin{tabular}[t]{@{}c@{}}\rule{0pt}{12pt}11.1s\\[-.2pt]{\fontsize{5.8}{6.4}\selectfont \textcolor{BrickRed}{$ \blacktriangledown\,19\% $}}\\[-.2pt]7.1s\\[-.2pt]{\fontsize{5.8}{6.4}\selectfont \textcolor{ForestGreen}{$ \blacktriangle\,16\% $}}\end{tabular} & \hspace*{1pt}\begin{tabular}[t]{@{}c@{}}\rule{0pt}{12pt}8.8s\\[-.2pt]{\fontsize{5.8}{6.4}\selectfont \textcolor{gray}{$ \blacktriangleright\,1\% $}}\\[-.2pt]5.7s\\[-.2pt]{\fontsize{5.8}{6.4}\selectfont \textcolor{ForestGreen}{$ \blacktriangle\,19\% $}}\end{tabular} & \hspace*{1pt}\begin{tabular}[t]{@{}c@{}}\rule{0pt}{12pt}13.5s\\[-.2pt]{\fontsize{5.8}{6.4}\selectfont \textcolor{ForestGreen}{$ \blacktriangle\,15\% $}}\\[-.2pt]\colorbox{BenchBest}{\strut 7.5s}\\[-.2pt]{\fontsize{5.8}{6.4}\selectfont \textcolor{ForestGreen}{$ \blacktriangle\,1.7\times $}}\end{tabular} & \hspace*{1pt}\begin{tabular}[t]{@{}c@{}}\rule{0pt}{12pt}\colorbox{BenchBest}{\strut 11.6s}\\[-.2pt]{\fontsize{5.8}{6.4}\selectfont \textcolor{ForestGreen}{$ \blacktriangle\,1.6\times $}}\\[-.2pt]\colorbox{BenchBest}{\strut 9.4s}\\[-.2pt]{\fontsize{5.8}{6.4}\selectfont \textcolor{ForestGreen}{$ \blacktriangle\,1.8\times $}}\end{tabular} & \hspace*{1pt}\begin{tabular}[t]{@{}c@{}}\rule{0pt}{12pt}1m\,34s\\[-.2pt]{\fontsize{5.8}{6.4}\selectfont \textcolor{BrickRed}{$ \blacktriangledown\,2.2\times $}}\\[-.2pt]1m\,24.9s\\[-.2pt]{\fontsize{5.8}{6.4}\selectfont \textcolor{BrickRed}{$ \blacktriangledown\,2.1\times $}}\end{tabular} & \hspace*{1pt}\begin{tabular}[t]{@{}c@{}}\rule{0pt}{12pt}1m\,33s\\[-.2pt]{\fontsize{5.8}{6.4}\selectfont \textcolor{BrickRed}{$ \blacktriangledown\,2.2\times $}}\\[-.2pt]1m\,22.3s\\[-.2pt]{\fontsize{5.8}{6.4}\selectfont \textcolor{BrickRed}{$ \blacktriangledown\,2.0\times $}}\end{tabular} & \textcolor{gray}{\ensuremath{\cdot}} & \hspace*{1pt}\begin{tabular}[t]{@{}c@{}}\rule{0pt}{12pt}1m\,44.1s\\[-.2pt]{\fontsize{5.8}{6.4}\selectfont \textcolor{ForestGreen}{$ \blacktriangle\,2.8\times $}}\\[-.2pt]1m\,38.1s\\[-.2pt]{\fontsize{5.8}{6.4}\selectfont \textcolor{ForestGreen}{$ \blacktriangle\,2.9\times $}}\end{tabular} & \hspace*{1pt}\begin{tabular}[t]{@{}c@{}}\rule{0pt}{12pt}5m\,11.1s\\[-.2pt]{\fontsize{5.8}{6.4}\selectfont \textcolor{BrickRed}{$ \blacktriangledown\,8.9\times $}}\\[-.2pt]1m\,29.9s\\[-.2pt]{\fontsize{5.8}{6.4}\selectfont \textcolor{BrickRed}{$ \blacktriangledown\,2.6\times $}}\end{tabular} & \hspace*{1pt}\begin{tabular}[t]{@{}c@{}}\rule{0pt}{12pt}4m\,41.1s\\[-.2pt]{\fontsize{5.8}{6.4}\selectfont \textcolor{BrickRed}{$ \blacktriangledown\,7.8\times $}}\\[-.2pt]1m\,27.4s\\[-.2pt]{\fontsize{5.8}{6.4}\selectfont \textcolor{BrickRed}{$ \blacktriangledown\,2.6\times $}}\end{tabular} & \textcolor{gray}{\ensuremath{\cdot}} & \hspace*{1pt}\begin{tabular}[t]{@{}c@{}}\rule{0pt}{12pt}10m\,23.5s\\[-.2pt]{\fontsize{5.8}{6.4}\selectfont \textcolor{BrickRed}{$ \blacktriangledown\,10.4\times $}}\\[-.2pt]1m\,47.3s\\[-.2pt]{\fontsize{5.8}{6.4}\selectfont \textcolor{BrickRed}{$ \blacktriangledown\,1.8\times $}}\end{tabular}\\[3pt]\hline
\BenchmarkSourceLabel{bench:aws-2-k20-w15:row:mpf-bounded-scalar-block}{R4}\shortstack[l]{{\color{gray}R4}} & \begin{tabular}[t]{@{}l@{}}\rule{0pt}{12pt}MPF / Bounded\\Scalar-DP / Block\end{tabular} & \hspace*{1pt}\begin{tabular}[t]{@{}c@{}}\rule{0pt}{12pt}3s\\[-.2pt]{\fontsize{5.8}{6.4}\selectfont \textcolor{ForestGreen}{$ \blacktriangle\,6.6\times $}}\\[-.2pt]2.3s\\[-.2pt]{\fontsize{5.8}{6.4}\selectfont \textcolor{ForestGreen}{$ \blacktriangle\,8.1\times $}}\end{tabular} & \hspace*{1pt}\begin{tabular}[t]{@{}c@{}}\rule{0pt}{12pt}2.8s\\[-.2pt]{\fontsize{5.8}{6.4}\selectfont \textcolor{ForestGreen}{$ \blacktriangle\,6.5\times $}}\\[-.2pt]2s\\[-.2pt]{\fontsize{5.8}{6.4}\selectfont \textcolor{ForestGreen}{$ \blacktriangle\,8.2\times $}}\end{tabular} & \hspace*{1pt}\begin{tabular}[t]{@{}c@{}}\rule{0pt}{12pt}5.3s\\[-.2pt]{\fontsize{5.8}{6.4}\selectfont \textcolor{ForestGreen}{$ \blacktriangle\,8.0\times $}}\\[-.2pt]3.3s\\[-.2pt]{\fontsize{5.8}{6.4}\selectfont \textcolor{ForestGreen}{$ \blacktriangle\,11.8\times $}}\end{tabular} & \hspace*{1pt}\begin{tabular}[t]{@{}c@{}}\rule{0pt}{12pt}\colorbox{BenchBest}{\strut 4.2s}\\[-.2pt]{\fontsize{5.8}{6.4}\selectfont \textcolor{ForestGreen}{$ \blacktriangle\,92.8\times $}}\\[-.2pt]3.3s\\[-.2pt]{\fontsize{5.8}{6.4}\selectfont \textcolor{ForestGreen}{$ \blacktriangle\,106.7\times $}}\end{tabular} & \hspace*{1pt}\begin{tabular}[t]{@{}c@{}}\rule{0pt}{12pt}9.5s\\[-.2pt]{\fontsize{5.8}{6.4}\selectfont \textcolor{gray}{$ \blacktriangleright\,2\% $}}\\[-.2pt]5.3s\\[-.2pt]{\fontsize{5.8}{6.4}\selectfont \textcolor{ForestGreen}{$ \blacktriangle\,1.6\times $}}\end{tabular} & \hspace*{1pt}\begin{tabular}[t]{@{}c@{}}\rule{0pt}{12pt}8.1s\\[-.2pt]{\fontsize{5.8}{6.4}\selectfont \textcolor{ForestGreen}{$ \blacktriangle\,7\% $}}\\[-.2pt]4.7s\\[-.2pt]{\fontsize{5.8}{6.4}\selectfont \textcolor{ForestGreen}{$ \blacktriangle\,1.5\times $}}\end{tabular} & \hspace*{1pt}\begin{tabular}[t]{@{}c@{}}\rule{0pt}{12pt}13.7s\\[-.2pt]{\fontsize{5.8}{6.4}\selectfont \textcolor{ForestGreen}{$ \blacktriangle\,14\% $}}\\[-.2pt]7.5s\\[-.2pt]{\fontsize{5.8}{6.4}\selectfont \textcolor{ForestGreen}{$ \blacktriangle\,1.7\times $}}\end{tabular} & \hspace*{1pt}\begin{tabular}[t]{@{}c@{}}\rule{0pt}{12pt}12.8s\\[-.2pt]{\fontsize{5.8}{6.4}\selectfont \textcolor{ForestGreen}{$ \blacktriangle\,1.4\times $}}\\[-.2pt]10.6s\\[-.2pt]{\fontsize{5.8}{6.4}\selectfont \textcolor{ForestGreen}{$ \blacktriangle\,1.6\times $}}\end{tabular} & \hspace*{1pt}\begin{tabular}[t]{@{}c@{}}\rule{0pt}{12pt}1m\,32.9s\\[-.2pt]{\fontsize{5.8}{6.4}\selectfont \textcolor{BrickRed}{$ \blacktriangledown\,2.2\times $}}\\[-.2pt]1m\,23.4s\\[-.2pt]{\fontsize{5.8}{6.4}\selectfont \textcolor{BrickRed}{$ \blacktriangledown\,2.0\times $}}\end{tabular} & \hspace*{1pt}\begin{tabular}[t]{@{}c@{}}\rule{0pt}{12pt}1m\,31.2s\\[-.2pt]{\fontsize{5.8}{6.4}\selectfont \textcolor{BrickRed}{$ \blacktriangledown\,2.1\times $}}\\[-.2pt]1m\,21s\\[-.2pt]{\fontsize{5.8}{6.4}\selectfont \textcolor{BrickRed}{$ \blacktriangledown\,2.0\times $}}\end{tabular} & \textcolor{gray}{\ensuremath{\cdot}} & \hspace*{1pt}\begin{tabular}[t]{@{}c@{}}\rule{0pt}{12pt}1m\,43.1s\\[-.2pt]{\fontsize{5.8}{6.4}\selectfont \textcolor{ForestGreen}{$ \blacktriangle\,2.8\times $}}\\[-.2pt]1m\,36.4s\\[-.2pt]{\fontsize{5.8}{6.4}\selectfont \textcolor{ForestGreen}{$ \blacktriangle\,3.0\times $}}\end{tabular} & \hspace*{1pt}\begin{tabular}[t]{@{}c@{}}\rule{0pt}{12pt}5m\,11.8s\\[-.2pt]{\fontsize{5.8}{6.4}\selectfont \textcolor{BrickRed}{$ \blacktriangledown\,8.9\times $}}\\[-.2pt]1m\,30.8s\\[-.2pt]{\fontsize{5.8}{6.4}\selectfont \textcolor{BrickRed}{$ \blacktriangledown\,2.7\times $}}\end{tabular} & \hspace*{1pt}\begin{tabular}[t]{@{}c@{}}\rule{0pt}{12pt}4m\,49.9s\\[-.2pt]{\fontsize{5.8}{6.4}\selectfont \textcolor{BrickRed}{$ \blacktriangledown\,8.1\times $}}\\[-.2pt]1m\,33.6s\\[-.2pt]{\fontsize{5.8}{6.4}\selectfont \textcolor{BrickRed}{$ \blacktriangledown\,2.7\times $}}\end{tabular} & \textcolor{gray}{\ensuremath{\cdot}} & \hspace*{1pt}\begin{tabular}[t]{@{}c@{}}\rule{0pt}{12pt}10m\,26.2s\\[-.2pt]{\fontsize{5.8}{6.4}\selectfont \textcolor{BrickRed}{$ \blacktriangledown\,10.4\times $}}\\[-.2pt]1m\,48.4s\\[-.2pt]{\fontsize{5.8}{6.4}\selectfont \textcolor{BrickRed}{$ \blacktriangledown\,1.9\times $}}\end{tabular}\\[3pt]\hline
\BenchmarkSourceLabel{bench:aws-2-k20-w15:row:mpf-bounded-simd-ksw2}{R5}\shortstack[l]{{\color{gray}R5}} & \begin{tabular}[t]{@{}l@{}}\rule{0pt}{12pt}MPF / Bounded\\SIMD-DP / KSW2\end{tabular} & \hspace*{1pt}\begin{tabular}[t]{@{}c@{}}\rule{0pt}{12pt}4.6s\\[-.2pt]{\fontsize{5.8}{6.4}\selectfont \textcolor{ForestGreen}{$ \blacktriangle\,4.3\times $}}\\[-.2pt]2.2s\\[-.2pt]{\fontsize{5.8}{6.4}\selectfont \textcolor{ForestGreen}{$ \blacktriangle\,8.3\times $}}\end{tabular} & \hspace*{1pt}\begin{tabular}[t]{@{}c@{}}\rule{0pt}{12pt}4.1s\\[-.2pt]{\fontsize{5.8}{6.4}\selectfont \textcolor{ForestGreen}{$ \blacktriangle\,4.5\times $}}\\[-.2pt]1.9s\\[-.2pt]{\fontsize{5.8}{6.4}\selectfont \textcolor{ForestGreen}{$ \blacktriangle\,8.5\times $}}\end{tabular} & \hspace*{1pt}\begin{tabular}[t]{@{}c@{}}\rule{0pt}{12pt}9.9s\\[-.2pt]{\fontsize{5.8}{6.4}\selectfont \textcolor{ForestGreen}{$ \blacktriangle\,4.3\times $}}\\[-.2pt]3.3s\\[-.2pt]{\fontsize{5.8}{6.4}\selectfont \textcolor{ForestGreen}{$ \blacktriangle\,12.1\times $}}\end{tabular} & \hspace*{1pt}\begin{tabular}[t]{@{}c@{}}\rule{0pt}{12pt}4.7s\\[-.2pt]{\fontsize{5.8}{6.4}\selectfont \textcolor{ForestGreen}{$ \blacktriangle\,83.2\times $}}\\[-.2pt]\colorbox{BenchBest}{\strut 1.9s}\\[-.2pt]{\fontsize{5.8}{6.4}\selectfont \textcolor{ForestGreen}{$ \blacktriangle\,184.9\times $}}\end{tabular} & \hspace*{1pt}\begin{tabular}[t]{@{}c@{}}\rule{0pt}{12pt}11.5s\\[-.2pt]{\fontsize{5.8}{6.4}\selectfont \textcolor{BrickRed}{$ \blacktriangledown\,23\% $}}\\[-.2pt]7.4s\\[-.2pt]{\fontsize{5.8}{6.4}\selectfont \textcolor{ForestGreen}{$ \blacktriangle\,12\% $}}\end{tabular} & \hspace*{1pt}\begin{tabular}[t]{@{}c@{}}\rule{0pt}{12pt}10.8s\\[-.2pt]{\fontsize{5.8}{6.4}\selectfont \textcolor{BrickRed}{$ \blacktriangledown\,23\% $}}\\[-.2pt]7.1s\\[-.2pt]{\fontsize{5.8}{6.4}\selectfont \textcolor{gray}{$ \blacktriangleright\,1\% $}}\end{tabular} & \hspace*{1pt}\begin{tabular}[t]{@{}c@{}}\rule{0pt}{12pt}15.4s\\[-.2pt]{\fontsize{5.8}{6.4}\selectfont \textcolor{gray}{$ \blacktriangleright\,2\% $}}\\[-.2pt]8.6s\\[-.2pt]{\fontsize{5.8}{6.4}\selectfont \textcolor{ForestGreen}{$ \blacktriangle\,1.5\times $}}\end{tabular} & \hspace*{1pt}\begin{tabular}[t]{@{}c@{}}\rule{0pt}{12pt}12.5s\\[-.2pt]{\fontsize{5.8}{6.4}\selectfont \textcolor{ForestGreen}{$ \blacktriangle\,1.5\times $}}\\[-.2pt]9.7s\\[-.2pt]{\fontsize{5.8}{6.4}\selectfont \textcolor{ForestGreen}{$ \blacktriangle\,1.8\times $}}\end{tabular} & \hspace*{1pt}\begin{tabular}[t]{@{}c@{}}\rule{0pt}{12pt}1m\,34.1s\\[-.2pt]{\fontsize{5.8}{6.4}\selectfont \textcolor{BrickRed}{$ \blacktriangledown\,2.2\times $}}\\[-.2pt]1m\,24s\\[-.2pt]{\fontsize{5.8}{6.4}\selectfont \textcolor{BrickRed}{$ \blacktriangledown\,2.0\times $}}\end{tabular} & \hspace*{1pt}\begin{tabular}[t]{@{}c@{}}\rule{0pt}{12pt}1m\,32.2s\\[-.2pt]{\fontsize{5.8}{6.4}\selectfont \textcolor{BrickRed}{$ \blacktriangledown\,2.2\times $}}\\[-.2pt]1m\,20s\\[-.2pt]{\fontsize{5.8}{6.4}\selectfont \textcolor{BrickRed}{$ \blacktriangledown\,1.9\times $}}\end{tabular} & \textcolor{gray}{\ensuremath{\cdot}} & \hspace*{1pt}\begin{tabular}[t]{@{}c@{}}\rule{0pt}{12pt}1m\,46.2s\\[-.2pt]{\fontsize{5.8}{6.4}\selectfont \textcolor{ForestGreen}{$ \blacktriangle\,2.7\times $}}\\[-.2pt]1m\,35.2s\\[-.2pt]{\fontsize{5.8}{6.4}\selectfont \textcolor{ForestGreen}{$ \blacktriangle\,3.0\times $}}\end{tabular} & \hspace*{1pt}\begin{tabular}[t]{@{}c@{}}\rule{0pt}{12pt}5m\,10.2s\\[-.2pt]{\fontsize{5.8}{6.4}\selectfont \textcolor{BrickRed}{$ \blacktriangledown\,8.9\times $}}\\[-.2pt]1m\,27.3s\\[-.2pt]{\fontsize{5.8}{6.4}\selectfont \textcolor{BrickRed}{$ \blacktriangledown\,2.6\times $}}\end{tabular} & \hspace*{1pt}\begin{tabular}[t]{@{}c@{}}\rule{0pt}{12pt}4m\,45.4s\\[-.2pt]{\fontsize{5.8}{6.4}\selectfont \textcolor{BrickRed}{$ \blacktriangledown\,7.9\times $}}\\[-.2pt]1m\,30.5s\\[-.2pt]{\fontsize{5.8}{6.4}\selectfont \textcolor{BrickRed}{$ \blacktriangledown\,2.6\times $}}\end{tabular} & \textcolor{gray}{\ensuremath{\cdot}} & \hspace*{1pt}\begin{tabular}[t]{@{}c@{}}\rule{0pt}{12pt}10m\,25.5s\\[-.2pt]{\fontsize{5.8}{6.4}\selectfont \textcolor{BrickRed}{$ \blacktriangledown\,10.4\times $}}\\[-.2pt]1m\,44.7s\\[-.2pt]{\fontsize{5.8}{6.4}\selectfont \textcolor{BrickRed}{$ \blacktriangledown\,1.8\times $}}\end{tabular}\\[3pt]\hline
\BenchmarkSourceLabel{bench:aws-2-k20-w15:row:mpf-bounded-scalar-ksw2}{R6}\shortstack[l]{{\color{gray}R6}} & \begin{tabular}[t]{@{}l@{}}\rule{0pt}{12pt}MPF / Bounded\\Scalar-DP / KSW2\end{tabular} & \hspace*{1pt}\begin{tabular}[t]{@{}c@{}}\rule{0pt}{12pt}4.6s\\[-.2pt]{\fontsize{5.8}{6.4}\selectfont \textcolor{ForestGreen}{$ \blacktriangle\,4.3\times $}}\\[-.2pt]2.2s\\[-.2pt]{\fontsize{5.8}{6.4}\selectfont \textcolor{ForestGreen}{$ \blacktriangle\,8.3\times $}}\end{tabular} & \hspace*{1pt}\begin{tabular}[t]{@{}c@{}}\rule{0pt}{12pt}4.1s\\[-.2pt]{\fontsize{5.8}{6.4}\selectfont \textcolor{ForestGreen}{$ \blacktriangle\,4.4\times $}}\\[-.2pt]1.9s\\[-.2pt]{\fontsize{5.8}{6.4}\selectfont \textcolor{ForestGreen}{$ \blacktriangle\,8.6\times $}}\end{tabular} & \hspace*{1pt}\begin{tabular}[t]{@{}c@{}}\rule{0pt}{12pt}9.7s\\[-.2pt]{\fontsize{5.8}{6.4}\selectfont \textcolor{ForestGreen}{$ \blacktriangle\,4.4\times $}}\\[-.2pt]\colorbox{BenchBest}{\strut 3.1s}\\[-.2pt]{\fontsize{5.8}{6.4}\selectfont \textcolor{ForestGreen}{$ \blacktriangle\,12.6\times $}}\end{tabular} & \hspace*{1pt}\begin{tabular}[t]{@{}c@{}}\rule{0pt}{12pt}5.8s\\[-.2pt]{\fontsize{5.8}{6.4}\selectfont \textcolor{ForestGreen}{$ \blacktriangle\,67.2\times $}}\\[-.2pt]3s\\[-.2pt]{\fontsize{5.8}{6.4}\selectfont \textcolor{ForestGreen}{$ \blacktriangle\,119.5\times $}}\end{tabular} & \hspace*{1pt}\begin{tabular}[t]{@{}c@{}}\rule{0pt}{12pt}9.5s\\[-.2pt]{\fontsize{5.8}{6.4}\selectfont \textcolor{gray}{$ \blacktriangleright\,2\% $}}\\[-.2pt]5.3s\\[-.2pt]{\fontsize{5.8}{6.4}\selectfont \textcolor{ForestGreen}{$ \blacktriangle\,1.6\times $}}\end{tabular} & \hspace*{1pt}\begin{tabular}[t]{@{}c@{}}\rule{0pt}{12pt}8.3s\\[-.2pt]{\fontsize{5.8}{6.4}\selectfont \textcolor{gray}{$ \blacktriangleright\,5\% $}}\\[-.2pt]4.6s\\[-.2pt]{\fontsize{5.8}{6.4}\selectfont \textcolor{ForestGreen}{$ \blacktriangle\,1.5\times $}}\end{tabular} & \hspace*{1pt}\begin{tabular}[t]{@{}c@{}}\rule{0pt}{12pt}14.2s\\[-.2pt]{\fontsize{5.8}{6.4}\selectfont \textcolor{ForestGreen}{$ \blacktriangle\,10\% $}}\\[-.2pt]7.5s\\[-.2pt]{\fontsize{5.8}{6.4}\selectfont \textcolor{ForestGreen}{$ \blacktriangle\,1.7\times $}}\end{tabular} & \hspace*{1pt}\begin{tabular}[t]{@{}c@{}}\rule{0pt}{12pt}13.3s\\[-.2pt]{\fontsize{5.8}{6.4}\selectfont \textcolor{ForestGreen}{$ \blacktriangle\,28\% $}}\\[-.2pt]10.5s\\[-.2pt]{\fontsize{5.8}{6.4}\selectfont \textcolor{ForestGreen}{$ \blacktriangle\,1.6\times $}}\end{tabular} & \hspace*{1pt}\begin{tabular}[t]{@{}c@{}}\rule{0pt}{12pt}1m\,27.2s\\[-.2pt]{\fontsize{5.8}{6.4}\selectfont \textcolor{BrickRed}{$ \blacktriangledown\,2.1\times $}}\\[-.2pt]1m\,17s\\[-.2pt]{\fontsize{5.8}{6.4}\selectfont \textcolor{BrickRed}{$ \blacktriangledown\,1.9\times $}}\end{tabular} & \hspace*{1pt}\begin{tabular}[t]{@{}c@{}}\rule{0pt}{12pt}1m\,39.6s\\[-.2pt]{\fontsize{5.8}{6.4}\selectfont \textcolor{BrickRed}{$ \blacktriangledown\,2.3\times $}}\\[-.2pt]1m\,26.7s\\[-.2pt]{\fontsize{5.8}{6.4}\selectfont \textcolor{BrickRed}{$ \blacktriangledown\,2.1\times $}}\end{tabular} & \textcolor{gray}{\ensuremath{\cdot}} & \hspace*{1pt}\begin{tabular}[t]{@{}c@{}}\rule{0pt}{12pt}1m\,46.4s\\[-.2pt]{\fontsize{5.8}{6.4}\selectfont \textcolor{ForestGreen}{$ \blacktriangle\,2.7\times $}}\\[-.2pt]1m\,36.3s\\[-.2pt]{\fontsize{5.8}{6.4}\selectfont \textcolor{ForestGreen}{$ \blacktriangle\,3.0\times $}}\end{tabular} & \hspace*{1pt}\begin{tabular}[t]{@{}c@{}}\rule{0pt}{12pt}5m\,10.4s\\[-.2pt]{\fontsize{5.8}{6.4}\selectfont \textcolor{BrickRed}{$ \blacktriangledown\,8.9\times $}}\\[-.2pt]1m\,27.2s\\[-.2pt]{\fontsize{5.8}{6.4}\selectfont \textcolor{BrickRed}{$ \blacktriangledown\,2.6\times $}}\end{tabular} & \hspace*{1pt}\begin{tabular}[t]{@{}c@{}}\rule{0pt}{12pt}4m\,48.1s\\[-.2pt]{\fontsize{5.8}{6.4}\selectfont \textcolor{BrickRed}{$ \blacktriangledown\,8.0\times $}}\\[-.2pt]1m\,29.7s\\[-.2pt]{\fontsize{5.8}{6.4}\selectfont \textcolor{BrickRed}{$ \blacktriangledown\,2.6\times $}}\end{tabular} & \textcolor{gray}{\ensuremath{\cdot}} & \hspace*{1pt}\begin{tabular}[t]{@{}c@{}}\rule{0pt}{12pt}10m\,29.2s\\[-.2pt]{\fontsize{5.8}{6.4}\selectfont \textcolor{BrickRed}{$ \blacktriangledown\,10.5\times $}}\\[-.2pt]1m\,47.6s\\[-.2pt]{\fontsize{5.8}{6.4}\selectfont \textcolor{BrickRed}{$ \blacktriangledown\,1.8\times $}}\end{tabular}\\[3pt]\hline
\BenchmarkSourceLabel{bench:aws-2-k20-w15:row:mpf-bounded-rmq-ksw2}{R7}\shortstack[l]{{\color{gray}R7}} & \begin{tabular}[t]{@{}l@{}}\rule{0pt}{12pt}MPF / Bounded\\RMQ / KSW2\end{tabular} & \hspace*{1pt}\begin{tabular}[t]{@{}c@{}}\rule{0pt}{12pt}4.5s\\[-.2pt]{\fontsize{5.8}{6.4}\selectfont \textcolor{ForestGreen}{$ \blacktriangle\,4.3\times $}}\\[-.2pt]\colorbox{BenchBest}{\strut 2.2s}\\[-.2pt]{\fontsize{5.8}{6.4}\selectfont \textcolor{ForestGreen}{$ \blacktriangle\,8.4\times $}}\end{tabular} & \hspace*{1pt}\begin{tabular}[t]{@{}c@{}}\rule{0pt}{12pt}4.2s\\[-.2pt]{\fontsize{5.8}{6.4}\selectfont \textcolor{ForestGreen}{$ \blacktriangle\,4.3\times $}}\\[-.2pt]2s\\[-.2pt]{\fontsize{5.8}{6.4}\selectfont \textcolor{ForestGreen}{$ \blacktriangle\,8.4\times $}}\end{tabular} & \hspace*{1pt}\begin{tabular}[t]{@{}c@{}}\rule{0pt}{12pt}9.7s\\[-.2pt]{\fontsize{5.8}{6.4}\selectfont \textcolor{ForestGreen}{$ \blacktriangle\,4.4\times $}}\\[-.2pt]3.2s\\[-.2pt]{\fontsize{5.8}{6.4}\selectfont \textcolor{ForestGreen}{$ \blacktriangle\,12.4\times $}}\end{tabular} & \hspace*{1pt}\begin{tabular}[t]{@{}c@{}}\rule{0pt}{12pt}6.2s\\[-.2pt]{\fontsize{5.8}{6.4}\selectfont \textcolor{ForestGreen}{$ \blacktriangle\,63.0\times $}}\\[-.2pt]3.4s\\[-.2pt]{\fontsize{5.8}{6.4}\selectfont \textcolor{ForestGreen}{$ \blacktriangle\,104.6\times $}}\end{tabular} & \hspace*{1pt}\begin{tabular}[t]{@{}c@{}}\rule{0pt}{12pt}9.4s\\[-.2pt]{\fontsize{5.8}{6.4}\selectfont \textcolor{gray}{$ \blacktriangleright\,1\% $}}\\[-.2pt]5.3s\\[-.2pt]{\fontsize{5.8}{6.4}\selectfont \textcolor{ForestGreen}{$ \blacktriangle\,1.6\times $}}\end{tabular} & \hspace*{1pt}\begin{tabular}[t]{@{}c@{}}\rule{0pt}{12pt}8.1s\\[-.2pt]{\fontsize{5.8}{6.4}\selectfont \textcolor{ForestGreen}{$ \blacktriangle\,8\% $}}\\[-.2pt]4.6s\\[-.2pt]{\fontsize{5.8}{6.4}\selectfont \textcolor{ForestGreen}{$ \blacktriangle\,1.5\times $}}\end{tabular} & \hspace*{1pt}\begin{tabular}[t]{@{}c@{}}\rule{0pt}{12pt}15.6s\\[-.2pt]{\fontsize{5.8}{6.4}\selectfont \textcolor{gray}{$ \blacktriangleright\,1\% $}}\\[-.2pt]9.1s\\[-.2pt]{\fontsize{5.8}{6.4}\selectfont \textcolor{ForestGreen}{$ \blacktriangle\,28\% $}}\end{tabular} & \hspace*{1pt}\begin{tabular}[t]{@{}c@{}}\rule{0pt}{12pt}12.6s\\[-.2pt]{\fontsize{5.8}{6.4}\selectfont \textcolor{ForestGreen}{$ \blacktriangle\,1.5\times $}}\\[-.2pt]9.8s\\[-.2pt]{\fontsize{5.8}{6.4}\selectfont \textcolor{ForestGreen}{$ \blacktriangle\,1.7\times $}}\end{tabular} & \hspace*{1pt}\begin{tabular}[t]{@{}c@{}}\rule{0pt}{12pt}1m\,32.9s\\[-.2pt]{\fontsize{5.8}{6.4}\selectfont \textcolor{BrickRed}{$ \blacktriangledown\,2.2\times $}}\\[-.2pt]1m\,23.1s\\[-.2pt]{\fontsize{5.8}{6.4}\selectfont \textcolor{BrickRed}{$ \blacktriangledown\,2.0\times $}}\end{tabular} & \hspace*{1pt}\begin{tabular}[t]{@{}c@{}}\rule{0pt}{12pt}1m\,41.2s\\[-.2pt]{\fontsize{5.8}{6.4}\selectfont \textcolor{BrickRed}{$ \blacktriangledown\,2.4\times $}}\\[-.2pt]1m\,28.4s\\[-.2pt]{\fontsize{5.8}{6.4}\selectfont \textcolor{BrickRed}{$ \blacktriangledown\,2.2\times $}}\end{tabular} & \textcolor{gray}{\ensuremath{\cdot}} & \hspace*{1pt}\begin{tabular}[t]{@{}c@{}}\rule{0pt}{12pt}1m\,46.6s\\[-.2pt]{\fontsize{5.8}{6.4}\selectfont \textcolor{ForestGreen}{$ \blacktriangle\,2.7\times $}}\\[-.2pt]1m\,35.5s\\[-.2pt]{\fontsize{5.8}{6.4}\selectfont \textcolor{ForestGreen}{$ \blacktriangle\,3.0\times $}}\end{tabular} & \hspace*{1pt}\begin{tabular}[t]{@{}c@{}}\rule{0pt}{12pt}5m\,9.8s\\[-.2pt]{\fontsize{5.8}{6.4}\selectfont \textcolor{BrickRed}{$ \blacktriangledown\,8.9\times $}}\\[-.2pt]1m\,27.1s\\[-.2pt]{\fontsize{5.8}{6.4}\selectfont \textcolor{BrickRed}{$ \blacktriangledown\,2.6\times $}}\end{tabular} & \hspace*{1pt}\begin{tabular}[t]{@{}c@{}}\rule{0pt}{12pt}4m\,52.9s\\[-.2pt]{\fontsize{5.8}{6.4}\selectfont \textcolor{BrickRed}{$ \blacktriangledown\,8.2\times $}}\\[-.2pt]1m\,32.9s\\[-.2pt]{\fontsize{5.8}{6.4}\selectfont \textcolor{BrickRed}{$ \blacktriangledown\,2.7\times $}}\end{tabular} & \textcolor{gray}{\ensuremath{\cdot}} & \hspace*{1pt}\begin{tabular}[t]{@{}c@{}}\rule{0pt}{12pt}10m\,33.7s\\[-.2pt]{\fontsize{5.8}{6.4}\selectfont \textcolor{BrickRed}{$ \blacktriangledown\,10.5\times $}}\\[-.2pt]1m\,55.3s\\[-.2pt]{\fontsize{5.8}{6.4}\selectfont \textcolor{BrickRed}{$ \blacktriangledown\,2.0\times $}}\end{tabular}\\[3pt]\hline
\BenchmarkSourceLabel{bench:aws-2-k20-w15:row:mps-bounded-scalar-block}{R8}\shortstack[l]{{\color{gray}R8}} & \begin{tabular}[t]{@{}l@{}}\rule{0pt}{12pt}MPS / Bounded\\Scalar-DP / Block\end{tabular} & \hspace*{1pt}\begin{tabular}[t]{@{}c@{}}\rule{0pt}{12pt}6.6s\\[-.2pt]{\fontsize{5.8}{6.4}\selectfont \textcolor{ForestGreen}{$ \blacktriangle\,3.0\times $}}\\[-.2pt]5.9s\\[-.2pt]{\fontsize{5.8}{6.4}\selectfont \textcolor{ForestGreen}{$ \blacktriangle\,3.2\times $}}\end{tabular} & \hspace*{1pt}\begin{tabular}[t]{@{}c@{}}\rule{0pt}{12pt}6.6s\\[-.2pt]{\fontsize{5.8}{6.4}\selectfont \textcolor{ForestGreen}{$ \blacktriangle\,2.7\times $}}\\[-.2pt]5.9s\\[-.2pt]{\fontsize{5.8}{6.4}\selectfont \textcolor{ForestGreen}{$ \blacktriangle\,2.8\times $}}\end{tabular} & \hspace*{1pt}\begin{tabular}[t]{@{}c@{}}\rule{0pt}{12pt}12.8s\\[-.2pt]{\fontsize{5.8}{6.4}\selectfont \textcolor{ForestGreen}{$ \blacktriangle\,3.3\times $}}\\[-.2pt]10.5s\\[-.2pt]{\fontsize{5.8}{6.4}\selectfont \textcolor{ForestGreen}{$ \blacktriangle\,3.8\times $}}\end{tabular} & \hspace*{1pt}\begin{tabular}[t]{@{}c@{}}\rule{0pt}{12pt}6.3s\\[-.2pt]{\fontsize{5.8}{6.4}\selectfont \textcolor{ForestGreen}{$ \blacktriangle\,61.6\times $}}\\[-.2pt]5.4s\\[-.2pt]{\fontsize{5.8}{6.4}\selectfont \textcolor{ForestGreen}{$ \blacktriangle\,65.7\times $}}\end{tabular} & \hspace*{1pt}\begin{tabular}[t]{@{}c@{}}\rule{0pt}{12pt}4m\,37.1s\\[-.2pt]{\fontsize{5.8}{6.4}\selectfont \textcolor{BrickRed}{$ \blacktriangledown\,29.7\times $}}\\[-.2pt]20s\\[-.2pt]{\fontsize{5.8}{6.4}\selectfont \textcolor{BrickRed}{$ \blacktriangledown\,2.4\times $}}\end{tabular} & \hspace*{1pt}\begin{tabular}[t]{@{}c@{}}\rule{0pt}{12pt}3m\,55.5s\\[-.2pt]{\fontsize{5.8}{6.4}\selectfont \textcolor{BrickRed}{$ \blacktriangledown\,27.0\times $}}\\[-.2pt]19.7s\\[-.2pt]{\fontsize{5.8}{6.4}\selectfont \textcolor{BrickRed}{$ \blacktriangledown\,2.8\times $}}\end{tabular} & \textcolor{gray}{\ensuremath{\cdot}} & \textcolor{gray}{\ensuremath{\cdot}} & \hspace*{1pt}\begin{tabular}[t]{@{}c@{}}\rule{0pt}{12pt}1m\,42.1s\\[-.2pt]{\fontsize{5.8}{6.4}\selectfont \textcolor{BrickRed}{$ \blacktriangledown\,2.4\times $}}\\[-.2pt]1m\,20s\\[-.2pt]{\fontsize{5.8}{6.4}\selectfont \textcolor{BrickRed}{$ \blacktriangledown\,1.9\times $}}\end{tabular} & \hspace*{1pt}\begin{tabular}[t]{@{}c@{}}\rule{0pt}{12pt}1m\,37.4s\\[-.2pt]{\fontsize{5.8}{6.4}\selectfont \textcolor{BrickRed}{$ \blacktriangledown\,2.3\times $}}\\[-.2pt]1m\,18s\\[-.2pt]{\fontsize{5.8}{6.4}\selectfont \textcolor{BrickRed}{$ \blacktriangledown\,1.9\times $}}\end{tabular} & \textcolor{gray}{\ensuremath{\cdot}} & \hspace*{1pt}\begin{tabular}[t]{@{}c@{}}\rule{0pt}{12pt}1m\,56.9s\\[-.2pt]{\fontsize{5.8}{6.4}\selectfont \textcolor{ForestGreen}{$ \blacktriangle\,2.5\times $}}\\[-.2pt]1m\,39.6s\\[-.2pt]{\fontsize{5.8}{6.4}\selectfont \textcolor{ForestGreen}{$ \blacktriangle\,2.9\times $}}\end{tabular} & \hspace*{1pt}\begin{tabular}[t]{@{}c@{}}\rule{0pt}{12pt}5m\,28.5s\\[-.2pt]{\fontsize{5.8}{6.4}\selectfont \textcolor{BrickRed}{$ \blacktriangledown\,9.4\times $}}\\[-.2pt]1m\,29.3s\\[-.2pt]{\fontsize{5.8}{6.4}\selectfont \textcolor{BrickRed}{$ \blacktriangledown\,2.6\times $}}\end{tabular} & \hspace*{1pt}\begin{tabular}[t]{@{}c@{}}\rule{0pt}{12pt}4m\,45.1s\\[-.2pt]{\fontsize{5.8}{6.4}\selectfont \textcolor{BrickRed}{$ \blacktriangledown\,7.9\times $}}\\[-.2pt]1m\,22.8s\\[-.2pt]{\fontsize{5.8}{6.4}\selectfont \textcolor{BrickRed}{$ \blacktriangledown\,2.4\times $}}\end{tabular} & \textcolor{gray}{\ensuremath{\cdot}} & \hspace*{1pt}\begin{tabular}[t]{@{}c@{}}\rule{0pt}{12pt}10m\,46.5s\\[-.2pt]{\fontsize{5.8}{6.4}\selectfont \textcolor{BrickRed}{$ \blacktriangledown\,10.8\times $}}\\[-.2pt]1m\,52.9s\\[-.2pt]{\fontsize{5.8}{6.4}\selectfont \textcolor{BrickRed}{$ \blacktriangledown\,1.9\times $}}\end{tabular}\\[3pt]\hline
\end{tabular}\par}
\smallskip{\footnotesize\noindent Each cell: A and B times, $\cdot$: no measurement. ---\textsuperscript{g}: GPU unavailable. n/a: no comparison.\par
\noindent Blue: one lowest A and one lowest B measurement for each column for each configuration.\par
\noindent Gray: minimap2. Yellow: my preferred MPF / Bounded / SIMD-DP / Block configuration.\par
\noindent MPF/MPS: minimizers per fragment/sequence. SIMD-ref: naive SIMD scan. Bounded: bounded SIMD algorithm.\par
}
\clearpage
\noindent{\large\bfseries Table \thetable: (E): AWS: 2 Gbp, k=20, w=15\hfill 2/7}\par\smallskip
{\fontsize{7.1}{8.1}\selectfont\setlength{\tabcolsep}{1pt}\setlength{\fboxsep}{2pt}\renewcommand{\arraystretch}{1.12}
\noindent\begin{tabular}{@{}p{14mm}p{27mm}|p{13.55mm}|p{13.55mm}|p{13.55mm}|p{13.55mm}|p{13.55mm}|p{13.55mm}|p{13.55mm}|p{13.55mm}|p{13.55mm}|p{13.55mm}|p{13.55mm}|p{13.55mm}|p{13.55mm}|p{13.55mm}|p{13.55mm}|p{13.55mm}|@{}}\hline
Engine & Execution & \multicolumn{8}{c|}{\textbf{Index}} & \multicolumn{8}{c|}{\textbf{No index}}\\
 &  & \multicolumn{8}{c|}{\textbf{2 Gbp}} & \multicolumn{8}{c|}{\textbf{2 Gbp}}\\
 &  & \multicolumn{8}{c|}{\textbf{k=20, w=15}} & \multicolumn{8}{c|}{\textbf{k=20, w=15}}\\
 &  & \multicolumn{4}{c|}{\textbf{DNA4}} & \multicolumn{4}{c|}{\textbf{DNA5}} & \multicolumn{4}{c|}{\textbf{DNA4}} & \multicolumn{4}{c|}{\textbf{DNA5}}\\
 &  & \multicolumn{1}{c|}{\textbf{M7A}} & \multicolumn{1}{c|}{\textbf{M7I}} & \multicolumn{1}{c|}{\textbf{M4}} & \multicolumn{1}{c|}{\textbf{C7G}} & \multicolumn{1}{c|}{\textbf{M7A}} & \multicolumn{1}{c|}{\textbf{M7I}} & \multicolumn{1}{c|}{\textbf{M4}} & \multicolumn{1}{c|}{\textbf{C7G}} & \multicolumn{1}{c|}{\textbf{M7A}} & \multicolumn{1}{c|}{\textbf{M7I}} & \multicolumn{1}{c|}{\textbf{M4}} & \multicolumn{1}{c|}{\textbf{C7G}} & \multicolumn{1}{c|}{\textbf{M7A}} & \multicolumn{1}{c|}{\textbf{M7I}} & \multicolumn{1}{c|}{\textbf{M4}} & \multicolumn{1}{c|}{\textbf{C7G}}\\
 & & {\color{gray}C1} & {\color{gray}C2} & {\color{gray}C3} & {\color{gray}C4} & {\color{gray}C5} & {\color{gray}C6} & {\color{gray}C7} & {\color{gray}C8} & {\color{gray}C9} & {\color{gray}C10} & {\color{gray}C11} & {\color{gray}C12} & {\color{gray}C13} & {\color{gray}C14} & {\color{gray}C15} & {\color{gray}C16}\\
\hline
\rowcolor{BenchBaseline}[0.5pt][0.5pt]\shortstack[l]{{\color{gray}R1}\\\textbf{minimap2}} & \textbf{Baseline} & \hspace*{1pt}\begin{tabular}[t]{@{}c@{}}\rule{0pt}{12pt}19.6s\\[-.2pt]{\fontsize{5.8}{6.4}\selectfont \textcolor{BrickRed}{$ \blacktriangledown\,6.7\times $}}\\[-.2pt]18.6s\\[-.2pt]{\fontsize{5.8}{6.4}\selectfont \textcolor{BrickRed}{$ \blacktriangledown\,8.4\times $}}\end{tabular} & \hspace*{1pt}\begin{tabular}[t]{@{}c@{}}\rule{0pt}{12pt}18.2s\\[-.2pt]{\fontsize{5.8}{6.4}\selectfont \textcolor{BrickRed}{$ \blacktriangledown\,6.9\times $}}\\[-.2pt]16.5s\\[-.2pt]{\fontsize{5.8}{6.4}\selectfont \textcolor{BrickRed}{$ \blacktriangledown\,8.6\times $}}\end{tabular} & \hspace*{1pt}\begin{tabular}[t]{@{}c@{}}\rule{0pt}{12pt}42.5s\\[-.2pt]{\fontsize{5.8}{6.4}\selectfont \textcolor{BrickRed}{$ \blacktriangledown\,8.2\times $}}\\[-.2pt]39.4s\\[-.2pt]{\fontsize{5.8}{6.4}\selectfont \textcolor{BrickRed}{$ \blacktriangledown\,12.6\times $}}\end{tabular} & \hspace*{1pt}\begin{tabular}[t]{@{}c@{}}\rule{0pt}{12pt}6m\,29.2s\\[-.2pt]{\fontsize{5.8}{6.4}\selectfont \textcolor{BrickRed}{$ \blacktriangledown\,92.8\times $}}\\[-.2pt]5m\,56.9s\\[-.2pt]{\fontsize{5.8}{6.4}\selectfont \textcolor{BrickRed}{$ \blacktriangledown\,184.9\times $}}\end{tabular} & \hspace*{1pt}\begin{tabular}[t]{@{}c@{}}\rule{0pt}{12pt}9.3s\\[-.2pt]{\fontsize{5.8}{6.4}\selectfont \textcolor{BrickRed}{$ \blacktriangledown\,24\% $}}\\[-.2pt]8.4s\\[-.2pt]{\fontsize{5.8}{6.4}\selectfont \textcolor{BrickRed}{$ \blacktriangledown\,1.6\times $}}\end{tabular} & \hspace*{1pt}\begin{tabular}[t]{@{}c@{}}\rule{0pt}{12pt}8.7s\\[-.2pt]{\fontsize{5.8}{6.4}\selectfont \textcolor{BrickRed}{$ \blacktriangledown\,1.5\times $}}\\[-.2pt]7s\\[-.2pt]{\fontsize{5.8}{6.4}\selectfont \textcolor{BrickRed}{$ \blacktriangledown\,1.7\times $}}\end{tabular} & \hspace*{1pt}\begin{tabular}[t]{@{}c@{}}\rule{0pt}{12pt}15.8s\\[-.2pt]{\fontsize{5.8}{6.4}\selectfont \textcolor{BrickRed}{$ \blacktriangledown\,26\% $}}\\[-.2pt]12.7s\\[-.2pt]{\fontsize{5.8}{6.4}\selectfont \textcolor{BrickRed}{$ \blacktriangledown\,1.7\times $}}\end{tabular} & \hspace*{1pt}\begin{tabular}[t]{@{}c@{}}\rule{0pt}{12pt}18.5s\\[-.2pt]{\fontsize{5.8}{6.4}\selectfont \textcolor{BrickRed}{$ \blacktriangledown\,1.6\times $}}\\[-.2pt]17s\\[-.2pt]{\fontsize{5.8}{6.4}\selectfont \textcolor{BrickRed}{$ \blacktriangledown\,1.8\times $}}\end{tabular} & \hspace*{1pt}\begin{tabular}[t]{@{}c@{}}\rule{0pt}{12pt}42.4s\\[-.2pt]{\fontsize{5.8}{6.4}\selectfont \textcolor{ForestGreen}{$ \blacktriangle\,2.0\times $}}\\[-.2pt]41.4s\\[-.2pt]{\fontsize{5.8}{6.4}\selectfont \textcolor{ForestGreen}{$ \blacktriangle\,1.9\times $}}\end{tabular} & \hspace*{1pt}\begin{tabular}[t]{@{}c@{}}\rule{0pt}{12pt}42.9s\\[-.2pt]{\fontsize{5.8}{6.4}\selectfont \textcolor{ForestGreen}{$ \blacktriangle\,2.0\times $}}\\[-.2pt]41.1s\\[-.2pt]{\fontsize{5.8}{6.4}\selectfont \textcolor{ForestGreen}{$ \blacktriangle\,1.8\times $}}\end{tabular} & \hspace*{1pt}\begin{tabular}[t]{@{}c@{}}\rule{0pt}{12pt}1m\,26.9s\\[-.2pt]{\fontsize{5.8}{6.4}\selectfont \textcolor{ForestGreen}{$ \blacktriangle\,1.9\times $}}\\[-.2pt]1m\,23.7s\\[-.2pt]{\fontsize{5.8}{6.4}\selectfont \textcolor{ForestGreen}{$ \blacktriangle\,1.7\times $}}\end{tabular} & \hspace*{1pt}\begin{tabular}[t]{@{}c@{}}\rule{0pt}{12pt}4m\,51.6s\\[-.2pt]{\fontsize{5.8}{6.4}\selectfont \textcolor{BrickRed}{$ \blacktriangledown\,2.9\times $}}\\[-.2pt]4m\,46.6s\\[-.2pt]{\fontsize{5.8}{6.4}\selectfont \textcolor{BrickRed}{$ \blacktriangledown\,3.0\times $}}\end{tabular} & \hspace*{1pt}\begin{tabular}[t]{@{}c@{}}\rule{0pt}{12pt}35s\\[-.2pt]{\fontsize{5.8}{6.4}\selectfont \textcolor{ForestGreen}{$ \blacktriangle\,8.6\times $}}\\[-.2pt]34s\\[-.2pt]{\fontsize{5.8}{6.4}\selectfont \textcolor{ForestGreen}{$ \blacktriangle\,2.3\times $}}\end{tabular} & \hspace*{1pt}\begin{tabular}[t]{@{}c@{}}\rule{0pt}{12pt}35.9s\\[-.2pt]{\fontsize{5.8}{6.4}\selectfont \textcolor{ForestGreen}{$ \blacktriangle\,6.7\times $}}\\[-.2pt]34.2s\\[-.2pt]{\fontsize{5.8}{6.4}\selectfont \textcolor{ForestGreen}{$ \blacktriangle\,2.1\times $}}\end{tabular} & \hspace*{1pt}\begin{tabular}[t]{@{}c@{}}\rule{0pt}{12pt}1m\,14.1s\\[-.2pt]{\fontsize{5.8}{6.4}\selectfont \textcolor{ForestGreen}{$ \blacktriangle\,6.9\times $}}\\[-.2pt]1m\,11.1s\\[-.2pt]{\fontsize{5.8}{6.4}\selectfont \textcolor{ForestGreen}{$ \blacktriangle\,2.2\times $}}\end{tabular} & \hspace*{1pt}\begin{tabular}[t]{@{}c@{}}\rule{0pt}{12pt}1m\,0.1s\\[-.2pt]{\fontsize{5.8}{6.4}\selectfont \textcolor{ForestGreen}{$ \blacktriangle\,10.4\times $}}\\[-.2pt]58.6s\\[-.2pt]{\fontsize{5.8}{6.4}\selectfont \textcolor{ForestGreen}{$ \blacktriangle\,1.8\times $}}\end{tabular}\\[3pt]\hline
\BenchmarkSourceLabel{bench:aws-2-k20-w15:row:mps-bounded-rmq-block}{R9}\shortstack[l]{{\color{gray}R9}\\\textbf{BioseqDB}} & \begin{tabular}[t]{@{}l@{}}\rule{0pt}{12pt}MPS / Bounded\\RMQ / Block\end{tabular} & \hspace*{1pt}\begin{tabular}[t]{@{}c@{}}\rule{0pt}{12pt}6.6s\\[-.2pt]{\fontsize{5.8}{6.4}\selectfont \textcolor{ForestGreen}{$ \blacktriangle\,3.0\times $}}\\[-.2pt]5.9s\\[-.2pt]{\fontsize{5.8}{6.4}\selectfont \textcolor{ForestGreen}{$ \blacktriangle\,3.2\times $}}\end{tabular} & \hspace*{1pt}\begin{tabular}[t]{@{}c@{}}\rule{0pt}{12pt}6.5s\\[-.2pt]{\fontsize{5.8}{6.4}\selectfont \textcolor{ForestGreen}{$ \blacktriangle\,2.8\times $}}\\[-.2pt]5.8s\\[-.2pt]{\fontsize{5.8}{6.4}\selectfont \textcolor{ForestGreen}{$ \blacktriangle\,2.8\times $}}\end{tabular} & \hspace*{1pt}\begin{tabular}[t]{@{}c@{}}\rule{0pt}{12pt}12.6s\\[-.2pt]{\fontsize{5.8}{6.4}\selectfont \textcolor{ForestGreen}{$ \blacktriangle\,3.4\times $}}\\[-.2pt]10.3s\\[-.2pt]{\fontsize{5.8}{6.4}\selectfont \textcolor{ForestGreen}{$ \blacktriangle\,3.8\times $}}\end{tabular} & \hspace*{1pt}\begin{tabular}[t]{@{}c@{}}\rule{0pt}{12pt}6.4s\\[-.2pt]{\fontsize{5.8}{6.4}\selectfont \textcolor{ForestGreen}{$ \blacktriangle\,61.0\times $}}\\[-.2pt]5.5s\\[-.2pt]{\fontsize{5.8}{6.4}\selectfont \textcolor{ForestGreen}{$ \blacktriangle\,65.2\times $}}\end{tabular} & \hspace*{1pt}\begin{tabular}[t]{@{}c@{}}\rule{0pt}{12pt}4m\,37.8s\\[-.2pt]{\fontsize{5.8}{6.4}\selectfont \textcolor{BrickRed}{$ \blacktriangledown\,29.8\times $}}\\[-.2pt]20.1s\\[-.2pt]{\fontsize{5.8}{6.4}\selectfont \textcolor{BrickRed}{$ \blacktriangledown\,2.4\times $}}\end{tabular} & \hspace*{1pt}\begin{tabular}[t]{@{}c@{}}\rule{0pt}{12pt}3m\,55.4s\\[-.2pt]{\fontsize{5.8}{6.4}\selectfont \textcolor{BrickRed}{$ \blacktriangledown\,26.9\times $}}\\[-.2pt]19.8s\\[-.2pt]{\fontsize{5.8}{6.4}\selectfont \textcolor{BrickRed}{$ \blacktriangledown\,2.8\times $}}\end{tabular} & \textcolor{gray}{\ensuremath{\cdot}} & \textcolor{gray}{\ensuremath{\cdot}} & \hspace*{1pt}\begin{tabular}[t]{@{}c@{}}\rule{0pt}{12pt}1m\,43.2s\\[-.2pt]{\fontsize{5.8}{6.4}\selectfont \textcolor{BrickRed}{$ \blacktriangledown\,2.4\times $}}\\[-.2pt]1m\,21s\\[-.2pt]{\fontsize{5.8}{6.4}\selectfont \textcolor{BrickRed}{$ \blacktriangledown\,2.0\times $}}\end{tabular} & \hspace*{1pt}\begin{tabular}[t]{@{}c@{}}\rule{0pt}{12pt}1m\,41.8s\\[-.2pt]{\fontsize{5.8}{6.4}\selectfont \textcolor{BrickRed}{$ \blacktriangledown\,2.4\times $}}\\[-.2pt]1m\,21.9s\\[-.2pt]{\fontsize{5.8}{6.4}\selectfont \textcolor{BrickRed}{$ \blacktriangledown\,2.0\times $}}\end{tabular} & \hspace*{1pt}\begin{tabular}[t]{@{}c@{}}\rule{0pt}{12pt}2m\,55s\\[-.2pt]{\fontsize{5.8}{6.4}\selectfont \textcolor{BrickRed}{$ \blacktriangledown\,2.0\times $}}\\[-.2pt]2m\,21s\\[-.2pt]{\fontsize{5.8}{6.4}\selectfont \textcolor{BrickRed}{$ \blacktriangledown\,1.7\times $}}\end{tabular} & \hspace*{1pt}\begin{tabular}[t]{@{}c@{}}\rule{0pt}{12pt}1m\,57.1s\\[-.2pt]{\fontsize{5.8}{6.4}\selectfont \textcolor{ForestGreen}{$ \blacktriangle\,2.5\times $}}\\[-.2pt]1m\,39.5s\\[-.2pt]{\fontsize{5.8}{6.4}\selectfont \textcolor{ForestGreen}{$ \blacktriangle\,2.9\times $}}\end{tabular} & \hspace*{1pt}\begin{tabular}[t]{@{}c@{}}\rule{0pt}{12pt}5m\,28.1s\\[-.2pt]{\fontsize{5.8}{6.4}\selectfont \textcolor{BrickRed}{$ \blacktriangledown\,9.4\times $}}\\[-.2pt]1m\,28.8s\\[-.2pt]{\fontsize{5.8}{6.4}\selectfont \textcolor{BrickRed}{$ \blacktriangledown\,2.6\times $}}\end{tabular} & \hspace*{1pt}\begin{tabular}[t]{@{}c@{}}\rule{0pt}{12pt}4m\,44.5s\\[-.2pt]{\fontsize{5.8}{6.4}\selectfont \textcolor{BrickRed}{$ \blacktriangledown\,7.9\times $}}\\[-.2pt]1m\,22.4s\\[-.2pt]{\fontsize{5.8}{6.4}\selectfont \textcolor{BrickRed}{$ \blacktriangledown\,2.4\times $}}\end{tabular} & \textcolor{gray}{\ensuremath{\cdot}} & \hspace*{1pt}\begin{tabular}[t]{@{}c@{}}\rule{0pt}{12pt}10m\,46.5s\\[-.2pt]{\fontsize{5.8}{6.4}\selectfont \textcolor{BrickRed}{$ \blacktriangledown\,10.8\times $}}\\[-.2pt]1m\,51.9s\\[-.2pt]{\fontsize{5.8}{6.4}\selectfont \textcolor{BrickRed}{$ \blacktriangledown\,1.9\times $}}\end{tabular}\\[3pt]\hline
\BenchmarkSourceLabel{bench:aws-2-k20-w15:row:mps-bounded-simd-block}{R10}\shortstack[l]{{\color{gray}R10}} & \begin{tabular}[t]{@{}l@{}}\rule{0pt}{12pt}MPS / Bounded\\SIMD-DP / Block\end{tabular} & \hspace*{1pt}\begin{tabular}[t]{@{}c@{}}\rule{0pt}{12pt}6.5s\\[-.2pt]{\fontsize{5.8}{6.4}\selectfont \textcolor{ForestGreen}{$ \blacktriangle\,3.0\times $}}\\[-.2pt]5.8s\\[-.2pt]{\fontsize{5.8}{6.4}\selectfont \textcolor{ForestGreen}{$ \blacktriangle\,3.2\times $}}\end{tabular} & \hspace*{1pt}\begin{tabular}[t]{@{}c@{}}\rule{0pt}{12pt}6.5s\\[-.2pt]{\fontsize{5.8}{6.4}\selectfont \textcolor{ForestGreen}{$ \blacktriangle\,2.8\times $}}\\[-.2pt]5.8s\\[-.2pt]{\fontsize{5.8}{6.4}\selectfont \textcolor{ForestGreen}{$ \blacktriangle\,2.8\times $}}\end{tabular} & \hspace*{1pt}\begin{tabular}[t]{@{}c@{}}\rule{0pt}{12pt}16.8s\\[-.2pt]{\fontsize{5.8}{6.4}\selectfont \textcolor{ForestGreen}{$ \blacktriangle\,2.5\times $}}\\[-.2pt]14.2s\\[-.2pt]{\fontsize{5.8}{6.4}\selectfont \textcolor{ForestGreen}{$ \blacktriangle\,2.8\times $}}\end{tabular} & \hspace*{1pt}\begin{tabular}[t]{@{}c@{}}\rule{0pt}{12pt}6.4s\\[-.2pt]{\fontsize{5.8}{6.4}\selectfont \textcolor{ForestGreen}{$ \blacktriangle\,60.5\times $}}\\[-.2pt]5.5s\\[-.2pt]{\fontsize{5.8}{6.4}\selectfont \textcolor{ForestGreen}{$ \blacktriangle\,64.7\times $}}\end{tabular} & \hspace*{1pt}\begin{tabular}[t]{@{}c@{}}\rule{0pt}{12pt}4m\,39s\\[-.2pt]{\fontsize{5.8}{6.4}\selectfont \textcolor{BrickRed}{$ \blacktriangledown\,29.9\times $}}\\[-.2pt]20.9s\\[-.2pt]{\fontsize{5.8}{6.4}\selectfont \textcolor{BrickRed}{$ \blacktriangledown\,2.5\times $}}\end{tabular} & \hspace*{1pt}\begin{tabular}[t]{@{}c@{}}\rule{0pt}{12pt}3m\,55.7s\\[-.2pt]{\fontsize{5.8}{6.4}\selectfont \textcolor{BrickRed}{$ \blacktriangledown\,27.0\times $}}\\[-.2pt]20.1s\\[-.2pt]{\fontsize{5.8}{6.4}\selectfont \textcolor{BrickRed}{$ \blacktriangledown\,2.9\times $}}\end{tabular} & \hspace*{1pt}\begin{tabular}[t]{@{}c@{}}\rule{0pt}{12pt}7m\,43.9s\\[-.2pt]{\fontsize{5.8}{6.4}\selectfont \textcolor{BrickRed}{$ \blacktriangledown\,29.3\times $}}\\[-.2pt]35.6s\\[-.2pt]{\fontsize{5.8}{6.4}\selectfont \textcolor{BrickRed}{$ \blacktriangledown\,2.8\times $}}\end{tabular} & \textcolor{gray}{\ensuremath{\cdot}} & \hspace*{1pt}\begin{tabular}[t]{@{}c@{}}\rule{0pt}{12pt}1m\,54.2s\\[-.2pt]{\fontsize{5.8}{6.4}\selectfont \textcolor{BrickRed}{$ \blacktriangledown\,2.7\times $}}\\[-.2pt]1m\,29.4s\\[-.2pt]{\fontsize{5.8}{6.4}\selectfont \textcolor{BrickRed}{$ \blacktriangledown\,2.2\times $}}\end{tabular} & \hspace*{1pt}\begin{tabular}[t]{@{}c@{}}\rule{0pt}{12pt}1m\,40.7s\\[-.2pt]{\fontsize{5.8}{6.4}\selectfont \textcolor{BrickRed}{$ \blacktriangledown\,2.4\times $}}\\[-.2pt]1m\,20.8s\\[-.2pt]{\fontsize{5.8}{6.4}\selectfont \textcolor{BrickRed}{$ \blacktriangledown\,2.0\times $}}\end{tabular} & \textcolor{gray}{\ensuremath{\cdot}} & \hspace*{1pt}\begin{tabular}[t]{@{}c@{}}\rule{0pt}{12pt}1m\,51.5s\\[-.2pt]{\fontsize{5.8}{6.4}\selectfont \textcolor{ForestGreen}{$ \blacktriangle\,2.6\times $}}\\[-.2pt]\colorbox{BenchBest}{\strut 1m\,35.2s}\\[-.2pt]{\fontsize{5.8}{6.4}\selectfont \textcolor{ForestGreen}{$ \blacktriangle\,3.0\times $}}\end{tabular} & \hspace*{1pt}\begin{tabular}[t]{@{}c@{}}\rule{0pt}{12pt}5m\,38.6s\\[-.2pt]{\fontsize{5.8}{6.4}\selectfont \textcolor{BrickRed}{$ \blacktriangledown\,9.7\times $}}\\[-.2pt]1m\,39.5s\\[-.2pt]{\fontsize{5.8}{6.4}\selectfont \textcolor{BrickRed}{$ \blacktriangledown\,2.9\times $}}\end{tabular} & \hspace*{1pt}\begin{tabular}[t]{@{}c@{}}\rule{0pt}{12pt}4m\,46.5s\\[-.2pt]{\fontsize{5.8}{6.4}\selectfont \textcolor{BrickRed}{$ \blacktriangledown\,8.0\times $}}\\[-.2pt]1m\,23.5s\\[-.2pt]{\fontsize{5.8}{6.4}\selectfont \textcolor{BrickRed}{$ \blacktriangledown\,2.4\times $}}\end{tabular} & \hspace*{1pt}\begin{tabular}[t]{@{}c@{}}\rule{0pt}{12pt}8m\,50.5s\\[-.2pt]{\fontsize{5.8}{6.4}\selectfont \textcolor{BrickRed}{$ \blacktriangledown\,7.2\times $}}\\[-.2pt]2m\,35s\\[-.2pt]{\fontsize{5.8}{6.4}\selectfont \textcolor{BrickRed}{$ \blacktriangledown\,2.2\times $}}\end{tabular} & \hspace*{1pt}\begin{tabular}[t]{@{}c@{}}\rule{0pt}{12pt}10m\,48.8s\\[-.2pt]{\fontsize{5.8}{6.4}\selectfont \textcolor{BrickRed}{$ \blacktriangledown\,10.8\times $}}\\[-.2pt]1m\,54.4s\\[-.2pt]{\fontsize{5.8}{6.4}\selectfont \textcolor{BrickRed}{$ \blacktriangledown\,2.0\times $}}\end{tabular}\\[3pt]\hline
\BenchmarkSourceLabel{bench:aws-2-k20-w15:row:mps-bounded-rmq-ksw2}{R11}\shortstack[l]{{\color{gray}R11}} & \begin{tabular}[t]{@{}l@{}}\rule{0pt}{12pt}MPS / Bounded\\RMQ / KSW2\end{tabular} & \hspace*{1pt}\begin{tabular}[t]{@{}c@{}}\rule{0pt}{12pt}8.2s\\[-.2pt]{\fontsize{5.8}{6.4}\selectfont \textcolor{ForestGreen}{$ \blacktriangle\,2.4\times $}}\\[-.2pt]5.7s\\[-.2pt]{\fontsize{5.8}{6.4}\selectfont \textcolor{ForestGreen}{$ \blacktriangle\,3.3\times $}}\end{tabular} & \hspace*{1pt}\begin{tabular}[t]{@{}c@{}}\rule{0pt}{12pt}8s\\[-.2pt]{\fontsize{5.8}{6.4}\selectfont \textcolor{ForestGreen}{$ \blacktriangle\,2.3\times $}}\\[-.2pt]5.6s\\[-.2pt]{\fontsize{5.8}{6.4}\selectfont \textcolor{ForestGreen}{$ \blacktriangle\,2.9\times $}}\end{tabular} & \hspace*{1pt}\begin{tabular}[t]{@{}c@{}}\rule{0pt}{12pt}17.4s\\[-.2pt]{\fontsize{5.8}{6.4}\selectfont \textcolor{ForestGreen}{$ \blacktriangle\,2.4\times $}}\\[-.2pt]10.1s\\[-.2pt]{\fontsize{5.8}{6.4}\selectfont \textcolor{ForestGreen}{$ \blacktriangle\,3.9\times $}}\end{tabular} & \hspace*{1pt}\begin{tabular}[t]{@{}c@{}}\rule{0pt}{12pt}8.3s\\[-.2pt]{\fontsize{5.8}{6.4}\selectfont \textcolor{ForestGreen}{$ \blacktriangle\,46.6\times $}}\\[-.2pt]5.3s\\[-.2pt]{\fontsize{5.8}{6.4}\selectfont \textcolor{ForestGreen}{$ \blacktriangle\,67.2\times $}}\end{tabular} & \hspace*{1pt}\begin{tabular}[t]{@{}c@{}}\rule{0pt}{12pt}4m\,39.5s\\[-.2pt]{\fontsize{5.8}{6.4}\selectfont \textcolor{BrickRed}{$ \blacktriangledown\,30.0\times $}}\\[-.2pt]20.2s\\[-.2pt]{\fontsize{5.8}{6.4}\selectfont \textcolor{BrickRed}{$ \blacktriangledown\,2.4\times $}}\end{tabular} & \hspace*{1pt}\begin{tabular}[t]{@{}c@{}}\rule{0pt}{12pt}3m\,57.2s\\[-.2pt]{\fontsize{5.8}{6.4}\selectfont \textcolor{BrickRed}{$ \blacktriangledown\,27.2\times $}}\\[-.2pt]19.8s\\[-.2pt]{\fontsize{5.8}{6.4}\selectfont \textcolor{BrickRed}{$ \blacktriangledown\,2.8\times $}}\end{tabular} & \textcolor{gray}{\ensuremath{\cdot}} & \textcolor{gray}{\ensuremath{\cdot}} & \hspace*{1pt}\begin{tabular}[t]{@{}c@{}}\rule{0pt}{12pt}1m\,48.4s\\[-.2pt]{\fontsize{5.8}{6.4}\selectfont \textcolor{BrickRed}{$ \blacktriangledown\,2.6\times $}}\\[-.2pt]1m\,24.7s\\[-.2pt]{\fontsize{5.8}{6.4}\selectfont \textcolor{BrickRed}{$ \blacktriangledown\,2.0\times $}}\end{tabular} & \hspace*{1pt}\begin{tabular}[t]{@{}c@{}}\rule{0pt}{12pt}1m\,39.1s\\[-.2pt]{\fontsize{5.8}{6.4}\selectfont \textcolor{BrickRed}{$ \blacktriangledown\,2.3\times $}}\\[-.2pt]1m\,18s\\[-.2pt]{\fontsize{5.8}{6.4}\selectfont \textcolor{BrickRed}{$ \blacktriangledown\,1.9\times $}}\end{tabular} & \textcolor{gray}{\ensuremath{\cdot}} & \hspace*{1pt}\begin{tabular}[t]{@{}c@{}}\rule{0pt}{12pt}2m\,0.6s\\[-.2pt]{\fontsize{5.8}{6.4}\selectfont \textcolor{ForestGreen}{$ \blacktriangle\,2.4\times $}}\\[-.2pt]1m\,39.5s\\[-.2pt]{\fontsize{5.8}{6.4}\selectfont \textcolor{ForestGreen}{$ \blacktriangle\,2.9\times $}}\end{tabular} & \hspace*{1pt}\begin{tabular}[t]{@{}c@{}}\rule{0pt}{12pt}5m\,38.2s\\[-.2pt]{\fontsize{5.8}{6.4}\selectfont \textcolor{BrickRed}{$ \blacktriangledown\,9.7\times $}}\\[-.2pt]1m\,37.1s\\[-.2pt]{\fontsize{5.8}{6.4}\selectfont \textcolor{BrickRed}{$ \blacktriangledown\,2.9\times $}}\end{tabular} & \hspace*{1pt}\begin{tabular}[t]{@{}c@{}}\rule{0pt}{12pt}4m\,45.8s\\[-.2pt]{\fontsize{5.8}{6.4}\selectfont \textcolor{BrickRed}{$ \blacktriangledown\,8.0\times $}}\\[-.2pt]1m\,22.4s\\[-.2pt]{\fontsize{5.8}{6.4}\selectfont \textcolor{BrickRed}{$ \blacktriangledown\,2.4\times $}}\end{tabular} & \textcolor{gray}{\ensuremath{\cdot}} & \hspace*{1pt}\begin{tabular}[t]{@{}c@{}}\rule{0pt}{12pt}10m\,55.2s\\[-.2pt]{\fontsize{5.8}{6.4}\selectfont \textcolor{BrickRed}{$ \blacktriangledown\,10.9\times $}}\\[-.2pt]1m\,53.8s\\[-.2pt]{\fontsize{5.8}{6.4}\selectfont \textcolor{BrickRed}{$ \blacktriangledown\,1.9\times $}}\end{tabular}\\[3pt]\hline
\BenchmarkSourceLabel{bench:aws-2-k20-w15:row:mps-bounded-simd-ksw2}{R12}\shortstack[l]{{\color{gray}R12}} & \begin{tabular}[t]{@{}l@{}}\rule{0pt}{12pt}MPS / Bounded\\SIMD-DP / KSW2\end{tabular} & \hspace*{1pt}\begin{tabular}[t]{@{}c@{}}\rule{0pt}{12pt}8.4s\\[-.2pt]{\fontsize{5.8}{6.4}\selectfont \textcolor{ForestGreen}{$ \blacktriangle\,2.3\times $}}\\[-.2pt]5.7s\\[-.2pt]{\fontsize{5.8}{6.4}\selectfont \textcolor{ForestGreen}{$ \blacktriangle\,3.3\times $}}\end{tabular} & \hspace*{1pt}\begin{tabular}[t]{@{}c@{}}\rule{0pt}{12pt}8s\\[-.2pt]{\fontsize{5.8}{6.4}\selectfont \textcolor{ForestGreen}{$ \blacktriangle\,2.3\times $}}\\[-.2pt]5.6s\\[-.2pt]{\fontsize{5.8}{6.4}\selectfont \textcolor{ForestGreen}{$ \blacktriangle\,2.9\times $}}\end{tabular} & \hspace*{1pt}\begin{tabular}[t]{@{}c@{}}\rule{0pt}{12pt}17.5s\\[-.2pt]{\fontsize{5.8}{6.4}\selectfont \textcolor{ForestGreen}{$ \blacktriangle\,2.4\times $}}\\[-.2pt]10.1s\\[-.2pt]{\fontsize{5.8}{6.4}\selectfont \textcolor{ForestGreen}{$ \blacktriangle\,3.9\times $}}\end{tabular} & \hspace*{1pt}\begin{tabular}[t]{@{}c@{}}\rule{0pt}{12pt}8.2s\\[-.2pt]{\fontsize{5.8}{6.4}\selectfont \textcolor{ForestGreen}{$ \blacktriangle\,47.3\times $}}\\[-.2pt]5.3s\\[-.2pt]{\fontsize{5.8}{6.4}\selectfont \textcolor{ForestGreen}{$ \blacktriangle\,67.7\times $}}\end{tabular} & \hspace*{1pt}\begin{tabular}[t]{@{}c@{}}\rule{0pt}{12pt}4m\,39.1s\\[-.2pt]{\fontsize{5.8}{6.4}\selectfont \textcolor{BrickRed}{$ \blacktriangledown\,29.9\times $}}\\[-.2pt]20.1s\\[-.2pt]{\fontsize{5.8}{6.4}\selectfont \textcolor{BrickRed}{$ \blacktriangledown\,2.4\times $}}\end{tabular} & \hspace*{1pt}\begin{tabular}[t]{@{}c@{}}\rule{0pt}{12pt}3m\,56.9s\\[-.2pt]{\fontsize{5.8}{6.4}\selectfont \textcolor{BrickRed}{$ \blacktriangledown\,27.1\times $}}\\[-.2pt]19.4s\\[-.2pt]{\fontsize{5.8}{6.4}\selectfont \textcolor{BrickRed}{$ \blacktriangledown\,2.8\times $}}\end{tabular} & \hspace*{1pt}\begin{tabular}[t]{@{}c@{}}\rule{0pt}{12pt}7m\,26s\\[-.2pt]{\fontsize{5.8}{6.4}\selectfont \textcolor{BrickRed}{$ \blacktriangledown\,28.2\times $}}\\[-.2pt]34.2s\\[-.2pt]{\fontsize{5.8}{6.4}\selectfont \textcolor{BrickRed}{$ \blacktriangledown\,2.7\times $}}\end{tabular} & \textcolor{gray}{\ensuremath{\cdot}} & \hspace*{1pt}\begin{tabular}[t]{@{}c@{}}\rule{0pt}{12pt}1m\,43.1s\\[-.2pt]{\fontsize{5.8}{6.4}\selectfont \textcolor{BrickRed}{$ \blacktriangledown\,2.4\times $}}\\[-.2pt]1m\,19.2s\\[-.2pt]{\fontsize{5.8}{6.4}\selectfont \textcolor{BrickRed}{$ \blacktriangledown\,1.9\times $}}\end{tabular} & \hspace*{1pt}\begin{tabular}[t]{@{}c@{}}\rule{0pt}{12pt}1m\,39.2s\\[-.2pt]{\fontsize{5.8}{6.4}\selectfont \textcolor{BrickRed}{$ \blacktriangledown\,2.3\times $}}\\[-.2pt]1m\,17.9s\\[-.2pt]{\fontsize{5.8}{6.4}\selectfont \textcolor{BrickRed}{$ \blacktriangledown\,1.9\times $}}\end{tabular} & \textcolor{gray}{\ensuremath{\cdot}} & \hspace*{1pt}\begin{tabular}[t]{@{}c@{}}\rule{0pt}{12pt}1m\,59.9s\\[-.2pt]{\fontsize{5.8}{6.4}\selectfont \textcolor{ForestGreen}{$ \blacktriangle\,2.4\times $}}\\[-.2pt]1m\,38.2s\\[-.2pt]{\fontsize{5.8}{6.4}\selectfont \textcolor{ForestGreen}{$ \blacktriangle\,2.9\times $}}\end{tabular} & \hspace*{1pt}\begin{tabular}[t]{@{}c@{}}\rule{0pt}{12pt}5m\,30.8s\\[-.2pt]{\fontsize{5.8}{6.4}\selectfont \textcolor{BrickRed}{$ \blacktriangledown\,9.5\times $}}\\[-.2pt]1m\,29.4s\\[-.2pt]{\fontsize{5.8}{6.4}\selectfont \textcolor{BrickRed}{$ \blacktriangledown\,2.6\times $}}\end{tabular} & \hspace*{1pt}\begin{tabular}[t]{@{}c@{}}\rule{0pt}{12pt}4m\,45.4s\\[-.2pt]{\fontsize{5.8}{6.4}\selectfont \textcolor{BrickRed}{$ \blacktriangledown\,7.9\times $}}\\[-.2pt]1m\,22.1s\\[-.2pt]{\fontsize{5.8}{6.4}\selectfont \textcolor{BrickRed}{$ \blacktriangledown\,2.4\times $}}\end{tabular} & \textcolor{gray}{\ensuremath{\cdot}} & \hspace*{1pt}\begin{tabular}[t]{@{}c@{}}\rule{0pt}{12pt}10m\,48.3s\\[-.2pt]{\fontsize{5.8}{6.4}\selectfont \textcolor{BrickRed}{$ \blacktriangledown\,10.8\times $}}\\[-.2pt]1m\,48.9s\\[-.2pt]{\fontsize{5.8}{6.4}\selectfont \textcolor{BrickRed}{$ \blacktriangledown\,1.9\times $}}\end{tabular}\\[3pt]\hline
\BenchmarkSourceLabel{bench:aws-2-k20-w15:row:mps-bounded-scalar-ksw2}{R13}\shortstack[l]{{\color{gray}R13}} & \begin{tabular}[t]{@{}l@{}}\rule{0pt}{12pt}MPS / Bounded\\Scalar-DP / KSW2\end{tabular} & \hspace*{1pt}\begin{tabular}[t]{@{}c@{}}\rule{0pt}{12pt}8.2s\\[-.2pt]{\fontsize{5.8}{6.4}\selectfont \textcolor{ForestGreen}{$ \blacktriangle\,2.4\times $}}\\[-.2pt]5.7s\\[-.2pt]{\fontsize{5.8}{6.4}\selectfont \textcolor{ForestGreen}{$ \blacktriangle\,3.3\times $}}\end{tabular} & \hspace*{1pt}\begin{tabular}[t]{@{}c@{}}\rule{0pt}{12pt}8.1s\\[-.2pt]{\fontsize{5.8}{6.4}\selectfont \textcolor{ForestGreen}{$ \blacktriangle\,2.3\times $}}\\[-.2pt]5.6s\\[-.2pt]{\fontsize{5.8}{6.4}\selectfont \textcolor{ForestGreen}{$ \blacktriangle\,2.9\times $}}\end{tabular} & \hspace*{1pt}\begin{tabular}[t]{@{}c@{}}\rule{0pt}{12pt}17.7s\\[-.2pt]{\fontsize{5.8}{6.4}\selectfont \textcolor{ForestGreen}{$ \blacktriangle\,2.4\times $}}\\[-.2pt]10.1s\\[-.2pt]{\fontsize{5.8}{6.4}\selectfont \textcolor{ForestGreen}{$ \blacktriangle\,3.9\times $}}\end{tabular} & \hspace*{1pt}\begin{tabular}[t]{@{}c@{}}\rule{0pt}{12pt}8.4s\\[-.2pt]{\fontsize{5.8}{6.4}\selectfont \textcolor{ForestGreen}{$ \blacktriangle\,46.5\times $}}\\[-.2pt]5.3s\\[-.2pt]{\fontsize{5.8}{6.4}\selectfont \textcolor{ForestGreen}{$ \blacktriangle\,66.8\times $}}\end{tabular} & \hspace*{1pt}\begin{tabular}[t]{@{}c@{}}\rule{0pt}{12pt}4m\,39.7s\\[-.2pt]{\fontsize{5.8}{6.4}\selectfont \textcolor{BrickRed}{$ \blacktriangledown\,30.0\times $}}\\[-.2pt]20.4s\\[-.2pt]{\fontsize{5.8}{6.4}\selectfont \textcolor{BrickRed}{$ \blacktriangledown\,2.4\times $}}\end{tabular} & \hspace*{1pt}\begin{tabular}[t]{@{}c@{}}\rule{0pt}{12pt}3m\,56.6s\\[-.2pt]{\fontsize{5.8}{6.4}\selectfont \textcolor{BrickRed}{$ \blacktriangledown\,27.1\times $}}\\[-.2pt]19.5s\\[-.2pt]{\fontsize{5.8}{6.4}\selectfont \textcolor{BrickRed}{$ \blacktriangledown\,2.8\times $}}\end{tabular} & \textcolor{gray}{\ensuremath{\cdot}} & \textcolor{gray}{\ensuremath{\cdot}} & \hspace*{1pt}\begin{tabular}[t]{@{}c@{}}\rule{0pt}{12pt}1m\,50.8s\\[-.2pt]{\fontsize{5.8}{6.4}\selectfont \textcolor{BrickRed}{$ \blacktriangledown\,2.6\times $}}\\[-.2pt]1m\,27.1s\\[-.2pt]{\fontsize{5.8}{6.4}\selectfont \textcolor{BrickRed}{$ \blacktriangledown\,2.1\times $}}\end{tabular} & \hspace*{1pt}\begin{tabular}[t]{@{}c@{}}\rule{0pt}{12pt}1m\,43.6s\\[-.2pt]{\fontsize{5.8}{6.4}\selectfont \textcolor{BrickRed}{$ \blacktriangledown\,2.4\times $}}\\[-.2pt]1m\,21.8s\\[-.2pt]{\fontsize{5.8}{6.4}\selectfont \textcolor{BrickRed}{$ \blacktriangledown\,2.0\times $}}\end{tabular} & \textcolor{gray}{\ensuremath{\cdot}} & \hspace*{1pt}\begin{tabular}[t]{@{}c@{}}\rule{0pt}{12pt}2m\,0.5s\\[-.2pt]{\fontsize{5.8}{6.4}\selectfont \textcolor{ForestGreen}{$ \blacktriangle\,2.4\times $}}\\[-.2pt]1m\,38.7s\\[-.2pt]{\fontsize{5.8}{6.4}\selectfont \textcolor{ForestGreen}{$ \blacktriangle\,2.9\times $}}\end{tabular} & \hspace*{1pt}\begin{tabular}[t]{@{}c@{}}\rule{0pt}{12pt}5m\,29.2s\\[-.2pt]{\fontsize{5.8}{6.4}\selectfont \textcolor{BrickRed}{$ \blacktriangledown\,9.4\times $}}\\[-.2pt]1m\,28s\\[-.2pt]{\fontsize{5.8}{6.4}\selectfont \textcolor{BrickRed}{$ \blacktriangledown\,2.6\times $}}\end{tabular} & \textcolor{gray}{\ensuremath{\cdot}} & \textcolor{gray}{\ensuremath{\cdot}} & \hspace*{1pt}\begin{tabular}[t]{@{}c@{}}\rule{0pt}{12pt}10m\,47.5s\\[-.2pt]{\fontsize{5.8}{6.4}\selectfont \textcolor{BrickRed}{$ \blacktriangledown\,10.8\times $}}\\[-.2pt]1m\,51.8s\\[-.2pt]{\fontsize{5.8}{6.4}\selectfont \textcolor{BrickRed}{$ \blacktriangledown\,1.9\times $}}\end{tabular}\\[3pt]\hline
\BenchmarkSourceLabel{bench:aws-2-k20-w15:row:mps-twostack-simd-block}{R14}\shortstack[l]{{\color{gray}R14}} & \begin{tabular}[t]{@{}l@{}}\rule{0pt}{12pt}MPS / Two-stack\\SIMD-DP / Block\end{tabular} & \hspace*{1pt}\begin{tabular}[t]{@{}c@{}}\rule{0pt}{12pt}39.5s\\[-.2pt]{\fontsize{5.8}{6.4}\selectfont \textcolor{BrickRed}{$ \blacktriangledown\,2.0\times $}}\\[-.2pt]6.8s\\[-.2pt]{\fontsize{5.8}{6.4}\selectfont \textcolor{ForestGreen}{$ \blacktriangle\,2.7\times $}}\end{tabular} & \hspace*{1pt}\begin{tabular}[t]{@{}c@{}}\rule{0pt}{12pt}35.7s\\[-.2pt]{\fontsize{5.8}{6.4}\selectfont \textcolor{BrickRed}{$ \blacktriangledown\,2.0\times $}}\\[-.2pt]7.2s\\[-.2pt]{\fontsize{5.8}{6.4}\selectfont \textcolor{ForestGreen}{$ \blacktriangle\,2.3\times $}}\end{tabular} & \textcolor{gray}{\ensuremath{\cdot}} & \hspace*{1pt}\begin{tabular}[t]{@{}c@{}}\rule{0pt}{12pt}33.9s\\[-.2pt]{\fontsize{5.8}{6.4}\selectfont \textcolor{ForestGreen}{$ \blacktriangle\,11.5\times $}}\\[-.2pt]6.9s\\[-.2pt]{\fontsize{5.8}{6.4}\selectfont \textcolor{ForestGreen}{$ \blacktriangle\,51.8\times $}}\end{tabular} & \hspace*{1pt}\begin{tabular}[t]{@{}c@{}}\rule{0pt}{12pt}4m\,25.9s\\[-.2pt]{\fontsize{5.8}{6.4}\selectfont \textcolor{BrickRed}{$ \blacktriangledown\,28.5\times $}}\\[-.2pt]13.3s\\[-.2pt]{\fontsize{5.8}{6.4}\selectfont \textcolor{BrickRed}{$ \blacktriangledown\,1.6\times $}}\end{tabular} & \hspace*{1pt}\begin{tabular}[t]{@{}c@{}}\rule{0pt}{12pt}3m\,42s\\[-.2pt]{\fontsize{5.8}{6.4}\selectfont \textcolor{BrickRed}{$ \blacktriangledown\,25.4\times $}}\\[-.2pt]12s\\[-.2pt]{\fontsize{5.8}{6.4}\selectfont \textcolor{BrickRed}{$ \blacktriangledown\,1.7\times $}}\end{tabular} & \textcolor{gray}{\ensuremath{\cdot}} & \textcolor{gray}{\ensuremath{\cdot}} & \textcolor{gray}{\ensuremath{\cdot}} & \textcolor{gray}{\ensuremath{\cdot}} & \textcolor{gray}{\ensuremath{\cdot}} & \textcolor{gray}{\ensuremath{\cdot}} & \textcolor{gray}{\ensuremath{\cdot}} & \textcolor{gray}{\ensuremath{\cdot}} & \textcolor{gray}{\ensuremath{\cdot}} & \textcolor{gray}{\ensuremath{\cdot}}\\[3pt]\hline
\BenchmarkSourceLabel{bench:aws-2-k20-w15:row:mps-simdref-simd-block}{R15}\shortstack[l]{{\color{gray}R15}} & \begin{tabular}[t]{@{}l@{}}\rule{0pt}{12pt}MPS / SIMD-ref\\SIMD-DP / Block\end{tabular} & \hspace*{1pt}\begin{tabular}[t]{@{}c@{}}\rule{0pt}{12pt}46s\\[-.2pt]{\fontsize{5.8}{6.4}\selectfont \textcolor{BrickRed}{$ \blacktriangledown\,2.3\times $}}\\[-.2pt]8.9s\\[-.2pt]{\fontsize{5.8}{6.4}\selectfont \textcolor{ForestGreen}{$ \blacktriangle\,2.1\times $}}\end{tabular} & \hspace*{1pt}\begin{tabular}[t]{@{}c@{}}\rule{0pt}{12pt}40.5s\\[-.2pt]{\fontsize{5.8}{6.4}\selectfont \textcolor{BrickRed}{$ \blacktriangledown\,2.2\times $}}\\[-.2pt]8.8s\\[-.2pt]{\fontsize{5.8}{6.4}\selectfont \textcolor{ForestGreen}{$ \blacktriangle\,1.9\times $}}\end{tabular} & \textcolor{gray}{\ensuremath{\cdot}} & \hspace*{1pt}\begin{tabular}[t]{@{}c@{}}\rule{0pt}{12pt}34.8s\\[-.2pt]{\fontsize{5.8}{6.4}\selectfont \textcolor{ForestGreen}{$ \blacktriangle\,11.2\times $}}\\[-.2pt]8.2s\\[-.2pt]{\fontsize{5.8}{6.4}\selectfont \textcolor{ForestGreen}{$ \blacktriangle\,43.6\times $}}\end{tabular} & \hspace*{1pt}\begin{tabular}[t]{@{}c@{}}\rule{0pt}{12pt}4m\,48.1s\\[-.2pt]{\fontsize{5.8}{6.4}\selectfont \textcolor{BrickRed}{$ \blacktriangledown\,30.9\times $}}\\[-.2pt]30.2s\\[-.2pt]{\fontsize{5.8}{6.4}\selectfont \textcolor{BrickRed}{$ \blacktriangledown\,3.6\times $}}\end{tabular} & \hspace*{1pt}\begin{tabular}[t]{@{}c@{}}\rule{0pt}{12pt}4m\,3.1s\\[-.2pt]{\fontsize{5.8}{6.4}\selectfont \textcolor{BrickRed}{$ \blacktriangledown\,27.8\times $}}\\[-.2pt]28.3s\\[-.2pt]{\fontsize{5.8}{6.4}\selectfont \textcolor{BrickRed}{$ \blacktriangledown\,4.0\times $}}\end{tabular} & \textcolor{gray}{\ensuremath{\cdot}} & \textcolor{gray}{\ensuremath{\cdot}} & \textcolor{gray}{\ensuremath{\cdot}} & \textcolor{gray}{\ensuremath{\cdot}} & \textcolor{gray}{\ensuremath{\cdot}} & \hspace*{1pt}\begin{tabular}[t]{@{}c@{}}\rule{0pt}{12pt}2m\,54.8s\\[-.2pt]{\fontsize{5.8}{6.4}\selectfont \textcolor{ForestGreen}{$ \blacktriangle\,1.7\times $}}\\[-.2pt]2m\,26.4s\\[-.2pt]{\fontsize{5.8}{6.4}\selectfont \textcolor{ForestGreen}{$ \blacktriangle\,2.0\times $}}\end{tabular} & \textcolor{gray}{\ensuremath{\cdot}} & \textcolor{gray}{\ensuremath{\cdot}} & \textcolor{gray}{\ensuremath{\cdot}} & \textcolor{gray}{\ensuremath{\cdot}}\\[3pt]\hline
\end{tabular}\par}
\smallskip{\footnotesize\noindent Each cell: A and B times, $\cdot$: no measurement. ---\textsuperscript{g}: GPU unavailable. n/a: no comparison.\par
\noindent Blue: one lowest A and one lowest B measurement for each column for each configuration.\par
\noindent Gray: minimap2. Yellow: my preferred MPF / Bounded / SIMD-DP / Block configuration.\par
\noindent MPF/MPS: minimizers per fragment/sequence. SIMD-ref: naive SIMD scan. Bounded: bounded SIMD algorithm.\par
}
\clearpage
\noindent{\large\bfseries Table \thetable: (E): AWS: 2 Gbp, k=20, w=15\hfill 3/7}\par\smallskip
{\fontsize{7.1}{8.1}\selectfont\setlength{\tabcolsep}{1pt}\setlength{\fboxsep}{2pt}\renewcommand{\arraystretch}{1.12}
\noindent\begin{tabular}{@{}p{14mm}p{27mm}|p{13.55mm}|p{13.55mm}|p{13.55mm}|p{13.55mm}|p{13.55mm}|p{13.55mm}|p{13.55mm}|p{13.55mm}|p{13.55mm}|p{13.55mm}|p{13.55mm}|p{13.55mm}|p{13.55mm}|p{13.55mm}|p{13.55mm}|p{13.55mm}|@{}}\hline
Engine & Execution & \multicolumn{8}{c|}{\textbf{Index}} & \multicolumn{8}{c|}{\textbf{No index}}\\
 &  & \multicolumn{8}{c|}{\textbf{2 Gbp}} & \multicolumn{8}{c|}{\textbf{2 Gbp}}\\
 &  & \multicolumn{8}{c|}{\textbf{k=20, w=15}} & \multicolumn{8}{c|}{\textbf{k=20, w=15}}\\
 &  & \multicolumn{4}{c|}{\textbf{DNA4}} & \multicolumn{4}{c|}{\textbf{DNA5}} & \multicolumn{4}{c|}{\textbf{DNA4}} & \multicolumn{4}{c|}{\textbf{DNA5}}\\
 &  & \multicolumn{1}{c|}{\textbf{M7A}} & \multicolumn{1}{c|}{\textbf{M7I}} & \multicolumn{1}{c|}{\textbf{M4}} & \multicolumn{1}{c|}{\textbf{C7G}} & \multicolumn{1}{c|}{\textbf{M7A}} & \multicolumn{1}{c|}{\textbf{M7I}} & \multicolumn{1}{c|}{\textbf{M4}} & \multicolumn{1}{c|}{\textbf{C7G}} & \multicolumn{1}{c|}{\textbf{M7A}} & \multicolumn{1}{c|}{\textbf{M7I}} & \multicolumn{1}{c|}{\textbf{M4}} & \multicolumn{1}{c|}{\textbf{C7G}} & \multicolumn{1}{c|}{\textbf{M7A}} & \multicolumn{1}{c|}{\textbf{M7I}} & \multicolumn{1}{c|}{\textbf{M4}} & \multicolumn{1}{c|}{\textbf{C7G}}\\
 & & {\color{gray}C1} & {\color{gray}C2} & {\color{gray}C3} & {\color{gray}C4} & {\color{gray}C5} & {\color{gray}C6} & {\color{gray}C7} & {\color{gray}C8} & {\color{gray}C9} & {\color{gray}C10} & {\color{gray}C11} & {\color{gray}C12} & {\color{gray}C13} & {\color{gray}C14} & {\color{gray}C15} & {\color{gray}C16}\\
\hline
\rowcolor{BenchBaseline}[0.5pt][0.5pt]\shortstack[l]{{\color{gray}R1}\\\textbf{minimap2}} & \textbf{Baseline} & \hspace*{1pt}\begin{tabular}[t]{@{}c@{}}\rule{0pt}{12pt}19.6s\\[-.2pt]{\fontsize{5.8}{6.4}\selectfont \textcolor{BrickRed}{$ \blacktriangledown\,6.7\times $}}\\[-.2pt]18.6s\\[-.2pt]{\fontsize{5.8}{6.4}\selectfont \textcolor{BrickRed}{$ \blacktriangledown\,8.4\times $}}\end{tabular} & \hspace*{1pt}\begin{tabular}[t]{@{}c@{}}\rule{0pt}{12pt}18.2s\\[-.2pt]{\fontsize{5.8}{6.4}\selectfont \textcolor{BrickRed}{$ \blacktriangledown\,6.9\times $}}\\[-.2pt]16.5s\\[-.2pt]{\fontsize{5.8}{6.4}\selectfont \textcolor{BrickRed}{$ \blacktriangledown\,8.6\times $}}\end{tabular} & \hspace*{1pt}\begin{tabular}[t]{@{}c@{}}\rule{0pt}{12pt}42.5s\\[-.2pt]{\fontsize{5.8}{6.4}\selectfont \textcolor{BrickRed}{$ \blacktriangledown\,8.2\times $}}\\[-.2pt]39.4s\\[-.2pt]{\fontsize{5.8}{6.4}\selectfont \textcolor{BrickRed}{$ \blacktriangledown\,12.6\times $}}\end{tabular} & \hspace*{1pt}\begin{tabular}[t]{@{}c@{}}\rule{0pt}{12pt}6m\,29.2s\\[-.2pt]{\fontsize{5.8}{6.4}\selectfont \textcolor{BrickRed}{$ \blacktriangledown\,92.8\times $}}\\[-.2pt]5m\,56.9s\\[-.2pt]{\fontsize{5.8}{6.4}\selectfont \textcolor{BrickRed}{$ \blacktriangledown\,184.9\times $}}\end{tabular} & \hspace*{1pt}\begin{tabular}[t]{@{}c@{}}\rule{0pt}{12pt}9.3s\\[-.2pt]{\fontsize{5.8}{6.4}\selectfont \textcolor{BrickRed}{$ \blacktriangledown\,24\% $}}\\[-.2pt]8.4s\\[-.2pt]{\fontsize{5.8}{6.4}\selectfont \textcolor{BrickRed}{$ \blacktriangledown\,1.6\times $}}\end{tabular} & \hspace*{1pt}\begin{tabular}[t]{@{}c@{}}\rule{0pt}{12pt}8.7s\\[-.2pt]{\fontsize{5.8}{6.4}\selectfont \textcolor{BrickRed}{$ \blacktriangledown\,1.5\times $}}\\[-.2pt]7s\\[-.2pt]{\fontsize{5.8}{6.4}\selectfont \textcolor{BrickRed}{$ \blacktriangledown\,1.7\times $}}\end{tabular} & \hspace*{1pt}\begin{tabular}[t]{@{}c@{}}\rule{0pt}{12pt}15.8s\\[-.2pt]{\fontsize{5.8}{6.4}\selectfont \textcolor{BrickRed}{$ \blacktriangledown\,26\% $}}\\[-.2pt]12.7s\\[-.2pt]{\fontsize{5.8}{6.4}\selectfont \textcolor{BrickRed}{$ \blacktriangledown\,1.7\times $}}\end{tabular} & \hspace*{1pt}\begin{tabular}[t]{@{}c@{}}\rule{0pt}{12pt}18.5s\\[-.2pt]{\fontsize{5.8}{6.4}\selectfont \textcolor{BrickRed}{$ \blacktriangledown\,1.6\times $}}\\[-.2pt]17s\\[-.2pt]{\fontsize{5.8}{6.4}\selectfont \textcolor{BrickRed}{$ \blacktriangledown\,1.8\times $}}\end{tabular} & \hspace*{1pt}\begin{tabular}[t]{@{}c@{}}\rule{0pt}{12pt}42.4s\\[-.2pt]{\fontsize{5.8}{6.4}\selectfont \textcolor{ForestGreen}{$ \blacktriangle\,2.0\times $}}\\[-.2pt]41.4s\\[-.2pt]{\fontsize{5.8}{6.4}\selectfont \textcolor{ForestGreen}{$ \blacktriangle\,1.9\times $}}\end{tabular} & \hspace*{1pt}\begin{tabular}[t]{@{}c@{}}\rule{0pt}{12pt}42.9s\\[-.2pt]{\fontsize{5.8}{6.4}\selectfont \textcolor{ForestGreen}{$ \blacktriangle\,2.0\times $}}\\[-.2pt]41.1s\\[-.2pt]{\fontsize{5.8}{6.4}\selectfont \textcolor{ForestGreen}{$ \blacktriangle\,1.8\times $}}\end{tabular} & \hspace*{1pt}\begin{tabular}[t]{@{}c@{}}\rule{0pt}{12pt}1m\,26.9s\\[-.2pt]{\fontsize{5.8}{6.4}\selectfont \textcolor{ForestGreen}{$ \blacktriangle\,1.9\times $}}\\[-.2pt]1m\,23.7s\\[-.2pt]{\fontsize{5.8}{6.4}\selectfont \textcolor{ForestGreen}{$ \blacktriangle\,1.7\times $}}\end{tabular} & \hspace*{1pt}\begin{tabular}[t]{@{}c@{}}\rule{0pt}{12pt}4m\,51.6s\\[-.2pt]{\fontsize{5.8}{6.4}\selectfont \textcolor{BrickRed}{$ \blacktriangledown\,2.9\times $}}\\[-.2pt]4m\,46.6s\\[-.2pt]{\fontsize{5.8}{6.4}\selectfont \textcolor{BrickRed}{$ \blacktriangledown\,3.0\times $}}\end{tabular} & \hspace*{1pt}\begin{tabular}[t]{@{}c@{}}\rule{0pt}{12pt}35s\\[-.2pt]{\fontsize{5.8}{6.4}\selectfont \textcolor{ForestGreen}{$ \blacktriangle\,8.6\times $}}\\[-.2pt]34s\\[-.2pt]{\fontsize{5.8}{6.4}\selectfont \textcolor{ForestGreen}{$ \blacktriangle\,2.3\times $}}\end{tabular} & \hspace*{1pt}\begin{tabular}[t]{@{}c@{}}\rule{0pt}{12pt}35.9s\\[-.2pt]{\fontsize{5.8}{6.4}\selectfont \textcolor{ForestGreen}{$ \blacktriangle\,6.7\times $}}\\[-.2pt]34.2s\\[-.2pt]{\fontsize{5.8}{6.4}\selectfont \textcolor{ForestGreen}{$ \blacktriangle\,2.1\times $}}\end{tabular} & \hspace*{1pt}\begin{tabular}[t]{@{}c@{}}\rule{0pt}{12pt}1m\,14.1s\\[-.2pt]{\fontsize{5.8}{6.4}\selectfont \textcolor{ForestGreen}{$ \blacktriangle\,6.9\times $}}\\[-.2pt]1m\,11.1s\\[-.2pt]{\fontsize{5.8}{6.4}\selectfont \textcolor{ForestGreen}{$ \blacktriangle\,2.2\times $}}\end{tabular} & \hspace*{1pt}\begin{tabular}[t]{@{}c@{}}\rule{0pt}{12pt}1m\,0.1s\\[-.2pt]{\fontsize{5.8}{6.4}\selectfont \textcolor{ForestGreen}{$ \blacktriangle\,10.4\times $}}\\[-.2pt]58.6s\\[-.2pt]{\fontsize{5.8}{6.4}\selectfont \textcolor{ForestGreen}{$ \blacktriangle\,1.8\times $}}\end{tabular}\\[3pt]\hline
\BenchmarkSourceLabel{bench:aws-2-k20-w15:row:mps-twostack-rmq-block}{R16}\shortstack[l]{{\color{gray}R16}\\\textbf{BioseqDB}} & \begin{tabular}[t]{@{}l@{}}\rule{0pt}{12pt}MPS / Two-stack\\RMQ / Block\end{tabular} & \hspace*{1pt}\begin{tabular}[t]{@{}c@{}}\rule{0pt}{12pt}39.9s\\[-.2pt]{\fontsize{5.8}{6.4}\selectfont \textcolor{BrickRed}{$ \blacktriangledown\,2.0\times $}}\\[-.2pt]6.8s\\[-.2pt]{\fontsize{5.8}{6.4}\selectfont \textcolor{ForestGreen}{$ \blacktriangle\,2.7\times $}}\end{tabular} & \hspace*{1pt}\begin{tabular}[t]{@{}c@{}}\rule{0pt}{12pt}35.5s\\[-.2pt]{\fontsize{5.8}{6.4}\selectfont \textcolor{BrickRed}{$ \blacktriangledown\,1.9\times $}}\\[-.2pt]7s\\[-.2pt]{\fontsize{5.8}{6.4}\selectfont \textcolor{ForestGreen}{$ \blacktriangle\,2.4\times $}}\end{tabular} & \textcolor{gray}{\ensuremath{\cdot}} & \textcolor{gray}{\ensuremath{\cdot}} & \hspace*{1pt}\begin{tabular}[t]{@{}c@{}}\rule{0pt}{12pt}4m\,24.9s\\[-.2pt]{\fontsize{5.8}{6.4}\selectfont \textcolor{BrickRed}{$ \blacktriangledown\,28.4\times $}}\\[-.2pt]12.5s\\[-.2pt]{\fontsize{5.8}{6.4}\selectfont \textcolor{BrickRed}{$ \blacktriangledown\,1.5\times $}}\end{tabular} & \hspace*{1pt}\begin{tabular}[t]{@{}c@{}}\rule{0pt}{12pt}3m\,41.4s\\[-.2pt]{\fontsize{5.8}{6.4}\selectfont \textcolor{BrickRed}{$ \blacktriangledown\,25.3\times $}}\\[-.2pt]11.5s\\[-.2pt]{\fontsize{5.8}{6.4}\selectfont \textcolor{BrickRed}{$ \blacktriangledown\,1.6\times $}}\end{tabular} & \textcolor{gray}{\ensuremath{\cdot}} & \textcolor{gray}{\ensuremath{\cdot}} & \textcolor{gray}{\ensuremath{\cdot}} & \textcolor{gray}{\ensuremath{\cdot}} & \textcolor{gray}{\ensuremath{\cdot}} & \textcolor{gray}{\ensuremath{\cdot}} & \textcolor{gray}{\ensuremath{\cdot}} & \textcolor{gray}{\ensuremath{\cdot}} & \textcolor{gray}{\ensuremath{\cdot}} & \textcolor{gray}{\ensuremath{\cdot}}\\[3pt]\hline
\BenchmarkSourceLabel{bench:aws-2-k20-w15:row:mps-twostack-scalar-block}{R17}\shortstack[l]{{\color{gray}R17}} & \begin{tabular}[t]{@{}l@{}}\rule{0pt}{12pt}MPS / Two-stack\\Scalar-DP / Block\end{tabular} & \hspace*{1pt}\begin{tabular}[t]{@{}c@{}}\rule{0pt}{12pt}39.8s\\[-.2pt]{\fontsize{5.8}{6.4}\selectfont \textcolor{BrickRed}{$ \blacktriangledown\,2.0\times $}}\\[-.2pt]6.8s\\[-.2pt]{\fontsize{5.8}{6.4}\selectfont \textcolor{ForestGreen}{$ \blacktriangle\,2.7\times $}}\end{tabular} & \hspace*{1pt}\begin{tabular}[t]{@{}c@{}}\rule{0pt}{12pt}35.8s\\[-.2pt]{\fontsize{5.8}{6.4}\selectfont \textcolor{BrickRed}{$ \blacktriangledown\,2.0\times $}}\\[-.2pt]6.5s\\[-.2pt]{\fontsize{5.8}{6.4}\selectfont \textcolor{ForestGreen}{$ \blacktriangle\,2.5\times $}}\end{tabular} & \textcolor{gray}{\ensuremath{\cdot}} & \textcolor{gray}{\ensuremath{\cdot}} & \hspace*{1pt}\begin{tabular}[t]{@{}c@{}}\rule{0pt}{12pt}4m\,25.7s\\[-.2pt]{\fontsize{5.8}{6.4}\selectfont \textcolor{BrickRed}{$ \blacktriangledown\,28.5\times $}}\\[-.2pt]12.7s\\[-.2pt]{\fontsize{5.8}{6.4}\selectfont \textcolor{BrickRed}{$ \blacktriangledown\,1.5\times $}}\end{tabular} & \hspace*{1pt}\begin{tabular}[t]{@{}c@{}}\rule{0pt}{12pt}3m\,41.4s\\[-.2pt]{\fontsize{5.8}{6.4}\selectfont \textcolor{BrickRed}{$ \blacktriangledown\,25.3\times $}}\\[-.2pt]11.5s\\[-.2pt]{\fontsize{5.8}{6.4}\selectfont \textcolor{BrickRed}{$ \blacktriangledown\,1.6\times $}}\end{tabular} & \textcolor{gray}{\ensuremath{\cdot}} & \textcolor{gray}{\ensuremath{\cdot}} & \textcolor{gray}{\ensuremath{\cdot}} & \textcolor{gray}{\ensuremath{\cdot}} & \textcolor{gray}{\ensuremath{\cdot}} & \textcolor{gray}{\ensuremath{\cdot}} & \textcolor{gray}{\ensuremath{\cdot}} & \textcolor{gray}{\ensuremath{\cdot}} & \textcolor{gray}{\ensuremath{\cdot}} & \textcolor{gray}{\ensuremath{\cdot}}\\[3pt]\hline
\BenchmarkSourceLabel{bench:aws-2-k20-w15:row:mps-scalar-simd-block}{R18}\shortstack[l]{{\color{gray}R18}} & \begin{tabular}[t]{@{}l@{}}\rule{0pt}{12pt}MPS / Scalar\\SIMD-DP / Block\end{tabular} & \hspace*{1pt}\begin{tabular}[t]{@{}c@{}}\rule{0pt}{12pt}43.8s\\[-.2pt]{\fontsize{5.8}{6.4}\selectfont \textcolor{BrickRed}{$ \blacktriangledown\,2.2\times $}}\\[-.2pt]8.6s\\[-.2pt]{\fontsize{5.8}{6.4}\selectfont \textcolor{ForestGreen}{$ \blacktriangle\,2.2\times $}}\end{tabular} & \hspace*{1pt}\begin{tabular}[t]{@{}c@{}}\rule{0pt}{12pt}38s\\[-.2pt]{\fontsize{5.8}{6.4}\selectfont \textcolor{BrickRed}{$ \blacktriangledown\,2.1\times $}}\\[-.2pt]8.7s\\[-.2pt]{\fontsize{5.8}{6.4}\selectfont \textcolor{ForestGreen}{$ \blacktriangle\,1.9\times $}}\end{tabular} & \textcolor{gray}{\ensuremath{\cdot}} & \hspace*{1pt}\begin{tabular}[t]{@{}c@{}}\rule{0pt}{12pt}36.6s\\[-.2pt]{\fontsize{5.8}{6.4}\selectfont \textcolor{ForestGreen}{$ \blacktriangle\,10.6\times $}}\\[-.2pt]8.2s\\[-.2pt]{\fontsize{5.8}{6.4}\selectfont \textcolor{ForestGreen}{$ \blacktriangle\,43.4\times $}}\end{tabular} & \hspace*{1pt}\begin{tabular}[t]{@{}c@{}}\rule{0pt}{12pt}4m\,40.3s\\[-.2pt]{\fontsize{5.8}{6.4}\selectfont \textcolor{BrickRed}{$ \blacktriangledown\,30.1\times $}}\\[-.2pt]23.1s\\[-.2pt]{\fontsize{5.8}{6.4}\selectfont \textcolor{BrickRed}{$ \blacktriangledown\,2.8\times $}}\end{tabular} & \hspace*{1pt}\begin{tabular}[t]{@{}c@{}}\rule{0pt}{12pt}4m\,0.1s\\[-.2pt]{\fontsize{5.8}{6.4}\selectfont \textcolor{BrickRed}{$ \blacktriangledown\,27.5\times $}}\\[-.2pt]22.3s\\[-.2pt]{\fontsize{5.8}{6.4}\selectfont \textcolor{BrickRed}{$ \blacktriangledown\,3.2\times $}}\end{tabular} & \textcolor{gray}{\ensuremath{\cdot}} & \textcolor{gray}{\ensuremath{\cdot}} & \hspace*{1pt}\begin{tabular}[t]{@{}c@{}}\rule{0pt}{12pt}2m\,52.9s\\[-.2pt]{\fontsize{5.8}{6.4}\selectfont \textcolor{BrickRed}{$ \blacktriangledown\,4.1\times $}}\\[-.2pt]2m\,16.2s\\[-.2pt]{\fontsize{5.8}{6.4}\selectfont \textcolor{BrickRed}{$ \blacktriangledown\,3.3\times $}}\end{tabular} & \hspace*{1pt}\begin{tabular}[t]{@{}c@{}}\rule{0pt}{12pt}2m\,34.3s\\[-.2pt]{\fontsize{5.8}{6.4}\selectfont \textcolor{BrickRed}{$ \blacktriangledown\,3.6\times $}}\\[-.2pt]2m\,5.6s\\[-.2pt]{\fontsize{5.8}{6.4}\selectfont \textcolor{BrickRed}{$ \blacktriangledown\,3.1\times $}}\end{tabular} & \textcolor{gray}{\ensuremath{\cdot}} & \hspace*{1pt}\begin{tabular}[t]{@{}c@{}}\rule{0pt}{12pt}2m\,43.8s\\[-.2pt]{\fontsize{5.8}{6.4}\selectfont \textcolor{ForestGreen}{$ \blacktriangle\,1.8\times $}}\\[-.2pt]2m\,16.9s\\[-.2pt]{\fontsize{5.8}{6.4}\selectfont \textcolor{ForestGreen}{$ \blacktriangle\,2.1\times $}}\end{tabular} & \hspace*{1pt}\begin{tabular}[t]{@{}c@{}}\rule{0pt}{12pt}6m\,48.5s\\[-.2pt]{\fontsize{5.8}{6.4}\selectfont \textcolor{BrickRed}{$ \blacktriangledown\,11.7\times $}}\\[-.2pt]2m\,32s\\[-.2pt]{\fontsize{5.8}{6.4}\selectfont \textcolor{BrickRed}{$ \blacktriangledown\,4.5\times $}}\end{tabular} & \textcolor{gray}{\ensuremath{\cdot}} & \textcolor{gray}{\ensuremath{\cdot}} & \textcolor{gray}{\ensuremath{\cdot}}\\[3pt]\hline
\BenchmarkSourceLabel{bench:aws-2-k20-w15:row:mps-scalar-rmq-block}{R19}\shortstack[l]{{\color{gray}R19}} & \begin{tabular}[t]{@{}l@{}}\rule{0pt}{12pt}MPS / Scalar\\RMQ / Block\end{tabular} & \hspace*{1pt}\begin{tabular}[t]{@{}c@{}}\rule{0pt}{12pt}43.6s\\[-.2pt]{\fontsize{5.8}{6.4}\selectfont \textcolor{BrickRed}{$ \blacktriangledown\,2.2\times $}}\\[-.2pt]8.8s\\[-.2pt]{\fontsize{5.8}{6.4}\selectfont \textcolor{ForestGreen}{$ \blacktriangle\,2.1\times $}}\end{tabular} & \hspace*{1pt}\begin{tabular}[t]{@{}c@{}}\rule{0pt}{12pt}37.4s\\[-.2pt]{\fontsize{5.8}{6.4}\selectfont \textcolor{BrickRed}{$ \blacktriangledown\,2.1\times $}}\\[-.2pt]8.8s\\[-.2pt]{\fontsize{5.8}{6.4}\selectfont \textcolor{ForestGreen}{$ \blacktriangle\,1.9\times $}}\end{tabular} & \textcolor{gray}{\ensuremath{\cdot}} & \textcolor{gray}{\ensuremath{\cdot}} & \hspace*{1pt}\begin{tabular}[t]{@{}c@{}}\rule{0pt}{12pt}4m\,39.7s\\[-.2pt]{\fontsize{5.8}{6.4}\selectfont \textcolor{BrickRed}{$ \blacktriangledown\,30.0\times $}}\\[-.2pt]22s\\[-.2pt]{\fontsize{5.8}{6.4}\selectfont \textcolor{BrickRed}{$ \blacktriangledown\,2.6\times $}}\end{tabular} & \hspace*{1pt}\begin{tabular}[t]{@{}c@{}}\rule{0pt}{12pt}3m\,57.3s\\[-.2pt]{\fontsize{5.8}{6.4}\selectfont \textcolor{BrickRed}{$ \blacktriangledown\,27.2\times $}}\\[-.2pt]21.1s\\[-.2pt]{\fontsize{5.8}{6.4}\selectfont \textcolor{BrickRed}{$ \blacktriangledown\,3.0\times $}}\end{tabular} & \textcolor{gray}{\ensuremath{\cdot}} & \textcolor{gray}{\ensuremath{\cdot}} & \hspace*{1pt}\begin{tabular}[t]{@{}c@{}}\rule{0pt}{12pt}2m\,53.3s\\[-.2pt]{\fontsize{5.8}{6.4}\selectfont \textcolor{BrickRed}{$ \blacktriangledown\,4.1\times $}}\\[-.2pt]2m\,17.7s\\[-.2pt]{\fontsize{5.8}{6.4}\selectfont \textcolor{BrickRed}{$ \blacktriangledown\,3.3\times $}}\end{tabular} & \hspace*{1pt}\begin{tabular}[t]{@{}c@{}}\rule{0pt}{12pt}2m\,48s\\[-.2pt]{\fontsize{5.8}{6.4}\selectfont \textcolor{BrickRed}{$ \blacktriangledown\,3.9\times $}}\\[-.2pt]2m\,18.9s\\[-.2pt]{\fontsize{5.8}{6.4}\selectfont \textcolor{BrickRed}{$ \blacktriangledown\,3.4\times $}}\end{tabular} & \hspace*{1pt}\begin{tabular}[t]{@{}c@{}}\rule{0pt}{12pt}3m\,24s\\[-.2pt]{\fontsize{5.8}{6.4}\selectfont \textcolor{BrickRed}{$ \blacktriangledown\,2.3\times $}}\\[-.2pt]2m\,47.8s\\[-.2pt]{\fontsize{5.8}{6.4}\selectfont \textcolor{BrickRed}{$ \blacktriangledown\,2.0\times $}}\end{tabular} & \textcolor{gray}{\ensuremath{\cdot}} & \hspace*{1pt}\begin{tabular}[t]{@{}c@{}}\rule{0pt}{12pt}6m\,48.4s\\[-.2pt]{\fontsize{5.8}{6.4}\selectfont \textcolor{BrickRed}{$ \blacktriangledown\,11.7\times $}}\\[-.2pt]2m\,33s\\[-.2pt]{\fontsize{5.8}{6.4}\selectfont \textcolor{BrickRed}{$ \blacktriangledown\,4.5\times $}}\end{tabular} & \textcolor{gray}{\ensuremath{\cdot}} & \textcolor{gray}{\ensuremath{\cdot}} & \textcolor{gray}{\ensuremath{\cdot}}\\[3pt]\hline
\BenchmarkSourceLabel{bench:aws-2-k20-w15:row:mps-scalar-scalar-block}{R20}\shortstack[l]{{\color{gray}R20}} & \begin{tabular}[t]{@{}l@{}}\rule{0pt}{12pt}MPS / Scalar\\Scalar-DP / Block\end{tabular} & \hspace*{1pt}\begin{tabular}[t]{@{}c@{}}\rule{0pt}{12pt}43.3s\\[-.2pt]{\fontsize{5.8}{6.4}\selectfont \textcolor{BrickRed}{$ \blacktriangledown\,2.2\times $}}\\[-.2pt]8.7s\\[-.2pt]{\fontsize{5.8}{6.4}\selectfont \textcolor{ForestGreen}{$ \blacktriangle\,2.2\times $}}\end{tabular} & \hspace*{1pt}\begin{tabular}[t]{@{}c@{}}\rule{0pt}{12pt}37.5s\\[-.2pt]{\fontsize{5.8}{6.4}\selectfont \textcolor{BrickRed}{$ \blacktriangledown\,2.1\times $}}\\[-.2pt]8.9s\\[-.2pt]{\fontsize{5.8}{6.4}\selectfont \textcolor{ForestGreen}{$ \blacktriangle\,1.9\times $}}\end{tabular} & \textcolor{gray}{\ensuremath{\cdot}} & \textcolor{gray}{\ensuremath{\cdot}} & \hspace*{1pt}\begin{tabular}[t]{@{}c@{}}\rule{0pt}{12pt}4m\,39.8s\\[-.2pt]{\fontsize{5.8}{6.4}\selectfont \textcolor{BrickRed}{$ \blacktriangledown\,30.0\times $}}\\[-.2pt]22s\\[-.2pt]{\fontsize{5.8}{6.4}\selectfont \textcolor{BrickRed}{$ \blacktriangledown\,2.6\times $}}\end{tabular} & \hspace*{1pt}\begin{tabular}[t]{@{}c@{}}\rule{0pt}{12pt}3m\,57.2s\\[-.2pt]{\fontsize{5.8}{6.4}\selectfont \textcolor{BrickRed}{$ \blacktriangledown\,27.1\times $}}\\[-.2pt]20.8s\\[-.2pt]{\fontsize{5.8}{6.4}\selectfont \textcolor{BrickRed}{$ \blacktriangledown\,3.0\times $}}\end{tabular} & \textcolor{gray}{\ensuremath{\cdot}} & \textcolor{gray}{\ensuremath{\cdot}} & \hspace*{1pt}\begin{tabular}[t]{@{}c@{}}\rule{0pt}{12pt}2m\,57.9s\\[-.2pt]{\fontsize{5.8}{6.4}\selectfont \textcolor{BrickRed}{$ \blacktriangledown\,4.2\times $}}\\[-.2pt]2m\,19.7s\\[-.2pt]{\fontsize{5.8}{6.4}\selectfont \textcolor{BrickRed}{$ \blacktriangledown\,3.4\times $}}\end{tabular} & \textcolor{gray}{\ensuremath{\cdot}} & \textcolor{gray}{\ensuremath{\cdot}} & \textcolor{gray}{\ensuremath{\cdot}} & \textcolor{gray}{\ensuremath{\cdot}} & \textcolor{gray}{\ensuremath{\cdot}} & \textcolor{gray}{\ensuremath{\cdot}} & \textcolor{gray}{\ensuremath{\cdot}}\\[3pt]\hline
\BenchmarkSourceLabel{bench:aws-2-k20-w15:row:mps-scalar-simd-ksw2}{R21}\shortstack[l]{{\color{gray}R21}} & \begin{tabular}[t]{@{}l@{}}\rule{0pt}{12pt}MPS / Scalar\\SIMD-DP / KSW2\end{tabular} & \hspace*{1pt}\begin{tabular}[t]{@{}c@{}}\rule{0pt}{12pt}45.1s\\[-.2pt]{\fontsize{5.8}{6.4}\selectfont \textcolor{BrickRed}{$ \blacktriangledown\,2.3\times $}}\\[-.2pt]8.4s\\[-.2pt]{\fontsize{5.8}{6.4}\selectfont \textcolor{ForestGreen}{$ \blacktriangle\,2.2\times $}}\end{tabular} & \hspace*{1pt}\begin{tabular}[t]{@{}c@{}}\rule{0pt}{12pt}37.9s\\[-.2pt]{\fontsize{5.8}{6.4}\selectfont \textcolor{BrickRed}{$ \blacktriangledown\,2.1\times $}}\\[-.2pt]8.2s\\[-.2pt]{\fontsize{5.8}{6.4}\selectfont \textcolor{ForestGreen}{$ \blacktriangle\,2.0\times $}}\end{tabular} & \hspace*{1pt}\begin{tabular}[t]{@{}c@{}}\rule{0pt}{12pt}1m\,6.6s\\[-.2pt]{\fontsize{5.8}{6.4}\selectfont \textcolor{BrickRed}{$ \blacktriangledown\,1.6\times $}}\\[-.2pt]14.5s\\[-.2pt]{\fontsize{5.8}{6.4}\selectfont \textcolor{ForestGreen}{$ \blacktriangle\,2.7\times $}}\end{tabular} & \hspace*{1pt}\begin{tabular}[t]{@{}c@{}}\rule{0pt}{12pt}40.5s\\[-.2pt]{\fontsize{5.8}{6.4}\selectfont \textcolor{ForestGreen}{$ \blacktriangle\,9.6\times $}}\\[-.2pt]8.3s\\[-.2pt]{\fontsize{5.8}{6.4}\selectfont \textcolor{ForestGreen}{$ \blacktriangle\,43.2\times $}}\end{tabular} & \hspace*{1pt}\begin{tabular}[t]{@{}c@{}}\rule{0pt}{12pt}4m\,40.4s\\[-.2pt]{\fontsize{5.8}{6.4}\selectfont \textcolor{BrickRed}{$ \blacktriangledown\,30.1\times $}}\\[-.2pt]21.8s\\[-.2pt]{\fontsize{5.8}{6.4}\selectfont \textcolor{BrickRed}{$ \blacktriangledown\,2.6\times $}}\end{tabular} & \hspace*{1pt}\begin{tabular}[t]{@{}c@{}}\rule{0pt}{12pt}3m\,58.8s\\[-.2pt]{\fontsize{5.8}{6.4}\selectfont \textcolor{BrickRed}{$ \blacktriangledown\,27.3\times $}}\\[-.2pt]20.9s\\[-.2pt]{\fontsize{5.8}{6.4}\selectfont \textcolor{BrickRed}{$ \blacktriangledown\,3.0\times $}}\end{tabular} & \textcolor{gray}{\ensuremath{\cdot}} & \textcolor{gray}{\ensuremath{\cdot}} & \hspace*{1pt}\begin{tabular}[t]{@{}c@{}}\rule{0pt}{12pt}2m\,54.8s\\[-.2pt]{\fontsize{5.8}{6.4}\selectfont \textcolor{BrickRed}{$ \blacktriangledown\,4.1\times $}}\\[-.2pt]2m\,18.1s\\[-.2pt]{\fontsize{5.8}{6.4}\selectfont \textcolor{BrickRed}{$ \blacktriangledown\,3.3\times $}}\end{tabular} & \hspace*{1pt}\begin{tabular}[t]{@{}c@{}}\rule{0pt}{12pt}2m\,39.7s\\[-.2pt]{\fontsize{5.8}{6.4}\selectfont \textcolor{BrickRed}{$ \blacktriangledown\,3.7\times $}}\\[-.2pt]2m\,8.9s\\[-.2pt]{\fontsize{5.8}{6.4}\selectfont \textcolor{BrickRed}{$ \blacktriangledown\,3.1\times $}}\end{tabular} & \hspace*{1pt}\begin{tabular}[t]{@{}c@{}}\rule{0pt}{12pt}4m\,27.6s\\[-.2pt]{\fontsize{5.8}{6.4}\selectfont \textcolor{BrickRed}{$ \blacktriangledown\,3.1\times $}}\\[-.2pt]3m\,34.6s\\[-.2pt]{\fontsize{5.8}{6.4}\selectfont \textcolor{BrickRed}{$ \blacktriangledown\,2.6\times $}}\end{tabular} & \textcolor{gray}{\ensuremath{\cdot}} & \hspace*{1pt}\begin{tabular}[t]{@{}c@{}}\rule{0pt}{12pt}6m\,52.2s\\[-.2pt]{\fontsize{5.8}{6.4}\selectfont \textcolor{BrickRed}{$ \blacktriangledown\,11.8\times $}}\\[-.2pt]2m\,34.3s\\[-.2pt]{\fontsize{5.8}{6.4}\selectfont \textcolor{BrickRed}{$ \blacktriangledown\,4.5\times $}}\end{tabular} & \textcolor{gray}{\ensuremath{\cdot}} & \textcolor{gray}{\ensuremath{\cdot}} & \textcolor{gray}{\ensuremath{\cdot}}\\[3pt]\hline
\BenchmarkSourceLabel{bench:aws-2-k20-w15:row:mps-simdref-rmq-block}{R22}\shortstack[l]{{\color{gray}R22}} & \begin{tabular}[t]{@{}l@{}}\rule{0pt}{12pt}MPS / SIMD-ref\\RMQ / Block\end{tabular} & \hspace*{1pt}\begin{tabular}[t]{@{}c@{}}\rule{0pt}{12pt}46.5s\\[-.2pt]{\fontsize{5.8}{6.4}\selectfont \textcolor{BrickRed}{$ \blacktriangledown\,2.4\times $}}\\[-.2pt]9.2s\\[-.2pt]{\fontsize{5.8}{6.4}\selectfont \textcolor{ForestGreen}{$ \blacktriangle\,2.0\times $}}\end{tabular} & \hspace*{1pt}\begin{tabular}[t]{@{}c@{}}\rule{0pt}{12pt}38.3s\\[-.2pt]{\fontsize{5.8}{6.4}\selectfont \textcolor{BrickRed}{$ \blacktriangledown\,2.1\times $}}\\[-.2pt]8.8s\\[-.2pt]{\fontsize{5.8}{6.4}\selectfont \textcolor{ForestGreen}{$ \blacktriangle\,1.9\times $}}\end{tabular} & \textcolor{gray}{\ensuremath{\cdot}} & \textcolor{gray}{\ensuremath{\cdot}} & \hspace*{1pt}\begin{tabular}[t]{@{}c@{}}\rule{0pt}{12pt}4m\,47s\\[-.2pt]{\fontsize{5.8}{6.4}\selectfont \textcolor{BrickRed}{$ \blacktriangledown\,30.8\times $}}\\[-.2pt]28.9s\\[-.2pt]{\fontsize{5.8}{6.4}\selectfont \textcolor{BrickRed}{$ \blacktriangledown\,3.4\times $}}\end{tabular} & \hspace*{1pt}\begin{tabular}[t]{@{}c@{}}\rule{0pt}{12pt}4m\,2.7s\\[-.2pt]{\fontsize{5.8}{6.4}\selectfont \textcolor{BrickRed}{$ \blacktriangledown\,27.8\times $}}\\[-.2pt]28.4s\\[-.2pt]{\fontsize{5.8}{6.4}\selectfont \textcolor{BrickRed}{$ \blacktriangledown\,4.0\times $}}\end{tabular} & \textcolor{gray}{\ensuremath{\cdot}} & \textcolor{gray}{\ensuremath{\cdot}} & \textcolor{gray}{\ensuremath{\cdot}} & \textcolor{gray}{\ensuremath{\cdot}} & \hspace*{1pt}\begin{tabular}[t]{@{}c@{}}\rule{0pt}{12pt}3m\,51.4s\\[-.2pt]{\fontsize{5.8}{6.4}\selectfont \textcolor{BrickRed}{$ \blacktriangledown\,2.7\times $}}\\[-.2pt]3m\,1.9s\\[-.2pt]{\fontsize{5.8}{6.4}\selectfont \textcolor{BrickRed}{$ \blacktriangledown\,2.2\times $}}\end{tabular} & \textcolor{gray}{\ensuremath{\cdot}} & \textcolor{gray}{\ensuremath{\cdot}} & \textcolor{gray}{\ensuremath{\cdot}} & \textcolor{gray}{\ensuremath{\cdot}} & \textcolor{gray}{\ensuremath{\cdot}}\\[3pt]\hline
\end{tabular}\par}
\smallskip{\footnotesize\noindent Each cell: A and B times, $\cdot$: no measurement. ---\textsuperscript{g}: GPU unavailable. n/a: no comparison.\par
\noindent Blue: one lowest A and one lowest B measurement for each column for each configuration.\par
\noindent Gray: minimap2. Yellow: my preferred MPF / Bounded / SIMD-DP / Block configuration.\par
\noindent MPF/MPS: minimizers per fragment/sequence. SIMD-ref: naive SIMD scan. Bounded: bounded SIMD algorithm.\par
}
\clearpage
\noindent{\large\bfseries Table \thetable: (E): AWS: 2 Gbp, k=20, w=15\hfill 4/7}\par\smallskip
{\fontsize{7.1}{8.1}\selectfont\setlength{\tabcolsep}{1pt}\setlength{\fboxsep}{2pt}\renewcommand{\arraystretch}{1.12}
\noindent\begin{tabular}{@{}p{14mm}p{27mm}|p{13.55mm}|p{13.55mm}|p{13.55mm}|p{13.55mm}|p{13.55mm}|p{13.55mm}|p{13.55mm}|p{13.55mm}|p{13.55mm}|p{13.55mm}|p{13.55mm}|p{13.55mm}|p{13.55mm}|p{13.55mm}|p{13.55mm}|p{13.55mm}|@{}}\hline
Engine & Execution & \multicolumn{8}{c|}{\textbf{Index}} & \multicolumn{8}{c|}{\textbf{No index}}\\
 &  & \multicolumn{8}{c|}{\textbf{2 Gbp}} & \multicolumn{8}{c|}{\textbf{2 Gbp}}\\
 &  & \multicolumn{8}{c|}{\textbf{k=20, w=15}} & \multicolumn{8}{c|}{\textbf{k=20, w=15}}\\
 &  & \multicolumn{4}{c|}{\textbf{DNA4}} & \multicolumn{4}{c|}{\textbf{DNA5}} & \multicolumn{4}{c|}{\textbf{DNA4}} & \multicolumn{4}{c|}{\textbf{DNA5}}\\
 &  & \multicolumn{1}{c|}{\textbf{M7A}} & \multicolumn{1}{c|}{\textbf{M7I}} & \multicolumn{1}{c|}{\textbf{M4}} & \multicolumn{1}{c|}{\textbf{C7G}} & \multicolumn{1}{c|}{\textbf{M7A}} & \multicolumn{1}{c|}{\textbf{M7I}} & \multicolumn{1}{c|}{\textbf{M4}} & \multicolumn{1}{c|}{\textbf{C7G}} & \multicolumn{1}{c|}{\textbf{M7A}} & \multicolumn{1}{c|}{\textbf{M7I}} & \multicolumn{1}{c|}{\textbf{M4}} & \multicolumn{1}{c|}{\textbf{C7G}} & \multicolumn{1}{c|}{\textbf{M7A}} & \multicolumn{1}{c|}{\textbf{M7I}} & \multicolumn{1}{c|}{\textbf{M4}} & \multicolumn{1}{c|}{\textbf{C7G}}\\
 & & {\color{gray}C1} & {\color{gray}C2} & {\color{gray}C3} & {\color{gray}C4} & {\color{gray}C5} & {\color{gray}C6} & {\color{gray}C7} & {\color{gray}C8} & {\color{gray}C9} & {\color{gray}C10} & {\color{gray}C11} & {\color{gray}C12} & {\color{gray}C13} & {\color{gray}C14} & {\color{gray}C15} & {\color{gray}C16}\\
\hline
\rowcolor{BenchBaseline}[0.5pt][0.5pt]\shortstack[l]{{\color{gray}R1}\\\textbf{minimap2}} & \textbf{Baseline} & \hspace*{1pt}\begin{tabular}[t]{@{}c@{}}\rule{0pt}{12pt}19.6s\\[-.2pt]{\fontsize{5.8}{6.4}\selectfont \textcolor{BrickRed}{$ \blacktriangledown\,6.7\times $}}\\[-.2pt]18.6s\\[-.2pt]{\fontsize{5.8}{6.4}\selectfont \textcolor{BrickRed}{$ \blacktriangledown\,8.4\times $}}\end{tabular} & \hspace*{1pt}\begin{tabular}[t]{@{}c@{}}\rule{0pt}{12pt}18.2s\\[-.2pt]{\fontsize{5.8}{6.4}\selectfont \textcolor{BrickRed}{$ \blacktriangledown\,6.9\times $}}\\[-.2pt]16.5s\\[-.2pt]{\fontsize{5.8}{6.4}\selectfont \textcolor{BrickRed}{$ \blacktriangledown\,8.6\times $}}\end{tabular} & \hspace*{1pt}\begin{tabular}[t]{@{}c@{}}\rule{0pt}{12pt}42.5s\\[-.2pt]{\fontsize{5.8}{6.4}\selectfont \textcolor{BrickRed}{$ \blacktriangledown\,8.2\times $}}\\[-.2pt]39.4s\\[-.2pt]{\fontsize{5.8}{6.4}\selectfont \textcolor{BrickRed}{$ \blacktriangledown\,12.6\times $}}\end{tabular} & \hspace*{1pt}\begin{tabular}[t]{@{}c@{}}\rule{0pt}{12pt}6m\,29.2s\\[-.2pt]{\fontsize{5.8}{6.4}\selectfont \textcolor{BrickRed}{$ \blacktriangledown\,92.8\times $}}\\[-.2pt]5m\,56.9s\\[-.2pt]{\fontsize{5.8}{6.4}\selectfont \textcolor{BrickRed}{$ \blacktriangledown\,184.9\times $}}\end{tabular} & \hspace*{1pt}\begin{tabular}[t]{@{}c@{}}\rule{0pt}{12pt}9.3s\\[-.2pt]{\fontsize{5.8}{6.4}\selectfont \textcolor{BrickRed}{$ \blacktriangledown\,24\% $}}\\[-.2pt]8.4s\\[-.2pt]{\fontsize{5.8}{6.4}\selectfont \textcolor{BrickRed}{$ \blacktriangledown\,1.6\times $}}\end{tabular} & \hspace*{1pt}\begin{tabular}[t]{@{}c@{}}\rule{0pt}{12pt}8.7s\\[-.2pt]{\fontsize{5.8}{6.4}\selectfont \textcolor{BrickRed}{$ \blacktriangledown\,1.5\times $}}\\[-.2pt]7s\\[-.2pt]{\fontsize{5.8}{6.4}\selectfont \textcolor{BrickRed}{$ \blacktriangledown\,1.7\times $}}\end{tabular} & \hspace*{1pt}\begin{tabular}[t]{@{}c@{}}\rule{0pt}{12pt}15.8s\\[-.2pt]{\fontsize{5.8}{6.4}\selectfont \textcolor{BrickRed}{$ \blacktriangledown\,26\% $}}\\[-.2pt]12.7s\\[-.2pt]{\fontsize{5.8}{6.4}\selectfont \textcolor{BrickRed}{$ \blacktriangledown\,1.7\times $}}\end{tabular} & \hspace*{1pt}\begin{tabular}[t]{@{}c@{}}\rule{0pt}{12pt}18.5s\\[-.2pt]{\fontsize{5.8}{6.4}\selectfont \textcolor{BrickRed}{$ \blacktriangledown\,1.6\times $}}\\[-.2pt]17s\\[-.2pt]{\fontsize{5.8}{6.4}\selectfont \textcolor{BrickRed}{$ \blacktriangledown\,1.8\times $}}\end{tabular} & \hspace*{1pt}\begin{tabular}[t]{@{}c@{}}\rule{0pt}{12pt}42.4s\\[-.2pt]{\fontsize{5.8}{6.4}\selectfont \textcolor{ForestGreen}{$ \blacktriangle\,2.0\times $}}\\[-.2pt]41.4s\\[-.2pt]{\fontsize{5.8}{6.4}\selectfont \textcolor{ForestGreen}{$ \blacktriangle\,1.9\times $}}\end{tabular} & \hspace*{1pt}\begin{tabular}[t]{@{}c@{}}\rule{0pt}{12pt}42.9s\\[-.2pt]{\fontsize{5.8}{6.4}\selectfont \textcolor{ForestGreen}{$ \blacktriangle\,2.0\times $}}\\[-.2pt]41.1s\\[-.2pt]{\fontsize{5.8}{6.4}\selectfont \textcolor{ForestGreen}{$ \blacktriangle\,1.8\times $}}\end{tabular} & \hspace*{1pt}\begin{tabular}[t]{@{}c@{}}\rule{0pt}{12pt}1m\,26.9s\\[-.2pt]{\fontsize{5.8}{6.4}\selectfont \textcolor{ForestGreen}{$ \blacktriangle\,1.9\times $}}\\[-.2pt]1m\,23.7s\\[-.2pt]{\fontsize{5.8}{6.4}\selectfont \textcolor{ForestGreen}{$ \blacktriangle\,1.7\times $}}\end{tabular} & \hspace*{1pt}\begin{tabular}[t]{@{}c@{}}\rule{0pt}{12pt}4m\,51.6s\\[-.2pt]{\fontsize{5.8}{6.4}\selectfont \textcolor{BrickRed}{$ \blacktriangledown\,2.9\times $}}\\[-.2pt]4m\,46.6s\\[-.2pt]{\fontsize{5.8}{6.4}\selectfont \textcolor{BrickRed}{$ \blacktriangledown\,3.0\times $}}\end{tabular} & \hspace*{1pt}\begin{tabular}[t]{@{}c@{}}\rule{0pt}{12pt}35s\\[-.2pt]{\fontsize{5.8}{6.4}\selectfont \textcolor{ForestGreen}{$ \blacktriangle\,8.6\times $}}\\[-.2pt]34s\\[-.2pt]{\fontsize{5.8}{6.4}\selectfont \textcolor{ForestGreen}{$ \blacktriangle\,2.3\times $}}\end{tabular} & \hspace*{1pt}\begin{tabular}[t]{@{}c@{}}\rule{0pt}{12pt}35.9s\\[-.2pt]{\fontsize{5.8}{6.4}\selectfont \textcolor{ForestGreen}{$ \blacktriangle\,6.7\times $}}\\[-.2pt]34.2s\\[-.2pt]{\fontsize{5.8}{6.4}\selectfont \textcolor{ForestGreen}{$ \blacktriangle\,2.1\times $}}\end{tabular} & \hspace*{1pt}\begin{tabular}[t]{@{}c@{}}\rule{0pt}{12pt}1m\,14.1s\\[-.2pt]{\fontsize{5.8}{6.4}\selectfont \textcolor{ForestGreen}{$ \blacktriangle\,6.9\times $}}\\[-.2pt]1m\,11.1s\\[-.2pt]{\fontsize{5.8}{6.4}\selectfont \textcolor{ForestGreen}{$ \blacktriangle\,2.2\times $}}\end{tabular} & \hspace*{1pt}\begin{tabular}[t]{@{}c@{}}\rule{0pt}{12pt}1m\,0.1s\\[-.2pt]{\fontsize{5.8}{6.4}\selectfont \textcolor{ForestGreen}{$ \blacktriangle\,10.4\times $}}\\[-.2pt]58.6s\\[-.2pt]{\fontsize{5.8}{6.4}\selectfont \textcolor{ForestGreen}{$ \blacktriangle\,1.8\times $}}\end{tabular}\\[3pt]\hline
\BenchmarkSourceLabel{bench:aws-2-k20-w15:row:mps-simdref-scalar-block}{R23}\shortstack[l]{{\color{gray}R23}\\\textbf{BioseqDB}} & \begin{tabular}[t]{@{}l@{}}\rule{0pt}{12pt}MPS / SIMD-ref\\Scalar-DP / Block\end{tabular} & \hspace*{1pt}\begin{tabular}[t]{@{}c@{}}\rule{0pt}{12pt}46.3s\\[-.2pt]{\fontsize{5.8}{6.4}\selectfont \textcolor{BrickRed}{$ \blacktriangledown\,2.4\times $}}\\[-.2pt]9.2s\\[-.2pt]{\fontsize{5.8}{6.4}\selectfont \textcolor{ForestGreen}{$ \blacktriangle\,2.0\times $}}\end{tabular} & \hspace*{1pt}\begin{tabular}[t]{@{}c@{}}\rule{0pt}{12pt}38.4s\\[-.2pt]{\fontsize{5.8}{6.4}\selectfont \textcolor{BrickRed}{$ \blacktriangledown\,2.1\times $}}\\[-.2pt]8.8s\\[-.2pt]{\fontsize{5.8}{6.4}\selectfont \textcolor{ForestGreen}{$ \blacktriangle\,1.9\times $}}\end{tabular} & \textcolor{gray}{\ensuremath{\cdot}} & \textcolor{gray}{\ensuremath{\cdot}} & \hspace*{1pt}\begin{tabular}[t]{@{}c@{}}\rule{0pt}{12pt}4m\,46.5s\\[-.2pt]{\fontsize{5.8}{6.4}\selectfont \textcolor{BrickRed}{$ \blacktriangledown\,30.7\times $}}\\[-.2pt]28.7s\\[-.2pt]{\fontsize{5.8}{6.4}\selectfont \textcolor{BrickRed}{$ \blacktriangledown\,3.4\times $}}\end{tabular} & \hspace*{1pt}\begin{tabular}[t]{@{}c@{}}\rule{0pt}{12pt}4m\,0.9s\\[-.2pt]{\fontsize{5.8}{6.4}\selectfont \textcolor{BrickRed}{$ \blacktriangledown\,27.6\times $}}\\[-.2pt]26.6s\\[-.2pt]{\fontsize{5.8}{6.4}\selectfont \textcolor{BrickRed}{$ \blacktriangledown\,3.8\times $}}\end{tabular} & \textcolor{gray}{\ensuremath{\cdot}} & \textcolor{gray}{\ensuremath{\cdot}} & \textcolor{gray}{\ensuremath{\cdot}} & \textcolor{gray}{\ensuremath{\cdot}} & \textcolor{gray}{\ensuremath{\cdot}} & \textcolor{gray}{\ensuremath{\cdot}} & \textcolor{gray}{\ensuremath{\cdot}} & \textcolor{gray}{\ensuremath{\cdot}} & \textcolor{gray}{\ensuremath{\cdot}} & \textcolor{gray}{\ensuremath{\cdot}}\\[3pt]\hline
\BenchmarkSourceLabel{bench:aws-2-k20-w15:row:mps-scalar-scalar-ksw2}{R24}\shortstack[l]{{\color{gray}R24}} & \begin{tabular}[t]{@{}l@{}}\rule{0pt}{12pt}MPS / Scalar\\Scalar-DP / KSW2\end{tabular} & \hspace*{1pt}\begin{tabular}[t]{@{}c@{}}\rule{0pt}{12pt}45.2s\\[-.2pt]{\fontsize{5.8}{6.4}\selectfont \textcolor{BrickRed}{$ \blacktriangledown\,2.3\times $}}\\[-.2pt]8.7s\\[-.2pt]{\fontsize{5.8}{6.4}\selectfont \textcolor{ForestGreen}{$ \blacktriangle\,2.1\times $}}\end{tabular} & \hspace*{1pt}\begin{tabular}[t]{@{}c@{}}\rule{0pt}{12pt}39.2s\\[-.2pt]{\fontsize{5.8}{6.4}\selectfont \textcolor{BrickRed}{$ \blacktriangledown\,2.2\times $}}\\[-.2pt]8.8s\\[-.2pt]{\fontsize{5.8}{6.4}\selectfont \textcolor{ForestGreen}{$ \blacktriangle\,1.9\times $}}\end{tabular} & \hspace*{1pt}\begin{tabular}[t]{@{}c@{}}\rule{0pt}{12pt}1m\,6.9s\\[-.2pt]{\fontsize{5.8}{6.4}\selectfont \textcolor{BrickRed}{$ \blacktriangledown\,1.6\times $}}\\[-.2pt]14.8s\\[-.2pt]{\fontsize{5.8}{6.4}\selectfont \textcolor{ForestGreen}{$ \blacktriangle\,2.7\times $}}\end{tabular} & \textcolor{gray}{\ensuremath{\cdot}} & \hspace*{1pt}\begin{tabular}[t]{@{}c@{}}\rule{0pt}{12pt}4m\,40.9s\\[-.2pt]{\fontsize{5.8}{6.4}\selectfont \textcolor{BrickRed}{$ \blacktriangledown\,30.1\times $}}\\[-.2pt]21.8s\\[-.2pt]{\fontsize{5.8}{6.4}\selectfont \textcolor{BrickRed}{$ \blacktriangledown\,2.6\times $}}\end{tabular} & \hspace*{1pt}\begin{tabular}[t]{@{}c@{}}\rule{0pt}{12pt}3m\,58.9s\\[-.2pt]{\fontsize{5.8}{6.4}\selectfont \textcolor{BrickRed}{$ \blacktriangledown\,27.3\times $}}\\[-.2pt]20.7s\\[-.2pt]{\fontsize{5.8}{6.4}\selectfont \textcolor{BrickRed}{$ \blacktriangledown\,2.9\times $}}\end{tabular} & \textcolor{gray}{\ensuremath{\cdot}} & \textcolor{gray}{\ensuremath{\cdot}} & \hspace*{1pt}\begin{tabular}[t]{@{}c@{}}\rule{0pt}{12pt}2m\,57.6s\\[-.2pt]{\fontsize{5.8}{6.4}\selectfont \textcolor{BrickRed}{$ \blacktriangledown\,4.2\times $}}\\[-.2pt]2m\,18.2s\\[-.2pt]{\fontsize{5.8}{6.4}\selectfont \textcolor{BrickRed}{$ \blacktriangledown\,3.3\times $}}\end{tabular} & \textcolor{gray}{\ensuremath{\cdot}} & \textcolor{gray}{\ensuremath{\cdot}} & \textcolor{gray}{\ensuremath{\cdot}} & \textcolor{gray}{\ensuremath{\cdot}} & \textcolor{gray}{\ensuremath{\cdot}} & \textcolor{gray}{\ensuremath{\cdot}} & \textcolor{gray}{\ensuremath{\cdot}}\\[3pt]\hline
\BenchmarkSourceLabel{bench:aws-2-k20-w15:row:mps-twostack-scalar-ksw2}{R25}\shortstack[l]{{\color{gray}R25}} & \begin{tabular}[t]{@{}l@{}}\rule{0pt}{12pt}MPS / Two-stack\\Scalar-DP / KSW2\end{tabular} & \hspace*{1pt}\begin{tabular}[t]{@{}c@{}}\rule{0pt}{12pt}43.6s\\[-.2pt]{\fontsize{5.8}{6.4}\selectfont \textcolor{BrickRed}{$ \blacktriangledown\,2.2\times $}}\\[-.2pt]6.8s\\[-.2pt]{\fontsize{5.8}{6.4}\selectfont \textcolor{ForestGreen}{$ \blacktriangle\,2.7\times $}}\end{tabular} & \hspace*{1pt}\begin{tabular}[t]{@{}c@{}}\rule{0pt}{12pt}39.5s\\[-.2pt]{\fontsize{5.8}{6.4}\selectfont \textcolor{BrickRed}{$ \blacktriangledown\,2.2\times $}}\\[-.2pt]7.2s\\[-.2pt]{\fontsize{5.8}{6.4}\selectfont \textcolor{ForestGreen}{$ \blacktriangle\,2.3\times $}}\end{tabular} & \textcolor{gray}{\ensuremath{\cdot}} & \textcolor{gray}{\ensuremath{\cdot}} & \hspace*{1pt}\begin{tabular}[t]{@{}c@{}}\rule{0pt}{12pt}4m\,31.6s\\[-.2pt]{\fontsize{5.8}{6.4}\selectfont \textcolor{BrickRed}{$ \blacktriangledown\,29.1\times $}}\\[-.2pt]12.4s\\[-.2pt]{\fontsize{5.8}{6.4}\selectfont \textcolor{BrickRed}{$ \blacktriangledown\,1.5\times $}}\end{tabular} & \hspace*{1pt}\begin{tabular}[t]{@{}c@{}}\rule{0pt}{12pt}3m\,44.6s\\[-.2pt]{\fontsize{5.8}{6.4}\selectfont \textcolor{BrickRed}{$ \blacktriangledown\,25.7\times $}}\\[-.2pt]11.3s\\[-.2pt]{\fontsize{5.8}{6.4}\selectfont \textcolor{BrickRed}{$ \blacktriangledown\,1.6\times $}}\end{tabular} & \textcolor{gray}{\ensuremath{\cdot}} & \textcolor{gray}{\ensuremath{\cdot}} & \textcolor{gray}{\ensuremath{\cdot}} & \textcolor{gray}{\ensuremath{\cdot}} & \textcolor{gray}{\ensuremath{\cdot}} & \textcolor{gray}{\ensuremath{\cdot}} & \textcolor{gray}{\ensuremath{\cdot}} & \textcolor{gray}{\ensuremath{\cdot}} & \textcolor{gray}{\ensuremath{\cdot}} & \textcolor{gray}{\ensuremath{\cdot}}\\[3pt]\hline
\BenchmarkSourceLabel{bench:aws-2-k20-w15:row:mps-twostack-simd-ksw2}{R26}\shortstack[l]{{\color{gray}R26}} & \begin{tabular}[t]{@{}l@{}}\rule{0pt}{12pt}MPS / Two-stack\\SIMD-DP / KSW2\end{tabular} & \hspace*{1pt}\begin{tabular}[t]{@{}c@{}}\rule{0pt}{12pt}43.7s\\[-.2pt]{\fontsize{5.8}{6.4}\selectfont \textcolor{BrickRed}{$ \blacktriangledown\,2.2\times $}}\\[-.2pt]6.8s\\[-.2pt]{\fontsize{5.8}{6.4}\selectfont \textcolor{ForestGreen}{$ \blacktriangle\,2.7\times $}}\end{tabular} & \hspace*{1pt}\begin{tabular}[t]{@{}c@{}}\rule{0pt}{12pt}39.6s\\[-.2pt]{\fontsize{5.8}{6.4}\selectfont \textcolor{BrickRed}{$ \blacktriangledown\,2.2\times $}}\\[-.2pt]7.2s\\[-.2pt]{\fontsize{5.8}{6.4}\selectfont \textcolor{ForestGreen}{$ \blacktriangle\,2.3\times $}}\end{tabular} & \textcolor{gray}{\ensuremath{\cdot}} & \textcolor{gray}{\ensuremath{\cdot}} & \hspace*{1pt}\begin{tabular}[t]{@{}c@{}}\rule{0pt}{12pt}4m\,30s\\[-.2pt]{\fontsize{5.8}{6.4}\selectfont \textcolor{BrickRed}{$ \blacktriangledown\,29.0\times $}}\\[-.2pt]12.4s\\[-.2pt]{\fontsize{5.8}{6.4}\selectfont \textcolor{BrickRed}{$ \blacktriangledown\,1.5\times $}}\end{tabular} & \hspace*{1pt}\begin{tabular}[t]{@{}c@{}}\rule{0pt}{12pt}3m\,45.3s\\[-.2pt]{\fontsize{5.8}{6.4}\selectfont \textcolor{BrickRed}{$ \blacktriangledown\,25.8\times $}}\\[-.2pt]11.5s\\[-.2pt]{\fontsize{5.8}{6.4}\selectfont \textcolor{BrickRed}{$ \blacktriangledown\,1.6\times $}}\end{tabular} & \textcolor{gray}{\ensuremath{\cdot}} & \textcolor{gray}{\ensuremath{\cdot}} & \textcolor{gray}{\ensuremath{\cdot}} & \textcolor{gray}{\ensuremath{\cdot}} & \textcolor{gray}{\ensuremath{\cdot}} & \textcolor{gray}{\ensuremath{\cdot}} & \textcolor{gray}{\ensuremath{\cdot}} & \textcolor{gray}{\ensuremath{\cdot}} & \textcolor{gray}{\ensuremath{\cdot}} & \textcolor{gray}{\ensuremath{\cdot}}\\[3pt]\hline
\BenchmarkSourceLabel{bench:aws-2-k20-w15:row:mps-twostack-rmq-ksw2}{R27}\shortstack[l]{{\color{gray}R27}} & \begin{tabular}[t]{@{}l@{}}\rule{0pt}{12pt}MPS / Two-stack\\RMQ / KSW2\end{tabular} & \hspace*{1pt}\begin{tabular}[t]{@{}c@{}}\rule{0pt}{12pt}43.4s\\[-.2pt]{\fontsize{5.8}{6.4}\selectfont \textcolor{BrickRed}{$ \blacktriangledown\,2.2\times $}}\\[-.2pt]6.8s\\[-.2pt]{\fontsize{5.8}{6.4}\selectfont \textcolor{ForestGreen}{$ \blacktriangle\,2.7\times $}}\end{tabular} & \hspace*{1pt}\begin{tabular}[t]{@{}c@{}}\rule{0pt}{12pt}39.6s\\[-.2pt]{\fontsize{5.8}{6.4}\selectfont \textcolor{BrickRed}{$ \blacktriangledown\,2.2\times $}}\\[-.2pt]6.7s\\[-.2pt]{\fontsize{5.8}{6.4}\selectfont \textcolor{ForestGreen}{$ \blacktriangle\,2.4\times $}}\end{tabular} & \textcolor{gray}{\ensuremath{\cdot}} & \textcolor{gray}{\ensuremath{\cdot}} & \hspace*{1pt}\begin{tabular}[t]{@{}c@{}}\rule{0pt}{12pt}4m\,29.8s\\[-.2pt]{\fontsize{5.8}{6.4}\selectfont \textcolor{BrickRed}{$ \blacktriangledown\,28.9\times $}}\\[-.2pt]12.6s\\[-.2pt]{\fontsize{5.8}{6.4}\selectfont \textcolor{BrickRed}{$ \blacktriangledown\,1.5\times $}}\end{tabular} & \hspace*{1pt}\begin{tabular}[t]{@{}c@{}}\rule{0pt}{12pt}3m\,45.3s\\[-.2pt]{\fontsize{5.8}{6.4}\selectfont \textcolor{BrickRed}{$ \blacktriangledown\,25.8\times $}}\\[-.2pt]11.5s\\[-.2pt]{\fontsize{5.8}{6.4}\selectfont \textcolor{BrickRed}{$ \blacktriangledown\,1.6\times $}}\end{tabular} & \textcolor{gray}{\ensuremath{\cdot}} & \textcolor{gray}{\ensuremath{\cdot}} & \textcolor{gray}{\ensuremath{\cdot}} & \textcolor{gray}{\ensuremath{\cdot}} & \textcolor{gray}{\ensuremath{\cdot}} & \textcolor{gray}{\ensuremath{\cdot}} & \textcolor{gray}{\ensuremath{\cdot}} & \textcolor{gray}{\ensuremath{\cdot}} & \textcolor{gray}{\ensuremath{\cdot}} & \textcolor{gray}{\ensuremath{\cdot}}\\[3pt]\hline
\BenchmarkSourceLabel{bench:aws-2-k20-w15:row:mps-scalar-rmq-ksw2}{R28}\shortstack[l]{{\color{gray}R28}} & \begin{tabular}[t]{@{}l@{}}\rule{0pt}{12pt}MPS / Scalar\\RMQ / KSW2\end{tabular} & \hspace*{1pt}\begin{tabular}[t]{@{}c@{}}\rule{0pt}{12pt}45.4s\\[-.2pt]{\fontsize{5.8}{6.4}\selectfont \textcolor{BrickRed}{$ \blacktriangledown\,2.3\times $}}\\[-.2pt]8.8s\\[-.2pt]{\fontsize{5.8}{6.4}\selectfont \textcolor{ForestGreen}{$ \blacktriangle\,2.1\times $}}\end{tabular} & \hspace*{1pt}\begin{tabular}[t]{@{}c@{}}\rule{0pt}{12pt}39.7s\\[-.2pt]{\fontsize{5.8}{6.4}\selectfont \textcolor{BrickRed}{$ \blacktriangledown\,2.2\times $}}\\[-.2pt]9s\\[-.2pt]{\fontsize{5.8}{6.4}\selectfont \textcolor{ForestGreen}{$ \blacktriangle\,1.8\times $}}\end{tabular} & \textcolor{gray}{\ensuremath{\cdot}} & \textcolor{gray}{\ensuremath{\cdot}} & \hspace*{1pt}\begin{tabular}[t]{@{}c@{}}\rule{0pt}{12pt}4m\,41.5s\\[-.2pt]{\fontsize{5.8}{6.4}\selectfont \textcolor{BrickRed}{$ \blacktriangledown\,30.2\times $}}\\[-.2pt]22.4s\\[-.2pt]{\fontsize{5.8}{6.4}\selectfont \textcolor{BrickRed}{$ \blacktriangledown\,2.7\times $}}\end{tabular} & \hspace*{1pt}\begin{tabular}[t]{@{}c@{}}\rule{0pt}{12pt}4m\,0.1s\\[-.2pt]{\fontsize{5.8}{6.4}\selectfont \textcolor{BrickRed}{$ \blacktriangledown\,27.5\times $}}\\[-.2pt]21.1s\\[-.2pt]{\fontsize{5.8}{6.4}\selectfont \textcolor{BrickRed}{$ \blacktriangledown\,3.0\times $}}\end{tabular} & \textcolor{gray}{\ensuremath{\cdot}} & \textcolor{gray}{\ensuremath{\cdot}} & \hspace*{1pt}\begin{tabular}[t]{@{}c@{}}\rule{0pt}{12pt}2m\,52.9s\\[-.2pt]{\fontsize{5.8}{6.4}\selectfont \textcolor{BrickRed}{$ \blacktriangledown\,4.1\times $}}\\[-.2pt]2m\,15.4s\\[-.2pt]{\fontsize{5.8}{6.4}\selectfont \textcolor{BrickRed}{$ \blacktriangledown\,3.3\times $}}\end{tabular} & \hspace*{1pt}\begin{tabular}[t]{@{}c@{}}\rule{0pt}{12pt}2m\,47.4s\\[-.2pt]{\fontsize{5.8}{6.4}\selectfont \textcolor{BrickRed}{$ \blacktriangledown\,3.9\times $}}\\[-.2pt]2m\,15.8s\\[-.2pt]{\fontsize{5.8}{6.4}\selectfont \textcolor{BrickRed}{$ \blacktriangledown\,3.3\times $}}\end{tabular} & \textcolor{gray}{\ensuremath{\cdot}} & \textcolor{gray}{\ensuremath{\cdot}} & \textcolor{gray}{\ensuremath{\cdot}} & \textcolor{gray}{\ensuremath{\cdot}} & \textcolor{gray}{\ensuremath{\cdot}} & \textcolor{gray}{\ensuremath{\cdot}}\\[3pt]\hline
\BenchmarkSourceLabel{bench:aws-2-k20-w15:row:mps-simdref-scalar-ksw2}{R29}\shortstack[l]{{\color{gray}R29}} & \begin{tabular}[t]{@{}l@{}}\rule{0pt}{12pt}MPS / SIMD-ref\\Scalar-DP / KSW2\end{tabular} & \hspace*{1pt}\begin{tabular}[t]{@{}c@{}}\rule{0pt}{12pt}47.8s\\[-.2pt]{\fontsize{5.8}{6.4}\selectfont \textcolor{BrickRed}{$ \blacktriangledown\,2.4\times $}}\\[-.2pt]9.1s\\[-.2pt]{\fontsize{5.8}{6.4}\selectfont \textcolor{ForestGreen}{$ \blacktriangle\,2.0\times $}}\end{tabular} & \hspace*{1pt}\begin{tabular}[t]{@{}c@{}}\rule{0pt}{12pt}39.9s\\[-.2pt]{\fontsize{5.8}{6.4}\selectfont \textcolor{BrickRed}{$ \blacktriangledown\,2.2\times $}}\\[-.2pt]8.7s\\[-.2pt]{\fontsize{5.8}{6.4}\selectfont \textcolor{ForestGreen}{$ \blacktriangle\,1.9\times $}}\end{tabular} & \textcolor{gray}{\ensuremath{\cdot}} & \textcolor{gray}{\ensuremath{\cdot}} & \hspace*{1pt}\begin{tabular}[t]{@{}c@{}}\rule{0pt}{12pt}4m\,48.7s\\[-.2pt]{\fontsize{5.8}{6.4}\selectfont \textcolor{BrickRed}{$ \blacktriangledown\,31.0\times $}}\\[-.2pt]28.7s\\[-.2pt]{\fontsize{5.8}{6.4}\selectfont \textcolor{BrickRed}{$ \blacktriangledown\,3.4\times $}}\end{tabular} & \hspace*{1pt}\begin{tabular}[t]{@{}c@{}}\rule{0pt}{12pt}4m\,3.3s\\[-.2pt]{\fontsize{5.8}{6.4}\selectfont \textcolor{BrickRed}{$ \blacktriangledown\,27.8\times $}}\\[-.2pt]27s\\[-.2pt]{\fontsize{5.8}{6.4}\selectfont \textcolor{BrickRed}{$ \blacktriangledown\,3.8\times $}}\end{tabular} & \textcolor{gray}{\ensuremath{\cdot}} & \textcolor{gray}{\ensuremath{\cdot}} & \textcolor{gray}{\ensuremath{\cdot}} & \textcolor{gray}{\ensuremath{\cdot}} & \textcolor{gray}{\ensuremath{\cdot}} & \textcolor{gray}{\ensuremath{\cdot}} & \textcolor{gray}{\ensuremath{\cdot}} & \textcolor{gray}{\ensuremath{\cdot}} & \textcolor{gray}{\ensuremath{\cdot}} & \textcolor{gray}{\ensuremath{\cdot}}\\[3pt]\hline
\end{tabular}\par}
\smallskip{\footnotesize\noindent Each cell: A and B times, $\cdot$: no measurement. ---\textsuperscript{g}: GPU unavailable. n/a: no comparison.\par
\noindent Blue: one lowest A and one lowest B measurement for each column for each configuration.\par
\noindent Gray: minimap2. Yellow: my preferred MPF / Bounded / SIMD-DP / Block configuration.\par
\noindent MPF/MPS: minimizers per fragment/sequence. SIMD-ref: naive SIMD scan. Bounded: bounded SIMD algorithm.\par
}
\clearpage
\noindent{\large\bfseries Table \thetable: (E): AWS: 2 Gbp, k=20, w=15\hfill 5/7}\par\smallskip
{\fontsize{7.1}{8.1}\selectfont\setlength{\tabcolsep}{1pt}\setlength{\fboxsep}{2pt}\renewcommand{\arraystretch}{1.12}
\noindent\begin{tabular}{@{}p{14mm}p{27mm}|p{13.55mm}|p{13.55mm}|p{13.55mm}|p{13.55mm}|p{13.55mm}|p{13.55mm}|p{13.55mm}|p{13.55mm}|p{13.55mm}|p{13.55mm}|p{13.55mm}|p{13.55mm}|p{13.55mm}|p{13.55mm}|p{13.55mm}|p{13.55mm}|@{}}\hline
Engine & Execution & \multicolumn{8}{c|}{\textbf{Index}} & \multicolumn{8}{c|}{\textbf{No index}}\\
 &  & \multicolumn{8}{c|}{\textbf{2 Gbp}} & \multicolumn{8}{c|}{\textbf{2 Gbp}}\\
 &  & \multicolumn{8}{c|}{\textbf{k=20, w=15}} & \multicolumn{8}{c|}{\textbf{k=20, w=15}}\\
 &  & \multicolumn{4}{c|}{\textbf{DNA4}} & \multicolumn{4}{c|}{\textbf{DNA5}} & \multicolumn{4}{c|}{\textbf{DNA4}} & \multicolumn{4}{c|}{\textbf{DNA5}}\\
 &  & \multicolumn{1}{c|}{\textbf{M7A}} & \multicolumn{1}{c|}{\textbf{M7I}} & \multicolumn{1}{c|}{\textbf{M4}} & \multicolumn{1}{c|}{\textbf{C7G}} & \multicolumn{1}{c|}{\textbf{M7A}} & \multicolumn{1}{c|}{\textbf{M7I}} & \multicolumn{1}{c|}{\textbf{M4}} & \multicolumn{1}{c|}{\textbf{C7G}} & \multicolumn{1}{c|}{\textbf{M7A}} & \multicolumn{1}{c|}{\textbf{M7I}} & \multicolumn{1}{c|}{\textbf{M4}} & \multicolumn{1}{c|}{\textbf{C7G}} & \multicolumn{1}{c|}{\textbf{M7A}} & \multicolumn{1}{c|}{\textbf{M7I}} & \multicolumn{1}{c|}{\textbf{M4}} & \multicolumn{1}{c|}{\textbf{C7G}}\\
 & & {\color{gray}C1} & {\color{gray}C2} & {\color{gray}C3} & {\color{gray}C4} & {\color{gray}C5} & {\color{gray}C6} & {\color{gray}C7} & {\color{gray}C8} & {\color{gray}C9} & {\color{gray}C10} & {\color{gray}C11} & {\color{gray}C12} & {\color{gray}C13} & {\color{gray}C14} & {\color{gray}C15} & {\color{gray}C16}\\
\hline
\rowcolor{BenchBaseline}[0.5pt][0.5pt]\shortstack[l]{{\color{gray}R1}\\\textbf{minimap2}} & \textbf{Baseline} & \hspace*{1pt}\begin{tabular}[t]{@{}c@{}}\rule{0pt}{12pt}19.6s\\[-.2pt]{\fontsize{5.8}{6.4}\selectfont \textcolor{BrickRed}{$ \blacktriangledown\,6.7\times $}}\\[-.2pt]18.6s\\[-.2pt]{\fontsize{5.8}{6.4}\selectfont \textcolor{BrickRed}{$ \blacktriangledown\,8.4\times $}}\end{tabular} & \hspace*{1pt}\begin{tabular}[t]{@{}c@{}}\rule{0pt}{12pt}18.2s\\[-.2pt]{\fontsize{5.8}{6.4}\selectfont \textcolor{BrickRed}{$ \blacktriangledown\,6.9\times $}}\\[-.2pt]16.5s\\[-.2pt]{\fontsize{5.8}{6.4}\selectfont \textcolor{BrickRed}{$ \blacktriangledown\,8.6\times $}}\end{tabular} & \hspace*{1pt}\begin{tabular}[t]{@{}c@{}}\rule{0pt}{12pt}42.5s\\[-.2pt]{\fontsize{5.8}{6.4}\selectfont \textcolor{BrickRed}{$ \blacktriangledown\,8.2\times $}}\\[-.2pt]39.4s\\[-.2pt]{\fontsize{5.8}{6.4}\selectfont \textcolor{BrickRed}{$ \blacktriangledown\,12.6\times $}}\end{tabular} & \hspace*{1pt}\begin{tabular}[t]{@{}c@{}}\rule{0pt}{12pt}6m\,29.2s\\[-.2pt]{\fontsize{5.8}{6.4}\selectfont \textcolor{BrickRed}{$ \blacktriangledown\,92.8\times $}}\\[-.2pt]5m\,56.9s\\[-.2pt]{\fontsize{5.8}{6.4}\selectfont \textcolor{BrickRed}{$ \blacktriangledown\,184.9\times $}}\end{tabular} & \hspace*{1pt}\begin{tabular}[t]{@{}c@{}}\rule{0pt}{12pt}9.3s\\[-.2pt]{\fontsize{5.8}{6.4}\selectfont \textcolor{BrickRed}{$ \blacktriangledown\,24\% $}}\\[-.2pt]8.4s\\[-.2pt]{\fontsize{5.8}{6.4}\selectfont \textcolor{BrickRed}{$ \blacktriangledown\,1.6\times $}}\end{tabular} & \hspace*{1pt}\begin{tabular}[t]{@{}c@{}}\rule{0pt}{12pt}8.7s\\[-.2pt]{\fontsize{5.8}{6.4}\selectfont \textcolor{BrickRed}{$ \blacktriangledown\,1.5\times $}}\\[-.2pt]7s\\[-.2pt]{\fontsize{5.8}{6.4}\selectfont \textcolor{BrickRed}{$ \blacktriangledown\,1.7\times $}}\end{tabular} & \hspace*{1pt}\begin{tabular}[t]{@{}c@{}}\rule{0pt}{12pt}15.8s\\[-.2pt]{\fontsize{5.8}{6.4}\selectfont \textcolor{BrickRed}{$ \blacktriangledown\,26\% $}}\\[-.2pt]12.7s\\[-.2pt]{\fontsize{5.8}{6.4}\selectfont \textcolor{BrickRed}{$ \blacktriangledown\,1.7\times $}}\end{tabular} & \hspace*{1pt}\begin{tabular}[t]{@{}c@{}}\rule{0pt}{12pt}18.5s\\[-.2pt]{\fontsize{5.8}{6.4}\selectfont \textcolor{BrickRed}{$ \blacktriangledown\,1.6\times $}}\\[-.2pt]17s\\[-.2pt]{\fontsize{5.8}{6.4}\selectfont \textcolor{BrickRed}{$ \blacktriangledown\,1.8\times $}}\end{tabular} & \hspace*{1pt}\begin{tabular}[t]{@{}c@{}}\rule{0pt}{12pt}42.4s\\[-.2pt]{\fontsize{5.8}{6.4}\selectfont \textcolor{ForestGreen}{$ \blacktriangle\,2.0\times $}}\\[-.2pt]41.4s\\[-.2pt]{\fontsize{5.8}{6.4}\selectfont \textcolor{ForestGreen}{$ \blacktriangle\,1.9\times $}}\end{tabular} & \hspace*{1pt}\begin{tabular}[t]{@{}c@{}}\rule{0pt}{12pt}42.9s\\[-.2pt]{\fontsize{5.8}{6.4}\selectfont \textcolor{ForestGreen}{$ \blacktriangle\,2.0\times $}}\\[-.2pt]41.1s\\[-.2pt]{\fontsize{5.8}{6.4}\selectfont \textcolor{ForestGreen}{$ \blacktriangle\,1.8\times $}}\end{tabular} & \hspace*{1pt}\begin{tabular}[t]{@{}c@{}}\rule{0pt}{12pt}1m\,26.9s\\[-.2pt]{\fontsize{5.8}{6.4}\selectfont \textcolor{ForestGreen}{$ \blacktriangle\,1.9\times $}}\\[-.2pt]1m\,23.7s\\[-.2pt]{\fontsize{5.8}{6.4}\selectfont \textcolor{ForestGreen}{$ \blacktriangle\,1.7\times $}}\end{tabular} & \hspace*{1pt}\begin{tabular}[t]{@{}c@{}}\rule{0pt}{12pt}4m\,51.6s\\[-.2pt]{\fontsize{5.8}{6.4}\selectfont \textcolor{BrickRed}{$ \blacktriangledown\,2.9\times $}}\\[-.2pt]4m\,46.6s\\[-.2pt]{\fontsize{5.8}{6.4}\selectfont \textcolor{BrickRed}{$ \blacktriangledown\,3.0\times $}}\end{tabular} & \hspace*{1pt}\begin{tabular}[t]{@{}c@{}}\rule{0pt}{12pt}35s\\[-.2pt]{\fontsize{5.8}{6.4}\selectfont \textcolor{ForestGreen}{$ \blacktriangle\,8.6\times $}}\\[-.2pt]34s\\[-.2pt]{\fontsize{5.8}{6.4}\selectfont \textcolor{ForestGreen}{$ \blacktriangle\,2.3\times $}}\end{tabular} & \hspace*{1pt}\begin{tabular}[t]{@{}c@{}}\rule{0pt}{12pt}35.9s\\[-.2pt]{\fontsize{5.8}{6.4}\selectfont \textcolor{ForestGreen}{$ \blacktriangle\,6.7\times $}}\\[-.2pt]34.2s\\[-.2pt]{\fontsize{5.8}{6.4}\selectfont \textcolor{ForestGreen}{$ \blacktriangle\,2.1\times $}}\end{tabular} & \hspace*{1pt}\begin{tabular}[t]{@{}c@{}}\rule{0pt}{12pt}1m\,14.1s\\[-.2pt]{\fontsize{5.8}{6.4}\selectfont \textcolor{ForestGreen}{$ \blacktriangle\,6.9\times $}}\\[-.2pt]1m\,11.1s\\[-.2pt]{\fontsize{5.8}{6.4}\selectfont \textcolor{ForestGreen}{$ \blacktriangle\,2.2\times $}}\end{tabular} & \hspace*{1pt}\begin{tabular}[t]{@{}c@{}}\rule{0pt}{12pt}1m\,0.1s\\[-.2pt]{\fontsize{5.8}{6.4}\selectfont \textcolor{ForestGreen}{$ \blacktriangle\,10.4\times $}}\\[-.2pt]58.6s\\[-.2pt]{\fontsize{5.8}{6.4}\selectfont \textcolor{ForestGreen}{$ \blacktriangle\,1.8\times $}}\end{tabular}\\[3pt]\hline
\BenchmarkSourceLabel{bench:aws-2-k20-w15:row:mps-simdref-rmq-ksw2}{R30}\shortstack[l]{{\color{gray}R30}\\\textbf{BioseqDB}} & \begin{tabular}[t]{@{}l@{}}\rule{0pt}{12pt}MPS / SIMD-ref\\RMQ / KSW2\end{tabular} & \hspace*{1pt}\begin{tabular}[t]{@{}c@{}}\rule{0pt}{12pt}48.4s\\[-.2pt]{\fontsize{5.8}{6.4}\selectfont \textcolor{BrickRed}{$ \blacktriangledown\,2.5\times $}}\\[-.2pt]9.1s\\[-.2pt]{\fontsize{5.8}{6.4}\selectfont \textcolor{ForestGreen}{$ \blacktriangle\,2.0\times $}}\end{tabular} & \hspace*{1pt}\begin{tabular}[t]{@{}c@{}}\rule{0pt}{12pt}40.4s\\[-.2pt]{\fontsize{5.8}{6.4}\selectfont \textcolor{BrickRed}{$ \blacktriangledown\,2.2\times $}}\\[-.2pt]8.8s\\[-.2pt]{\fontsize{5.8}{6.4}\selectfont \textcolor{ForestGreen}{$ \blacktriangle\,1.9\times $}}\end{tabular} & \textcolor{gray}{\ensuremath{\cdot}} & \textcolor{gray}{\ensuremath{\cdot}} & \hspace*{1pt}\begin{tabular}[t]{@{}c@{}}\rule{0pt}{12pt}4m\,48.1s\\[-.2pt]{\fontsize{5.8}{6.4}\selectfont \textcolor{BrickRed}{$ \blacktriangledown\,30.9\times $}}\\[-.2pt]28.9s\\[-.2pt]{\fontsize{5.8}{6.4}\selectfont \textcolor{BrickRed}{$ \blacktriangledown\,3.4\times $}}\end{tabular} & \hspace*{1pt}\begin{tabular}[t]{@{}c@{}}\rule{0pt}{12pt}4m\,3.5s\\[-.2pt]{\fontsize{5.8}{6.4}\selectfont \textcolor{BrickRed}{$ \blacktriangledown\,27.9\times $}}\\[-.2pt]27.1s\\[-.2pt]{\fontsize{5.8}{6.4}\selectfont \textcolor{BrickRed}{$ \blacktriangledown\,3.9\times $}}\end{tabular} & \textcolor{gray}{\ensuremath{\cdot}} & \textcolor{gray}{\ensuremath{\cdot}} & \textcolor{gray}{\ensuremath{\cdot}} & \textcolor{gray}{\ensuremath{\cdot}} & \textcolor{gray}{\ensuremath{\cdot}} & \textcolor{gray}{\ensuremath{\cdot}} & \textcolor{gray}{\ensuremath{\cdot}} & \textcolor{gray}{\ensuremath{\cdot}} & \textcolor{gray}{\ensuremath{\cdot}} & \textcolor{gray}{\ensuremath{\cdot}}\\[3pt]\hline
\BenchmarkSourceLabel{bench:aws-2-k20-w15:row:mps-simdref-simd-ksw2}{R31}\shortstack[l]{{\color{gray}R31}} & \begin{tabular}[t]{@{}l@{}}\rule{0pt}{12pt}MPS / SIMD-ref\\SIMD-DP / KSW2\end{tabular} & \hspace*{1pt}\begin{tabular}[t]{@{}c@{}}\rule{0pt}{12pt}47.5s\\[-.2pt]{\fontsize{5.8}{6.4}\selectfont \textcolor{BrickRed}{$ \blacktriangledown\,2.4\times $}}\\[-.2pt]8.9s\\[-.2pt]{\fontsize{5.8}{6.4}\selectfont \textcolor{ForestGreen}{$ \blacktriangle\,2.1\times $}}\end{tabular} & \hspace*{1pt}\begin{tabular}[t]{@{}c@{}}\rule{0pt}{12pt}40.7s\\[-.2pt]{\fontsize{5.8}{6.4}\selectfont \textcolor{BrickRed}{$ \blacktriangledown\,2.2\times $}}\\[-.2pt]8.9s\\[-.2pt]{\fontsize{5.8}{6.4}\selectfont \textcolor{ForestGreen}{$ \blacktriangle\,1.9\times $}}\end{tabular} & \hspace*{1pt}\begin{tabular}[t]{@{}c@{}}\rule{0pt}{12pt}1m\,6.7s\\[-.2pt]{\fontsize{5.8}{6.4}\selectfont \textcolor{BrickRed}{$ \blacktriangledown\,1.6\times $}}\\[-.2pt]14.4s\\[-.2pt]{\fontsize{5.8}{6.4}\selectfont \textcolor{ForestGreen}{$ \blacktriangle\,2.7\times $}}\end{tabular} & \textcolor{gray}{\ensuremath{\cdot}} & \hspace*{1pt}\begin{tabular}[t]{@{}c@{}}\rule{0pt}{12pt}4m\,47.7s\\[-.2pt]{\fontsize{5.8}{6.4}\selectfont \textcolor{BrickRed}{$ \blacktriangledown\,30.9\times $}}\\[-.2pt]28.5s\\[-.2pt]{\fontsize{5.8}{6.4}\selectfont \textcolor{BrickRed}{$ \blacktriangledown\,3.4\times $}}\end{tabular} & \hspace*{1pt}\begin{tabular}[t]{@{}c@{}}\rule{0pt}{12pt}4m\,3.9s\\[-.2pt]{\fontsize{5.8}{6.4}\selectfont \textcolor{BrickRed}{$ \blacktriangledown\,27.9\times $}}\\[-.2pt]27.7s\\[-.2pt]{\fontsize{5.8}{6.4}\selectfont \textcolor{BrickRed}{$ \blacktriangledown\,3.9\times $}}\end{tabular} & \textcolor{gray}{\ensuremath{\cdot}} & \textcolor{gray}{\ensuremath{\cdot}} & \textcolor{gray}{\ensuremath{\cdot}} & \textcolor{gray}{\ensuremath{\cdot}} & \textcolor{gray}{\ensuremath{\cdot}} & \textcolor{gray}{\ensuremath{\cdot}} & \textcolor{gray}{\ensuremath{\cdot}} & \textcolor{gray}{\ensuremath{\cdot}} & \textcolor{gray}{\ensuremath{\cdot}} & \textcolor{gray}{\ensuremath{\cdot}}\\[3pt]\hline
\BenchmarkSourceLabel{bench:aws-2-k20-w15:row:mpf-twostack-rmq-block}{R32}\shortstack[l]{{\color{gray}R32}} & \begin{tabular}[t]{@{}l@{}}\rule{0pt}{12pt}MPF / Two-stack\\RMQ / Block\end{tabular} & \hspace*{1pt}\begin{tabular}[t]{@{}c@{}}\rule{0pt}{12pt}4m\,37.6s\\[-.2pt]{\fontsize{5.8}{6.4}\selectfont \textcolor{BrickRed}{$ \blacktriangledown\,14.2\times $}}\\[-.2pt]4m\,37.1s\\[-.2pt]{\fontsize{5.8}{6.4}\selectfont \textcolor{BrickRed}{$ \blacktriangledown\,14.9\times $}}\end{tabular} & \hspace*{1pt}\begin{tabular}[t]{@{}c@{}}\rule{0pt}{12pt}7m\,30s\\[-.2pt]{\fontsize{5.8}{6.4}\selectfont \textcolor{BrickRed}{$ \blacktriangledown\,24.7\times $}}\\[-.2pt]7m\,29.2s\\[-.2pt]{\fontsize{5.8}{6.4}\selectfont \textcolor{BrickRed}{$ \blacktriangledown\,27.2\times $}}\end{tabular} & \hspace*{1pt}\begin{tabular}[t]{@{}c@{}}\rule{0pt}{12pt}17m\,24.5s\\[-.2pt]{\fontsize{5.8}{6.4}\selectfont \textcolor{BrickRed}{$ \blacktriangledown\,24.6\times $}}\\[-.2pt]17m\,22.9s\\[-.2pt]{\fontsize{5.8}{6.4}\selectfont \textcolor{BrickRed}{$ \blacktriangledown\,26.5\times $}}\end{tabular} & \textcolor{gray}{\ensuremath{\cdot}} & \hspace*{1pt}\begin{tabular}[t]{@{}c@{}}\rule{0pt}{12pt}3m\,23s\\[-.2pt]{\fontsize{5.8}{6.4}\selectfont \textcolor{BrickRed}{$ \blacktriangledown\,21.8\times $}}\\[-.2pt]3m\,21.6s\\[-.2pt]{\fontsize{5.8}{6.4}\selectfont \textcolor{BrickRed}{$ \blacktriangledown\,24.0\times $}}\end{tabular} & \hspace*{1pt}\begin{tabular}[t]{@{}c@{}}\rule{0pt}{12pt}5m\,46.9s\\[-.2pt]{\fontsize{5.8}{6.4}\selectfont \textcolor{BrickRed}{$ \blacktriangledown\,39.7\times $}}\\[-.2pt]5m\,45.4s\\[-.2pt]{\fontsize{5.8}{6.4}\selectfont \textcolor{BrickRed}{$ \blacktriangledown\,49.2\times $}}\end{tabular} & \hspace*{1pt}\begin{tabular}[t]{@{}c@{}}\rule{0pt}{12pt}11m\,57.4s\\[-.2pt]{\fontsize{5.8}{6.4}\selectfont \textcolor{BrickRed}{$ \blacktriangledown\,45.4\times $}}\\[-.2pt]11m\,54.8s\\[-.2pt]{\fontsize{5.8}{6.4}\selectfont \textcolor{BrickRed}{$ \blacktriangledown\,56.1\times $}}\end{tabular} & \textcolor{gray}{\ensuremath{\cdot}} & \hspace*{1pt}\begin{tabular}[t]{@{}c@{}}\rule{0pt}{12pt}1m\,54.6s\\[-.2pt]{\fontsize{5.8}{6.4}\selectfont \textcolor{BrickRed}{$ \blacktriangledown\,2.7\times $}}\\[-.2pt]1m\,19s\\[-.2pt]{\fontsize{5.8}{6.4}\selectfont \textcolor{BrickRed}{$ \blacktriangledown\,1.9\times $}}\end{tabular} & \hspace*{1pt}\begin{tabular}[t]{@{}c@{}}\rule{0pt}{12pt}1m\,45.2s\\[-.2pt]{\fontsize{5.8}{6.4}\selectfont \textcolor{BrickRed}{$ \blacktriangledown\,2.5\times $}}\\[-.2pt]1m\,14.4s\\[-.2pt]{\fontsize{5.8}{6.4}\selectfont \textcolor{BrickRed}{$ \blacktriangledown\,1.8\times $}}\end{tabular} & \textcolor{gray}{\ensuremath{\cdot}} & \textcolor{gray}{\ensuremath{\cdot}} & \hspace*{1pt}\begin{tabular}[t]{@{}c@{}}\rule{0pt}{12pt}5m\,31.5s\\[-.2pt]{\fontsize{5.8}{6.4}\selectfont \textcolor{BrickRed}{$ \blacktriangledown\,9.5\times $}}\\[-.2pt]1m\,18.5s\\[-.2pt]{\fontsize{5.8}{6.4}\selectfont \textcolor{BrickRed}{$ \blacktriangledown\,2.3\times $}}\end{tabular} & \hspace*{1pt}\begin{tabular}[t]{@{}c@{}}\rule{0pt}{12pt}4m\,59.4s\\[-.2pt]{\fontsize{5.8}{6.4}\selectfont \textcolor{BrickRed}{$ \blacktriangledown\,8.3\times $}}\\[-.2pt]1m\,13.4s\\[-.2pt]{\fontsize{5.8}{6.4}\selectfont \textcolor{BrickRed}{$ \blacktriangledown\,2.1\times $}}\end{tabular} & \textcolor{gray}{\ensuremath{\cdot}} & \textcolor{gray}{\ensuremath{\cdot}}\\[3pt]\hline
\BenchmarkSourceLabel{bench:aws-2-k20-w15:row:mpf-twostack-rmq-ksw2}{R33}\shortstack[l]{{\color{gray}R33}} & \begin{tabular}[t]{@{}l@{}}\rule{0pt}{12pt}MPF / Two-stack\\RMQ / KSW2\end{tabular} & \hspace*{1pt}\begin{tabular}[t]{@{}c@{}}\rule{0pt}{12pt}4m\,34.2s\\[-.2pt]{\fontsize{5.8}{6.4}\selectfont \textcolor{BrickRed}{$ \blacktriangledown\,14.0\times $}}\\[-.2pt]4m\,32.2s\\[-.2pt]{\fontsize{5.8}{6.4}\selectfont \textcolor{BrickRed}{$ \blacktriangledown\,14.6\times $}}\end{tabular} & \hspace*{1pt}\begin{tabular}[t]{@{}c@{}}\rule{0pt}{12pt}6m\,52.3s\\[-.2pt]{\fontsize{5.8}{6.4}\selectfont \textcolor{BrickRed}{$ \blacktriangledown\,22.7\times $}}\\[-.2pt]6m\,49.6s\\[-.2pt]{\fontsize{5.8}{6.4}\selectfont \textcolor{BrickRed}{$ \blacktriangledown\,24.8\times $}}\end{tabular} & \hspace*{1pt}\begin{tabular}[t]{@{}c@{}}\rule{0pt}{12pt}17m\,29.3s\\[-.2pt]{\fontsize{5.8}{6.4}\selectfont \textcolor{BrickRed}{$ \blacktriangledown\,24.7\times $}}\\[-.2pt]17m\,23.1s\\[-.2pt]{\fontsize{5.8}{6.4}\selectfont \textcolor{BrickRed}{$ \blacktriangledown\,26.5\times $}}\end{tabular} & \textcolor{gray}{\ensuremath{\cdot}} & \hspace*{1pt}\begin{tabular}[t]{@{}c@{}}\rule{0pt}{12pt}3m\,23.9s\\[-.2pt]{\fontsize{5.8}{6.4}\selectfont \textcolor{BrickRed}{$ \blacktriangledown\,21.9\times $}}\\[-.2pt]3m\,21.3s\\[-.2pt]{\fontsize{5.8}{6.4}\selectfont \textcolor{BrickRed}{$ \blacktriangledown\,24.0\times $}}\end{tabular} & \hspace*{1pt}\begin{tabular}[t]{@{}c@{}}\rule{0pt}{12pt}5m\,43s\\[-.2pt]{\fontsize{5.8}{6.4}\selectfont \textcolor{BrickRed}{$ \blacktriangledown\,39.3\times $}}\\[-.2pt]5m\,39.6s\\[-.2pt]{\fontsize{5.8}{6.4}\selectfont \textcolor{BrickRed}{$ \blacktriangledown\,48.3\times $}}\end{tabular} & \hspace*{1pt}\begin{tabular}[t]{@{}c@{}}\rule{0pt}{12pt}12m\,2.7s\\[-.2pt]{\fontsize{5.8}{6.4}\selectfont \textcolor{BrickRed}{$ \blacktriangledown\,45.7\times $}}\\[-.2pt]11m\,55.6s\\[-.2pt]{\fontsize{5.8}{6.4}\selectfont \textcolor{BrickRed}{$ \blacktriangledown\,56.2\times $}}\end{tabular} & \textcolor{gray}{\ensuremath{\cdot}} & \hspace*{1pt}\begin{tabular}[t]{@{}c@{}}\rule{0pt}{12pt}1m\,59.6s\\[-.2pt]{\fontsize{5.8}{6.4}\selectfont \textcolor{BrickRed}{$ \blacktriangledown\,2.8\times $}}\\[-.2pt]1m\,19.6s\\[-.2pt]{\fontsize{5.8}{6.4}\selectfont \textcolor{BrickRed}{$ \blacktriangledown\,1.9\times $}}\end{tabular} & \hspace*{1pt}\begin{tabular}[t]{@{}c@{}}\rule{0pt}{12pt}1m\,48.5s\\[-.2pt]{\fontsize{5.8}{6.4}\selectfont \textcolor{BrickRed}{$ \blacktriangledown\,2.5\times $}}\\[-.2pt]1m\,14.5s\\[-.2pt]{\fontsize{5.8}{6.4}\selectfont \textcolor{BrickRed}{$ \blacktriangledown\,1.8\times $}}\end{tabular} & \textcolor{gray}{\ensuremath{\cdot}} & \textcolor{gray}{\ensuremath{\cdot}} & \hspace*{1pt}\begin{tabular}[t]{@{}c@{}}\rule{0pt}{12pt}5m\,35.5s\\[-.2pt]{\fontsize{5.8}{6.4}\selectfont \textcolor{BrickRed}{$ \blacktriangledown\,9.6\times $}}\\[-.2pt]1m\,17.7s\\[-.2pt]{\fontsize{5.8}{6.4}\selectfont \textcolor{BrickRed}{$ \blacktriangledown\,2.3\times $}}\end{tabular} & \hspace*{1pt}\begin{tabular}[t]{@{}c@{}}\rule{0pt}{12pt}5m\,2s\\[-.2pt]{\fontsize{5.8}{6.4}\selectfont \textcolor{BrickRed}{$ \blacktriangledown\,8.4\times $}}\\[-.2pt]1m\,12.5s\\[-.2pt]{\fontsize{5.8}{6.4}\selectfont \textcolor{BrickRed}{$ \blacktriangledown\,2.1\times $}}\end{tabular} & \textcolor{gray}{\ensuremath{\cdot}} & \textcolor{gray}{\ensuremath{\cdot}}\\[3pt]\hline
\BenchmarkSourceLabel{bench:aws-2-k20-w15:row:mpf-twostack-simd-block}{R34}\shortstack[l]{{\color{gray}R34}} & \begin{tabular}[t]{@{}l@{}}\rule{0pt}{12pt}MPF / Two-stack\\SIMD-DP / Block\end{tabular} & \hspace*{1pt}\begin{tabular}[t]{@{}c@{}}\rule{0pt}{12pt}4m\,36.1s\\[-.2pt]{\fontsize{5.8}{6.4}\selectfont \textcolor{BrickRed}{$ \blacktriangledown\,14.1\times $}}\\[-.2pt]4m\,35.5s\\[-.2pt]{\fontsize{5.8}{6.4}\selectfont \textcolor{BrickRed}{$ \blacktriangledown\,14.8\times $}}\end{tabular} & \hspace*{1pt}\begin{tabular}[t]{@{}c@{}}\rule{0pt}{12pt}7m\,3.9s\\[-.2pt]{\fontsize{5.8}{6.4}\selectfont \textcolor{BrickRed}{$ \blacktriangledown\,23.3\times $}}\\[-.2pt]7m\,3s\\[-.2pt]{\fontsize{5.8}{6.4}\selectfont \textcolor{BrickRed}{$ \blacktriangledown\,25.6\times $}}\end{tabular} & \hspace*{1pt}\begin{tabular}[t]{@{}c@{}}\rule{0pt}{12pt}17m\,28.7s\\[-.2pt]{\fontsize{5.8}{6.4}\selectfont \textcolor{BrickRed}{$ \blacktriangledown\,24.7\times $}}\\[-.2pt]17m\,27.2s\\[-.2pt]{\fontsize{5.8}{6.4}\selectfont \textcolor{BrickRed}{$ \blacktriangledown\,26.6\times $}}\end{tabular} & \hspace*{1pt}\begin{tabular}[t]{@{}c@{}}\rule{0pt}{12pt}5m\,27s\\[-.2pt]{\fontsize{5.8}{6.4}\selectfont \textcolor{ForestGreen}{$ \blacktriangle\,16\% $}}\\[-.2pt]5m\,26s\\[-.2pt]{\fontsize{5.8}{6.4}\selectfont \textcolor{ForestGreen}{$ \blacktriangle\,9\% $}}\end{tabular} & \hspace*{1pt}\begin{tabular}[t]{@{}c@{}}\rule{0pt}{12pt}3m\,24.2s\\[-.2pt]{\fontsize{5.8}{6.4}\selectfont \textcolor{BrickRed}{$ \blacktriangledown\,21.9\times $}}\\[-.2pt]3m\,22.7s\\[-.2pt]{\fontsize{5.8}{6.4}\selectfont \textcolor{BrickRed}{$ \blacktriangledown\,24.2\times $}}\end{tabular} & \hspace*{1pt}\begin{tabular}[t]{@{}c@{}}\rule{0pt}{12pt}5m\,44.1s\\[-.2pt]{\fontsize{5.8}{6.4}\selectfont \textcolor{BrickRed}{$ \blacktriangledown\,39.4\times $}}\\[-.2pt]5m\,42.7s\\[-.2pt]{\fontsize{5.8}{6.4}\selectfont \textcolor{BrickRed}{$ \blacktriangledown\,48.8\times $}}\end{tabular} & \hspace*{1pt}\begin{tabular}[t]{@{}c@{}}\rule{0pt}{12pt}11m\,58.9s\\[-.2pt]{\fontsize{5.8}{6.4}\selectfont \textcolor{BrickRed}{$ \blacktriangledown\,45.5\times $}}\\[-.2pt]11m\,56.2s\\[-.2pt]{\fontsize{5.8}{6.4}\selectfont \textcolor{BrickRed}{$ \blacktriangledown\,56.2\times $}}\end{tabular} & \hspace*{1pt}\begin{tabular}[t]{@{}c@{}}\rule{0pt}{12pt}4m\,0.4s\\[-.2pt]{\fontsize{5.8}{6.4}\selectfont \textcolor{BrickRed}{$ \blacktriangledown\,13.0\times $}}\\[-.2pt]3m\,59s\\[-.2pt]{\fontsize{5.8}{6.4}\selectfont \textcolor{BrickRed}{$ \blacktriangledown\,14.0\times $}}\end{tabular} & \hspace*{1pt}\begin{tabular}[t]{@{}c@{}}\rule{0pt}{12pt}1m\,57.3s\\[-.2pt]{\fontsize{5.8}{6.4}\selectfont \textcolor{BrickRed}{$ \blacktriangledown\,2.8\times $}}\\[-.2pt]1m\,21.6s\\[-.2pt]{\fontsize{5.8}{6.4}\selectfont \textcolor{BrickRed}{$ \blacktriangledown\,2.0\times $}}\end{tabular} & \hspace*{1pt}\begin{tabular}[t]{@{}c@{}}\rule{0pt}{12pt}1m\,45.1s\\[-.2pt]{\fontsize{5.8}{6.4}\selectfont \textcolor{BrickRed}{$ \blacktriangledown\,2.5\times $}}\\[-.2pt]1m\,15.1s\\[-.2pt]{\fontsize{5.8}{6.4}\selectfont \textcolor{BrickRed}{$ \blacktriangledown\,1.8\times $}}\end{tabular} & \textcolor{gray}{\ensuremath{\cdot}} & \textcolor{gray}{\ensuremath{\cdot}} & \hspace*{1pt}\begin{tabular}[t]{@{}c@{}}\rule{0pt}{12pt}5m\,30.7s\\[-.2pt]{\fontsize{5.8}{6.4}\selectfont \textcolor{BrickRed}{$ \blacktriangledown\,9.5\times $}}\\[-.2pt]1m\,17.9s\\[-.2pt]{\fontsize{5.8}{6.4}\selectfont \textcolor{BrickRed}{$ \blacktriangledown\,2.3\times $}}\end{tabular} & \hspace*{1pt}\begin{tabular}[t]{@{}c@{}}\rule{0pt}{12pt}4m\,57.2s\\[-.2pt]{\fontsize{5.8}{6.4}\selectfont \textcolor{BrickRed}{$ \blacktriangledown\,8.3\times $}}\\[-.2pt]1m\,13.2s\\[-.2pt]{\fontsize{5.8}{6.4}\selectfont \textcolor{BrickRed}{$ \blacktriangledown\,2.1\times $}}\end{tabular} & \textcolor{gray}{\ensuremath{\cdot}} & \textcolor{gray}{\ensuremath{\cdot}}\\[3pt]\hline
\BenchmarkSourceLabel{bench:aws-2-k20-w15:row:mpf-twostack-simd-ksw2}{R35}\shortstack[l]{{\color{gray}R35}} & \begin{tabular}[t]{@{}l@{}}\rule{0pt}{12pt}MPF / Two-stack\\SIMD-DP / KSW2\end{tabular} & \hspace*{1pt}\begin{tabular}[t]{@{}c@{}}\rule{0pt}{12pt}4m\,37.8s\\[-.2pt]{\fontsize{5.8}{6.4}\selectfont \textcolor{BrickRed}{$ \blacktriangledown\,14.2\times $}}\\[-.2pt]4m\,35.9s\\[-.2pt]{\fontsize{5.8}{6.4}\selectfont \textcolor{BrickRed}{$ \blacktriangledown\,14.8\times $}}\end{tabular} & \hspace*{1pt}\begin{tabular}[t]{@{}c@{}}\rule{0pt}{12pt}7m\,37.6s\\[-.2pt]{\fontsize{5.8}{6.4}\selectfont \textcolor{BrickRed}{$ \blacktriangledown\,25.1\times $}}\\[-.2pt]7m\,34.4s\\[-.2pt]{\fontsize{5.8}{6.4}\selectfont \textcolor{BrickRed}{$ \blacktriangledown\,27.5\times $}}\end{tabular} & \hspace*{1pt}\begin{tabular}[t]{@{}c@{}}\rule{0pt}{12pt}17m\,30.1s\\[-.2pt]{\fontsize{5.8}{6.4}\selectfont \textcolor{BrickRed}{$ \blacktriangledown\,24.7\times $}}\\[-.2pt]17m\,23.9s\\[-.2pt]{\fontsize{5.8}{6.4}\selectfont \textcolor{BrickRed}{$ \blacktriangledown\,26.5\times $}}\end{tabular} & \textcolor{gray}{\ensuremath{\cdot}} & \hspace*{1pt}\begin{tabular}[t]{@{}c@{}}\rule{0pt}{12pt}3m\,24.7s\\[-.2pt]{\fontsize{5.8}{6.4}\selectfont \textcolor{BrickRed}{$ \blacktriangledown\,22.0\times $}}\\[-.2pt]3m\,22.1s\\[-.2pt]{\fontsize{5.8}{6.4}\selectfont \textcolor{BrickRed}{$ \blacktriangledown\,24.1\times $}}\end{tabular} & \hspace*{1pt}\begin{tabular}[t]{@{}c@{}}\rule{0pt}{12pt}5m\,44.4s\\[-.2pt]{\fontsize{5.8}{6.4}\selectfont \textcolor{BrickRed}{$ \blacktriangledown\,39.4\times $}}\\[-.2pt]5m\,40.9s\\[-.2pt]{\fontsize{5.8}{6.4}\selectfont \textcolor{BrickRed}{$ \blacktriangledown\,48.5\times $}}\end{tabular} & \hspace*{1pt}\begin{tabular}[t]{@{}c@{}}\rule{0pt}{12pt}12m\,4s\\[-.2pt]{\fontsize{5.8}{6.4}\selectfont \textcolor{BrickRed}{$ \blacktriangledown\,45.8\times $}}\\[-.2pt]11m\,56.8s\\[-.2pt]{\fontsize{5.8}{6.4}\selectfont \textcolor{BrickRed}{$ \blacktriangledown\,56.3\times $}}\end{tabular} & \textcolor{gray}{\ensuremath{\cdot}} & \hspace*{1pt}\begin{tabular}[t]{@{}c@{}}\rule{0pt}{12pt}1m\,58.9s\\[-.2pt]{\fontsize{5.8}{6.4}\selectfont \textcolor{BrickRed}{$ \blacktriangledown\,2.8\times $}}\\[-.2pt]1m\,18.9s\\[-.2pt]{\fontsize{5.8}{6.4}\selectfont \textcolor{BrickRed}{$ \blacktriangledown\,1.9\times $}}\end{tabular} & \hspace*{1pt}\begin{tabular}[t]{@{}c@{}}\rule{0pt}{12pt}1m\,47.5s\\[-.2pt]{\fontsize{5.8}{6.4}\selectfont \textcolor{BrickRed}{$ \blacktriangledown\,2.5\times $}}\\[-.2pt]1m\,13.8s\\[-.2pt]{\fontsize{5.8}{6.4}\selectfont \textcolor{BrickRed}{$ \blacktriangledown\,1.8\times $}}\end{tabular} & \textcolor{gray}{\ensuremath{\cdot}} & \textcolor{gray}{\ensuremath{\cdot}} & \hspace*{1pt}\begin{tabular}[t]{@{}c@{}}\rule{0pt}{12pt}5m\,34.8s\\[-.2pt]{\fontsize{5.8}{6.4}\selectfont \textcolor{BrickRed}{$ \blacktriangledown\,9.6\times $}}\\[-.2pt]1m\,17.6s\\[-.2pt]{\fontsize{5.8}{6.4}\selectfont \textcolor{BrickRed}{$ \blacktriangledown\,2.3\times $}}\end{tabular} & \hspace*{1pt}\begin{tabular}[t]{@{}c@{}}\rule{0pt}{12pt}5m\,3.8s\\[-.2pt]{\fontsize{5.8}{6.4}\selectfont \textcolor{BrickRed}{$ \blacktriangledown\,8.5\times $}}\\[-.2pt]1m\,12.8s\\[-.2pt]{\fontsize{5.8}{6.4}\selectfont \textcolor{BrickRed}{$ \blacktriangledown\,2.1\times $}}\end{tabular} & \textcolor{gray}{\ensuremath{\cdot}} & \textcolor{gray}{\ensuremath{\cdot}}\\[3pt]\hline
\BenchmarkSourceLabel{bench:aws-2-k20-w15:row:mpf-twostack-scalar-ksw2}{R36}\shortstack[l]{{\color{gray}R36}} & \begin{tabular}[t]{@{}l@{}}\rule{0pt}{12pt}MPF / Two-stack\\Scalar-DP / KSW2\end{tabular} & \hspace*{1pt}\begin{tabular}[t]{@{}c@{}}\rule{0pt}{12pt}4m\,38.1s\\[-.2pt]{\fontsize{5.8}{6.4}\selectfont \textcolor{BrickRed}{$ \blacktriangledown\,14.2\times $}}\\[-.2pt]4m\,36.2s\\[-.2pt]{\fontsize{5.8}{6.4}\selectfont \textcolor{BrickRed}{$ \blacktriangledown\,14.8\times $}}\end{tabular} & \hspace*{1pt}\begin{tabular}[t]{@{}c@{}}\rule{0pt}{12pt}7m\,34.7s\\[-.2pt]{\fontsize{5.8}{6.4}\selectfont \textcolor{BrickRed}{$ \blacktriangledown\,25.0\times $}}\\[-.2pt]7m\,31.6s\\[-.2pt]{\fontsize{5.8}{6.4}\selectfont \textcolor{BrickRed}{$ \blacktriangledown\,27.3\times $}}\end{tabular} & \hspace*{1pt}\begin{tabular}[t]{@{}c@{}}\rule{0pt}{12pt}17m\,28.2s\\[-.2pt]{\fontsize{5.8}{6.4}\selectfont \textcolor{BrickRed}{$ \blacktriangledown\,24.7\times $}}\\[-.2pt]17m\,22.1s\\[-.2pt]{\fontsize{5.8}{6.4}\selectfont \textcolor{BrickRed}{$ \blacktriangledown\,26.5\times $}}\end{tabular} & \textcolor{gray}{\ensuremath{\cdot}} & \hspace*{1pt}\begin{tabular}[t]{@{}c@{}}\rule{0pt}{12pt}3m\,24.7s\\[-.2pt]{\fontsize{5.8}{6.4}\selectfont \textcolor{BrickRed}{$ \blacktriangledown\,22.0\times $}}\\[-.2pt]3m\,22.1s\\[-.2pt]{\fontsize{5.8}{6.4}\selectfont \textcolor{BrickRed}{$ \blacktriangledown\,24.1\times $}}\end{tabular} & \hspace*{1pt}\begin{tabular}[t]{@{}c@{}}\rule{0pt}{12pt}5m\,43.5s\\[-.2pt]{\fontsize{5.8}{6.4}\selectfont \textcolor{BrickRed}{$ \blacktriangledown\,39.3\times $}}\\[-.2pt]5m\,40.1s\\[-.2pt]{\fontsize{5.8}{6.4}\selectfont \textcolor{BrickRed}{$ \blacktriangledown\,48.4\times $}}\end{tabular} & \hspace*{1pt}\begin{tabular}[t]{@{}c@{}}\rule{0pt}{12pt}12m\,2.2s\\[-.2pt]{\fontsize{5.8}{6.4}\selectfont \textcolor{BrickRed}{$ \blacktriangledown\,45.7\times $}}\\[-.2pt]11m\,55s\\[-.2pt]{\fontsize{5.8}{6.4}\selectfont \textcolor{BrickRed}{$ \blacktriangledown\,56.1\times $}}\end{tabular} & \textcolor{gray}{\ensuremath{\cdot}} & \hspace*{1pt}\begin{tabular}[t]{@{}c@{}}\rule{0pt}{12pt}1m\,59.1s\\[-.2pt]{\fontsize{5.8}{6.4}\selectfont \textcolor{BrickRed}{$ \blacktriangledown\,2.8\times $}}\\[-.2pt]1m\,19.3s\\[-.2pt]{\fontsize{5.8}{6.4}\selectfont \textcolor{BrickRed}{$ \blacktriangledown\,1.9\times $}}\end{tabular} & \hspace*{1pt}\begin{tabular}[t]{@{}c@{}}\rule{0pt}{12pt}1m\,47.1s\\[-.2pt]{\fontsize{5.8}{6.4}\selectfont \textcolor{BrickRed}{$ \blacktriangledown\,2.5\times $}}\\[-.2pt]1m\,13.3s\\[-.2pt]{\fontsize{5.8}{6.4}\selectfont \textcolor{BrickRed}{$ \blacktriangledown\,1.8\times $}}\end{tabular} & \textcolor{gray}{\ensuremath{\cdot}} & \textcolor{gray}{\ensuremath{\cdot}} & \hspace*{1pt}\begin{tabular}[t]{@{}c@{}}\rule{0pt}{12pt}5m\,34.9s\\[-.2pt]{\fontsize{5.8}{6.4}\selectfont \textcolor{BrickRed}{$ \blacktriangledown\,9.6\times $}}\\[-.2pt]1m\,17.2s\\[-.2pt]{\fontsize{5.8}{6.4}\selectfont \textcolor{BrickRed}{$ \blacktriangledown\,2.3\times $}}\end{tabular} & \hspace*{1pt}\begin{tabular}[t]{@{}c@{}}\rule{0pt}{12pt}4m\,58.5s\\[-.2pt]{\fontsize{5.8}{6.4}\selectfont \textcolor{BrickRed}{$ \blacktriangledown\,8.3\times $}}\\[-.2pt]1m\,10.5s\\[-.2pt]{\fontsize{5.8}{6.4}\selectfont \textcolor{BrickRed}{$ \blacktriangledown\,2.1\times $}}\end{tabular} & \textcolor{gray}{\ensuremath{\cdot}} & \textcolor{gray}{\ensuremath{\cdot}}\\[3pt]\hline
\end{tabular}\par}
\smallskip{\footnotesize\noindent Each cell: A and B times, $\cdot$: no measurement. ---\textsuperscript{g}: GPU unavailable. n/a: no comparison.\par
\noindent Blue: one lowest A and one lowest B measurement for each column for each configuration.\par
\noindent Gray: minimap2. Yellow: my preferred MPF / Bounded / SIMD-DP / Block configuration.\par
\noindent MPF/MPS: minimizers per fragment/sequence. SIMD-ref: naive SIMD scan. Bounded: bounded SIMD algorithm.\par
}
\clearpage
\noindent{\large\bfseries Table \thetable: (E): AWS: 2 Gbp, k=20, w=15\hfill 6/7}\par\smallskip
{\fontsize{7.1}{8.1}\selectfont\setlength{\tabcolsep}{1pt}\setlength{\fboxsep}{2pt}\renewcommand{\arraystretch}{1.12}
\noindent\begin{tabular}{@{}p{14mm}p{27mm}|p{13.55mm}|p{13.55mm}|p{13.55mm}|p{13.55mm}|p{13.55mm}|p{13.55mm}|p{13.55mm}|p{13.55mm}|p{13.55mm}|p{13.55mm}|p{13.55mm}|p{13.55mm}|p{13.55mm}|p{13.55mm}|p{13.55mm}|p{13.55mm}|@{}}\hline
Engine & Execution & \multicolumn{8}{c|}{\textbf{Index}} & \multicolumn{8}{c|}{\textbf{No index}}\\
 &  & \multicolumn{8}{c|}{\textbf{2 Gbp}} & \multicolumn{8}{c|}{\textbf{2 Gbp}}\\
 &  & \multicolumn{8}{c|}{\textbf{k=20, w=15}} & \multicolumn{8}{c|}{\textbf{k=20, w=15}}\\
 &  & \multicolumn{4}{c|}{\textbf{DNA4}} & \multicolumn{4}{c|}{\textbf{DNA5}} & \multicolumn{4}{c|}{\textbf{DNA4}} & \multicolumn{4}{c|}{\textbf{DNA5}}\\
 &  & \multicolumn{1}{c|}{\textbf{M7A}} & \multicolumn{1}{c|}{\textbf{M7I}} & \multicolumn{1}{c|}{\textbf{M4}} & \multicolumn{1}{c|}{\textbf{C7G}} & \multicolumn{1}{c|}{\textbf{M7A}} & \multicolumn{1}{c|}{\textbf{M7I}} & \multicolumn{1}{c|}{\textbf{M4}} & \multicolumn{1}{c|}{\textbf{C7G}} & \multicolumn{1}{c|}{\textbf{M7A}} & \multicolumn{1}{c|}{\textbf{M7I}} & \multicolumn{1}{c|}{\textbf{M4}} & \multicolumn{1}{c|}{\textbf{C7G}} & \multicolumn{1}{c|}{\textbf{M7A}} & \multicolumn{1}{c|}{\textbf{M7I}} & \multicolumn{1}{c|}{\textbf{M4}} & \multicolumn{1}{c|}{\textbf{C7G}}\\
 & & {\color{gray}C1} & {\color{gray}C2} & {\color{gray}C3} & {\color{gray}C4} & {\color{gray}C5} & {\color{gray}C6} & {\color{gray}C7} & {\color{gray}C8} & {\color{gray}C9} & {\color{gray}C10} & {\color{gray}C11} & {\color{gray}C12} & {\color{gray}C13} & {\color{gray}C14} & {\color{gray}C15} & {\color{gray}C16}\\
\hline
\rowcolor{BenchBaseline}[0.5pt][0.5pt]\shortstack[l]{{\color{gray}R1}\\\textbf{minimap2}} & \textbf{Baseline} & \hspace*{1pt}\begin{tabular}[t]{@{}c@{}}\rule{0pt}{12pt}19.6s\\[-.2pt]{\fontsize{5.8}{6.4}\selectfont \textcolor{BrickRed}{$ \blacktriangledown\,6.7\times $}}\\[-.2pt]18.6s\\[-.2pt]{\fontsize{5.8}{6.4}\selectfont \textcolor{BrickRed}{$ \blacktriangledown\,8.4\times $}}\end{tabular} & \hspace*{1pt}\begin{tabular}[t]{@{}c@{}}\rule{0pt}{12pt}18.2s\\[-.2pt]{\fontsize{5.8}{6.4}\selectfont \textcolor{BrickRed}{$ \blacktriangledown\,6.9\times $}}\\[-.2pt]16.5s\\[-.2pt]{\fontsize{5.8}{6.4}\selectfont \textcolor{BrickRed}{$ \blacktriangledown\,8.6\times $}}\end{tabular} & \hspace*{1pt}\begin{tabular}[t]{@{}c@{}}\rule{0pt}{12pt}42.5s\\[-.2pt]{\fontsize{5.8}{6.4}\selectfont \textcolor{BrickRed}{$ \blacktriangledown\,8.2\times $}}\\[-.2pt]39.4s\\[-.2pt]{\fontsize{5.8}{6.4}\selectfont \textcolor{BrickRed}{$ \blacktriangledown\,12.6\times $}}\end{tabular} & \hspace*{1pt}\begin{tabular}[t]{@{}c@{}}\rule{0pt}{12pt}6m\,29.2s\\[-.2pt]{\fontsize{5.8}{6.4}\selectfont \textcolor{BrickRed}{$ \blacktriangledown\,92.8\times $}}\\[-.2pt]5m\,56.9s\\[-.2pt]{\fontsize{5.8}{6.4}\selectfont \textcolor{BrickRed}{$ \blacktriangledown\,184.9\times $}}\end{tabular} & \hspace*{1pt}\begin{tabular}[t]{@{}c@{}}\rule{0pt}{12pt}9.3s\\[-.2pt]{\fontsize{5.8}{6.4}\selectfont \textcolor{BrickRed}{$ \blacktriangledown\,24\% $}}\\[-.2pt]8.4s\\[-.2pt]{\fontsize{5.8}{6.4}\selectfont \textcolor{BrickRed}{$ \blacktriangledown\,1.6\times $}}\end{tabular} & \hspace*{1pt}\begin{tabular}[t]{@{}c@{}}\rule{0pt}{12pt}8.7s\\[-.2pt]{\fontsize{5.8}{6.4}\selectfont \textcolor{BrickRed}{$ \blacktriangledown\,1.5\times $}}\\[-.2pt]7s\\[-.2pt]{\fontsize{5.8}{6.4}\selectfont \textcolor{BrickRed}{$ \blacktriangledown\,1.7\times $}}\end{tabular} & \hspace*{1pt}\begin{tabular}[t]{@{}c@{}}\rule{0pt}{12pt}15.8s\\[-.2pt]{\fontsize{5.8}{6.4}\selectfont \textcolor{BrickRed}{$ \blacktriangledown\,26\% $}}\\[-.2pt]12.7s\\[-.2pt]{\fontsize{5.8}{6.4}\selectfont \textcolor{BrickRed}{$ \blacktriangledown\,1.7\times $}}\end{tabular} & \hspace*{1pt}\begin{tabular}[t]{@{}c@{}}\rule{0pt}{12pt}18.5s\\[-.2pt]{\fontsize{5.8}{6.4}\selectfont \textcolor{BrickRed}{$ \blacktriangledown\,1.6\times $}}\\[-.2pt]17s\\[-.2pt]{\fontsize{5.8}{6.4}\selectfont \textcolor{BrickRed}{$ \blacktriangledown\,1.8\times $}}\end{tabular} & \hspace*{1pt}\begin{tabular}[t]{@{}c@{}}\rule{0pt}{12pt}42.4s\\[-.2pt]{\fontsize{5.8}{6.4}\selectfont \textcolor{ForestGreen}{$ \blacktriangle\,2.0\times $}}\\[-.2pt]41.4s\\[-.2pt]{\fontsize{5.8}{6.4}\selectfont \textcolor{ForestGreen}{$ \blacktriangle\,1.9\times $}}\end{tabular} & \hspace*{1pt}\begin{tabular}[t]{@{}c@{}}\rule{0pt}{12pt}42.9s\\[-.2pt]{\fontsize{5.8}{6.4}\selectfont \textcolor{ForestGreen}{$ \blacktriangle\,2.0\times $}}\\[-.2pt]41.1s\\[-.2pt]{\fontsize{5.8}{6.4}\selectfont \textcolor{ForestGreen}{$ \blacktriangle\,1.8\times $}}\end{tabular} & \hspace*{1pt}\begin{tabular}[t]{@{}c@{}}\rule{0pt}{12pt}1m\,26.9s\\[-.2pt]{\fontsize{5.8}{6.4}\selectfont \textcolor{ForestGreen}{$ \blacktriangle\,1.9\times $}}\\[-.2pt]1m\,23.7s\\[-.2pt]{\fontsize{5.8}{6.4}\selectfont \textcolor{ForestGreen}{$ \blacktriangle\,1.7\times $}}\end{tabular} & \hspace*{1pt}\begin{tabular}[t]{@{}c@{}}\rule{0pt}{12pt}4m\,51.6s\\[-.2pt]{\fontsize{5.8}{6.4}\selectfont \textcolor{BrickRed}{$ \blacktriangledown\,2.9\times $}}\\[-.2pt]4m\,46.6s\\[-.2pt]{\fontsize{5.8}{6.4}\selectfont \textcolor{BrickRed}{$ \blacktriangledown\,3.0\times $}}\end{tabular} & \hspace*{1pt}\begin{tabular}[t]{@{}c@{}}\rule{0pt}{12pt}35s\\[-.2pt]{\fontsize{5.8}{6.4}\selectfont \textcolor{ForestGreen}{$ \blacktriangle\,8.6\times $}}\\[-.2pt]34s\\[-.2pt]{\fontsize{5.8}{6.4}\selectfont \textcolor{ForestGreen}{$ \blacktriangle\,2.3\times $}}\end{tabular} & \hspace*{1pt}\begin{tabular}[t]{@{}c@{}}\rule{0pt}{12pt}35.9s\\[-.2pt]{\fontsize{5.8}{6.4}\selectfont \textcolor{ForestGreen}{$ \blacktriangle\,6.7\times $}}\\[-.2pt]34.2s\\[-.2pt]{\fontsize{5.8}{6.4}\selectfont \textcolor{ForestGreen}{$ \blacktriangle\,2.1\times $}}\end{tabular} & \hspace*{1pt}\begin{tabular}[t]{@{}c@{}}\rule{0pt}{12pt}1m\,14.1s\\[-.2pt]{\fontsize{5.8}{6.4}\selectfont \textcolor{ForestGreen}{$ \blacktriangle\,6.9\times $}}\\[-.2pt]1m\,11.1s\\[-.2pt]{\fontsize{5.8}{6.4}\selectfont \textcolor{ForestGreen}{$ \blacktriangle\,2.2\times $}}\end{tabular} & \hspace*{1pt}\begin{tabular}[t]{@{}c@{}}\rule{0pt}{12pt}1m\,0.1s\\[-.2pt]{\fontsize{5.8}{6.4}\selectfont \textcolor{ForestGreen}{$ \blacktriangle\,10.4\times $}}\\[-.2pt]58.6s\\[-.2pt]{\fontsize{5.8}{6.4}\selectfont \textcolor{ForestGreen}{$ \blacktriangle\,1.8\times $}}\end{tabular}\\[3pt]\hline
\BenchmarkSourceLabel{bench:aws-2-k20-w15:row:mpf-twostack-scalar-block}{R37}\shortstack[l]{{\color{gray}R37}\\\textbf{BioseqDB}} & \begin{tabular}[t]{@{}l@{}}\rule{0pt}{12pt}MPF / Two-stack\\Scalar-DP / Block\end{tabular} & \hspace*{1pt}\begin{tabular}[t]{@{}c@{}}\rule{0pt}{12pt}4m\,36s\\[-.2pt]{\fontsize{5.8}{6.4}\selectfont \textcolor{BrickRed}{$ \blacktriangledown\,14.1\times $}}\\[-.2pt]4m\,35.5s\\[-.2pt]{\fontsize{5.8}{6.4}\selectfont \textcolor{BrickRed}{$ \blacktriangledown\,14.8\times $}}\end{tabular} & \hspace*{1pt}\begin{tabular}[t]{@{}c@{}}\rule{0pt}{12pt}6m\,53.8s\\[-.2pt]{\fontsize{5.8}{6.4}\selectfont \textcolor{BrickRed}{$ \blacktriangledown\,22.7\times $}}\\[-.2pt]6m\,52.9s\\[-.2pt]{\fontsize{5.8}{6.4}\selectfont \textcolor{BrickRed}{$ \blacktriangledown\,25.0\times $}}\end{tabular} & \hspace*{1pt}\begin{tabular}[t]{@{}c@{}}\rule{0pt}{12pt}17m\,25.1s\\[-.2pt]{\fontsize{5.8}{6.4}\selectfont \textcolor{BrickRed}{$ \blacktriangledown\,24.6\times $}}\\[-.2pt]17m\,23.5s\\[-.2pt]{\fontsize{5.8}{6.4}\selectfont \textcolor{BrickRed}{$ \blacktriangledown\,26.5\times $}}\end{tabular} & \textcolor{gray}{\ensuremath{\cdot}} & \hspace*{1pt}\begin{tabular}[t]{@{}c@{}}\rule{0pt}{12pt}3m\,24.8s\\[-.2pt]{\fontsize{5.8}{6.4}\selectfont \textcolor{BrickRed}{$ \blacktriangledown\,22.0\times $}}\\[-.2pt]3m\,23.4s\\[-.2pt]{\fontsize{5.8}{6.4}\selectfont \textcolor{BrickRed}{$ \blacktriangledown\,24.3\times $}}\end{tabular} & \hspace*{1pt}\begin{tabular}[t]{@{}c@{}}\rule{0pt}{12pt}5m\,38s\\[-.2pt]{\fontsize{5.8}{6.4}\selectfont \textcolor{BrickRed}{$ \blacktriangledown\,38.7\times $}}\\[-.2pt]5m\,36.6s\\[-.2pt]{\fontsize{5.8}{6.4}\selectfont \textcolor{BrickRed}{$ \blacktriangledown\,47.9\times $}}\end{tabular} & \hspace*{1pt}\begin{tabular}[t]{@{}c@{}}\rule{0pt}{12pt}11m\,57.4s\\[-.2pt]{\fontsize{5.8}{6.4}\selectfont \textcolor{BrickRed}{$ \blacktriangledown\,45.4\times $}}\\[-.2pt]11m\,54.9s\\[-.2pt]{\fontsize{5.8}{6.4}\selectfont \textcolor{BrickRed}{$ \blacktriangledown\,56.1\times $}}\end{tabular} & \textcolor{gray}{\ensuremath{\cdot}} & \hspace*{1pt}\begin{tabular}[t]{@{}c@{}}\rule{0pt}{12pt}1m\,55.2s\\[-.2pt]{\fontsize{5.8}{6.4}\selectfont \textcolor{BrickRed}{$ \blacktriangledown\,2.7\times $}}\\[-.2pt]1m\,18.8s\\[-.2pt]{\fontsize{5.8}{6.4}\selectfont \textcolor{BrickRed}{$ \blacktriangledown\,1.9\times $}}\end{tabular} & \hspace*{1pt}\begin{tabular}[t]{@{}c@{}}\rule{0pt}{12pt}1m\,43.5s\\[-.2pt]{\fontsize{5.8}{6.4}\selectfont \textcolor{BrickRed}{$ \blacktriangledown\,2.4\times $}}\\[-.2pt]1m\,13.7s\\[-.2pt]{\fontsize{5.8}{6.4}\selectfont \textcolor{BrickRed}{$ \blacktriangledown\,1.8\times $}}\end{tabular} & \textcolor{gray}{\ensuremath{\cdot}} & \textcolor{gray}{\ensuremath{\cdot}} & \hspace*{1pt}\begin{tabular}[t]{@{}c@{}}\rule{0pt}{12pt}5m\,30.7s\\[-.2pt]{\fontsize{5.8}{6.4}\selectfont \textcolor{BrickRed}{$ \blacktriangledown\,9.5\times $}}\\[-.2pt]1m\,17.1s\\[-.2pt]{\fontsize{5.8}{6.4}\selectfont \textcolor{BrickRed}{$ \blacktriangledown\,2.3\times $}}\end{tabular} & \hspace*{1pt}\begin{tabular}[t]{@{}c@{}}\rule{0pt}{12pt}4m\,55s\\[-.2pt]{\fontsize{5.8}{6.4}\selectfont \textcolor{BrickRed}{$ \blacktriangledown\,8.2\times $}}\\[-.2pt]1m\,11.7s\\[-.2pt]{\fontsize{5.8}{6.4}\selectfont \textcolor{BrickRed}{$ \blacktriangledown\,2.1\times $}}\end{tabular} & \textcolor{gray}{\ensuremath{\cdot}} & \textcolor{gray}{\ensuremath{\cdot}}\\[3pt]\hline
\BenchmarkSourceLabel{bench:aws-2-k20-w15:row:mpf-scalar-scalar-block}{R38}\shortstack[l]{{\color{gray}R38}} & \begin{tabular}[t]{@{}l@{}}\rule{0pt}{12pt}MPF / Scalar\\Scalar-DP / Block\end{tabular} & \hspace*{1pt}\begin{tabular}[t]{@{}c@{}}\rule{0pt}{12pt}5m\,39.5s\\[-.2pt]{\fontsize{5.8}{6.4}\selectfont \textcolor{BrickRed}{$ \blacktriangledown\,17.3\times $}}\\[-.2pt]5m\,39s\\[-.2pt]{\fontsize{5.8}{6.4}\selectfont \textcolor{BrickRed}{$ \blacktriangledown\,18.2\times $}}\end{tabular} & \hspace*{1pt}\begin{tabular}[t]{@{}c@{}}\rule{0pt}{12pt}12m\,20.4s\\[-.2pt]{\fontsize{5.8}{6.4}\selectfont \textcolor{BrickRed}{$ \blacktriangledown\,40.7\times $}}\\[-.2pt]12m\,19.8s\\[-.2pt]{\fontsize{5.8}{6.4}\selectfont \textcolor{BrickRed}{$ \blacktriangledown\,44.8\times $}}\end{tabular} & \hspace*{1pt}\begin{tabular}[t]{@{}c@{}}\rule{0pt}{12pt}22m\,35.5s\\[-.2pt]{\fontsize{5.8}{6.4}\selectfont \textcolor{BrickRed}{$ \blacktriangledown\,31.9\times $}}\\[-.2pt]22m\,34.5s\\[-.2pt]{\fontsize{5.8}{6.4}\selectfont \textcolor{BrickRed}{$ \blacktriangledown\,34.4\times $}}\end{tabular} & \hspace*{1pt}\begin{tabular}[t]{@{}c@{}}\rule{0pt}{12pt}12m\,31.2s\\[-.2pt]{\fontsize{5.8}{6.4}\selectfont \textcolor{BrickRed}{$ \blacktriangledown\,1.9\times $}}\\[-.2pt]12m\,30.6s\\[-.2pt]{\fontsize{5.8}{6.4}\selectfont \textcolor{BrickRed}{$ \blacktriangledown\,2.1\times $}}\end{tabular} & \hspace*{1pt}\begin{tabular}[t]{@{}c@{}}\rule{0pt}{12pt}7m\,58.4s\\[-.2pt]{\fontsize{5.8}{6.4}\selectfont \textcolor{BrickRed}{$ \blacktriangledown\,51.3\times $}}\\[-.2pt]7m\,54.7s\\[-.2pt]{\fontsize{5.8}{6.4}\selectfont \textcolor{BrickRed}{$ \blacktriangledown\,56.6\times $}}\end{tabular} & \hspace*{1pt}\begin{tabular}[t]{@{}c@{}}\rule{0pt}{12pt}17m\,10s\\[-.2pt]{\fontsize{5.8}{6.4}\selectfont \textcolor{BrickRed}{$ \blacktriangledown\,117.9\times $}}\\[-.2pt]17m\,7.3s\\[-.2pt]{\fontsize{5.8}{6.4}\selectfont \textcolor{BrickRed}{$ \blacktriangledown\,146.2\times $}}\end{tabular} & \hspace*{1pt}\begin{tabular}[t]{@{}c@{}}\rule{0pt}{12pt}29m\,8.1s\\[-.2pt]{\fontsize{5.8}{6.4}\selectfont \textcolor{BrickRed}{$ \blacktriangledown\,110.6\times $}}\\[-.2pt]29m\,2.5s\\[-.2pt]{\fontsize{5.8}{6.4}\selectfont \textcolor{BrickRed}{$ \blacktriangledown\,136.8\times $}}\end{tabular} & \hspace*{1pt}\begin{tabular}[t]{@{}c@{}}\rule{0pt}{12pt}15m\,18.5s\\[-.2pt]{\fontsize{5.8}{6.4}\selectfont \textcolor{BrickRed}{$ \blacktriangledown\,49.7\times $}}\\[-.2pt]15m\,15.3s\\[-.2pt]{\fontsize{5.8}{6.4}\selectfont \textcolor{BrickRed}{$ \blacktriangledown\,53.8\times $}}\end{tabular} & \hspace*{1pt}\begin{tabular}[t]{@{}c@{}}\rule{0pt}{12pt}2m\,53.1s\\[-.2pt]{\fontsize{5.8}{6.4}\selectfont \textcolor{BrickRed}{$ \blacktriangledown\,4.1\times $}}\\[-.2pt]2m\,14.9s\\[-.2pt]{\fontsize{5.8}{6.4}\selectfont \textcolor{BrickRed}{$ \blacktriangledown\,3.3\times $}}\end{tabular} & \hspace*{1pt}\begin{tabular}[t]{@{}c@{}}\rule{0pt}{12pt}2m\,43.4s\\[-.2pt]{\fontsize{5.8}{6.4}\selectfont \textcolor{BrickRed}{$ \blacktriangledown\,3.8\times $}}\\[-.2pt]2m\,13.1s\\[-.2pt]{\fontsize{5.8}{6.4}\selectfont \textcolor{BrickRed}{$ \blacktriangledown\,3.2\times $}}\end{tabular} & \textcolor{gray}{\ensuremath{\cdot}} & \textcolor{gray}{\ensuremath{\cdot}} & \hspace*{1pt}\begin{tabular}[t]{@{}c@{}}\rule{0pt}{12pt}6m\,44s\\[-.2pt]{\fontsize{5.8}{6.4}\selectfont \textcolor{BrickRed}{$ \blacktriangledown\,11.6\times $}}\\[-.2pt]2m\,28s\\[-.2pt]{\fontsize{5.8}{6.4}\selectfont \textcolor{BrickRed}{$ \blacktriangledown\,4.4\times $}}\end{tabular} & \hspace*{1pt}\begin{tabular}[t]{@{}c@{}}\rule{0pt}{12pt}6m\,4.7s\\[-.2pt]{\fontsize{5.8}{6.4}\selectfont \textcolor{BrickRed}{$ \blacktriangledown\,10.2\times $}}\\[-.2pt]2m\,28s\\[-.2pt]{\fontsize{5.8}{6.4}\selectfont \textcolor{BrickRed}{$ \blacktriangledown\,4.3\times $}}\end{tabular} & \textcolor{gray}{\ensuremath{\cdot}} & \textcolor{gray}{\ensuremath{\cdot}}\\[3pt]\hline
\BenchmarkSourceLabel{bench:aws-2-k20-w15:row:mpf-scalar-scalar-ksw2}{R39}\shortstack[l]{{\color{gray}R39}} & \begin{tabular}[t]{@{}l@{}}\rule{0pt}{12pt}MPF / Scalar\\Scalar-DP / KSW2\end{tabular} & \hspace*{1pt}\begin{tabular}[t]{@{}c@{}}\rule{0pt}{12pt}5m\,39.7s\\[-.2pt]{\fontsize{5.8}{6.4}\selectfont \textcolor{BrickRed}{$ \blacktriangledown\,17.3\times $}}\\[-.2pt]5m\,37.2s\\[-.2pt]{\fontsize{5.8}{6.4}\selectfont \textcolor{BrickRed}{$ \blacktriangledown\,18.1\times $}}\end{tabular} & \hspace*{1pt}\begin{tabular}[t]{@{}c@{}}\rule{0pt}{12pt}12m\,8.9s\\[-.2pt]{\fontsize{5.8}{6.4}\selectfont \textcolor{BrickRed}{$ \blacktriangledown\,40.0\times $}}\\[-.2pt]12m\,7.2s\\[-.2pt]{\fontsize{5.8}{6.4}\selectfont \textcolor{BrickRed}{$ \blacktriangledown\,44.0\times $}}\end{tabular} & \hspace*{1pt}\begin{tabular}[t]{@{}c@{}}\rule{0pt}{12pt}22m\,21.6s\\[-.2pt]{\fontsize{5.8}{6.4}\selectfont \textcolor{BrickRed}{$ \blacktriangledown\,31.6\times $}}\\[-.2pt]22m\,15.8s\\[-.2pt]{\fontsize{5.8}{6.4}\selectfont \textcolor{BrickRed}{$ \blacktriangledown\,33.9\times $}}\end{tabular} & \hspace*{1pt}\begin{tabular}[t]{@{}c@{}}\rule{0pt}{12pt}10m\,39.3s\\[-.2pt]{\fontsize{5.8}{6.4}\selectfont \textcolor{BrickRed}{$ \blacktriangledown\,1.6\times $}}\\[-.2pt]10m\,36.8s\\[-.2pt]{\fontsize{5.8}{6.4}\selectfont \textcolor{BrickRed}{$ \blacktriangledown\,1.8\times $}}\end{tabular} & \hspace*{1pt}\begin{tabular}[t]{@{}c@{}}\rule{0pt}{12pt}7m\,46.7s\\[-.2pt]{\fontsize{5.8}{6.4}\selectfont \textcolor{BrickRed}{$ \blacktriangledown\,50.1\times $}}\\[-.2pt]7m\,43s\\[-.2pt]{\fontsize{5.8}{6.4}\selectfont \textcolor{BrickRed}{$ \blacktriangledown\,55.2\times $}}\end{tabular} & \hspace*{1pt}\begin{tabular}[t]{@{}c@{}}\rule{0pt}{12pt}16m\,59.7s\\[-.2pt]{\fontsize{5.8}{6.4}\selectfont \textcolor{BrickRed}{$ \blacktriangledown\,116.7\times $}}\\[-.2pt]16m\,56.8s\\[-.2pt]{\fontsize{5.8}{6.4}\selectfont \textcolor{BrickRed}{$ \blacktriangledown\,144.7\times $}}\end{tabular} & \hspace*{1pt}\begin{tabular}[t]{@{}c@{}}\rule{0pt}{12pt}28m\,56.8s\\[-.2pt]{\fontsize{5.8}{6.4}\selectfont \textcolor{BrickRed}{$ \blacktriangledown\,109.8\times $}}\\[-.2pt]28m\,50.8s\\[-.2pt]{\fontsize{5.8}{6.4}\selectfont \textcolor{BrickRed}{$ \blacktriangledown\,135.9\times $}}\end{tabular} & \hspace*{1pt}\begin{tabular}[t]{@{}c@{}}\rule{0pt}{12pt}15m\,4.5s\\[-.2pt]{\fontsize{5.8}{6.4}\selectfont \textcolor{BrickRed}{$ \blacktriangledown\,48.9\times $}}\\[-.2pt]15m\,0.9s\\[-.2pt]{\fontsize{5.8}{6.4}\selectfont \textcolor{BrickRed}{$ \blacktriangledown\,52.9\times $}}\end{tabular} & \hspace*{1pt}\begin{tabular}[t]{@{}c@{}}\rule{0pt}{12pt}2m\,57.1s\\[-.2pt]{\fontsize{5.8}{6.4}\selectfont \textcolor{BrickRed}{$ \blacktriangledown\,4.2\times $}}\\[-.2pt]2m\,16.9s\\[-.2pt]{\fontsize{5.8}{6.4}\selectfont \textcolor{BrickRed}{$ \blacktriangledown\,3.3\times $}}\end{tabular} & \hspace*{1pt}\begin{tabular}[t]{@{}c@{}}\rule{0pt}{12pt}3m\,14.4s\\[-.2pt]{\fontsize{5.8}{6.4}\selectfont \textcolor{BrickRed}{$ \blacktriangledown\,4.5\times $}}\\[-.2pt]2m\,40.8s\\[-.2pt]{\fontsize{5.8}{6.4}\selectfont \textcolor{BrickRed}{$ \blacktriangledown\,3.9\times $}}\end{tabular} & \textcolor{gray}{\ensuremath{\cdot}} & \textcolor{gray}{\ensuremath{\cdot}} & \hspace*{1pt}\begin{tabular}[t]{@{}c@{}}\rule{0pt}{12pt}6m\,47s\\[-.2pt]{\fontsize{5.8}{6.4}\selectfont \textcolor{BrickRed}{$ \blacktriangledown\,11.6\times $}}\\[-.2pt]2m\,28.7s\\[-.2pt]{\fontsize{5.8}{6.4}\selectfont \textcolor{BrickRed}{$ \blacktriangledown\,4.4\times $}}\end{tabular} & \hspace*{1pt}\begin{tabular}[t]{@{}c@{}}\rule{0pt}{12pt}6m\,3s\\[-.2pt]{\fontsize{5.8}{6.4}\selectfont \textcolor{BrickRed}{$ \blacktriangledown\,10.1\times $}}\\[-.2pt]2m\,24.7s\\[-.2pt]{\fontsize{5.8}{6.4}\selectfont \textcolor{BrickRed}{$ \blacktriangledown\,4.2\times $}}\end{tabular} & \textcolor{gray}{\ensuremath{\cdot}} & \textcolor{gray}{\ensuremath{\cdot}}\\[3pt]\hline
\BenchmarkSourceLabel{bench:aws-2-k20-w15:row:mpf-scalar-simd-block}{R40}\shortstack[l]{{\color{gray}R40}} & \begin{tabular}[t]{@{}l@{}}\rule{0pt}{12pt}MPF / Scalar\\SIMD-DP / Block\end{tabular} & \hspace*{1pt}\begin{tabular}[t]{@{}c@{}}\rule{0pt}{12pt}5m\,39.7s\\[-.2pt]{\fontsize{5.8}{6.4}\selectfont \textcolor{BrickRed}{$ \blacktriangledown\,17.3\times $}}\\[-.2pt]5m\,39.2s\\[-.2pt]{\fontsize{5.8}{6.4}\selectfont \textcolor{BrickRed}{$ \blacktriangledown\,18.2\times $}}\end{tabular} & \hspace*{1pt}\begin{tabular}[t]{@{}c@{}}\rule{0pt}{12pt}12m\,5.2s\\[-.2pt]{\fontsize{5.8}{6.4}\selectfont \textcolor{BrickRed}{$ \blacktriangledown\,39.8\times $}}\\[-.2pt]12m\,4.6s\\[-.2pt]{\fontsize{5.8}{6.4}\selectfont \textcolor{BrickRed}{$ \blacktriangledown\,43.9\times $}}\end{tabular} & \hspace*{1pt}\begin{tabular}[t]{@{}c@{}}\rule{0pt}{12pt}22m\,32.6s\\[-.2pt]{\fontsize{5.8}{6.4}\selectfont \textcolor{BrickRed}{$ \blacktriangledown\,31.8\times $}}\\[-.2pt]22m\,31.6s\\[-.2pt]{\fontsize{5.8}{6.4}\selectfont \textcolor{BrickRed}{$ \blacktriangledown\,34.3\times $}}\end{tabular} & \hspace*{1pt}\begin{tabular}[t]{@{}c@{}}\rule{0pt}{12pt}12m\,48.6s\\[-.2pt]{\fontsize{5.8}{6.4}\selectfont \textcolor{BrickRed}{$ \blacktriangledown\,2.0\times $}}\\[-.2pt]12m\,47.9s\\[-.2pt]{\fontsize{5.8}{6.4}\selectfont \textcolor{BrickRed}{$ \blacktriangledown\,2.2\times $}}\end{tabular} & \hspace*{1pt}\begin{tabular}[t]{@{}c@{}}\rule{0pt}{12pt}7m\,52s\\[-.2pt]{\fontsize{5.8}{6.4}\selectfont \textcolor{BrickRed}{$ \blacktriangledown\,50.6\times $}}\\[-.2pt]7m\,48.3s\\[-.2pt]{\fontsize{5.8}{6.4}\selectfont \textcolor{BrickRed}{$ \blacktriangledown\,55.9\times $}}\end{tabular} & \hspace*{1pt}\begin{tabular}[t]{@{}c@{}}\rule{0pt}{12pt}16m\,56.1s\\[-.2pt]{\fontsize{5.8}{6.4}\selectfont \textcolor{BrickRed}{$ \blacktriangledown\,116.3\times $}}\\[-.2pt]16m\,53.4s\\[-.2pt]{\fontsize{5.8}{6.4}\selectfont \textcolor{BrickRed}{$ \blacktriangledown\,144.2\times $}}\end{tabular} & \hspace*{1pt}\begin{tabular}[t]{@{}c@{}}\rule{0pt}{12pt}29m\,42.8s\\[-.2pt]{\fontsize{5.8}{6.4}\selectfont \textcolor{BrickRed}{$ \blacktriangledown\,112.8\times $}}\\[-.2pt]29m\,37.4s\\[-.2pt]{\fontsize{5.8}{6.4}\selectfont \textcolor{BrickRed}{$ \blacktriangledown\,139.5\times $}}\end{tabular} & \hspace*{1pt}\begin{tabular}[t]{@{}c@{}}\rule{0pt}{12pt}15m\,28s\\[-.2pt]{\fontsize{5.8}{6.4}\selectfont \textcolor{BrickRed}{$ \blacktriangledown\,50.2\times $}}\\[-.2pt]15m\,24.8s\\[-.2pt]{\fontsize{5.8}{6.4}\selectfont \textcolor{BrickRed}{$ \blacktriangledown\,54.3\times $}}\end{tabular} & \hspace*{1pt}\begin{tabular}[t]{@{}c@{}}\rule{0pt}{12pt}2m\,45.2s\\[-.2pt]{\fontsize{5.8}{6.4}\selectfont \textcolor{BrickRed}{$ \blacktriangledown\,3.9\times $}}\\[-.2pt]2m\,8.9s\\[-.2pt]{\fontsize{5.8}{6.4}\selectfont \textcolor{BrickRed}{$ \blacktriangledown\,3.1\times $}}\end{tabular} & \hspace*{1pt}\begin{tabular}[t]{@{}c@{}}\rule{0pt}{12pt}2m\,44.3s\\[-.2pt]{\fontsize{5.8}{6.4}\selectfont \textcolor{BrickRed}{$ \blacktriangledown\,3.8\times $}}\\[-.2pt]2m\,14.4s\\[-.2pt]{\fontsize{5.8}{6.4}\selectfont \textcolor{BrickRed}{$ \blacktriangledown\,3.3\times $}}\end{tabular} & \textcolor{gray}{\ensuremath{\cdot}} & \hspace*{1pt}\begin{tabular}[t]{@{}c@{}}\rule{0pt}{12pt}2m\,41s\\[-.2pt]{\fontsize{5.8}{6.4}\selectfont \textcolor{ForestGreen}{$ \blacktriangle\,1.8\times $}}\\[-.2pt]2m\,13.3s\\[-.2pt]{\fontsize{5.8}{6.4}\selectfont \textcolor{ForestGreen}{$ \blacktriangle\,2.2\times $}}\end{tabular} & \hspace*{1pt}\begin{tabular}[t]{@{}c@{}}\rule{0pt}{12pt}6m\,49.3s\\[-.2pt]{\fontsize{5.8}{6.4}\selectfont \textcolor{BrickRed}{$ \blacktriangledown\,11.7\times $}}\\[-.2pt]2m\,33s\\[-.2pt]{\fontsize{5.8}{6.4}\selectfont \textcolor{BrickRed}{$ \blacktriangledown\,4.5\times $}}\end{tabular} & \hspace*{1pt}\begin{tabular}[t]{@{}c@{}}\rule{0pt}{12pt}6m\,44.7s\\[-.2pt]{\fontsize{5.8}{6.4}\selectfont \textcolor{BrickRed}{$ \blacktriangledown\,11.3\times $}}\\[-.2pt]2m\,51.6s\\[-.2pt]{\fontsize{5.8}{6.4}\selectfont \textcolor{BrickRed}{$ \blacktriangledown\,5.0\times $}}\end{tabular} & \textcolor{gray}{\ensuremath{\cdot}} & \hspace*{1pt}\begin{tabular}[t]{@{}c@{}}\rule{0pt}{12pt}11m\,45.1s\\[-.2pt]{\fontsize{5.8}{6.4}\selectfont \textcolor{BrickRed}{$ \blacktriangledown\,11.7\times $}}\\[-.2pt]2m\,31.9s\\[-.2pt]{\fontsize{5.8}{6.4}\selectfont \textcolor{BrickRed}{$ \blacktriangledown\,2.6\times $}}\end{tabular}\\[3pt]\hline
\BenchmarkSourceLabel{bench:aws-2-k20-w15:row:mpf-scalar-rmq-block}{R41}\shortstack[l]{{\color{gray}R41}} & \begin{tabular}[t]{@{}l@{}}\rule{0pt}{12pt}MPF / Scalar\\RMQ / Block\end{tabular} & \hspace*{1pt}\begin{tabular}[t]{@{}c@{}}\rule{0pt}{12pt}5m\,39.9s\\[-.2pt]{\fontsize{5.8}{6.4}\selectfont \textcolor{BrickRed}{$ \blacktriangledown\,17.3\times $}}\\[-.2pt]5m\,39.5s\\[-.2pt]{\fontsize{5.8}{6.4}\selectfont \textcolor{BrickRed}{$ \blacktriangledown\,18.2\times $}}\end{tabular} & \hspace*{1pt}\begin{tabular}[t]{@{}c@{}}\rule{0pt}{12pt}12m\,19.6s\\[-.2pt]{\fontsize{5.8}{6.4}\selectfont \textcolor{BrickRed}{$ \blacktriangledown\,40.6\times $}}\\[-.2pt]12m\,19.1s\\[-.2pt]{\fontsize{5.8}{6.4}\selectfont \textcolor{BrickRed}{$ \blacktriangledown\,44.8\times $}}\end{tabular} & \hspace*{1pt}\begin{tabular}[t]{@{}c@{}}\rule{0pt}{12pt}22m\,41.7s\\[-.2pt]{\fontsize{5.8}{6.4}\selectfont \textcolor{BrickRed}{$ \blacktriangledown\,32.0\times $}}\\[-.2pt]22m\,40.7s\\[-.2pt]{\fontsize{5.8}{6.4}\selectfont \textcolor{BrickRed}{$ \blacktriangledown\,34.6\times $}}\end{tabular} & \hspace*{1pt}\begin{tabular}[t]{@{}c@{}}\rule{0pt}{12pt}12m\,57s\\[-.2pt]{\fontsize{5.8}{6.4}\selectfont \textcolor{BrickRed}{$ \blacktriangledown\,2.0\times $}}\\[-.2pt]12m\,56.4s\\[-.2pt]{\fontsize{5.8}{6.4}\selectfont \textcolor{BrickRed}{$ \blacktriangledown\,2.2\times $}}\end{tabular} & \hspace*{1pt}\begin{tabular}[t]{@{}c@{}}\rule{0pt}{12pt}7m\,52.2s\\[-.2pt]{\fontsize{5.8}{6.4}\selectfont \textcolor{BrickRed}{$ \blacktriangledown\,50.7\times $}}\\[-.2pt]7m\,48.5s\\[-.2pt]{\fontsize{5.8}{6.4}\selectfont \textcolor{BrickRed}{$ \blacktriangledown\,55.9\times $}}\end{tabular} & \hspace*{1pt}\begin{tabular}[t]{@{}c@{}}\rule{0pt}{12pt}17m\,9.4s\\[-.2pt]{\fontsize{5.8}{6.4}\selectfont \textcolor{BrickRed}{$ \blacktriangledown\,117.8\times $}}\\[-.2pt]17m\,6.8s\\[-.2pt]{\fontsize{5.8}{6.4}\selectfont \textcolor{BrickRed}{$ \blacktriangledown\,146.1\times $}}\end{tabular} & \hspace*{1pt}\begin{tabular}[t]{@{}c@{}}\rule{0pt}{12pt}21m\,1.2s\\[-.2pt]{\fontsize{5.8}{6.4}\selectfont \textcolor{BrickRed}{$ \blacktriangledown\,79.8\times $}}\\[-.2pt]20m\,58.9s\\[-.2pt]{\fontsize{5.8}{6.4}\selectfont \textcolor{BrickRed}{$ \blacktriangledown\,98.8\times $}}\end{tabular} & \hspace*{1pt}\begin{tabular}[t]{@{}c@{}}\rule{0pt}{12pt}15m\,4.2s\\[-.2pt]{\fontsize{5.8}{6.4}\selectfont \textcolor{BrickRed}{$ \blacktriangledown\,48.9\times $}}\\[-.2pt]15m\,1.3s\\[-.2pt]{\fontsize{5.8}{6.4}\selectfont \textcolor{BrickRed}{$ \blacktriangledown\,52.9\times $}}\end{tabular} & \hspace*{1pt}\begin{tabular}[t]{@{}c@{}}\rule{0pt}{12pt}2m\,49.9s\\[-.2pt]{\fontsize{5.8}{6.4}\selectfont \textcolor{BrickRed}{$ \blacktriangledown\,4.0\times $}}\\[-.2pt]2m\,13.4s\\[-.2pt]{\fontsize{5.8}{6.4}\selectfont \textcolor{BrickRed}{$ \blacktriangledown\,3.2\times $}}\end{tabular} & \hspace*{1pt}\begin{tabular}[t]{@{}c@{}}\rule{0pt}{12pt}2m\,29.9s\\[-.2pt]{\fontsize{5.8}{6.4}\selectfont \textcolor{BrickRed}{$ \blacktriangledown\,3.5\times $}}\\[-.2pt]2m\,1.1s\\[-.2pt]{\fontsize{5.8}{6.4}\selectfont \textcolor{BrickRed}{$ \blacktriangledown\,2.9\times $}}\end{tabular} & \textcolor{gray}{\ensuremath{\cdot}} & \textcolor{gray}{\ensuremath{\cdot}} & \hspace*{1pt}\begin{tabular}[t]{@{}c@{}}\rule{0pt}{12pt}6m\,44.4s\\[-.2pt]{\fontsize{5.8}{6.4}\selectfont \textcolor{BrickRed}{$ \blacktriangledown\,11.6\times $}}\\[-.2pt]2m\,28.2s\\[-.2pt]{\fontsize{5.8}{6.4}\selectfont \textcolor{BrickRed}{$ \blacktriangledown\,4.4\times $}}\end{tabular} & \hspace*{1pt}\begin{tabular}[t]{@{}c@{}}\rule{0pt}{12pt}6m\,31.8s\\[-.2pt]{\fontsize{5.8}{6.4}\selectfont \textcolor{BrickRed}{$ \blacktriangledown\,10.9\times $}}\\[-.2pt]2m\,38.1s\\[-.2pt]{\fontsize{5.8}{6.4}\selectfont \textcolor{BrickRed}{$ \blacktriangledown\,4.6\times $}}\end{tabular} & \textcolor{gray}{\ensuremath{\cdot}} & \textcolor{gray}{\ensuremath{\cdot}}\\[3pt]\hline
\BenchmarkSourceLabel{bench:aws-2-k20-w15:row:mpf-simdref-scalar-ksw2}{R42}\shortstack[l]{{\color{gray}R42}} & \begin{tabular}[t]{@{}l@{}}\rule{0pt}{12pt}MPF / SIMD-ref\\Scalar-DP / KSW2\end{tabular} & \hspace*{1pt}\begin{tabular}[t]{@{}c@{}}\rule{0pt}{12pt}5m\,40.6s\\[-.2pt]{\fontsize{5.8}{6.4}\selectfont \textcolor{BrickRed}{$ \blacktriangledown\,17.4\times $}}\\[-.2pt]5m\,36.9s\\[-.2pt]{\fontsize{5.8}{6.4}\selectfont \textcolor{BrickRed}{$ \blacktriangledown\,18.1\times $}}\end{tabular} & \hspace*{1pt}\begin{tabular}[t]{@{}c@{}}\rule{0pt}{12pt}12m\,16.8s\\[-.2pt]{\fontsize{5.8}{6.4}\selectfont \textcolor{BrickRed}{$ \blacktriangledown\,40.5\times $}}\\[-.2pt]12m\,13.7s\\[-.2pt]{\fontsize{5.8}{6.4}\selectfont \textcolor{BrickRed}{$ \blacktriangledown\,44.4\times $}}\end{tabular} & \hspace*{1pt}\begin{tabular}[t]{@{}c@{}}\rule{0pt}{12pt}22m\,52.9s\\[-.2pt]{\fontsize{5.8}{6.4}\selectfont \textcolor{BrickRed}{$ \blacktriangledown\,32.3\times $}}\\[-.2pt]22m\,46.8s\\[-.2pt]{\fontsize{5.8}{6.4}\selectfont \textcolor{BrickRed}{$ \blacktriangledown\,34.7\times $}}\end{tabular} & \hspace*{1pt}\begin{tabular}[t]{@{}c@{}}\rule{0pt}{12pt}12m\,8.2s\\[-.2pt]{\fontsize{5.8}{6.4}\selectfont \textcolor{BrickRed}{$ \blacktriangledown\,1.9\times $}}\\[-.2pt]12m\,4.8s\\[-.2pt]{\fontsize{5.8}{6.4}\selectfont \textcolor{BrickRed}{$ \blacktriangledown\,2.0\times $}}\end{tabular} & \hspace*{1pt}\begin{tabular}[t]{@{}c@{}}\rule{0pt}{12pt}6m\,5.2s\\[-.2pt]{\fontsize{5.8}{6.4}\selectfont \textcolor{BrickRed}{$ \blacktriangledown\,39.2\times $}}\\[-.2pt]6m\,1.5s\\[-.2pt]{\fontsize{5.8}{6.4}\selectfont \textcolor{BrickRed}{$ \blacktriangledown\,43.1\times $}}\end{tabular} & \hspace*{1pt}\begin{tabular}[t]{@{}c@{}}\rule{0pt}{12pt}13m\,19.6s\\[-.2pt]{\fontsize{5.8}{6.4}\selectfont \textcolor{BrickRed}{$ \blacktriangledown\,91.5\times $}}\\[-.2pt]13m\,16.8s\\[-.2pt]{\fontsize{5.8}{6.4}\selectfont \textcolor{BrickRed}{$ \blacktriangledown\,113.4\times $}}\end{tabular} & \hspace*{1pt}\begin{tabular}[t]{@{}c@{}}\rule{0pt}{12pt}24m\,19.3s\\[-.2pt]{\fontsize{5.8}{6.4}\selectfont \textcolor{BrickRed}{$ \blacktriangledown\,92.3\times $}}\\[-.2pt]24m\,13.1s\\[-.2pt]{\fontsize{5.8}{6.4}\selectfont \textcolor{BrickRed}{$ \blacktriangledown\,114.1\times $}}\end{tabular} & \hspace*{1pt}\begin{tabular}[t]{@{}c@{}}\rule{0pt}{12pt}11m\,40.3s\\[-.2pt]{\fontsize{5.8}{6.4}\selectfont \textcolor{BrickRed}{$ \blacktriangledown\,37.9\times $}}\\[-.2pt]11m\,36.8s\\[-.2pt]{\fontsize{5.8}{6.4}\selectfont \textcolor{BrickRed}{$ \blacktriangledown\,40.9\times $}}\end{tabular} & \hspace*{1pt}\begin{tabular}[t]{@{}c@{}}\rule{0pt}{12pt}2m\,42s\\[-.2pt]{\fontsize{5.8}{6.4}\selectfont \textcolor{BrickRed}{$ \blacktriangledown\,3.8\times $}}\\[-.2pt]2m\,2.1s\\[-.2pt]{\fontsize{5.8}{6.4}\selectfont \textcolor{BrickRed}{$ \blacktriangledown\,2.9\times $}}\end{tabular} & \hspace*{1pt}\begin{tabular}[t]{@{}c@{}}\rule{0pt}{12pt}2m\,30.1s\\[-.2pt]{\fontsize{5.8}{6.4}\selectfont \textcolor{BrickRed}{$ \blacktriangledown\,3.5\times $}}\\[-.2pt]1m\,57.8s\\[-.2pt]{\fontsize{5.8}{6.4}\selectfont \textcolor{BrickRed}{$ \blacktriangledown\,2.9\times $}}\end{tabular} & \textcolor{gray}{\ensuremath{\cdot}} & \textcolor{gray}{\ensuremath{\cdot}} & \hspace*{1pt}\begin{tabular}[t]{@{}c@{}}\rule{0pt}{12pt}6m\,17.7s\\[-.2pt]{\fontsize{5.8}{6.4}\selectfont \textcolor{BrickRed}{$ \blacktriangledown\,10.8\times $}}\\[-.2pt]2m\,16.5s\\[-.2pt]{\fontsize{5.8}{6.4}\selectfont \textcolor{BrickRed}{$ \blacktriangledown\,4.0\times $}}\end{tabular} & \hspace*{1pt}\begin{tabular}[t]{@{}c@{}}\rule{0pt}{12pt}6m\,21.8s\\[-.2pt]{\fontsize{5.8}{6.4}\selectfont \textcolor{BrickRed}{$ \blacktriangledown\,10.6\times $}}\\[-.2pt]2m\,44.5s\\[-.2pt]{\fontsize{5.8}{6.4}\selectfont \textcolor{BrickRed}{$ \blacktriangledown\,4.8\times $}}\end{tabular} & \textcolor{gray}{\ensuremath{\cdot}} & \textcolor{gray}{\ensuremath{\cdot}}\\[3pt]\hline
\BenchmarkSourceLabel{bench:aws-2-k20-w15:row:mpf-scalar-simd-ksw2}{R43}\shortstack[l]{{\color{gray}R43}} & \begin{tabular}[t]{@{}l@{}}\rule{0pt}{12pt}MPF / Scalar\\SIMD-DP / KSW2\end{tabular} & \hspace*{1pt}\begin{tabular}[t]{@{}c@{}}\rule{0pt}{12pt}5m\,42.8s\\[-.2pt]{\fontsize{5.8}{6.4}\selectfont \textcolor{BrickRed}{$ \blacktriangledown\,17.5\times $}}\\[-.2pt]5m\,40.6s\\[-.2pt]{\fontsize{5.8}{6.4}\selectfont \textcolor{BrickRed}{$ \blacktriangledown\,18.3\times $}}\end{tabular} & \hspace*{1pt}\begin{tabular}[t]{@{}c@{}}\rule{0pt}{12pt}12m\,19.8s\\[-.2pt]{\fontsize{5.8}{6.4}\selectfont \textcolor{BrickRed}{$ \blacktriangledown\,40.6\times $}}\\[-.2pt]12m\,18s\\[-.2pt]{\fontsize{5.8}{6.4}\selectfont \textcolor{BrickRed}{$ \blacktriangledown\,44.7\times $}}\end{tabular} & \hspace*{1pt}\begin{tabular}[t]{@{}c@{}}\rule{0pt}{12pt}22m\,43.4s\\[-.2pt]{\fontsize{5.8}{6.4}\selectfont \textcolor{BrickRed}{$ \blacktriangledown\,32.1\times $}}\\[-.2pt]22m\,39.5s\\[-.2pt]{\fontsize{5.8}{6.4}\selectfont \textcolor{BrickRed}{$ \blacktriangledown\,34.5\times $}}\end{tabular} & \hspace*{1pt}\begin{tabular}[t]{@{}c@{}}\rule{0pt}{12pt}11m\,32.1s\\[-.2pt]{\fontsize{5.8}{6.4}\selectfont \textcolor{BrickRed}{$ \blacktriangledown\,1.8\times $}}\\[-.2pt]11m\,29.4s\\[-.2pt]{\fontsize{5.8}{6.4}\selectfont \textcolor{BrickRed}{$ \blacktriangledown\,1.9\times $}}\end{tabular} & \hspace*{1pt}\begin{tabular}[t]{@{}c@{}}\rule{0pt}{12pt}8m\,18.5s\\[-.2pt]{\fontsize{5.8}{6.4}\selectfont \textcolor{BrickRed}{$ \blacktriangledown\,53.5\times $}}\\[-.2pt]8m\,14.8s\\[-.2pt]{\fontsize{5.8}{6.4}\selectfont \textcolor{BrickRed}{$ \blacktriangledown\,59.0\times $}}\end{tabular} & \hspace*{1pt}\begin{tabular}[t]{@{}c@{}}\rule{0pt}{12pt}16m\,50.3s\\[-.2pt]{\fontsize{5.8}{6.4}\selectfont \textcolor{BrickRed}{$ \blacktriangledown\,115.6\times $}}\\[-.2pt]16m\,47.4s\\[-.2pt]{\fontsize{5.8}{6.4}\selectfont \textcolor{BrickRed}{$ \blacktriangledown\,143.4\times $}}\end{tabular} & \hspace*{1pt}\begin{tabular}[t]{@{}c@{}}\rule{0pt}{12pt}29m\,22.4s\\[-.2pt]{\fontsize{5.8}{6.4}\selectfont \textcolor{BrickRed}{$ \blacktriangledown\,111.5\times $}}\\[-.2pt]29m\,16.4s\\[-.2pt]{\fontsize{5.8}{6.4}\selectfont \textcolor{BrickRed}{$ \blacktriangledown\,137.9\times $}}\end{tabular} & \hspace*{1pt}\begin{tabular}[t]{@{}c@{}}\rule{0pt}{12pt}15m\,17.5s\\[-.2pt]{\fontsize{5.8}{6.4}\selectfont \textcolor{BrickRed}{$ \blacktriangledown\,49.6\times $}}\\[-.2pt]15m\,13.9s\\[-.2pt]{\fontsize{5.8}{6.4}\selectfont \textcolor{BrickRed}{$ \blacktriangledown\,53.7\times $}}\end{tabular} & \hspace*{1pt}\begin{tabular}[t]{@{}c@{}}\rule{0pt}{12pt}2m\,52.8s\\[-.2pt]{\fontsize{5.8}{6.4}\selectfont \textcolor{BrickRed}{$ \blacktriangledown\,4.1\times $}}\\[-.2pt]2m\,15.7s\\[-.2pt]{\fontsize{5.8}{6.4}\selectfont \textcolor{BrickRed}{$ \blacktriangledown\,3.3\times $}}\end{tabular} & \hspace*{1pt}\begin{tabular}[t]{@{}c@{}}\rule{0pt}{12pt}2m\,32.2s\\[-.2pt]{\fontsize{5.8}{6.4}\selectfont \textcolor{BrickRed}{$ \blacktriangledown\,3.6\times $}}\\[-.2pt]2m\,0.8s\\[-.2pt]{\fontsize{5.8}{6.4}\selectfont \textcolor{BrickRed}{$ \blacktriangledown\,2.9\times $}}\end{tabular} & \textcolor{gray}{\ensuremath{\cdot}} & \textcolor{gray}{\ensuremath{\cdot}} & \hspace*{1pt}\begin{tabular}[t]{@{}c@{}}\rule{0pt}{12pt}6m\,48.6s\\[-.2pt]{\fontsize{5.8}{6.4}\selectfont \textcolor{BrickRed}{$ \blacktriangledown\,11.7\times $}}\\[-.2pt]2m\,28.4s\\[-.2pt]{\fontsize{5.8}{6.4}\selectfont \textcolor{BrickRed}{$ \blacktriangledown\,4.4\times $}}\end{tabular} & \hspace*{1pt}\begin{tabular}[t]{@{}c@{}}\rule{0pt}{12pt}6m\,49.3s\\[-.2pt]{\fontsize{5.8}{6.4}\selectfont \textcolor{BrickRed}{$ \blacktriangledown\,11.4\times $}}\\[-.2pt]2m\,53.7s\\[-.2pt]{\fontsize{5.8}{6.4}\selectfont \textcolor{BrickRed}{$ \blacktriangledown\,5.1\times $}}\end{tabular} & \textcolor{gray}{\ensuremath{\cdot}} & \textcolor{gray}{\ensuremath{\cdot}}\\[3pt]\hline
\end{tabular}\par}
\smallskip{\footnotesize\noindent Each cell: A and B times, $\cdot$: no measurement. ---\textsuperscript{g}: GPU unavailable. n/a: no comparison.\par
\noindent Blue: one lowest A and one lowest B measurement for each column for each configuration.\par
\noindent Gray: minimap2. Yellow: my preferred MPF / Bounded / SIMD-DP / Block configuration.\par
\noindent MPF/MPS: minimizers per fragment/sequence. SIMD-ref: naive SIMD scan. Bounded: bounded SIMD algorithm.\par
}
\clearpage
\noindent{\large\bfseries Table \thetable: (E): AWS: 2 Gbp, k=20, w=15\hfill 7/7}\par\smallskip
{\fontsize{7.1}{8.1}\selectfont\setlength{\tabcolsep}{1pt}\setlength{\fboxsep}{2pt}\renewcommand{\arraystretch}{1.12}
\noindent\begin{tabular}{@{}p{14mm}p{27mm}|p{13.55mm}|p{13.55mm}|p{13.55mm}|p{13.55mm}|p{13.55mm}|p{13.55mm}|p{13.55mm}|p{13.55mm}|p{13.55mm}|p{13.55mm}|p{13.55mm}|p{13.55mm}|p{13.55mm}|p{13.55mm}|p{13.55mm}|p{13.55mm}|@{}}\hline
Engine & Execution & \multicolumn{8}{c|}{\textbf{Index}} & \multicolumn{8}{c|}{\textbf{No index}}\\
 &  & \multicolumn{8}{c|}{\textbf{2 Gbp}} & \multicolumn{8}{c|}{\textbf{2 Gbp}}\\
 &  & \multicolumn{8}{c|}{\textbf{k=20, w=15}} & \multicolumn{8}{c|}{\textbf{k=20, w=15}}\\
 &  & \multicolumn{4}{c|}{\textbf{DNA4}} & \multicolumn{4}{c|}{\textbf{DNA5}} & \multicolumn{4}{c|}{\textbf{DNA4}} & \multicolumn{4}{c|}{\textbf{DNA5}}\\
 &  & \multicolumn{1}{c|}{\textbf{M7A}} & \multicolumn{1}{c|}{\textbf{M7I}} & \multicolumn{1}{c|}{\textbf{M4}} & \multicolumn{1}{c|}{\textbf{C7G}} & \multicolumn{1}{c|}{\textbf{M7A}} & \multicolumn{1}{c|}{\textbf{M7I}} & \multicolumn{1}{c|}{\textbf{M4}} & \multicolumn{1}{c|}{\textbf{C7G}} & \multicolumn{1}{c|}{\textbf{M7A}} & \multicolumn{1}{c|}{\textbf{M7I}} & \multicolumn{1}{c|}{\textbf{M4}} & \multicolumn{1}{c|}{\textbf{C7G}} & \multicolumn{1}{c|}{\textbf{M7A}} & \multicolumn{1}{c|}{\textbf{M7I}} & \multicolumn{1}{c|}{\textbf{M4}} & \multicolumn{1}{c|}{\textbf{C7G}}\\
 & & {\color{gray}C1} & {\color{gray}C2} & {\color{gray}C3} & {\color{gray}C4} & {\color{gray}C5} & {\color{gray}C6} & {\color{gray}C7} & {\color{gray}C8} & {\color{gray}C9} & {\color{gray}C10} & {\color{gray}C11} & {\color{gray}C12} & {\color{gray}C13} & {\color{gray}C14} & {\color{gray}C15} & {\color{gray}C16}\\
\hline
\rowcolor{BenchBaseline}[0.5pt][0.5pt]\shortstack[l]{{\color{gray}R1}\\\textbf{minimap2}} & \textbf{Baseline} & \hspace*{1pt}\begin{tabular}[t]{@{}c@{}}\rule{0pt}{12pt}19.6s\\[-.2pt]{\fontsize{5.8}{6.4}\selectfont \textcolor{BrickRed}{$ \blacktriangledown\,6.7\times $}}\\[-.2pt]18.6s\\[-.2pt]{\fontsize{5.8}{6.4}\selectfont \textcolor{BrickRed}{$ \blacktriangledown\,8.4\times $}}\end{tabular} & \hspace*{1pt}\begin{tabular}[t]{@{}c@{}}\rule{0pt}{12pt}18.2s\\[-.2pt]{\fontsize{5.8}{6.4}\selectfont \textcolor{BrickRed}{$ \blacktriangledown\,6.9\times $}}\\[-.2pt]16.5s\\[-.2pt]{\fontsize{5.8}{6.4}\selectfont \textcolor{BrickRed}{$ \blacktriangledown\,8.6\times $}}\end{tabular} & \hspace*{1pt}\begin{tabular}[t]{@{}c@{}}\rule{0pt}{12pt}42.5s\\[-.2pt]{\fontsize{5.8}{6.4}\selectfont \textcolor{BrickRed}{$ \blacktriangledown\,8.2\times $}}\\[-.2pt]39.4s\\[-.2pt]{\fontsize{5.8}{6.4}\selectfont \textcolor{BrickRed}{$ \blacktriangledown\,12.6\times $}}\end{tabular} & \hspace*{1pt}\begin{tabular}[t]{@{}c@{}}\rule{0pt}{12pt}6m\,29.2s\\[-.2pt]{\fontsize{5.8}{6.4}\selectfont \textcolor{BrickRed}{$ \blacktriangledown\,92.8\times $}}\\[-.2pt]5m\,56.9s\\[-.2pt]{\fontsize{5.8}{6.4}\selectfont \textcolor{BrickRed}{$ \blacktriangledown\,184.9\times $}}\end{tabular} & \hspace*{1pt}\begin{tabular}[t]{@{}c@{}}\rule{0pt}{12pt}9.3s\\[-.2pt]{\fontsize{5.8}{6.4}\selectfont \textcolor{BrickRed}{$ \blacktriangledown\,24\% $}}\\[-.2pt]8.4s\\[-.2pt]{\fontsize{5.8}{6.4}\selectfont \textcolor{BrickRed}{$ \blacktriangledown\,1.6\times $}}\end{tabular} & \hspace*{1pt}\begin{tabular}[t]{@{}c@{}}\rule{0pt}{12pt}8.7s\\[-.2pt]{\fontsize{5.8}{6.4}\selectfont \textcolor{BrickRed}{$ \blacktriangledown\,1.5\times $}}\\[-.2pt]7s\\[-.2pt]{\fontsize{5.8}{6.4}\selectfont \textcolor{BrickRed}{$ \blacktriangledown\,1.7\times $}}\end{tabular} & \hspace*{1pt}\begin{tabular}[t]{@{}c@{}}\rule{0pt}{12pt}15.8s\\[-.2pt]{\fontsize{5.8}{6.4}\selectfont \textcolor{BrickRed}{$ \blacktriangledown\,26\% $}}\\[-.2pt]12.7s\\[-.2pt]{\fontsize{5.8}{6.4}\selectfont \textcolor{BrickRed}{$ \blacktriangledown\,1.7\times $}}\end{tabular} & \hspace*{1pt}\begin{tabular}[t]{@{}c@{}}\rule{0pt}{12pt}18.5s\\[-.2pt]{\fontsize{5.8}{6.4}\selectfont \textcolor{BrickRed}{$ \blacktriangledown\,1.6\times $}}\\[-.2pt]17s\\[-.2pt]{\fontsize{5.8}{6.4}\selectfont \textcolor{BrickRed}{$ \blacktriangledown\,1.8\times $}}\end{tabular} & \hspace*{1pt}\begin{tabular}[t]{@{}c@{}}\rule{0pt}{12pt}42.4s\\[-.2pt]{\fontsize{5.8}{6.4}\selectfont \textcolor{ForestGreen}{$ \blacktriangle\,2.0\times $}}\\[-.2pt]41.4s\\[-.2pt]{\fontsize{5.8}{6.4}\selectfont \textcolor{ForestGreen}{$ \blacktriangle\,1.9\times $}}\end{tabular} & \hspace*{1pt}\begin{tabular}[t]{@{}c@{}}\rule{0pt}{12pt}42.9s\\[-.2pt]{\fontsize{5.8}{6.4}\selectfont \textcolor{ForestGreen}{$ \blacktriangle\,2.0\times $}}\\[-.2pt]41.1s\\[-.2pt]{\fontsize{5.8}{6.4}\selectfont \textcolor{ForestGreen}{$ \blacktriangle\,1.8\times $}}\end{tabular} & \hspace*{1pt}\begin{tabular}[t]{@{}c@{}}\rule{0pt}{12pt}1m\,26.9s\\[-.2pt]{\fontsize{5.8}{6.4}\selectfont \textcolor{ForestGreen}{$ \blacktriangle\,1.9\times $}}\\[-.2pt]1m\,23.7s\\[-.2pt]{\fontsize{5.8}{6.4}\selectfont \textcolor{ForestGreen}{$ \blacktriangle\,1.7\times $}}\end{tabular} & \hspace*{1pt}\begin{tabular}[t]{@{}c@{}}\rule{0pt}{12pt}4m\,51.6s\\[-.2pt]{\fontsize{5.8}{6.4}\selectfont \textcolor{BrickRed}{$ \blacktriangledown\,2.9\times $}}\\[-.2pt]4m\,46.6s\\[-.2pt]{\fontsize{5.8}{6.4}\selectfont \textcolor{BrickRed}{$ \blacktriangledown\,3.0\times $}}\end{tabular} & \hspace*{1pt}\begin{tabular}[t]{@{}c@{}}\rule{0pt}{12pt}35s\\[-.2pt]{\fontsize{5.8}{6.4}\selectfont \textcolor{ForestGreen}{$ \blacktriangle\,8.6\times $}}\\[-.2pt]34s\\[-.2pt]{\fontsize{5.8}{6.4}\selectfont \textcolor{ForestGreen}{$ \blacktriangle\,2.3\times $}}\end{tabular} & \hspace*{1pt}\begin{tabular}[t]{@{}c@{}}\rule{0pt}{12pt}35.9s\\[-.2pt]{\fontsize{5.8}{6.4}\selectfont \textcolor{ForestGreen}{$ \blacktriangle\,6.7\times $}}\\[-.2pt]34.2s\\[-.2pt]{\fontsize{5.8}{6.4}\selectfont \textcolor{ForestGreen}{$ \blacktriangle\,2.1\times $}}\end{tabular} & \hspace*{1pt}\begin{tabular}[t]{@{}c@{}}\rule{0pt}{12pt}1m\,14.1s\\[-.2pt]{\fontsize{5.8}{6.4}\selectfont \textcolor{ForestGreen}{$ \blacktriangle\,6.9\times $}}\\[-.2pt]1m\,11.1s\\[-.2pt]{\fontsize{5.8}{6.4}\selectfont \textcolor{ForestGreen}{$ \blacktriangle\,2.2\times $}}\end{tabular} & \hspace*{1pt}\begin{tabular}[t]{@{}c@{}}\rule{0pt}{12pt}1m\,0.1s\\[-.2pt]{\fontsize{5.8}{6.4}\selectfont \textcolor{ForestGreen}{$ \blacktriangle\,10.4\times $}}\\[-.2pt]58.6s\\[-.2pt]{\fontsize{5.8}{6.4}\selectfont \textcolor{ForestGreen}{$ \blacktriangle\,1.8\times $}}\end{tabular}\\[3pt]\hline
\BenchmarkSourceLabel{bench:aws-2-k20-w15:row:mpf-simdref-rmq-block}{R44}\shortstack[l]{{\color{gray}R44}\\\textbf{BioseqDB}} & \begin{tabular}[t]{@{}l@{}}\rule{0pt}{12pt}MPF / SIMD-ref\\RMQ / Block\end{tabular} & \hspace*{1pt}\begin{tabular}[t]{@{}c@{}}\rule{0pt}{12pt}5m\,44s\\[-.2pt]{\fontsize{5.8}{6.4}\selectfont \textcolor{BrickRed}{$ \blacktriangledown\,17.6\times $}}\\[-.2pt]5m\,43.3s\\[-.2pt]{\fontsize{5.8}{6.4}\selectfont \textcolor{BrickRed}{$ \blacktriangledown\,18.4\times $}}\end{tabular} & \hspace*{1pt}\begin{tabular}[t]{@{}c@{}}\rule{0pt}{12pt}12m\,19.6s\\[-.2pt]{\fontsize{5.8}{6.4}\selectfont \textcolor{BrickRed}{$ \blacktriangledown\,40.6\times $}}\\[-.2pt]12m\,18.9s\\[-.2pt]{\fontsize{5.8}{6.4}\selectfont \textcolor{BrickRed}{$ \blacktriangledown\,44.7\times $}}\end{tabular} & \hspace*{1pt}\begin{tabular}[t]{@{}c@{}}\rule{0pt}{12pt}23m\,12.3s\\[-.2pt]{\fontsize{5.8}{6.4}\selectfont \textcolor{BrickRed}{$ \blacktriangledown\,32.7\times $}}\\[-.2pt]23m\,11s\\[-.2pt]{\fontsize{5.8}{6.4}\selectfont \textcolor{BrickRed}{$ \blacktriangledown\,35.3\times $}}\end{tabular} & \hspace*{1pt}\begin{tabular}[t]{@{}c@{}}\rule{0pt}{12pt}11m\,15.2s\\[-.2pt]{\fontsize{5.8}{6.4}\selectfont \textcolor{BrickRed}{$ \blacktriangledown\,1.7\times $}}\\[-.2pt]11m\,14.4s\\[-.2pt]{\fontsize{5.8}{6.4}\selectfont \textcolor{BrickRed}{$ \blacktriangledown\,1.9\times $}}\end{tabular} & \hspace*{1pt}\begin{tabular}[t]{@{}c@{}}\rule{0pt}{12pt}6m\,4.6s\\[-.2pt]{\fontsize{5.8}{6.4}\selectfont \textcolor{BrickRed}{$ \blacktriangledown\,39.1\times $}}\\[-.2pt]6m\,1s\\[-.2pt]{\fontsize{5.8}{6.4}\selectfont \textcolor{BrickRed}{$ \blacktriangledown\,43.1\times $}}\end{tabular} & \hspace*{1pt}\begin{tabular}[t]{@{}c@{}}\rule{0pt}{12pt}13m\,18.8s\\[-.2pt]{\fontsize{5.8}{6.4}\selectfont \textcolor{BrickRed}{$ \blacktriangledown\,91.4\times $}}\\[-.2pt]13m\,16.2s\\[-.2pt]{\fontsize{5.8}{6.4}\selectfont \textcolor{BrickRed}{$ \blacktriangledown\,113.3\times $}}\end{tabular} & \hspace*{1pt}\begin{tabular}[t]{@{}c@{}}\rule{0pt}{12pt}23m\,19.7s\\[-.2pt]{\fontsize{5.8}{6.4}\selectfont \textcolor{BrickRed}{$ \blacktriangledown\,88.5\times $}}\\[-.2pt]23m\,16.1s\\[-.2pt]{\fontsize{5.8}{6.4}\selectfont \textcolor{BrickRed}{$ \blacktriangledown\,109.6\times $}}\end{tabular} & \hspace*{1pt}\begin{tabular}[t]{@{}c@{}}\rule{0pt}{12pt}12m\,38s\\[-.2pt]{\fontsize{5.8}{6.4}\selectfont \textcolor{BrickRed}{$ \blacktriangledown\,41.0\times $}}\\[-.2pt]12m\,35s\\[-.2pt]{\fontsize{5.8}{6.4}\selectfont \textcolor{BrickRed}{$ \blacktriangledown\,44.4\times $}}\end{tabular} & \hspace*{1pt}\begin{tabular}[t]{@{}c@{}}\rule{0pt}{12pt}2m\,40.3s\\[-.2pt]{\fontsize{5.8}{6.4}\selectfont \textcolor{BrickRed}{$ \blacktriangledown\,3.8\times $}}\\[-.2pt]2m\,2.4s\\[-.2pt]{\fontsize{5.8}{6.4}\selectfont \textcolor{BrickRed}{$ \blacktriangledown\,3.0\times $}}\end{tabular} & \hspace*{1pt}\begin{tabular}[t]{@{}c@{}}\rule{0pt}{12pt}2m\,36.1s\\[-.2pt]{\fontsize{5.8}{6.4}\selectfont \textcolor{BrickRed}{$ \blacktriangledown\,3.6\times $}}\\[-.2pt]2m\,4.7s\\[-.2pt]{\fontsize{5.8}{6.4}\selectfont \textcolor{BrickRed}{$ \blacktriangledown\,3.0\times $}}\end{tabular} & \textcolor{gray}{\ensuremath{\cdot}} & \textcolor{gray}{\ensuremath{\cdot}} & \hspace*{1pt}\begin{tabular}[t]{@{}c@{}}\rule{0pt}{12pt}6m\,15.1s\\[-.2pt]{\fontsize{5.8}{6.4}\selectfont \textcolor{BrickRed}{$ \blacktriangledown\,10.7\times $}}\\[-.2pt]2m\,16.3s\\[-.2pt]{\fontsize{5.8}{6.4}\selectfont \textcolor{BrickRed}{$ \blacktriangledown\,4.0\times $}}\end{tabular} & \hspace*{1pt}\begin{tabular}[t]{@{}c@{}}\rule{0pt}{12pt}6m\,12.9s\\[-.2pt]{\fontsize{5.8}{6.4}\selectfont \textcolor{BrickRed}{$ \blacktriangledown\,10.4\times $}}\\[-.2pt]2m\,40.4s\\[-.2pt]{\fontsize{5.8}{6.4}\selectfont \textcolor{BrickRed}{$ \blacktriangledown\,4.7\times $}}\end{tabular} & \textcolor{gray}{\ensuremath{\cdot}} & \textcolor{gray}{\ensuremath{\cdot}}\\[3pt]\hline
\BenchmarkSourceLabel{bench:aws-2-k20-w15:row:mpf-scalar-rmq-ksw2}{R45}\shortstack[l]{{\color{gray}R45}} & \begin{tabular}[t]{@{}l@{}}\rule{0pt}{12pt}MPF / Scalar\\RMQ / KSW2\end{tabular} & \hspace*{1pt}\begin{tabular}[t]{@{}c@{}}\rule{0pt}{12pt}5m\,46.6s\\[-.2pt]{\fontsize{5.8}{6.4}\selectfont \textcolor{BrickRed}{$ \blacktriangledown\,17.7\times $}}\\[-.2pt]5m\,44.5s\\[-.2pt]{\fontsize{5.8}{6.4}\selectfont \textcolor{BrickRed}{$ \blacktriangledown\,18.5\times $}}\end{tabular} & \hspace*{1pt}\begin{tabular}[t]{@{}c@{}}\rule{0pt}{12pt}12m\,26.1s\\[-.2pt]{\fontsize{5.8}{6.4}\selectfont \textcolor{BrickRed}{$ \blacktriangledown\,41.0\times $}}\\[-.2pt]12m\,24.3s\\[-.2pt]{\fontsize{5.8}{6.4}\selectfont \textcolor{BrickRed}{$ \blacktriangledown\,45.1\times $}}\end{tabular} & \hspace*{1pt}\begin{tabular}[t]{@{}c@{}}\rule{0pt}{12pt}22m\,22.5s\\[-.2pt]{\fontsize{5.8}{6.4}\selectfont \textcolor{BrickRed}{$ \blacktriangledown\,31.6\times $}}\\[-.2pt]22m\,16.8s\\[-.2pt]{\fontsize{5.8}{6.4}\selectfont \textcolor{BrickRed}{$ \blacktriangledown\,33.9\times $}}\end{tabular} & \hspace*{1pt}\begin{tabular}[t]{@{}c@{}}\rule{0pt}{12pt}11m\,45.8s\\[-.2pt]{\fontsize{5.8}{6.4}\selectfont \textcolor{BrickRed}{$ \blacktriangledown\,1.8\times $}}\\[-.2pt]11m\,43.3s\\[-.2pt]{\fontsize{5.8}{6.4}\selectfont \textcolor{BrickRed}{$ \blacktriangledown\,2.0\times $}}\end{tabular} & \hspace*{1pt}\begin{tabular}[t]{@{}c@{}}\rule{0pt}{12pt}7m\,49.2s\\[-.2pt]{\fontsize{5.8}{6.4}\selectfont \textcolor{BrickRed}{$ \blacktriangledown\,50.3\times $}}\\[-.2pt]7m\,45.6s\\[-.2pt]{\fontsize{5.8}{6.4}\selectfont \textcolor{BrickRed}{$ \blacktriangledown\,55.5\times $}}\end{tabular} & \hspace*{1pt}\begin{tabular}[t]{@{}c@{}}\rule{0pt}{12pt}17m\,0.3s\\[-.2pt]{\fontsize{5.8}{6.4}\selectfont \textcolor{BrickRed}{$ \blacktriangledown\,116.8\times $}}\\[-.2pt]16m\,57.4s\\[-.2pt]{\fontsize{5.8}{6.4}\selectfont \textcolor{BrickRed}{$ \blacktriangledown\,144.8\times $}}\end{tabular} & \hspace*{1pt}\begin{tabular}[t]{@{}c@{}}\rule{0pt}{12pt}29m\,26.3s\\[-.2pt]{\fontsize{5.8}{6.4}\selectfont \textcolor{BrickRed}{$ \blacktriangledown\,111.7\times $}}\\[-.2pt]29m\,20.3s\\[-.2pt]{\fontsize{5.8}{6.4}\selectfont \textcolor{BrickRed}{$ \blacktriangledown\,138.2\times $}}\end{tabular} & \hspace*{1pt}\begin{tabular}[t]{@{}c@{}}\rule{0pt}{12pt}15m\,2.9s\\[-.2pt]{\fontsize{5.8}{6.4}\selectfont \textcolor{BrickRed}{$ \blacktriangledown\,48.8\times $}}\\[-.2pt]14m\,59.6s\\[-.2pt]{\fontsize{5.8}{6.4}\selectfont \textcolor{BrickRed}{$ \blacktriangledown\,52.8\times $}}\end{tabular} & \hspace*{1pt}\begin{tabular}[t]{@{}c@{}}\rule{0pt}{12pt}2m\,53.6s\\[-.2pt]{\fontsize{5.8}{6.4}\selectfont \textcolor{BrickRed}{$ \blacktriangledown\,4.1\times $}}\\[-.2pt]2m\,15.9s\\[-.2pt]{\fontsize{5.8}{6.4}\selectfont \textcolor{BrickRed}{$ \blacktriangledown\,3.3\times $}}\end{tabular} & \hspace*{1pt}\begin{tabular}[t]{@{}c@{}}\rule{0pt}{12pt}2m\,44s\\[-.2pt]{\fontsize{5.8}{6.4}\selectfont \textcolor{BrickRed}{$ \blacktriangledown\,3.8\times $}}\\[-.2pt]2m\,12.1s\\[-.2pt]{\fontsize{5.8}{6.4}\selectfont \textcolor{BrickRed}{$ \blacktriangledown\,3.2\times $}}\end{tabular} & \textcolor{gray}{\ensuremath{\cdot}} & \textcolor{gray}{\ensuremath{\cdot}} & \hspace*{1pt}\begin{tabular}[t]{@{}c@{}}\rule{0pt}{12pt}6m\,46.2s\\[-.2pt]{\fontsize{5.8}{6.4}\selectfont \textcolor{BrickRed}{$ \blacktriangledown\,11.6\times $}}\\[-.2pt]2m\,28.4s\\[-.2pt]{\fontsize{5.8}{6.4}\selectfont \textcolor{BrickRed}{$ \blacktriangledown\,4.4\times $}}\end{tabular} & \hspace*{1pt}\begin{tabular}[t]{@{}c@{}}\rule{0pt}{12pt}6m\,33.5s\\[-.2pt]{\fontsize{5.8}{6.4}\selectfont \textcolor{BrickRed}{$ \blacktriangledown\,11.0\times $}}\\[-.2pt]2m\,49.4s\\[-.2pt]{\fontsize{5.8}{6.4}\selectfont \textcolor{BrickRed}{$ \blacktriangledown\,5.0\times $}}\end{tabular} & \textcolor{gray}{\ensuremath{\cdot}} & \textcolor{gray}{\ensuremath{\cdot}}\\[3pt]\hline
\BenchmarkSourceLabel{bench:aws-2-k20-w15:row:mpf-simdref-simd-ksw2}{R46}\shortstack[l]{{\color{gray}R46}} & \begin{tabular}[t]{@{}l@{}}\rule{0pt}{12pt}MPF / SIMD-ref\\SIMD-DP / KSW2\end{tabular} & \hspace*{1pt}\begin{tabular}[t]{@{}c@{}}\rule{0pt}{12pt}5m\,50.5s\\[-.2pt]{\fontsize{5.8}{6.4}\selectfont \textcolor{BrickRed}{$ \blacktriangledown\,17.9\times $}}\\[-.2pt]5m\,46.9s\\[-.2pt]{\fontsize{5.8}{6.4}\selectfont \textcolor{BrickRed}{$ \blacktriangledown\,18.6\times $}}\end{tabular} & \hspace*{1pt}\begin{tabular}[t]{@{}c@{}}\rule{0pt}{12pt}12m\,30.9s\\[-.2pt]{\fontsize{5.8}{6.4}\selectfont \textcolor{BrickRed}{$ \blacktriangledown\,41.3\times $}}\\[-.2pt]12m\,27.6s\\[-.2pt]{\fontsize{5.8}{6.4}\selectfont \textcolor{BrickRed}{$ \blacktriangledown\,45.3\times $}}\end{tabular} & \hspace*{1pt}\begin{tabular}[t]{@{}c@{}}\rule{0pt}{12pt}22m\,56s\\[-.2pt]{\fontsize{5.8}{6.4}\selectfont \textcolor{BrickRed}{$ \blacktriangledown\,32.4\times $}}\\[-.2pt]22m\,50s\\[-.2pt]{\fontsize{5.8}{6.4}\selectfont \textcolor{BrickRed}{$ \blacktriangledown\,34.8\times $}}\end{tabular} & \hspace*{1pt}\begin{tabular}[t]{@{}c@{}}\rule{0pt}{12pt}11m\,32.8s\\[-.2pt]{\fontsize{5.8}{6.4}\selectfont \textcolor{BrickRed}{$ \blacktriangledown\,1.8\times $}}\\[-.2pt]11m\,29.4s\\[-.2pt]{\fontsize{5.8}{6.4}\selectfont \textcolor{BrickRed}{$ \blacktriangledown\,1.9\times $}}\end{tabular} & \hspace*{1pt}\begin{tabular}[t]{@{}c@{}}\rule{0pt}{12pt}6m\,4.7s\\[-.2pt]{\fontsize{5.8}{6.4}\selectfont \textcolor{BrickRed}{$ \blacktriangledown\,39.1\times $}}\\[-.2pt]6m\,1.1s\\[-.2pt]{\fontsize{5.8}{6.4}\selectfont \textcolor{BrickRed}{$ \blacktriangledown\,43.1\times $}}\end{tabular} & \hspace*{1pt}\begin{tabular}[t]{@{}c@{}}\rule{0pt}{12pt}13m\,22.2s\\[-.2pt]{\fontsize{5.8}{6.4}\selectfont \textcolor{BrickRed}{$ \blacktriangledown\,91.8\times $}}\\[-.2pt]13m\,19.2s\\[-.2pt]{\fontsize{5.8}{6.4}\selectfont \textcolor{BrickRed}{$ \blacktriangledown\,113.7\times $}}\end{tabular} & \hspace*{1pt}\begin{tabular}[t]{@{}c@{}}\rule{0pt}{12pt}25m\,24.5s\\[-.2pt]{\fontsize{5.8}{6.4}\selectfont \textcolor{BrickRed}{$ \blacktriangledown\,96.4\times $}}\\[-.2pt]25m\,18.3s\\[-.2pt]{\fontsize{5.8}{6.4}\selectfont \textcolor{BrickRed}{$ \blacktriangledown\,119.2\times $}}\end{tabular} & \hspace*{1pt}\begin{tabular}[t]{@{}c@{}}\rule{0pt}{12pt}12m\,52.2s\\[-.2pt]{\fontsize{5.8}{6.4}\selectfont \textcolor{BrickRed}{$ \blacktriangledown\,41.8\times $}}\\[-.2pt]12m\,48.8s\\[-.2pt]{\fontsize{5.8}{6.4}\selectfont \textcolor{BrickRed}{$ \blacktriangledown\,45.2\times $}}\end{tabular} & \hspace*{1pt}\begin{tabular}[t]{@{}c@{}}\rule{0pt}{12pt}2m\,40.4s\\[-.2pt]{\fontsize{5.8}{6.4}\selectfont \textcolor{BrickRed}{$ \blacktriangledown\,3.8\times $}}\\[-.2pt]2m\,0.9s\\[-.2pt]{\fontsize{5.8}{6.4}\selectfont \textcolor{BrickRed}{$ \blacktriangledown\,2.9\times $}}\end{tabular} & \hspace*{1pt}\begin{tabular}[t]{@{}c@{}}\rule{0pt}{12pt}2m\,41.5s\\[-.2pt]{\fontsize{5.8}{6.4}\selectfont \textcolor{BrickRed}{$ \blacktriangledown\,3.8\times $}}\\[-.2pt]2m\,8.1s\\[-.2pt]{\fontsize{5.8}{6.4}\selectfont \textcolor{BrickRed}{$ \blacktriangledown\,3.1\times $}}\end{tabular} & \textcolor{gray}{\ensuremath{\cdot}} & \textcolor{gray}{\ensuremath{\cdot}} & \hspace*{1pt}\begin{tabular}[t]{@{}c@{}}\rule{0pt}{12pt}6m\,16.9s\\[-.2pt]{\fontsize{5.8}{6.4}\selectfont \textcolor{BrickRed}{$ \blacktriangledown\,10.8\times $}}\\[-.2pt]2m\,16s\\[-.2pt]{\fontsize{5.8}{6.4}\selectfont \textcolor{BrickRed}{$ \blacktriangledown\,4.0\times $}}\end{tabular} & \hspace*{1pt}\begin{tabular}[t]{@{}c@{}}\rule{0pt}{12pt}6m\,11s\\[-.2pt]{\fontsize{5.8}{6.4}\selectfont \textcolor{BrickRed}{$ \blacktriangledown\,10.3\times $}}\\[-.2pt]2m\,39.5s\\[-.2pt]{\fontsize{5.8}{6.4}\selectfont \textcolor{BrickRed}{$ \blacktriangledown\,4.7\times $}}\end{tabular} & \textcolor{gray}{\ensuremath{\cdot}} & \textcolor{gray}{\ensuremath{\cdot}}\\[3pt]\hline
\BenchmarkSourceLabel{bench:aws-2-k20-w15:row:mpf-simdref-rmq-ksw2}{R47}\shortstack[l]{{\color{gray}R47}} & \begin{tabular}[t]{@{}l@{}}\rule{0pt}{12pt}MPF / SIMD-ref\\RMQ / KSW2\end{tabular} & \hspace*{1pt}\begin{tabular}[t]{@{}c@{}}\rule{0pt}{12pt}5m\,50.6s\\[-.2pt]{\fontsize{5.8}{6.4}\selectfont \textcolor{BrickRed}{$ \blacktriangledown\,17.9\times $}}\\[-.2pt]5m\,47s\\[-.2pt]{\fontsize{5.8}{6.4}\selectfont \textcolor{BrickRed}{$ \blacktriangledown\,18.6\times $}}\end{tabular} & \hspace*{1pt}\begin{tabular}[t]{@{}c@{}}\rule{0pt}{12pt}12m\,12.8s\\[-.2pt]{\fontsize{5.8}{6.4}\selectfont \textcolor{BrickRed}{$ \blacktriangledown\,40.3\times $}}\\[-.2pt]12m\,9.8s\\[-.2pt]{\fontsize{5.8}{6.4}\selectfont \textcolor{BrickRed}{$ \blacktriangledown\,44.2\times $}}\end{tabular} & \hspace*{1pt}\begin{tabular}[t]{@{}c@{}}\rule{0pt}{12pt}23m\,7.9s\\[-.2pt]{\fontsize{5.8}{6.4}\selectfont \textcolor{BrickRed}{$ \blacktriangledown\,32.6\times $}}\\[-.2pt]23m\,1.9s\\[-.2pt]{\fontsize{5.8}{6.4}\selectfont \textcolor{BrickRed}{$ \blacktriangledown\,35.1\times $}}\end{tabular} & \hspace*{1pt}\begin{tabular}[t]{@{}c@{}}\rule{0pt}{12pt}10m\,52.1s\\[-.2pt]{\fontsize{5.8}{6.4}\selectfont \textcolor{BrickRed}{$ \blacktriangledown\,1.7\times $}}\\[-.2pt]10m\,49.7s\\[-.2pt]{\fontsize{5.8}{6.4}\selectfont \textcolor{BrickRed}{$ \blacktriangledown\,1.8\times $}}\end{tabular} & \hspace*{1pt}\begin{tabular}[t]{@{}c@{}}\rule{0pt}{12pt}6m\,6.7s\\[-.2pt]{\fontsize{5.8}{6.4}\selectfont \textcolor{BrickRed}{$ \blacktriangledown\,39.3\times $}}\\[-.2pt]6m\,3s\\[-.2pt]{\fontsize{5.8}{6.4}\selectfont \textcolor{BrickRed}{$ \blacktriangledown\,43.3\times $}}\end{tabular} & \hspace*{1pt}\begin{tabular}[t]{@{}c@{}}\rule{0pt}{12pt}13m\,19.1s\\[-.2pt]{\fontsize{5.8}{6.4}\selectfont \textcolor{BrickRed}{$ \blacktriangledown\,91.5\times $}}\\[-.2pt]13m\,16.2s\\[-.2pt]{\fontsize{5.8}{6.4}\selectfont \textcolor{BrickRed}{$ \blacktriangledown\,113.3\times $}}\end{tabular} & \hspace*{1pt}\begin{tabular}[t]{@{}c@{}}\rule{0pt}{12pt}24m\,33.5s\\[-.2pt]{\fontsize{5.8}{6.4}\selectfont \textcolor{BrickRed}{$ \blacktriangledown\,93.2\times $}}\\[-.2pt]24m\,27.3s\\[-.2pt]{\fontsize{5.8}{6.4}\selectfont \textcolor{BrickRed}{$ \blacktriangledown\,115.2\times $}}\end{tabular} & \hspace*{1pt}\begin{tabular}[t]{@{}c@{}}\rule{0pt}{12pt}12m\,1.2s\\[-.2pt]{\fontsize{5.8}{6.4}\selectfont \textcolor{BrickRed}{$ \blacktriangledown\,39.0\times $}}\\[-.2pt]11m\,57.8s\\[-.2pt]{\fontsize{5.8}{6.4}\selectfont \textcolor{BrickRed}{$ \blacktriangledown\,42.2\times $}}\end{tabular} & \hspace*{1pt}\begin{tabular}[t]{@{}c@{}}\rule{0pt}{12pt}2m\,41.8s\\[-.2pt]{\fontsize{5.8}{6.4}\selectfont \textcolor{BrickRed}{$ \blacktriangledown\,3.8\times $}}\\[-.2pt]2m\,2.1s\\[-.2pt]{\fontsize{5.8}{6.4}\selectfont \textcolor{BrickRed}{$ \blacktriangledown\,2.9\times $}}\end{tabular} & \hspace*{1pt}\begin{tabular}[t]{@{}c@{}}\rule{0pt}{12pt}2m\,37.2s\\[-.2pt]{\fontsize{5.8}{6.4}\selectfont \textcolor{BrickRed}{$ \blacktriangledown\,3.7\times $}}\\[-.2pt]2m\,3.8s\\[-.2pt]{\fontsize{5.8}{6.4}\selectfont \textcolor{BrickRed}{$ \blacktriangledown\,3.0\times $}}\end{tabular} & \textcolor{gray}{\ensuremath{\cdot}} & \textcolor{gray}{\ensuremath{\cdot}} & \hspace*{1pt}\begin{tabular}[t]{@{}c@{}}\rule{0pt}{12pt}6m\,18.7s\\[-.2pt]{\fontsize{5.8}{6.4}\selectfont \textcolor{BrickRed}{$ \blacktriangledown\,10.8\times $}}\\[-.2pt]2m\,18s\\[-.2pt]{\fontsize{5.8}{6.4}\selectfont \textcolor{BrickRed}{$ \blacktriangledown\,4.1\times $}}\end{tabular} & \hspace*{1pt}\begin{tabular}[t]{@{}c@{}}\rule{0pt}{12pt}6m\,21.7s\\[-.2pt]{\fontsize{5.8}{6.4}\selectfont \textcolor{BrickRed}{$ \blacktriangledown\,10.6\times $}}\\[-.2pt]2m\,44.6s\\[-.2pt]{\fontsize{5.8}{6.4}\selectfont \textcolor{BrickRed}{$ \blacktriangledown\,4.8\times $}}\end{tabular} & \textcolor{gray}{\ensuremath{\cdot}} & \textcolor{gray}{\ensuremath{\cdot}}\\[3pt]\hline
\BenchmarkSourceLabel{bench:aws-2-k20-w15:row:mpf-simdref-simd-block}{R48}\shortstack[l]{{\color{gray}R48}} & \begin{tabular}[t]{@{}l@{}}\rule{0pt}{12pt}MPF / SIMD-ref\\SIMD-DP / Block\end{tabular} & \hspace*{1pt}\begin{tabular}[t]{@{}c@{}}\rule{0pt}{12pt}5m\,50.7s\\[-.2pt]{\fontsize{5.8}{6.4}\selectfont \textcolor{BrickRed}{$ \blacktriangledown\,17.9\times $}}\\[-.2pt]5m\,49.9s\\[-.2pt]{\fontsize{5.8}{6.4}\selectfont \textcolor{BrickRed}{$ \blacktriangledown\,18.8\times $}}\end{tabular} & \hspace*{1pt}\begin{tabular}[t]{@{}c@{}}\rule{0pt}{12pt}12m\,26.8s\\[-.2pt]{\fontsize{5.8}{6.4}\selectfont \textcolor{BrickRed}{$ \blacktriangledown\,41.0\times $}}\\[-.2pt]12m\,25.9s\\[-.2pt]{\fontsize{5.8}{6.4}\selectfont \textcolor{BrickRed}{$ \blacktriangledown\,45.2\times $}}\end{tabular} & \hspace*{1pt}\begin{tabular}[t]{@{}c@{}}\rule{0pt}{12pt}22m\,58.7s\\[-.2pt]{\fontsize{5.8}{6.4}\selectfont \textcolor{BrickRed}{$ \blacktriangledown\,32.4\times $}}\\[-.2pt]22m\,57.3s\\[-.2pt]{\fontsize{5.8}{6.4}\selectfont \textcolor{BrickRed}{$ \blacktriangledown\,35.0\times $}}\end{tabular} & \hspace*{1pt}\begin{tabular}[t]{@{}c@{}}\rule{0pt}{12pt}13m\,2.2s\\[-.2pt]{\fontsize{5.8}{6.4}\selectfont \textcolor{BrickRed}{$ \blacktriangledown\,2.0\times $}}\\[-.2pt]13m\,1.5s\\[-.2pt]{\fontsize{5.8}{6.4}\selectfont \textcolor{BrickRed}{$ \blacktriangledown\,2.2\times $}}\end{tabular} & \hspace*{1pt}\begin{tabular}[t]{@{}c@{}}\rule{0pt}{12pt}6m\,8.9s\\[-.2pt]{\fontsize{5.8}{6.4}\selectfont \textcolor{BrickRed}{$ \blacktriangledown\,39.6\times $}}\\[-.2pt]6m\,5.3s\\[-.2pt]{\fontsize{5.8}{6.4}\selectfont \textcolor{BrickRed}{$ \blacktriangledown\,43.6\times $}}\end{tabular} & \hspace*{1pt}\begin{tabular}[t]{@{}c@{}}\rule{0pt}{12pt}13m\,31.2s\\[-.2pt]{\fontsize{5.8}{6.4}\selectfont \textcolor{BrickRed}{$ \blacktriangledown\,92.8\times $}}\\[-.2pt]13m\,28.4s\\[-.2pt]{\fontsize{5.8}{6.4}\selectfont \textcolor{BrickRed}{$ \blacktriangledown\,115.0\times $}}\end{tabular} & \hspace*{1pt}\begin{tabular}[t]{@{}c@{}}\rule{0pt}{12pt}24m\,58.1s\\[-.2pt]{\fontsize{5.8}{6.4}\selectfont \textcolor{BrickRed}{$ \blacktriangledown\,94.7\times $}}\\[-.2pt]24m\,52.4s\\[-.2pt]{\fontsize{5.8}{6.4}\selectfont \textcolor{BrickRed}{$ \blacktriangledown\,117.1\times $}}\end{tabular} & \hspace*{1pt}\begin{tabular}[t]{@{}c@{}}\rule{0pt}{12pt}13m\,3.8s\\[-.2pt]{\fontsize{5.8}{6.4}\selectfont \textcolor{BrickRed}{$ \blacktriangledown\,42.4\times $}}\\[-.2pt]13m\,0.9s\\[-.2pt]{\fontsize{5.8}{6.4}\selectfont \textcolor{BrickRed}{$ \blacktriangledown\,45.9\times $}}\end{tabular} & \hspace*{1pt}\begin{tabular}[t]{@{}c@{}}\rule{0pt}{12pt}2m\,38.4s\\[-.2pt]{\fontsize{5.8}{6.4}\selectfont \textcolor{BrickRed}{$ \blacktriangledown\,3.7\times $}}\\[-.2pt]2m\,0.1s\\[-.2pt]{\fontsize{5.8}{6.4}\selectfont \textcolor{BrickRed}{$ \blacktriangledown\,2.9\times $}}\end{tabular} & \hspace*{1pt}\begin{tabular}[t]{@{}c@{}}\rule{0pt}{12pt}2m\,35.4s\\[-.2pt]{\fontsize{5.8}{6.4}\selectfont \textcolor{BrickRed}{$ \blacktriangledown\,3.6\times $}}\\[-.2pt]2m\,4.3s\\[-.2pt]{\fontsize{5.8}{6.4}\selectfont \textcolor{BrickRed}{$ \blacktriangledown\,3.0\times $}}\end{tabular} & \textcolor{gray}{\ensuremath{\cdot}} & \hspace*{1pt}\begin{tabular}[t]{@{}c@{}}\rule{0pt}{12pt}2m\,49.5s\\[-.2pt]{\fontsize{5.8}{6.4}\selectfont \textcolor{ForestGreen}{$ \blacktriangle\,1.7\times $}}\\[-.2pt]2m\,21.9s\\[-.2pt]{\fontsize{5.8}{6.4}\selectfont \textcolor{ForestGreen}{$ \blacktriangle\,2.0\times $}}\end{tabular} & \hspace*{1pt}\begin{tabular}[t]{@{}c@{}}\rule{0pt}{12pt}6m\,15.5s\\[-.2pt]{\fontsize{5.8}{6.4}\selectfont \textcolor{BrickRed}{$ \blacktriangledown\,10.7\times $}}\\[-.2pt]2m\,16.4s\\[-.2pt]{\fontsize{5.8}{6.4}\selectfont \textcolor{BrickRed}{$ \blacktriangledown\,4.0\times $}}\end{tabular} & \hspace*{1pt}\begin{tabular}[t]{@{}c@{}}\rule{0pt}{12pt}5m\,52.6s\\[-.2pt]{\fontsize{5.8}{6.4}\selectfont \textcolor{BrickRed}{$ \blacktriangledown\,9.8\times $}}\\[-.2pt]2m\,26.6s\\[-.2pt]{\fontsize{5.8}{6.4}\selectfont \textcolor{BrickRed}{$ \blacktriangledown\,4.3\times $}}\end{tabular} & \textcolor{gray}{\ensuremath{\cdot}} & \hspace*{1pt}\begin{tabular}[t]{@{}c@{}}\rule{0pt}{12pt}12m\,1.9s\\[-.2pt]{\fontsize{5.8}{6.4}\selectfont \textcolor{BrickRed}{$ \blacktriangledown\,12.0\times $}}\\[-.2pt]3m\,7.8s\\[-.2pt]{\fontsize{5.8}{6.4}\selectfont \textcolor{BrickRed}{$ \blacktriangledown\,3.2\times $}}\end{tabular}\\[3pt]\hline
\BenchmarkSourceLabel{bench:aws-2-k20-w15:row:mpf-simdref-scalar-block}{R49}\shortstack[l]{{\color{gray}R49}} & \begin{tabular}[t]{@{}l@{}}\rule{0pt}{12pt}MPF / SIMD-ref\\Scalar-DP / Block\end{tabular} & \hspace*{1pt}\begin{tabular}[t]{@{}c@{}}\rule{0pt}{12pt}5m\,51.9s\\[-.2pt]{\fontsize{5.8}{6.4}\selectfont \textcolor{BrickRed}{$ \blacktriangledown\,18.0\times $}}\\[-.2pt]5m\,51.1s\\[-.2pt]{\fontsize{5.8}{6.4}\selectfont \textcolor{BrickRed}{$ \blacktriangledown\,18.9\times $}}\end{tabular} & \hspace*{1pt}\begin{tabular}[t]{@{}c@{}}\rule{0pt}{12pt}12m\,10.6s\\[-.2pt]{\fontsize{5.8}{6.4}\selectfont \textcolor{BrickRed}{$ \blacktriangledown\,40.1\times $}}\\[-.2pt]12m\,9.8s\\[-.2pt]{\fontsize{5.8}{6.4}\selectfont \textcolor{BrickRed}{$ \blacktriangledown\,44.2\times $}}\end{tabular} & \hspace*{1pt}\begin{tabular}[t]{@{}c@{}}\rule{0pt}{12pt}23m\,42.4s\\[-.2pt]{\fontsize{5.8}{6.4}\selectfont \textcolor{BrickRed}{$ \blacktriangledown\,33.5\times $}}\\[-.2pt]23m\,41s\\[-.2pt]{\fontsize{5.8}{6.4}\selectfont \textcolor{BrickRed}{$ \blacktriangledown\,36.1\times $}}\end{tabular} & \hspace*{1pt}\begin{tabular}[t]{@{}c@{}}\rule{0pt}{12pt}11m\,56.4s\\[-.2pt]{\fontsize{5.8}{6.4}\selectfont \textcolor{BrickRed}{$ \blacktriangledown\,1.8\times $}}\\[-.2pt]11m\,55.6s\\[-.2pt]{\fontsize{5.8}{6.4}\selectfont \textcolor{BrickRed}{$ \blacktriangledown\,2.0\times $}}\end{tabular} & \hspace*{1pt}\begin{tabular}[t]{@{}c@{}}\rule{0pt}{12pt}6m\,3.8s\\[-.2pt]{\fontsize{5.8}{6.4}\selectfont \textcolor{BrickRed}{$ \blacktriangledown\,39.0\times $}}\\[-.2pt]6m\,0.1s\\[-.2pt]{\fontsize{5.8}{6.4}\selectfont \textcolor{BrickRed}{$ \blacktriangledown\,43.0\times $}}\end{tabular} & \hspace*{1pt}\begin{tabular}[t]{@{}c@{}}\rule{0pt}{12pt}13m\,15.2s\\[-.2pt]{\fontsize{5.8}{6.4}\selectfont \textcolor{BrickRed}{$ \blacktriangledown\,91.0\times $}}\\[-.2pt]13m\,12.6s\\[-.2pt]{\fontsize{5.8}{6.4}\selectfont \textcolor{BrickRed}{$ \blacktriangledown\,112.8\times $}}\end{tabular} & \hspace*{1pt}\begin{tabular}[t]{@{}c@{}}\rule{0pt}{12pt}24m\,21.9s\\[-.2pt]{\fontsize{5.8}{6.4}\selectfont \textcolor{BrickRed}{$ \blacktriangledown\,92.5\times $}}\\[-.2pt]24m\,16.2s\\[-.2pt]{\fontsize{5.8}{6.4}\selectfont \textcolor{BrickRed}{$ \blacktriangledown\,114.3\times $}}\end{tabular} & \hspace*{1pt}\begin{tabular}[t]{@{}c@{}}\rule{0pt}{12pt}11m\,53.1s\\[-.2pt]{\fontsize{5.8}{6.4}\selectfont \textcolor{BrickRed}{$ \blacktriangledown\,38.6\times $}}\\[-.2pt]11m\,50s\\[-.2pt]{\fontsize{5.8}{6.4}\selectfont \textcolor{BrickRed}{$ \blacktriangledown\,41.7\times $}}\end{tabular} & \hspace*{1pt}\begin{tabular}[t]{@{}c@{}}\rule{0pt}{12pt}2m\,40s\\[-.2pt]{\fontsize{5.8}{6.4}\selectfont \textcolor{BrickRed}{$ \blacktriangledown\,3.8\times $}}\\[-.2pt]2m\,1.7s\\[-.2pt]{\fontsize{5.8}{6.4}\selectfont \textcolor{BrickRed}{$ \blacktriangledown\,2.9\times $}}\end{tabular} & \hspace*{1pt}\begin{tabular}[t]{@{}c@{}}\rule{0pt}{12pt}2m\,29.9s\\[-.2pt]{\fontsize{5.8}{6.4}\selectfont \textcolor{BrickRed}{$ \blacktriangledown\,3.5\times $}}\\[-.2pt]1m\,58.9s\\[-.2pt]{\fontsize{5.8}{6.4}\selectfont \textcolor{BrickRed}{$ \blacktriangledown\,2.9\times $}}\end{tabular} & \textcolor{gray}{\ensuremath{\cdot}} & \textcolor{gray}{\ensuremath{\cdot}} & \hspace*{1pt}\begin{tabular}[t]{@{}c@{}}\rule{0pt}{12pt}6m\,18.4s\\[-.2pt]{\fontsize{5.8}{6.4}\selectfont \textcolor{BrickRed}{$ \blacktriangledown\,10.8\times $}}\\[-.2pt]2m\,19.3s\\[-.2pt]{\fontsize{5.8}{6.4}\selectfont \textcolor{BrickRed}{$ \blacktriangledown\,4.1\times $}}\end{tabular} & \hspace*{1pt}\begin{tabular}[t]{@{}c@{}}\rule{0pt}{12pt}6m\,19.9s\\[-.2pt]{\fontsize{5.8}{6.4}\selectfont \textcolor{BrickRed}{$ \blacktriangledown\,10.6\times $}}\\[-.2pt]2m\,45.3s\\[-.2pt]{\fontsize{5.8}{6.4}\selectfont \textcolor{BrickRed}{$ \blacktriangledown\,4.8\times $}}\end{tabular} & \textcolor{gray}{\ensuremath{\cdot}} & \textcolor{gray}{\ensuremath{\cdot}}\\[3pt]\hline
\BenchmarkSourceLabel{bench:aws-2-k20-w15:row:vulkan-unavailable}{R50}\shortstack[l]{{\color{gray}R50}\\\rule{0pt}{12pt}\textbf{BioseqDB}} & Vulkan (all modes) & \textcolor{gray}{---\textsuperscript{g}} & \textcolor{gray}{---\textsuperscript{g}} & \textcolor{gray}{---\textsuperscript{g}} & \textcolor{gray}{---\textsuperscript{g}} & \textcolor{gray}{---\textsuperscript{g}} & \textcolor{gray}{---\textsuperscript{g}} & \textcolor{gray}{---\textsuperscript{g}} & \textcolor{gray}{---\textsuperscript{g}} & \textcolor{gray}{---\textsuperscript{g}} & \textcolor{gray}{---\textsuperscript{g}} & \textcolor{gray}{---\textsuperscript{g}} & \textcolor{gray}{---\textsuperscript{g}} & \textcolor{gray}{---\textsuperscript{g}} & \textcolor{gray}{---\textsuperscript{g}} & \textcolor{gray}{---\textsuperscript{g}} & \textcolor{gray}{---\textsuperscript{g}}\\[3pt]\hline
\end{tabular}\par}
\smallskip{\footnotesize\noindent Each cell: A and B times, $\cdot$: no measurement. ---\textsuperscript{g}: GPU unavailable. n/a: no comparison.\par
\noindent Blue: one lowest A and one lowest B measurement for each column for each configuration.\par
\noindent Gray: minimap2. Yellow: my preferred MPF / Bounded / SIMD-DP / Block configuration.\par
\noindent MPF/MPS: minimizers per fragment/sequence. SIMD-ref: naive SIMD scan. Bounded: bounded SIMD algorithm.\par
}
\end{landscape}
\restoregeometry

\clearpage\subsection{Comparison with other aligners}\label{sec:external-aligner-results}
The following adjustments were needed to run the compared aligners:
\begin{itemize}
\item \textbf{Minimap2.}
A single FASTA record must fit in 32-bit integer, so I split the 4 Gbp reference into two records of about 2 Gbp each. 
Returned positions are converted back to coordinates in the original reference. The 2 Gbp reference didn't require any spliting, so I left it as-is.
\item \textbf{Winnowmap.} I re-enabled saving its native index and changed loading to rebuild the weighting Bloom filter from the saved list of frequent k-mers produced by Meryl (Meryl is the tool Winnomap uses for extracting k-mers).
The 4 Gbp references used index batches of 2 Mbp, because otherwise Winnowmap failed even with sufficient SWAP. This means that Winnowmap had far better conditions to minimap2.
Sadly comparison wasn't possible otherwise, beccause I encountered repeatable memory errors.
\item \textbf{blastn.} BLAST used sequences unchanged. It has no second branch and scores N symbol differently - as $-3$ rather than $-1$, but it had no impact on results.
\item \textbf{LexicMap.} It uses 28-bit coordinates meaning that the sequences was required to be split into 268Mbp parts. 
On machines with less than 10 GiB RAM, the 4 Gbp reference was indexed one record per batch, because otherwise it produced an error. LexicMap by default replaces N with A in target, which also didn't affect the quality of results.
\end{itemize}
\newgeometry{left=10mm,right=10mm,top=12mm,bottom=12mm}
\begin{landscape}
\clearpage
\refstepcounter{table}\label{tab:bench:external-4-k15-w10}\label{tab:external-aligners}
\noindent{\large\bfseries (F): Table \thetable: aligner comparison, 4 Gbp, k=15, w=10\hfill 1/4}\par\smallskip
{\fontsize{7.1}{8.1}\selectfont\setlength{\tabcolsep}{1pt}\setlength{\fboxsep}{2pt}\renewcommand{\arraystretch}{1.12}
\noindent\begin{tabular}{@{}p{15mm}p{26mm}|p{13.55mm}|p{13.55mm}|p{13.55mm}|p{13.55mm}|p{13.55mm}|p{13.55mm}|p{13.55mm}|p{13.55mm}|p{13.55mm}|p{13.55mm}|p{13.55mm}|p{13.55mm}|p{13.55mm}|p{13.55mm}|p{13.55mm}|p{13.55mm}|@{}}\hline
Aligner & Configuration & \multicolumn{8}{c|}{\textbf{Index}} & \multicolumn{8}{c|}{\textbf{No index}}\\
 &  & \multicolumn{8}{c|}{\textbf{4 Gbp}} & \multicolumn{8}{c|}{\textbf{4 Gbp}}\\
 &  & \multicolumn{8}{c|}{\textbf{k=15, w=10}} & \multicolumn{8}{c|}{\textbf{k=15, w=10}}\\
 &  & \multicolumn{4}{c|}{\textbf{DNA4}} & \multicolumn{4}{c|}{\textbf{DNA5}} & \multicolumn{4}{c|}{\textbf{DNA4}} & \multicolumn{4}{c|}{\textbf{DNA5}}\\
 &  & \multicolumn{1}{c|}{\textbf{M7A}} & \multicolumn{1}{c|}{\textbf{M7I}} & \multicolumn{1}{c|}{\textbf{M4}} & \multicolumn{1}{c|}{\textbf{C7G}} & \multicolumn{1}{c|}{\textbf{M7A}} & \multicolumn{1}{c|}{\textbf{M7I}} & \multicolumn{1}{c|}{\textbf{M4}} & \multicolumn{1}{c|}{\textbf{C7G}} & \multicolumn{1}{c|}{\textbf{M7A}} & \multicolumn{1}{c|}{\textbf{M7I}} & \multicolumn{1}{c|}{\textbf{M4}} & \multicolumn{1}{c|}{\textbf{C7G}} & \multicolumn{1}{c|}{\textbf{M7A}} & \multicolumn{1}{c|}{\textbf{M7I}} & \multicolumn{1}{c|}{\textbf{M4}} & \multicolumn{1}{c|}{\textbf{C7G}}\\
 & & \BenchmarkSourceLabel{bench:external-4-k15-w10:column:index-4-k15-w10-dna4-m7a}{C1}{\color{gray}C1} & \BenchmarkSourceLabel{bench:external-4-k15-w10:column:index-4-k15-w10-dna4-m7i}{C2}{\color{gray}C2} & \BenchmarkSourceLabel{bench:external-4-k15-w10:column:index-4-k15-w10-dna4-m4}{C3}{\color{gray}C3} & \BenchmarkSourceLabel{bench:external-4-k15-w10:column:index-4-k15-w10-dna4-c7g}{C4}{\color{gray}C4} & \BenchmarkSourceLabel{bench:external-4-k15-w10:column:index-4-k15-w10-dna5-m7a}{C5}{\color{gray}C5} & \BenchmarkSourceLabel{bench:external-4-k15-w10:column:index-4-k15-w10-dna5-m7i}{C6}{\color{gray}C6} & \BenchmarkSourceLabel{bench:external-4-k15-w10:column:index-4-k15-w10-dna5-m4}{C7}{\color{gray}C7} & \BenchmarkSourceLabel{bench:external-4-k15-w10:column:index-4-k15-w10-dna5-c7g}{C8}{\color{gray}C8} & \BenchmarkSourceLabel{bench:external-4-k15-w10:column:noindex-4-k15-w10-dna4-m7a}{C9}{\color{gray}C9} & \BenchmarkSourceLabel{bench:external-4-k15-w10:column:noindex-4-k15-w10-dna4-m7i}{C10}{\color{gray}C10} & \BenchmarkSourceLabel{bench:external-4-k15-w10:column:noindex-4-k15-w10-dna4-m4}{C11}{\color{gray}C11} & \BenchmarkSourceLabel{bench:external-4-k15-w10:column:noindex-4-k15-w10-dna4-c7g}{C12}{\color{gray}C12} & \BenchmarkSourceLabel{bench:external-4-k15-w10:column:noindex-4-k15-w10-dna5-m7a}{C13}{\color{gray}C13} & \BenchmarkSourceLabel{bench:external-4-k15-w10:column:noindex-4-k15-w10-dna5-m7i}{C14}{\color{gray}C14} & \BenchmarkSourceLabel{bench:external-4-k15-w10:column:noindex-4-k15-w10-dna5-m4}{C15}{\color{gray}C15} & \BenchmarkSourceLabel{bench:external-4-k15-w10:column:noindex-4-k15-w10-dna5-c7g}{C16}{\color{gray}C16}\\\hline
\rowcolor{BenchBaseline}[0.5pt][0.5pt]\BenchmarkSourceLabel{bench:external-4-k15-w10:row:minimap2}{R1}\shortstack[l]{{\color{gray}R1}\\minimap2} & Baseline & \hspace*{1pt}\begin{tabular}[t]{@{}c@{}}\rule{0pt}{12pt}36.2s\\[-.2pt]{\fontsize{5.8}{6.4}\selectfont \textcolor{BrickRed}{$ \blacktriangledown\,3.6\times $}}\\[-.2pt]35s\\[-.2pt]{\fontsize{5.8}{6.4}\selectfont \textcolor{BrickRed}{$ \blacktriangledown\,3.9\times $}}\end{tabular} & \hspace*{1pt}\begin{tabular}[t]{@{}c@{}}\rule{0pt}{12pt}38.1s\\[-.2pt]{\fontsize{5.8}{6.4}\selectfont \textcolor{BrickRed}{$ \blacktriangledown\,2.9\times $}}\\[-.2pt]36s\\[-.2pt]{\fontsize{5.8}{6.4}\selectfont \textcolor{BrickRed}{$ \blacktriangledown\,3.0\times $}}\end{tabular} & \hspace*{1pt}\begin{tabular}[t]{@{}c@{}}\rule{0pt}{12pt}1m\,32.4s\\[-.2pt]{\fontsize{5.8}{6.4}\selectfont \textcolor{BrickRed}{$ \blacktriangledown\,9.8\times $}}\\[-.2pt]1m\,28.5s\\[-.2pt]{\fontsize{5.8}{6.4}\selectfont \textcolor{BrickRed}{$ \blacktriangledown\,12.7\times $}}\end{tabular} & \hspace*{1pt}\begin{tabular}[t]{@{}c@{}}\rule{0pt}{12pt}6m\,31.5s\\[-.2pt]{\fontsize{5.8}{6.4}\selectfont \textcolor{BrickRed}{$ \blacktriangledown\,11.2\times $}\textsuperscript{g}}\\[-.2pt]6m\,26.3s\\[-.2pt]{\fontsize{5.8}{6.4}\selectfont \textcolor{BrickRed}{$ \blacktriangledown\,11.4\times $}}\end{tabular} & \hspace*{1pt}\begin{tabular}[t]{@{}c@{}}\rule{0pt}{12pt}31.6s\\[-.2pt]{\fontsize{5.8}{6.4}\selectfont \textcolor{BrickRed}{$ \blacktriangledown\,1.4\times $}}\\[-.2pt]30.3s\\[-.2pt]{\fontsize{5.8}{6.4}\selectfont \textcolor{BrickRed}{$ \blacktriangledown\,1.6\times $}}\end{tabular} & \hspace*{1pt}\begin{tabular}[t]{@{}c@{}}\rule{0pt}{12pt}34.1s\\[-.2pt]{\fontsize{5.8}{6.4}\selectfont \textcolor{BrickRed}{$ \blacktriangledown\,1.4\times $}}\\[-.2pt]32s\\[-.2pt]{\fontsize{5.8}{6.4}\selectfont \textcolor{BrickRed}{$ \blacktriangledown\,1.5\times $}}\end{tabular} & \hspace*{1pt}\begin{tabular}[t]{@{}c@{}}\rule{0pt}{12pt}1m\,9.5s\\[-.2pt]{\fontsize{5.8}{6.4}\selectfont \textcolor{BrickRed}{$ \blacktriangledown\,28\% $}}\\[-.2pt]1m\,6.1s\\[-.2pt]{\fontsize{5.8}{6.4}\selectfont \textcolor{BrickRed}{$ \blacktriangledown\,1.4\times $}}\end{tabular} & \hspace*{1pt}\begin{tabular}[t]{@{}c@{}}\rule{0pt}{12pt}3m\,1s\\[-.2pt]{\fontsize{5.8}{6.4}\selectfont \textcolor{BrickRed}{$ \blacktriangledown\,3.6\times $}\textsuperscript{g}}\\[-.2pt]2m\,51.5s\\[-.2pt]{\fontsize{5.8}{6.4}\selectfont \textcolor{BrickRed}{$ \blacktriangledown\,3.7\times $}}\end{tabular} & \hspace*{1pt}\begin{tabular}[t]{@{}c@{}}\rule{0pt}{12pt}\colorbox{BenchBest}{\strut 2m\,55.5s}\\[-.2pt]{\fontsize{5.8}{6.4}\selectfont \textcolor{ForestGreen}{$ \blacktriangle\,22\% $}}\\[-.2pt]\colorbox{BenchBest}{\strut 2m\,54s}\\[-.2pt]{\fontsize{5.8}{6.4}\selectfont \textcolor{ForestGreen}{$ \blacktriangle\,19\% $}}\end{tabular} & \hspace*{1pt}\begin{tabular}[t]{@{}c@{}}\rule{0pt}{12pt}\colorbox{BenchBest}{\strut 2m\,50.4s}\\[-.2pt]{\fontsize{5.8}{6.4}\selectfont \textcolor{ForestGreen}{$ \blacktriangle\,28\% $}}\\[-.2pt]\colorbox{BenchBest}{\strut 2m\,48.2s}\\[-.2pt]{\fontsize{5.8}{6.4}\selectfont \textcolor{ForestGreen}{$ \blacktriangle\,25\% $}}\end{tabular} & \hspace*{1pt}\begin{tabular}[t]{@{}c@{}}\rule{0pt}{12pt}\colorbox{BenchBest}{\strut 5m\,35.7s}\\[-.2pt]{\fontsize{5.8}{6.4}\selectfont \textcolor{ForestGreen}{$ \blacktriangle\,18\% $}}\\[-.2pt]\colorbox{BenchBest}{\strut 5m\,31.9s}\\[-.2pt]{\fontsize{5.8}{6.4}\selectfont \textcolor{ForestGreen}{$ \blacktriangle\,15\% $}}\end{tabular} & \hspace*{1pt}\begin{tabular}[t]{@{}c@{}}\rule{0pt}{12pt}44m\,20.5s\\[-.2pt]{\fontsize{5.8}{6.4}\selectfont \textcolor{gray}{n/a}\textsuperscript{g}}\\[-.2pt]44m\,15.6s\\[-.2pt]{\fontsize{5.8}{6.4}\selectfont \textcolor{gray}{n/a}}\end{tabular} & \hspace*{1pt}\begin{tabular}[t]{@{}c@{}}\rule{0pt}{12pt}\colorbox{BenchBest}{\strut 1m\,47.1s}\\[-.2pt]{\fontsize{5.8}{6.4}\selectfont \textcolor{ForestGreen}{$ \blacktriangle\,5.5\times $}}\\[-.2pt]\colorbox{BenchBest}{\strut 1m\,45.8s}\\[-.2pt]{\fontsize{5.8}{6.4}\selectfont \textcolor{ForestGreen}{$ \blacktriangle\,2.1\times $}}\end{tabular} & \hspace*{1pt}\begin{tabular}[t]{@{}c@{}}\rule{0pt}{12pt}\colorbox{BenchBest}{\strut 1m\,47.2s}\\[-.2pt]{\fontsize{5.8}{6.4}\selectfont \textcolor{ForestGreen}{$ \blacktriangle\,4.7\times $}}\\[-.2pt]\colorbox{BenchBest}{\strut 1m\,45.1s}\\[-.2pt]{\fontsize{5.8}{6.4}\selectfont \textcolor{ForestGreen}{$ \blacktriangle\,2.3\times $}}\end{tabular} & \hspace*{1pt}\begin{tabular}[t]{@{}c@{}}\rule{0pt}{12pt}\colorbox{BenchBest}{\strut 3m\,40.3s}\\[-.2pt]{\fontsize{5.8}{6.4}\selectfont \textcolor{ForestGreen}{$ \blacktriangle\,4.2\times $}}\\[-.2pt]\colorbox{BenchBest}{\strut 3m\,36.6s}\\[-.2pt]{\fontsize{5.8}{6.4}\selectfont \textcolor{ForestGreen}{$ \blacktriangle\,1.7\times $}}\end{tabular} & \hspace*{1pt}\begin{tabular}[t]{@{}c@{}}\rule{0pt}{12pt}29m\,43.4s\\[-.2pt]{\fontsize{5.8}{6.4}\selectfont \textcolor{gray}{n/a}\textsuperscript{g}}\\[-.2pt]29m\,38.5s\\[-.2pt]{\fontsize{5.8}{6.4}\selectfont \textcolor{gray}{n/a}}\end{tabular}\\[3pt]\hline
\rowcolor{BenchPreferred}[0.5pt][0.5pt]\BenchmarkSourceLabel{bench:external-4-k15-w10:row:bioseq}{R2}\shortstack[l]{{\color{gray}R2}\\BioseqDB} & \begin{tabular}[t]{@{}l@{}}\rule{0pt}{12pt}MPF / Bounded\\SIMD-DP / Block\end{tabular} & \hspace*{1pt}\begin{tabular}[t]{@{}c@{}}\rule{0pt}{12pt}\colorbox{BenchBest}{\strut 9.9s}\\[-.2pt]{\fontsize{5.8}{6.4}\selectfont \textcolor{ForestGreen}{$ \blacktriangle\,3.6\times $}}\\[-.2pt]\colorbox{BenchBest}{\strut 9s}\\[-.2pt]{\fontsize{5.8}{6.4}\selectfont \textcolor{ForestGreen}{$ \blacktriangle\,3.9\times $}}\end{tabular} & \hspace*{1pt}\begin{tabular}[t]{@{}c@{}}\rule{0pt}{12pt}\colorbox{BenchBest}{\strut 13s}\\[-.2pt]{\fontsize{5.8}{6.4}\selectfont \textcolor{ForestGreen}{$ \blacktriangle\,2.9\times $}}\\[-.2pt]\colorbox{BenchBest}{\strut 11.9s}\\[-.2pt]{\fontsize{5.8}{6.4}\selectfont \textcolor{ForestGreen}{$ \blacktriangle\,3.0\times $}}\end{tabular} & \hspace*{1pt}\begin{tabular}[t]{@{}c@{}}\rule{0pt}{12pt}\colorbox{BenchBest}{\strut 9.4s}\\[-.2pt]{\fontsize{5.8}{6.4}\selectfont \textcolor{ForestGreen}{$ \blacktriangle\,9.8\times $}}\\[-.2pt]\colorbox{BenchBest}{\strut 7s}\\[-.2pt]{\fontsize{5.8}{6.4}\selectfont \textcolor{ForestGreen}{$ \blacktriangle\,12.7\times $}}\end{tabular} & \hspace*{1pt}\begin{tabular}[t]{@{}c@{}}\rule{0pt}{12pt}\colorbox{BenchBest}{\strut 35s}\\[-.2pt]{\fontsize{5.8}{6.4}\selectfont \textcolor{ForestGreen}{$ \blacktriangle\,11.2\times $}}\\[-.2pt]\colorbox{BenchBest}{\strut 33.8s}\\[-.2pt]{\fontsize{5.8}{6.4}\selectfont \textcolor{ForestGreen}{$ \blacktriangle\,11.4\times $}}\end{tabular} & \hspace*{1pt}\begin{tabular}[t]{@{}c@{}}\rule{0pt}{12pt}\colorbox{BenchBest}{\strut 23.1s}\\[-.2pt]{\fontsize{5.8}{6.4}\selectfont \textcolor{ForestGreen}{$ \blacktriangle\,27\% $}}\\[-.2pt]\colorbox{BenchBest}{\strut 18.6s}\\[-.2pt]{\fontsize{5.8}{6.4}\selectfont \textcolor{ForestGreen}{$ \blacktriangle\,1.6\times $}}\end{tabular} & \hspace*{1pt}\begin{tabular}[t]{@{}c@{}}\rule{0pt}{12pt}\colorbox{BenchBest}{\strut 25.2s}\\[-.2pt]{\fontsize{5.8}{6.4}\selectfont \textcolor{ForestGreen}{$ \blacktriangle\,26\% $}}\\[-.2pt]\colorbox{BenchBest}{\strut 21.2s}\\[-.2pt]{\fontsize{5.8}{6.4}\selectfont \textcolor{ForestGreen}{$ \blacktriangle\,1.5\times $}}\end{tabular} & \hspace*{1pt}\begin{tabular}[t]{@{}c@{}}\rule{0pt}{12pt}\colorbox{BenchBest}{\strut 54.4s}\\[-.2pt]{\fontsize{5.8}{6.4}\selectfont \textcolor{ForestGreen}{$ \blacktriangle\,22\% $}}\\[-.2pt]\colorbox{BenchBest}{\strut 46.7s}\\[-.2pt]{\fontsize{5.8}{6.4}\selectfont \textcolor{ForestGreen}{$ \blacktriangle\,29\% $}}\end{tabular} & \hspace*{1pt}\begin{tabular}[t]{@{}c@{}}\rule{0pt}{12pt}\colorbox{BenchBest}{\strut 49.8s}\\[-.2pt]{\fontsize{5.8}{6.4}\selectfont \textcolor{ForestGreen}{$ \blacktriangle\,3.6\times $}}\\[-.2pt]\colorbox{BenchBest}{\strut 46.4s}\\[-.2pt]{\fontsize{5.8}{6.4}\selectfont \textcolor{ForestGreen}{$ \blacktriangle\,3.7\times $}}\end{tabular} & \hspace*{1pt}\begin{tabular}[t]{@{}c@{}}\rule{0pt}{12pt}3m\,44.5s\\[-.2pt]{\fontsize{5.8}{6.4}\selectfont \textcolor{BrickRed}{$ \blacktriangledown\,28\% $}}\\[-.2pt]3m\,35.9s\\[-.2pt]{\fontsize{5.8}{6.4}\selectfont \textcolor{BrickRed}{$ \blacktriangledown\,24\% $}}\end{tabular} & \hspace*{1pt}\begin{tabular}[t]{@{}c@{}}\rule{0pt}{12pt}3m\,57.4s\\[-.2pt]{\fontsize{5.8}{6.4}\selectfont \textcolor{BrickRed}{$ \blacktriangledown\,1.4\times $}}\\[-.2pt]3m\,45.7s\\[-.2pt]{\fontsize{5.8}{6.4}\selectfont \textcolor{BrickRed}{$ \blacktriangledown\,1.3\times $}}\end{tabular} & \hspace*{1pt}\begin{tabular}[t]{@{}c@{}}\rule{0pt}{12pt}6m\,50.6s\\[-.2pt]{\fontsize{5.8}{6.4}\selectfont \textcolor{BrickRed}{$ \blacktriangledown\,22\% $}}\\[-.2pt]6m\,30.7s\\[-.2pt]{\fontsize{5.8}{6.4}\selectfont \textcolor{BrickRed}{$ \blacktriangledown\,18\% $}}\end{tabular} & \textcolor{gray}{\ensuremath{\cdot}} & \hspace*{1pt}\begin{tabular}[t]{@{}c@{}}\rule{0pt}{12pt}9m\,45.5s\\[-.2pt]{\fontsize{5.8}{6.4}\selectfont \textcolor{BrickRed}{$ \blacktriangledown\,5.5\times $}}\\[-.2pt]3m\,41.5s\\[-.2pt]{\fontsize{5.8}{6.4}\selectfont \textcolor{BrickRed}{$ \blacktriangledown\,2.1\times $}}\end{tabular} & \hspace*{1pt}\begin{tabular}[t]{@{}c@{}}\rule{0pt}{12pt}8m\,25.6s\\[-.2pt]{\fontsize{5.8}{6.4}\selectfont \textcolor{BrickRed}{$ \blacktriangledown\,4.7\times $}}\\[-.2pt]3m\,56.7s\\[-.2pt]{\fontsize{5.8}{6.4}\selectfont \textcolor{BrickRed}{$ \blacktriangledown\,2.3\times $}}\end{tabular} & \hspace*{1pt}\begin{tabular}[t]{@{}c@{}}\rule{0pt}{12pt}15m\,24.9s\\[-.2pt]{\fontsize{5.8}{6.4}\selectfont \textcolor{BrickRed}{$ \blacktriangledown\,4.2\times $}}\\[-.2pt]6m\,0.2s\\[-.2pt]{\fontsize{5.8}{6.4}\selectfont \textcolor{BrickRed}{$ \blacktriangledown\,1.7\times $}}\end{tabular} & \textcolor{gray}{\ensuremath{\cdot}}\\[3pt]\hline
\BenchmarkSourceLabel{bench:external-4-k15-w10:row:blastn}{R3}\shortstack[l]{{\color{gray}R3}\\blastn} & Native\textsuperscript{*} & \hspace*{1pt}\begin{tabular}[t]{@{}c@{}}\rule{0pt}{12pt}11m\,14.7s\\[-.2pt]{\fontsize{5.8}{6.4}\selectfont \textcolor{BrickRed}{$ \blacktriangledown\,18.6\times $}\textsuperscript{c}}\\[-.2pt]11m\,9.5s\\[-.2pt]{\fontsize{5.8}{6.4}\selectfont \textcolor{BrickRed}{$ \blacktriangledown\,19.1\times $}}\end{tabular} & \hspace*{1pt}\begin{tabular}[t]{@{}c@{}}\rule{0pt}{12pt}16m\,40.6s\\[-.2pt]{\fontsize{5.8}{6.4}\selectfont \textcolor{BrickRed}{$ \blacktriangledown\,26.3\times $}\textsuperscript{c}}\\[-.2pt]16m\,29.1s\\[-.2pt]{\fontsize{5.8}{6.4}\selectfont \textcolor{BrickRed}{$ \blacktriangledown\,27.5\times $}}\end{tabular} & \hspace*{1pt}\begin{tabular}[t]{@{}c@{}}\rule{0pt}{12pt}22m\,6.7s\\[-.2pt]{\fontsize{5.8}{6.4}\selectfont \textcolor{BrickRed}{$ \blacktriangledown\,14.4\times $}\textsuperscript{c}}\\[-.2pt]21m\,51.5s\\[-.2pt]{\fontsize{5.8}{6.4}\selectfont \textcolor{BrickRed}{$ \blacktriangledown\,14.8\times $}}\end{tabular} & \hspace*{1pt}\begin{tabular}[t]{@{}c@{}}\rule{0pt}{12pt}12m\,25.4s\\[-.2pt]{\fontsize{5.8}{6.4}\selectfont \textcolor{BrickRed}{$ \blacktriangledown\,1.9\times $}\textsuperscript{c}}\\[-.2pt]12m\,19.3s\\[-.2pt]{\fontsize{5.8}{6.4}\selectfont \textcolor{BrickRed}{$ \blacktriangledown\,1.9\times $}}\end{tabular} & \hspace*{1pt}\begin{tabular}[t]{@{}c@{}}\rule{0pt}{12pt}11m\,40.1s\\[-.2pt]{\fontsize{5.8}{6.4}\selectfont \textcolor{BrickRed}{$ \blacktriangledown\,22.2\times $}\textsuperscript{a,c}}\\[-.2pt]9m\,50.3s\\[-.2pt]{\fontsize{5.8}{6.4}\selectfont \textcolor{BrickRed}{$ \blacktriangledown\,19.5\times $}}\end{tabular} & \hspace*{1pt}\begin{tabular}[t]{@{}c@{}}\rule{0pt}{12pt}16m\,52.8s\\[-.2pt]{\fontsize{5.8}{6.4}\selectfont \textcolor{BrickRed}{$ \blacktriangledown\,29.7\times $}\textsuperscript{a,c}}\\[-.2pt]13m\,48.2s\\[-.2pt]{\fontsize{5.8}{6.4}\selectfont \textcolor{BrickRed}{$ \blacktriangledown\,25.9\times $}}\end{tabular} & \hspace*{1pt}\begin{tabular}[t]{@{}c@{}}\rule{0pt}{12pt}23m\,30.3s\\[-.2pt]{\fontsize{5.8}{6.4}\selectfont \textcolor{BrickRed}{$ \blacktriangledown\,20.3\times $}\textsuperscript{a,c}}\\[-.2pt]18m\,49.3s\\[-.2pt]{\fontsize{5.8}{6.4}\selectfont \textcolor{BrickRed}{$ \blacktriangledown\,17.1\times $}}\end{tabular} & \hspace*{1pt}\begin{tabular}[t]{@{}c@{}}\rule{0pt}{12pt}12m\,50.3s\\[-.2pt]{\fontsize{5.8}{6.4}\selectfont \textcolor{BrickRed}{$ \blacktriangledown\,4.3\times $}\textsuperscript{a,c}}\\[-.2pt]10m\,46.6s\\[-.2pt]{\fontsize{5.8}{6.4}\selectfont \textcolor{BrickRed}{$ \blacktriangledown\,3.8\times $}}\end{tabular} & \hspace*{1pt}\begin{tabular}[t]{@{}c@{}}\rule{0pt}{12pt}11m\,28s\\[-.2pt]{\fontsize{5.8}{6.4}\selectfont \textcolor{BrickRed}{$ \blacktriangledown\,3.9\times $}\textsuperscript{c}}\\[-.2pt]11m\,22.8s\\[-.2pt]{\fontsize{5.8}{6.4}\selectfont \textcolor{BrickRed}{$ \blacktriangledown\,3.9\times $}}\end{tabular} & \hspace*{1pt}\begin{tabular}[t]{@{}c@{}}\rule{0pt}{12pt}16m\,55s\\[-.2pt]{\fontsize{5.8}{6.4}\selectfont \textcolor{BrickRed}{$ \blacktriangledown\,6.0\times $}\textsuperscript{c}}\\[-.2pt]16m\,43.4s\\[-.2pt]{\fontsize{5.8}{6.4}\selectfont \textcolor{BrickRed}{$ \blacktriangledown\,6.0\times $}}\end{tabular} & \hspace*{1pt}\begin{tabular}[t]{@{}c@{}}\rule{0pt}{12pt}22m\,36.2s\\[-.2pt]{\fontsize{5.8}{6.4}\selectfont \textcolor{BrickRed}{$ \blacktriangledown\,4.0\times $}\textsuperscript{c}}\\[-.2pt]22m\,20.8s\\[-.2pt]{\fontsize{5.8}{6.4}\selectfont \textcolor{BrickRed}{$ \blacktriangledown\,4.0\times $}}\end{tabular} & \hspace*{1pt}\begin{tabular}[t]{@{}c@{}}\rule{0pt}{12pt}13m\,1.6s\\[-.2pt]{\fontsize{5.8}{6.4}\selectfont \textcolor{ForestGreen}{$ \blacktriangle\,3.4\times $}\textsuperscript{c}}\\[-.2pt]12m\,55.8s\\[-.2pt]{\fontsize{5.8}{6.4}\selectfont \textcolor{ForestGreen}{$ \blacktriangle\,3.4\times $}}\end{tabular} & \hspace*{1pt}\begin{tabular}[t]{@{}c@{}}\rule{0pt}{12pt}12m\,18.6s\\[-.2pt]{\fontsize{5.8}{6.4}\selectfont \textcolor{BrickRed}{$ \blacktriangledown\,6.9\times $}\textsuperscript{a,c}}\\[-.2pt]10m\,29.2s\\[-.2pt]{\fontsize{5.8}{6.4}\selectfont \textcolor{BrickRed}{$ \blacktriangledown\,5.9\times $}}\end{tabular} & \hspace*{1pt}\begin{tabular}[t]{@{}c@{}}\rule{0pt}{12pt}17m\,25.7s\\[-.2pt]{\fontsize{5.8}{6.4}\selectfont \textcolor{BrickRed}{$ \blacktriangledown\,9.8\times $}\textsuperscript{a,c}}\\[-.2pt]14m\,22.8s\\[-.2pt]{\fontsize{5.8}{6.4}\selectfont \textcolor{BrickRed}{$ \blacktriangledown\,8.2\times $}}\end{tabular} & \hspace*{1pt}\begin{tabular}[t]{@{}c@{}}\rule{0pt}{12pt}24m\,29.4s\\[-.2pt]{\fontsize{5.8}{6.4}\selectfont \textcolor{BrickRed}{$ \blacktriangledown\,6.7\times $}\textsuperscript{a,c}}\\[-.2pt]19m\,49.6s\\[-.2pt]{\fontsize{5.8}{6.4}\selectfont \textcolor{BrickRed}{$ \blacktriangledown\,5.5\times $}}\end{tabular} & \hspace*{1pt}\begin{tabular}[t]{@{}c@{}}\rule{0pt}{12pt}13m\,43.9s\\[-.2pt]{\fontsize{5.8}{6.4}\selectfont \textcolor{ForestGreen}{$ \blacktriangle\,2.2\times $}\textsuperscript{a,c}}\\[-.2pt]11m\,39.5s\\[-.2pt]{\fontsize{5.8}{6.4}\selectfont \textcolor{ForestGreen}{$ \blacktriangle\,2.5\times $}}\end{tabular}\\[3pt]\hline
\BenchmarkSourceLabel{bench:external-4-k15-w10:row:winnowmap}{R4}\shortstack[l]{{\color{gray}R4}\\winnowmap} & Native\textsuperscript{*} & \hspace*{1pt}\begin{tabular}[t]{@{}c@{}}\rule{0pt}{12pt}52.9s\\[-.2pt]{\fontsize{5.8}{6.4}\selectfont \textcolor{BrickRed}{$ \blacktriangledown\,1.5\times $}\textsuperscript{c,e}}\\[-.2pt]50.6s\\[-.2pt]{\fontsize{5.8}{6.4}\selectfont \textcolor{BrickRed}{$ \blacktriangledown\,1.4\times $}}\end{tabular} & \hspace*{1pt}\begin{tabular}[t]{@{}c@{}}\rule{0pt}{12pt}59.4s\\[-.2pt]{\fontsize{5.8}{6.4}\selectfont \textcolor{BrickRed}{$ \blacktriangledown\,1.6\times $}\textsuperscript{c,e}}\\[-.2pt]55.2s\\[-.2pt]{\fontsize{5.8}{6.4}\selectfont \textcolor{BrickRed}{$ \blacktriangledown\,1.5\times $}}\end{tabular} & \hspace*{1pt}\begin{tabular}[t]{@{}c@{}}\rule{0pt}{12pt}1m\,49.1s\\[-.2pt]{\fontsize{5.8}{6.4}\selectfont \textcolor{BrickRed}{$ \blacktriangledown\,18\% $}\textsuperscript{c,e}}\\[-.2pt]1m\,41.3s\\[-.2pt]{\fontsize{5.8}{6.4}\selectfont \textcolor{BrickRed}{$ \blacktriangledown\,14\% $}}\end{tabular} & \hspace*{1pt}\begin{tabular}[t]{@{}c@{}}\rule{0pt}{12pt}1m\,16.1s\\[-.2pt]{\fontsize{5.8}{6.4}\selectfont \textcolor{ForestGreen}{$ \blacktriangle\,5.1\times $}\textsuperscript{c,d}}\\[-.2pt]1m\,8.2s\\[-.2pt]{\fontsize{5.8}{6.4}\selectfont \textcolor{ForestGreen}{$ \blacktriangle\,5.7\times $}}\end{tabular} & \hspace*{1pt}\begin{tabular}[t]{@{}c@{}}\rule{0pt}{12pt}31.6s\\[-.2pt]{\fontsize{5.8}{6.4}\selectfont \textcolor{gray}{$ \blacktriangleright\,0\% $}\textsuperscript{c,e}}\\[-.2pt]29.4s\\[-.2pt]{\fontsize{5.8}{6.4}\selectfont \textcolor{gray}{$ \blacktriangleright\,3\% $}}\end{tabular} & \hspace*{1pt}\begin{tabular}[t]{@{}c@{}}\rule{0pt}{12pt}44.2s\\[-.2pt]{\fontsize{5.8}{6.4}\selectfont \textcolor{BrickRed}{$ \blacktriangledown\,30\% $}\textsuperscript{c,e}}\\[-.2pt]40.3s\\[-.2pt]{\fontsize{5.8}{6.4}\selectfont \textcolor{BrickRed}{$ \blacktriangledown\,26\% $}}\end{tabular} & \hspace*{1pt}\begin{tabular}[t]{@{}c@{}}\rule{0pt}{12pt}59s\\[-.2pt]{\fontsize{5.8}{6.4}\selectfont \textcolor{ForestGreen}{$ \blacktriangle\,15\% $}\textsuperscript{c,e}}\\[-.2pt]51.8s\\[-.2pt]{\fontsize{5.8}{6.4}\selectfont \textcolor{ForestGreen}{$ \blacktriangle\,22\% $}}\end{tabular} & \hspace*{1pt}\begin{tabular}[t]{@{}c@{}}\rule{0pt}{12pt}1m\,8s\\[-.2pt]{\fontsize{5.8}{6.4}\selectfont \textcolor{ForestGreen}{$ \blacktriangle\,2.7\times $}\textsuperscript{c,d}}\\[-.2pt]1m\,0.3s\\[-.2pt]{\fontsize{5.8}{6.4}\selectfont \textcolor{ForestGreen}{$ \blacktriangle\,2.8\times $}}\end{tabular} & \hspace*{1pt}\begin{tabular}[t]{@{}c@{}}\rule{0pt}{12pt}8m\,6s\\[-.2pt]{\fontsize{5.8}{6.4}\selectfont \textcolor{BrickRed}{$ \blacktriangledown\,2.8\times $}\textsuperscript{c,e}}\\[-.2pt]8m\,3.7s\\[-.2pt]{\fontsize{5.8}{6.4}\selectfont \textcolor{BrickRed}{$ \blacktriangledown\,2.8\times $}}\end{tabular} & \hspace*{1pt}\begin{tabular}[t]{@{}c@{}}\rule{0pt}{12pt}7m\,59.6s\\[-.2pt]{\fontsize{5.8}{6.4}\selectfont \textcolor{BrickRed}{$ \blacktriangledown\,2.8\times $}\textsuperscript{c,e}}\\[-.2pt]7m\,55.3s\\[-.2pt]{\fontsize{5.8}{6.4}\selectfont \textcolor{BrickRed}{$ \blacktriangledown\,2.8\times $}}\end{tabular} & \hspace*{1pt}\begin{tabular}[t]{@{}c@{}}\rule{0pt}{12pt}15m\,50.3s\\[-.2pt]{\fontsize{5.8}{6.4}\selectfont \textcolor{BrickRed}{$ \blacktriangledown\,2.8\times $}\textsuperscript{c,e}}\\[-.2pt]15m\,42.3s\\[-.2pt]{\fontsize{5.8}{6.4}\selectfont \textcolor{BrickRed}{$ \blacktriangledown\,2.8\times $}}\end{tabular} & \hspace*{1pt}\begin{tabular}[t]{@{}c@{}}\rule{0pt}{12pt}\colorbox{BenchBest}{\strut 7m\,57s}\\[-.2pt]{\fontsize{5.8}{6.4}\selectfont \textcolor{ForestGreen}{$ \blacktriangle\,5.6\times $}\textsuperscript{c,d}}\\[-.2pt]\colorbox{BenchBest}{\strut 7m\,48.9s}\\[-.2pt]{\fontsize{5.8}{6.4}\selectfont \textcolor{ForestGreen}{$ \blacktriangle\,5.7\times $}}\end{tabular} & \hspace*{1pt}\begin{tabular}[t]{@{}c@{}}\rule{0pt}{12pt}6m\,51.2s\\[-.2pt]{\fontsize{5.8}{6.4}\selectfont \textcolor{BrickRed}{$ \blacktriangledown\,3.8\times $}\textsuperscript{c,e}}\\[-.2pt]6m\,49.1s\\[-.2pt]{\fontsize{5.8}{6.4}\selectfont \textcolor{BrickRed}{$ \blacktriangledown\,3.9\times $}}\end{tabular} & \hspace*{1pt}\begin{tabular}[t]{@{}c@{}}\rule{0pt}{12pt}6m\,41.1s\\[-.2pt]{\fontsize{5.8}{6.4}\selectfont \textcolor{BrickRed}{$ \blacktriangledown\,3.7\times $}\textsuperscript{c,e}}\\[-.2pt]6m\,37.2s\\[-.2pt]{\fontsize{5.8}{6.4}\selectfont \textcolor{BrickRed}{$ \blacktriangledown\,3.8\times $}}\end{tabular} & \hspace*{1pt}\begin{tabular}[t]{@{}c@{}}\rule{0pt}{12pt}12m\,42.6s\\[-.2pt]{\fontsize{5.8}{6.4}\selectfont \textcolor{BrickRed}{$ \blacktriangledown\,3.5\times $}\textsuperscript{c,e}}\\[-.2pt]12m\,35.1s\\[-.2pt]{\fontsize{5.8}{6.4}\selectfont \textcolor{BrickRed}{$ \blacktriangledown\,3.5\times $}}\end{tabular} & \hspace*{1pt}\begin{tabular}[t]{@{}c@{}}\rule{0pt}{12pt}\colorbox{BenchBest}{\strut 7m\,10.6s}\\[-.2pt]{\fontsize{5.8}{6.4}\selectfont \textcolor{ForestGreen}{$ \blacktriangle\,4.1\times $}\textsuperscript{c,d}}\\[-.2pt]\colorbox{BenchBest}{\strut 7m\,3.1s}\\[-.2pt]{\fontsize{5.8}{6.4}\selectfont \textcolor{ForestGreen}{$ \blacktriangle\,4.2\times $}}\end{tabular}\\[3pt]\hline
\BenchmarkSourceLabel{bench:external-4-k15-w10:row:lexicmap}{R5}\shortstack[l]{{\color{gray}R5}\\lexicmap} & Native\textsuperscript{*} & \hspace*{1pt}\begin{tabular}[t]{@{}c@{}}\rule{0pt}{12pt}4m\,7.1s\\[-.2pt]{\fontsize{5.8}{6.4}\selectfont \textcolor{BrickRed}{$ \blacktriangledown\,6.8\times $}\textsuperscript{c}}\\[-.2pt]4m\,7s\\[-.2pt]{\fontsize{5.8}{6.4}\selectfont \textcolor{BrickRed}{$ \blacktriangledown\,7.0\times $}}\end{tabular} & \hspace*{1pt}\begin{tabular}[t]{@{}c@{}}\rule{0pt}{12pt}6m\,29.9s\\[-.2pt]{\fontsize{5.8}{6.4}\selectfont \textcolor{BrickRed}{$ \blacktriangledown\,10.2\times $}\textsuperscript{c}}\\[-.2pt]6m\,29.7s\\[-.2pt]{\fontsize{5.8}{6.4}\selectfont \textcolor{BrickRed}{$ \blacktriangledown\,10.8\times $}}\end{tabular} & \hspace*{1pt}\begin{tabular}[t]{@{}c@{}}\rule{0pt}{12pt}9m\,24.6s\\[-.2pt]{\fontsize{5.8}{6.4}\selectfont \textcolor{BrickRed}{$ \blacktriangledown\,6.1\times $}\textsuperscript{c}}\\[-.2pt]9m\,24.4s\\[-.2pt]{\fontsize{5.8}{6.4}\selectfont \textcolor{BrickRed}{$ \blacktriangledown\,6.4\times $}}\end{tabular} & \hspace*{1pt}\begin{tabular}[t]{@{}c@{}}\rule{0pt}{12pt}7m\,16.7s\\[-.2pt]{\fontsize{5.8}{6.4}\selectfont \textcolor{BrickRed}{$ \blacktriangledown\,12\% $}\textsuperscript{c}}\\[-.2pt]7m\,16.4s\\[-.2pt]{\fontsize{5.8}{6.4}\selectfont \textcolor{BrickRed}{$ \blacktriangledown\,13\% $}}\end{tabular} & \hspace*{1pt}\begin{tabular}[t]{@{}c@{}}\rule{0pt}{12pt}4m\,7.8s\\[-.2pt]{\fontsize{5.8}{6.4}\selectfont \textcolor{BrickRed}{$ \blacktriangledown\,7.9\times $}\textsuperscript{b,c}}\\[-.2pt]4m\,7.5s\\[-.2pt]{\fontsize{5.8}{6.4}\selectfont \textcolor{BrickRed}{$ \blacktriangledown\,8.2\times $}}\end{tabular} & \hspace*{1pt}\begin{tabular}[t]{@{}c@{}}\rule{0pt}{12pt}6m\,29.2s\\[-.2pt]{\fontsize{5.8}{6.4}\selectfont \textcolor{BrickRed}{$ \blacktriangledown\,11.4\times $}\textsuperscript{b,c}}\\[-.2pt]6m\,28.8s\\[-.2pt]{\fontsize{5.8}{6.4}\selectfont \textcolor{BrickRed}{$ \blacktriangledown\,12.1\times $}}\end{tabular} & \hspace*{1pt}\begin{tabular}[t]{@{}c@{}}\rule{0pt}{12pt}11m\,26.4s\\[-.2pt]{\fontsize{5.8}{6.4}\selectfont \textcolor{BrickRed}{$ \blacktriangledown\,9.9\times $}\textsuperscript{b,c}}\\[-.2pt]11m\,25.5s\\[-.2pt]{\fontsize{5.8}{6.4}\selectfont \textcolor{BrickRed}{$ \blacktriangledown\,10.4\times $}}\end{tabular} & \hspace*{1pt}\begin{tabular}[t]{@{}c@{}}\rule{0pt}{12pt}7m\,15.7s\\[-.2pt]{\fontsize{5.8}{6.4}\selectfont \textcolor{BrickRed}{$ \blacktriangledown\,2.4\times $}\textsuperscript{b,c}}\\[-.2pt]7m\,13s\\[-.2pt]{\fontsize{5.8}{6.4}\selectfont \textcolor{BrickRed}{$ \blacktriangledown\,2.5\times $}}\end{tabular} & \hspace*{1pt}\begin{tabular}[t]{@{}c@{}}\rule{0pt}{12pt}7m\,34.2s\\[-.2pt]{\fontsize{5.8}{6.4}\selectfont \textcolor{BrickRed}{$ \blacktriangledown\,2.6\times $}\textsuperscript{c}}\\[-.2pt]7m\,34.1s\\[-.2pt]{\fontsize{5.8}{6.4}\selectfont \textcolor{BrickRed}{$ \blacktriangledown\,2.6\times $}}\end{tabular} & \hspace*{1pt}\begin{tabular}[t]{@{}c@{}}\rule{0pt}{12pt}11m\,38.9s\\[-.2pt]{\fontsize{5.8}{6.4}\selectfont \textcolor{BrickRed}{$ \blacktriangledown\,4.1\times $}\textsuperscript{c}}\\[-.2pt]11m\,38.7s\\[-.2pt]{\fontsize{5.8}{6.4}\selectfont \textcolor{BrickRed}{$ \blacktriangledown\,4.2\times $}}\end{tabular} & \hspace*{1pt}\begin{tabular}[t]{@{}c@{}}\rule{0pt}{12pt}17m\,11.5s\\[-.2pt]{\fontsize{5.8}{6.4}\selectfont \textcolor{BrickRed}{$ \blacktriangledown\,3.1\times $}\textsuperscript{c}}\\[-.2pt]17m\,11.2s\\[-.2pt]{\fontsize{5.8}{6.4}\selectfont \textcolor{BrickRed}{$ \blacktriangledown\,3.1\times $}}\end{tabular} & \hspace*{1pt}\begin{tabular}[t]{@{}c@{}}\rule{0pt}{12pt}15m\,28s\\[-.2pt]{\fontsize{5.8}{6.4}\selectfont \textcolor{ForestGreen}{$ \blacktriangle\,2.9\times $}\textsuperscript{c}}\\[-.2pt]15m\,27.2s\\[-.2pt]{\fontsize{5.8}{6.4}\selectfont \textcolor{ForestGreen}{$ \blacktriangle\,2.9\times $}}\end{tabular} & \hspace*{1pt}\begin{tabular}[t]{@{}c@{}}\rule{0pt}{12pt}7m\,48.7s\\[-.2pt]{\fontsize{5.8}{6.4}\selectfont \textcolor{BrickRed}{$ \blacktriangledown\,4.4\times $}\textsuperscript{b,c}}\\[-.2pt]7m\,48.4s\\[-.2pt]{\fontsize{5.8}{6.4}\selectfont \textcolor{BrickRed}{$ \blacktriangledown\,4.4\times $}}\end{tabular} & \hspace*{1pt}\begin{tabular}[t]{@{}c@{}}\rule{0pt}{12pt}11m\,34.4s\\[-.2pt]{\fontsize{5.8}{6.4}\selectfont \textcolor{BrickRed}{$ \blacktriangledown\,6.5\times $}\textsuperscript{b,c}}\\[-.2pt]11m\,34s\\[-.2pt]{\fontsize{5.8}{6.4}\selectfont \textcolor{BrickRed}{$ \blacktriangledown\,6.6\times $}}\end{tabular} & \hspace*{1pt}\begin{tabular}[t]{@{}c@{}}\rule{0pt}{12pt}17m\,36.1s\\[-.2pt]{\fontsize{5.8}{6.4}\selectfont \textcolor{BrickRed}{$ \blacktriangledown\,4.8\times $}\textsuperscript{b,c}}\\[-.2pt]17m\,35.3s\\[-.2pt]{\fontsize{5.8}{6.4}\selectfont \textcolor{BrickRed}{$ \blacktriangledown\,4.9\times $}}\end{tabular} & \hspace*{1pt}\begin{tabular}[t]{@{}c@{}}\rule{0pt}{12pt}16m\,6s\\[-.2pt]{\fontsize{5.8}{6.4}\selectfont \textcolor{ForestGreen}{$ \blacktriangle\,1.8\times $}\textsuperscript{b,c}}\\[-.2pt]16m\,3.7s\\[-.2pt]{\fontsize{5.8}{6.4}\selectfont \textcolor{ForestGreen}{$ \blacktriangle\,1.8\times $}}\end{tabular}\\[3pt]\hline
\end{tabular}\par}
\smallskip{\footnotesize\noindent Each cell: A, comparison, B, comparison. $\cdot$: no measurement. n/a: no comparison.\par
\noindent Blue: one lowest A and one lowest B measurement for each column for each configuration.\par
\noindent Gray: minimap2. Yellow: my preferred MPF / Bounded / SIMD-DP / Block configuration.\par
\noindent Comparisons use minimap2. The minimap2 row compares with BioseqDB.\par
\noindent\textsuperscript{*} Each algorithm except BLAST required intensive splitting, sometimes not equal. The splitting was the biggest possible to allow algorithm not to crash.\par
\noindent\textsuperscript{a}BLAST scores N as -3, compared with -1 in minimap2.\par
\noindent\textsuperscript{b}LexicMap stores reference N as A.\par
\noindent\textsuperscript{c}The 4 Gbp reference is split into two input records.\par
\noindent\textsuperscript{d}Winnowmap uses overlapping 250 Mb records and 499M index batches. Merging results can affect MAPQ.\par
\noindent\textsuperscript{e}Winnowmap merges results from 1999M index batches.\par

}
\clearpage
\refstepcounter{table}\label{tab:bench:external-4-k20-w15}
\noindent{\large\bfseries (G): Table \thetable: aligner comparison, 4 Gbp, k=20, w=15\hfill 2/4}\par\smallskip
{\fontsize{7.1}{8.1}\selectfont\setlength{\tabcolsep}{1pt}\setlength{\fboxsep}{2pt}\renewcommand{\arraystretch}{1.12}
\noindent\begin{tabular}{@{}p{15mm}p{26mm}|p{13.55mm}|p{13.55mm}|p{13.55mm}|p{13.55mm}|p{13.55mm}|p{13.55mm}|p{13.55mm}|p{13.55mm}|p{13.55mm}|p{13.55mm}|p{13.55mm}|p{13.55mm}|p{13.55mm}|p{13.55mm}|p{13.55mm}|p{13.55mm}|@{}}\hline
Aligner & Configuration & \multicolumn{8}{c|}{\textbf{Index}} & \multicolumn{8}{c|}{\textbf{No index}}\\
 &  & \multicolumn{8}{c|}{\textbf{4 Gbp}} & \multicolumn{8}{c|}{\textbf{4 Gbp}}\\
 &  & \multicolumn{8}{c|}{\textbf{k=20, w=15}} & \multicolumn{8}{c|}{\textbf{k=20, w=15}}\\
 &  & \multicolumn{4}{c|}{\textbf{DNA4}} & \multicolumn{4}{c|}{\textbf{DNA5}} & \multicolumn{4}{c|}{\textbf{DNA4}} & \multicolumn{4}{c|}{\textbf{DNA5}}\\
 &  & \multicolumn{1}{c|}{\textbf{M7A}} & \multicolumn{1}{c|}{\textbf{M7I}} & \multicolumn{1}{c|}{\textbf{M4}} & \multicolumn{1}{c|}{\textbf{C7G}} & \multicolumn{1}{c|}{\textbf{M7A}} & \multicolumn{1}{c|}{\textbf{M7I}} & \multicolumn{1}{c|}{\textbf{M4}} & \multicolumn{1}{c|}{\textbf{C7G}} & \multicolumn{1}{c|}{\textbf{M7A}} & \multicolumn{1}{c|}{\textbf{M7I}} & \multicolumn{1}{c|}{\textbf{M4}} & \multicolumn{1}{c|}{\textbf{C7G}} & \multicolumn{1}{c|}{\textbf{M7A}} & \multicolumn{1}{c|}{\textbf{M7I}} & \multicolumn{1}{c|}{\textbf{M4}} & \multicolumn{1}{c|}{\textbf{C7G}}\\
 & & \BenchmarkSourceLabel{bench:external-4-k20-w15:column:index-4-k20-w15-dna4-m7a}{C1}{\color{gray}C1} & \BenchmarkSourceLabel{bench:external-4-k20-w15:column:index-4-k20-w15-dna4-m7i}{C2}{\color{gray}C2} & \BenchmarkSourceLabel{bench:external-4-k20-w15:column:index-4-k20-w15-dna4-m4}{C3}{\color{gray}C3} & \BenchmarkSourceLabel{bench:external-4-k20-w15:column:index-4-k20-w15-dna4-c7g}{C4}{\color{gray}C4} & \BenchmarkSourceLabel{bench:external-4-k20-w15:column:index-4-k20-w15-dna5-m7a}{C5}{\color{gray}C5} & \BenchmarkSourceLabel{bench:external-4-k20-w15:column:index-4-k20-w15-dna5-m7i}{C6}{\color{gray}C6} & \BenchmarkSourceLabel{bench:external-4-k20-w15:column:index-4-k20-w15-dna5-m4}{C7}{\color{gray}C7} & \BenchmarkSourceLabel{bench:external-4-k20-w15:column:index-4-k20-w15-dna5-c7g}{C8}{\color{gray}C8} & \BenchmarkSourceLabel{bench:external-4-k20-w15:column:noindex-4-k20-w15-dna4-m7a}{C9}{\color{gray}C9} & \BenchmarkSourceLabel{bench:external-4-k20-w15:column:noindex-4-k20-w15-dna4-m7i}{C10}{\color{gray}C10} & \BenchmarkSourceLabel{bench:external-4-k20-w15:column:noindex-4-k20-w15-dna4-m4}{C11}{\color{gray}C11} & \BenchmarkSourceLabel{bench:external-4-k20-w15:column:noindex-4-k20-w15-dna4-c7g}{C12}{\color{gray}C12} & \BenchmarkSourceLabel{bench:external-4-k20-w15:column:noindex-4-k20-w15-dna5-m7a}{C13}{\color{gray}C13} & \BenchmarkSourceLabel{bench:external-4-k20-w15:column:noindex-4-k20-w15-dna5-m7i}{C14}{\color{gray}C14} & \BenchmarkSourceLabel{bench:external-4-k20-w15:column:noindex-4-k20-w15-dna5-m4}{C15}{\color{gray}C15} & \BenchmarkSourceLabel{bench:external-4-k20-w15:column:noindex-4-k20-w15-dna5-c7g}{C16}{\color{gray}C16}\\\hline
\rowcolor{BenchBaseline}[0.5pt][0.5pt]\BenchmarkSourceLabel{bench:external-4-k20-w15:row:minimap2}{R1}\shortstack[l]{{\color{gray}R1}\\minimap2} & Baseline & \hspace*{1pt}\begin{tabular}[t]{@{}c@{}}\rule{0pt}{12pt}2m\,48.1s\\[-.2pt]{\fontsize{5.8}{6.4}\selectfont \textcolor{BrickRed}{$ \blacktriangledown\,18.3\times $}}\\[-.2pt]2m\,41.7s\\[-.2pt]{\fontsize{5.8}{6.4}\selectfont \textcolor{BrickRed}{$ \blacktriangledown\,19.0\times $}}\end{tabular} & \hspace*{1pt}\begin{tabular}[t]{@{}c@{}}\rule{0pt}{12pt}2m\,50.3s\\[-.2pt]{\fontsize{5.8}{6.4}\selectfont \textcolor{BrickRed}{$ \blacktriangledown\,15.0\times $}}\\[-.2pt]2m\,43s\\[-.2pt]{\fontsize{5.8}{6.4}\selectfont \textcolor{BrickRed}{$ \blacktriangledown\,15.3\times $}}\end{tabular} & \hspace*{1pt}\begin{tabular}[t]{@{}c@{}}\rule{0pt}{12pt}4m\,54.3s\\[-.2pt]{\fontsize{5.8}{6.4}\selectfont \textcolor{BrickRed}{$ \blacktriangledown\,43.2\times $}}\\[-.2pt]4m\,39s\\[-.2pt]{\fontsize{5.8}{6.4}\selectfont \textcolor{BrickRed}{$ \blacktriangledown\,54.6\times $}}\end{tabular} & \hspace*{1pt}\begin{tabular}[t]{@{}c@{}}\rule{0pt}{12pt}25m\,5.5s\\[-.2pt]{\fontsize{5.8}{6.4}\selectfont \textcolor{BrickRed}{$ \blacktriangledown\,38.5\times $}\textsuperscript{g}}\\[-.2pt]25m\,0.7s\\[-.2pt]{\fontsize{5.8}{6.4}\selectfont \textcolor{BrickRed}{$ \blacktriangledown\,39.5\times $}}\end{tabular} & \hspace*{1pt}\begin{tabular}[t]{@{}c@{}}\rule{0pt}{12pt}33.9s\\[-.2pt]{\fontsize{5.8}{6.4}\selectfont \textcolor{BrickRed}{$ \blacktriangledown\,1.9\times $}}\\[-.2pt]32.9s\\[-.2pt]{\fontsize{5.8}{6.4}\selectfont \textcolor{BrickRed}{$ \blacktriangledown\,2.1\times $}}\end{tabular} & \hspace*{1pt}\begin{tabular}[t]{@{}c@{}}\rule{0pt}{12pt}34.4s\\[-.2pt]{\fontsize{5.8}{6.4}\selectfont \textcolor{BrickRed}{$ \blacktriangledown\,2.6\times $}}\\[-.2pt]32.7s\\[-.2pt]{\fontsize{5.8}{6.4}\selectfont \textcolor{BrickRed}{$ \blacktriangledown\,2.9\times $}}\end{tabular} & \hspace*{1pt}\begin{tabular}[t]{@{}c@{}}\rule{0pt}{12pt}1m\,8.2s\\[-.2pt]{\fontsize{5.8}{6.4}\selectfont \textcolor{BrickRed}{$ \blacktriangledown\,3.7\times $}}\\[-.2pt]1m\,5.1s\\[-.2pt]{\fontsize{5.8}{6.4}\selectfont \textcolor{BrickRed}{$ \blacktriangledown\,4.6\times $}}\end{tabular} & \hspace*{1pt}\begin{tabular}[t]{@{}c@{}}\rule{0pt}{12pt}3m\,9.8s\\[-.2pt]{\fontsize{5.8}{6.4}\selectfont \textcolor{BrickRed}{$ \blacktriangledown\,3.3\times $}\textsuperscript{g}}\\[-.2pt]3m\,3.7s\\[-.2pt]{\fontsize{5.8}{6.4}\selectfont \textcolor{BrickRed}{$ \blacktriangledown\,3.4\times $}}\end{tabular} & \hspace*{1pt}\begin{tabular}[t]{@{}c@{}}\rule{0pt}{12pt}3m\,28.3s\\[-.2pt]{\fontsize{5.8}{6.4}\selectfont \textcolor{BrickRed}{$ \blacktriangledown\,1.3\times $}}\\[-.2pt]3m\,22.8s\\[-.2pt]{\fontsize{5.8}{6.4}\selectfont \textcolor{BrickRed}{$ \blacktriangledown\,1.4\times $}}\end{tabular} & \hspace*{1pt}\begin{tabular}[t]{@{}c@{}}\rule{0pt}{12pt}3m\,53.2s\\[-.2pt]{\fontsize{5.8}{6.4}\selectfont \textcolor{BrickRed}{$ \blacktriangledown\,1.4\times $}}\\[-.2pt]3m\,43.5s\\[-.2pt]{\fontsize{5.8}{6.4}\selectfont \textcolor{BrickRed}{$ \blacktriangledown\,1.4\times $}}\end{tabular} & \hspace*{1pt}\begin{tabular}[t]{@{}c@{}}\rule{0pt}{12pt}7m\,14.6s\\[-.2pt]{\fontsize{5.8}{6.4}\selectfont \textcolor{BrickRed}{$ \blacktriangledown\,1.4\times $}}\\[-.2pt]6m\,57.2s\\[-.2pt]{\fontsize{5.8}{6.4}\selectfont \textcolor{BrickRed}{$ \blacktriangledown\,1.5\times $}}\end{tabular} & \hspace*{1pt}\begin{tabular}[t]{@{}c@{}}\rule{0pt}{12pt}51m\,22s\\[-.2pt]{\fontsize{5.8}{6.4}\selectfont \textcolor{gray}{n/a}\textsuperscript{g}}\\[-.2pt]51m\,17.2s\\[-.2pt]{\fontsize{5.8}{6.4}\selectfont \textcolor{gray}{n/a}}\end{tabular} & \hspace*{1pt}\begin{tabular}[t]{@{}c@{}}\rule{0pt}{12pt}\colorbox{BenchBest}{\strut 1m\,10.3s}\\[-.2pt]{\fontsize{5.8}{6.4}\selectfont \textcolor{ForestGreen}{$ \blacktriangle\,7.6\times $}}\\[-.2pt]\colorbox{BenchBest}{\strut 1m\,9s}\\[-.2pt]{\fontsize{5.8}{6.4}\selectfont \textcolor{ForestGreen}{$ \blacktriangle\,2.4\times $}}\end{tabular} & \hspace*{1pt}\begin{tabular}[t]{@{}c@{}}\rule{0pt}{12pt}\colorbox{BenchBest}{\strut 1m\,8.6s}\\[-.2pt]{\fontsize{5.8}{6.4}\selectfont \textcolor{ForestGreen}{$ \blacktriangle\,6.4\times $}}\\[-.2pt]\colorbox{BenchBest}{\strut 1m\,6.9s}\\[-.2pt]{\fontsize{5.8}{6.4}\selectfont \textcolor{ForestGreen}{$ \blacktriangle\,2.5\times $}}\end{tabular} & \hspace*{1pt}\begin{tabular}[t]{@{}c@{}}\rule{0pt}{12pt}\colorbox{BenchBest}{\strut 2m\,21.5s}\\[-.2pt]{\fontsize{5.8}{6.4}\selectfont \textcolor{ForestGreen}{$ \blacktriangle\,6.2\times $}}\\[-.2pt]\colorbox{BenchBest}{\strut 2m\,18.4s}\\[-.2pt]{\fontsize{5.8}{6.4}\selectfont \textcolor{ForestGreen}{$ \blacktriangle\,2.2\times $}}\end{tabular} & \hspace*{1pt}\begin{tabular}[t]{@{}c@{}}\rule{0pt}{12pt}15m\,56.8s\\[-.2pt]{\fontsize{5.8}{6.4}\selectfont \textcolor{gray}{n/a}\textsuperscript{g}}\\[-.2pt]15m\,51.9s\\[-.2pt]{\fontsize{5.8}{6.4}\selectfont \textcolor{gray}{n/a}}\end{tabular}\\[3pt]\hline
\rowcolor{BenchPreferred}[0.5pt][0.5pt]\BenchmarkSourceLabel{bench:external-4-k20-w15:row:bioseq}{R2}\shortstack[l]{{\color{gray}R2}\\BioseqDB} & \begin{tabular}[t]{@{}l@{}}\rule{0pt}{12pt}MPF / Bounded\\SIMD-DP / Block\end{tabular} & \hspace*{1pt}\begin{tabular}[t]{@{}c@{}}\rule{0pt}{12pt}\colorbox{BenchBest}{\strut 9.2s}\\[-.2pt]{\fontsize{5.8}{6.4}\selectfont \textcolor{ForestGreen}{$ \blacktriangle\,18.3\times $}}\\[-.2pt]\colorbox{BenchBest}{\strut 8.5s}\\[-.2pt]{\fontsize{5.8}{6.4}\selectfont \textcolor{ForestGreen}{$ \blacktriangle\,19.0\times $}}\end{tabular} & \hspace*{1pt}\begin{tabular}[t]{@{}c@{}}\rule{0pt}{12pt}\colorbox{BenchBest}{\strut 11.4s}\\[-.2pt]{\fontsize{5.8}{6.4}\selectfont \textcolor{ForestGreen}{$ \blacktriangle\,15.0\times $}}\\[-.2pt]\colorbox{BenchBest}{\strut 10.7s}\\[-.2pt]{\fontsize{5.8}{6.4}\selectfont \textcolor{ForestGreen}{$ \blacktriangle\,15.3\times $}}\end{tabular} & \hspace*{1pt}\begin{tabular}[t]{@{}c@{}}\rule{0pt}{12pt}\colorbox{BenchBest}{\strut 6.8s}\\[-.2pt]{\fontsize{5.8}{6.4}\selectfont \textcolor{ForestGreen}{$ \blacktriangle\,43.2\times $}}\\[-.2pt]\colorbox{BenchBest}{\strut 5.1s}\\[-.2pt]{\fontsize{5.8}{6.4}\selectfont \textcolor{ForestGreen}{$ \blacktriangle\,54.6\times $}}\end{tabular} & \hspace*{1pt}\begin{tabular}[t]{@{}c@{}}\rule{0pt}{12pt}\colorbox{BenchBest}{\strut 39.1s}\\[-.2pt]{\fontsize{5.8}{6.4}\selectfont \textcolor{ForestGreen}{$ \blacktriangle\,38.5\times $}}\\[-.2pt]\colorbox{BenchBest}{\strut 38s}\\[-.2pt]{\fontsize{5.8}{6.4}\selectfont \textcolor{ForestGreen}{$ \blacktriangle\,39.5\times $}}\end{tabular} & \hspace*{1pt}\begin{tabular}[t]{@{}c@{}}\rule{0pt}{12pt}\colorbox{BenchBest}{\strut 18.2s}\\[-.2pt]{\fontsize{5.8}{6.4}\selectfont \textcolor{ForestGreen}{$ \blacktriangle\,1.9\times $}}\\[-.2pt]\colorbox{BenchBest}{\strut 16s}\\[-.2pt]{\fontsize{5.8}{6.4}\selectfont \textcolor{ForestGreen}{$ \blacktriangle\,2.1\times $}}\end{tabular} & \hspace*{1pt}\begin{tabular}[t]{@{}c@{}}\rule{0pt}{12pt}\colorbox{BenchBest}{\strut 13.2s}\\[-.2pt]{\fontsize{5.8}{6.4}\selectfont \textcolor{ForestGreen}{$ \blacktriangle\,2.6\times $}}\\[-.2pt]\colorbox{BenchBest}{\strut 11.4s}\\[-.2pt]{\fontsize{5.8}{6.4}\selectfont \textcolor{ForestGreen}{$ \blacktriangle\,2.9\times $}}\end{tabular} & \hspace*{1pt}\begin{tabular}[t]{@{}c@{}}\rule{0pt}{12pt}\colorbox{BenchBest}{\strut 18.6s}\\[-.2pt]{\fontsize{5.8}{6.4}\selectfont \textcolor{ForestGreen}{$ \blacktriangle\,3.7\times $}}\\[-.2pt]\colorbox{BenchBest}{\strut 14.3s}\\[-.2pt]{\fontsize{5.8}{6.4}\selectfont \textcolor{ForestGreen}{$ \blacktriangle\,4.6\times $}}\end{tabular} & \hspace*{1pt}\begin{tabular}[t]{@{}c@{}}\rule{0pt}{12pt}57.2s\\[-.2pt]{\fontsize{5.8}{6.4}\selectfont \textcolor{ForestGreen}{$ \blacktriangle\,3.3\times $}}\\[-.2pt]53.9s\\[-.2pt]{\fontsize{5.8}{6.4}\selectfont \textcolor{ForestGreen}{$ \blacktriangle\,3.4\times $}}\end{tabular} & \hspace*{1pt}\begin{tabular}[t]{@{}c@{}}\rule{0pt}{12pt}\colorbox{BenchBest}{\strut 2m\,38.5s}\\[-.2pt]{\fontsize{5.8}{6.4}\selectfont \textcolor{ForestGreen}{$ \blacktriangle\,24\% $}}\\[-.2pt]\colorbox{BenchBest}{\strut 2m\,29.8s}\\[-.2pt]{\fontsize{5.8}{6.4}\selectfont \textcolor{ForestGreen}{$ \blacktriangle\,26\% $}}\end{tabular} & \hspace*{1pt}\begin{tabular}[t]{@{}c@{}}\rule{0pt}{12pt}\colorbox{BenchBest}{\strut 2m\,49s}\\[-.2pt]{\fontsize{5.8}{6.4}\selectfont \textcolor{ForestGreen}{$ \blacktriangle\,28\% $}}\\[-.2pt]\colorbox{BenchBest}{\strut 2m\,37s}\\[-.2pt]{\fontsize{5.8}{6.4}\selectfont \textcolor{ForestGreen}{$ \blacktriangle\,30\% $}}\end{tabular} & \hspace*{1pt}\begin{tabular}[t]{@{}c@{}}\rule{0pt}{12pt}\colorbox{BenchBest}{\strut 5m\,3.3s}\\[-.2pt]{\fontsize{5.8}{6.4}\selectfont \textcolor{ForestGreen}{$ \blacktriangle\,1.4\times $}}\\[-.2pt]\colorbox{BenchBest}{\strut 4m\,42.7s}\\[-.2pt]{\fontsize{5.8}{6.4}\selectfont \textcolor{ForestGreen}{$ \blacktriangle\,1.5\times $}}\end{tabular} & \textcolor{gray}{\ensuremath{\cdot}} & \hspace*{1pt}\begin{tabular}[t]{@{}c@{}}\rule{0pt}{12pt}8m\,52.3s\\[-.2pt]{\fontsize{5.8}{6.4}\selectfont \textcolor{BrickRed}{$ \blacktriangledown\,7.6\times $}}\\[-.2pt]2m\,45.5s\\[-.2pt]{\fontsize{5.8}{6.4}\selectfont \textcolor{BrickRed}{$ \blacktriangledown\,2.4\times $}}\end{tabular} & \hspace*{1pt}\begin{tabular}[t]{@{}c@{}}\rule{0pt}{12pt}7m\,22.2s\\[-.2pt]{\fontsize{5.8}{6.4}\selectfont \textcolor{BrickRed}{$ \blacktriangledown\,6.4\times $}}\\[-.2pt]2m\,48.2s\\[-.2pt]{\fontsize{5.8}{6.4}\selectfont \textcolor{BrickRed}{$ \blacktriangledown\,2.5\times $}}\end{tabular} & \hspace*{1pt}\begin{tabular}[t]{@{}c@{}}\rule{0pt}{12pt}14m\,43.3s\\[-.2pt]{\fontsize{5.8}{6.4}\selectfont \textcolor{BrickRed}{$ \blacktriangledown\,6.2\times $}}\\[-.2pt]5m\,9.5s\\[-.2pt]{\fontsize{5.8}{6.4}\selectfont \textcolor{BrickRed}{$ \blacktriangledown\,2.2\times $}}\end{tabular} & \textcolor{gray}{\ensuremath{\cdot}}\\[3pt]\hline
\BenchmarkSourceLabel{bench:external-4-k20-w15:row:blastn}{R3}\shortstack[l]{{\color{gray}R3}\\blastn} & Native\textsuperscript{*} & \hspace*{1pt}\begin{tabular}[t]{@{}c@{}}\rule{0pt}{12pt}4m\,35.4s\\[-.2pt]{\fontsize{5.8}{6.4}\selectfont \textcolor{BrickRed}{$ \blacktriangledown\,1.6\times $}\textsuperscript{c}}\\[-.2pt]4m\,30.3s\\[-.2pt]{\fontsize{5.8}{6.4}\selectfont \textcolor{BrickRed}{$ \blacktriangledown\,1.7\times $}}\end{tabular} & \hspace*{1pt}\begin{tabular}[t]{@{}c@{}}\rule{0pt}{12pt}6m\,21.2s\\[-.2pt]{\fontsize{5.8}{6.4}\selectfont \textcolor{BrickRed}{$ \blacktriangledown\,2.2\times $}\textsuperscript{c}}\\[-.2pt]6m\,10.8s\\[-.2pt]{\fontsize{5.8}{6.4}\selectfont \textcolor{BrickRed}{$ \blacktriangledown\,2.3\times $}}\end{tabular} & \hspace*{1pt}\begin{tabular}[t]{@{}c@{}}\rule{0pt}{12pt}8m\,55.3s\\[-.2pt]{\fontsize{5.8}{6.4}\selectfont \textcolor{BrickRed}{$ \blacktriangledown\,1.8\times $}\textsuperscript{c}}\\[-.2pt]8m\,41.3s\\[-.2pt]{\fontsize{5.8}{6.4}\selectfont \textcolor{BrickRed}{$ \blacktriangledown\,1.9\times $}}\end{tabular} & \hspace*{1pt}\begin{tabular}[t]{@{}c@{}}\rule{0pt}{12pt}5m\,9.4s\\[-.2pt]{\fontsize{5.8}{6.4}\selectfont \textcolor{ForestGreen}{$ \blacktriangle\,4.9\times $}\textsuperscript{c}}\\[-.2pt]5m\,4.2s\\[-.2pt]{\fontsize{5.8}{6.4}\selectfont \textcolor{ForestGreen}{$ \blacktriangle\,4.9\times $}}\end{tabular} & \hspace*{1pt}\begin{tabular}[t]{@{}c@{}}\rule{0pt}{12pt}5m\,39.4s\\[-.2pt]{\fontsize{5.8}{6.4}\selectfont \textcolor{BrickRed}{$ \blacktriangledown\,10.0\times $}\textsuperscript{a,c}}\\[-.2pt]4m\,13.6s\\[-.2pt]{\fontsize{5.8}{6.4}\selectfont \textcolor{BrickRed}{$ \blacktriangledown\,7.7\times $}}\end{tabular} & \hspace*{1pt}\begin{tabular}[t]{@{}c@{}}\rule{0pt}{12pt}7m\,32.4s\\[-.2pt]{\fontsize{5.8}{6.4}\selectfont \textcolor{BrickRed}{$ \blacktriangledown\,13.1\times $}\textsuperscript{a,c}}\\[-.2pt]5m\,27.1s\\[-.2pt]{\fontsize{5.8}{6.4}\selectfont \textcolor{BrickRed}{$ \blacktriangledown\,10.0\times $}}\end{tabular} & \hspace*{1pt}\begin{tabular}[t]{@{}c@{}}\rule{0pt}{12pt}11m\,29.3s\\[-.2pt]{\fontsize{5.8}{6.4}\selectfont \textcolor{BrickRed}{$ \blacktriangledown\,10.1\times $}\textsuperscript{a,c}}\\[-.2pt]8m\,4s\\[-.2pt]{\fontsize{5.8}{6.4}\selectfont \textcolor{BrickRed}{$ \blacktriangledown\,7.4\times $}}\end{tabular} & \hspace*{1pt}\begin{tabular}[t]{@{}c@{}}\rule{0pt}{12pt}6m\,10.2s\\[-.2pt]{\fontsize{5.8}{6.4}\selectfont \textcolor{BrickRed}{$ \blacktriangledown\,2.0\times $}\textsuperscript{a,c}}\\[-.2pt]4m\,39s\\[-.2pt]{\fontsize{5.8}{6.4}\selectfont \textcolor{BrickRed}{$ \blacktriangledown\,1.5\times $}}\end{tabular} & \hspace*{1pt}\begin{tabular}[t]{@{}c@{}}\rule{0pt}{12pt}4m\,58.3s\\[-.2pt]{\fontsize{5.8}{6.4}\selectfont \textcolor{BrickRed}{$ \blacktriangledown\,1.4\times $}\textsuperscript{c}}\\[-.2pt]4m\,53.3s\\[-.2pt]{\fontsize{5.8}{6.4}\selectfont \textcolor{BrickRed}{$ \blacktriangledown\,1.4\times $}}\end{tabular} & \hspace*{1pt}\begin{tabular}[t]{@{}c@{}}\rule{0pt}{12pt}6m\,40.8s\\[-.2pt]{\fontsize{5.8}{6.4}\selectfont \textcolor{BrickRed}{$ \blacktriangledown\,1.7\times $}\textsuperscript{c}}\\[-.2pt]6m\,30.3s\\[-.2pt]{\fontsize{5.8}{6.4}\selectfont \textcolor{BrickRed}{$ \blacktriangledown\,1.7\times $}}\end{tabular} & \hspace*{1pt}\begin{tabular}[t]{@{}c@{}}\rule{0pt}{12pt}9m\,34s\\[-.2pt]{\fontsize{5.8}{6.4}\selectfont \textcolor{BrickRed}{$ \blacktriangledown\,1.3\times $}\textsuperscript{c}}\\[-.2pt]9m\,19.8s\\[-.2pt]{\fontsize{5.8}{6.4}\selectfont \textcolor{BrickRed}{$ \blacktriangledown\,1.3\times $}}\end{tabular} & \hspace*{1pt}\begin{tabular}[t]{@{}c@{}}\rule{0pt}{12pt}\colorbox{BenchBest}{\strut 5m\,40.5s}\\[-.2pt]{\fontsize{5.8}{6.4}\selectfont \textcolor{ForestGreen}{$ \blacktriangle\,9.1\times $}\textsuperscript{c}}\\[-.2pt]\colorbox{BenchBest}{\strut 5m\,35.1s}\\[-.2pt]{\fontsize{5.8}{6.4}\selectfont \textcolor{ForestGreen}{$ \blacktriangle\,9.2\times $}}\end{tabular} & \hspace*{1pt}\begin{tabular}[t]{@{}c@{}}\rule{0pt}{12pt}6m\,15.6s\\[-.2pt]{\fontsize{5.8}{6.4}\selectfont \textcolor{BrickRed}{$ \blacktriangledown\,5.3\times $}\textsuperscript{a,c}}\\[-.2pt]4m\,49.2s\\[-.2pt]{\fontsize{5.8}{6.4}\selectfont \textcolor{BrickRed}{$ \blacktriangledown\,4.2\times $}}\end{tabular} & \hspace*{1pt}\begin{tabular}[t]{@{}c@{}}\rule{0pt}{12pt}8m\,11.1s\\[-.2pt]{\fontsize{5.8}{6.4}\selectfont \textcolor{BrickRed}{$ \blacktriangledown\,7.2\times $}\textsuperscript{a,c}}\\[-.2pt]6m\,2.7s\\[-.2pt]{\fontsize{5.8}{6.4}\selectfont \textcolor{BrickRed}{$ \blacktriangledown\,5.4\times $}}\end{tabular} & \hspace*{1pt}\begin{tabular}[t]{@{}c@{}}\rule{0pt}{12pt}12m\,29.6s\\[-.2pt]{\fontsize{5.8}{6.4}\selectfont \textcolor{BrickRed}{$ \blacktriangledown\,5.3\times $}\textsuperscript{a,c}}\\[-.2pt]9m\,3.4s\\[-.2pt]{\fontsize{5.8}{6.4}\selectfont \textcolor{BrickRed}{$ \blacktriangledown\,3.9\times $}}\end{tabular} & \hspace*{1pt}\begin{tabular}[t]{@{}c@{}}\rule{0pt}{12pt}\colorbox{BenchBest}{\strut 7m\,4.7s}\\[-.2pt]{\fontsize{5.8}{6.4}\selectfont \textcolor{ForestGreen}{$ \blacktriangle\,2.3\times $}\textsuperscript{a,c}}\\[-.2pt]\colorbox{BenchBest}{\strut 5m\,34.7s}\\[-.2pt]{\fontsize{5.8}{6.4}\selectfont \textcolor{ForestGreen}{$ \blacktriangle\,2.8\times $}}\end{tabular}\\[3pt]\hline
\BenchmarkSourceLabel{bench:external-4-k20-w15:row:winnowmap}{R4}\shortstack[l]{{\color{gray}R4}\\winnowmap} & Native\textsuperscript{*} & \hspace*{1pt}\begin{tabular}[t]{@{}c@{}}\rule{0pt}{12pt}55.6s\\[-.2pt]{\fontsize{5.8}{6.4}\selectfont \textcolor{ForestGreen}{$ \blacktriangle\,3.0\times $}\textsuperscript{c,e}}\\[-.2pt]53.8s\\[-.2pt]{\fontsize{5.8}{6.4}\selectfont \textcolor{ForestGreen}{$ \blacktriangle\,3.0\times $}}\end{tabular} & \hspace*{1pt}\begin{tabular}[t]{@{}c@{}}\rule{0pt}{12pt}59.6s\\[-.2pt]{\fontsize{5.8}{6.4}\selectfont \textcolor{ForestGreen}{$ \blacktriangle\,2.9\times $}\textsuperscript{c,e}}\\[-.2pt]56.2s\\[-.2pt]{\fontsize{5.8}{6.4}\selectfont \textcolor{ForestGreen}{$ \blacktriangle\,2.9\times $}}\end{tabular} & \hspace*{1pt}\begin{tabular}[t]{@{}c@{}}\rule{0pt}{12pt}1m\,43s\\[-.2pt]{\fontsize{5.8}{6.4}\selectfont \textcolor{ForestGreen}{$ \blacktriangle\,2.9\times $}\textsuperscript{c,e}}\\[-.2pt]1m\,36.3s\\[-.2pt]{\fontsize{5.8}{6.4}\selectfont \textcolor{ForestGreen}{$ \blacktriangle\,2.9\times $}}\end{tabular} & \hspace*{1pt}\begin{tabular}[t]{@{}c@{}}\rule{0pt}{12pt}1m\,0.7s\\[-.2pt]{\fontsize{5.8}{6.4}\selectfont \textcolor{ForestGreen}{$ \blacktriangle\,24.8\times $}\textsuperscript{c,d}}\\[-.2pt]54.1s\\[-.2pt]{\fontsize{5.8}{6.4}\selectfont \textcolor{ForestGreen}{$ \blacktriangle\,27.7\times $}}\end{tabular} & \hspace*{1pt}\begin{tabular}[t]{@{}c@{}}\rule{0pt}{12pt}24.6s\\[-.2pt]{\fontsize{5.8}{6.4}\selectfont \textcolor{ForestGreen}{$ \blacktriangle\,28\% $}\textsuperscript{c,e}}\\[-.2pt]22.8s\\[-.2pt]{\fontsize{5.8}{6.4}\selectfont \textcolor{ForestGreen}{$ \blacktriangle\,1.4\times $}}\end{tabular} & \hspace*{1pt}\begin{tabular}[t]{@{}c@{}}\rule{0pt}{12pt}23.7s\\[-.2pt]{\fontsize{5.8}{6.4}\selectfont \textcolor{ForestGreen}{$ \blacktriangle\,1.4\times $}\textsuperscript{c,e}}\\[-.2pt]20.4s\\[-.2pt]{\fontsize{5.8}{6.4}\selectfont \textcolor{ForestGreen}{$ \blacktriangle\,1.6\times $}}\end{tabular} & \hspace*{1pt}\begin{tabular}[t]{@{}c@{}}\rule{0pt}{12pt}36.1s\\[-.2pt]{\fontsize{5.8}{6.4}\selectfont \textcolor{ForestGreen}{$ \blacktriangle\,1.9\times $}\textsuperscript{c,e}}\\[-.2pt]29.6s\\[-.2pt]{\fontsize{5.8}{6.4}\selectfont \textcolor{ForestGreen}{$ \blacktriangle\,2.2\times $}}\end{tabular} & \hspace*{1pt}\begin{tabular}[t]{@{}c@{}}\rule{0pt}{12pt}\colorbox{BenchBest}{\strut 41.4s}\\[-.2pt]{\fontsize{5.8}{6.4}\selectfont \textcolor{ForestGreen}{$ \blacktriangle\,4.6\times $}\textsuperscript{c,d}}\\[-.2pt]\colorbox{BenchBest}{\strut 34.7s}\\[-.2pt]{\fontsize{5.8}{6.4}\selectfont \textcolor{ForestGreen}{$ \blacktriangle\,5.3\times $}}\end{tabular} & \hspace*{1pt}\begin{tabular}[t]{@{}c@{}}\rule{0pt}{12pt}10m\,41.8s\\[-.2pt]{\fontsize{5.8}{6.4}\selectfont \textcolor{BrickRed}{$ \blacktriangledown\,3.1\times $}\textsuperscript{c,e}}\\[-.2pt]10m\,40s\\[-.2pt]{\fontsize{5.8}{6.4}\selectfont \textcolor{BrickRed}{$ \blacktriangledown\,3.2\times $}}\end{tabular} & \hspace*{1pt}\begin{tabular}[t]{@{}c@{}}\rule{0pt}{12pt}11m\,28.6s\\[-.2pt]{\fontsize{5.8}{6.4}\selectfont \textcolor{BrickRed}{$ \blacktriangledown\,3.0\times $}\textsuperscript{c,e}}\\[-.2pt]11m\,25.2s\\[-.2pt]{\fontsize{5.8}{6.4}\selectfont \textcolor{BrickRed}{$ \blacktriangledown\,3.1\times $}}\end{tabular} & \hspace*{1pt}\begin{tabular}[t]{@{}c@{}}\rule{0pt}{12pt}33m\,6.8s\\[-.2pt]{\fontsize{5.8}{6.4}\selectfont \textcolor{BrickRed}{$ \blacktriangledown\,4.6\times $}\textsuperscript{c,e}}\\[-.2pt]33m\,0.1s\\[-.2pt]{\fontsize{5.8}{6.4}\selectfont \textcolor{BrickRed}{$ \blacktriangledown\,4.7\times $}}\end{tabular} & \hspace*{1pt}\begin{tabular}[t]{@{}c@{}}\rule{0pt}{12pt}12m\,28.4s\\[-.2pt]{\fontsize{5.8}{6.4}\selectfont \textcolor{ForestGreen}{$ \blacktriangle\,4.1\times $}\textsuperscript{c,d}}\\[-.2pt]12m\,21.9s\\[-.2pt]{\fontsize{5.8}{6.4}\selectfont \textcolor{ForestGreen}{$ \blacktriangle\,4.1\times $}}\end{tabular} & \hspace*{1pt}\begin{tabular}[t]{@{}c@{}}\rule{0pt}{12pt}8m\,28.7s\\[-.2pt]{\fontsize{5.8}{6.4}\selectfont \textcolor{BrickRed}{$ \blacktriangledown\,7.2\times $}\textsuperscript{c,e}}\\[-.2pt]8m\,26.9s\\[-.2pt]{\fontsize{5.8}{6.4}\selectfont \textcolor{BrickRed}{$ \blacktriangledown\,7.3\times $}}\end{tabular} & \hspace*{1pt}\begin{tabular}[t]{@{}c@{}}\rule{0pt}{12pt}8m\,59.5s\\[-.2pt]{\fontsize{5.8}{6.4}\selectfont \textcolor{BrickRed}{$ \blacktriangledown\,7.9\times $}\textsuperscript{c,e}}\\[-.2pt]8m\,56.1s\\[-.2pt]{\fontsize{5.8}{6.4}\selectfont \textcolor{BrickRed}{$ \blacktriangledown\,8.0\times $}}\end{tabular} & \hspace*{1pt}\begin{tabular}[t]{@{}c@{}}\rule{0pt}{12pt}24m\,58.6s\\[-.2pt]{\fontsize{5.8}{6.4}\selectfont \textcolor{BrickRed}{$ \blacktriangledown\,10.6\times $}\textsuperscript{c,e}}\\[-.2pt]24m\,52s\\[-.2pt]{\fontsize{5.8}{6.4}\selectfont \textcolor{BrickRed}{$ \blacktriangledown\,10.8\times $}}\end{tabular} & \hspace*{1pt}\begin{tabular}[t]{@{}c@{}}\rule{0pt}{12pt}10m\,19s\\[-.2pt]{\fontsize{5.8}{6.4}\selectfont \textcolor{ForestGreen}{$ \blacktriangle\,1.5\times $}\textsuperscript{c,d}}\\[-.2pt]10m\,12.1s\\[-.2pt]{\fontsize{5.8}{6.4}\selectfont \textcolor{ForestGreen}{$ \blacktriangle\,1.6\times $}}\end{tabular}\\[3pt]\hline
\BenchmarkSourceLabel{bench:external-4-k20-w15:row:lexicmap}{R5}\shortstack[l]{{\color{gray}R5}\\lexicmap} & Native\textsuperscript{*} & \hspace*{1pt}\begin{tabular}[t]{@{}c@{}}\rule{0pt}{12pt}56.9s\\[-.2pt]{\fontsize{5.8}{6.4}\selectfont \textcolor{ForestGreen}{$ \blacktriangle\,3.0\times $}\textsuperscript{c}}\\[-.2pt]56.7s\\[-.2pt]{\fontsize{5.8}{6.4}\selectfont \textcolor{ForestGreen}{$ \blacktriangle\,2.9\times $}}\end{tabular} & \hspace*{1pt}\begin{tabular}[t]{@{}c@{}}\rule{0pt}{12pt}1m\,21.3s\\[-.2pt]{\fontsize{5.8}{6.4}\selectfont \textcolor{ForestGreen}{$ \blacktriangle\,2.1\times $}\textsuperscript{c}}\\[-.2pt]1m\,21.1s\\[-.2pt]{\fontsize{5.8}{6.4}\selectfont \textcolor{ForestGreen}{$ \blacktriangle\,2.0\times $}}\end{tabular} & \hspace*{1pt}\begin{tabular}[t]{@{}c@{}}\rule{0pt}{12pt}2m\,1.2s\\[-.2pt]{\fontsize{5.8}{6.4}\selectfont \textcolor{ForestGreen}{$ \blacktriangle\,2.4\times $}\textsuperscript{c}}\\[-.2pt]2m\,0.9s\\[-.2pt]{\fontsize{5.8}{6.4}\selectfont \textcolor{ForestGreen}{$ \blacktriangle\,2.3\times $}}\end{tabular} & \hspace*{1pt}\begin{tabular}[t]{@{}c@{}}\rule{0pt}{12pt}1m\,16.5s\\[-.2pt]{\fontsize{5.8}{6.4}\selectfont \textcolor{ForestGreen}{$ \blacktriangle\,19.7\times $}\textsuperscript{c}}\\[-.2pt]1m\,16.1s\\[-.2pt]{\fontsize{5.8}{6.4}\selectfont \textcolor{ForestGreen}{$ \blacktriangle\,19.7\times $}}\end{tabular} & \hspace*{1pt}\begin{tabular}[t]{@{}c@{}}\rule{0pt}{12pt}56.4s\\[-.2pt]{\fontsize{5.8}{6.4}\selectfont \textcolor{BrickRed}{$ \blacktriangledown\,1.7\times $}\textsuperscript{b,c}}\\[-.2pt]55.8s\\[-.2pt]{\fontsize{5.8}{6.4}\selectfont \textcolor{BrickRed}{$ \blacktriangledown\,1.7\times $}}\end{tabular} & \hspace*{1pt}\begin{tabular}[t]{@{}c@{}}\rule{0pt}{12pt}1m\,20.1s\\[-.2pt]{\fontsize{5.8}{6.4}\selectfont \textcolor{BrickRed}{$ \blacktriangledown\,2.3\times $}\textsuperscript{b,c}}\\[-.2pt]1m\,19.6s\\[-.2pt]{\fontsize{5.8}{6.4}\selectfont \textcolor{BrickRed}{$ \blacktriangledown\,2.4\times $}}\end{tabular} & \hspace*{1pt}\begin{tabular}[t]{@{}c@{}}\rule{0pt}{12pt}2m\,9.7s\\[-.2pt]{\fontsize{5.8}{6.4}\selectfont \textcolor{BrickRed}{$ \blacktriangledown\,1.9\times $}\textsuperscript{b,c}}\\[-.2pt]2m\,8.5s\\[-.2pt]{\fontsize{5.8}{6.4}\selectfont \textcolor{BrickRed}{$ \blacktriangledown\,2.0\times $}}\end{tabular} & \hspace*{1pt}\begin{tabular}[t]{@{}c@{}}\rule{0pt}{12pt}1m\,14.6s\\[-.2pt]{\fontsize{5.8}{6.4}\selectfont \textcolor{ForestGreen}{$ \blacktriangle\,2.5\times $}\textsuperscript{b,c}}\\[-.2pt]1m\,12.8s\\[-.2pt]{\fontsize{5.8}{6.4}\selectfont \textcolor{ForestGreen}{$ \blacktriangle\,2.5\times $}}\end{tabular} & \hspace*{1pt}\begin{tabular}[t]{@{}c@{}}\rule{0pt}{12pt}4m\,52.5s\\[-.2pt]{\fontsize{5.8}{6.4}\selectfont \textcolor{BrickRed}{$ \blacktriangledown\,1.4\times $}\textsuperscript{c}}\\[-.2pt]4m\,52.3s\\[-.2pt]{\fontsize{5.8}{6.4}\selectfont \textcolor{BrickRed}{$ \blacktriangledown\,1.4\times $}}\end{tabular} & \hspace*{1pt}\begin{tabular}[t]{@{}c@{}}\rule{0pt}{12pt}6m\,45.3s\\[-.2pt]{\fontsize{5.8}{6.4}\selectfont \textcolor{BrickRed}{$ \blacktriangledown\,1.7\times $}\textsuperscript{c}}\\[-.2pt]6m\,45.1s\\[-.2pt]{\fontsize{5.8}{6.4}\selectfont \textcolor{BrickRed}{$ \blacktriangledown\,1.8\times $}}\end{tabular} & \hspace*{1pt}\begin{tabular}[t]{@{}c@{}}\rule{0pt}{12pt}10m\,57.1s\\[-.2pt]{\fontsize{5.8}{6.4}\selectfont \textcolor{BrickRed}{$ \blacktriangledown\,1.5\times $}\textsuperscript{c}}\\[-.2pt]10m\,56.8s\\[-.2pt]{\fontsize{5.8}{6.4}\selectfont \textcolor{BrickRed}{$ \blacktriangledown\,1.6\times $}}\end{tabular} & \hspace*{1pt}\begin{tabular}[t]{@{}c@{}}\rule{0pt}{12pt}9m\,46.1s\\[-.2pt]{\fontsize{5.8}{6.4}\selectfont \textcolor{ForestGreen}{$ \blacktriangle\,5.3\times $}\textsuperscript{c}}\\[-.2pt]9m\,45.6s\\[-.2pt]{\fontsize{5.8}{6.4}\selectfont \textcolor{ForestGreen}{$ \blacktriangle\,5.3\times $}}\end{tabular} & \hspace*{1pt}\begin{tabular}[t]{@{}c@{}}\rule{0pt}{12pt}4m\,52.2s\\[-.2pt]{\fontsize{5.8}{6.4}\selectfont \textcolor{BrickRed}{$ \blacktriangledown\,4.2\times $}\textsuperscript{b,c}}\\[-.2pt]4m\,51.8s\\[-.2pt]{\fontsize{5.8}{6.4}\selectfont \textcolor{BrickRed}{$ \blacktriangledown\,4.2\times $}}\end{tabular} & \hspace*{1pt}\begin{tabular}[t]{@{}c@{}}\rule{0pt}{12pt}6m\,46.8s\\[-.2pt]{\fontsize{5.8}{6.4}\selectfont \textcolor{BrickRed}{$ \blacktriangledown\,5.9\times $}\textsuperscript{b,c}}\\[-.2pt]6m\,46.3s\\[-.2pt]{\fontsize{5.8}{6.4}\selectfont \textcolor{BrickRed}{$ \blacktriangledown\,6.1\times $}}\end{tabular} & \hspace*{1pt}\begin{tabular}[t]{@{}c@{}}\rule{0pt}{12pt}10m\,14.1s\\[-.2pt]{\fontsize{5.8}{6.4}\selectfont \textcolor{BrickRed}{$ \blacktriangledown\,4.3\times $}\textsuperscript{b,c}}\\[-.2pt]10m\,13.1s\\[-.2pt]{\fontsize{5.8}{6.4}\selectfont \textcolor{BrickRed}{$ \blacktriangledown\,4.4\times $}}\end{tabular} & \hspace*{1pt}\begin{tabular}[t]{@{}c@{}}\rule{0pt}{12pt}9m\,49.2s\\[-.2pt]{\fontsize{5.8}{6.4}\selectfont \textcolor{ForestGreen}{$ \blacktriangle\,1.6\times $}\textsuperscript{b,c}}\\[-.2pt]9m\,47.7s\\[-.2pt]{\fontsize{5.8}{6.4}\selectfont \textcolor{ForestGreen}{$ \blacktriangle\,1.6\times $}}\end{tabular}\\[3pt]\hline
\end{tabular}\par}
\smallskip{\footnotesize\noindent Each cell: A, comparison, B, comparison. $\cdot$: no measurement. n/a: no comparison.\par
\noindent Blue: one lowest A and one lowest B measurement for each column for each configuration.\par
\noindent Gray: minimap2. Yellow: my preferred MPF / Bounded / SIMD-DP / Block configuration.\par
\noindent Comparisons use minimap2. The minimap2 row compares with BioseqDB.\par
\noindent\textsuperscript{*} Each algorithm except BLAST required intensive splitting, sometimes not equal. The splitting was the biggest possible to allow algorithm not to crash.\par
\noindent\textsuperscript{a}BLAST scores N as -3, compared with -1 in minimap2.\par
\noindent\textsuperscript{b}LexicMap stores reference N as A.\par
\noindent\textsuperscript{c}The 4 Gbp reference is split into two input records.\par
\noindent\textsuperscript{d}Winnowmap uses overlapping 250 Mb records and 499M index batches. Merging results can affect MAPQ.\par
\noindent\textsuperscript{e}Winnowmap merges results from 1999M index batches.\par

}
\clearpage
\refstepcounter{table}\label{tab:bench:external-2-k15-w10}
\noindent{\large\bfseries (H): Table \thetable: aligner comparison, 2 Gbp, k=15, w=10\hfill 3/4}\par\smallskip
{\fontsize{7.1}{8.1}\selectfont\setlength{\tabcolsep}{1pt}\setlength{\fboxsep}{2pt}\renewcommand{\arraystretch}{1.12}
\noindent\begin{tabular}{@{}p{15mm}p{26mm}|p{13.55mm}|p{13.55mm}|p{13.55mm}|p{13.55mm}|p{13.55mm}|p{13.55mm}|p{13.55mm}|p{13.55mm}|p{13.55mm}|p{13.55mm}|p{13.55mm}|p{13.55mm}|p{13.55mm}|p{13.55mm}|p{13.55mm}|p{13.55mm}|@{}}\hline
Aligner & Configuration & \multicolumn{8}{c|}{\textbf{Index}} & \multicolumn{8}{c|}{\textbf{No index}}\\
 &  & \multicolumn{8}{c|}{\textbf{2 Gbp}} & \multicolumn{8}{c|}{\textbf{2 Gbp}}\\
 &  & \multicolumn{8}{c|}{\textbf{k=15, w=10}} & \multicolumn{8}{c|}{\textbf{k=15, w=10}}\\
 &  & \multicolumn{4}{c|}{\textbf{DNA4}} & \multicolumn{4}{c|}{\textbf{DNA5}} & \multicolumn{4}{c|}{\textbf{DNA4}} & \multicolumn{4}{c|}{\textbf{DNA5}}\\
 &  & \multicolumn{1}{c|}{\textbf{M7A}} & \multicolumn{1}{c|}{\textbf{M7I}} & \multicolumn{1}{c|}{\textbf{M4}} & \multicolumn{1}{c|}{\textbf{C7G}} & \multicolumn{1}{c|}{\textbf{M7A}} & \multicolumn{1}{c|}{\textbf{M7I}} & \multicolumn{1}{c|}{\textbf{M4}} & \multicolumn{1}{c|}{\textbf{C7G}} & \multicolumn{1}{c|}{\textbf{M7A}} & \multicolumn{1}{c|}{\textbf{M7I}} & \multicolumn{1}{c|}{\textbf{M4}} & \multicolumn{1}{c|}{\textbf{C7G}} & \multicolumn{1}{c|}{\textbf{M7A}} & \multicolumn{1}{c|}{\textbf{M7I}} & \multicolumn{1}{c|}{\textbf{M4}} & \multicolumn{1}{c|}{\textbf{C7G}}\\
 & & \BenchmarkSourceLabel{bench:external-2-k15-w10:column:index-2-k15-w10-dna4-m7a}{C1}{\color{gray}C1} & \BenchmarkSourceLabel{bench:external-2-k15-w10:column:index-2-k15-w10-dna4-m7i}{C2}{\color{gray}C2} & \BenchmarkSourceLabel{bench:external-2-k15-w10:column:index-2-k15-w10-dna4-m4}{C3}{\color{gray}C3} & \BenchmarkSourceLabel{bench:external-2-k15-w10:column:index-2-k15-w10-dna4-c7g}{C4}{\color{gray}C4} & \BenchmarkSourceLabel{bench:external-2-k15-w10:column:index-2-k15-w10-dna5-m7a}{C5}{\color{gray}C5} & \BenchmarkSourceLabel{bench:external-2-k15-w10:column:index-2-k15-w10-dna5-m7i}{C6}{\color{gray}C6} & \BenchmarkSourceLabel{bench:external-2-k15-w10:column:index-2-k15-w10-dna5-m4}{C7}{\color{gray}C7} & \BenchmarkSourceLabel{bench:external-2-k15-w10:column:index-2-k15-w10-dna5-c7g}{C8}{\color{gray}C8} & \BenchmarkSourceLabel{bench:external-2-k15-w10:column:noindex-2-k15-w10-dna4-m7a}{C9}{\color{gray}C9} & \BenchmarkSourceLabel{bench:external-2-k15-w10:column:noindex-2-k15-w10-dna4-m7i}{C10}{\color{gray}C10} & \BenchmarkSourceLabel{bench:external-2-k15-w10:column:noindex-2-k15-w10-dna4-m4}{C11}{\color{gray}C11} & \BenchmarkSourceLabel{bench:external-2-k15-w10:column:noindex-2-k15-w10-dna4-c7g}{C12}{\color{gray}C12} & \BenchmarkSourceLabel{bench:external-2-k15-w10:column:noindex-2-k15-w10-dna5-m7a}{C13}{\color{gray}C13} & \BenchmarkSourceLabel{bench:external-2-k15-w10:column:noindex-2-k15-w10-dna5-m7i}{C14}{\color{gray}C14} & \BenchmarkSourceLabel{bench:external-2-k15-w10:column:noindex-2-k15-w10-dna5-m4}{C15}{\color{gray}C15} & \BenchmarkSourceLabel{bench:external-2-k15-w10:column:noindex-2-k15-w10-dna5-c7g}{C16}{\color{gray}C16}\\\hline
\rowcolor{BenchBaseline}[0.5pt][0.5pt]\BenchmarkSourceLabel{bench:external-2-k15-w10:row:minimap2}{R1}\shortstack[l]{{\color{gray}R1}\\minimap2} & Baseline & \hspace*{1pt}\begin{tabular}[t]{@{}c@{}}\rule{0pt}{12pt}11.2s\\[-.2pt]{\fontsize{5.8}{6.4}\selectfont \textcolor{BrickRed}{$ \blacktriangledown\,3.0\times $}}\\[-.2pt]10s\\[-.2pt]{\fontsize{5.8}{6.4}\selectfont \textcolor{BrickRed}{$ \blacktriangledown\,3.4\times $}}\end{tabular} & \hspace*{1pt}\begin{tabular}[t]{@{}c@{}}\rule{0pt}{12pt}10.7s\\[-.2pt]{\fontsize{5.8}{6.4}\selectfont \textcolor{BrickRed}{$ \blacktriangledown\,3.0\times $}}\\[-.2pt]8.6s\\[-.2pt]{\fontsize{5.8}{6.4}\selectfont \textcolor{BrickRed}{$ \blacktriangledown\,3.2\times $}}\end{tabular} & \hspace*{1pt}\begin{tabular}[t]{@{}c@{}}\rule{0pt}{12pt}22.5s\\[-.2pt]{\fontsize{5.8}{6.4}\selectfont \textcolor{BrickRed}{$ \blacktriangledown\,3.3\times $}}\\[-.2pt]18.8s\\[-.2pt]{\fontsize{5.8}{6.4}\selectfont \textcolor{BrickRed}{$ \blacktriangledown\,4.2\times $}}\end{tabular} & \hspace*{1pt}\begin{tabular}[t]{@{}c@{}}\rule{0pt}{12pt}1m\,13.1s\\[-.2pt]{\fontsize{5.8}{6.4}\selectfont \textcolor{BrickRed}{$ \blacktriangledown\,8.3\times $}\textsuperscript{g}}\\[-.2pt]1m\,5.7s\\[-.2pt]{\fontsize{5.8}{6.4}\selectfont \textcolor{BrickRed}{$ \blacktriangledown\,8.3\times $}}\end{tabular} & \hspace*{1pt}\begin{tabular}[t]{@{}c@{}}\rule{0pt}{12pt}9.9s\\[-.2pt]{\fontsize{5.8}{6.4}\selectfont \textcolor{BrickRed}{$ \blacktriangledown\,21\% $}}\\[-.2pt]8.8s\\[-.2pt]{\fontsize{5.8}{6.4}\selectfont \textcolor{BrickRed}{$ \blacktriangledown\,1.5\times $}}\end{tabular} & \hspace*{1pt}\begin{tabular}[t]{@{}c@{}}\rule{0pt}{12pt}9.5s\\[-.2pt]{\fontsize{5.8}{6.4}\selectfont \textcolor{BrickRed}{$ \blacktriangledown\,1.4\times $}}\\[-.2pt]7.6s\\[-.2pt]{\fontsize{5.8}{6.4}\selectfont \textcolor{BrickRed}{$ \blacktriangledown\,1.6\times $}}\end{tabular} & \hspace*{1pt}\begin{tabular}[t]{@{}c@{}}\rule{0pt}{12pt}16.5s\\[-.2pt]{\fontsize{5.8}{6.4}\selectfont \textcolor{BrickRed}{$ \blacktriangledown\,20\% $}}\\[-.2pt]12.9s\\[-.2pt]{\fontsize{5.8}{6.4}\selectfont \textcolor{BrickRed}{$ \blacktriangledown\,1.5\times $}}\end{tabular} & \hspace*{1pt}\begin{tabular}[t]{@{}c@{}}\rule{0pt}{12pt}28.2s\\[-.2pt]{\fontsize{5.8}{6.4}\selectfont \textcolor{BrickRed}{$ \blacktriangledown\,1.4\times $}\textsuperscript{g}}\\[-.2pt]26.2s\\[-.2pt]{\fontsize{5.8}{6.4}\selectfont \textcolor{BrickRed}{$ \blacktriangledown\,1.6\times $}}\end{tabular} & \hspace*{1pt}\begin{tabular}[t]{@{}c@{}}\rule{0pt}{12pt}\colorbox{BenchBest}{\strut 53s}\\[-.2pt]{\fontsize{5.8}{6.4}\selectfont \textcolor{ForestGreen}{$ \blacktriangle\,2.2\times $}}\\[-.2pt]\colorbox{BenchBest}{\strut 51.9s}\\[-.2pt]{\fontsize{5.8}{6.4}\selectfont \textcolor{ForestGreen}{$ \blacktriangle\,2.1\times $}}\end{tabular} & \hspace*{1pt}\begin{tabular}[t]{@{}c@{}}\rule{0pt}{12pt}\colorbox{BenchBest}{\strut 56.6s}\\[-.2pt]{\fontsize{5.8}{6.4}\selectfont \textcolor{ForestGreen}{$ \blacktriangle\,2.1\times $}}\\[-.2pt]\colorbox{BenchBest}{\strut 54.6s}\\[-.2pt]{\fontsize{5.8}{6.4}\selectfont \textcolor{ForestGreen}{$ \blacktriangle\,1.9\times $}}\end{tabular} & \hspace*{1pt}\begin{tabular}[t]{@{}c@{}}\rule{0pt}{12pt}\colorbox{BenchBest}{\strut 1m\,55.5s}\\[-.2pt]{\fontsize{5.8}{6.4}\selectfont \textcolor{ForestGreen}{$ \blacktriangle\,1.9\times $}}\\[-.2pt]\colorbox{BenchBest}{\strut 1m\,51.8s}\\[-.2pt]{\fontsize{5.8}{6.4}\selectfont \textcolor{ForestGreen}{$ \blacktriangle\,1.8\times $}}\end{tabular} & \hspace*{1pt}\begin{tabular}[t]{@{}c@{}}\rule{0pt}{12pt}6m\,1.3s\\[-.2pt]{\fontsize{5.8}{6.4}\selectfont \textcolor{BrickRed}{$ \blacktriangledown\,2.9\times $}\textsuperscript{g}}\\[-.2pt]5m\,55.2s\\[-.2pt]{\fontsize{5.8}{6.4}\selectfont \textcolor{BrickRed}{$ \blacktriangledown\,3.0\times $}}\end{tabular} & \hspace*{1pt}\begin{tabular}[t]{@{}c@{}}\rule{0pt}{12pt}\colorbox{BenchBest}{\strut 48.1s}\\[-.2pt]{\fontsize{5.8}{6.4}\selectfont \textcolor{ForestGreen}{$ \blacktriangle\,7.3\times $}}\\[-.2pt]\colorbox{BenchBest}{\strut 47s}\\[-.2pt]{\fontsize{5.8}{6.4}\selectfont \textcolor{ForestGreen}{$ \blacktriangle\,2.4\times $}}\end{tabular} & \hspace*{1pt}\begin{tabular}[t]{@{}c@{}}\rule{0pt}{12pt}\colorbox{BenchBest}{\strut 44.5s}\\[-.2pt]{\fontsize{5.8}{6.4}\selectfont \textcolor{ForestGreen}{$ \blacktriangle\,6.5\times $}}\\[-.2pt]\colorbox{BenchBest}{\strut 42.6s}\\[-.2pt]{\fontsize{5.8}{6.4}\selectfont \textcolor{ForestGreen}{$ \blacktriangle\,2.7\times $}}\end{tabular} & \hspace*{1pt}\begin{tabular}[t]{@{}c@{}}\rule{0pt}{12pt}\colorbox{BenchBest}{\strut 1m\,30.2s}\\[-.2pt]{\fontsize{5.8}{6.4}\selectfont \textcolor{ForestGreen}{$ \blacktriangle\,6.4\times $}}\\[-.2pt]\colorbox{BenchBest}{\strut 1m\,26.7s}\\[-.2pt]{\fontsize{5.8}{6.4}\selectfont \textcolor{ForestGreen}{$ \blacktriangle\,2.3\times $}}\end{tabular} & \hspace*{1pt}\begin{tabular}[t]{@{}c@{}}\rule{0pt}{12pt}\colorbox{BenchBest}{\strut 2m\,46.4s}\\[-.2pt]{\fontsize{5.8}{6.4}\selectfont \textcolor{ForestGreen}{$ \blacktriangle\,4.1\times $}\textsuperscript{g}}\\[-.2pt]2m\,42.7s\\[-.2pt]{\fontsize{5.8}{6.4}\selectfont \textcolor{BrickRed}{$ \blacktriangledown\,16\% $}}\end{tabular}\\[3pt]\hline
\rowcolor{BenchPreferred}[0.5pt][0.5pt]\BenchmarkSourceLabel{bench:external-2-k15-w10:row:bioseq}{R2}\shortstack[l]{{\color{gray}R2}\\BioseqDB} & \begin{tabular}[t]{@{}l@{}}\rule{0pt}{12pt}MPF / Bounded\\SIMD-DP / Block\end{tabular} & \hspace*{1pt}\begin{tabular}[t]{@{}c@{}}\rule{0pt}{12pt}\colorbox{BenchBest}{\strut 3.7s}\\[-.2pt]{\fontsize{5.8}{6.4}\selectfont \textcolor{ForestGreen}{$ \blacktriangle\,3.0\times $}}\\[-.2pt]\colorbox{BenchBest}{\strut 3s}\\[-.2pt]{\fontsize{5.8}{6.4}\selectfont \textcolor{ForestGreen}{$ \blacktriangle\,3.4\times $}}\end{tabular} & \hspace*{1pt}\begin{tabular}[t]{@{}c@{}}\rule{0pt}{12pt}\colorbox{BenchBest}{\strut 3.5s}\\[-.2pt]{\fontsize{5.8}{6.4}\selectfont \textcolor{ForestGreen}{$ \blacktriangle\,3.0\times $}}\\[-.2pt]\colorbox{BenchBest}{\strut 2.7s}\\[-.2pt]{\fontsize{5.8}{6.4}\selectfont \textcolor{ForestGreen}{$ \blacktriangle\,3.2\times $}}\end{tabular} & \hspace*{1pt}\begin{tabular}[t]{@{}c@{}}\rule{0pt}{12pt}\colorbox{BenchBest}{\strut 6.7s}\\[-.2pt]{\fontsize{5.8}{6.4}\selectfont \textcolor{ForestGreen}{$ \blacktriangle\,3.3\times $}}\\[-.2pt]\colorbox{BenchBest}{\strut 4.5s}\\[-.2pt]{\fontsize{5.8}{6.4}\selectfont \textcolor{ForestGreen}{$ \blacktriangle\,4.2\times $}}\end{tabular} & \hspace*{1pt}\begin{tabular}[t]{@{}c@{}}\rule{0pt}{12pt}\colorbox{BenchBest}{\strut 8.8s}\\[-.2pt]{\fontsize{5.8}{6.4}\selectfont \textcolor{ForestGreen}{$ \blacktriangle\,8.3\times $}\textsuperscript{g}}\\[-.2pt]\colorbox{BenchBest}{\strut 7.9s}\\[-.2pt]{\fontsize{5.8}{6.4}\selectfont \textcolor{ForestGreen}{$ \blacktriangle\,8.3\times $}}\end{tabular} & \hspace*{1pt}\begin{tabular}[t]{@{}c@{}}\rule{0pt}{12pt}\colorbox{BenchBest}{\strut 8.2s}\\[-.2pt]{\fontsize{5.8}{6.4}\selectfont \textcolor{ForestGreen}{$ \blacktriangle\,17\% $}}\\[-.2pt]\colorbox{BenchBest}{\strut 5.8s}\\[-.2pt]{\fontsize{5.8}{6.4}\selectfont \textcolor{ForestGreen}{$ \blacktriangle\,1.5\times $}}\end{tabular} & \hspace*{1pt}\begin{tabular}[t]{@{}c@{}}\rule{0pt}{12pt}\colorbox{BenchBest}{\strut 6.6s}\\[-.2pt]{\fontsize{5.8}{6.4}\selectfont \textcolor{ForestGreen}{$ \blacktriangle\,1.4\times $}}\\[-.2pt]\colorbox{BenchBest}{\strut 4.6s}\\[-.2pt]{\fontsize{5.8}{6.4}\selectfont \textcolor{ForestGreen}{$ \blacktriangle\,1.6\times $}}\end{tabular} & \hspace*{1pt}\begin{tabular}[t]{@{}c@{}}\rule{0pt}{12pt}\colorbox{BenchBest}{\strut 13.7s}\\[-.2pt]{\fontsize{5.8}{6.4}\selectfont \textcolor{ForestGreen}{$ \blacktriangle\,17\% $}}\\[-.2pt]\colorbox{BenchBest}{\strut 8.8s}\\[-.2pt]{\fontsize{5.8}{6.4}\selectfont \textcolor{ForestGreen}{$ \blacktriangle\,1.5\times $}}\end{tabular} & \hspace*{1pt}\begin{tabular}[t]{@{}c@{}}\rule{0pt}{12pt}\colorbox{BenchBest}{\strut 20.2s}\\[-.2pt]{\fontsize{5.8}{6.4}\selectfont \textcolor{ForestGreen}{$ \blacktriangle\,28\% $}}\\[-.2pt]\colorbox{BenchBest}{\strut 16.8s}\\[-.2pt]{\fontsize{5.8}{6.4}\selectfont \textcolor{ForestGreen}{$ \blacktriangle\,1.6\times $}}\end{tabular} & \hspace*{1pt}\begin{tabular}[t]{@{}c@{}}\rule{0pt}{12pt}1m\,56.4s\\[-.2pt]{\fontsize{5.8}{6.4}\selectfont \textcolor{BrickRed}{$ \blacktriangledown\,2.2\times $}}\\[-.2pt]1m\,48.1s\\[-.2pt]{\fontsize{5.8}{6.4}\selectfont \textcolor{BrickRed}{$ \blacktriangledown\,2.1\times $}}\end{tabular} & \hspace*{1pt}\begin{tabular}[t]{@{}c@{}}\rule{0pt}{12pt}1m\,56.2s\\[-.2pt]{\fontsize{5.8}{6.4}\selectfont \textcolor{BrickRed}{$ \blacktriangledown\,2.1\times $}}\\[-.2pt]1m\,46.2s\\[-.2pt]{\fontsize{5.8}{6.4}\selectfont \textcolor{BrickRed}{$ \blacktriangledown\,1.9\times $}}\end{tabular} & \hspace*{1pt}\begin{tabular}[t]{@{}c@{}}\rule{0pt}{12pt}3m\,42.9s\\[-.2pt]{\fontsize{5.8}{6.4}\selectfont \textcolor{BrickRed}{$ \blacktriangledown\,1.9\times $}}\\[-.2pt]3m\,25s\\[-.2pt]{\fontsize{5.8}{6.4}\selectfont \textcolor{BrickRed}{$ \blacktriangledown\,1.8\times $}}\end{tabular} & \hspace*{1pt}\begin{tabular}[t]{@{}c@{}}\rule{0pt}{12pt}\colorbox{BenchBest}{\strut 2m\,5.6s}\\[-.2pt]{\fontsize{5.8}{6.4}\selectfont \textcolor{ForestGreen}{$ \blacktriangle\,2.9\times $}\textsuperscript{g}}\\[-.2pt]\colorbox{BenchBest}{\strut 1m\,59.7s}\\[-.2pt]{\fontsize{5.8}{6.4}\selectfont \textcolor{ForestGreen}{$ \blacktriangle\,3.0\times $}}\end{tabular} & \hspace*{1pt}\begin{tabular}[t]{@{}c@{}}\rule{0pt}{12pt}5m\,52.1s\\[-.2pt]{\fontsize{5.8}{6.4}\selectfont \textcolor{BrickRed}{$ \blacktriangledown\,7.3\times $}}\\[-.2pt]1m\,53.5s\\[-.2pt]{\fontsize{5.8}{6.4}\selectfont \textcolor{BrickRed}{$ \blacktriangledown\,2.4\times $}}\end{tabular} & \hspace*{1pt}\begin{tabular}[t]{@{}c@{}}\rule{0pt}{12pt}4m\,49.2s\\[-.2pt]{\fontsize{5.8}{6.4}\selectfont \textcolor{BrickRed}{$ \blacktriangledown\,6.5\times $}}\\[-.2pt]1m\,56.8s\\[-.2pt]{\fontsize{5.8}{6.4}\selectfont \textcolor{BrickRed}{$ \blacktriangledown\,2.7\times $}}\end{tabular} & \hspace*{1pt}\begin{tabular}[t]{@{}c@{}}\rule{0pt}{12pt}9m\,39.5s\\[-.2pt]{\fontsize{5.8}{6.4}\selectfont \textcolor{BrickRed}{$ \blacktriangledown\,6.4\times $}}\\[-.2pt]3m\,18.1s\\[-.2pt]{\fontsize{5.8}{6.4}\selectfont \textcolor{BrickRed}{$ \blacktriangledown\,2.3\times $}}\end{tabular} & \hspace*{1pt}\begin{tabular}[t]{@{}c@{}}\rule{0pt}{12pt}11m\,24.8s\\[-.2pt]{\fontsize{5.8}{6.4}\selectfont \textcolor{BrickRed}{$ \blacktriangledown\,4.1\times $}}\\[-.2pt]\colorbox{BenchBest}{\strut 2m\,20.4s}\\[-.2pt]{\fontsize{5.8}{6.4}\selectfont \textcolor{ForestGreen}{$ \blacktriangle\,14\% $}}\end{tabular}\\[3pt]\hline
\BenchmarkSourceLabel{bench:external-2-k15-w10:row:blastn}{R3}\shortstack[l]{{\color{gray}R3}\\blastn} & Native\textsuperscript{*} & \hspace*{1pt}\begin{tabular}[t]{@{}c@{}}\rule{0pt}{12pt}5m\,36.2s\\[-.2pt]{\fontsize{5.8}{6.4}\selectfont \textcolor{BrickRed}{$ \blacktriangledown\,30.1\times $}\textsuperscript{f}}\\[-.2pt]5m\,30.7s\\[-.2pt]{\fontsize{5.8}{6.4}\selectfont \textcolor{BrickRed}{$ \blacktriangledown\,33.0\times $}}\end{tabular} & \hspace*{1pt}\begin{tabular}[t]{@{}c@{}}\rule{0pt}{12pt}8m\,14.3s\\[-.2pt]{\fontsize{5.8}{6.4}\selectfont \textcolor{BrickRed}{$ \blacktriangledown\,46.4\times $}\textsuperscript{f}}\\[-.2pt]8m\,2.4s\\[-.2pt]{\fontsize{5.8}{6.4}\selectfont \textcolor{BrickRed}{$ \blacktriangledown\,56.0\times $}}\end{tabular} & \hspace*{1pt}\begin{tabular}[t]{@{}c@{}}\rule{0pt}{12pt}10m\,51.3s\\[-.2pt]{\fontsize{5.8}{6.4}\selectfont \textcolor{BrickRed}{$ \blacktriangledown\,28.9\times $}}\\[-.2pt]10m\,35.3s\\[-.2pt]{\fontsize{5.8}{6.4}\selectfont \textcolor{BrickRed}{$ \blacktriangledown\,33.8\times $}}\end{tabular} & \hspace*{1pt}\begin{tabular}[t]{@{}c@{}}\rule{0pt}{12pt}5m\,54.4s\\[-.2pt]{\fontsize{5.8}{6.4}\selectfont \textcolor{BrickRed}{$ \blacktriangledown\,4.9\times $}\textsuperscript{f}}\\[-.2pt]5m\,48.6s\\[-.2pt]{\fontsize{5.8}{6.4}\selectfont \textcolor{BrickRed}{$ \blacktriangledown\,5.3\times $}}\end{tabular} & \hspace*{1pt}\begin{tabular}[t]{@{}c@{}}\rule{0pt}{12pt}6m\,37.9s\\[-.2pt]{\fontsize{5.8}{6.4}\selectfont \textcolor{BrickRed}{$ \blacktriangledown\,40.2\times $}\textsuperscript{a}}\\[-.2pt]4m\,57.9s\\[-.2pt]{\fontsize{5.8}{6.4}\selectfont \textcolor{BrickRed}{$ \blacktriangledown\,33.9\times $}}\end{tabular} & \hspace*{1pt}\begin{tabular}[t]{@{}c@{}}\rule{0pt}{12pt}9m\,32.3s\\[-.2pt]{\fontsize{5.8}{6.4}\selectfont \textcolor{BrickRed}{$ \blacktriangledown\,60.3\times $}\textsuperscript{a}}\\[-.2pt]6m\,18s\\[-.2pt]{\fontsize{5.8}{6.4}\selectfont \textcolor{BrickRed}{$ \blacktriangledown\,49.6\times $}}\end{tabular} & \hspace*{1pt}\begin{tabular}[t]{@{}c@{}}\rule{0pt}{12pt}13m\,27.4s\\[-.2pt]{\fontsize{5.8}{6.4}\selectfont \textcolor{BrickRed}{$ \blacktriangledown\,49.0\times $}\textsuperscript{a}}\\[-.2pt]8m\,52.1s\\[-.2pt]{\fontsize{5.8}{6.4}\selectfont \textcolor{BrickRed}{$ \blacktriangledown\,41.2\times $}}\end{tabular} & \hspace*{1pt}\begin{tabular}[t]{@{}c@{}}\rule{0pt}{12pt}6m\,53.6s\\[-.2pt]{\fontsize{5.8}{6.4}\selectfont \textcolor{BrickRed}{$ \blacktriangledown\,14.7\times $}\textsuperscript{a}}\\[-.2pt]4m\,51.8s\\[-.2pt]{\fontsize{5.8}{6.4}\selectfont \textcolor{BrickRed}{$ \blacktriangledown\,11.1\times $}}\end{tabular} & \hspace*{1pt}\begin{tabular}[t]{@{}c@{}}\rule{0pt}{12pt}5m\,39.9s\\[-.2pt]{\fontsize{5.8}{6.4}\selectfont \textcolor{BrickRed}{$ \blacktriangledown\,6.4\times $}}\\[-.2pt]5m\,34.5s\\[-.2pt]{\fontsize{5.8}{6.4}\selectfont \textcolor{BrickRed}{$ \blacktriangledown\,6.4\times $}}\end{tabular} & \hspace*{1pt}\begin{tabular}[t]{@{}c@{}}\rule{0pt}{12pt}8m\,22.6s\\[-.2pt]{\fontsize{5.8}{6.4}\selectfont \textcolor{BrickRed}{$ \blacktriangledown\,8.9\times $}}\\[-.2pt]8m\,10.6s\\[-.2pt]{\fontsize{5.8}{6.4}\selectfont \textcolor{BrickRed}{$ \blacktriangledown\,9.0\times $}}\end{tabular} & \hspace*{1pt}\begin{tabular}[t]{@{}c@{}}\rule{0pt}{12pt}11m\,6.9s\\[-.2pt]{\fontsize{5.8}{6.4}\selectfont \textcolor{BrickRed}{$ \blacktriangledown\,5.8\times $}}\\[-.2pt]10m\,50.9s\\[-.2pt]{\fontsize{5.8}{6.4}\selectfont \textcolor{BrickRed}{$ \blacktriangledown\,5.8\times $}}\end{tabular} & \hspace*{1pt}\begin{tabular}[t]{@{}c@{}}\rule{0pt}{12pt}6m\,6.4s\\[-.2pt]{\fontsize{5.8}{6.4}\selectfont \textcolor{gray}{$ \blacktriangleright\,1\% $}}\\[-.2pt]6m\,0.6s\\[-.2pt]{\fontsize{5.8}{6.4}\selectfont \textcolor{gray}{$ \blacktriangleright\,2\% $}}\end{tabular} & \hspace*{1pt}\begin{tabular}[t]{@{}c@{}}\rule{0pt}{12pt}6m\,51.8s\\[-.2pt]{\fontsize{5.8}{6.4}\selectfont \textcolor{BrickRed}{$ \blacktriangledown\,8.6\times $}\textsuperscript{a}}\\[-.2pt]5m\,12.2s\\[-.2pt]{\fontsize{5.8}{6.4}\selectfont \textcolor{BrickRed}{$ \blacktriangledown\,6.6\times $}}\end{tabular} & \hspace*{1pt}\begin{tabular}[t]{@{}c@{}}\rule{0pt}{12pt}9m\,48.4s\\[-.2pt]{\fontsize{5.8}{6.4}\selectfont \textcolor{BrickRed}{$ \blacktriangledown\,13.2\times $}\textsuperscript{a}}\\[-.2pt]6m\,35.3s\\[-.2pt]{\fontsize{5.8}{6.4}\selectfont \textcolor{BrickRed}{$ \blacktriangledown\,9.3\times $}}\end{tabular} & \hspace*{1pt}\begin{tabular}[t]{@{}c@{}}\rule{0pt}{12pt}13m\,55.3s\\[-.2pt]{\fontsize{5.8}{6.4}\selectfont \textcolor{BrickRed}{$ \blacktriangledown\,9.3\times $}\textsuperscript{a}}\\[-.2pt]9m\,23.3s\\[-.2pt]{\fontsize{5.8}{6.4}\selectfont \textcolor{BrickRed}{$ \blacktriangledown\,6.5\times $}}\end{tabular} & \hspace*{1pt}\begin{tabular}[t]{@{}c@{}}\rule{0pt}{12pt}7m\,14.5s\\[-.2pt]{\fontsize{5.8}{6.4}\selectfont \textcolor{BrickRed}{$ \blacktriangledown\,2.6\times $}\textsuperscript{a}}\\[-.2pt]5m\,12.1s\\[-.2pt]{\fontsize{5.8}{6.4}\selectfont \textcolor{BrickRed}{$ \blacktriangledown\,1.9\times $}}\end{tabular}\\[3pt]\hline
\BenchmarkSourceLabel{bench:external-2-k15-w10:row:winnowmap}{R4}\shortstack[l]{{\color{gray}R4}\\winnowmap} & Native\textsuperscript{*} & \hspace*{1pt}\begin{tabular}[t]{@{}c@{}}\rule{0pt}{12pt}12.9s\\[-.2pt]{\fontsize{5.8}{6.4}\selectfont \textcolor{BrickRed}{$ \blacktriangledown\,15\% $}}\\[-.2pt]11.1s\\[-.2pt]{\fontsize{5.8}{6.4}\selectfont \textcolor{BrickRed}{$ \blacktriangledown\,11\% $}}\end{tabular} & \hspace*{1pt}\begin{tabular}[t]{@{}c@{}}\rule{0pt}{12pt}13.5s\\[-.2pt]{\fontsize{5.8}{6.4}\selectfont \textcolor{BrickRed}{$ \blacktriangledown\,27\% $}}\\[-.2pt]10.3s\\[-.2pt]{\fontsize{5.8}{6.4}\selectfont \textcolor{BrickRed}{$ \blacktriangledown\,20\% $}}\end{tabular} & \hspace*{1pt}\begin{tabular}[t]{@{}c@{}}\rule{0pt}{12pt}23.2s\\[-.2pt]{\fontsize{5.8}{6.4}\selectfont \textcolor{gray}{$ \blacktriangleright\,3\% $}}\\[-.2pt]17.1s\\[-.2pt]{\fontsize{5.8}{6.4}\selectfont \textcolor{ForestGreen}{$ \blacktriangle\,9\% $}}\end{tabular} & \hspace*{1pt}\begin{tabular}[t]{@{}c@{}}\rule{0pt}{12pt}44.3s\\[-.2pt]{\fontsize{5.8}{6.4}\selectfont \textcolor{ForestGreen}{$ \blacktriangle\,1.7\times $}\textsuperscript{d}}\\[-.2pt]37.8s\\[-.2pt]{\fontsize{5.8}{6.4}\selectfont \textcolor{ForestGreen}{$ \blacktriangle\,1.7\times $}}\end{tabular} & \hspace*{1pt}\begin{tabular}[t]{@{}c@{}}\rule{0pt}{12pt}12.3s\\[-.2pt]{\fontsize{5.8}{6.4}\selectfont \textcolor{BrickRed}{$ \blacktriangledown\,24\% $}}\\[-.2pt]10.6s\\[-.2pt]{\fontsize{5.8}{6.4}\selectfont \textcolor{BrickRed}{$ \blacktriangledown\,21\% $}}\end{tabular} & \hspace*{1pt}\begin{tabular}[t]{@{}c@{}}\rule{0pt}{12pt}13s\\[-.2pt]{\fontsize{5.8}{6.4}\selectfont \textcolor{BrickRed}{$ \blacktriangledown\,1.4\times $}}\\[-.2pt]10s\\[-.2pt]{\fontsize{5.8}{6.4}\selectfont \textcolor{BrickRed}{$ \blacktriangledown\,1.3\times $}}\end{tabular} & \hspace*{1pt}\begin{tabular}[t]{@{}c@{}}\rule{0pt}{12pt}23.8s\\[-.2pt]{\fontsize{5.8}{6.4}\selectfont \textcolor{BrickRed}{$ \blacktriangledown\,1.4\times $}}\\[-.2pt]17.6s\\[-.2pt]{\fontsize{5.8}{6.4}\selectfont \textcolor{BrickRed}{$ \blacktriangledown\,1.4\times $}}\end{tabular} & \hspace*{1pt}\begin{tabular}[t]{@{}c@{}}\rule{0pt}{12pt}37.3s\\[-.2pt]{\fontsize{5.8}{6.4}\selectfont \textcolor{BrickRed}{$ \blacktriangledown\,1.3\times $}\textsuperscript{d}}\\[-.2pt]31.4s\\[-.2pt]{\fontsize{5.8}{6.4}\selectfont \textcolor{BrickRed}{$ \blacktriangledown\,20\% $}}\end{tabular} & \hspace*{1pt}\begin{tabular}[t]{@{}c@{}}\rule{0pt}{12pt}3m\,33.2s\\[-.2pt]{\fontsize{5.8}{6.4}\selectfont \textcolor{BrickRed}{$ \blacktriangledown\,4.0\times $}}\\[-.2pt]3m\,31.5s\\[-.2pt]{\fontsize{5.8}{6.4}\selectfont \textcolor{BrickRed}{$ \blacktriangledown\,4.1\times $}}\end{tabular} & \hspace*{1pt}\begin{tabular}[t]{@{}c@{}}\rule{0pt}{12pt}3m\,22.6s\\[-.2pt]{\fontsize{5.8}{6.4}\selectfont \textcolor{BrickRed}{$ \blacktriangledown\,3.6\times $}}\\[-.2pt]3m\,19.3s\\[-.2pt]{\fontsize{5.8}{6.4}\selectfont \textcolor{BrickRed}{$ \blacktriangledown\,3.7\times $}}\end{tabular} & \hspace*{1pt}\begin{tabular}[t]{@{}c@{}}\rule{0pt}{12pt}6m\,21.3s\\[-.2pt]{\fontsize{5.8}{6.4}\selectfont \textcolor{BrickRed}{$ \blacktriangledown\,3.3\times $}}\\[-.2pt]6m\,15s\\[-.2pt]{\fontsize{5.8}{6.4}\selectfont \textcolor{BrickRed}{$ \blacktriangledown\,3.4\times $}}\end{tabular} & \hspace*{1pt}\begin{tabular}[t]{@{}c@{}}\rule{0pt}{12pt}3m\,56.6s\\[-.2pt]{\fontsize{5.8}{6.4}\selectfont \textcolor{ForestGreen}{$ \blacktriangle\,1.5\times $}\textsuperscript{d}}\\[-.2pt]3m\,50.3s\\[-.2pt]{\fontsize{5.8}{6.4}\selectfont \textcolor{ForestGreen}{$ \blacktriangle\,1.5\times $}}\end{tabular} & \hspace*{1pt}\begin{tabular}[t]{@{}c@{}}\rule{0pt}{12pt}3m\,7.4s\\[-.2pt]{\fontsize{5.8}{6.4}\selectfont \textcolor{BrickRed}{$ \blacktriangledown\,3.9\times $}}\\[-.2pt]3m\,5.8s\\[-.2pt]{\fontsize{5.8}{6.4}\selectfont \textcolor{BrickRed}{$ \blacktriangledown\,4.0\times $}}\end{tabular} & \hspace*{1pt}\begin{tabular}[t]{@{}c@{}}\rule{0pt}{12pt}3m\,2.2s\\[-.2pt]{\fontsize{5.8}{6.4}\selectfont \textcolor{BrickRed}{$ \blacktriangledown\,4.1\times $}}\\[-.2pt]2m\,59.3s\\[-.2pt]{\fontsize{5.8}{6.4}\selectfont \textcolor{BrickRed}{$ \blacktriangledown\,4.2\times $}}\end{tabular} & \hspace*{1pt}\begin{tabular}[t]{@{}c@{}}\rule{0pt}{12pt}6m\,35.9s\\[-.2pt]{\fontsize{5.8}{6.4}\selectfont \textcolor{BrickRed}{$ \blacktriangledown\,4.4\times $}}\\[-.2pt]6m\,29.8s\\[-.2pt]{\fontsize{5.8}{6.4}\selectfont \textcolor{BrickRed}{$ \blacktriangledown\,4.5\times $}}\end{tabular} & \hspace*{1pt}\begin{tabular}[t]{@{}c@{}}\rule{0pt}{12pt}3m\,33.9s\\[-.2pt]{\fontsize{5.8}{6.4}\selectfont \textcolor{BrickRed}{$ \blacktriangledown\,29\% $}\textsuperscript{d}}\\[-.2pt]3m\,28.1s\\[-.2pt]{\fontsize{5.8}{6.4}\selectfont \textcolor{BrickRed}{$ \blacktriangledown\,28\% $}}\end{tabular}\\[3pt]\hline
\BenchmarkSourceLabel{bench:external-2-k15-w10:row:lexicmap}{R5}\shortstack[l]{{\color{gray}R5}\\lexicmap} & Native\textsuperscript{*} & \hspace*{1pt}\begin{tabular}[t]{@{}c@{}}\rule{0pt}{12pt}2m\,14.1s\\[-.2pt]{\fontsize{5.8}{6.4}\selectfont \textcolor{BrickRed}{$ \blacktriangledown\,12.0\times $}}\\[-.2pt]2m\,14s\\[-.2pt]{\fontsize{5.8}{6.4}\selectfont \textcolor{BrickRed}{$ \blacktriangledown\,13.4\times $}}\end{tabular} & \hspace*{1pt}\begin{tabular}[t]{@{}c@{}}\rule{0pt}{12pt}3m\,26.7s\\[-.2pt]{\fontsize{5.8}{6.4}\selectfont \textcolor{BrickRed}{$ \blacktriangledown\,19.4\times $}}\\[-.2pt]3m\,26.6s\\[-.2pt]{\fontsize{5.8}{6.4}\selectfont \textcolor{BrickRed}{$ \blacktriangledown\,24.0\times $}}\end{tabular} & \hspace*{1pt}\begin{tabular}[t]{@{}c@{}}\rule{0pt}{12pt}4m\,40.1s\\[-.2pt]{\fontsize{5.8}{6.4}\selectfont \textcolor{BrickRed}{$ \blacktriangledown\,12.4\times $}}\\[-.2pt]4m\,39.9s\\[-.2pt]{\fontsize{5.8}{6.4}\selectfont \textcolor{BrickRed}{$ \blacktriangledown\,14.9\times $}}\end{tabular} & \hspace*{1pt}\begin{tabular}[t]{@{}c@{}}\rule{0pt}{12pt}3m\,24.3s\\[-.2pt]{\fontsize{5.8}{6.4}\selectfont \textcolor{BrickRed}{$ \blacktriangledown\,2.8\times $}}\\[-.2pt]3m\,24.2s\\[-.2pt]{\fontsize{5.8}{6.4}\selectfont \textcolor{BrickRed}{$ \blacktriangledown\,3.1\times $}}\end{tabular} & \hspace*{1pt}\begin{tabular}[t]{@{}c@{}}\rule{0pt}{12pt}2m\,12.9s\\[-.2pt]{\fontsize{5.8}{6.4}\selectfont \textcolor{BrickRed}{$ \blacktriangledown\,13.4\times $}\textsuperscript{b}}\\[-.2pt]2m\,12.5s\\[-.2pt]{\fontsize{5.8}{6.4}\selectfont \textcolor{BrickRed}{$ \blacktriangledown\,15.1\times $}}\end{tabular} & \hspace*{1pt}\begin{tabular}[t]{@{}c@{}}\rule{0pt}{12pt}3m\,24.1s\\[-.2pt]{\fontsize{5.8}{6.4}\selectfont \textcolor{BrickRed}{$ \blacktriangledown\,21.5\times $}\textsuperscript{b}}\\[-.2pt]3m\,23.6s\\[-.2pt]{\fontsize{5.8}{6.4}\selectfont \textcolor{BrickRed}{$ \blacktriangledown\,26.7\times $}}\end{tabular} & \hspace*{1pt}\begin{tabular}[t]{@{}c@{}}\rule{0pt}{12pt}4m\,56.4s\\[-.2pt]{\fontsize{5.8}{6.4}\selectfont \textcolor{BrickRed}{$ \blacktriangledown\,18.0\times $}\textsuperscript{b}}\\[-.2pt]4m\,55.6s\\[-.2pt]{\fontsize{5.8}{6.4}\selectfont \textcolor{BrickRed}{$ \blacktriangledown\,22.9\times $}}\end{tabular} & \hspace*{1pt}\begin{tabular}[t]{@{}c@{}}\rule{0pt}{12pt}3m\,22.8s\\[-.2pt]{\fontsize{5.8}{6.4}\selectfont \textcolor{BrickRed}{$ \blacktriangledown\,7.2\times $}\textsuperscript{b}}\\[-.2pt]3m\,22.5s\\[-.2pt]{\fontsize{5.8}{6.4}\selectfont \textcolor{BrickRed}{$ \blacktriangledown\,7.7\times $}}\end{tabular} & \hspace*{1pt}\begin{tabular}[t]{@{}c@{}}\rule{0pt}{12pt}3m\,18.9s\\[-.2pt]{\fontsize{5.8}{6.4}\selectfont \textcolor{BrickRed}{$ \blacktriangledown\,3.8\times $}}\\[-.2pt]3m\,18.7s\\[-.2pt]{\fontsize{5.8}{6.4}\selectfont \textcolor{BrickRed}{$ \blacktriangledown\,3.8\times $}}\end{tabular} & \hspace*{1pt}\begin{tabular}[t]{@{}c@{}}\rule{0pt}{12pt}5m\,10.5s\\[-.2pt]{\fontsize{5.8}{6.4}\selectfont \textcolor{BrickRed}{$ \blacktriangledown\,5.5\times $}}\\[-.2pt]5m\,10.3s\\[-.2pt]{\fontsize{5.8}{6.4}\selectfont \textcolor{BrickRed}{$ \blacktriangledown\,5.7\times $}}\end{tabular} & \hspace*{1pt}\begin{tabular}[t]{@{}c@{}}\rule{0pt}{12pt}8m\,19.4s\\[-.2pt]{\fontsize{5.8}{6.4}\selectfont \textcolor{BrickRed}{$ \blacktriangledown\,4.3\times $}}\\[-.2pt]8m\,19.2s\\[-.2pt]{\fontsize{5.8}{6.4}\selectfont \textcolor{BrickRed}{$ \blacktriangledown\,4.5\times $}}\end{tabular} & \hspace*{1pt}\begin{tabular}[t]{@{}c@{}}\rule{0pt}{12pt}5m\,43.1s\\[-.2pt]{\fontsize{5.8}{6.4}\selectfont \textcolor{ForestGreen}{$ \blacktriangle\,5\% $}}\\[-.2pt]5m\,43s\\[-.2pt]{\fontsize{5.8}{6.4}\selectfont \textcolor{gray}{$ \blacktriangleright\,3\% $}}\end{tabular} & \hspace*{1pt}\begin{tabular}[t]{@{}c@{}}\rule{0pt}{12pt}3m\,16.5s\\[-.2pt]{\fontsize{5.8}{6.4}\selectfont \textcolor{BrickRed}{$ \blacktriangledown\,4.1\times $}\textsuperscript{b}}\\[-.2pt]3m\,16.1s\\[-.2pt]{\fontsize{5.8}{6.4}\selectfont \textcolor{BrickRed}{$ \blacktriangledown\,4.2\times $}}\end{tabular} & \hspace*{1pt}\begin{tabular}[t]{@{}c@{}}\rule{0pt}{12pt}5m\,5s\\[-.2pt]{\fontsize{5.8}{6.4}\selectfont \textcolor{BrickRed}{$ \blacktriangledown\,6.9\times $}\textsuperscript{b}}\\[-.2pt]5m\,4.5s\\[-.2pt]{\fontsize{5.8}{6.4}\selectfont \textcolor{BrickRed}{$ \blacktriangledown\,7.2\times $}}\end{tabular} & \hspace*{1pt}\begin{tabular}[t]{@{}c@{}}\rule{0pt}{12pt}7m\,40.8s\\[-.2pt]{\fontsize{5.8}{6.4}\selectfont \textcolor{BrickRed}{$ \blacktriangledown\,5.1\times $}\textsuperscript{b}}\\[-.2pt]7m\,39.9s\\[-.2pt]{\fontsize{5.8}{6.4}\selectfont \textcolor{BrickRed}{$ \blacktriangledown\,5.3\times $}}\end{tabular} & \hspace*{1pt}\begin{tabular}[t]{@{}c@{}}\rule{0pt}{12pt}5m\,52.3s\\[-.2pt]{\fontsize{5.8}{6.4}\selectfont \textcolor{BrickRed}{$ \blacktriangledown\,2.1\times $}\textsuperscript{b}}\\[-.2pt]5m\,51.9s\\[-.2pt]{\fontsize{5.8}{6.4}\selectfont \textcolor{BrickRed}{$ \blacktriangledown\,2.2\times $}}\end{tabular}\\[3pt]\hline
\end{tabular}\par}
\smallskip{\footnotesize\noindent Each cell: A, comparison, B, comparison. $\cdot$: no measurement. n/a: no comparison.\par
\noindent Blue: one lowest A and one lowest B measurement for each column for each configuration.\par
\noindent Gray: minimap2. Yellow: my preferred MPF / Bounded / SIMD-DP / Block configuration.\par
\noindent Comparisons use minimap2. The minimap2 row compares with BioseqDB.\par
\noindent\textsuperscript{*} Each algorithm except BLAST required intensive splitting, sometimes not equal. The splitting was the biggest possible to allow algorithm not to crash.\par
\noindent\textsuperscript{a}BLAST scores N as -3, compared with -1 in minimap2.\par
\noindent\textsuperscript{b}LexicMap stores reference N as A.\par
\noindent\textsuperscript{d}Winnowmap uses overlapping 250 Mb records and 499M index batches. Merging results can affect MAPQ.\par
\noindent\textsuperscript{f}Swap settings were not recorded.\par

}
\clearpage
\refstepcounter{table}\label{tab:bench:external-2-k20-w15}
\noindent{\large\bfseries (I): Table \thetable: aligner comparison, 2 Gbp, k=20, w=15\hfill 4/4}\par\smallskip
{\fontsize{7.1}{8.1}\selectfont\setlength{\tabcolsep}{1pt}\setlength{\fboxsep}{2pt}\renewcommand{\arraystretch}{1.12}
\noindent\begin{tabular}{@{}p{15mm}p{26mm}|p{13.55mm}|p{13.55mm}|p{13.55mm}|p{13.55mm}|p{13.55mm}|p{13.55mm}|p{13.55mm}|p{13.55mm}|p{13.55mm}|p{13.55mm}|p{13.55mm}|p{13.55mm}|p{13.55mm}|p{13.55mm}|p{13.55mm}|p{13.55mm}|@{}}\hline
Aligner & Configuration & \multicolumn{8}{c|}{\textbf{Index}} & \multicolumn{8}{c|}{\textbf{No index}}\\
 &  & \multicolumn{8}{c|}{\textbf{2 Gbp}} & \multicolumn{8}{c|}{\textbf{2 Gbp}}\\
 &  & \multicolumn{8}{c|}{\textbf{k=20, w=15}} & \multicolumn{8}{c|}{\textbf{k=20, w=15}}\\
 &  & \multicolumn{4}{c|}{\textbf{DNA4}} & \multicolumn{4}{c|}{\textbf{DNA5}} & \multicolumn{4}{c|}{\textbf{DNA4}} & \multicolumn{4}{c|}{\textbf{DNA5}}\\
 &  & \multicolumn{1}{c|}{\textbf{M7A}} & \multicolumn{1}{c|}{\textbf{M7I}} & \multicolumn{1}{c|}{\textbf{M4}} & \multicolumn{1}{c|}{\textbf{C7G}} & \multicolumn{1}{c|}{\textbf{M7A}} & \multicolumn{1}{c|}{\textbf{M7I}} & \multicolumn{1}{c|}{\textbf{M4}} & \multicolumn{1}{c|}{\textbf{C7G}} & \multicolumn{1}{c|}{\textbf{M7A}} & \multicolumn{1}{c|}{\textbf{M7I}} & \multicolumn{1}{c|}{\textbf{M4}} & \multicolumn{1}{c|}{\textbf{C7G}} & \multicolumn{1}{c|}{\textbf{M7A}} & \multicolumn{1}{c|}{\textbf{M7I}} & \multicolumn{1}{c|}{\textbf{M4}} & \multicolumn{1}{c|}{\textbf{C7G}}\\
 & & \BenchmarkSourceLabel{bench:external-2-k20-w15:column:index-2-k20-w15-dna4-m7a}{C1}{\color{gray}C1} & \BenchmarkSourceLabel{bench:external-2-k20-w15:column:index-2-k20-w15-dna4-m7i}{C2}{\color{gray}C2} & \BenchmarkSourceLabel{bench:external-2-k20-w15:column:index-2-k20-w15-dna4-m4}{C3}{\color{gray}C3} & \BenchmarkSourceLabel{bench:external-2-k20-w15:column:index-2-k20-w15-dna4-c7g}{C4}{\color{gray}C4} & \BenchmarkSourceLabel{bench:external-2-k20-w15:column:index-2-k20-w15-dna5-m7a}{C5}{\color{gray}C5} & \BenchmarkSourceLabel{bench:external-2-k20-w15:column:index-2-k20-w15-dna5-m7i}{C6}{\color{gray}C6} & \BenchmarkSourceLabel{bench:external-2-k20-w15:column:index-2-k20-w15-dna5-m4}{C7}{\color{gray}C7} & \BenchmarkSourceLabel{bench:external-2-k20-w15:column:index-2-k20-w15-dna5-c7g}{C8}{\color{gray}C8} & \BenchmarkSourceLabel{bench:external-2-k20-w15:column:noindex-2-k20-w15-dna4-m7a}{C9}{\color{gray}C9} & \BenchmarkSourceLabel{bench:external-2-k20-w15:column:noindex-2-k20-w15-dna4-m7i}{C10}{\color{gray}C10} & \BenchmarkSourceLabel{bench:external-2-k20-w15:column:noindex-2-k20-w15-dna4-m4}{C11}{\color{gray}C11} & \BenchmarkSourceLabel{bench:external-2-k20-w15:column:noindex-2-k20-w15-dna4-c7g}{C12}{\color{gray}C12} & \BenchmarkSourceLabel{bench:external-2-k20-w15:column:noindex-2-k20-w15-dna5-m7a}{C13}{\color{gray}C13} & \BenchmarkSourceLabel{bench:external-2-k20-w15:column:noindex-2-k20-w15-dna5-m7i}{C14}{\color{gray}C14} & \BenchmarkSourceLabel{bench:external-2-k20-w15:column:noindex-2-k20-w15-dna5-m4}{C15}{\color{gray}C15} & \BenchmarkSourceLabel{bench:external-2-k20-w15:column:noindex-2-k20-w15-dna5-c7g}{C16}{\color{gray}C16}\\\hline
\rowcolor{BenchBaseline}[0.5pt][0.5pt]\BenchmarkSourceLabel{bench:external-2-k20-w15:row:minimap2}{R1}\shortstack[l]{{\color{gray}R1}\\minimap2} & Baseline & \hspace*{1pt}\begin{tabular}[t]{@{}c@{}}\rule{0pt}{12pt}19.6s\\[-.2pt]{\fontsize{5.8}{6.4}\selectfont \textcolor{BrickRed}{$ \blacktriangledown\,6.6\times $}}\\[-.2pt]18.6s\\[-.2pt]{\fontsize{5.8}{6.4}\selectfont \textcolor{BrickRed}{$ \blacktriangledown\,8.1\times $}}\end{tabular} & \hspace*{1pt}\begin{tabular}[t]{@{}c@{}}\rule{0pt}{12pt}18.2s\\[-.2pt]{\fontsize{5.8}{6.4}\selectfont \textcolor{BrickRed}{$ \blacktriangledown\,6.9\times $}}\\[-.2pt]16.5s\\[-.2pt]{\fontsize{5.8}{6.4}\selectfont \textcolor{BrickRed}{$ \blacktriangledown\,8.6\times $}}\end{tabular} & \hspace*{1pt}\begin{tabular}[t]{@{}c@{}}\rule{0pt}{12pt}42.5s\\[-.2pt]{\fontsize{5.8}{6.4}\selectfont \textcolor{BrickRed}{$ \blacktriangledown\,7.9\times $}}\\[-.2pt]39.4s\\[-.2pt]{\fontsize{5.8}{6.4}\selectfont \textcolor{BrickRed}{$ \blacktriangledown\,11.7\times $}}\end{tabular} & \hspace*{1pt}\begin{tabular}[t]{@{}c@{}}\rule{0pt}{12pt}6m\,29.2s\\[-.2pt]{\fontsize{5.8}{6.4}\selectfont \textcolor{BrickRed}{$ \blacktriangledown\,86.7\times $}}\\[-.2pt]5m\,56.9s\\[-.2pt]{\fontsize{5.8}{6.4}\selectfont \textcolor{BrickRed}{$ \blacktriangledown\,97.9\times $}}\end{tabular} & \hspace*{1pt}\begin{tabular}[t]{@{}c@{}}\rule{0pt}{12pt}9.3s\\[-.2pt]{\fontsize{5.8}{6.4}\selectfont \textcolor{BrickRed}{$ \blacktriangledown\,24\% $}}\\[-.2pt]8.4s\\[-.2pt]{\fontsize{5.8}{6.4}\selectfont \textcolor{BrickRed}{$ \blacktriangledown\,1.6\times $}}\end{tabular} & \hspace*{1pt}\begin{tabular}[t]{@{}c@{}}\rule{0pt}{12pt}8.7s\\[-.2pt]{\fontsize{5.8}{6.4}\selectfont \textcolor{BrickRed}{$ \blacktriangledown\,1.5\times $}}\\[-.2pt]7s\\[-.2pt]{\fontsize{5.8}{6.4}\selectfont \textcolor{BrickRed}{$ \blacktriangledown\,1.7\times $}}\end{tabular} & \hspace*{1pt}\begin{tabular}[t]{@{}c@{}}\rule{0pt}{12pt}15.8s\\[-.2pt]{\fontsize{5.8}{6.4}\selectfont \textcolor{BrickRed}{$ \blacktriangledown\,26\% $}}\\[-.2pt]12.7s\\[-.2pt]{\fontsize{5.8}{6.4}\selectfont \textcolor{BrickRed}{$ \blacktriangledown\,1.6\times $}}\end{tabular} & \hspace*{1pt}\begin{tabular}[t]{@{}c@{}}\rule{0pt}{12pt}18.5s\\[-.2pt]{\fontsize{5.8}{6.4}\selectfont \textcolor{BrickRed}{$ \blacktriangledown\,1.4\times $}\textsuperscript{g}}\\[-.2pt]17s\\[-.2pt]{\fontsize{5.8}{6.4}\selectfont \textcolor{BrickRed}{$ \blacktriangledown\,1.6\times $}}\end{tabular} & \hspace*{1pt}\begin{tabular}[t]{@{}c@{}}\rule{0pt}{12pt}\colorbox{BenchBest}{\strut 42.4s}\\[-.2pt]{\fontsize{5.8}{6.4}\selectfont \textcolor{ForestGreen}{$ \blacktriangle\,2.0\times $}}\\[-.2pt]\colorbox{BenchBest}{\strut 41.4s}\\[-.2pt]{\fontsize{5.8}{6.4}\selectfont \textcolor{ForestGreen}{$ \blacktriangle\,1.9\times $}}\end{tabular} & \hspace*{1pt}\begin{tabular}[t]{@{}c@{}}\rule{0pt}{12pt}\colorbox{BenchBest}{\strut 42.9s}\\[-.2pt]{\fontsize{5.8}{6.4}\selectfont \textcolor{ForestGreen}{$ \blacktriangle\,2.0\times $}}\\[-.2pt]\colorbox{BenchBest}{\strut 41.1s}\\[-.2pt]{\fontsize{5.8}{6.4}\selectfont \textcolor{ForestGreen}{$ \blacktriangle\,1.8\times $}}\end{tabular} & \hspace*{1pt}\begin{tabular}[t]{@{}c@{}}\rule{0pt}{12pt}\colorbox{BenchBest}{\strut 1m\,26.9s}\\[-.2pt]{\fontsize{5.8}{6.4}\selectfont \textcolor{ForestGreen}{$ \blacktriangle\,1.9\times $}}\\[-.2pt]\colorbox{BenchBest}{\strut 1m\,23.7s}\\[-.2pt]{\fontsize{5.8}{6.4}\selectfont \textcolor{ForestGreen}{$ \blacktriangle\,1.7\times $}}\end{tabular} & \hspace*{1pt}\begin{tabular}[t]{@{}c@{}}\rule{0pt}{12pt}4m\,51.6s\\[-.2pt]{\fontsize{5.8}{6.4}\selectfont \textcolor{BrickRed}{$ \blacktriangledown\,2.9\times $}\textsuperscript{g}}\\[-.2pt]4m\,46.6s\\[-.2pt]{\fontsize{5.8}{6.4}\selectfont \textcolor{BrickRed}{$ \blacktriangledown\,3.0\times $}}\end{tabular} & \hspace*{1pt}\begin{tabular}[t]{@{}c@{}}\rule{0pt}{12pt}\colorbox{BenchBest}{\strut 35s}\\[-.2pt]{\fontsize{5.8}{6.4}\selectfont \textcolor{ForestGreen}{$ \blacktriangle\,8.6\times $}}\\[-.2pt]\colorbox{BenchBest}{\strut 34s}\\[-.2pt]{\fontsize{5.8}{6.4}\selectfont \textcolor{ForestGreen}{$ \blacktriangle\,2.5\times $}}\end{tabular} & \hspace*{1pt}\begin{tabular}[t]{@{}c@{}}\rule{0pt}{12pt}\colorbox{BenchBest}{\strut 35.9s}\\[-.2pt]{\fontsize{5.8}{6.4}\selectfont \textcolor{ForestGreen}{$ \blacktriangle\,6.7\times $}}\\[-.2pt]\colorbox{BenchBest}{\strut 34.2s}\\[-.2pt]{\fontsize{5.8}{6.4}\selectfont \textcolor{ForestGreen}{$ \blacktriangle\,2.5\times $}}\end{tabular} & \hspace*{1pt}\begin{tabular}[t]{@{}c@{}}\rule{0pt}{12pt}\colorbox{BenchBest}{\strut 1m\,14.1s}\\[-.2pt]{\fontsize{5.8}{6.4}\selectfont \textcolor{ForestGreen}{$ \blacktriangle\,6.9\times $}}\\[-.2pt]\colorbox{BenchBest}{\strut 1m\,11.1s}\\[-.2pt]{\fontsize{5.8}{6.4}\selectfont \textcolor{ForestGreen}{$ \blacktriangle\,2.2\times $}}\end{tabular} & \hspace*{1pt}\begin{tabular}[t]{@{}c@{}}\rule{0pt}{12pt}\colorbox{BenchBest}{\strut 1m\,0.1s}\\[-.2pt]{\fontsize{5.8}{6.4}\selectfont \textcolor{ForestGreen}{$ \blacktriangle\,10.4\times $}\textsuperscript{g}}\\[-.2pt]\colorbox{BenchBest}{\strut 58.6s}\\[-.2pt]{\fontsize{5.8}{6.4}\selectfont \textcolor{ForestGreen}{$ \blacktriangle\,1.9\times $}}\end{tabular}\\[3pt]\hline
\rowcolor{BenchPreferred}[0.5pt][0.5pt]\BenchmarkSourceLabel{bench:external-2-k20-w15:row:bioseq}{R2}\shortstack[l]{{\color{gray}R2}\\BioseqDB} & \begin{tabular}[t]{@{}l@{}}\rule{0pt}{12pt}MPF / Bounded\\SIMD-DP / Block\end{tabular} & \hspace*{1pt}\begin{tabular}[t]{@{}c@{}}\rule{0pt}{12pt}\colorbox{BenchBest}{\strut 3s}\\[-.2pt]{\fontsize{5.8}{6.4}\selectfont \textcolor{ForestGreen}{$ \blacktriangle\,6.6\times $}}\\[-.2pt]\colorbox{BenchBest}{\strut 2.3s}\\[-.2pt]{\fontsize{5.8}{6.4}\selectfont \textcolor{ForestGreen}{$ \blacktriangle\,8.1\times $}}\end{tabular} & \hspace*{1pt}\begin{tabular}[t]{@{}c@{}}\rule{0pt}{12pt}\colorbox{BenchBest}{\strut 2.6s}\\[-.2pt]{\fontsize{5.8}{6.4}\selectfont \textcolor{ForestGreen}{$ \blacktriangle\,6.9\times $}}\\[-.2pt]\colorbox{BenchBest}{\strut 1.9s}\\[-.2pt]{\fontsize{5.8}{6.4}\selectfont \textcolor{ForestGreen}{$ \blacktriangle\,8.6\times $}}\end{tabular} & \hspace*{1pt}\begin{tabular}[t]{@{}c@{}}\rule{0pt}{12pt}\colorbox{BenchBest}{\strut 5.4s}\\[-.2pt]{\fontsize{5.8}{6.4}\selectfont \textcolor{ForestGreen}{$ \blacktriangle\,7.9\times $}}\\[-.2pt]\colorbox{BenchBest}{\strut 3.4s}\\[-.2pt]{\fontsize{5.8}{6.4}\selectfont \textcolor{ForestGreen}{$ \blacktriangle\,11.7\times $}}\end{tabular} & \hspace*{1pt}\begin{tabular}[t]{@{}c@{}}\rule{0pt}{12pt}\colorbox{BenchBest}{\strut 4.5s}\\[-.2pt]{\fontsize{5.8}{6.4}\selectfont \textcolor{ForestGreen}{$ \blacktriangle\,86.7\times $}\textsuperscript{g}}\\[-.2pt]\colorbox{BenchBest}{\strut 3.6s}\\[-.2pt]{\fontsize{5.8}{6.4}\selectfont \textcolor{ForestGreen}{$ \blacktriangle\,97.9\times $}}\end{tabular} & \hspace*{1pt}\begin{tabular}[t]{@{}c@{}}\rule{0pt}{12pt}\colorbox{BenchBest}{\strut 7.5s}\\[-.2pt]{\fontsize{5.8}{6.4}\selectfont \textcolor{ForestGreen}{$ \blacktriangle\,19\% $}}\\[-.2pt]\colorbox{BenchBest}{\strut 5.2s}\\[-.2pt]{\fontsize{5.8}{6.4}\selectfont \textcolor{ForestGreen}{$ \blacktriangle\,1.6\times $}}\end{tabular} & \hspace*{1pt}\begin{tabular}[t]{@{}c@{}}\rule{0pt}{12pt}\colorbox{BenchBest}{\strut 5.8s}\\[-.2pt]{\fontsize{5.8}{6.4}\selectfont \textcolor{ForestGreen}{$ \blacktriangle\,1.5\times $}}\\[-.2pt]\colorbox{BenchBest}{\strut 4s}\\[-.2pt]{\fontsize{5.8}{6.4}\selectfont \textcolor{ForestGreen}{$ \blacktriangle\,1.7\times $}}\end{tabular} & \hspace*{1pt}\begin{tabular}[t]{@{}c@{}}\rule{0pt}{12pt}\colorbox{BenchBest}{\strut 12.6s}\\[-.2pt]{\fontsize{5.8}{6.4}\selectfont \textcolor{ForestGreen}{$ \blacktriangle\,21\% $}}\\[-.2pt]\colorbox{BenchBest}{\strut 7.9s}\\[-.2pt]{\fontsize{5.8}{6.4}\selectfont \textcolor{ForestGreen}{$ \blacktriangle\,1.6\times $}}\end{tabular} & \hspace*{1pt}\begin{tabular}[t]{@{}c@{}}\rule{0pt}{12pt}\colorbox{BenchBest}{\strut 12.9s}\\[-.2pt]{\fontsize{5.8}{6.4}\selectfont \textcolor{ForestGreen}{$ \blacktriangle\,1.4\times $}}\\[-.2pt]\colorbox{BenchBest}{\strut 10.6s}\\[-.2pt]{\fontsize{5.8}{6.4}\selectfont \textcolor{ForestGreen}{$ \blacktriangle\,1.6\times $}}\end{tabular} & \hspace*{1pt}\begin{tabular}[t]{@{}c@{}}\rule{0pt}{12pt}1m\,24.8s\\[-.2pt]{\fontsize{5.8}{6.4}\selectfont \textcolor{BrickRed}{$ \blacktriangledown\,2.0\times $}}\\[-.2pt]1m\,16.9s\\[-.2pt]{\fontsize{5.8}{6.4}\selectfont \textcolor{BrickRed}{$ \blacktriangledown\,1.9\times $}}\end{tabular} & \hspace*{1pt}\begin{tabular}[t]{@{}c@{}}\rule{0pt}{12pt}1m\,25.4s\\[-.2pt]{\fontsize{5.8}{6.4}\selectfont \textcolor{BrickRed}{$ \blacktriangledown\,2.0\times $}}\\[-.2pt]1m\,15.7s\\[-.2pt]{\fontsize{5.8}{6.4}\selectfont \textcolor{BrickRed}{$ \blacktriangledown\,1.8\times $}}\end{tabular} & \hspace*{1pt}\begin{tabular}[t]{@{}c@{}}\rule{0pt}{12pt}2m\,44s\\[-.2pt]{\fontsize{5.8}{6.4}\selectfont \textcolor{BrickRed}{$ \blacktriangledown\,1.9\times $}}\\[-.2pt]2m\,26.2s\\[-.2pt]{\fontsize{5.8}{6.4}\selectfont \textcolor{BrickRed}{$ \blacktriangledown\,1.7\times $}}\end{tabular} & \hspace*{1pt}\begin{tabular}[t]{@{}c@{}}\rule{0pt}{12pt}\colorbox{BenchBest}{\strut 1m\,42s}\\[-.2pt]{\fontsize{5.8}{6.4}\selectfont \textcolor{ForestGreen}{$ \blacktriangle\,2.9\times $}}\\[-.2pt]\colorbox{BenchBest}{\strut 1m\,35.7s}\\[-.2pt]{\fontsize{5.8}{6.4}\selectfont \textcolor{ForestGreen}{$ \blacktriangle\,3.0\times $}}\end{tabular} & \hspace*{1pt}\begin{tabular}[t]{@{}c@{}}\rule{0pt}{12pt}5m\,1.5s\\[-.2pt]{\fontsize{5.8}{6.4}\selectfont \textcolor{BrickRed}{$ \blacktriangledown\,8.6\times $}}\\[-.2pt]1m\,24.2s\\[-.2pt]{\fontsize{5.8}{6.4}\selectfont \textcolor{BrickRed}{$ \blacktriangledown\,2.5\times $}}\end{tabular} & \hspace*{1pt}\begin{tabular}[t]{@{}c@{}}\rule{0pt}{12pt}4m\,0.6s\\[-.2pt]{\fontsize{5.8}{6.4}\selectfont \textcolor{BrickRed}{$ \blacktriangledown\,6.7\times $}}\\[-.2pt]1m\,24.3s\\[-.2pt]{\fontsize{5.8}{6.4}\selectfont \textcolor{BrickRed}{$ \blacktriangledown\,2.5\times $}}\end{tabular} & \hspace*{1pt}\begin{tabular}[t]{@{}c@{}}\rule{0pt}{12pt}8m\,31.7s\\[-.2pt]{\fontsize{5.8}{6.4}\selectfont \textcolor{BrickRed}{$ \blacktriangledown\,6.9\times $}}\\[-.2pt]2m\,34.1s\\[-.2pt]{\fontsize{5.8}{6.4}\selectfont \textcolor{BrickRed}{$ \blacktriangledown\,2.2\times $}}\end{tabular} & \hspace*{1pt}\begin{tabular}[t]{@{}c@{}}\rule{0pt}{12pt}10m\,27.6s\\[-.2pt]{\fontsize{5.8}{6.4}\selectfont \textcolor{BrickRed}{$ \blacktriangledown\,10.4\times $}}\\[-.2pt]1m\,51.4s\\[-.2pt]{\fontsize{5.8}{6.4}\selectfont \textcolor{BrickRed}{$ \blacktriangledown\,1.9\times $}}\end{tabular}\\[3pt]\hline
\BenchmarkSourceLabel{bench:external-2-k20-w15:row:blastn}{R3}\shortstack[l]{{\color{gray}R3}\\blastn} & Native\textsuperscript{*} & \hspace*{1pt}\begin{tabular}[t]{@{}c@{}}\rule{0pt}{12pt}2m\,19.5s\\[-.2pt]{\fontsize{5.8}{6.4}\selectfont \textcolor{BrickRed}{$ \blacktriangledown\,7.1\times $}}\\[-.2pt]2m\,14.5s\\[-.2pt]{\fontsize{5.8}{6.4}\selectfont \textcolor{BrickRed}{$ \blacktriangledown\,7.2\times $}}\end{tabular} & \hspace*{1pt}\begin{tabular}[t]{@{}c@{}}\rule{0pt}{12pt}3m\,12.3s\\[-.2pt]{\fontsize{5.8}{6.4}\selectfont \textcolor{BrickRed}{$ \blacktriangledown\,10.6\times $}}\\[-.2pt]3m\,1.8s\\[-.2pt]{\fontsize{5.8}{6.4}\selectfont \textcolor{BrickRed}{$ \blacktriangledown\,11.0\times $}}\end{tabular} & \hspace*{1pt}\begin{tabular}[t]{@{}c@{}}\rule{0pt}{12pt}4m\,30.1s\\[-.2pt]{\fontsize{5.8}{6.4}\selectfont \textcolor{BrickRed}{$ \blacktriangledown\,6.4\times $}}\\[-.2pt]4m\,16s\\[-.2pt]{\fontsize{5.8}{6.4}\selectfont \textcolor{BrickRed}{$ \blacktriangledown\,6.5\times $}}\end{tabular} & \hspace*{1pt}\begin{tabular}[t]{@{}c@{}}\rule{0pt}{12pt}2m\,27.9s\\[-.2pt]{\fontsize{5.8}{6.4}\selectfont \textcolor{ForestGreen}{$ \blacktriangle\,2.6\times $}}\\[-.2pt]2m\,22.7s\\[-.2pt]{\fontsize{5.8}{6.4}\selectfont \textcolor{ForestGreen}{$ \blacktriangle\,2.5\times $}}\end{tabular} & \hspace*{1pt}\begin{tabular}[t]{@{}c@{}}\rule{0pt}{12pt}3m\,27.6s\\[-.2pt]{\fontsize{5.8}{6.4}\selectfont \textcolor{BrickRed}{$ \blacktriangledown\,22.3\times $}\textsuperscript{a}}\\[-.2pt]2m\,22.7s\\[-.2pt]{\fontsize{5.8}{6.4}\selectfont \textcolor{BrickRed}{$ \blacktriangledown\,17.0\times $}}\end{tabular} & \hspace*{1pt}\begin{tabular}[t]{@{}c@{}}\rule{0pt}{12pt}4m\,37.9s\\[-.2pt]{\fontsize{5.8}{6.4}\selectfont \textcolor{BrickRed}{$ \blacktriangledown\,31.8\times $}\textsuperscript{a}}\\[-.2pt]2m\,47s\\[-.2pt]{\fontsize{5.8}{6.4}\selectfont \textcolor{BrickRed}{$ \blacktriangledown\,23.8\times $}}\end{tabular} & \hspace*{1pt}\begin{tabular}[t]{@{}c@{}}\rule{0pt}{12pt}7m\,14.2s\\[-.2pt]{\fontsize{5.8}{6.4}\selectfont \textcolor{BrickRed}{$ \blacktriangledown\,27.5\times $}\textsuperscript{a}}\\[-.2pt]4m\,26.5s\\[-.2pt]{\fontsize{5.8}{6.4}\selectfont \textcolor{BrickRed}{$ \blacktriangledown\,20.9\times $}}\end{tabular} & \hspace*{1pt}\begin{tabular}[t]{@{}c@{}}\rule{0pt}{12pt}3m\,30.4s\\[-.2pt]{\fontsize{5.8}{6.4}\selectfont \textcolor{BrickRed}{$ \blacktriangledown\,11.4\times $}\textsuperscript{a}}\\[-.2pt]2m\,14.4s\\[-.2pt]{\fontsize{5.8}{6.4}\selectfont \textcolor{BrickRed}{$ \blacktriangledown\,7.9\times $}}\end{tabular} & \hspace*{1pt}\begin{tabular}[t]{@{}c@{}}\rule{0pt}{12pt}2m\,28.5s\\[-.2pt]{\fontsize{5.8}{6.4}\selectfont \textcolor{BrickRed}{$ \blacktriangledown\,3.5\times $}}\\[-.2pt]2m\,23.4s\\[-.2pt]{\fontsize{5.8}{6.4}\selectfont \textcolor{BrickRed}{$ \blacktriangledown\,3.5\times $}}\end{tabular} & \hspace*{1pt}\begin{tabular}[t]{@{}c@{}}\rule{0pt}{12pt}3m\,20.9s\\[-.2pt]{\fontsize{5.8}{6.4}\selectfont \textcolor{BrickRed}{$ \blacktriangledown\,4.7\times $}}\\[-.2pt]3m\,10.4s\\[-.2pt]{\fontsize{5.8}{6.4}\selectfont \textcolor{BrickRed}{$ \blacktriangledown\,4.6\times $}}\end{tabular} & \hspace*{1pt}\begin{tabular}[t]{@{}c@{}}\rule{0pt}{12pt}4m\,46.7s\\[-.2pt]{\fontsize{5.8}{6.4}\selectfont \textcolor{BrickRed}{$ \blacktriangledown\,3.3\times $}}\\[-.2pt]4m\,32.6s\\[-.2pt]{\fontsize{5.8}{6.4}\selectfont \textcolor{BrickRed}{$ \blacktriangledown\,3.3\times $}}\end{tabular} & \hspace*{1pt}\begin{tabular}[t]{@{}c@{}}\rule{0pt}{12pt}2m\,38.7s\\[-.2pt]{\fontsize{5.8}{6.4}\selectfont \textcolor{ForestGreen}{$ \blacktriangle\,1.8\times $}}\\[-.2pt]2m\,33.5s\\[-.2pt]{\fontsize{5.8}{6.4}\selectfont \textcolor{ForestGreen}{$ \blacktriangle\,1.9\times $}}\end{tabular} & \hspace*{1pt}\begin{tabular}[t]{@{}c@{}}\rule{0pt}{12pt}3m\,43.6s\\[-.2pt]{\fontsize{5.8}{6.4}\selectfont \textcolor{BrickRed}{$ \blacktriangledown\,6.4\times $}\textsuperscript{a}}\\[-.2pt]2m\,37.7s\\[-.2pt]{\fontsize{5.8}{6.4}\selectfont \textcolor{BrickRed}{$ \blacktriangledown\,4.6\times $}}\end{tabular} & \hspace*{1pt}\begin{tabular}[t]{@{}c@{}}\rule{0pt}{12pt}4m\,51.4s\\[-.2pt]{\fontsize{5.8}{6.4}\selectfont \textcolor{BrickRed}{$ \blacktriangledown\,8.1\times $}\textsuperscript{a}}\\[-.2pt]2m\,59.6s\\[-.2pt]{\fontsize{5.8}{6.4}\selectfont \textcolor{BrickRed}{$ \blacktriangledown\,5.3\times $}}\end{tabular} & \hspace*{1pt}\begin{tabular}[t]{@{}c@{}}\rule{0pt}{12pt}7m\,39.4s\\[-.2pt]{\fontsize{5.8}{6.4}\selectfont \textcolor{BrickRed}{$ \blacktriangledown\,6.2\times $}\textsuperscript{a}}\\[-.2pt]4m\,49.6s\\[-.2pt]{\fontsize{5.8}{6.4}\selectfont \textcolor{BrickRed}{$ \blacktriangledown\,4.1\times $}}\end{tabular} & \hspace*{1pt}\begin{tabular}[t]{@{}c@{}}\rule{0pt}{12pt}3m\,50.8s\\[-.2pt]{\fontsize{5.8}{6.4}\selectfont \textcolor{BrickRed}{$ \blacktriangledown\,3.8\times $}\textsuperscript{a}}\\[-.2pt]2m\,36.2s\\[-.2pt]{\fontsize{5.8}{6.4}\selectfont \textcolor{BrickRed}{$ \blacktriangledown\,2.7\times $}}\end{tabular}\\[3pt]\hline
\BenchmarkSourceLabel{bench:external-2-k20-w15:row:winnowmap}{R4}\shortstack[l]{{\color{gray}R4}\\winnowmap} & Native\textsuperscript{*} & \hspace*{1pt}\begin{tabular}[t]{@{}c@{}}\rule{0pt}{12pt}28.4s\\[-.2pt]{\fontsize{5.8}{6.4}\selectfont \textcolor{BrickRed}{$ \blacktriangledown\,1.4\times $}}\\[-.2pt]26.9s\\[-.2pt]{\fontsize{5.8}{6.4}\selectfont \textcolor{BrickRed}{$ \blacktriangledown\,1.4\times $}}\end{tabular} & \hspace*{1pt}\begin{tabular}[t]{@{}c@{}}\rule{0pt}{12pt}27.6s\\[-.2pt]{\fontsize{5.8}{6.4}\selectfont \textcolor{BrickRed}{$ \blacktriangledown\,1.5\times $}}\\[-.2pt]24.8s\\[-.2pt]{\fontsize{5.8}{6.4}\selectfont \textcolor{BrickRed}{$ \blacktriangledown\,1.5\times $}}\end{tabular} & \hspace*{1pt}\begin{tabular}[t]{@{}c@{}}\rule{0pt}{12pt}1m\,2.8s\\[-.2pt]{\fontsize{5.8}{6.4}\selectfont \textcolor{BrickRed}{$ \blacktriangledown\,1.5\times $}}\\[-.2pt]55s\\[-.2pt]{\fontsize{5.8}{6.4}\selectfont \textcolor{BrickRed}{$ \blacktriangledown\,1.4\times $}}\end{tabular} & \hspace*{1pt}\begin{tabular}[t]{@{}c@{}}\rule{0pt}{12pt}8m\,35.8s\\[-.2pt]{\fontsize{5.8}{6.4}\selectfont \textcolor{BrickRed}{$ \blacktriangledown\,1.3\times $}}\\[-.2pt]8m\,4.4s\\[-.2pt]{\fontsize{5.8}{6.4}\selectfont \textcolor{BrickRed}{$ \blacktriangledown\,1.4\times $}}\end{tabular} & \hspace*{1pt}\begin{tabular}[t]{@{}c@{}}\rule{0pt}{12pt}10.8s\\[-.2pt]{\fontsize{5.8}{6.4}\selectfont \textcolor{BrickRed}{$ \blacktriangledown\,16\% $}}\\[-.2pt]9.4s\\[-.2pt]{\fontsize{5.8}{6.4}\selectfont \textcolor{BrickRed}{$ \blacktriangledown\,12\% $}}\end{tabular} & \hspace*{1pt}\begin{tabular}[t]{@{}c@{}}\rule{0pt}{12pt}10.9s\\[-.2pt]{\fontsize{5.8}{6.4}\selectfont \textcolor{BrickRed}{$ \blacktriangledown\,24\% $}}\\[-.2pt]8s\\[-.2pt]{\fontsize{5.8}{6.4}\selectfont \textcolor{BrickRed}{$ \blacktriangledown\,14\% $}}\end{tabular} & \hspace*{1pt}\begin{tabular}[t]{@{}c@{}}\rule{0pt}{12pt}20.1s\\[-.2pt]{\fontsize{5.8}{6.4}\selectfont \textcolor{BrickRed}{$ \blacktriangledown\,27\% $}}\\[-.2pt]14.5s\\[-.2pt]{\fontsize{5.8}{6.4}\selectfont \textcolor{BrickRed}{$ \blacktriangledown\,14\% $}}\end{tabular} & \hspace*{1pt}\begin{tabular}[t]{@{}c@{}}\rule{0pt}{12pt}21s\\[-.2pt]{\fontsize{5.8}{6.4}\selectfont \textcolor{BrickRed}{$ \blacktriangledown\,14\% $}}\\[-.2pt]17.1s\\[-.2pt]{\fontsize{5.8}{6.4}\selectfont \textcolor{gray}{$ \blacktriangleright\,0\% $}}\end{tabular} & \hspace*{1pt}\begin{tabular}[t]{@{}c@{}}\rule{0pt}{12pt}4m\,16.8s\\[-.2pt]{\fontsize{5.8}{6.4}\selectfont \textcolor{BrickRed}{$ \blacktriangledown\,6.1\times $}}\\[-.2pt]4m\,15.2s\\[-.2pt]{\fontsize{5.8}{6.4}\selectfont \textcolor{BrickRed}{$ \blacktriangledown\,6.2\times $}}\end{tabular} & \hspace*{1pt}\begin{tabular}[t]{@{}c@{}}\rule{0pt}{12pt}5m\,0.6s\\[-.2pt]{\fontsize{5.8}{6.4}\selectfont \textcolor{BrickRed}{$ \blacktriangledown\,7.0\times $}}\\[-.2pt]4m\,57.9s\\[-.2pt]{\fontsize{5.8}{6.4}\selectfont \textcolor{BrickRed}{$ \blacktriangledown\,7.2\times $}}\end{tabular} & \hspace*{1pt}\begin{tabular}[t]{@{}c@{}}\rule{0pt}{12pt}10m\,23.6s\\[-.2pt]{\fontsize{5.8}{6.4}\selectfont \textcolor{BrickRed}{$ \blacktriangledown\,7.2\times $}}\\[-.2pt]10m\,18s\\[-.2pt]{\fontsize{5.8}{6.4}\selectfont \textcolor{BrickRed}{$ \blacktriangledown\,7.4\times $}}\end{tabular} & \hspace*{1pt}\begin{tabular}[t]{@{}c@{}}\rule{0pt}{12pt}16m\,47.8s\\[-.2pt]{\fontsize{5.8}{6.4}\selectfont \textcolor{BrickRed}{$ \blacktriangledown\,3.5\times $}}\\[-.2pt]16m\,11s\\[-.2pt]{\fontsize{5.8}{6.4}\selectfont \textcolor{BrickRed}{$ \blacktriangledown\,3.4\times $}}\end{tabular} & \hspace*{1pt}\begin{tabular}[t]{@{}c@{}}\rule{0pt}{12pt}3m\,21.4s\\[-.2pt]{\fontsize{5.8}{6.4}\selectfont \textcolor{BrickRed}{$ \blacktriangledown\,5.8\times $}}\\[-.2pt]3m\,19.9s\\[-.2pt]{\fontsize{5.8}{6.4}\selectfont \textcolor{BrickRed}{$ \blacktriangledown\,5.9\times $}}\end{tabular} & \hspace*{1pt}\begin{tabular}[t]{@{}c@{}}\rule{0pt}{12pt}4m\,0.5s\\[-.2pt]{\fontsize{5.8}{6.4}\selectfont \textcolor{BrickRed}{$ \blacktriangledown\,6.7\times $}}\\[-.2pt]3m\,57.6s\\[-.2pt]{\fontsize{5.8}{6.4}\selectfont \textcolor{BrickRed}{$ \blacktriangledown\,7.0\times $}}\end{tabular} & \hspace*{1pt}\begin{tabular}[t]{@{}c@{}}\rule{0pt}{12pt}8m\,56.2s\\[-.2pt]{\fontsize{5.8}{6.4}\selectfont \textcolor{BrickRed}{$ \blacktriangledown\,7.2\times $}}\\[-.2pt]8m\,50.6s\\[-.2pt]{\fontsize{5.8}{6.4}\selectfont \textcolor{BrickRed}{$ \blacktriangledown\,7.5\times $}}\end{tabular} & \hspace*{1pt}\begin{tabular}[t]{@{}c@{}}\rule{0pt}{12pt}4m\,43.2s\\[-.2pt]{\fontsize{5.8}{6.4}\selectfont \textcolor{BrickRed}{$ \blacktriangledown\,4.7\times $}}\\[-.2pt]4m\,39.3s\\[-.2pt]{\fontsize{5.8}{6.4}\selectfont \textcolor{BrickRed}{$ \blacktriangledown\,4.8\times $}}\end{tabular}\\[3pt]\hline
\BenchmarkSourceLabel{bench:external-2-k20-w15:row:lexicmap}{R5}\shortstack[l]{{\color{gray}R5}\\lexicmap} & Native\textsuperscript{*} & \hspace*{1pt}\begin{tabular}[t]{@{}c@{}}\rule{0pt}{12pt}38.8s\\[-.2pt]{\fontsize{5.8}{6.4}\selectfont \textcolor{BrickRed}{$ \blacktriangledown\,2.0\times $}}\\[-.2pt]38.6s\\[-.2pt]{\fontsize{5.8}{6.4}\selectfont \textcolor{BrickRed}{$ \blacktriangledown\,2.1\times $}}\end{tabular} & \hspace*{1pt}\begin{tabular}[t]{@{}c@{}}\rule{0pt}{12pt}54.3s\\[-.2pt]{\fontsize{5.8}{6.4}\selectfont \textcolor{BrickRed}{$ \blacktriangledown\,3.0\times $}}\\[-.2pt]54s\\[-.2pt]{\fontsize{5.8}{6.4}\selectfont \textcolor{BrickRed}{$ \blacktriangledown\,3.3\times $}}\end{tabular} & \hspace*{1pt}\begin{tabular}[t]{@{}c@{}}\rule{0pt}{12pt}1m\,29.8s\\[-.2pt]{\fontsize{5.8}{6.4}\selectfont \textcolor{BrickRed}{$ \blacktriangledown\,2.1\times $}}\\[-.2pt]1m\,29.4s\\[-.2pt]{\fontsize{5.8}{6.4}\selectfont \textcolor{BrickRed}{$ \blacktriangledown\,2.3\times $}}\end{tabular} & \hspace*{1pt}\begin{tabular}[t]{@{}c@{}}\rule{0pt}{12pt}49.9s\\[-.2pt]{\fontsize{5.8}{6.4}\selectfont \textcolor{ForestGreen}{$ \blacktriangle\,7.8\times $}}\\[-.2pt]49.7s\\[-.2pt]{\fontsize{5.8}{6.4}\selectfont \textcolor{ForestGreen}{$ \blacktriangle\,7.2\times $}}\end{tabular} & \hspace*{1pt}\begin{tabular}[t]{@{}c@{}}\rule{0pt}{12pt}37.8s\\[-.2pt]{\fontsize{5.8}{6.4}\selectfont \textcolor{BrickRed}{$ \blacktriangledown\,4.1\times $}\textsuperscript{b}}\\[-.2pt]37.3s\\[-.2pt]{\fontsize{5.8}{6.4}\selectfont \textcolor{BrickRed}{$ \blacktriangledown\,4.4\times $}}\end{tabular} & \hspace*{1pt}\begin{tabular}[t]{@{}c@{}}\rule{0pt}{12pt}53.5s\\[-.2pt]{\fontsize{5.8}{6.4}\selectfont \textcolor{BrickRed}{$ \blacktriangledown\,6.1\times $}\textsuperscript{b}}\\[-.2pt]52.8s\\[-.2pt]{\fontsize{5.8}{6.4}\selectfont \textcolor{BrickRed}{$ \blacktriangledown\,7.5\times $}}\end{tabular} & \hspace*{1pt}\begin{tabular}[t]{@{}c@{}}\rule{0pt}{12pt}1m\,18.6s\\[-.2pt]{\fontsize{5.8}{6.4}\selectfont \textcolor{BrickRed}{$ \blacktriangledown\,5.0\times $}\textsuperscript{b}}\\[-.2pt]1m\,17.5s\\[-.2pt]{\fontsize{5.8}{6.4}\selectfont \textcolor{BrickRed}{$ \blacktriangledown\,6.1\times $}}\end{tabular} & \hspace*{1pt}\begin{tabular}[t]{@{}c@{}}\rule{0pt}{12pt}49.7s\\[-.2pt]{\fontsize{5.8}{6.4}\selectfont \textcolor{BrickRed}{$ \blacktriangledown\,2.7\times $}\textsuperscript{b}}\\[-.2pt]49.3s\\[-.2pt]{\fontsize{5.8}{6.4}\selectfont \textcolor{BrickRed}{$ \blacktriangledown\,2.9\times $}}\end{tabular} & \hspace*{1pt}\begin{tabular}[t]{@{}c@{}}\rule{0pt}{12pt}1m\,41.8s\\[-.2pt]{\fontsize{5.8}{6.4}\selectfont \textcolor{BrickRed}{$ \blacktriangledown\,2.4\times $}}\\[-.2pt]1m\,41.7s\\[-.2pt]{\fontsize{5.8}{6.4}\selectfont \textcolor{BrickRed}{$ \blacktriangledown\,2.5\times $}}\end{tabular} & \hspace*{1pt}\begin{tabular}[t]{@{}c@{}}\rule{0pt}{12pt}2m\,35.6s\\[-.2pt]{\fontsize{5.8}{6.4}\selectfont \textcolor{BrickRed}{$ \blacktriangledown\,3.6\times $}}\\[-.2pt]2m\,35.4s\\[-.2pt]{\fontsize{5.8}{6.4}\selectfont \textcolor{BrickRed}{$ \blacktriangledown\,3.8\times $}}\end{tabular} & \hspace*{1pt}\begin{tabular}[t]{@{}c@{}}\rule{0pt}{12pt}4m\,6.7s\\[-.2pt]{\fontsize{5.8}{6.4}\selectfont \textcolor{BrickRed}{$ \blacktriangledown\,2.8\times $}}\\[-.2pt]4m\,6.3s\\[-.2pt]{\fontsize{5.8}{6.4}\selectfont \textcolor{BrickRed}{$ \blacktriangledown\,2.9\times $}}\end{tabular} & \hspace*{1pt}\begin{tabular}[t]{@{}c@{}}\rule{0pt}{12pt}3m\,25.5s\\[-.2pt]{\fontsize{5.8}{6.4}\selectfont \textcolor{ForestGreen}{$ \blacktriangle\,30\% $}}\\[-.2pt]3m\,25.3s\\[-.2pt]{\fontsize{5.8}{6.4}\selectfont \textcolor{ForestGreen}{$ \blacktriangle\,28\% $}}\end{tabular} & \hspace*{1pt}\begin{tabular}[t]{@{}c@{}}\rule{0pt}{12pt}1m\,43.4s\\[-.2pt]{\fontsize{5.8}{6.4}\selectfont \textcolor{BrickRed}{$ \blacktriangledown\,3.0\times $}\textsuperscript{b}}\\[-.2pt]1m\,42.6s\\[-.2pt]{\fontsize{5.8}{6.4}\selectfont \textcolor{BrickRed}{$ \blacktriangledown\,3.0\times $}}\end{tabular} & \hspace*{1pt}\begin{tabular}[t]{@{}c@{}}\rule{0pt}{12pt}2m\,37.7s\\[-.2pt]{\fontsize{5.8}{6.4}\selectfont \textcolor{BrickRed}{$ \blacktriangledown\,4.4\times $}\textsuperscript{b}}\\[-.2pt]2m\,37.1s\\[-.2pt]{\fontsize{5.8}{6.4}\selectfont \textcolor{BrickRed}{$ \blacktriangledown\,4.6\times $}}\end{tabular} & \hspace*{1pt}\begin{tabular}[t]{@{}c@{}}\rule{0pt}{12pt}3m\,45.3s\\[-.2pt]{\fontsize{5.8}{6.4}\selectfont \textcolor{BrickRed}{$ \blacktriangledown\,3.0\times $}\textsuperscript{b}}\\[-.2pt]3m\,43.9s\\[-.2pt]{\fontsize{5.8}{6.4}\selectfont \textcolor{BrickRed}{$ \blacktriangledown\,3.2\times $}}\end{tabular} & \hspace*{1pt}\begin{tabular}[t]{@{}c@{}}\rule{0pt}{12pt}3m\,16.3s\\[-.2pt]{\fontsize{5.8}{6.4}\selectfont \textcolor{BrickRed}{$ \blacktriangledown\,3.3\times $}\textsuperscript{b}}\\[-.2pt]3m\,15.8s\\[-.2pt]{\fontsize{5.8}{6.4}\selectfont \textcolor{BrickRed}{$ \blacktriangledown\,3.3\times $}}\end{tabular}\\[3pt]\hline
\end{tabular}\par}
\smallskip{\footnotesize\noindent Each cell: A, comparison, B, comparison. $\cdot$: no measurement. n/a: no comparison.\par
\noindent Blue: one lowest A and one lowest B measurement for each column for each configuration.\par
\noindent Gray: minimap2. Yellow: my preferred MPF / Bounded / SIMD-DP / Block configuration.\par
\noindent Comparisons use minimap2. The minimap2 row compares with BioseqDB.\par
\noindent\textsuperscript{*} Each algorithm except BLAST required intensive splitting, sometimes not equal. The splitting was the biggest possible to allow algorithm not to crash.\par
\noindent\textsuperscript{a}BLAST scores N as -3, compared with -1 in minimap2.\par
\noindent\textsuperscript{b}LexicMap stores reference N as A.\par

}
\end{landscape}
\restoregeometry
\clearpage\subsection{Vulkan on a Vast.ai GPU instance}\label{sec:vastai-gpu-results}
\begin{table}[H]\centering\small
\caption{Vast.ai machine specs}\label{tab:vastai-machine}
\begin{tabularx}{\linewidth}{@{}l>{\raggedright\arraybackslash}X@{}}\toprule
Property & Value\\\midrule
CPU & AMD EPYC 7282 16-Core Processor\\[3pt]
CPU quota & 7.68 cores\\[3pt]
Memory limit & 120.8 GiB\\[3pt]
GPU & NVIDIA GeForce RTX 5090\\[3pt]
GPU memory & 32,607 MiB\\[3pt]
OS & Ubuntu 24.04\\[3pt]
NVIDIA driver & 595.84\\[3pt]
Vulkan API & 1.4.329\\[3pt]
\bottomrule\end{tabularx}\end{table}
\newgeometry{left=10mm,right=10mm,top=12mm,bottom=12mm}
\begin{landscape}
\clearpage
\refstepcounter{table}\label{tab:vastai-gpu}\label{tab:bench:vastai}
\noindent{\large\bfseries (J): Table \thetable: Vast.ai GPU, Vulkan vs minimap2\hfill 1/1}\par\smallskip
{\fontsize{7.1}{8.1}\selectfont\setlength{\tabcolsep}{1pt}\setlength{\fboxsep}{2pt}\renewcommand{\arraystretch}{1.12}
\noindent\begin{tabular}{@{}p{14mm}p{27mm}|p{13.55mm}|p{13.55mm}|p{13.55mm}|p{13.55mm}|@{}}\hline
Engine & Execution & \multicolumn{4}{c|}{\textbf{Index}}\\
 &  & \multicolumn{2}{c|}{\textbf{4 Gbp}} & \multicolumn{2}{c|}{\textbf{2 Gbp}}\\
 &  & \multicolumn{1}{c|}{\textbf{\shortstack{k=15\\w=10}}} & \multicolumn{1}{c|}{\textbf{\shortstack{k=20\\w=15}}} & \multicolumn{1}{c|}{\textbf{\shortstack{k=15\\w=10}}} & \multicolumn{1}{c|}{\textbf{\shortstack{k=20\\w=15}}}\\
 &  & \multicolumn{1}{c|}{\textbf{DNA4}} & \multicolumn{1}{c|}{\textbf{DNA4}} & \multicolumn{1}{c|}{\textbf{DNA4}} & \multicolumn{1}{c|}{\textbf{DNA4}}\\
 & & \BenchmarkSourceLabel{bench:vastai:column:index-4-k15-w10-dna4-vastai}{C1}{\color{gray}C1} & \BenchmarkSourceLabel{bench:vastai:column:index-4-k20-w15-dna4-vastai}{C2}{\color{gray}C2} & \BenchmarkSourceLabel{bench:vastai:column:index-2-k15-w10-dna4-vastai}{C3}{\color{gray}C3} & \BenchmarkSourceLabel{bench:vastai:column:index-2-k20-w15-dna4-vastai}{C4}{\color{gray}C4}\\\hline
\rowcolor{BenchBaseline}[0.5pt][0.5pt]\BenchmarkSourceLabel{bench:vastai:row:minimap2}{R1}\shortstack[l]{{\color{gray}R1}\\minimap2} & \textbf{Baseline} & \hspace*{1pt}\begin{tabular}[t]{@{}c@{}}\rule{0pt}{12pt}\colorbox{BenchBest}{\strut 19.5s}\\[-.2pt]{\fontsize{5.8}{6.4}\selectfont \textcolor{ForestGreen}{$ \blacktriangle\,21.1\times $}}\\[-.2pt]\colorbox{BenchBest}{\strut 18s}\\[-.2pt]{\fontsize{5.8}{6.4}\selectfont \textcolor{ForestGreen}{$ \blacktriangle\,19.3\times $}}\end{tabular} & \hspace*{1pt}\begin{tabular}[t]{@{}c@{}}\rule{0pt}{12pt}\colorbox{BenchBest}{\strut 42.5s}\\[-.2pt]{\fontsize{5.8}{6.4}\selectfont \textcolor{ForestGreen}{$ \blacktriangle\,7.6\times $}}\\[-.2pt]\colorbox{BenchBest}{\strut 41.5s}\\[-.2pt]{\fontsize{5.8}{6.4}\selectfont \textcolor{ForestGreen}{$ \blacktriangle\,6.3\times $}}\end{tabular} & \hspace*{1pt}\begin{tabular}[t]{@{}c@{}}\rule{0pt}{12pt}\colorbox{BenchBest}{\strut 14.1s}\\[-.2pt]{\fontsize{5.8}{6.4}\selectfont \textcolor{ForestGreen}{$ \blacktriangle\,7.7\times $}}\\[-.2pt]\colorbox{BenchBest}{\strut 13s}\\[-.2pt]{\fontsize{5.8}{6.4}\selectfont \textcolor{ForestGreen}{$ \blacktriangle\,4.5\times $}}\end{tabular} & \hspace*{1pt}\begin{tabular}[t]{@{}c@{}}\rule{0pt}{12pt}\colorbox{BenchBest}{\strut 1m\,34.6s}\\[-.2pt]{\fontsize{5.8}{6.4}\selectfont \textcolor{gray}{$ \blacktriangleright\,4\% $}}\\[-.2pt]1m\,33.6s\\[-.2pt]{\fontsize{5.8}{6.4}\selectfont \textcolor{BrickRed}{$ \blacktriangledown\,2.1\times $}}\end{tabular}\\[3pt]\hline
\BenchmarkSourceLabel{bench:vastai:row:mps-vulkan-vulkan-ksw2}{R2}\shortstack[l]{{\color{gray}R2}\\BioseqDB} & \begin{tabular}[t]{@{}l@{}}\rule{0pt}{12pt}MPS / Vulkan\\Vulkan-DP / KSW2\end{tabular} & \textcolor{BrickRed}{\ensuremath{\times}}\textsuperscript{2} & \textcolor{BrickRed}{\ensuremath{\times}}\textsuperscript{2} & \hspace*{1pt}\begin{tabular}[t]{@{}c@{}}\rule{0pt}{12pt}1m\,54.9s\\[-.2pt]{\fontsize{5.8}{6.4}\selectfont \textcolor{BrickRed}{$ \blacktriangledown\,8.2\times $}}\\[-.2pt]1m\,2.3s\\[-.2pt]{\fontsize{5.8}{6.4}\selectfont \textcolor{BrickRed}{$ \blacktriangledown\,4.8\times $}}\end{tabular} & \hspace*{1pt}\begin{tabular}[t]{@{}c@{}}\rule{0pt}{12pt}1m\,38.5s\\[-.2pt]{\fontsize{5.8}{6.4}\selectfont \textcolor{gray}{$ \blacktriangleright\,4\% $}}\\[-.2pt]\colorbox{BenchBest}{\strut 44.1s}\\[-.2pt]{\fontsize{5.8}{6.4}\selectfont \textcolor{ForestGreen}{$ \blacktriangle\,2.1\times $}}\end{tabular}\\[3pt]\hline
\BenchmarkSourceLabel{bench:vastai:row:mps-vulkan-vulkan-block}{R3}\shortstack[l]{{\color{gray}R3}\\BioseqDB} & \begin{tabular}[t]{@{}l@{}}\rule{0pt}{12pt}MPS / Vulkan\\Vulkan-DP / Block\end{tabular} & \textcolor{BrickRed}{\ensuremath{\times}}\textsuperscript{2} & \textcolor{BrickRed}{\ensuremath{\times}}\textsuperscript{1} & \hspace*{1pt}\begin{tabular}[t]{@{}c@{}}\rule{0pt}{12pt}1m\,48.7s\\[-.2pt]{\fontsize{5.8}{6.4}\selectfont \textcolor{BrickRed}{$ \blacktriangledown\,7.7\times $}}\\[-.2pt]58.9s\\[-.2pt]{\fontsize{5.8}{6.4}\selectfont \textcolor{BrickRed}{$ \blacktriangledown\,4.5\times $}}\end{tabular} & \hspace*{1pt}\begin{tabular}[t]{@{}c@{}}\rule{0pt}{12pt}1m\,38.5s\\[-.2pt]{\fontsize{5.8}{6.4}\selectfont \textcolor{gray}{$ \blacktriangleright\,4\% $}}\\[-.2pt]47s\\[-.2pt]{\fontsize{5.8}{6.4}\selectfont \textcolor{ForestGreen}{$ \blacktriangle\,2.0\times $}}\end{tabular}\\[3pt]\hline
\BenchmarkSourceLabel{bench:vastai:row:mpf-vulkan-vulkan-ksw2}{R4}\shortstack[l]{{\color{gray}R4}\\BioseqDB} & \begin{tabular}[t]{@{}l@{}}\rule{0pt}{12pt}MPF / Vulkan\\Vulkan-DP / KSW2\end{tabular} & \hspace*{1pt}\begin{tabular}[t]{@{}c@{}}\rule{0pt}{12pt}6m\,50.8s\\[-.2pt]{\fontsize{5.8}{6.4}\selectfont \textcolor{BrickRed}{$ \blacktriangledown\,21.1\times $}}\\[-.2pt]5m\,48.5s\\[-.2pt]{\fontsize{5.8}{6.4}\selectfont \textcolor{BrickRed}{$ \blacktriangledown\,19.3\times $}}\end{tabular} & \hspace*{1pt}\begin{tabular}[t]{@{}c@{}}\rule{0pt}{12pt}5m\,24.2s\\[-.2pt]{\fontsize{5.8}{6.4}\selectfont \textcolor{BrickRed}{$ \blacktriangledown\,7.6\times $}}\\[-.2pt]4m\,22.1s\\[-.2pt]{\fontsize{5.8}{6.4}\selectfont \textcolor{BrickRed}{$ \blacktriangledown\,6.3\times $}}\end{tabular} & \hspace*{1pt}\begin{tabular}[t]{@{}c@{}}\rule{0pt}{12pt}3m\,38.5s\\[-.2pt]{\fontsize{5.8}{6.4}\selectfont \textcolor{BrickRed}{$ \blacktriangledown\,15.5\times $}}\\[-.2pt]2m\,49.8s\\[-.2pt]{\fontsize{5.8}{6.4}\selectfont \textcolor{BrickRed}{$ \blacktriangledown\,13.0\times $}}\end{tabular} & \hspace*{1pt}\begin{tabular}[t]{@{}c@{}}\rule{0pt}{12pt}2m\,58.3s\\[-.2pt]{\fontsize{5.8}{6.4}\selectfont \textcolor{BrickRed}{$ \blacktriangledown\,1.9\times $}}\\[-.2pt]2m\,9.4s\\[-.2pt]{\fontsize{5.8}{6.4}\selectfont \textcolor{BrickRed}{$ \blacktriangledown\,1.4\times $}}\end{tabular}\\[3pt]\hline
\BenchmarkSourceLabel{bench:vastai:row:mpf-vulkan-vulkan-block}{R5}\shortstack[l]{{\color{gray}R5}\\BioseqDB} & \begin{tabular}[t]{@{}l@{}}\rule{0pt}{12pt}MPF / Vulkan\\Vulkan-DP / Block\end{tabular} & \hspace*{1pt}\begin{tabular}[t]{@{}c@{}}\rule{0pt}{12pt}7m\,7.5s\\[-.2pt]{\fontsize{5.8}{6.4}\selectfont \textcolor{BrickRed}{$ \blacktriangledown\,21.9\times $}}\\[-.2pt]6m\,8.1s\\[-.2pt]{\fontsize{5.8}{6.4}\selectfont \textcolor{BrickRed}{$ \blacktriangledown\,20.4\times $}}\end{tabular} & \hspace*{1pt}\begin{tabular}[t]{@{}c@{}}\rule{0pt}{12pt}5m\,48s\\[-.2pt]{\fontsize{5.8}{6.4}\selectfont \textcolor{BrickRed}{$ \blacktriangledown\,8.2\times $}}\\[-.2pt]4m\,48s\\[-.2pt]{\fontsize{5.8}{6.4}\selectfont \textcolor{BrickRed}{$ \blacktriangledown\,6.9\times $}}\end{tabular} & \hspace*{1pt}\begin{tabular}[t]{@{}c@{}}\rule{0pt}{12pt}4m\,37s\\[-.2pt]{\fontsize{5.8}{6.4}\selectfont \textcolor{BrickRed}{$ \blacktriangledown\,19.7\times $}}\\[-.2pt]3m\,47s\\[-.2pt]{\fontsize{5.8}{6.4}\selectfont \textcolor{BrickRed}{$ \blacktriangledown\,17.4\times $}}\end{tabular} & \hspace*{1pt}\begin{tabular}[t]{@{}c@{}}\rule{0pt}{12pt}3m\,8.4s\\[-.2pt]{\fontsize{5.8}{6.4}\selectfont \textcolor{BrickRed}{$ \blacktriangledown\,2.0\times $}}\\[-.2pt]2m\,21.2s\\[-.2pt]{\fontsize{5.8}{6.4}\selectfont \textcolor{BrickRed}{$ \blacktriangledown\,1.5\times $}}\end{tabular}\\[3pt]\hline
\end{tabular}\par}
\smallskip{\footnotesize\noindent Each cell: A, comparison, B, comparison. Light gray = minimap2.\par
\noindent Blue: one lowest A and one lowest B measurement for each column for each configuration.\par
\noindent MPF/MPS = minimizers per fragment/sequence. Vulkan rows compare with minimap2 on the same machine.\par
\noindent $\cdot$ = no result. n/a: no comparison. \textcolor{BrickRed}{$\bigcirc$} = setup rejected. \textcolor{BrickRed}{$\times$} = execution failed.\par
\noindent\textsuperscript{1}The GPU index exceeds the available memory budget.\par
\noindent\textsuperscript{2}The MPS control buffer exceeds the Vulkan buffer limit.\par
}
\end{landscape}
\restoregeometry

